% ===== tex/DM2022_0.tex =====


\section{Введение}
\begin{frame}
\frametitle{О чём речь}
{\large
 В этом курсе рассматриваются некоторые элементарные понятия
\em{дискретной}, или \em{конечной}, математики, то есть математики,
прежде всего изучающей конечные множества и различные структуры на
них. Это означает, что понятия континуума, предела или
непрерывности не являются предметом изучения, хотя могут
использоваться как вспомогательные средства. Дискретная математика
имеет широкий спектр приложений, прежде всего в областях,
связанных с информационными технологиями и компьютерами. В самом
первоначальном, ныне редко используемом названии компьютера~---
<<электронная цифровая вычислительная машина>>~--- слово
<<цифровая>> указывает на принципиально дискретный характер работы
данного устройства. 
Современные компьютеры неотделимы от
современных программистов, для которых и
создан этот курс.
}

\end{frame}



\subsection{История развития курса}
\begin{frame}
\frametitle{История развития курса (1/3)}

Текущая реализация курса 
является результатом
достаточно длительного развития.

\smallskip
Авторский курс <<Дискретная математика>>  был поставлен на кафедре 
<<Прик\-ладная математика>> Ленинградского политехнического института осенью 1986 года,
читается непрерывно более 37 лет в различных университетах 
и из года в год совершенствуется, сохраняя при этом
преемственность.
В курсе меняется всё: 
распределение нагрузки по учебному плану, 
объём и сложность рассматриваемых вопросов,
состав преподавательского корпуса, применяемые методики преподавания,
учебные материалы. В то же время
основные цели, задачи  и конкурентные преимущества курса 
остаются постоянными.       

\smallskip 
Перечислим основные вехи развития курса.

\smallskip 
В 2000 году издательство <<Питер>> выпустило в свет \em{ учебное
пособие} Ф.~А.~Новикова <<Дискретная математика для
программистов>>. Книга оказалась удачной, о чём свидетельствовали
тиражи и письма читателей. Поэтому в 2004 году было выпущено второе, а в 2008 году~--- третье
издание с тем же названием, которые также пользовались широким спросом. Эти учебные пособия есть в библиотеке Политехнического университета. 
%Резонно
%предположить, что объём, стиль и содержание книги отвечают
%возрастающим потребностям студентов, изучающих дискретную
%математику, и инженеров, применяющих математические модели в своей
%работе. 

\end{frame}

\begin{frame}
\frametitle{История развития курса (2/3)}

%\vspace{-2pt}
В 2010 году было выпущено
новое издание, \em{ учебник}
<<Дискретная математика>>, созданный на основе учебного пособия
<<Дискретная математика для программистов>>. 
Содержа\-ние  учебника было приведено в соответствие с
новым федеральным государственным образовательным стандартом
высшего про\-фес\-сио\-нально\-го образования 
по направлению подготовки <<прикладная математика и информатика>>. Учебник получил рекомендацию Учебно-методического объединения по университетскому политехническому образованию.

\smallskip
В 2012 году
вышло в свет второе издание учебника <<Дискретная математика>>, 
в 2017 году вышло \em{ третье издание},
в значительной степени исправленное и дополненное. Это издание без изменений допечатывалось до 2022 года.
 

\smallskip
Все эти книги имеют 
единый план и стиль, но каждый следующий вариант 
более точен в деталях и подробностях,
содержит меньше ошибок,
включает большее количество понятий, фактов и примеров,
а потому имеет несколько больший объём,
консервативно расширяя и дополняя
предыдущие издания.
Благодаря внимательности студентов, 
количество ошибок в последних изданиях существенно меньше, 
чем в первых изданиях, поэтому желательно использовать последние издания.
\end{frame}

\begin{frame}
\frametitle{История развития курса (3/3)}
%\vspace{-2pt}
По мере появления технических
возможностей в курсе используются 
средства электронного обучения~--- 
сейчас это слайды, 
размещенные на используемой платформе дистанционного
обучения (в 2024 году это MS Teams).
В текущей реализации курса слайды и третье 
издание учебника созданы на основе единой базы текстов.
Это позволяет использовать как слайды, так и учебник
на аудиторных занятиях, для самостоятельной работы студентов
и для подготовки к экзамену.
По замечаниям студентов в курс постоянно
вносятся исправления и улучшения. 
Исправления в слайды вносятся оперативно,
а в учебник исправления не вносятся,
поэтому в случае расхождений верить следует последним версиям слайдов. 
     
\smallskip  
Слайды и учебник имеют общую  мелкую рубрикацию: параграф помещается на
нескольких страницах (или слайдах) и может быть прочтён за один
присест, не поднимая головы (или не отрывая взгляда от экрана).
Как в слайдах, так и в книге текст насыщен формулами и иллюстрациями.
Заголовки, формулы, алгоритмы и картинки~--- 
это существенно необходимые строительные блоки курса,
а поясняющие слова~--- неизбежный связующий раствор.
Для облегчения восприятия, поясняющий текст в слайдах 
несколько сокращён по сравнению с книгой, но с сохранением канвы изложения.

\end{frame}


\subsection{Цели курса}
\begin{frame}
\frametitle{Цели курса}
В целом курс преследует три основные цели.

%\vspace{-4pt}
\begin{enumerate}
\item
Познакомить студентов с максимально широким кругом понятий дискретной математики.
Тем самым у студентов
формируется \em{ терминологический запас}, необходимый для самостоятельного
изучения специальной математической и тео\-ре\-тико\-программистской литературы.
\item
Сообщить студентам необходимые \em{ конкретные сведения} из дискретной
математики, предусматриваемые Федеральным
государственным образовательным стандартом. Разборы доказательств приведённых утверждений и разнообразных примеров позволяют студентам овладеть методами дискретной математики, наиболее
употребительными при  решении практических задач.
\item
Пополнить запас примеров \em{ нетривиальных алгоритмов}. 
Изучение алгоритмов решения
типовых задач дискретной математики и способов представления математических
объектов в программах абсолютно необходимо практикующему программисту,
поскольку позволяет уменьшить трудозатраты на <<изобретение велосипеда>>.
\end{enumerate}

\end{frame}


\subsection{Структура курса}
\begin{frame}
\frametitle{Структура курса (1/2)}
%\vspace{-2pt}
Курс состоит из 10 глав (\em{ тем}), которые можно разделить на три группы. 

\smallskip
Первая группа, несколько б$\acute{\text{о}}$льшая по
объёму (темы~1, 3 и 5), содержит самые 
\em{ общие сведения} из основных
разделов дискретной математики. Доля алгоритмов в темах первой
группы несколько меньше, а доля определений несколько
больше, чем во второй и третьей группах.

\smallskip
Во второй группе
(темы~2,~4~и~6) собраны различные \em{ специальные разделы}, которые
можно включать или не включать в учебный курс в зависимости от
конкретных обстоятельств. 

\smallskip
Третья группа глав (темы~7--10) целиком
посвящена \em{ теории графов}
(в частности, алгоритмам над графами), 
поскольку теория графов наиболее широко применяется 
при построении дискретных математических моделей.

\smallskip
Главы (темы) делятся на \em{ разделы}, которые, в свою очередь, делятся на
\em{ параграфы}. 


\smallskip
Каждый раздел по объёму соответствует экзаменационному вопросу.
Параграфы в пределах раздела упорядочены по сложности: сначала
приводятся основные определения, а затем рассматриваются более
сложные вопросы. 



\end{frame}


\begin{frame}
\frametitle{Структура курса (2/2)}

\smallskip
В курсе имеются общие справочные механизмы, призванные облегчить работу с материалом,
особенно в случае выборочного изучения.
На \em{ титульном слайде} приведён адрес электронной почты
автора, по которому всегда можно обратиться с замечаниями и вопросами. 
\em{ Содержание} доведено до нижнего 
уровня рубрикации; элементы содержания 
работают как гиперссылки. 


В конце курса приведены \em{ указатель имён} 
и актуальный \em{ предметный указатель},
которым автор советует пользоваться для 
быстрого поиска определяющих вхождений терминов.


Каждому разделу 
предпослана таблица, в которой выписаны 
\em{ ключевые термины и обозначения}, вводимые в каждом параграфе. Фактически, это план ответа 
на экзаменационный вопрос, соответствующий разделу. 
В этой таблице знаком \RS ~отмечены
параграфы, которые нельзя опустить 
без ущерба для понимания остального материала,
а знаком \BS ~отмечены
параграфы, которые, напротив, можно опустить 
при первом чтении, поскольку они
содержат материал хотя и познавательный, но не обязательный по образовательному стандарту.  
В конце большинства разделов приводятся
диаграммы UML, которые в компактной 
графической форме представляют \em{ онтологию}
понятий этого раздела. Способ чтения
этих диаграмм указан ниже 
 в разделе <<Онтология дискретной математики>>.
 

\end{frame}
%\subsection{План 2019 года}
%\begin{frame}
%\frametitle{План 2019 года}
%
%\begin{enumerate}
%\item \em{ Множества и отношения}. Исключая характеристические функции, четыре лекции, восемь вопросов, 12, 16,
%19 и 26 февраля.
%\item 
%\em{ Алгебраические структуры}. Оставляем для самостоятельного изучения.
%\item 
%\em{ Булевы функции}. Без сокращений, три лекции, шесть вопросов, 2, 5 и 12 марта.
%\item
%\em{ Логические исчисления}. Без сокращений, три лекции, шесть вопросов, \\ 16, 19 и 23 марта.
%\item \em{ Комбинаторный анализ}.
%За исключением формул обращения, три лекции, шесть вопросов, 26 и 30 марта, 2 апреля. 
% \item \em{ Кодирование}.
%Оставляем для самостоятельного изучения.  
%\item  \em{ Графы}. Без сокращений, две лекции, пять вопросов,
%6 и 9 апреля.  
%\item  \em{ Связность}. Исключая потоки в сетях, две лекции, четыре вопроса, 13 и 16 апреля.    
%\item  \em{ Деревья}. Сокращая деревья сортировки, но включая двоичные кучи и деревья отрезков, три лекции, пять вопросов, 20, 23 и 27 апреля. 
%\item  \em{ Циклы, независимость и раскраска}. Без сокращения, три лекции, восемь вопросов, \\
%7, 14 и 21 мая.
%\end{enumerate}
%
%\end{frame}



\subsection{Используемые обозначения}
\begin{frame}
\frametitle{Используемые обозначения (1/8)}
Предметом курса является дискретная математика, но стиль изложения алгоритмический.
Предполагается, что студенты не испытывают затруднений в понимании
текстов программ на языке высокого уровня. Именно поэтому в
качестве нотации для записи алгоритмов используется \em{
некоторый\/} неспецифицированный алгоритмический язык (\em{
псевдокод\/}).

\smallskip
Ключевые слова псевдокода выделены полужирным шрифтом, идентификаторы
объeктов записываются курсивом и часто состоят из одной буквы, 
как это принято для математических
объектов, имена процедур записываются прямым шрифтом, 
как названия математических функций,
примечания отделяются двойной косой чертой //.

\smallskip
В языке используются общеупотребительные структуры
управления: \em{ ветвления} в краткой ({\bf if} ... {\bf then} ... {\bf
end~if}), полной ({\bf if} ... {\bf then} ... {\bf else} ... {\bf
end~if}) и
 расширенной ({\bf if} ... {\bf then} ... 
 {\bf elseif} ... {\bf then} ...
  {\bf else} ... {\bf
end~if}) формах, \em{ циклы} со счётчиком ({\bf for} ... {\bf from} ...
{\bf to} ... {\bf do}
 ... {\bf end~for}), с постусловием
({\bf~repeat}~...~{\bf~until}), с предусловием \linebreak[3]({\bf
while}~...~{\bf~do}~...~{\bfseries end~while}) и по множеству
({\bf for} ... {\bf do} ... {\bf end~for}), а также переходы ({\bf
go to} ...).
\end{frame}



\begin{frame}
\frametitle{Используемые обозначения (2/8)}
Для обозначения \em{ присваивания} используется традиционный знак
$\Assign$, вызов процедуры или функции ключевыми словами не
выделяется, выход из процедуры вместе с возвратом значения
обозначается ключевым словом {\bf return}.


\smallskip
Подразумевается строгая статическая \em{ типизация}, даже приведения не
используются, за исключением разыменования, которое
подразумевается везде, где оно нужно по здравому программистскому
смыслу. 


\smallskip
В то же время объявления локальных переменных почти всегда
опускаются, если их тип ясен из контекста.


\smallskip
Из структур данных считаются доступными массивы {\bf array} [...]
{\bf of} ..., \\структуры {\bf struct \{ } ... {\bf \}} и указатели $\uparrow$.
Указатель может иметь специальное значение
{\bf nil}, не указывающее ни на один объект.


\smallskip
<<Лишние>> разделители систематически опускаются, функции могут
возвращать несколько значений, и вообще, разрешаются любые вольности,
которые позволяют сделать тексты программ более краткими и читабельными.
\end{frame}


\begin{frame}
\frametitle{Используемые обозначения (3/8)}

%\vspace{-2pt}
Особого упоминания заслуживают три естественных оператора, аналоги
которых в явном виде редко встречаются в <<настоящих>> языках
программирования.
Оператор

\vspace{-2pt}
\begin{algorithmic}
\STATE \textbf{select} $m\in M$
\end{algorithmic}

\vspace{-3pt}
означает выбор \em{ произвольного} элемента $m$ из множества $M$.
Этот оператор часто необходим в <<переборных>> алгоритмах.
Оператор

\vspace{-2pt}
\begin{algorithmic}
\STATE \textbf{yield} $x$
\end{algorithmic}

\vspace{-3pt}
означает \em{ возврат} значения $x$, но при этом 
\em{ выполнение} функции \em{ не прекращается\/}, а продолжается со следующего оператора. Этот
оператор позволяет очень просто записать <<генерирующие>>
алгоритмы, результатом которых является некоторое заранее
неизвестное множество значений. Оператор

\vspace{-2pt}
\begin{algorithmic}
\STATE \textbf{next for} $x$
\end{algorithmic}

\vspace{-2pt}
означает прекращение выполнения текущего фрагмента программы и
продолжение выполнения \em{ со следующего шага цикла} по переменной $x$.
Этот оператор должен находиться в теле цикла по $x$ (на любом
уровне вложенности). 
\end{frame}



\begin{frame}
\frametitle{Используемые обозначения (4/8)}
Операторы \textbf{select}, \textbf{yield} и
\textbf{next for} не являются чем-то необычным и новым: первый из них
легко моделируется использованием датчика псевдослучайных чисел,
второй~--- присоединением элемента к результирующему множеству,
а~третий~--- переходом по метке. Однако указанные операторы сокращают запись и~повышают её наглядность.

\smallskip
В программах широко используются математические обозначения,
которые являются самоочевидными в конкретном контексте. Например,
конструкция

\begin{algorithmic}
\FOR{$x\in M$}
\STATE $P(x)$
\ENDFOR
\end{algorithmic}

означает применение процедуры $P$ ко всем
элементам множества $M$, а конструкция

\begin{algorithmic}
\STATE \textbf{$\sgn (x) \Assign$ if $x>0$ then $1$ elseif 
$x<0$ then $-1$ else $0$ end if}
\end{algorithmic}

означает определение функции <<знак числа>>.

\end{frame}



\begin{frame}
\frametitle{Используемые обозначения (5/8)}

Поскольку этот курс излагается на <<программно-математическом>>
языке, в нём не только математические обозначения используются в
программах, но и некоторые программистские обозначения
используются в математическом тексте.

\medskip
Прежде всего обсудим способы введения обозначений объектов и
операций. 

\medskip
Если обозначение объекта или операции считается
\em{ глобальным} (в пределах данного курса или в более широком контексте), то используется знак $\Def$,
который следует читать <<по определению есть>>. При этом слева от
знака $\Def$ может стоять обозначение объекта, а справа~---
выражение, определяющее этот объект. Например, формула
$$
\mathbb{R}_+\Def \Set{x\in\mathbb{R}}{x>0}
$$
определяет множество положительных вещественных чисел.



\end{frame}



\begin{frame}
\frametitle{Используемые обозначения (6/8)}

В случае определения операции в левой части стоит \em{
выражение\/}. В этом случае левую часть следует неформально
понимать как заголовок процедуры выполнения определяемой операции,
а правую часть~--- как тело этой процедуры. Например, формула
$$
A = B \Def A \subset B \And B \subset A
$$
означает следующее: <<чтобы вычислить значение выражения
$A = B$, нужно вычислить значения выражений $A \subset B$ и $B
\subset A$, а затем вычислить конъюнкцию этих значений>>.

\medskip
Отметим широкое использование символа присваивания
$\Assign$, который в математическом контексте можно читать как
<<положим>>. Если в левой части такого присваивания стоит простая
переменная, а в правой части~--- выражение, определяющее конкретный
объект, то это имеет очевидный смысл введения \em{ локального}
(в пределах доказательства и т.~п.)
обозначения. Например, запись $Z'\Assign Z\cup Z_{e_k} \setminus \{e_k\}$ означает:
<<положим, что $Z'$ содержит все элементы $Z$ и $Z_{e_k}$ за исключением $e_k$>>.

\end{frame}


\begin{frame}
\frametitle{Используемые обозначения (7/8)}

Наряду с
обычным для математики <<статическим>> использованием знака
$\Assign$, когда выражению в левой части придается постоянное в данном контексте
значение (определение константы), знак присваивания используется и
в <<динамическом>>, программистском смысле. Например, пусть в
графе $G$ есть вершина $v$. 
Тогда формулу $G\Assign G-v$
следует читать и понимать так: <<удалим в~графе $G$ вершину $v$>>.

\medskip
Одной из задач курса является выработка у студентов навыка чтения
математических текстов. Поэтому, начиная с самого первого слайда,
без дополнительных объяснений интенсивно используются язык
исчисления предикатов и другие общепринятые математические
обозначения. При этом стиль записи формул совершенно свободный и
неформальный. Например, вместо строгой формулы 
$\A{k}{(k<n)\Impl P(k)}$ может быть написано $\A{k<n}{P(k)}$
или даже $\forall\,k<n~P(k)$ в предположении, что читатель
знает или догадается, какие синтаксические упрощения используются.
\end{frame}



\begin{frame}
\frametitle{Используемые обозначения (8/8)}

В первых главах формулировки теорем
дублируются, то~есть приводятся на естественном языке и на языке
формул. Тем самым на примерах объясняется используемый язык
формул. Но в последних главах книги и в доказательствах
использование естественного языка сведено к минимуму. Это
существенно сокращает объём текста, но требует внимания при
чтении.

\medskip
Для краткости в тексте часто используются обороты, которые
подразумевают восстановление читателем опущенных слов по
контексту. Например, если символ $A$ означает множество, то вместо
аккуратного оборота <<где объект $x$ является
элементом множества $A$>>, может быть использована более короткая
фраза <<где $x$ принадлежит $A$>> или даже формула <<где $x\in A$>>.
В целом используемые обозначения строго следуют классическим
образцам, принятым в учебной математической и программистской
литературе. Непривычным может показаться только совместное
использование этих обозначений. Но такое
взаимопроникновение обозначений продиктовано основной задачей
курса.
\end{frame}


\subsection{Стандартные математические символы}
\begin{frame}
\frametitle{Стандартные математические символы (1/6)}
Некоторые обозначения являются глобальными, то есть означают одно
и то же в любом контексте (в этих слайдах).  Таких стандартных обозначений
не так уж много. Они собраны в следующих таблицах.

{\bf Метаобозначения}

\begin{tabular}{l|l|l}
\hline
 Обозначение &  Смысл & Примеры \\
\hline
   $\Def$    &  По определению есть &  $|\emptyset|\Def 0$\\
   $\Assign$   &  Положим равным &    $c\Assign (a+b)/2$ \\
   $=$   &  Равно &    $1= 1, 2\times 2 =4$ \\
\hline
\end{tabular}

\vspace{0.5\baselineskip}
{\bf Числовые множества}

\begin{tabular}{l|l|l}
\hline
 Обозначение &  Название & Примеры \\
\hline
   $\mathbb{N}$    &  Натуральные числа &  $1,2,3$\\
   $\mathbb{Z}$    &  Целые числа &  $0,1,-1,2,-2$\\
   $\mathbb{Q}$    &  Рациональные числа &
   $\displaystyle\frac{1}{2},\frac{1}{3}, \frac{553}{113}, 2.718281828$\\
   $\mathbb{R}$    &  Вещественные числа &  $1, 1.5, \sqrt{2},\pi, e$ \\
\hline
\end{tabular}
\end{frame}


\begin{frame}
\frametitle{Стандартные математические символы (2/6)}

{\bf Совокупности}

\begin{tabular}{p{25mm}|p{75mm}|p{40mm}}
\hline
 Обозначение &  Применение  & Примеры \\
\hline
   $\{\Tuple{a}{1}{n}\}$
   &  Неупорядоченное множество различных элементов
   &  $\{1,2,3\}$\\
   $(\Tuple{a}{1}{n})$
   &  Упорядоченная последовательность однородных элементов
   &  $(0,1,0,1)$\\
   $\langle A;B,C \rangle$
   &   Упорядоченная последовательность разнородных элементов
   &  $\langle \mathbb{R};+,\ast \rangle$\\
   $\langle \Tuple{a}{1}{n} \rangle$
   &  Упорядоченная последовательность составных элементов
   &  $\langle a\to 01, b\to 10 \rangle$\\
\hline
\end{tabular}

\end{frame}

\begin{frame}
\frametitle{Стандартные математические символы (3/6)}
{\bf Операции с множествами}

\begin{tabular}{p{22mm}|p{70mm}|p{60mm}}
\hline
 Обозначение &  Прочтение  & Примеры \\
\hline
   $a\in A$
   &  {\raggedright Элемент $a$ принадлежит множеству $A$}
   &  $1\in\{1,2,3\}$ \\
   $a\not\in A$
   &  Элемент $a$ не принадлежит множеству~$A$
   &  $4\not\in\{1,2,3\}$, $\sqrt{2}\not\in\mathbb{Q}$ \\
   $A\subset B$
   & Множество $A$ является подмножеством множества $B$
   &  $\{2,3\}\subset\{1,2,3\}$ \\
   $A\cup B$
   &   Объединение множеств $A$ и $B$
   &  $\{1,2\}\cup \{2,3\}=\{1,2,3\}$ \\
   $A\cap B$
   &   Пересечение множеств $A$ и $B$
   &  $\{1,2\}\cap \{2,3\}=\{2\}$ \\
   $A\setminus B$
   &   Разность множеств $A$ и $B$
   &  $\{1,2\}\setminus \{2,3\}=\{1\}$ \\
   $A\vartriangle B$
   & Симметрическая разность множеств $A$~и~$B$
   &  $\{1,2\}\vartriangle \{2,3\}=\{1,3\}$ \\
   $A\times B$
   &   Прямое произведение множеств $A$ и $B$  
   &  $\{1,2\}\times \{2,3\}=$\hfill $ = \{(1,2),(1,3),(2,2),(2,3)\}$ \\
   $\emptyset$
   &    Пустое множество
   &  $\{\ \}$ \\
   $|A|$
   &    Мощность множества $A$
   &  $|\{1,2\}| = 2$ \\
\hline
\end{tabular}

\end{frame}


\begin{frame}
\frametitle{Стандартные математические символы (4/6)}

\vspace{-2pt}
{\bf Логические обозначения}

\begin{tabular}{p{25mm}|p{50mm}|p{55mm}}
\hline
 Обозначение &  Название  & Прочтение \\
\hline
   $\Not P$
   &  Отрицание
   &  Не $P$ \\
   $P\Or Q$
   &  Дизъюнкция
   &  $P$ или $Q$ \\
   $P\And Q$
   &  Конъюнкция
   &  $P$ и $Q$ \\
   $P\Impl Q$
   &  Импликация
   &  Если $P$, то $Q$ \\
   $\A{x}{P(x)}$
   &  Квантор всеобщности
   &  Для всех $x$ выполнено $P(x)$ \\
   $\E{x}{P(x)}$
   &  Квантор существования
   &  Существует $x$, такой, что $P(x)$ \\
\hline
\end{tabular}


{\bf Стандартные числовые функции}

\begin{tabular}{p{25mm}|p{50mm}|p{70mm}}
\hline
 Обозначение &  Прочтение & Примечание  \\
\hline
   $\lfloor x \rfloor$
   &  Целая часть числа
   & $\A{x\in \Bbb{R}}{x>0 \Impl \lfloor x \rfloor \le x}$
   \\
   $\log_a x $
   &  Логарифм по основанию $a$ 
   & $\log x \Assign \log_2 x$; $\ln x \Assign \log_e x$
   \\
   $m \DIV n $
   &  Неполное частное 
   & $m = (m \DIV n) * n + m \MOD n$
   \\
   $m \MOD n $
   &  Остаток от деления 
   & $m = (m \DIV n) * n + m \MOD n$
   \\
   $ n! $
   &  Факториал 
   & $n! \Def \prod\limits_{i=1}^n i\ , 0!\Def 1$
   \\
\hline
\end{tabular}

\end{frame}


\begin{frame}
\frametitle{Стандартные математические символы (5/6)}

\vspace{-2pt}
{\bf Отношения}

\begin{tabular}{p{25mm}|p{85mm}}
\hline
 Обозначение &  Название  \\
\hline
   $R\circ S$
   &  Композиция отношений \\
   $R^n$
   &  Степень отношения \\
%   $\boxed{R}$
%   &  Матрица отношения \\
   $\equiv$
   &  Отношение эквивалентности \\
   $\prec$
   &  Отношение частичного порядка \\
   $<$
   &  Отношение строгого линейного порядка \\
   $\leq$
   &  Отношение нестрогого линейного порядка \\
\hline
\end{tabular}

\medskip
{\bf Групповые операции, обозначаемые знаками}

\begin{tabular}{p{20mm}|p{65mm}|p{50mm}}
\hline
 Операция &  Прочтение & Примечание  \\
\hline
   $\sum_{i=1}^{n}{a_i}$
   &  Сумма элементов $\Tuple{a}{1}{n}$
   & $\sum_{i=1}^{n}{a_i} = a_1+\ldots+a_n$\\
   $\prod_{i=1}^{n}{a_i}$
   &  Произведение элементов  $\Tuple{a}{1}{n}$
   & $\prod_{i=1}^{n}{a_i} = a_1\times\ldots\times a_n$\\
   $\BigOr_{i=1}^{n}{a_i}$
   &  Дизъюнкция элементов $\Tuple{a}{1}{n}$
   & $\BigOr_{i=1}^{n}{a_i} = a_1\Or\ldots\Or a_n$\\
   $\bigwedge_{i=1}^{n}{a_i}$
   &  Конъюнкция элементов  $\Tuple{a}{1}{n}$
   & $\bigwedge_{i=1}^{n}{a_i} = a_1\And\ldots\And a_n$\\   
\hline
\end{tabular}

\end{frame}


\begin{frame}
\frametitle{Стандартные математические символы (6/6)}
{\bf Групповые операции, обозначаемые именами}

\begin{tabular}{p{28mm}|p{70mm}|p{40mm}}
\hline
 Операция &  Прочтение & Примечание  \\
\hline
   $\min (\Tuple{a}{1}{n})$
   &  Минимальный из элементов $\Tuple{a}{1}{n}$
   & $\min(a,b)\leq a$\\
   $\max  (\Tuple{a}{1}{n})$
   &  Максимальный из элементов $\Tuple{a}{1}{n}$
   & $\max(a,b)\geq  b$\\
\hline
\end{tabular}

\smallskip
{\bf Функции и отображения}

\begin{tabular}{p{25mm}|p{50mm}|p{60mm}}
\hline
 Обозначение &  Прочтение & Примечание  \\
\hline
   $\Map{f}{A}{B}$
   &  Функция $f$ из $A$ в $B$
   & $f\subset A\times B$\\
   $b=f(a)$
   &  $b$ является значением функции $f$ для аргумента $a$
   & $a\mapsto b$\\
   $f\circ g$
   &  Суперпозиция функций $f$ и $g$
   & $(f\circ g)(x)= f(g(x))$ \\
   $f^{-1}$
   &  Обратная функция к $f$
   & $\Map{f^{-1}}{B}{A}$\\
   $\Dom f$
   &  Множество определения функции $\Map{f}{A}{B}$
   & $\forall a\!\in\!\Dom f (\exists b\!\in\! B (b\!=\!f(a)))$\\
   $\Im f$
   &  Множество значений функции $\Map{f}{A}{B}$
   & $\forall b\!\in\!\Im f (\exists a\!\in\! A (b\!=\!f(a)))$\\
\hline
\end{tabular}

\end{frame}



\subsection{Выделения в тексте}


\begin{frame}
\frametitle{Выделения в тексте (1/4)}


В учебнике и в слайдах имеются следующие основные виды текстов: 
\em{ определения}, \em{ теоремы}, \em{ леммы}
и \em{ следствия}, \em{ доказательства} и \em{ обоснования}, 
\em{ замечания}, \em{ отступления}, 
\em{ алгоритмы} и
\em{ примеры}. 
Фактически, обычный связующий текст сведен к минимуму в целях
сокращения объёма.

\smallskip
 \em{ Определения\/} не выделяются никакими 
специальными словами,
поскольку составляют львиную долю основного текста книги. Вместо
этого {\it курсивом\/}, а в слайдах ещё и 
{\color{red} красным цветом}, выделяются
\odmkey{определяющие~вхождения}{Вхождение
(occurrence)!определяющее (defining)} терминов, а текст,
соседствующий с термином, и есть определение. 

\smallskip
\em{ Формулировки\/} теорем, лемм и следствий в соответствии с
общепринятой в математической литературе практикой выделены {\it
курсивом\/}. При этом формулировке предшествует соответствующее
ключевое слово: <<теорема>>, <<лемма>> или <<следствие>>. Ключевые слова выделяются 
{\color{brown}коричневым цветом}. Как
правило, утверждения не нумеруются, за исключением случаев
вхождения нескольких однородных утверждений в один параграф. Для
ссылок на утверждения используются либо номера параграфов, в
которых утверж\-дения сформулированы, либо собственные имена
утверждений. 

\end{frame}

\begin{frame}
\frametitle{Выделения в тексте (2/4)}
Дело в том, что в~данном курсе теоремами оформлены
как простые утверждения, являю\-щиеся непосредственными
следствиями определений, так и глубокие факты, действительно
заслуживающие статуса теоремы. В последнем случае приводится
собственное имя (название теоремы), которое либо выносится в
название параграфа (или раздела), либо указывается в квадратных скобках
после слова <<теорема>>.

\smallskip
\em{ Доказательства\/} относятся к предшествующим утверждениям, а
\em{ обоснования\/}~---~к предшествующим алгоритмам. В курсе
нет недоказанных утверж\-де\-ний и приведённых без
достаточного обоснования алгоритмов. Из этого правила имеется
несколько исключений, то есть утверждений, доказательства которых
не приводятся, поскольку автор считает их технически слишком
сложными для выбранного уровня изложения. Отсутствие 
сложного доказательства всегда явно отмечается в тексте. Иногда же
доказательство опускается, потому что утверждение легко может быть
обосновано студентами самостоятельно. Такие случаи не отмечаются в
тексте. Доказательства и обоснования начинаются с соответствующего
ключевого слова и
заканчиваются специальным значком \qedsymbol. 
При этом, если доказательство (или обоснование) 
не помещается на одной странице,
то оно прерывается
 значком \qedsymbol, 
а на следующей странице возобновляется с
пометкой <<(продолжение)>>.

\end{frame}




\begin{frame}
\frametitle{Выделения в тексте (3/4)}
Если формулировка
теоремы содержит несколько утверждений или если доказательство
состоит из нескольких шагов (частей), 
то доказательство делится на абзацы, а
в начале каждого абзаца указывается
 {\sffamily\small\color{brown}[в~скобках]},
что именно в нём доказывается.

\smallskip
\em{ Замечания\/} и \em{ отступления\/} также начинаются с
соответствующего ключевого слова: <<замечание>> или
<<отступление>>. Назначение этих абзацев различно. Заме\-чание
сообщает некоторые специальные или дополнительные сведения, прямо
относящиеся к основному материалу курса. Отступление не связано
непосредственно с основным материалом, его назначение~---
расширить кругозор читателя и показать связь основного материала с
вопросами, лежащими за пределами курса.

%\fussy

\begin{remark}
Иногда {\it курсив} используется не для выделения
определяющих вхождений терминов и формулировок утверждений, а для
того, чтобы сделать \em{ эмфатическое} ударение на каком-то слове в
тексте. Эмфатические ударения, в отличие от определяющих вхождений, в слайдах выделяются {\color{blue} синим цветом}. 
\end{remark}

\begin{digress}
Использование отступлений необходимо, в противном случае
у читателя может сложиться ошибочное впечатление о
замкнутости изучаемого предмета и его оторванности
от других областей знаний.
\end{digress}
\end{frame}




\begin{frame}
\frametitle{Выделения в тексте (4/4)}
%\vspace{-3pt}
\em{ Алгоритмы\/}, как уже было сказано, записаны на псевдокоде. 
Как правило, перед текстом
алгоритма на естественном языке указываются его назначение, а
также входные и выходные данные. После
алгоритма приводятся его обоснование и иногда пример протокола
выполнения.

\vspace{-2pt}
\begin{algorithmic}
\REQUIRE текст алгоритма.
\ENSURE 1, если алгоритм понят, 0~--- если не понят.
\WHILE{текст алгоритма на экране}
   \STATE {смотреть на текст алгоритма}
   \COMMENT {\color{gray}  по возможности вдумчиво}
   \IF{понятно}
   \STATE {\textbf{return} 1}
   \ENDIF    
\ENDWHILE
\STATE \textbf{return} 0
\end{algorithmic}

\vspace{-2pt}
\begin{examples} Как правило, примеры приводятся непосредственно вслед за
определением понятия, поэтому не используется никаких связующих
слов, поясняющих, к чему относятся примеры. 
\end{examples}
\end{frame}

\subsection{Онтология дискретной математики}
\begin{frame}
\frametitle{Онтология дискретной математики (1/8)}
\begin{center}
\begin{columns}[t] 
\begin{column}[T]{8.5cm}
Новинкой реализации курса в 2018 году явилось применение понятия \odmkey{онтологии}{Онтология (ontology)}, то есть формального описания некоторой области знаний с помощью концептуальной схемы, точнее
говоря \odmkey{диаграммы}{Диаграмма (diagram)}. Применение диаграмм
оказало положительный эффект на понимание
материала студентами,
поэтому в реализации курса в 2019 году применение диаграмм было несколько расширено. 
Мы описываем онтологии в графической форме с помощью \em{ унифицированного языка моделирования UML}. 

Прежде всего, в онтологию входят \em{ понятия}. Например, само слово <<онтология>> --- это пример понятия.
Понятия изображаются прямоугольниками, и им даются имена.
\end{column}
\begin{column}[T]{6cm} 
\includegraphics[width=5.5cm]{0_1.png}
\end{column}
\end{columns}
\end{center}

\end{frame}

\begin{frame}
\frametitle{Онтология дискретной математики (2/8)}
\begin{center}
\begin{columns}[t] 
\begin{column}[T]{8.5cm}
Далее в онтологии указываются различные отношения между понятиями.

\smallskip
Очень часто встречаются иерархические отношения типа \em{ часть--целое}. Например, онтология состоит из многих элементов.

\smallskip
Понятия могут иметь \em{ атрибуты}. Например, у онтологии есть автор, который ее составил, а у элементов есть имена.  

\smallskip
Понятия могут быть \em{ обобщающими} и \em{ специализированными}. Например, мы рассматриваем в онтологиях только два частных случая элементов: сущности и отношения. 
\end{column}
\begin{column}[T]{6cm} 
\includegraphics[width=5.5cm]{0_4.png}
\end{column}
\end{columns}
\end{center}

\end{frame}
\begin{frame}
\frametitle{Онтология дискретной математики (3/8)}
\begin{center}
\begin{columns}[t] 
\begin{column}[T]{8.5cm}
Между понятиями могут быть и любые другие отношения, которым обычно даются имена, чтобы было понятно, что это за отношение. 

\smallskip
Можно указывать направление чтения имени отношения, и тогда диаграммы можно читать вслух. 
Например: \em{ бинарное отношение между сущностями}.  

\smallskip
Показывая различные отношения между понятиями в онтологии, мы получаем возможность наглядно, без лишних слов, показать целостную картину. 
\end{column}
\begin{column}[T]{6cm} 
\includegraphics[width=5.5cm]{0_6.png}
\end{column}
\end{columns}
\end{center}

\end{frame}
\begin{frame}
\frametitle{Онтология дискретной математики (4/8)}
\begin{center}
\begin{columns}[t] 
\begin{column}[T]{8.5cm}

Еще одно важное обозначение --- это \em{ зависимость} между понятиями, которая возникает, когда зависимое понятие определяется через независимое понятие. Например, принимая определение: 
<<\odmkey{линия}{Линия (line)} имеет длину, но не имеет площади>>, тем самым мы постулируем, что понятия длины и площади первичны, а понятие линии можно определить через них. Разумеется, в онтологии не должно быть циклов из пунктирных стрелочек. Если такие циклы появляются, то это порочный круг из определений.

%\smallskip
%.
\end{column}
\begin{column}[T]{6cm} 
\includegraphics[width=5.5cm]{0_7.png}
\end{column}
\end{columns}
\end{center}
\end{frame}

\begin{frame}
\frametitle{Онтология дискретной математики (5/8)}
\begin{center}
\begin{columns}[t] 
\begin{column}[T]{8.5cm}

Понятия могут быть не только пассивными сущностями,
то есть множествами объектов,
но и активными сущностями, то есть процессами обработки объектов. Процесс мы изображаем 
в виде сущности со
стереотипом <<activity>>.   

\smallskip
Наконец, последнее  обозначение --- это \em{ поток информации} (пунктирная стрелка со стереотипом <<flow>>). Мы используем поток информации, чтобы показать, какие сущности являются аргументами (то есть, входами) и результатами
(то есть, выходами) процесса.    

\smallskip
Например, найти площадь любой фигуры 
можно с помощью процесса 
численного интегрирования.    

\end{column}
\begin{column}[T]{6cm} 
\includegraphics[width=5.5cm]{0_8.png}
\end{column}
\end{columns}
\end{center}

\end{frame}

\begin{frame}
\frametitle{Онтология дискретной математики (6/8)}
\begin{center}
\begin{columns}[t] 
\begin{column}[T]{7.5cm}
В 2020 году 
концепция использования 
онтологий в образовании получила дальнейшее
развитие.

\smallskip
Новая, можно сказать революционная идея
состоит в том, чтобы структурировать онтологии 
не только по группам понятий,
но и по аспектам, или же граням,
отражаемым на диаграммах.
Отсюда возникает несколько типов диаграмм.

\smallskip
Предлагается использовать три грани онтологии.

\smallskip
На первой грани с помощью отношений 
обобщения, агрегации, композиции и 
ассоциации
показывается общая иерархия понятий 
и их составных частей.  
\end{column}
\begin{column}[T]{7cm} 
\includegraphics[width=7cm]{0_A.jpg}
\end{column}
\end{columns}
\end{center}

\end{frame}
\begin{frame}
\frametitle{Онтология дискретной математики (7/8)}
\begin{center}
\begin{columns}[t] 
\begin{column}[T]{8cm}
На второй грани предлагается показывать 
не структурные, а причинно-следственные и иные 
\em{ зависимости} между понятиями.

\smallskip
Например, если некоторое понятие определяется через
другие понятия, и тем cамым зависит от них,
проводится стрелка зависимости 
(без стереотипа) от определяемого понятия
к определяющему понятию.

\smallskip
Если некоторое конкретное понятие
является одним из возможных конкретных воплощений
некоторого более абстрактного понятия,
то от конкретного понятия
проводится стрелка зависимости к абстрактному
понятию со стереотипом <<manifest>>.

\end{column}
\begin{column}[T]{6.5cm} 
\includegraphics[width=6.5cm]{0_B.jpg}
\end{column}
\end{columns}
\end{center}

\end{frame}


\begin{frame}
\frametitle{Онтология дискретной математики (8/8)}
\begin{center}
\begin{columns}[t] 
\begin{column}[T]{6.5cm}
\smallskip
Аналогично 
проводится стрелка зависимости от экземпляра 
класса к классу
 со стереотипом <<instance of>>.
 
 \begin{remark}
 Эти и другие используемые стереотипы
 являются стандартными стереотипами 
 языка UML~2.
 \end{remark}

На третьей грани показываются
алгоритмические аспекты, то есть
аргументы и результаты действий в нотации диаграмм деятельности UML.

\smallskip
Разнесение онтологии по трём граням 
резко повышает наглядность каждой
из диаграмм и информативность
онтологии в целом.
\end{column}
\begin{column}[T]{8cm} 
\includegraphics[width=8cm]{0_C.jpg}
\end{column}
\end{columns}
\end{center}

\end{frame}




\subsection{Подведение итогов}
\begin{frame}
\frametitle{Подведение итогов (1/6)}

В курсе традиционно используются три вида занятий: 
\em{ лекции}, \em{ упражнения} и 
\em{ лабораторные работы}. 
С 2022 года в программу курса 
добавлена компьютерная система 
\em{ автоматической} проверки
правильности решений учебных задач,
что позволяет значительно интенсифицировать
учебный процесс.  
Объём занятий каждого вида 
определяется отдельно для каждой реализации курса в зависимости от специфических факторов семестра и университета.

\smallskip
Порядок и условия получения \em{ зачётов} как по упражнениям, так и по лабораторным определяют преподаватели,
ведущие упражнения и лабораторные. Получение зачётов не
является \em{ предварительным} условием для получения экзаменационной оценки.

\smallskip
На лекциях в семестре ведётся учёт посещаемости, успеваемости, 
решённых задач
и других измеримых показателей (\em{портфолио}). Для этого применяются
подписанные \em{ бумажные листочки}, которые надобно сдавать
в конце лекции, если лекция проводится очно,
 или \em{ встроенный чат},
если лекция проводятся средствами телекоммуникаций.
Учёт решённых задач в системе автоматической проверки
ведётся средствами этой системы. 

    
Текущие значения показателей 
(то есть рейтинг по портфолио) доступны всем участникам 
учебного процесса,
включая студентов, преподавателей и руководство 
высшей школы.

\end{frame}

\begin{frame}
\frametitle{Подведение итогов (2/6)}

Предусматривается две различных формы экзамена:
традиционный \em{ экзамен в сессию} и 
\em{ экзамен по  портфолио},
то есть по сумме результатов, показанных
на протяжении всего семестра. 

\smallskip
\em{ Экзамен в сессию} проводится  по расписанию деканата. 
Если экзамен не сдан c первой попытки,
то вторая попытка сдать экзамен предоставляется 
на дополнительной сессии, а третья и последняя попытка~--- 
на заседании предметной комиссии.
Пересдача экзамена <<в коридоре на бегу>> не предусматривается. 

\smallskip
На традиционном экзамене необходимо \em{ письменно} ответить на два или три вопроса,
каждый из которых является одним разделом курса.
Ответ должен иметь форму развёрнутого плана.
Порядок изложения и буквальная точность формулировок несущественны.
Важно знать и понимать, о чём идёт речь в вопросе, и записать
\em{ исчерпывающий план ответа}.
План должен содержать, по меньшей мере, заголовки параграфов, и может содержать список определяемых в вопросе (в разделе) понятий, 
формулировки доказываемых утверждений и спецификации алгоритмов.
При предъявлении плана ответа на вопрос, студент должен показать способность устно конкретизировать план в любом месте.
%: дать развёрнутое 
%определение, привести пример, доказать утверждение, обосновать алгоритм.  

\end{frame}

\begin{frame}
\frametitle{Подведение итогов (3/6)}

\em{ Оценка за ответ} определяется полнотой 
материала, включённого студентом в план ответа.
Допускается <<забыть>> не более одного параграфа в разделе (вопросе).
Параграфы, помеченные знаком \RS, опускать не разрешается. Если к концу периода времени, отведённого на подготовку, не написаны полные планы ответа
по всем предложенным вопросам, то ответ в целом не принимается и экзамен считается проваленным. 

\smallskip
Часть материала курса, возможно, \em{ не выносится на экзамен}
в конкретной реализации курса, но оставлена в слайдах для перекрёстных ссылок и ради совпадения нумерации с учебником.
Каждый такой случай особо оговаривается на лекциях.
Заголовки опущенных параграфов и разделов не нужно включать в план ответа.

\smallskip
Выбор экзаменационного билета \em{ не является случайным}.
Билеты выбираются индивидуально,  
<<в открытую>>, по очереди, строго по рейтингу, установленному по результатам работы 
студентов в семестре.
Активные студенты с высоким рейтингом выбирают
и отвечают те вопросы, которые хотят, 
а плохо успевающим студентам достаются оставшиеся билеты.
Как именно вопросы группируюся по билетам, является секретом, который открывается непосредственно на экзамене.
\end{frame}

\begin{frame}
\frametitle{Подведение итогов (4/6)}

\em{ Оценка за экзамен} определяется как округлённое среднее оценок за ответы на отдельные вопросы плюс бонус за портфолио в размере одного процента рейтинга, при этом для вопросов используется пятибалльная система оценок.
\begin{enumerate}
\item \em{ Посредственно, E = 1} --- план написан, но студент к плану ничего не может добавить. 
\item \em{ Удовлетворительно, D = 2} --- студент знает определения основных понятий, может привести примеры. 
\item \em{ Хорошо, C = 3} --- студент может без ошибок сформулировать теоремы и входы--выходы алгоритмов.
\item \em{ Очень хорошо, B = 4} --- студент может доказывать теоремы и объяснять алгоритмы. 
\item \em{ Отлично, A = 5} --- студент знает всё или <<почти всё>>.
\end{enumerate}
Округлённая оценка прямо транслируется в итоговую 
оценку: $5 \Impl$ отлично, $4 \Impl$ хорошо,
 $3 \Impl$ удовлетворительно, остальное неудовлетворительно, экзамен не сдан. 
\begin{examples}
$(1+3+5)/3+0,3=$ удовлетворительно, $(3+4+5)/3+0,7=$ отлично, 
$(1+1+5)/3 +0,1=$ неудовлетворительно.  
\end{examples} 

\em{ На традиционном экзамене запрещается пользоваться любыми источниками,
кроме собственных знаний.} 
\end{frame}

\begin{frame}
\frametitle{Подведение итогов (5/6)}

В случае дистанционного обучения экзамен, как правило,
проводится \em{ по портфолио}, то есть 
\em{ по совокупности содеянного} в семестре 
(<<автоматом>>). 
Для этого достаточно за весь семестр набрать определённое число {\bf баллов}.
Баллы начисляются \em{ субъективно} по отчуждаемым артефактам.   
\begin{enumerate}
\item \em{ Посещение лекции}, отмеченное сданным <<листочком>>~--- $1$.
\item \em{ Правильный вопрос}, заданный на лекции, или иное полезное действие~--- $1..3$.
\item \em{ Исправленная ошибка} в слайдах
~--- $1..5$.
\item \em{ Прослушанный МООК} по теме курса,~--- $5..10$ 
\item \em{ Чудесный подвиг} в дискретной математике~--- $10$ и более.
\end{enumerate}

Вес баллов определяется динамически для 
каждого студенческого потока. 
В конце семестра, когда лекции закончены, вычисляются среднее
$x$ и стандартное отклонение $y$ суммы набранных баллов по всему потоку.
По умолчанию число баллов более $x+y$ ранжируется
как оценка
<<отлично>>, число баллов от $x$ до $x+y$ ранжируется
как оценка 
<<хорошо>>,
число баллов от $x-y$ до $x$ ранжируется
как оценка 
<<удовлетворительно>>, 
а число баллов менее $x-y$ ранжируется
как оценка
<<неудовлетворительно>>.
В эти формулы могут быть добавлены поправочные коэффициенты по усмотрению преподавателя.
\end{frame}

\begin{frame}
\frametitle{Подведение итогов (6/6)}  

Начиная с весеннего семестра 2022 года отдельно учитываются успехи в работе 
с программами автоматической проверки решения задач.
При этом баллы, автоматически выставленные программой, 
суммируются и ранжируются по тому же принципу, как и 
баллы, выставленные преподавателем.
Итоговая оценка складывается из двух оценок:
оценки, выставленной преподавателем,
и оценки, выставленной программой.
Выбор способа совместного учёта двух оценок 
остаётся за преподавателем.

\smallskip
\begin{remark}
Для дальнейшего 
использования программ автоматического
оценивания решений учебных задач очень
важно экспериментально подтвердить
высокую корреляцию оценок, выставленных вручную и автоматически.  
\end{remark}


\smallskip
В случае экзамена <<по портфолио>> возможно 
<<добирать>> баллы 
в период подготовки к экзамену
за счёт работы со слайдами и за
счёт <<чудесных подвигов>>.
Если студент не набрал необходимых 
баллов к сроку сдачи экзамена,
то аттестация по портфолио 
проведена быть не может~--- студент должен сдавать
традиционный экзамен. 


Возможно совмещение двух форм экзамена: студенты,
имеющие наивысшие рейтинги к моменту экзамена 
получают  <<автоматом>> оценку <<отлично>> по портфолио, 
увеличивая шансы остальных на получение удачного билета. 
\end{frame}

% ===== tex/DM2022_1.tex =====

\begin{frame}
\section{1. Множества и отношения}
\frametitle{1. Множества и отношения}
{\large
Понятия <<множество>>, <<отношение>>, <<функция>> и близкие к ним
составляют основной словарь дискретной (равно как и
<<непрерывной>>) математики. Именно эти базовые понятия
рассматриваются в первой главе, тем самым закладывается
необходимая основа для дальнейших построений. Особенность
изложения состоит в том, что здесь рассматриваются почти
исключительно \em{конечные\/} множества, а~тонкие и сложные
вопросы, связанные с рассмотрением бесконечных множеств,
излагаются с <<программистской>> точки зрения. С другой стороны,
значительное внимание уделяется <<представлению>> множеств в
программах. Эти вопросы не имеют прямого отношения к собственно
теории множеств в её классическом виде, но очень важны для практикующего программиста.
}
\end{frame}

\begin{frame}
\subsection{1.1. Множества}
\label{1.1.}
\frametitle{1.1. Множества (1/2)}
{\large\linespread{1.5}\selectfont 

При построении доступной для рационального анализа картины мира
часто используется термин <<объект>> для обозначения некой
сущности, отделимой от остальных. Выделение объектов~--- это 
не более чем произвольный акт нашего сознания. В одной и той же ситуации
объекты могут быть выделены по-разному, в зависимости от точки зрения,
целей анализа и других обстоятельств. Но
как бы то ни было, выделение объектов и их совокупностей~---
естественный (или даже единственно возможный) способ организации
нашего мышления, поэтому неудивительно, что он лежит в основе
главного инструмента описания точного знания~--- математики.
}
\end{frame}

\begin{frame}
\frametitle{1.1. Множества (2/2)}

\begin{tabular}{p{12mm}|p{40mm}|p{92mm}}
  № & Название параграфа
  & Ключевые термины и обозначения \\
  \hline
\DMref{1.1.1.}{1.1.1.} & Элементы и множества 
  & 
  $\in$~--- отношение принадлежности;
  $\emptyset$~--- пустое множество;
  $\Bbb{N}$~--- множество натуральных чисел;
  $\Bbb{Z}$~--- множество целых чисел;
  $\Bbb{R}$~--- множество вещественных чисел;
  семейство; класс 
   \\
\DMref{1.1.2.}{1.1.2.} & Задание множеств 
  & 
{
 $\{\Tuple{a}{1}{n}\}$~--- перечисление элементов;

 $\Set{x}{P(x)}$~--- характеристический предикат;

 $\Set{x}{\dots\textbf{yield}\ x}$~--- порождающая процедура;

 $m..n$~--- диапазон целых чисел;

  $\Bbb{N}_0$~--- натуральные числа с нулём
}
  \\
\DMref{1.1.3.}{1.1.3.} \RS & Парадокс Рассела &
Конструктивизм, аксиома регулярности\\
\DMref{1.1.4.}{1.1.4.}  & Мультимножества &
{$\widehat{X} = \left[ x_1^{a_1}, \dots, x_n^{a_n} \right] =\langle
x_1, \dots, x_1 ; \dots ; x_n,
\dots, x_n \rangle =$ 

$= \langle a_1(x_1), \dots, a_n(x_n)
\rangle$; носитель; показатель; состав} 
\\
\DMref{1.1.5.}{1.1.5.} & Конечные последовательности &
{Алфавит; слово; язык; префикс; постфикс; конкатенация}\\   
\end{tabular}
\end{frame}

\begin{frame}

\label{1.1.1.}\subsubsection{1.1.1. Элементы и множества }\frametitle{1.1.1. Элементы и множества (1/2)}

\odmkey{Множество}{Множество
(set)}~--- это любая определённая совокупность объектов. Объекты,
из которых составлено множество, называются его
\odmkey{элементами}{Элемент (element)!множества (of set)}.
Элементы множества различны и отличимы друг от друга. Как
множествам, так и элементам можно давать имена или присваивать
символьные обозначения.

\begin{examples}
Множество $S$ слайдов в данной презентации. Множество ${\Bbb N}$
\odmkey{натуральных}{Число (number)!натуральное (natural)} чисел
\{1, 2, 3,~\dots\}.
Множество ${\Bbb P}$ \odmkey{простых}{Число
(number)!простое (prime)} чисел \{2, 3, 5, 7, 11,~\dots\}.
Множество $\Bbb Z$ \odmkey{целых}{Число
(number)!целое (integer)} чисел \{\dots,~$-2$, $-1$, 0, 1,
2,~\dots\}. Множество ${\Bbb R}$ \odmkey{вещественных}{Число
(number)!вещественное (real)} чисел. Множество $A$ различных
символов на этом слайде.
\end{examples}


\smallskip
Отношение 
\odmkey{принадлежности}{Принадлежность (membership)} элемента $x$ множеству $M$
записывают так:
$ x\in M$ (пишут $x\notin M$, если элемент 
не принадлежит множеству).

\smallskip
Множество, не содержащее элементов, называется
\odmkey{пустым}{Множество (set)!пустое (empty)}
и обозначает так: $\emptyset$.

\smallskip
Множества как объекты могут быть элементами других множеств.
Множество, элементами которого являются множества, иногда называют
\odmkey{семейством}{Семейство (family)!множеств (of sets)}.
Обычно семейства обозначают <<рукописными>> буквами ($\cal M$),
множества~--- прописными буквами ($M$), а~элементы~--- строчными буквами ($m$).

\end{frame}


\begin{frame}
\frametitle{1.1.1. Элементы и множества (2/2)}

\begin{digress}
Понятия множества, элемента и принадлежности, которые на первый
взгляд представ\-ляются интуитивно ясными, при ближайшем
рассмотрении такую ясность утрачивают. Во-первых, даже возможность различить
элементы при более глубоком рассмотрении представляет некоторую
проблему. Например, символы <<а>> и <<{\it а}>>, которые
встречаются на предыдущем слайде,~--- это один элемент множества $A$
или два разных элемента? Или же, например, два вхождения символа
<<о>> в слово <<множество>>~--- графически они неразличимы
невооружённым глазом, но это символы букв, которые обозначают
звуки, а читаются эти две гласные по-разному: первая под
ударением, а вторая~--- безударная. 
Во-вторых, проблематична
возможность (без дополнительных усилий) указать, принадлежит ли
данный элемент данному множеству. Например, является ли число
869584769215374850678576637 простым?
\end{digress}

\medskip
Совокупность объектов, которая \em{не} является множеством,
называется \odmkey{классом}{Класс (class)} 
(со сравнительно меньшей ответственностью за
определённость, отличимость и бесповторность элементов).

\begin{example}
Буквы латинского алфавита образуют множество,
а совокупность слов английского языка лучше считать
классом, ввиду наличия диалектов, жаргона, омонимии  и так далее.  
\end{example}

\end{frame}

\begin{frame}
\label{1.1.2.}\subsubsection{1.1.2. Задание множеств }\frametitle{1.1.2. Задание множеств (1/2)}

Способы задания множеств:
\begin{tabbing}
\hspace{5mm}\= \odmkey{характеристический предикат}{Множество
(set)!заданное (defined by)!характеристическим предикатом
(characteristic predicate)}:$\quad$\= \kill \>\odmkey{перечисление
элементов}{Множество (set)!заданное (defined by)!перечислением
элементов (enumeration)}: \> 
$M\Assign\{a,b,c,\dots,z\};$\\[0pt]
\> \odmkey{характеристический предикат}{Множество (set)!заданное
(defined by)!характеристическим предикатом (characteristic
predicate)}: \>
$M\Assign\Set{x}{P(x)};$ \\[0pt]
\>\odmkey{порождающая процедура} {Множество (set)!заданное
(defined by)!порождающей процедурой (generating procedure)}: \>
$M\Assign\Set{x}{x\Assign f}$.

\end{tabbing}

\begin{examples}

$M_{9}\Assign \{1,2,3,4,5,6,7,8,9\}$.%

$M_{9}\Assign
\Set{n}{n\in\Bbb N\And n < 10}$.%

$M_{9}\Assign
\Set{n}{\textbf{for $n$ from $1$
to $9$ do yield $n$ end for}}$.

\end{examples}

При задании множеств перечислением обозначения элементов иногда
снабжают индексами и указывают множество, из которого берутся
индексы. В частности, запись $\left\{a_i\right\}_{i=1}^{k}$
означает то же, что $\{a_1,\dots,a_k\}$.

Знак многоточия ($\dots$), который употребляется
при задании множеств, является сокращённой формой записи, в
которой порождающая процедура для множества индексов считается
очевидной.

\begin{example}
$\{a_1, a_2, a_3,\dots\} = \left\{a_i\right\}_{i=1}^{\infty}$.
\end{example}
\end{frame}

\begin{frame}
\frametitle{1.1.2. Задание множеств (2/2)}

Множество целых чисел в диапазоне от $m$ до $n$:
$
m..n\Def\Set{k\in\Bbb Z}{m\leq k\And k\leq n}=$
$=\Set{k\in \Bbb Z}{\textbf{for $k$ from $m$ to $n$ do yield $k$ end~for}}.
$

\smallskip
Перечислением элементов можно задавать только
\odmkey{конечные}{Множество (set)!конечное (finite)} множества.
\odmkey{Бесконечные}{Множество (set)!бесконечное (infinite)}
множества задаются характеристическим предикатом или порождающей процедурой.

\begin{example}
${\Bbb N}\Def \Set{n}{n\Assign 0; \textbf{while true do $n\Assign n+1;$
yield $n$ end~while}}$.
\end{example}

\smallskip
Если  $x\notin A$,
то элемент $x$ можно
\odmkey{добавить}{Операция (operation)!добавления элемента
(of element adding)} в множество $A$:
$$A+x\Def\Set{y}{y\in A\Or y=x}.$$

Если 
$x\in A$, то элемент $x$ можно \odmkey{удалить}{Операция
(operation)!удаления элемента (of element removal)} из множества
$A$:
$$
A-x\Def\Set{y}{y\in A\And y\neq x}.
$$

\begin{examples}
$M_9 \Assign \emptyset +1+2+3+4+5+6+7+8+9$.


Множество натуральных чисел с нулём: $\Bbb{N}_0\Def \Bbb{N}+0$. 
\end{examples}

\end{frame}

\begin{frame}
\label{1.1.3.}\subsubsection{1.1.3. Парадокс Рассела }\frametitle{1.1.3. Парадокс Рассела (1/2)}
Возможность задания множеств характеристическим предикатом
зависит от предиката.
Использование некоторых предикатов для этой цели
может приводить к
противоречиям.
Например,
все рассмотренные в примерах множества не
содержат себя в качестве элемента,
то есть множества, не содержащие себя в качестве 
элемента, заведомо существуют. 
Рассмотрим множество $Y$ всех множеств,
не содержащих себя в качестве элемента:
$$Y\Def\Set{X}{X\notin X}.$$

Если множество $Y$ существует, то мы должны иметь
возможность ответить на следующий вопрос: $Y\!\in Y$? Пусть
$Y\!\in Y$, тогда $Y\!\not\in Y$. Пусть $Y\!\not\in Y$, тогда
$Y\!\in Y$. Имеем неустранимое логическое противоречие~--- \odmkey{парадокс}{Парадокс (paradox)}. 
Три способа избежать этого 
парадокса.

\begin{enumerate}

\item Ограничить используемые характеристические предикаты видом\\
$
P(x) = x\in A\And Q(x),
$
где $A$~--- известное, заведомо существующее множество
(\odmkey{универсум}{Универсум (universe)}). Обычно при этом
используют обозначение $\Set{x\in A}{Q(x)}$. Для $Y$ универсум не
указан, а потому $Y$ множеством не является. 
\end{enumerate}
\end{frame}

\begin{frame}
\frametitle{1.1.3. Парадокс Рассела (2/2)}
\begin{enumerate}
\item[2.] Теория типов.
Объекты имеют тип~0, множества элементов типа 0 имеют тип~1,
множества элементов типа 0 и 1~--- тип~2 и~т.~д. Класс $Y$ не имеет типа
и потому множеством не является. 
\item[3.] Явный запрет принадлежности
множества самому себе: $X\in X$~--- недопустимый предикат. При
аксиоматическом построении теории множеств соответствующая аксиома
называется \odmkey{аксиомой регулярности}{Аксиома
(axiom)!регулярности (of regularity)}.
\end{enumerate}

\medskip
При этом, однако, нет полной уверенности в том, что не обнаружатся другие противоречия. Полноценным выходом из ситуации являлось бы 
\em{аксиоматическое} построение теории множеств и доказательство непротиворечивости построенной формальной теории. 

\smallskip
Однако исследование парадоксов и~непротиворечивости систем аксиом является технически трудной задачей и~уводит далеко в сторону от программистской практики, для которой важнейшими являются конечные множества. Поэтому здесь мы ограничиваемся 
\em{наивной теорией множеств}.

\end{frame}

\label{1.1.4.}\subsubsection{1.1.4. Мультимножества}
\begin{frame}
\frametitle{1.1.4. Мультимножества}

\odmkey{Мультимножеством}{Мультимножество (multiset)}
$\widehat{X}$ 
называется совокупность элементов множества $X=\{\Tuple{x}{1}{n}\}$, в которую элемент $x_i$
входит $a_i$ раз, $a_i\ge 0$.
Мультимножество обозначается одним из следующих способов:
$$
\widehat{X} = \left[ x_1^{a_1}, \dots, x_n^{a_n} \right] = \langle
x_1, \dots, x_1 ; \dots ; x_n,
\dots, x_n \rangle = \langle a_1(x_1), \dots, a_n(x_n)
\rangle.
$$
\begin{example}
Пусть $X=\{a,b,c\}$. Тогда $\widehat{X} = \left[ a^0, b^3,
c^4\right] = \langle b,b,b;c,c,c,c \rangle = \langle 0(a),
3(b), 4(c) \rangle$.
\end{example}

\smallskip
Пусть {$\widehat{X}\!=\!\langle a_1(x_1), \dots, a_n(x_n)
\rangle$}~--- мультимножество над множеством
$X\!=\!\{\Tuple{x}{1}{n}\}$. Тогда число $a_i$ называется
\odmkey{показателем}{Показатель элемента (index of element)}
элемента $x_i$, множество $X$~--- \odmkey{носителем}{Носитель
(support)!мультимножества (of multiset)}, число $m=a_1+\ldots+a_n$~---
\odmkey{мощностью}{Мощность (cardinality)!мультимножества (of multiset)}, 
а  $\underline{\widehat{X}}=\Set{x_i\in
X}{a_i>0}$ называется \odmkey{составом}{Состав
(compound)!мультимножества (of multiset)} мультимножества
$\widehat{X}$.
\begin{example}
Пусть $\widehat{X} = \left[ a^0, b^3, c^4\right]$~---
мультимножество над множеством $X=\{a,b,c\}$. Тогда
$\underline{\widehat{X}}=\{b,c\}$.
\end{example}

\smallskip
Мультимножество $\widehat{X}=\langle a_1(x_1), \dots, a_n(x_n) \rangle$
над множеством $X=\{\Tuple{x}{1}{n}\}$ называется
\odmkey{индикатором}{Индикатор (indicator)}, если $\A{i\in 1..n}{a_i=0\Or a_i=1}$.

\end{frame}

\begin{frame}
\label{1.1.5.}\subsubsection{1.1.5. Конечные последовательности }\frametitle{1.1.5. Конечные последовательности (1/3)}
%\vspace{-4pt}
К понятию
мультимножества тесно примыкает широко используемое понятие 
(конечной) последовательности.

\smallskip
Пусть задано  множество $X =\{\Tuple{x}{1}{n}\}$.
Рассмотрим
мультимножество 
$$\widehat{X}=\langle a_1(x_1), \dots, a_n(x_n)\rangle $$ над $X$,
$\sum_{i=1}^n a_i=m$.
Элементы мультимножества можно \em{пронумеровать}, назначив
им номера из отрезка $1..m$. Это можно сделать многими 
различными способами.      
Мультимно\-жест\-во с конкретной нумерацией
называется \odmkey{последовательностью}{Последовательность (sequence)} длины $m$ 
над множеством элементов $X$. 
Другими словами, если задана последовательность, то всегда можно сказать,
какой элемент стоит на $i$-м месте в последовательности.  
Отсюда вытекает общепринятый способ
записи последовательностей: выписать элементы мультимножества слева направо
в порядке  возрастания номеров.

\begin{example}
Пусть $X=\{a, b\}$ --- исходное множество элементов, 
$\widehat{X}=\langle 2(a),2(b)\rangle$ --- мультимножество,
тогда  $aabb$, $abab$, $abba$, $baab$, $baba$, $bbaa$ --- все 6 возможных  
в данном случае последовательностей.   
\end{example} 

\smallskip   
\begin{remark}
Конечную последовательность чисел часто называют \odmkey{вектором}{Вектор (vector)}.
\end{remark}
\end{frame}

\begin{frame}
\frametitle{1.1.5. Конечные последовательности (2/3)}

\begin{digress}
Одномерный массив \textbf{$A : $ array $[1..n]$ of $X$} является
программным представлением последовательности 
длины $n$ над
множеством $X$. При этом множество $X$ называется 
\odmkey{типом}{Тип (type)} (элементов) массива,
а выражение $A[i]$ доставляет $i$-й элемент последовательности.
Выражение $A[i]$ вычисляется очень эффективно, при условии,
что все элементы массива (последовательности) занимают 
одинаковый объём памяти компьютера.     
\end{digress}

\medskip
Пусть теперь элементами исходного множества 
являются \odmkey{буквы}{Буква (letter)}, 
или \odmkey{символы}{Символ (symbol)}.
В таком случае множество букв называется \odmkey{алфавитом}{Алфавит (alphabet)}
и обычно обозначается $A$. 

Последовательность букв называется \odmkey{словом}{Слово (word)}, 
или \odmkey{цепочкой}{Цепочка (string)}, в данном алфавите $A$.
Одна буква может входить в слово несколько раз, 
каждый отдельный экземпляр буквы в слове 
называется \odmkey{вхождением}{Вхождение (occurrence)} буквы в слово. 

Если слово $\alpha = a_1\dots a_k, a_i\in A$, 
то количество букв в слове называется 
\odmkey{длиной}{Длина (length)!слова (of word)} слова: 
$|\alpha| = |a_1\dots a_k| = k$.
\odmkey{Пустое слово}{Слово (word)!пустое (empty)} обозначается 
$\varepsilon$; $|\varepsilon|\Def 0$. 

\smallskip
Множество слов в алфавите ${A}$, включая пустое слово, обозначается~${A^*}$.
Пустое слово не содержит букв алфавита, но считается словом:
$\varepsilon \not\in A$, но $\varepsilon\in A^*$. 
Если пустое слово исключается из рассмотрения, 
то множество слов в алфавите $A$ обозначается~$A^+$. 

\end{frame}


\begin{frame}
\frametitle{1.1.5. Конечные последовательности (3/3)}
Для слов определена операция 
\odmkey{конкатенации}{Операция (operation)!конкатенации (of concatenation)}, 
или \odmkey{сцепления}{Операция (operation)!сцепления (of concatenation)}.
Обычно операция конкатенации никак не обозначается, а сцепляемые слова 
просто записываются подряд. 

\begin{example}
Конкатенация слов <<пара>> и <<план>> образует слово <<параплан>>.
\end{example}

Если $\alpha=\alpha_1\alpha_2$, то $\alpha_1$ 
называется \odmkey{началом}{Слово (word)!начало (prefix)}, 
или \odmkey{префиксом}{Префикс (prefix)}, 
слова $\alpha$, а $\alpha_2$~--- \odmkey{окончанием}{Слово (word)!окончание (postfix)}, или \odmkey{постфиксом}{Постфикс (postfix)}, слова $\alpha$. 

\begin{example}
Слова <<д>> и <<до>> являются префиксами, а слова <<ом>> и <<м>> являются постфиксами слова <<дом>>.
\end{example}

\begin{remark}
Каждое слово является своим постфиксом и префиксом.
\end{remark}

%\medskip
Слова $\alpha$ и $\beta$ могут \odmkey{пересекаться}
{Операция (operation)!пересечения (of intersection)}.
При пересечении либо одно слово является частью другого, 
$\alpha=\gamma\beta\delta$,
либо постфикс одного слова является префиксом другого, 
$\alpha=\delta\gamma, \beta=\gamma\zeta$.
Для пары слов $\alpha$ и $\beta$ тройка слов 
$\gamma$, $\delta$, $\zeta$ определяется однозначно,
если выбирать наибольшее пересечение.

\begin{example}
Пусть $\alpha=abb$, $\beta=bbc$. Тогда можно
взять $\gamma\Assign b$, $\delta\Assign ab$, $\zeta\Assign bc$.
Однако наибольшее пересечение даёт единственная тройка 
$\gamma\Assign bb$, $\delta\Assign a$, $\zeta\Assign c$.
\end{example} 

\medskip
Произвольное множество $L$ слов в~некотором алфавите называется 
\odmkey{языком}{Язык (language)} в этом алфавите;~$L\subset A^*$. 


\end{frame}

\begin{frame}
\frametitle{Онтология наивной теории множеств}
\begin{center}
\includegraphics[height=7.5cm]{1_1_A.jpg}
\end{center}
\end{frame}


\begin{frame}
\frametitle{Онтология наивной теории множеств}
\begin{center}
\includegraphics[height=7.5cm]{1_1_B.jpg}
\end{center}
\end{frame}

\begin{frame}
\frametitle{Онтология наивной теории множеств}
\begin{center}
\includegraphics[width=14.5cm]{1_1_C.jpg}
\end{center}
\end{frame}

\begin{frame}
\frametitle{1.2. Алгебра подмножеств (1/3)}
\subsection{1.2. Алгебра подмножеств}
\label{1.2.}
Самого по себе понятия множества ещё недостаточно~--- нужно определить
способы конструирования новых множеств из уже имеющихся,
то есть определить 
\odmkey{операции}{Операция (operation)} над множествами.

\medskip
Рассматриваемый совместно набор операций, определённых на некотором множестве,
образует \odmkey{алгебру}{Алгебра (algebra)}. При этом рассматриваемое множество объектов называют 
\odmkey{носителем}{Носитель (support)} алгебры.

\begin{example}
\odmkey{Арифметика}{Арифметика (arithmetic)}, изучаемая в начальной школе, --- это алгебра,
включающая операции сложения ($+$), умножения ($\times$), вычитания
($-$) и деления с остатком ($/$),
причём носителем является множество натуральных чисел и ноль. 
\end{example} 

\medskip
Наряду с операциями, результаты которых принадлежат носителю,
часто бывают необходимы и такие операции,
результаты которых не принадлежат носителю.    

\begin{examples}
В арифметике, наряду с арифметическими операциями,
используются \odmkey{операции сравнения}{Операция (operation)!сравнения (of comparison)}: больше ($>$), меньше ($<$) и другие. Аргументами этих операций являются натуральные числа, а результатами --- истинностные значения.
Операция вычитания может дать отрицательное число,
которое не считается натуральным.
\end{examples}

\end{frame}

\begin{frame}
\frametitle{1.2. Алгебра подмножеств (2/3)}
\begin{tabular}{p{12mm}|p{40mm}|p{92mm}}
  № & Название параграфа
  & Ключевые термины и обозначения \\
  \hline
\DMref{1.2.1.}{1.2.1.} & Сравнение множеств 
  & 
  $\subset$~--- включение множеств;
  $=$~--- равенство множеств;
  собственные и несобственные подмножества 
   \\
\DMref{1.2.2.}{1.2.2.} & Равномощные множества 
  & 
 $\hookrightarrow$~--- вложение множеств;
 $\sim$~--- взаимно-одно\-значное соответствие; равномощность
  \\
\DMref{1.2.3.}{1.2.3.} \RS & Конечные и бесконечные множества &
$|A|=\infty$\\
\DMref{1.2.4.}{1.2.4.}  & Счётные и несчётные множества &
$T(n)$~--- треугольные числа;
алгоритм перечисления конечных подмножеств натурального ряда 
\\
\DMref{1.2.5.}{1.2.5.} & Мощность конечного множества &
Принцип Дирихле\\   
\DMref{1.2.6.}{1.2.6.} & Операции над множествами 
  & 
  $\cup$~--- объединение;
  $\cap$~--- пересечение;
  $\setminus$~--- разность;
  $\vartriangle$~--- симметрическая разность;
 дополнение, диаграмма Венна 

\end{tabular}
\end{frame}

\begin{frame}
\frametitle{1.2. Алгебра подмножеств (3/3)}
\begin{tabular}{p{12mm}|p{40mm}|p{92mm}}
  № & Название параграфа
  & Ключевые термины и обозначения \\
  \hline
     \\
\DMref{1.2.7.}{1.2.7.} & Разбиения и покрытия 
  & 
{Дизъюнктное семейство, 

теорема Кантора--Бернштейна}\\
\DMref{1.2.8.}{1.2.8.} \RS & Булеан &
$2^M$~--- множество подмножеств; 
алгебра подмножеств, парадокс Кантора\\
\DMref{1.2.9.}{1.2.9.}  & Свойства операций над множествами &
Идемпотентность, коммутативность, ассоциативность, дистрибутивность, поглощение, инволютивность,
законы де Моргана
\end{tabular}
\end{frame}

\begin{frame}
\label{1.2.1.}\subsubsection{1.2.1. Сравнение множеств }\frametitle{1.2.1. Сравнение множеств (1/2)}
%\vspace{-4pt}
Множество $A$ \odmkey{содержится}{Содержится (contain)} в множестве $B$ (множество $B$
\odmkey{включает}{Включение (inclusion)!множеств (of sets)}
множество $A$), если каждый элемент множества $A$ есть элемент
множества $B$:
$A\subset B \Def x\in A \Impl x\in B.$
При этом говорят, что
$A$~--- 
\odmkey{подмножество}{Подмножество (subset)} $B$, $B$~---
\odmkey{надмножество}{Надмножество (superset)} $A$. 

По определению $\A{M}{\emptyset \subset M}$.


Два множества \odmkey{равны}{Равенство (equality)!множеств (of
sets)}, если они являются подмножествами друг друга:
$$A=B\Def A\subset B \And B\subset A.$$

\begin{theorem}
Включение множеств обладает следующими свойствами:
%\begin{enumerate}[label=\arabic*)]
%\item
%$\A{A}{A\subset A}$;%
%\item
%$\A{A,B}{A\subset B \And B\subset A \Impl A=B}$;%
%\item
%$\A{A,B,C}{A\subset B \And B\subset C \Impl A\subset C}$.
%\end{enumerate}
\begin{enumerate}
\item[1)]
%\\{\color{blue} 1)}
$\A{A}{A\subset A}$;%
%\\{\color{blue} 2)}
\item[2)]
$\A{A,B}{A\subset B \And B\subset A \Impl A=B}$;%
%\\{\color{blue} 3)}
\item[3)]
$\A{A,B,C}{A\subset B \And B\subset C \Impl A\subset C}$.
\end{enumerate}
\end{theorem}

\begin{proof}
\proofsection{\color{blue}1)} Имеем $x\in A\Impl x\in A$ и, значит, $A\subset A$. 
\proofsection{\color{blue}2)} По определению. 
\proofsection{\color{blue}3)} Если $x\in A\Impl x\in B$ 
и $x\in B\Impl x\in C$, то $x\in A\Impl x\in C$
и $A\subset C$.
\end{proof}
\end{frame}

\begin{frame}
\frametitle{1.2.1. Сравнение множеств (2/2)}
Если $A \subset B$ и $A \ne B$, $A \ne \emptyset$, то $A$ называется
\odmkey{собственным}{Подмножество (subset)!собственное (proper)}
подмножеством $B$,
а $B$~---
\odmkey{собственным}{Надмножество (superset)!собственное (proper)}
надмножеством $A$. Таким образом,
само множество и пустое множество считаются 
\em{несобственными} подмножествами.

\begin{example}
Любое одноэлементное множество имеет два различных несобственных подмножества и не имеет собственных подмножеств.
\end{example}

\smallskip
\begin{remark}
В этом курсе для обозначения включения собственных подмножеств
используется знак $\SubsetNE$,
а для несобственных~---~$\subseteq$.
В других источниках иногда для обозначения
собственного включения 
используют знак $\subset$,
оставляя знак $\subseteq$ 
для обозначения общего случая включения,
которое может быть несобственным. 
\end{remark}

\smallskip
\begin{corollary}
$A\SubsetNE B\subseteq C\Impl A\SubsetNE C$.
\end{corollary}

\smallskip
\begin{remark}
Нетрудно видеть, что понятие подмножества множества $X$
равнообъёмно понятию индикатора над множеством $X$. Действительно,
если $X=\{\Tuple{x}{1}{n}\}$~---  множество, а $Y$~--- любое
подмножество множества $X$, $Y\subset X$, то существует
единственный индикатор $\widehat{Y}=\langle
a_1(x_1),\dots,a_n(x_n)\rangle$ над множеством $X$, такой, что $$
\A{i\in 1..n}{a_i=1 \Equiv x_i\in Y}. $$
\end{remark}


\end{frame}

\begin{frame}
\label{1.2.2.}\subsubsection{1.2.2. Равномощные множества }\frametitle{1.2.2. Равномощные множества (1/3)}
%\vspace{-2pt}

Говорят, что установлено 
\odmkey{соот\-вет\-ст\-вие} {Соответствие (correspondence)} из 
множества $A$ в множество $B$, если каждому
элементу множества $A$  поставлен в соответствие один и только
один элемент множества $B$,
причём различным элементам множества $A$ поставлены в соответствие различные элементы множества $B$. 
В этом случае 
говорят также, что
множество $A$ \odmkey{вложено}{Вложение
(injection)!множества (of set)} в множество $B$, и 
используют обозначение $A\hookrightarrow B$. 
Если элементу $a\in A$ соответствует
элемент $b\in B$, то данное обстоятельство обозначают следующим образом: $a\mapsto b$.

\smallskip
\begin{example}
$\Bbb{N}\hookrightarrow\Bbb{Q}$. 
Соответствие: $n\mapsto\frac{n}{1}$.
\end{example} 

\medskip
\begin{remark}
Весьма общее понятие соответствия трактуется здесь с
сугубо программистских позиций.
Если сказано, что установлено соответствие
из множества $A$ в множество 
$B$, то подразумевается, что задан способ
по любому элементу $a\in A$ однозначно определить соответствующий ему элемент
$b\in B$.  Способ может быть любым, если возможность его применения
не вызывает сомнений. Например, это может быть массив
$\mbox{\textbf{array} } [A] \mbox{ \textbf{of} } B$, или процедура
типа $\mbox{\textbf{proc} } (A) : B$, или же выражение, зависящее
от переменной типа $A$ и доставляющее значение типа $B$.
\end{remark}


\end{frame}


\begin{frame}
\frametitle{1.2.2. Равномощные множества (2/3)}

%\vspace{-2pt}
Всякий элемент любого множества очевидно соответствует
сам себе, поэтому  $$A\subset B\Impl A\hookrightarrow B.$$

\begin{example}
Поскольку $\Bbb{N}\subset\Bbb{Q}\subset\Bbb{R}$,
имеем $\Bbb{N}\hookrightarrow \Bbb{Q}\hookrightarrow \Bbb{R}$.
\end{example}

\medskip
Говорят, что между множествами $A$ и $B$ установлено
\odmkey{взаимно-однозначное соот\-вет\-ст\-вие} {Соответствие (correspondence)! взаимно-однозначное (one-to-one correspondence)}, если каждому
элементу $A$  поставлен в соответствие один и только
один элемент $B$, и для каждого элемента $B$
один и только один элемент $A$ поставлен в соответствие
этому элементу $B$. 
В этом случае 
говорят также, что
множества $A$ и $B$ \em{изоморфны}, и 
используют обозначение $A\sim
B$. 
Если элементу $a\in A$ взаимно-однозначно соответствует
элемент $b\in B$, то данное обстоятельство обозначают так: $a \leftrightarrow b$.

\begin{examples}

Соответствие $n\leftrightarrow 2n$ устанавливает взаимно-однозначное соответствие
между множеством натуральных чисел $\mathbb{N}$
и множеством чётных натуральных чисел $2\mathbb{N}$;
$\mathbb{N}\sim 2\mathbb{N}$.

Соответствие $n\leftrightarrow n^2$ устанавливает взаимно-однозначное соответствие
между множеством натуральных чисел 
и множеством полных квадратов. 
\end{examples}
\end{frame}


\begin{frame}
\frametitle{1.2.2. Равномощные множества (3/3)}

\vspace{-2pt}
Если между двумя множествами, $A$ и $B$, может быть установлено
взаимно-одно\-знач\-ное соответствие, то говорят, что множества
имеют одинаковую \odmkey{мощность}{Мощность (cardinality)!множества (of set)}, или что множества
\odmkey{равномощны}{Отношение (relation)!равномощности (of
equivalence)}, и записывают это так: $|A|=|B|$. 

\begin{theorem}
Равномощность множеств обладает следующими свойствами:
\begin{enumerate}
\item[1)]
$\A{A}{|A|=|A|}$;%
\item[2)]
$\A{A,B}{|A|=|B|  \Impl |B|=|A|}$;%
\item[3)]
$\A{A,B,C}{|A|=|B| \And |B|=|C| \Impl |A|=|C|}$.
\end{enumerate}
\end{theorem}
\begin{proof}
\proofsection{\color{blue}1)} достаточно рассмотреть соответствие $a\leftrightarrow a$, 
где $a\in A$; 
\proofsection{\color{blue}2)} ввиду взаимной однозначности соответствия $a\leftrightarrow b$; 
\proofsection{\color{blue}3)} соответствие
устанавливается с использованием элемента $b\in B$:
$a\leftrightarrow b\leftrightarrow c$, где ${a\in A}, {b\in B}, {c\in C}$.
\end{proof}

%\smallskip
Пусть $A\hookrightarrow B$. Обозначим 
$A^{\hookrightarrow^B}\Def\Set{b\in B}{\E{a\in A}{a\mapsto b}}$. Заметим, что 
$A^{\hookrightarrow^B}\subset B$.
\begin{remark}
Если $B$ ясно из контекста, то его 
можно опустить: $A^{\hookrightarrow}= A^{\hookrightarrow^B}$. 
\end{remark}
\smallskip
\begin{lemma}
$A\hookrightarrow B \Impl |A|=|A^{\hookrightarrow}|$.
\end{lemma}  
\begin{proof}
Достаточно положить $a\leftrightarrow b\Assign 
a\mapsto b$.
\end{proof}
\end{frame}



\begin{frame}
\label{1.2.3.}\subsubsection{1.2.3. Конечные и бесконечные множества }\frametitle{1.2.3. Конечные и бесконечные множества (1/2)}
Вопрос о том, чем \em{конечное} отличается от \em{бесконечного},
может поставить в тупик.
Здесь мы рассматриваем применительно к множествам один
из возможных ответов на этот вопрос.
С античных времен известен принцип: <<часть меньше целого>>.
%Оказывается, этот принцип можно использовать в качестве
%характеристического свойства конечных множеств.
Оказывается, для бесконечных множеств этот принцип не имеет места.

\medskip
Множество $A$ называется
\odmkey{конечным}{Множество (set)!конечное (finite)}, если
у него нет равномощного собственного подмножества:
$$
\A{B}{(B\subset A\And|B|=|A|)\Impl (B=A)}.
$$
Для конечного множества $A$ используется запись $|A|<\infty$.
Все остальные множества называются
\odmkey{бесконечными}{Множество (set)!бесконечное (infinite)}:
$$
\E{B}{B\subset A \And |B|=|A| \And B\neq A},
$$
то есть бесконечное множество равномощно
некоторому своему собственному подмножеству.
Для бесконечного множества $A$ используется запись
$|A|=\infty$.
\end{frame}

\begin{frame}
\frametitle{1.2.3. Конечные и бесконечные множества (2/2)}
\begin{example}
Множество натуральных чисел бесконечно, $|\mathbb{N}|=\infty$,
поскольку оно равномощно своему собственному подмножеству
чётных чисел.
\end{example}
\begin{theorem}
Множество, имеющее бесконечное подмножество,
бесконечно:
$$
(B\subset A \And |B|=\infty) \Impl (|A|=\infty).
$$
\end{theorem}
\begin{proof}
Множество $B$ бесконечно, значит $B\sim C$, причём $C\subset B$. Обозначим это соответствие
$x\mapsto x'$. Построим соответствие между множеством $A$ и его
собственным подмножеством $D$:
$
\textbf{$x\mapsto$ if $x\in B$ then $x'$ else $x$ end if}.
$
\end{proof}
\begin{center}
\begin{columns}[t] % contents are top vertically aligned
\begin{column}[T]{4cm} % alternative top-align that's better for graphics
\includegraphics[width=4cm]{f1_2_3_2_2.jpg}
\end{column}
\begin{column}[T]{10cm}
Другими словами, на элементах из $B$ мы пользуемся заданным
соответствием, а~остальным элементам сопоставляем их самих. 
Это взаимно-одно\-значное соответ\-ствие между
множеством $A$ и его собственным подмножеством $D$, и, значит,
$|A|=\infty$.
\end{column}
\end{columns}
\end{center}

\vspace{-8pt}
\begin{corollary}
Все подмножества конечного множества конечны.
\end{corollary}
\begin{proof}
Пусть $\E{A\subset B}{|A|=\infty 
\And |B|<\infty}$.
Но по теореме $|B|=\infty$, так как
$|A|=\infty \And A\subset B$. Противоречие.
\end{proof}
\end{frame}

\begin{frame}
\label{1.2.4.}\subsubsection{1.2.4. Счётные и несчётные множества }\frametitle{1.2.4. Счётные и несчётные множества (1/5)}
Бесконечные множества, равномощные множеству натуральных чисел, 
называются \odmkey{счёт\-ными}
{Множество (set)!счётное (countable)}.
Счётные множества отличаются тем, 
что для каждого элемента можно указать \\
взаимно-однозначно соответствующее 
натуральное число, то есть \em{номер}. 
Другими словами, множество счётно, если его элементы можно
 \em{перенумеровать}.

\begin{example}
Множество целых чисел $\Bbb{Z}$ счётно. 
Соответствие задаётся последовательностью: $0, 1, -1, 2, -2, \dots$.
\end{example} 

\medskip 
Фактическая нумерация элементов счётного множества не всегда столь проста, 
иногда требуются более изощрённые построения.

\begin{example}
Множество точек на плоскости с натуральными координатами счётно. 
\end{example}

Для построения нумерации здесь используются так называемые 
\odmkey{треугольные}{Число (number)!треугольное (triangular)} числа $T(k)$, вычисляемые по формуле

%\vspace{-12pt}
$$
T(k) = k(k+1)/2 = \sum\limits_{i=1}^k i = T(k-1)+k. 
$$

\end{frame}

\begin{frame}
\frametitle{1.2.4. Счётные и несчётные множества (2/5)} 

На рисунке слева показано начало одной из многих возможных нумераций
с помощью треугольных чисел. 
Заметим, что для всех точек $(m,n)$, 
лежащих на каждой <<обратной диагонали>>, 
сумма $m+n$ постоянна, а точки, лежащие на вертикальной оси $m=1$, 
имеют номера, являющиеся треугольными числами (см. рисунок справа). 

\begin{center}
\includegraphics[height=50mm]{f1_2_4_2_5.jpg}
\end{center}
\end{frame}

\begin{frame}
\frametitle{ 1.2.4. Счётные и несчётные множества (3/5)}


Вернемся теперь к установлению взаимно-однозначного соответствия 
между точками плоскости и натуральными числами. 

Точка с координатами $(m,n)$ получает номер $k$, вычисляемый следующим образом:\\
$k\Assign T(m+n-1)-m+1=(m+n)(m+n-1)/2-m+1$.

Если точка имеет номер $k$ в данной нумерации, 
то её координаты $(m,n)$ вычисляются следующим образом:

\textbf{$i:=1; T:=1;$ while $T<k$ do $i\Assign i+1; T\Assign T+i$ 
end while};\\ \textbf{$m:=T-k+1; n:=i-m+1$} 

\begin{columns}[t] % contents are top vertically aligned
\begin{column}[T]{3cm} % alternative top-align that's 
\begin{remark}
Почему числа $T(n$) называются треугольными, объяснено на рисунке.
\end{remark}
\end{column}
\begin{column}[T]{12cm}
\begin{center}
\includegraphics[height=30mm]{f1_2_4_3_5.jpg}
\end{center}
\end{column}
\end{columns} 
\end{frame}

\begin{frame}
\frametitle{1.2.4. Счётные и несчётные множества (4/5)}
Не все бесконечные множества являются счётными. 
Существуют \odmkey{несчётные}{Множество (set)!несчётное (countless)} множества.

\begin{theorem}
Множество всех подмножеств натурального ряда несчётно.
\end{theorem}


\begin{proof}
Допустим, что множество всех подмножеств натурального ряда счётно. 
Тогда их можно перенумеровать, и каждое множество чисел получит свой номер. 
Обозначим $N(X)$~--- номер множества $X$. 
Может статься, что некоторые множества чисел содержат 
свой номер в качестве элемента, 
а другие не содержат. 
Рассмотрим числовое множество $Y$, 
которое состоит из номеров множеств, не содержащих свой номер в качестве элемента:
$$
Y = \Set{x\in\Bbb{N}}{ \E{X\subset\Bbb{N}}{x=N(X)\And x\not\in X}}.
$$ 

Пусть множество $Y$ имеет номер $y=N(Y)$. 
Тогда, если $y\not\in Y$, то множество $Y$ не содержит свой номер, 
а значит, по своему определению, $y\in Y$. 
Обратно: если $y\in Y$, то $y\not\in Y$. 
Противоречие, 
и значит нумерация всех числовых множеств невозможна.
\end{proof}


\end{frame}

%\frame[shrink]{
\begin{frame}
\frametitle{ 1.2.4. Счётные и несчётные множества (5/5)}

\vspace{-2pt}
Следующий алгоритм перечисляет все \em{конечные} подмножества натурального
ряда, откуда следует, что множество всех 
конечных подмножеств натурального ряда счётно.  

\vspace{-4pt}
\begin{algorithmic}
\STATE{$n\Assign 0$}
\COMMENT{\small\color{gray} добавляемое число }
\STATE{$Q\leftarrow \emptyset$}
\COMMENT{\small\color{gray} очередь множеств содержит пустое множество }
\STATE{\textbf{yield $\emptyset$}}
\COMMENT{\small\color{gray} первое подмножество пустое}
\WHILE{\textbf{true}}
\STATE{$S\leftarrow Q$}
\COMMENT{\small\color{gray} извлекаем первое множество из очереди }
\STATE{\textbf{if $S=\emptyset$ then $n\Assign n+1$ end if}}
\COMMENT{\small\color{gray} следующее добавляемое число }
\STATE{$Q\leftarrow S$}
\COMMENT{\small\color{gray} возвращаем извлечённое множество в очередь }
\STATE{$S\Assign S + n $}
\COMMENT{\small\color{gray} добавляем число $n$ в множество $S$ }
\STATE{$Q\leftarrow S$}
\COMMENT{\small\color{gray} добавляем новое множество в очередь }
\STATE{\textbf{yield $S$}}
\COMMENT{\small\color{gray} очередное множество }
\ENDWHILE
\end{algorithmic}

\vspace{-2pt}
\begin{ground}  Множество всех конечных подмножеств отрезка $1..n$ 
состоит из двух равномощных частей: 
множества всех подмножеств отрезка $1..(n-1)$ 
и множества этих же подмножеств, в каждое из которых добавлено число $n$.
\end{ground}     

\end{frame}
%}


\begin{frame}
\label{1.2.5.}\subsubsection{1.2.5. Мощность конечного множества }\frametitle{1.2.5. Мощность конечного множества (1/5)}

\begin{theorem}
Любое непустое конечное множество равномощно
некоторому отрезку натурального ряда:
$
\A{A}{A\neq\emptyset\And|A|<\infty\Impl\E{k\in\mathbb{N}}{|A|=|1..k|}}.
$
\end{theorem}
\begin{proof}
Рассмотрим следующую программу:
%\vspace{-\baselineskip}
\begin{algorithmic}
\STATE $i\Assign 0$
\COMMENT{\small\color{gray} счётчик элементов }
\WHILE{$A\neq\emptyset$}
\STATE \textbf{select $x\in A$}
\COMMENT{\small\color{gray} выбираем элемент }
\STATE $i\Assign i+1$
\COMMENT{\small\color{gray} увеличиваем счётчик }
\STATE $x\mapsto i$
\COMMENT{\small\color{gray} ставим элементу в соответствие его номер }
\STATE $A\Assign A-x$
\COMMENT{\small\color{gray} удаляем элемент из множества }
\ENDWHILE
\end{algorithmic}
Если эта программа не заканчивает работу,
то она даёт соответствие $B\sim \mathbb{N}$
для некоторого множества $B\subset A$,
что невозможно
ввиду конечности $A$. Значит, процедура заканчивает работу
при~$i=k$. Но в этом случае построено взаимно-однозначное
соответствие
$A\sim 1..k$.
\end{proof}
\end{frame}

\begin{frame}
\frametitle{1.2.5. Мощность конечного множества (2/5)}
\begin{lemma}
Любое непустое подмножество множества натуральных чисел со\-дер\-жит
наименьший элемент.
\end{lemma}
\begin{proof}
Пусть $A$~---  произвольное подмножество множества натуральных
чисел, конечное или бесконечное. Рассмотрим задание множества $A$
следую\-щей порождающей процедурой:

\begin{algorithmic}
\STATE $n\Assign 0$
\COMMENT{\small\color{gray} переменная для натуральных чисел }
\WHILE{\textbf{true}}
\STATE $n\Assign n+1$
\COMMENT{\small\color{gray} следующее натуральное число }
\IF{$n\in A$}
\STATE \textbf{yield $n$}
\COMMENT{\small\color{gray} число входит в множество $A$ }
\ENDIF   
\ENDWHILE
\end{algorithmic}

Ясно, что множество $A$ действительно порождается этой процедурой,
причём тот элемент, который порождается первым, является
наименьшим.
\end{proof}

\end{frame}

\begin{frame}
\frametitle{1.2.5. Мощность конечного множества (3/5)}
%\vspace{-2pt}
\begin{theorem}
Любой отрезок натурального ряда конечен:
$$
\A{n\in\mathbb{N}}{|1..n|<\infty}.
$$
\end{theorem}
\begin{proof}
От противного. Пусть существуют бесконечные отрезки натурального
ряда. Рассмотрим наименьшее $n$ такое, что $|1..n|=\infty$. Тогда
отрезок $1..n$ равномощен некоторому своему собственному
подмножеству $A$: $|1..n|=|A|$, $A\subsetneq 1..n$, то есть существует взаимно-однозначное соответствие $1..n \sim A$. 
Пусть при этом соответствии $n\leftrightarrow i$.
Может случиться так, что $n\in A$ и $n\ne i$.
Тогда <<подправим>> соответствие так, что
$n\leftrightarrow n$ и $i \leftrightarrow i$. 
Рассмотрим теперь соответствие между отрезком $1..(n-1)$ и его
собственным подмножеством $A-i$, задаваемое соответствием между
отрезком $1..n$ и множеством $A$. Это соответствие является взаимно-однозначным, а
значит, отрезок $1..(n-1)$ изоморфен своему собственному
подмножеству и~является бесконечным, что противоречит выбору $n$.
\end{proof}

\end{frame}

\begin{frame}
\frametitle{1.2.5. Мощность конечного множества (4/5)}
\begin{corollary}
Различные отрезки натурального ряда неравномощны:
$$
n\neq m \Impl |1..n|\neq|1..m|.
$$
\end{corollary}
\begin{proof}
Пусть для определённости $n>m$. 
Тогда $1..m \subset 1..n$ и $1..m \neq  1..n$. Если $|1..n| = |1..m|$, 
то $|1..n| = \infty$, что противоречит теореме.
\end{proof}

\medskip
\begin{remark}
Это следствие часто излагают в следующей форме: если $n\neq m$, то $n$ кроликов невозможно рассадить в $m$ ящиков так, чтобы в каждом ящике было по одному кролику. Действительно, если бы рассадить кроликов было возможно, то было бы установлено взаимно-однозначное соответствие между различными отрезками натурального ряда. Это утверждение называется \odmkey{принцип Дирихле}{Принцип (principle)!Дирихле (Dirichlet)}.
\end{remark}

\medskip
Говорят, что конечное множество $A$ имеет
\odmkey{мощность}{Мощность (cardinality)!множества (of set)!конечного (finite)} $k$ (обозначения: $|A|=k$, $\card A =k$, $\#A =k$),
если оно равномощно отрезку $1..k$:
$$
|A|=k\Def A\sim 1..k.
$$
\end{frame}

\begin{frame}
\frametitle{1.2.5. Мощность конечного множества (5/5)}


%\vspace{-4pt}
\begin{remark}
Таким образом, если множество $A$ конечно, $|A|=k$, то элементы
$A$ всегда можно \odmkey{перенумеровать}{Нумерация множества
(numbering of set)}, то есть поставить им в соответствие номера из
отрезка $1..k$ с помощью некоторой процедуры. Наличие такой
процедуры подразумевается, когда употребляется запись
$A=\{\Tuple{a}{1}{k}\}$.
\end{remark}

\medskip
По определению $|\emptyset|\Def 0$.

\medskip
\begin{digress}
Определение мощности конечного множества как 
натурального числа позволяет сравнивать конечные множества по мощности естественным и привычным образом, например, 
констатировать, что   $A\subset B \Impl |A|\leq|B|$.
Однако указанное определение мощности для
конечных множеств не распространяется на бесконечные множества,
поэтому сравнивать по мощности бесконечные
множества надобно с большой осторожностью.
Для этого вводятся 
\odmkey{кардинальные числа}{Число (number)!кардинальное  (cardinal)} 
и другие \odmkey{трансфинитные}{Трансфинитный (transfinite)} построения.
Точное определение соответствующих понятий выходит за намеченные здесь рамки наивной теории множеств, 
поскольку эти понятия \em{никогда} не используются в программировании.   
\end{digress}

\end{frame}

\begin{frame}
\label{1.2.6.}\subsubsection{1.2.6. Операции над множествами }\frametitle{1.2.6. Операции над множествами (1/4)}
Обычно рассматриваются следующие
операции над множествами:

\begin{tabbing}
\hspace{5mm}\= \odmkey{симметрическая}
{Разность (difference)!симметрическая
(symmetric)}: \=
$A \vartriangle B \Def (A \cup B) \setminus (A \cap B) $; =\kill
\>\odmkey{объединение}{Объединение (union)!множеств (of sets)}: 
\>{$A\cup B \Def \Set{x}{x\in A \lor x\in B}$}; \\[1pt]
\>\odmkey{пересечение}{Пересечение (intersection)!множеств (of sets)!}: 
\>{$A\cap B \Def \Set{x}{x\in A \And x\in B}$}; \\[1pt]
\>\odmkey{разность}{Разность (difference)!множеств (of sets)}:
\>{$A\setminus B \Def \Set{x}{x\in A \And x\notin B}$}; \\[1pt]
\>\odmkey{симметрическая}{Разность (difference)!симметрическая
(symmetric)} 
\>$A \vartriangle B \Def (A \cup B) \setminus (A \cap B) $; \\[1pt]
\>\odmkey{разность}{Разность (difference)!симметрическая
(symmetric)}: 
%\>$A \vartriangle B \Def (A \cup B) \setminus (A \cap B) $; \\[3pt]
\>{$A \vartriangle B \Def \Set{x}{x\in A \not\equiv x\in B}$;} \\[1pt]
\>\odmkey{дополнение}{Дополнение (complement)!множества (of a set)}:
\>{\footnotesize$\overline{A} \Def \Set{x}{x\notin A}$}.
\end{tabbing}

Результаты операций объединения, пересечения, разности и симметрической разности множеств
всегда определены и являются множествами. Операция дополнения подразумевает, что задан некоторый
универсум $U$:
$\overline{A}=U\setminus A.$
В противном случае операция дополнения не определена.

\begin{example}
Пусть $A\Assign \{1,2,3\}$, $B\Assign \{3,4,5\}$. 
Тогда $A\cap B = \{3\}$, 
$A\cup B = \{1,2,3,4,5\}$,  $A\setminus B = \{1,2\}$,
$A\vartriangle B = \{1,2,4,5\}$. 
\end{example}
\end{frame}

\begin{frame}
\frametitle{1.2.6. Операции над множествами (2/4)}

\vspace{-2pt}
На рисунке приведены \odmkey{диаграммы Венна}{Диаграмма
(diagram)!Венна (John Venn)},
иллюстрирующие операции над множествами. Сами
исходные множества изображаются фигурами (в данном случае
кругами), а результат выделяется графически  (в данном случае для
выделения использована заливка цветом).

\vspace{-6pt}
\begin{figure}[h]
\begin{center}
\includegraphics[width=120mm]{f1_2_6_2_4.jpg}
\end{center}
\end{figure}

\vspace{-4pt}
В некоторых случаях операции над множествами выразимы через комбинации других операций.

\begin{tabular}{| c | c | c | c |}
\hline & $ \cup , \cap $ & $ \cup , \setminus $ & $ \cup , \vartriangle $  \\[-2pt]
\hline
$A\cup B =$  & $A\cup B$       & $ A\cup B$           &    $ A\cup B$     \\
$A\cap B =$  & $A\cap B$  & {$A\setminus (A\setminus B)$}
& {$(A\cup B) \vartriangle (A \vartriangle B)$} \\
$A\setminus B =$  &     ---        & $ A \setminus B$     
& {$(A\cup B) \vartriangle B$}\\
$A\vartriangle B =$ &    ---      & {$(A\setminus B)\cup (B\setminus A)$}
& $A \vartriangle B$\\
\hline
\end{tabular}


\end{frame}


\begin{frame}
\frametitle{ 1.2.6. Операции над множествами (3/4)}
%\vspace{-4pt}

\begin{tabular}{| c | c | c | c |}
\hline
& $ \cap ,\setminus $ & $ \cap , \vartriangle $ & $ \setminus , \vartriangle $ \\
\hline
$A\cup B =$  & --- & {$(A\vartriangle B) \vartriangle (A\cap B)$}&
{$(A\vartriangle B)\vartriangle (A\setminus (A\setminus B))$} \\
$A\cap B =$  & $ A \cap B $   & $ A \cap B $ 
& {$A\setminus (A\setminus B)$} \\
$A\setminus B =$ & $A \setminus B$ & {$A \vartriangle (A \cap B)$} 
& $ A\setminus B$  \\
$A\vartriangle B =$ & --- & $A  \vartriangle B$ & $A  \vartriangle B$ \\
\hline
\end{tabular}


\medskip
В то же время, в четырёх случаях выражения не существует. 

Разность и симметрическая разность невыразимы через объединение и пересечение. 
Действительно, если $A=B\ne\emptyset$, 
то  $A\setminus B = \emptyset$,
$A\vartriangle B = \emptyset$, 
при этом любые комбинации объединения и пересечения непустого множества 
с самим собой дают это же самое множество. 

Объединение и симметрическая разность невыразимы через пересечение и разность,
поскольку результат пересечения и разности всегда является подмножеством первого аргумента, 
и поэтому значение любого выражения, 
составленного из пересечений и разностей, 
является подмножеством одного из исходных аргументов.

\end{frame}



\begin{frame}
\frametitle{1.2.6. Операции над множествами (4/4)}

Если множества $A$ и $B$ конечны, то из определений и диаграмм Венна
ясно, что
\begin{align*}
|A\cup B| &= |A|+|B| - |A\cap B|, \\
|A\setminus B| &= |A| - |A\cap B|, \\
|A\vartriangle B| &= |A|+|B| - 2|A\cap B|.
\end{align*}
%$ |A\cup B| = |A|+|B| - |A\cap B|$,
%$|A\setminus B| = |A| - |A\cap B|$,
%$|A\vartriangle B| = |A|+|B| - 2|A\cap B|$.

\medskip
Операции пересечения и объединения двух множеств допускают следующее об\-общение.
Пусть задано семейство множеств $\left\{A_i\right\}_{i\in I}$. Тогда


$$
\bigcup_{i\in I} A_i \Def \Set{x}{\E{i\in I}{x\in A_i}},
\quad
\bigcap_{i\in I} A_i \Def \Set{x}{\A{i\in I}{x\in A_i}}.
$$

%\vspace{-4pt}
\begin{remark}
Если операция ассоциативна 
(см. \DMref{п.~1.2.9}{1.2.9.} и 
\DMref{2.1.4}{2.1.4.}), то неважно,
как расставлены скобки в выражении с этой операцией.
В таком случае можно считать,
что операция допускает \em{переменное} количество 
аргументов, два или более, 
и считать операцию 
\odmkey{групповой}{Операция (operation)!групповая (group)}.
Примерами групповых операций являются
введённые объединение $\bigcup$ и пересечение $\bigcap$,
а также хорошо известные сумма  
$\sum$ и произведение $\prod$,
дизъюнкция $\BigOr$ и конъюнкция $\bigwedge$ 
(\DMref{п.~3.2.2}{3.2.2.}).
\end{remark}

\end{frame}

\begin{frame}
\label{1.2.7.}\subsubsection{1.2.7. Разбиения и покрытия }\frametitle{1.2.7. Разбиения и покрытия (1/5)}
Пусть $\cal E=\left\{E_i\right\}_{i\in I}$~---
семейство подмножеств множества $M$; ${E_i\subset M}$.
Семейство $\cal E$ называется \odmkey{покрытием}{Покрытие
(covering)} множества $M$, если каждый элемент $M$ принадлежит
хотя бы одному из $E_i$:
$$
\A{x\in M}{\E{i \in I}{x\in E_i}}.
$$
Семейство $\cal E$ называется
\odmkey{дизъюнктным}{Семейство (family)!дизъюнктное (disjunct)}, если элементы
этого семейства попарно не пересекаются, то~есть каждый элемент
$M$ принадлежит не более чем одному из множеств $E_i$:
$$ \A{i,j\in I}{i \ne j \Impl E_i \cap E_j = \emptyset}. $$
Дизъюнктное покрытие называется \odmkey{разбиением}{Разбиение
(partition)} множества $M$. Элементы разбиения, то есть
подмножества множества $M$, часто называют \odmkey{блоками}{Блок
(block)!разбиения (of partition)} разбиения.

\begin{theorem}
Если $\cal E = \left\{E_i\right\}_{i\in I}$ есть дизъюнктное
семейство подмножеств множества $M$, то существует разбиение
$\cal B = \left\{B_i\right\}_{i\in I}$ множества $M$ такое, что
каждый элемент дизъюнктного семейства $\cal E$ является
подмножеством блока разбиения~$\cal B$: $\A{i\in I}{E_i\subset B_i}.$
\end{theorem}
\end{frame}

\begin{frame}
\frametitle{1.2.7. Разбиения и покрытия (2/5)}

\begin{proof}
Выберем произвольный элемент $i_0\in I$
и построим семейство \\$\cal B = \left\{B_i\right\}_{i\in I}$
следующим образом:
$$
B_{i_0}\Assign M\setminus\bigcup\limits_{i\in I-i_0} E_i, \qquad
\A{i\in I-i_0}{B_i\Assign E_i}.
$$
Семейство $\cal B$ по построению является дизъюнктным покрытием,
то есть разбиением.
Ясно, что $E_{i_0}\subset B_{i_0}$,
а для остальных элементов требуемое включение имеет место
по построению.
\end{proof}

\begin{example}
Пусть $M\Assign \{1,2,3\}$. Тогда
элементы дизъюнктного семейства\\ 
$\cal E = \{E_1=\{1\}, E_2=\{2\}\}$
являются подмножествами блоков
разбиения \\$\cal B =\{B_1=\{1, 3\},B_2=\{2\}\}$.
\end{example}

\begin{corollary}
Если все множества $E_i$ конечны и множество 
индексов $I$ конечно, а семейство 
$\cal{E}=\left\{E_i\right\}_{i\in I}$ дизъюнктно, то 
$$\displaystyle{\left\vert \bigcup_{ i\in I} E_i \right\vert = \sum_{i\in I} \vert E_i \vert.}$$
\end{corollary} 
\end{frame}

\begin{frame}
\frametitle{1.2.7. Разбиения и покрытия (3/5)}
%\vspace{-2pt}
\begin{theorem}
Если $A\hookrightarrow B$ и $A\hookleftarrow B$,
то $A\sim B$.
\end{theorem}
%%%%%%% старое доказательство
\vspace{-2pt}
\begin{proof} Определим 
операцию ${}^*$ так:
$\A{C\subset A}{C^* \Assign A\setminus 
(B\setminus C^{\hookrightarrow})^{\hookleftarrow}}$.
Заметим, что $D\subset C\subset A\Impl D^{\hookrightarrow}\subset C^{\hookrightarrow}\subset B\Impl$
$B\setminus C^{\hookrightarrow}\subset B\setminus D^{\hookrightarrow}\Impl$
$(B\setminus C^{\hookrightarrow})^{\hookleftarrow}
\subset (B\setminus D^{\hookrightarrow})^{\hookleftarrow}\Impl$
$A\setminus (B\setminus D^{\hookrightarrow})^{\hookleftarrow}\subset A\setminus (B\setminus C^{\hookrightarrow})^{\hookleftarrow}\Impl$
$D^*\subset C^*\subset A$,
поскольку $\A{X,Y,Z}{X\subset Y\subset Z\Impl
Z\setminus Y\subset Z\setminus X}$
(в терминах \DMref{п.~1.8.6}{1.8.6.} операция ${}^*$ монотонна 
по включению).
Заметим также, что $(C\cup D)^* = C^* \cup D^*$, 
поскольку 

\vspace{-12pt}
$$\A{X,Y,Z,W}
{X\setminus (Y\setminus (Z\cup W))=
(X\setminus (Y\setminus Z))\cup
(X\setminus (Y\setminus W))}.$$

\vspace{-2pt}
В терминах \DMref{п.~1.7.6}{1.7.6.} операция ${}^*$~--- гомоморфизм.
Пусть $\cal{M}\Assign
\Set{C\!\subset\! A}{C\!\subset\! C^*}$. Семейство 
$\cal{M}$ не пустое:
$\emptyset\in\cal{M}$.
Рассмотрим множество $S\Assign\bigcup_{C\in\cal{M}} C$. Имеем $S^* =\left(\bigcup_{C\in\cal{M}} C\right)^* =$ $= \bigcup_{C\in\cal{M}} C^*$, значит $S\subset S^*$ и $S\in \cal{M}$.     
По свойству монотонности получаем 
$S^*\subset {S^*}^*$, так что $S^*\in \cal{M}$, а значит, $S^*\subset S$. 
Таким образом, 
$S= S^* = A \setminus (B \setminus S^{\hookrightarrow})^{\hookleftarrow}$, 
откуда \\$A\setminus S = (B\setminus S^{\hookrightarrow})^{\hookleftarrow}$,
поскольку $\A{X,Y,Z}
{X= Y\setminus Z \And Z\subset Y \Impl
Z=Y\setminus X}$.
Теперь рассмотрим разбиение множества
$A = (A\setminus S)\cup S$ и множества
$B = (B\setminus S^{\hookrightarrow})\cup S^{\hookrightarrow}$.
Определим 
требуемое взаимно-однозначное соответствие 
на подмножествах $S$ и $S^{\hookrightarrow}$
с помощью вложения
$A\hookrightarrow B$,
а на подмножествах $A\setminus S$ 
и $B\setminus S^{\hookrightarrow}$
с помощью вложения
$A\hookleftarrow B$. 
\end{proof}
\end{frame}

\begin{frame}
\frametitle{1.2.7. Разбиения и покрытия (4/5)}
%\vspace{-2pt}
%%%%%%%%%%%%%% Козаков 2020
Приведём доказательство этой важной теоремы
без опоры на понятие семейства.
\begin{proof}
Введём обозначения:
$\ A_0 \Assign A$,
$\ B_0 \Assign B$, 
$\ B_1 \Assign A^{\hookrightarrow}$,
$\ A_1 \Assign B^{\hookleftarrow}$,
и так далее, 
$B_{i+1} \Assign A_i^{\hookrightarrow}$,
$A_{i+1} \Assign B_i^{\hookleftarrow}$. 
В этих обозначениях
$A_0 \supset A_1 \supset A_2 \supset \dots$,\\ 
$B_0 \supset B_1 \supset B_2 \supset \dots$,
причём 
$A_0 \sim B_1 \sim A_2 \sim  \dots$, 
$B_0 \sim A_1 \sim B_2 \sim \dots$,\\
то есть
$A_i \sim A_{i+2}$, 
$B_i \sim B_{i+2}$.
Осталось доказать, что $A_0 \sim A_1$.
Обозначим $C_i\Assign A_i\setminus A_{i+1}$.

%\medskip
\begin{columns}[t]
\begin{column}[T]{71mm}
\includegraphics[width=70mm]{f1_2_7_4_5.jpg}
\end{column}
\begin{column}[T]{82mm}
Имеем
$C_i=A_i\setminus A_{i+1} \sim 
A_{i+2}\setminus A_{i+1+2}= C_{i+2}$.
Обозначим $C\Assign \bigcap_i A_i$~--- 
<<сердцевина>>, возможно, пустая. 
По построению $A_0 = C_0 \cup C_1 \cup C_2 \cup \dots \cup C $ и 
$A_1 =  C_1 \cup C_2 \cup \dots \cup C $.
Пусть вложение
$A\hookrightarrow B$ элементу $a\in A$ 
сопоставляет элемент $f(a)\in B$, а
вложение
$A\hookleftarrow B$ элементу $b\in B$ 
сопоставляет элемент $g(b)\in A$.
\end{column}
\end{columns}


Тогда построим 
соответствие $h$ из $A_0$ в $A_1$:
\textbf{$h(a) \Assign $ if $a\in C_{2n}$ then 
$g(f(a))$ else $a$ end if}.
Это соответствие взаимно-однозначно.
В самом деле, $h(C_{2n}) = C_{2n+2}$, $h(C_{2n+1}) = C_{2n+1}$, и $h(C) = C$. Таким образом, $h(A_0) = h(C_0 \cup C_1 \cup C_2 \cup \dots \cup C) = 
C_1 \cup C_2 \cup \dots  \cup C = A_1$,
и значит $A_0 \sim A_1$.
\end{proof}
\end{frame}

\begin{frame}
\frametitle{1.2.7. Разбиения и покрытия (5/5)}

\vspace{-2pt}
\begin{remark}
Доказанное утверждение в отечественной традиции 
называется \odmkey{теоремой Кантора--Бернштейна}
{Теорема (theorem)!Кантора--Бернштейна (Schr$\ddot{o}$der--Bernstein)}, а в зарубежной традиции~--- теоремой Шрёдера--Бернштейна.
\end{remark}

\begin{corollary}
Если $A\subset B \subset C$ и $A\sim C$, то $B\sim C$. 
\end{corollary}

\begin{proof}
$B\subset C \Impl B\hookrightarrow C$, 
$A\subset B \And A\sim C \Impl B\hookleftarrow C$ 
\end{proof}

\begin{digress}
Для случая конечных множеств 
теорема Кантора--Бернштейна
представляется интуитивно 
очевидной. Однако доказательство теоремы для общего случая
оказывается сравнительно
сложным.
Это не случайно, при трансфинитных построениях
интуиция часто подводит, 
поскольку рассуждения не опираются на повседневный опыт. В таких случаях формализация оказывается более надёжным инструментом по сравнению с интуицией. 
\end{digress}

\vspace{-2pt}
\begin{columns}[t] 
\begin{column}[T]{3.5cm}
\vspace{-12pt}
\begin{center}
\begin{figure}
\includegraphics[height=3cm]{cantor_pic.png}
\end{figure}
\end{center}
\end{column}
\begin{column}[T]{12cm}

\vspace{-2pt}
\begin{example} 
Пусть $A$~--- множество точек квадрата, а $B$~--- множество точек вписанного в него круга.
С одной стороны, поскольку круг вписан в квадрат, имеем $B \hookrightarrow A$. 
С другой стороны, исходный квадрат можно сжать в 
$\sqrt{2}/2$ раз (биекция), получив квадрат $A'$,
и вписать его в круг так, что $A \sim A' \hookrightarrow B$. Получаем $A\sim B$, причём соответствие хорошо видно на рисунке: белые точки соответстствуют сами себе, а чёрные точки сопоставляет друг другу \em{гомотетия} $1/\sqrt{2}$.   
\end{example}
\end{column}
\end{columns}

\end{frame}


\begin{frame}
\label{1.2.8.}\subsubsection{1.2.8. Булеан }
\frametitle{1.2.8. Булеан (1/4)}

Множество всех подмножеств множества $M$ называется
\odmkey{булеаном}{Булеан (power set)} множества $M$ и
обозначается~$2^M$:
$$
2^M\Def\Set{A}{A\subset M}.
$$
\begin{theorem} Если множество $M$ конечно, то
$\displaystyle{| 2^M | = 2^{| M |}}$.
\end{theorem}
\begin{proof}
Индукция по $\vert M\vert$. База: если $\vert M\vert = 0$, то $M =
\emptyset$ и $2^{\emptyset} = \{\emptyset\}$. Пусть
$\A{M}{\vert M\vert<k \Impl \vert 2^M\vert = 2^{\vert
M\vert}}$. Рассмотрим $M = \{\Tuple{a}{1}{k}\}$, $\vert M\vert =
k$. Положим
$
{\cal M}_1\Assign\Set{X\in 2^M}{a_k\notin X}$
и 
${\cal M}_2\Assign\Set{X\in 2^M}{a_k\in X}.
$
Имеем $2^M = {\cal M}_1 \cup {\cal M}_2$,\\${\cal M}_1 \cap {\cal
M}_2 = \emptyset$, ${\cal M}_1 \sim {\cal M}_2$ за счёт
взаимно-однозначного соответствия $X\leftrightarrow X+a_k$ и ${\cal
M}_1\sim 2^{\{\Tuple{a}{1}{k-1}\}}$. По индукционному
предположению $\vert 2^{\{\Tuple{a}{1}{k-1}\}}\vert=2^{k-1}$, и,
значит, $\vert{\cal M}_1\vert = \vert{\cal M}_2\vert = 2^{k-1}$.
Следовательно, $\vert 2^M\vert = \vert {\cal M}_1\vert + \vert
{\cal M}_2\vert = 2^{k-1} + 2^{k-1} = 2^k $.
\end{proof}

\begin{examples}
$2^\emptyset = \{\emptyset\}$;
$|2^\emptyset| = 2^{|\emptyset|}=2^0=1$.\\
$2^{\{1\}} = \{\emptyset, \{1\}\}$;
$|2^{\{1\}}| = 2^{|{\{1\}}|}=2^1=2$.
\end{examples}

\medskip
Применение
операций к подмножествам универсума не выводит за его пределы.

Множество всех подмножеств множества $U$ с операциями пересечения,
объединения, разности и дополнения образует \odmkey{алгебру
подмножеств}{Алгебра (algebra)!подмножеств (of subsets)} множества
$U$.
\end{frame}

\begin{frame}
\frametitle{1.2.8. Булеан (2/4)}
В булеане $2^M$ естественно выделяются семейства 
подмножеств $\cal{C}^k (M)$, состоящие из подмножеств множества $M$, 
имеющих одинаковую мощность $k$, то есть \\$\cal{C}^k (M) =\Set{S\subset M}{ k=|S| }$. Такие семейства называются 
\odmkey{семействами равномощных подмножеств}
{Семейство (family)! равномощных подмножеств (of equipotent subsets)}. 
Заметим, что семейства равномощных подмножеств, 
за исключением \\$\cal{C}^0 (M)$, являются покрытиями множества $M$, 
а $\cal{C}^1 (M)$~--- семейство одноэлементных подмножеств множества $M$~---
является разбиением. Кроме того, семейства $\cal{C}^k (M)$ не пересекаются. 
Если $|M|=n$, то очевидно, что

\vspace{-8pt}
$$
2^M=\bigcup_{k=0}^n \cal{C}^k (M)\quad   \text{ и }
\quad     2^n=\sum_{k=0}^n |\cal{C}^k (M) |.
$$

\vspace{-2pt}
\begin{remark}
В \DMref{п.~5.1.5.}{5.1.5.} показано, что мощность семейства равномощных 
подмножеств $\cal{C}^k (M)$ является биномиальным коэффициентом, 
$|\cal{C}^k (M)|=C(n,k)$. 
\end{remark}

\begin{example}
Множество всех подмножеств множества $\{1, 2, 3\}$
содержит одно несобственное подмножество $\emptyset$,
семейство из трёх одноэлементных подмножеств 
$\{1\}, \{2\}, \{3\}$,
семейство из трёх двухэлементных подмножеств
$\{1,2\}, \{2,3\}, \{1,3\}$
и одно несобственное подмножество $\{1,2,3\}$. 
\end{example} 

\end{frame}

\begin{frame}
\frametitle{1.2.8. Булеан (3/4)}

Неравномощные множества можно сравнивать.
Если множества $A$ и  $B$ неравномощны, $|A| \ne |B|$,
и при этом в множестве $B$ существует собственное 
подмножество $C$, $C\SubsetNE B$, такое что $|A|=|C|$,
то говорят, что множество $B$ 
\odmkey{мощнее}{Мощность (cardinality)! строго больше (strictly bigger)} множества
$A$ и записывают это так: $|A|<|B|$.
Более формально:
$$
|A|<|B| \Def A\hookrightarrow B \And |A|\ne|B|.
$$


\begin{theorem} \NB{ [Кантора]}
Булеан множества мощнее множества: $|A|<|2^{A}|$.
\end{theorem}
\begin{proof}
От противного. Пусть $A$~--- некоторое множество и 
существует взаимно\-однозначное
соответствие $A \stackrel{f}{\sim} 2^A$.
Положим $Y\Assign \Set{x\in A}{x\notin f(x)}$.
Имеем $Y\in 2^A$, значит $\E{y\in A}{f(y)=Y}$.
Рассмотрим элемент $y$ и множество $Y$.
Если $y\in Y$, то $y\notin Y$, а если $y \notin Y$, то
$y\in Y$~--- противоречие. Следовательно,
соответствие $f$ не существует, 
а значит множество и его булеан неравномощны.   
Но множество $A$ равномощно множеству одноэлементных
подмножеств, которое является собственным
подмножеством булеана, $A\sim \cal{C}^1 (A)$.
\end{proof}

\begin{remark}
Иногда используют обозначение $|A|\le|B|$,
полагая $|A|\le|B| \Def A\hookrightarrow B$. 
\end{remark}
\end{frame}

\begin{frame}
\frametitle{1.2.8. Булеан (4/4)}

%\vspace{-2pt}
\begin{corollary}
Не существует множества всех множеств.
\end{corollary}

\begin{proof} От противного.
Пусть множество \em{всех} множеств существует, 
обозначим его $U$.
Тогда $U\sim \cal{C}^1(U)\subset 2^U$, и значит
$U\hookrightarrow 2^U$. 
С другой стороны, все элементы булеана 
также являются множествами,
$\A{X\in 2^U}{X\in U}$, то есть 
$2^U\subset U$, и значит $2^U\hookrightarrow U$.
По теореме Кантора--Бернштейна имеем $U\sim 2^U$, что противоречит 
 $|U|<|2^U|$.  
\end{proof}

%\medskip
\begin{remark}
Это следствие известно как \odmkey{парадокс Кантора}{Парадокс (paradox)!Кантора (Cantor)}.
\end{remark}

\begin{example} 
Множества всех 
\em{ конечных} множеств также
не существует. Действительно,
пусть такое множество $M$ существует.
Тогда $\A{A}{A\hookrightarrow M}$: $a\mapsto \{a\}$,
$a\in A$, $\{a\}\in M$, и значит $|A|\le|M|$.
Если взять $A\Assign 2^M$, то
$|2^M|\le|M|$,
что противоречит теореме Кантора.
\end{example}

%\medskip
\begin{digress}
При обобщении любой теории неизбежно 
возникают трудные вопросы, 
на которые надобно ответить 
для обоснования обобщения.   
Разобранный парадокс Кантора является примером 
ответа на такой вопрос в наивной теории множеств. 
Другие примеры трудных вопросов: существует ли надмножество всех множеств? Существует ли множество максимальной мощности? Существует ли множество $Y$, такое что
$|X|<|Y|<|2^X|$ для бесконечного 
множества $X$? (\odmkey{обобщённая континуум-гипотеза}{Обобщённая континуум-гипотеза \\(generalized continuum-hypothesis)}). При поиске ответа на трудные вопросы следует оценить,
окупит ли найденный ответ затраты на решение вопроса. 
\end{digress}

\end{frame}

\subsubsection{1.2.9. Свойства операций над множествами}\label{1.2.9.}
\begin{frame}
\frametitle{1.2.9. Свойства операций над множествами (1/4)}
\begin{theorem} Пусть $A, B, C \subset U\ne \emptyset$. Тогда
операции над множествами обладают \\следующими свойствами:
\begin{tabbing}
11)
\= выражение для разности: $qquad$\=
$A \setminus B = A \cap \overline{B}$.\kill 
1) 
\> \odmkey{идемпотентность}{Идемпотентность (idempotency)!пересечения
(of intersection)}:
\index{Идемпотентность (idempotency)!объединения (of union)} \>
$A \cup A = A, \quad A \cap A = A$; \\
2)
\>\odmkey{коммутативность}{Коммутативность
(commutativity)!пересечения
(of intersection)}:
\index{Коммутативность (commutativity)!объединения (of union)}\> 
$A \cup B = B \cup A, \quad A \cap B = B \cap A$;\\
3)
\>\odmkey{ассоциативность}{Ассоциативность
(associativity)!пересечения
(of intersection)}:
\index{Ассоциативность (associativity)!объединения (of union)}\> 
$A \cup (B \cup C) = (A \cup B) \cup C$,\\
\> \> $A \cap (B \cap C) =
(A \cap B) \cap C$;\\
4)
\>\odmkey{дистрибутивность}
{Дистрибутивность (distributivity)!пересечения относительно
объединения (of intersection over union)}:
\index{Дистрибутивность (distributivity)!объединения относительно пересечения (of union over intersection)} \>
$A \cup (B \cap C) = (A \cup B) \cap (A \cup C)$, \\
\> \> $A \cap (B \cup C) = (A \cap B) \cup (A \cap C)$;\\
5)
\>\odmkey{поглощение}{Поглощение (absorption)}: 
\>$(A \cap B) \cup A = A, \quad (A \cup B) \cap A = A$;\\
6)
\> свойства нуля:\> 
$A \cup \emptyset =A,\quad A \cap \emptyset = \emptyset$;\\
7)
\> свойства единицы:\>
$A \cup U = U, \quad A \cap U = A$;\\
8)
\> \odmkey{инволютивность}{Инволютивность (involutivity)!дополнения
(of a complement)} дополнения:\>
$\overline{\overline{A}} = A$;\\
9)
\>\odmkey{законы де Моргана}{Законы де Моргана (de Morgan laws)}\>
$\overline{A \cap B} = \overline{A} \cup \overline{B}, \quad
\overline{A \cup B} = \overline{A} \cap \overline{B}$;\\
10)
\> свойства дополнения:\>
$A \cup \overline{A} = U, \quad A \cap \overline{A} = \emptyset$;\\
11)
\> выражение для разности:\>
$A \setminus B = A \cap \overline{B}$.
\end{tabbing}
\end{theorem}


\end{frame}


\begin{frame}
\frametitle{1.2.9. Свойства операций над множествами (2/4)}
\begin{proof}
Для первого равенства:
$
x\in A\cup A \Equiv x\in A \Or x\in A \Equiv x\in A.
$
Аналогично и для остальных
равенств.
\end{proof} 

\begin{digress}
Чтобы наглядно усмотреть 
справедливость указанных равенств, очень полезно нарисовать соответствующие диаграммы 
Венна.
\end{digress}

\medskip
Введённые операции над множествами~--- 
объединение, пересечение, разность и 
симметрическая разность~---
примечательны тем, что
определены для любых множеств,
поскольку элементы множества-результата
берутся из множеств-операндов.


\begin{remark}
Если наложить ограничение,
что в результат операции над множествами
входят только элементы операндов,
то набор из четырёх введённых операций 
оказывается исчерпывающим.
\end{remark}

\medskip
Однако в алгебре подмножеств любого множества $U$
%,
%то есть в булеане $2^U$, 
возможно определять такие операции,
в результаты которых входят элементы
множества $U$, не входящие в операнды.

\begin{example}
Результат операции дополнения 
$\overline{A}\Def\Set{a\in U}{a\not\in A}$
состоит из элементов,
которые не входят в операнд. 
\end{example}


\end{frame}


\begin{frame}
\frametitle{1.2.9. Свойства операций над множествами (3/4)}

Рассмотрим ещё две операции: $\mid$ и 
$\downarrow$ над  подмножествами множества  
$U$:
$$
A\mid B \Def \Set{x\in U}{x\not\in A \Or
x\not\in B },\quad
A\downarrow B \Def \Set{x\in U}{x\not\in A \And
x\not\in B }.
$$

\begin{remark}
В алгебре подмножеств эти операции обычно не
рассматриваются, а потому не именуются.
Аналогичные операции в алгебре логики
называются \odmkey{штрих Шеффера}{Штрих Шеффера (Scheffer stroke)} и 
\odmkey{стрелка Пирса}{Стрелка Пирса (Pierce arrow)}, соответственно.
\end{remark}  

Нетрудно показать, что введённые здесь операции
выражаются через привычные операции над множествами:
$$
A\mid B = U\setminus (A\cap B)=\overline{A\cap B}=\overline{A}\cup\overline{B}, \quad
A\downarrow B = U\setminus (A\cup B) =\overline{A\cup B}=\overline{A}\cap\overline{B}.
$$   

Это наблюдение позволяет заметить,
что некоторые нетривиальные свойства привычных 
операций выполняются и для вновь введённых 
операций, например, законы де Моргана
$$A\mid B =\overline{\overline{A} \downarrow\overline{B}},
\quad
A\downarrow B =\overline{\overline{A} \mid\overline{B}}.$$

\vspace{-2pt}
\begin{remark}
Если определена инволютивная операция дополнения,
то операции, для которых выполняются 
законы де Моргана, называются
\odmkey{двойственными}{Двойственность (duality)}.
Двойственными являются пересечение и объединение,
штрих Шеффера и стрелка Пирса.
\end{remark} 

\end{frame}


\begin{frame}
\frametitle{1.2.9. Свойства операций над множествами (4/4)}

Все традиционные операции над множествами непосредственно выражаются как через штрих Шеффера, так и через стрелку Пирса 
(см. также \DMref{п.~3.6.2}{3.6.2.}):
$$
\overline{A} = \overline{A}\cup \overline{A} =
  A\mid A = \overline{A}\cap \overline{A} = A\downarrow A,
$$   
$$
A\cup B = (A\mid A)\mid(B\mid B) =
 (A\downarrow B)\downarrow (A\downarrow B),
$$   
$$
A\cap B = (A\mid B)\mid(A\mid B) =
 (A\downarrow A)\downarrow (B\downarrow B),
$$
$$
A\setminus B = A\cap\overline{B} =
 (A\mid (B\mid B))\mid (A\mid (B\mid B)) =
 (A\downarrow A)\downarrow B,
$$
$$
A\triangle B =
(A\setminus B)\cup(B\setminus A)=
(((A\!\downarrow\! A)\!\downarrow\! B)
\!\downarrow\! 
((B\!\downarrow\! B)\!\downarrow\! A)) \!\downarrow\!
(((A\!\downarrow \!A)\!\downarrow\! B) \!\downarrow\! 
((B\!\downarrow\! B)\!\downarrow\! A)).  $$

\begin{digress}
Отмеченные свойства бывают присущи различным операциям, в частности~--- арифметическим операциям с числами.
Однако операции над множествами отличаются
от числовых операций в следующем
существенном аспекте. 
Образно говорят, 
что бинарная операция $\diamond$
обладает памятью, если по результату 
операции и одному из операндов можно 
восстановить второй операнд.
Другими словами,  
если уравнение $a\diamond x = b$
имеет однозначное решение,
выражаемое через $a$ и $b$.
Например, в арифметике, 
если $a+x=b$, то $x=b-a$.
Однако, если 
$A\cup X =B$, то в общем случае можно
утверждать, что $B\setminus A\subset X \subset B$, но не более того. 
\end{digress}
    
\end{frame}

\begin{frame}
\frametitle{Онтология наивной теории множеств}
\begin{center}
\includegraphics[height=7.8cm]{1_2_A.jpg}
\end{center}
\end{frame}

\begin{frame}
\frametitle{Онтология наивной теории множеств}
\begin{center}
\includegraphics[height=7.8cm]{1_2_C.jpg}
\end{center}
\end{frame}

\begin{frame}
\frametitle{1.3. Представление множеств в программах (1/3)}
\subsection{1.3. Представление множеств в программах}
Термин <<представление>> применительно к программированию означает
следующее.

\medskip
Представить в программе какой-либо объект (в данном
случае множество)~--- это значит описать в терминах системы
программирования структуру данных, используемую для хранения
информации о представляемом объекте, и алгоритмы над выбранными
структурами данных, которые реализуют присущие данному объекту
операции. 

\medskip
Следует подчеркнуть, что, как правило, один и тот же объект может
быть представлен многими разными способами, причём нельзя указать
способ, который является наилучшим для всех возможных случаев. В
одних случаях выгодно использовать одно представление, а в
других~--- другое. Умение выбрать наилучшее для данного случая представление
является основой искусства практического программирования. Хороший
программист отличается тем, что он знает \em{много} разных
способов представления и умело выбирает наиболее подходящий.

\medskip
Целью этого раздела является обсуждение 
трёх существенно различных представлений множеств в программах.

\end{frame}

\begin{frame}
\frametitle{1.3. Представление множеств в программах (2/3)}
\begin{tabular}{p{12mm}|p{40mm}|p{92mm}}
  № & Название параграфа
  & Ключевые термины и обозначения \\
  \hline
\DMref{1.3.1.}{1.3.1.} & Битовые шкалы 
  & 
  $\And$~--- конъюнкция;
  $\Or$~--- дизъюнкция;
  $+_2$~--- сложение по модулю два;
  $\Not$~--- инвертирование; бит 
   \\
\DMref{1.3.2.}{1.3.2.} & Представление натуральных чисел 
  & 
 система счисления; основание;
 цифра
  \\
\DMref{1.3.3.}{1.3.3.} & Перебор подмножеств множества &
$B(i)$~--- двоичный код числа $i$;
$Q(i)$~--- число битов, в которых 
различаются $B(i-1)$ и $B(i)$ \\
\DMref{1.3.4.}{1.3.4.} \RS & Алгоритм построения кода Грея &
бинарный код Грея;
левый зеркальный код Грея; рекуррентные вычисления 
\\
\DMref{1.3.5.}{1.3.5.} & Представление множеств списками &
Алгоритм слияния
\end{tabular}
\end{frame}

\begin{frame}
\frametitle{1.3. Представление множеств в программах (3/3)}
\begin{tabular}{p{12mm}|p{40mm}|p{92mm}}
  № & Название параграфа
  & Ключевые термины и обозначения \\
  \hline
\DMref{1.3.6.}{1.3.6.} & Проверка включения слиянием 
  & 
  \\
\DMref{1.3.7.}{1.3.7.} & Вычисление объединения слиянием 
  & 
\\
\DMref{1.3.8.}{1.3.8.} & Вычисление пересечения слиянием 
  & 
\\
\DMref{1.3.9.}{1.3.9.} \RS & Представление множеств итераторами &
Первичные операции; алгоритм пересечения итераторов;
алгоритм вычитания итераторов;
алгоритм дизъюнктного объединения итераторов;
оператор структурного перехода
\end{tabular}
\end{frame}


\begin{frame}
\label{1.3.1.}\subsubsection{1.3.1. Битовые шкалы }\frametitle{1.3.1. Битовые шкалы (1/3)}

Булевский массив \textbf{array} $[1..n]$ \textbf{of} $0..1$ называют 
\odmkey{битовой шкалой}{Битовая шкала (bit scale)} (или \odmkey{двоичным вектором}{Вектор (vector)!двоичный (binary)}, 
или \odmkey{машинным словом}{Слово (word)!машинное (machine)}, или \odmkey{двоичным кодом}{Код (code)!двоичный (binary)}) длины $n$. 
Элементы битовой шкалы называют \odmkey{разрядами}{Разряд (place)} или \odmkey{битами}{Бит (bit)}. 
Значения разрядов обозначают цифрами $0$ и $1$
и считают, что к ним применимы следующие операции:%бинарные операции: 
%\odmkey{конъюнкция}{}, или \odmkey{логическое умножение}{} (обозначается $\And$);
%\odmkey{дизъюнкция}{}, или \odmkey{логическое сложение}{} (обозначается $\Or$); 
%\odmkey{сложение по модулю 2}{} или \odmkey{исключающее или}{} 
%(обозначается $+_2$, причём $1 +_2 1=0$); а также унарная операция 
%\odmkey{логическое отрицание}{}, или \odmkey{инвертирование}{} (обозначается $\Not$).

\medskip
\begin{tabular}{| l | c |}
\hline Название операции & Обозначение \\\hline
\odmkey{конъюнкция}{Конъюнкция (conjunction)}, или \odmkey{логическое умножение}{Умножение (multiplication)!логическое (logical)} & $\And$ \\\hline
\odmkey{дизъюнкция}{Дизъюнкция (disjunction)}, или \odmkey{логическое сложение}{Сложение (addition)!логическое (logical)} & $\Or$ \\\hline
\odmkey{сложение по модулю 2}{Сложение (addition)!по модулю 2 (by modulo 2)}, или \odmkey{исключающее или}{Исключающее ИЛИ (exclusive or)} & $+_2$ \\\hline
\odmkey{логическое отрицание}{Логическое отрицание (logical negation)}, или \odmkey{инвертирование}{Инвертирование (inversing)} & $\Not$ \\\hline
\end{tabular}

\medskip
\begin{remark}
В большинстве компьютеров для некоторых характерных значений
$n$ (сейчас наиболее распространены случаи $n=32$, $n=64$)  время выполнения поразрядных операций не зависит от длины машинного слова $n$ (выполняется за один такт). 
\end{remark} 

\end{frame}

\begin{frame}
\frametitle{1.3.1. Битовые шкалы (2/3)}
%\vspace{-2pt} 
Пусть задано конечное множество $U=\{\Tuple{u}{1}{n}\}$.
Подмножество $A$ множества $U$ представляется кодом\odmkey{}{Код
(code)!множества (of set)}  \textbf{$C:$ array $[1..n]$ of $0..1$}, в
котором:
$$ \A{i\in
1..n}{C[i]=u_i\in A}. $$ 
Тогда
код пересечения множеств $A$ и $B$ есть поразрядная конъюнкция
кода множества $A$ и кода множества $B$,
код объединения множеств $A$ и $B$ есть поразрядная дизъюнкция
кода множества $A$ и кода множества $B$,
код симметрической разности множеств $A$ и $B$ есть поразрядное сложение по модулю 2
кода множества $A$ и кода множества $B$,
код дополнения множества $A$ есть поразрядное инвертирование
кода множества $A$.

\medskip
Для вычисления разности множеств обычно  пользуются равенством $A \setminus B = A \cap \overline{B}$
(\DMref{п.~1.2.9}{1.2.9.}), а включение множеств проверяют, пользуясь свойством
$A\subset B\Equiv A=A\cap B$.
Для проверки поразрядного равенства шкал в компьютерах также есть машинная команда.    
Таким образом, в этом представлении \em{все} операции
над множествами 
выполняются весьма эффективно.


\end{frame}

\begin{frame}
\frametitle{1.3.1. Битовые шкалы (3/3)} 
%\vspace{-4pt}
\begin{remark}
Если мощность множества превосходит размер машинного слова, но не очень велика,
то для представления подмножеств используются массивы битовых шкал. В этом случае
операции над множествами реализуются с помощью циклов по элементам массива.
\end{remark}

%\medskip
Эту же идею можно использовать для представления мультимножеств.

Если $\widehat{X} =\left[x_1^{a_1},\dots,x_n^{a_n}\right]$~---
мультимножество над множеством $X=\{\Tuple{x}{1}{n}\}$, то
массив \textbf{$A:$ array $[1..n]$ of integer}, где $\A{i\in
1..n}{A[i]=a_i}$, является представлением мультимножества
$\widehat{X}$.

\begin{remark}
Представление индикатора совпадает с представлением его состава.
\end{remark}

%\medskip
\begin{theorem}
Существует $2^n$ различных битовых шкал длины $n$.
\end{theorem} 
\begin{proof}
Рассмотрим любое множество $M$ из $n$ элементов и перенумеруем элементы. 
Тогда каждое подмножество множества $M$ представляется единственным образом битовой шкалой, 
и каждая битовая шкала длины $n$ является представлением 
определённого подмножества множества $M$.
Таким образом, множество битовых шкал длины $n$ и 
множество всех подмножеств $n$-элементного множества равномощны, 
откуда по первой теореме \DMref{п.~1.2.8}{1.2.8.}
 следует требуемый результат. 
\end{proof}
\end{frame}

\begin{frame}
\label{1.3.2.}\subsubsection{1.3.2. Представление натуральных чисел}
\frametitle{1.3.2. Представление натуральных чисел (1/8)}
%\vspace{-2pt}
\begin{lemma}
\NB{ [1]} $\A{k\in\mathbb{N}}{\E{ !\, m\ge 0}{2^m \le k<2^{m+1}}}.$
%Для любого натурального числа $k$ 
%существует единственное целое неотрицательное число
%$m\ge 0$ такое, что $2^m \le k<2^{m+1}$.
\end{lemma}
\begin{proof}
Положим $m \Assign \lfloor \log k \rfloor $.
%($\log x\Def \log_2 x$, см. 0.2.3).
\end{proof}

\medskip
\begin{lemma} \NB{ [2]}   					
$\sum_{i=0}^m 2^i = 2^{m+1}-1.$
\end{lemma}
\begin{proof} Cуммa геометрической прогрессии.	
\end{proof}

\medskip
\begin{theorem}
Для любого натурального числа $k$ существует единственное множество \\
$B(k)=\{\Tuple{m}{1}{s} \}$ 
такое, что $k=\sum\limits_{i=1}^s 2^{m_i}.$ 
\end{theorem} 
\begin{remark}
$\A{i\in 1..s}{m_i\in \Bbb{N}_0}$.
\end{remark}
\begin{proof} 
Индукция по числу $m$, которое существует для любого натурального $k$ 
по лемме 1.
При $m=0$ имеем $k=1=2^0$. Пусть для всех чисел
$j<2^m$ существует единственное множество 
$B(j)=\{\Tuple{m}{1}{s}\}$ такое,
что $j=\sum_{i=1}^s 2^{m_i}$. 
Рассмотрим число $k$ такое, что 
$2^m\!\le\!k<2^{m\!+\!1}$.
Если $k\!=\!2^m$, то множество
$B(2^m)\!=\!\{m\}$ из одного элемента является требуемым.
Если же число $k>2^m$, то рассмотрим число $j\Assign k-2^m$. Имеем $j<2^m$. 
По индукционному предположению 
$\E{!B(j)=\{\Tuple{m}{1}{s}\}}{j=\sum_{i=1}^s 2^{m_i}}$.
Но тогда добавим элемент $m$ и 
получим требуемое множество $B(k)= \{\Tuple{m}{1}{s},m\}$.
\end{proof}
\end{frame}

\begin{frame}
\frametitle{ 1.3.2. Представление натуральных чисел (2/8)}
\vspace{-2pt}
\begin{corollary} \NB{ [1]}
Любое натуральное число $k$ такое, что $2^m \le k < 2^{m+1}$, 
можно представить в виде линейной комбинации степеней числа 2:
$k=\sum_{i=0}^m a_i 2^i$, 
где коэффициенты $a_i$ являются двоичными цифрами 0, 1, причём $a_m=1$.
\end{corollary} 
\begin{proof}
% Рассмотрим число $k$ такое, что $2^m \le k<2^{m+1}$. 
По теореме $\E{!\ B(k)=\{\Tuple{m}{1}{s} \}}{k=\sum_{i=1}^s 2^{m_i}}$. 
Не ограничивая общности можно считать, что $m_1 < \ldots < m_s$. 
По лемме 2  
$m_s=m$, иначе суммы степеней не достанет до числа $k$. 
Положим коэффициент $a_i \Assign 1$, если $i \in B(k)=\{\Tuple{m}{1}{s} \}$, 
и положим $a_i \Assign 0$, если $i\notin B(k)$.
\end{proof}
\begin{corollary} \NB{ [2]} 
Существует взаимно-однозначное соответствие между битовыми шкалами длины $n$ 
и целыми неотрицательными числами $k$: $0 \le k \le 2^n-1$. 
\end{corollary}

\begin{proof} 
Одно из возможных соответствий можно установить следующим образом.
Поставим степени числа $2$ в соответствие разрядам двоичной шкалы
\textbf{$B:$ array$[1..n]$ of $0..1$}. 
Пусть $B[1] \mapsto 2^{n-1}$,  
$B[2] \mapsto 2^{n-2}$, 
и так далее, 
$B[n-1] \mapsto 2^1$, 
$B[n] \mapsto 2^0$. 
Сами значения разрядов шкалы будем считать коэффициентами линейной комбинации из следствия 1, причём $a_i=B[n-i]$. 
Тогда всякому неотрицательному числу
из указанного диапазона соответствует 
линейная комбинация и 
битовая шкала,
и обратно,
всякой битовой шкале соответствует число.
\end{proof}

\end{frame}


\begin{frame}
\frametitle{1.3.2. Представление натуральных чисел (3/8)}
Указанный  способ интерпретации двоичных шкал называется 
\odmkey{двоичным представлением}{Представление (representation)!двоичное (binary)}, или \odmkey{двоичным кодом}{Код (code)!двоичный (binary)} числа. 
При этом число 0 представляется кодом, состоящим из нулей; число $2^k$ 
представляется кодом, содержащим ровно одну единицу в $(k+1)$-ом (считая справа)
разряде; число $2^k-1$ представляется кодом, состоящим из $k$ единиц, считая справа. 
Следующий алгоритм строит двоичный код числа. 
В этом разделе двоичный код числа обозначается буквой $B_n$, 
причём, если длина кода $n$ ясна из контекста, то нижний индекс опускается. 

\begin{algorithmic}
\REQUIRE натуральное число $k$:  $0\le k \le 2^n-1$.
\ENSURE битовая шкала \textbf{$B$: array$[1..n]$ of $0..1$}.
%~---
%двоичное представление числа $k$.
\FOR{\textbf{$i$ from $n$ downto $1$}}
\STATE  $B[i] \Assign k \MOD 2$
\COMMENT {\small\color{gray} очередной бит}
\STATE  $k \Assign k \DIV 2$
\COMMENT {\small\color{gray} оставшееся число}
\ENDFOR
\end{algorithmic}
%\end{algorithm}

%\vspace{-0.25\baselineskip}
\begin{ground}
Операция $k \MOD 2$ в любом случае дает 0 или 1, так что $B$ 
действительно является битовой шкалой длины $n$. 
Заметим, что 
$\A{a\in \mathbb{N}}{a= 2(a \DIV 2 )+(a \MOD 2)}$.
\end{ground}
\end{frame}

\begin{frame}
\frametitle{ 1.3.2. Представление натуральных чисел (4/8)}
\begin{ground} (Окончание)
Таким образом, если вырезка первых $n\!-\!1$ разрядов из битовой шкалы
\textbf{ $B_{n\!-\!1}$:~array~$[1..(n\!-\!1)]$~of~$0..1$} является двоичным кодом числа
$ (k \DIV 2)$, то и вся шкала $B_n$ является двоичным кодом числа $k$, 
и последний разряд определён правильно. 
Далее, применяя это рассуждение к числу $(k \DIV 2)$, 
к числу $((k \DIV 2) \DIV 2) = k \DIV 2^2$ и т.~д., 
получаем, что все левые отрезки 
битовой шкалы \textbf{$B_{n-i}$ : array $[1..n-i]$ of $0..1$} 
являются двоичными кодами чисел $k \DIV 2^i$, 
а все правые отрезки битовой шкалы \textbf{array $[n-i+1..n]$ of $0..1$} 
являются двоичными кодами чисел $k \MOD 2^i$.
\end{ground}


\begin{example}
Построение двоичного кода числа <<десять>>. 


\begin{tabular}{|l|l|l|l|l||c|}
\hline
$i$ & $k$ & $k \DIV 2$ & $k \MOD 2$ & $B [i..n]$ &	$\sum\limits_{j=i}^n B[j] 2^{n-j}$ \\ 
\hline
4 & 10 & 5 & 0 & $B[4..4] = ~~~0$ & 0 \\
3 &	5 &	2 & 1 & $B[3..4] = ~~10$ & 2 \\
2 & 2 & 1 & 0 & $B[2..4] = ~010$ & 2 \\
1 &	1 & 0 & 1 & $B[1..4] = 1010$ &	10 \\
\hline
\end{tabular}

\end{example}
\end{frame}

\begin{frame}
\frametitle{ 1.3.2. Представление натуральных чисел (5/8)}
\begin{remark}
Обратный алгоритм вычисления числа по его двоичному коду и связанные с этим полезные формулы см. в \DMref{п.~3.6.1}{3.6.1.}.
\end{remark}

\medskip
\begin{digress}
Здесь считается, 
что коэффициент $a_i$ при $2^{n-i}$ хранится в разряде с индексом 
$n-i$, т.~е. в первом разряде хранится коэффициент при старшей степени, 
а в последнем разряде хранится коэффициент при младшей степени. 
Такой порядок называется <<от старшего к младшему>> (big-endian по-английски) 
и используется во многих случаях, например, при обычной записи чисел в десятичной системе, 
в протоколе TCP/IP, в процессорах компаний IBM и Motorola. 
Однако применяются и другие порядки хранения коэффициентов. 
Так, в процессорах архитектуры X86 компании Intel применяется противоположный порядок --- little-endian.
Аналогичным способом можно представлять отрицательные целые, 
используя разряд знака,
а также рациональные числа,  используя отрицательные степени
числа 2 
($2^{-1}=\frac{1}{2}$, $2^{-2}=\frac{1}{4}$, $2^{-3}=\frac{1}{8}, \dots$)  
в качестве значений разрядов дробной части числа.
Для представления отрицательных чисел можно использовать 
дополнительный разряд знака и~т.~д.
Важно осознать, что числа \em{возможно} представлять 
в программах битовыми шкалами, 
и такое представление оказывается вполне удовлетворительным с практической точки зрения.
\end{digress}

\end{frame}

\begin{frame}
\frametitle{ 1.3.2. Представление натуральных чисел (6/8)}
%\vspace{-2pt}
\begin{remark}
Все утверждения этого параграфа 
обобщаются на случай позиционной
системы счисления с \odmkey{основанием}{Основание (radix)} $b$. 
\end{remark}

\medskip
\begin{digress}
Известно множество  различных представлений чисел в программах. 
\odmkey{Представления чисел}{Представление (representation)!чисел (of numbers)}, или \odmkey{системы счисления}{Система счисления (number language, notation, numeration, number system)}, 
то есть правила записи чисел с помощью системы знаков, 
которые в этом случае называют \odmkey{цифрами}{Цифра (digit)}, принято делить на три класса. 
К первому классу относятся \odmkey{позиционные}{Система счисления (number language, notation, numeration, number system)!позиционная (positional)} системы счисления с определённым основанием, 
в которых число представляется как линейная комбинация степеней основания 
(например: двоичная, десятичная). 
Второй класс образуют \odmkey{смешанные}{Система счисления (number language, notation, numeration, number system)!смешанная (mixed)} системы счисления, 
в которых число также представляется линейной комбинацией, 
но вместо степеней основания используется иной возрастающий ряд чисел 
(например: исчисление времени --- секунды, минуты, часы; 
\odmkey{факториальная}{Система счисления (number language, notation, numeration, number system)!факториальная (factorial)} система счисления; \odmkey{фибоначчиева}{Система счисления (number language, notation, numeration, number system)!фибоначчиева (fibonacci)} система счисления). 
В системах третьего класса
значение цифры в числе определяется не позицией в записи, 
а иным способом (например: римская система счисления, \odmkey{биномиальная}{Система счисления (number language, notation, numeration, number system)!биномиальная (binomial)} система счисления, 
система \odmkey{остаточных классов}{Система счисления (number language, notation, numeration, number system)! остаточных классов (of residual classes)}). 
\end{digress}


\end{frame}

\begin{frame}
\frametitle{ 1.3.2. Представление натуральных чисел (7/8)}

\begin{digress} (Продолжение)
Двоичный код в программах является самым распространённым, 
но не потому, что этот код <<самый лучший>>. 

Например, двоичная система кажется самой экономной с точки зрения расхода памяти, 
но это не так.
Действительно, 
пусть $b$~--- основание системы счисления,
тогда для записи \em{всех} не более чем $k$-разрядных
чисел в этой системе потребуется $s=kb$ знаков
(экземпляров цифр). Например, для записи всех
трёхразрядных десятичных чисел от $0$ до $999$ потребуется $30$ знаков. 
\odmkey{Экономичностью системы счисления}{Экономичность системы счисления (radix efficiency)}
называется максимальное количество чисел $N$,
которые возможно записать при заданном количестве знаков $s$. В системе с основанием $b$ с помощью
$k$ разрядов можно записать $N=b^k =b^{s/b}$ чисел.    Например, тридцатью знаками
в десятичной системе можно записать $10^3=1000$ чисел,  
а в двоичной системе можно записать $2^{15}=32768$ чисел. Двоичная система явно экономичнее десятичной. 
Средствами математического анализа нетрудно
найти максимум выражения $b^{s/b}$ при заданном числе знаков $s$. Максимум достигается 
при $b=e$. 
Число $3$ ближе к числу $e$, чем число $2$, поэтому троичная система несколько экономичнее двоичной. 
\end{digress}     

\end{frame}

\begin{frame}
\frametitle{ 1.3.2. Представление натуральных чисел (8/8)} 
%\vspace{-2pt}
\begin{digress} (окончание) 
Арифметические операции над двоичными числами кажутся достаточно эффективными. 
Сложение выполняется за время $O(n)$, 
а умножение~--- 
 за $O(n^2)$, если умножать <<в столбик>>,  
и несколько быстрее, если применять современные изощрённые алгоритмы. 
Умножение труднее сложения, это кажется законом природы, 
но это не так --- в системах остаточных классов 
обе эти операции выполняются за линейное время. 
Вывод из этих наблюдений состоит в том, что в каждой математической модели для конкретной предметной области необходимо анализировать фактические требования к системе счисления 
и выбирать наиболее подходящую, а не довольствоваться той, которая предоставлена системой программирования по умолчанию.
\end{digress}


\medskip
{\large
Таким образом, между множеством всех битовых шкал длины $n$, 
множеством всех подмножеств множества $\{\Tuple{a}{1}{n}\}$ 
и множеством всех целых
чисел из диапазона $0..(2^n-1)$ 
установлены \em{взаимно-однозначные соответствия\/}, что позволяет представлять подмножества числами, а также представлять как числа, так и множества битовыми шкалами. 
}
\end{frame}

\begin{frame}
\label{1.3.3.}\subsubsection{1.3.3. Перебор подмножеств множества }\frametitle{1.3.3. Перебор подмножеств множества (1/4)}
%\vspace{-4pt}
Во многих переборных алгоритмах требуется последовательно
рассмотреть все подмножества заданного $n$-элементного множества, 
то есть требуется рассмотреть все возможные битовые шкалы длины $n$.
Пусть $B(i)$~--- это функция, доставляющая двоичный код числа $i$. 
Тогда следующий \odmkey{алгоритм перечисления подмножеств}{Алгоритм (algorithm)!генерации всех подмножеств (generating of all subsets)}  перечисляет все подмножества 
$n$-элементного множества.
%\begin{algorithm}
%\caption{Генерация всех подмножеств $n$-элементного множества}
%\index{Алгоритм (algorithm)!генерации всех подмножеств (generation
%of all subsets)} \label{simplesets}

%\vspace{-4pt}
\begin{algorithmic}
\REQUIRE число $n \geq 0$~--- мощность множества.
\ENSURE последовательность кодов подмножеств $B(i)$.
\STATE \textbf{for $i$ from $0$ to $2^n-1$ do yield $B(i)$ end for}
\end{algorithmic}
%\end{algorithm}

\begin{ground}
Алгоритм выдает $2^n$ различных целых чисел, следовательно, $2^n$
различных кодов. Таким образом, все подмножества
генерируются, причём ровно по одному разу.
\end{ground}

Недостаток этого алгоритма состоит в том, что порядок генерации подмножеств
никак не связан с их составом. Например, вслед за подмножеством с кодом 0111,
состоящим из трёх элементов,
будет перечислено подмножество с кодом 1000,
состоящее из одного элемента.
\end{frame}

\begin{frame}
\frametitle{1.3.3. Перебор подмножеств  множества (2/4)}
Рассмотрим $B(0),B(1),\dots$~--- последовательность подмножеств, 
которую строит предыдущий алгоритм. 
Введем функцию $Q(i)$, равную количеству разрядов, в которых коды 
$B(i)$ и $B(i-1)$ отличаются.
% (в п. 6.3.3. показано, что эта функция является расстоянием на множестве битовых шкал длины n). На рис. 1.3. 
Ниже показано начало графика этой функции, для наглядности точки графика соединены линиями, а под значениями аргументов указаны соответствующие битовые шкалы. 
\begin{figure}[h]
\begin{center}
\includegraphics[height=55mm]{f1_3_3_2_4.jpg}
\end{center}
\end{figure}
\end{frame}


\begin{frame}
\frametitle{ 1.3.3. Перебор  подмножеств  множества (3/4)}
%\vspace{-2pt}
\begin{theorem}
Если $k<n$, то в последовательности битовых шкал 
$B(0),\dots, B(2^n-1)$,  $k$ правых разрядов кода $B(i)$ при $i<2^k$ 
совпадают с $k$ правыми разрядами кода $B(2^k+i):$
$\A{k\!\in\!\mathbb{N}}{\A{i\!<\!2^k}{B(i)[n\!-\!k\!+\!1..n]=B(2^k\!+\!i)[n\!-\!k\!+\!1..n]}}.$
\end{theorem}

\begin{proof}
Непосредственно следует из наблюдения, 
состоящего в том, что $k$ правых разрядов двоичного кода числа $n$ 
являются двоичным кодом числа $n \MOD 2^k$.
\end{proof}

%\vspace{\baselineskip}
\medskip
\begin{corollary} \NB{ [1]}
$\forall k,m,i\!\in\!\mathbb{N}\left( i\!<\!2^k\!\Impl\!B(i)[n\!-\!k\!+\!1..n]\!=\! 
B(m2^k\!+\!i)[n\!-\!k\!+\!1..n] \right)$.
\end{corollary}
\begin{proof}
$\A{m,i\!\in\!\mathbb{N}}
{ i\!<\!2^k\!\Impl\!i\!=\!(m2^k\!+\!i)\MOD 2^k}.$ 
\end{proof}

%\vspace{\baselineskip}
\medskip
\begin{corollary} \NB{ [2]}
$\A{k\in\mathbb{N}}{\A{m \in 1..(2^k-1)}{Q(m)=Q(2^k+m)}}$.
\end{corollary} 
\begin{proof} 
$\A{i<2^k}{B(i)[n-k+1..n]=B(2^k+i)[n-k+1..n]}$.
\end{proof}

\medskip
\begin{corollary} \NB{ [3]}
$\A{k\!\in\!\mathbb{N}}{\A{m\!\in\!0..(2^k\!-\!1)}{Q(2^k\!+\!m)\!=\!Q(2^k\!-\!m)}}$.
\end{corollary}

\begin{proof}
Индукция по числу $k$. Для $k=1, 2, 3$  выполнено. 
%(см. рисунок на предыдущем слайде).
Индукционное предположение: пусть выполнено для $k-1$, 
рассмотрим $k$.
Пусть $m=2^{k-1}+l$. 
Тогда $Q(2^k+m) = $ $=Q(m) =Q(2^{k-1}\!+\!l) = Q(2^{k-1}\!-\!l) = Q(2^k\!-\!(2^{k-1}\!+\!l) ) = Q(2^k\!-\!m)$.
\end{proof}

\end{frame}


\begin{frame}
\frametitle{ 1.3.3. Перебор  подмножеств  множества (4/4)}

\begin{corollary} \NB{ [4]}
$\A{i,j}{i+j=2^n-1 \Impl B(i)= \Not B(j)} $.
\end{corollary} 

\begin{proof}
Индукция по числу $n$. Для $n=1$ имеем: 
$B(0)\!=\!0\!=\!\Not 1\!=\!\Not B(1)$. 
Для $n=2$ имеем: 
$B(0)\!=\!00\!=\!\Not 11\!=\!\Not B(3)$, $B(1)\!=\!01\!=\!\Not 10\!=\!\Not B(2)$. 
Индукционное предположение: 
пусть следствие выполнено для $n$, 
рассмотрим $n\!+\!1$. Пусть $i\!+\!j\!=\!2^{n+1}\!-\!1$ и для определённости $i\!<\!j$. 
Тогда $0\!\le\!i\!<\!2^n\!\le\!j\!\le\!2^{n+1}\!-\!1$. 
Положим $k\Assign j\!-\!2^n$. 
Тогда $i+k=i+j-2^n=2^{n+1}-1-2^n=2^n-1$. 
Рассмотрим правые $n$ битов соответствующих кодов. 
По теореме имеем: $B(k)[2..n+1]= B(2^n+k)[2..n+1]=B(j)[2..n+1]$, 
и по индукционному предположению $B(k)[2..n+1]=\Not B(i)[2..n+1]$. 
Но $B(i)[1]=0$ и $B(j)[1]=1$, откуда $B(i)=\Not B(j)$.
\end{proof}

%\vspace{\baselineskip}
\medskip
\begin{corollary} \NB{ [5]}
Значение функции $Q(m)$ равно 
количеству нулей на конце кода $B(m)$, увеличенному на 1.
\end{corollary}


\begin{proof}
Код нечётного числа имеет в последнем разряде $1$ и именно 
в этом разряде отличается от кода предшествующего чётного числа. 
Коды чётных чисел имеют на конце столько нулей, 
какова степень числа $2$, на которую данное число делится нацело.
\end{proof} 

\end{frame}

\begin{frame}
\label{1.3.4.}\subsubsection{1.3.4. Алгоритм построения кода Грея }
\frametitle{1.3.4. Алгоритм построения кода Грея (1/7)}
%\vspace{-2pt}
\odmkey{Бинарный код Грея}{Код (code)!Грея (Gray)} $n$-элементного множества~--- 
%\footnote(Фрэнк Греу (1887--1969) )
это последовательность всех подмножеств этого множества (то есть битовых шкал), 
в которой любые два соседних подмножества в последовательности, 
в том числе первое и последнее подмножество, 
различаются в точности одним элементом (разрядом). 

\medskip
Последовательности битовых шкал с указанным свойством существуют. 
Действительно,  
для $n=1$ искомая последовательность кодов есть $0, 1$. 
Пусть известна искомая последовательность кодов $G_1, G_2,\dots, G_{2^k}$
для $n=k$. Тогда построим из этой последовательности новую следующим образом: 
сначала выпишем все коды известной последовательности, 
прибавляя слева разряд $0$, а затем эти же коды в обратном порядке, 
но прибавляя слева разряд $1$. Имеем последовательность
$0G_1, 0G_2,\dots,0G_{2^k}, 1G_{2^k},\dots, 1G_2,1G_1$. 
В этой последовательности $2^{k+1}$ кодов, 
они все различны, и любые два соседних, в том числе первый и последний, 
различаются ровно в одном разряде по построению. 

\medskip
Далее приведён \odmkey{алгоритм построения бинарного кода Грея}{Алгоритм (algorithm)!построения бинарного кода Грея (generating of binary Grey code)}  
$n$-элементного множества.
\end{frame}

\begin{frame}
\frametitle{1.3.4. Алгоритм построения кода Грея (2/7)}


\vspace{-4pt}
\begin{algorithmic}
\REQUIRE число $n \geq 0$~--- мощность множества.
\ENSURE последовательность кодов подмножеств $B$.
\STATE \textbf{$B:$ array $[1..n]$ of $0..1$}
\COMMENT{\small\color{gray} битовая шкала}
\STATE{\textbf{for $i$ from $1$ to $n$ do  $B[i]\Assign 0$ end for}}
\COMMENT{\small\color{gray} инициализация }
\STATE \textbf{yield $B$}
\COMMENT{\small\color{gray} пустое множество }
\FOR{$i$ \textbf{from} 1 \textbf{to} $2^n-1$}
\STATE $p \Assign n + 1- Q (i)$
\COMMENT{\small\color{gray} определение номера элемента }
\STATE $B[p] \Assign 1 - B[p]$
\COMMENT{\small\color{gray} добавление или удаление элемента }
\STATE \textbf{yield} $B$
\COMMENT{\small\color{gray} очередное подмножество }
\ENDFOR
\end{algorithmic}
%\end{algorithm}

\vspace{-4pt}
Функция $Q$: 
%с помощью которой определяется номер разряда,
%подлежащего изменению, возвращает количество нулей на конце
%двоичной записи числа $i$, увеличенное на 1.

%\begin{algorithm}\caption{Функция $Q$ определения номера изменяемого разряда}
%\label{gray1}
\vspace{-4pt}
\begin{algorithmic}
\REQUIRE $i$~--- номер подмножества.
\ENSURE номер изменяемого разряда.
\STATE $q \Assign 1$; $j  \Assign i$
\STATE\textbf{while $j \MOD 2 = 0$ do 
$j \Assign j \DIV 2$; $q \Assign q + 1$ end while}
\STATE \textbf{return} $q$
\end{algorithmic}
%\end{algorithm}
\end{frame}

\begin{frame}
\frametitle{1.3.4. Алгоритм построения кода Грея (3/7)}
%\vspace{-4pt}
\begin{remark}
Функция $Q$ может быть вычислена достаточно эффективно. 
Действительно, результат операции $j \MOD 2$~--- это младший разряд двоичного кода, 
который можно вычислить за один такт, например, 
взяв конъюнкцию кода числа $j$ и кода числа $1$. Результат операции 
$j \DIV 2$~--- это сдвинутый вправо на один разряд код числа $j$. В большинстве компьютеров сдвиги кодов влево (умножение на~$2$) 
и вправо (деление на~$2$) также выполняются за один такт.
\end{remark}
\begin{ground}
На нулевом шаге алгоритм выдаёт правильное подмножество $B$ (пустое).
Пусть за первые $2^k-1$ шагов алгоритм выдал последовательность значений $B$. 
При этом изменения происходили только в правых $k$ разрядах, 
а все левые $n-k$ разрядов оставались нулевыми: $B[1]=B[2]=\ldots=B[n-k]=0$. 
На $2^k$-м шаге разряд $B[n-k]$ изменяет своё значение с $0$ на $1$. 
После этого повторяется последовательность изменений значений $B[n-k+1..n]$ 
в обратном порядке, поскольку $Q(2^k+m)=Q(2^k-m)$ для $0\le m \le 2^k-1$. 
\end{ground}
\end{frame}

\begin{frame}
\frametitle{1.3.4. Алгоритм построения кода Грея (4/7)}
\begin{example}
Протокол выполнения алгоритма для $n=3$.
\begin{center}
\begin{tabular}{p{10mm}|p{10mm}|p{5mm}p{5mm}p{5mm}}
$i$ & $p$ & & $B$ & \\
\hline
&   & 0 & 0 & 0 \\
1 & 1 & 0 & 0 & 1 \\
2 & 2 & 0 & 1 & 1 \\
3 & 1 & 0 & 1 & 0 \\
4 & 3 & 1 & 1 & 0 \\
5 & 1 & 1 & 1 & 1 \\
6 & 2 & 1 & 0 & 1 \\
7 & 1 & 1 & 0 & 0 \\
\hline
\end{tabular}
\end{center}
\end{example}
%{\small
\begin{remark}
Построенная последовательность отнюдь не единственная последовательность, 
обладающая указанным свойством. 
Действительно, можно начать не с последовательности $0, 1$, 
а с последовательности $1, 0$, 
можно отражать не вправо, а влево, 
можно приписывать разряды не слева, а справа и~т.~д. 
Поэтому по\-строенный пример кода точнее называть 
\odmkey{левым зеркальным кодом Грея}{Код (code)!Грея  (Gray)!левый зеркальный (left mirror)}.
\end{remark}
%}
\end{frame}

\begin{frame}
\frametitle{ 1.3.4. Алгоритм построения кода Грея (5/7)}
Вычислять коды Грея можно различными способами.
В частности, левые зеркальные коды обладают свойством, позволяющим прямо вычислить код по его номеру.
Пусть $B(i)$~--- двоичный код числа $i$, 
а $G(i)$~--- $i$-ый левый зеркальный код Грея.
\begin{example}
Рассмотрим двоичные коды, результаты некоторых операций с кодами и зеркальный код Грея для $n=2$.

\begin{center}
\begin{tabular}{|c|c|c|c|c|}
\hline
$i$	& $B(i)$ & $B(i) \DIV 2$ & $B(i) +_2  (B(i) \DIV 2)$ &	$G(i)$ \\
\hline
0 & 00 & 00 & 00 & 00 \\
1 & 01 & 00 & 01 & 01 \\
2 & 10 & 01 & 11 & 11 \\
3 & 11 & 01 & 10 & 10 \\
\hline
\end{tabular}
\end{center}

Можно заметить, что значение выражения 
в четвёртом столбце и код Грея совпадают. Насколько случайным является это обстоятельство? 
\end{example}
\end{frame}

\begin{frame}
\frametitle{ 1.3.4. Алгоритм построения кода Грея (6/7)}
\begin{theorem}
$G(i)= B(i) +_2  (B(i)  \DIV 2 )$.  
\end{theorem}

\begin{proof}
Индукция по $n$. 
Для $n=1$ утверждение теоремы тривиально, 
для $n=2$ проверено в предыдущем примере. 
Пусть утверждение теоремы верно для битовых шкал длины $n$, 
рассмотрим $G_{n+1} (i)$. Возможны два случая.

\proofsection{$0\le i < 2^n$} 
Заметим, что $B_{n+1}(i)[1]=0$. 
По алгоритму имеем
$G_{n+1} (i)=  0\cdot G_n (i)$. 
По индукционному предположению 
имеем $G_n (i)= B_n (i) +_2  (B_n (i)  \DIV 2 )$, 
откуда\\
$G_{n+1} (i) =  0\cdot G_n (i) =  
0\cdot (B_n (i) +_2  (B_n (i) \DIV  2 ) )=
0 \cdot B_n (i) +_2  (0\cdot B_n (i)  \DIV 2 )=$\\ 
$=B_{n+1} (i) +_2  (B_{n+1} (i)  \DIV 2 )$,
где знак $\cdot$ обозначает конкатенацию.

\proofsection{$2^n \le i <2^{n+1}$} 
Заметим, что $B_{n+1} (i)[1]=1$. Положим $j\Assign i-2^n$ 
и рассмотрим $k\Assign 2^{n+1}-i-1 = 2^{n+1}-2^n-j-1 = 2^n-j-1$. 
Тогда по алгоритму построения кода Грея 
$G_{n+1} (i)=  1\cdot G_n (k)$. При этом $j+k=2^n-1$ 
и выполняется условие следствия 4. 
Заметим также, что $A +_2 B = \Not A+_2 \Not B$ для любых битовых шкал $A$ и $B$. 
По индукционному предположению 
$G_n (k)= B_n (k) +_2  (B_n (k)  \DIV 2 )$, 
получаем: 
$G_{n+1} (i)=$\\$=  1\cdot G_n (k) = 1\cdot (B_n (k) +_2  (B_n (k)  \DIV 2 ) )=  
1\cdot(\Not B_n (j) +_2  (\Not B_n (j)  \DIV 2 ) ) 
=$\\
$= 1\cdot (B_n (j) +_2  (B_n (j)  \DIV 2 ) )= 
1\cdot (B_n (i) +_2  (B_n (i)  \DIV 2 ) )  = B_{n+1} (i) +_2  (B_{n+1} (i)  \DIV 2 )$.
\end{proof}

\end{frame}

\begin{frame}
\frametitle{ 1.3.4. Алгоритм построения кода Грея (7/7)}

\vspace{-2pt}
Операции с битовыми шкалами можно выполнять не только параллельно,
сразу для всех разрядов, но и последовательно,
вычисляя биты один за другим и используя полученные значения
при вычислении следующих битов. Такие вычисления называются
\odmkey{рекуррентными}{Рекуррентные вычисления (recurrent computing)}. Иногда это бывает весьма полезно.

\begin{example}
Сумма $Z$ натуральных чисел $X$ и $Y$, представленных битовыми шкалами длины $n$,
вычисляется так:   
$p\Assign 0$; \textbf{for} $i$ \textbf{from}
$n$ \textbf{downto} $1$ \textbf{do} 
$q\Assign X[i]+Y[i]+p$; $Z[i]\Assign q \MOD 2$; $p\Assign q \DIV 2$ \textbf{end for} 
\end{example}

Применяя идею рекуррентного вычисления битов, 
легко вывести полезные следствия. 

\begin{corollary} \NB{ [1]}
$\A{1<i\le n}{G[i]=B[i] +_2 B[i-1]}$, $G[1]=B[1]$.
\end{corollary}
\begin{proof} 
$\A{k}{(B(k)\DIV 2 )[1]= 0 \And (B(k) \DIV 2 )[i] = B(k)[i - 1]}$.
\end{proof}

\begin{corollary} \NB{ [2]}
$\A{1<i\le n}{B[i]=G[i] +_2 B[i-1]}$, $G[1]=B[1]$.
\end{corollary}
\begin{proof}
$G[i]+_2B[i - 1] = B[i] +_2 B[i - 1] +_2 B[i - 1] = B[i]$.
\end{proof}

\begin{corollary} \NB{ [3]}
$\A{i}{1<i\le n \Impl B[i]= \sum_{j=1}^{i} G[j]}$.
\end{corollary}

\vspace{-1pt}
\begin{remark}
Здесь и далее $\sum$ означает сложение по модулю 2, если результат двоичный.
\end{remark}
\begin{proof}
По следствию 2
имеем $B[i]=G[i] +_2 B[i-1]= G[i]+_2 G[i-1] +_2 B[i-2]$ и т.~д. 
Учитывая, что $B[0]=0$ по определению, 
%(это возникающий слева <<ниоткуда>> разряд при сдвиге кода вправо),
получаем требуемое.
\end{proof}


\end{frame}


\begin{frame}
\label{1.3.5.}\subsubsection{1.3.5. Представление множеств списками }\frametitle{1.3.5. Представление множеств списками (1/2)}

Если универсум очень велик (или бесконечен), а
рассмат\-риваемые подмножества универсума не очень велики, то
представление с помощью битовых шкал не является эффективным с
точки зрения экономии памяти. В этом случае множества обычно
представляются \odmkey{списками элементов}{Список (list)!элементов (of elements)}. 
Элемент списка при этом
представляется записью с двумя полями: информационным и указателем
на следующий элемент. Весь список задаётся указателем на первый
элемент.

\begin{algorithmic}
\STATE elem = \textbf{struct \{ }
\STATE $\quad i\colon$ info;
\COMMENT{\small\color{gray} информационное поле}
\STATE $\quad n \colon \uparrow$ elem
\COMMENT{\small\color{gray} указатель на следующий элемент }
\STATE \textbf{\} }
\end{algorithmic}

При таком представлении трудоёмкость операции $\in$ составит $O(n)$, а
трудоёмкость операций $\subset$, $\cap$, $\cup$ составит $O(nm)$,
где $n$ и $m$~--- мощности участвующих в~операции множеств.
\end{frame}


\begin{frame}
\frametitle{1.3.5. Представление множеств списками (2/2)}

%\addvspace{-3pt}
\begin{remark}
Используемый здесь и далее символ $O$ в выражениях вида
$f(n)=O(g(n))$ читается <<О~большое>>
и означает, что функция $f$
асимптотически ограничена по сравнению с $g$ (ещё говорят, что $f$
<<имеет тот же порядок величины>>, что и $g$):
$$
f(n)=O(g(n))\Def \E{n_0, c > 0}{\A{n>n_0}{f(n)\leq c g(n)}}.
$$
\end{remark}

\medskip
Если элементы в списках упорядочить, например, по возрастанию
значения поля $i$, то трудоёмкость всех операций составит
$O(n+m)$. Эффективная реализация операций над множествами,
представленными в виде упорядоченных списков, основана на весьма
общем алгоритме, известном как алгоритм \odmkey{слияния}{Алгоритм
(algorithm)!слияния (merge)}. Алгоритм слияния параллельно
просматривает два множества, представленных упорядоченными
списками, причём на каждом шаге продвижение происходит в том
множестве, в котором текущий элемент меньше.
\end{frame}

\begin{frame}
\label{1.3.6.}\subsubsection{1.3.6. Проверка включения слиянием }\frametitle{1.3.6. Проверка включения слиянием (1/2)}

%\vspace{-3pt}
Рассмотрим \odmkey{алгоритм слияния}{Алгоритм (algorithm)!проверки включения слиянием (checking
inclusion by merge)},
который определяет, является ли множество~$A$
подмножеством множества $B$.
%\begin{algorithm}
%\caption{Проверка включения слиянием} \label{subset}
%\index{Алгоритм (algorithm)!}

\vspace{-3pt}
\begin{algorithmic}
\REQUIRE %проверяемые 
множества $A$ и $B$, которые заданы указателями $a$ и $b$.
\ENSURE {\bf true}, если $A \subset B$, в противном случае {\bf false}.
\STATE $pa \Assign a; pb \Assign b$
\COMMENT{\small\color{gray} инициализация }
\WHILE{$pa \not=\textbf{nil} \And pb \not= \textbf{nil}$}
\IF{$pa.i < pb.i$}
\STATE \textbf{return false}
\COMMENT {\small\color{gray} элемент $A$ отсутствует в  $B$ }
\ELSIF{$pa.i > pb.i$}
\STATE $pb \Assign pb.n$
\COMMENT{\small\color{gray} элемент  $A$, может быть, присутствует в $B$ }
\ELSE
\STATE $pa \Assign pa.n$
\COMMENT{\small\color{gray} здесь $pa.i = pb.i$, то~есть }
\STATE $pb \Assign pb.n$
\COMMENT{\small\color{gray} элемент $A$ точно присутствует в $B$ }
\ENDIF
\ENDWHILE
\STATE \textbf{return $pa=$ nil}
\COMMENT{\small\color{gray} {\bf true}, если $A$ исчерпано, {\bf false} в противном случае }
\end{algorithmic}
%\end{algorithm}
\end{frame}

\begin{frame}
\frametitle{1.3.6. Проверка включения слиянием (2/2)}

\begin{ground}
На каждом шаге основного цикла возможна одна из трёх ситуаций:
текущий элемент множества $A$ меньше, больше или равен текущему
элементу множества $B$. 

\smallskip
В первом случае текущий элемент множества
$A$ заведомо меньше, чем текущий и все последующие элементы
множества $B$, а потому он не содержится в множестве $B$ и можно
завершить выполнение алгоритма. 

\smallskip
Во втором случае происходит
продвижение по множеству $B$ в надежде отыскать элемент,
совпадающий с текущим элементом множества $A$. 

\smallskip
В третьем случае
найдены совпадающие элементы и происходит продвижение сразу в
обоих множествах. 

\smallskip
По завершении основно\-го~цикла возможны два
варианта:\\ либо $pa=\textbf{nil}$, либо 
$pa\not=\textbf{nil}$. Первый вариант означает, что для всех элементов
множества $A$ удалось найти совпадающие элементы в множестве $B$.
Второй вариант означает, что множество $B$ закончилось раньше, то~есть не для
всех элементов множества $A$ удалось найти совпадающие элементы в
множестве~$B$.
\end{ground}

\end{frame}

\begin{frame}
\label{1.3.7.}\subsubsection{1.3.7. Вычисление объединения слиянием }\frametitle{1.3.7. Вычисление объединения слиянием (1/3)}
\vspace{-2pt}
Рассмотрим \odmkey{алгоритм слияния}
{Алгоритм (algorithm)!вычисления объединения слиянием (finding union by merge)}, который вычисляет объединение двух множеств,
представленных упорядоченными списками.
%\begin{algorithm}
%\caption{Вычисление объединения слиянием} \label{union}
%\index{Алгоритм (algorithm)!вычисления объединения слиянием (union
%by merge)}
\vspace{-4pt}
%{\small
\begin{algorithmic}
\REQUIRE  множества $A$ и $B$, заданные указателями $a$ и $b$.
\ENSURE объединение $C=A\cup B$, заданное указателем $c$.
\STATE $pa \Assign a; pb \Assign b;  c\Assign$ \textbf{nil}; $e\Assign$ \textbf{nil}
\COMMENT{\small\color{gray} инициализация }
\WHILE{$pa \not=$ \textbf{nil} $\And pb \not=$ \textbf{nil}}
\IF{$pa.i < pb.i$}
\STATE $d \Assign pa.i;  pa \Assign pa.n$
\COMMENT {\small\color{gray} добавляем элемент из $A$ }
\ELSIF{$pa.i > pb.i$}
\STATE $d \Assign pb.i; pb \Assign pb.n$
\COMMENT {\small\color{gray} добавляем элемент из $B$ }
\ELSE
\STATE $d \Assign pa.i$ 
\COMMENT{\small\color{gray} здесь $pa.i = pb.i$, и можно взять любой }
\STATE $pa \Assign pa.n; pb \Assign pb.n$ 
\COMMENT{\small\color{gray} продвижение в обоих }
\ENDIF
\STATE $\Append(c,e,d)$
\COMMENT{\small\color{gray} добавление элемента $d$ в конец $e$ списка $c$ }
\ENDWHILE
\end{algorithmic}
%} 
\end{frame}

\begin{frame}
\frametitle{1.3.7. Вычисление объединения слиянием (2/3) }


\begin{algorithmic}
%\IF{  \COMMENT{\small\color{gray} нужно добавить в результат оставшиеся элементы $B$ }
%\ENDIF
\STATE $p \Assign$\textbf{nil} 
\COMMENT{\small\color{gray} указатель <<хвоста>> }
\STATE \textbf{if $pa \not=$ nil then $p \Assign pa$
end if}
\COMMENT{\small\color{gray} нужно добавить в результат оставшиеся элементы $A$ }
\STATE \textbf{if $pb \not=$ nil tnen $p \Assign pb$ end if}
\COMMENT{\small\color{gray} нужно добавить в результат оставшиеся элементы $B$ }
\WHILE{$p \not=$ \textbf{nil}}
\STATE $\Append(c,e,p.i)$ \COMMENT{\small\color{gray} добавление элемента $p.i$ в конец списка $c$ }
\STATE $p \Assign p.n$ \COMMENT{\small\color{gray} продвижение указателя <<хвоста>> }
\ENDWHILE
\end{algorithmic}
%\end{algorithm}


\begin{ground}
На каждом шаге основного цикла возможна одна из трёх ситуаций:
текущий элемент множества $A$ меньше, больше или равен текущему
элементу множества $B$. В первом случае в результирующий список
добавляется текущий элемент множества $A$ и происходит продвижение
в этом множестве, во втором~--- аналогичная операция производится
с множеством $B$, а в третьем случае найдены совпадающие элементы
и происходит продвижение сразу в обоих множествах.  По завершении
основного цикла один из указателей, $pa$ и $pb$ (но не оба
вместе!), может быть не равен \textbf{nil}. В этом случае остаток
соответствующего множества без проверки добавляется в результат.
\end{ground}

\end{frame}

\begin{frame}
\frametitle{1.3.7. Вычисление объединения слиянием (3/3) }
Вспомогательная процедура $\Append(c,e,d)$ присоединяет
элемент $d$ к хвосту $e$ списка $c$.
%\begin{algorithm}
%\caption{ Процедура Append~--- присоединение
%элемента в конец списка }
%\label{append}
%\vspace{-0.5cm}
\begin{algorithmic}
\REQUIRE указатель $c$ на первый элемент списка, указатель $e$ на последний
элемент списка, добавляемый элемент $d$.
\ENSURE указатель $c$ на первый элемент списка, указатель $e$ на последний
элемент списка.
\STATE $q \Assign$ new(elem); $q.i \Assign d; q.n \Assign$ \textbf{nil}
\COMMENT{\small\color{gray} новый элемент списка }
\STATE\textbf{if $c =$ nil then $c \Assign q$
else $e.n \Assign q$ end if}
\STATE $e \Assign q$
\end{algorithmic}
%\end{algorithm}


\begin{ground}
До первого вызова функции $\Append$ имеем: $c=\textbf{nil}$ и
$e=\textbf{nil}$. После первого вызова  указатель $c$ указывает на
первый элемент списка, а указатель $e$~--- на последний элемент.
Если указатели $c$ и $e$ не пустые, то после очередного вызова
функции $\Append$ это свойство сохраняется. 
\end{ground}

\end{frame}

\begin{frame}
\label{1.3.8.}\subsubsection{1.3.8. Вычисление пересечения слиянием }\frametitle{1.3.8. Вычисление пересечения слиянием (1/2)}
\vspace{-2pt}
Рассмотрим \odmkey{алгоритм слияния}{Алгоритм (algorithm)!вычисления пересечения слиянием (finding intersection by merge)},
который вычисляет пересечение двух множеств, представленных
упорядоченными списками.
%\begin{algorithm}
%\caption{Вычисление пересечения слиянием} \label{intersect}
%\index{Алгоритм (algorithm)!вычисления пересечения слиянием
%(intersection by merge)}
%\addvspace{-4.3pt}%
%\end{algorithm}%
\vspace{-2pt}

\begin{algorithmic}
\REQUIRE  множества $A$ и $B$, заданные указателями $a$ и $b$.
\ENSURE пересечение $C=A\cap B$, заданное указателем $c$.
\STATE $pa \Assign a;~~ pb \Assign b;~~ c \Assign\textbf{nil};~~ e \Assign$ \textbf{nil}
\COMMENT{\small\color{gray} инициализация }
\WHILE{$pa \not=$ \textbf{nil} $\And pb \not=$ \textbf{nil}}
\IF{$pa.i < pb.i$}
\STATE $pa \Assign pa.n$
\COMMENT{\small\color{gray} элемент множества $A$ не принадлежит пересечению }
\ELSIF{$pa.i > pb.i$}
\STATE $pb \Assign pb.n$
\COMMENT{\small\color{gray} элемент множества $B$ не принадлежит пересечению }
\ELSE
\STATE
\COMMENT {\small\color{gray} здесь $pa.i = pb.i$~--- данный элемент принадлежит пересечению }
\STATE $\Append(c,e,pa.i)$ 
\COMMENT{\small\color{gray} добавление элемента }
\STATE $pa \Assign pa.n; pb \Assign pb.n$ 
\COMMENT{\small\color{gray} продвижение указателей }  
\ENDIF
\ENDWHILE
\end{algorithmic}

\end{frame}

\begin{frame}
\frametitle{1.3.8. Вычисление пересечения слиянием (2/2)}

\begin{ground}
На каждом шаге основного цикла возможна одна из трёх ситуаций:
текущий элемент множества $A$ меньше, больше или равен текущему
элементу множества $B$. В первом случае текущий элемент множества
$A$ не принадлежит пересечению, он пропускается, и происходит
продвижение в этом множестве. Во втором случае то же самое производится с
множеством $B$. В третьем случае найдены совпадающие элементы,
один экземпляр элемента добавляется в результат, и происходит
продвижение сразу в обоих множествах. Таким образом, в~результат
попадают все совпадающие элементы обоих множеств, причём  ровно
один раз.
\end{ground}

\end{frame}

\begin{frame}
\subsubsection{1.3.9. Представление множеств итераторами }
\label{1.3.9.}\frametitle{1.3.9. Представление множеств итераторами (1/8)}
%\vspace{-4pt}
Если
структуры данных представления предполагаются известными и используются при
программировании операций, то:
\begin{enumerate}
\item[1)] возникают затруднения при совместном использовании
нескольких \em{альтернативных} представлений для одного типа объектов; 
\item[2)] необходимо заранее определять набор операций,
которым открыт доступ к внутренней структуре данных
объектов (такие операции часто называют
\odmkey{первичными}{Операция (operation)!первичная (primary)}).
\end{enumerate}

\begin{digress}
Парадигма объектно-ориентированного программирования предполагает
совместное согласованное определение структур данных и процедур
доступа к ним. 
Другими словами, в объектно-ориентированном программировании
первичные операции класса объектов заранее фиксированы в форме \odmkey{методов}{Объектно-ориентированное программирование (object-oriented programming)!метод  (method)},
а представление данных для объектов класса  
единственно и фиксировано в форме \odmkey{свойств}{Объектно-ориентированное программирование (object-oriented programming)!cвойство (property)}. 
Парадигма объектно-ориентированного программирования~--- весьма
полезная, но отнюдь не универсальная парадигма
программирования.
\end{digress}
\end{frame}

\begin{frame}
\frametitle{1.3.9. Представление множеств итераторами (2/8) }
Для множеств понятие
принадлежности является первичным и не может быть
запрограммировано без знания структуры данных для хранения
множеств.
Остальные операции над множествами можно запрограммировать независимо от
представления множеств. 
Для представления множества
достаточно
\odmkey{итератора}{Итератор (iterator)}~---
процедуры, которая перебирает элементы множества и делает с ними
то, что нужно. 

\medskip
Наличие
итератора $X$ позволяет написать цикл

\begin{algorithmic}
\FOR{$x\in X$}
\STATE $S(x)$
\COMMENT{\color{gray} где $S$~--- любая процедура, зависящая от $x$ }
\ENDFOR
\end{algorithmic}
и этот цикл означает применение оператора $S$ ко всем элементам
множества $X$ по одному разу в некотором неопределённом порядке.

\medskip
Итераторы можно реализовать в любом императивном 
языке программирования. Для этого достаточно 
иметь возможность использовать процедурные параметры.


\end{frame}


\begin{frame}
\frametitle{1.3.9. Представление множеств итераторами (3/8) }


\begin{example} Итератор для множества, представленного списком, где
$p$~--- указатель на первый элемент.
%\vspace{-\baselineskip}


\begin{algorithmic}
\REQUIRE указатель $p$ на первый элемент списка, 
процедура $S$ обработки очередного элемента списка.
\ENSURE применение процедуры $S$ ко всем элементам списка
\WHILE{$p \ne $ \textbf{nil}}
\STATE {$S(p.i)$}
\COMMENT{\small\color{gray} очередной элемент}
\STATE $p\Assign p.n$
\COMMENT{\small\color{gray} продвижение указателя }
\ENDWHILE
\end{algorithmic}
\end{example}

\begin{remark}
В доказательствах цикл \textbf{for $x\in X$ do $S(x)$ end for}
используется и в том случае, когда $|X| =\infty$.
Это не означает
реального процесса исполнения тела цикла в дискретные моменты времени,
а означает всего лишь мысленную возможность
применить оператор $S$ ко всем элементам
множества $X$ независимо от его мощности.
\end{remark}
\end{frame}

\begin{frame}
\frametitle{1.3.9. Представление множеств итераторами (4/8)}
В таких обозначениях итератор пересечения
множеств $X\cap Y$ может быть задан следующим  \odmkey{алгоритмом пересечения итераторов}
{Алгоритм (algorithm)!пересечения итераторов (intersection
of iterators)}:

\begin{algorithmic}
\FOR{$x\in X$}
\FOR{$y\in Y$}
\IF{$x=y$}
\STATE $S(x)$
\COMMENT{\small\color{gray} тело цикла }
\STATE \textbf{next for $x$}
\COMMENT{\small\color{gray} следующий $x$ }
\ENDIF
\ENDFOR
\ENDFOR
\end{algorithmic}

\begin{remark}
Оператор \textbf{next for $x$}~--- это 
\odmkey{оператор структурного перехода}{Оператор (statement)!структурного перехода
(structural go to)},
выполнение которого в данном случае означает 
прерывание
выполнения цикла по $y$ и переход к следующему шагу
в цикле по $x$.
\end{remark}
\end{frame}

\begin{frame}
\frametitle{1.3.9. Представление множеств итераторами (5/8)}
Аналогично итератор разности множеств $X\setminus Y$
может быть
задан следующим \odmkey{алгоритмом вычитания итераторов}{Алгоритм (algorithm)!вычитания итераторов (subtraction of iterators)} :
%
%\begin{algorithm}
%\caption{Итератор разности множеств} \label{iter9} \index{Алгоритм
%(algorithm)!вычитания итераторов (substraction of iterators)}
%\vspace{-\baselineskip}


\begin{algorithmic}
\FOR{$x\in X$}
\FOR{$y\in Y$}
\IF{$x=y$}
\STATE \textbf{next for $x$}
\COMMENT{\small\color{gray} следующий $x$ }
\ENDIF
\ENDFOR
\STATE $S(x)$
\COMMENT{\small\color{gray} тело цикла }
\ENDFOR
\end{algorithmic}
%\end{algorithm}

\end{frame}

\begin{frame}
\frametitle{\large 1.3.9. Представление множеств итераторами (6/8) }

Заметим, что $A\cup B = (A\setminus B)\cup B$ и $(A\setminus
B)\cap B = \emptyset$. Поэтому достаточно рассмотреть \odmkey{алгоритм 
объединения дизъюнктных множеств}{Алгоритм (algorithm)!объединения дизъюнктных итераторов
(union of disjunctive iterators)} $X\cup Y$, где ${X\cap Y =\emptyset}$.

\begin{algorithmic}
\FOR{$x\in X$}
\STATE $S(x)$
\COMMENT{\small\color{gray} тело цикла }
\ENDFOR
\FOR{$y\in Y$}
\STATE $S(y)$
\COMMENT{\small\color{gray} тело цикла }
\ENDFOR
\end{algorithmic}
%\end{algorithm}

Стоит обратить внимание на то, что
итератор пересечения реализуется вложенными циклами,
а итератор дизъюнктного объединения~---
последовательностью циклов.


\end{frame}

\begin{frame}
\frametitle{1.3.9. Представление множеств итераторами (7/8) }

%\vspace{-2pt}
Сложность
алгоритмов для операций с множествами
при использовании 
представления, не зависящего от реализации,
составляет в худшем случае $O(nm)$, где
$m$ и $n$~--- мощности участвующих множеств,
в то время как для представления упорядоченными списками
сложность в худшем случае равна $O(n+m)$.


\begin{remark}
На практике почти всегда оказывается так, что самая изощрённая
реализация операции, не зависящая от представления, оказывается
менее эффективной по сравнению с самой прямолинейной реализацией,
ориентированной на конкретное представление.
\end{remark}


\begin{digress}
При современном состоянии аппаратной базы компьютеров
оказывается, что во многих задачах можно пренебречь некоторым
выигрышем в эффективности, например различием между $O(nm)$ и
$O(n+m)$. 
Другими словами, нынешние компьютеры настолько быстры,
что часто можно не гнаться за самым эффективным представлением,
ограничившись <<достаточно эффективным>>.
\end{digress}
\end{frame}

\begin{frame}
\frametitle{1.3.9. Представление множеств итераторами (8/8) }

\begin{example}
В развитом пакете программ для решения задач линейной алгебры 
используются матрицы различных видов 
(разреженные, треугольные, блочно-диагональные и т.~д.), 
причём виды матриц настолько отличаются, что для их представления удобно использовать 
\em{различные} структуры данных (одномерные и двумерные массивы, 
списки ненулевых элементов и т.~д.). 
В таком случае нерационально программировать 
для каждой структуры данных свои функции, реализующие операции с матрицами. 
Тогда можно реализовать для каждой структуры данных 
первичные операции (скажем, чтение и запись значения элемента матрицы) 
и дальше программировать требуемые функции, 
которые уже будут работать с первичными операциями, 
независимо от структур данных, выбранных для представления матриц.

%\medskip
%{\color{red} Контрольный вопрос: как построить 
%итератор симметрической разности?}
\end{example}
\end{frame}

\begin{frame}
\frametitle{Онтология наивной теории множеств}
\begin{center}
\includegraphics[width=12cm]{1_3_B.jpg}
\end{center}
\end{frame}

\begin{frame}
\frametitle{Онтология наивной теории множеств}
\begin{center}
\includegraphics[width=12cm]{1_3_C.jpg}
\end{center}
\end{frame}


\begin{frame}
\frametitle{1.4. Отношения (1/2)}
\subsection{1.4. Отношения}
\label{1.4.}

Обычное широко используемое понятие <<отношение>>
имеет вполне точное математическое значение,
которое вводится в этом разделе.

Формально отношения являются множествами, 
и тему можно было бы назвать
<<Множества, в том числе отношения>>.
Однако с помощью отношений вводят в рассмотрение
целый ряд понятий и обозначений,
которые находят, может быть, даже больше
применений, чем рассмотренные операции с множествами.
Поэтому понятие <<отношение>> вполне правомерно
поставить в один ряд с понятием <<множество>>.

\medskip
\begin{tabular}{p{13mm}|p{40mm}|p{88mm}}
  № & Название параграфа
  & Ключевые термины и обозначения \\
  \hline
\DMref{1.4.1.}{1.4.1.} \RS & Упорядоченные пары и наборы 
  & 
  {$(\Tuple{a}{1}{n})$~--- кортеж;
  
  $|(\Tuple{a}{1}{n})|$~--- длина кортежа}
  \\
\DMref{1.4.2.}{1.4.2.} \RS & Прямое произведение множеств 
  & 
{$A\times B$~--- декартово произведение множеств;

$A^k$~--- степень множества}
  \\
\DMref{1.4.3.}{1.4.3.} \BS & Размеченное объединение множеств &
$A\oplus B$~--- прямая сумма множеств \\
 \end{tabular}   
\end{frame}

\begin{frame}
\frametitle{1.4. Отношения (2/2)}

\begin{tabular}{p{13mm}|p{30mm}|p{100mm}}
  № & Назв. параграфа
  & Ключевые термины и обозначения \\
  \hline
\DMref{1.4.4.}{1.4.4.} & Бинарные отношения &
{График отношения, область отправления, область прибытия; $\Dom$~--- область определения,
 $\Im$~--- область значений; инфиксная форма записи; 
$R^{-1}$~--- обратное отношение; $\overline{R}$~---дополнение отношения; $U$~---универсальное отношение; 
 тождественное отношение} 
\\
\DMref{1.4.5.}{1.4.5.} & Композиция отношений &
{$\circ$~--- операция композиции}\\  
\DMref{1.4.6.}{1.4.6.} \RS & Свойства отношений &
{Рефлексивность, антирефлексивность, 
симметричность, антисимметричность,
 транзитивность, линейность}\\  
\DMref{1.4.7.}{1.4.7.} & Степень отношения и циклы &
{Длина цикла, треугольник, петля, ациклическое отношение}\\  
\DMref{1.4.8.}{1.4.8.} \RS & Ядро отношения &
{$\ker $~--- операция взятия ядра отношения }\\  \DMref{1.4.9.}{1.4.9.} & Представление отношений в программах &
{Булевая матрица, транспонирование, инвертирование, вычитание, умножение, дизъюнкция }
\end{tabular}  

\end{frame}


\begin{frame}
\label{1.4.1.}\subsubsection{1.4.1. Упорядоченные пары и наборы }\frametitle{1.4.1. Упорядоченные пары и наборы (1/3)}
Если $a$ и $b$~--- объекты, то через $(a,b)$ обозначим
\odmkey{упорядоченную пару}{Упорядоченная пара (ordered couple)}.
\odmkey{Равенство}{Равенство (equality)!упорядоченных пар (of
ordered pairs)} упорядоченных пар определяется следующим
образом:
$${(a,b)=(c,d)\Def a=c \And b=d}.$$ 
Вообще говоря,
${(a,b)\ne(b,a)}$, если $a\ne b$.

\begin{remark}
Формально, упорядоченные пары можно определить как множества, если
определить их так: 
${(a,b)\Def\{\{a\},\{a,b\}\}}.$
Таким образом,
введённое понятие упорядоченной пары не выводит рассмотрение за
пределы наивной теории множеств.
\end{remark}

\medskip
Аналогично упорядоченным парам можно рассматривать упорядоченные
тройки, четвёрки и~т.\;д. В общем случае подобные объекты
называются \odmkey{$n$-ками}{$n$-ка ($n$-tuple)}, \odmkey{кортежами}{Кортеж (tuple)},
\odmkey{наборами}{Набор (tuple)} или (конечными)
\odmkey{последовательностями}{Последовательность
(sequence)!конечная (finite)}. Упорядоченный набор из $n$
элементов обозначается $\left(\Tuple{a}{1}{n}\right)$.

\begin{remark} Обычно элементы последовательности считают
однотипными, то есть принадлежащими одному множеству,
хотя, вообще говоря, это не обязательно.  
\end{remark}

\end{frame}

\begin{frame}
\frametitle{1.4.1. Упорядоченные пары и наборы (2/3)}


Набор
$\left(\Tuple{a}{1}{n}\right)$ можно определить рекурсивно,
используя понятие упорядоченной пары:
$$
\left(\Tuple{a}{1}{n}\right) \Def \left(\left(\Tuple{a}{1}{n-1}\right),a_n\right).
$$
Количество элементов в наборе называется его \odmkey{длиной}{Длина
(length)!набора (of tuple)} и обозначается следующим образом:
$\left\vert\left(\Tuple{a}{1}{n}\right)\right\vert$.


\medskip
\begin{remark}
Наиболее естественным представлением
в программе $n$-ки с элементами из множества
$A$ является массив \textbf{array $[1..n]$ of $A$}.
\end{remark}

\medskip
\begin{example}
Место в зрительном зале обычно 
определяется парой натуральных чисел: номер ряда и 
номер места в ряду.
\end{example}

\end{frame}

\begin{frame}
\frametitle{1.4.1. Упорядоченные пары и наборы (3/3)}
\begin{theorem} Два набора одной длины равны, если
равны их соответствующие эле\-менты:
$$
\A{n}{\left(\Tuple{a}{1}{n}\right)=\left(\Tuple{b}{1}{n}\right)
\Equiv \BigAnd\limits_{i=1}^{n} a_i=b_i}.
$$
\end{theorem}
Отсюда следует, что порядок <<отщепления>> элементов в рекурсивном
определении кортежа не имеет значения.

\begin{proof}
Индукция по $n$. База: при $n=2$  по определению равенства упорядоченных пар.
Пусть теперь
$$
\left(\Tuple{a}{1}{n-1}\right)=\left(\Tuple{b}{1}{n-1}\right)
\Equiv \BigAnd\limits_{i=1}^{n-1} a_i=b_i.
$$
Тогда
$\left(\Tuple{a}{1}{n}\right)=\left(\Tuple{b}{1}{n}\right) \Equiv$\\
$\Equiv\left(\left(\Tuple{a}{1}{n-1}\right),a_n\right)=
\left(\left(\Tuple{b}{1}{n-1}\right),b_n\right) \Equiv$\\
$\Equiv \left(\Tuple{a}{1}{n-1}\right)=\left(\Tuple{b}{1}{n-1}\right)
\And a_n=b_n \Equiv \BigAnd\limits_{i=1}^{n} a_i=b_i$.
\end{proof}



\end{frame}

\begin{frame}
\label{1.4.2.}\subsubsection{1.4.2. Прямое произведение множеств }\frametitle{1.4.2. Прямое произведение множеств (1/2)}

%\vspace{-4pt}
%Пусть $A$ и $B$~--- два множества.
\odmkey{Прямым (декартовым)
произведением}{Произведение (product)!прямое (direct)}
\odmkey{}{Произведение (product)!декартово (Cartesian)}
двух множеств $A$ и $B$ называется множество всех
упорядоченных пар, в которых первый элемент принадлежит множеству $A$, а
второй принадлежит множеству $B$:
$$ A\times B \Def \Set{(a,b)}{a\in A \And b\in B}. $$

%\vspace{-4pt}
\begin{remark}
Точка на плоскости может быть задана упорядоченной парой
координат, т.~е. двумя точками на координатных осях. Таким
образом, $\Bbb R^2 = \Bbb R \times \Bbb R$. Своим появлением метод
координат обязан Декарту,
%\footnotemark\footnotetext{\;Рене Декарт
%(1596--1650).}, 
отсюда и название <<декартово произведение>>.
\end{remark}

\begin{theorem} 
Если множества $A$ и $B$ конечны, то
$|A\times B|=|A|\,|B|$.
\end{theorem}
\begin{proof}
Первый компонент упорядоченной пары можно выбрать $\vert A\vert$
способами, второй~--- $\vert B\vert$ способами, причём независимо от
выбора первого элемента. Таким образом, всего
имеется $\vert A\vert\vert B\vert$ различных упорядоченных пар.
\end{proof}

Понятие прямого произведения допускает обобщение. 

\vspace{-4pt}
$$
A_1 \times \dots \times A_n \Def \Set{(\Tuple{a}{1}{n})} {a_1\in
A_1 \And \dots \And a_n\in A_n}.
$$

\end{frame}

\begin{frame}
\frametitle{1.4.2. Прямое произведение множеств (2/2)}
%\vspace{-2pt}
В прямом произведении $A_1 \times \dots \times A_n$ множества $A_i$ не обязательно различны.

\medskip
\begin{corollary} 
$\vert A_1 \times \cdots \times A_n\vert = 
\vert A_1 \vert \cdot \dots \cdot \vert A_n\vert$.
\end{corollary}

\medskip
Прямое произведение нескольких множеств можно вычислять в произвольном порядке.

\begin{lemma}
$(A\times B)\times C \sim A\times (B\times C)$.
\end{lemma}
\begin{proof}
$(A\times B)\times C = 
\Set{((a,b),c)}{(a,b)\in A\times B \And c\in C} =$ \\
$= \Set{(a,b,c)}{a\in A \And b\in B \And c\in C} \ToS{(a,b,c)\leftrightarrow(a,(b,c))}$ \\ 
$\sim \Set{(a,(b,c))}{a\in A \And (b,c)\in B\times C} =  A\times (B\times C).$
\end{proof}

\smallskip
\odmkey{Степенью}{Степень (degree)!множества (of set)} множества
$A$ называется его $n$-кратное прямое произведение самого на себя.
Обозначение:
$
A^n\Def\underbrace{A \times \dots \times A}_{\text{$n$ раз}}.
$
Соответственно, $A^1\Def A$, $A^2\Def A\times A$ и вообще $$A^n\Def
A^{n-1}\times A.$$

\smallskip
\begin{corollary} $|A^n|=|A|^n$.
\end{corollary}
%
%
\end{frame}

\begin{frame}
\label{1.4.3.}\subsubsection{1.4.3. Размеченное объединение множеств }\frametitle{1.4.3. Размеченное объединение множеств (1/3)}
Пусть $\Tuple{X}{1}{n}$~--- семейство множеств, 
возможно имеющих общие элементы. 
\odmkey{Размечен\-ным объединением}{Объединение (union)!множеств (of sets)!размеченное (tagged)} 
(употребляются также термины \odmkey{дизъюнктное объединение множеств}{Объединение (union)!множеств (of sets)!дизъюнктивное (disjunctive)} и \odmkey{прямая сумма множеств}{Прямая сумма множеств (direct sum of sets)}) 
$n$ множеств $\Tuple{X}{1}{n}$
называется множество всех упорядоченных пар, в которых первый элемент принадлежит множеству $X_i$, а второй элемент~--- 
натуральное 
число $i$:
%\vspace{-\baselineskip}
$$X_1\oplus X_2\oplus {\dots} \oplus X_n\Def\Set{(x,i)}{ 1\le i \le n \And x \in X_i }.$$
%\vspace{-0.25\baselineskip}
Операция размеченного объединения приписывает каждому элементу нового множества индекс, 
по которому можно понять, из какого конкретно множества элемент попал в результат, 
в отличие от операции объединения, в которой эта информация теряется.

\begin{remark}
Размеченное объединение множеств $\Tuple{X}{1}{n}$
есть объединение декартовых произведений  $X_i\times\{i\}$:

\vspace{-12pt}
$$X_1\oplus X_2\oplus {\dots} \oplus X_n=\bigcup_{i=1}^{n} \left(X_i\times\{i\}\right).$$  

Эта операция считается 
\em{вариаргументной}, поэтому скобки в записи 
$X_1\oplus X_2\oplus {\dots} \oplus X_n$ опускаются,
и порядок аргументов также не имеет значения.
\end{remark}
\end{frame}

\begin{frame}
\frametitle{1.4.3. Размеченное объединение множеств (2/3)}
\begin{theorem}
Для конечных множеств 
$\left|X_1\oplus X_2\oplus\dots\oplus X_n\right| = 
\sum\limits_{i=1}^n |X_i|$.
\end{theorem}

\begin{proof}
Множества пар $X_i\times \{i\}$ 
с различными метками $i$ не имеют общих элементов, значит:

\vspace{-\baselineskip}
$$
\left\vert \bigcup\limits_i X_i\times \{ i \} \right\vert = 
\sum_i \left\vert X_i\times \{ i \} \right\vert = 
\sum_i \vert X_i \vert \  \vert \{ i \} \vert = \sum_i \vert X_i\vert.
$$

\vspace{-0.25\baselineskip}
Учитывая замечание, получаем: 
$|X_1\oplus X_2\oplus\dots\oplus X_n |=|X_1 |+|X_2 |+\dots+|X_n|$. 						
\end{proof}

%{\small
\begin{example}
Пусть $S$ и $U$ --- множества двоичных последовательностей длины $n$. 
Ясно, что они могут иметь неотличимые элементы. 
Пусть для $S$ задано отображение в целые числа, 
которое трактует элемент $S$ как двоичное представление числа со знаком;  
для $U$ задано отображение в натуральные числа, 
трактующее его элемент как двоичное представление числа без знака. 
Тогда для размеченного объединения множеств 
$X=S\oplus U$ 
можно построить отображение в целые числа, 
которое в зависимости от метки элемента будет переводить его 
в целое число по соответствующему правилу.
\end{example}
%}
\end{frame}

\begin{frame}

\frametitle{1.4.3. Размеченное объединение множеств (3/3)}
%\vspace{-0.2\baselineskip}
\begin{digress}
В программировании операция размеченного
объединения поддерживается структурой данных,
используемой для хранения значения одного
из нескольких типов. 
Текущим является тип значения, записанного последним.
В языках семейства Си объединение типов задается ключевым словом \textbf{union}. 
При этом информация о текущем типе нигде не сохраняется, 
и за безошибочное обращение со значением 
полностью отвечает программист. 
Размеченное объединение \textbf{record case} 
в языке Паскаль имеет специальное поле 
(\odmkey{тег}{Тег (tag)}),
хранящее тип текущего значения. 
Анализ тега позволяет программисту избегать ошибок.
На первый взгляд то же поведение можно обеспечить, 
если взять тип прямого произведения 
(\textbf{struct} в языках семейства Си),
завести в нем поля необходимых типов и 
поле тега. 
Преимуществом размеченного объединения 
\textbf{record case} перед типом прямого произведения
\textbf{struct} является выигрыш по памяти. 
В случае размеченного объединения достаточно выделить память 
под наиболее объёмное значение (и под поле тега), 
а в случае прямого произведения необходимо выделить 
память под значения всех типов в сумме.
\end{digress}


\end{frame}

\begin{frame}
\label{1.4.4.}\subsubsection{1.4.4. Бинарные отношения }\frametitle{1.4.4. Бинарные отношения (1/4)}
\odmkey{Бинарным отношением}{Отношение (relation)!бинарное (binary)}
между множествами $A$ и~$B$
называется тройка $\langle  A, B, R \rangle$,
где $R$~--- подмножество
прямого произведения $A$ и $B$:
$ R \subset A \times B. $
Множество $R$ называется \odmkey{графиком отношения}{График (graph)!отношения (of
relation)}, $A$ называется \odmkey{областью
отправления}{Область отправления (source range)!отношения (of
relation)}, а $B$~--- \odmkey{областью прибытия}{Область прибытия
(target range)!отношения (of relation)}.
Если множества $A$ и $B$
определены контекстом, то часто просто говорят, что задано
отношение $R$. При этом для краткости отношение обозначают тем же
символом, что и график.

\medskip
\begin{remark}
Принято говорить <<отношение \em{между} множествами>>,
хотя порядок множеств имеет значение
и отношение между $A$ и $B$ совсем не то же самое,
что отношение между $B$ и $A$.
Поэтому иногда, чтобы подчеркнуть это обстоятельство, употребляют оборот
<<отношение $R$ \em{из} множества $A$ \em{в} множество $B$>>.
\end{remark}

\medskip
Среди элементов множеств $A$ и $B$ могут быть такие,
которые входят в одну или несколько пар, 
определяющих отношение  $R$, а могут быть такие, которые не
входят ни в одну из пар, определяющих отношение  $R$. 

\end{frame}

\begin{frame}
\frametitle{1.4.4. Бинарные отношения (2/4)}
%\vspace{-0.2\baselineskip}
Множество
элементов области отправления, которые присутствуют в парах
отношения, называется \odmkey{областью определения}{Область
определения (domain)!отношения (of relation)} отношения $R$ и
обозначается $\Dom R$,
а множество элементов области прибытия,
которые присутствуют в парах отношения, называется
\odmkey{областью значений}{Область значений (codomain, value
range)!отношения (of relation)} отношения $R$ и обозначается $\Im
R$:\\
$\Dom R \Def \Set{a\in A}{\E{b\in B}{(a,b)\in R}},\quad \Dom R \subset A$\\
$\Im R \Def \Set{b\in B}{\E{a\in A}{(a,b)\in R}},\quad \Im R \subset B$.

\medskip
Если $A=B$ ($R \subset A^2 $), то говорят, что $R$ есть \odmkey{отношение
на множестве}{Отношение (relation)!на множестве (on set)}~$A$.
Для бинарных отношений обычно используется
\odmkey{инфиксная форма записи}{Форма записи (notation)!инфиксная (infix)}  :
$aRb \Def (a,b)\in R \subset A \times B.$
Инфиксная форма позволяет более кратко записывать некоторые  формы
утверж\-дений относительно отношений:
$
aRbRc\!\Def\! (a,b)\!\in\! R\! \And\! (b,c)\!\in\!R.
$

\medskip
\begin{example}
Пусть задано множество $U$. Тогда $\in$ (принадлежность)~---  
отношение между элементами множества $U$ и элементами булеана $2^U$, 
а $\subset$ (включение) и $=$ (равенство)~--- отношения на $2^U$. 
Хорошо известны отношения $=,<,\leq,>,\geq,\not=$ между числами.
\end{example}
\end{frame}

\begin{frame}
\frametitle{1.4.4. Бинарные отношения (3/4)}
%\vspace{-2pt}
Пусть $R$~--- отношение между множествами $A$ и $B$: $R \subset A \times B$.
Введём следующие понятия:
\begin{tabbing}
\odmkey{Тождественное}{Отношение!тождественное (identical)}
отношение: \= $R^{-1}$\=\kill

\odmkey{Обратное}{Отношение (relation)!обратное (converse)}
отношение: \> $R^{-1}$ \> $\Def \Set{(b,a)}{(a,b)\in R}\subset B\times A$. \\


\odmkey{Дополнение}{Отношение (relation)!дополнительное (complement)} отношения:
\> $\overline{R}$ \> $\Def \Set{(a,b)}{(a,b) \notin R}\subset A \times B$. \\


\odmkey{Универсальное}{Отношение (relation)!универсальное (universal)} отношение:
\> $U$\>$\Def \Set{(a,b)}{a\!\in\!A\!\And\!b\!\in\!B}= A\!\times\!B$.
\end{tabbing}

\begin{columns}[t]
\begin{column}[T]{6cm}
\begin{example} Графики отношений\\ 
$R=\Set{(x,y)}{4>x-y>1}$,\\
$x\in 1..5$, $y\in 1..4$ и\\  
$R^{-1}=\Set{(x,y)}{4>y-x>1}$,\\ $x\in 1..4$, $y\in 1..5$.
\end{example}
\end{column}
\begin{column}[T]{8cm}
\vspace{-12pt}
\begin{figure}[h]
\begin{center}
\includegraphics[height=45mm]{f1_4_4_3_4.jpg}
\end{center}
\end{figure}
\end{column}
\end{columns}


\end{frame}

\begin{frame}
\frametitle{1.4.4. Бинарные отношения (4/4)}
%\vspace{-4pt}
На любом множестве можно определить
\odmkey{тождественное}{Отношение (relation)!тождественное (identical)} отношение: \\
$I\Def \Set{(a,a)}{a\in A} \subset A^2$. Обычно для обозначения 
тождественного отношения используют знак равенства $=$.

\begin{examples}
Равенство множеств и равенство чисел~--- тождественные отношения.
\end{examples}

\begin{remark}
График тождественного отношения на множестве $A$ иногда называют
\odmkey{диагональю}{Диагональ (diagonal)!прямого произведения (of
direct product)} в $A\times A$.
\end{remark}

\medskip
Обобщённое понятие отношения:
\odmkey{$n$-местное}{Отношение (relation)!$n$-местное ($n$-place)}
(\odmkey{$n$-арное}{Отношение (relation)!$n$-арное ($n$-ary)})
отношение $R$~--- это подмножество прямого произведения $n$
множеств, то есть множество 
\odmkey{кортежей}{Кортеж (tuple)}:

\vspace{-12pt}
$$
R\subset A_1\times\dots\times A_n \Def
\Set{(\Tuple{a}{1}{n})}{a_1\in A_1\And\dots\And a_n\in A_n}.
$$

\vspace{-2pt}
Обычно вместо $(\Tuple{a}{1}{n})\in R$ пишут $R(\Tuple{a}{1}{n})$.
Число $n$, то есть длину кортежей отношения, называют 
\odmkey{вместимостью}{Вместимость отношения (arity of relation)} или \odmkey{арностью}{Арность отношения (arity of relation)} от английского слова arity.

\begin{digress}
Многоместные отношения используются, например, в теории баз данных. Само
название <<реляционная>> база данных происходит от слова relation
(отношение).
\end{digress}

%\vspace{-3pt}
Далее рассматриваются только двуместные (бинарные) отношения, при этом слово
<<бинарные>> опускается.
\end{frame}

\begin{frame}
\label{1.4.5.}\subsubsection{1.4.5. Композиция отношений }\frametitle{1.4.5. Композиция отношений (1/3) }
%\vspace{-2pt}
Пусть $R_1\subset A\times C$~--- отношение из $A$ в $C$, а
$R_2\subset C\times B$~--- отношение из $C$ в $B$.
\odmkey{Композицией}{Композиция (composition)!отношений (of
relations)} двух
отношений $R_1$ и $R_2$ называется отношение $R\subset A\times B$
из $A$ в $B$, определяемое следующим образом:
$$
R = R_1\circ R_2\Def \Set{(a,b)}{a\in A \And b\in B \And
\E{c\in C}{aR_1c \And cR_2b}}.
$$
Другими словами,
$
a(R_1\circ R_2)b \Equiv \E{c\in C}{aR_1c\And cR_2b}.
$

\begin{theorem} Композиция отношений ассоциативна, то
есть
$$
\forall R_1\subset A\times B, R_2\subset B\times C, R_3\subset C\times D \left(
(R_1\circ R_2)\circ R_3 = R_1\circ(R_2\circ R_3)
\right).$$
\end{theorem}
\begin{proof}
\begin{align*}
a(R_1\circ R_2)\circ R_3 d
&\Equiv \E{c\in C}{a R_1\circ R_2 c \And c R_3 d} \Equiv \\
&\Equiv \E{c\in C}{(\E{b\in B}{a R_1 b \And b R_2 c}) \And c R_3 d} \Equiv \\
&\Equiv \E{b\in B, c\in C}{a R_1 b \And b R_2 c \And c R_3 d} \Equiv \\
&\Equiv \E{b\in B}{a R_1 b \And (\E{c\in C}{b R_2 c\And c R_3 d})} \Equiv \\
&\Equiv \E{b\in B}{a R_1 b \And b R_2\circ R_3 d} \Equiv a
R_1\circ(R_2\circ R_3) d. \hspace{30pt}\qed\end{align*}
\renewcommand{\qed}{\vspace{-\baselineskip}}\end{proof}
\end{frame}

\begin{frame}
\frametitle{1.4.5. Композиция отношений (2/3) }


Композиция отношений на множестве $A$ является отношением на
множестве $A$.

\begin{example}
Композиция отношений $<$ и $\leq$  на множестве вещественных чисел
совпадает с отношением $<$: $<\circ \leq\ = \ <$.
\end{example}

\begin{remark}
В общем случае композиция отношений не коммутативна, то есть\\
$R_1\circ R_2 \ne R_2\circ R_1$.
\end{remark}

\begin{example}
На множестве людей определены отношения
кровного родства и супружества. Отношения <<кровные родственники супругов>> и
<<супруги кровных родственников>> различны.
\end{example}

\medskip
Из определения композиции непосредственно следует, 
что если $S_1 \subset R_1  \And S_2 \subset R_2$,\\ 
то $S_1\circ S_2 \subset R_1 \circ R_2$. 

\end{frame}

\begin{frame}
\frametitle{1.4.5. Композиция отношений (3/3) }
\begin{columns}[t]
\begin{column}[T]{5cm}
\begin{example} Геометрическая иллюстрация композиции отношений.\\
$R_1\subset A\times C$,\\
$R_2\subset C\times B$,\\
$R=R_1\circ R_2\subset A\times B$.\\
Нетрудно заметить, что\\
$\Dom R\subset \Dom R_1$ 
и \\$\Im R\subset \Im R_2$.\\
Ясно также, что если\\
$\Im R_1 \cap \Dom R_2 = \emptyset$,\\
то $|R|=0$.
\end{example}
\end{column}
\begin{column}[T]{10cm} 
\begin{figure}[h]
\begin{center}
\includegraphics[height=70mm]{f1_4_5_3_3.jpg}
\end{center}
\end{figure}
\end{column}
\end{columns}

\end{frame}

\begin{frame}
\label{1.4.6.}\subsubsection{1.4.6. Свойства отношений }\frametitle{1.4.6. Свойства отношений (1/4)}

\vspace{-2pt}
Пусть $R\subset A^2$. Тогда
отношение $R$ называется

\vspace{-5pt}
\begin{tabbing}
\hspace{5mm}\=\odmkey{антирефлексивным}{Отношение (relation)!антирефлексивное
(antireflexive)},\quad \=\kill
\>\odmkey{рефлексивным}{Отношение (relation)!рефлексивное (reflexive)},
\>если $\A{a\in A}{aRa}$;\\
\>\odmkey{антирефлексивным}{Отношение (relation)!антирефлексивное (antireflexive)},
\>если $\A{a\in A}{\Not aRa}$;\\
\>\odmkey{симметричным}{Отношение (relation)!симметричное (symmentic)},
\>если $\A{a,b\in A}{aRb \Impl bRa}$;\\
\>\odmkey{антисимметричным}{Отношение (relation)!антисимметричное (antisymmetric)},
\>если $\A{a,b\in A}{aRb \And bRa \Impl a=b}$;\\
\>\odmkey{транзитивным}{Отношение (relation)!транзитивное (transitive)},
\>если $\A{a,b,c\in A}{aRb \And bRc \Impl aRc}$;\\
\>\odmkey{линейным}{Отношение (relation)!линейное (linear)},
\>если $\A{a,b\in A}{a = b \Or aRb \Or bRa}$.
\end{tabbing}

\vspace{-5pt}
\begin{remark}
Иногда линейное отношение $R$ называют \odmkey{полным}{Отношение
(relation)!полное (complete)}. 
Такой термин является оправданным,
поскольку для любых двух различных элементов, $a$ и $b$, либо пара
$(a,b)$, либо пара $(b,a)$ принадлежит отношению $R$. 
Отношение,
не обладающее свойством полноты, при этом вполне естественно
называть \odmkey{частичным}{Отношение (relation)!частичное
(partial)}. Однако использование термина <<полное отношение>>
может порождать терминологическую путаницу. Дело в
том, что словосочетание <<вполне упорядоченное
множество>> 
оказывается созвучным словосочетанию <<множество с полным
порядком>>, хотя эти словосочетания обозначают существенно
различные объекты. Поэтому здесь 
термин <<линейное отношение>> считается антонимом термину <<частичное
отношение>>.
\end{remark}
\end{frame}

\begin{frame}
\frametitle{1.4.6. Свойства отношений (2/4)}

\begin{theorem} Пусть $R \subset A\times A$~{\upshape---} отношение на $A$. Тогда:

%\vspace{-4pt}
\begin{tabbing}
\=$1.$ $R$ антисимметрично$\quad$\= \kill
\>$1.$ $R$ рефлексивно
\> $\Equiv I\subset R${\upshape;} \\
\>$2.$ $R$ симметрично
\>$\Equiv R=R^{-1}${\upshape;} \\
\>$3.$ $R$ транзитивно
\>$\Equiv R\circ R\subset R${\upshape;} \\
\>$4.$ $R$ антисимметрично
\>$\Equiv R \cap R^{-1}\subset I${\upshape;} \\
\>$5.$ $R$ антирефлексивно
\>$\Equiv R \cap I = \emptyset${\upshape;} \\
\>$6.$ $R$ линейно \>$\Equiv R\cup I\cup R^{-1} = U$.
\end{tabbing}
\end{theorem}

%\vspace{-6pt}

\begin{proof}
Используя определения свойств отношений, имеем:
%%\settowidth{\leftmargini}{1. $\Equiv.$}
%%\addtolength{\leftmargini}{\labelsep}
\proofsection{1. $\Equiv$}
$\A{a\in A}{aRa}
\Equiv\A{a\in A}{(a,a)\in R}
\Equiv I\subset R$.
\proofsection{2. $\Equiv$}
$\A{a,b\in A}{aRb\Impl bRa}\Equiv
\A{a,b\in A}{(a,b)\in R \Impl (b,a) \in  R} \Equiv$\\
$\Equiv \A{a,b \in A}{\left((a,b) \in R \Impl(b,a)\in R\right)
\And\left((b,a)\in R \Impl(a,b)\in R\right)}\Equiv$ \\
$\Equiv \A{a,b \in  A}{\left((a,b) \in  R \Impl(a,b) \in  R^{-1}\right)
\And \left((a,b) \in  R^{-1} \Impl(a,b) \in  R\right)}\Equiv$\\
$\Equiv R\subset R^{-1}  \And  R^{-1}\subset R \Equiv R=R^{-1}$.
\end{proof}
\end{frame}

\begin{frame}
\frametitle{1.4.6. Свойства отношений (3/4)}
\begin{proof}(Продолжение)
\proofsection{3. $\Equiv$}
$\A{a,b,c  \in  A}{aRb \And  bRc\Impl aRc}\Equiv$\\
$\Equiv\A{a,b,c  \in  A}{(a,b) \in  R  \And  (b,c) \in  R\Impl (a,c) \in  R}\Equiv$\\
$\Equiv\A{a,c  \in  A}{\E{b \in  A}{(a,b) \in  R  \And  (b,c) \in  R}\Impl (a,c) \in  R}\Equiv$\\
$\Equiv\A{a,c  \in  A}{(a,c) \in  R\circ R\Impl (a,c) \in  R}\Equiv$\\
$\Equiv\A{a,c  \in  A}{a R\circ R c\Impl aRc}\Equiv R\circ R \subset
R$.
\proofsection{4. $\Impl$} От противного. $R\cap
R^{-1}\not\subset I\Impl
\E{a\ne b}{aRb \And  aR^{-1}b}\Impl$\\
$\Impl\E{a\ne b}{aRb \And  bRa}$.%
\proofsection{4. $\Impll$} $R \cap 
R^{- 1} \subset  I \Impl (aRb \And  aR^{- 1}b\Impl a = b) \Impl (aRb \And 
bRa \Impl  a =b )$. 
\proofsection{5. $\Equiv$}
$R\cap I=\emptyset\Equiv\neg\E{a \in  A}{aRa}\Equiv\A{a \in  A}{\neg aRa}$.
\proofsection{6. $\Equiv$}
$\A{a,b \in  A}{a=b\Or aRb\Or bRa}\Equiv$\\
$\Equiv\A{a,b \in  A}{(a,b) \in  I \Or (a,b) \in  R \Or 
(a,b) \in  R^{-1}}\Equiv$\\
$\Equiv U\subset I\cup R\cup R^{-1}
\Equiv U=R\cup I\cup R^{-1}$.
\end{proof}

\end{frame}

\begin{frame}
\frametitle{1.4.6. Свойства отношений (4/4)}

\begin{columns}[t] % contents are top vertically aligned
\begin{column}[T]{10cm}
\begin{example}
Пусть на множестве $ M=\{1,2,3,4\}$
задано бинарное отношение $R$,
график которого показан справа.
Какими из перечисленных свойств обладает это отношение?
Отношение $R$ не является ни рефлексивным, ни симметричным, 
ни транзитивным, ни линейным, ни антирефлексивным,
ни антисимметричным.
\end{example}
\end{column}
\begin{column}[T]{4.5cm}
\vspace{-12pt} 
\includegraphics[height=3.2cm]{f1_4_6_4_4.jpg}
\end{column}
\end{columns}

\medskip
Свойства отношений не являются 
полностью независимыми, 
в некоторых случаях они несовместны.
\begin{example}
Отношение не может одновременно быть рефлексивным и 
антирефлексивным.
\end{example} 

\begin{theorem}
Если отношение линейно, симметрично
и транзитивно, то оно рефлексивно.
\end{theorem}
\begin{proof}
Рассмотрим $R\subset M^2$. Пусть $a\in M, b\in M, a\ne b$.
Тогда, в силу линейности, $aRb \Or bRa$,
ввиду симметричности
$aRb \And bRa$ и по транзитивности $aRa$.
\end{proof}

\begin{remark} Однако из симметричности и транзитивности рефлексивность не вытекает.
Контрпример~--- пустое отношение,
которое симметрично и транзитивно,
но не линейно, а потому и не рефлексивно. 
\end{remark}

\end{frame}



\begin{frame}
\label{1.4.7.}\subsubsection{1.4.7. Степень отношения и циклы }\frametitle{1.4.7. Степень отношения и циклы (1/4)}
Пусть $R$~--- отношение на множестве $A$.
\odmkey{Степенью отношения}{Степень (degree)!отношения (of relation)}
$R$ на множестве $A$ называется его $n$-кратная
композиция с самим собой. Обозначение:
%\vspace{-12pt}
$$
R^n\Def\underbrace{R \circ \dots \circ R}_{\text{$n$ раз}}.
$$
Соответственно, $R^0\Def I$, $R^1\Def R$, $R^2\Def R\circ R$
и вообще $R^n\Def
R^{n-1}\circ R$.

\medskip
Если $(a,b)\in R^n$, то это значит, что 
$\E{\Tuple{c}{1}{n-1} \in A}{a\,R\,c_1\, R\dots R\,c_{n-1}\,R\,b}$.
Последо\-ва\-тельность промежуточных элементов, 
если их индивидуальные имена не важны, обозначим для краткости 
$\langle a\,R\dots R\,b\rangle$. 
Количество промежуточных элементов на $1$ меньше степени отношения. 
Заметим, что не все элементы в последовательности элементов обязательно различны. 
Вполне может быть, что $(a,a)\in R^n$. 
В таком случае последовательность $\langle a\,R\dots R\,a\rangle$ 
называется \odmkey{циклом}{Цикл (cycle)}. 
Количество вхождений отношения $R$ в запись цикла 
(т.~е. степень отношения) называется \odmkey{длиной}{Длина (length)!цикла (of cycle)} цикла.
\end{frame}

\begin{frame}
\frametitle{1.4.7. Степень отношения и циклы (2/4)}
\begin{examples}  
Простейший цикл длины $2$ образует симметричная пара $(a,b)$,
такая что \\$aRb \And bRa$. Тройка элементов $(a,b,c)$, 
такая что $(a,b)\in R \And (b,c)\in R \And (c,a)\in R$, 
образует цикл длины $3$ и называется \odmkey{треугольником}{Треугольник (triangle)} и т.~д.
\end{examples}

\medskip
\begin{remark}
Строго говоря, рефлексивный элемент $a$, 
такой что $(a,a)\in R=R^1$, также образует цикл длины $1$ 
и называется \odmkey{петлёй}{Петля (loop)}. 
Однако такой цикл является тривиальным и неинтересным 
с точки зрения свойств отношения в целом. 
Действительно, если $aRa$, то $\A{n}{aR^n a}$. 
По петле можно <<накручивать>> сколь угодно длинные циклы, 
и это не даёт никакой новой информации. 
Поэтому петли обычно не считают циклами.
\end{remark}

\medskip
Отношение $R$ на множестве $A$ ($R\subset A^2$) называется \odmkey{ациклическим}{Отношение (relation)!ациклическое (acyclic)}, 
если оно не содержит циклов, 
т.~е. последовательностей $\langle a\,R\dots R\,a\rangle$. 
Свойство ацикличности тесно связано с другими свойствами отношений.
\end{frame}

\begin{frame}
\frametitle{1.4.7. Степень отношения и циклы (3/4)}
\begin{lemma} \NB{ [1]}
Если отношение $R$ на множестве $A$ антисимметрично и транзитивно, то оно ациклично.
\end{lemma} 
\begin{proof}
От противного. Пусть $\langle a\,R\,b\,R\dots R\,a\rangle$ --- цикл. Тогда, с одной стороны, 
$aRb$, но, с другой стороны, $bRa$ по транзитивности, что противоречит антисимметричности.
\end{proof}

\medskip
\begin{lemma} \NB{  [2]}
Если отношение $R$ на множестве $A$ ациклично, 
то оно не симметрично и не содержит симметричных пар $aRb \And bRa$. 
\end{lemma}
\begin{proof}
Действительно, каждая такая пара образует цикл длины 2.
\end{proof}

\medskip
Вообще говоря, допустимо рассматривать произвольные степени отношения, 
то есть последовательности элементов и, в частности, циклы произвольной длины. 
Однако это не имеет особого смысла, как показывает следующая теорема.


%\begin{proof}
%Обозначим $c_0\Assign a$, $c_k\Assign b$. По определению
%$$(a,b)\in R^k\Impl \E{\Tuple{c}{1}{k-1}\in A}{c_0Rc_1Rc_2R
%\dots Rc_{k-1}Rc_k}.$$ Далее, если $\vert A\vert = n \And k\ge
%n$, то $\E{i,j}{c_i=c_j \And i\ne j}$ и, значит,
%$$c_0Rc_1R \dots Rc_iRc_{j+1}R \dots Rc_{k-1}Rc_k,$$ то есть
%$(a,b)\in R^{k-(j-i)}$. Обозначим через $d$ процедуру, которая
%вычисляет пару чисел $(i,j)$. Существование требуемого числа $k$
%(степени отношения $R$) обеспечивается следующим построением:
%\begin{algorithmic}
%\WHILE{$k \geq n$}
%\STATE $c_0\Assign a$; $c_k\Assign b$
%\STATE $(i,j) \Assign d()$
%\STATE $k \Assign k-(j-i)$
%\ENDWHILE\hfill\qed
%\end{algorithmic}
%\renewcommand{\qed}{\vspace{-\baselineskip}}
%\end{proof}
%
%\vspace{-\baselineskip}

\end{frame}

\begin{frame}
\frametitle{1.4.7. Степень отношения и циклы (4/4)}
\begin{theorem} Если некоторая пара $(a,b)$ принадлежит какой-либо
степени отношения $R$ на множестве $A$ мощности $n$,
то эта пара принадлежит и
некоторой степени $R$ не выше $n-1$:
$$R\subset A^2 \And |A|=n \Impl
\A{a,b\in A}{\E{k}{a R^k b}\Impl\E{k < n}{a R^k b}}.
$$
\end{theorem}

\begin{proof}
Обозначим $c_0\Assign a$, $c_k\Assign b$. По определению
$$(a,b)\in R^k\Impl \E{\Tuple{c}{1}{k-1}\in A}{c_0Rc_1Rc_2R
\dots Rc_{k-1}Rc_k}.$$ 
Далее, если $\vert A\vert = n \And k\ge
n$, то по принципу Дирихле (\DMref{п.~1.2.5}{1.2.5.}) 
$\E{i,j}{c_i=c_j \And i\ne j}$, 
то есть цикл $\langle c_i R \dots Rc_j \rangle$, который можно удалить, получив последовательность меньшей длины
$c_0Rc_1R \dots Rc_iRc_{j+1}R \dots Rc_{k-1}Rc_k,$ 
то есть
$(a,b)\in R^{k-(j-i)}$.
Обозначим через $d$ процедуру, которая
вычисляет пару чисел $(i,j)$. Существование требуемого числа $k$
(степени отношения $R$) обеспечивается следующим построением:\\
\textbf{while $k \geq n$ do
$(i,j) \Assign d();$
$k \Assign k-(j-i)$
end while}
\end{proof}

\begin{corollary}
$ R\subset A^2\And\vert A\vert = n \Impl \bigcup_{i=1}^{\infty}
R^i = \bigcup_{i=1}^{n-1} R^i. $
\end{corollary}
\end{frame}


\begin{frame}
\label{1.4.8.}\subsubsection{1.4.8. Ядро отношения }\frametitle{1.4.8. Ядро отношения (1/3)}

Если $R\subset A\times B$~--- отношение между множествами $A$ и
$B$, то композиция $R\circ R^{-1}$ называется \odmkey{ядром}{Ядро
(kernel)!отношения (of relation)} отношения~$R$ и обозначается
$\ker R$. Другими словами,
$$a_1 \ker R\, a_2 \Equiv \E{b\in B}{a_1 R\, b \And a_2 R\, b}.$$

\begin{examples}
\begin{enumerate}
\item 
Пусть отношение $N_1$ <<находиться на расстоянии не более 1>> на множестве
целых чисел определено следующим образом:
$N_1\Assign \Set{(n,m)}{|n-m|\leq 1}.$
Тогда ядром этого отношения является отношение $N_2$
<<находиться на рас\-сто\-я\-нии не более~2>>:
$N_2\Assign \Set{(n,m)}{|n-m|\leq 2}.$
\item 
Ядром отношения <<быть знакомыми>> является отношение <<быть знакомыми или \\иметь общего знакомого>>. 
\item 
Ядро отношения $\in$ между $M$ и $2^M$ универсально:
$\A{a,b\in M}{\{a,b\}\in 2^M}$.
\item 
Ядро отношения $\subset$ на $2^M$ также универсально.
\end{enumerate}
\end{examples}

Ядро отношения~$R$ между $A$ и
$B$ является отношением на $A$:
$$R\subset A\times B \Impl \ker R\subset A^2.$$

\end{frame}

\begin{frame}
\frametitle{1.4.8. Ядро отношения (2/3)}
\begin{examples} (продолжение)
{\color{blue} 5.}
Рассмотрим отношение <<тесного включения>> на $2^U$:
$$X\SubsetOne Y \Def X\subset Y \And |X|+1=|Y|.$$
Ядром отношения $\SubsetOne$ является отношение
<<отличаться не более чем одним элементом>>:
$\ker \SubsetOne = \Set{X,Y\in 2^U}{|X|=|Y| \And |X\cap Y|\geq |X|-1}$.
\end{examples}





\smallskip
\begin{remark}
Термин <<ядро>> и обозначение $\ker$ используются также
и для других объектов. Их не следует путать~--- из контекста
обычно ясно, в каком смысле используется термин <<ядро>>.
\end{remark}

\smallskip
\begin{theorem}
Ядро любого отношения рефлексивно и симметрично на области опре\-де\-ления.
\end{theorem}
\begin{proof}
\proofsection{Рефлексивность} 
Пусть $a\in \Dom R$. Тогда 
$\E{b\in B}{aRb} \Impl$\\ $\Impl \E{b\in B}{aRb \And bR^{-1}a}\Impl (a,a)\in
R\circ R^{-1} \Impl a\ker R\, a$.
\proofsection{Симметричность}
Пусть $a,b\in \Dom R \And a\ker R\, b$.  Тогда
$\E{c\in B}{aRc\And c R^{-1} b} \Impl$  
$\Impl\E{c\in B}{bRc \And cR^{-1}a}$
$\Impl (b,a)\in R\circ R^{-1} \Impl b\ker R\, a$.
\end{proof}
\end{frame}



\begin{frame}
\frametitle{1.4.8. Ядро отношения (3/3)}
\begin{example}
Пусть на множестве $M=\{1,2,3,4,5,6\}$ 
задано бинарное отношение \\$R=\Set{(x,y)}{xy=6}$.
Обратное отношение совпадает с исходным, а ядро рефлексивно и симметрично, хотя и не линейно.
\begin{figure}[h]
\begin{center}
\includegraphics[height=60mm]{f1_4_8_3_3.jpg}
\end{center}
\end{figure}
\end{example} 
\end{frame}

\begin{frame}
\label{1.4.9.}\subsubsection{1.4.9. Представление отношений в программах }\frametitle{1.4.9. Представление отношений в программах (1/3)}
Пусть $R$~--- отношение на $A$, $R\subset A^2$ и $|A|=n$. 
Перенумеруем элементы множества~$A$, $A=\{\Tuple{a}{1}{n}\}$. 
Тогда отношение $R$ можно представить 
\odmkey{булевой}{Матрица (matrix)!булева (boolean)} матрицей \\
\textbf{$\boxed{R}:$ array $[1..n,1..n]$ of $0..1$}, 
где $\A{i,j\in 1..n}{\boxed{R}\,~[i,j]=a_i R a_j}$.

\begin{remark}
Термин <<булева матрица>> означает, что элементы матрицы~---
булевы значения и операции над ними выполняются по
соответствующим правилам. В
частности, если $\boxed{A}$ и~$\boxed{B}$~---  булевы матрицы, то
операции  над ними  определяются  следующим  образом:

\begin{tabbing}
\hspace{5mm}\= \odmkey{транспонирование}{Матрица (matrix)!инверсия
(inversion)}:,\quad \=\kill

\>\odmkey{транспонирование}{Матрица (matrix)!транспонирование (transposition)}:
\>$\boxed{A}^{\mathrm T} [i,j] \Def \boxed{A}~[j,i]$;\\[4pt]

\>\odmkey{инвертирование}{Матрица (matrix)!инверсия (inversion)}:
\>$\overline{\boxed{A\vphantom{\overline{A}}}}~[i,j]\Def 1-\boxed{A}~[i,j]$;\\

\>\odmkey{вычитание}{Матрица (matrix)!вычитание (subtraction)}:
\>$\left(\boxed{A}-\boxed{B}\right)[i,j]\Def
\boxed{A}~[i,j]\And\overline{\boxed{B\vphantom{\overline{B}}}}~[i,j]$;\\


\>\odmkey{умножение}{Матрица (matrix)!умножение (multiplication)}:
\>$\left(\boxed{A}\times\boxed{B}\right)[i,j]\Def\BigOr\limits_{k=1}^{n}
\left(\boxed{A}~[i,k]\And\boxed{B}~[k,j]\right)$;\\

\>\odmkey{дизъюнкция}{Матрица (matrix)!дизъюнкция (disjunction)}: 
\>$\left(\boxed{A}\Or\boxed{B}\right)[i,j]\Def
\boxed{A}~[i,j]\Or\boxed{B}~[i,j]$.
\end{tabbing}

\end{remark}
\end{frame}

\begin{frame}
\frametitle{1.4.9. Представление отношений в программах (2/3)}
%\setcounter{vartheorem}{0}
\begin{remark}
Представление отношений с помощью булевых матриц~---
это только один из возможных способов представления отношений.
Другие варианты рассматриваются
при обсуждении представления графов.
\end{remark}

В следующих утверждениях предполагается, что все рассматриваемые отношения
определены на множестве $A$, причём $|A|=n$.
\begin{theorem} \NB{  [1]}
$\boxed{R^{-1}}={\boxed{R\vphantom{R^{-1}}}}^{\;{\mathrm T}}$.
\end{theorem}
\begin{proof}
$(b,a)\!\in\! R^{-1}\!\Equiv\!(a,b)\!\in\! R \Equiv \boxed{R}~
[a,b]\!=\!1\!\Equiv\!
\boxed{R\vphantom{R^{-1}}}^{\;{\mathrm T}}[b,a]\!=\!1$.
\end{proof}

%\vspace{-.5\baselineskip}

В универсальное отношение $U$ входят все пары элементов, поэтому
все элементы матрицы универсального отношения $\boxed{U}$ равны 1.

\begin{theorem} \NB{  [2]}
$\boxed{\overline{R}}=\boxed{U}-\boxed{R}$\;.
\end{theorem}
\begin{proof}
$(a,b)\in \overline{R}\Equiv(a,b)\not\in R \Equiv
\boxed{R}~[a,b]=0\Equiv$\\ $\Equiv\boxed{U}~[a,b]-\boxed{R}~[a,b] =1
\Equiv \left(\boxed{U}- \boxed{R}\right)[a,b]=1$.
\end{proof}

%\vspace{-.5\baselineskip}

\begin{corollary}
$\boxed{\overline{R}}=\overline{\boxed{R\vphantom{\overline{R}}}}$\;.
\end{corollary}
\end{frame}

\begin{frame}
\frametitle{1.4.9. Представление отношений в программах (3/3)}
\begin{theorem} \NB{  [3]}
$\boxed{R_1\circ R_2} = \boxed{R_1} \times \boxed{R_2}$\;.
\end{theorem}

\vspace{.5\baselineskip}
\begin{proof}
$(a,b)\in R_1\circ R_2\Equiv\E{c\in A}{aR_{1}c\And cR_{2}b} \Equiv\newline\Equiv
\E{c\in A}{\boxed{R_1}~[a,c] = 1\And \boxed{R_2}~[c,b] = 1}\Equiv\newline\Equiv
\E{c\in A}{\left(\boxed{R_1}~[a,c]\And \boxed{R_2}~[c,b]\right) =
1}\Equiv\newline\Equiv \left(\BigOr\limits_{k=1}^n\boxed{R_1}~[a,k]\And
\boxed{R_2}~[k,b]\right) = 1\Equiv \left(\boxed{R_1}\times
\boxed{R_2}\right)~[a,b]~=~1$.
\end{proof}

%\vspace{-.5\baselineskip}
\vspace{.5\baselineskip}
\begin{corollary}
$
\boxed{R^k}=\boxed{R}^{\,k}.
$
\end{corollary}

\vspace{.5\baselineskip}
\begin{theorem} \NB{  [4]}
$\boxed{R_1\cup R_2}=\boxed{R_1}\Or\boxed{R_2}$\;.
\end{theorem}

\vspace{.5\baselineskip}
\begin{proof}
$(a,b)\in R_1\cup R_2\Equiv aR_1b\vee aR_2b\Equiv$\\
$\Equiv\boxed{R_1}~[a,b]=1\Or\boxed{R_2}~[a,b]=1\Equiv
\left(\boxed{R_1}\vee\boxed{R_2}\right)~[a,b]=1$.
\end{proof}
%
\end{frame}

\begin{frame}
\frametitle{Онтология наивной теории множеств}
\begin{center}
\includegraphics[height=7.5cm]{1_4_A.jpg}
\end{center}
\end{frame}

\begin{frame}
\frametitle{Онтология наивной теории множеств}
\begin{center}
\includegraphics[height=7.5cm]{1_4_C.jpg}
\end{center}
\end{frame}

\begin{frame}
\frametitle{1.5. Замыкание и сокращение отношений (1/2)}
\subsection{1.5. Замыкание и сокращение отношений}
\label{1.5.}

\odmkey{Замыкание}{Замыкание (closure)}
является весьма общим математическим понятием. 
Неформально говоря, замкнутость
означает, что многократное выполнение допустимых шагов не выводит
за определённые границы. Например,
предел сходящейся последовательности чисел из
\em{замкнутого\/} интервала $[a,b]$ обязательно
принадлежит этому интервалу, а предел сходящейся
последовательности чисел из
открытого (то есть \em{не\/} замкнутого) интервала $(a,b)$
может лежать вне этого интервала.
В некоторых случаях можно <<расширить>>
незамкнутый объект так,
чтобы он оказался замкнутым.

\medskip
\odmkey{Сокращение}{Сокращение (restriction)!отношения (of relation)}, или \odmkey{редукция}{Редукция (reduction)!отношения (of relation)},
в некотором смысле противоположно замыканию. 
Если в замыкание включается всё, что только можно получить из исходного объекта, 
то в сокращение включается минимум информации из исходного объекта, 
но так, чтобы можно было получить всё то, 
что несёт в себе исходный объект. 
Другими словами, замыкание редукции некоторого объекта 
совпадает с замыканием исходного объекта. 

\end{frame}

\begin{frame}
\frametitle{1.5. Замыкание и сокращение отношений (2/2)}

\begin{tabular}{p{11mm}|p{41mm}|p{92mm}}
  № & Название параграфа
  & Ключевые термины и обозначения \\
  \hline
\DMref{1.5.1.}{1.5.1.} & Транзитивное и рефлексивное замыкание &
{$R^+$~--- объединение положительных степеней $R$;

$R^*$~--- объединение неотрицательных степеней $R$} 
\\
\DMref{1.5.2.}{1.5.2.}\RS & Алгоритм Уоршалла &
{ }\\  
\DMref{1.5.3.}{1.5.3.} & Транзитивное сокращение &
{Редукция, минимальное представление}\\  
\DMref{1.5.4.}{1.5.4.} & Диаграммы Хассе &
{ }\\  
\DMref{1.5.5.}{1.5.5.}\BS & Симметричное замыкание и сокращение &
{ }\\  
\end{tabular}  

\end{frame}

\begin{frame}
\label{1.5.1.}\subsubsection{1.5.1. Транзитивное и рефлексивное замыкание }
\frametitle{1.5.1. Транзитивное и рефлексивное замыкание (1/2)}

Пусть $R$ и $R'$~--- отношения на множестве $M$. Отношение $R'$ называется
\odmkey{замыканием}{Замыкание (closure)!отношения (of relation)} отношения $R$
относительно свойства $C$, если

1) $R'$ обладает свойством $C$: $\quad C(R')$;

2) $R'$ является надмножеством $R$: $\quad R\subset R'$;

3) $R'$ является наименьшим:  
$\quad C(R'')\And R\subset R''\Impl R'\subset R''$.

Пусть $R^+$~--- объединение положительных степеней $R$, а
$R^*$~--- объединение не\-от\-ри\-ца\-тель\-ных степеней $R$:

\vspace{-8pt}
$$
R^+\Def\bigcup_{i=1}^{\infty}R^i,\qquad
R^*\Def\bigcup_{i=0}^{\infty}R^i.
$$

\vspace{-2pt}
\begin{theorem}
$R^+$~--- транзитивное замыкание $R$.
\end{theorem}
\begin{proof} Проверим выполнение свойств замыкания при условии,
что свойство $C$~--- это транзитивность. 
\proofsection{$C(R^+)$}
Пусть $aR^{+}b\And bR^{+}c$. Тогда $\E{n}{aR^{n}b}\And
\E{m}{bR^{m}c}$. Следовательно,\\
$\E{\Tuple{c}{1}{n-1}}{aRc_1R\dots Rc_{n-1}Rb}$ и
$\E{\Tuple{d}{1}{m-1}}{bRd_1R\dots Rd_{m-1}Rc}$. Таким образом,
$aRc_1R \dots Rc_{n-1}RbRd_1R \dots Rd_{m-1}Rc\Impl
aR^{n+m}c\Impl aR^+c$. 
\end{proof}
\end{frame}

\begin{frame}
\frametitle{1.5.1. Транзитивное и рефлексивное замыкание (2/2)}
\begin{proof} (Продолжение)
\proofsection{$R\subset R^+$} $R=R^1\Impl
R\subset\bigcup\limits_{i=1}^\infty R^i\Impl R\subset R^+$.
\proofsection{$C(R'')\And R\subset R''\Impl R^+\subset R''$} $a
R^{+} b \Impl \E{k}{a R^k b}{\Impl}$\\
$\Impl\E{\Tuple{c}{1}{k-1}}{a R
c_1 R \dots R c_{k-1}Rb}$; $R \subset R''\Impl a R''c_1
R''\dots R'' c_{k-1} R'' b \Impl a R'' b$. Таким образом,
$R^+\subset R''$.
\end{proof}

\smallskip
\begin{corollary}
$R^*$~--- рефлексивное замыкание $R^+$.

Таким образом, $R^*$~--- рефлексивное и транзитивное 
замыкание $R$.
\end{corollary}

\smallskip
Вычислить транзитивное замыкание заданного
отношения можно следующим образом:
$$
R^+=\bigcup_{i=1}^{\infty} R^i=\bigcup_{i=1}^{n-1}R^i\Impl
\boxed{R^+}=\BigOr_{i=1}^{n-1}\boxed{R^i}=\BigOr_{i=1}^{n-1}\boxed{R}^{\,i}.
$$
Такое вычисление имеет сложность $O(n^4)$.

\begin{remark}
Здесь $\BigOr$ означает групповую дизъюнкцию
(см. \DMref{п. 3.2.2}{3.2.2.}).
\end{remark}

\end{frame}

\begin{frame}
\label{1.5.2.}\subsubsection{1.5.2. Алгоритм Уоршалла }\frametitle{1.5.2. Алгоритм Уоршалла (1/2)}

\vspace{-2pt}
Следующий \odmkey{алгоритм Уоршалла}{Алгоритм (algorithm)!Уоршалла (Warshall)} вычисляет транзитивное 
замыкание  за
$O(n^3)$ шагов.

\vspace{-3pt}
\begin{algorithmic}
\REQUIRE матрица отношения $R$: \textbf{array$[1..n,1..n]$ of $0..1$}. 
\ENSURE матрица замыкания $T$: \textbf{array$[1..n,1..n]$ of $0..1$}. 
\STATE $T \Assign R$ 
\FOR{\textbf{$i$ from $1$ to $n$}} 
\FOR{\textbf{$j$ from $1$ to $n$}} 
\FOR{\textbf{$k$ from $1$ to $n$}} 
\STATE $T[j,k] \Assign T[j,k] \Or T[j,i]\And T[i,k]$ 
\ENDFOR 
\ENDFOR
\ENDFOR
\end{algorithmic}
%%\end{algorithm}

\vspace{-3pt}
\begin{ground} До начала цикла $T[j,k]=1$, если
элементы $j$ и $k$ непосредственно связаны отношением $R$
без промежуточных элементов.   
На каждом шаге основного цикла (по $i$) в транзитивное замыкание
добавляются такие пары элементов с номерами $j$ и $k$ (то~есть
$T[j,k]\Assign 1$), для которых существует последовательность
промежуточных элементов с номерами в диапазоне $1..i$,
связанных отношением $R$ (то есть $j\,R\,\dots\,R\,k$).  
\end{ground}

\end{frame}

\begin{frame}
\frametitle{1.5.2. Алгоритм Уоршалла  (2/2)}
\vspace{-4pt}
\begin{ground} (Окончание)
Последовательность
промежуточных элементов с номерами из диапазона $1..i$
(эта последовательность отнюдь не обязательно включает \em{все} элементы с номерами из диапазона $1..i$),
связанных отношением~$R$, для элементов с номерами $j$ и $k$
существует в одном из двух случаев: либо уже существует
последовательность промежуточных элементов с номерами из диапазона
$1..i$ для пары элементов с номерами $j$ и $k$, либо
существуют \em{две\/} последовательности элементов с номерами из
диапазона $1..i$~--- одна для пары элементов с номерами
$j$ и $i$, а вторая для пары элементов с номерами $i$ и $k$. Именно
это отражено в операторе 
$T[j,k]\!\Assign\!T[j,k]\!\Or\!T[j,i]\!\And\!T[i,k]$. 
После окончания основного цикла в качестве промежуточных
рассмотрены все элементы, и,~таким образом, построено транзитивное
замыкание.
Заметим, что в начале каждого шага цикла по $i$ матрица $T$ 
содержит все пары с промежуточными элементами с 
номерами из диапазона $1..(i-1)$, 
а в конце этого шага по $i$~--- 
все пары с промежуточными элементами с 
номерами из диапазона $1..i$. 
\end{ground}

\end{frame}

\begin{frame}
\label{1.5.3.}\subsubsection{1.5.3. Транзитивное сокращение }\frametitle{1.5.3. Транзитивное сокращение (1/6)} 
%\vspace{-4pt}
\odmkey{Транзитивным сокращением}{Сокращение (restriction)!отношения (of relation)!транзитивное (transitive)} 
(или \odmkey{транзитивной редукцией}{Редукция (reduction)!отношения (of relation)!транзитивная (transitive)},
или \odmkey{минимальным представлением}{Представление (representation)!отношения (of relation)!минимальное  (minimal)}) 
отношения $R$ на множестве $A$ 
называется наименьшее (по количеству элементов)
отношение, обозначаемое $R^{-}$, такое, что 
${R^{-}}^+ = R^+$. 
Транзитивное сокращение содержит минимально необходимые данные, 
по которым, используя транзитивность, 
можно восстановить транзитивное замыкание
отношения в полном объёме.

\begin{example}
На рисунке представлены (слева направо) графики отношения, 
его транзитивного замыкания и транзитивной редукции, 
определённые на множестве $1..5$.
\begin{figure}[h]
\begin{center}
\includegraphics[height=40mm]{f1_5_3_1_6.jpg}
\end{center}
\end{figure}
\end{example} 
\end{frame}

\begin{frame}
\frametitle{1.5.3. Транзитивное сокращение (2/6)} 

Задача отыскания транзитивного сокращения заданного 
отношения $R$ является практически значимой и не является тривиальной. 

Действительно, обозначим $\cal{R}(R)=\Set{S \subset A^2}{S^+ = R^+ }$~--- 
семейство отношений, транзитивные замыкания которых совпадают 
с транзитивным замыканием исходного отношения. 
Ясно, что $\cal{R}(R)\ne \emptyset$, 
поскольку $R\in\cal{R}(R)$, и ясно, что семейство отношений $\cal{R}(R)$
может быть весьма велико, грубая оценка сверху $|\cal{R}(R)|\le 2^{|A|^2}$, 
поскольку состав семейства $\cal R(R)$ не очевиден. 
Таким образом, налицо экстремальная задача: 
найти глобальный минимум в обширном пространстве поиска.
К счастью, некоторые свойства исходного отношения 
упрощают нахождение его транзитивного сокращения.


\begin{remark}
Полезно сравнить содержание этого параграфа
 с \DMref{п.~8.5.3.}{8.5.3.}
Исторически идея транзитивной редукции возникла применительно к ориентированным графам 
как средство построить минимальный по числу дуг орграф,
имеющий заданное отношение достижимости. 
В такой постановке задача транзитивной редукции особенно наглядна. 
Поэтому алгоритмы отыскания минимального 
представления удобнее формулировать на языке графов, 
хотя они применимы к любым отношениям.
\end{remark}

\end{frame}

\begin{frame}
\frametitle{1.5.3. Транзитивное сокращение (3/6)} 
\begin{lemma} \NB{ [1]}
Пусть $S,\, T \subset A^2$~--- конечные ациклические отношения 
на множестве $A$ и $S^+=T^+$. Тогда если 
$\E{(a,b)}{(a,b) \in S \And (a,b)\not\in T}$, то $(S\setminus\{(a,b)\} )^+=S^+$.
\end{lemma} 

\begin{proof}
Пусть $(a,b)\in S \And (a,b)\not\in T$. 
Поскольку $S^+=T^+$, существует последовательность 
промежуточных элементов $\langle a T\dots T b\rangle$. 
Обозначим один из этих промежуточных элементов $c$. 
Тогда существуют последовательности 
$\langle a T\dots T c\rangle$ и $\langle c T\dots T b\rangle$, 
но в силу того, что $S^+=T^+$, существуют и 
последовательности  
$\langle a S\dots S c\rangle$ и $\langle c S\dots S b\rangle$. 
Если предположить, что $b \in \langle a S\dots S c\rangle$, 
получаем две последовательности $\langle b S\dots S c\rangle$ и $\langle c S\dots S b\rangle$, 
то есть цикл, что противоречит ацикличности $S$. 
Аналогично, если предположить, 
что $a \in \langle c S\dots S b\rangle$, 
получаем две последовательности 
$\langle a S\dots S c\rangle$ и $\langle c S\dots S a\rangle$, 
то есть цикл, что также противоречит ацикличности $S$. 
Таким образом, пара $(a,b)$ для построения транзитивного 
замыкания $S$ избыточна, что и доказывает требуемое равенство $(S\setminus\{(a,b)\})^+=S^+$. 	
\end{proof}

\medskip
\begin{corollary}
Если отношение $R$ ациклично и конечно,
то его транзитивное сокращение единственно.
\end{corollary}
\begin{proof}
От противного. Пусть  
%у отношения $R$ есть два
%транзитивных сокращения  
$\E{S,T}{S=R^- \And T=R^- \And S\ne T}$.
Тогда по лемме любое из сокращений можно ещё уменьшить,
что противоречит их выбору.
\end{proof}
\end{frame}

\begin{frame}
\frametitle{1.5.3. Транзитивное сокращение (4/6)} 
\begin{lemma} \NB{  [2]}
Если отношение $R$ ациклично и конечно, то семейство $\cal{R}(R)$ 
замкнуто относительно пересечения, то есть
$\A{S, T\in \cal{R}(R)}{S \cap T \in \cal{R}(R)}$.
\end{lemma}

\begin{proof}
Пусть $S,T \in \cal{R}(R)$. 
Обозначим\\
$\{(a_1,b_1), (a_2,b_2), \dots , (a_k,b_k)\} \Assign S\setminus T = S\setminus (S\cap T)$. 
Последовательно применяя лемму 1 имеем:
$(S\setminus \{(a_1,b_1)\}\setminus \{(a_2,b_2)\}\setminus\dots\setminus \{(a_k,b_k)\})^+
=S^+=R^+$.\\
Но $S\setminus\{(a_1,b_1 )\}\setminus \{(a_2,b_2 ) \}\setminus\dots\setminus
\{(a_k,b_k ) \}= S\cap T$, значит $(S\cap T)^+= S^+=R^+$\\ и $(S\cap T)\in \cal{R}(R)$.
\end{proof}


\begin{theorem}
Если отношение $R$ ациклично и конечно, 
то транзитивная редукция $R^{-}$  
единственна и определяется по формуле

\vspace{-12pt}
$$R^{-} = \bigcap_{S\in\cal{R}(R)} S .$$
\end{theorem}

\vspace{-6pt}
\begin{proof}
Отношение $R$ конечно, поэтому и семейство $\cal{R}(R)$ конечно,
следовательно, $R^-$ всегда определено и единственно.
По лемме 2 имеем $R^-\in\cal{R}(R)$. Наконец, $\A{S\in\cal{R}(R)}{|R^-|\le|S|}$. 
\end{proof}
\end{frame}

\begin{frame}
\frametitle{1.5.3. Транзитивное сокращение (5/6)} 

\begin{corollary}
Если отношение $R$ ациклично и конечно, 
то транзитивное сокращение $R^{-}$ обладает следующими свойствами:
\begin{enumerate}
\item[1)] $R^{-} \subset R$; 
\item[2)] $\A{(a,b)\in R^{-}}{\Not\exists\langle a R^{-}\dots R^{-}b\rangle}$;
\item[3)] $R^{-} =R \setminus (R\circ R^+)$. 
\end{enumerate}
\end{corollary}

\begin{proof}
%\proofsection{1} 
\proofsection{\color{blue} 1)} $R^-=\bigcap_{S\in\cal{R}(R)} S \subseteq R$, поскольку $R\in \cal{R}(R)$;
\proofsection{\color{blue} 2)} если $\exists \langle aR^-\dots R^- b\rangle$, 
то $(R^-\setminus \{(a,b)\} )^+ = (R^-)^+$, 
и значит, пару $(a,b)$ можно исключить из транзитивного сокращения, 
что противоречит его минимальности;
\proofsection{\color{blue} 3)} cледует из того, 
что $(a,b)\in R\circ R^+ \Equiv \exists \langle aR\dots Rb\rangle$. 
\end{proof}

\begin{remark}
Транзитивное замыкание определено для любого отношения, 
единственно и содержит исходное отношение. 
Если выполнены условия теоремы, 
то транзитивное сокращение также определено, 
единственно и содержится в исходном отношении. 
Однако если условия теоремы не выполнены, 
то для транзитивного сокращения все условия  могут нарушаться.
\end{remark}

\end{frame}

\begin{frame}
\frametitle{1.5.3. Транзитивное сокращение (6/6)} 

\begin{example}
На рисунке слева представлен график отношения на множестве $\{1,2,3\}$, 
которое содержит четыре элемента (два цикла длины $2$) 
и имеет в качестве транзитивного замыкания универсальное отношение. 
При этом ни одну из пар нельзя исключить из отношения так, 
чтобы транзитивное замыкание сохранилось. 
Однако справа на этом же рисунке представлены графики двух различных отношений, 
у которых транзитивное замыкание также является универсальным отношением, 
но при этом в них всего по три элемента. 
Обратите внимание, что ни одно из транзитивных сокращений 
не является подмножеством исходного отношения! 
\begin{figure}[h]
\begin{center}
\includegraphics[height=35mm]{f1_5_3_6_6.jpg}
\end{center}
\end{figure}
\end{example} 


\end{frame}

\begin{frame}
\label{1.5.4.}\subsubsection{1.5.4. Диаграммы Хассе }\frametitle{1.5.4. Диаграммы Хассе (1/6)}
В самом общем случае \odmkey{диаграммой Хассе}{Диаграмма (diagram)!Хассе (Hasse)}  
для отношения $R$ на множестве $A$ (обозначение $\cal{H}(R)$) 
называется отношение
$$ \cal{H}(R)\Def \Set{(a,b)\in R
\And a\ne b}{\Not\E{c\in A}{a\ne c \And c\ne b \And aRc \And cRb}}.$$
Другими словами, в диаграмму Хассе заданного отношения
включаются только такие пары элементов,
которые являются <<соседними>>.  
\begin{example}
Рассмотрим отношение <<тесного включения>> на $2^U$ (см. пример 5 \DMref{п.~1.4.8}{1.4.8.}):
$$X\SubsetOne Y \Def X\subset Y \And |X|+1=|Y|.$$
Ядром отношения $\SubsetOne$ является отношение
<<отличаться не более чем одним элементом>>:
$$
\ker \SubsetOne = \Set{X,Y\in 2^U}{|X|=|Y| \And |X\cap Y|\geq |X|-1}.
$$
При этом отношение <<тесного включения>>  
является диаграммой Хассе для отношения включения на булеане $2^U$:
$
\cal{H}(\subset) = \SubsetOne.
$
\end{example}
\end{frame}


\begin{frame}
\frametitle{1.5.4. Диаграммы Хассе (2/6)}
Диаграмма Хассе, подобно транзитивному сокращению, 
более экономно представляет исходное отношение. 
Однако диаграмма Хассе не тождественна транзитивному сокращению 
хотя бы потому, что диаграмма Хассе обязательно 
является подмножеством исходного отношения, 
в то время как транзитивное сокращение 
может даже не пересекаться с исходным отношением. 
Более того, в некоторых случаях диаграммы Хассе 
может быть \em{недостаточно} для восстановления 
самого отношения или его транзитивного замыкания.
\begin{examples} На следующем рисунке: 
1) все отношения
совпадают со своими диаграммами Хассе; 
2) два отношения справа 
являются транзитивными сокращениями друг для друга; 
3) для универсального отношения на трёх элементах диаграмма Хассе пуста.

\begin{figure}[h]
\begin{center}
\includegraphics[height=25mm]{1_5_3_2.jpg}
\end{center}
\end{figure} 
\end{examples}
\end{frame}

\begin{frame}
\frametitle{1.5.4. Диаграммы Хассе (3/6)}

\begin{theorem}
Если конечное отношение $R$ антисимметрично и транзитивно, 
то $\cal{H}(R)=R^-$.
\end{theorem}
\begin{proof}
Антисимметричное транзитивное отношение ациклично (\DMref{п.~1.4.7}{1.4.7.}, лемма~1), 
значит, действует теорема \DMref{п.~1.5.3}{1.5.3.} и следствия к ней. 
По первому следствию \\$R^{-}\subset R$, 
а по второму следствию
$\A{(a,b)\in R^{-}}{\Not\exists {\langle aR^- \dots R^- b\rangle}}$, 
и тем более\\ $\A{(a,b)\in R^{-}}{\Not\E{c\in A}{a\,R\,c\,R\,b}}$. 
Таким образом, для транзитивного сокращения в данном случае 
выполняются условия определения диаграммы Хассе.
\end{proof}

\medskip
\begin{remark}
Антисимметричное и транзитивное отношение~---
это отношение \em{порядка}, \\
см.~\DMref{п.~1.8.2}{1.8.2.}. 
\end{remark}


Мы видим, что для отношения порядка
на конечном множестве   
диаграмма Хассе существует и единственна.
По этой причине обычно 
диаграммы Хассе рассматривают применительно
только к упорядоченным множествам.  



\begin{remark}
Для отношений на бесконечных множествах диаграмма Хассе может не существовать, 
даже если эти множества упорядочены.
\end{remark}

\begin{examples}
Множество вещественных чисел $\mathbb{R}$ и множество рациональных чисел $\mathbb{Q}$ 
с обычным отношением порядка не имеют диаграмм Хассе: 
между любыми двумя элементами всегда можно вставить третий. 
\end{examples}
\end{frame}

\begin{frame}
\frametitle{1.5.4. Диаграммы Хассе (4/6)}
Если отношение не является антисимметричным 
(и, тем самым, не является отношением порядка),
то диаграмма Хассе может не быть определена
разумным образом,
даже если множество конечно.
\begin{example}
Пусть $U=A^2$~--- универсальное отношение на 
конечном множестве $A$, $|A|>2$. Тогда
$\cal{H}(U)=\emptyset$,
поскольку $\A{a, b \in A}{\E{c\in A}{aUc \And cUb}}$. 
\end{example}

\medskip
Пустая диаграмма Хассе, разумеется, не может служить средством
редукции: информация об исходном отношении
теряется безвозвратно.   
Однако в некоторых случаях возможно построить
отношение, подобное по своим полезным свойствам диаграмме Хассе, но не удовлетворяющее формальному
определению.
\begin{example}
Пусть $U=A^2$~--- универсальное отношение на 
конечном множестве \\$A=\{\Tuple{a}{1}{n}\}$, $n>2$. Тогда положим
$W\Assign\{(a_1,a_2), (a_2, a_3),\ldots,
(a_{n-1},a_n), (a_n, a_1)\}$.
Имеем $W\subset U$, 
$W^+ = U^+ = U$ и 
$\A{(a, b) \in W}{\Not \E{c\in A}{aWc \And cWb}}$. 
\end{example}

\begin{remark}
На языке теории графов более подробно этот вопрос обсуждается в 
\DMref{п.~8.5.2}{8.5.2.}.
\end{remark}    

\medskip
Для любого 
%\em{ациклического } 
отношения на \em{конечном} множестве диаграмма Хассе существует, 
что показывает следующий алгоритм,
очевидным образом вытекающий из предыдущих формул.
\end{frame}

\begin{frame}
\frametitle{1.5.4. Диаграммы Хассе (5/6)}


%Алгоритм 1.12. Построение диаграммы Хассе.
\vspace{-6pt}
\begin{algorithmic}
\REQUIRE отношение
%, заданное матрицей 
$R$: \textbf{array} $[1..n, 1..n]$ \textbf{of} $0..1$.
\ENSURE диаграмма Хассе
%, заданная матрицей 
$H$: \textbf{array}$[1..n, 1..n]$ \textbf{of} $0..1$.
\FOR{$i$ \textbf{from} $1$ \textbf{to} $n$} 
\FOR{$j$ \textbf{from} $1$ \textbf{to} $n$} 
\IF{$ R[i, j] = 1$}  
\STATE {$H[i, j] \Assign 1$}
\COMMENT{\small\color{gray} возможно, это соседи }
\FOR{$k$ \textbf{from} $1$ \textbf{to} $n$} 
\IF{$ R[i, k] = 1 \And R[k, j] = 1$}  
\STATE{$H[i, j] \Assign 0;$ 
\textbf{exit for} $k$}
\COMMENT{\small\color{gray} найден промежуточный элемент }
%            \STATE{}
%            \COMMENT{\small\color{gray} можно больше не искать }				     
\ENDIF
\ENDFOR
\ELSE
\STATE{$H[i, j]\Assign 0$}
\COMMENT{\small\color{gray} диаграмма Хассе --- подмножество }
\ENDIF
\ENDFOR 
\ENDFOR
\end{algorithmic}

\end{frame}

\begin{frame}
\frametitle{1.5.4. Диаграммы Хассе (6/6)}
Особенно удобным является графическое представление, 
из-за которого диаграммы Хассе и получили название диаграмм. 
При этом элементы множества $A$ изображаются в определённых фигурах, 
например, в кругах или овалах, 
а элементы самой диаграммы,~т.~е. упорядоченные пары $(a,b)$, 
изображаются в виде стрелочки, ведущей от $a$ к $b$. 
\begin{columns}[t] % contents are top vertically aligned
\begin{column}[T]{7cm} % each column can also be its own environment
\vspace{-10pt}
\begin{figure}[h]
\begin{center}
\includegraphics[width=70mm]{f1_5_4_6_6.jpg}
\end{center}
\end{figure}
\end{column}
\begin{column}[T]{6cm}
\begin{example}
На рисунке приведён пример диаграммы Хассе 
для отношения включения на множестве $2^{\{1,2,3\}}$. 
\end{example}

Диаграммы Хассе  часто используют для представления 
упорядоченных множеств (см. 
\DMref{п.~1.8.2}{1.8.2.}) и решёток 
(см. \DMref{п.~2.6.1}{2.6.1.}).
\end{column}
\end{columns}

\end{frame}

\label{1.5.5.}\subsubsection{1.5.5. Симметричное замыкание и сокращение}
\begin{frame}
\frametitle{1.5.5. Симметричное замыкание и сокращение (1/2)}

Симметричное отношение совпадает со своим обратным (\DMref{п.~1.4.6}{1.4.6.}).
Это характеристическое свойство симметричности 
позволяет легко построить
замыкание и сокращение 
любого отношения относительно свойства симметричности.

\begin{lemma}     
Операция взятия
обратного отношения инволютивна, 
то есть $\left(R^{-1}\right)^{-1}=R$.
\end{lemma}

\begin{lemma}     
Операция взятия
обратного отношения уважает объединение,\\ 
то есть $\left(R_1 \cup R_2 \right)^{-1}=
(R_1)^{-1}\cup (R_2)^{-1}$.
\end{lemma}


\begin{theorem}
Если $R\subset M^2$~--- отношение 
на множестве $M$, то 
отношение $R\cup R^{-1}$
является симметричным замыканием 
$R$.   
\end{theorem}
\begin{proof}
Положим $S\Assign R\cup R^{-1}$, и пусть $S'$~---
 некоторое симметричное отношение, 
 являющееся надмножеством $R$.
Проверим выполнение свойств замыкания. 
\proofsection{$S$ симметрично} Имеем
$S^{-1}=\left(R\cup R^{-1}\right)^{-1}
=R^{-1} \cup \left(R^{-1}\right)^{-1} =
R^{-1} \cup R = S$ 
и по второму пункту теоремы
\DMref{п.~1.4.6}{1.4.6.} отношение $S$
симметрично.
\proofsection{$R\subset S$} По построению.
\proofsection{$S\subset S'$}
От противного.
Пусть $aSb$, но при этом $\Not aS'b$.
Если $aRb$, то $aS'b$~--- противоречие.
Если же $\Not aRb$, то $aR^{-1}b$  и
$bRa$, а значит $bS'a$ и $aS'b$~--- опять противоречие.
\end{proof}
\end{frame}

\begin{frame}
\frametitle{1.5.5. Симметричное замыкание и сокращение (2/2)}
\odmkey{Симметричным сокращением}{Сокращение (restriction)!отношения (of relation)!симметричное (symmetric)} отношения называется
такое минимальное подмножество отношения,
симметричное замыкание которого совпадает
с симметричным замыканием исходного отношения.
\begin{corollary}
Если $R\subset M^2$~--- отношение 
на множестве $M=\{\Tuple{a}{1}{n}\}$, то 
отношение $(R\setminus R^{-1})\cup
\Set{(a_i, a_j)\in R\cap R^{-1}}{i \geq j}$
является симметричным сокращением 
$R$.   
\end{corollary}
\begin{digress}
Замыкание отношений относительно 
антирефлексивности и антисимметричности невозможно определить непротиворечивым образом.
Действительно, если отношение не
антирефлексивно или не антисимметрично, то есть, если отношение
содержит хотя бы одну пару элементов,
нарушающую требуемое свойство, 
то любое надмножество исходного отношения
также не будет иметь требуемого свойства.

\smallskip
Замыкание  
относительно линейности вполне возможно,
но при этом само по себе малосодержательно.
Действительно, если есть 
пары элементов $a, b$, которые не участвуют 
в отношении, то есть $\Not aRb \And \Not bRa$,
то можно \em{произвольным} образом
доопределить отношение, и оно станет линейным.
Такое доопределение
не даёт ничего полезного.
С другой стороны,
алгоритм топологической сортировки
(\DMref{п.~1.8.4.}{1.8.4.})~---
это, фактически, построение линейного замыкания отношения частичного порядка,
то есть антисимметричного транзитивного отношения.     
\end{digress}
\end{frame}


\begin{frame}
\frametitle{Онтология наивной теории множеств}
\begin{center}
\includegraphics[height=7.5cm]{ODM1_5.png}
\end{center}
\end{frame}


\begin{frame}
\frametitle{1.6. Функции (1/2)}
\subsection{1.6. Функции}
\label{1.6.}
{\large
Понятие <<функция>> является одним из основополагающих в
математике. В~данном случае подразумеваются прежде всего функции,
отображающие одно конечное множество объектов в другое конечное
множество. Термин <<отображение>> мы считаем синонимом термина
<<функция>>, но различаем контексты употребления. Термин
<<функция>> используется в общем случае и в тех случаях, ког\-да
элементы отображаемых множеств не имеют структуры, или она не
важна. Термин <<отображение>> применяется, напротив, только в тех
случаях, ког\-да структура отображаемого множества имеет
первостепенное значение. Например, мы говорим <<булева функция>>,
но <<отображение групп>>. Такое предпочтение слову <<функция>>
оказывается в расчёте на постоянное сопоставление читателем
математического понятия функции с понятием функции в языках
программирования.
}
\end{frame}

\begin{frame}
\frametitle{1.6. Функции (2/2)}

\begin{tabular}{p{11mm}|p{41mm}|p{92mm}}
  № & Название параграфа
  & Ключевые термины и обозначения \\
  \hline
\DMref{1.6.1.}{1.6.1.} & Функциональные отношения &
{Функция, аргумент, значение, область отправления, область прибытия; $\Dom$~--- область определения,
 $\Im$~--- область значений; префиксная форма записи, многозначные отношения, тотальная и частичная функция, сужение и продолжение функции, преобразование, $n$-местная функция} 
\\
\DMref{1.6.2.}{1.6.2.}\RS & Инъекция, сюръекция, биекция & {Взаимно-однозначная функция, обратная функция}\\  
\DMref{1.6.3.}{1.6.3.} & Образы и прообразы &
{Переход к образам, переход к прообразам}\\ 
\DMref{1.6.4.}{1.6.4.} & Суперпозиция функций &
{$\circ$~--- операция суперпозиции}\\  
\DMref{1.6.5.}{1.6.5.} & Степень функции &
{ Заключительно периодическая последовательность}\\
\DMref{1.6.6.}{1.6.6.} & Представление функций в программах &
{ Оператор возврата \textbf{return}} \\ 
\end{tabular}  

\end{frame}


\begin{frame}
\label{1.6.1.}\subsubsection{1.6.1. Функциональные отношения }\frametitle{1.6.1. Функциональные отношения (1/3)}

%\vspace{-0.2\baselineskip}
Пусть $f$~--- отношение между множествами $A$ и $B$ такое, что
\vspace{-0.8\baselineskip}

$$
\A{a,b,c}{(a,b)\in f \And (a,c)\in f \Impl b=c}.
$$

\vspace{-0.2\baselineskip}

Такое свойство отношения называется
\odmkey{однозначностью}{Отношение (relation)!однозначное (single-value)}, или
\odmkey{функциональностью}{Отношение (relation)!функциональное (functional)},
а само отношение
называется \odmkey{функцией}{Функция (function)}
из $A$ в
$B$, причём для записи используется одна из следующих форм:

\vspace{-1\baselineskip}
$$
\Map{f}{A}{B}\qquad\text{или}\qquad A\To{f} B.
$$

\vspace{-0.2\baselineskip}

Если $\Map{f}{A}{B}$, то обычно используется
\odmkey{префиксная}{Запись (notation)!префиксная (prefix)} форма
записи:
$$
b=f(a)\Def (a,b)\in f.
$$

Если $b=f(a)$, то $a$ называют \odmkey{аргументом}{Аргумент
(argument)!функции (of function)}, а $b$~---
\odmkey{значением}{Значение (value)!функции (of function)}
функции. 
Если из контекста ясно, о какой функции идет речь, то
иногда используется обозначение $a \mapsto b$. 

Поскольку функция
является отношением, для функции $\Map{f}{A}{B}$
можно использовать
уже введённые понятия: множество $A$ называется \odmkey{областью
отправления}{Область отправления (source range)!функции (of
function)}, множество $B$~--- \odmkey{областью прибытия}{Область
прибытия (target range)!функции (of function)}
функции. 

\smallskip
Отношения, не обладающие свойством
однозначности, называют 
\odmkey{многозначными}{Функция (function)!многозначная (multifunction)}. 
\end{frame}

\begin{frame}
\frametitle{1.6.1. Функциональные отношения (2/3)}
Также используют термины\\ 
\odmkey{область определения}{Область определения
(domain)!функции (of function)} функции:
$\Dom f\Def\Set{a\in A}{\E{b\in B}{b=f(a)}}$;\\ 
\odmkey{область значений}{Область значений (codomain, 
range)!функции (of function)} функции:  $\Im f\Def\Set{b\in
B}{\E{a\in A}{b=f(a)}}$.

\medskip
Область определения является подмножеством области отправления,
$\Dom f \subset A$, а область значений является подмножеством
области прибытия, $\Im f \subset B$. 
Если $\Dom f=A$, то функция
называется \odmkey{тотальной}{Функция (function)!тотальная
(total)}, а если $\Dom f\ne A$~--- \odmkey{частичной}{Функция
(function)!частичная (partial)}. 

\medskip
\odmkey{Сужением}{Сужение функции (restriction of function)} функции $\Map{f}{A}{B}$ на
множество $M\subset A$ называется функция
$f|_{M}$, определяемая
следующим образом:
$$
f|_M\Def\Set{(a,b)}{(a,b)\in f\And a\in M}.
$$
Функция $f$ называется \odmkey{продолжением}{Продолжение (extension)!функции (of function)} функции $g$, если $g$ является сужением
$f$.

\medskip
{\large
Далее, если явно не оговорено противное, речь идёт о тотальных
функциях, поэтому области отправления и  определения совпадают.
}
\end{frame}

\begin{frame}
\frametitle{1.6.1. Функциональные отношения (3/3)}
\begin{remark}
Любому многозначному отношению 
$F\subset A\times B$ 
соответствует однозначная 
функция $\Map{f}{A}{2^B}$,
достаточно положить
$f(a)\Assign\Set{b\in B}{(a,b)\in F}$.
\end{remark}

\medskip
Тотальная функция
$\Map{f}{M}{M}$ называется
\odmkey{преобразованием}{Преобразование (transformation)} над~$M$.

\medskip
\begin{digress}
В математике, как правило, рассматриваются
тотальные функции, а частичные функции объявлены
<<ущербными>>. Такое соглашение 
упрощает определения и доказательства. 
Про\-грам\-мис\-ты же имеют дело
с частичными функциями, которые определены
не для всех допустимых значений.
Например, математическая функция~$y=x^2$ определена для всех $x$,
а стандартная функция языка программирования 
$\text{sqr}(x)=x*x$
даёт правильный результат отнюдь не для всех возможных
значений аргумента.
\end{digress}

\medskip
Функция $\Map{f}{A_1\times {\dots} \times A_n}{B}$ называется функцией
\odmkey{$n$ аргументов}
{Функция (function)!$n$~аргументов ($n$ arguments)},
или \odmkey{$n$-мест\-ной}{Функция (function)!$n$-местная ($n$-place)} функцией.
Такая функция отображает кортеж
$(a_1, \dots , a_n) \in A_1\times {\dots} \times A_n$ в элемент
$b\in B$,
$b=f((a_1,\dots , a_n))$.
При записи лишние скобки обычно опускают: 
$b=f(a_1,\dots , a_n)$.
\end{frame}

\begin{frame}
\label{1.6.2.}\subsubsection{1.6.2. Инъекция, сюръекция и биекция }\frametitle{1.6.2. Инъекция, сюръекция и биекция (1/5)}

Пусть $\Map{f}{A}{B}$. Тогда (тотальная) функция $f$
называется:


\odmkey{инъективной}{Функция (function)!инъективная (injective)}, или \odmkey{инъекцией}{Инъекция (injection)},
если $b\!=\!f(a)\!\And\!b\!=\!f(c) \Impl a\!=\!c$;\\ 
\odmkey{сюръективной}{Функция (function)!сюръективная  (surjective)}, или \odmkey{сюръекцией}{Сюръекция (surjection)},
если $\A{b\in B}{\E{a\in A}{b=f(a)}}$;\\
\odmkey{биективной}{Функция (function)!биективная (bijective)}, или \odmkey{биекцией}{Биекция (bijection)},
если она инъективная и сюръективная.

\begin{remark}
Биективную функцию также называют
\odmkey{взаимно однозначной}{Функция
(function)!взаимно-однозначная (one-to-one)}. Введенное в
\DMref{п.~1.2.2}{1.2.2.} вза\-им\-но-од\-но\-знач\-ное соответствие есть
биекция.
\end{remark}
%

%%\begin{remark}
\medskip
Понятия функциональности, инъективности, сюръективности и тотальности применимы к
\em{любым} отношениям, а не только к функциям. 
Таким образом, получаем
четыре свойства отношения $f\subset A\times B$, которые для
удобства запоминания определений можно свести в следующую таблицу:

%\vspace{-1.2\baselineskip}
\begin{center}
%{\sffamily\small
%\def\arraystretch{1.1}
\begin{tabular}{c|c}
%  $A$  & $B$   \\
\hline
{Функциональность} & {Инъективность}\\
$(a,b)\in f\And (a,c)\in f \Impl b = c$ & $(a,b)\in f\And (c,b)\in f \Impl a = c$\\
\hline
{Тотальность} & {Сюръективность}\\
$\A{a\in A}{\E{b\in B}{(a,b)\in f}}$ & $\A{b\in B}{\E{a\in A}{(a,b)\in f}}$\\
\hline
\end{tabular}
\end{center}
\end{frame}

\begin{frame}
\frametitle{1.6.2. Инъекция, сюръекция и биекция (2/5)}
\begin{theorem} Если $\Map{f}{A}{B}$~--- биекция,
то отношение $f^{-1}\subset B \times A$
(\odmkey{обратная}{Функция (function)!обратная (inverse)} функция)
является биекцией.
\end{theorem}
\begin{proof}
Поскольку $f$~--- биекция, имеем
$b_1 = f(a)\And b_2 = f(a)\Impl b_1=b_2,$ 
$b = f(a_1)\And b = f(a_2)\Impl a_1 = a_2,$ 
$\A{b \in  B}{\E{a \in  A}{b=f(a)}}.$
\proofsection{$f^{-1}$~--- функция} Имеем
$f^{-1}\Def\Set{(b,a)}{a \in  A \And b \in  B\And b=f(a)}$.\\
Пусть $a_1 = f^{-1}(b)\And a_2 = f^{-1}(b)$.
Тогда $b = f(a_1)\And b = f(a_2)\Impl a_1 = a_2$.
\proofsection{$f^{-1}$~--- инъекция} Пусть
$a = f^{-1}(b_1)\And a=f^{-1}(b_2)$. \\Тогда
$b_1 =f(a)\And b_2 = f(a)\Impl b_1 = b_2$.
\proofsection{$f^{-1}$~--- сюръекция} От противного. Пусть
$\E{a \in  A}{\neg\E{b \in  B}{a =  f^{-1}(b)}}$.\\
Тогда $\E{a \in  A}{\A{b \in  B}{a \not=  f^{-1}(b)}}$. Обозначим этот
элемент $a_0$. Поскольку $f$ тотальна, имеем
$\A{b}{a_0\ne f^{-1}(b)}\Impl
\A{b}{b\ne f(a_0)}\Impl {a_0\notin \Dom f = A} $.
\end{proof}

\begin{corollary}
Если $\Map{f}{A}{B}$~{\upshape---} инъективная функция, то
отношение ${f^{-1}\!\!\subset\! B\!\times\! A}$ является частичной функцией
$\Map{f^{-1}}{B}{A}$.
\end{corollary}
\begin{remark}
Важно осознавать,
что хотя отношение $f^{-1}$ существует всегда,
оно отнюдь не всегда является функцией.
\end{remark}

\end{frame} 

\begin{frame}
\frametitle{1.6.2. Инъекция, сюръекция и биекция (3/5)}

Пусть задано некоторое отношение $R$ 
из множества $A$ в множество $B$ и 
пара $(a,b)$ принадлежит отношению $R$, 
$(a,b)\in R\subset A\times B$.
В этом случае говорят, что элементы 
\odmkey{связаны}{Связь (link)}
отношением $R$. Графически связь 
изображают обычно в виде стрелочки, проведенной от элемента множества 
$A$ к элементу множества $B$.

Число связей, в которых участвует элемент,
уместно назвать его 
\odmkey{степенью}{Степень (degree)!элемента (of element)}. 
Свойства отношения можно определить исходя из ограничений на степени элементов, 
как показано в следующей таблице.

\begin{center}

\begin{tabular}{|c|c|c|}
\hline
&{Элементы в множестве $A$} & {Элементы в множестве $B$}\\
\hline
Cтепень не более 1 & {Однозначность} & {Инъективность}\\
\hline
Cтепень не менее 1 & {Тотальность} & {Сюръективность}\\
\hline
\end{tabular}
\end{center}  

Свойства функциональности, тотальности, инъективности и сюръективности независимы,
и могут быть или не быть присущи отношению
в любых комбинациях. Таким образом, 
всего возможно $2^4=16$ комбинаций свойств.
\end{frame} 

\begin{frame}
\frametitle{1.6.2. Инъекция, сюръекция и биекция (4/5)}

На  рисунке представлены примеры отношений для всех возможных 
комбинаций свойств.

\begin{figure}[h]
\begin{center}
\includegraphics[height=70mm]{f1_6_2_4_5.jpg}
\end{center}
\end{figure}
\end{frame}

\begin{frame}
\frametitle{1.6.2. Инъекция, сюръекция и биекция (5/5)}
Если область отправления $A$ и
область прибытия $B$ равномощны,
$|A|=|B|$, в частности, если $A=B$,
то не все комбинации свойств возможны,
как показано на рисунке.

\begin{figure}[h]
\begin{center}
\includegraphics[height=70mm]{f1_6_2_5_5.jpg}
\end{center}
\end{figure}
\end{frame}

%\begin{frame}
%\frametitle{1.6.2. Инъекция, сюръекция и биекция (5/6)}
%На  рисунке представлены примеры возможных 
%тотальных многозначных отношений.
%
%\begin{figure}[h]
%\begin{center}
%\includegraphics[height=70mm]{1_6_4.jpg}
%\end{center}
%\end{figure}
%\end{frame}
%\begin{frame}
%\frametitle{1.6.2. Инъекция, сюръекция и биекция (6/6)}
%На  рисунке представлены примеры  возможных 
%тотальных однозначных функций.
%
%\begin{figure}[h]
%\begin{center}
%\includegraphics[height=70mm]{1_6_5.jpg}
%\end{center}
%\end{figure}
%\end{frame}



\begin{frame}
\label{1.6.3.}\subsubsection{1.6.3. Образы и прообразы }\frametitle{1.6.3. Образы и прообразы (1/2)}

Пусть $\Map{f}{A}{B}$, и пусть
$A_1\subset A$, $B_1\subset B$. Тогда множество
$$
F(A_1)\Def\Set{b\in B}{\E{a\in A_1}{b=f(a)}}
$$
называется \odmkey{образом}{Образ (image)}
множества $A_1$ (при отображении $f$),
а множество
$$
F^{-1}(B_1)\Def\Set{a\in A}{\E{b\in B_1}{b=f(a)}}
$$
называется \odmkey{прообразом}{Прообраз (preimage)} множества
$B_1$ (при отображении $f$).

Введём обозначения $f_A\Assign\Dom
f$, $f_B\Assign\Im f$. Заметим, что $F$ является отношением между
множествами $2^{f_A}$ и $2^{f_B}$:
$$
F\Def\Set{(A_1,B_1)}{A_1\subset A\And B_1\subset B \And B_1=F(A_1)},
$$
а $F^{-1}$ является обратным отношением между
множествами $2^{f_B}$ и $2^{f_A}$:
$$
F^{-1}\Def\Set{(B_1,A_1)}{A_1\subset A\And B_1\subset B \And B_1=F(A_1)}.
$$

Заметим, что если $f$ --- функция, то для любого множества образ и прообраз определены однозначно. Из этого наблюдения немедленно следует теорема.   
\end{frame}

\begin{frame}
\frametitle{1.6.3. Образы и прообразы (2/2)}

\begin{theorem}
Если $\Map{f}{A}{B}$~{\upshape---} функция, то
$\Map{F}{2^{f_A}}{2^{f_B}}$ (\odmkey{переход к образам}{Функция
(function)!перехода к образам (image)})
и \\$\Map{F^{-1}}{2^{f_B}}{2^{f_A}}$ (\odmkey{переход к
прообразам}{Функция (function)!перехода к прообразам (preimage)})~{\upshape---} тоже функции.
\end{theorem}



%%%%%%%%%%%%%%%%%%%%%%%%% Изменения 2012
\begin{remark}
Фактически это означает, что функцию можно применять не только к
элементам области определения, но и к любым её подмножествам. На
практике в целях упрощения записи для обозначения применения
функции к  подмножеству иногда используют ту же самую букву, что и~для
применения функции к отдельному элементу.
\end{remark}

\medskip
Пусть задана функция $\Map{f}{M}{M}$,
а функции $\Map{F, F^{-1}}{2^M}{2^M}$ обозначают функции перехода
к образам и прообразам соответственно. Тогда нетрудно показать, что
$$
\A{X,Y \subset M}{F(X\cap Y) \subset F(X) \cap F(Y)}
$$
и
$$
\A{X,Y \subset M}{F(X\cup Y) \subset F(X) \cup F(Y)}.
$$
Обычно по определению полагают, что $F(\emptyset) = F^{-1}(\emptyset) = \emptyset$.
%%%%%%%%%%%%%%%%%%%%%%%%% Конец изменений 2012


\end{frame}

\begin{frame}
\label{1.6.4.}\subsubsection{1.6.4. Суперпозиция функций}
\frametitle{1.6.4. Суперпозиция функций}

Поскольку функции являются частным случаем отношений, для них
определена композиция. Композиция функций называется
\odmkey{суперпозицией}{Суперпозиция (superposition)}. Для
обозначения суперпозиции применяют тот же знак $\circ$, но
операнды записывают в обратном порядке: если $\Map{f}{A}{B}$ и
$\Map{g}{B}{C}$, то суперпозиция функций $f$ и $g$ записывается
так:  $g\circ f$. Этот способ записи суперпозиции функций
объясняется тем, что обозначение функции принято писать слева от
списка аргументов:
$$
(f\circ g)(x) \Def f(g(x)).
$$

\begin{theorem} Суперпозиция функций является функцией:
$$
{\Map{f}{A}{B}}\:\And\:{\Map{g}{B}{C}}\Impl\Map{g\circ f}{A}{C}.
$$
\end{theorem}
\begin{proof}
Пусть $b=(g\circ f)(a)$ и $c =(g\circ f)(a)$. Тогда $b =
g(f(a))$ и $c =g(f(a))$. Но $f$ и $g$~--- функции, поэтому $b=c$.
\end{proof}

%\vspace{-.5\baselineskip}

\begin{corollary}
Суперпозиция биекций является биекцией.
\end{corollary}
\end{frame}

\begin{frame}
\label{1.6.5.}\subsubsection{1.6.5. Степень функции }\frametitle{1.6.5. Степень функции (1/3)}

%%%%%%%%%%%%%%%%%%%%%%%%% Изменения 2012
Поскольку функция является отношением, можно ввести степень
функции относительно суперпозиции. \odmkey{Степенью}{Степень
(degree)!функции (of function)} $n$ функции $f$ (обозначение
$f^n$) называется $n$-кратная суперпозиция функции $f$.
Соответственно, $f^0(x) = x$, $f^1(x) = f(x)$ и вообще $f^n = f
\circ f^{n-1}$. При записи суперпозиций функций знак $\circ$ часто опускают. 


\begin{remark} В математическом анализе часто используют такое же
обозначение для степени \em{значения} функции относительно умножения.
Например, $\sin^2(x) = \sin (x) \cdot \sin (x)$. Следует различать по
контексту, в каком смысле используется степень функции. В частности, в
данном контексте  $\sin^2(x) = \sin(\sin(x))$.
\end{remark}
%%%%%%%%%%%%%%%%%%%%%%%%% Конец изменений 2012

\medskip
Пусть $f$~--- любая функция. 
Раcсмотрим последовательность степеней: 
$f^0=x$, $f^1=f$, $f^2=f f,\dots, f^n=ff^{n-1}$.
Заметим, что порядок сомножителей в определении степени несуществен.

\begin{lemma} $f f^n =  f^n f$.
\end{lemma}
\begin{proof}
По индукции. База: $f f$ =  $f f$. Пусть $f f^n =  f^n f$. 
Тогда $f^{n+1} f = (f f^n) f =$ $= f (f^n f) = f (f f^n)= f f^{n+1}$. 
\end{proof}
\end{frame}

\begin{frame}
\frametitle{1.6.5. Степень функции (2/3)}

Порядок вычисления степеней функции также несуществен, важна только сама степень.

\begin{theorem}
$f^n f^m = f^{n+m}$.
\end{theorem}
\begin{proof}
$f^n f^m = (f f^{n-1})( f^m) = (f^{n-1} f)( f^m)= (f^{n-1})(f f^m) = $\\
$=(f^{n-1})( f^{m+1}) = {\dots} = (f^0)(f^{n+m}) = f^{n+m}$.
\end{proof}

\medskip
Хотя композиция отношений в общем случае не коммутативна, 
но степени функции коммутируют.

\begin{corollary}
$ f^n f^m= f^m f^n$.
\end{corollary}

\begin{proof}
$f^n f^m= f^{n+m} = f^{m+n} = f^m f^n$.
\end{proof}

\medskip
Бесконечная последовательность элементов $\left( a_i \right)_{i=1}^\infty$ называется
\odmkey{заключительно периодической}
{Последовательность (sequence)!заключительно периодическая (eventially periodic)}, если 
$\E{m, k\in\mathbb{N}}{\A{i\ge m}{a_{i}=a_{i+k}}}$.
Число $k$ называется \odmkey{периодом}{Период (period)}. Если число
$m=1$, то последовательность называется просто \odmkey{периодической}{Последовательность (sequence)!периодическая (periodic)}.

\medskip
Пусть задана функция $\Map{f}{M}{M}$, $|M|=n$, $\Dom f = M$.
Во многих приложениях используются 
последовательности $\left(f^i(a)\right)_{i=1}^\infty$ для некоторого $a\in M$.

\end{frame} 


\begin{frame}
\frametitle{1.6.5. Степень функции (3/3)}
\vspace{-2pt}
\begin{theorem}
Если функция на конечном множестве тотальна, то
последовательность значений степеней этой функции для 
любого аргумента заключительно периодическая.
$$\Map{f}{M}{M}\And |M|=n \And \Dom f = M 
\Impl$$
$$ \A{a\in M}{\E{m, k\in \mathbb{N}}{\A{i\ge m}{f^{i}(a)=f^{i+k}(a)}}}.$$  
\end{theorem}

\begin{proof}
Рассмотрим последовательность 
$\left(f^i(a)\right)_{i=1}^\infty$ для некоторого $a\in M$.
Эта последовательность существует, поскольку функция $f$ тотальна.
По принципу Дирихле в этой последовательности есть равные элементы.
Пусть $i$ и $j$, $i<j$ --- наименьшая пара чисел, такая что
$f^i(a) = f^j(a)$. Тогда положим $m\Assign i$, 
$k\Assign j-i$. 
\end{proof}

\begin{example}
Рассмотрим числа, записанные четырьмя \em{цифрами}. 
Отсортируем цифры в порядке невозрастания и в порядке неубывания,
и вычтем меньшее число из большего. Получится число,
записанное четырьмя цифрами.   
Например, взяв число $0231$, имеем $3210 - 0123 = 3087$.
Таким образом, определена тотальная функция \\
$\Map{f}{0..9999}{0..9999}$.
Удивительный факт: если исходное число содержит различные цифры, 
то не более чем за семь применений функции $f$
получается число $6174$, которое
далее периодически повторяется.
Число $6174$ называется 
\odmkey{постоянной Капрекара}
{Постоянная (constant)!Капрекар (Kaprekar)}. 
\end{example}
%}

\end{frame}



\begin{frame}
\label{1.6.6.}\subsubsection{1.6.6. Представление функций в программах }\frametitle{1.6.6. Представление функций в программах (1/2)}

Пусть $\Map{f}{A}{B}$,
множество $A$ конечно и не очень велико: $\vert A\vert = n$. Наиболее общим
представлением такой функции является массив \textbf{array} [A] \textbf{of} B,
где A~--- тип данных, значения которого представляют элементы множества $A$,
а B~--- тип данных, значения которого представляют элементы множества $B$.
Функция нескольких аргументов представляется многомерным массивом.
\smallskip
\begin{remark}
Если среда программирования допускает массивы только с натуральными
индексами, то элементы множества $A$ нумеруются
(то~есть $A = \{\Tuple{a}{1}{n}\}$)
и функция представляется с помощью массива \textbf{array $[1..n]$ of $B$}.
Таким образом, последовательности являются частным случаем функций. 
\end{remark}




\begin{digress}
Представление функции с помощью массива является эффективным по времени,
поскольку реализация массивов в большинстве случаев обеспечивает
вычисление значения функции для заданного значения аргумента (индекса)
за постоянное время, не зависящее от размера массива
и значения индекса.
\end{digress}

\end{frame}



\begin{frame}
\frametitle{1.6.6. Представление функций в программах (2/2)}



Если множество $A$ велико или бесконечно, то для представления функций
используется особый вид процедур, возвращающих единственное значение для
заданного значения аргумента. Обычно такие процедуры также
называются <<функциями>>.
В некоторых языках программирования определения функций вводятся ключевым
словом \textbf{function}. Многоместные функции представляются с помощью нескольких
формальных параметров в определении функции.
Свойство функциональности обеспечивается
\odmkey{оператором возврата}{Оператор (statement)!возврата (return)},
часто обозначаемым ключевым словом
\textbf{return}, 
который прекращает выполнение тела функции и одновременно возвращает
значение.
\begin{digress}
В языке программирования Фортран и в некоторых других языках вызов
функции и~обращение к массиву синтаксически неразличимы, что
подчеркивает родственность этих понятий.
\end{digress}

\end{frame}

\begin{frame}
\frametitle{Онтология наивной теории множеств}
\begin{center}
\includegraphics[height=8cm]{ODM1_6.png}
\end{center}
\end{frame}


\begin{frame}
\frametitle{1.7. Отношения эквивалентности (1/2)}
\subsection{1.7. Отношения эквивалентности}
\label{1.7.} 
Различные встречающиеся на
практике отношения могут обладать (или не обладать) теми или иными
свойствами. Свойства, введённые в \DMref{п.~1.4.6}{1.4.6.},
встречаются особенно часто (потому им и даны специальные
названия). Более того, оказывается, что некоторые устойчивые
комбинации этих свойств встречаются настолько часто, что
заслуживают отдельного названия и специального изучения. 

\smallskip
Здесь
рассматриваются классы отношений, обладающих определённым набором
свойств. Такое <<абстрактное>> изучение классов отношений обладает
тем преимуществом, что, один раз установив некоторые следствия из
наличия у отношения определённого набора свойств, далее эти
следствия можно автоматически распространить на все конкретные
отношения, которые обладают данным набором свойств. Рассмотрение
начинается отношениями эквивалентности (в~этом разделе) и
продолжается отношениями порядка (в следующем разделе).
\end{frame}

\begin{frame}
\frametitle{1.7. Отношения эквивалентности (2/2)}

\begin{tabular}{p{11mm}|p{41mm}|p{92mm}}
  № & Название параграфа
  & Ключевые термины и обозначения \\
  \hline
\DMref{1.7.1.}{1.7.1.}\RS  & Классы эквивалентности &
{Отношение эквивалентности, измельчение разбиения } 
\\
\DMref{1.7.2.}{1.7.2.} & Ядро функционального отношения & { }\\  
\DMref{1.7.3.}{1.7.3.}\RS & Фактормножество &
{$\nat R(x)$~--- натуральное отображение, факторизация, индуцированная функция, разложение функции}\\  
\DMref{1.7.4.}{1.7.4} & Множество уровня &
{ }\\  
\DMref{1.7.5.}{1.7.5.}\BS & Толерантность  &
{ Класс толерантности}\\
\DMref{1.7.6.}{1.7.6.}\RS & Гомоморфизмы и изоморфизмы &
{ Теорема о гомоморфизме}\\  
\end{tabular}  

\end{frame}


\begin{frame}
\label{1.7.1.}\subsubsection{1.7.1. Классы эквивалентности }\frametitle{1.7.1. Классы эквивалентности (1/6)}
%\vspace{-4pt}
Рефлексивное симметричное
транзитивное отношение называется
\odmkey{отношением эквивалентности}{Отношение (relation)!эквивалентности (of
equivalence)}. Обычно отношение эквивалентности обозначают знаком
$\equiv$. 
%\vspace{-0.25\baselineskip}

%\vspace{-2pt}
\begin{examples}
\begin{enumerate}
\item Отношения равенства чисел, равенства множеств. \\[-2pt]
\item Отношение равномощности множеств. \\[-2pt]
\item Отношение между мультимножествами <<иметь одинаковый состав>>. \\[-2pt]
\item Отношения между конечными последовательностями: 
<<быть последовательностями над одним и
тем же мультимножеством>> 
и <<быть последовательностями над мультимножествами одного состава>>. 
При этом первое из этих отношений является собственным подмножеством второго.\\
\end{enumerate}
\end{examples}

\vspace{-2pt}
Пусть $\equiv$~--- отношение эквивалентности на множестве $M$
 и $x\in M$. 
 Подмножество элементов множества $M$, эквивалентных $x$,
называется \odmkey{классом эквивалентности}{Класс
(class)!эквивалентности (of equivalence)}
для $x$: $$[x]_{\equiv}\Def\Set{y\in M}{y\equiv x}.$$
Если отношение подразумевается, то нижний индекс опускают.

\end{frame}

\begin{frame}
\frametitle{1.7.1. Классы эквивалентности (2/6)}
\begin{lemma} \NB{[1]}
$ \A{a\in M}{[a]\not= \emptyset}$.
\end{lemma}
\begin{proof}
$a\equiv a\Impl a\in [a]$.
\end{proof}

\medskip
\begin{lemma} \NB{[2]}
$a\equiv b\Impl [a] = [b]$.
\end{lemma}
\begin{proof}
Пусть $a\equiv b$.
Тогда $x\in [a]\Equiv x\equiv a \Equiv x\equiv b\Equiv x\in [b]$.
\end{proof}

\medskip
\begin{lemma} \NB{[3]}
$a\not\equiv b\Impl [a]\cap [b] = \emptyset$.
\end{lemma}
\begin{proof}
От противного: $[a]\cap [b]\not= \emptyset \Impl$ $\E{c\in M}{c\in
[a]\cap[b]}\Impl$ \\$\Impl {c\in [a]} \And c\in [b] \Impl$
$c \equiv a\And c\equiv b\Impl$ $a\equiv c\And
c\equiv b\Impl$ $a\equiv b$, противоречие.
\end{proof}

\medskip
\begin{theorem}
Если $\equiv$~{\upshape---} отношение эквивалентности на множестве
$M$, то классы эквивалентности по этому отношению образуют
разбиение множества $M$, причём среди элементов разбиения нет
пустых. И обратно, всякое разбиение $\cal B =\{B_i\}$ множества
$M$, не содержащее пустых элементов, определяет отношение
эквивалентности на множестве $M$, классами эквивалентности
которого являются элементы разбиения.
\end{theorem}
\end{frame}

\begin{frame}
\frametitle{1.7.1. Классы эквивалентности (3/6)}
\begin{proof}
\proofsection{$\Impl$} Построение требуемого разбиения
обеспечивается следующим алгоритмом.

%\vspace{-12pt}
\begin{algorithmic}
\REQUIRE множество $M$, отношение эквивалентности $\equiv$ на $M$.
\ENSURE разбиение $\cal B$ множества $M$ на классы эквивалентности.
\STATE $\cal B\Assign\emptyset$
\COMMENT{\small\color{gray} вначале разбиение пусто }
\FOR{$a\in M$}
\FOR{$B\in \cal B$}
\STATE\textbf{select} $b\in B$
\COMMENT{\small\color{gray}  берём представителя из $B$ }
\IF{$a\equiv b$}
\STATE $B\Assign B+a$
\COMMENT{\small\color{gray}  пополняем класс эквивалентности }
\STATE \textbf{next for $a$}
\COMMENT{\small\color{gray}  переходим к новому элементу из $M$ }
\ENDIF
\ENDFOR
\STATE $\cal{B}\Assign \cal{B}\cup\left\{{\{a\}}\right\}$
\COMMENT{\small\color{gray}  вводим новый класс эквивалентности }
\ENDFOR
\end{algorithmic}
\end{proof}
\end{frame}

\begin{frame}
\frametitle{1.7.1. Классы эквивалентности (4/6)}

\begin{proof} (Продолжение)
\proofsection{$\Impll$} Положим $a\equiv
b\Assign\E{i}{a\in B_i\And b\in B_i}$. Тогда отношение $\equiv$
рефлексивно, поскольку $M = \bigcup B_i\Impl\A{a\in M}{\E{i}{a\in
B_i}}$ и $a\in B_i\Impl {a\in B_i\And} a\in B_i\Impl a\equiv
a$. 

Отношение $\equiv$ симметрично, поскольку \\${a\equiv
b}\Impl$ $\E{i}{a\in B_i\And b\in B_i}\Impl$
$\E{i}{b\in B_i\And a\in B_i}\Impl b\equiv a$. 

Отношение $\equiv$
транзитивно, поскольку 
\\$a\equiv b\And b\equiv c\Impl [a] =
[b]\And [b] = [c]\Impl [a] = [c]\Impl {a\equiv c}$.
\end{proof}

\medskip
Пусть $\cal{B}_1, \cal{B}_2$ --- два разбиения множества $M$. 
Говорят, что разбиение $\cal{B}_1$ 
есть \odmkey{измельчение}{Измельчение разбиения (refinement of partition)} разбиения $\cal{B}_2$, 
если каждый элемент разбиения $\cal{B}_2$ 
является объединением некоторых элементов разбиения $\cal{B}_1$. 
Измельчение разбиений порождает измельчение отношений эквивалентности.

Предельным случаем измельчения отношения эквивалентности 
является отношение равенства, 
когда каждый элемент эквивалентен (равен) только самому себе.

\end{frame}

\begin{frame}
\frametitle{1.7.1. Классы эквивалентности (5/6)}


\begin{examples}
\begin{enumerate}
\item На рисунке слева приведён график отношения \\ 
$\Set{(i,j)}{i\in 1..6 \And j\in 1..6 \And (i+j)\MOD 2 = 1 }$.
Далее приведены два разбиения: по столбцам и по строкам~--- и
справа приведено их общее измельчение. 
\end{enumerate}

%\vspace{-4pt}
\begin{figure}[h]
\begin{center}
\includegraphics[width=150mm]{f1_7_1_5_6.jpg}
\end{center}
\end{figure}
\end{examples}
\end{frame}

\begin{frame}
\frametitle{1.7.1. Классы эквивалентности (6/6)}
\begin{enumerate}
\item[2.] Упорядоченную пару точек 
(\odmkey{направленный отрезок}{Направленный отрезок (directed interval)} 
часто называют \odmkey{вектором}{Вектор (vector)}. 
Говорят, что векторы \odmkey{коллинеарны}{Коллинеарность (co-linear)}, 
если они лежат на параллельных прямых или на одной прямой. 
Легко видеть, что коллинеарность векторов является отношением эквивалентности. 
Тогда классы отношения эквивалентности состоят из всех коллинеарных векторов. 
Эти классы можно измельчить, различая 
сонаправленные и противонаправленные векторы. 
\item[3.] Отношение между конечными последовательностями 
<<быть последовательностями над одним и тем же мультимножеством>> 
является измельчением отношения 
<<быть последовательностями над мультимножествами одного состава>>. 
\item[4.] Отношение равенства множеств является
измельчением отношения равномощности множеств.
\end{enumerate}


\end{frame}

\begin{frame}
\label{1.7.2.}\subsubsection{1.7.2. Ядро функционального отношения }\frametitle{1.7.2. Ядро функционального отношения (1/2)}

%\vspace{-2pt}
Всякая функция $f$, будучи отношением, имеет ядро
$f\circ f^{-1}$, которое
является отношением на области определения функции.

%\medskip
\begin{remark}
Даже если $f$~--- функция, $f^{-1}$ отнюдь не обязательно
функция, поэтому здесь $\circ$~--- знак композиции отношений,
а не суперпозиции функций.
\end{remark}

\begin{theorem} Ядро функционального отношения
является отношением эквивалентности
на множестве определения.
\end{theorem}
\begin{proof}
Пусть $\Map{f}{A}{B}$. Можно считать, что ${\Dom f = A}$ и $\Im f = B$.
\proofsection{Рефлексивность} Пусть $a\in A$. Тогда $\E{b\in
B}{b=f(a)}\Impl a\in f^{-1}(b)\Impl$\\
$\Impl {(a,b)\in f\And (b,a)\in
f^{-1}}\Impl (a,a)\in f\circ f^{-1}$.
\proofsection{Симметричность} Пусть $(a_1,a_2)\in f\circ f^{-1}$.
Тогда
$\E{b}{b=f(a_1)\And a_2\in f^{-1}(b)}$,\\ 
$a_1\in f^{-1}(b)\And b=f(a_2)$ и 
$ b=f(a_2)\And a_1\in f^{-1}(b)$, откуда
${(a_2,a_1)\in f\circ f^{-1}}$.\proofsection{Транзитивность} Пусть
$(a_1,a_2)\in f\circ f^{-1}$ и $(a_2,a_3)\in f\circ f^{-1}$. Это
означает, что \\ $\E{b_1\in B}{b_1=f(a_1)\And a_2\in f^{-1}(b_1)}$ и
$\E{b_2\in B}{b_2=f(a_2)\And a_3\in f^{-1}(b_2)}$. Тогда
$b_1=f(a_1)\And b_1=f(a_2)\And b_2=f(a_2)\And b_2=f(a_3)\Impl
b_1=b_2$. Положим $b\Assign b_1$.
% (или $b\Assign b_2$). 
Тогда
$b=f(a_1)\And b=f(a_3)\Impl b=f(a_1)\And a_3\in f^{-1}(b)
\Impl{(a_1,a_3)\in f\circ f^{-1}}$.
\end{proof}
\end{frame}

\begin{frame}
\frametitle{1.7.2. Ядро функционального отношения (2/2)}
\begin{remark}
Термин <<ядро>> применительно к функции $f$
и обозначение $\ker f$ обычно резервируются для
других целей,
поэтому в формулировке предыдущей теоремы вместо
слова <<функция>> употреблён несколько
тяжеловесный оборот <<функциональное отношение>>.
\end{remark}

\begin{columns}[t]
\begin{column}[T]{5cm}
\begin{example}
Функциональное
отношение $R\subset 1..3\times 1..3$,
$y-1=(x-2)^2$, обратное отношение и
ядро  отношения. 
\end{example}
\end{column}
\begin{column}[T]{10cm}
\vspace{-12pt}
\begin{center}
\includegraphics[height=25mm]{f1_7_2_2_2.jpg}
\end{center}
\end{column}
\end{columns}


\begin{corollary}
Ядром биекции является тождественное отношение. 
\end{corollary}


\begin{columns}[t]
\begin{column}[T]{5cm}
\begin{example}
Биективное
отношение $R\subset 1..3\times 1..3$,\\
$x+y=4$, обратное отношение и
ядро  отношения. 
\end{example}
\end{column}
\begin{column}[T]{10cm}
\vspace{-12pt}
\begin{center}
\includegraphics[height=25mm]{f1_7_2_2_2_.jpg}
\end{center}
\end{column}
\end{columns}
\end{frame}

\begin{frame}
\label{1.7.3.}\subsubsection{1.7.3. Фактормножество }\frametitle{1.7.3. Фактормножество (1/3)}

Если $R$~--- отношение эквивалентности на
множестве $M$, то множество классов эквивалентности называется
\odmkey{фактормножеством}{Фактормножество (factor set)}
множества $M$ относительно
эквивалентности $R$ и обозначается $M/R$:
$
M/R\Def\{[x]_{R}\}_{x\in M}.
$
Фактормножество является подмножеством булеана: $M/R\subset 2^M$.


\begin{columns}[t] 
\begin{column}[T]{11cm} 
\begin{example}
Рассмотрим множество $A=\{1,2,3,4\}$ 
и отношение эквивалентности 
$R=\{(1,1)$, $(2,2)$, $(3,3)$, $(4,4)$,  
$(1,2)$, $(2,1)$, $(3,4)$, $(4,3)\}$.
Элементы исходного отношения 
обозначены закрашенными кружками,
элементы исходного множества обозначены квадратиками, 
а элементы фактормножества обозначены цветом
квадратиков. В данном случае в фактормножестве два элемента, поэтому использованы два цвета. 
\end{example}
\end{column}
\begin{column}[T]{4cm} 
\vspace{-12pt}
\begin{center}
\includegraphics[height=35mm]{f1_7_3_1_3.jpg}
\end{center}
\end{column}
\end{columns}

\begin{theorem}
Если $\Map{f}{A}{B}$ функция, то $|\Dom f/\left(f\circ
f^{-1}\right)|=|\Im f|$. 
\end{theorem}
\begin{proof}
Ядро функционального отношения является 
эквивалентностью на области определения,
и каждому классу эквивалентности взаимно однозначно соответствует тот элемент области значений,
в который отображаются элементы класса.  
\end{proof}

\end{frame}

\begin{frame}
\frametitle{1.7.3. Фактормножество (2/3)}
Функция  $\Map{\nat R}{M}{M/R}$ называется
\odmkey{отождествлением}{Функция (function)!отождествления
(identification)} 
(или \odmkey{натуральным отображением}{Натуральное отображение (natural representation)}, 
или \odmkey{канонической проекцией}{Проекция (projection)!каноническая (canonical)}) 
и определяется следующим образом:
$$
\nat R (x)\Def[x]_{R}.
$$

Использование фактормножества вместо множества называется 
\odmkey{факторизацией}{Факторизация (factorization)}.

\medskip
Пусть задана функция $\Map{f}{A}{B}$. 
Тогда $\ker f$~--- это эквивалентность, 
которая определяет фактормножество $A/ \ker f$
с элементами $[x]_{\ker f}$  
и функцию отождествления  $$\Map{\nat \ker f}{A}{A/\ker f}.
$$ 
Определим \odmkey{индуцированную функцию}{Функция (function)!индуцированная (induced)} 
из фактормножества в множество $B$:
$$  f/{\ker f}\Assign 
\Set{(s,b)}{s= [x]_{\ker f}   \And b=f(x)  }.$$ 
Тогда исходная функция может быть 
\odmkey{факторизована}{Факторизация (factorization)!функции (of function)}, 
или \em{разложена}, в суперпозицию
$$f=(\nat \ker f )\circ(f/{\ker f}).$$
Заметим, что при этом сначала применяется отождествление, 
которое сюръективно, а потом индуцированная функция, которая инъективна.
\end{frame}

\begin{frame}
\frametitle{1.7.3. Фактормножество (3/3)}
\begin{examples}
\NB{1.} Пусть $A$~--- множество студентов, 
а $B$~--- множество проводимых занятий (в данный момент времени). 
Отношение <<посещаемости>> между студентами и занятиями функционально: 
студент может посещать (или прогуливать) только одно занятие. 
Значит, ядро этого функционального отношения является эквивалентностью. 
Классы эквивалентности по отношению посещаемости называются группами. 
Студентов одной группы объединяет общий учебный план. 
При составлении расписания занятий обычно применяется 
факторизация по группам (а также по курсам и факультетам). 
На стене вывешивается таблица с расписанием, 
которая представляет индуцированную функцию. 
Функцию отождествления студент должен выполнить сам, 
и таким образом, удаётся компактно представить довольно 
сложное отношение (количество элементов в исходных 
множествах исчисляется тысячами).

\NB{2.} Рассмотрим отношение между точками на координатной плоскости <<лежать на одинаковом расстоянии от заданного объекта>>. Для любого объекта это отношение является отношением эквивалетности. В частности, классами эквивалентности, то есть элементами фактормножества по отношению <<лежать на одинаковом расстоянии от начала координат>>,
являются окружности с центром в начале координат. 
\end{examples}
\end{frame}

\subsubsection{1.7.4. Множества уровня }
\label{1.7.4.}
\begin{frame}
\frametitle{1.7.4. Множества уровня }

Пусть $\Map{f}{A}{B}$~--- функция и $ f\circ f^{-1}$~--- отношение
эквивалентности на $\Dom f$. 
Рассмотрим фактормножество $\Dom
f/\left(f\circ f^{-1}\right)$. 
Класс эквивалентности 
(элемент
фактормножества)
~--- это подмножество элементов $A$, которые
отображаются в один и тот же элемент $b\in B$. Такие множества
называются \odmkey{множествами уровня}{Множество (set)!уровня
(level)} функции~$f$.

%\medskip
%Напомним, что $|\Dom f/\left(f\circ
%f^{-1}\right)|=|\Im f|$. 
%
%\medskip
\begin{example}
Множество уровня 0 для функции $\sin(x)$~--- это $\pi\mathbb{Z}$.
\end{example}

Неформально множество уровня функции~$f$,
соответствующее значению $c$,~--- это множество решений уравнения
$f(x)=c$.

%
%\medskip
\begin{digress}
Термин <<множество уровня>> опирается на очевидные геометрические
ассоциации. Если нарисовать график функции одной переменной
в обычной декартовой системе координат
и провести горизонтальную прямую на определённом \em{уровне},
то абсциссы точек пересечения с графиком функции дадут множество уровня
функции.
\end{digress}

\begin{example}
Пусть задано множество $M$, $\vert M\vert = n$.
Рассмотрим функциональное отношение (мощность) $P$ между булеаном $2^M$ и
множеством неотрицательных целых
чисел:\\
$P\subset 2^M\times 0..n$,
где
$P\Assign\Set{(X,k)}{X\subset M\And k\in 0..n \And \vert X\vert = k}$.
Ядром  является отношение равномощности,
а множествами уровня~--- семейства равномощных подмножеств.
\end{example}

\end{frame}


\begin{frame}
\label{1.7.5.}\subsubsection{1.7.5. Толерантность }\frametitle{1.7.5. Толерантность (1/2)}

Рефлексивное симметричное отношение называется \odmkey{толерантностью}{Отношение (relation)!толерантности (tolerance)}. 
Отношение эквивалентности 
является частным случаем отношения толерантности.
В отличие от отношения эквивалентности, 
дающего разбиение множества элементов, 
на котором оно определено, 
отношение толерантности даёт покрытие этого множества. 
Действительно, 
пусть $T\subset A^2$~--- отношение толерантности. 
Определим \odmkey{класс толерантности}{Класс (class)!толерантности (of tolerance)} 
для элемента $x$ по отношению толерантности $T$ 
как множество $\Set{y\in A}{x\,T\,y}$. 
Тогда семейство классов толерантности образует покрытие множества $A$.


На содержательном уровне толерантность означает следующее: 
любой объект толерантен самому себе (рефлексивность), 
а толерантность двух объектов не зависит от того, 
в каком порядке они рассматриваются (симметричность). 
Но если один объект толерантен с другим, а этот другой~--- 
с третьим, то это вовсе не значит, 
что все три объекта толерантны между собой.
Другими словами, толерантность не предполагает транзитивность.
\end{frame}



\begin{frame}
\frametitle{1.7.5. Толерантность (2/2)}
\begin{examples}
\begin{enumerate}
\item 
Отношение знакомства на множестве людей.
Иванов может быть знаком с Петровым, Петров~---
с Сидоровым, но при этом Иванов и Сидоров 
могут не быть знакомы между собой.
\item 
Отношение похожести между поисковыми запросами 
в системах информационного поиска.
\item 
Отношение на множестве слов, которое задаётся 
как наличие хотя бы одной общей буквы. 
В этом отношении находятся, например, 
пересекающиеся слова кроссворда.
\end{enumerate}
\end{examples}
\end{frame}

\begin{frame}
\frametitle{Онтология наивной теории множеств}

\begin{center}
\includegraphics[height=7.5cm]{ODM1_7_1.png}
\end{center}
\end{frame}

\begin{frame}
\label{1.7.6.}\subsubsection{1.7.6. Гомоморфизмы и изоморфизмы }\frametitle{1.7.6. Гомоморфизмы и изоморфизмы (1/5)}
Пусть $R$~--- $n$-местное отношение на множестве $X$, причём 
$R\subset X^n$; 
а $S$~--- $n$-местное отношение на множестве $Y$, причём 
$S\subset Y^n$, $n>1$.
Тогда функция $\Map{f}{X}{Y}$ называется \odmkey{гомоморфизмом}{Гомоморфизм (homomorphism)},
если
$$
\A{\Tuple{x}{1}{n}\in X}{R(x_1, \dots, x_n) = S(f(x_1), \dots, f(x_n))}.
$$

\begin{example}
Пусть $R$~--- некоторое отношение эквивалентности
на множестве $A$, а $S$~--- тождественное отношение 
на фактор-множестве $A/R$. 
Тогда функция отождествления \\$\Map{\nat R}{A}{A/R}$ ~---
это гомоморфизм.    
\end{example}


\begin{remark}
Образно говорят, что гомоморфизм из множества в множество
<<уважает>> соответствующие отношения на этих множествах.
\end{remark}

\smallskip
Функция может быть гомоморфизмом и уважать одни отношения
на множествах, и при этом не уважать другие отношения
на этих же множествах.

\begin{example}
Рассмотрим функцию $n\DIV 2 + 1$ на множестве натуральных чисел.
Эта функция уважает отношение порядка $\le$ на множестве
натуральных чисел, поскольку $n\le m \Equiv (n\DIV 2 +1)\le (m\DIV 2 +1)$,
но не уважает отношение неравенства $\ne$, поскольку
$4\ne 5$, но $4\DIV 2 +1 = 3 = 5\DIV 2 +1 $.   
\end{example}

\end{frame}

\begin{frame}
\frametitle{1.7.6. Гомоморфизмы и изоморфизмы (2/5)}

Биективный гомоморфизм называется \odmkey{изоморфизмом}{Изоморфизм (isomorphism)}. 
Любая биекция $\Map{f}{X}{Y}$ уважает тождественные отношения 
$\stackrel{X}{=}$ и $\stackrel{Y}{=}$
на множествах $X$ и $Y$: $x_1 \stackrel{X}{=} x_2 \Equiv f(x_1) \stackrel{Y}{=} f(x_2)$.

\begin{example}
Взаимно-однозначное соответствие между множествами является изоморфизмом.
Равномощные множества изоморфны относительно 
тождественных отношений.
\end{example}

\smallskip
Пусть $X$ и $Y$ --- множества с определёнными на них $n$-местными отношениями $R$ и $S$ 
соответственно и $\Map{f}{X}{Y}$ --- изоморфизм.

\begin{theorem}
Если $\Map{f}{X}{Y}$ --- изоморфизм, то $\Map{f^{-1}}{Y}{X}$ --- тоже изоморфизм.
\end{theorem}
\begin{proof}
Пусть $x_1, \dots, x_n \in X$. 
Обозначим $y_1 \Assign f(x_1), \dots, y_n \Assign f(x_n)$,\\ 
где $y_1, \dots, y_n \in Y$. 
Поскольку $f$~--- биекция, имеем $x_1 = f^{-1}(y_1), \dots, x_n = f^{-1}(y_n)$. 
Тогда
$S(y_1, \dots, y_n) = S(f(x_1), \dots, f(x_n)) = R(x_1, \dots, x_n)= R(f^{-1}(y_1), \dots, f^{-1}(y_n))$.
\end{proof}

Если $\Map{f}{X}{Y}$ --- изоморфизм, то множества $X$ и $Y$ 
называют \odmkey{изоморфными}{Изоморфные (isomorphic)!множества (sets)} 
и этот факт обозначают так: $X \stackrel{f}{\sim}Y$. 
Если $f$ ясно из контекста, 
то пишут просто $X \sim Y$.

\begin{remark}
Не всякая биекция является изоморфизмом.
\end{remark}

\begin{example}
Рассмотрим функцию изменения знака $\Map{f}{\Bbb{Z}}{\Bbb{Z}}$,
$f(x)\Assign -x$. Это биекция, которая не уважает, например, отношение порядка $<$. 
\end{example}
\end{frame}


\begin{frame}
\frametitle{1.7.6. Гомоморфизмы и изоморфизмы (3/5)}

Даже если множества равномощны, $|X|=|Y|$,
отношения $R\subset X^n$ и $S\subset Y^n$
могут быть таковы, что не существует
изоморфизма, уважающего эти отношения.

\begin{example}
Пусть $|X|=|Y|$. Положим $R\Assign\stackrel{X}{=}$,
$S\Assign\stackrel{Y}{\ne}$. 
Тогда \\ $\A{\Map{f}{X}{Y}}{x_1 = x_2 
\Impl \Not\left(f(x_1)\ne f( x_2)\right)}$. 
\end{example}

\smallskip
В то же время,
если  задано $n$-местное отношение $R$ на множестве $X$
и биекция $\Map{f}{X}{Y}$,
то на множестве $Y$ всегда 
возможно задать $n$-местное отношение $S$ так,
чтобы $X\sim Y$. Для этого достаточно определить
$$S\Assign \Set{(\Tuple{y}{1}{n})\in Y^n}{R(f^{-1}(y_1),\dots,f^{-1}(y_n))}.$$

\smallskip 
Таким образом, если задано семейство множеств $\cal{X}=\{\Tuple{X}{1}{n}\}$
и на каждом множестве $X_i \in \cal{X}$ задано $n$-местное отношение $R_i\subset X_i^n$,
то некоторые пары множеств могут оказаться изоморфными: $X_i \sim X_j$,
а другие~--- нет. Всё зависит от заданных отношений. Семейство пар изоморфных множеств образует отношение 
\odmkey{изоморфности}{Отношение (relation)!изоморфности (isomorphy)} на 
семействе множеств $\cal{X}$.  

\begin{remark}
Неравномощные множества не могут быть изоморфны.
\end{remark}

\end{frame}

\begin{frame}
\frametitle{1.7.6. Гомоморфизмы и изоморфизмы (4/5)}


\begin{theorem}
Изоморфность является отношением эквивалентности.
\end{theorem}

\begin{proof}
\proofsection{Рефлексивность}
Функция $x\mapsto x$ --- требуемый изоморфизм.
\proofsection{Симметричность}
Следует из предыдущей теоремы.
\proofsection{Транзитивность}
Пусть $X$, $Y$ и $Z$ --- три множества 
с заданными на них $n$-местными отношениями $R_x$, $R_y$ и $R_z$ соответственно,
$\Map{f}{X}{Y}$ и $\Map{g}{Y}{Z}$ --- изоморфизмы. 
Тогда 
$
R_x(x_1, \dots, x_n)\!=\!R_y(f(x_1), \dots, f(x_n))\!=\!R_z(g (f(x_1)), \dots, g(f(x_n))).
$
\end{proof}

\medskip
Всякий гомоморфизм индуцирует связанный 
изоморфизм. 

\begin{theorem} \NB{ [О гомоморфизме]}
Если $\Map{f}{X}{Y}$ --- гомоморфизм, то $ \Dom{f}/\ker{f} \sim  \Im{f}$.
\end{theorem}

\begin{proof}
\proofsection{Биекция} 
Рассмотрим индуцированную функцию $\Map{g=f/\ker{f}}{\Dom{f}/\ker{f}}{Y}$.\\
Функция $g$ инъективна по построению. 
Функция $\Map{h}{\Dom{f}/\ker{f}}{\Im{f}}$, 
полученная из $g$ сокращением области прибытия, сюръективна, 
при этом так же как и $g$ инъективна, а значит, биективна.
\end{proof}
\end{frame}

\begin{frame}
\frametitle{1.7.6. Гомоморфизмы и изоморфизмы (5/5)}

\begin{proof} (Продолжение)
\proofsection{Гомоморфизм} Пусть $R$~--- исходное 
$n$-местное отношение на множестве $X$. 
Оно индуцирует новое $n$-местное отношение $R_f$ на множестве $\Dom{f}/\ker{f}$:\\
$R_f([x_1], \dots, [x_n]) \Assign R(x_1, \dots, x_n)$.
Представим функцию $f$ в виде 
$$f =(\nat{\ker{f}}) \circ (f/\ker{f})= (\nat{\ker{f}}) \circ h 
.$$ 
Имеем \\
$R_f([x_1], \dots, [x_n]) =  R(x_1, \dots, x_n)= S(f(x_1), \dots, f(x_n))= S(h([x_1]), 
\dots,h([x_n]))$.
\end{proof}

\begin{example}
Рассмотрим множество слов $A^{+}$, 
заданных над алфавитом $A$, и отношение равенства длин слов на нём. 
Также рассмотрим множество $S$ конечных мультимножеств букв из алфавита $A$ 
и отношение равномощности на нём. 
Пусть $\Map{f}{A^{+}}{S}$~--- функция, 
отображающая слово в мультимножество составляющих его букв. 
Она является гомоморфизмом относительно описанных отношений. 
Тогда по теореме о гомоморфизме множество всех множеств слов, 
составленных из одних и тех же букв, 
изоморфно множеству мультимножеств этих букв.
\end{example}

\begin{remark}
Другие примеры гомоморфизмов рассмотрены в 
\DMref{п.~2.1.6}{2.1.6.}.
\end{remark}

\end{frame}

\begin{frame}
\frametitle{Онтология наивной теории множеств}
\begin{remark}
Мощности классов эквивалентности по ядру гомоморфизма в некотором смысле характеризуют его отклонение от инъективности. Если все мощности равны $1$, то гомоморфизм инъективен.
\end{remark}

\begin{center}
\includegraphics[width=12cm]{ODM1_7_2.png}
\end{center}
\end{frame}

\begin{frame}
\frametitle{1.8. Отношения порядка (1/2)}
\subsection{1.8. Отношения порядка}
\label{1.8.}

В этом разделе определяются и описываются различные варианты отношений порядка.
Отношение порядка позволяет \em{сравнивать\/} между собой различные
элементы одного множества. Фактически, интуитивные представления
об отношениях порядка можно считать общеизвестными, и такие представления  в этом курсе уже использованы при описании
работы с упорядоченными списками.
Однако при построении математических моделей
интуитивных представлений может оказаться недостаточно, 
и необходимо более детальное описание свойств
отношений  порядка.

\bigskip
\begin{tabular}{p{11mm}|p{42mm}|p{90mm}}
  № & Название параграфа
  & Ключевые термины и обозначения \\
  \hline
\DMref{1.8.1.}{1.8.1.}\BS  & Предпорядок  &
{Отношение строгого предпорядка } 
\\
\DMref{1.8.2.}{1.8.2.}\RS & Частичный и линейный порядок & 
{ Нестрогий и строгий порядок; ЧУМ; цепочка множеств; возрастающая последовательность, высота множества}\\  
\end{tabular}  

\end{frame}

\begin{frame}
\frametitle{1.8. Отношения порядка (2/2)}

\begin{tabular}{p{11mm}|p{45mm}|p{88mm}}
  № & Название параграфа
  & Ключевые термины и обозначения \\
  \hline
\DMref{1.8.3.}{1.8.3.} & Минимальные и наименьшие элементы &
{Максимальные и наибольшие элементы; двойственность}
\\  
\DMref{1.8.4.}{1.8.4.} & Алгоритм топологической сортировки &
{ }\\  
\DMref{1.8.5.}{1.8.5.} & Верхние и нижние границы  &
{ $\sup$~--- супремум, $\inf$~---инфимум; линейно полное множество}\\
\DMref{1.8.6.}{1.8.6.} & Монотонные и непрерывные функции &
{ Возрастающая, убывающая, строго возрастающая и убывающая функция; непрерывность по Скотту}\\  
\DMref{1.8.7.}{1.8.7.}\BS & Наименьшая неподвижная точка функции &
{ Теорема о наименьшей неподвижной точке функции}\\  
\DMref{1.8.8.}{1.8.8.} & Вполне упорядоченные множества &
{ изоморфизм упорядоченных множеств}\\  
\DMref{1.8.9.}{1.8.9.}\RS & Индукция &
{ Аксиома выбора}\\  
\end{tabular}  

\end{frame}

\begin{frame}
\label{1.8.1.}\subsubsection{1.8.1. Предпорядок }\frametitle{1.8.1. Предпорядок (1/2)}
Рефлексивное транзитивное отношение называется отношением 
\odmkey{предпорядка}{Отношение (relation)!предпорядка (preorder)} 
(иногда используют термин \odmkey{квазипорядок}{Отношение (relation)!квазипорядка (quasiorder)}). 
Отношение предпорядка обычно обозначают знаком 
$\sqsubseteq$.

\smallskip
Отношение предпорядка, с одной стороны, является обобщением эквивалентности, 
а с другой стороны, является обобщением нестрогого частичного порядка 
(см. следующий параграф). 
Предпорядок, будучи транзитивным, совпадает со своим транзитивным замыканием:
$\sqsubseteq^+  = \sqsubseteq$. 

\smallskip
\begin{theorem}
Если $R$~--- отношение предпорядка, то 
$\A{n\ge 1}{\A{k>0}{R^{n+k} \subset R^n}}$.
\end{theorem}
\begin{proof} 
По теореме \DMref{п.~1.4.6}{1.4.6.} $R^2\subset R$, и значит, 
$R^3=R^2\circ R\subset R\circ R=R^2$, и так далее, 
имеем $\A{n>0}{ R^{n+1}\subset R^n}$. 
Откуда по транзитивности отношения $\sqsubseteq$ 
имеем $\A{n\ge 1}{\A{k>0}{R^{n+k} \subset R^n}}$
\end{proof}

\smallskip
\begin{corollary}
Если $R$~--- отношение предпорядка, то $R^*=R^+=R$.
\end{corollary}
\begin{proof}
По формулам параграфа \DMref{1.5.1.}{1.5.1.}
\end{proof}
\end{frame}

\begin{frame}
\frametitle{1.8.1. Предпорядок (2/2)}
Антирефлексивное транзитивное отношение называется отношением 
\odmkey{строгого предпорядка}{Отношение (relation)!предпорядка (preorder)!строгого (strict)} 
и обозначается знаком $\sqsubset$.

\smallskip 
\begin{example}
В <<Повести временных лет>> сказано:
<<И были три брата: один по имени Кий, 
другой~--- Щек и третий~--- Хорив, а сестра их была Лыбедь>>. 
Отношения на множестве людей <<$x$~брат~$y$>> и <<$x$~сестра~$y$>> являются строгими предпорядками.
На рисунке показаны диаграммы этих отношений.

\begin{figure}[h]
\begin{center}
\includegraphics[height=40mm]{f1_8_1_2_2.jpg}
\end{center}
\end{figure}


\end{example}

\end{frame}

\begin{frame}
\label{1.8.2.}\subsubsection{1.8.2. Частичный и линейный порядок }\frametitle{1.8.2. Частичный и линейный порядок (1/7)}
%\vspace{-0.2\baselineskip}
Антисимметричное транзитивное отношение называется отношением
\odmkey{порядка}{Отношение (relation)!порядка (order)}. Отношение
порядка может быть рефлексивным, и~тогда оно называется отношением
\odmkey{нестрогого порядка}{Отношение (relation)!порядка
(order)!нестрогого (non-strict)}.
Отношение порядка может быть антирефлексивным, и~тогда оно называется
отношением \odmkey{строгого порядка}{Отношение (relation)!порядка
(order)!строгого (strict)}.
Отношение порядка может быть линейным, и
тогда оно называется отношением \odmkey{линейного
порядка}{Отношение (relation)!порядка
(order)!линейного (linear)}.
Отношение порядка может не обладать свойством линейности, и тогда оно
называется отношением
\odmkey{частичного порядка}{Отношение (relation)!порядка
(order)!частичного (partial)}.
Обычно отношение строгого порядка (линейного или частичного)
обозначается знаком $\prec$, а отношение нестрогого порядка~--- знаком
$\preceq$. 


\begin{remark}
Для \em{строгого} порядка свойство антисимметричности
можно определить так:
$\A{a,b}{\Not(a\prec b)\Or \neg(b\prec a)}$,
что эквивалентно
$\A{a,b} {\Not(a\prec b \And b\prec a)}$, а если порядок ещё и \em{линейный\/}, то антисимметричность задаётся формулой
$\A{a,b}{a\prec b \Equiv \neg(b\prec a)}$.
\end{remark}

\begin{examples}
Отношение $<$ на множестве вещественных чисел
является отношением строгого линейного порядка.
Отношение $\le$ на множестве вещественных чисел является отношением
нестрогого линейного порядка.
Отношение $\subseteq$ на булеане $2^M$~--- 
нестрогий частичный порядок, 
а отношение $\SubsetNE$~---  
строгий частичный порядок. 
\end{examples}

\end{frame}

\begin{frame}
\frametitle{1.8.2. Частичный и линейный порядок (2/7)}

\vspace{-2pt}
Множество, на котором задано отношение частичного порядка,
называется \odmkey{частично упорядоченным}{Множество (set)!частично
упорядоченное (partially ordered)}.
Множество, на котором задано отношение линейного порядка,
называется
\odmkey{линейно упорядоченным}{Множество (set)!линейно упорядоченное
(linearly ordered)}.

\begin{example}
Множество вещественных чисел упорядочено линейно, а булеан упорядочен частично.
\end{example}

\begin{theorem} 
Если $\prec$~--- отношение частичного порядка, то обратное отношение $\succ$ также является отношением частичного порядка.
\end{theorem}
\begin{proof} 
\proofsection{Антисимметричность}
$\A{a, b}{(a \prec b) \And (b \prec a) \Impl a = b} \Equiv$\\
$\Equiv \A{a, b}{(b \prec a) \And (a \prec b) \Impl a = b} \Equiv$
$ \A{a, b}{(a \succ b) \And (b \succ a) \Impl a = b}$.
\proofsection{Транзитивность}
$\A{a, b, c}{(a \prec b) \And (b \prec c) \Impl a \prec c} \Equiv$\\
$ \A{a, b, c}{(b \succ a) \And (c \succ b) \Impl c \succ a} \Equiv$
$ \A{a, b, c}{(c \succ b) \And (b \succ a) \Impl c \succ a}$.
\end{proof}

\begin{corollary}
Если $\prec$ является отношением частичного порядка,
то отношения $\prec \setminus I$ и $\succ \setminus I$ являются 
отношениями строгого частичного порядка,
и отношения $\prec \cup I$ и $\succ \cup I$ являются отношениями 
нестрогого частичного порядка.
\end{corollary}
\begin{proof}
По теореме \DMref{п.~1.4.6.}{1.4.6.}
\end{proof}
\end{frame}

\begin{frame}
\frametitle{1.8.2. Частичный и линейный порядок (3/7)}
\begin{lemma}
Отношение $R$ антисимметрично и линейно тогда и только тогда,
когда дополнение отношения $\overline{R}$ совпадает
с обратным отношением $R^{-1}$ везде, 
за исключением диагонали.
\end{lemma}
\begin{proof}
По условиям леммы, исключим диагональ, 
то есть рассмотрим пары \em{различных} элементов $a$ и $b$.
\proofsection{Необходимость}
В указанных условиях свойство антисимметричности гласит:\\ 
$\A{a,b}{\Not\left(aRb \And bRa\right)}$,
а свойство линейности гласит: $\A{a,b}{aRb \Or bRa}$, 
так что из двух пар
$(a,b)$ и $(b,a)$ в точности одна принадлежит отношению 
$R$, а другая пара принадлежит как дополнению, так и 
обратному отношению.
\proofsection{Достаточность}
Если дополнение отношения $R$ совпадает с обратным отношением,
то это означает, что $\A{a,b}{\Not(aRb)\Equiv bRa}$,
то есть отношение $R$ антисимметрично. 
Линей\-ность докажем от противного. 
Пусть $\E{a,b}{\Not(aRb \Or bRa)}$. 
Тогда $\E{a,b}{\Not aRb \And \Not bRa}$,
и по уже доказанному, $\E{a,b}{ bRa \And aRb}$,
что противоречит антисимметричности.  
\end{proof}


\end{frame}

\begin{frame}
\frametitle{1.8.2. Частичный и линейный порядок (4/7)}
\begin{corollary}
Если отношение $R$ антисимметрично и линейно, то 
$$\A{a,b}{a \neq b \Impl (aRb \Equiv \Not (b R a))}.$$  
\end{corollary}
\smallskip
\begin{theorem}
Если $\prec$ является отношением линейного порядка, 
то его дополнение $\not\prec$ также является отношением линейного порядка.
\end{theorem}
\begin{proof}
По лемме $\A{a,b}{aRb \Equiv \Not (b R a)}$.
\proofsection{Линейность}
$\A{a,b}{a=b \Or a\prec b \Or b\prec a} \Equiv 
\A{a,b}{a=b \Or b\not\prec a \Or a\not\prec b}\Equiv$
\\ $\Equiv \A{a,b}{a=b \Or a\not\prec b \Or b\not\prec a}$.
\proofsection{Антисимметричность}
$\A{a,b}{a\prec b \And b\prec a \Impl a=b } \Equiv$\\ 
$\Equiv \A{a,b}{b\not\prec a \And a\not\prec b \Impl a=b } \Equiv 
\A{a,b}{a\not\prec b \And b\not\prec a \Impl a=b }$.
\proofsection{Транзитивность}
$\A{a,b,c}{c\prec b \And b\prec a  \Impl c\prec a} \Equiv$\\
$\Equiv \A{a,b,c}{b\! \not\prec\! c \And a\!\not\prec\! b \Impl a\!\not\prec\! c} \Equiv
\A{a,b,c}{a\! \not\prec\! b \And b\!\not\prec\! c \Impl a\!\not\prec\! c}$.
\end{proof}

\end{frame}

\begin{frame}
\frametitle{1.8.2. Частичный и линейный порядок (5/7)}


\begin{corollary}
Если $\preceq$ является отношением нестрогого линейного порядка, 
то дополнительное отношение $\succ$ является отношением строгого линейного порядка. 
Обратно, если $\prec$ является отношением строгого линейного порядка, 
то дополнительное отношение $\succeq$ является отношением нестрогого линейного порядка.
\end{corollary}
\begin{proof}
Достаточно проверить,
в какое отношение входит 
диагональ.
\end{proof}

\medskip
\begin{example}
Отношение $<$ является отношением строгого линейного порядка, 
а его дополнение $\geq$ является отношением 
нестрогого линейного порядка на \em{любом} подмножестве
\odmkey{числовой оси}{Числовая ось (number line)}, то есть любом подмножестве 
(конечном или бесконечном) множества чисел~---
целых, рациональных и вещественных.
\end{example}


\begin{remark}
Линейность существенна: дополнение частичного порядка не 
обязательно частичный порядок.
\end{remark}


\begin{example}
Пусть есть три конфеты: мятная, яблочная и лимонная. Пусть мятная конфета не менее вкусная, чем лимонная,
и  любая конфета не менее вкусная,
чем она сама. Рассмотрим дополнительное отношение
<<менее вкусная>>. Тогда яблочная конфета менее вкусная, чем лимонная, и в то же время,
лимонная менее вкусная, чем яблочная~--- стало быть, 
это совсем не отношение порядка.
\end{example}
\end{frame}

\begin{frame}
\frametitle{1.8.2. Частичный и линейный порядок (6/7)}


\begin{example}
Отношение $\subseteq$ на булеане является 
отношением нестрогого частичного порядка, 
но его дополнение не является отношением строгого частичного порядка. 
Более того, оно вообще не является отношением порядка,
поскольку не антисимметрично и не транзитивно.
На рисунке показаны графики отношения $\subseteq$ и его дополнения
на булеане $2^{\{1,2,3\}}$.

\begin{center}
\includegraphics[height=50mm]{f1_8_2_6_7.jpg}
\end{center}
\end{example}
\end{frame}

\begin{frame}
\frametitle{1.8.2. Частичный и линейный порядок (7/7)}
\vspace{-2pt}
Даже если частично упорядоченное множество не является линейно
упорядоченным, в нем могут найтись линейно упорядоченные подмножества.

\begin{example}
Семейство подмножеств $\{C_i\}$ множества $A$ называется
\odmkey{цепочкой}{Цепочка (chain)} в~$A$, если\\
$\A{i}{C_i\subset C_{i+1} \And C_i \ne C_{i+1}}$. Любая цепочка в $A$
образует линейно упорядоченное подмножество
частично упорядоченного множества $2^A$.
\end{example}
Линейно упорядоченное подмножество $\{a_1, \dots, a_n, \dots\}$
можно рассматривать как последовательность
$\langle a_1, \dots, a_n, \dots \rangle$,
в которой $\A{i}{a_i\prec a_{i+1} \And a_i \ne a_{i+1}}$.
Такие последовательности называются
\odmkey{строго монотонно возрастающими}
{Последовательность (sequence)!строго монотонно возрастающая
(strictly monotone increasing)}, или просто
\odmkey{возрастающими}{Последовательность (sequence)! возрастающая
(increasing)}. 

\medskip
Последовательности могут быть конечными или бесконечными.
Если последовательность конечная, количество элементов в ней называется
\odmkey{длиной}{Длина (length)!последовательности (of sequence)}.

\medskip
Говорят, что частично упорядоченное множество имеет конечную
\odmkey{высоту}{Высота (height)!множества (of set)} $k$, если любая возрастающая последовательность имеет длину не более $k$.
%%%%%%%%%%%%%%%%%%%%%%%%% Конец изменений 2012
\end{frame}

\begin{frame}
\frametitle{Онтология наивной теории множеств}
\begin{center}
\includegraphics[width=12cm]{ODM1_8_1.png}
\end{center}
\end{frame}

\begin{frame}
\label{1.8.3.}\subsubsection{1.8.3. Минимальные и наименьшие элементы }\frametitle{1.8.3. Минимальные и наименьшие элементы (1/3) }
Элемент $x$ множества $M$ с отношением порядка
$\prec$ называется \odmkey{минимальным}{Элемент
(element)!минимальный (minimal)}, если в множестве $M$ не
существует элементов, меньших, чем элемент $x$:
$$\neg\E{y \in M}{y \prec x \And y \ne x},$$
иначе говоря, 
$$\A{y \in M}{\neg (y \prec x) \Or y = x}.$$

\begin{example}
Пустое множество $\emptyset$ является минимальным элементом
булеана по включению.
\end{example}

\medskip
\begin{theorem}
Во всяком конечном непустом частично упорядоченном множестве существует
минимальный элемент.
\end{theorem}
\begin{proof}
От противного. Пусть $\neg(\E{x\in M}{\neg\E{y\in M}{y\prec x\And
y\ne x}})$. Тогда $\A{x \in  M}{\E{y \in M}{y \prec x}} \Impl 
\E{{(u_i)}^\infty_{i=1}}{\A{i}{u_{i+1} \prec u_i \And  u_{i+1} \ne u_i}}.$ 
Поскольку $\vert M\vert < \infty$, по принципу Дирихле имеем $\E{i,j}{i<j\And
u_i=u_j}$. Но по транзитивности
$$u_i\succ u_{i+1}\succ \dots\succ u_j\Impl u_{i+1}\succ u_j=u_i.$$
Таким образом, $u_{i+1}\prec u_i\And u_{i+1}\succ u_i\Impl
u_{i+1}=u_i$~--- противоречие.
\end{proof}

\end{frame}

\begin{frame}
\frametitle{1.8.3. Минимальные и наименьшие элементы (2/3) }
%\vspace{-2pt}
\begin{remark}
Линейно упорядоченное конечное множество содержит ровно один минимальный
элемент, а в произвольном конечном частично упорядоченном множестве их
может быть несколько.
\end{remark}

\begin{theorem}
Всякий частичный порядок на конечном множестве может быть дополнен до линейного.
\end{theorem}
\begin{proof}
Алгоритм топологической сортировки 
(\DMref{п.~1.8.4}{1.8.4.}).
\end{proof}
\begin{remark}
В данном случае слова <<может быть дополнен>> означают,
что существует отношение линейного порядка, которое является надмножеством
заданного отношения частичного порядка.
\end{remark}

%%%%%%%%%%%%%%%%%%%%%%%%% Изменения 2012
\medskip
Пусть $M$ --- частично упорядоченное множество с отношением
порядка $\prec$, а $X$~--- его подмножество: $X\subset M$. Элемент
$a$ называется \odmkey{наименьшим}{Элемент (element)!наименьший
(least)} в множестве $X$ (обозначение $a = \min X$), если $a \in X
\And \A{x\in X}{x \ne a \Impl a \prec x}$.

\medskip
Очевидно, что наименьший элемент, если он существует, является
минимальным. Но элемент может быть минимальным, не будучи
наименьшим. 

\end{frame}

\begin{frame}
\frametitle{1.8.3. Минимальные и наименьшие элементы (3/3) }
%\vspace{-2pt}
Более того, конечное частично упорядоченное множество,
имея минимальные элементы, может не иметь наименьшего элемента.

\begin{example}
Пусть в множестве $\{a, b, c, d\}$ задано отношение порядка $\{(a,
b), (c, d)\}$, то есть $a\prec b$ и $c \prec d$. Тогда элементы
$a$ и $c$ минимальные, но наименьшего элемента не существует.
\end{example}


\begin{theorem}
Наименьший элемент, если он существует, является единственным.
\end{theorem}
\begin{proof}
Пусть существуют два наименьших элемента, $a$ и $b$.
Тогда $a \prec b$, поскольку $a$ наименьший, и, аналогично,
$b \prec a$, поскольку $b$  наименьший.
По антисимметричности, $a = b$.
\end{proof}

\medskip
В частично упорядоченном множестве с отношением порядка $\prec$
%наряду с~минимальными и наименьшими элементами 
можно определить
\odmkey{максимальные}{Элемент (element)!максимальный (maximal)} и~\odmkey{наибольшие}{Элемент (element)!наибольший (greatest)}
элементы, используя обратное отношение $\succ$.

\begin{remark}
Используя теоремы \DMref{п.~1.8.2}{1.8.2.}, нетрудно показать, что все утверждения относительно минимальных и наименьших элементов имеют место и относительно максимальных и~наибольших элементов соответственно,
то есть имеет место \odmkey{двойственность}{Двойственность (duality)}.
\end{remark}
%%%%%%%%%%%%%%%%%%%%%%%%% Конец изменений 2012



\end{frame}

\label{1.8.4.}\subsubsection{1.8.4. Алгоритм топологической сортировки}
\begin{frame}
\frametitle{1.8.4. Алгоритм топологической сортировки (1/3)}


\vspace{-0.2\baselineskip}
Алгоритм дополнения частичного порядка до линейного на конечном
множестве. \index{Алгоритм (algorithm)!топологической сортировки
(topological sorting)}
%\begin{algorithm}
%\caption{Алгоритм топологической сортировки}
%\label{knuth}

\vspace{-0.2\baselineskip}
\begin{algorithmic}
\REQUIRE конечное частично упорядоченное множество $U$.
\ENSURE линейно упорядоченное множество $W$.
\WHILE{$U\ne\emptyset$}
\STATE  $m\Assign \min(U)$
\COMMENT{ \small\color{gray} минимальный элемент }
\STATE \textbf{yield} $m$
\COMMENT{ \small\color{gray}  включаем элемент $m$ в множество $W$ }
\STATE $U\Assign U - m$
\COMMENT{ \small\color{gray}  элемент $m$ более не рассматривается }
\ENDWHILE
\end{algorithmic}
%\end{algorithm}

\vspace{-0.2\baselineskip}
\begin{ground}
Всякая процедура, генерирующая объекты с помощью оператора
\textbf{yield}, определяет
линейный порядок на множестве своих результатов.
Действительно, последо\-ва\-тельность, в которой объекты
генерируются во время работы процедуры~--- это линей\-ный порядок.
\end{ground}

\begin{remark}
Используемая в алгоритме
функция $\min$ возвращает минимальный элемент,
существование которого га\-ран\-тируется теоремой.
При удалении этого элемента из множества
частичный порядок сохраняется.
Даже если оставшиеся элементы окажутся несравнимы,
алгоритм работает правильно, поскольку все элементы
минимальны.  
\end{remark}
\end{frame}

\begin{frame}
\frametitle{1.8.4. Алгоритм топологической сортировки (2/3)}

\vspace{-4pt}
\begin{algorithmic}
\REQUIRE конечное частично упорядоченное множество $U$.
\ENSURE минимальный элемент $m$.
\STATE \textbf{select $m\in U$}
\COMMENT{\small\color{gray} выбираем  кандидата в минимальные }
\FOR{$u\in U$}
\IF{$u\prec m$}
\STATE  $m\Assign u$
\COMMENT{\small\color{gray} меняем кандидата в минимальные }
\ENDIF
\ENDFOR
\STATE \textbf{return} $m$
\COMMENT{\small\color{gray} элемент $m$ минимальный }
\end{algorithmic}

\vspace{-4pt}
\begin{ground}
Рассмотрим последовательность элементов, которые присваивались $m$ во время работы $\min$. Эта последовательность не может быть бесконечной, и она линейно упорядочена. Рассмотрим последний элемент этой последовательности. Этот элемент $m$~--- минимальный. Действительно, от противного, пусть существует элемент $u$ такой, что $u\prec m \And u\neq m$ и $u$ не входит в последовательность. Тогда по транзитивности $u$ меньше любого члена последовательности, и он был бы включен в~последовательность в момент своего рассмотрения.
\end{ground}
\end{frame}

\begin{frame}
\frametitle{1.8.4. Алгоритм топологической сортировки (3/3)}

\begin{remark}
Если отношение порядка представлено матрицей,
то функцию $\min$ можно реализовать, например, так~---
найти в матрице отношения первую строку, содержащую только нули.
\end{remark}

\medskip
Дополнение частичного порядка до линейного
в общем случае не единственно.

\begin{example}
На рисунке слева показана диаграмма отношения частичного порядка,
а справа показаны два дополнения из четырёх возможных дополнений  этого порядка до линейного.
\begin{figure}[h]
\begin{center}
\includegraphics[width=140mm]{f1_8_4_3_3.jpg}
\end{center}
\end{figure} 
\end{example}
\end{frame}

\begin{frame}
\label{1.8.5.}\subsubsection{1.8.5. Верхние и нижние границы }\frametitle{1.8.5. Верхние и нижние границы (1/2)}
\vspace{-2pt}
Пусть $X\subset M$~--- подмножество упорядоченного
множества $M$ с отношением~порядка $\prec$. 
Элемент $m\in M$
называется \odmkey{нижней границей}{Граница (bound)!нижняя (lower)}
для подмножества $X$, если $$\A{{x\in X}}{m\ne x \Impl m\prec x}.$$ 
Элемент $m\in
M$ называется \odmkey{верхней границей}{Граница (bound)!верхняя
(upper)} для подмножества $X$, 
если $$\A{x\in X}{m\ne x \Impl x\prec m}.$$
Верхние и нижние границы не обязаны существовать для любого
множества, и если существуют, то не всегда единственны. 

%\medskip
Если
существует наибольшая нижняя граница, то она называется
\odmkey{инфимумом}{Инфимум (infimum)} и~обозначается $\inf(X)$.
Если существует наименьшая верхняя граница, то она называется
\odmkey{супремумом}{Супремум (supremum)} и обозначается $\sup(X)$.

%\medskip
\begin{example}
Рассмотрим множество рациональных чисел $\mathbb{Q}$
с обычным отношением порядка~$<$ и его подмножество
$X\Assign\Set{x\in \mathbb{Q}}{x>0 \And x^2>2}$.
Тогда $\frac{17}{12}\in X$, а $\frac{17}{13}$ является одной
из нижних границ множества $X$.
Инфимума это множество не имеет. 
Если в качестве  $X$ мы возьмем подмножество действительных, а не рациональных чисел, то 
инфимумом будет $\sqrt2$. 
\end{example}
\end{frame}

\begin{frame}
\frametitle{1.8.5. Верхние и нижние границы (2/2)}

\vspace{-2pt}
\begin{theorem}
Пусть $M$~--- частично упорядоченное множество, а~$X$~---
любое его подмножество. Тогда:

\vspace{-1pt}
\begin{enumerate}
\item[1)] Если существует наименьший элемент $a = \min X$, то~$a = \inf X$;
\item[2)] если существует инфимум $a = \inf X$ и $a\in X$, то~$a = \min X$.
\end{enumerate}
\end{theorem}
\begin{proof}

\vspace{-1pt}
\proofsection{\color{blue}1)} $\A{x \in X}{a \prec x}$, следовательно, $a$~---
нижняя граница $X$. Пусть $b$~--- любая нижняя граница $X$, тогда
$b \prec a$, так как $a \in X$. 
Следовательно, $a$~--- наибольшая нижняя граница;
\proofsection{\color{blue}2)} $\A{x \in X}{a \prec x}$, поскольку $a$~---
нижняя граница $X$, и значит, это наименьший элемент, поскольку $a \in X$.
\end{proof}

\vspace{-1pt}
\begin{corollary}
Пусть $M$ --- частично упорядоченное множество, а $X$~--- любое его подмножество. Тогда:

\vspace{-2pt} 
\begin{enumerate}
\item[1)] если существует наибольший элемент $a\! =\! \max X$, то $a\! =\! \sup X$;
\item[2)] если существует супремум $a = \sup X$ и $a \in X$, то $a = \max X$.
\end{enumerate}
\end{corollary}

\vspace{-2pt}
Говорят, что частично упорядоченное множество
\odmkey{линейно полно}{Множество (set)!линейно полное (completely ordered)},
если любое его линейно упорядоченное подмножество имеет супремум.

\vspace{-1pt}
\begin{example}
Если множество $X$ конечно, то $2^X$ --- линейно полно относительно~$\subset$.
\end{example}

\end{frame}


\begin{frame}
\label{1.8.6}\label{1.8.6.}\subsubsection{1.8.6. Монотонные и непрерывные функции }\frametitle{1.8.6. Монотонные и непрерывные функции (1/4)}

Пусть $A$ и $B$~--- упорядоченные
множества и $\Map{f}{A}{B}$.
Тогда 
функция $f$ называется\\
\odmkey{монотонно возрастающей}{Функция (function)!монотонно возрастающая
(monotone increasing)},
если\\
$
\phantom{xxxxxxxxx} a_1\prec a_2 \And a_1\neq a_2 \Impl f(a_1)\prec f(a_2) \Or f(a_1)= f(a_2);
$\\
\odmkey{монотонно убывающей}{Функция (function)!монотонно убывающая
(monotone decreasing)}, если\\ 
$
\phantom{xxxxxxxxx} a_1\prec a_2 \And a_1\neq a_2 \Impl f(a_2)\prec f(a_1) \Or f(a_1)= f(a_2);
$\\
\odmkey{строго монотонно возрастающей}{Функция (function)!строго
монотонно возрастающая (strictly monotone increasing)}, если\\
$\phantom{xxxxxxxxx} 
a_1\prec a_2 \And a_1\neq a_2 \Impl f(a_1)\prec f(a_2)
\And f(a_1)\neq f(a_2);$\\
\odmkey{строго монотонно
убывающей}{Функция (function)!строго монотонно убывающая (strictly
monotone decreasing)}, если\\
$\phantom{xxxxxxxxx} 
a_1\prec a_2 \And a_1\neq a_2 \Impl f(a_2)\prec f(a_1) \And f(a_1)\neq f(a_2).
$

Монотонно возрастаю\-щие и убывающие
функции называются \odmkey{монотонными}{Функция
(function)!монотонная (monotone)}, а строго монотонно возрастающие
и убывающие функции называются \odmkey{строго монотонными}{Функция
(function)!строго монотонная
(strictly monotone)}. 

Очевидно, что строго монотонная функция монотонна,
но обратное неверно.

\begin{remark}
Монотонная функция~--- это гомоморфизм, уважающий порядок.
\end{remark}

\end{frame}

\begin{frame}
\frametitle{\large 1.8.6. Монотонные и непрерывные функции (2/4)}


\begin{examples}
\begin{enumerate}

%{\color{blue} 1.} 
\item Тождественная функция $y=x$ для чисел является строго
монотонно возрастающей, а знак числа $\sign(x)\Def$
\textbf{if $x<0$ then $-1$ elseif $x>0$ then $1$ else $0$ end if}
является монотонно возрастающей функцией.

%{\color{blue} 2.} 

\item Пусть $2^M$~--- булеан над
линейно упорядоченным конечным множеством $M$, а~частичный порядок
на булеане~--- это включение. Тогда функция 
из алгоритма топологической сортировки, 
доставляющая
минимальный элемент, является монотонно убывающей, но не строго
монотонной.

%{\color{blue} 3.} 

\item Пусть порядок на булеане $2^M$~--- это
собственное включение. Тогда функция, которая сопоставляет
подмножеству $X\subset M$ его мощность,  является
строго монотонно возрастающей.

%{\color{blue} 4.}

\item Пусть $\Map{f}{M}{M}$ --- некоторая функция. Тогда функция
перехода к образам \\$\Map{F}{2^M}{2^M}$ 
является монотонно возрастающей по включению.
\end{enumerate}
\end{examples}


\end{frame}

\begin{frame}
\frametitle{1.8.6. Монотонные и непрерывные функции (3/4)}

\begin{theorem}
Суперпозиция одинаково монотонных функций монотонна в том же смысле.
\end{theorem}

\begin{proof}
Пусть $a \prec b$ и $g$ монотонно возрастает, тогда $g(a) \prec g(b)$.
Но $f$ также монотонно возрастает, и значит, $f(g(a)) \prec f(g(b))$. Для других случаев монотонности доказательство аналогично.
\end{proof}

\smallskip
Функция $\Map{f}{M}{M}$, где $M$ --- частично упорядоченное множество, называется \odmkey{непрерывной}{Функция (function)! непрерывная по Скотту (Scott-continuous)}
(по Скотту),
если для любой возрастающей последовательности\\
$a_1 \prec \dots \prec a_n \prec \dots$,
для которой существует супремум, выполнено равенство
$$f(\sup\{a_i\}) = \sup\{f(a_i)\}.$$

\begin{remark}
Учитывая теоремы \DMref{п.~1.8.2}{1.8.2.} и последнее замечание \DMref{п.~1.8.3}{1.8.3.},
непрерывность можно определить через строго убывающие
и имеющие инфимум последовательности.
\end{remark}

\end{frame}

\begin{frame}
\frametitle{1.8.6. Монотонные и непрерывные функции (4/4)}

\begin{theorem}
Непрерывная функция монотонно возрастает.
\end{theorem}
\begin{proof}
Пусть $a \prec b$. Тогда $\sup\{a, b\} = b$.
Имеем: \\ $f(b) = f(\sup\{a, b\}) = \sup\{f(a), f(b)\}$.
Следовательно, $f(a) \prec f(b)$.
\end{proof}

%\vspace{-.5\baselineskip}

\begin{theorem}
Суперпозиция непрерывных функций непрерывна.
\end{theorem}
\begin{proof}
Пусть строго монотонно возрастающая последовательность $\{a_n\}$
имеет супремум $a=\sup\{a_n\}$. Тогда\\
$f(g(a)) = f(g(\sup\{a_n\})) = f(\sup\{g(a_n)\}) = \sup\{f(g(a_n))\}.$
\end{proof}

%\vspace{-.5\baselineskip}

\end{frame}

\begin{frame}
\label{1.8.7.}\subsubsection{1.8.7. Наименьшая неподвижная точка функции }\frametitle{1.8.7. Наименьшая неподвижная точка функции (1/4) }


Рассмотрим функцию $\Map{f}{M}{M}$. Элемент $x \in M$ называется
\odmkey{неподвижной точкой}{Неподвижная точка (fixed point)}
функции $f$, если $f(x) = x$.

\medskip
\begin{remark}
Неподвижная точка функции $f$ --- это то же,
что и \odmkey{корень}{Корень уравнения (root of equation)}
уравнения $f(x) = x$.
\end{remark}

\begin{example}
Пусть $\Map{f}{X}{X}$. Рассмотрим функцию перехода к образам
$\Map{F}{2^X}{2^X}$.
Тогда множество 
$A \Assign \bigcap \Set{Y\!\subset\!X}{F(Y)\!\subset\!Y}$ 
(то~есть пересечение всех таких подмножеств $Y$
множества $X$, образы которых являются подмножествами аргументов)
является неподвижной точкой функции $F$. Покажем, что $F(A)
\subset A$. Действительно, \\$F(A) = F\left(\bigcap \Set{Y\subset
X}{F(Y)\subset Y}\right) \subset \cap F(Y) = A$. Отсюда по
монотонности $F$ имеем $F(F(A)) \subset F(A)$, то есть $F(A)$
удовлетворяет условию для множеств $Y$, и значит,\\ $A \subset F(A)$.
Имеем $A = F(A)$, и значит, $A$ --- неподвижная точка функции
перехода к образам $F$.
\end{example}

\medskip
\begin{remark}
Если $A= \bigcap \Set{Y\subset X}{F(Y)\subset Y} = \emptyset$, то
$F(\emptyset) = \emptyset$~--- пустое множество является неподвижной
точкой функции перехода к образам по определению.
\end{remark}
\end{frame}

\begin{frame}
\frametitle{1.8.7. Наименьшая неподвижная точка функции (2/4) }

Вообще говоря, функция $\Map{f}{X}{X}$ может не иметь неподвижных точек, иметь одну неподвижную точку или иметь их несколько,\vadjust{\penalty-1000} может быть, даже бесконечно много. Если множество $X$ упорядочено
и функция $\Map{f}{X}{X}$ имеет неподвижные точки,
то среди них можно выделить
\odmkey{наименьшую неподвижную точку}{Неподвижная точка (fixed point)!наименьшая (least)}.
Однако, даже если функция имеет неподвижные точки, среди них может не быть наименьшей просто потому, что не всякое, даже конечное, частично упорядоченное множество имеет наименьший элемент.

\begin{example}
Наименьшей неподвижной точкой для функции перехода к образам
$\Map{F}{2^X}{2^X}$ является множество $A = \bigcap \Set{Y\!\subset\!X}{F(Y)\!\subset\!Y}$, введённое в~предыдущем примере. Действительно,
пусть $\E{B \subset X}{F(B)=B \And \Not A\subset B}$. Тогда
$F(B)\subset B$, и значит, множество $B$ является одним из
множеств $Y$. Получается, что пересечение не является подмножеством
одного из пересекаемых множеств~--- противоречие.
\end{example}
\end{frame}

\begin{frame}
\frametitle{1.8.7. Наименьшая неподвижная точка функции (3/4) }

\begin{theorem} \NB{[о наименьшей неподвижной точке функции]}
Если множество $X$ частично упорядочено, линейно полно, имеет
наименьший элемент $\Bl{0} = \min X$ и функция $\Map{f}{X}{X}$
непрерывна и тотальна, $\Dom f = X$, то элемент
$a = \sup\Set{f^n(\Bl{0})}{n\ge 0}$ является наименьшей
неподвижной точкой функции $f$.
\end{theorem}
\begin{proof} Покажем сначала, что элемент $a$ существует.
Поскольку функция $f$ тотальна, элементы $f^n(\Bl{0})$
определены для любого $n\ge 0$.
Рассмотрим подмножество \\ $Y := \Set{f^n(\Bl{0})}{n \ge 0}$.
Имеем $f^0(\Bl{0})= \Bl{0} \prec f^1(\Bl{0})$,
так как $\Bl{0}$ --- наименьший элемент.
Поскольку функция $f$ непрерывна, она и монотонна, а значит, \\
$f^1(\Bl{0})= f(f^0(\Bl{0})) \prec f^2(\Bl{0}) = f(f^1(\Bl{0}))$,
$f^2(\Bl{0})= f(f^1(\Bl{0})) \prec f^3(\Bl{0}) = f(f^2(\Bl{0}))$
и вообще $f^n(\Bl{0}) \prec f^{n+1}(\Bl{0})$
для любого $n \ge 0$.
Таким образом, подмножество $Y$ линейно упорядочено,
а значит, элемент $a = \sup\Set{f^n(\Bl{0})}{n \ge 0}$ определён по условию теоремы.
\end{proof}
\end{frame}

\begin{frame}
\frametitle{1.8.7. Наименьшая неподвижная точка функции (4/4) }
\begin{proof} (Продолжение)
Покажем теперь, что элемент $a$ действительно является неподвижной точкой функции $f$, то есть $f(a) = a$. Имеем
\begin{align*}
f(a) &= f(\sup\Set{f^n(\Bl{0})}{n \ge 0})=
\sup\Set{f(f^n(\Bl{0}))}{n \ge 0} = \\
&= \sup\Set{f^n(\Bl{0})}{n \ge 1}
= \sup\{\Bl{0}, a\} = a.
\end{align*}
Покажем теперь, что $a$ --- наименьшая неподвижная точка.
Пусть $b$~--- любая неподвижная точка.
Поскольку $f(b)=b$, имеем $\A{n\ge 0}{f^n(b) =b}$. Но $\Bl{0}\prec b$, 
и функция $f$ монотонна, следовательно, $\A{n\ge 0}{f^n(\Bl{0})\prec b}$,
то есть $b$ --- верхняя граница множества
$\Set{f^n(\Bl{0})}{n \ge 0}$, а значит, $a \prec b$.
\end{proof}

\begin{remark}
Требование непрерывности в формулировке теоремы нужно только для того, чтобы гарантировать существование супремумов. Вообще говоря, если супремумы заведомо существуют, то от функции требуется только монотонность.
\end{remark}

\begin{digress}
Фактически, приведенная теорема о наименьшей неподвижной точке функции
служит обоснованием одного из случаев, когда для решения уравнения $x=f(x)$ правомерно применять метод итераций: $x_{i+1}:=f(x_i)$.
\end{digress}
%%%%%%%%%%%%%%%%%%%%%%%%% Конец изменений 2012

\end{frame}

\begin{frame}
\label{1.8.8.}\subsubsection{1.8.8. Вполне упорядоченные множества }
\frametitle{1.8.8. Вполне упорядоченные множества (1/3)}
\vspace{-2pt}
Частично упорядоченное множество $X$ называется
\odmkey{вполне упорядоченным}{Множество (set)!вполне упорядоченное
(well-ordered)},
если любое его непустое подмножество имеет наименьший элемент.
В частности, для любых двух элементов $a,b\in X$ один из
них обязан быть наименьшим в подмножестве 
$\{a,b\}$,
а значит, 
вполне упорядоченное
множество упорядочено линейно.

%\medskip
\begin{remark}
Не надо путать вполне упорядоченные множества и множества,
на которых определён линейный  порядок.
Линейно упорядоченное множество может не быть
вполне упорядоченным.
\end{remark}

%\medskip
\begin{examples}
{\color{blue} 1.}
Всякое конечное линейно упорядоченное множество
вполне упорядочено.
\\{\color{blue} 2.}
Множество натуральных чисел $\mathbb{N}$ с обычным
отношением порядка вполне упорядочено.
\\{\color{blue} 3.}
Множество рациональных чисел $\mathbb{Q}$ с обычным
отношением порядка является линейно упорядоченным, но не является
вполне упорядоченным, так как множество\\
$X\Assign \Set{x\in \mathbb{Q}}{x^2>2}$ 
(равно как и само множество $\mathbb{Q}$)
не имеет наименьшего элемента.

{\color{blue} 4.}
Множество вещественных чисел $\mathbb{R}$
с обычным отношением порядка не является
вполне упорядоченным, так как множество $X\Assign \Set{x\in \mathbb{R}}{x>0}$ 
(равно как и само множество $\mathbb{R}$)
не имеет наименьшего элемента.

\end{examples}
\end{frame}

\begin{frame}
\frametitle{1.8.8. Вполне упорядоченные множества (2/3)}

Говорят, что два линейно упорядоченных множества $A$ и $B$
\odmkey{изоморфны}{Изоморфизм (isomorphism)!линейно упорядоченных
множеств (of linear ordered sets)} (обозначение $A\sim B$), если
между ними существует взаимно-однозначное соответствие,
сохраняющее порядок:

\vspace{-12pt}
$$
A \sim B \Def |A| = |B| \And  (\A{a_1,a_2 \in A}
{a_1 \prec a_2 \And a_1 
\mapsto b_1 \And  a_2 \mapsto  b_2 \Impl b_1 \prec b_2}).
$$


%\vspace{-12pt}
\begin{theorem}
Линейно упорядоченные множества 
$A$ и $B$
изоморфны тогда и только тогда,
когда между ними существует строго монотонное
вза\-им\-но-од\-но\-знач\-ное соответствие.
\end{theorem}

\begin{proof}
Пусть $A$ и $B$~--- линейно упорядоченные множества с отношениями
порядка $<_A$ и $<_B$.
\proofsection{$\Longrightarrow$}
Если $A\sim B$, то $\E{\Map{f}{A}{B}}{a_1 <_A a_2 
\Impl f(a_1) <_B f(a_2)}$, 
причём $f$~--- биекция, то есть требуемое строго монотонно возрастающее 
взаимно-однозначное соответствие.
\proofsection{$\Longleftarrow$}
Если $\Map{f}{A}{B}$~--- строго монотонно возрастающее 
взаимно-однозначное соответствие, то
оно является требуемым изоморфизмом, и $A\sim B$.
Если же $\Map{f}{A}{B}$~--- строго монотонно убывающее 
взаимно-однозначное соответствие, то
достаточно вместо отношения $<_B$ рассмотреть обратное отношение $>_B$.   
\end{proof}
\end{frame}

\begin{frame}
\frametitle{1.8.8. Вполне упорядоченные множества (3/3)}
\begin{remark}
Понятие \odmkey{изоморфизма}{Изоморфизм (isomorphism)!вполне
упорядоченных множеств (of well-ordered sets)} применимо и к
вполне упорядоченным множествам, поскольку они упорядочены линейно.
\end{remark}

\medskip
\begin{example}
Вполне упорядоченные множества натуральных чисел и чётных чисел с
обычным порядком $<$ изоморфны, поскольку
соответствие $n\mapsto 2n$ является строго
монотонно возрастающим.
\end{example}
\end{frame}

\begin{frame}
\label{1.8.9.}\subsubsection{1.8.9. Индукция }\frametitle{1.8.9. Индукция (1/2)}
Важность вполне упорядоченных множеств
определяется тем, что для них можно обобщить принцип
\odmkey{индукции}{Принцип (principle)!индукции (of induction)},
применяемый обычно для множества натуральных чисел.
\begin{theorem}
\NB{[Индукция по вполне упорядоченному множеству]}
Пусть
$X$~--- вполне упорядоченное множество, $x_0$~--- его наименьший элемент, а $P$~---
некоторый предикат, зависящий от элементов $X$. Тогда если
$$
P(x_0)\And\A{x_1\in X}{\left(\A{x\in X}{x\prec x_1 \Impl P(x)}\right) \Impl P(x_1)},
$$
то
$$
\A{x\in X}{P(x)}.
$$
\end{theorem}
\begin{proof}
Обозначим $\overline{P}\Assign\Set{x\in X}{\Not P(x)}$, $\overline{P}\subset X$.
Далее от противного. Пусть $\overline{P}\neq\emptyset$.
Поскольку $X$ вполне упорядочено, 
$\overline{P}$ имеет
наименьший элемент, обозначим его $x_1$. Тогда
$\A{x\in X}{x\prec x_1 \Impl P(x)}$ по выбору $x_1$, и значит,
$P(x_1)$ по условию теоремы, что противоречит выбору $x_1$.
\end{proof}

\end{frame}

\begin{frame}
\frametitle{1.8.9. Индукция (2/2)}
\begin{remark}
Обычная математическая индукция соответствует индукции
по вполне упорядоченному множеству  $\mathbb{N}$.
\end{remark}

\medskip
Индукция по вполне упорядоченному множеству
сильнее обычного принципа математической индукции
в силу следующей замечательной теоремы.
\begin{theorem}
Любое множество может быть вполне упорядочено.
\end{theorem}
Данное утверждение эквивалентно так называемой \odmkey{аксиоме
выбора}{Аксиома (axiom)!выбора (of choice)} в теории множеств и
само может быть принято за аксиому. 
%Мы опускаем доказательство
%эквивалентности этой теоремы аксиоме выбора.

\medskip
\begin{example}
Рассмотрим множество положительных рациональных чисел \\
$
\mathbb{Q}_+\Def\Set{\frac{m}{n}}{m,n\in\mathbb{N}},
$
где $m$ и $n$ взаимно просты.
Определим отношение $\prec$ на множестве
$\mathbb{Q}_+$
следующим образом:
$
\frac{m_1}{n_1}\prec\frac{m_2}{n_2}\Def m_1<m_2 \Or \left( m_1=m_2 \And n_1<n_2\right)\!,
$
где $<$~--- обычное отношение порядка на $\mathbb{N}$.
Нетрудно видеть, что множество $\mathbb{Q}_+$ с  отношением $\prec$
является вполне упорядоченным,
в то время как с обычным отношением порядка $<$
оно таковым не является.
\end{example}
\end{frame}

\begin{frame}
\frametitle{Онтология наивной теории множеств}
\begin{center}
\includegraphics[width=12cm]{ODM1_8_2.png}
\end{center}
\end{frame}


\subsection{1.9. Характеристические функции }
\label{1.9.}
\begin{frame}
\frametitle{1.9. Характеристические функции (1/2)}
%\vspace{-3pt}
В математике нередко бывает так,
что один и тот же объект описывается различными средствами, на разных языках.
Например, в планиметрии окружность определяется как геометрическое место точек,
равноудалённых от одной точки, 
называемой центром окружности, 
а в аналитической геометрии окружностью называется множество пар $(x, y)$,
удовлетворяющих уравнению $(x-a)^2+(y-b)^2 = r^2$.
Речь идёт об одном и том же объекте, 
обладающем присущими ему свойствами. 
Но в одном случае используется язык геометрических образов, 
а в другом~--- алгебраических формул.
При решении геометрических задач на построение 
с помощью циркуля и линейки удобнее использовать язык планиметрии, 
а при решении задач на вычисление координат точек пересечения окружностей 
удобнее использовать язык аналитической геометрии. 
Подобным образом содержание наивной теории множеств 
может быть передано разными средствами: 
через рассмотренное понятие совокупности элементов 
или через понятие \odmkey{характеристической функции}{Функция (function)!характеристическая (characteristic)}, 
рассматриваемое в этом разделе. 
Понятие характеристической функции 
наталкивает на различные обобщения,
которые оказываются полезными во многих приложениях. 

\end{frame}

\begin{frame}
\frametitle{1.9. Характеристические функции (2/2)}

\begin{tabular}{p{11mm}|p{45mm}|p{88mm}}
  № & Название параграфа
  & Ключевые термины и обозначения \\
  \hline
\DMref{1.9.1.}{1.9.1.} & Классификация характеристических функций &
{Множество характеристик; бинарная характеристическая функция, классификатор, неотрицательно целочисленная характеристическая функция, нечёткая характеристическая функция }
\\  
\DMref{1.9.2.}{1.9.2.} \RS & Характеристические функции множеств &
{ $1_A$~--- характеристическая функция множества $A$}\\  
\DMref{1.9.3.}{1.9.3.} & Характеристические функции отношений &
{ Предикат}\\
\DMref{1.9.4.}{1.9.4.} & Характеристические функции мультимножеств &
{ Операции над мультимножествами}\\  
\DMref{1.9.5.}{1.9.5.} & Классификаторы &
{Фасетный классификатор, иерархический классификатор}\\  
\DMref{1.9.6.}{1.9.6.} & Нечёткие множества &
{ Нечёткое дополнение, пересечение и объединение}\\  
\end{tabular}  

\end{frame}

\begin{frame}
\label{1.9.1.}\subsubsection{1.9.1. Классификация характеристических функций }\frametitle{1.9.1. Классификация характеристических функций (1/8)}

Пусть задано непустое множество $U$, которое называется множеством элементов,
и непустое частично-упорядоченное множество $\mathrm{X}$, 
которое называется множеством \odmkey{характеристик}{Множество (set)!характеристик (of characteristics)}.
Тогда всякая функция 
$\Map{\mu}{U}{\mathrm{X}}$ называется 
\odmkey{характеристической функцией}{Характеристическая функция (characteristic function)}, поскольку она однозначно сопоставляет каждому 
элементу $ u\in U$ характеристику $\chi \in \mathrm{X}$.

Характеристические функции используются в различных областях математики,
и при этом применяется различная терминология, в зависимости 
от природы множества характеристик.   


Если $|\mathrm{X}|=1$,
то все элементы имеют одну характеристику,
и характеристические функции неинформативны,
а потому не используются.

\smallskip
Наиболее употребительными являются следующие четыре случая:

\begin{enumerate}
\item[1)] $\mathrm{X} = \{0,1\}$~---
\odmkey{бинарные характеристические функции}{Характеристическая функция (characteristic function)!бинарная (binary)};
\item[2)] $\mathrm{X} =\{\Tuple{\chi}{1}{n}\}$~---
\odmkey{классификаторы}{Классификатор (classifier)};
\item[3)] $\mathrm{X} = \Bbb{N}_0$~--- 
\odmkey{неотрицательно целочисленные функции}{Характеристическая функция (characteristic function)!неотрицательно целочисленная \\(integer-valued non-negative)};
\item[4)] $\mathrm{X} = [0;1]$~---
\odmkey{нечёткие характеристические функции}{Характеристическая функция (characteristic function)!нечёткая (fuzzy)}.
\end{enumerate}

\end{frame}

\begin{frame}
\frametitle{1.9.1. Классификация характеристических функций (2/8)}

Наиболее употребительны бинарные характеристические функции.

\begin{example}
Всякому подмножеству $A$ множества $U$ можно взаимно-однозначно сопоставить
тотальную функцию $\Map{1_A}{U}{0..1}$, где 
$1_A(a)\Def \textbf{ if } 
a\in A \textbf{ then } 1 \textbf{ else } 0 
\textbf{ end~if}$.
Функция $1_A$ называется \odmkey{характеристической функцией множества}{Характеристическая функция (characteristic function)!множества (of set)} $A$.
\end{example}

\medskip
Элементы множества $U$ могут иметь сложную структуру.
В частности, элементами множества $U$ могут  быть кортежи, то есть $U=U_1\times\ldots\times U_n$. 
Это означает, что       
характеристические функции (как бинарные, так и иные) не 
обязаны иметь только один аргумент, а могут иметь несколько аргументов.  


\begin{example}
Всякому бинарному отношению $R$ между множествами $A$ и $B$ ($R\subset A\times B$)
можно взаимно-однозначно сопоставить тотальную функцию
$\Map{1_R}{A\times B}{0..1}$, которая определяется так:
$1_R(a,b)\Def$ \textbf{if $aRb$ then $1$ else $0$ end~if}
и называется \odmkey{характеристической функцией отношения}{Характеристическая функция (characteristic function)!отношения (of relation)} $R$.
\end{example}

\begin{remark}
Бинарные характеристические функции активно применяют
в математической логике (глава 4).   
В этом случае элементы множества $\mathrm{X}$ 
называют \odmkey{истинностными значениями}{Истинностное значение (truth value)}, 
а характеристические  функции называют \odmkey{предикатами}{Предикат (predicate)}.
\end{remark}

\end{frame}

\begin{frame}
\frametitle{ 1.9.1. Классификация характеристических функций (3/8)}

В приложениях очень часто используются характеристические функции
с конечным набором характеристик.
Эти функции называют \odmkey{классификаторами}{Классификатор (classifier)}, 
поскольку они приписывают каждому элементу определённую характеристику, то есть относят элемент к определённому классу.

\begin{columns}[t]
\begin{column}[T]{8cm}
\begin{example}
Оси координат разбивают плоскость на четыре \odmkey{квадранта}{Квадрант (quadrant)}. Квадранты 
по традиции нумеруются римскими цифрами. Для каждой точки плоскости, за исключением точек на осях,  можно определить, к какому квадранту она относится, 
т.~е. задать классификатор на $\Bbb{R}^2$. 
\end{example}
\end{column}
\begin{column}[T]{3cm} 
\begin{center}
\vspace{-1.5\baselineskip}
\begin{figure}[t]
\includegraphics[height=20mm]{f1_9_1_3_8.jpg}
\end{figure}\end{center}
\end{column}
\begin{column}[T]{1cm}
$\mathrm{I}:$ \\
$\mathrm{II}:$ \\
$\mathrm{III}:$ \\
$\mathrm{IV}:$
\end{column}
\begin{column}[T]{3cm}
$x > 0, y > 0$ \\
$x < 0, y > 0$ \\
$x < 0, y < 0$ \\
$x > 0, y < 0$ \\
\end{column}
\end{columns}

\medskip
\begin{remark} В предыдущем примере
точки на осях не отнесены ни к одному из квадрантов,
то есть функция $\mu$ не тотальная. Такие классификаторы называются \odmkey{неполными}{Классификатор (classifier)!неполный (incomplete)} и требуют оговорок или доопределения при использовании.  
\end{remark}

\end{frame}

\begin{frame}
\frametitle{1.9.1. Классификация характеристических функций (4/8)}

Пусть теперь областью значений характеристических функций являются 
натуральные числа и $0$. 
Тогда имеется 
взаимно однозначное соответствие 
между множеством всех мультимножеств
над множеством $U$ 
и множеством всех характеристических функций 
из $U$ в $\Bbb{N}_0$. Действительно, 
пусть $\widehat{X} = \left[ u_1^{a_1}, \dots, u_n^{a_n} \right]$~--- 
мультимножество над множеством $U$. 
Тогда мультимножеству $\widehat{X}$
взаимно однозначно соответствует 
характеристическая функция
$\Map{\mu_{\widehat{X}}}{U}{\Bbb{N}_0}$,
определяемая следующим образом:
$\mu_{\widehat{X}}(u_i)\Assign a_i$.

\medskip
Взаимно однозначное соответствие 
между множеством мультимножеств над $U$ 
и множеством характеристических функций
из $U$ в $\Bbb{N}_0$  обозначим$ f_1$. 
Сужение (\DMref{п.~1.6.1}{1.6.1.})
соответствия $f_1$
на множество индикаторов устанавливает 
взаимно однозначное соответствие $f_2$ между множеством индикаторов и
множеством бинарных характеристических функций.
Взаимно однозначное соответствие
между булеаном над множеством $U$
и множеством индикаторов над множеством $U$
(\DMref{п.~1.2.1}{1.2.1.}) обозначим $g_2$,
а взаимно однозначное соответствие
между булеаном над множеством $U$
и множеством характеристических функций
подмножеств множества $U$
обозначим $h_2$.
Таким образом, имеют место взаимно однозначные
соответствия и включения, представленные на следующем рисунке.  
\end{frame}

\begin{frame}
\frametitle{ 1.9.1. Классификация характеристических функций (5/8)}


\begin{center}
\includegraphics[width=110mm]{f1_9_1_5_8.jpg}
\end{center}
\end{frame}

\begin{frame}
\frametitle{ 1.9.1. Классификация характеристических функций (6/8)}

\begin{remark}
Множество индикаторов является
подмножеством множества мультимно\-жеств,
а множество бинарных характеристических функций
является подмножеством множества неотрицательно целочисленных характеристических функций. 
\end{remark}

\medskip
\begin{digress}
Объекты, обозначенные на предыдущем рисунке пунктирными линиями и вопросительными знаками, не определяются и не используются в наивной теории множеств. 
Это не означает, что в этих объектах
зашифрована какая-то загадка. 
Это означает, что язык наивной теории множеств,
как и всякий формальный язык, имеет 
ограниченную сферу применения. 
В подавляющем большинстве случаев
языка наивной теории множеств оказывается
достаточно для построения
практически удобных математических моделей,
и в этих случаях всё равно, на каком языке 
формулировать модели: на языке множеств, на языке индикаторов или
на языке характеристических функций~---
поскольку эти языки эквивалентны.
Если же языка теории множеств оказывается недостаточно, то не следует его коверкать,
изобретая доморощенные <<обобщения>>. 
Лучше выбрать другой язык с более широкой сферой 
применения, подходящий к данной задаче. 
\end{digress}


\end{frame}

\begin{frame}
\frametitle{ 1.9.1. Классификация характеристических функций (7/8)}

Отношение принадлежности элемента множеству 
в наивной теории множеств не имеет 
никакой неопределённости: либо элемент принадлежит множеству,
либо не принадлежит, и третьего не дано.
Но в жизни такой определённости может не быть.   

\begin{example}
Мы пользуемся понятием <<высокий человек>>
и вправе рассматривать множество высоких
людей, а также его дополнение~--- множество низких людей.
Однако если посмотреть на картинку, становится понятно, что такой классификации  недостаточно, бинарная характеристическая функция  
оказывается слишком грубым классификатором.


\begin{center}
\includegraphics[height=4cm]{1_9_1_(6).jpg}
\end{center}

\end{example}

\end{frame}

\begin{frame}
\frametitle{ 1.9.1. Классификация характеристических функций (8/8)}

Можно ввести трёхзначный классификатор  и 
множество людей среднего роста,
однако некоторая неудовлетворённость при этом остаётся,
и появляются выражения <<несколько ниже среднего>>, <<чуть выше среднего>>.

В таких случаях используют  
\odmkey{нечёткую характеристическую
функцию}{Функция (function)!характеристическая (characteristic)!нечёткая (fuzzy)} с множеством значений из отрезка $[0;1]$.
Значение этой функции трактуется как степень
принадлежности элемента множеству.

Множество, определяемое характеристической функцией 
с множеством значений $[0; 1]$, называется \odmkey{нечётким}{Нечёткое множество (fuzzy set)}.
\begin{digress}
Впервые термин \odmkey{нечёткая логика}{Логика (logic)!нечёткая (fuzzy)}
 (fuzzy logic) был введён Заде 
%\footnote{Лотфи Заде (род. 1921)} 
%Случилось 
%это 
в 1965 году. Бытует мнение, что толчком для создания новой теории послужил спор Заде со своим другом о том, чья жена привлекательнее. К единому мнению спорящие, естественно, так и не пришли. Но это спровоцировало Заде сформулировать концепцию выражения нечётких понятий 
(<<привлекательность>>, <<молодость>> и пр.) в числовой форме.
\end{digress}


\end{frame}

\begin{frame}
\label{1.9.2.}\subsubsection{1.9.2. Характеристические функции множеств }\frametitle{1.9.2. Характеристические функции множеств (1/5)}
В этом параграфе рассматриваются обычные множества и соответствующие им бинарные характеристические функции.
Из определения 
$1_A(a)\Def$ \textbf{if $a\in A$ then $1$ else $0$ end~if}
характеристической функции $\Map{1_A}{U}{0..1}$ подмножества $A\subset U$ 
легко вывести, что $1_\emptyset =0$ и $1_U =1$.   


Характеристические функции результатов операций над множествами выражаются 
через характеристические функции операндов с помощью арифметических операций.

\begin{theorem}
Пусть $\Map{1_A}{U}{0..1}$~--- характеристическая
функция множества $A\subset U$ и
$\Map{1_B}{U}{0..1}$~--- характеристическая
функция множества $B\subset U$. Тогда:
\begin{enumerate}
\item[1)] $1_{\overline{A}}=1-1_A$;
\item[2)] $1_{A\cap B}=\min(1_A,1_B)=1_A 1_B$;
\item[3)] $1_{A\cup B}=\max(1_A,1_B)=1_A+1_B - 1_A 1_B$;
\item[4)] $1_{A\setminus B}=1_A(1-1_B)=1_A- 1_A 1_B$;
\item[5)] $1_{A\triangle B}=1_A+1_B - 2\cdot 1_A 1_B$.
\end{enumerate} 
\end{theorem}
\end{frame}

\begin{frame}
\frametitle{ 1.9.2. Характеристические функции множеств (2/5)}

\begin{proof}
Возьмём произвольный элемент $x\in U$ и вычислим значения проверяемых выражений,  рассматривая все возможные случаи принадлежности элемента $x$. 

\proofsection{1)} Возможны два случая, разобранные в таблице.

\begin{tabular}{|ccc|}
\hline
$1_A(x)$ & $1_{\overline{A}}(x)$  & $1 -1_A(x)$\\
\hline
$0$ & $1$ & $1$ \\ 
$1$ & $0$ & $0$ \\ 
\hline         
\end{tabular}

\medskip
\proofsection{2)} Возможны четыре случая, разобранные в таблице.

\begin{tabular}{|ccccc|}
\hline
$1_A(x)$ & $1_B(x)$  & $1_{A\cap B}(x)$ & $\min(1_A(x),1_B(x))$ & $1_A(x) 1_B(x)$ \\
\hline
$0$ & $0$ & $0$ & $0$ & $0$ \\ 
$0$ & $1$ & $0$ & $0$ & $0$ \\ 
$1$ & $0$ & $0$ & $0$ & $0$ \\ 
$1$ & $1$ & $1$ & $1$ & $1$ \\ 
\hline         
\end{tabular}

\medskip
Остальные утверждения доказываются аналогично.
\end{proof}
\end{frame}

\begin{frame}
\frametitle{ 1.9.2. Характеристические функции множеств (3/5)}

\begin{theorem}
\begin{enumerate}
\item
$A \subset B \Equiv  \A{x\in U}{1_A(x) \leq 1_B(x)}$.
\item
$A = B  \Equiv \A{x\in U}{1_A(x) = 1_B(x)}$.
\end{enumerate}
\end{theorem}
\begin{proof}
\proofsection{1} Возможны четыре случая, разобранные в таблице.

\begin{tabular}{|cccc|}
\hline
$1_A(x)$ & $1_B(x)$  & $x\in A \Impl x\in B$ & $1_A(x) \leq 1_B(x)$ \\
\hline
$0$ & $0$ & $1$ & $1$ \\ 
$0$ & $1$ & $1$ & $1$ \\ 
$1$ & $0$ & $0$ & $0$ \\ 
$1$ & $1$ & $1$ & $1$ \\ 
\hline         
\end{tabular}

\medskip
\proofsection{2} По определению, $A = B \Equiv A \subset B \And B \subset A$. 
Имеем: $A = B \Equiv$\\ $\Equiv\A{x\in U}{1_A(x) \leq 1_B(x) \And 1_B(x) \leq 1_A(x)}$
$\Equiv \A{x\in U}{1_A(x) = 1_B(x)}$
\end{proof}
\end{frame}

\begin{frame}
\frametitle{1.9.2. Характеристические функции множеств (4/5)}
Характеристические функции часто оказываются 
полезны при решении задач в алгебре множеств.

\smallskip
\begin{example}
Требуется доказать равенство $(A \triangle B) \triangle C = A \triangle (B \triangle C)$.
\smallskip
Выразим левую и правую часть равенства в терминах характеристических функций.

\smallskip
$(A \triangle B) \triangle C \mapsto 1_A + 1_B - 2\cdot 1_A \cdot 1_B + 1_C - 2\cdot (1_A + 1_B - 2 \cdot 1_A \cdot 1_B)\cdot 1_C  =$ \\
$= 1_A + 1_B -2\cdot 1_A\cdot 1_B 
+ 1_C - 2\cdot 1_A \cdot 1_C - 2\cdot 1_B\cdot  1_C + 4\cdot 1_A\cdot 1_B\cdot 1_C $.

\smallskip
$A \triangle (B \triangle C) \mapsto 1_A + (1_B + 1_C - 2\cdot 1_B\cdot 1_C) - 2\cdot 1_A \cdot(1_B + 1_C - 2\cdot 1_B\cdot 1_C) =$\\
$= 1_A + 1_B + 1_C - 2\cdot 1_B\cdot 1_C - 
2\cdot 1_A\cdot 1_B -2\cdot 1_A\cdot 1_C 
+ 4\cdot 1_A\cdot 1_B \cdot 1_C$.

\smallskip
Результаты совпали, равенство доказано.
Таким образом, доказательство получилось в один шаг, тогда как при решении задачи эквивалентными преобразованиями по равенствам 
\DMref{п.~1.2.9}{1.2.9.} требуется больше десяти шагов.
\end{example}



\end{frame}
\begin{frame}
\frametitle{1.9.2. Характеристические функции множеств (5/5)}

\begin{remark}
Нетрудно показать, что все равенства, перечисленные в 
\DMref{п.~1.2.9}{1.2.9.}, выполняются для характеристических функций. Например, свойство идемпотентности: $A\cup A = A$, так как $\max(1_A, 1_A) = 1_A$.
Таким образом, вычисления в терминах характеристических функций множеств сводимы к операциям наивной теории множеств и обратно.
\end{remark}

\smallskip
\begin{digress}
Арифметические вычисления обычному человеку выполнять проще, легче и привычнее всего, поэтому решение
в терминах характеристических функций оказывается наглядным и коротким. На самом деле оно эквивалентно
решению в терминах операций наивной теории множеств,
просто многие арифметические операции выполняются
<<в уме>> и не выписываются в выкладках явно. 
\end{digress}

\end{frame}

\begin{frame}
\label{1.9.3.}\subsubsection{1.9.3. Характеристические функции отношений }\frametitle{1.9.3. Характеристические функции отношений (1/3)}
В этом параграфе рассматриваются 
бинарные отношения на множестве и соответствующие им бинарные характеристические функции двух аргументов, используя определение 
$1_R(a,b)\Def$ \textbf{if $aRb$ then $1$ else $0$ end~if}
характеристической функции $\Map{1_R}{A\times A}{0..1}$ отношения $R\subset A^2$.   

\smallskip
Характеристические функции результатов некоторых операций над отношениями выражаются 
через характеристические функции операндов.

\smallskip
\begin{theorem}
Пусть $\Map{1_R}{A^2}{0..1}$~--- характеристическая
функция бинарного отношения $R\subset A^2$,
$A=\{\Tuple{a}{1}{n}\}$, $x, y\in A$. Тогда:
\begin{enumerate}
\item[1)] $1_{\overline{R}}(x,y) = 1-1_R(x,y)$;
\item[2)]$1_{R\circ R} (x, y) = \BigOr\limits_{i=1}^n
1_{R}(x, a_i)1_{R}(a_i,y)$;
\item[3)]$1_{\ker R} (x, y) = \BigOr\limits_{i=1}^n
1_{R}(x, a_i)1_{R}(y,a_i)$;
\item[4)] $1_{R^{-1}}(x,y) = 1_R(y,x)$.
\end{enumerate} 
\end{theorem}
\end{frame}

\begin{frame}
\frametitle{1.9.3. Характеристические функции отношений (2/3)}

\begin{proof}
Рассмотрим произвольную пару 
$(x,y)\in A^2$ и вычислим значения проверяемых
выражений в различных случаях.

\proofsection{1)}   Возможны два случая, разобранные в таблице. 

\medskip
\begin{tabular}{|ccc|}
\hline
$1_R(x,y)$ & $1_{\overline{R}}(x,y)$  & $1 -1_R(x,y)$\\
\hline
$0$ & $1$ & $1$ \\ 
$1$ & $0$ & $0$ \\ 
\hline         
\end{tabular}

\medskip
\proofsection{2)}   Возможны два случая, разобранные в таблице. 

\medskip
\begin{tabular}{|ccc|}
\hline
$\E{i\in 1..n}{xRa_i\And a_iRy}$ & 
$1_{R\circ R}(x,y)$  & $\BigOr\limits_{i=1}^n
1_{R}(x, a_i)1_{R}(a_i,y)$\\
\hline
$0$ & $0$ & $0$ \\ 
$1$ & $1$ & $1$ \\ 
\hline         
\end{tabular}

\end{proof}

\end{frame}

\begin{frame}
\frametitle{\large 1.9.3. Характеристические функции отношений (3/3)}
\begin{proof} (Продолжение)

\proofsection{3)}   Возможны два случая, разобранные в таблице. 

\medskip
\begin{tabular}{|ccc|}
\hline
$\E{i\in 1..n}{x R a_i\And y R a_i}$ & 
$1_{\ker R}(x,y)$  & $\BigOr\limits_{i=1}^n
1_{R}(x, a_i)1_{R}(y,a_i)$\\
\hline
$0$ & $0$ & $0$ \\ 
$1$ & $1$ & $1$ \\ 
\hline         
\end{tabular}

\medskip
\proofsection{4)}   В любом случае
$(x,y)\in R \Equiv (y,x)\in R^{-1}$. 
\end{proof}

\smallskip
\begin{remark}
Полезно сравнить эту теорему с теоремами
\DMref{п.~1.4.9}{1.4.9.}.
\end{remark}

\smallskip
\begin{digress}
Операции с характеристическими функциями 
отношений не удаётся свести к простым арифметическим
манипуляциям так, как это сделано для характеристических функций множеств.
Поэтому на практике характеристические функции отношений в том виде, как они представлены в этом параграфе, используются редко.
Чаще характеристические функции отношений 
(под названием 
предикатов) трактуются средствами математической логики (см. главу 4).
\end{digress}

\end{frame}

\begin{frame}
\label{1.9.4.}\subsubsection{1.9.4. Характеристические функции мультимножеств }\frametitle{1.9.4. Характеристические функции мультимножеств (1/2)}
Пусть $\widehat{X} = \left[ u_1^{a_1}, \dots, u_n^{a_n} \right]$~--- мультимножество над множеством
$U = \{\Tuple{u}{1}{n}\}$, где $a_i$~--- показатель элемента  $u_i$.
\odmkey{Характеристической функцией мультимножества}{Характеристическая функция (characteristic function)!мультимножества (of multiset)} 
$\widehat{X}$ называется тотальная функция  
$\Map{\mu_{\widehat{X}}}{U}{\Bbb{N}_0}$, такая, что
$\mu_{\widehat{X}}(u_i)\Def a_i$. 

\medskip
Пользуясь характеристическими функциями, для 
мультимножеств легко ввести те же операции,
что и для множеств. Пусть $A, B$~--- 
мультимножества над множеством $U$, 
а $\mu_A, \mu_B$~--- 
их характеристические функции, соответственно. Тогда: 

\odmkey{мощность мультимножества}{Мощность (cardinality)!мультимножества (of multiset)}:
$|A| \Def \sum\limits_{i=1}^n \mu_A (u_i)$;

\odmkey{включение мультимножеств}{Включение (inclusion)!мультимножеств (of multisets)}:
$A \subset B \Def  \A{u\in U}{\mu_A(u) \leq \mu_B(u)}$;

\odmkey{равенство мультимножеств}{Равенство (equality)!мультимножеств (of multisets)}:
$A = B  \Def \A{u\in U}{\mu_A(u) = \mu_B(u)}$;

\odmkey{пересечение мультимножеств}{Пересечение (intersection)!мультимножеств (of multisets)}:
$  \mu_{A\cap B}(u_i) \Def \min(\mu_A(u_i), \mu_B(u_i))$;

\odmkey{объединение мультимножеств}{Объединение (union)!мультимножеств (of multisets)}:
$  \mu_{A\cup B}(u_i) \Def \max(\mu_A(u_i), \mu_B(u_i))$.
\end{frame}

\begin{frame}
\frametitle{1.9.4. Характеристические функции мультимножеств (2/2)}

\begin{remark}
Учитывая теоремы \DMref{п.~1.9.2}{1.9.2.} легко показать,
что введённые таким образом операции 
над мультимножествами согласованы 
с одноимёнными операциями над множествами,
в частности, совпадают при применении к индикаторам.
Однако в общем случае 
операции над мультимножествами
не обязаны обладать такими же свойстами,
как операции над множествами.
Например, равенство $\mu_{A\cap B} = 
\mu_A \cdot \mu_B$ всегда выполняется
для индикаторов, 
но не всегда выполняется для  
мультимножеств в общем случае.
\end{remark}

\medskip
Для мультимножеств легко определить и другие полезные операции,
например, \odmkey{арифметическую сумму мультимножеств}{Арифметическая сумма мультимножеств (arithmetic sum of multisets)} 
$$\mu_{A + B} (u_i) \Def \mu_{A} (u_i) + \mu_{B} (u_i) ,$$
получая равенства, для которых нет аналогов в наивной теории множеств.
Например:
$$A+B=(A\cup B)+(A\cap B).$$


%%%%%%%%%%%%%%%%%%%%%%%%%%%%%%%%%%%%%%%%
% дальнейшее в этом параграфе считаю слишком сырым 
% ФАН 25.06.2015
%
%Объединение и арифметическая сумма определяются аналогично:
%$$C = A\cup B \Def \Set{x_i^{c_i}}{c_i = \max(a_i, b_i) },$$
%$$C = A + B \Def \Set{x_i^{c_i}}{c_i = a_i + b_i }.$$
%
%Очевидным образом характеристические функции результатов операций над мультимножествами выражаются 
%через характеристические функции операндов:
%\begin{enumerate}
%\item[1)] $f_{A\cap B} = \min(f_A, f_B)$;
%\item[2)] $f_{A\cup B} = \max(f_A, f_B)$;
%\item[3)] $f_{A+B} = f_A + f_B$.
%\end{enumerate} 

\end{frame}

%\begin{frame}
%\frametitle{\large 1.9.4. Характеристические функции мультимножеств (3/3)}
%\begin{digress} 
%Введённые таким образом операции лежат в основе тропической математики. Пусть есть множество $\mathbb{R}_{max} = \mathbb{R} \cup \{-\infty\}$ и две операции на нём:
%\begin{enumerate}
%\item[1)]$x \oplus y \Def max\{x, y\}$;
%\item[2)]$x \odot y \Def x+y$.
%\end{enumerate}
%Тогда, интерпретируя характеристическую функцию, как отображение $f: X \to \mathbb{R}_{max}$, получаем выражения для характеристической функции объединения и арифметической суммы мультимножеств через новые операции:
%\begin{enumerate}
%\item[1)]$f_{A\cup B} = f_A \oplus f_B$
%\item[2)]$f_{A+B} = f_A \odot f_B$.
%\end{enumerate}
%\end{digress}
%\end{frame}

\begin{frame}
\label{1.9.5.}\subsubsection{1.9.5. Классификаторы }\frametitle{1.9.5. Классификаторы (1/4)}

Характеристическая функция  $\Map{\mu}{U}{\mathrm{X}}$, 
где $\mathrm{X}=\{\Tuple{\chi}{1}{n}\}$~--- конечное множество,
называется \odmkey{классификатором}{Классификатор (classifier)}.

\medskip 
Классификатор задаёт семейство 
подмножеств $\cal E$ множества $U$ ($\cal E = {\{E_i\}}_{i=1}^n$), 
в котором все элементы одного подмножества объединены одной характеристикой:\\ 
$\A{x,y \in E_i}{\mu(x)=\mu(y)=\chi_i}$
(семейство множеств уровня, 
\DMref{п.~1.7.4}{1.7.4.}).
Поскольку классификатор~--- тотальная функция, это семейство является дизъюнктным покрытием, т.~е. разбиением (\DMref{п.~1.2.7}{1.2.7.}) и задаёт отношение
эквивалентности на множестве $U$ (\DMref{п.~1.7.1}{1.7.1.}). Таким образом, 
классификаторы~--- ещё один способ описания
часто используемой концепции разбиений, имеющий множество приложений.
Например,
\odmkey{дискриминантный анализ}{Анализ (analysis)!дискриминантный (discriminant)} в статистике,
\odmkey{распознавание образов}{Распознавание образов (pattern recognition)} в искусственном интеллекте, \odmkey{систематика}{Систематика (systematics)} живых организмов в биологии.  
Особый интерес представляет случай, когда множество $U$ очень велико
или даже бесконечно. При этом одним классификатором трудно обойтись:
либо классы оказываются слишком большими, либо их оказывается слишком много.
\end{frame}

\begin{frame}
\frametitle{1.9.5. Классификаторы (2/4)}


В этом случае на практике применяют один из двух основных методов классификации: фасетную классификацию или иерархическую классификацию.

\medskip
При \odmkey{фасетной классификации}{Классификатор (classifier)!фасетный (faceted)}
на множестве $U$ определяются независимо несколько классификаторов:
$\Tuple{\mu}{1}{k}$, каждый из которых задает своё разбиение,
$\Map{\mu_i}{U}{\mathrm{X_i}}$.
Эти классификаторы иногда называют признаками, или свойствами.
\odmkey{Фасетный классификатор}{Классификатор (classifier)!фасетный
(faceted)}~---
это общее измельчение всех разбиений, 
$\Map{\mu}{U}{\mathrm{X}_1\times\ldots\times\mathrm{X}_k}$.

\begin{example}
В системах просмотра фильмов через Интернет
каталог доступных фильмов обычно устроен как фасетный классификатор:
можно указать жанр фильма, страну производства фильма,
год выпуска и т.~д. Обычно для точного поиска оказывается 
достаточным иметь около десяти независимых признаков.   
\end{example}

При фасетной классификации класс каждого объекта определяется как неупорядоченное \em{множество}
характеристик $\{\Tuple{\chi}{1}{k}\}$, причём все объекты имеют одинаковое количество характеристик.

\begin{remark}
Термин <<фасетный>> происходит от заимствованного слова <<фасет>>, что означает скошенную грань. Поэтому фасетный
классификатор было бы уместно называть \em{многогранным}.
\end{remark} 

\end{frame}

\begin{frame}
\frametitle{1.9.5. Классификаторы (3/4)}

\odmkey{Иерархический классификатор}{Классификатор (classifier)!иерархический (hierarchical)}    
также задаётся несколькими классификаторами. 

\smallskip
Сначала определяется классификатор верхнего уровня.
Затем на всех блоках разбиения классификатора первого уровня определяются
свои классификаторы второго уровня, и так далее.
Уровней измельчения разбиений может быть сколько угодно.
Такие разбиения образуют иерархию, т.~е. неупорядоченное ориентированное дерево 
(\DMref{п.~9.2.1}{9.2.1.}). 

\smallskip
При измельчении блоков используемые классификаторы и глубина 
измельчения могут совпадать, а могут и не совпадать с другими блоками 
того же уровня. В первом случае дерево классификации имеет 
постоянную
\odmkey{ширину ветвления}{Ширина ветвления (branch-width)} и является \odmkey{выровненным}{Дерево (tree)!классификации (classification)!выровненное (aligned)},
а во втором случае дерево может иметь произвольную форму. 

\smallskip
При иерархической классификации класс каждого объекта определяется как \em{последовательность}
уточняющих характеристик $(\Tuple{\chi}{1}{k})$ (путь в дереве от корня),
причём разные объекты могут иметь различное количество характеристик.
%в котором все элементы одного подмножества объединены одной характеристикой: 
%$\A{x,y \in E_i}{\mu(x)=\mu(y)=\chi_i}$
%(семейство множеств уровня, \DMref{п.~1.7.4).}{1.7.4).}
%Поскольку классификатор~--- тотальная функция, это семейство является дизъюнктным покрытием, т.~е. разбиением (п. 1.2.7) и задаёт отношение
%эквивалентности на множестве $U$ (п.~1.7.1).
%\end{frame}
%
%\begin{frame}
%\frametitle{1.9.5. Классификаторы (2/3)}
%Пусть на множестве $U$ определены два классификатора $\mu$ и $\lambda$. 
%Они задают разбиения $\cal E={\{E_i\}}_{i=1}^n$ и $\cal B={\{B_i\}}_{i=1}^n$ соответственно,
%причём разбиение $\cal B$ является измельчением
%разбиения $\cal E$, т.~е. каждый блок $\cal E$ является объединением блоков $\cal B$  
%(п\DMref{п.~1.7.1,}{1.7.1,}5.4.1). 
%Уровней измельчения разбиений может быть сколько угодно.
%Такие разбиения образуют иерархию, т.~е. неупорядоченное ориентированное дерево (п. 9.2.1). 
%%%%%%%%%%%%%%%%%%%%%%%%%%%%%%%%%%%%%%%%%%%
%%% текст и картинка очень сырые
%проиллюстрированное на картинке.
%В корне этой иерархии находится множество $U$, а в каждом узле – подмножество своего родителя. 
%На каждой ступени располагаются элементы только одного разбиения.
%Причём чем ниже ступень иерархии, тем мельче разбиение (на $(i+1)$-й ступени находится измельчение $i$-й ступени).
%
%\begin{center}
%\includegraphics[height=2.5cm]{1_9_5_(2).jpg}
%\end{center}
\end{frame}

\begin{frame}
\frametitle{1.9.5. Классификаторы (4/4)}
\begin{example}
\odmkey{Универсальная десятичная классификация}{Классификация (classification)!универсальная десятичная (universal decimal)} (УДК)~--- международная система классификации информации,
используемая в России. 
Вся информация делится на 10 блоков, каждому из которых соответствует цифра от 0 до 9. 
Каждый блок в свою очередь делится ещё на 10 подблоков, каждому из которых ставится в соответствие цифра от 0 до 9, и~т.~д. 
Таким образом каждому блоку соответствует код, составленный из кода своего предка и своей цифры. 
Так, например, блоку «Религии Индийского субконтинента» соответствует код 23. 
\end{example}
\begin{center}
\includegraphics[height=4cm]{1_9_5_(3).jpg}
\end{center}
\end{frame}

\begin{frame}
\label{1.9.6.}\subsubsection{1.9.6. Нечёткие множества }\frametitle{1.9.6. Нечёткие множества (1/5)}

Нечёткое множество $A$ описывается нечёткой характеристической функцией вида \\ $\Map{\mu_A}{U}{[0; 1]}$.
Для нечётких множеств определены следующие операции:
\begin{enumerate}
\item[1)]
\odmkey{Нечёткое дополнение}{Дополнение (complement)!нечёткое (fuzzy)}
$\mu_{\overline{A}}\Def1-\mu_A$;
\item[2)] 
\odmkey{Нечёткое пересечение}{Пересечение (intersection)!нечёткое (fuzzy)}
$\mu_{A\cap B}\Def\min(\mu_A,\mu_B)$;
\item[3)] 
\odmkey{Нечёткое объединение}{Объединение (union)!нечёткое (fuzzy)}
$\mu_{A\cup B}\Def\max(\mu_A,\mu_B)$;
\end{enumerate} 
\begin{example}
На картинке проиллюстрировано пересечение двух нечётких множеств, а именно множеств старых и молодых людей.
\end{example}
\begin{center}
\includegraphics[height=3.5cm]{1_9_6_(1).jpg} 
\end{center}

\end{frame}

\begin{frame}
\frametitle{1.9.6. Нечёткие множества (2/5)}
Некоторые свойства операций над обычными множествами также сохраняются и для нечётких множеств.

\smallskip
\begin{theorem}
Для операций над нечёткими множествами выполняются законы де Моргана:
$\overline{A \cup B} = \overline{A} \cap \overline{B}$,
$\overline{A \cap B} = \overline{A} \cup \overline{B}$.
\end{theorem}

\begin{proof}
В терминах характеристических функций первое утверждение может быть записано в виде 
$1 - \max(\mu_A, \mu_B) = \min(1 - \mu_A, 1 - \mu_B)$ и легко проверено с учётом двух случаев: 
$f_A(x) > f_B(x)$ и $f_A(x) \le f_B(x)$.

Второе утверждение доказывается аналогично.
\end{proof}

\smallskip
\begin{theorem}
Нечёткие объединение и пересечение взаимно дистрибутивны:\\
$A \cup (B \cap C) = (A \cup B) \cap (A \cup C)$, 
$A \cap (B \cup C) = (A \cap B) \cup (A \cap C)$.
\end{theorem}
\begin{proof}
В терминах характеристических функций первое утверждение может быть записано в виде
$\max(\mu_A, \min(\mu_B, \mu_C)) = \min(\max(\mu_A, \mu_B), \max(\mu_A, \mu_C))$ и проверено с учётом шести следующих случаев.
\end{proof}


\end{frame}

\begin{frame}
\frametitle{1.9.6. Нечёткие множества (3/5)}

\begin{proof} (Продолжение)

\begin{tabular}{|ccc|}
\hline
случай & $\max(\mu_A, \min(\mu_B, \mu_C))$ 
& $\min(\max(\mu_A, \mu_B), \max(\mu_A, \mu_C))$ \\
$\mu_A \ge \mu_B \ge \mu_C$ & $\mu_A$ & $\mu_A$ \\ 
$\mu_A \ge \mu_C \ge \mu_B$ & $\mu_A$ & $\mu_A$ \\ 
$\mu_B \ge \mu_A \ge \mu_C$ & $\mu_A$ & $\mu_A$ \\ 
$\mu_B \ge \mu_C \ge \mu_A$ & $\mu_C$ & $\mu_C$ \\  
$\mu_C \ge \mu_A \ge \mu_B$ & $\mu_A$ & $\mu_A$ \\  
$\mu_C \ge \mu_B \ge \mu_A$ & $\mu_B$ & $\mu_B$ \\ 
\hline
\end{tabular}

Второе утверждение доказывается аналогично.
\end{proof}

\medskip
\begin{remark}
Не все свойства операций над обычными множествами сохраняются и для нечётких множеств.
В частности,
идемпотентность,
коммутативность,
ассоциативность,
дистрибутивность,
поглощение,
инволютивность и
законы де Моргана выполняются,
но, например, свойство дополнения не выполняется, 
так как $A \cup \overline{A} = U$ в терминах характеристических функций записывается в виде
$\max(\mu_A, 1 - \mu_A) = 1$, что, очевидно, неверно для всех 
$\Map{\mu_A}{U}{[0; 1]}$, отличных от констант 0 и 1.
\end{remark}


\end{frame}


\begin{frame}
\frametitle{1.9.6. Нечёткие множества (4/5)}
Для нечётких множеств определяются следующие алгебраические операции, которых нет в наивной теории множеств:
\begin{tabbing}
\hspace{0mm}\=\odmkey{Алгебраическое умножение}{Нечёткое множество (fuzzy set)!алгебраическое умножение (algebraic multiplication)},\quad \=\kill
\>\odmkey{Алгебраическое умножение}{Нечёткое множество (fuzzy set)!алгебраическое умножение (algebraic multiplication)}:
\>$\mu_{AB} = \mu_A\mu_B$;\\
\>\odmkey{Алгебраическая сумма}{Нечёткое множество (fuzzy set)!алгебраическая сумма (algebraic sum)}:
\>$\mu_{A+B} = \mu_A + \mu_B$;\\
\>\odmkey{Абсолютная разность}{Нечёткое множество (fuzzy set)!абсолютная разность (absolute difference)}:
\> $\mu_{|A-B|} = |\mu_A - \mu_B|$.
%;\\
%\>\odmkey{Выпуклая комбинация}{Выпуклая комбинация нечётких множеств}:
%\> $\mu_G(x) = (\mu_C(x) - \mu_B(x))/(\mu_A(x)-\mu_B(x))$.
\end{tabbing}

\begin{remark}
Можно заметить, что для обычных множеств пересечение и алгебраическое умножение, а также объединение и алгебраическая сумма являются эквивалентными понятиями.
\end{remark}
\end{frame}

%
%\begin{frame}
%\label{1.9.3.}\subsubsection{1.9.3. Свойства нечётких множеств }\frametitle{1.9.3. Свойства нечётких множеств (1/?)}
%{\color{blue} TO DO. Чрезвычайно важное указание.
%Я не знаю, как поступить с последующим текстом.
%Возможно он полезен и даже необходим.
%Тогда его придётся выправить.
%Но сначала требуется получить ответ на вопрос: зачем?
%Более конкретно:
%
%1) зачем мы рассматриваем множества вещественных чисел
%и гладкие характеристические функции? Понятно, что в 
%математическом анализе это необходимо; возможно, это полезно
%при переходе к вероятности; но в дискретной математике это не нужно;
%
%2) понятия выпуклости, ограниченности, евклидова пространства и 
%метрической окрестности нельзя считать известными, их придётся определить;
%
%3) что это даёт в дискретных моделях?    
%
%\vspace{\baselineskip}
%Возможно, я чего-то не вижу. Покажите! 
%    
%ФАН}
%
%\end{frame}
%
%
%\begin{frame}
%\label{1.9.3.}\subsubsection{1.9.3. Свойства нечётких множеств }\frametitle{1.9.3. Свойства нечётких множеств (1/7)}
%Будем предполагать, что $X$ - вещественное евклидово пространство. 
%Нечёткое множество $A$ \odmkey{выпукло}{выпуклое нечёткое множество} тогда и только тогда, когда множества 
%$G_a = \Set{x}{f_A(x) \geq a}$ выпуклы для всех $a \in (0; 1]$. Это утверждение можно переписать в более удобной для использования форме: нечёткое множество $A$ выпукло тогда и только тогда, когда  $\A {x_1, x_2 \in X, a\in [0; 1]}{f_A[ax_1 + (1 - a)x_2] \geq min[f_A(x_1), f_A(x_2)]}$.
%
%\begin{remark}
%Важно заметить, что это определение не требует, чтобы $f_A$ была выпуклой функцией. Это проиллюстрировано на картинке ниже:
%\end{remark}
%\begin{center}
%\includegraphics[scale=0.4]{1_9_3_(1).jpg} 
%\end{center}
%
%\end{frame}
%
%\begin{frame}
%\frametitle{1.9.3. Свойства нечётких множеств (2/7)}
%\begin{theorem}
% Если A и B выпуклы, то $A\cap B$ выпукло.
%\end{theorem}
%\begin{proof}
%$C = A\cap B$. Тогда $f_C[ax_1+(1-a)x_2] = min[f_A[ax_1+(1-a)x_2], f_B[ax_1+(1-a)x_2]]$. 
%
%Так как $A$ и $B$ выпуклы, то:
%
%$f_A[ax_1+(1-a)x_2] \geq min[f_A(x_1), f_A(x_2)]$, $f_B[ax_1+(1-a)x_2] \geq min[f_B(x_1), f_B(x_2)]$.
%
%Следовательно, $f_C[ax_1+(1-a)x_2] \geq min[min[f_A(x_1), f_A(x_2)], min[f_B(x_1), f_B(x_2)]] \Equiv f_C[ax_1+(1-a)x_2] \geq min[min[f_A(x_1), f_B(x_1)], min[f_A(x_2), f_B(x_2)]] \Equiv f_C[ax_1+(1-a)x_2] \geq min[f_C(x_1), f_C(x_2)]$.
%\end{proof}
%\end{frame}
%
%\begin{frame}
%\frametitle{1.9.3. Свойства нечётких множеств (3/7)}
%Нечёткое множество $A$ называется \odmkey{ограниченным}{органиченное нечёткое множество} тогда и только тогда, когда все множества $G_a = \Set{x}{f_A(x) \geq a}$ ограничены для любых $a > 0$.
%
%\begin{lemma}
%Пусть $A$~--- нечёткое ограниченное множество, и пусть $M = sup_xf_A(x)$.
%Тогда существует по крайней мере одна точка $x_0$, на которой $M$ существенно достигается, то есть 
%для каждого $e > 0$ точки каждой сферической окрестности $x_0$ содержатся в множестве  $Q(e) = \Set{x}{f_A(x) \geq M - e}$.
%\end{lemma}
%
%\begin{proof}
%Рассмотрим последовательность вложенных ограниченных множеств $G_1$, $G_2$, ..., где $G_n = \Set{x}{f_A(x) \geq M - M/(n+1)}$. Заметим, что $G_n$ непусто для всех конечных $n$ вследствие определения $M$.
%Возьмём за $x_n$ произвольную точку из $G_n$. Тогда $x_1$, $x_2$, ... - последовательность точек в замкнутом ограниченном множесте $G_1$. По теореме Больцана-Вейерштрасса эта последовательность должна иметь хотя бы одну предельную точку $x_0$ в $G_1$. 
%\end{proof}
%\end{frame}
%
%\begin{frame}
%\frametitle{1.9.3. Свойства нечётких множеств (4/7)}
%\begin{proof}
%(Продолжение) Следовательно, каждая сферическая окрестность $x_0$ будет содержать бесконечно много точек из последовательности $x_1$, $x_2$, ..., и более конкретно, из последовательности $x_{N+1}$, $x_{N+2}$, ..., где $N \geq M/e$. 
%Так как точки этой подпоследовательности попадают в $Q(e) = \Set{x}{f_A(x) \geq M - e}$, лемма доказана.
%\end{proof}
%\begin{remark}
%Здесь и далее используется понятие метрики, такое как сферическая (метрическая) окрестность. В пункте 6.6.3 это понятие будет рассмотрено подробно, однако для краткости изложения будем использовать его в этом разделе, так как его значение в данном контексте интуитивно понятно.
%\end{remark}
%\end{frame}
%
%\begin{frame}
%\frametitle{1.9.3. Свойства нечётких множеств (5/7)}
%Нечёткое множество $A$ называется \odmkey{строго выпуклым}{строго выпуклое нечёткое множество},
%
%если для любого $a\in (0, 1]$, все множества $G_a$ строго выпуклы 
%
%(т.~е., если $\A{x_1,x_2 \in G_a,x_m \in (x_1, x_2)}{x_m \in G_a}$).
%
%Нечёткое множество $A$ называется \odmkey{сильно выпуклым}{сильно выпуклое нечёткое множество}, если 
%$\A{a\in (0,1)}{\A{x_1, x_2}{f_A[ax_1+(1-a)x_2] > min[f_A(x_1), f_a(x_2)]}}$.
%\begin{digress}
%Строгая выпуклость не означает сильной, и наоборот. Также, очевидно, если $A$ и $B$ ограничены, то их объединение и пересечение тоже ограничены. 
%\end{digress}
%\end{frame}
%
%\begin{frame}
%\frametitle{1.9.3. Свойства нечётких множеств (6/7)}
%Пусть $A$~--- выпуклое нечёткое множество и точка $M = sup_xf_A(x)$. Если $A$ ограничено, тогда, как показано выше, либо $M$ достижима для какой-то точки
%$x_0$, либо существует хотя бы одна точка $x_0$, в которой $M$ существенно достигается, то есть 
%для каждого $e > 0$ точки каждой сферической окрестности $x_0$ содержатся в множестве $Q(e) = \Set{x}{M-f_A(x)\leq e}$. В частности, если $A$ сильно выпуклое, и $x_0$ достигается, то $x_0$ является уникальной. Более того, пусть $ C(A)$~--- множество всех точек в $X$, в котором $M$ существенно достижима. Это множество будет называться \odmkey{ядром}{ядро нечёткого множество} $А$. 
%
%\begin{theorem}
%Если $A$~--- выпуклое нечёткое множество, то его ядро~--- выпуклое множество.
%\end{theorem}
%\begin{proof}
%Будет достаточно показать, что если $M$ существенно достижима в $x_0$ и $x_1$, $x_1\neq x_0$, то она также существенно достижима для любого $x = ax_0+(1-a)x_1$, $0\leq a\leq 1$.
%Возьмём $P$ - цилиндр радиуса $e$, с осью, проходящей через $x_0$ и $x_1$. Пусть $x_0'$ точка~--- сферы радиуса $e$ с центром в $x_0$, и $x_1'$~--- точка сферы радиуса $e$ с центром в $x_1$, для которого выполнены неравенства: $f_A(x_0') \geq M - e$, $f_A(x_1') \geq M - e$. 
%\end{proof}
%\end{frame}
%
%\begin{frame}
%\frametitle{1.9.3. Свойства нечётких множеств (7/7)}
%\begin{proof}
%(Продолжение) 
%
%Тогда из выпулости $A$: $\A{u\in [x_0', x_1']}{f_A(u)\geq M - e}$.
%
%Более того, из выпуклости $P$: $\A{u\in [x_0', x_1']}{u\in P}$.
%
%Так как $[x_0', x_1']\in P$, $\A{x\in [x_0, x_1]}{distance(x, [x_0', x_1']) \leq e}$.
%
%Следовательно, сфера радиуса $e$ с центром $x$ будет содержать хотя бы одну точку из $[x_0', x_1']$ и, следовательно, будет содержать хотя бы одну точку $w : f_A(w) \geq M - e$. 
%
%Это означает, что $M$ существенно достижима в $x$.
%\end{proof}
%
%\begin{corollary}
%Если $X = E^1$, и $A$ сильно выпуклое, то точка, в которой $M$ существенно достижима, уникальна.
%\end{corollary}
%\end{frame}



\begin{frame}
\frametitle{1.9.6. Нечёткие множества (5/5)}
\begin{digress}
Элементами обычной логики являются простые высказывания~--- утверждения, которые либо истинны, либо ложны (см. 
\DMref{п.~4.1.1}{4.1.1.}).
Для построения такой логики вполне достаточно
наивной теории множеств с чётким понятием принадлежности.
Однако, в реальных задачах часто приходится сталкиваться с ситуациями, где этого недостаточно, где невозможно однозначно ответить на вопрос принадлежит ли какой-либо объект множеству или нет, верно некоторое утверждение или нет. 
Возможно лишь \em{оценить} степень  принадлежности элемента или степень достоверности утверждения. Тогда необходимо использование нечёткой логики, построенной на понятиях нечётких множеств.
Нечёткая логика тесно связана с вероятностью и
используется во многих приложениях,
в частности, в искусственном интеллекте.
\end{digress}
\end{frame}

\begin{frame}
\frametitle{Онтология наивной теории множеств}
\begin{center}
\includegraphics[width=14cm]{ODM1_9.png}
\end{center}

%{\color{red} Контрольный вопрос:
%что целесообразно добавить к содержанию раздела?}
\end{frame}

%\begin{frame}
%\frametitle{Предметный указатель}
%
%
%\end{frame}





% ===== tex/DM2022_2.tex =====
\section{2. Алгебраические структуры}

\begin{frame}
\frametitle{2. Алгебраические структуры}


Алгебраические методы находят самое широкое применение при
формализации различных предметных областей. Грубо говоря, при
построе\-нии модели предметной области всё начинается с введения
подходящих обозна\-чений для операций и отношений с последующим
исследованием их свойств. Вла\-дение алгебраической терминологией,
таким образом, входит в арсенал средств, необходимых для
абстрактного моделирования, предшествующего практическому
программированию задач конкретной предметной области. Материал
этой главы, помимо введения в терминологию общей алгебры, содержит
некоторое количество примеров конкретных алгебраических структур.
Вначале бегло рассматриваются классические структуры, которые
обычно включаются в курсы общей алгебры, а 
также элементарное введение в теорию чисел, 
необходимую в некоторых современных приложениях.
Затем обсуждаются
некоторые более специальные структуры, наиболее часто применяемые
в конкретных программных построениях. 
Особого внимания заслуживает
понятие матроида, обсуждаемое в этой главе, поскольку матроиды
тесно связаны с важнейшим классом эффективных алгоритмов,
называемых \odmkey{жадными}{Алгоритм (algorithm)!жадный (greedy)}.
\end{frame}

\begin{frame}
\frametitle{2.1. Алгебры и морфизмы (1/2)}
\subsection{2.1. Алгебры и морфизмы}
\label{2.1.}

Словом <<алгебра>>  называют, вообще говоря,  раздел
математики, посвящённый изучению систем операций 
над элементами некоторого множества.
В такой самой общей постановке область применения алгебраических методов 
оказывается практически неограниченной,
поскольку в любой предметной области
явно выделяются определённые множества объектов
и характерные операции с этими объектами.
Чтобы подчеркнуть весьма общий характер алгебраических методов,
слово <<алгебра>> в этом смысле употребляют вместе с подходящим прилагательным:
<<общая алгебра>>, <<универсальная алгебра>>, <<абстрактная алгебра>>. 

\medskip
При этом конкретные объекты, изучаемые в алгебре, 
также называют словосочетанием, в которое входит слово алгебра.
Например, <<алгебра подмножеств>> 
(\DMref{п.~1.2.8}{1.2.8.}), <<булева алгебра>> (\DMref{п.~2.6.6}{2.6.6.}),
<<алгебра булевых функций>> 
(\DMref{п.~3.2.4}{3.2.4.}).
Для краткости и без риска возникновения
недоразумений слово <<алгебра>> без дополнительных слов
в этом курсе используется как
родовое понятие для обозначения разнообразных алгебраических
структур, как конкретных, так и абстрактных. 

\end{frame}

\begin{frame}
\frametitle{2.1. Алгебры и морфизмы (2/2)}
\raggedbottom
\begin{tabular}{p{11mm}|p{37mm}|p{95mm}}
  № & Название параграфа
  & Ключевые термины и обозначения \\
  \hline
\DMref{2.1.1.}{2.1.1.}\RS
 & Операции и их носители &
{$n$-арная операция; префиксная, инфиксная и постфиксная формы записи; 
универсальная алгебра $\Alg{M}{\Sigma}$;
носитель, сигнатура и тип алгебры; модель;
многоосновная алгебра} 
\\
{\DMref{2.1.2.}{2.1.2.}} & Замыкания и подалгебры & {замкнутое подмножество носителя, замыкание относительно сигнатуры}\\  
{\DMref{2.1.3.}{2.1.3.}} & Система образующих &
{конечно-порождённая алгебра; 
определяющие соотношения, ранг, свободная алгебра}\\ 
{\DMref{2.1.4.}{2.1.4.}}\RS & Свойства операций &
{ассоциативность, коммутативность, дистрибутивность,
поглощение, идемпотентность, нейтральный элемент,
сократимость, обратный элемент, разрешимость}\\  
\DMref{2.1.5.}{2.1.5.}\RS & Таблицы Кэли &
{латинский квадрат}\\ 
{\DMref{2.1.6.}{2.1.6.}} & Гомоморфизмы и изоморфизмы
алгебр &
{гомоморфизм, изоморфизм, мономорфизм, эпиморфизм, эндоморфизм, автоморфизм; коммутативная диаграмма;
теория категорий}\\
\end{tabular}  


\end{frame}

\begin{frame}
\label{2.1.1.}
\subsubsection{2.1.1. Операции и их носители }
\frametitle{2.1.1. Операции и их носители (1/2)}

Всюду определённая (тотальная) функция $\Map{\varphi}{M^n}{M}$
называется \odmkey{$n$-ар\-ной 
операцией}{Операция (operation)!$n$-арная ($n$-ary)} на $M$. Если
операция $\varphi$~--- бинарная,
то обычно используют \odmkey{инфиксную форму записи}{Форма записи (notation)!инфиксная (infix)} $a \varphi b$ вместо $\varphi
(a,b)$ или $a * b$, где $*$~--- знак операции.
Если же операция унарная, 
то используют \odmkey{префиксную форму записи}{Форма записи (notation)!префиксная (prefix)}  
$\varphi a$ или \odmkey{постфиксную форму записи}{Форма записи (notation)!постфиксная (postfix)}  $a\varphi$. 

\smallskip
\begin{remark}
Всякая $n$-арная операция $\Map{\varphi}{M^n}{M}$ 
задаёт $(n+1)$-арное отношение $\Phi \subset M^{n+1}$ на множестве $M$:
$\Phi\Assign \Set{(\Tuple{x}{1}{n}, y)\in M^{n+1}}{y=\varphi(\Tuple{x}{1}{n})}$.
\end{remark}

\smallskip
Множество $M$
вместе с набором операций $\Sigma =\{\Tuple{\varphi}{1}{m}\}$,
${\Map{\varphi_i}{M^{n_i}}{M}}$, где $n_i$~--- арность операции
$\varphi_i$, называется \odmkey {алгебраической
структурой}{Структура алгебраическая (algebraic structure)},
\odmkey{универсальной алгеброй} {Алгебра (algebra)!универсальная
(universal)}
 или просто
\odmkey{алгеброй}{Алгебра (algebra)}. При этом множество $M$
называется
\odmkey{основным (несущим)}{Множество (set)!несущее (support)}
множеством или \odmkey{основой (носителем)}{Носитель (support)};\index{Основа (base)}
 вектор арностей $(n_1, \dots,n_m)$ называется \odmkey{типом}{Тип (type)}; множество операций
$\Sigma$ называется \odmkey{сигнатурой}{Сигнатура (signature)}. Запись:
$\Alg{M}{\Sigma}$ или $\Alg{M}{\Tuple{\varphi}{1}{m}}$. Операции
$\varphi_i$ \odmkey{конечноместны} {Операция (operation)!конечноместная
(finitude)},
сигнатура $\Sigma$ конечна. Носитель не~обязательно конечен, но обязательно не
пуст.

\smallskip
Если в качестве $\varphi_i$ допускаются не только функции, но и
отношения, то множество $M$ вместе с набором операций и отношений
называется \odmkey{моделью}{Модель (model)}. 

\end{frame}

\begin{frame}
\frametitle{2.1.1. Операции и их носители (2/2)}

\begin{remark}
Далее для обозначения алгебры 
используется прописная рукописная буква, а для обозначения её
носителя~--- соответствующая обычная прописная буква:
$\mathcal{A}=\Alg{A}{\Sigma}$. Такое соглашение позволяет
использовать вольности в обозначениях, не вводя явно две буквы для
алгебры и для носителя, а подразумевая одну из них по умолчанию.
Например, выражение <<алгебра $A$>> означает алгебру $\mathcal{A}$
с носителем $A$, а запись
$a\in\mathcal{A}$ означает элемент $a$ из носителя $A$ алгебры
$\mathcal{A}$.
\end{remark}

\medskip
В приложениях обычно
используется следующее обобщение понятия алгебры.\\ Пусть
$M=\{\Tuple{M}{1}{n}\}$~--- множество \odmkey{основ}{Основа
(base)}, $\Sigma=\{\Tuple{\varphi}{1}{m}\}$~--- сигнатура, \\причём
$\varphi_i\colon M_{i_1}\times {\dots} \times M_{i_{n_i}}\to M_{i_0}$.
Тогда $\Alg{M}{\Sigma}$ называется \odmkey{многоосновной}{Алгебра
(algebra)!многоосновная (multi-based)} алгеброй. 
Другими
словами, многоосновная алгебра имеет несколько носителей,
а~операция сигнатуры действует из прямого произведения некоторых
носителей в~некоторый носитель.

\smallskip
\begin{remark}
Использование многоосновных алгебр позволяет
считать любую модель многоосновной алгеброй.
Для этого достаточно ввести в число основ
истинностные значения (\DMref{п.~3.1.1.}{3.1.1.})
$\{0,1\}$ и считать, что отношение 
$\varphi \subset M\times M$ является операцией\\
$\Map{\varphi}{M\times M}{\{0,1\}}$. 
\end{remark}

\end{frame}

\begin{frame}
\label{2.1.2.}
\subsubsection{2.1.2. Замыкания и подалгебры }
\frametitle{2.1.2. Замыкания и подалгебры (1/3)}

Если
$
\A{\Tuple{x}{1}{n} \in X}{\varphi(\Tuple{x}{1}{n}) \in X}$,
то подмножество носителя $X \subset M$ называется \odmkey{замкнутым}
{Множество (set)!замкнутое (closed)}
относительно операции $\varphi$. 
Если $X$ замкнуто относительно всех $\varphi\in\Sigma$, то
$\Alg{X}{\Sigma_X}$ называется \odmkey {подалгеброй}{Подалгебра (subalgebra)}
$\Alg{M}{\Sigma}$, где $\Sigma_X = \{\varphi^X_{i}\}, \quad
\varphi^X_{i} = \varphi_i \big\vert_{{X}^{k}}, \quad k = n_i$.

\begin{examples}
\begin{enumerate}
\item
Алгебра $\Alg{\Bbb R}{+,\cdot}$~--- \odmkey{поле
вещественных чисел}{Поле (field)!веществественных чисел (of real
numbers)}.
 Тип~--- $(2,2)$. Все конечные
подмножества, кроме $\{0\}$, не замкнуты относительно сложения, и
все конечные подмножества, кроме $\{0\}$, $\{1\}$,
$\{0,1\}$, $\{0,1,-1\}$ и $\{1,-1\}$,
 не замкнуты относительно умножения.
 \odmkey{Поле рациональных чисел}{Поле (field)!рациональных
чисел (of rational numbers)} $\Alg{\Bbb Q}{+,\cdot}$ образует
подалгебру.
 \odmkey{Кольцо целых чисел}{Кольцо (ring)!целых чисел (of integer numbers)}
 $\Alg{\Bbb Z}{+,\cdot}$ образует подалгебру алгебры рациональных
 и, соответственно, вещественных чисел.
\item
Алгебра $\Alg{2^M}{\cup,\cap,^-}$~--- \odmkey{алгебра
подмножеств}{Алгебра (algebra)!подмножеств (of subsets)} над
множеством $M$. Тип~--- $(2,2,1)$. При этом $\A{X\subset
M}{\Alg{2^X}{\cup,\cap,^-}}$~--- её подалгебра.
\item
Алгебра
гладких функций $\Alg{\Set{f}{\Map{f}{\Bbb R}{\Bbb
R}}}{\frac{d}{dx}}$, где $\frac{d}{dx}$~--- операция
дифференцирования. Множество полиномов одной переменной $x$
образует подалгебру, которая обозначается $\Bbb R[x]$.
\end{enumerate}
\end{examples}

\end{frame}

\begin{frame}
\frametitle{2.1.2. Замыкания и подалгебры (2/3)}

\begin{theorem} Непустое пересечение подалгебр некоторой алгебры
образует подалгебру этой же алгебры.
\end{theorem}
\begin{proof}
Пусть $\Alg{X_i}{\Sigma_{X_i}}$~--- подалгебра  $\Alg{M}{\Sigma}$.
Обозначим $X\Assign\bigcap X_i$. Тогда\\
$\A{i}{\A{j}{\varphi_j^{X_i}(\Tuple{x}{1}{n_j})\in X_i}} 
\Impl\A{j}{\varphi_j^{X}(\Tuple{x}{1}{n_j})\in X}$.
Другими словами, любые операции любой подалгебры на её носителе не выводят за его пределы, а мы рассматриваем пересечение  носителей подалгебр.
\end{proof}

\odmkey{Замыканием множества}{Замыкание (closure)!множества (of a
set)} $X \subset M$ относительно сигнатуры $\Sigma$ (обозначается
$[X]_{\Sigma}$) называется множество всех элементов (включая сами
элементы множества $X$), которые можно получить, применяя операции
из $\Sigma$. Если сигнатура подразумевается, её можно не
указывать.


\begin{theorem}
Замыкание обладает следующими свойствами:
\begin{enumerate}
\item[1)] $X \subset Y \Impl [X]\subset [Y]$;\\
\item[2)] $X \subset [X]$;\\
\item[3)] $[[X]]=[X]$;\\
\item[4)] $[X] \cup [Y] \subset [X \cup Y]$.
\end{enumerate}
\end{theorem}

\end{frame}

\begin{frame}
\frametitle{2.1.2. Замыкания и подалгебры (3/3)}
\begin{proof}
\proofsection{\color{blue}1)} Очевидно.
\proofsection{\color{blue}2)} По определению.
\proofsection{\color{blue}3)} По определению.
\proofsection{\color{blue}4)} В
$[X \cup Y]$ входят не только те элементы, которые
можно получить из элементов $X$  и из элементов $Y$, но и те элементы,
которые можно получить из комбинаций этих элементов.   
\end{proof}

\smallskip
\begin{example}
В алгебре натуральных чисел  $\Alg{\Bbb
N}{+,\cdot}$ (в арифметике) замыканием числа 2 являются чётные числа, то~есть
$[\{2\}]=\Set{n\in \Bbb N}{n=2k\And k\in\Bbb N}$.
\end{example}

\smallskip
Пусть $\mathcal{A}=\Alg{A}{\Sigma}$~--- некоторая алгебра и
$X_1,\dots,X_n\subset A$~--- некоторые подмножества носителя, а
$\varphi\in\Sigma$~--- одна из операций алгебры.
Тогда используется следующее соглашение об обозначениях:
$\ \ \varphi(X_1,\dots,X_n)\Def \Set{\varphi(x_1,\dots,x_n)}{x_1\in
X_1\And\dots\And x_n\in X_n},
$ \\
то есть алгебраические операции можно применять не только к
отдельным элементам, но и к множествам (подмножествам носителя),
получая, соответственно, не
отдельные элементы, а множества (подмножества носителя), аналогично \DMref{п.~1.6.3}{1.6.3.}.
\end{frame}

\begin{frame}
\label{2.1.3.}\subsubsection{2.1.3. Система образующих }\frametitle{2.1.3. Система образующих (1/3)}

Множество $S\subset M$ называется \odmkey{системой образующих}
{Система (system)!образующих (generating)} 
или 
\odmkey{порождающим множеством}
{Множество (set)!порождающее (generating)} алгебры
$\Alg{M}{\Sigma}$, если $[S]_{\Sigma} = M$. Другими словами, любой элемент носителя можно породить, 
применяя операции к образующим. 
Если алгебра имеет
конечную систему образующих, то алгебра называется
\odmkey{конечно-порождённой}{Алгебра (algebra)!конечно-порожденная
 (finitely generated)}.
Бесконечные алгебры могут иметь конечные системы образующих.

\begin{examples}
\begin{enumerate}
\item Алгебра \odmkey{натуральных чисел}{Число (number)!натуральное (natural)}
$\Alg{\Bbb N}{+}$ с операцией сложения имеет конечную систему образующих
$1\in\Bbb N$.
\item Алгебра $\Alg{\mathbb{N}}{*}$ 
с тем же носителем, но с операцией умножения не является конечно\-порождённой. 
Системой образующих этой алгебры
является множество простых чисел, бесконечность которого установлена в \DMref{п.~2.4.4}{2.4.4.}.
\item
Пустое множество и множество одноэлементных
подмножеств $\cal{C}^1(M)$ (\DMref{п.~1.2.8}{1.2.8.}) являются 
системой образующих булеана $2^M$ 
в алгебре подмножеств для любого множества $M$. 
\item
Пустое множество и любая полная цепочка 
подмножеств (\DMref{п.~5.1.4}{5.1.4.}) являются 
системой образующих булеана $2^M$ 
в алгебре подмножеств для любого множества $M$. 
\end{enumerate}
\end{examples}

\end{frame}

\begin{frame}
\frametitle{2.1.3. Система образующих (2/3)}
Последние два примера показывают, 
что порождающее множество не
обязательно единственное. Например,
любое надмножество порождающего 
множества является порождающим,
сам носитель является порождающим множеством.
Порождающее множество называется
минимальным, если его подмножества не являются порождающими.

Мощность наименьшего 
порождающего множества называется
\odmkey{рангом}{Ранг (rank)!алгебры (of algebra)} алгебры.

\begin{example}
Алгебра натуральных чисел
$\Alg{\Bbb N}{+}$ имеет ранг $1$,
а алгебра подмножеств множества $M$
имеет ранг 
$|M|$.
\end{example}

\smallskip
Выражение элемента носителя через образующие 
не обязательно единственно. Другими словами,
разные выражения, построенные 
из образующих с помощью операций 
алгебр, могут иметь одно значение.

\begin{example}
В алгебре натуральных чисел 
$\Alg{\Bbb N}{+, \cdot}$ система образующих
состоит из одного элемента $\{1\}$,
при этом $(1+1)+(1+1)=4=(1+1)\cdot(1+1)$.  
\end{example} 

\smallskip
Равенство различных выражений над
образующими называют 
\odmkey{ограничением}{Ограничение (constraint)}
(а также \odmkey{определяющим соотношением}{Соотношение (relationship)!определяющее (defining)}, см. \DMref{п.~2.2.2}{2.2.2.}). 
Алгебра, в которой нет ограничений,
то есть любой элемент носителя выражается через
образующие единственным образом,
называется \odmkey{свободной}{Алгебра (algebra)!свободная (free)}.
  
\end{frame}

\begin{frame}
\frametitle{2.1.3. Система образующих (3/3)}
\begin{example}
Рассмотрим алгебру натуральных чисел 
$\Alg{\Bbb N}{\plusplus}$, где единственная
унарная операция $\plusplus$ добавляет
к аргументу $1$. В этой алгебре система образующих
состоит из одного элемента $\{1\}$,
при этом алгебра является свободной,
поскольку каждое натуральное число
определяется просто количеством применений
операции $\plusplus$.
\end{example}

\smallskip
Пусть заданы набор функциональных символов $\Phi =
\{\Tuple{\varphi}{1}{m}\}$, служащих обозначениями функций
некоторой\vadjust{\penalty-1000} сигнатуры $\Sigma$ типа $N=(\Tuple{n}{1}{m})$,  и
множество переменных $V=\{x_1,x_2,\dots\}$.
 Определим множество
\odmkey{термов}{Терм (term)} $T$ индуктивным образом:
\begin{enumerate}
\item[1)] $V \subset T$;
\item[2)] $\Tuple{t}{1}{n_i} \in T \And \varphi_i \in \Phi \Impl
\varphi_i (\Tuple{t}{1}{n_i}) \in T$.
\end{enumerate}

Алгебра $\Alg{T}{\Phi}$ называется \odmkey{свободной алгеброй
термов}{Алгебра (algebra)!термов свободная (free algebra of
terms)}. Носителем этой алгебры является множество термов, то~есть
формальных выражений, построенных с помощью знаков операций
сигнатуры $\Sigma$. Заметим, что множество $V$ является системой
образующих свободной алгебры термов.

\begin{remark}
Свободная алгебра термов имеет
бесконечный носитель при любой сигнатуре.
Алгебра с конечным носителем не может быть свободной.
\end{remark}

\end{frame}

\begin{frame}
\label{2.1.4.}\subsubsection{2.1.4. Свойства операций }\frametitle{2.1.4. Свойства операций (1/8)}
Некоторые часто встречающиеся свойства операций имеют специальные
названия. 
Пусть задана алгебра $\Alg{M}{\Sigma}$ и $a,b,c\in M$;
$\circ,\diamond\in\Sigma$; $\circ,\Map{\diamond}{M\times M}{M}$.
Тогда свойства операций определяются следующим образом:
\begin{enumerate}
\item[(А1)]
\odmkey{ассоциативность}{Операция (operation)!ассоциативная
(associative)}: $\ \A{a,b,c\in M}{(a \circ b)\circ c = a \circ(b \circ c)};$
\item[(А2)]
\odmkey{коммутативность}{Операция (operation)!коммутативная
(commutative)}: $\ \A{a,b\in M}{a \circ b = b \circ a}$;
\item[(А3)]
\odmkey{дистрибутивность}{Операция (operation)!дистрибутивная
(distributive)} $\diamond$ относительно $\circ$ 
\\слева: 
$ \ \A{a,b,c\in M}{a \diamond (b \circ c)=(a \diamond b) \circ (a \diamond c)};$ \\
справа:
$\ \A{a,b,c\in M}{(a \circ b) \diamond c =(a \diamond c) \circ (b \diamond c)};$
\item[(А4)]
\odmkey{поглощение}{Поглощение (absorption)} ($\diamond$
поглощает $\circ$)\\ 
справа: $\ \A{a,b\in M}{(a \circ b)\diamond a=a};$\\
слева: $\ \A{a,b\in M}{a  \diamond (a \circ b)=a};$
\item[(А5)]
\odmkey{идемпотентность}{Операция (operation)!идемпотентная
(idempotent)}: $\ \A{a\in M}{a \circ a = a}.$
%\end{tabbing}
\end{enumerate}

\begin{remark}
Если операция $\diamond$ коммутативна,
то говорят просто о дистрибутивности или поглощении, не добавляя уточнения слева или справа.
\end{remark}
\end{frame}

\begin{frame}
\frametitle{2.1.4. Свойства операций (2/8)}

\begin{examples}
\begin{enumerate}
\item 
Ассоциативные операции: сложение и умножение чисел,
объединение и пересечение множеств, 
композиция отношений.
Неассоциативные операции: возведение чисел в степень, вычитание
множеств.\\[-6pt]
\item 
Коммутативные операции: сложение и умножение
чисел, объединение и пересечение множеств. Некоммутативные
операции: умножение матриц, композиция отношений, возведение в
степень.\\[-6pt]
\item 
Дистрибутивные операции: умножение относительно
сложения чисел. Недистри\-бу\-тивные операции: возведение в степень
дистрибутивно относительно умножения справа, но не слева:
$\bigl((ab)^c=a^cb^c, a^{bc}\ne a^ba^c\bigr)$.\\[-6pt]
\item
Пересечение
поглощает объединение, объединение поглощает пересечение. Сложе\-ние
и умножение не поглощают друг друга.\\[-6pt]
\item 
Идемпотентные операции:
наибольший общий делитель натуральных чисел, объединение и
пересечение множеств, минимум и максимум чисел. Неидемпотентные операции: сложение и
умножение чисел.
\end{enumerate}
\end{examples}

\end{frame}

\begin{frame}
\frametitle{2.1.4. Свойства операций (3/8)}

Свойства операций не являются полностью независимыми~--- некоторые комбинации\\ свойств влекут 
друг друга. 

\begin{theorem} \NB{[1]}
Если обе операции поглощают 
друг друга, то они идемпотентны.
\end{theorem}

\begin{proof}
$a=(a \circ c)\diamond a=(a\circ (a\diamond b))\diamond a =a \diamond a$, $c=a\diamond b$\\
$a=(a \diamond c)\circ a=(a\diamond (a\circ b))\circ a =a \circ a$, $c=a\circ b$.
\end{proof}

\smallskip
Продолжаем перечисление свойств операций: 
\begin{enumerate}
\item[(А6)]
\odmkey{нейтральный элемент}{Элемент (element)!нейтральный
(neutral)}
\\левый: 
$\E{e\in M}{\A{a\in M}{e\circ a =a}}$ \\
правый:
$\E{e\in M}{\A{a\in M}{a\circ e =a}}$
\end{enumerate}

\begin{theorem} \NB{[2]}
Если существуют левый и правый нейтральные элементы, то они равны.
\end{theorem}

\begin{proof}
Пусть $e_l$ и $e_r$ --- соответственно левый и правый нейтральные элементы относительно операции $\circ$. 
Тогда  $e_l  \circ  e_r  = e_l  \And  e_l  \circ  e_r  = e_r  \Impl e_l=e_r$.
\end{proof}

\smallskip
Элемент $e \in M$ называется 
\odmkey{нейтральным}{Элемент (element)!нейтральный (neutral)!относительно операции (for operation)} относительно операции $\circ$, если он нейтральный слева и справа,  
$\A{a\in M}{ e  \circ  a = a = a  \circ  e}$.

\begin{example}
Ноль является нейтральным элементом по сложению целых, рациональных и вещественных чисел.
\end{example}
\end{frame}

\begin{frame}
\frametitle{2.1.4. Свойства операций (4/8)}


\begin{theorem} \NB{[3]} Если нейтральный элемент 
%относительно бинарной операции $\circ$  
существует, 
то он единственен.
\end{theorem}
\begin{proof}
Пусть существуют два нейтральных
элемента: $e_1, e_2$. \\ Тогда $e_1\circ
e_2=e_1\And e_1 \circ e_2= e_2\Impl e_1=e_2$.
\end{proof}

\smallskip
Продолжаем перечисление свойств операций: 
\begin{enumerate}
\item[(А7)]
\odmkey{сократимость}{Операция (operation)!сократимая (cancellation)}
\\ 
слева : 
$\ \A{a,x,y \in M}{a\circ x = a\circ y \Impl x= y}$
\\
справа: 
$ \ \A{a,x,y \in M}{x\circ a = y\circ a \Impl x= y}$
\end{enumerate}

\begin{example}
Сложение и умножение натуральных чисел~--- сократимые операции слева и справа,
а объединение и пересечение множеств несократимы ни слева, ни справа.
\end{example}

\begin{theorem} \NB{[4]}
Если операция $\circ$ ассоциативна и сократима справа, и относительно неё существует левый нейтральный элемент, то существует и правый нейтральный элемент.
\end{theorem}
\begin{proof}
Пусть $e$ --- левый нейтральный элемент относительно операции $\circ$. 
Тогда $\A{a\in M}{a\circ a  =  a  \circ (e  \circ  a)  =   (a  \circ  e)  \circ  a   \Impl a  \circ  e = a}$
\end{proof}
\begin{corollary}
Если операция $\circ$ ассоциативна и сократима слева, и относительно неё существует правый нейтральный элемент, то существует и левый нейтральный элемент.
\end{corollary}

\end{frame}

\begin{frame}
\frametitle{2.1.4. Свойства операций (5/8)}
\begin{remark}
Свойство ассоциативности существенно. Рассмотрим
операцию вычитания множеств, которая не является ассоциативной 
$A\setminus (B\setminus C)\neq (A\setminus B)\setminus C$. Эта операция имеет правый нейтральный элемент
$A\setminus\emptyset=A$,
но не имеет левого нейтрального элемента и не сократима 
ни слева, ни справа.
Таким образом, только в случае сократимой ассоциативной операции можно говорить про нейтральный элемент, не указывая, левый он или правый.
\end{remark}

\begin{example}
Операция умножения на множестве 
$\Bbb{N}$ ассоциативна и сократима, и число $1$ является нейтральным элементом.
При этом, операция возведения в степень не ассоциативна,
сократима справа 
$\A{a,x,y \in \Bbb{N}}{x^a = y^a \Impl x= y}$,
но не слева 
$\Not(1^x = 1^y \Impl x= y)$,
и
множество $\Bbb{N}$ относительно операции возведения в степень имеет правый нейтральный элемент
$\A{n\in \Bbb{N}}{n^1 =n}$, 
но не имеет левого нейтрального элемента\\
$\Not\E{x\in \Bbb{N}}{\A{n\in \Bbb{N}}{x^n =n}}.$ 
\end{example}

\smallskip
Продолжаем перечисление свойств операций:
\begin{enumerate}
\item[(А8)] \odmkey{обратный элемент}{Элемент (element)!обратный (inverse)}
относительно операции $\circ$
с нейтральным элементом $e$ \\
слева: $\ a^{-1} \circ a=e$,\\ 
справа: $\ a \circ a^{-1}=e$.
\end{enumerate}

\end{frame}

\begin{frame}
\frametitle{2.1.4. Свойства операций (6/8)}
\begin{theorem} \NB{[5]}
Если операция $\circ$ ассоциативна, и существует левый обратный элемент, то существует и правый обратный элемент, и они равны между собой.
\end{theorem}

\begin{proof}
Пусть $a_1$~--- левый обратный элемент относительно ассоциативной бинарной операции $\circ$ 
для элемента $a$, $e$~--- нейтральный элемент (по ассоциативности он единственный).
Тогда
$a=e\circ a = a\circ e = a \circ  (a_1 \circ a) =
 (a \circ a_1)\circ  a  \Impl a \circ a_1=e$.
\end{proof}
\begin{remark}
В случае ассоциативной операции, можно говорить про обратный элемент, не указывая, левый он или правый.
\end{remark}

\begin{theorem} \NB{[6]}
Обратный элемент относительно ассоциативной операции единственен.
\end{theorem}
\begin{proof}
Пусть $a\circ a^{-1}=a^{-1}\circ a=e\And a\circ b = b\circ a=e$. \\Тогда
$a^{-1}={a^{-1}\circ e}={a^{-1}\circ(a\circ b)}=
{(a^{-1}\circ a)\circ b}={e\circ b=b}$.
\end{proof} 

\begin{example} 
Каждое вещественное число $a$
относительно операции сложения $+$ имеет обратный
элемент $a^{-1}= - a$.
\end{example}

\smallskip
Завершаем перечисление свойств операций:
\begin{enumerate}
\item[(А9)]
\odmkey{разрешимость}{Операция (operation)!разрешимая (resolvable)}:
$\ \A{a,b\in M}{\E{!\ x, y \in M}{a\circ x = y\circ a = b}}$
\end{enumerate}

Другими словами, операция разрешима,
если по результату операции и одному из операндов,
можно найти другой операнд.


\end{frame}

\begin{frame}
\frametitle{2.1.4. Свойства операций (7/8)}
\begin{example}
Операция сложения в целых числа разрешима,
а операция умножения~--- нет, потому что 
$\A{x\in\Bbb{Z}}{0\ast x =0}$.  
\end{example}

\begin{remark}
Разрешимость~--- весьма полезное 
свойство, поскольку позволяет решать уравнения
$a\circ x = b$.
Разрешимые операции называют также
\odmkey{операциями с памятью}{Операция (operation)!с памятью (with memory)}.
\end{remark}


Рассмотренные четыре свойства, а именно: 
существование нейтрального элемента, обратимость,
сократимость и разрешимость очень тесно
связаны друг с другом, однако не тождественны 
и не сводятся друг к другу,
что показывают следующие теоремы.
 


\begin{theorem} \NB{[7]}
Если операция разрешима и существует нейтральный элемент, то существует и обратный элемент.
\end{theorem} 

\begin{proof}
$\A{a\in M}{\E{!\ x,y}{a\circ x = y\circ a = e}}$,
то есть $x$ и $y$~--- левый и правый обратные элементы. 
\end{proof}

\smallskip
\begin{theorem} \NB{[8]}
Из обратимости ассоциативной операции следует сократимость.
\end{theorem}
\begin{proof} 
Рассмотрим операцию $\circ$, 
 имеющую левый обратный элемент.\\
Пусть $a,x,y \in M$ и $a\circ x = a\circ y$. 
Домножим обе части равенства слева на $a^{-1}$, 
получим $x = y$. Аналогично для сократимости справа.
\end{proof}
Условие ассоциативности в этой теореме существенно!
\end{frame}

\begin{frame}
\frametitle{2.1.4. Свойства операций (8/8)}
\begin{example}
Неассоциативная операция вычитания множеств, обладающая обратимостью (обратным к элементу является он сам), не является сократимой. 
Так из $A\setminus B = A\setminus C$ 
не следует $B = C$, а из $B\setminus A = C\setminus A$ также не следует $B = C$.
\end{example}
	


\begin{theorem} \NB{[9]}
Из разрешимости операции следует сократимость.
\end{theorem}
\begin{proof}
Рассмотрим разрешимую операцию $\circ$.
Пусть $a, x, y \in M$ и \\
$a\circ x = a\circ y = c$. 
Из разрешимости следует, что уравнение $a\circ b = c$ имеет единственное решение, а значит $x = y$.  Аналогично для сократимости справа.
\end{proof}	
 
\begin{remark}
Из сократимости не следует ни обратимости, ни разрешимости.
\end{remark}	


\begin{example}
На множестве вещественных чисел $\Bbb{R}$ сложение  оказывается обратимым  и разрешимым, но на множестве неотрицательных
вещественных чисел $\Bbb{R}_{+}$ сложение обладает свойством сократимости, 
поскольку 
$a+x=b+x \Impl a=b$, но обратного элемента нет ни у какого элемента кроме нуля, и уравнения вида 
$a + x = b$ не разрешимы при $a > b$.
\end{example}

\smallskip
\begin{theorem} \NB{[10]}
Для ассоциативной операции из обратимости следует разрешимость.
\end{theorem}
\begin{proof}
Фиксируем $a,b \in M$. При $x = a^{-1}\circ b$ имеем $a \circ x = b$. Аналогично для уравнений вида $y \circ a = b$ ($y = b\circ a^{-1}$).
\end{proof}

%%\smallskip
%%\begin{theorem} \NB{[11]}
%%??? Из разрешимости и сократимости следует обратимость операции и существование нейтрального элемента для неё.
%%\end{theorem}
%%\begin{proof}
%%Положим $b\Assign a$ в определении 
%%разрешимости операции $\circ$.\\
%%Имеем, 
%%$\A{a}{\E{!\ x, y}{a\circ x = a \And y \circ a = a}}$, 
%%.
%%
%%%Тогда 
%% => a * x = a * e & y * a = a * e. Из сократимости: x = e = y (единственность нейтрального элемента доказана в теореме из "2.1.4. Свойства операции (2/7)").
%%Для операции с разрешимостью и нейтральным элементом, для каждого элемента существуют левый и правый обратные элементы (формулировка теоремы из "2.1.4. Свойства операции (6/7)")" .
%\end{proof}
\end{frame}


\begin{frame}
\subsubsection{2.1.5. Таблицы Кэли }\label{2.1.5.}
\frametitle{2.1.5. Таблицы Кэли (1/6)}
Бинарную операцию $\ast$, 
заданную на конечном множестве $A$, 
удобно представлять квадратной матрицей
\textbf{$K_{\ast}:$ array $[A, A]$ of $A$}. 
Обычно эту матрицу изображают в виде таблицы,
строки и столбцы которой помечены элементами 
множества $A$, 
причём \\$\A{a,b\in A}{K_{\ast}[a,b]=a\ast b}$.
Такая таблица называется \odmkey{таблицей Кэли}{Таблица (table)!Кэли (Cayley)}.
		
\begin{example} 
Рассмотрим 
булеан множества $\lbrace 1, 2 \rbrace$ с операциями пересечения и разности множеств. 
Их таблицы Кэли выглядят следующим образом:
		
\begin{center}
\begin{tabular}{|c||c|c|c|c|}
\hline
$\cap$ & $\emptyset$ & $\{1\}$ & $\{2\}$ & 
$\{1,2\}$ \\ \hline \hline
$\emptyset$ & $\emptyset$ & $\emptyset$ & $\emptyset$ & $\emptyset$ \\ \hline
$\{1\}$ & $\emptyset$ & $\{1\}$  & $\emptyset$ & $\{1\}$\\ \hline
$\{2\}$ & $\emptyset$ & $\emptyset$ & $\{2\}$ & $\{2\}$\\ \hline
$\{1,2\}$ & $\emptyset$ & $\{1\}$ & $\{2\}$ & $\{1,2\}$ \\ \hline
\end{tabular}
\hspace{2cm}
\begin{tabular}{|c||c|c|c|c|}
\hline
$\setminus$ & $\emptyset$ & $\{1\}$ & $\{2\}$ & 
$\{1,2\}$ \\ \hline \hline
$\emptyset$ & $\emptyset$ & $\emptyset$ & $\emptyset$ & $\emptyset$ \\ \hline
$\{1\}$ & $\{1\}$ & $\emptyset$  & $\{1\}$ & $\emptyset$\\ \hline
$\{2\}$ & $\{2\}$ & $\{2\}$ & $\emptyset$ & $\emptyset$\\ \hline
$\{1,2\}$ & $\{1,2\}$ & $\{2\}$ & $\{1\}$ & $\emptyset$ \\ \hline
\end{tabular}
\end{center}
\end{example}

\medskip				
\begin{remark}
Далее $K[a,*]$ означает строку 
таблицы Кэли, помеченную элементом $a$,
а $K[*,b]$ означает столбец, помеченный элементом $b$.
\end{remark}
\end{frame}
	
\begin{frame}
\frametitle{2.1.5 Таблицы Кэли (2/6)}
Некоторые свойства операций наглядно 
проявляются в свойствах таблицы Кэли.

\medskip

		
\begin{theorem} Таблица Кэли для коммутативной операции $\ast$ симметрична относительно главной диагонали.
\end{theorem}
\begin{proof} $a\ast b = b\ast a \Equiv 
K_{\ast}[a,b] = K_{\ast}[b,a]$.
\end{proof}


\medskip
\begin{theorem} Если существует элемент $a$, такой, что  $\A{b\in A}{K[a,b] = b}$, то $a$~--- левый нейтральный элемент. Если существует элемент $a$, такой, что  $\A{b\in A}{K[b,a] = b}$, то $a$~--- правый нейтральный элемент. 
\end{theorem}

\begin{proof}
Если $K[a,b] = b$ для любого $b$, то $a\ast b = b$ для любого $b$. Значит, $a$~--- левый нейтральный элемент. Аналогично для правого.
\end{proof}

\begin{remark}
Напомним, что если существуют левый и правый нейтральные элементы, то они совпадают, 
и если нейтральный элемент существует,
то он единственный.
\end{remark}

\medskip
\begin{example}
Из таблиц $K_{\cap}$ и $K_{\setminus}$ хорошо видно,
что пересечение коммутативно, а разность~--- нет,
что $\{1, 2\}$ является нейтральным элементом для пересечения, а разность не имеет нейтрального элемента.
\end{example}
\end{frame}
	
\begin{frame}
\frametitle{2.1.5 Таблицы Кэли (3/6)}
\begin{theorem} Если в каждой строке (столбце) таблицы Кэли 
присутствует нейтральный элемент, то для любого элемента существует правый (левый) обратный.   
\end{theorem}

\begin{proof} Пусть $e$~--- нейтральный элемент.
Если $K[a,b] = e$, то $b$ является правым обратным для $a$, а $a$ является левым обратным для $b$.
\end{proof}


\begin{remark}
Левый и правый обратные элементы вовсе не обязаны совпадать.
\end{remark}
\begin{example} Рассмотрим три операции, обозначаемые $+, \ast, \wedge $, заданные над множеством $A = \lbrace e, a, b, c\rbrace$. При данных таблицах Кэли для этих операций, элемент $e$ является нейтральным для всех трёх операций, 
для любого элемента $A$ существует обратный по операции $+$ (он сам), существуют левый и правый обратные по операции $\ast$, причём частично не совпадающие. 
По операции $\wedge$ для элемента $a$ нет ни левого, ни правого обратных.

\begin{center}
\begin{tabular}{|l|l|l|l|l|}
\hline
$+$ & $e$ & $a$ & $b$ & $c$ \\ \hline
$e$ & $e$ & $a$ & $b$ & $c$ \\ \hline
$a$ & $a$ & $e$ & $b$ & $c$ \\ \hline
$b$ & $b$ & $a$ & $e$ & $c$ \\ \hline
$c$ & $c$ & $a$ & $b$ & $e$ \\ \hline
\end{tabular}
\hspace{1cm}
\begin{tabular}{|l|l|l|l|l|}
\hline
$\ast$ & $e$ & $a$ & $b$ & $c$ \\ \hline
$e$ & $e$ & $a$ & $b$ & $c$ \\ \hline
$a$ & $a$ & $b$ & $c$ & $e$ \\ \hline
$b$ & $b$ & $e$ & $c$ & $a$ \\ \hline
$c$ & $c$ & $b$ & $e$ & $a$ \\ \hline
\end{tabular}
\hspace{1cm}
\begin{tabular}{|l|l|l|l|l|}
\hline
$\wedge$ & $e$ & $a$ & $b$ & $c$ \\ \hline
$e$ & $e$ & $a$ & $b$ & $c$ \\ \hline
$a$ & $a$ & $b$ & $b$ & $a$ \\ \hline
$b$ & $b$ & $c$ & $e$ & $b$ \\ \hline
$c$ & $c$ & $a$ & $b$ & $e$ \\ \hline
\end{tabular}
\end{center}
\end{example}
\end{frame}
	
\begin{frame}
\frametitle{2.1.5 Таблицы Кэли (4/6)}		
\begin{example}
Из таблицы $K_{\cap}$ хорошо видно,
что пересечение необратимо. Разность
также необратима уже потому, что нет нейтрального 
элемента.
\end{example}

\begin{theorem}
В каждой строке (столбце) таблицы Кэли операции с левой (правой) сократимостью не повторяются элементы.
\end{theorem}
\begin{proof}
Если в строке не повторяются элементы,
то $K[a,b]= K[a,c] \Equiv b=c$. Значит $a\circ b=a\circ c \Equiv b=c$ (левая сократимость). Аналогично для правой сократимости.
\end{proof}

\begin{columns}[t]
\begin{column}{10cm}

\vspace{-2\baselineskip}
\begin{example}
 В приведённой справа таблицe Кэли
 в строках не повторяются элементы и
 эта операция сократима слева, но не справа.
\end{example}
\end{column}
\begin{column}{5cm}
\begin{tabular}{|c||c|c|}
\hline
$\circ$ & $a$ & $b$  \\ \hline \hline
$a$ & $a$ & $b$  \\ \hline
$b$ & $a$ & $b$  \\ \hline
\end{tabular}
\end{column}
\end{columns}				
\begin{theorem} Диагональ таблицы Кэли для идемпотентной операции совпадает с множеством элементов.
\end{theorem}
		
\begin{proof} Если операция идемпотентна, то  на пересечении строки и столбца элемента $a$, лежащем на диагонали, стоит $a$.
\end{proof}
		
\begin{example} Пересечение множеств идемпотентно, а разность множеств --- нет.
\end{example}

\end{frame}
	
\begin{frame}
\frametitle{2.1.5 Таблицы Кэли (5/6)}

\begin{theorem}
 В каждой строке и столбце таблицы Кэли разрешимой операции есть все элементы, и они не повторяются.
\end{theorem}

\begin{proof}
Условие теоремы означает,
что для любых $a, b$ существуют единственные $c, d$, такие что $ K[a,c]=b=K[d,a]$,
значит  $a \circ c=b=d\circ a$ (разрешимость).
\end{proof}

\begin{digress}
Квадратная таблица размера $n\times n$, 
в которой содержится $n$ различных 
элементов, причём в каждой строке 
и в каждом столбце все элементы содержатся ровно
по одному разу, называется 
\odmkey{латинским квадратом}{Квадрат (square)!латинский (Latin)}.
Таким образом, таблица Кэли разрешимой 
операции является латинским квадратом.
Латинские квадраты встречаются 
в самых неожиданных разделах математики,
см., например, \DMref{п.~8.3.3}{8.3.3.}.  
\end{digress}

\smallskip
Отношение между операциями и таблицами Кэли не является 
однозначным. Все таблицы, которые получаются
друг из друга одновременной перестановкой строк и столбцов
(то есть перестановкой элементов несущего множества)
задают одну и ту же операцию и являются 
\odmkey{изоморфными}{Изоморфизм (isomorphism)!таблиц Кэли 
(of Cayley tables)}. Более формально
$$
K_1 \sim K_2 \Equiv \E{\Map{f}{M}{M}}{\A{a,b\in M}
{f(K_1[a,b])= K_2[f(a),f(b)]}},
$$ 
причём $f$ является перестановкой множества $M$.
 
\end{frame}
	
\begin{frame}
\frametitle{2.1.5 Таблицы Кэли (6/6)}
Другими словами, таблицы изоморфны, если между ними существует биективный гомоморфизм (\DMref{п.~1.7.6}{1.7.6.},
\DMref{п.~2.1.6}{2.1.6.}). 

\smallskip
Для конечного множества $|M|=n$ существует $n^{n^2}$ различных
таблиц Кэли. Однако, поскольку отношение изоморфности таблиц Кэли очевидно является эквивалентностью, 
и классы эквивалентности по этому отношению 
взаимно однозначно соответствуют операциям на множестве $M$,
существенно различных операций, то есть операций, обладающих
различными наборами свойств, на самом деле меньше.   
		
\begin{example}
Рассмотрим случай $M=\{a,b\}, n=2$. Легко понять, что существует  16 таблиц Кэли, но классов эквивалентности всего лишь 10. Ниже
приведены все 16 таблиц, распределённых на 10 столбцов, 
соответствующих классам изоморфности таблицы. Таблицы в столбцах получаются одна из другой одновременной перестановкой
строк и столбцов.  
\begin{center}
\begin{tabular}{|cc|cc|cc|cc|cc|cc|cc|cc|cc|cc|}
\hline
$a$ & $a$ & $a$ & $a$ & $a$ & $b$ & $a$ & $a$& $a$ & $b$ &
$a$ & $b$ & $b$ & $a$ & $b$ & $b$ & $b$ & $a$& $b$ & $b$ \\ 
$a$ & $a$ & $a$ & $b$ & $a$ & $a$ & $b$ & $b$& $a$ & $b$ &
$b$ & $a$ & $a$ & $b$ & $b$ & $a$ & $b$ & $b$& $b$ & $b$ \\ \hline
    &     & $b$ & $a$ & $a$ & $a$ & $b$ & $b$& $b$ & $a$ &
    &     &     &     & $a$ & $b$ & $b$ & $b$&     &     \\
    &     & $a$ & $a$ & $b$ & $a$ & $a$ & $a$& $b$  & $a$ &
    &     &     &     & $b$ & $b$ & $a$ & $b$&     &     \\
\hline 
\end{tabular}
\end{center}
\end{example}
		

%\medskip		
%\TODO{{\color{blue}(Важно! Про изоморфизм.)}
%Вообще говоря, всего существует $n^{n^2}$ таблиц Кэли. Но среди них есть изоморфные,
%то есть задающие одну и ту же операцию.
%Их матрицы получаются {\color{blue}(как получаются изоморфные матрицы? Сколько их?)}.
%{\color{blue} Пример.} Ниже приведены все
%{\color{blue} 2? 8? 10?} попарно неизоморфные 
%таблицы Кэли для $n=2$ и под каждой 
%таблицей указано,какими свойствами 
%из набора свойств А1--А9 обладает эта операция
%}

\end{frame}

%	\begin{frame}
%		\frametitle{2.1.5 Таблицы Кэли (7/8)}
%		
%		Так как в операции ассоциативности участвуют три элемента, мы не можем по таблице определить выполнение этого свойства.Тем не менее, \odmkey{тест ассоциативности Лайта} \; может определить ассоциативность с меньшими усилиями, чем грубый перебор. Ниже приведена сама процедура, а которой мы рассматриваем алгебру $(A, \cdot)$.
%		
%		Выберем $a \in A$ и рассмотрим две новые операции: $x*y=x \cdot (a \cdot y)$ и $x \star y=(x \cdot a) \cdot y$
%		Затем необходимо построить и сравнить таблицы Кэли для этих двух операций. Если они совпали, то $x \cdot (a \cdot y) = (x \cdot a) \cdot y$ для всех $x$ и $y$. Указанная процедура повторяется для всех элементов $A$.
%		
%		К сожалению, эта процедура всё равно имеет сложность $O(n^3)$ в худшем случае (что, впрочем, случается достаточно редко) для алгебры на $n$ элементах.
%		
%	\end{frame}
	
%	
%	\begin{frame}
%		\frametitle{2.2.6 Квазигруппы (?/?)}
%		
%		\begin{theorem} Таблица Кэли для квазигруппы является латинским квадратом.
%		\end{theorem}
%		
%		\begin{proof}
%			Допустим, что в какой-то строке (случай столбца аналогичен) найдутся два совпадающих элемента. Это значит, что $\exist a,b,c \in A: ab = ac$. Значит, $a^{-1}ab = a^{-1}ac$, то есть $b = c$.% бред
%		\end{proof}
%		
%		
%		\begin{example} Рассмотрим $A = \lbrace 1, -1, i, -i\rbrace$, с заданной на нём операцией умножение. Данная алгебра является квазигруппой и её таблица Кэли - латинский квадрат.
%		\end{example}
%		
%		\begin{table}
%			\begin{tabular}{|l|l|l|l|l|} \hline
%				\vdot & 1 & -1 & $i$ & $-i$  \\ \hline
%				1 & 1 & -1 & $i$ & $-i$  \\ \hline
%				-1 & -1 & 1 & $-i$ & $i$  \\ \hline
%				$i$ & $i$ & $-i$ & -1 & 1  \\ \hline
%				$-i$ & $-i$ & $i$ & 1 & -1  \\ \hline
%			\end{tabular}
%		\end{table}
%		
%		%\begin{remark} Данное условие является необходимым. но не достаточным. Так, например множество $A = \lbrace a, b, c, d\rbrace$ с заданной на нём ниже приведённой таблицей Кэли операцией $*$ квазигруппой не является (например, тут нет нейтрального элемента).
%		%\end{remark}
%		
%		\vspace{3mm}
%		
%		%\begin{minipage}[b]{0.50\columnwidth}
%		%\begin{center}
%		%\begin{table}
%		%\begin{tabular}{|l|l|l|l|l|} \hline
%		%\vdot & 1 & -1 & $i$ & $-i$  \\ \hline
%		%1 & 1 & -1 & $i$ & $-i$  \\ \hline
%		%-1 & -1 & 1 & $-i$ & $i$  \\ \hline
%		%$i$ & $i$ & $-i$ & -1 & 1  \\ \hline
%		%$-i$ & $-i$ & $i$ & 1 & -1  \\ \hline
%		%\end{tabular}
%		%\end{table}
%		%\end{center}
%		%\end{minipage}
%		%\hfill
%		%\begin{minipage}[b]{0.50\columnwidth}
%		%\begin{center}
%		%\begin{table}
%		%\centering
%		%\begin{tabular}{|l|l|l|l|l|} \hline
%		%* & $a$ & $b$ & $c$ & $d$  \\ \hline
%		%$a$ & $a$ & $c$ & $d$ & $b$  \\ \hline
%		%$b$ & $d$ & $b$ & $a$ & $c$  \\ \hline
%		%$c$ & $b$ & $d$ & $c$ & $a$  \\ \hline
%		%$d$ & $c$ & $a$ & $b$ & $d$  \\ \hline
%		%\end{tabular}
%		%\end{table}
%		%\end{center}
%		%\end{minipage}
%		
%	\end{frame}
%	

\begin{frame}
\label{2.1.6.}\subsubsection{2.1.6. Гомоморфизмы и изоморфизмы алгебр }\frametitle{2.1.6. Гомоморфизмы и изоморфизмы алгебр (1/7)}
Понятие гомоморфизма является одним из ключевых
понятий алгебры. Каждая алгебраическая структура определённым
образом выделяет класс <<разумных>> отображений между объектами с
данной структурой, согласованных с операция\-ми этой структуры.
Поскольку операция~--- частный случай отношения, 
согласованность означает выполнение основного свойства гомоморфизма  
(\DMref{п.~1.7.6}{1.7.6.}).

\smallskip
Применительно к алгебраическим структурам 
гомоморфизм определяется для алгебр одного типа. 
Пусть 
${\cal A} = \Alg{A}{\Tuple{\varphi}{1}{m}}$ и 
${\cal B}=\Alg{B}{\Tuple{\psi}{1}{m}}$~--- две алгебры одного типа. 
Если функция $\Map{f}{A}{B}$ является гомоморфизмом для всех пар $\varphi_i$ и $\psi_i$, то есть

\vspace{-10pt}
$$
\A{i\in 1..m}{f(\varphi_i(\Tuple{a}{1}{n_i}))=\psi_i(f(a_1),\dots,f(a_{n_i}))},
$$

то говорят, что $f$~--- \odmkey{гомоморфизм}{Гомоморфизм
(homomorphism)} из ${\cal A}$ в ${\cal B}$.

\smallskip
\begin{examples}
\begin{enumerate}
\item
Пусть ${\cal A}=\Alg{\Bbb N}{+}$, ${\cal B}=\Alg{\Bbb
N_{10}}{+_{10}}$, где ${\Bbb N_{10}}=\{0,1,2,3,4,5,6,7,8,9\}$, а
$+_{10}$~--- сложение по модулю $10$. Тогда $f: a\mapsto (a \mod
10)$~--- гомоморфизм из ${\cal A}$~в~${\cal B}$.
\item
Пусть ${\cal A}=\Alg{\Bbb N}{\cdot}$, ${\cal B}=\Alg{\Bbb
N^2}{\cdot}$, где ${\Bbb N^2}=\{1,4,9, \ldots\}$, а
$\cdot$~--- умножение. \\ Тогда $f: a\mapsto a^2$~--- гомоморфизм из ${\cal A}$~в~${\cal B}$.
\end{enumerate}
\end{examples}

\end{frame}


\begin{frame}
\frametitle{ 2.1.6. Гомоморфизмы и изоморфизмы алгебр (2/7)}

Биективный гомоморфизм
называется \odmkey{изоморфизмом}{Изоморфизм (isomorphism)}, \DMref{п.~1.7.6}{1.7.6.}.

\begin{examples}
\begin{enumerate}
\item Биекция $n\mapsto 2n$ является изоморфизмом между
множеством натуральных чисел и множеством чётных чисел с операцией сложения.
\item Пусть $|M|=n$, 
${\cal A}=\Alg{2^M}{\cap,\cup}$, ${\cal B}=\Alg{2^M}{\cup^*, \cup^*}$, 
где $A\cup^* B \Assign \overline{A}\cap\overline{B}$,\\
$A\cap^* B \Assign \overline{A}\cup\overline{B}$. 
Тогда $f: A\mapsto \overline{A}$~--- 
изоморфизм между ${\cal A}$~и~${\cal B}$.
\end{enumerate}
\end{examples}

\begin{remark}
Сложение является операцией на множестве чётных чисел
(сумма чётных чисел~--- чётное число), но
не является операцией на множестве нечётных чисел.
\end{remark}

\smallskip
Гомоморфизмы, обладающие дополнительными свойствами, имеют
специальные названия.
Инъективный гомоморфизм называется
\odmkey{мономорфизмом}{Мономорфизм (monomorphism)}. 
Сюръективный гомоморфизм называется
\odmkey{эпиморфизмом}{Эпиморфизм (epimorphism)}.
Другими словами, изоморфизм является
 одновременно мономорфизмом и эпиморфизмом.
Если $A=B$, то гомоморфизм называется \odmkey
{эндоморфизмом}{Эндоморфизм (endomorphism)},
а изоморфизм называется \odmkey{автоморфизмом}{Автоморфизм (automorphism)}.
Общее соотношение между различными видами гомоморфизмов представлено на рисунке на следующем слайде.

\end{frame}

\begin{frame}
\frametitle{2.1.6. Гомоморфизмы и изоморфизмы алгебр (3/7)}
\begin{center}
\includegraphics[width=110mm]{2_1_5_1_.jpg}
\end{center}

\begin{remark}
Прочерки на рисунке означают, что таким гомоморфизмам
не дают специальных названий.
\end{remark}

\begin{examples}
\NB{1.} Гомоморфизм $n\mapsto 2n$, сохраняющий отношение порядка, является  мономорфизмом и эпиморфизмом, следовательно, и изоморфизмом.

\NB{2.} Пусть $A^{+}$~--- множество непустых слов в алфавите $A$
и  $\vert\alpha\vert$~--- длина слова  $\alpha\in A^{+}$. 
Длина  является гомоморфизмом относительно 
операции конкатенации на множестве $A^{+}$ и операции $+$ 
на множестве $\Bbb N$:
$\A{\alpha, \beta \in A^+}{|\alpha\beta| = |\alpha| + |\beta|}$. Этот гомоморфизм является эпиморфизмом.
\end{examples}


\end{frame}
\begin{frame}
\frametitle{2.1.6. Гомоморфизмы и изоморфизмы алгебр (4/7)}
Пусть ${\cal A} = \Alg{A}{\varphi}$, ${\cal B} = \Alg{B}{\psi}$,
тип~$=(1)$ и $\Map{f}{A}{B}$. Действие функций $\varphi,\psi,f$ можно
изобразить с помощью следующей диаграммы:

\vspace{-12pt}
\begin{figure}[h]
\begin{center}
\includegraphics[height=30mm]{2_1_5_2.jpg}
\end{center}
\end{figure}

\vspace{-4pt}
Пусть $f$~--- гомоморфизм. Тогда,
если взять конкретное значение
$a\in A$ и двигаться по двум различным путям на диаграмме, то
получится один и тот же элемент $b\in B$ (так~как
$f(\varphi(a))=\psi(f(a))$). В таком случае диаграмма называется
\odmkey{коммутативной}{Диаграмма (diagram)!коммутативная (commutative)}.
 Коммутатив\-ной
диаграмма называется потому, что условие
гомоморфизма можно переписать так:
$
f \circ \varphi  = \psi \circ f,
$
где $\circ$~--- суперпозиция функций.
Коммутативные диаграммы широко используются
в абстрактной алгебре и теории категорий.

\end{frame}

\begin{frame}
\frametitle{2.1.6. Гомоморфизмы и изоморфизмы алгебр (5/7)}

Если $\Map{f}{A}{B}$~--- изоморфизм, то алгебры ${\cal A}$ и
${\cal B}$ называют \odmkey{изоморфными}{Изоморфизм
(isomorphism)!алгебр (of algebras)} и обозначают так: ${\cal
A}\ToS{f}{\cal B}$. Если $f$ ясно из контекста или просто неважно
в конкретном рассуждении, то пишут $\cal A \sim \cal B$.

\medskip
Из теорем \DMref{п.~1.7.6}{1.7.6.} непосредственно следует, 
что отношение изоморфизма на множестве однотипных алгебр является
эквивалентностью.

\begin{remark}
Напомним, что изоморфизм алгебр  уважает \em{все} операции.
\end{remark} 

\begin{examples}
\begin{enumerate}
\item Пусть ${\cal A}=\Alg{\Bbb N}{+}$,
${\cal B}=\Alg{\Set{n}{n=2k,k\in \Bbb N}}{+}$~---
\odmkey{чётные}{Число (number)!чётное (even)} числа.
Тогда ${\cal A}\ToS{\times
2}{\cal B}$.
\item ${\cal A}=\Alg{2^M}{\cap,\cup}\ToS{f}{\cal B}=\Alg{2^M}{\cup,\cap}$, где
$f (X)=\overline X$.
\item ${\cal A}=\Alg{{\Bbb R}_+}{\cdot}\ToS{f}{\cal B}=\Alg{\Bbb R}{+}$,
где $f(x) =\log_a x$.
\item
Пусть $F$ - множество характеристических функций всех подмножеств универсума $U$. Тогда алгебра 
$\Alg{2^U}{\cup, \cap}$ изоморфна алгебре
 $\Alg{F}{ \max, \min}$ (\DMref{п.~1.9.2}{1.9.2.})
\end{enumerate}
\end{examples}

\end{frame}

\begin{frame}
\frametitle{2.1.6. Гомоморфизмы и изоморфизмы алгебр (6/7)}

Понятие изоморфизма является одним из центральных понятий,
оправдывающих применимость алгебраических методов в различных
областях. 

\smallskip
Действительно, пусть $\cal A =\Alg{A}{\Sigma_\varphi}$,
$\cal B=\Alg{B}{\Sigma_\psi}$ и $\cal A\ToS{f}\cal B$. Пусть в
алгебре $\cal A$ установлено свойство $\Phi_1=\Phi_2$, где
$\Phi_1$ и $\Phi_2$~--- некоторые формулы в сигнатуре
$\Sigma_\varphi$. Поскольку алгебры $\cal A$ и $\cal B$ изоморфны,
отсюда немедленно следует, что в алгебре $\cal B$ справедливо
свойство $\Psi_1=\Psi_2$, где $\Psi_1$ и $\Psi_2$~--- формулы,
полученные из формул $\Phi_1$ и $\Phi_2$ заменой операций из
сигнатуры $\Sigma_\varphi$ соответствующими операциями из
сигнатуры $\Sigma_\psi$. Таким образом, достаточно установить
некоторое свойство в одной алгебре, и оно автоматически
распространяется на все изоморфные алгебры.

\smallskip 
Алгебраические
структуры принято рассматривать \em{с точностью до изоморфизма}.

\smallskip 
Понятие изоморфизма применяется не только в алгебре, но и
практически во всех областях математики. Например, в этом курсе
рассматривается изоморфизм множеств 
(\DMref{п.~1.2.2}{1.2.2.}), 
линейно
упорядоченных множеств 
(\DMref{п.~1.8.8}{1.8.8.}), 
графов
 (\DMref{п.~7.1.6}{7.1.6.}).
\end{frame}

\begin{frame}
\frametitle{2.1.6. Гомоморфизмы и изоморфизмы алгебр (7/7)}


\begin{digress}
Термин <<морфизм>>, вынесенный в название раздела,
используется как собирательное понятие, подобно тому,
как в объектно-ориентированном программировании
используется термин <<абстрактный класс>>.
Напомним, что абстрактный класс~---  это класс,
который не имеет непосредственных экземпляров,
но является обобщением (предком) некоторых конкретных классов,
которые могут иметь непосредственные экземпляры.
Понятие морфизма является одним из ключевых понятий математики.
 
В этом разделе дано определение нескольких подклассов морфизмов, 
применимое к алгебраическим структурам, то есть множествам с операциями. 
В \DMref{п.~1.7.6}{1.7.6.} рассмотрены подклассы морфизмов, применимых к множествам с отношениями.
Более общее определение морфизмов, 
пригодное для самых разных математических объектов 
(множеств, графов, топологических пространств и т.д.), 
даёт 
\odmkey{теория категорий}{Теория (theory)!категорий (of category)}. 
\end{digress}

\end{frame}

\begin{frame}
\frametitle{Онтология общей алгебры}
\begin{center}
\includegraphics[height=7.8cm]{ODM2_1.png}
\end{center}
\end{frame}


\begin{frame}
\frametitle{2.2. Алгебры с одной операцией (1/3)}
\subsection{2.2. Алгебры с одной операцией}
\label{2.2.}
%\label{sect2-3}
Естественно начать изучение алгебраических структур с наиболее
простых. Самой простой структурой является алгебра с одной унарной
операцией, или \odmkey{унар}{Унар (unary algebra)}. 
\begin{example}
Алгебра $\Alg{\Bbb{N}}{\plusplus}$ (см. \DMref{п.~2.1.3.}{2.1.3.}) является унаром.
\end{example}
%\medskip
%\TODO{ Надо рассмотреть для полноты картины}
%
Здесь этот случай не рассматривается подробно, хотя он достаточно содержателен. 
\begin{example}
Любая тотальная функция на конечном множестве
образует унар, см. \DMref{п.~1.6.5}{1.6.5.}.
\end{example}
\medskip
Следующим по порядку сложности является случай
алгебры с одной бинарной операцией
$\Map{\ast}{M\times M}{M}$, для которой не 
предполагается никаких специальных свойств. Алгебра с такой операцией называется \odmkey{группоид}{Группоид (groupoid)} или
\odmkey{магма}{Магма (magma)}. Однако их общность не позволяет развить сколько-нибудь содержательную теорию, поэтому здесь рассматриваются только группоиды 
со специальными свойствами операции. 
Таким группоидам дают индивидуальные названия в соответствии 
с исторической традицией.

\medskip
В рассматриваемых алгебрах операция обычно обозначается знаком $*$ и 
называется \\ <<ум\-ножение>>, но это не более
чем упрощение речи. Выполнение свойств 
арифметического умножения <<по умолчанию>>
не предполагается. 
\end{frame}
\begin{frame}
\frametitle{2.2. Алгебры с одной операцией  (2/3)}
\raggedbottom
\begin{tabular}{p{12mm}|p{45mm}|p{86mm}}
  № & Название параграфа
  & Ключевые термины и обозначения \\
  \hline
  & & \\
 \DMref{2.2.1.}{2.2.1.}\RS
 & Полугруппы и квазигруппы &
{циклическая полугруппа } 
\\
{\DMref{2.2.2.}{2.2.2.}} & Определяющие соотношения & \\  
{\DMref{2.2.3.}{2.2.3.}} \BS & Системы подстановок термов &
{отношение редукции $\leadsto$; терминирующая система; редуцируемый терм; конфлюэнтная система;
нормальная форма; свойство Чёрча--Россера
}\\ 
{\DMref{2.2.4.}{2.2.4.}} \BS & Алгоритм Кнута--Бендикса &
{}\\  
\end{tabular}  

\end{frame}
\begin{frame}
\frametitle{2.2. Алгебры с одной операцией  (3/3)}
\raggedbottom
\begin{tabular}{p{12mm}|p{45mm}|p{86mm}}
  № & Название параграфа
  & Ключевые термины и обозначения \\
  \hline
  & & \\
{\DMref{2.2.5.}{2.2.5.}}\RS & Моноиды и лупы &
{подстановка, замена, преобразование над множеством}\\
%\TODO{ и лупы ?}  }\\
\DMref{2.2.6.}{2.2.6.} \RS  & Группы &
{ порядок группы; абелева группа; }\\ 
\DMref{2.2.7.}{2.2.7.}   & Действие группы на множестве &
{ орбита}\\
%\TODO{ идеал} } \\
\DMref{2.2.8.}{2.2.8.}   & Группа перестановок &
{ симметрическая группа; перестановка; подстановка } 
\\
\DMref{2.2.9.}{2.2.9.} \BS  & Перестановочные матрицы &
{ скалярное произведение } 
\end{tabular}  

\end{frame}

\begin{frame}
\label{2.2.1.}\subsubsection{2.2.1. Полугруппы и квазигруппы}
\frametitle{2.2.1. Полугруппы и квазигруппы (1/2)}

%\odmkey{Полугруппа}{Полугруппа (semigroup)}~--- это алгебра с одной
%ассоциативной бинарной операцией
%$
%a \ast (b \ast c) = (a \ast b) \ast c.
%$

%\vspace{-2pt}
Если для операции $\Map{\ast}{M\times M}{M}$ выполнено свойство 
{\color{blue}(A1)} (то есть она является ассоциативной), то $\Alg{M}{\ast}$ называется 
\odmkey{полугруппой}{Полугруппа (semi-group)}. 

\begin{examples}
\begin{enumerate}
\item Множество непустых слов $A^+$ в алфавите $A$ образует
полугруппу от\-но\-си\-тель\-но операции \odmkey{конкатенации}{Операция
(operation)!конкатенации (of concatenation)}
(см. \DMref{п.~1.1.5}{1.1.5.}).
\item
 Всякое множество тотальных функций
одного аргумента $\Set{f}{\Map{f}{M}{M}}$, 
зам\-кнутое
относительно суперпозиции, является полугруппой.
\end{enumerate}
\end{examples}

Если в полугруппе существует система образующих, состоящая из
одного элемента, то такая полугруппа называется \odmkey
{циклической}{Полугруппа (semi-group)!циклическая (cyclic)}.

\begin{example}
$\Alg{\Bbb N}{+}$ является циклической полугруппой, поскольку
$\{1\}$ является системой образую\-щих.
\end{example}
Пусть $P\! =\! \Alg{M}{\ast}$~--- полугруппа с конечной системой
образующих $A$, $A\subset M$,\\ $A = \{\Tuple{a}{1}{n}\}$. Тогда
$
\A{x\in M}{\E{\Tuple{y}{1}{k}\in A}{x=y_1\ast\ldots\ast y_k}}.
$
Если опустить обозначение операции $\ast$, то всякий элемент $a\in
M$ можно представить как слово $\alpha$ в алфавите~$A$, то~есть существует функция
$\Map{f}{A^+}{M}$. 
\end{frame}

\begin{frame}
\frametitle{2.2.1. Полугруппы и квазигруппы (2/2)}
Свойство ассоциативности очень широко распространено,
поэтому полугруппы являются основой для построения 
более сложных 
алгебр с одной бинарной операцией, в частности, 
моноидов и групп (см. \DMref{п.~2.2.6}{2.2.6.},
\DMref{п.~2.2.5}{2.2.5.}).
Однако существуют возможность опеределить
более сложные алгебры, начиная не с ассоциативности, 
а с других свойств. 

\smallskip   
Если для операции $\Map{\ast}{M\times M}{M}$ выполнено свойство 
{\color{blue}(A9)} (то есть она является разрешимой),
 то $\Alg{M}{\ast}$ называется 
\odmkey{квазигруппой}{Квазигруппа (quasi-group)}.

\begin{example}
Множество целых чисел 
с операцией вычитания 
образует квазигруппу $\Alg{\Bbb{Z}}{-}$,
поскольку вычитание не ассоциативно
($0=(3-2)-1 \ne 3-(2-1)=2$) и в то же время уравнения
$a-x=b$ и $x-a=b$ легко разрешимы. 
\end{example} 

Квазигруппа $\Alg{M}{\ast}$ называется
\odmkey{полностью антисимметричной}{Квазигруппа (quasi-group)!полностью антисимметричная (totally anti-symmetric)}, если
\begin{itemize}
\item $a\ast b = b\ast a\Impl a=b$ и
\item $(a\ast  b)\ast c = (a \ast c)\ast b\Impl b=c$.
\end{itemize}

\begin{digress}
Построение полностью антисимметричных квазигрупп
считается значительным достижением математики XXI века,
поскольку с помощью этих квазигрупп 
возможно эффективное кодирование с исправлением ошибок.
\end{digress}
\end{frame}

\begin{frame}
\label{2.2.2.}\subsubsection{2.2.2. Определяющие соотношения }\frametitle{2.2.2. Определяющие соотношения (1/5)}
\begin{remark}
Функция, которая сопоставляет слову в 
алфавите образующих элемент носителя
полугруппы, обязательно сюръективна, но не обязательно инъективна!
\end{remark}

\smallskip
Пусть задана полугруппа с конечной системой образующих $A$.
Обозначим через $\bar\alpha$ элемент, определяемый
словом $\alpha$. Слова $\alpha$ и $\beta$ могут быть различными,
но соответствую\-щие им элементы $\bar \alpha$ и $\bar \beta$
могут быть равны: $\bar \alpha = \bar \beta$, $\alpha, \beta \in
A^+$, $\bar \alpha, \bar \beta \in M$. Такие равенства называются
\odmkey{определяющими соотношениями}{Соотношение (relationship)!определяющее (defining)}. 
Если в полугруппе нет
определяющих соотношений, и все различные слова, составленные из
образующих, суть различные элементы носителя, то полугруппа
называется \odmkey{свободной}{Полугруппа (semi-group)!свободная
(free)}. 

\smallskip
Всякая конечно-порождённая полугруппа может быть получена
из свободной введением определяющих соотношений. 

\smallskip
Пусть в
конечно-порождённой полугруппе некоторый элемент $\bar\alpha$
определяется словом $\alpha$, причём $\alpha=\beta\gamma\delta$ и
задано определяющее соотношение $\gamma=\sigma$. Тогда вхождение
слова $\gamma$ в слово $\alpha$ можно заменить словом $\sigma$,
причём полученное слово $\beta\sigma\delta$ будет определять тот
же самый элемент $\bar\alpha$. Такое преобразование слова будем
называть \em{применением} определяющего соотношения.  

\end{frame}

\begin{frame}
\frametitle{2.2.2. Определяющие соотношения (2/5)}
Два слова в
алфавите $A$ считаются равными, если одно из другого получается
применением определяющих соотношений, то есть слова равны, если
равны определяемые ими элементы. Отношение равенства слов в
полугруппе с определяющими соотношениями является отношением
эквивалентности. Классы эквивалентности по этому отношению
соответствуют элементам полугруппы.
%\addvspace{-3pt}
\begin{examples}
%\addvspace{-3pt}
\begin{enumerate}
\item 
В полугруппе $\Alg{\Bbb N}{+}$ имеется конечная система
образующих $\{1\}$. Другими словами, каждое натуральное число
можно представить как последовательность знаков 1. Очевидно, что
различные слова в алфавите $\{1\}$ суть различные элементы
носителя, то~есть эта полугруппа свободна.
\item 
Пусть полугруппа
$\mathcal{P}$ задана системой двух образующих $\{a,b\}$ и~\hbox{двумя}
опре\-де\-ляю\-щими соотношениями: $aa=a$ и $bb=b$. \hbox{Следующий}
эле\-мен\-тар\-ный алгоритм определяет, равны ли два слова, $\alpha$ и
$\beta$,~в~полугруп\-пе~$\mathcal{P}$.
\end{enumerate}
\end{examples}


\end{frame}

\begin{frame}
\frametitle{2.2.2. Определяющие соотношения (3/5)}
%\vspace{-6pt}
\begin{algorithmic}
\REQUIRE  слова \textbf{$\alpha:$ array $[1..s]$ of $\{a,b\}$} и
\textbf{$\beta:$ array $[1..t]$ of $\{a,b\}$};
\ENSURE  значение выражения $\alpha=\beta$ в полугруппе $\mathcal{P}$.
\STATE \textbf{ if $\alpha[1]\neq\beta[1]$ then
  return false end if}
  \COMMENT{\small\color{gray} первые буквы не совпадают }
\STATE $N_\alpha:=0$
\COMMENT{\small\color{gray}  счётчик изменений буквы в $\alpha$ }
\FOR{\textbf{$i$ from $2$ to $s$}}
\IF{$\alpha[i]\neq\alpha[i-1]$}
  \STATE $N_\alpha:=N_\alpha+1$
  \COMMENT{\small\color{gray}  буква изменилась }
\ENDIF
\ENDFOR
\STATE $N_\beta:=0$
\COMMENT{\small\color{gray}  счётчик изменений буквы в $\beta$ }
\FOR{\textbf{$i$ from $2$ to $t$}}
\IF{$\beta[i]\neq\beta[i-1]$}
  \STATE $N_\beta:=N_\beta+1$
  \COMMENT{\small\color{gray}  буква изменилась }
\ENDIF
\ENDFOR
\STATE \textbf{return $N_\alpha=N_\beta$}
\COMMENT{\small\color{gray}  слова равны, если значения счётчиков одинаковы }
\end{algorithmic}



%\odmkey{}{Теорема (theorem)!Маркова--Поста
%(Markov--Post)}


\end{frame}

\begin{frame}
\frametitle{2.2.2. Определяющие соотношения (4/5)}
В предыдущем примере определяющие соотношения таковы, что
равенство любых слов как элементов легко проверяется. Однако это
не всегда так.

\medskip
\begin{theorem} \NB{[Марков--Пост]}
%\addtocounter{footnote}{-1} \footnotetext{\;Андрей Андреевич
%Марков (1903--1979).} \addtocounter{footnote}{1}
%\footnotetext{\;Эмиль Леон Пост (1897--1954).} 
Существует
полугруппа, в которой проблема распознавания равенства слов
алгоритмически неразрешима.
\end{theorem}
\begin{proof}
Без доказательства.
\end{proof}

%\addvspace{-3pt}
\begin{remark}
Термин <<алгоритмическая неразрешимость>> принадлежит
теории алгоритмов, которая  рассматривается в этом курсе
в главе 4,
а потому строгое определение этого термина \em{здесь} не приводится.
Неформально можно сказать, что
проблема называется
\odmkey{алгоритмически неразрешимой}{Проблема (problem)!алгоритмически
неразрешимая \\  (algorithmically unsolvable)},
если  не существует программы (алгоритма) её решения (см. \DMref{п.~4.5.7}{4.5.7.}).
Здесь словосочетание <<не существует>> означает не то, что
программа ещё не составлена,
 а то, что требуемая программа не может быть составлена в принципе.
%В теории алгоритмов строго установлено, что алгоритмически
%неразрешимые проблемы существуют. Приведём простое наблюдение,
%которое не является строгим доказательством последнего
%утверждения, но может служить косвенным намёком на причины
%существования алгоритмически неразрешимых проблем. Пусть в
%конечном алфавите $A$ задано слово $\alpha$, $\alpha\in A^*$ и
%язык $L$, $L\subset A^*$. Проблема: определить, принадлежит ли
%слово $\alpha$ языку $L$? Обозначим через $P_L(\alpha)$ программу,
%которая для заданного слова $\alpha\in A^*$ выдаёт 1, если
%$\alpha\in L$,  и выдаёт 0 в противном случае. Множество всех
%программ не более чем счётно (просто потому, что всякая
%программа~--- это слово в конечном алфавите). В то же время
%множество всех языков несчётно, потому что это множество всех
%подмножеств счётного множества. Таким образом, заведомо существуют
%языки, для которых не существует программы, распознающей все слова
%этого языка. Приведённое наблюдение не является доказательством,
%потому что мы не определили, как именно задаются объекты $\alpha$,
%$L$ и $P_L(\alpha)$.
\end{remark}
\end{frame}

\begin{frame}
\frametitle{2.2.2. Определяющие соотношения (5/5)}
\begin{example}
Григорий Самуилович~Цейтин
нашёл легко описываемый пример полугруппы, в которой проблема
равенства слов алгоритмически неразрешима. 

В полугруппе c 
пятью образующими $\{a,b,c,d,e\}$ и семью определяющими соотношениями:
$ac =ca$, 
$ad =da$, 
$bc =cb$, 
$bd =db$, 
$abac =abace$, 
$eca =ae$,
$edb =be$~---   
невозможно составить программу, которая
бы для любых заданных слов в алфавите $\{a,b,c,d,e\}$ проверяла,
можно ли преобразовать одно слово в~другое применением указанных
определяющих соотношений.
\end{example}

\smallskip
{\large
Некоторые программисты полагают, что если в условиях задачи всё
дискретно и конечно, то для решения такой задачи программу можно
составить в любом случае (используя метод <<полного перебора>>).

Предшествующие теорема и пример показывают, что это мнение
ошибочно.
}
\end{frame}

%_{•}
%\begin{frame}
%\frametitle{Математические знаки (2/3)}
%$\leftrightarrow$ leftrightarrow
%$\leq$ leq (relation)
%$\lfloor$ lfloor (delimiter)
%$\lhd$ (left-pointing arrow head)
%$\ll$ ll (relation)
%$\lnot$ lnot
%$\longleftarrow$ longleftarrow
%$\longleftrightarrow$ longleftrightarrow
%$\longmapsto$ longmapsto
%$\longrightarrow$ longrightarrow
%$\lor$ lor
%$\mapsto$ mapsto
%$\mho$ mho 
%$\mid$ mid (relation)
%$\models$ models (relation)
%$\mp$ mp (binary operation)
%$\mu$ mu
%$\nabla$ nabla
%$\natural$ natural
%$\ne$ ne
%$\nearrow$ nearrow
%$\neg$ neg
%$\neq$ neq (relation)
%$\ni$ ni (relation)
%$\not$
%Overstrike a following operator with a /, as in $\not=$.
%$\notin$ $\ni$ $\nu$ nu
%$\nwarrow$ nwarrow
%$\odot$ odot (binary operation)
%$\oint$ oint
%$\Omega$ Omega $\omega$ omega
%$\ominus$ ominus (binary operation)
%$\oplus$ oplus (binary operation)
%$\oslash$ oslash (binary operation)
%$\otimes$ otimes (binary operation)
%$\owns$ owns
%$\parallel$ parallel (relation)
%$\partial$ partial
%$\perp$ perp (relation)
%$\phi$ phi
%$\Pi$ Pi $\pi$ pi
%$\pm$ pm (binary operation)
%$\prec$ prec (relation)
%$\preceq$ preceq (relation)
%$\prime$ prime
%$\prod$ prod
%$\propto$ propto (relation)
%$\Psi$ Psi $\psi$ psi
%$\rangle$ rangle (delimiter)
%$\rbrace$ rbrace (delimiter)
%$\rbrack$ rbrack (delimiter)
%$\rceil$ rceil (delimiter)
%$\Re$ Re
%$\rfloor$ rfloor
%$\rhd$ (binary operation)
%$\rho$ rho
%$\Rightarrow$ Rightarrow
%$\rightarrow$ rightarrow
%$\rightharpoondown$ rightharpoondown
%$\rightharpoonup$ rightharpoonup
%$\rightleftharpoons$ rightleftharpoons
%$\searrow$ searrow
%$\setminus$ setminus (binary operation)
%$\sharp$ sharp
%$\Sigma$ Sigma $\sigma$ sigma
%$\sim$ sim (relation) $\simeq$ simeq (relation)
%$\smallint$ smallint
%$\smile$ smile (relation)
%$\spadesuit$ spadesuit
%$\sqcap$ sqcap (binary operation)
%$\sqcup$ sqcup (binary operation)
%$\sqsubset$ (relation)
%$\sqsubseteq$ sqsubseteq (relation)
%$\sqsupset$ (relation)
%$\sqsupseteq$ sqsupseteq (relation)
%$\star$ star (binary operation)
%$\subset$ subset (relation)
%$\subseteq$ subseteq (relation)
%$\succ$ succ (relation)
%\end{frame}
%
%\begin{frame}
%\frametitle{Математические знаки (3/3)}
%$\succeq$ succeq (relation)
%$\sum$ sum
%$\supset$ supset (relation)
%$\supseteq$ supseteq (relation)
%$\surd$ surd
%$\swarrow$ swarrow
%$\tau$ tau
%$\theta$ theta
%$\times$ times (binary operation)
%$\to$ to
%$\top$ top
%$\triangle$ triangle
%$\triangleleft$ triangleleft (binary operation)
%$\triangleright$ triangleright (binary operation)
%$\unlhd$ 
%left-pointing arrowhead with line under (binary operation)
%$\unrhd$
%right-pointing arrowhead with line under (binary operation)
%$\Uparrow$ Uparrow (delimiter)
%$\uparrow$ uparrow (delimiter)
%$\Updownarrow$ Updownarrow (delimiter)
%$\updownarrow$ updownarrow (delimiter)
%$\uplus$ uplus (binary operation)
%$\Upsilon$ Upsilon
%$\upsilon$ upsilon
%$\varepsilon$ varepsilon
%$\varphi$ varphi
%$\varpi$ varpi
%$\varrho$ varrho
%$\varsigma$ varsigma
%$\vartheta$ vartheta
%$\vdash$ vdash (relation)
%$\vee$ vee (binary operation)
%$\Vert$ Vert (delimiter)
%$\vert$ vert (delimiter)
%$\wedge$ wedge (binary operation)
%$\wp$ wp $\wr$ wr (binary operation)
%$\Xi$ Xi $\xi$ xi
%$\zeta$ zeta
%\end{frame}

\begin{frame}

\label{2.2.3.}\subsubsection{2.2.3. Системы подстановок термов }\frametitle{2.2.3. Системы подстановок термов (1/4)}

Слова в полугруппе~--- это выражения, 
построенные с помощью одной ассоциативной 
бинарной операции. Как указано в предыдущем параграфе, 
зная правила преобразований выражений (определяющие соотношения)
даже в таком простом случае, как полугруппа, не всегда 
возможно решить проблему равенства выражений. 
Тем более трудной является проверка равенства более сложных выражений. 
В данном параграфе рассматриваются 
\em{системы подстановок термов}, 
применение которых позволяет в некоторых случаях это сделать.

\smallskip
\odmkey{Системой подстановок термов}{Система (system)!подстановок термов (of term rewriting)} 
(или \odmkey{системой правил переписывания}
{Система (system)!правил переписывания (of rewriting rules)}) 
называется пара $\Alg{A}{\leadsto}$, 
где $A$~--- множество, элементы которого 
обычно называют \odmkey{термами}{Терм (term)}, 
а $\leadsto$~--- отношение на множестве $A$, 
 именуемое \odmkey{отношением переписывания}{Отношение (relation)!переписывания (of rewriting)},
или \odmkey{отношением редукции}{Отношение (relation)!редукции (of reduction)}.
Если $a\leadsto b$, 
то говорят, что терм $a$ переписывается в терм $b$ \odmkey{за один шаг}{Переписывание термов (term rewriting)!за один шаг (one-step)}.
Говорят, что терм $a$  \odmkey{переписывается}{Переписывание термов (term rewriting)} в терм $b$, 
если существует последовательность термов  $\Tuple{t}{0}{n}$, 
такая, что $a=t_0$, $t_n=b$ и $t_0\leadsto t_1 \leadsto \dots \leadsto t_n $.

\smallskip
\begin{remark}
Обычно считают, что терм переписывается в себя, и в качестве
отношения редукции рассматривают
рефлексивное транзитивное замыкание, обозначаемое 
$\Reduce$.
\end{remark}

\end{frame}

\begin{frame}
\frametitle{2.2.3. Системы подстановок термов (2/4)}


Говорят, что терм $x\in A$
может быть \odmkey{редуцирован}{Терм (term)!редуцируемый (reducable)}, если $\E{y \in A, y\ne x}{x \Reduce y}$. 
В противном случае терм $x$ называется \odmkey{не редуцируемым}{Терм (term)!не редуцируемый (non-reducable)}.
Система подстановок термов называется \odmkey{терминирующей}{Система (system)!подстановок термов (of term rewriting)!терминирующая (terminating)}
(или \odmkey{остановочной}{Система (system) !подстановок термов (of term rewriting)!остановочная (stopping)}), 
если не существует \em{бесконечной} последовательности термов $x_1, x_2, \dots$, 
такой, что $x_1 \leadsto x_2 \leadsto \dots$.

\begin{remark}
В терминирующей системе любой терм можно 
переписать в некоторый нередуцируемый терм.
\end{remark}

\begin{example}
В примере 2 из \DMref{п.~2.2.2}{2.2.2.} множество термов $A=\{a,b\}^*$ состоит из всех слов
в алфавите $\{a,b\}$, 
а система $aa\leadsto a$ и $bb\leadsto b$ является 
терминирующей.
 Нередуцируемыми являются все слова, в которых нет двух одинаковых букв подряд. 
\end{example}

\medskip
Утверждение о том, что если любой терм можно переписать в некоторый нередуцируемый, то система терминирующая~--- неверно.

\begin{example}
Рассмотрим систему $A = \{a, b, c\}, a \leadsto b, b \leadsto a, a \leadsto c, b \leadsto c$. Любой терм в такой системе можно переписать в терм $c$, который нередуцируем. 
Однако последовательность $a\leadsto b\leadsto a\leadsto b\leadsto\dots$, очевидно, бесконечна.
\end{example}

\end{frame}

\begin{frame}
\frametitle{2.2.3. Системы подстановок термов (3/4)}

Система подстановок термов называется \odmkey{конфлюэнтной}{Система (system)!подстановок термов (of term rewriting)!конфлюэнтная (confluent)}, 
если \\
$\A{ x\!\in\! A }{\E{ u, v\!\in\!A}{x\!\Reduce\!u \And x\!\Reduce\!v}\Impl
 \E {y\!\in\!A}{u\!\Reduce\!y \And v\!\Reduce\!y}}$.

\begin{example}
Система $A = \{a, b, c, d\}, a \leadsto b, b \leadsto a, a \leadsto c, b \leadsto d$
не конфлюэнтна, а система предыдущего примера~--- конфлюэнтна.
\end{example}

\smallskip
Терм $y$ называется \odmkey{нормальной формой}{Форма (form)!нормальная (normal)! терма (term)}
терма $x$, если $x\Reduce y$ и терм $y$ нередуцируемый.
В терминирующей системе все термы имеют нормальные формы,
однако не обязательно единственные.

\begin{example}
В системе 
$A = \{a, b, c, d\}, a \leadsto b, b \leadsto a, a \leadsto c, b \leadsto d$
оба терма $a$ и $b$ имеют по две нормальные формы $c$ и $d$.
\end{example}

\smallskip
Говорят, что термы $x$ и  $y$ имеют 
\odmkey{общую нормальную форму}{Форма (form)!общая нормальная (general normal)! терма (term)}\\ (и пишут $x \downarrow y$), 
если $\E{z\in A}{x \Reduce z \And y \Reduce z}$, причём терм
$z$ нередуцируемый.

Пусть $\stackrel{*}{\leftrightarrow}$~--- рефлексивное симметричное
транзитивное замыкание отношения $\leadsto$. 
Тогда система подстановок термов обладает 
\odmkey{свойством Чёрча-Россера}{Свойство (property)!Чёрча--Россера (Church--Rosser)}, 
если  $x \stackrel{*}{\leftrightarrow} y$ 
тогда и только тогда, когда $x \downarrow y$.



\end{frame}

\begin{frame}

\frametitle{2.2.3. Системы подстановок термов (4/4)}

В системе подстановок термов со свойством Чёрча-Россера все термы 
имеют единственную нормальную форму.
Другими словами, система подстановок термов 
обладает свойством Чёрча-Россера, если термы преобразуются друг в друга по правилам переписывания
тогда и только тогда, когда у них есть общая нормальная форма.   


Свойство Чёрча-Россера~--- очень сильное, редкое и полезное 
свойство. Действительно, в случае наличия этого свойства
для проверки равенства термов достаточно свести 
обе части проверяемого равенства к нормальной форме.
Если результаты совпадут, то равенство имеет место,
в противном случае не имеет места.
Однако проверка 
 свойства 
Чёрча-Россера является алгоритмически неразрешимой задачей.   

\begin{theorem}
Терминирующая система подстановок термов 
обладает свойством\\ Чёрча-Россера тогда и только тогда, 
когда она конфлюэнтна.
\end{theorem}
\begin{proof}
Без доказательства.
\end{proof}

Теорема имеет важное значение, поскольку конфлюэнтность системы в некотором смысле проще проверять, чем свойство Чёрча-Россера, 
и, что важнее, проще построить систему с таким свойством.

\end{frame}


\begin{frame}

\label{2.2.4.}\subsubsection{2.2.4. Алгоритм Кнута--Бендикса }\frametitle{2.2.4. Алгоритм Кнута--Бендикса (1/7)}

Известен \odmkey{алгоритм Кнута--Бендикса}{Алгоритм (algorithm)!Кнута--Бендикса (Knuth--Bendix)},
который по заданной системе подстановок термов (в случае успешного завершения!) 
строит эквивалентную терминирующую конфлюэнтную систему подстановок термов.
Для простоты алгоритм излагается на примере полугрупп 
с конечной системой образующих (см.~\DMref{п.~2.1.3}{2.1.3.}) и конечным набором определяющих соотношений (см.~\DMref{п.~2.2.2}{2.2.2.}). В такой постановке термами являются слова в 
алфавите образующих (см.~\DMref{п.~1.1.5}{1.1.5.}).

\smallskip
Пусть задан конечный алфавит $A$ (образующие) 
и конечный набор пар слов \\ $R=\Set{(\lambda,\rho)}{\lambda\in A^*\And\rho\in A^*}$
(определяющие соотношения). На словах 
определён линейный порядок (cм.~\DMref{п.~1.8.2}{1.8.2.},
не ограничивая общности этот порядок можно считать лексикографическим). 



\end{frame}

\begin{frame}

\frametitle{2.2.4. Алгоритм Кнута--Бендикса (2/7)}


Основная идея алгоритма вытекает из следующих наблюдений.

{\color{blue}\bf 1.} Если использовать определяющее соотношение
<<в обе стороны>> $\alpha\leadsto \beta $ и $\beta\leadsto \alpha$,
то система заведомо не будет терминирующей, а 
потому такое использование надобно запретить и <<ориентировать>> 
соотношения в одну сторону.   

{\color{blue}\bf 2.} 
Определяющее соотношение $\alpha = \beta $
следует ориентировать так, чтобы $|\alpha| > |\beta| $. 
Действительно, в таком случае в цепочке преобразований длина слова всё время уменьшается, 
а значит любая цепочка обрывается и система является терминирующей. 

{\color{blue}\bf 3.} Термы в обеих частях правил должны быть нередуцируемыми
с помощью других правил. В противном случае надо провести редукции 
и упростить исходные правила.

{\color{blue}\bf 4.} Если в системе есть тривиальные правила, заданные исходно или полученные в результате редукции, то их следует исключить.

{\color{blue}\bf 5.} Нарушение конфлюэнтности может происходить
только в случае наличия правил с пересекающимися левыми частями.
Значит, такие правила необходимо устранить, 
заменяя их эквивалентными правилами. 
\end{frame}

\begin{frame}

\frametitle{2.2.4. Алгоритм Кнута--Бендикса (3/7)}


Важной операцией в алгоритме является выделение пересекающихся пар слов.  
Пересече\-ние слов может быть двух видов: 
либо суффикс одного слова есть префикс другого (слова $\alpha\beta$ и $\beta\gamma$), 
либо одно слово~--- подстрока другого слова ($\alpha\beta\gamma$ и $\beta$).
В обоих случаях при пересечении однозначно 
выделяются три слова $\alpha, \beta, \gamma$,
которые не могут быть одновременно  пустыми.  


\begin{example}
В строках $abab$ и $baba$ имеем $\alpha=a$, $\beta=bab$ и $\gamma=a$. 
\end{example}

\smallskip 
Основной структурой данных является множество пар слов $R=\{(\lambda,\rho)\}$, 
то есть набор правил переписывания $\lambda\leadsto \rho$.
Вначале это исходные определяющие соотношения, а в конце,
возможно, конфлюэнтная система. Алгоритм может не закончить работу,
если исходная система не обладает свойством Чёрча--Россера.


\end{frame}

\begin{frame}

\frametitle{2.2.4. Алгоритм Кнута--Бендикса (4/7)}

Вспомогательная процедура $\Clear$ проводит <<чистку>>
множества пар слов $R$, а именно: применяет 
все возможные редукции к каждому правилу, и если правило редуцируется к 
тривиальному, то удаляет его.

\begin{example}
Если $R=\{aa\leadsto bb, a\leadsto b\}$, то после применения процедуры $\Clear$ 
имеем $R=\{a\leadsto b\}$.
\end{example}       

\medskip
Вспомогательная процедура $\Conf$ находит
в $R$ два правила, 
левые части которых пересекаются, и возвращает
соответствующие слова $\alpha, \beta, \gamma$, 
а также результаты редукции слова $\alpha\beta\gamma$
по первому правилу ($\rho_1$) и по второму правилу ($\rho_2$).


\begin{example}
$\text{Conf} (\{abab\leadsto c, baba\leadsto d\})=$\\
$=(\alpha=a, \beta=bab,\gamma=a,  \rho_1 =ca,
 \rho_2=ad)$.
\end{example}       


\end{frame}

\begin{frame}%[shrink=20] %чтобы алгоритм попал на один слайд
\frametitle{2.2.4. Алгоритм Кнута--Бендикса (5/7)}

%\vspace{-4pt}
\begin{algorithmic}
\REQUIRE $R$~--- множество определяющих соотношений.
\ENSURE $R$~--- множество пар слов, 
входящих в отношение переписывания искомой системы.
\FOR{ $(\lambda,\rho)\in R$}
  \STATE{\textbf{if $|\lambda|<|\rho|$ then $\lambda :=: \rho$ end if } }
  \STATE
   \COMMENT{\small\color{gray} меняем местами, так чтобы левая часть была больше правой}
\ENDFOR
\WHILE {\textbf{true}}
\STATE {$\Clear$}
\COMMENT{\small\color{gray} чистка текущего множества правил }
\STATE {$(\alpha, \beta,\gamma,  \rho_1,
 \rho_2)\Assign \Conf$}
%\STATE  
\COMMENT{\small\color{gray} ищем пару правил c пересекающимися левыми частями }
\STATE{\textbf{if $\alpha=\beta=\gamma=\varepsilon$ then stop end if}}
%\STATE
\COMMENT{\small\color{gray} нет конфликтующих правил --- система построена }
\IF {$\rho_1 \neq \rho_2$} 
\STATE $R \Assign R + (\max(\rho_1,\rho_2),\min(\rho_1,\rho_2))$ 
\COMMENT {\small\color{gray} добавляем новую пару}
\ENDIF
\ENDWHILE
\end{algorithmic}
\end{frame}

\begin{frame}
\frametitle{2.2.4. Алгоритм Кнута--Бендикса (6/7)}
%\vspace{-0.3\baselineskip}
\begin{ground}
Построенный алгоритм строит конфлюэнтную и терминирующую систему, что следует
из наблюдений, лежащих в основе алгоритма. Кроме того,  построенная система $R^\prime$  
является следствием исходного набора соотношений $R$. Действительно, возьмём любые
$\alpha$ и $\beta$ такие, что $\alpha\stackrel{*}{\leadsto}_R\beta$, и рассмотрим
последовательность определяющих соотношений (правил), 
при применении которых $\alpha\stackrel{*}{\leadsto}_R\beta$.
Если какое-то из правил не входит в $R^\prime$, то значит, 
оно было <<удалено>> в процедуре $\Clear$,
то есть при помощи применения некоторых редукций из $R^\prime$ данное правило сводится к тривиальному. Тогда
применение этого правила можно 
заменить на применение некоторого количества правил из $R^\prime$.
Таким образом, $\alpha\stackrel{*}{\leadsto}_{R^\prime}\beta$.
\end{ground}


\end{frame}

\begin{frame}
\frametitle{2.2.4. Алгоритм Кнута--Бендикса (7/7)}
\begin{example}
Пусть $X = \{a, b\}$, $R = \{aa = \varepsilon, ab = \varepsilon\}$.
В таблице отображены состояния множества $R$  после каждой итерации основного цикла в алгоритме и результаты выполнения процедур $\Clear$ и $\Conf$.

\begin{center}
\begin{tabular}{l|c|c}
\hline 
$R$ & $\Clear$ & $\Conf$ \\
\hline
$aa \leadsto \varepsilon, ab \leadsto \varepsilon$ & $\emptyset$ & $b\leadsto a$  \\ 
\hline
$aa \leadsto \varepsilon, ab \leadsto \varepsilon, b\leadsto a$ & 
$ab \leadsto \varepsilon$ & $\emptyset$  \\ 
\hline
$aa \leadsto \varepsilon, b \leadsto a$ & $\emptyset$  & $\emptyset$ \\
\hline
\end{tabular}
\end{center}

\end{example}




\begin{example}
Пусть $X = \{a, b\}$, $R = \{abab \leadsto ab, baba \leadsto ba\}$. 
При применении алгоритма на каждом шаге получаем тривиальные 
правила переписывания $aba \leadsto aba$ или $ababa \leadsto ababa$, и
 искомую конфлюэнтную систему подстановок термов 
не удаётся построить. 
\end{example}



%%%%%%%%%%%%%%%%%%%%%%%%%%%%
% Это замечание и пример внушают мне сомнения
% я их не проверял, и не видно что это пример из теории групп. ФАН 27.06.2015 
%%%%%%%%%%%%%%%%%%%%%%%%%%%%%
%Алгоритм Кнута--Бендикса можно применять
%не только к системам подстановок слов, 
%но и к системам подстановок более сложно устроенных термов.
%\begin{example}
%Рассмотрим произвольное множество $X$ и функции 
% $f: X \times X \rightarrow X$, $g: X \rightarrow X$
%такие, что:
%
%\vspace{-\baselineskip}
%$$f(1, x) = x, \quad f(g(x), x) = 1, \quad f(f(x,y), z) = f(x, f(y, z)).$$
%
%%\vspace{-0.5\baselineskip}
%После применения алгоритма Кнута-Бендикса получим следующие правила переписывания.
%
%\begin{tabular}{l|l||l|l}
%\hline 
%$R1$ &  $f(1, x)\leadsto x$ & $R6$ & $f(x,1)\leadsto x$ \\
%\hline 
%$R2$ & $f(g(x), x)\leadsto 1$ & $R7$ & $g(g(x))\leadsto x$ \\
%\hline 
%$R3$ & $f(f(x, y),z)\leadsto f(x,f(y,z))$ & $R8$ & $f(x, g(x))\leadsto 1$ \\
%\hline 
%$R4$ & $f(g(x),f(x,y))\leadsto y$ & $R9$ & $f(x, f(g(x),y))\leadsto y$ \\
%\hline 
%$R5$ & $g(1)\leadsto 1$ & $R10$ & $g(f(x,y))\leadsto f(g(y),g(x))$ \\
%\hline 
%\end{tabular}
%
%Установим равенство двух термов
%
%\vspace{-\baselineskip} 
%$$g(f(f(g(a),a),f(b,g(b))))= f(b,f(g(f(a,b)),a)),$$
%
%%\vspace{-0.5\baselineskip}
%приведя их к нормальной форме.
%
%$g(f(f(g(a),a),f(b,g(b)))) \leadsto_{R2} g(f(1,f(b,g(b)))) \leadsto_{R8} g(f(1,1)) \leadsto_{R1} g(1) \leadsto_{R5} 1$.
%
%$f(b,f(g(f(a,b)),a)) \leadsto_{R10} f(b,f(f(g(b), g(a)),a)) \leadsto_{R3} f(b,f(g(b), f(g(a),a))) \leadsto_{R2} f(b,f(g(b), 1)) \leadsto_{R6} f(b,g(b)) \leadsto_{R8} 1$.
%\end{example}

Таким образом, приведённый алгоритм не является универсальным.
Если алгоритм Кнута\-Бендикса успешно завершается, то 
в данной полугруппе можно решить проблему равенства слов.
Если же алгоритм не завершается,
то это ничего не означает.   

\begin{remark}
Применение описанного алгоритма к полугруппе 
из \DMref{п.~2.2.2}{2.2.2.}. 
не приводит к построению конфлюэнтной терминирующей
 системы подстановок термов.
\end{remark}

\end{frame}

\begin{frame}
\label{2.2.5.}\subsubsection{2.2.5. Моноиды и лупы }
\frametitle{2.2.5. Моноиды и лупы (1/3)}
%\vspace{-2pt}

%\odmkey{Моноид}{Моноид (monoid)}~--- это полугруппа с
%\odmkey{единицей}{Единица (identity)!моноида (of monoid)}:
%$
%\E{e}{\A{a}{a \ast e = e \ast a = a}}.
%$
%\begin{remark}
%Единицу также называют
%\odmkey{нейтральным элементом}{Элемент (element)!нейтральный (neutral)}.
%\end{remark}

Если для бинарной операции 
$\Map{\ast}{M\times M}{M}$ 
выполнены свойства {\color{blue}(A1)} и 
{\color{blue}(A6)}, то алгебра
 $\Alg{M}{\ast}$ называется \odmkey{моноидом}{Моноид (monoid)}.
 Иначе говоря, моноид~--- это полугруппа  
 с нейтральным элементом, который называют также 
 \odmkey{единицей}{Единица (identity)!моноида (of monoid)}.

\begin{examples}
\begin{enumerate}
\item Множество слов $A^*$ в алфавите $A$ вместе с операцией
конкатенации и пустым словом $\varepsilon$ образует моноид.
\item Пусть $T$~--- множество термов над множеством переменных $V$ и
сигнатурой~$\Sigma$. \\ \odmkey{Подстановкой}{Подстановка
(substitution)}, или \odmkey{заменой переменных}{Замена переменной
(substitution)}, называется множество пар
$
\sigma=[t_i/ v_i]_{i\in I},
$
где $t_i$~--- термы, а $v_i$~--- переменные. Результатом
применения подстановки $\sigma$ к терму $t$ 
(обозначается
$t\sigma$) 
называется терм, который получается заменой \em{всех}
вхождений переменных $v_i$ соответствующими термами $t_i$.
\odmkey{Композицией}{Композиция (composition)!подстановок (of
substitutions)} подстановок $\sigma_1=[t_i/ v_i]_{i\in I}$ и
$\sigma_2=[t_j/ v_j]_{j\in J}$ называется
подстановка $\sigma_1\!\circ\!\sigma_2\Assign [t_k/
v_k]_{k\in K}$, где $K=I\!\cup\!J$, \\ а
$t_k=$ \textbf{if} $k\in I$ \textbf{then} $t_i\sigma_2$ \textbf{else} 
$t_j$.
Множество подстановок образует моноид относительно композиции,
причём тождественная подстановка $[v_i/ v_i]$ является
единицей.
\end{enumerate}
\end{examples}
\end{frame}


\begin{frame}
\frametitle{2.2.5. Моноиды и лупы (2/3)}

Функция $\Map{f}{M}{M}$ 
называется \odmkey{преобразованием над множеством}{Преобразование (transformation)!над множеством (on set)} $M$.
Всякое множество преобразований, замкнутое относительно суперпозиции,
образует моноид (тождественное преобразование является нейтральным элементом).  

\begin{theorem} Всякий моноид над $M$ изоморфен некоторому моноиду
преобразований над $M$.
\end{theorem}
\begin{proof}
Пусть $\mathcal{M}=\Alg{M}{\ast}$~--- моноид над $M$
с единицей
$e$. Построим \\ $\cal F=\Alg{F}{\circ}$~--- моноид
преобразований над $M$, где \\
$F{\Assign}\Set{\Map{f_m}{M}{M}}{f_m(x)\Assign 
x \ast m}_{m\in M}$, 
а~$\circ$~--- суперпозиция функций. 
\\ Рассмотрим функцию $\Map{h}{M}{F}$,
$h(m)\Assign f_m$. 

Тогда $\cal F$~--- моноид, поскольку
суперпозиция функций ассоциативна, $f_e$~--- тождественная функция
(так~как $f_e(x)=x\ast e=x$) и $F$ замкнуто относительно $\circ$,
так~как \\$f_a \circ f_b=f_{a\ast b}\colon$ $f_{a \ast b}(x)=x \ast
(a\ast b)= (x\ast a)\ast b=f_a(x)\ast b=f_b(f_a(x))$. Далее,
$h$~--- гомоморфизм, так~как 
$h(a\ast b) = f_{a\ast b} = f_a\!\circ\! f_b = 
h(a) \circ  h(b)$. И наконец, $h$~--- биекция, поскольку $h$~---
сюръекция по построению, и  $h$~--- инъекция 
(так~как \\
$(f_a(e) = e \ast  a= a\And f_b(e) =  e \ast  b = b)\Impl (a \ne b\Impl
f_a \ne f_b)$).
\end{proof}
\end{frame}


\begin{frame}
\frametitle{2.2.5. Моноиды и лупы (3/3)}

Если для бинарной операции 
$\Map{\ast}{M\times M}{M}$ 
выполнены свойства {\color{blue}(A6)} и 
{\color{blue}(A9)}, то алгебра
 $\Alg{M}{\ast}$ называется \odmkey{лупой}{Лупа (loop)}.
 Иначе говоря, лупа~--- это квазигруппа  
 с нейтральным элементом.

\begin{digress}
Таблица Кэли (\DMref{п.~2.1.5}{2.1.5.}) для квазигруппы непрменно является
латинским квадратом (\DMref{п.~8.3.3}{8.3.3.}).
Для лупы этот латинский квадрат характерен тем,
что в нём есть строка, совпадающая с заголовками столбцов и
есть столбец, совпадающий с заголовками строк. 
Наименьшая неассоциативная лупа имеет порядок 5,
 её таблица Кэли 
приведена ниже.
\begin{center}
\begin{tabular}{|c||c|c|c|c|c|}
\hline
$\ast$ & $1$ & $2$ & $3$ & $4$ & $5$ \\ \hline \hline
$1$ & $1$ & $2$ & $3$ & $4$ & $5$ \\ \hline
$2$ & $2$ & $1$ & $5$ & $3$ & $4$ \\ \hline
$3$ & $3$ & $4$ & $1$ & $5$ & $2$ \\ \hline
$4$ & $4$ & $5$ & $2$ & $1$ & $3$ \\ \hline
$5$ & $5$ & $3$ & $4$ & $2$ & $1$ \\ \hline
\end{tabular}
\end{center}
\end{digress} 

В \DMref{п.~2.1.4}{2.1.4.} показано,
что из разрешимости (свойство {\color{blue}(A9)})
и существования нейтрального элемента 
(свойство {\color{blue}(A6)})
следует обратимость (свойство {\color{blue}(A8)}).
Поэтому в лупе для любого элемента определён обратный. 

\end{frame}

\begin{frame}
\label{2.2.6.}\subsubsection{2.2.6. Группы }\frametitle{2.2.6. Группы (1/4)}

Теория групп, некоторые положения которой рассматриваются в этом
разделе, является ярчайшим примером мощи алгебраических методов,
позволяющих получить из немногих базовых понятий и аксиом
множество важнейших следствий.

Если для бинарной операции 
$\Map{\ast}{M\times M}{M}$ 
 выполнены свойства {\color{blue}(A1)}, 
 {\color{blue}(A6)} и {\color{blue}(A8)},
 то алгебра $\Alg{M}{\ast}$ называется 
 \odmkey{группой}{Группа (group)}. Иначе говоря, группа~--- это моноид, в котором для любого элемента существует обратный.

%\odmkey{Группа}{Группа (group)}~--- это моноид, в котором
%$
%\A{a}{\E{a^{-1}}{a \ast a^{-1}  \!=\! a^{-1} \ast a  \!=\! e}}.
%$
%Элемент $a^{-1}$ называется \odmkey{обратным}{Элемент (element)!обратный
%(inverse)},
%а операция $\ast$ называется умножением.

\begin{remark}
Так как операция умножения в группе не обязательно коммутативна,
то, вообще говоря, свойства $a^{-1} \ast a =e$ и $a \ast
a^{-1}=e$ неравносильны. Но теоремы \DMref{п.~2.1.4}{2.1.4.} показывают, что в условиях ассоциативности
не требуется различать обратные слева и
обратные справа элементы. 
Введённая аксиома требует существования
двустороннего обратного элемента. Можно показать, что она следует
из более слабых аксиом существования левой единицы и левого обратного элемента.
\end{remark}

\smallskip
Группу можно определить и как \em{ассоциативную лупу}. Нетрудно
и очень полезно проверить, что приведённые определения эквивалентны.

\smallskip
Обычно группу обозначают буквой $\mathcal{G}\Assign\Alg{G}{\ast}$. Мощность носителя $G$ группы 
$\mathcal{G}$ называется
\odmkey{порядком группы}{Порядок (order)!группы (of group)} и
обозначается $|G|$.


\end{frame}

\begin{frame}[shrink]
\frametitle{2.2.6. Группы (2/4)}

\begin{examples}
\begin{enumerate}
\item Множество невырожденных квадратных матриц порядка $n$ образует
группу относительно операции умножения матриц. Единицей группы
является единичная матрица. Обратным элементом является обратная матрица.
\item Множество перестановок, то~есть множество
взаимно однозначных функций 
\\$\Map{f}{M}{M}$, является группой
относительно операции суперпозиции. Единицей группы является
тождественная функция, а обратным элементом~--- обратная функция.
\item Множество поворотов правильного треугольника относительно его центра на\\ 
$0^\circ, 120^\circ, 240^\circ$ образует группу поворотов, 
переводящих треугольник в себя. Ниже приведена <<таблица умножения>> группы
и геометрическая иллюстрация.
\begin{columns}[t] % contents are top vertically aligned
\begin{column}[T]{7cm} % each column can also be its own environment
\begin{tabular}{|c|c|c|c|}
\hline
                  & $0^\circ$ & $120^\circ$ & $240^\circ$\\
\hline
$0^\circ$         & $0^\circ$ & $120^\circ$ & $240^\circ$\\
$120^\circ$       & $120^\circ$ & $240^\circ$ & $0^\circ$\\
$240^\circ$       & $240^\circ$ & $0^\circ$ & $120^\circ$\\              
\hline
\end{tabular}

\end{column}
\begin{column}[T]{4cm} % alternative top-align that's better for graphics
\includegraphics[height=20mm]{2_2_6.jpg}
\end{column}
\end{columns}
\end{enumerate}
\end{examples}
\end{frame}

\begin{frame}
\frametitle{2.2.6. Группы (3/4)}


\begin{theorem} В любой группе выполняются следующие соотношения:
\begin{enumerate}
\item[1)] ${(a\ast b)}^{-1}=b^{-1}\ast a^{-1}$;
\item[2)] ${(a^{-1})}^{-1} = a$.
\end{enumerate}
\end{theorem}


\begin{proof}
\proofsection{\color{blue}1)}
$(a  \ast  b)  \ast (b^{-1}  \ast  a^{-1})  = a  \ast (b  \ast  b^{-1})  \ast  a^{-1}  = a  \ast  e  \ast 
a^{-1}  = a  \ast  a^{-1}  = e$.
\proofsection{\color{blue}2)}
$a^{-1}  \ast  a  = e = a^{-1} \ast (a^{-1})^{-1}   \Impl {(a^{-1})}^{-1} = a$.
\end{proof}


\begin{theorem} В любой группе выполняется свойство сократимости, то есть:
\begin{enumerate}
\item[1)] $a \ast b = a \ast c \Impl b = c$;
\item[2)] $b \ast a = c \ast a \Impl b = c$;
\end{enumerate}
\end{theorem}


\begin{proof}
\proofsection{\color{blue}1)}
$b  =  e   \ast  b   =  (a^{-1}  \ast  a)  \ast  b   =  a^{-1}  \ast (a  \ast  b)  = a^{-1}  \ast (a  \ast  c)  = 
(a^{-1}  \ast  a)  \ast  c  =  e  \ast  c   =  c$.
\proofsection{\color{blue}2)}
$b  =  b   \ast  e   =  b   \ast  (a  \ast  a^{-1})   =  (b   \ast  a)  \ast  a^{-1}  = (c   \ast  a)  \ast  a^{-1}  = 
c   \ast (a   \ast  a^{-1})   =  c   \ast  e  =  c$.
\end{proof}


\begin{theorem} В группе можно однозначно решить уравнение $a \ast x = b$
({решение: } $x = a^{-1} \ast b$).
\end{theorem}
\begin{proof}
$a \ast x = b\Impl a^{-1} \ast (a\ast x)=a^{-1} \ast b \Impl(a^{-1}\ast a) \ast x=a^{-1} \ast b \Impl$
$\Impl  e \ast
x =a^{-1} \ast b \Impl x=a^{-1}\ast b$.
\end{proof}

\end{frame}

\begin{frame}
\frametitle{2.2.6. Группы (4/4)}
%\vspace{-2pt}

%\odmkey{Коммутативная}{Группа (group)!коммутативная (commutative)}
%группа,  в которой
%$
%\A{a,b}{a\! \ast\! b\! =\! b\! \ast\! a},
%$
Если в группе выполнено свойство {\color{blue}(А2)},
то есть операция группы коммутативна, то группа
называется \odmkey{абелевой}{Группа (group)!абелева (abelian)}.
 В абелевых группах 
 обычно приняты
следующие обозначения: 
операция в такой группе обозначается $+$,
обратный элемент к $a$ обозначается $-a$, 
нейтральный элемент группы
обозначается $0$ и называется \odmkey{нулем}{Нуль!группы (group
identity)}.
% или \odmkey{нейтральным элементом}{Элемент
%(element)!нейтральный (neutral)}.
\begin{examples}
\begin{enumerate}
\item
 $\Alg{\Bbb Z}{+}$~--- множество целых чисел образует абелеву
группу относительно сложения. 
Нейтральным элементом группы
является число~0. Обратным элементом является число с
противоположным знаком: $x^{-1}\Assign -x$. 
\item 
$\Alg{\Bbb Q_+}{\cdot}$~--- множество положительных рациональных чисел
образует абелеву группу относительно умножения. 
Нейтральным
элементом группы является число~1. \\Обратным элементом является
обратное число: $(m/n)^{-1}\Assign n/m$.
\item 
$\Alg{2^M}{\vartriangle}$~--- булеан образует абелеву группу
относительно симметрической разности. 
Нейтральным элементом группы
является пустое множество $\emptyset$. Обратным элементом
к~элементу $x$ является он сам: $x^{-1}\Assign x$.
\end{enumerate}
\end{examples}
\end{frame}

\begin{frame}
\label{2.2.7.}\subsubsection{2.2.7. Действие группы на множестве }\frametitle{2.2.7. Действие группы на множестве (1/3)}

Пусть $S$~--- множество, а $G$~--- группа 
с операцией $\ast$ и нейтральным элементом $e$. 

\smallskip
\odmkey{ Действием группы $G$ на множестве $S$ }{Действие группы (group action)!на множестве (on set)} 
называется функция
$\Map{\phi}{G\times S}{S}$ такая,
что  $$ \A{s \in  S}{\phi(e,s)=s} \text{ и } 
\A{g_1, g_2 \in  G, s \in  S }{\phi(g_1 \ast  g_2, s)=
\phi(g_1, \phi(g_2,s))}.$$
Для краткости записи принято опускать знак операции $\ast$ 
и имя отображения $\phi$, обозначая $\phi(g,s)$ просто $gs$.
В таких обозначениях $$ \A{s\in S}{es=s}\text{ и }
\A{g_1, g_2\in G, s\in S }{(g_1g_2)s=
g_1(g_2 s)}.$$   
В этом случае также говорят,
что группа $G$ \odmkey{ действует}{Действие группы (group action)!на множестве (on set)}  на множестве $S$ или что
$S$ есть $G$-множество.

\smallskip
Если $s\in S$, то множество $Gs \Def \Set{t\in S}{t=gs \And g\in G}$,
то есть множества всех элементов из $S$ вида $gs$, где $g$
<<пробегает>> носитель группы $G$, называется
\odmkey{орбитой}{Орбита элемента (element orbit)} (или \odmkey{траекторией}{Траектория элемента (element trajectory)}) элемента $s$
в $G$-множестве $S$.  


\end{frame}

\begin{frame}
\frametitle{2.2.7. Действие группы на множестве (2/3)} 

\begin{example}
Пусть $S_1$~--- единичная окружность, а $\varphi\in [0;2\pi]$~---
некоторый угол. Определим действие группы $\Alg{\Bbb{Z}}{+}$ на 
множестве $S_1$ следующим образом: число $n\in Z$ 
действует на точку $x\in S_1$, как поворот на угол $n\varphi$.
Положительные углы отсчитываются против часовой стрелки, отрицательные~---
по часовой стрелке.    
Тогда если угол $\varphi$ соизмерим с числом $\pi$,
то все орбиты состоят из конечного числа точек.
Если же отношение $\varphi /\pi$ иррационально,
то все орбиты бесконечны.
\end{example}

\smallskip
\begin{lemma}
Если $G$~--- группа, то $\A{g\in G}{Gg=G}$.
\end{lemma}

\begin{proof}
Ясно, что $\A{g\in G}{Gg\subset G}$. Пусть
$g_1\in G$. \\ Тогда $\A{g\!\in\!G}{\E{g_2\!\in\!G}{g_2g=g_1}}$ 
по теореме \DMref{п.~2.2.6}{2.2.6.}, 
$g_1\in Gg$ и $G\subset Gg$.  
\end{proof}


\end{frame}

\begin{frame}
\frametitle{2.2.7. Действие группы на множестве (3/3)} 

\begin{theorem}
Если группа $G$ действует на множестве $S$,
то орбиты не пусты и образуют разбиение 
множества $S$.
\end{theorem}

\begin{proof}
Каждый элемент принадлежит своей орбите, 
так как из определения действия группы $G$ 
на множестве $S$ имеем $\A{s\in S}{es=s}$, поэтому орбиты не пусты
и образуют покрытие (\DMref{п.~1.2.7}{1.2.7.}).
Пусть две орбиты $Gs_1$ и $Gs_2$ пересекаются, то есть\\
 $\E{s\in S}{s \in Gs_1 \And s\in Gs_2}$. 
 Тогда $\E{g_1\in G}{s=g_1s_1}$ и по лемме $Gs=Gg_1s_1= Gs_1$.
Аналогично $\E{g_2\in G}{s=g_2s_2}$ и $Gs=Gg_2s_2= Gs_2$,
откуда $Gs_1 =Gs_2$.    
\end{proof}

\smallskip
Таким образом, действие группы на множестве определяет
отношение эквивалентности на этом множестве (\DMref{п.~1.7.1}{1.7.1.}).

\smallskip
\begin{example}
Каждая группа $G$ действует на самой себе с помощью умножения.
Это действие имеет всего одну орбиту. 
\end{example}

\smallskip
Задать действие группы на множестве~--- всё равно, что 
рассматривать элементы группы как преобразования множества. 
\end{frame}

\begin{frame}
\label{2.2.8.}\subsubsection{2.2.8. Группа перестановок }\frametitle{2.2.8. Группа перестановок (1/4)}
%\subsection{}\label{3.1.6}
В этом параграфе рассматривается одна из важнейших групп,
называемая группой \odmkey{перестановок}{Группа
(group)!перестановок (of permutations)}, или
\odmkey{симметрической}{Группа (group)!симметрическая (symmetric)}
группой. 

Биективная функция $\Map{f}{X}{X}$ называется
\odmkey{перестановкой}{Перестановка (permutation)} множества~$X$.
\begin{remark}
Если множество $X$ конечно ($\vert X\vert = n$), то, не
ограничивая общности, можно считать, что $X=1..n$. В этом случае
перестановку $\Map{f}{1..n}{1..n}$ удобно задавать таблицей из
двух строк. В первой строке~--- значения аргументов, во второй~---
соответствующие значения функции. Такая таблица называется
\odmkey{подстановкой}{Подстановка (substitution)}.
В сущности, перестановка и~подстановка~--- синонимы.
\end{remark}
\begin{example}
$$
f =
\begin{vmatrix}
1 & 2 & 3 & 4 & 5 \\
5 & 2 & 1 & 4 & 3
\end{vmatrix},\qquad
g =
\begin{vmatrix}
1 & 2 & 3 & 4 & 5 \\
4 & 1 & 2 & 3 & 5
\end{vmatrix}.
$$
\end{example}
\odmkey{Произведением}{Произведение (product)!перестановок (of
permutations)} перестановок $f$ и $g$ (обозначается $fg$)
называется их суперпозиция $g\circ f$.
\end{frame}

\begin{frame}
\frametitle{2.2.8. Группа перестановок (2/4)}
\begin{example}
$$
fg =
\begin{vmatrix}
1 & 2 & 3 & 4 & 5 \\
5 & 1 & 4 & 3 & 2
\end{vmatrix}.
$$
\end{example}


\odmkey{Тождественная}{Перестановка (permutation)!тождественная
(identity)} перестановка~--- это перестановка $e$, такая, что
$e(x)=x$.
\begin{example}
$$
e =
\begin{vmatrix}
1 & 2 & 3 & 4 & 5 \\
1 & 2 & 3 & 4 & 5
\end{vmatrix}.
$$
\end{example}
\odmkey{Обратная}{Перестановка (permutation)!обратная (inverse)}
перестановка~--- это
обратная функция, которая всегда существует, поскольку
перестановка является биекцией.
\begin{remark}
Таблицу обратной подстановки можно получить, если просто поменять
местами строки таблицы исходной подстановки.
\end{remark}
\begin{example}
$$
f =
\begin{vmatrix}
1 & 2 & 3 & 4 & 5 \\
3 & 4 & 2 & 1 & 5
\end{vmatrix},\qquad
f^{-1} =
\begin{vmatrix}
3 & 4 & 2 & 1 & 5\\
1 & 2 & 3 & 4 & 5
\end{vmatrix} =
\begin{vmatrix}
1 & 2 & 3 & 4 & 5\\
4 & 3 & 1 & 2 & 5
\end{vmatrix}.
$$
\end{example}
\end{frame}

\begin{frame}
\frametitle{2.2.8. Группа перестановок (3/4)}
Таким образом, поскольку суперпозиция функций ассоциативна,
а единичный и~обратный элементы существуют,
множество перестановок $n$-элементного множества
образует \odmkey{группу перестановок}{Группа (group)!перестановок (of permutations)} 
относительно операции суперпозиции. Эта группа
называется также
 \odmkey{симметрической группой}{Группа (group)!симметрическая (symmetric)}
порядка
$n$ и обозначается $S_n$.
\smallskip
\begin{theorem}
Всякая конечная группа $G$ порядка $n$ 
изоморфна некоторой подгруппе симметрической группы $S_n$.
\end{theorem}

\begin{proof} 
Пусть $G$~--- группа, $|G|=n$, нейтральный элемент $e$. 
Рассмотрим действие группы $G$ на множестве $G$ с помощью умножения
(\DMref{п.~2.2.7}{2.2.7.}). Каждый элемент $g\in G$ задаёт подстановку $f_g \in S_n$.
Рассмотрим множество 
$$H \Assign\Set{\Map{f_g}{G}{G}}{f_g(x) \Assign x\ast g \And g\in G}\subset S_n.$$
Ясно, что отображение $g\mapsto f_g$ биективно и уважает умножение,
так что $G\sim H$.
Покажем, что $H$~--- подгруппа группы $S_n$. Действительно,
$H$ замкнуто относительно $\circ$,
так~как \\ $f_{g\ast h}(x)=x\ast
(g\ast h)= (x\ast g)\ast h=f_g(x)\ast h=f_h(f_g(x))$;
 $f_e$~--- тождественная функция,
так~как $f_e(x)=x\ast e=x$,   
$f_{g^{-1}}$~--- обратная функция,
так~как $f_g\left(f_{g^{-1}}(x)\right)=(x\ast g^{-1})\ast g=x$.
\end{proof}

\end{frame}

\begin{frame}
\frametitle{2.2.8. Группа перестановок (4/4)}
\begin{example}
Группа поворотов правильного треугольника (см. пример 3 в \DMref{п.~2.2.6}{2.2.6.})
изоморфна следующей подгруппе группы $S_3$:


$$
e=\begin{vmatrix}
1 & 2 & 3 \\
1 & 2 & 3
\end{vmatrix},\qquad
f=\begin{vmatrix}
1 & 2 & 3 \\
2 & 3 & 1
\end{vmatrix},\qquad
f^{-1}=\begin{vmatrix}
1 & 2 & 3 \\
3 & 1 & 2
\end{vmatrix}.
$$
\end{example}

\medskip
\begin{remark}
Не следует думать, что описание подгрупп симметрической группы
является простой задачей. Это реально только для небольших $n$.
\end{remark}
\end{frame}

\begin{frame}
\label{2.2.9.}\subsubsection{2.2.9. Перестановочные матрицы }\frametitle{2.2.9. Перестановочные матрицы (1/3)}

Если $\Map{f}{1..n}{1..n}$~--- перестановка, то 
\odmkey{матрицей подстановки}{Матрица (matrix)!подстановки (of substitution)} или \odmkey{перестановочной матрицей}{Матрица (matrix)!перестановочная (of permutation)}
называется квадратная булева матрица \\
 \textbf{$P_f:$ array $[1..n,1..n]$ of $0..1$},
такая, что $P_f[i,j]\Assign (j=f(i))$. 

\smallskip
Нетрудно видеть, что в каждой строке и в каждом столбце 
перестановочной матрицы содержится ровно одна единица, 
остальные элементы~---  нули.
Более того, всякой булевой матрице $P$, 
у которой в каждой строке и 
в каждом столбце 
содержится ровно одна единица,
соответствует единственная перестановка $f$, 
которая получится, если матрицу $P_f$ умножить справа на
вектор-столбец $(1, \dots, n )$.

\medskip
\begin{example}
$f =
\begin{vmatrix}
1 & 2 & 3 & 4 & 5 \\
3 & 4 & 2 & 1 & 5
\end{vmatrix},\qquad
P_f =
\begin{pmatrix}
0 & 0 & 1 & 0 & 0\\
0 & 0 & 0 & 1 & 0\\
0 & 1 & 0 & 0 & 0\\
1 & 0 & 0 & 0 & 0\\
0 & 0 & 0 & 0 & 1
\end{pmatrix}$.
\end{example}

\begin{remark}
Хотя перестановочные матрицы булевы, умножение 
матриц может производиться с помощью обычных арифметических, а не логических
операций.
\end{remark}
\end{frame}

\begin{frame}
\frametitle{2.2.9. Перестановочные матрицы (2/3)}

Введём обозначение: $\delta_i^n$~--- булевский вектор длины $n$,
в котором в $i$-м разряде находится единица, а в остальных разрядах находятся нули.  
В этих обозначениях для каждой строки перестановочной матрицы
 $P_f[i,1..n]\!=\! \delta^n_{f(i)}$,
а для каждого столбца $P_f[1..n,j]\!=\! \delta^n_{f^{-1}(j)}$.

\smallskip
\begin{lemma}
Множество перестановочных матриц одного порядка 
образует некоммутативную группу относительно операции умножения матриц.
\end{lemma}



\begin{proof} Произведение
перестановочных матриц является перестановочной матрицей.
Действительно, поскольку
\odmkey{скалярное произведение}{Произведение (product)!скалярное (scalar)} 

\vspace{-8pt}
$$\langle\delta^n_i\mid \delta^n_j\rangle\Def \sum\nolimits_{k=1}^n \delta_i^n[k]\cdot\delta_j^n[k]= \textbf{ if } i=j \textbf{ then } 1 
\textbf{ else } 0 \textbf{ end~if },$$

\vspace{-2pt}
то произведение булевых матриц является булевой матрицей.
Также $P_f\times\delta^n_i = \delta^n_{f(i)}$, 
то есть каждая строка произведения содержит ровно одну единицу, 
и аналогично для столбцов, то есть произведение 
является перестановочной матрицей.
Далее, умножение перестановочных матриц ассоциативно в силу
ассоциативности умножения квадратных матриц.
Единичная матрица является перестановочной и 
служит нейтральным элементом. Транспонирование
перестановочной матрицы даёт перестановочную матрицу, которая
является обратной, то есть $P_f P_f^{\mathrm{T}}= I_n$,
где $I_n$~--- единичная матрица.    
\end{proof}  


\end{frame}

\begin{frame}
\frametitle{2.2.9. Перестановочные матрицы (3/3)}

Обозначим группу перестановочных матриц порядка $n$ как $P_n$.

\smallskip
\begin{theorem}
Группа перестановочных матриц $P_n$ 
изоморфна симметрической группе $S_n$. 
\end{theorem}
\begin{proof}
В \DMref{п.~5.1.4}{5.1.4.} показано, что $|S_n|=n!$,
и по правилу произведения (см.~\DMref{п.~5.1.1}{5.1.1.}) имеем $|P_n|= n!$,
поскольку единицу в первой строке можно выбрать $n$
способами, единицу во второй строке можно выбрать $n-1$ способами
и т.~д. Отображение $f\mapsto P_f$, определённое в начале параграфа,
является биекцией. Осталось показать, что это отображение
является гомоморфизмом, то есть $P_{fg}=P_f P_g$.
Имеем: 
$(P_f P_g )[i,j]=1 \Equiv $
$\Equiv\sum_{k=1}^n P_f[i,k]\cdot P_g[k,j]=1  \Equiv 
\left\langle \delta^n_{f(i)} \mid \delta^n_{g^{-1}(j)} \right\rangle =1 \Equiv f(i)=g^{-1}(j) \Equiv$ 
$\Equiv g(f(i))=j \Equiv 
P_{fg}[i,g(f(i))]=1
$. 
\end{proof}

\medskip
Значение перестановочных матриц состоит в том, 
что это эффективное \em{представление}
перестановки, позволяющее свести применение перестановки
как функции к операциям с булевыми матрицами.
В частности, если $(\Tuple{a}{1}{n})$~---
любой числовой $n$-вектор (не обязательно перестановка $1..n$),
а $P_f$~--- перестановочная матрица для перестановки $f$,
то $P_f\times(\Tuple{a}{1}{n})^T = (a_{f(1)}, \dots, a_{f(n)})^T$.    


\end{frame}


\begin{frame}
\frametitle{Онтология общей алгебры}
\begin{center}
\includegraphics[height=7.5cm]{ODM2_2.png}
\end{center}
\end{frame}


\subsection{2.3. Алгебры с двумя операциями}
\label{2.3.}
\begin{frame}
\frametitle{2.3. Алгебры с двумя операциями}


В этом разделе мы обращаемся к объектам,
знакомым читателю со
школы. Среди алгебр с двумя
операциями наиболее важными являются кольца и поля,
а основными примерами колец и полей являются множества целых,
рациональных и~вещественных чисел с операциями сложения и
умножения.

\begin{tabular}{p{12mm}|p{45mm}|p{86mm}}
  № & Название параграфа
  & Ключевые термины и обозначения \\
  \hline
  & & \\
{\DMref{2.3.1.}{2.3.1.}}\RS & Кольца &
{коммутативные и с единицей}\\
\DMref{2.3.2.}{2.3.2.}  & Области целостности и тела &
{делители нуля}\\ 
\DMref{2.3.3.}{2.3.3.}   & Поля & \\
\DMref{2.3.4.}{2.3.4.}  \RS  & Тропическая математика &
\end{tabular}  
\end{frame}

\begin{frame}
\label{2.3.1.}\subsubsection{2.3.1. Кольца }\frametitle{2.3.1. Кольца (1/3)}


Пусть $M$~--- множество, на котором заданы
две бинарные операции: 
$\Map{+}{M\times M}{M}$ и
$\Map{\ast}{M\times M}{M}$
(они называются сложением и умножением соответственно). Если при этом для сложения выполнены свойства {\color{blue}(A1)}, {\color{blue}(A2)},
{\color{blue}(A6)} и {\color{blue}(A8)}, 
для умножения – {\color{blue}(А1)}, и операции связаны между собой дистрибутивным законом {\color{blue}(A3)},
то алгебра $\Alg{M}{+,\ast}$
называется \odmkey{кольцом}{Кольцо (ring)}. 
Другими словами, кольцо~--- это абелева группа по сложению и
полугруппа по умножению, причём 
умножение дистрибутивно относительно сложения слева и справа. 

\begin{remark}
Умножение в кольце не обязательно коммутативно,
поэтому необходимо требовать дистрибутивность слева и справа.
\end{remark}
\begin{digress}
В любой алгебре, в которой действуют 
несколько операций, обязательно должны 
присутствовать законы (аксиомы),
которые связывают различные операции. 
\end{digress}
%\begin{tabbing}
%%2.2\= $\E{0 \in M}{\A{a}{a + 0 = 0 + a = a}}\qquad$\=существует нуль;\kill
%\begin{enumerate}
%\item[1)] $(a + b) + c = a + (b + c)$~---
%сложение ассоциативно;
%
%\item[2)] $\E{0 \in M}{\A{a}{a + 0 = 0 + a = a}}$~---
%существует нуль; 
%
%\item[3)] $\A{a}{\E{-a}{a + (-a) = 0}}$~---
% существует обратный элемент; 
% 
%\item[4)] $a + b = b + a $~---
%сложение коммутативно, то~есть 
% кольцо~--- абелева группа по сложению; 
%
%\item[5)] $a \ast (b \ast c) = (a \ast b) \ast c$~---
% умножение ассоциативно, то~есть 
% кольцо~--- полугруппа по умножению; 
% 
%\item[6)] $a \ast (b + c) = (a \ast b) + (a \ast c)$, 
%$(a + b) \ast c = (a \ast c) + (b \ast c)$~---
% умножение дистрибутивно относительно
%сложения слева и справа. 
%\end{enumerate}
%\end{tabbing}
Кольцо называется
\odmkey{коммутативным}{Кольцо (ring)!коммутативное (commutative)}, если для умножения 
выполнено свойство {\color{blue}(A2)},
то есть умножение коммутативно. 
Кольцо называется кольцом
\odmkey{с единицей}{Кольцо (ring)!с единицей (with identity)}, если для умножения 
выполнено свойство {\color{blue}(A6)}, то~есть кольцо 
 c единицей~--- моноид по умножению.

\begin{remark}
В кольце нейтральный элемент по сложению называется нулём и обозначается $0$, а нейтральный
элемент по умножению называется единицей 
и обозначается $1$.
\end{remark} 
 
\end{frame}

\begin{frame}
\frametitle{2.3.1. Кольца (2/3)}

\begin{theorem} В кольце выполняются следующие соотношения:
\begin{columns}[t]
\begin{column}[T]{7cm}
\begin{enumerate}
\item[1)] $0 \ast a = a\ast 0 = 0$;%
\item[2)] $a \ast (-b) = (-a) \ast b = -(a \ast b)$;%
\item[3)] $(-a) \ast (-b) = a \ast b$;%
\end{enumerate}
\end{column}
\begin{column}[T]{7cm}
\begin{enumerate}
\item[4)] $(-a) = a \ast (-1)$;%
\item[5)] $-(a +  b)=(-a) + (-b)$;%
\item[6)] $a \not= 0 \Impl{(a^{-1})}^{-1} = a$.%
\end{enumerate}
\end{column}
\end{columns}
\end{theorem}
\begin{proof}
\proofsection{\color{blue}1)}
$0\ast a=(0 + 0)\ast a=(0\ast a) + (0\ast a)\Impl
 -(0\ast a) + (0\ast a) = {-(0\ast a)}  +$ $+{\bigl((0\ast a)}  + {(0\ast
a)\bigr)} = \bigl(-(0\ast a)  + (0\ast a)\bigr) + (0\ast
a) \Impl
 {0 =0 + (0\ast a)} ={0\ast a}$.
\proofsection{\color{blue}2)}
$\bigl(a\ast(-b)\bigr) + (a\ast b)=a\ast(-b +  b)=a\ast
0=0$, \\
$(a\ast b) + \bigl((-a)\ast b\bigr)= \bigl(a + (-a)\bigr)\ast
b=0\ast b=0$. 
\proofsection{\color{blue}3)}
$(-a)\ast(-b)=-\bigl(a\ast(-b)\bigr)=-\bigl(-(a\ast b)\bigr)=a\ast
b$. 
\proofsection{\color{blue}4)} $\bigl(a\ast (-1)\bigr) +
a=\bigl(a\ast(-1)\bigr) + (a \ast 1)= a\ast(-1 + 1)=a\ast 0=0$.
\proofsection{\color{blue}5)} $(a +  b) +  \bigl((-a) +  (-b)\bigr)=(a +  b) +
\bigl((-b) +  (-a)\bigr) ={a +  \bigl(b +  (-b)\bigr) +
(-a)} =$\\$= {a  + 0  + (-a)} = {a +  (-a)=0}$. 
\proofsection{\color{blue}6)}
$a^{-1}\ast a=1$.
\end{proof}

\end{frame}

\begin{frame}
\frametitle{2.3.1. Кольца (3/3)}

\begin{examples}
\begin{enumerate}
\item Целые числа $\Alg{\Bbb Z}{+,*}$~--- коммутативное кольцо с единицей. 
\item Множество классов вычетов (\DMref{п.~2.4.6}{2.4.6.})
$\Alg{\Bbb Z_n}{+,*}$~--- коммутативное
кольцо с единицей при $n\ge 2$. 
\item \odmkey{Машинная арифметика}{Арифметика (arithmetic)!машинная (machine)} целых чисел
$\Alg{\Bbb Z_{2^{15}}}{+,*}$~--- коммутативное кольцо с единицей.
\item Множество числовых векторов одного размера,
в которых сложение и умножение выполняются покомпонентно,~--- коммутативное кольцо с единицей.
%; нулем является нулевой вектор,
%а единицей является единичный вектор. 
\item Множество квадратных числовых матриц одного порядка~--- некоммутативное кольцо с единицей.
%;
%единицей является единичная матрица. 
%\item Множество многочленов одной переменной~--- коммутативное кольцо с единицей. 
\end{enumerate}

Если множество $M$ содержит единственный элемент,
то всегда выполняются все свойства операций.
Такой неинтересный случай называется 
\odmkey{вырожденным}{Кольцо (ring)!вырожденный случай (degenerate case)}.

\begin{corollary}
В невырожденном кольце с единицей $0\ne 1$.
\end{corollary}
\end{examples}

\end{frame}

\begin{frame}
\label{2.3.2.}\subsubsection{2.3.2. Области целостности и тела }\frametitle{2.3.2. Области целостности и тела (1/2)}
%\subsection{}
Если в кольце для некоторых ненулевых элементов $x$, $y$
выполняется равенство $x \ast y = 0$, то $x$ называется
\odmkey{левым}{Делитель нуля (zero divisor)!левый (left)}, а
$y$~--- \odmkey{правым}{Делитель нуля (zero divisor)!правый
(right)} \odmkey{делителем нуля}{Делитель нуля (zero divisor)}.
\begin{example}
В машинной арифметике $\Alg{\Bbb Z_{2^{15}}}{+,*}$ имеем
$256*128= 2^8 * 2^7  =2^{15} =\nobreak 0$.
\end{example}

Заметим, что если в кольце нет делителей нуля, то $\A{x\ne 0, y\ne\nobreak
0}{x\ast y \ne 0}$. В~группе $\A{a,b,c}{\bigl(a \ast b = a \ast c
\Impl b = c\bigr) \And \bigl(b \ast a = c \ast a \Impl b =
c\bigr)}$, однако в~произвольном кольце это не так.
\begin{example}
В машинной арифметике $\Alg{\Bbb Z_{2^{15}}}{+,*}$ имеем
$2^8 * 2^7 =0$ и $0 * 2^7 =0$, \\но $0\ne 2^8=256$.
\end{example}
\begin{theorem}
$\forall a\ne 0$ $\bigl(a \ast b = a \ast c \Impl b = c\bigr) \And
\bigl(b \ast a = c \ast a \Impl b = c\bigr)\Equiv$ \\
$\Equiv \A{x \ne 0, y\ne 0}{x\ast y \ne 0}.$
\end{theorem}
\begin{proof}
\proofsection{$\Impl$} От противного. Пусть $\E{x\ne 0,y\ne
0}{x\ast y=0}$. \\Тогда $x\not=0\And x\ast y ={0 \And x \ast
0} =0 \Impl y=0$. 
\proofsection{$\Impll$} $0=(a \ast b) + 
\bigl(-(a \ast b)\bigr)= (a \ast b) + \bigl(-(a \ast c)\bigr)=(a \ast
b) + \bigl(a \ast (-c)\bigr) = a \ast \bigl(b + (-c)\bigr)$,
$a \ast \bigl(b + (-c)\bigr) = 0 \And a \not= 0 \Impl b + (-c) = 0\Impl
b = c$.
\end{proof}
\end{frame}

\begin{frame}
\frametitle{2.3.2. Области целостности и тела (2/2)}


Коммутативное кольцо с единицей, не имеющее делителей нуля,
называется \odmkey{областью целостности}{Область (domain)!целостности (integral)}.
\begin{example}
Целые числа $\Alg{\Bbb Z}{+,*}$ являются областью целостности, а
машинная арифметика $\Alg{\Bbb Z_{2^{15}}}{+,*}$~--- не является.
\end{example}

\medskip
Кольцо с единицей, в котором для всех
ненулевых элементов существуют обратные по умножению,
называется \odmkey{телом}{Тело (skew-field)}. 

\begin{theorem}
Тело не имеет делителей нуля.
\end{theorem}

\begin{proof}
От противного. Пусть $x\ast y =0$, $x\ne 0$, $y\ne 0$. 
Тогда $0=x^{-1}\ast 0= x^{-1}\ast( x\ast y) 
=(x^{-1}\ast x)\ast y = 1\ast y = y$.
\end{proof}

\begin{corollary}
Коммутативное тело является областью целостности.
\end{corollary}

\begin{example}
Множество чисел вида $a+b\sqrt{2}$, 
где $a,b\in \Bbb{Q}$, образует коммутативное тело 
относительно обычных операций сложения и умножения.
\end{example}
\end{frame}


\begin{frame}
\label{2.3.3.}
\subsubsection{2.3.3. Поля }\frametitle{2.3.3. Поля (1/3)}

Пусть $M$~--- множество, на котором заданы
две бинарные операции: 
$\Map{+}{M\times M}{M}$ (сложение) и
$\Map{\ast}{M\times M}{M}$
(умножение). Если при этом для сложения выполнены свойства {\color{blue}(A1)}, {\color{blue}(A2)},
{\color{blue}(A6)} и {\color{blue}(A8)} c нейтральным элементом $0$, 
а для умножения также выполнены свойства {\color{blue}(A1)}, {\color{blue}(A2)},
{\color{blue}(A6)} и {\color{blue}(A8)}, но на множестве $M-0$, и операции связаны между собой дистрибутивным законом {\color{blue}(A3)},
то алгебра $\Alg{M}{+,\ast}$
называется \odmkey{полем}{Поле (field)}.
 
Другими словами, поле~--- это абелева группа по сложению и
абелева группа по умножению, исключая $0$, причём 
умножение дистрибутивно относительно сложения. 

Учитывая введённые ранее понятия, можно констатировать также,
что поле~--- это коммутативное кольцо, которое является коммутативным телом и, стало быть, областью целостности.  
%
%\odmkey{Поле}{Поле (field)}~--- это множество $M$ c двумя
%бинарными операциями $ + $ и $\ast$, такими, что
%
%\begin{enumerate}
%\item[1)] $(a + b) + c = a + (b+ c)$~---
%сложение ассоциативно; 
%
%\item[2)] $\E{0 \in M}{a + 0 = 0 + a = a}$~---
% существует нуль;
%
%\item[3)] $\A{a}{\E{-a}{a + (-a) = 0}} $~---
%существует обратный элемент по сложению; 
%
%\item[4)] $a + b = b + a$~---
% сложение коммутативно, то~есть поле~--- абелева 
% группа по сложению;
% 
%\item[5)] $a \ast (b \ast c) = (a \ast b) \ast c$~---
%умножение ассоциативно; 
%
%\item[6)] $\E{1 \in M}{a \ast 1 = 1 \ast a = a}$~---
%существует единица;
%
%\item[7)] $\A{a \ne 0}{\E{a^{-1}}{a^{-1} \ast a=1}}$~---
%существует обратный элемент по умножению; 
%
%\item[8)] $a \ast b = b \ast a $~---
%умножение коммутативно, то~есть ненулевые элементы 
%образуют абелеву группу по умножению;
%
%\item[9)] $a \ast (b + c) = (a \ast b) + (a \ast c)$~---
%умножение дистрибутивно относительно сложения.
%%\end{tabbing}
%\end{enumerate}

\begin{examples}
\begin{enumerate}
\item $\Alg{\Bbb R}{+,*}$~--- поле вещественных чисел. 

\item
$\Alg{\Bbb Q}{+,*}$~--- поле рациональных чисел.

\end{enumerate}
\end{examples}

\end{frame}

\begin{frame}
\frametitle{2.3.3. Поля (2/3)}
\begin{examples}
\begin{enumerate}

\item[3.]  Пусть $E_2\Def\{0,1\}$. Определим операции $\Map{ +
,\cdot}{E_2\times E_2}{E_2}$ следующим образом: $0\cdot 0=0\cdot 1
= 1\cdot 0 = 0$, $1\cdot 1 = 1$ (конъюнкция), $0  +  0 = 1 + 1 =
0$, $0 +  1 = 1  + 0= 1$ (сложение по модулю 2). Тогда $\Bbb
Z_2\Def\Alg{E_2}{ + ,\cdot}$ является полем и называется
\odmkey{двоичной арифметикой}{Арифметика (arithmetic)!двоичная (binary)}.
 В двоичной арифметике
нуль~--- это 0, единица~--- это 1, $-1 = 1^{-1} = 1$, а
$0^{-1}$~--- как всегда, не определён. 

\item[4.] Если в условиях
предыдущего примера взять в качестве сложения дизъюнкцию, то есть
определить операции следующим образом: $0\cdot 0=0\cdot 1 = 1\cdot
0 = 0$, $1\cdot 1 = 1$ (конъюнкция), $0+0=0$, $0+1=1+0=1+1=1$
(дизъюнкция), то структура $\Alg{E_2}{ \Or ,\cdot}$ не является
полем, поскольку у элемента 1 нет обратного по сложению.
\item[5.] Множество классов вычетов по простому модулю
является полем (\DMref{п.~2.5.4}{2.5.4.}).
\end{enumerate}
\end{examples}


\end{frame}

\begin{frame}
\frametitle{2.3.3. Поля (3/3)}
%
%
%Таким образом, поле~--- это коммутативное кольцо с единицей,
%в котором каждый элемент,
%кроме нуля,
%имеет обратный элемент по умножению.
%
%\medskip
%\begin{theorem} Поле является областью целостности:
% $a \ast b = 0 \Impl a = 0 \Or b = 0$.
%\end{theorem}
%\begin{proof}
%$a\ast b=0\And a\not =0\,{\Impl}\,b =1\ast b={(a^{-1}\ast
%a)\ast b}={a^{-1}\ast\,(a\ast b)}={a^{-1}\ast0\,=0}$,
%$a\ast b =0\And b\not = 0\,{\Impl}\,a =1\ast
%a =(b^{-1}\ast b)\ast a ={b^{-1}\ast (b\ast a)} ={b^{-1}\ast
%(a\ast b)} ={b^{-1}\ast 0} =0$.
%\end{proof}
%
%\medskip
\begin{theorem} Если $a\ne 0$,
то в поле единственным образом разрешимо уравнение
$a \ast x  +  b = 0$ (решение: $x = -(a^{-1}) \ast b$).
\end{theorem}
\begin{proof}
$a\ast x +  b =0 \Impl a\ast x +  b + (-b) = 0 +
(-b)\Impl$\\
$\Impl a\ast x + \bigl(b +  (-b)\bigr) =-b \Impl
a\ast x+0 =-b \Impl a\ast x =-b\Impl$  \\ $\Impl a^{-1}\ast
(a\ast x) =a^{-1}\ast (-b)\Impl (a^{-1}\ast a)\ast
x = -(a^{-1}\ast b)\Impl 1\ast x = -(a^{-1}\ast
b) \Impl$ $\Impl x =-(a^{-1}\ast b)$.
\end{proof}


\begin{digress}
В неявном виде понятие поля возникло  давно,
в процессе выявления свойств основных операций
арифметики: сложения, вычитания, умножения и деления.
Со временем область применения этого понятия
значительно расширилась. Например, найденное Галуа необходимое
и достаточное условие разрешимости уравнения 
с одним неизвестным в радикалах
позволило доказать неразрешимость задач древности
о квадратуре круга, трисекции угла и 
удвоении куба с помощью циркуля и линейки.
В настоящее время модулярная арифметика 
и вычисления в конечных полях составляют основу криптографических алгоритмов, 
и тем самым в некоторой мере всех информационных технологий.     
\end{digress}

\end{frame}

\begin{frame}\label{2.3.4.}
\subsubsection{2.3.4. Тропическая математика}
\frametitle{2.3.4. Тропическая математика (1/5)}

Пока среди рассмотренных примеров операций в алгебраических структурах
не использовалось свойство {\color{blue}(A5)}~--- идемпотентность. Это не удивительно, поскольку 
те операции, которые рассматриваются чаше всего,
а именно сложение и умножение чисел, \em{не} идемпотенты.
Однако во второй половине XX века и особенно в XXI веке
обнаружилось, что в целом ряде важнейших приложений
удобно и целесообразно применение алгебр с идемпотентными
операциями.

\smallskip
\odmkey{Тропическая}{Математика (mathematics)!тропическая (tropical)} (правильнее говорить 
\odmkey{идемпотентная}{Математика (mathematics)!идемпотентная (idempotent)}) математика 
изучает теорию и приложения различных систем
с идемпотентными операциями.

\medskip
Пусть $M$~--- множество, на котором заданы
две бинарные операции: 
$$\Map{\oplus}{M\times M}{M},\quad 
\Map{\odot}{M\times M}{M},$$
которые по традиции называют <<сложение>> и <<умножение>>, соответственно, но обозначают специальными значками в кружках,
чтобы не путать с обычным сложением и умножением, поскольку часто употребляют тропические операции в одном контексте с обычными. 

\end{frame}

\begin{frame}
\frametitle{2.3.4. Тропическая математика (2/5)} 
Если при этом для сложения выполнены свойства 
{\color{blue}(A1)}, {\color{blue}(A2)},
{\color{blue}(A6)} c нейтральным элементом $\Bl{0}$ и {\color{blue}(A5)}, 
а для умножения также выполнены свойства {\color{blue}(A1)}, {\color{blue}(A2)},
{\color{blue}(A6)} c нейтральным элементом $\Bl{1}$ и {\color{blue}(A8)}, но на множестве $M-\Bl{0}$, и операции связаны между собой дистрибутивным законом {\color{blue}(A3)},
а также если $\A{x\in M}{x\odot \Bl{0}=\Bl{0}}$,
то алгебра $\Alg{M}{\oplus,\odot}$
называется \odmkey{идемпотентным полуполем}{Полуполе (semifield)}.


Другими словами, полуполе~--- это идемпотентный 
моноид по сложению и
абелева группа по умножению, исключая $0$, причём 
умножение дистрибутивно относительно сложения. 
Учитывая введённые ранее понятия, можно констатировать также,
что полуполе отличается от поля отсутствием обратимости 
сложения, которое идемпотентно. Это внешне небольшое отличие
радикально меняет свойства алгебры.  

\odmkey{Полуполе}{Полуполе (semifield)}~---
 это множество $M$ c двумя
бинарными операциями $ \oplus $ и $\odot$, такими, что

\begin{enumerate}
\item[1)] $\A{a,b,c \in M}{(a \oplus b) \oplus c = a \oplus (b\oplus c)}$~---
сложение ассоциативно; 

\item[2)] $\A{a,b \in M}{a \oplus b = b \oplus a}$~---
 сложение коммутативно;

\item[3)] $\E{\Bl{0} \in M}{\A{a\in M}{a \oplus \Bl{0} = \Bl{0}}}$~---
 существует нуль;

\item[4)] $\A{a}{a \oplus a = a}$~---
сложение идемпотентно; 
 
\item[5)] $\A{a,b,c \in M}{a \odot (b \odot c) = (a \odot b) \odot c}$~---
умножение ассоциативно; 
\end{enumerate}
\end{frame}

\begin{frame}
\frametitle{2.3.4. Тропическая математика (3/5)} 
\begin{enumerate}

\item[6)] $\A{a,b\in M}{a \odot b = b \odot a }$~---
умножение коммутативно;
\item[7)] $\E{\Bl{1} \in M}{\A{a\in M}{a \odot \Bl{1} =  a}}$~---
существует единица;
\item[8)] $\A{a \ne \Bl{0}}{\E{a^{-1}}{a^{-1} \odot a=\Bl{1}}}$~---
существует обратный элемент по умножению; 
\item[9)] $a \odot (b \oplus c) = (a \odot b) \oplus (a \odot c)$~---
умножение дистрибутивно относительно сложения;
\item[10)] $\A{x\in M}{x\odot \Bl{0}=\Bl{0}}$~---
умножение на нуль даёт нуль.
\end{enumerate}

\smallskip
Полуполя существуют, в том числе и числовые. Ниже приведены 
четыре важнейших примера и указаны традиционные обозначения
этих алгебр. В этих примерах $\Bbb{R}$~--- множество вещественных чисел,
$\Bbb{R}_+$~--- множество положительных вещественных чисел,
$-\infty, +\infty$~--- специальные дополнительные элементы, обладающие тем свойством, что \\$\A{x\in \Bbb{R}}{\max(x,-\infty)=x, \min(x,+\infty)= x}$. Знак $+$ означает обычное сложение, знак $\times$ означает здесь не декартово произведение,
а обычное умножение, как в школе.
$$
R_{\max,+}\Def\Alg{\Bbb{R}\cup\{-\infty\}}{\max,+},\quad
R_{\min,+}\Def\Alg{\Bbb{R}\cup\{+\infty\}}{\min,+},
$$
$$
R_{\max,\times}\Def\Alg{\Bbb{R_+}\cup\{0\}}{\max,\times},\quad
R_{\min,\times}\Def\Alg{\Bbb{R_+}\cup\{+\infty\}}{\min,\times},
$$
\end{frame}

\begin{frame}
\frametitle{2.3.4. Тропическая математика (4/5)} 
Нейтральные и обратные элементы определяются следующим образом:\\
Для $R_{\max,+}$ положим $\Bl{0}\Def -\infty$,
$\Bl{1}\Def 0$, $x^{-1}\Def -x$.\\ 
Для $R_{\min,+}$ положим $\Bl{0}\Def +\infty$,
$\Bl{1}\Def 0$, $x^{-1}\Def -x$.\\ 
Для $R_{\max,\times}$ положим $\Bl{0}\Def 0$,
$\Bl{1}\Def 1$, $x^{-1}\Def 1/x$.\\ 
Для $R_{\min,\times}$ положим $\Bl{0}\Def 0$,
$\Bl{1}\Def 1$, $x^{-1}\Def 1/x$. 

Ясно, что при таких определениях
все аксиомы тропического полуполя выполнены.

\begin{theorem}
Алгебры $R_{\max,+}$, $R_{\min,+}$, $R_{\max,\times}$, $R_{\max,+}$ изоморфны.
\end{theorem}


\begin{proof}
\vspace{-12pt}
\begin{figure}[h]
\begin{center}
\includegraphics[height=35mm]{2_3_4.jpg}
\end{center}
\end{figure}
\end{proof}
\end{frame}




\begin{frame}
\frametitle{2.3.4. Тропическая математика (5/5)}

Поскольку умножение ассоциативно, можно обычным образом определить возведение в степень $x^n = nx$, которое оказывается просто $n$-кратным сложением.
  

Если аксиома 8 не выполнена, то есть не существует обратного
элемента поумножению, то полуполе называют 
\odmkey{полукольцом}{Полукольцо (semiring)}.

\begin{example}
В полукольце $R_{\max, \min}\Def \Alg{\Bbb{R}\cup \{-\infty, +\infty\}}{\max , \min} $ нулём служит $-\infty$,
а единицей~--- $+\infty$.
\end{example}

\smallskip
Тропическая математика, которая кажется праздным умствованием,
имеет неожиданно широкую область применения. Дело в том,
что тропические аналоги традиционных формул и алгоритмов,
оказываются существенно проще. Для примера
приведём одну теорему, подтверждающую эту мысль.

\begin{theorem} {\NB{[биномиальная теорема]}}
$(a\oplus b)^n = a^n \oplus b^n $.
\end{theorem}   

\begin{proof}
Если $a \geq b$, то $(a\oplus b)^n = a^n$.
Если $a \leq b$, то $(a\oplus b)^n = b^n$.
В любом случае то $(a\oplus b)^n = \max(a^n, b^n) = a^n \oplus b^n$.  
\end{proof}
 

\begin{digress}
Одной из ранних работ явилась работа Клини
1956 года об идемпотентных полукольцах,
а эпитет <<тропическая>>, согласно легенде,
эта отрасль математики
получила в беседе двух математиков,
обсуждавших название нового направления
на пляже в Бразилии. 
\end{digress} 
\end{frame}


\begin{frame}
\frametitle{Онтология общей алгебры}
\begin{center}
\includegraphics[width=9cm]{ODM2_3.png}
\end{center}
\end{frame}




\begin{frame}
\frametitle{2.4. Элементарная теория чисел (1/4)}
\subsection{2.4. Элементарная теория чисел}

Во всех предметных областях и языках программирования используются целые числа и операции с ними (то есть, \odmkey{арифметика}{Арифметика (arithmetic)}).

\medskip
Традиционно в арифметике используют четыре операции:
сложение, умножение, вычитание и деление.   
Алгебра $\left\langle {\Bbb{Z};+,\ast} \right\rangle$
с операциями сложения и умножения является кольцом,
и этого достаточно для описания свойств сложения и умножения.
Множество целых чисел ${\mathbb Z}$ с операцией сложения 
образует абелеву группу  
(\DMref{п.~2.2.6.}{2.2.6.}), при этом 
$\A {x}{ \E {y} {x + y = 0}}$, 
то есть у любого элемента 
этого множества существует обратный по сложению элемент,
что в достаточной мере описывает свойства операции вычитания. 
Множество ${\mathbb Z}$ с операцией умножения 
образует полугруппу (моноид) 
(\DMref{п.~2.2.5}{2.2.5.}), причём обратный элемент по умножению существует лишь для чисел $1$
и $-1$. В связи с этим операция, обратная к умножению~--- деление~--- является в 
некотором смысле особенной и требует специфического рассмотрения.

\medskip
Раздел математики, изучающий целые числа и операции с ними, называется
\odmkey{теорией чисел}{Теория (theory)!чисел (of numbers)} или \odmkey{высшей арифметикой}{Арифметика (arithmetic)!высшая (higher)}.

\end{frame}

\begin{frame}
\frametitle{2.4. Элементарная теория чисел (2/4)}

Та часть теории чисел, в которой целые числа изучаются
без использования методов других разделов математики,
называется \odmkey{элементарной теорией чисел}{Теория (theory)!чисел (of numbers)!элементарная (elementary)}.  
Многие задачи элементарной теории чисел, такие как делимость целых чисел,
вычисление наибольшего общего делителя и наименьшего общего кратного, 
разложение числа на простые множители, 
вычисления по модулю, нашли в последнее время ряд важнейших приложений
в информационных технологиях.
В этом разделе бегло рассматривается элементарная теория чисел,
вводятся обозначения и устанавливаются основные факты,
необходимые в приложениях. 
Сами приложения рассматриваются в других курсах 
по информационным технологиям. 

\begin{remark}
В изложении рассматриваются натуральные числа $\Bbb{N}$ и ноль  
(напомним обозначение $\Bbb{N}_0\Def \Bbb{N}+0$) 
только ради упрощения формулировок и доказательств.  
Все определения, утверждения и алгоритмы этого раздела легко распространяются 
на множество целых чисел $\Bbb{Z}$ за счёт учёта знака числа.
\end{remark}

 
В соответствии с традициями, в теории чисел 
знак операции умножения $\ast$ опускается, а запись $x+(-y)$ упрощается до $x-y$.
\end{frame}

\begin{frame}
\frametitle{2.4. Элементарная теория чисел (3/4)}
\begin{tabular}{p{12mm}|p{40mm}|p{90mm}}
  № & Название параграфа
  & Ключевые термины и обозначения \\
  \hline
  & & \\
\DMref{2.4.1.}{2.4.1.}\RS
 & Делимость чисел &
{делимое, делитель, частное, неполное частное, остаток; собственные и несобственные делители,
признаки делимости, совершенные числа} 
\\
{\DMref{2.4.2.}{2.4.2.}} & Наибольший общий делитель
и наименьшее общее кратное & {взаимная простота; алгоритм Евклида}\\  
{\DMref{2.4.3.}{2.4.3.}} & Линейное представление 
наибольшего общего делителя &
{ диофантовы уравнения; коэффициенты Безу }\\ 
{\DMref{2.4.4.}{2.4.4.}}\RS & Простые числа &
{простое число, составное число; решето Эратосфена;
 \TODO{ алгоритмы разложения на множители}}\\  
\end{tabular}  

\end{frame}

\begin{frame}
\frametitle{2.4. Элементарная теория чисел (4/4)}
\begin{tabular}{p{12mm}|p{40mm}|p{90mm}}
  № & Название параграфа
  & Ключевые термины и обозначения \\
  \hline
  & & \\
{\DMref{2.4.5.}{2.4.5.}} & Сравнения &
{сравнение по модулю $\Mod{a}{b}{m}$}\\
\DMref{2.4.6.}{2.4.6.} & Системы вычетов &
{классы вычетов, полная система вычетов } \\
{\DMref{2.4.7.}{2.4.7.}}\RS & Китайская теорема об остатках &
{}\\
{\DMref{2.4.8.}{2.4.8.}} & Вычисления в остаточных классах &
{модулярная арифметика}\\
{\DMref{2.4.9.}{2.4.9.}} & Функция Эйлера &
{аддитивные и мультипликативные функции}
\end{tabular}  

\end{frame}


\begin{frame}
\label{2.4.1.}\subsubsection{2.4.1. Делимость чисел }\frametitle{2.4.1. Делимость чисел (1/5)}

Число $a\in \Bbb{N}$ \odmkey{делится на}{Делится на (divisible by)} 
число $b\in \Bbb{N}$ (обозначение $b|a$),  
если $\E{q\in \Bbb{N}}{a=bq}$.
Число $a$ называется \odmkey{кратным}{Число (number)!кратное (multiple)} 
числу $b$, число $b$ называется \odmkey{делителем}{Делитель (divisor)} числа $a$,
а число $q$ называется \odmkey{частным}{Частное (quotient)} от деления  $a$ на $b$. 


\medskip
\begin{theorem}
$\A{a,b \in \Bbb{N}}{\E{ q\in\Bbb{N}_0,r \in \Bbb{N}_0}{ a = bq + r \And 0 \le r < b}}$.
\end{theorem}
\begin{proof} Следующий элементарный \odmkey{алгоритм деления вычитанием}{Алгоритм (algorithm)!деления вычитанием (of division by subtraction)}
находит числа $q$ и $r$.
$q\Assign 0;\ r\Assign a;$
\textbf{while  $r\ge b$ do $r\Assign r-b;$ $ q\Assign q+1$ end while}.
\end{proof}

\medskip
\begin{corollary}
Если $a = bq + r \And 0 \le r < b$, то числа $q$ и $r$ единственны.
\end{corollary} 
\begin{proof}
Пусть $a=bq_1+r_1, 0 \le r_1 < b, a=bq_2+r_2, 0 \le r_2 < b$. \\
Тогда $0 = (bq_1 + r_1) - (bq_2 + r_2) = b(q_1-q_2)+(r_1-r_2)$.
Если $q_1 \ne q_2$, то $|b(q_1-q_2)|\ge b$, но $|r_1-r_2|<b$,
и равенство невозможно. 
Следовательно, $q_1 = q_2$, откуда $r_1-r_2=0$ и $r_1=r_2$.
\end{proof}

\medskip
Число $q$  
называется \odmkey{неполным частным}{Неполное частное (quotient)}, 
а число $r$ называется 
\odmkey{остатком}{Остаток от деления (remainder)} от деления числа $a$ на число $b$.
Если $b|a$, то $r=0$. 

\end{frame}

\begin{frame}
\frametitle{2.4.1. Делимость чисел (2/5)}

%\vspace{-0.25\baselineskip}

\begin{remark}
Для вычисления неполного частного и остатка (деление с остатком) 
обычно определяются специальные операции.
Здесь они обозначаются
$a\DIV b$ и $a\MOD b$ соответственно.
\end{remark} 

Делить с остатком можно и отрицательные числа. Действительно,
пусть $a, b, q, r \in \Bbb{N}_0$, $a=bq+r$, $0\le r<b$.
Тогда если $b|a$, то $r=0$ и $-a=b(-q)$, а в противном случае
$-a=-bq-r=b(-(q+1))+(b-r)$, причём остаток $0\le b-r<b$. 

Множество делителей числа $a\in\Bbb{N}$ обозначается 
$$D(a)\Def\Set{b\in\Bbb{N}}{b|a}.$$
Очевидно, что любое число делится на 1 и на себя. Эти делители называются 
\odmkey{тривиальными}{Делитель (divisor)!тривиальный (trivial)}. Прочие делители, если они есть,
называются \odmkey{нетривиальными}{Делитель (divisor)!тривиальный (trivial)}, или \odmkey{собственными}{Делитель (divisor)!собственный (proper)}.


\end{frame}

\begin{frame}
\frametitle{2.4.1. Делимость чисел (3/5)}
\begin{digress} 
Определить делители произвольного числа не так просто.
Для некоторых делителей известны \odmkey{признаки делимости}{Признаки (criteria)!делимости (divisibility)},
то есть сравнительно простые алгоритмы, позволяющие 
ответить на вопрос, является ли число $b$ делителем числа $a$.
Читателю, несомненно, известны со школы признаки делимости на 2, 3, 5, 10.

Все признаки делимости следуют общей схеме: cтроится убывающая 
последовательность $\Tuple{a}{0}{n}$, такая что $a_0=a$ и
$b|a_i\Equiv b|a_{i+1}$. Однако,
все признаки делимости индивидуальны для каждого делителя,
существенно зависят от системы счисления и построены 
далеко не для всех делителей.
В реальных компьютерных вычислениях признаки делимости не используются.
\end{digress}

\medskip
В общем случае деление в компьютере является сравнительно трудоёмкой операцией.
В большинстве арифметических устройств компьютеров деление выполняется
примерно в два раза дольше умножения, 
а умножение выполняется примерно в два раза дольше сложения.

\end{frame}

\begin{frame}
\frametitle{2.4.1. Делимость чисел (4/5)}

\odmkey{Отношение делимости}{Отношение (relation)!делимости (of divisibility)} рефлексивно 
$\A{a \in \Bbb{N}}{a|a}$,
транзитивно \\ $\A{a, b, c \in \Bbb{N}}{a|b \And b|c \Impl a|c}$ и 
антисимметрично $\A{a,b \in \Bbb{N}}{a|b \And b|a \Impl a=b}$.
Таким образом, делимость чисел~--- 
это нестрогий частичный порядок на множестве $\Bbb{N}$.

\begin{example}
На рисунке  представлена 
диаграмма Хассе отношения делимости  
для чисел от 1 до 12. Хорошо видно, 
что делимость имеет <<нерегулярный>> характер.
               
\begin{center}
\includegraphics[height=48mm]{2_4_1.jpg}
\end{center}
\end{example}

%\vspace{\baselineskip}

\end{frame}

\begin{frame}
\frametitle{2.4.1. Делимость чисел (5/5)}

\begin{digress}
С отношением делимости связано множество
любопытных и поучительных числовых казусов и фокусов.
Например, ещё в древности были обнаружены так называемые
\odmkey{совершенные}{Число (number)!совершенное (perfect)} числа, то есть числа, 
равные сумме своих делителей, считая 1, но не считая само число. 
Первые совершенные числа:

$6 = 1+2+3$,

$28  = 1+2+4+7+14$, 

$496=1+2+4+8+16+31+62+124+248$,\\...

Все обнаруженные совершенные числа являются чётными и треугольными.
До сих пор неизвестно, конечно ли множество совершенных чисел и
существуют ли нечётные совершенные числа.  
\end{digress}

\end{frame}

\begin{frame}
\label{2.4.2.}\subsubsection{2.4.2. Наибольший общий делитель и наименьшее общее кратное}\frametitle{2.4.2. Наибольший общий делитель 
и наименьшее общее кратное (1/6)}

Натуральное число $c$ такое, что $c|a$ и 
$c|b$, называется \odmkey{общим делителем}{Делитель (divisor)} чисел $a$ и $b$.\\
Множество общих делителей чисел $a$ и $b$ 
обозначим $D(a,b)$.
Наибольшее из таких чисел называется \odmkey{наибольшим общим делителем}{Наибольший общий делитель (greatest common divisor)} чисел $a$ и $b$.
Обозначение: $\gcd(a,b)$, НОД$(a,b)$ 
или просто $(a,b)$.

Для обозначения наибольшего общего 
делителя здесь используется функция $\gcd(a,b)$.
Если $\gcd(a,b)=1$, то числа $a$ и $b$ 
называются \odmkey{взаимно простыми}{Числа (numbers)!взаимно простые (coprime)}. Обозначение: $a \perp b$.

\begin{columns}[t] % contents are top vertically aligned
\begin{column}[T]{6cm} % each column can also be its own environment
\begin{example} Из определения вытекает, что отношение взаимной простоты
антирефлексивно (для всех чисел, кроме единицы) 
 и симметрично.
На рисунке представлена часть графика отношения взаимной простоты.
\end{example}
\end{column}
\begin{column}[T]{8cm} % alternative top-align that's better for graphics
\includegraphics[height=45mm]{2_4_2.jpg}
\end{column}
\end{columns}

\end{frame}

\begin{frame}
\frametitle{2.4.2. Наибольший общий делитель 
и наименьшее общее кратное (2/6)}

\begin{lemma}
Если $b|a$, то множество общих делителей чисел $a$ и $b$ совпадает с множеством делителей $b$: $b|a \Impl D(a,b) = D(b)$.
\end{lemma}

\begin{proof}
Любой общий делитель чисел $a$ и $b$, очевидно, является делителем $b$. Любой делитель $b$ делит также и $a$ в силу транзитивности.
\end{proof}


\begin{corollary}
$b|a \Impl \gcd(a,b)=b$.
\end{corollary}

\medskip
\begin{lemma}
Если $a = bq+r$, то множество общих делителей чисел $a$ и $b$ совпадает с множеством общих делителей $b$ и $r$: 
$a = bq+r \Impl D(a,b) = D(b,r).$
\end{lemma}


\begin{proof}
Пусть $x \in D(a,b)$. Тогда $\E{k,l}{a=kx \And b=lx}$ и $kx=lxq+r$, \\ $r=(k-lq)x$, следовательно $x|r$, $x \in D(b,r)$.
Пусть $y \in D(b,r)$. Тогда
\\ $\E{m,n}{b=my \And r=ny}$ и $a=myq+ny=(mq+n)y$, следовательно, $y|a$, $y \in D(a,b)$.
Таким образом, $D(a,b) = D(b,r)$.
\end{proof}

\begin{corollary}
$a=bq+r \Impl \gcd(a,b)=\gcd(b,r)$.
\end{corollary}

\begin{example}
$60=7\cdot 8 +4$,  $\gcd(60,8)=\gcd(8,4)=4$, $\gcd(60,7)=\gcd(7,4)=1$. 
\end{example}

\end{frame}

\begin{frame}
\frametitle{2.4.2. Наибольший общий делитель 
и наименьшее общее кратное (3/6)}

Доказанные утверждения лежат в основе алгоритма нахождения наибольшего 
общего делителя двух натуральных чисел $a$ и $b$, известного как \odmkey{алгоритм Евклида}{Алгоритм (algorithm)!Евклида (Euclid)}. В силу симметричности $\gcd(a,b)$, не ограничивая общности,
можно считать, что $a\ge b$.


\begin{algorithmic}
\REQUIRE $a,b$ --- натуральные числа, $a \ge b$.
\ENSURE значение $\gcd(a,b)$ в переменной $a$.
%\IF {$a<b$}
%\STATE $a :=: b$ \COMMENT{поменяли числа местами}
%\ENDIF
\WHILE{$b > 0$}
\STATE {$r \Assign a \MOD b;\ a \Assign b;\ b \Assign r$}
\ENDWHILE
%\RETURN {$a$}
\end{algorithmic}

\vspace{-2pt}
\begin{ground}
На каждом шаге алгоритма 
пара $(a,b)$ замещается парой $(b,r)$,
где $r$ --- остаток от деления $a$ на $b$.
Согласно следствию, 
после каждой итерации цикла искомый наибольший общий делитель 
остаётся неизменным, т.~е. является инвариантом алгоритма. 
Когда $a$ станет кратно $b$, $r$ будет равно нулю. Это означает, что искомый наибольший общий делитель найден.
Цикл обязательно завершится, потому что числа $a$ и $b$ не могут уменьшаться бесконечно.
\end{ground}

\begin{example}
$\gcd(616,286) = \gcd(286,44) = \gcd(44,22) = 22$.
\end{example}

\end{frame}

\begin{frame}
\frametitle{2.4.2. Наибольший общий делитель 
и наименьшее общее кратное (4/6)}

\begin{digress}
Алгоритм Евклида даёт повод сделать некоторые замечания относительно
\odmkey{стиля программирования}{Стиль программирования (coding standard)}, то есть
набора необязательных правил и соглашений, 
направленных на \em{улучшение} кода.
В эту категорию входят использование 
\odmkey{дисциплины имён}{Дисциплина имён (naming convention)},     
\odmkey{отступов}{Отступы (tabbing)},
\odmkey{комментариев}{Комментарии (comment)} и 
множества других полезных приёмов.
Важно понимать, что использование того или иного 
стиля программирования улучшает код не абсолютно,
а относительно тех целей, которые преследуют выбранные правила.               
Рассмотрим в качестве примера ещё одну реализацию алгоритма Евклида.

\begin{algorithmic}
\REQUIRE $a,b$ --- натуральные числа.
\ENSURE значение $\gcd(a,b)$ есть сумма $a + b$.
\WHILE{$a\neq0 \And b\neq0$}
\STATE \textbf{if $a>b$ then 
$a \Assign a \MOD b$ else $ b \Assign b \MOD a$ end if}
\ENDWHILE
\end{algorithmic}
\end{digress}
\end{frame}

\begin{frame}
\frametitle{2.4.2. Наибольший общий делитель 
и наименьшее общее кратное (5/6)}

\begin{digress} (Продолжение)
С одной стороны, этот вариант несколько эффективнее по времени,
поскольку на каждом шаге цикла выполняется одно присваивание
вместо трёх,
и по памяти, поскольку используются
две переменные вместо трёх. С другой стороны, этот код чуть длиннее,
значительно дальше от доказанных математических формул 
и потому требует существенного дополнительного обоснования,
а самое главное, в этом варианте нарушается 
<<правило хорошего тона>>, гласящее, что каждый объект 
(в частности, переменная) в программе 
при изменении состояния (значения) не должен менять
свой <<смысл>> или свою <<ответственность>>.
\end{digress}

\smallskip
\begin{remark}
Антирефлексивное и симметричное отношение
называют иногда отношением 
\odmkey{строгой толерантности}{Отношение (relation)!толерантности (of tolerance)!строгой (strict)}
по аналогии с отношением 
толерантности, см. \DMref{п.~1.7.5.}{1.7.5.} 
\end{remark}

\begin{example}
Отношение взаимной простоты чисел является отношением
строгой толерантности.
\end{example}

\end{frame}

\begin{frame}
\frametitle{2.4.2. Наибольший общий делитель 
и наименьшее общее кратное (6/6)}

Число $d$ 
называется \odmkey{общим кратным}{Общее кратное (common multiple)} чисел $a$ и $b$, если $a|d$ и 
$b|d$.
Множество общих кратных  $K(a,b)\Def\Set{d\!\in\!\Bbb{N}}{a|d \And b|d}$.
Наименьшее из таких чисел называется \odmkey{наименьшим общим кратным }{Наименьшее общее кратное (least common multiple)} $a$ и $b$.
Обозначение: $\lcm(a,b)$, НОК$(a,b)$ или $[a,b]$.

\begin{remark}
В отличие от конечного множества $D(a,b)$, множество $K(a,b)$ является бесконечным. Однако $K(a,b)\subset \Bbb{N}$, 
поэтому $K(a,b)$ имеет наименьший элемент 
и определение наименьшего общего кратного корректно.
\end{remark}

\begin{theorem}
Для любых натуральных чисел $a$, $b$ и $m$ справедливо:
\begin{enumerate}
\item[1)] $a|m \And b|m \Impl \lcm(a,b)|m$;
\item[2)] $\gcd(a,b)\lcm(a,b)=ab$.
\end{enumerate}
\end{theorem}

\begin{proof}
Пусть $d \Assign \gcd(a,b)$, тогда $a=a_1d, b=b_1d$. Заметим, что $a_1 \perp b_1$. Действительно, от противного. Пусть для определенности $a_1|b_1$, то есть существует такой $q$, что $q a_1=b_1$.
Значит, $b=b_1, d=q a_1, d=q a$. Очевидно, что $a>d$, значит $d$~--- не НОД. Противоречие.
Пусть теперь $m = ak$. Тогда $ak/b = a_1k/b_1 \in \Bbb{N}$, 
следовательно $b_1|k$, т.е. $\E{ t \in \Bbb{N}}{k = b_1 t}$. 
Отсюда $m=ak=ab_1t=abt/d$.
С~другой стороны,  любое число вида $abt/d\in K(a,b)$.
Наименьшее общее кратное достигается при $t = 1$, т.е. равно 
$ab/d$. 
\end{proof}
\end{frame}

\begin{frame}
\label{2.4.3.}\subsubsection{2.4.3. Линейное представление наибольшего общего делителя}
\frametitle{2.4.3. Линейное представление наибольшего общего делителя (1/4)}

Рассмотрим протокол применения алгоритма 
Евклида в двух экивалентных формах.
\begin{columns}[t]
\begin{column}{7cm}
$a=bq_0+r_1$

$b=r_1q_1+r_2$

$r_1=r_2q_2+r_3$

$\dots$

$r_{n-2}=r_{n-1}q_{n-1}+r_n$
\end{column}
\begin{column}{7cm}
$r_1=a - bq_0$

$r_2=b - r_1q_1$

$r_3=r_1 -r_2q_2$

$\dots$

$r_n =r_{n-2} - r_{n-1}q_{n-1}$
\end{column}
\end{columns}

\medskip
При этом $r_n$ является последним ненулевым остатком, то есть $r_{n-1}=r_nq_n$.

\medskip
Подставляя полученные равенства снизу вверх, 
получаем
$$\gcd(a,b) = r_n= r_{n-2}- r_{n-1}q_{n-1} = 
r_{n-2}- (r_{n-3}- r_{n-2}q_{n-2})q_{n-1} = 
\dots = ax+by.$$

Эти наблюдения позволяют сформулировать теорему.
\begin{theorem}
$\A{a,b \in \Bbb{N}}{\E{x,y\in \Bbb{Z}}{\gcd(a,b)=ax+by}}.$
\end{theorem}

Такое представление наибольшего общего делителя
называется \odmkey{соотношением Безу}{Соотношение (identity)!Безу (Bezout)}, или \em{линейным представлением}, 
а числа $x$ и $y$ называются
\odmkey{коэффициентами Безу}{Коэффициент (coefficient)!Безу (Bezout)}
\end{frame} 

\begin{frame}
\frametitle{2.4.3. Линейное представление наибольшего общего делителя (2/4)}

Как найти коэффициенты Безу?

\medskip
Алгоритм Евклида строит последовательность
чисел $a>b>r_1>r_2>\dots>r_n>0$. Оказывается, 
не только последний остаток $r_n=\gcd(a,b)$,
но и все числа в этой последовательности
выражаются в виде линейных комбинаций
чисел $a$ и $b$, то есть $r_k=ax_k+by_k$. 
Коэффициенты разложения некоторого числа 
в линейную комбинацию чисел $a, b$, то есть
выражение $ax+by$, запишем в векторной форме $(x,y)$.
В таких обозначениях $a=(1,0)$, $b=(0,1)$.
Операции с векторами, как обычно, выполняются покомпонентно.

\medskip
Рассмотрим ещё раз протокол работы алгоритма Евклида,
подставляя формулы сверху вниз и в векторной форме:
 
\begin{columns}[t]
%\begin{column}{3cm}
%$r_1=a - bq_0$
%
%$r_2=b - r_1q_1$
%
%$r_3=r_1 -r_2q_2$
%
%
%$\dots$
%
%\end{column}
\begin{column}{8cm}
$r_1= a\cdot 1 +b(-q_0)$

$r_2= a(-q_1)+b(1+q_0q_1)$

$r_3= a(1+q_1q_2)+b(-q_0-q_2+q_0q_1q_2)$

$\dots$

\end{column}
\begin{column}{7cm}
$(1,-q_0) = (1,0)-q_0(0,1)$

$(-q_1,1+q_0 q_1)= (0,1)-q_1(1,-q_0)$

$(1+q_1 q_2, -q_0-q_2-q_0 q_1q_2)=$\\$=(1,-q_0)-q_2(-q_1,1+q_0 q_1)$

$\dots$

\end{column}

\end{columns}

\end{frame}

\begin{frame}
\frametitle{2.4.3. Линейное представление наибольшего общего делителя (3/4)}
\begin{lemma}
В указанных обозначениях
$
(x_k, y_k) = (x_{k-2}, y_{k-2}) - q_{k-1}(x_{k-1}, y_{k-1})
$.
\end{lemma}

\begin{proof}
По индукции. База индукции проверена выше.
Индукционный переход.
Пусть утверждение верно для всех чисел, меньших $k$. Рассмотрим $(x_{k},y_{k})$.
Имеем 
\begin{align*}
r_{k}&=ax_k+by_k =r_{k-2} - q_{k-1}r_{k-1} = \\
&=ax_{k-2}+by_{k-2} - q_{k-1}(ax_{k-1}+by_{k-1})=\\
&=a(x_{k-2}-q_{k-1}x_{k-1})+b(y_{k-2}-q_{k-1}y_{k-1}) .
\end{align*}
и значит $x_k=x_{k-2}-q_{k-1}x_{k-1}$, 
$y_k=y_{k-2}-q_{k-1}y_{k-1}$.
\end{proof}

%\begin{remark}
%Можно показать, что коэффициенты последовательно
%вычисляются друг через друга по реккурентной формуле
%$z_{i+1} = z_{i-1}- q_{n+1-i}z_i$, причём
%$z_0=0, z_1=1$. 
%Два последних элемента в этой последовательности
%дают коэффициенты Безу,
%то есть $\gcd(a,b)= az_{n+1}+bz_{n+2}$
%\end{remark}

\vspace{12pt}
\begin{columns}[t]
\begin{column}{7cm}

\vspace{-42pt}
\begin{example}
Найдём линейное представление
наибольшего общего делителя чисел $616$ и $286$.
Получаем\\
$22=-6\cdot 616 + 13\cdot 286$
\end{example}
\end{column}
\begin{column}{7cm}
\begin{tabular}{r|r|r|r|r}
$i$ & $-1$ &$0$ &$1$ &$2$  \\
\hline
$r$ & $616$ &$286$ &$44$ &$22$  \\
\hline
$q$ &  &$2$ &$6$ &$2$ \\
\hline
$x$ & $1$ &$0$ &$1$ &$-6$ \\
\hline
$y$ & $0$ &$1$ &$-2$ &$13$ \\
\hline
\end{tabular}
\end{column}
\end{columns}

\end{frame}

\begin{frame}
\frametitle{2.4.3. Линейное представление наибольшего общего делителя (4/4)}
Линейные уравнения, в которых все числа целые,
как известные коэффициенты, так и неизвестные величины, называются \odmkey{диофантовыми}{Уравнение (equation)!диофантово (diophantine)}.

\medskip
Мы уже установили, что уравнение
$ax+by=\gcd(a,b)$
всегда разрешимо в целых числах и решениями
оказываются коэффициенты Безу.

\medskip
\begin{theorem}
Уравнение с целыми коэффициентами 
$ax+by=c$
разрешимо в целых числах тогда и только тогда, когда
$\gcd(a,b)|c$.
\end{theorem}

\begin{proof}
Обозначим $d=\gcd(a,b)$. 
Левая часть $ax+by$ делится на $d$,
значит уравнение разрешимо тогда и только тогда,
когда $d|c$. Пусть $c=de$. Уже доказано,
что уравнение $ax+by=d$ всегда имеет решение.
Обозначим это решение $x_0, y_0$.
Тогда $x_1=x_0e, y_1=y_0e$~--- решение
уравнения $ax+by=c$.
\end{proof}

\medskip
\begin{remark}
Решения уравнения $ax+by=c$ 
описываются формулами
$$ x= x_1-\frac{b}{d}k,\quad
y= y_1+\frac{a}{d}k, \quad k\in \Bbb{Z},$$
где $x_1, y_1$~--- любое частное решение.
\end{remark}
\end{frame}

\begin{frame}
\label{2.4.4.}\subsubsection{2.4.4. Простые числа }\frametitle{2.4.4. Простые числа (1/7)}

Натуральное число $a > 1$ называется \odmkey{простым}{Число (number)!простое (prime)}, если оно не имеет собственных 
делителей. В противном случае число $a$ называется \odmkey{составным}{Число (number)!составное (composite)}.

\begin{remark}
Число 1 имеет только один делитель, поэтому его не относят 
ни к простым, ни к составным.
\end{remark}

Множество простых чисел обозначается $\Bbb{P}$.

\begin{theorem}
Для любого натурального числа $a > 1$ его наименьший делитель, 
отличный от 1, является простым числом.
\end{theorem}

\begin{proof}
От противного. Пусть этот делитель составной, тогда у него тоже есть делитель, меньший его и отличный от 1. Но тогда он делит и $a$, а значит, исходный делитель не был наименьшим. 
\end{proof}

\begin{corollary}
Для составного числа $a > 1$ его наименьший делитель, отличный от 1, 
не превосходит $\sqrt{a}$.
\end{corollary}

\begin{proof}
Пусть $p > 1$ --- наименьший делитель $a$. Тогда $a=pq$, где $q \geq p$. Следовательно, $a \geq p^2$.
\end{proof}
\end{frame}

\begin{frame}
\frametitle{2.4.4. Простые числа (2/7)}
\begin{theorem}
Простых чисел бесконечно много.
\end{theorem}

\begin{proof} 
От противного. Пусть  $\vert \Bbb{P}\vert=n$. 
Тогда  $\Bbb{P}=\{\Tuple{p}{1}{n}\}$. 
Cоставим число $p \Assign p_1p_2 \dots p_n + 1$. Это число не делится ни 
на одно из чисел $\Tuple{p}{1}{n}$, так как даёт остаток 1 при 
делении на каждое из них, и не совпадает ни с одним из них. Следовательно, 
число $p$ само простое. Противоречие.
\end{proof}

\medskip
Для построения множества простых чисел, не превосходящих данного числа $n$, используют простой способ, называемый \odmkey{решетом Эратосфена}{Решето Эратосфена (Sieve of Eratosthenes)}.
\end{frame}

\begin{frame}
\frametitle{2.4.4. Простые числа (3/7)}


\begin{algorithmic}
\REQUIRE натуральное число $n > 1$.
\ENSURE множество простых чисел, не превосходящих $n$.
\STATE{$B \Assign \left\{ 2,3,4,5 \ldots ,n \right\}$}
\WHILE{$B \ne \emptyset$}
	\STATE {$x \Assign \min(B)$} 
	\COMMENT {\color{gray}\small  наименьший элемент $B$}
	\STATE {\textbf{yield} $x$} \COMMENT {\color{gray}\small выдать $x$}
	\STATE {$B \Assign B - x$} \COMMENT {\color{gray}\small удалить $x$ из $B$}
	\STATE {$y \Assign x^2$}
	\WHILE {$y \le n$}
		\STATE {$B \Assign B - y$; $y \Assign y + x$}
		\COMMENT {\color{gray}\small удалить из $B$ все числа, кратные $x$}
		
	\ENDWHILE
\ENDWHILE

\end{algorithmic}


\begin{ground}
На каждом шаге алгоритма минимальный элемент $x \in B$~--- простое число. 
В противном случае 
у него был бы простой делитель $p < x$, 
и $x$ был бы удалён из $B$ ранее, как кратный $p$. Для 
простого $x$ достаточно удалить из $B$ все кратные ему, начиная с $x^2$, т.к. все меньшие составные, 
кратные $x$, были удалены из $B$ ранее.
\end{ground}

\end{frame}

\begin{frame}
\frametitle{2.4.4. Простые числа (4/7)}

\begin{example} Найдём с помощью этого алгоритма все простые числа, не 
превосходящие 100. В правой нижней ячейке  список простых чисел.
\end{example}


\begin{table}[htbp]
\begin{tabular}{|p{5pt}|p{185pt}|p{205pt}|}
\hline
$x$ & 
Вычеркнутые составные числа& 
Полученные | Оставшиеся числа \\
\hline
2& 
4,6,8,10,12,14,16,18,20,22,24,26,28,\par  30,32,34,36,38,40,42,44,46,48,50,52,\par
54,56,58,60,62,64,66,68,70,72,74,76,\par 78,80,82,84,86,88,90,92,94,96,98,100& 
2 | 3,5,7,9,11,13,15,17,19,21,23,25,27,29,31,\par 33,35,37,39,41,43,45,47,49,51,53,55,57,59,\par  61,63,65,67,69,71,73,75,77,79,81,83,85,87,\par
89,91,93,95,97,99 \\
\hline
3& 
9,15,21,27,33,39,45,51,57,63, \par 69,75,81,87,93,99& 
2,3 | 5,7,11,13,17,19,23,25,29,31, \par 35,37,41,43,47,49,53,55,59,61,65, \par 67,71,73,77,79,83,85,89,91,95,97 \\
\hline
5& 
25,35,55,65,85,95& 
2,3,5 | 7,11,13,17,19,23,29, \par 31,37,41,43,47,49,53,59,61,67, \par 71,73,77,79,83,89,91,97 \\
\hline
7& 
49,77,91& 
2,3,5,7 | 11,13,17,19,23,29,31,37,41,43, \par 47,53,59,61,67,71,73,79,83,89,97 \\
\hline
\end{tabular}
\end{table}


\end{frame}

\begin{frame}
\frametitle{2.4.4. Простые числа (5/7)}

\begin{lemma}
 $\A{a,b,c \in \Bbb{N}}{b \perp c\Impl\gcd(ac,b) = \gcd(a,b)}$.
\end{lemma}

\begin{proof}
Рассмотрим $\gcd(ac,b)$. Так как $b \perp c$, то $\gcd(ac,b)$ не имеет общих делителей с $c$, отсюда $\gcd(ac,b)=\gcd(a,b)$
\end{proof}

\smallskip
\begin{corollary}
 $\A{a,b,c \in \Bbb{N}}{ b \perp c \And b|ac \Impl b|a}$.
\end{corollary}

\smallskip
\begin{lemma} {\NB{[Лемма Евклида]}}
$\A{a,b \in \Bbb{N}, p \in \Bbb{P}}{p|ab \Impl p|a \Or p|b}$.
\end{lemma}

\begin{proof}
 От противного. Пусть ни $a$, ни $b$ не делятся на $p$. Тогда из простоты $p$ 
следует, что $a \perp p$, $b \perp p$. По лемме $\gcd(ab,p) = \gcd(a,p) = 1$, $ab \perp p$~--- противоречие.
\end{proof}

\smallskip
\begin{corollary}
$\A{a_1, a_2, \dots ,a_n \in \Bbb{N}, p \in \Bbb{P}}{p|a_1a_2 \dots a_n
\Impl\E{i}{p|a_i}}$.
\end{corollary}

\smallskip
\begin{theorem} {\NB{[Основная теорема арифметики]}}
Любое натуральное число, большее 1, можно представить в виде произведения простых множителей, причём единственным образом.
\end{theorem}



\end{frame}

\begin{frame}
\frametitle{2.4.4. Простые числа (6/7)}

\begin{proof}
\proofsection{Существование} От противного. 
Пусть $1< n \in \Bbb{N}$~--- наименьшее 
натуральное число, не представимое в виде произведения нескольких простых 
чисел.  Тогда если
число $n$~--- простое, то оно является своим разложением. 
Если же число $n$ --- составное, то $n=n_1n_2$, где $n_1 <n$, $n_2 <n$, причём каждое из $n_1$, $n_2$ имеет разложение, а значит, можно разложить и исходное число $n$. Противоречие.

\proofsection{Единственность} Пусть $n$ имеет два различных 
разложения на простые множители. Рас\-смотрим любое число $p$ в одном из разложений.
Тогда произведение чисел второго разложения делится на $p$ и по следствию из леммы Евклида одно из простых чисел второго разложения совпадает с $p$. 
Сократим оба разложения на $p$ и применим к оставшимся разложениям такое же рассуждение. Так по одному сократим все числа в первом разложении, 
найдя и сократив равные им во втором. Процесс закончится ввиду конечности разложений.
\end{proof}





\end{frame}

\begin{frame}
\frametitle{2.4.4. Простые числа (7/7)}

\begin{corollary}
Пусть $a = p_{1}^{a_{1}}p_{2}^{a_{2}}\cdots p_{n}^{a_{n}}$, $b = p_{1}^{b_{1}}p_{2}^{b_{2}}\cdots p_{n}^{b_{n}}$, $p_{1},p_{2},\ldots ,p_{n}\in \mathbb P$, \\$a_{1},a_{2},\ldots ,a_{n},b_{1},b_{2},\ldots ,b_{n}\ge 0$. Тогда $\gcd (a,b) = p_{1}^{c_{1}}p_{2}^{c_{2}}\cdots p_{n}^{c_{n}}$, $\lcm (a,b) = p_{1}^{d_{1}}p_{2}^{d_{2}}\cdots p_{n}^{d_{n}}$, где $c_{i} = \min (a_{i},b_{i})$, $d_{i} = \max (a_{i},b_{i})$.
\end{corollary}

\medskip
\begin{example}
$24 = 2^3 3^1, 36 = 2^2 3^2, \gcd(24,36)=2^2 3^1 = 12, \lcm(24,36)=2^3 3^2 =72$. 
\end{example}


\medskip
\begin{digress}
Разложение натурального числа на простые множители 
является одной из древнейших математических задач,
которая остаётся актуальной и сегодня.
Для сходной задачи проверки простоты числа
уже давно найдены достаточно эффективные алгоритмы
(см.~решето Эратосфена), а для разложения числа на множители пока не найдено
достаточно эффективных с практической точки зрения алгоритмов,
но при этом не доказано, что таких алгоритмов не существует.  
Между тем, очень многие современные информационные технологии,
например, шифрование с открытым ключом, критически 
зависят от наличия или отсутствия этих алгоритмов.
\end{digress}

\end{frame}

\begin{frame}
\label{2.4.5.}\subsubsection{2.4.5. Сравнения }\frametitle{2.4.5. Сравнения (1/3)}

Пусть числа $a, b, m \in \Bbb{N}$. Говорят, что число
$a$ \odmkey{сравнимо с}{Сравнение (comparison)} числом $b$ \odmkey{по модулю}{Сравнение (comparison)!по модулю (modulo)} $m$, если оба 
числа имеют одинаковый остаток при делении на $m$, то есть
$a=pm+r$, $b=qm+r$.
Обозначение: $\Mod{a}{b}{m}$.

\smallskip
\begin{theorem} Пусть $a, b\in \Bbb{N}, a>b$.
Число $a$ сравнимо с числом $b$ по модулю $m$ тогда и только тогда, когда $m$ 
является делителем $a-b$:

\vspace{-0.75\baselineskip}
$$\Mod{a}{b}{m} \Equiv m|(a-b).$$
\end{theorem}


\begin{proof}
\proofsection{$\Mod{a}{b}{m} \Impl m|(a-b)$}
Пусть $a=pm+r$, $b=qm+r$.\\ 
Тогда $a-b = (p - q)m + (r - r)= km$ и $m|(a-b)$. 
\proofsection{$ m|(a-b) \Impl  \Mod{a}{b}{m}$}
Пусть $a-b=km$. \\Тогда 
$\E{p,q, r_a, r_b \in \Bbb{N}_0}{a=pm+r_a \And b=qm+r_b \And r_a<m \And r_b<m}$. \\
Имеем $a-b = (p - q)m + (r_a - r_b)$, откуда $r_a-r_b=0$  и $r_a=r_b=r$. 
\end{proof}

\end{frame}

\begin{frame}
\frametitle{2.4.5. Сравнения (2/3)}

\begin{corollary}
Отношение $\Mod{a}{b}{m}$ является отношением эквивалентности на множестве 
$\Bbb{N}$.
\end{corollary}
\begin{proof}
\proofsection{Рефлексивность} Остаток определён однозначно.
\proofsection{Симметричность} Определение сравнимости не зависит от порядка.
\proofsection{Транзитивность} В силу транзитивности равенства остатков.
\end{proof}


\smallskip
Рассмотрим основные свойства введённого отношения сравнимости.

\begin{lemma} \NB{ [1]}
Пусть $\Mod{a}{b}{m}$ и $\Mod{c}{d}{m}$.
Тогда $\Mod{(a+c)}{(b+d)}{m}$.
\end{lemma}

\begin{proof}
Пусть $a=q_a m +r_1$, $b=q_b m +r_1$, $c=q_c m +r_2$, $d=q_d m +r_2$.
Тогда $a+c = q_a m +r_1 + q_c m +r_2=  q_{ac} m +r$ и 
$b+d = q_b m +r_1 + q_d m +r_2=  q_{bd} m +r$. 
Если $r =r_1+r_2\ge m$, то положим $r\Assign r-m$, 
$q_{ac} \Assign q_{ac}+1$, $q_{bd} \Assign q_{bd}+1$.
\end{proof}

\smallskip
\begin{remark}
Отношение сравнения по модулю 
легко может быть определено не только для натуральных,
но и для всех целых чисел.
Это обстоятельство используется 
нами в случае необходимости без дополнительных оговорок.
\end{remark}
\end{frame}

\begin{frame}
\frametitle{2.4.5. Сравнения (3/3)}
\begin{lemma} \NB{ [2]}
Пусть $\Mod{a}{b}{m}$ и $\Mod{c}{d}{m}$.
Тогда $\Mod{ac}{bd}{m}$.
\end{lemma}

\begin{proof}
Не ограничивая общности, можно считать, что 
$a>b, c>d$.
Так как $m|(a - b)$ и $m|(c - d)$, 
то $a-b = km$, $c-d = lm$ для некоторых целых $k,l$. Имеем $ac = (km + b)(lm + d) = (klm+bl+dk)m +bd$. 
Из этого равенства следует, что остатки $ac$ и $bd$ при делении на $m$ равны, а значит, $\Mod{ac}{bd}{m}$.
\end{proof}
\smallskip
\begin{lemma}  \NB{ [3]}
Если $\Mod{a}{b}{m}$, $n \in \Bbb{N}$,
то $\Mod{a^n}{b^n}{m}$.
\end{lemma}

\begin{proof}
Индукция по $n$. База при $n = 1$ верна по условию леммы. 
Индукционный переход по лемме 2. 
\end{proof}

\smallskip
Из доказанных лемм непосредственно следует теорема.

\begin{theorem}
Пусть $f(x) = a_nx^n + a_{n-1}x^{n-1} + {\dots} + a_0$ --- многочлен с натуральными коэффициентами. Тогда если $\Mod{a}{b}{m}$, то $\Mod{f(a)}{f(b)}{m}$.
\end{theorem}

\end{frame}

\begin{frame}
\label{2.4.6.}\subsubsection{2.4.6. Системы вычетов }\frametitle{2.4.6. Системы вычетов (1/5)}


Поскольку деление с остатком определено для всех целых чисел,
отношение сравнимости по модулю также распространяется на всё множество
целых чисел.

\smallskip
\odmkey{Классом вычетов по модулю}{Класс (class)!вычетов по модулю (of residue modulo)}  $m$ 
(обозначение $[a]_m$)
называется класс эквивалентности (см.~\DMref{п.~1.7.1}{1.7.1.}.) 
по отношению сравнимости:

\vspace{-12pt}
$$
[a]_m \Def \Set{b\in\Bbb{Z}}{\Mod{b}{a}{m}},\quad a\in \Bbb{Z}, m\in \Bbb{N}.
$$ 


\begin{remark}
 Слово <<вычет>> употребляется как синоним слова <<остаток>>.
Всего классов вычетов по модулю $m$ ровно $m$ 
и один класс вычетов составляют все 
числа, дающие один и тот же остаток при делении на $m$.
\end{remark}

\smallskip
Определим над классами вычетов операции сложения $+_m$ и умножения $\ast_m$:
$$[x]_m +_m [y]_m \Def [(x + y)]_m,\quad  [x]_m \ast_m [y]_m \Def [(xy)]_m.$$
Множество классов вычетов с введёнными операциями обозначается $\Bbb{Z}_m$.

\smallskip
\begin{remark}
Остаток от деления на модуль $m$
является эпиморфизмом (см. \DMref{2.1.6}{2.1.6.})
из множества целых чисел $\Bbb{Z}$
на множество классов вычетов $\Bbb{Z}_m$
как по сложению, так и по умножению. 
\end{remark}
\end{frame}

\begin{frame}
\frametitle{2.4.6. Системы вычетов (2/5)}

\begin{lemma} \NB{ [1]}
Множество классов вычетов $\Bbb{Z}_m$ 
является коммутативным кольцом с единицей.
\end{lemma}

\begin{proof}
В качестве нуля нужно взять $[0]_m$, 
а в качестве единицы~--- $[1]_m$. 
Свойства из определения кольца (\DMref{п.~2.3.1}{2.3.1.}) выполнены в силу 
арифметических свойств сравнений.
\end{proof}


\smallskip
\odmkey{Полной системой вычетов по модулю }
{Система (system)!вычетов по модулю (of residue modulo)!полная (complete)} $m$ называют множество целых чисел, содержащее ровно по 
одному элементу из каждого класса вычетов по модулю $m$.

\begin{example}
$\{7,10,-1,3,101\}$~-- полная 
система вычетов по модулю 5.
\end{example}

\smallskip
\begin{remark}
Любое целое число сравнимо
с одним и только с одним элементом любой полной системы вычетов по любому модулю. 
\end{remark}

\smallskip
Множество $\{ 0, 1, 2, \dots , m - 1\}$ называется системой \odmkey{наименьших неотрицательных вычетов}{Система (system)!наименьших неотрицательных вычетов (smallest non-negative residues)} по модулю $m$. По принципу Дирихле эта система полна.

\begin{theorem} \NB{ [о полных системах вычетов]}
Пусть $\{\Tuple{a}{1}{m}\}$~--- полная система вычетов по модулю $m$, $x, y$ --- целые числа, $x \perp m$. Тогда $\{\Tuple{y+xa}{1}{m}\}$~--- полная система вычетов по модулю $m$.
\end{theorem}
\end{frame}

\begin{frame}
\frametitle{2.4.6. Системы вычетов (3/5)}

\begin{proof}
Покажем, что числа $\Tuple{y+xa}{1}{m}$ 
дают различные остатки при делении на число $m$. 
Тогда, поскольку чисел ровно $m$, 
они образуют полную систему вычетов. 
Предположим противное, пусть $y+xa_i$ и $y+xa_j$ дают 
одинаковые остатки при делении на $m$. 
Тогда $\Mod{xa_i + y}{xa_j + y}{m}$ или 
$m|x(a_i - a_j)$. 
Так как $x \perp m$, $m|(a_i - a_j)$ или 
$\Mod{a_i}{a_j}{m}$, 
что противоречит тому, что исходная система вычетов 
$\{\Tuple{a}{1}{m}\}$ 
является полной системой вычетов по модулю $m$.
\end{proof}


\begin{corollary}
$\Bbb{Z}_p$ является полем тогда и только тогда,
когда $p \in \Bbb{P}$.
\end{corollary}

\begin{proof}
\proofsection{$p \in \Bbb{P} \Impl \Bbb{Z}_p$~--- поле}
Поскольку по лемме $\Bbb{Z}_p$ 
является кольцом (\DMref{п.~2.3.1}{2.3.1.}), 
достаточно установить наличие у каждого вычета 
$[x]_p$
обратного по введённой операции умножения. 
Для этого рассмотрим любую 
полную систему вычетов $\{\Tuple{a}{1}{p}\}$, умноженную на $x$,
$\{\Tuple{xa}{1}{p}\}$. По 
теореме это будет полная система вычетов, 
то есть,\\ $\E{i\in 1..p}{\Mod{xa_i}{1}{p}}$. 
Класс $[a_i]_p$ и будет искомым обратным элементом для $[x]_p$.
\proofsection{$\Bbb{Z}_p$~--- поле 
$\Impl p \in \Bbb{P}$} От противного. 
Пусть $p = k l, 1<k\le l <p$.
Тогда $[k]_p \ne [0]_p$ и $[l]_p \ne [0]_p$,
но $[kl]_p = [p]_p = [0]_p$,
и значит $[k]_p$ и $[l]_p$
являются ненулевыми делителями нуля, что невозможно в поле.  
\end{proof}

\end{frame}

\begin{frame}
\frametitle{2.4.6. Системы вычетов (4/5)}

\begin{lemma} \NB{ [2]}
Если $\Mod{a}{b}{m}$, то $\gcd(a,m) = \gcd(b,m)$.
\end{lemma}

\begin{proof}
$\Mod{a}{b}{m} \Impl b = a + mk$, $k \in \Bbb{Z} \Impl \gcd(a,m) = \gcd(b,m)$.
Последнее вытекает из леммы \DMref{п.~2.4.2}{2.4.2.}.
\end{proof}

\medskip
\odmkey{Приведённой системой вычетов по модулю }
{Система (system)!вычетов по модулю (of residue modulo)!приведённая (reduced)} $m$ 
называют множество целых чисел, взаимно простых 
с числом $m$, и содержащихся в различных 
классах вычетов по модулю $m$.
Число элементов в приведённой системе вычетов равно 
количеству чисел, взаимно простых с $m$ и меньших $m$ (функция Эйлера, см. \DMref{п.~2.4.9}{2.4.9.}).  
Любое целое число, взаимно простое с $m$, 
сравнимо по модулю $m$
с одним и только с одним элементом любой приведённой системы вычетов по модулю $m$. 

\begin{example}
$ \{7, -1, 3, 101\}$~--- приведённая система вычетов по модулю 5.
\end{example}

\medskip
\begin{theorem}
Если $\{\Tuple{a}{1}{n}\}$~--- приведённая система вычетов 
по модулю $m$, $x \in \Bbb{Z}$, $x \perp m$, 
то $\{\Tuple{xa}{1}{n}\}$~--- приведённая система вычетов по модулю $m$.
\end{theorem}

\begin{proof}
Аналогично случаю полной системы вычетов.
\end{proof}
\end{frame}

\begin{frame}
\frametitle{2.4.6. Системы вычетов (5/5)}

\begin{corollary}
Приведённая система вычетов образует
абелеву группу по умножению.
\end{corollary}

\begin{proof}
Ассоциативность, коммутативность и наличие
нейтрального элемента очевидны в любой системе вычетов. Осталось показать, что 
в приведённой системе вычетов для любого элемента
найдётся обратный.
Пусть $\{\Tuple{a}{1}{n}\}$~--- приведённая 
система вычетов по модулю $m$ и $x \perp m$.
По теореме система $\{\Tuple{xa}{1}{n}\}$ также приведённая,
и значит    
$\E{i\in 1..n}{\Mod{xa_i}{1}{m}}$.
Положим $[x]_m^{-1}\Assign [a_i]_m$. 
\end{proof}


\end{frame}

\begin{frame}
\label{2.4.7.}\subsubsection{2.4.7. Китайская теорема об остатках}
\frametitle{2.4.7. Китайская теорема об остатках (1/2)}

\begin{remark}
\odmkey{Китайская теорема об остатках}{Китайская теорема об остатках (chinese remainder theorem)} 
была сформулирована в Китае, предположительно, 
в 3 веке н.э., китайским математиком Сунь Цзы, откуда и получила свое 
название.
\end{remark}

\begin{theorem}
Для любых попарно взаимно простых натуральных чисел $\Tuple{a}{1}{n}$ 
и для любых целых чисел $\Tuple{r}{1}{n}$ таких, 
что $\A{i\in 1..n}{0 \le r_i < a_i}$ существует такое число $s$, что деление его на каждое из чисел $a_i$ даёт в остатке число $r_i$.
\end{theorem}
\begin{proof}
Индукция по $n$. База: при $n = 1$ возьмём $s\Assign r_1$.
Индукционный переход. 
Пусть  $\E{t\in \Bbb{N}}{\A{i\in 1..(n-1)}{t=q_ia_i+r_i}}$.
Рассмотрим систему наименьших неотрицательных вычетов
$\{0, \dots, a_n-1\}$ по модулю $a_n$ и применим теорему о полных системах вычетов,   
полагая $x \Assign a_1 a_2 \dots a_{n-1}$ и $y\Assign t$. 
Тогда система чисел $s_1 = t$, $s_2 = t +x$, $\dots$, $s_{a_n} = t + (a_n - 1)x$
является полной системой вычетов по модулю $a_n$,
причём $\A{i\in 1..(n-1)}{\A{j\in 1..a_n}{\Mod{s_j}{t}{a_i}}}$ по построению.
Ввиду полноты системы среди чисел $\Tuple{s}{1}{a_n}$ 
существует такое число $s$, которое при делении на $a_n$ даст в остатке $r_n$. 
\end{proof}

\begin{corollary}
Пусть $m \Assign m_1 m_2 \dots m_n$, $\A{i,j\in 1..n, i\ne j}{m_i\perp m_j}$, 
$x, a\in \Bbb{Z}$.\\ 
Тогда $\Mod{x}{a}{m} \Equiv \A{i \in 1.. n}{\Mod{x}{a}{m_i}}$.
\end{corollary}
\end{frame}

\begin{frame}
\frametitle{2.4.7. Китайская теорема об остатках (2/2) }

В терминах общей алгебры китайская теорема
об остатках может быть сформулирована следующим 
образом.

\begin{theorem}
Если числа $\Tuple{m}{1}{k}$ попарно взаимно просты и
$M=m_1\cdot\ldots\cdot  m_k$, то
отображение 
$\ \Bbb{Z}_M \to \Bbb{Z}_{m_1} \times \dots
\times \Bbb{Z}_{m_k}$,
которое задаётся соответствием\\ 
$x\mapsto (x\MOD m_1, \dots, x\MOD m_k)$,
является изоморфизмом колец.
\end{theorem} 

\begin{proof}
Операции сложения и умножения 
на наборах остатков 
$(r_1, \dots, r_k)$
опеределены покомпонентно, то есть
$(\Tuple{r}{1}{k})+
(\Tuple{s}{1}{k})\Def (r_1+s_1, \dots, r_k+s_k)$,\\
$(\Tuple{r}{1}{k})\ast
(\Tuple{s}{1}{k})\Def (r_1\ast s_1, \dots, r_k\ast s_k)$.
Из этого, и из лемм \DMref{п.~2.4.5}{2.4.5.}
следует, что указанное в теореме отображение~---
гомоморфизм. 
По принципу Дирихле сюръективное 
отображение равномощных множеств биективно.
Сюръективность проверена в предыдущей теореме.
Осталось проверить равномощность.
По следствию к теореме \DMref{1.4.2}{1.4.2.}
$\left\vert \Bbb{Z}_{m_1} \times \dots \times
\Bbb{Z}_{m_k}\right\vert
= m_1\cdot\ldots\cdot  m_k = M =
\left\vert \Bbb{Z}_{M}\right\vert$.
\end{proof}


\begin{remark}
Пусть $M = m_1m_2 \dots m_k$, причём числа $m_1, m_2 \dots m_k$ --- попарно взаимно просты. Тогда китайская теорема об остатках позволяет
заменить вычисления в кольце $\Bbb{Z}_M$
вычислениями в кольце векторов 
$\Bbb{Z}_{m_1} \times \dots \times
\Bbb{Z}_{m_k}$.
\end{remark}
\end{frame}


\begin{frame}
\label{2.4.8.}\subsubsection{2.4.8. Вычисления в остаточных классах }\frametitle{2.4.8. Вычисления в остаточных классах (1/2)}

Пусть $m \Assign m_1 m_2 \dots m_n$, $\A{i,j\in 1..n, i\ne j}{m_i\perp m_j}$.
По следствию к китайской теореме об остатках любое число $a\in 0..(m-1)$
однозначно определяется остатками $\Tuple{a}{1}{n}$ от деления
числа $a$ на числа $\Tuple{m}{1}{n}$.
Такая непозиционная система счисления называется 
\odmkey{системой остаточных классов }{Система (system)!остаточных классов (of residue number)}, 
или \odmkey{модулярной арифметикой}{Арифметика (arithmetic)!модулярная (modular)}. 
Набор чисел $\Tuple{m}{1}{n}$ называется 
\odmkey{основанием}{Основание (base)!системы остаточных классов (system of residue number)} системы остаточных классов.\\
Обозначение: $a=(\Tuple{a}{1}{n})$. 

В системе остаточных классов операции сложения и умножения 
выполняются покомпонентно, то есть, 
если $a=(\Tuple{a}{1}{n})$, $b=(\Tuple{b}{1}{n})$, 
$a + b = c$, $ab = d$, то $c=(\Tuple{c}{1}{n})$, $d=(\Tuple{d}{1}{n})$, 
где $c_i \Assign (a_i + b_i)\MOD m_i $ и $d_i\Assign (a_i  b_i)\MOD m_i $ соответственно.

\begin{example}
Пусть система остаточных классов задана своим основанием $ (3,4,5)$. 
Вычислим в этой системе сумму и произведение чисел $7$ и $8$.
Имеем: $7 = (7 \MOD 3, 7 \MOD 4, 7 \MOD 5) = (1, 3, 2)$, $8 = (2,0,3)$. 

Тогда сложение: 
$7 + 8 = ((1+2) \MOD 3, (3+0) \MOD 4, (2 + 3) \MOD 5) = (0, 3, 0)$
и  $15 = (15 \MOD 3, 15 \MOD 4, 15 \MOD 5) = (0, 3 ,0)$.

Умножение: 
$7 \cdot 8 = ((1\cdot 2) \MOD 3, 
(3\cdot0) \MOD 4, (2\cdot 3) \MOD 5) = (2, 0, 1)$
и  $56 = (56 \MOD 3, 56 \MOD 4, 56 \MOD 5) = (2, 0 ,1)$.
\end{example}

\end{frame}


\begin{frame}
\frametitle{2.4.8. Вычисления в остаточных классах (2/2)}

\begin{digress}
Вычисления в остаточных классах в некоторых специальных применениях
используются наряду или даже вместо вычислений 
в обычной позиционной двоичной системе счисления.
 
\smallskip
Одним из главных преимуществ вычислений в системах остаточных классов 
является отсутствие переноса разрядов при сложении и умножении,
поскольку
достаточно сложить или умножить  представляющие векторы поэлементно.
В специализированном процессоре это можно сделать параллельно, 
а поэтому \em{время выполнения обеих операций 
не зависит от количества чисел в основании системы остаточных классов}.
Сами числа в основании системы остаточных классов можно выбирать небольшими,
и аппаратные операции с ними можно выполнять очень быстро.  

\smallskip
Ограничение диапазона представления чисел не является существенным.
Например, если взять основания $3, 5, 7, 11, 13$,
то для хранения остатков потребуется $2+3+3+4+4=16$ бит.
При этом диапазон представления составит $0..15014$,
что вполне сопоставимо с диапазоном $0..65535$
для 16-битных двоичных чисел.

\smallskip
Действительно существенным недостатком
 вычислений в остаточных классах
являются медленные алгоритмы
перевода из позиционной системы счисления в систему остаточных классов 
и обратно, а также деления с остатком.
Эффективные алгоритмы для этих операций неизвестны.
\end{digress}

\end{frame}



\begin{frame}
\label{2.4.9.}\subsubsection{2.4.9. Функция Эйлера }\frametitle{2.4.9. Функция Эйлера (1/5)}

% Мне эта преамбула показалась лишней. ФАН 29.04.2015 
%Функция и теорема Эйлера обобщают знаменитую теорему Ферма, 
%которая является одной из ключевых теорем в теории чисел, 
%и поэтому занимают одно из центральных мест в этой области математики.

\odmkey{Функцией Эйлера}{Функция (function)!Эйлера (Euler)} от натурального аргумента $n$ 
называется количество  чисел  
меньших $n$ и взаимно простых с $n$. 
Стандартным обозначением для функции Эйлера является $\varphi (n)$.
Из определения ясно, что $n>1\Impl 1\le \varphi(n)<n$.
По определению $\varphi(1)\Assign 1$.

\begin{example}
$\varphi(1)=1, \varphi(2)=1, \varphi(3)=2, \varphi(4)=2, \varphi(5)=4, \varphi(6)=2$.
\end{example}

\begin{remark}
Приведённая система вычетов по модулю $n$ содержит $\varphi(n)$ элементов.
\end{remark}

\smallskip
Вообще говоря, числовая функция $f$ называется 
\odmkey{аддитивной}{Функция (function)!аддитивная (additive)}, если она уважает сложение,
т.~е. $f(x+y)=f(x)+f(y)$, и называется  
\odmkey{мультипликативной}{Функция (function)!мультипликативная (multiplicative)}, если она уважает умножение,
т.~е. $f(xy)=f(x)f(y)$.

\smallskip  
В теории чисел используются несколько отличающиеся определения
этих понятий. \\Функция $\Map{f}{\Bbb{N}}{\Bbb{N}}$ называется
\odmkey{аддитивной}{Функция (function)!аддитивная (additive)}, если
$$\A{x, y \in \Bbb{N}}{x \perp y \Impl f(x+y) = f(x) + f(y)}$$
и называется
\odmkey{мультипликативной}{Функция (function)!мультипликативная (multiplicative)}, если
$$\A{x, y \in \Bbb{N}}{x \perp y \Impl f(xy) = f(x)f(y)}.$$

\end{frame}

\begin{frame}
\frametitle{2.4.9. Функция Эйлера (2/5)}

\begin{examples}
Количество множителей в разложении числа $n$ на простые без учета кратности
 $\Omega (n)$ --- аддитивная функция. Число делителей числа $n$ (обозначается $d(n)$) и
сумма делителей числа $n$ (обозначается $\sigma (n)$)~--- мультипликативные функции. 
\end{examples}


\smallskip
Также в теории чисел выделяют 
\odmkey{вполне аддитивные}{Функция (function)!вполне аддитивная (completely additive)} и 
\odmkey{вполне мультипликативные}{Функция (function)!вполне мультипликативная (completely multiplicative)} функции, 
пренебрегая взаимной простотой в 
соответствующих определениях.

\begin{example}
Степенная функция $f(n) = n^a, a > 0$ вполне 
мультипликативна.
\end{example}


Зададимся мультипликативной функцией $f$. 
Заметим, что для подсчета её значения для произвольного $n$ достаточно 
посчитать её значения для чисел 
вида $p^a$, где $p \in \Bbb{P}$, \\$a \in \Bbb{N}_{0}$. 
Пусть $n$~--- произвольное натуральное число. 
Тогда по основной теореме арифметики (\DMref{п.~2.4.4}{2.4.4.}) 
число $n$ можно единственным образом разложить на простые 
множители $n = p_1^{a_1}p_2^{a_2} \dots p_k^{a_k}$, 
и из мультипликативности получим 
$$f(n)=f(p_1^{a_1}p_2^{a_2} \dots p_k^{a_k})=f(p_1^{a_1})f(p_2^{a_2}) \dots f(p_k^{a_k}).$$

\end{frame}

\begin{frame}
\frametitle{2.4.9. Функция Эйлера (3/5)}
\begin{theorem} \NB{[1]}
Функция Эйлера мультипликативна.
\end{theorem}

\begin{proof}
Пусть $a \perp b$. 
Рассмотрим множество $$U \Assign\Set{ ax + by}{ x \in \{ 1,2, \ldots ,b \}, 
y \in \{ 1,2, \ldots ,a \} },$$ 
состоящее из $ab$ чисел. 
Заметим, что все они дают разные остатки при делении на $ab$. 
Действительно, пусть $\Mod{ax_1 + by_1}{ax_2 + by_2}{ab}$. 
Тогда $\Mod{ax_1}{ax_2}{b}$ и \\$\Mod{by_1}{by_2}{a}$,
откуда $x_1 = x_2$ и $y_1 = y_2$. 

\smallskip
Далее, $ax+by \perp ab \Equiv x \perp b \And y \perp a$. 
Действительно, 
$ax + by \perp ab \Equiv$ \\
$ \Equiv ax + by \perp b \And ax + by \perp a \Equiv ax  \perp b \And  by \perp a \Equiv$
$x  \perp b \And y \perp a $.

\smallskip
Таким образом, множество элементов $U$, взаимно простых с $ab$, 
изоморфно множеству пар $(x,y)$ таких, 
что $x \in 1..b \And x \perp b$, $y \in 1..a \And y \perp a$. \\ 
Но $|\Set{ax+by=u\in U}{ax+by \perp ab}| = \varphi(ab)$,
$|\Set{x \in 1..b}{x\perp b}| = \varphi(b)$,\\
$|\Set{y \in 1..a}{y\perp a}| = \varphi(a)$.
Поэтому $\varphi (ab) = \varphi (a) \varphi (b)$.
\end{proof}

\end{frame}

\begin{frame}
\frametitle{2.4.9. Функция Эйлера (4/5)}
\begin{theorem} \NB{[2]}
Если $n=p_1^{a_1}p_2^{a_2} \dots p_k^{a_k}$, $p_i\in\Bbb{P}$, то
 $$\varphi (n) =  
 \left(p_1^{a_1} - p_1^{a_1 - 1}\right)
 \left(p_2^{a_2} - p_2^{a_2 - 1}\right) \dots 
 \left(p_k^{a_k} - p_k^{a_k - 1}\right) = 
 n \prod\limits_{i = 1}^k \left(1 - \frac{1}{p_i}\right).$$
\end{theorem}


\begin{proof}
Заметим, что $\varphi (p^{a}) = p^{a} - p^{a - 1}$
 для $a > 0$. 
Действительно, числами, не 
взаимно простыми со степенью простого числа, являются те и только те, 
которые на это простое число делятся. 
Эти числа имеют вид $1p, 
2p,3p, \ldots , p^{a−1}p$, 
а их ровно $p^{a - 1}$. Отсюда и из мультипликативности функции Эйлера и следует утверждение теоремы. 
\end{proof}


\smallskip
Cамое известное и важное свойство функции Эйлера выражается в теореме Эйлера.

\begin{theorem} \NB{ [Эйлера]}.
Если $a\perp n$, то $\Mod{a^{\varphi (n)}}{1}{n}$.
\end{theorem}

\begin{proof}
Рассмотрим приведённую систему вычетов $\{\Tuple{x}{1}{\varphi(n)}\}$ по модулю $n$. 
В ней ровно $\varphi (n)$ элементов. 
Домножим каждый элемент  этой системы на $a$. По теореме из \DMref{п.~2.4.6}{2.4.6.}. полученная система тоже является 
приведённой системой вычетов по модулю $n$.
\end{proof}

\end{frame}

\begin{frame}
\frametitle{2.4.9. Функция Эйлера (5/5)}

\vspace{-2pt}
\begin{proof} (Продолжение)
Таким образом, произведение всех полученных вычетов сравнимо с произведением всех исходных (поскольку новая система вычетов также приведённая, то она является 
некоторой перестановкой исходной системы).
Отсюда $\Mod{ax_1  ax_2 \dots ax_{\varphi (n)}}{x_1x_2 \dots x_{\varphi (n)}}{n}$. Поскольку $x_1x_2 \dots x_{\varphi (n)}$ взаимно просто с $n$, то на это произведение можно поделить обе части сравнения.
Получим $\Mod{a^{\varphi (n)}}{1}{n}$.
\end{proof}

\begin{corollary} \NB{ [малая теорема Ферма]} 
Пусть $p$~--- простое число. 
Тогда $\Mod{a^{p-1}}{1}{p}$ для любого $a$, не делящегося на $p$.
\end{corollary}


\begin{example}
Покажем, как теорему Эйлера можно применять в задаче, 
на первый взгляд ничего общего с ней не имеющей.
Пусть числа $m$ и $n$ взаимно просты, $n$ не делится ни на 2, ни на 5. 
Оценить длину периода дроби $x=\frac{m}{n}$.
Решение. Существует $t$ такое, что 
$(10^t - 1)x = y$, где $y$ --- целое число. 
Отсюда $x = \frac{y}{10^t - 1}$. Поскольку $\frac{m}{n}$ --- несократимая дробь, 
имеем $n|(10^t - 1)$. Заметим, что $t = \varphi (n)$ подходит по теореме Эйлера, а длина периода --- минимальное 
такое $t$, значит, длина периода не превосходит $\varphi (n)$.
\end{example}

\end{frame}



\begin{frame}
\frametitle{Онтология элементарной теории чисел}
\begin{center}
\includegraphics[height=7.5cm]{ODM2_4.png}
\end{center}
\end{frame}


\begin{frame}
\frametitle{2.5. Векторные пространства (1/2)}
\subsection{2.5. Векторные пространства}
\label{2.5.}

\begin{columns}[t]
\begin{column}[T]{6cm}
Понятие векторного пространства должно быть известно читателю из
курса средней школы и других математических курсов. Обычно это
понятие ассоциируется с геометрической интерпретацией векторов в
пространствах $\Bbb R^2$ и $\Bbb R^3$. В этом разделе даны и
другие примеры векторных пространств, которые используются в
последующих главах для решения задач, весьма далёких от
геометрической интерпретации.
\end{column}
\begin{column}[T]{8cm}
\begin{figure}[h]
\begin{center}
\includegraphics[height=70mm]{odm02-01_1}
\end{center}
\end{figure}
\end{column}
\end{columns}
\end{frame}

\begin{frame}
\frametitle{2.5. Векторные пространства (2/2)}

\begin{tabular}{p{11mm}|p{43mm}|p{89mm}}
  № & Название параграфа
  & Ключевые термины и обозначения \\
  \hline
\DMref{2.5.1}{2.5.1.}\RS &	
Пространство векторов над полем	& 
аддитивный и мультипликативный нейтральные элементы, векторное пространство, вектор, скаляр \\
\DMref{2.5.2}{2.5.2.} &	Линейные комбинации	& линейно независимое множество, линейно зависимое множество \\
\DMref{2.5.3}{2.5.3.} & Базис и размерность &	порождающее множество, множество образующих, базис, координаты, размерность векторного пространства  \\
\DMref{2.5.4}{2.5.4.} &	Конечные поля &	поле Галуа, сужение и расширение поля, порядок по сложению, характеристика \\
\DMref{2.5.5}{2.5.5.} &	Модули & левые и правые модули \\
\DMref{2.5.6}{2.5.6.}\BS & Алгебры Ли & линейный оператор, линейное пространство, $[\ , \ ]$~--- скобка Ли; антикоммутативность, кососимметричность, билинейность, нильпотентность, коммутатор, тождество Якоби
\end{tabular}  
\end{frame}

\begin{frame}
\label{2.5.1.}\subsubsection{2.5.1. Пространство векторов над полем }\frametitle{2.5.1. Пространство векторов над полем (1/3)}
%\subsection{Векторное пространство}\label{2.7.1}
%\index{Векторное пространство (vector space)} 
Пусть ${\cal F} =\langle F ; +,\cdot\rangle$~--- поле с операцией сложения $+$,
операцией умножения $\cdot$, \odmkey{аддитивным}{Элемент
(element)!нейтральный (neutral)!аддитивный (additive)} нейтральным
элементом (нулем)~0 и \odmkey{мультипликативным}{Элемент
(element)!нейтральный (neutral)!мультипликативный (multiplicative)} (единицей)~1. Пусть ${\cal V} = \langle
V;+\rangle$~---  абелева группа с операцией $+$ и нейтральным
элементом~$\Bl{0}$. Если существует операция $F \times V \to V$
(знак этой операции опускается), такая, что для любых $a,b\in F$ и
для любых $\Bl{x},\Bl{y}\in V$ выполняются соотношения:
%\begin{tabbing}
%2.2\=$a(\Bl{x} + \Bl{y}) = a\Bl{x} + a\Bl{y}$\kill
%\>1)\' $(a+b)\Bl{x} = a\Bl{x} + b\Bl{x},$   \\[3pt]
%\>2)\' $a(\Bl{x} + \Bl{y}) = a\Bl{x} + a\Bl{y},$  \\[3pt]
%\>3)\' $(a \cdot b)\Bl{x} = a(b\Bl{x}),$  \\[3pt]
%\>4)\' $1\Bl{x} = \Bl{x},$
%\end{tabbing}
\begin{enumerate}
\item[1)]$(a+b)\Bl{x} = a\Bl{x} + b\Bl{x}$,
\item[2)]$a(\Bl{x} + \Bl{y}) = a\Bl{x} + a\Bl{y}$,
\item[3)]$(a \cdot b)\Bl{x} = a(b\Bl{x})$,
\item[4)]$1\Bl{x} = \Bl{x}$,
\end{enumerate}
то ${\cal V}$ называется \odmkey{векторным
пространством}{Пространство (space)!векторное (vector)} над полем
${\cal F}$, элементы $F$ называются \odmkey{скалярами}{Скаляр
(scalar)}, элементы $V$ называются \odmkey{векторами}{Вектор
(vector)}, нейтральный элемент группы $V$ называется
\odmkey{нуль-вектором}{Нуль-вектор (nill vector)} и обозначается
$\Bl{0}$, а необозначенная операция ${F \times V \to V}$
называется \odmkey{умножением вектора на скаляр}{Умножение
(multiplication)!вектора на скаляр (of vector on scalar)}.
\end{frame}

\begin{frame}%[shrink]
\frametitle{2.5.1. Пространство векторов над полем (2/3)}
\begin{examples}
%\begin{enumerate}
{\color{blue} 1.}
Пусть $\cal F=\Alg{F}{+,\cdot}$~--- некоторое поле. Рассмотрим
множество кортежей $F^n$. Тогда ${\cal F}^n=\Alg{F^n}{+}$, где
$$
(\Tuple{a}{1}{n})+(\Tuple{b}{1}{n})\Def (a_1+b_1,\dots, a_n+b_n)
$$
 является абелевой группой, в которой
$$-(\Tuple{a}{1}{n})\Def
(\Tuple{-a}{1}{n}) \quad\text{и}\quad \Bl{0}\Def (\Tuple{0}{}{}).$$
 Положим
$a(\Tuple{a}{1}{n})\Def (\Tuple{a\cdot a}{1}{n})$.
 Тогда ${\cal
F}^n$ является векторным пространством над $\cal F$ для любого
(конечного)~$n$. В частности, ${\Bbb R}^n$ является векторным
пространством для любого $n$. Векторное пространство ${\Bbb R}^2$
(и ${\Bbb R}^3)$ имеет естественную геометрическую интерпретацию.



{\color{blue} 2.} Двоичная арифметика ${\Bbb
Z_2}=\Alg{E_2}{{+}_2,\cdot}$ является полем, а булеан
$\Alg{2^M}{\vartriangle}$ с симметрической разностью является
абелевой группой. Положим $1X\Def X$, ${0X \Def \emptyset}$.
Таким образом, булеан с симметрической разностью является
векторным пространством над двоичной арифметикой.
\end{examples}
\end{frame}

\begin{frame}
\frametitle{2.5.1. Пространство векторов над полем (3/3)}
\begin{examples}
{\color{blue} 3.} Пусть
$X=\{\Tuple{x}{1}{n}\}$~--- произвольное конечное множество, \\ а
$\mathcal{F}=\Alg{F}{+,\cdot}$~--- поле. Рассмотрим множество
$V_X$ всех формальных сумм вида
$
\sum_{x\in X} \lambda_x x,
$
где $\lambda_x\in F$.
Положим

\vspace{-12pt}
$$
\begin{array}{lcl}
  \sum\limits_{x\in X} \lambda_x x + \sum\limits_{x\in X} \mu_x x & \Def &
  \sum\limits_{x\in X} (\lambda_x+\mu_x) x, \\
  \alpha\sum\limits_{x\in X} \lambda_x x & \Def &
  \sum\limits_{x\in X} (\alpha\cdot\lambda_x) x,
\end{array}
$$
где $\alpha,\lambda_x,\mu_x\in F$.
Ясно, что $V_X$ является векторным пространством, а $X$ является
его множеством образующих. В таком
случае говорят, что $V_X$~--- это векторное пространство,
<<натянутое>> на множество $X$. 

{\color{blue} 4.} В условиях предыдущего
примера ($X=\{\Tuple{x}{1}{n}\}$~--- произвольное конечное
множество, а $\mathcal{F}=\Alg{F}{+,\cdot}$~--- поле) рассмотрим
множество $\Phi\Assign\Set{f}{\Map{f}{X}{F}}$. Положим

\vspace{-6pt}
$$
\begin{array}{l}
  (f+g)(x) \Def  f(x)+g(x), \quad \text{где}\quad f, g \in \Phi,\\
  (\alpha f)(x)  \Def  \alpha\cdot f(x),\quad
  \text{где}\quad \alpha\in F, f\in\Phi.
\end{array}
$$
Нетрудно видеть, что $\Phi$ является векторным
пространством.
%\end{enumerate}
\end{examples}

\end{frame}

\begin{frame}
\label{2.5.2.}\subsubsection{2.5.2. Линейные комбинации }\frametitle{2.5.2. Линейные комбинации (1/2)}

\begin{lemma}
{\NB{1)}} $\ \A{{\Bl x} \in V}{0{\Bl x} ={\Bl{0}}},\quad$
{\NB{2)}} $\ \A{a \in F}{a{\Bl{0}} = {\Bl{0}}}$.
\end{lemma}

\begin{proof}
\proofsection{1}
$0{\Bl x}=(1-1){\Bl x}=1{\Bl x}-1{\Bl x}={\Bl x}-{\Bl x}={\Bl 0}$.
\proofsection{2}
$a{\Bl 0}=a({\Bl 0}-{\Bl 0})=a{\Bl 0}-a{\Bl 0}=(a-a){\Bl 0}=
0{\Bl 0}={\Bl 0}$.
\end{proof}

Если ${\cal V}$~--- векторное пространство над полем ${\cal F}$,
$S$~--- некоторое множество векторов,
$S \subset V$, то конечная сумма вида
\vspace{-3pt}
$$
\sum_{i=1}^n{a_i\Bl{x}_i}, \quad  a_i \in F, \quad  {\Bl{x}_i} \in
S,
$$
называется
\odmkey{линейной комбинацией}{Линейная комбинация (linear combination)}
векторов из $S$.
Пусть $X = \{\Tuple{\Bl{x}}{1}{k}\}$~--- конечное множество
векторов. Если
\vspace{-3pt}
$$
\sum_{i=1}^k {{a_i}{\Bl{x}_i}} = \Bl{0} \Impl \A{i\in 1..k}{a_i=0},
$$
то множество $X$ называется \odmkey{линейно независимым}{Множество
(set)!линейно независимое \\ (linearly independent)}.
\end{frame}

\begin{frame}
\frametitle{2.5.2. Линейные комбинации (2/2)}

 В противном
случае множество $X$ называется \odmkey{линейно зависимым}{Множество
(set)!линейно зависимое (linearly dependent)}:
\vspace{-6pt}
$$
\E{\Tuple{a}{1}{k} \in F} {\BigOr_{i=1}^k{a_i \not= 0} \And
\sum_{i=1}^k{a_i{\Bl{x}_i} = \Bl{0}}}.
$$

\begin{theorem}
Линейно независимое множество векторов не содержит нуль-вектора.
\end{theorem}
\begin{proof}
От противного. \\ Пусть линейно независимое множество
$S =\{\Bl{x}_1,\Bl{x}_2, \dots, \Bl{x}_k\}$ содержит
нуль-вектор и пусть для определённости $\Bl{x}_1= \Bl{0}$. Положим
$a_1\Assign 1$, $a_2\Assign 0,\dots,a_k\Assign 0$. Имеем
$$\sum\limits_{i=1}^k a_i\Bl{x}_i=1\Bl{0}+0\Bl{x}_2+\dots
+0\Bl{x}_k = \Bl{0}+\Bl{0}+\ldots +\Bl{0} =\Bl{0},$$ что
противоречит линейной независимости $S$.
\end{proof}
\end{frame}

\begin{frame}
\label{2.5.3.}\subsubsection{2.5.3. Базис и размерность }\frametitle{2.5.3. Базис и размерность (1/4)}
%\subsection{Базис и размерность}
%\label{basedim} 
Подмножество векторов $S \in V$, такое что любой
элемент $V$ может быть представлен в виде линейной комбинации
элементов $S$, называется \odmkey{порождающим}{Множество
(set)!порождающее (generating)} множеством пространства ${\cal V}$
или \odmkey{множеством образующих}{Множество (set)!образующих (of
generators)}. Линейно независимое порождающее множество называется
\odmkey{базисом}{Базис (basis)!векторного пространства (of vector
space)} векторного пространства. \setcounter{theorem}{0}
\begin{theorem}
Пусть векторное пространство ${\cal V}$ имеет базис $B =
\{\Tuple{\Bl{e}}{1}{n}\}$.
Тогда каждый элемент векторного пространства имеет единственное
представление в данном базисе.
\end{theorem}

\begin{proof}
Пусть $\Bl{x}=\sum\limits_{i=1}^n a_i\Bl{e}_i$ и
$\Bl{x}=\sum\limits_{i=1}^n b_i\Bl{e}_i$. \\Тогда
$\Bl{0}=\Bl{x}-\Bl{x}=\sum\limits_{i=1}^n
a_i\Bl{e}_i -\sum\limits_{i=1}^n b_i\Bl{e}_i =
\sum\limits_{i=1}^n (a_i-b_i)\Bl{e}_i\Impl$ \\
$\Impl\A{i\in 1..n}{a_i-b_i}=0\Impl \A{i}{a_i=b_i}$.
\end{proof}

\smallskip
Коэффициенты разложения вектора в данном базисе называются его
\odmkey{координатами}{Координаты (coordinates)}.

\end{frame}

\begin{frame}
\frametitle{2.5.3. Базис и размерность (2/4)}
\begin{theorem}
Пусть $B = \{\Tuple{\Bl{e}}{1}{n}\}$~{\upshape---} базис, а $X =
\{\Tuple{\Bl{x}}{1}{m}\}$~{\upshape---} линейно независимое
множество векторов. Тогда $n \geq m$.
\end{theorem}

\begin{proof}
От противного. Пусть $n < m$ и $B$~--- базис. Тогда
$$
\E{\Tuple{a}{1}{n}\in F}{\Bl{x}_1=a_1\Bl{e}_1+\ldots
+a_n\Bl{e}_n}.
$$
Имеем $\Bl{x}_1\not=
\Bl{0}\Impl\bigvee\limits_{i=1}^n{a_i\not= 0}$. Пусть для
определённости $a_1\not= 0$. Тогда
$$
\Bl{e}_1 =
a_1^{-1}\Bl{x}_1-(a_1^{-1}a_2)\Bl{e}_2-{\dots}-(a_1^{-1}a_n)\Bl{e}_n.
$$
Так~как $B$ порождает ${\cal V}$, то и $\{\Bl{x}_1, \Bl{e}_2,
\dots, \Bl{e}_n\}$ тоже порождает ${\cal V}$. Аналогично \\
$\{\Bl{x}_1, \Bl{x}_2, \Bl{e}_3, \dots, \Bl{e}_n\}$ порождает
${\cal V}$. Продолжая процесс, получаем, что
$\{\Tuple{\Bl{x}}{1}{n}\}$ порождает ${\cal V}$. Следовательно,
$$
\Bl{x}_{n+1}=\sum\limits_{i=1}^n{b_i\Bl{x}_i} \Impl
\Bl{x}_{n+1}-\sum\limits_{i=1}^n{b_i\Bl{x}_i} = \Bl{0}.
$$
Таким образом, $X$~--- линейно зависимое множество, что
противоречит условию.
\end{proof}
\end{frame}

\begin{frame}
\frametitle{2.5.3. Базис и размерность (3/4)}
\begin{corollary}
Если в векторном пространстве ${\cal V}$ множество векторов
$X=\{\Tuple{\Bl{x}}{1}{n}\}$ линейно независимо
и векторы множества $Y=\{\Tuple{\Bl{y}}{1}{m}\}$ выражаются в виде
линейных комбинаций векторов множества $X$, то есть
$\A{i\in 1..m}{\Bl{y}_i = \sum_{j=1}^{n} a_{ij}\Bl{x}_j}$, причём
$m>n$, то множество $Y$~{\upshape---} линейно зависимое.
\end{corollary}
\begin{proof}
Рассмотрим подпространство ${\cal V}_X$  векторного пространства
${\cal V}$, порождаемое векторами множества $X$. В подпространстве
${\cal V}_X$ множество $X$~--- базис, а векторы множества $Y$~--- элементы.
 Пусть множество $Y$ линейно независимое, тогда
по теореме $n\ge m$, что противоречит условию.
\end{proof}

\medskip
\begin{theorem}
Если $B_1$ и $B_2$~{\upshape---} базисы векторного пространства
${\cal V}$, то $ \vert B_1 \vert = \vert B_2 \vert. $
\end{theorem}
\begin{proof}
$B_1$~--- базис, и  $B_2$~--- линейно независимое множество.
Следовательно, $\vert B_1 \vert \geq \vert B_2 \vert$. С другой
стороны, $B_2$~--- базис, и $B_1$~--- линейно независимое
множество. Следовательно, $\vert B_2 \vert \geq \vert B_1 \vert$.
Имеем $\vert B_1\vert=\vert B_2 \vert$.\vadjust{\vspace*{-.75\baselineskip}}
\end{proof}

\end{frame}

\begin{frame}
\frametitle{2.5.3. Базис и размерность (4/4)}

Если векторное пространство ${\cal V}$ имеет базис $B$, то
количество элементов в базисе называется
\odmkey{размерностью}{Размерность (dimension)!векторного
пространства (of vector space)} векторного пространства и
обозначается $\dim{\cal V}$. Векторное пространство, имеющее
конечный базис, называется \odmkey{конечномерным}{Пространство (space)!векторное (vector)!конечномерное (finite-dimensional)}.
Векторное пространство, не имеющее базиса (в смысле приведённого
определения), называется \odmkey{бесконечномерным}{Пространство (space)!векторное (vector)!бесконечномерное (infinite-dimensional)}.
\begin{examples}
\begin{enumerate}
\item Одноэлементные подмножества образуют базис булеана, $\dim 2^M=\vert
M\vert$.
\item Кортежи вида $(0,\dots,1,0,\dots,0)$ образуют базис пространства
${\cal F}^n$, $\dim{\cal F}^n=n$.
\end{enumerate}
\end{examples}
\begin{remark}
Если определить базис как максимальное линейно независимое
множество (то есть множество, которое нельзя расширить, не нарушив
свойства линейной независимости), то понятие базиса можно
распространить и на бесконечномерные пространства.
\end{remark}
\end{frame}

\begin{frame}
\label{2.5.4.}\subsubsection{2.5.4. Конечные поля }\frametitle{2.5.4. Конечные поля (1/5)}

\vspace{-2pt}
На примере понятия конечных полей здесь демонстрируется тесная связь различных разделов общей алгебры и теории чисел.
  
Поле $\Bbb{F}_n$, состоящее из конечного числа элементов $n$, называется 
\odmkey{конечным полем}{Поле (field)!конечное (finite)}, 
или \odmkey{полем Галуа}{Поле (field)!Галуа (Galois)}.
 Иногда используется обозначение $GF(n)$ (от англ. Galois Field). 


\begin{example}
Если $p$~--- простое число, то $\Bbb{Z}_p$~--- множество классов вычетов по модулю $p$~---  поле Галуа (\DMref{п.~2.4.6}{2.4.6.}).
\end{example}

Если подмножество $G$ элементов поля $F$ также является полем, то $G$ называется \odmkey{сужением}{Сужение (constriction)!поля (field)} поля $F$, а $F$~--- \odmkey{расширением}{Расширение (dilation)!поля (field)} поля $G$. 

\begin{example}
Поле $\Bbb{F}_4=\Alg{\{0,1,a,b\}}{+,\ast}$ является расширением поля $\Bbb{Z}_2=\Alg{\{0,1\}}{+,\ast}$, 
если определить операции следующим образом:
$$
\Bbb{Z}_2: 
\begin{array}{c|cc}
 + &  0 &  1 \\
\hline
 0 & 0 & 1 \\
 1 & 1 & 0 \\
\end{array},
\begin{array}{c|cc}
 \ast &  0 &  1 \\
\hline
 0 & 0 & 0 \\
 1 & 0 & 1 \\
\end{array};\ 
\Bbb{F}_4:  
\begin{array}{c|cccc}
 + &  0 &  1 &  a &  b \\
\hline
 0 & 0 & 1 & a & b \\
 1 & 1 & 0 & b & a \\
 a & a & b & 0 & 1 \\
 b & b & a & 1 & 0 \\
\end{array},
\begin{array}{c|cccc}
 \ast &  0 &  1 &  a &  b \\
\hline
 0 & 0 & 0 & 0 & 0 \\
 1 & 0 & 1 & a & b \\
 a & 0 & a & b & 1 \\
 b & 0 & b & 1 & a \\
\end{array}.
$$ 
\end{example}
\end{frame}

\begin{frame}
\frametitle{2.5.4. Конечные поля (2/5)}
 
\odmkey{Порядком по сложению}{} элемента 
$a$ из абелевой группы $F$ называется
\em{наименьшее} натуральное число $q$ такое, 
что $q a = \underbrace{a + a + \ldots + a}_{q\ \text{раз}} = 0$.
Если такого числа нет, то говорят, что порядок элемента равен бесконечности (также иногда говорят, что порядок равен нулю).
Не во всякой абелевой группе элементы имеют конечный порядок по сложению.

\begin{example}
В группе 
$\Alg{\Bbb{Z}}{+}$ все элементы, кроме 
$0$, не имеют порядка по сложению. 
В группе $\Alg{\Bbb{Z}}{+_2}$  
порядок по сложению элемента $0$ равен $1$, 
а порядок по сложению элемента $1$ равен $2$.
\end{example}


\begin{lemma}
В конечном поле у
любого элемента  существует порядок по сложению: \\ 
$\A{a\in\Bbb{F}_n}{\E{q\in \Bbb{N}}{q a =0}}$.
\end{lemma}

\begin{proof}
Ввиду конечности поля в 
бесконечной последовательности 
$a, 2a, 3a, \dots$ 
существуют повторяющиеся элементы. Пусть $q_1a = q_2a$, $q_2 > q_1$. 
Тогда $(q_2 - q_1)a = 0$, и следовательно порядок элемента $a$ не больше $q_2 - q_1$.
\end{proof}


\end{frame}

\begin{frame}
\frametitle{2.5.4. Конечные поля (3/5)}

\begin{remark}
Не следует путать операцию умножения элемента на число \\$q a = a + a + \ldots + a\ (q\ \text{раз})$, где $a \in \Bbb{F}_n$, 
$q \in \Bbb{N}$, с операцией умножения в конечном поле $a*b$, где $a, b \in \Bbb{F}_n$.
\end{remark}

\begin{theorem} \NB{[1]}
Любое конечное поле $\Bbb{F}_n$ 
является расширением поля, изоморфного $\Bbb{Z}_p$, где $p$~--- простое число.
\end{theorem}

\begin{proof}
Рассмотрим $1$~--- единицу по умножению
в поле $\Bbb{F}_n$. 
Пусть $m$~--- порядок $1$ по сложению. 
Тогда $m>1$, поскольку в любом поле $0\ne 1$.
Если число $m$~--- составное,
то есть $m = bc$ ($b > 1$, $c > 1$),
то $m1 = 0 \Impl (bc)1 = 0\Impl (b1)*(c1) = 0$. Так как $b < m$, $c<m$, 
имеем $b1\ne 0$, $c1\ne 0$. У нуля оказались ненулевые делители, что противоречит свойству полей (\DMref{п.~2.3.3}{2.3.3.}),
значит  $m$~--- простое число.
Но $1\in \Bbb{F}_n$ и  
$\A{k\in\Bbb{N}}{k1\in\Bbb{F}_n}$. 
Каждому элементу $k1$ ($k<m$) можно сопоставить класс вычетов $k$ из поля вычетов $\Bbb{Z}_m$. 
Таким образом, в поле $\Bbb{F}_n$ можно выделить множество 
$G = \{0, 1, 1+1, \ldots, (m-1)1\}$, 
замкнутое относительно операций поля, и  изоморфное полю $\Bbb{Z}_m$.
\end{proof}


\end{frame}

\begin{frame}
\frametitle{2.5.4. Конечные поля (4/5)}
%// следующая теорема нужна для доказательства теоремы 3, а также для записи элементов F_n в виде многочлена.

%Согласен.
% Прошу включить материал
% попадая в обозначения и стиль 
% по приведенному образцу

%\vspace{-2pt}
\begin{theorem} \NB{[2]}
Пусть $\Bbb{F}_n$~--- конечное поле, 
а $\Bbb{G}$~--- сужение поля $\Bbb{F}_n$
изоморфное полю $\Bbb{Z}_p$, где число $p$~--- простое. 
Тогда конечное поле $\Bbb{F}_n$~---
векторное пространство над своим подполем $\Bbb{G}$.
\end{theorem}



\begin{proof}
Элементы поля $\Bbb{F}_n$ являются векторами,
сложение $+$ является векторным сложением, элементы поля $\Bbb{G}$ являются скалярами (и одновременно векторами!), а операция умножения $\ast$ является умножением на скаляр.  
Из доказательства теоремы 1 видно,
что нейтральные элементы в полях $\Bbb{F}_n$ и
 $\Bbb{G}$ совпадают.
Таким образом, все четыре аксиомы векторного пространства (\DMref{п.~2.5.1}{2.5.1.})
выполняются, как частные случаи свойств операций поля.   
%Достаточно построить базис 
%$X=\{\Tuple{x}{1}{m}\}$ такой, что 
%$\A{x\in F}{\E{\Tuple{a}{1}{m} \in \Bbb{G}}{x=\sum_{i=1}^n a_i\ast x_i}}$.
%
%\begin{algorithmic}
%\ENSURE множество F, G.
%\REQUIRE базис F.
%\STATE{$H\Assign\emptyset$}
%\COMMENT{\color{gray}\small $H$ можно получить из построенных элементов базиса}
%\STATE {$i\Assign 0$} 
%\COMMENT {\color{gray}\small $i$~--- количество построенных элементов базиса}
%\WHILE{$F \ne H$}
%\STATE {$i \Assign i + 1$}
%\STATE\textbf{select $x_i \in F \setminus H$}
%\STATE\textbf{yield $x_i$}
%\STATE {$T \Assign H$}
%\COMMENT {\color{gray}\small сохраняем множество H}
%%  for x из HH do
%%    for a из G do
%%      H = H + {x + a * xi}
%%    end for
%%  end for
%\ENDWHILE
%\end{algorithmic}
%Обоснование.
%H строится так, что на каждом шаге H = {a1*x1 + ... +ai*xi | a1, ...,ai ∈ G}, H принадлежит F.
%Ввиду конечности F и постепенного увеличения H, алгоритм завершится (через m шагов).
%xi выбирается из F - H, то есть xi не выражается через x1, ..., xi-1, а значит, x1, ..., xm линейно независимы. H = F, то есть x1, ..., xm – система образующих, а, следовательно, базис. Теорема доказана.
%
\end{proof}

\smallskip
\begin{corollary}
Пусть $\Bbb{F}_n$~--- конечное поле, 
а $\Bbb{G}$~--- сужение поля $\Bbb{F}_n$
изоморфное полю $\Bbb{Z}_p$, \\где число $p$~--- простое. 
Тогда существует множество  
$X=\{\Tuple{x}{1}{m}\}$ такое, \\что 
$\A{x\in F}{\E{\Tuple{a}{1}{m} \in \Bbb{G}}{x=\sum_{i=1}^m a_i\ast x_i}}$.
\end{corollary}

\begin{proof}
В качестве $X$ можно взять любое порождающее линейно независимое
множество, то есть базис.  
\end{proof}
\end{frame}

\begin{frame}
\frametitle{2.5.4. Конечные поля (5/5)}
\begin{theorem} \NB{[3]}
В поле Галуа $\Bbb{F}_n$ количество элементов $n= p^k$, 
где $p\in \Bbb{P}$~--- простое число, $k\in \Bbb{N}$~--- натуральное число.
\end{theorem}

\begin{proof}
 Пусть $p$~--- порядок $1$ по сложению. 
В теореме 1 показано, что в поле $\Bbb{F}_n$ существует подполе 
$\Bbb{G} =\{0,1,1+1, \dots, (p-1)1 \}$  порядка $p$.
По теореме 2 поле $\Bbb{F}_n$ является векторным пространством над
полем $\Bbb{G}$ и имеет базис $X=\{\Tuple{x}{1}{k}\}$. 
Таким образом, каждый элемент $x\in\Bbb{F}_n$ 
\em{единственным образом} представляется в виде суммы  
 $x=\sum_{i=1}^k a_i\ast x_i$, и каждая сумма 
 представляет некоторый элемент. Таких сумм ровно $p^k$,
поскольку коэффициенты   
$\Tuple{a}{1}{k}$ принимают все $p$ 
возможных значений, а $k$~--- число векторов в базисе.
%\proofsection{2} Построим поле F = {a1 + a2*x + ... + ak*x^k | a1, ..., ak ∈ Z_p}, F – поле Галуа с p^k элементами.
\end{proof}

\medskip
Число $p$ называется \odmkey{характеристикой}{Характеристика поля (field characteristic)} поля $F_{p^k}$.

\begin{remark} 
Можно показать, что верно и обратное, 
то есть
для любого $p\in \Bbb{P}$ и $k\in \Bbb{N}$ 
существует поле Галуа $\Bbb{F}_{p^k}$.
\end{remark}



%Порядок по сложению любого элемента по сложению равен p.

%Элементы поля Fp^k можно представить через k коэффициентов разложения по базису. Эти коэффициенты являются элементами поля Fp. Их можно записать в виде упорядоченных наборов, коэффициентов полиномов степени k – 1 или разрядов числа в системе счисления с основанием p.
%Пример: F2^3= {(0,0,0), (0,0,1), (0,1,0), (0,1,1), (1,0,0), (1,0,1), (1,1,0), (1,1,1)} = {0, 1, x, x + 1, x^2, x^2 + 1, x^2 + x, x^2 + x+ 1} = {0, 1, 2, 3, 4, 5, 6, 7} (при x = 2); F3^2 = {(0,0), (0,1), (0,2), (1,0), (1,1), (1,2), (2,0), (2,1), (2,2) } = {0, 1, 2, x, x + 1, x + 2, 2*x, 2*x + 1, 2*x + 2} = {0, 1, 2, 3, 4, 5, 6, 7, 8} (при x = 3).
%Так как поле Fp^k – векторное пространство, для выполнения операции сложения на нём необходимо выполнить операцию сложения на поле Zp для каждого коэффициента полинома.
%
%\odmkey{Полиномом над конечным полем}{Полином!над конечным полем}
%$\Bbb{F}_n$ называется сумма вида 
%$a_0 + a_1 * x + \ldots + a_m *x^m$, 
%$a_i \in\Bbb{F}_n$.
%
%Класс вычетов по модулю многочлена g(x) (обозначение [f(x)]g(x)) – класс эквивалентности по отношению сравнимости.
%[f(x)]g(x) = (def) {h(x) ∈ Zp(x) | f(x) = h(x) (mod g(x))}, где g(x) ∈ Zp(x), Zp(x) – множество полиномов от x с коэффициентами из Zp, p – простое число.
%Кольцо вычетов по модулю многочлена g(x) – множество всех классов вычетов, где операции сложения и умножения выполняются по модулю g(x).
%Замечание.
%Данные определения являются аналогом вычетов по модулю числа. 
%// Я считаю, что неприводимый многочлен обязательно следует упомянуть, так как он необходим для определения операции умножения в этом поле. Он вполне относится к basic, ведь его определение аналогично определению простого числа.
%Неприводимый многочлен над полем – это многочлен, который нельзя разложить в данном поле на неконстантные многочлены (степени 1 или больше).
%Замечание.
%Если g(x) – это неприводимый многочлен над полем Fp степени k, то кольцо вычетов – это поле Fp^k, в котором операция умножения выполняется по модулю g(x).
%Замечание.
%Нахождение неприводимых многочленов над полем Fp обычно производится перебором, можно также использовать готовые таблицы.
%Пример.
%F2^2. Неприводимый многочлен x^2 + x + 1.
%0 + 0 = 0; 0 + 1 = 1; 0 + x = x; 0 + x + 1 = x + 1;
%1 + 0 = 1; 1 + 1 = 0; 1 + x = x + 1; 1 + (x + 1) = x;
%x + 0 = x; x + 1 = x + 1; x + x = 0; x + (x + 1) = 1;
%(x + 1) + 0 = x + 1; (x + 1) + 1 = x; (x + 1) + x = 1; (x + 1) + (x + 1) = 0;
%0 * 0 = 0; 0 * 1 = 0; 0 * x = 0; 0 * (x + 1) = 0;
%1 * 0 = 0; 1 * 1 = 1; 1 * x = x; 1 * (x + 1) = x;
%x * 0 = 0; x * 1 = x; x * x = x + 1; x * (x + 1) = 1;
%(x + 1) * 0 = 0; (x + 1) * 1 = x + 1; (x + 1) * x = 1; (x + 1) * (x + 1) = x;
%Пример.
%Поле F2 = {0, 1}. В этом случае операции сложения и умножения можно представить как дизъюнкцию и конъюнкцию соответственно (см. п. 3.1.4).

%Хвотовик неубедительный и не увлекательный. Может быть вернуть картинку с "квадратными" кодами?
%Без объяснений! ФАН 

%В кодировании конечные поля используются для построения циклических кодов и их частного случая, кода Боуза-Чоудхури-Хоквингема (БЧХ кода).
%Отступление.
%Конечные поля используются в теории Галуа для решения задач о том, какие фигуры можно построить циркулем и линейкой, а также какие алгебраические уравнения разрешимы относительно операций сложения, вычитания, умножения, деления и выделения корня.

%Не для всякого числа $n\in \Bbb{N}$ существует поле Галуа $\Bbb{F}_n$.
%
%
%
%
\end{frame}


\begin{frame}
\label{2.5.5.}\subsubsection{2.5.5. Модули }\frametitle{2.5.5. Модули (1/2)}

 Понятие модуля во многом аналогично понятию векторного
пространства, с той лишь разницей, что векторы умножаются не на
элементы из поля, а на элементы из произвольного кольца.
Пусть $\mathcal{R}= \Alg{R}{+,\ast}$~--- коммутативное кольцо с
операцией сложения $+$, операцией умножения $\ast$, нулевым
элементом $0$ и единичным элементом $1$. Абелева группа
$\mathcal{M}=\Alg{M}{+}$ называется \odmkey{модулем}{Модуль (module)} над
кольцом $\mathcal{R}$, если задана операция умножения на скаляр
$R\times M\rightarrow M$ такая, что:
\begin{enumerate}
\item[1)] $(a+b)\Bl{x} = a\Bl{x} + b\Bl{x}$;%
\item[2)] $a(\Bl{x} + \Bl{y}) = a\Bl{x} + a\Bl{y}$;%
\item[3)] $(a \ast b)\Bl{x} = a(b\Bl{x})$;%
\item[4)] $1 \Bl{x} = \Bl{x}$,
\end{enumerate}
где $a,b \in R$, а $\Bl{x}, \Bl{y} \in M$. Как и в случае
векторных пространств, элементы кольца $\mathcal{R}$ называются
\odmkey{скаля\-рами}{Скаляр (scalar)}, а элементы группы
$\mathcal{M}$~--- \odmkey{векторами}{Вектор (vector)}.

\end{frame}


\begin{frame}
\frametitle{2.5.5. Модули (2/2)}
\begin{examples}

\begin{enumerate}

\item Векторное пространство $\mathcal{V}$ 
над полем $\mathcal{F}$ является модулем над $\mathcal{F}$, так как поле по определению является кольцом.

\item  Множество векторов с целочисленными координатами в
$\mathbb{R}^n$ является модулем над кольцом $\mathbb{Z}$.

\item  Любая абелева группа есть модуль над кольцом $\mathbb{Z}$, 
если операцию умножения на скаляр $n\in \mathbb{Z}, n>0$ определить следующим образом:
$n\Bl{v}\Def\Bl{v} + \Bl{v} + \dots + \Bl{v}$, где в сумму входит $n$ слагаемых.
\end{enumerate}
\end{examples}

\smallskip
Если кольцо некоммутативное, то различают 
\odmkey{левые модули}{Модуль (module)!левый (left)}, в которых аксиомы записываются так,
как указано на предыдущем слайде, и  
\odmkey{правые модули}{Модуль (module)!правый (right)}, в которых третье
свойство имеет вид:
\begin{enumerate}
\item[3)] $(a \ast b)\Bl{x} = b(a\Bl{x})$.
\end{enumerate} 

\end{frame}

\begin{frame}
\subsubsection{2.5.6. Алгебры Ли}\label{2.5.6.}
\frametitle{2.5.6. Алгебры Ли (1/9)}

\odmkey{Алгеброй Ли}{Алгебра (algebra)!Ли (Lie)} называется векторное пространство $\mathcal{V}$ над полем $\mathcal{F}$ с бинарной операцией $\Map{[\;,\;]}{\mathcal{V} \times \mathcal{V}}{\mathcal{V}}$ 
(операция называется \odmkey{скобка Ли}
{Операция (operation)!скобка Ли (Lie bracket)}), для которой выполняются  следующие условия:

\begin{enumerate}
\item $\A{\Bl{x}\in V}{[\Bl{x}, \Bl{x}] = \Bl{0}}$~--- \odmkey{нильпотентность}{Нильпотентность (nilpotency)} второй степени; 
\item $\A{\Bl{x}, \Bl{y}, \Bl{z}\in V}{[\Bl{x},[\Bl{y}, \Bl{z}]] + [\Bl{y},[\Bl{z},\Bl{x}]] + 
[\Bl{z}, [\Bl{x}, \Bl{y}]] = \Bl{0}}$~--- 
\odmkey{тождество Якоби}{Тождество (identity)!Якоби  (Jacobi)}; 
\item $\A{\Bl{x}, \Bl{y}, \Bl{z} \in V, a, b \in F}{[a\Bl{x}+b\Bl{y}, \Bl{z}] = a[\Bl{x},\Bl{z}]
+b[\Bl{y},\Bl{z}]}$ \\
$\A{\Bl{x}, \Bl{y}, \Bl{z} \in V, a, b \in F}
{[\Bl{x} ,a\Bl{y}+b\Bl{z}] = a[\Bl{x},\Bl{y}]
+b[\Bl{x},\Bl{z}]}$~--- \odmkey{билинейность}{Операция (operation)!билинейная (bilinear)} 
скобки Ли.
\end{enumerate} 

\begin{theorem}
Скобка Ли \odmkey{антикоммутативна}{Операция (operation)!антикоммутативная (anticommutativе)}, то есть
$\A{\Bl{x}, \Bl{y} \in V}{[\Bl{x},\Bl{y}]= -[\Bl{y},\Bl{x}]}$.
\end{theorem}

\begin{proof}
$\Bl{0} = [\Bl{x}+\Bl{y}, \Bl{x}+\Bl{y} ]=
[\Bl{x},\Bl{x}] + [\Bl{x},\Bl{y}] +[\Bl{y},\Bl{x}]+ [\Bl{y},\Bl{y}] = [\Bl{x},\Bl{y}] +[\Bl{y},\Bl{x}].
$
\end{proof}

\begin{remark}
Антикоммутативную операцию называют также
\odmkey{кососимметричной}{Операция (operation)!кососимметричная (skew-symmetric)}.
\end{remark}
 
%Можно показать, что если характеристика поля $\mathcal{F}$ больше $2$, то из антикоммутативности следует тождество $[\Bl{x},\Bl{x}] = 0$.

\begin{theorem}
Если характеристика поля $\mathcal{F}$ не равна двум,
то из кососимметричности операции следует её нильпотентность.  
\end{theorem}

\begin{proof}
Имеем  $\Bl{0} = [\Bl{x} + \Bl{y}, \Bl{x}] + 
    [\Bl{x}, \Bl{x} + \Bl{y}] = 
    [\Bl{x}, \Bl{x}] + [\Bl{y}, \Bl{x}] + [\Bl{x}, \Bl{x}]
     + [\Bl{x}, \Bl{y}] =$\\
     $= [\Bl{x}, \Bl{x}] + 
     [\Bl{x}, \Bl{x}] = 2[\Bl{x}, \Bl{x}]$.
Так как характеристика поля $\mathcal{F}$ не равна двум, то из  равенства $\Bl{0} = 2[\Bl{x}, \Bl{x}]$ заключаем:
    $[\Bl{x}, \Bl{x}] = \Bl{0}$.
\end{proof}

\end{frame}

\begin{frame}
\frametitle{2.5.6. Алгебры Ли (2/9)}

\begin{example}
Трёхмерное векторное пространство $\Bbb{R}^3$ над полем $\Bbb{R}$ с операцией векторного произведения является алгеброй Ли. 
Напоминание: 
\odmkey{векторным произведением}{Произведение (product)!векторное (cross)}
 векторов $\Bl{a}$ и $\Bl{b}$ называется вектор 
 $\Bl{c}$, перпендикулярный обоим векторам $\Bl{a}$ и $\Bl{b}$,  длина которого 
 $|\Bl{c}| = |\Bl{a}| \cdot |\Bl{b}| \cdot \sin\alpha$, где $\alpha$ -- угол между $\Bl{a}$ и $\Bl{b}$, а одно из двух возможных направлений определяется по 
 \odmkey{правилу правой руки}{Правило (rule)!правой руки (right-hand)} так: если указательный палец правой руки направлен по вектору $\mathbf{a}$, средний палец направлен по вектору 
  $\Bl{b}$, то вектор $\Bl{c}$ будет сонаправлен с большим пальцем, который перпендикулярен двум другим.
Нетрудно проверить выполнение свойств 
скобки Ли для векторного произведения. 
  
\begin{enumerate}
\item Нильпотентность следует из того, что синус угла между коллениарными векторами равен нулю, а значит векторное произведение двух совпадающих векторов даёт нуль-вектор.
\item Тождество Якоби проще всего наглядно 
представить себе для трёх ортогональных единичных векторов $x$, $y$, $z$,
как на рисунке слева на следующем слайде: по правилу правой руки каждое слагаемое в тождестве Якоби оказывается нуль-вектором в силу уже доказанной нильпотентности. 
\end{enumerate}  
\end{example}
\end{frame}

\begin{frame}
\frametitle{2.5.6 Алгебры Ли (3/9)}
\begin{enumerate}
\item[3.] Билинейность можно усмотреть аналогично тождеству Якоби, как показано на рисунках в центре и справа для случая ортогональных единичных векторов и единичных коэффициентов.
\end{enumerate}  
\begin{center}
\includegraphics[width=155mm]{2_5_6.jpg}
\end{center}

\begin{digress}
Алгебры Ли играют весьма важную роль в современной алгебре и связаны со многими прикладными областями, в особенности в физике, такими как электромагнетизм, квантовая физика и многое другое, но здесь мы ограничиваемся лишь рассмотрением некоторых примеров в рамках изучения векторных пространств и других алгебраических структур.
\end{digress}

\end{frame}

\begin{frame}
\frametitle{2.5.6 Алгебры Ли (4/9)}

Начнём с теории групп.
Пусть $\mathcal{G}$~--- произвольная группа. 
Бинарная операция  
$$[x,y] \Def (yx)^{-1}(xy) = x^{-1}y^{-1}xy$$ 
называется \odmkey{коммутатором}{Коммутатор (commutator)}. 
Нетрудно видеть, что 
$$[x,y]^{-1}= ((yx)^{-1}xy)^{-1} = (xy)^{-1}(yx)=
y^{-1}x^{-1}yx.$$

\begin{remark}
Коммутатор показывает меру некоммутативности операции, поэтому так и назван. Ведь $[x, y] = e$ тогда и только тогда, когда элементы $x$ и $y$ коммутируют.
\end{remark}

Коммутатор в группе обладает свойствами нильпотентности
и кососимметричности,
$$
[x,x] = e, \quad [x,y]=[y,x]^{-1},
$$   
то есть коммутатор обладает некоторыми свойствами скобки Ли, именно поэтому используется сходное обозначение.
Но тождество Якоби и дистрибутивность коммутатора относительно групповой операции не всегда выполняются,
поэтому группа с коммутатором в общем случае 
не является алгеброй Ли. 

\end{frame}

\begin{frame}
\frametitle{2.5.6 Алгебры Ли (5/9)}


\begin{digress}
Заметим, что если группа абелева, 
то коммутатор любой пары элементов вырождается в нейтральный элемент, а потому тождество Якоби 
и дистрибутивность выполняются,
и можно было бы формально считать, что 
абелева группа является алгеброй Ли 
над двоичной арифметикой.
Впрочем, если операция даёт во всех случаях один результат,
то~есть является константой,
то формально выполняются любые <<законы>>, 
записанные как равенста выражений с этой операцией.
Таким <<методом>> можно строить сколь угодно вычурные
абстрактные теории, которые, однако, малосодержательны.  
\end{digress}      

\smallskip
Рассмотрим теперь векторные пространства.

Векторное пространство $\mathcal{A}$ с 
носителем $A$ над полем 
$\mathcal{F}$, снабжённое билинейной операцией 
$\Map{\cdot}{A \times A}{A}$, называется 
\odmkey{алгеброй над полем}{Алгебра (algebra)!над полем (over a field)} $\mathcal{F}$.
Билинейную операцию в алгебре над полем принято называть умножением. 
Часто это операцию никак не обозначают.

\begin{remark}
Операцию умножения векторов в алгебре над полем 
не следует путать с операцией умножения скаляров на 
векторы в этом же векторном пространстве,
а также с операцией умножения скаляров в поле,
хотя все эти операции никак не обозначаются.
Какое именно умножение использовано в формуле определяется
типами операндов. 
%Групповую операцию самого векторного пространства
%$\mathcal{A}$ обозначают $+$, а обратный элемент знаком
%$-$, поскольку векторное пространство~--- это абелева группа. 
\end{remark}

\end{frame}

\begin{frame}
\frametitle{2.5.6 Алгебры Ли (6/9)}

Из билинейности умножения векторов в алгебре над полем
немедленно следует, что для любых векторов 
$\Bl{x}, \Bl{y}, \Bl{z} \in A$ и для любых
скаляров $a, b  \in F$ выполняются следующие свойства:

\begin{enumerate}
    \item $(a\Bl{x} + b\Bl{y})\Bl{z} = a\Bl{x z} + b\Bl{y z}$;
    \item $\Bl{x}(a\Bl{y} + b\Bl{z}) = a\Bl{x y} + b\Bl{x z}$;
    \item $(a \Bl{x})(b \Bl{y}) = (a b)(\Bl{x y})$.
\end{enumerate}

Алгебра над полем называется ассоциативной, если операция умножения в ней ассоциативна, и называется алгеброй с единицей, если умножение имеет нейтральный элемент.
Напомним, что если алгебра над полем является конечнопорождённой,
то число элементов в минимальном порождающем множестве называется размерностью. 

\begin{example}
Комплексные числа являются двумерной ассоциативной алгеброй с единицей
над полем вещественных чисел.
\end{example}

\medskip
Пусть $\mathcal{A}$ --- ассоциативная алгебра над полем
 $\mathcal{F}$. Определим в ней бинарную операцию по формуле: 
 $[\Bl{x},\Bl{y}]\Def\Bl{xy-yx},$
это выражение также называется 
 \odmkey{коммутатором}{Коммутатор (commutator)}.

\end{frame}

\begin{frame}
\frametitle{2.5.6 Алгебры Ли (7/9)}

\begin{remark}
Коммутатор в алгебре над полем определяется
через введённую операцию умножения векторов,
которая никак не обозначается, и через 
взятие обратного элемента по сложению.
Групповая операция $+$ в определении
коммутатора в алгебре над полем 
никак не используется, в отличие от 
коммутатора в группе.
\end{remark}

\begin{theorem}
    Коммутатор в ассоциативной алгебре 
    над полем является скобкой Ли.
\end{theorem}

\begin{proof} Проверим свойства скобки Ли.
\proofsection{Нильпотентность}    
$[\Bl{x}, \Bl{x}] = \Bl{xx} - \Bl{xx} = \Bl{0}$.
\proofsection{Тождество Якоби}$[\Bl{x}, [\Bl{y}, \Bl{z}]] + [\Bl{y}, [\Bl{z}, \Bl{x}]] + [\Bl{z}, [\Bl{x}, \Bl{y}]] = \Bl{xyz} - \Bl{xzy} - \Bl{yzx} + \Bl{zyx} + \Bl{yzx} - \Bl{xyz} - \Bl{zxy} + \Bl{xzy} + \Bl{zxy} - \Bl{zyx} - \Bl{xyz} + \Bl{yxz} = \Bl{0}$.
\proofsection{Линейность слева} $[a  \Bl{x} + b   \Bl{y}, \Bl{z}] = (a  \Bl{x} + b   \Bl{y})\Bl{z} - \Bl{z}(a  \Bl{x} + b   \Bl{y}) = a  \Bl{xz} - a  \Bl{zx} + b   \Bl{yz} - b   \Bl{zy} = a [\Bl{x}, \Bl{z}] + b  [\Bl{y}, \Bl{z}]$.
\proofsection{Линейность справа} $[\Bl{x}, a  \Bl{y} + b   \Bl{z}] = \Bl{x}(a  \Bl{y} + b   \Bl{z}) - (a  \Bl{y} + b   \Bl{z})\Bl{x} = a  \Bl{xy} - a  \Bl{yx} + b   \Bl{xz} - b   \Bl{zx} = a  [\Bl{x}, \Bl{y}] + b  [\Bl{x}, \Bl{z}]$.
 \end{proof}
 
Таким образом, всякая ассоциативная алгебра над полем  обладает естественной структурой алгебры Ли над тем же полем.
\end{frame}

\begin{frame}
\frametitle{2.5.6 Алгебры Ли (8/9)}

Пусть $\mathcal{V}$~--- векторное пространство 
с носителем $V$ над полем $\mathcal{F}$ с носителем
$F$.
Функция $\Map{L}{\mathcal{V}}{\mathcal{V}}$ 
называется 
\odmkey{линейным оператором}{Оператор (operator)!линейный (linear)} над векторным пространством,
или \odmkey{линейным преобразованием}{Преобразование (transformation)!линейное (linear)}, если  
$\A{\Bl{x},\Bl{y}\in V, a,b\in F}{L(a\Bl{x}+b\Bl{y})= aL(\Bl{x})+bL(\Bl{y})}$. 
\begin{example}
Множество полиномов одной вещественной переменной
над полем вещественных чисел является векторным пространством. Операция дифференцирования
является преобразованием этого 
пространства, причём дифференцирование линейно.     
\end{example}  

\medskip
Множество линейных преобразований векторного пространства
$\mathcal{V}$ (обычно обозначается 
$\End(\mathcal{V}))$, поскольку 
такие преобразования являются эндоморфизмами,  
\DMref{2.1.6.}{2.1.6.})  --- также является векторным пространством над $\mathcal{F}$.
В этом пространстве есть естественная операция умножения~--- суперпозиция преобразований, 
что позволяет ввести коммутатор: 
$$[\Bl{x}, \Bl{y}] \Def \Bl{x}\circ \Bl{y} - \Bl{y}\circ \Bl{x}.$$


\smallskip
\begin{remark}
Линейные преобразования играют очень большую роль в
теории векторных пространств, 
поэтому иногда 
термин \odmkey{линейное пространство}{Пространство (space)!линейное (linear)} используют как синоним 
термина векторное пространство.
\end{remark}

\end{frame}

\begin{frame}
\frametitle{2.5.6 Алгебры Ли (9/9)}

\begin{theorem}
 Пространство линейных преобразований векторного пространства  является алгеброй Ли.
\end{theorem}

\begin{proof}
Заметим, что пространство линейных преобразований с введённой операцией умножения задает ассоциативную алгебру над полем $\mathcal{F}$. Тогда в силу предыдущей теоремы получаем, что коммутатор является скобкой Ли, а значит мы действительно получаем алгебру Ли.
\end{proof}


\begin{digress}
Если $\mathcal{V}$~--- $n$-мерное  векторное пространство над полем 
$\mathcal{F}$, то  
пространство линейных преобразований
векторного пространства  имеет размерность $n^2$ и может быть представлено множеством квадратных матриц размера $n\times n$.
Применению преобразования соответствует
умножение вектора на матрицу, 
а суперпозиции преобразований
соответствует умножение матриц.
Все эти понятия и операции очень важны 
с прикладной точки зрения, 
поскольку образуют предмет \em{линейной алгебры}~--- математического аппарата, едва ли не чаще всего используемого в приложениях.
Буквально во всех предметных областях
естественных наук и техники используются 
вычислительные методы линейной алгебры.
Получается, 
что на первый взгляд абстрактные 
и оторванные от жизни определения,
с которых начинается этот параграф,
оказываются наиболее общим математическим аппаратом,
описывающим прикладные аспекты   
вычислительной науки.
\end{digress}

\end{frame}

\begin{frame}
\frametitle{Онтология общей алгебры}
\begin{center}
\includegraphics[height=7.8cm]{ODM2_5.png}
\end{center}
\end{frame}


\subsection{2.6. Решётки}
\begin{frame}
\frametitle{2.6. Решётки}

%\label{sect2-6}
Решётки иногда называют <<структурами>>, но слово <<структура>>
перегружено, и мы не будем использовать его в этом значении.
Решётки сами по себе часто встречаются в разных программистских
задачах, но ещё важнее то, что понятие решётки непосредственно
подводит к понятию булевой алгебры, значение которой для основ
современной двоичной компьютерной техники трудно переоценить.

\medskip
\begin{tabular}{p{11mm}|p{43mm}|p{89mm}}
  № & Название параграфа
  & Ключевые термины и обозначения \\
  \hline
\DMref{2.6.1}{2.6.1.} & Дистрибутивные и ограниченные решётки &  принцип двойственности, нижняя и верхняя грань решётки \\
\DMref{2.6.2}{2.6.2.} & Решётки с дополнением & \\	
\DMref{2.6.3}{2.6.3.}\RS & Частичный порядок в решетке & нижняя и верхняя границы, нижняя и верхняя грани
\\
\DMref{2.6.4}{2.6.4.} & Полные решётки & \\	
\DMref{2.6.5}{2.6.5.}\BS & Полурешётки & \\	
\DMref{2.6.6}{2.6.6.}\RS &	Булева алгебра	& \\
\end{tabular}  

\end{frame}

\begin{frame}
\label{2.6.1.}\subsubsection{2.6.1. Дистрибутивные и ограниченные решётки }\frametitle{2.6.1. Дистрибутивные и ограниченные решётки (1/3)}

\odmkey{Решётка}{Решётка (lattice)}~--- это множество $M$ с двумя
бинарными операциями, $\cap$ и $\cup$, такими, что выполнены
следующие условия (аксиомы решётки, см. \DMref{п.~2.1.4.}{2.1.4.})):

\begin{enumerate}
\item[(А1)] ассоциативность:
$(a \cap b)\cap c = a \cap (b \cap c)$,
$(a \cup b)\cup c = a \cup (b \cup c)$ ; 
\item[(А2)] коммутативность:
   $a \cap b = b \cap a$,
   $a \cup b = b \cup a$ ; 
\item[(А4)] поглощение:
 $(a \cap b)\cup a = a$,
$(a \cup b)\cap a = a$ ;
\item[(А5)] идемпотентность: $a \cap a = a$,
   $a \cup a=a$  . 
\end{enumerate}

\begin{example}
Булеан любого множества $M$ 
образует решётку $\Alg{2^M}{\cap,\cup}$ относительно пересечения и объединения.
\end{example}


\begin{remark}
В каждой паре аксиом решётки одна формула пары получается из другой
подстановкой $\cup\leftrightarrow\cap$, то есть аксиомы в каждой паре
\odmkey{двойственны}{Аксиома (axiom)!двойственная (dual)}  друг другу. Отсюда следует, что любое утверждение,
истинное для произвольных элементов решётки, остаётся истинным
при подстановке $\cup\leftrightarrow\cap$. Этот полезный факт
называется \odmkey{принципом двойственности}{Принцип (principle)!двойственности (of duality)} для решёток.         
\end{remark}

\begin{example}
Булеан любого множества $M$ также
образует решётку $\Alg{2^M}{\cup, \cap}$ относительно объединения и пересечения,
поскольку объединение и пересечение двойственны!
\end{example}
\end{frame}

\begin{frame}%[shrink]
\frametitle{2.6.1. Дистрибутивные и ограниченные решётки (2/3)}

Решётка называется \odmkey{дистрибутивной}{Решётка
(lattice)!дистрибутивная (distributive)},
если  $$a \cup (b \cap c) = (a \cup b)\cap (a \cup c),$$
   $$a \cap (b \cup c) = (a \cap b)\cup (a \cap c).$$
\begin{example}
Обе решётки, образуемые булеаном любого множества, дистрибутивны.
\end{example}


Если в решётке $\E{0\in M}{\A{a}{0\cap a =
0}}$, то 0 называется \odmkey{нулем}{Нуль!решётки (lattice
identity)} (или \odmkey{нижней гранью}{Грань (bound)!решётки
(lattice)!нижняя (low)}) решётки. Если в решётке $\E{1\in
M}{\A{a}{1 \cup a = 1}}$, то 1 называется
\odmkey{единицей}{Единица (identity)!решётки (of lattice)} (или
\odmkey {верхней гранью}{Грань (bound)!решётки (lattice)!верхняя
(upper)})
 решётки. Решётка с верхней и
нижней гранями называется \odmkey
{ограниченной}{Решётка (lattice)!ограниченная (bounded)}.
%\setcounter{theorem}{0}

\begin{example}
В решётке  $\Alg{2^M}{\cap,\cup}$ пустое множество $\emptyset$
образует нижнюю грань решётки, а множество $M$ образует верхнюю грань решётки. 
\end{example}

\begin{theorem} Если нижняя (верхняя) грань существует, то она
единственна.
\end{theorem}
\begin{proof}
Пусть $0'$~--- \!ещё один нуль решётки. Тогда $0' =
0\cap 0' = 0' \cap 0 = 0$.
\end{proof}
\end{frame}


\begin{frame}
\frametitle{ 2.6.1. Дистрибутивные и ограниченные решётки (3/3)}

\begin{example}
Множество натуральных чисел, 
где  $a\cap b\Assign \gcd(a, b)$   
(см.~\DMref{п.~2.4.2}{2.4.2.}) и \\$a\cup b\Assign \lcm(a, b)$
 (см.~\DMref{п.~2.4.3}{2.4.3.}) является дистрибутивной решёткой с нижней гранью.
В этой решётке число $1$ является нижней гранью, 
верхней грани эта решётка не имеет.
Идемпотентность и коммутативность
очевидны.
Ассоциативность проверяется с использованием свойств делимости. 
Пусть $x=\gcd(\gcd(a, b),c)$, $y=\gcd(a,\gcd(b, c))$. Тогда \\
$x|\gcd(a, b)\And x|c \Impl x|a \And x|b \And x|c \Impl x|a\And x|\gcd(b,c)
\Impl x|y$. Аналогично в обратную сторону и для второй операции.
Поглощение вытекает из теоремы \DMref{п.~2.4.3}{2.4.3.} и следствия к лемме из \DMref{п.~2.4.2}{2.4.2.}.
Дистрибутивность проще всего показать,
опираясь на следствие из основной теоремы арифметики 
(\DMref{п.~2.4.4}{2.4.4.}). 
\end{example}

\medskip
\begin{theorem}
$a \cap b = b \Equiv a \cup b = a$.
\end{theorem}
\begin{proof}
\proofsection{$\Impl$}
$a\cup b=a\cup(a\cap b)=(a\cap b)\cup a = a$.
\proofsection{$\Impll$}
$a\cap b=(a\cup b)\cap b=(b\cup a)\cap b=b$.
\end{proof}
\begin{corollary}
$ 0 \cup a = 1 \cap a =  a$.
\end{corollary}
\end{frame}

\begin{frame}
\label{2.6.2.}\subsubsection{2.6.2. Решётки с дополнением }\frametitle{2.6.2. Решётки с дополнением (1/2)}
%\subsection{Решётки с дополнением}
%\label{complattice}
\vspace{-2pt} 
 В ограниченной решётке элемент $a'$ называется
\odmkey{дополнением}{Дополнение (complement)} элемента $a$, если
$a \cap a'=0$ и $a \cup a'=1$. Если дополнение существует у любого элемента,
 то ограниченная решётка
называется \odmkey{решёткой с дополнением}{Решётка (lattice)!с
дополнением (complement)}.
 Вообще говоря, дополнение не обязано
существовать и не обязано быть единственным.

\begin{example}
Множество всех \em{нечётких} подмножеств некоторого множества $M$
образует, как нетрудно проверить, ограниченную дистрибутивную решётку
относительно пересечения и объединения 
(см. \DMref{п.~1.9.6}{1.9.6.}).
Однако, хотя для нечётких множеств дополнение определено, 
требуемые в решётке свойства дополнения не выполняются:
$\min (\mu_A, 1-\mu_A) \ne 0$, $\max (\mu_A, 1-\mu_A) \ne 1$.         
\end{example}

\begin{theorem} \NB{[о свойствах дополнения]} В
ограниченной дистрибутивной решётке с~дополнением выполняется
следующее:
\begin{enumerate}
\item[1)] дополнение $a'$ единственно;%
\item[2)] дополнение инволютивно: $a'' = a$;%
\item[3)] грани дополняют друг друга: $1' = 0$, $0' = 1$;%
\item[4)] законы де Моргана: ${(a \cap b)}' = a' \cup b'$, ${(a \cup b)}' = a' \cap b'$.%
\end{enumerate}
\end{theorem}
\end{frame}

\begin{frame}
\frametitle{2.6.2. Решётки с дополнением (2/2)}
\begin{proof} 
\proofsection{\color{blue}1)}
 Пусть $x,y$~--- дополнения $a$. Тогда
$x=x\cap 1=x\cap(a\cup y)=$ \\ $={(x\cap a)\cup(x\cap y)} =
{0\cup(x\cap y)} ={x\cap y}  = y\cap x = 0\cup(y\cap x)=$ \\
$=(y\cap a)\cup(y\cap x)  ={y\cap(a\cup x)} = {y\cap 1=y}$;

\proofsection{\color{blue}2)} $(a\cup a'=1\Impl a'\cup a=1$, $a\cap a'=0\Impl
a'\cap a=0)\Impl a=a''$; 

\proofsection{\color{blue}3)} $(1\cap 0=0, 0'\cap
0=0)\Impl 1=0'$, $(1\cup 0=1, 1\cup 1'=1)\Impl 0=1'$;

\proofsection{\color{blue}4)}
$(a \cap  b) \cap (a'\cup b')$ 
$=(a \cap  b \cap  a')\cup(a \cap  b \cap  b')=$\\ 
$=(0 \cap  b)\cup(a \cap  0)$ $=0\cup 0=0$,\\ 
$(a \cap  b)\cup(a'\cup b')= (a\cup a'\cup b') \cap (b\cup
a'\cup b')=$\\ 
$=(1\cup b') \cap (1\cup a')$  $= 1 \cap  1=1$,\\ 
$(a\cup b) \cap (a' \cap  b')=(a \cap  a' \cap  b')\cup(b \cap  a' \cap  b')=$\\ 
 $=(0 \cap  b')\cup(a' \cap  0)$ $=0\cup 0=0$,\\ 
$(a\cup b)\cup(a' \cap  b')$  
$=(a\cup b\cup a') \cap (a\cup b\cup b')=$\\ $=(1\cup b) \cap (1\cup a)=
1 \cap  1=1$.
\end{proof}

%\addvspace{6pt}
\end{frame}

\begin{frame}
\label{2.6.3.}\subsubsection{2.6.3. Частичный порядок в решётке }\frametitle{2.6.3. Частичный порядок в решётке (1/3)}
%\subsection{Частичный порядок в решётке}
В любой решётке можно естественным образом ввести частичный
порядок, а~именно: $a \prec b \Def a \cap b = a$.
%\setcounter{theorem}{0}
\begin{theorem} Отношение $\prec$ является
отношением частичного порядка.
\end{theorem}
%\addvspace{-9pt}
\begin{proof}
\proofsection{Рефлексивность} $a\cap a=a\Impl a\prec a$.
\proofsection{Антисимметричность} $a\prec b\And b\prec a\Impl
a=a\cap b= b \cap a=b$. 
\proofsection{Транзитивность} $a\prec
b\And b\prec c\Impl  a\cap c=(a\cap b)\cap c =a\cap(b\cap
c)=a\cap b=a \Impl a\prec c$.
\end{proof}
%\addvspace{-9pt}

Наличие частичного порядка в решётке не случайно, это её
характеристическое свойство. Более того, обычно решётку
определяют, начиная с частичного порядка, следующим образом. 

\medskip
Пусть $M$~--- частично упорядоченное множество с частичным
порядком $\prec$. 
Напомним, (см.~\DMref{п.~1.8.5}{1.8.5.}),
что элемент $x$ называется \odmkey{нижней
границей}{Граница (bound)!нижняя (low)}
для $a$ и $b$, если $x \prec a \And x
\prec b$. Аналогично, $y$ называется
\odmkey{верхней границей}{Граница (bound)!верхняя (upper)}
 для $a$ и $b$, если $a \prec y \And b
\prec y$. \\ Здесь $a, b, x, y \in M$.
\end{frame}

\begin{frame}
\frametitle{2.6.3. Частичный порядок в решётке (2/3)}

Элемент $x$ называется 
\odmkey{нижней гранью}{Граница (bound)!нижняя (low)!наибольшая (greatest)}
 элементов $a$ и $b$, если $x$~--- нижняя
граница элементов $a$ и $b$ и для любой другой нижней границы $v$
элементов $a$ и $b$ выполняется $v \prec x$. Обозначение: $x =
\inf (a,b)$. 

\medskip
Аналогично, $y$ называется 
\odmkey{верхней гранью}{Граница (bound)!верхняя
(upper)!наименьшая (least)} элементов $a$ и $b$, если $y$~---
верхняя граница элементов $a$ и $b$ и для любой другой верхней
границы $u$ элементов $a$ и $b$ выполняется $y \prec u$.
Обозначение: $y = \sup(a,b)$.


\begin{lemma} Если нижняя (верхняя) грань любых двух элементов существует, 
то она единственна.
\end{lemma}
%\addvspace{-12pt}
\begin{proof}
$ x = \inf(a,b) \And y = \inf(a,b) \Impl y \prec x \And x \prec y \Impl 
 x = y $.
\end{proof}


\begin{theorem} Если в частично упорядоченном множестве для любых двух
элементов существуют нижняя и верхняя грани, то это множество
образует решётку относительно $\inf$ и $\sup$.
\end{theorem}
%\addvspace{-12pt}
\end{frame}

\begin{frame}
\frametitle{2.6.3. Частичный порядок в решётке (3/3)}
\begin{proof}
Пусть  $x \cap y\Assign \inf(x,y)$, $x \cup y\Assign \sup(x,y)$.
\proofsection{\color{blue}1}~$\inf(x,x) = x \Impl x \cap x = x$, 
$\sup(x,x) = x \Impl x \cup x = x$.
\proofsection{\color{blue}2}~$\inf(x,y) = \inf(y,x) \Impl x \cap y = y \cap x$,
$\sup(x,y) = \sup(y,x) \Impl x \cup y = y \cup x$.
\proofsection{\color{blue}3}~$\inf\bigl(x,\inf(y,z)\bigr) = \inf\bigl(\inf(x,y),z\bigr)$
$\Impl x \cap (y \cap z) = (x \cap y) \cap z$,\\ 
$\sup\bigl(x, \sup(y,z)\bigr) = \sup\bigl(\sup(x,y),z\bigr) \Impl 
x \cup (y \cup z) = (x \cup y) \cup z$.
\proofsection{\color{blue}4}~$ \sup\bigl(\inf(a, b), a\bigr) = a \Impl (a \cap
 b) \cup a = a $, $ \inf\bigl(\sup(a, b), a\bigr) = a \Impl
 (a \cup b) \cap a = a $.
\end{proof}

\smallskip
\begin{example}
Включение $\subset$ является естественным частичным порядком
на булеане $\Alg{2^M}{\cap,\cup}$, поскольку
$A\subset B\Equiv A\cap B = A$.
\end{example}

\end{frame}

\begin{frame}
\label{2.6.4.}\subsubsection{2.6.4. Полные решётки }\frametitle{2.6.4. Полные решётки (1/2)}

Частичный порядок можно естественным образом ввести в любой решётке,
но при этом не всякое подмножество носителя решётки
обязано иметь верхнюю и нижнюю грани.
Решётка называется \odmkey{полной}{Решётка (lattice)!полная (complete)},
если любое подмножество имеет супремум и инфимум.
Пол\-ная решётка $L$ является ограниченной, 
если положить : $\Bl{0} \Assign \inf L$,
$\Bl{1} \Assign \sup L$. Кроме того, ясно, что полная решётка является линейно полным 
(\DMref{п.~1.8.5}{1.8.5.}) частично упорядоченным множеством.

\smallskip
\begin{example}
Булеан является полной решёткой, поскольку если 
$X\subset 2^M$, то 
$$\sup X = \bigcup\limits_{Y\in X} Y, \quad
\inf X = \bigcap\limits_{Y\in X} Y.$$
\end{example}

\smallskip 
Свойство полноты решётки является более сильным, 
чем свойство линейной полноты в частично упорядоченных множествах, 
поэтому в полной решётке справедлива теорема 
о наименьшей неподвижной точке функции (\DMref{п.~1.8.7}{1.8.7.})
при более слабом предположении о свойствах функции.

\end{frame}

\begin{frame}
\frametitle{2.6.4. Полные решётки (2/2)}
\begin{theorem}
Если $L$ {\upshape ---} полная решётка, а функция $\Map{f}{L}{L}$ монотонна,
то функция $f$ имеет наименьшую неподвижную точку.
\end{theorem}
\begin{proof}
Рассмотрим множество $A \Assign \Set{x \in L}{f(x)\prec x}$ и
элемент $a  = \inf A$.
Имеем $\A{x \in A}{a \prec x}$, и~значит, $\A{x \in A}{f(a) \prec
f(x) \prec x}$, следовательно, $f(a)$ --- нижняя граница для
множества $A$ и $f(a) \prec a$. Далее, $f(f(a)) \prec f(a)$, и 
значит, $f(a)\in A$, откуда по определению элемента $a$ имеем $a
\prec f(a)$ и по антисимметричности $f(a)=a$. 
Таким образом, $a$ --- неподвижная точка функции
$f$. Пусть теперь $b$ --- любая неподвижная точка функции $f$, то
есть $f(b)=b$. По рефлексивности естественного частичного порядка
 $f(b) \prec b$. Тем
самым $b \in A$ и значит $a \prec b$.
\end{proof}
\end{frame}

\begin{frame}
\label{2.6.5.}\subsubsection{2.6.5. Полурешётки }
\frametitle{2.6.5. Полурешётки (1/4)}
%\subsection{Полурешётки}
Естественный частичный порядок в решётке можно определить
одним из двух способов:
$$a \prec b := a \cap b = a\quad \text{или}\quad a \prec b := a \cup b = b.$$
Из теоремы \DMref{п.~2.6.1}{2.6.1.} следует, что это одно и то же отношение. Таким образом, для определения частичного порядка достаточно иметь только одну подходящую операцию.

Множество $M$ c операцией $\Map{\cap}{M \times M}{M}$ называется
\odmkey{полурешёткой}{Полурешётка (semi-lattice)}, если выполнены
следующие условия (аксиомы полурешётки):
%\begin{tabbing}
%   3. \= $a \cap (b \cap c) = (a \cap b) \cap c$\quad \=ассоциативность.\kill
%    1. \> $a\cap a = a$ \>идемпотентность, \\
%    2. \> $a\cap b = b \cap a$ \> коммутативность, \\
%    3. \> $a \cap (b \cap c) = (a \cap b) \cap c$ \>ассоциативность.
%\end{tabbing}
\begin{enumerate}
\item[(А1)] ассоциативность: $\;a \cap (b \cap c) = (a \cap b) \cap c$ ;
\item[(А2)] коммутативность: $a\cap b = b \cap a$ ;
\item[(А5)] идемпотентность: $\,a\cap a = a$ .
\end{enumerate}

Такой набор свойств встречается достаточно часто, так что многие операции образуют полурешётку.

\begin{example}
Операции объединения и пересечения множеств образуют полурешётки.
\end{example}
\end{frame}

\begin{frame}
\frametitle{2.6.5. Полурешётки (2/4)}
Как видно из доказательства теоремы  
\DMref{п.~2.5.4}{2.5.4.}, аксиом
полурешётки достаточно для определения естественного частичного порядка
$a \prec b \Assign a\cap b = a$.

\begin{remark}
Для доказательства рефлексивности используется идемпотентность, для доказательства антисимметричности --- коммутативность, а транзитивность следует из ассоциативности.
\end{remark}

Функция $\Map{f}{M}{M}$, определённая на полурешётке $M$, называется
\odmkey{ди\-стри\-бу\-тивной}{Функция (function)!дистрибутивная (distributive)},
если $f(x\cap y) = f(x) \cap f(y)$.

\begin{remark}
В терминах \DMref{п.~2.1.5}{2.1.5.} дистрибутивная функция --- это гомоморфизм.
\end{remark}

Нетрудно видеть, что дистрибутивная функция $f$ монотонна:
$$
x \prec y \Impl x = x\cap y \Impl f(x) = f(x\cap y) =
f(x) \cap f(y) \Impl f(x) \prec f(y).
$$

\begin{example} Пересечение в булеане монотонно, а разность --- нет.
\end{example}

\smallskip 
Полурешётка называется \odmkey{ограниченной}{Полурешётка
(semi-lattice)!ограниченная (bounded)}, если у неё существуют
верхняя и нижняя грань.
\end{frame}

\begin{frame}
\frametitle{2.6.5. Полурешётки (3/4)}
\vspace{-2pt}
Поскольку в полурешётке определён естественный частичный порядок, 
опре\-делены и все другие понятия:
%%% Я забыл, как делаются маркированные списки. Новиков
%%% Далее идет маркированный список
\begin{itemize}
\item верхние и нижние границы,
\item верхние и нижние грани (\em{супремум} и \em{инфимум}),
\item максимальные и минимальные элементы,
\item наибольшие и наименьшие элементы,
\item линейная полнота,
\item конечная высота.
\end{itemize}

\begin{theorem}
Если $L$ {\upshape ---} ограниченная полурешётка конечной высоты, а
$\Map{f}{L}{L}$~{\upshape ---} монотонная функция, то функция $f$ имеет
наименьшую неподвижную точку, которая может быть получена методом
итераций из нижней грани.
\end{theorem}
\begin{proof}
В ограниченной полурешётке нижняя грань $\Bl{0}$
является наименьшим элементом. Поскольку полурешётка имеет конечную высоту, всякое линейно упорядоченное подмножество можно рассматривать как монотонную последовательность, которая заканчивается наибольшим элементом (супремумом). 
\end{proof}

\end{frame}

\begin{frame}
\frametitle{2.6.5. Полурешётки (4/4)}
\begin{proof} (Окончание) 
В частности, последовательность
$\langle f^0(\Bl{0}), f^1(\Bl{0}), f^2(\Bl{0}),\dots \rangle$
монотонна в силу монотонности функции $f$ и имеет супремум
$a = \sup\Set{f^n(\Bl{0})}{n\ge 0}$, 
который является наименьшей неподвижной точкой 
(см. доказательство теоремы о наименьшей неподвижной точке в \DMref{п.~1.8.7}{1.8.7.}).
\end{proof}

%\vspace*{-.5\baselineskip}

\begin{remark}
Можно показать, что множество неподвижных точек монотонной функции на ограниченной полурешётке конечной высоты само образует ограниченную полурешётку конечной высоты.
\end{remark}

\begin{example}
Рассмотрим конечное множество $A$ и булеан $2^A$. 
Ясно, что $2^A$ является
ограниченной полурешёткой конечной высоты относительно
пересечения. 
Рассмотрим функцию $\Map{f}{2^A}{2^A}$, где $f(X)
\Assign X+a$. Наименьшей неподвижной точкой этой функции является
$\{a\}$. 
Действительно, $f(\emptyset)=a$, $f(a)=a$,
$f(f(a))=a,\dots$. Более того, все подмножества $2^A$, содержащие
элемент $a$, являются неподвижными точками функции $f$ и~образуют
ограниченную полурешётку конечной высоты.
\end{example}
%\vspace*{-.5\baselineskip}


%%%%%%%%%%%%%%%%%%%%%%%%% Конец изменений 2012

%\addvspace{12pt}
\end{frame}

\begin{frame}
\label{2.6.6.}\subsubsection{2.6.6. Булевы алгебры }\frametitle{2.6.6. Булевы алгебры (1/2)}

%\subsection{}
%\label{boolalg} 
Дистрибутивная ограниченная решётка, в которой для
каждого элемента су\-ще\-ствует дополнение, называется \odmkey{булевой
алгеброй}{Алгебра (algebra)!булева (boolean)}. Свойства
булевой алгебры:
\begin{enumerate}
\item[1)] $a\cup a=a$, \;\;$a\cap a=a$~---  по определению решётки;
 
\item[2)] $a\cup b=b\cup a$, $a\cap b=b\cap a$~--- по определению решётки;

\item[3)] $a\cup(b\cup c)=(a\cup b)\cup c$, $a\cap(b\cap c)=(a\cap b)\cap c$~---
 по определению решётки; 
 
\item[4)] $(a\cap b)\cup a=a$, $(a\cup b)\cap a=a$~---
 по определению решётки; 
 
\item[5)] $a\cup(b\cap c)=(a\cup b)\cap(a\cup c)$, 
$a\cap(b\cup c)=(a\cap b)\cup(a\cap c)$~---
 по  дистрибутивности;
 
\item[6)] $a\cup 1=1$, $a\cap 0=0$~--- по ограниченности;

\item[7)] $a\cup 0=a$, $a\cap 1=a$~---
 по следствию из теоремы~2 \DMref{п.~2.6.1}{2.6.1.};
 
\item[8)] $a''=a $~--- по теореме \DMref{п.~2.6.2}{2.6.2.};

\item[9)] $(a\cap b)'=a'\cup b'$, $(a\cup b)'=a'\cap b'$ 
 по теореме \DMref{п.~2.6.2}{2.6.2.};
 
\item[10)] $a\cup a'=1$, $a\cap a'=0$~---  по определению.
%\end{tabbing}
\end{enumerate}
\end{frame}

\begin{frame}
\frametitle{2.6.6. Булевы алгебры (2/2)}

\begin{examples}
%\begin{enumerate}

{\color{blue}1.} $\Alg{2^M}{\cap, \cup, ^-}$~--- булева алгебра, причём
$M$~--- верхняя грань, $\emptyset$~--- нижняя грань, $\subset$~---
естественный частичный порядок. 

{\color{blue}2.} $\Alg{E_2}{\And, \Or,
\Not}$~--- булева алгебра, где $\And$~--- конъюнкция, $\Or$~---
дизъюнкция, $\Not$~--- отрицание, причём $1$~--- верхняя грань,
$0$~--- нижняя грань, а импликация $\Impl$~--- естественный
частичный порядок. 

{\color{blue}3.} Пусть $T $~--- некоторое
 конечное множество взаимно простых чисел.
Пусть $P$~--- множество всех возможных произведений различных
чисел из $T$, включая пустое произведение, по определению равное
1. Заметим, что $P\sim 2^T$. Тогда 
$\langle P; \gcd, \lcm, ()^\prime\rangle$~--- булева алгебра, 
где $\gcd$~---
наибольший общий делитель, 
$\lcm$~--- наименьшее общее кратное,~а
$
(n)^\prime\Def
\frac{1}{n}\prod\limits_{q\in T}q,
$
причём $\prod\limits_{q\in T} q$~--- верхняя грань,
 $1$~--- нижняя грань,
отношение <<$n|m$>>~--- естественный частичный порядок.
%\end{enumerate}
\end{examples}

\end{frame}


\begin{frame}
\frametitle{Онтология общей алгебры}
\begin{center}
\includegraphics[height=8cm]{ODM2_6.png}
\end{center}
\end{frame}

\subsection{2.7. Матроиды и жадные алгоритмы}
\begin{frame}
\frametitle{2.7. Матроиды и жадные алгоритмы (1/2)}

В этом разделе рассматривается следующая простая и
часто встречающаяся экстремальная задача. Задано конечное
множество $E$, элементам которого приписаны положительные веса, и
семейство $\mathcal{E}\subset 2^E$ подмножеств множества $E$.
Требуется найти в семействе $\mathcal{E}$ элемент (подмножество
$E$) максимального суммарного веса. Оказывается, что если
семейство $\mathcal{E}$ обладает особой структурой (является
матроидом), то для решения задачи достаточно применить очень
простой и эффективный алгоритм (\em{жадный алгоритм}). Ясное понимание
природы и области применимости жадных алгоритмов совершенно
необходимо каждому программисту. Матроиды, рассматриваемые в этом
разделе, вообще говоря, не являются алгебраическими структурами в
том смысле, который  придан этому понятию в первом разделе данной
главы. Однако матроиды имеют много общего с рассмотренными
алгебраическими структурами (в частности, с линейно независимыми
множествами в векторных пространствах и модулях) и изучаются
сходными методами.

\end{frame}

\begin{frame}
\frametitle{2.7. Матроиды и жадные алгоритмы (2/2)}


\medskip
\begin{tabular}{p{11mm}|p{43mm}|p{89mm}}
  № & Название параграфа
  & Ключевые термины и обозначения \\
  \hline
\DMref{2.7.1}{2.7.1.} & Матроиды & \\	
\DMref{2.7.2}{2.7.2.} & Максимальные независимые подмножества & ранг, неравенство полумодулярности \\
\DMref{2.7.3}{2.7.3.} & Базисы & \\	
\DMref{2.7.4}{2.7.4.} & Циклы в матроидах & \\
\DMref{2.7.5}{2.7.5.}\RS & Жадный алгоритм & весовая функция \\
\DMref{2.7.6}{2.7.6.} & Примеры применения матроидов & матроид трансверсалей, частичная трансверсаль
\end{tabular}  

\end{frame}

\begin{frame}
\label{2.7.1.}\subsubsection{2.7.1. Матроиды }\frametitle{2.7.1. Матроиды (1/2)}

\vspace{-2pt}
\odmkey{Матроидом}{Матроид (matroid)} $M=\lrle{E,\cal
E}$ называется пара, состоящая из конечного множества $E$, $\vert
E\vert=n$, и семейства его подмножеств $\cal E \subset 2^E$,
таких, что выполняются следующие три аксиомы: 
\begin{enumerate}
\item[$M_1\!:$]  $\emptyset\in \cal E$;
\item[$M_2\!:$] $A\in\cal E\And B\subset A\Impl B\in\cal E$;
\item[$M_3\!:$] $A,B\in\cal E\And\vert B\vert=\vert A\vert+1\Impl\E{e\in
B\setminus A}{A+e \in\cal E}$.
\end{enumerate}
 Элементы множества $\cal E$ называются
\odmkey{независимыми}{Множество (set)!независимое (independent)},
а элементы множества
$2^E\setminus\cal E$~--- \odmkey{зависимыми}{Множество
(set)!зависимое (dependent)} множествами.

\smallskip
\begin{remark}
Аксиома $M_1$ исключает из рассмотрения вырожденный случай $\cal E=\emptyset$.\\
Аксиома $M_3$ формулируется для \em{любых} независимых
множеств
$A, B$, в том числе допускается, что 
$A\cap B\ne \emptyset$ или даже $A\subset B$.
В условиях действия аксиомы $M_2$
условие в аксиоме $M_3$ до 
$\vert A\vert  < \vert B\vert$,
но принятое условие  $\vert B\vert=\vert A\vert+1$
удобнее использовать в доказательствах

\end{remark} 

\smallskip
\begin{digress}
Понятие матроида возникло
в результате исследования линейной независимости
в~векторных пространствах,
но нашло приложение в программировании.
\end{digress}

%\vspace{-0.5\baselineskip}
\end{frame}

\begin{frame}
\frametitle{2.7.1. Матроиды (2/2)}
\begin{example}
Семейство линейно независимых множеств векторов любого векторного
пространства является матроидом. Действительно, 
аксиомы $M_1$ и $M_2$ выполняются. 
Пусть $A{\Assign}\{\Tuple{\Bl a}{1}{m}\}$ и
$B\Assign\{\Tuple{\Bl b}{1}{m+1}\}$~--- линейно независимые
множества векторов. Если бы все векторы из множества $B$
выражались в виде линейной комбинации векторов из множества~$A$,
то множество $B$ было бы линейно зависимым по следствию к теореме
\DMref{п.~2.5.3}{2.5.3.}. 

Стало быть, среди векторов множества
$B$ есть по крайней мере один вектор $\Bl b$, который не входит в
множество $A$ и не выражается в виде линейной комбинации векторов
из множества $A$.
 Добавление вектора $\Bl b$ к множеству $A$
образует линейно независимое множество.
\end{example}

\end{frame}

\begin{frame}
\label{2.7.2.}
\subsubsection{2.7.2. Максимальные независимые подмножества }
\frametitle{2.7.2. Максимальные независимые подмножества (1/4)}

Пусть $M=\lrle{E,\cal E}$~--- матроид и $X\subset
E$~--- произвольное множество. \odmkey{Максимальным}{Множество
(set)!независимое (independent)!максимальное (maximal)}
независимым подмножеством множества $X$ называется множество $Y$,
такое, что
$$ Y\subset X\And Y\in\cal E\And \A{Z\in\cal E}
{Y\subset Z\subset X \Impl Z= Y}. $$
Множество максимальных независимых подмножеств
множества $X$ обозначим~$\underline{X}$:
$$
\underline{X}\Def \Set{Y\!\subset\!E}{Y\!\subset\! X\And Y\!\in\!\cal E\And
\A{Z\in\cal E} {Y\subset Z\subset X \Impl Z= Y}}.
$$
Рассмотрим следующее утверждение: 

\settowidth{\leftmargini}{$M_1\colon$}
\addtolength{\leftmargini}{\labelsep}
\begin{enumerate}
\item[$M_4\colon$]
$\quad \A{X}{Y\in\underline{X}\And Z\in\underline{X}
\Impl \vert Y\vert=\vert Z\vert},$
\end{enumerate}

то~есть все максимальные независимые подмножества данного
множества рав\-но\-мощны.

\begin{theorem}
Пусть $M=\lrle{E,\cal E}$ и выполнены аксиомы $M_1$ и $M_2$. Тогда
аксиома $M_3$ и~утверждение $M_4$ эквивалентны, то есть
$$
A,B\in\cal E\And
\vert B\vert=\vert A\vert+1\Impl
\E{e\in B\setminus A}{A+e \in\cal E},
$$
тогда и только тогда, когда
$
\A{X}{Y\in\underline{X} \And Z\in\underline{X}\Impl\vert
Y\vert=\vert Z\vert}.
$
\end{theorem}

\end{frame}

\begin{frame}
\frametitle{2.7.2. Максимальные независимые подмножества (2/4)}

\begin{proof}
\proofsection{$\Impl$} Пусть выполнены утверждения $M_1$, $M_2$,
$M_3$ (то~есть $M$~--- матроид). Покажем от противного, что
выполняется и $M_4$. Пусть $Y,Z\in\underline{X}$, $\vert
Y\vert\not=\vert Z\vert $ и для определённости $\vert Y\vert>\vert
Z\vert$. Возьмём $Y'\subset Y$, так что $\vert Y'\vert=\vert
Z\vert +1$. Тогда по свойству $M_3$ имеем \\
 $\E{e\in Y'\setminus
Z}{W\Assign Z+e\in\cal E}$. Таким образом, имеем $W\in\cal E$,
${Z\subset W}$, $W\subseteq X$, $Z\ne W$, что противоречит
предположению $Z\in\underline{X}$. 
\proofsection{$\Impll$} Пусть
выполнены утверждения $M_1$, $M_2$, $M_4$. Покажем от противного,
что выполняется и $M_3$. Возьмём $A,B\in\cal E$, так что $\vert
B\vert= \vert A\vert+1$. Допустим, что \\$\neg\E{e\in B\setminus
A}{A+e\in\cal E}$, то~есть $\A{e\in B\setminus A}{A+e\notin\cal
E}$. Рассмотрим  $C\Assign A\cup B$. Имеем $A\in\underline{C}$. Но
$B\in\cal E$, поэтому $\E{B'}{B\subset B'\And B'\in\cal E\And
B'\in \underline{C}}$. По условию $M_4$ имеем $\vert B'\vert =
\vert A\vert$. Но $\vert A\vert = \vert B'\vert \ge \vert B\vert =
\vert A\vert +1$~--- противоречие.
\end{proof}

\medskip
\begin{remark}
Таким образом, $M_1,M_2,M_4$~--- эквивалентная система аксиом матроида.
\end{remark}

\medskip
Поскольку в матроиде максимальные независимые подмножества любого множества имеют
одну и ту же мощность, возможно определить \odmkey{ранг}{Ранг (rank)!множества (of a set)} любого множества
как функцию:
$$r(A) \Def \max \{|B| \mid B \subset A \And B \in \mathcal{E} \}.$$
\end{frame}

\begin{frame}
\frametitle{2.7.2. Максимальные независимые подмножества (3/4)}

%
%{\textbf{Лемма:}}Дан матроид $M = \langle E,\varepsilon \rangle$⟩ и множество  $A \subset E$. Пусть $ B \subset A, B\in \varepsilon$, тогда существует $D:B \subset D \subset A,D \in \varepsilon,|D|=r(A)$.
%
%{\textbf{Доказательство:}}Пусть S — подмножество A такое, что $r(A)=|S|,S \in \varepsilon$ (по определению ранговой функции такое S всегда существует).
%
%Предположим, что лемма неверна и максимальное независимое подмножество, которое мы можем получить из B добавляя элементы из A — это C, причем $|C|<r(A)$. Тогда имеем: $C \in \varepsilon,S \in \varepsilon,|C|<|S|$, следовательно существует элемент $x \in S \backslash C:C \cup \{x\} \in \varepsilon$. Заметим также что $|C \cup x| = |C|+1 >|C|$ и $x \in A$, т.к.$ S \backslash C \subset A, B \subset C \cup\{x\}$. Итак пришли к противоречию, мы получили множество большее по мощности, чем C такое, что $B\subset C'\subset A,C' \in \varepsilon$, значит исходное предположение было не верно, и мы можем найти множество D удовлетворяющее необходимым условиям.

Из определения ранговой функции ясно, что ранг $r(A)$~--- это мощность максимального независимого
подмножества $B\subset A$, причём если $A\in \cal E$, то $B=A$, 
а если $A \not\in \cal E$  то $B$ не обязательно единственно.   

\begin{theorem}
Пусть $M = \langle E, \cal E \rangle$~--- матроид, 
и $\Map{r}{2^E}{\Bbb{N}}$~--- его ранговая функция. Тогда для любых множеств $A,B \subset 2^E $ выполняется следующее:
\begin{enumerate}
\item $0 \le r(A) \le |A|$;
\item $A \subset B \Impl r(A) \le r(B)$;
\item $r(A \cup B)+r(A \cap B)\le r(A)+r(B)$ (\odmkey{неравенство полумодулярности}{Неравенство (inequality)!полумодулярности (submodularity)}).
\end{enumerate}
\end{theorem}

\begin{proof}
\proofsection{1.} Очевидно из определения: максимальное независимое 
подмножество по мощности не может быть больше самого множества и меньше нуля.
\proofsection{2.} Пусть $C \subset A$~--- максимальное независимое подмножество. 
Так как $A \subset B$, то $C \subseteq B$~--- независимое подмножество. 
Поэтому $r(B) \ge |C|$ по определению, а значит $r(B) \ge r(A).$
\end{proof}
\end{frame}

\begin{frame}
\frametitle{2.7.2. Максимальные независимые подмножества (4/4)}

\begin{proof} (Продолжение)
\proofsection{3.} Рассмотрим множество $D_{1} \subset A \cap B \And D_{1} \in \cal E \And 
|D_{1}|=r(A \cap B)$, такое множество всегда существует по определению функции $r$. 
Дополним множество $D_{1}$ элементами из
множества $B \setminus D_{1}$ 
до множества $D_{2}$ так, что 
$|D_{2}| = r(B) \And D_{2} \in \cal E$.
Далее дополним $D_{2}$ элементами из 
$A \cup B \setminus D_{2}$ до множества $D_{3}$
 так, что $|D_{3}|=r(A \cup B) \And  D_{3} \in \cal E$. 
 Заметим, что на последнем шаге будут добавляться только элементы из $A$.  Действительно,
 пусть на третьем этапе мы взяли $x\in B$, 
 тогда $\{x\} \cup D_{2} \subset D_{3}\And
 D_{3} \in \cal E$, следовательно 
 $\{x\} \cup D_{2} \in \cal E$ (по определению матроида), а также $|\{x\} \cup D_{2}|=|D_{2}|+1=r(B)+1$, что невозможно по определению 
 $r$.\\
Заметим также, что
$(D_{3}\setminus D_{2})\cup D_{1} \subset A$, $(D_{3}\setminus D_{2})\cup D_{1} \in \cal E$
(по определению матроида), значит по определению ранговой функции:
$r(A)\ge |(D_{3} \setminus D_{2})\cup D_{1}|=|D_{3}|-|D_{2}|+|D_{1}|$
Заменяя мощности на ранги имеем:
$r(A)+r(B) \ge r(A \cup B)+r(A \cap B)$.
\end{proof}

\medskip
\TODO{Пример желателен}
\end{frame}

\begin{frame}\label{2.7.3.}
\subsubsection{2.7.3. Базисы матроида}
\frametitle{2.7.3. Базисы матроида(1/2)}


Максимальные независимые подмножества множества $E$ называются
\odmkey{базисами}{Базис (basis)!матроида (of matroid)} матроида
$M=\lrle{E,\cal E}$. Всякий матроид имеет базисы, что видно из
следующего алгоритма.

\begin{algorithmic}
\REQUIRE Матроид $M=\lrle{E,\cal E}$.
\ENSURE Множество $B\subset E$ элементов, образующих базис.
\STATE $B\Assign\emptyset$ 
\COMMENT{\color{gray}\small вначале базис пуст }
\FOR{$e\in E$}
\IF{$B+e\in\cal E$}
\STATE $B\Assign B+e$ 
\COMMENT{\color{gray}\small расширяем базис допустимым образом }
\ENDIF
\ENDFOR
\end{algorithmic}

\end{frame}

\begin{frame}
\frametitle{2.7.3. Базисы матроида(2/2)}
\begin{ground} Пусть $B_0=\emptyset$, $B_1,...,$ $B_k=B$~---
последовательность значений переменной $B$ в процессе работы
алгоритма. По построению $\A{i}{B_i\in\cal E}$. Пусть
$B {\not\in}\underline{E}$. Тогда
%$\E{B'}{B\subset B'\And B'\not=B\And B'\in\cal E}$.
$\exists\,B'\;(B\subset B'\And B' {\not=}B\And B'\in\cal E)$.
Возьмём $B''\subset B'$, так что $\vert B'' \vert = \vert B
\vert + 1$, $B\subset B''$ и $B''\in \cal E$. Рассмотрим $e\in
B'\setminus B$. Элемент $e\not\in B$, значит  
$\E{i}{(B_{i-1}+e)\notin \cal E}$. Но $e\in B'' \And B_{i-1}\subset
B''\Impl \left(B_{i-1}+e\right)\subset B''$. По аксиоме $M_2$ имеем
$\left( B_{i-1}+e\right)\in \cal E$~--- противоречие.
\end{ground}

\smallskip
\begin{corollary}
Все базисы матроида равномощны.
\end{corollary}

\medskip
Ранг носителя матроида~---
 это размерность его базиса. Ранговая функция $r(E)$ определяет размер максимального независимого подмножества носителя $Е$. 
\end{frame}

\begin{frame}
\subsubsection{2.7.4. Циклы в матроидах }\label{2.7.4.}
\frametitle{2.7.4. Циклы в матроидах (1/4)}

Минимальное \em{зависимое} подмножество $C$ множества $E$ называется \odmkey{циклом}{Цикл (cycle)!матроида (of matroid)} матроида \\$M=\lrle{E,\cal E}$. 
Множество циклов матроида обозначим $\cal C$.
Другими словами
$$
\cal C \Def \Set{C\subset E}{C \not\in \cal E 
\And (X\SubsetNE C \Impl X\in \cal E)}.
$$

\begin{theorem}
Пусть задан матроид $M=\lrle{E,\cal E}$ 
и его семейство циклов $\cal C$. Тогда
\begin{enumerate}
\item $\emptyset\not\in \cal C$;
\item $\A{C_1, C_2 \in \cal C}{C_1\! \neq\! C_2 \Impl 
(C_1\! \not\subset\! C_2 \And C_2\! \not\subset\! C_1)}$;
\item $\A{C_1, C_2 \in \cal C}{(C_1\! \neq\! C_2 \And \E{e\!\in\! E}{e\in C_1\!\cap\! C_2})\!\Impl\! 
(\E{C\in \cal C}{C\subseteq (C_1\!\cup\!C_2)\!\setminus\!\{e\}})} $.
\end{enumerate}
\end{theorem}
\begin{proof} Напомним, что выполнены аксиомы $M_1$, $M_2$, $M_3$.
\proofsection{1} По аксиоме $M_1$ имеем $\emptyset\in \cal E$ 
 и значит   $\emptyset\not\in \cal C$.
\proofsection{2} От противного. Если $C_1\SubsetNE C_2$,
то $C_1\in \cal E$. Противоречие. Аналогично $C_2 \not\subset C_1$.
\proofsection{3} От противного. Обозначим  
$D\Assign (C_1\cup C_2)\setminus \{e\}$, $A\Assign C_1\cap C_2$
и допустим, что $D\in \cal E$. Тогда 
из предыдущего пункта теоремы следует, что
$|C_1\setminus C_2|>0$ и $|C_2\setminus C_1|>0$,
а значит $|D| = |(C_1\cup C_2)| - |\{e\}|$
$=|C_1\setminus C_2|  + |C_2\setminus C_1| + |C_1\cap C_2| - |\{e\}|\geq$\\
$\geq 1+1 +|A| -1 > |A|$, причем $A\in \cal E$. Применив нужное число раз
аксиому  $M_3$ заключаем, что $\E{B\subset E}{|B|=|D| \And 
A\SubsetNE B \And B\in \cal E}$.  
\end{proof}
\end{frame}

\begin{frame}
\frametitle{2.7.4. Циклы в матроидах (2/4)}
\begin{proof} (продолжение)
Поскольку $C_1 \in \cal C \And B\in \cal E$,
имеем $C_1\not\subset B$ и значит 
$\E{b\in C_1\setminus A}{b\not\in B}$,
следовательно, в множестве $B$ лежит не более, чем 
$|C_1\setminus A|-1$ элементов из множества $C_1\setminus A$.
Аналогично, в множестве $B$ лежит не более, чем \\
$|C_2\setminus A|-1$ элементов из множества $C_2\setminus A$.
Tогда $|B|\leq |A| + |C_1\setminus A| - 1 + |C_2\setminus A| - 1 =$ $=|C_1\cup C_2|- 2 = |D|-1 $, что противоречит равенству $|B|=|D|$. 
\end{proof}
\medskip
Понятие цикла даёт третий способ определения матроида следующим образом.

 \odmkey{Матроид}{Матроид (matroid)}~--- это пара 
$\langle E, \cal C\rangle$, где $E$~--- носитель матроида, а $\cal C$~--- семейство непустых подмножеств $E$, называемое множеством циклов матроида, для которых выполняются следующие условия:

\begin{enumerate}
\item Ни один цикл не является подмножеством другого.

\item Если есть два различных цикла $C_1$ и $C_2$, 
имеющих общий элемент $e\in (C_{1}\cap C_{2}) $, то $(C_{1}\cup C_{2})\setminus \{e\}$ содержит цикл.
\end{enumerate}
 
\begin{theorem}
Множество циклов однозначно определяет матроид. 
\end{theorem} 
 
\begin{proof} Пусть задано семейство непустых циклов $\cal C$,
удовлетворяющих указанным условиям. Определим семейство 
$\cal E \Def \Set{X\subset E}{\Not\E{C\in \cal C}{C\subset X}}.
$ 
\end{proof} 
 \end{frame}

\begin{frame}
\frametitle{2.7.4. Циклы в матроидах (3/4)}
\begin{proof} (продолжение) Другими словами,
определим семейство всех множеств, не содержащих циклов.
Тогда семейство $\cal E$ образует матроид.
\proofsection{$M_1$}
Пустое множество не содержит ничего, 
поэтому $\emptyset \in \cal E$.
\proofsection{$M_2$}
Если множество не содержит циклов, то и любое его 
подмножество не содержит циклов, и вторая аксиома выполнена.
\proofsection{$M_3$}
От противного. Допустим, что найдутся множества
$X\in \cal E$, $Y\in \cal E$, $|X|<|Y|$, такие что аксиома 
$M_3$ не выполнена. 
Среди всех таких пар множеств $X$ и $Y$, 
выберем пару, в которой достигается $\min|X\cup Y|$, 
то есть пару с максимальным числом совпадающих 
элементов. Обозначим $\{\Tuple{y}{1}{n}\} \Assign Y\setminus X$.
Если $n=1$, то $X+y_1 = Y$ и аксиома $M_3$ выполнена, то есть можно считать, что
$n>1$. Заметим, что в силу предположения $\A{i \in 1..n}{X+y_i\not\in \cal E
\And \E{C_i\in \cal C}{C_i \subset X+y_i} \And y_i\in C_i}$,
причём все множества $C_i$ попарно различны.

Рассмотрим множество $C_1$. Имеем $y_1\in C_1\subset X+y_1$.
Но тогда $\E{x_1\in X\setminus Y}{x_1\in C_1}$,
поскольку $Y\in \cal E$. Рассмотрим теперь
множество $Z \Assign X - x_1 + y_1$.

Пусть $Z\not\in \cal E$. Тогда $\E{C^\prime\in \cal C}{y_1\in C^\prime \subset Z}$, и $\E{C^{\prime\prime} \in \cal C}{C^{\prime\prime} \subset (C_1 \cup C^{\prime})-y_1}\subset X$, поскольку $y_1 \in C_1 \cap C^\prime$, что противоречит условию $X\in \cal E$.  
\end{proof}
\end{frame}

\begin{frame}
\frametitle{2.7.4. Циклы в матроидах (4/4)}
\begin{proof} (Окончание)
Пусть $Z\in \cal E$. Заметим, что  $|Z\cup Y|<|X\cup Y|$,
поскольку $Z\Assign X - x_1+ y_1$. Но пара $X,Y$ выбрана
как минимальная, значит для пары $Z,Y$ 
аксиома $M_3$ выполняется  и cуществует некоторый
элемент $y_j \in Y$ такой, что $Z+y_j \in \cal E$.  
Рассмотрим множество $C_j$. Имеем $y_j\in C_j \subset X+y_j$.
Если $x_1 \not\in C_j$, то \\$C_j\subset X -x_1+ y_j 
\subset X-x_1+y_1 + y_j = Z+y_j$, что противоречит $Z+y_j \in \cal E$.
Следовательно, $x_1 \in C_1\cap C_j \And C_1\ne C_j$.  
Но тогда $\E{C\in \cal C}{C\subset (C_j\cup C_1)-x_1}$.
Имеем $C\subset (C_j\cup C_1)-x_1\subset (C_j-x_1)\cup (C_1-x_1)$
$\subset (X-x_1+y_j)\cup (X-x_1+y_1) = Z+y_j\in \cal E$,
что невозможно.
\end{proof}

\medskip
\begin{corollary} {\NB{1}}
Если $X$~--- любое независимое множество матроида $M=\lrle{E,\cal E}$,
а $e\in E$~--- произвольный элемент носителя,
то множество $X+e$ содержит не более одного цикла.
\end{corollary}

\medskip
\begin{corollary} {\NB{2}}
Если $B$~--- любой базис матроида $M=\lrle{E,\cal E}$,
а $e\in E\setminus B$~--- произвольный элемент носителя вне базиса,
то множество $B+e$ содержит в точности один цикл, включающий $e$.
\end{corollary}


\end{frame}

\begin{frame}\label{2.7.5.}
\subsubsection{2.7.5. Жадный алгоритм }\frametitle{2.7.5. Жадный алгоритм (1/5)}


Сформулируем более точно рассматриваемую
экстремальную задачу. Пусть имеются конечное множество $E$,
$\vert E\vert=n$, \odmkey{весовая}{Функция (function)!весовая (weight)}
функция $\Map{w}{E}{\Bbb R_+}$ и~семейство $\cal E\subset 2^E$.
Требуется найти $X\in\cal E$, такое, что
$w(X)=\max\limits_{Y\in\cal E}w(Y)$,
где $w(Z)\Def\sum\limits_{e\in Z\subset E}w(e)$.
Другими словами, необходимо выбрать в указанном семействе
 подмножество наибольшего веса.
Не ограничивая общности, можно считать, что $w(e_1)\geq\ldots \geq w(e_n)>0$. Рассмотрим следующий алгоритм.



\vspace{-2pt}
\begin{algorithmic}
\REQUIRE множество $E=\{\Tuple{e}{1}{n}\}$,
семейство его подмножеств $\cal E$
и весовая функция $w$, $w(e_1)\geq\ldots \geq w(e_n)>0$.
%Множество $E$ линейно упорядочено в порядке
%убывания весов элементов.
\ENSURE подмножество $X$, такое, что
$X\in \cal E$ и $w(X)\leq \max\limits_{Y\in \cal E} w(Y)$.
\vspace{-2pt}
\STATE $X\Assign\emptyset$
 \COMMENT{\color{gray}\small  вначале множество $X$ пусто }
\FOR{\textbf{$i$ from $1$ to $n$}}
 \IF{$X+e_i\in\cal E$}
   \STATE $X\Assign X+e_i$
   \COMMENT{\color{gray}\small  добавляем в $X$ первый подходящий элемент }
 \ENDIF
\ENDFOR
\end{algorithmic}


\end{frame}


\begin{frame}
\frametitle{2.7.5. Жадный алгоритм (2/5)}

Алгоритм такого типа называется \odmkey{жадным}{Алгоритм
(algorithm)!жадный (greedy)}. Совершенно очевидно, что по
построению окончательное множество $X\in\cal E$. Также очевидно,
что жадный алгоритм является чрезвычайно эффективным: количество
шагов составляет $O(n)$, то~есть жадный алгоритм является
\odmkey{линейным}{Алгоритм (algorithm)!линейный (linear)}. 
%(Не
%считая затрат на сортировку множества $E$ и проверку независимости
%$X+e_i\in\cal E$.) 
Возникает вопрос: в каких случаях жадный
алгоритм действительно решает задачу, поставленную в начале
параграфа? Другими словами: когда выгодно быть жадным?
%
%\vspace*{-.5\baselineskip}

\begin{example}
Пусть дана матрица {\color{blue}$(0)$}.
В задаче <<выбрать по одному элементу \em{из каждого столбца\/}
так, чтобы их сумма была максимальна>>
жадный алгоритм
выберет элементы {\color{blue}$(1)$},
которые действительно являются решением задачи.
В то же время, в задаче
<<выбрать по
одному элементу \em{из каждого столбца и из каждой строки\/} так,
чтобы их сумма была максимальна>>
жадный алгоритм выберет элементы {\color{blue}$(2)$},
которые не являются решением,
потому что {\color{blue}$(3)$}~--- лучшее решение . 
$$
{\color{blue}(0)} \begin{matrix}
7 & 5 & 1 \\
3 & 4 & 3 \\
2 & 3 & 1
\end{matrix}
\qquad
{\color{blue}(1)} \begin{matrix}
\boxed{7} & \boxed{5} & 1 \\
3 & 4 & \boxed{3} \\
2 & 3 & 1
\end{matrix}
\qquad
{\color{blue}(2)} 
\begin{matrix}
\boxed{7} & 5 & 1 \\
3 & \boxed{4} & 3 \\
2 & 3 & \boxed{1}
\end{matrix}
\qquad
{\color{blue}(3)} 
\begin{matrix}
\boxed{7} & 5 & 1 \\
3 & 4 & \boxed{3} \\
2 & \boxed{3} & 1
\end{matrix}
$$
\end{example}

\end{frame}


\begin{frame}
\frametitle{2.7.5. Жадный алгоритм (3/5)}
\begin{theorem} {\NB{[Радо--Эдмондса]}}
Если $M=\lrle{E,\cal E}$~{\upshape---} матроид, то для любой
функции $w$ жадный алгоритм находит независимое множество $X$ с
наибольшим весом; обратно, если же $M=\lrle{E,\cal E}$ не является
матроидом, то существует такая функция $w$, что множество $X$,
найденное жадным алгоритмом, не будет подмножеством наибольшего
веса.
\end{theorem}

\begin{proof}
Рассмотрим матроид $M=\lrle{E,\cal E}$ и
 множество $X=\{\Tuple{x}{1}{k}\}$, построенное
жадным алгоритмом. По построению $w(x_1)\geq\ldots\geq w(x_k)>0$ и
множество $X$~--- базис матроида $M$. 
Пусть теперь 
$Y =\{\Tuple{y}{1}{m}\}\in\cal E$~--- некоторое
независимое множество, $w(y_1)\geq\ldots\geq w(y_m)>0$. Имеем $m\leq k$. 
Покажем от противного, что\\ 
$\A{i \in 1..m}{w(y_i) \leq w(x_i)}$. 
Пусть $\E{i \in 1..m}{w(y_i) > w(x_i)}$.
Рассмотрим независимые множества $A=\{\Tuple{x}{1}{i-1}\}$ и
$B=\{\Tuple{y}{1}{i-1},y_i\}$. \\Имеем $\E{j\leq
i}{\{\Tuple{x}{1}{i-1},y_j\}}$~--- независимое
множество. Тогда \\$w(y_j)\geq w(y_i)>w(x_i)$, откуда
следует, что \\
$\E{p \leq i}{w(x_1) \geq \ldots \geq w(x_{p-1}) \geq w(y_j) \geq
 w(x_p)}$, что противоречит тому, что $x_p$~--- элемент
с~наибольшим весом, добавление которого к множеству
$\{\Tuple{x}{1}{p-1}\}$ не нарушает независимости. Следовательно,
$\A{i}{w(y_i)\leq w(x_i)}$, и значит,
$w(Y)\leq w(X)$. 
\end{proof}\end{frame}


\begin{frame}
\frametitle{2.7.5. Жадный алгоритм (4/5)}

\begin{proof} (Продолжение) 
 Пусть теперь
$M=\Alg{E}{\cal E}$ не является матроидом. Если нарушено условие
$M_2$, то~есть $\E{A,B}{A\subset B\in\cal E}$, но $A\not\in \cal
E$, то определим функцию $w$ следующим образом:
\textbf{$w(e)\Assign$ if $e\in A$ then $1$ else $0$ end~if}.
Тогда \\
$A\not\subset X\Impl w(X)<w(A)=w(B)$.
Если же условие $M_2$ выполнено, но нарушено условие $M_3$,
то $\E{A,B}{\vert A\vert=k\And \vert B\vert=k+1}$
и для любого элемента $e\in B\setminus A$ множество $A+e$~--- зависимое.
Пусть $p\Assign\vert A\cap B\vert$. Тогда $p<k$. Выберем число
$\varepsilon$ так, что
$0<\varepsilon<1/(k-p)$.
Определим функцию $w$: \\
$w(e)\Assign$ \textbf{if $e\in A$ then
$1+\varepsilon$ else if $e\in B\setminus A$ then $1$ else $0$ end~if}.
Заметим, что при таких весах жадный алгоритм сначала выберет все
элементы из $A$ и отбросит элементы из $B\setminus A$. В
результате будет выбрано множество $X$, вес которого меньше веса
множества $B$. Действительно, $w(X)\!=\! w(A)=k(1\!+ \varepsilon) = 
(k - p)(1 + \varepsilon) +p(1 + \varepsilon) <$ \\$< 
(k - p)(1 + 1/(k-p)) + p(1 + \varepsilon)  = 
(k - p+ 1) +p(1 + \varepsilon) = w(B).$
\end{proof}

\end{frame}


\begin{frame}
\frametitle{2.7.5. Жадный алгоритм (5/5)}

%\vspace{-5pt}
\begin{remark}
Тот факт, что семейство $\cal E$
является матроидом, означает,
что для решения поставленной
экстремальной задачи можно
применить жадный алгоритм,
однако из этого не следует,
что не может существовать ещё
более эффективного алгоритма.
С другой стороны,
если семейство $\cal E$ не
является матроидом, то это ещё
не значит, что жадный алгоритм
не найдёт правильного решения~---
всё зависит от свойств конкретной
функции $w$.
\end{remark}

\medskip
%\vspace*{-.25\baselineskip}

\begin{digress}
Жадные алгоритмы были исследованы сравнительно
недавно, но их значение в практике программирования чрезвычайно
велико. Если удаётся свести конкретную экстремальную задачу к
такой постановке, где множество допустимых вариантов 
является матроидом, то в большинстве
случаев следует сразу применять жадный алгоритм, поскольку он
достаточно эффективен в практическом смысле. Если же, наоборот,
оказывается, что множество допустимых вариантов не
образует матроида, то это <<плохой признак>>. Скорее всего, данная
задача окажется труднорешаемой. В этом случае целесообразно
 исследовать задачу, чтобы избежать бесплодных попыток
<<изобрести>> эффективный алгоритм там, где это на самом деле
невозможно.
\end{digress}

\end{frame}

\begin{frame}
\subsubsection{2.7.6. Примеры применения матроидов }
\label{2.7.6.}\frametitle{2.7.6. Примеры применения матроидов (1/5)}

Ниже приведены некоторые примеры матроидов, встречающихся
в рассмотренных областях дискретной математики. В последних главах курса
имеются дополнительные примеры матроидов, связанных с графами
(\DMref{п.~9.6.3}{9.6.3.}, \DMref{П.~10.5.1}{10.5.1.}). 

\medskip
\begin{examples}
\end{examples}
{\color{blue}1. }\odmkey{Свободные матроиды}{Матроид (matroid)!свободный
(free)}. Если $E$~--- произвольное конечное множество,
то $M=\lrle{E,2^E}$~--- матроид. Такой матроид
называется \odmkey{свободным}{Матроид (matroid)!свободный (free)}.
В свободном матроиде каждое множество независимое.
Более того, если $D\subset E$, то   
$\lrle{E,2^D}$~--- матроид над множеством $E$ 
для любого подмножества $D$.

\medskip
{\color{blue}2. }
\odmkey{Матроиды разбиений}{Матроид (matroid)!разбиений (of
partitions)}. Пусть $\{\Tuple{E}{1}{k}\}$~--- некоторое разбиение
множества $E$ на непустые множества. Другими словами,
$\bigcup\limits^k_{i=1}E_i=E$, $E_i\cap E_j=\emptyset$,
$E_j\not=\emptyset$. Положим $\cal E\Assign\Set{A\subseteq E}{1\ge
\vert A\cap E_i\vert}$, то~есть в независимое множество входит не
более чем один элемент каждого блока разбиения. Тогда
$M=\lrle{E,\cal E}$~--- матроид. 
Действительно, условие $M_2$
выполнено.
Если $A,B\in\cal E$ и $\vert B\vert=\vert A\vert+1$, то
$\E{i}{\vert E_i\cap A\vert=0\And\vert E_i\cap B\vert=1}$.
Обозначим $e\Assign E_i\cap B$. Тогда $A+e\in\cal E$. Значит,
выполнено условие~$M_3$. 
\end{frame}

\begin{frame}
\frametitle{2.7.6. Примеры применения матроидов (2/5)}

{\color{blue}3. }\em{Векторные пространства}.
Линейно независимые множества векторов любого векторного
пространства образуют матроид. Действительно, условия $M_1$ и
$M_2$, очевидно, выполнены для линейно независимых множеств
векторов. Условие $M_4$ выполнено по теореме
\DMref{п.~2.4.3}{2.4.3.}.

{\color{blue}4. }\odmkey{Матроид
трансверсалей}{Матроид (matroid)!трансверсалей (of transversals)}.
Пусть $E$~--- некоторое множество и $\cal E$~--- семейство
подмножеств этого множества (не обязательно различных), $\cal
E\subset 2^E$. Подмножество $A\subset E$ называется
\odmkey{частичной трансверсалью}{Трансверсаль
(transversal)!частичная (partial)} семейства $\cal E$, если $A$
содержит не более чем по одному элементу для каждого подмножества
$E_i$ семейства $\cal E$. Частичные трансверсали над $E$ образуют
матроид.

 Необходимо ясно понимать границы применимости жадных алгоритмов. 
Если семейство множеств допустимых решений
экстремальной задачи образует матроид, 
то жадный алгоритм в любом случае найдет максимальное 
допустимое множество. Другими словами, свойств матроида
достаточно для применимости жадного алгоритма.
Если же свойства матроида не выполняются, то жадный 
алгоритм найдет некоторое <<решение>>,
но оно может оказаться как максимальным, так и не максимальным.
Другими словами, свойства матроида не являются необходимыми для применимости жадного алгоритма.
\end{frame}

\begin{frame}
\frametitle{2.7.6. Примеры применения матроидов (3/5)}
\begin{examples}Рассмотрим серию примеров на тему известной задачи о рюкзаке. Имеется $n$ предметов, в наших примерах
это наборы конфет, каждый из которых имеет вес $w_i$ 
и стоимость $c_i$, а также рюкзак с грузоподъёмностью $K$.
Требуется выбрать множество предметов $X$ так, чтобы
их суммарная стоимость была максимальна, 
а суммарный вес не превышал $K$:

\vspace{-20pt}
$$
\sum_{i=1}^n c_i x_i \to \max,\quad \sum_{i=1}^n w_i x_i \le K,
\quad x_i \in E_2=\{0,1\}. 
$$
Не ограничивая общности и для удобства мы рассматриваем целые числа.    
\end{examples}

{\color{blue}1. } Пусть все конфеты упакованы 
в наборы одного веса, который можно принять за единицу,
$\A{i \in 1..n}{w_i =1}$. Тогда семейство
множеств, которые помещаются в рюкзак, образует матроид.
Действительно, пустое множество помещается в рюкзак,
если некоторое множество помещается в рюкзак,
то любое его подмножество помещается в рюкзак,
и если есть два множества $A$ и $B$, 
такие что $|A|+1=|B|$, то по принципу Дирихле
в множестве $B\setminus A$ можно взять любой
элемент и добавить в множество $A$, и после этого
оно будет помещаться в рюкзак. 
Легко видеть, что жадный алгоритм, добавляющий предметы 
в порядке невозрастания стоимости, находит оптимальное решение.
\end{frame}

\begin{frame}
\frametitle{2.7.6. Примеры применения матроидов (4/5)}

{\color{blue}2. } Пусть теперь конфеты упакованы 
в наборы \em{разного} веса, тогда исходные данные 
представляет собой множество пар 
$\{(c_i, w_i)\}_{i=1}^n$. В этом случае семейство
множеств, которые помещаются в рюкзак, вообще говоря, \em{не} образует матроид.
Хотя пустое множество помещается в рюкзак, и
если некоторое множество помещается в рюкзак,
то любое его подмножество помещается в рюкзак,
но нетрудно предъявить конкретный пример
двух множеств $A$ и $B$, 
таких что $|A|+1=|B|$, но  $A$ невозможно расширить 
за счёт $B$ так, чтобы оно помещалось в рюкзак.
Пусть $\{(c_i, w_i)\}_{i=1}^n = \{(800,1), (300,1),
 (780,3), (450,2)\}$, и пусть 
$A = \{(800,1),  (780,3)\}$,  $B = \{(800,1), (300,1),
 (450,2)\}$ и $K=4$.
В этом случае множество $A$ вообще невозможно расширить.
При решении задачи о рюкзаке обычно рассматривают цену за килограмм $c_i/w_i$, а  жадный алгоритм добавляет 
предметы в порядке невозрастания цены.
В этом примере цены суть $(800, 300, 260, 225)$.
Легко видеть, что жадный алгоритм, добавляющий предметы 
в порядке невозрастания цены, находит неоптимальное решение
$\{(800,1), (300,1), (450,2)\}$, в то время
как оптимальным является множество $\{(800,1), (780,3)\}$.
Идея добавлять предметы в порядке невозрастания стоимости
срабатывает в приведённом случае, но 
опровергается таким примером: $\{(c_i, w_i)\}_{i=1}^n = \{(800,4), (780,3), (450,2), (300,1)\}$.
\end{frame}

\begin{frame}
\frametitle{2.7.6. Примеры применения матроидов (5/5)}
{\color{blue}3. } 
Однако отсутствие свойств матроида не означает,
что жадный алгоритм не найдёт решения. Пусть
$\{(c_i, w_i)\}_{i=1}^n = \{(800,2), (600,2),
 (780,3), (450,2)\}$, и пусть цены более равномерные:
  $(400, 300, 260, 225)$. 
Множество 
$A = \{(780,3)\}$ при $K=4$ невозможно расширить, один килограмм
вместимости рюкзака остаётся неиспользованным при любом выборе
множества $B$.
Таким образом, семейство множеств наборов, 
помещающихся в рюкзак, не является матроидом в данном случае.
Однако жадный алгоритм находит оптимальное решение
$\{(800,2), (600,2) \}$,
причём как при добавлении предметов в порядке невозрастания цены, 
так и при добавлении предметов в порядке невозрастания стоимости.

\medskip
Таким образом, свойства матроида являются достаточными
для применимости жадного алгоритма, но не являются необходимыми.
Необходимые и достаточные условия применимости 
жадного алгоритма в общем случае неизвестны,
и если свойства матроида не выполняются, то
приходится индивидуально исследовать применимость
жадного алгоритма в каждом конкретном случае. 
\end{frame}

%\begin{frame}
%\frametitle{Онтология общей алгебры}
%\begin{center}
%\includegraphics[height=8cm]{ODM2_6.png}
%\end{center}
%\end{frame}


% ===== tex/DM2022_3.tex =====

\begin{frame}
\frametitle{3. Булевы функции (1/2)}
\section{3. Булевы функции}
{\large
Занимаясь исследованием законов мышления, 
английский математик \\Джордж Буль  применил в логике 
систему формальных обозначений и правил, близкую к математической. 
Так возникла \em{алгебра логики}, или \em{булева алгебра}. 
Значение логической алгебры долгое время игнорировалось, 
поскольку она не находила применения в науке и технике того времени. 
С появлением электронной техники операции, 
введённые Булем, оказались весьма полезны. 
Они изначально ориентированы на работу только с двумя сущностями~--- 
истина и ложь, что хорошо согласуется с двоичным кодом, 
который в современных компьютерах также 
представляется всего двумя сигналами~--- ноль и единица.
}
\end{frame}

\begin{frame}
\frametitle{3. Булевы функции (2/2)}
{\large
Данная глава имеет двойное назначение. С одной стороны,
помимо основных фактов из теории булевых функций здесь
затрагиваются весьма общие понятия, такие как \em{реализация функций
формулами}, \em{нормальные формы}, \em{двойственность}, \em{полнота}.
Эти понятия затруднительно описать исчерпывающим образом на
выбранном элементарном уровне изложения, но знакомство с ними
необходимо. Поэтому рассматриваются частные случаи
указанных понятий на простейшем примере булевых функций.
С другой стороны, материал этой главы служит для
<<накопления фактов>> и овладения базисной логической техникой,
что должно создать у читателя
необходимый эффект <<узнавания знакомого>> при
изучении довольно абстрактного и формального
материала следующей главы.
}
\end{frame}
\subsection{3.1. Элементарные булевы функции}
\label{3.1.}
\begin{frame}
\frametitle{3.1. Элементарные булевы функции }

Подобно тому как в классической математике знакомство с основами анализа
начинают с изучения \em{элементарных функций}
(рациональные функции вещественной переменной 
$x$, $e^x$, $\ln x$, прямые и обратные тригонометрические функции и
их суперпозиции),
изложение теории булевых функций естественно начать
с выделения, идентификации и изучения \em{элементарных
булевых функций} (дизъюнкция, конъюнкция и т.~д.).

\medskip
\begin{tabular}{p{12mm}|p{50mm}|p{80mm}}
  № & Название параграфа
  & Ключевые термины и обозначения \\
  \hline
\DMref{3.1.1.}{3.1.1.} & Функции алгебры логики  &
{ Булевы функции, истинностные значения, таблица истинности, установленный порядок}\\  
\DMref{3.1.2.}{3.1.2.} & Существенные и несущественные переменные &
{Введение переменной, удаление переменной }\\  
\DMref{3.1.3.}{3.1.3.} & Булевы функции одной переменной  &
{ Константы, $0$~--- нуль, тождественная функция, 
$\Not x$~--- отрицание, $1$~--- единица }\\  
\DMref{3.1.4.}{3.1.4.}\RS & Булевы функции двух переменных  &
{ $\And$~--- конъюнкция, $+_{2}$~--- сложение по модулю, $\mid$~--- штрих Шеффера}
\\  
\DMref{3.1.5.}{3.1.5.}\BS & Функции $k$-значной логики  &
{ }\\    
\end{tabular}  


\end{frame}

\begin{frame}
\label{3.1.1.}
\label{3.1.1}\subsubsection{3.1.1. Функции алгебры логики }\frametitle{3.1.1. Функции алгебры логики (1/3)}
%\subsection{Функции алгебры логики} \label{f-a-l}
Функции $\Map{f}{E^{n}_{2}}{E_2}$, где $E_2\Def \{0,1\}$,
называются \odmkey{функциями алгебры логики}{Функция
(function)!алгебры логики (of boolean algebra)}, или
\odmkey{булевыми функциями}{Функция (function)!булева (boolean)}
от $n$ переменных. 
Элементы множества $E_2$ называются 
\odmkey{истинностными значениями}{Истинностное значение (truth-value)}. 
Множество булевых функций от $n$ переменных
обозначим $P_n$, \\$P_n\Def \{f\mid f\colon E_2^{n}\to E_2\}$.


Булеву функцию от $n$ переменных можно задать \odmkey{таблицей
истинности}{Таблица (table)!истинности (of truth values)}$\colon$

%\vspace{-8pt}
\begin{center}
%{\sffamily\small
%\def\arraystretch{1.1}
\begin{tabular}{cccc|c}
%\hline
 $x_1$ &  $\ldots$& $x_{n-1}$& $x_n$ & $f(\Tuple{x}{1}{n})$ \\
\hline
  $0$    &  $\ldots$ &    $0$     &  $0$    &    $f(0,\ldots,0,0)$\\
  $0$    &  $\ldots$ &    $0$     &  $1$    &    $f(0,\ldots,0,1)$\\
  $0$    &  $\ldots$ &    $1$     &  $0$    &    $f(0,\ldots,1,0)$\\
$\ldots$ &  $\ldots$ &   $\ldots$ & $\ldots $  & $\ldots $\\
  $1$    &  $\ldots$ &     $1$     &   $1$    &    $f(1,\ldots,1,1)$\\
\hline
\end{tabular}
%}
\end{center}

Если число переменных равно $n$, то в таблице истинности имеется
$2^n$ строк, соответствующих всем различным комбинациям значений
переменных (см. \DMref{п. 5.1.2}{5.1.2.}). Следовательно, существует $2^{2^n}$ различных
столбцов, каждый из которых определяет булеву функцию от $n$
переменных. 
%Таким образом, $|P_n|=2^{2^n}$.


\end{frame}

\begin{frame}
\frametitle{3.1.1. Функции алгебры логики (2/3)}



Если нужно задать несколько (например, $k$) булевых функций от $n$
переменных, то это удобно сделать с помощью одной таблицы,
использовав несколько столбцов:

%\addvspace{6pt}
\begin{center}
\begin{tabular}{ccc|c|c|c}
\hline
 $x_1$ & $\ldots$ & $x_n$ & $f_1(\Tuple{x}{1}{n})$ & $\ldots$ & $f_k(\Tuple{x}{1}{n})$ \\
\hline
   0   & $\ldots$ &  0    & $f_1(0,\ldots,0)$       & $\ldots$ & $f_k(0,\ldots,0)$       \\
$\ldots$ &  $\ldots$ &   $\ldots$ & $\ldots $  & $\ldots$ & $\ldots $ \\
   1   & $\ldots$ &  1    & $f_1(1,\ldots,1)$       & $\ldots$ & $f_k(1,\ldots,1)$       \\
\hline
\end{tabular}
\end{center}

%\addvspace{6pt}
Такая таблица имеет $2^n$ строк, $n$ столбцов для значений
переменных и ещё $k$ столбцов для значений функций. 
Вообще говоря,
значения переменных можно не хранить, если принять соглашение о
перечислении наборов переменных в определённом порядке. Далее во всех
таблицах истинности  переменные всегда
перечисляются в лексикографическом порядке, а кортежи булевых
значений~--- в порядке возрастания целых чисел,
задаваемых
кортежами как двоичными шкалами,
что совпадает с~лексикографическим порядком на кортежах.
Такой
порядок здесь называется \odmkey{установленным}{Порядок
(order)!установленный (established)}.
\end{frame}

\begin{frame}
\frametitle{3.1.1. Функции алгебры логики (3/3)}
Из предшествующих наблюдений вытекает теорема.
\begin{theorem}
$\vert P_n \vert = 2^{2^n}.$
\end{theorem}

Если $k>2^n$ (что, как показывает предыдущая формула, не редкость), то
таблицу истинности удобно <<транспонировать>>, выписывая наборы
значений в столбцах, а значения функций~--- в строках:

%\addvspace{6pt}
\begin{center}

\begin{tabular}{c|ccc}
\hline
 $x_1$ & 0 & $\dots$ & 1  \\
$\dots$ & $\dots$ & $\dots$ & $\dots$  \\
$x_n$ & 0 & $\dots$ & 1  \\
\hline
   $f_1$   & $f_1(0,\dots,0)$ &  $\dots$    & $f_1(1,\dots,1)$  \\
   $\dots$ &  $\dots$         &  $\dots$    & $\dots $           \\
   $f_k$   & $f_k(0,\dots,0)$ &  $\dots$    & $f_k(1,\dots,1)$  \\
\hline
\end{tabular}
\end{center}

\addvspace{6pt}
Именно такой способ записи таблицы истинности использован в
пп.~\DMref{3.1.3}{3.1.3.} и~\DMref{3.1.4}{3.1.4.}.


\end{frame}

\begin{frame}
\label{3.1.2.}\subsubsection{3.1.2. Существенные и несущественные переменные }\frametitle{3.1.2. Существенные и несущественные переменные (1/3)}
%\subsection{Существенные и несущественные переменные}
%\label{s31}

Булева функция $f \in P_n$
\em{существенно зависит\/} от переменной $x_i$, если существует такой
набор значений $\Tuple{a}{1}{i-1},\Tuple{a}{i+1}{n}$, что
$$
f(\Tuple{a}{1}{i-1},0,\Tuple{a}{i+1}{n}) \not=
f(\Tuple{a}{1}{i-1},1,\Tuple{a}{i+1}{n}).
$$
В этом случае $x_i$ называют
\odmkey{существенной переменной}{Переменная (variable)!существенная
(essential)},
в противном случае $x_i$
называют \odmkey{{несущественной
(\odmkey{фиктивной}{Переменная (variable)!фиктивная (inessential)})
 переменной}}
{Переменная (variable)!несущественная (inessential)}.

\begin{example}
Пусть булевы функции $f_1(x_1,x_2)$ и
$f_2(x_1,x_2)$ заданы %следующей
таблицей истинности:

\addvspace{6pt}
\begin{center}
\def\arraystretch{1.2}
\begin{tabular}{p{10mm}p{10mm}|p{10mm}|p{10mm}}
$x_1$ & $x_2$ & $f_1$ & $f_2$ \\
\hline
   0  &   0   &   0   &   1   \\
   0  &   1   &   0   &   1   \\
   1  &   0   &   1   &   0   \\
   1  &   1   &   1   &   0   \\
\hline
\end{tabular}
\end{center}

Для этих функций переменная
$x_1$~--- существенная, а переменная $x_2$~--- несущественная.
\end{example}
\end{frame}

\begin{frame}
\frametitle{ 3.1.2. Существенные и несущественные переменные (2/3)}
\vspace{-2pt}
Пусть заданы две булевы функции,
$f_1(\Tuple{x}{1}{n-1})$ и $f_2(\Tuple{x}{1}{n-1},x_n)$,
и пусть переменная $x_n$~--- несущественная для функции $f_2$,
а при одинаковых значениях остальных переменных значения
функций совпадают:

\vspace{-16pt}
$$
\A{\Tuple{a}{1}{n-1},a_n}{f_1(\Tuple{a}{1}{n-1})=f_2(\Tuple{a}{1}{n-1},a_n)}.
$$

\vspace{-4pt}
В таком случае говорят, что $f_2$ получается из $f_1$
\em{введением}
несущественной переменной $x_n$,
а $f_1$ получается из $f_2$
\em{удалением}
несущественной переменной $x_n$.

\begin{remark}
Процедурно введение и удаление несущественных переменных
выполняются достаточно просто. 
Чтобы ввести несущественную
переменную, нужно продублировать каждую строку таблицы истинности,
добавить новый столбец и заполнить этот столбец чередующимися
значениями $0$ и $1$ (или $1$ и $0$, что несущественно).
Удаление несущественной переменной выполняется аналогично:
нужно отсортировать таблицу истинности так, чтобы она состояла из пар строк,
различающихся только в разряде несущественной переменной,
после чего удалить столбец несущественной переменной
и удалить каждую вторую строку таблицы.
\end{remark}
\end{frame}

\begin{frame}
\frametitle{ 3.1.2. Существенные и несущественные переменные (3/3)}
Всюду в дальнейшем булевы функции рассматриваются 
\em{с точностью до несущественных переменных}. Это позволяет считать, что все булевы функции (в~данной системе функций) зависят от одних и тех же переменных. Для этого в~данной системе булевых функций можно сначала удалить все несущественные переменные, а~потом добавить несущественные переменные так, чтобы уравнять количество переменных у всех функций.
\end{frame}

\begin{frame}
\label{3.1.3.}\subsubsection{3.1.3. Булевы функции одной переменной}
\frametitle{3.1.3. Булевы функции одной переменной}

%\subsection{Булевы функции одной переменной}\label{sec-3-1-3}

В следующей таблице собраны все булевы функции одной переменной.
Две из них, фактически, являются \em{константами}, поскольку их
значения не зависят от значения аргумента.

\addvspace{6pt}
%\begin{center}
{\sffamily\small
\def\arraystretch{1.2}
\begin{tabular}{l|r|c|c|l}
%\hline
\multicolumn{2}{r|}{Переменная $x$}         & $0$ & $1$ &            \\
\hline
Название      & Обозначение               &   &   & Несущественные   \\
\hline
Нуль          &           $0$               & $0$ & $0$ & $x$   \\
Тождественная &          $x$              & $0$ & $1$ &      \\
Отрицание     & $\neg x$, $\bar x$, $x'$, $\sim x$  & $1$ & $0$ &      \\
Единица       &           $1$               & $1$ & $1$ & $x$   \\
\hline
\end{tabular}
}
%\end{center}

\end{frame}

\begin{frame}
\label{3.1.4.}\subsubsection{3.1.4. Булевы функции двух переменных }\frametitle{3.1.4. Булевы функции двух переменных (1/2)}
В следующих двух таблицах собраны все булевы функции
двух переменных. Из них две являются константами,
четыре зависят от одной переменной,
и только десять существенно зависят от обеих переменных.

\begin{tabular}{l|r|c|c|c|c|l}
\multicolumn{2}{r|}{Переменная $x$}          & $0$ & $0$ & $1$ & $1$ & \\
%\hline
 \multicolumn{2}{r|}{Переменная $y$}         & $0$ & $1$ & $0$ & $1$ &   \\
\hline
Название             & Обозначение           &   &   &   &   & Несущественные   \\
\hline
Нуль                 &           $0$           & $0$ & $0$ & $0$ & $0$ & $x$, $y$   \\
Конъюнкция           &$\cdot$, $\And$, $\wedge$& $0$ & $0$ & $0$ & $1$ &       \\
                     &                       & $0$ & $0$ & $1$ & $0$ &         \\
                     &                       & $0$ & $0$ & $1$ & $1$ & $y$      \\
                     &                       & $0$ & $1$ & $0$ & $0$ &        \\
                     &                       & $0$ & $1$ & $0$ & $1$ & $x$     \\
Сложение по модулю 2 &$+$, ${+}_2$, $\not\equiv$, $\oplus$,
                       $\vartriangle$      & $0$ & $1$ & $1$ & $0$ &       \\
Дизъюнкция           &   $\Or$               & $0$ & $1$ & $1$ & $1$ &        \\
\hline
\end{tabular}

\end{frame}

\begin{frame}
\frametitle{3.1.4. Булевы функции двух переменных (2/2)}

\begin{tabular}{l|r|c|c|c|c|l}
\multicolumn{2}{r|}{Переменная $x$}          & $0$ & $0$ & $1$ & $1$ & \\
%\hline
 \multicolumn{2}{r|}{Переменная $y$}         & $0$ & $1$ & $0$ & $1$ &   \\
\hline
Название             & Обозначение           &   &   &   &   & Несущественные   \\
\hline
Стрелка Пирса
                     &   $\downarrow$        & $1$ & $0$ & $0$ & $0$ &       \\
Эквивалентность      &    $\equiv$    & $1$ & $0$ & $0$ & $1$ &       \\
                     &                       & $1$ & $0$ & $1$ & $0$ & $x$      \\
                     &                       & $1$ & $0$ & $1$ & $1$ &        \\
                     &                       & $1$ & $1$ & $0$ & $0$ & $y$       \\
%логическое следование&                       &   &   &   &   &        \\
Импликация           &$\rightarrow$,
                      $\Rightarrow$, $\supset$& $1$ & $1$ & $0$ & $1$ &        \\
Штрих Шеффера
                     &     $\mid$            & $1$ & $1$ & $1$ & $0$ &        \\
Единица              &         $1$             & $1$ & $1$ & $1$ & $1$ &  $x$, $y$   \\
\hline
\end{tabular}


\medskip
\begin{remark}
Пустоты в столбцах <<Название>> и <<Обозначение>> в предыдущих таблицах
означают, что булева функция редко используется,
а потому не имеет специального названия и обозначения.
\end{remark}

%%%%%%%%%%%%%%%%%%%%%%%%% Изменения 2012
\end{frame}

\begin{frame}
\label{3.1.5.}\subsubsection{3.1.5. Функции $k$-значной логики }\frametitle{3.1.5. Функции $k$-значной логики (1/3)}
%\subsection{Функции $k$-значной логики}

Рассмотрим множество $E_k = \{0, 1, \dots , k-1\}$ и множество функций
$P_{k,n}$, имеющих $n$ аргументов типа $E_k$ и возвращающих значение типа
$E_k$. Такие функции называются
\odmkey{функциями $k$-значной логики}{Функция (function)!$k$-значной логики ($k$-valued logic)}
и являются обобщением функций алгебры логики: $P_{2,n} = P_n$.
Функции $k$-значной логики можно задавать таблицами, подобными таблицам
истинности для обычной двузначной логики, причём в таблицах используются
элементы из множества $E_k$, а не из множества $E_2$.
Нетрудно видеть, что $\vert P_{k, n} \vert = k^{k^n}$.
Число функций $k$-значной логики растет ещё быстрее, чем число булевых
функций. Так, уже для трёхзначной логики аналог таблицы \DMref{п.~3.1.4}{3.1.4.}
для всех функций двух переменных содержит $3^9=19\,683$ строк и не может
быть помещён в книгу или на слайды.


\end{frame}

\begin{frame}
\frametitle{3.1.5. Функции $k$-значной логики (2/3)}

Функцию $k$-значной логики можно представить как систему функций обычной
двузначной логики следующим образом. Функция
$f(x_1, \dots, x_n) \in P_{k,n}$ представляется системой функций
$\{g_1(y_{11},  \dots,  y_{1m}, \dots,  y_{n1},  \dots,  y_{nm}),
 \dots,  g_m(y_{11},  \dots,  y_{1m},  \dots,  y_{n1},  \dots,  y_{nm})\}$,
где $m = \lceil \log k \rceil$. Здесь значения функции $f$ и
переменных $x_i$ принадлежат множеству $E_k$, а значения функций
$g_j$ и переменных $y_{ij}$ принадлежат множеству $E_2$. 

\smallskip
При этом,
если значение переменной $x_i$ имеет двоичный код $a_{i1} \dots
a_{im}$, то значение функции $f(x_1, \dots, x_n)$ имеет двоичный
код \\ $g_1(a_{11}, \dots, a_{1m}, \dots, a_{n1}, \dots, a_{nm})\dots
g_m(a_{11}, \dots, a_{1m}, \dots, a_{n1}, \dots, a_{nm})$. 
Другими
словами, функция $g_j$ вычисляет $j$-ю цифру в двоичном
представлении числа $f(x_1, \dots, x_n)$.

\smallskip
 Таким образом, хотя
$k$-значная логика может быть сведена к~двузначной логике, и
многие полезные утверждения двузначной логики распространяются на
случай $k$-значной логики, но вычислительные процедуры во всех
случаях являются существенно более трудоёмкими для функций
$k$-значной логики.

\end{frame}

\begin{frame}
\frametitle{3.1.5. Функции $k$-значной логики (3/3)}

\begin{columns}[t]
\begin{column}[T]{7cm}
\begin{example}
В трёхзначной логике 
$E_3=\{0,1,2\}$, 
 $k=3$ и $m=\lceil \log 3\rceil=2$. 
Рассмотрим функцию <<трёхзначная конъюнкция>>, 
которую можно определить, 
например, как минимум истинностных значений аргументов: 
$x_1 \And x_2 \Assign \min(x_1,x_2)$.
 
В таком случае таблица истинности имеет вид,
показанный справа.
\end{example}
\end{column}
\begin{column}[T]{7cm}
\begin{tabular}{|c|c|c|c|c|c|}
\hline
$x_1$&$y_{11} y_{12}$ & $x_2$ & $y_{21} y_{22}$ & $x_1 \And x_2$ & $g_1 g_2$ \\
\hline
0 & 00 & 0 & 00 & 0 & 00 \\
0 & 00 & 1 & 01 & 0 & 00 \\
0 & 00 & 2 & 10 & 0 & 00 \\
1 & 01 & 0 & 00 & 0 & 00 \\
1 & 01 & 1 & 01 & 1 & 01 \\
1 & 01 & 2 & 10 & 1 & 01 \\
2 & 10 & 0 & 00 & 0 & 00 \\
2 & 10 & 1 & 01 & 1 & 01 \\
2 & 10 & 2 & 10 & 2 & 10 \\
\hline
\end{tabular}
\end{column}
\end{columns}

\smallskip
Таким образом, 
конъюнкцию представляет следующая система функций:

$g_1(y_{11},y_{12},y_{21},y_{22})=y_{11}\And y_{21}$,  

$g_2(y_{11},y_{12},y_{21},y_{22})=y_{12}\And y_{22}\Or y_{12} 
\And y_{21}\Or y_{11}\And y_{22}$.


\smallskip
Функции $k$-значной логики далее в этом курсе не рассматриваются.
\end{frame}



\begin{frame}
\frametitle{3.2. Формулы}
\subsection{3.2. Формулы}\label{3.2.}

В этом разделе обсуждается целый ряд важных понятий,
которые часто считаются самоочевидными, а потому их объяснение опускается.
Такими понятиями являются, в частности, само понятие \em{формулы}, 
\em{реализация} функций формулами,
\em{интерпретация} формул для вычисления значений функций, 
\em{равносильность} формул,
\em{подстановка} и~\em{замена} в формулах.
Между тем программная реализация работы с формулами
требует учёта некоторых тонкостей, связанных с данными понятиями,
которые целесообразно явно указать и обсудить.

\medskip
\begin{tabular}{p{12mm}|p{65mm}|p{65mm}}
  № & Название параграфа
  & Ключевые термины и обозначения \\
  \hline
\DMref{3.2.1.}{3.2.1.}\RS & Реализация функций формулами  &
{ Базис, главная операция, подформула, алгоритм интерпретации }\\  
\DMref{3.2.2.}{3.2.2.} & Равносильные функции &
{Кратная операция, вариаргументная операция }\\  
\DMref{3.2.3.}{3.2.3.} & Подстановка и замена  &
{ Правило подстановки, правило замены }\\  
\DMref{3.2.4.}{3.2.4.} & Алгебра булевых функций  & {Алгебра Линденбаума--Тарского}\\
\end{tabular}  

\end{frame}

\begin{frame}
\label{3.2.1.}\subsubsection{3.2.1. Реализация функций формулами }\frametitle{3.2.1. Реализация функций формулами (1/8)}

Начнём обсуждение с привычного понятия формулы.
Вообще говоря, \odmkey{формула}{Формула (formula)}~--- 
это цепочка символов, имеющих заранее
оговорённый смысл.
Обычно формулы строятся по определённым правилам из
обозначений объектов и вспомогательных символов~--- разделителей.
Применительно к рассматриваемому случаю объектами являются
переменные и обозначения булевых функций, 
перечисленные в таблицах
пп.\DMref{~3.1.3}{3.1.3.} и~\DMref{3.1.4}{3.1.4.},
а разделителями~--- скобки <<(>>, <<)>> и запятая <<,>>.
Мы считаем, что обозначения всех объектов  определены и
отличимы друг от друга и от разделителей,
а правила составления формул известны.

\begin{remark}
Для обозначения объектов используются как отдельные символы, так и
группы символов, в том числе со специальным начертанием: верхние и
нижние индексы, наклон шрифта и т.~п. Допущение о синтаксической
различимости остаётся в силе для всех таких особенностей.
\end{remark}

\begin{example}
Операцию сложения чисел принято обозначать знаком
<<$+$>>, 
%$\Map{+}{\Bbb R\times\Bbb R}{\Bbb R}$,
а функцию <<синус>>~--- словом <<$\sin$>>,
которое записывается прямым шрифтом и считается одним символом.
%,$\Map{\sin}{\Bbb R}{\Bbb R}$
\end{example}
\end{frame}

\begin{frame}
\frametitle{3.2.1. Реализация функций формулами (2/8)}
\begin{digress}
В математических текстах часто используется \em{нелинейная}
форма записи формул, при которой формула \em{не} является цепочкой символов.
Примерами могут служить дроби, радикалы, суммы, интегралы.
Подобные нелинейные формы записи формул привычны и~удобны.
Исполь\-зо\-ва\-ние нелинейной записи формул не является принципиальным
расширением языка формул.
Можно ограничиться только линейными цепочками символов,
что блестяще подтверждает система TeX, с помощью которой
подготовлены эти слайды. 
\end{digress}


Пусть $F = \{\Tuple{f}{1}{m}\}$~--- некоторое множество булевых функций
от $n$ переменных.
\odmkey{Формулой}{Формула (formula)} $\cal F$ над $F$
называется выражение (цепочка символов) вида
$$
{\cal F}[F] = f(\Tuple{t}{1}{n}),
$$
где $f \in F$ и $t_i$~--- либо переменная, либо формула над $F$.
Множество $F$ называется \odmkey{базисом}{Базис (base)}, функция
$f$ называется \odmkey{главной}{Операция (operation)!главная
(principal)} (\odmkey{внешней}{Операция (operation)!внешняя
(outermost)}) операцией (функцией), а $t_i$ называются
\odmkey{подформулами}{Подформула (subformula)}.

\end{frame}

\begin{frame}
\frametitle{3.2.1. Реализация функций формулами (3/8)}

Если базис $F$
ясен из контекста, то его обозначение опускают.

Множество всех формул над базисом 
$F$ здесь обозначается $\mathfrak{F}[F]$.  


\smallskip
\begin{remark}
Обычно при записи формул, составленных из элементарных булевых функций,
обозначения бинарных булевых функций записываются как знаки операций в
инфиксной форме, отрицание записывается как знак унарной операции
в префиксной форме, тождественные функции записываются как константы~0 и~1,
устанавливается приоритет операций
($\neg$, $\And$, $\Or$, $\to$), и лишние
скобки опускаются.
\end{remark}

\smallskip
Всякой формуле ${\cal F}$ однозначно соответствует некоторая функция
$f$. Это соответствие задаётся приведённым на следующем слайде
\odmkey{алгоритмом интерпретации}{Алгоритм (algorithm)!интерпретации (interpretation)},
который позволяет вычислить
значение формулы ${\cal F}$ при заданных значениях переменных.

\smallskip
Алгоритм интерпретации реализован 
рекурсивной функцией Eval. 
\end{frame}

\begin{frame}
\frametitle{3.2.1. Реализация функций формулами (4/8)}
\begin{algorithmic}
\REQUIRE формула ${\cal F}$, множество $F$ функций базиса,
 значения переменных $\Tuple{x}{1}{n}$. 
\ENSURE значение формулы
${\cal F}$ на значениях $\Tuple{x}{1}{n}$ или значение
\textbf{fail}, если значение формулы не~может быть определено.
\STATE \textbf{if ${\cal F}$ = $'{x_i}'$ 
 return $x_i$ end if } 
  \COMMENT{\color{gray}\small значение переменной задано } 
 
\IF{${\cal F}$ = $'f(\Tuple{t}{1}{n})'$} 
  \STATE \textbf{if $f\not\in F$ return fail end if} 
  \COMMENT{\color{gray}\small функция не входит в базис } 
 \FOR{$i \in 1..n$} 
 \STATE $y_i\Assign$  Eval
 ($t_i,F,\Tuple{x}{1}{n}$) \COMMENT{\color{gray}\small значение $i$-го аргумента }
 \IF{$y_i$ = \textbf{fail}} 
   \STATE \textbf{ return fail }
   \COMMENT{\color{gray}\small аргумент не может быть вычислен } 
   \ENDIF
 \ENDFOR 
 \STATE \textbf{return} $f$($\Tuple{y}{1}{n}$) 
 \COMMENT{\color{gray}\small вычисленное значение главной операции } 
\ENDIF 
\STATE \textbf{return fail} \COMMENT{\color{gray}\small это не формула }
\end{algorithmic}
\end{frame}

\begin{frame}
\frametitle{3.2.1. Реализация функций формулами (5/8)}
%\begin{digress}

\vspace{-2pt}
Некоторые программистские замечания по поводу процедуры Eval.

\begin{remark}
\end{remark}
{\color{blue} 1.}
При программной реализации алгоритма
интерпретации формул важно учитывать, что в общем случае результат
вычисления значения формулы может быть не определён
(\textbf{fail}). Это имеет место, если формула построена
синтаксически неправильно или если в ней используются функции
(операции), способ вычисления которых не задан, то~есть они не
входят в базис. 
Таким образом, необходимо либо проверять
правильность формулы до начала работы алгоритма интерпретации,
либо предусматривать невозможность вычисления значения в самом
алгоритме.

%\item 
{\color{blue} 2.}
Это не единственный возможный алгоритм вычисления
значения формулы, 
более того, он не самый лучший. 
Для конкретных
классов формул, например, для некоторых классов формул,
реализующих булевы функции, известны более эффективные алгоритмы
(см., например, \DMref{п.~3.5.3}{3.5.3.})


%\item 

\end{frame}

\begin{frame}
\frametitle{3.2.1. Реализация функций формулами (6/8)}

{\color{blue} 3.}
Сама идея этого алгоритма:
<<сначала вычисляются значения аргументов, а потом значение
функции>>,~--- не является догмой.

 Например, можно построить такой
алгоритм интерпретации, который вычисляет значения только 
\em{некоторых\/} аргументов, а потом вычисляет значение формулы, в
результате чего получается новая \em{формула\/}, реализующая
функцию, которая зависит от меньшего числа аргументов. Такой
алгоритм интерпретации называется \odmkey{смешанными
вычислениями}{Смешанные вычисления \\(mixed computations)}.


%\item
{\color{blue} 4.}
Порядок вычисления аргументов \em{считается\/} неопределённым, то
есть предполагается, что базисные функции не имеют
\odmkey{побочных эффектов}{Побочный эффект (side-effect)}. Если же
базисные функции имеют побочные эффекты, то результат вычисления
значения функции может зависеть от порядка вычисления значений
аргументов. Побочные эффекты могут иметь различные причины.
Наиболее частая: использование глобальных переменных. Например,
если
$
f(x)\Def a\Assign a+1; \textbf{return}\ x+a,\quad 
g(x)\Def a\Assign a\ast 2; \textbf{return}\ x*a,
$
причём переменная $a$ глобальна, то (если $a\ne 6$)
$f(2)+g(3)\neq g(3)+f(2).$
%\end{enumerate}
%\end{digress}
\end{frame}

\begin{frame}
\frametitle{3.2.1. Реализация функций формулами (7/8)}
%\addvspace{-3pt} 
Если формула $\cal F$ и базис $F$ заданы,
то процедура Eval($\cal F,F,\Tuple{x}{1}{n}$) является некоторой
булевой функцией $f$ переменных $\Tuple{x}{1}{n}$. В этом случае
говорят, что формула ${\cal F}$ \odmkey{реализует}{Реализация функции (realization of function)} функцию $f$,
и записывают это так: 
$\func{\cal F} = f$. 


\begin{remark}
Для обозначения реализуемости применяют и другие приёмы. Выбранное
обозначение обладает тем достоинством, что согласовано
с другими обозначениями
в слайдах.
\end{remark}

Зная таблицы истинности для функций базиса, можно вычислить таблицу
истинности той функции, которую реализует данная формула.

%\begin{examples}
%\end{examples}
%\begin{enumerate}
%\item 

{\bf Примеры. }{\color{blue}\bf 1.}
$F_1 \Assign (x_1\wedge x_2)\vee ((x_1\wedge \overline x_2)\vee
(\overline x_1\wedge x_2))$

{\small
%\addvspace{6pt}
\def\arraystretch{0.9}
\begin{tabular}{p{5mm}|p{5mm}||p{15mm}|p{15mm}|c|p{15mm}|p{5mm}}
%\hline
$x_1$ & $x_2$ & $x_1\wedge \overline x_2$ & $\overline x_1\wedge x_2$ &
$(x_1\wedge \overline x_2)\vee (\overline x_1\wedge x_2)$ &
$x_1\wedge x_2$ & $F_1$ \\
\hline
0 & 0 & 0 & 0 & 0 & 0 & 0 \\
0 & 1 & 0 & 1 & 1 & 0 & 1 \\
1 & 0 & 1 & 0 & 1 & 0 & 1 \\
1 & 1 & 0 & 0 & 0 & 1 & 1 \\
\hline
\end{tabular}
}
%\addvspace{6pt}

Таким образом, формула $F_1$ реализует дизъюнкцию.

%\addvspace{12pt}

%\item 
\end{frame}

\begin{frame}
\frametitle{3.2.1. Реализация функций формулами (8/8)}

{\color{blue}\bf 2.}
$F_2 \Assign (x_1\wedge x_2)\to {x_1}$

%\addvspace{6pt}
%\begin{center}
\def\arraystretch{0.9}
\begin{tabular}{p{5mm}|p{5mm}||p{15mm}|p{10mm}}
%\hline
$x_1$ & $x_2$ & $x_1\wedge x_2$ & $F_2$ \\
\hline
0 & 0 & 0 & 1 \\
0 & 1 & 0 & 1 \\
1 & 0 & 0 & 1 \\
1 & 1 & 1 & 1 \\
\hline
\end{tabular}

%\addvspace{6pt}

Таким образом, формула $F_2$ реализует константу 1.

\addvspace{12pt}

%\item 
{\color{blue}\bf 3.}
$F_3 \Assign ((x_1\wedge x_2)+x_1)+x_2$

%\addvspace{6pt}
%\begin{center}
\def\arraystretch{0.9}
\begin{tabular}{p{10mm}p{10mm}||c|c|c}
%\hline
$x_1$ & $x_2$ & $x_1\wedge x_2$  & $(x_1\wedge x_2)+x_1$ &
$((x_1\wedge x_2)+x_1)+x_2$ \\
\hline
0 & 0 & 0 & 0 & 0 \\
0 & 1 & 0 & 0 & 1 \\
1 & 0 & 0 & 1 & 1 \\
1 & 1 & 1 & 0 & 1 \\
\hline
\end{tabular}

%\addvspace{6pt}

Таким образом, формула $F_3$ также реализует дизъюнкцию.
%\end{enumerate}
%\end{examples}

\end{frame}

\begin{frame}
\label{3.2.2.}\subsubsection{3.2.2. Равносильные формулы }\frametitle{3.2.2. Равносильные формулы (1/2)}
%\subsection{Равносильные формулы}\label{sec1-10}
\vspace{-0.2\baselineskip}
Функция может иметь несколько реализаций над данным базисом.
Формулы, реализующие одну и ту же функцию, называются
\odmkey{равносильными}{Формула (formula)!равносильная (equivalent)}: \\
${\cal F_1}  =  {\cal F_2}  \Def  \E{f}{\func{\cal F_1}  = 
f \And \func{\cal F_2}  =  f}$.
Отношение равносильности формул является эквивалентностью.
Имеют место, в частности, следующие равносильности. 

\vspace{-0.4\baselineskip}
\begin{tabbing}
15.5\= $a \Or (b \wedge c) = (a \Or b) \wedge (a \Or c)$,$\quad$\= \kill
\>1. $a \Or a = a$,
  $a \wedge a = a$.     \\[-2pt]
\>2. $a \Or b = b \Or a$,
  $a \wedge b = b \wedge a$. \\[-2pt]
\>3. $a \Or (b \Or c) = (a \Or b) \Or c$,
  $a \wedge (b \wedge c) = (a \wedge b) \wedge c$. \\[-2pt]
\>4. ($a \wedge b) \Or a = a$,
  $(a \Or b) \wedge a = a$. \\[-2pt]
\>5. $a \Or (b \wedge c) = (a \Or b) \wedge (a \Or c)$,
  $a \wedge (b \Or c) = (a \wedge b) \Or (a \wedge c)$. \\[-2pt]
\>6. $a \Or 1 = 1$,
  $a \wedge 0 = 0$. \\[-2pt]
\>7. $a \Or 0 = a$,
  $a \wedge 1 = a$. \\[-2pt]
\>8. $\neg\neg a = a$. \\[-2pt]
\>9. $\neg(a \wedge b) = \neg a \Or \neg b$,
  $\neg(a \Or b) = \neg a \wedge \neg b$. \\[-2pt]
\>10. $a \Or \neg a = 1$,
  $a \wedge \neg a = 0$.
\end{tabbing}

\vspace{-0.4\baselineskip}
Все указанные равносильности могут быть легко проверены построением соответствующих
таблиц истинности.
Таким образом, $\Alg{E_2}{\Or, \wedge, \neg}$~--- булева алгебра
(см. \DMref{п.~2.6.6}{2.6.6.}).
\end{frame}

\begin{frame}
\frametitle{3.2.2. Равносильные формулы (2/2)}

\begin{remark}
Если некоторая бинарная 
операция ассоциативна,
то в формуле, построенной с 
применением такой операции,
порядок выполнения операций 
не важен и скобки можно 
опустить. Другими словами,
ассоциативная бинарная
операция определяет 
\odmkey{групповую операцию}{Операция (operation)!групповая (group)},
применимую не только к паре, 
но и к любому набору аргументов.
В этом случае употребляют также термины 
\odmkey{кратная операция}{Операция (operation)!кратная (multiple)} и
\odmkey{вариаргументная операция}{Операция (operation)!вариаргументная (varyadic)}.
Если же операция вдобавок коммутативна,
то не только скобки можно опустить,
но и аргументы групповой операции можно записывать в любом порядке.  
Ввиду выполнения равносильностей 2 и 3,
можно определить кратные операции дизъюнкции и конъюнкции, для которых
используются следующие обозначения:


$$
\BigOr_{i=1}^{n} x_i \Def x_1 \Or \dots \Or x_n, \qquad
\bigwedge_{i=1}^{n} x_i \Def x_1 \wedge \dots \wedge x_n \Def x_1  \dots  x_n.
$$


Аналогичными свойствами обладают сложение и умножение чисел,
объединение и пересечение множеств
\end{remark}
\end{frame}

\begin{frame}
\label{3.2.3.}\subsubsection{3.2.3. Подстановка и замена }\frametitle{3.2.3. Подстановка и замена (1/4)}
%\subsection{Подстановка и замена}\label{sec--3-2-3}
Если в формулу $\cal F$ входит переменная $x$, то это
обстоятельство обозначается так: \\$\cal F(\dots x\dots)$.
Соответственно, запись $\cal F(\dots \cal G\dots)$ обозначает, что
в формулу $\cal F$ входит подформула $\cal G$. 

\smallskip
Вместо подформулы
(в частности, вместо переменной) в формулу можно подставить другую
формулу (в частности, переменную), в результате чего получится
новая правильно построенная формула. 

\smallskip
Если подстановка формулы $\cal G$
производится вместо \em{всех\/} вхождений заменяемой переменной
$x$ (или подформулы), то результат подстановки обозначается
следующим образом: $\cal F(\dots x\dots)[\cal G/x]$. Если же
подстановка производится вместо \em{некоторых\/} вхождений (в том
числе вместо одного), то результат подстановки  
не имеет общепринятого обозначения. 
В этом курсе принято дублировать знак подстановки 
и писать между знаками список номеров вхождений, подлежащих замене. При этом если номера не указываются, то подразумевается некоторое одно 
вхождение.



\end{frame}

\begin{frame}
\frametitle{3.2.3. Подстановка и замена (2/4)}
\begin{examples}
\begin{enumerate}
\item Замена всех вхождений переменной:
$ x\vee\neg x [y\wedge z/x] = (y\wedge z)\vee\neg(y\wedge z)$.
\item Замена всех вхождений подформулы:
$x\vee y\vee z [\neg x / y\vee z ] = x\vee\neg x$.
\item 
Замена первого вхождения переменной:
$ x\vee\neg x[y/1/x] = y\vee\neg x$.
\item 
 Замена первого вхождения подформулы:
$x\vee y\vee z [\neg x /1/ y\vee z ] = x\vee\neg x$.
\end{enumerate}
\end{examples}

Известны два правила: \odmkey{правило
 подстановки}{Правило (rule)!подстановки (of substitution)} и
\odmkey{правило замены}{Правило (rule)!замены (of
replacement)}, которые позволяют преобразовывать формулы с
сохранением равносильности.

\begin{theorem} {\NB{ [Правило подстановки]}}  Если в двух
равносильных формулах вместо всех вхождений некоторой переменной
$x$ подставить одну и ту же формулу, то получатся равносильные
формулы:
$$
\forall \cal G \big( \cal F_1(\dots x \dots) = \cal F_2(\dots x \dots)
\Impl\cal F_1(\dots x\dots)[\cal G /x] = \cal F_2(\dots x
\dots)[\cal G /x]\big).
$$
\end{theorem}

\end{frame}

\begin{frame}
\frametitle{3.2.3. Подстановка и замена (3/4)}

\begin{proof}
Чтобы доказать равносильность двух формул, нужно показать, что они
реализуют одну и ту же функцию. А это можно сделать, если взять
произвольный набор значений переменных и убедиться, что значения,
полученные при вычислении формул, совпадают.

Рассмотрим произвольный набор значений $a_1,\dots,a_x,\dots,a_n$
переменных \\ $x_1, \dots$, $x, \dots, x_n$. Обозначим
$a\Assign\Eval(\cal G,F,a_1,\dots,a_x,\dots,a_n)$. По определению
алгоритма интерпретации
$$
\Eval(\cal F_1[\cal G
/x],F,a_1,\dots,a_x,\dots,a_n) = \Eval(\cal F_1,F,a_1,\dots,a,\dots,a_n)
$$
и, аналогично,
$$
\Eval(\cal F_2[\cal G/x],F,a_1,\dots,a_x,\dots,a_n) =
\Eval(\cal F_2,F,a_1,\dots,a,\dots,a_n).
$$
Но $\cal F_1 = \cal F_2$  и,
значит,
$$
\Eval(\cal F_1,F,a_1,\dots,a_x,\dots,a_n) = \Eval(\cal
F_2,F,a_1,\dots,a_x,\dots,a_n),
$$
откуда
$
 \Eval(\cal F_1[\cal G
/x],F,a_1,\dots,a_x,\dots,a_n) = \Eval(\cal F_2[\cal G
/x],F,a_1,\dots,a_x,\dots,a_n).
$
\end{proof}


\end{frame}

\begin{frame}
\frametitle{3.2.3. Подстановка и замена (4/4)}

\begin{remark}
В правиле подстановки условие замены \em{всех} вхождений
существенно: например, $x\vee\neg x  = 1$ и $ x\vee\neg x
[y/x] = y\vee\neg y = 1$, но $ x\vee\neg x[y/1/x] = y\vee\neg x
\not= 1$.
\end{remark}

\begin{theorem} {\NB{ [Правило замены]}}
Если в формуле заменить некоторую подформулу равносильной формулой, то
получится формула, равносильная исходной:
$$
\A{\cal F(\dots \cal G_1 \dots)} {\cal G_1 = \cal G_2 \Impl \cal
F(\dots \cal G_1 \dots) = \cal F(\dots \cal G_1 \dots)[\cal
G_2//\cal G_1]}.
$$
\end{theorem}
\begin{proof}
Рассмотрим произвольный набор значений $a_1,\dots,a_n$ переменных \\
$x_1,\dots,x_n$. Имеем   $\cal G_1 = \cal G_2$ и, значит,
$
\Eval(\cal
G_1,F,a_1,\dots,a_n) = \Eval(\cal G_2,F,a_1,\dots,a_n),
$
откуда 
$
\Eval(\cal F[\cal G_1 //\cal G_1],F,a_1,\dots,a_n) =
\Eval(\cal F[\cal G_2 //\cal G_1],F,a_1,\dots,a_n).
$
\end{proof}

\medskip
Пусть $F = \{\Tuple{f}{1}{m}\}$ и $G = \{\Tuple{g}{1}{m}\}$.
Тогда говорят, что формулы ${\cal F}[F]$ и~${\cal G}[G]$ имеют
\em{одинаковое строение\/},
если ${\cal F}$ совпадает с результатами подстановки в формулу ${\cal G}$
функций $f_i$ вместо функций $g_i$:
$
\cal F[F] = \cal G[G][f_i^{}/g_i^{}]_{i=1}^{m}.
$
\end{frame}

\begin{frame}
\label{3.2.4.}\subsubsection{3.2.4. Алгебра булевых функций }\frametitle{3.2.4. Алгебра булевых функций (1/3)}

Булевы функции $\Or$, $\wedge$, $\neg$ (и любые
другие) могут рассматриваться как операции на множестве булевых функций,
$\Map{\Or, \wedge}{P_n\times P_n}{P_n}$, $\Map{\neg}{P_n}{P_n}$.
Действительно, пусть формулы $\cal F_1$ и $\cal F_2$ равносильны и реализуют
функцию $f$, а~формулы $\cal G_1$ и $\cal G_2$ равносильны и реализуют
функцию $g$:
$\func{\cal F_1} = f$,
$\func{\cal F_2} = f$,
$\func{\cal G_1} = g$,
$\func{\cal G_2} = g$.
Тогда, применяя правило замены нужное число раз, имеем
$$
{\cal F_1} \Or {\cal G_1} = {\cal F_2} \Or {\cal G_2}, \quad 
{\cal F_1} \wedge {\cal G_1} = {\cal F_2} \wedge {\cal G_2},\quad 
\neg{\cal F_1} = \neg{\cal F_2}.$$

Таким образом, если взять любые формулы $\cal F$ и $\cal G$,
реализующие функции $f$ и~$g$ соответственно, то каждая из формул
$ \cal F \wedge \cal G$, $\cal F\Or\cal G$ и $\neg \cal F$
реализует одну и ту же функцию, независимо от выбора
реализующих формул $\cal F$ и $\cal G$. Следовательно, функции,
которые реализуются соответствующими формулами, можно по
определению считать результатами применения соответствующих
операций. Другими словами, если
$\func\cal F = f$, $\func{\cal G} = g$,
то
$f\wedge g \Def \func(\cal F\wedge \cal G)$,
$f\Or g \Def \func(\cal F\Or \cal G)$, 
$\neg f \Def \func (\neg \cal F)$.
Алгебраическая структура $\Alg{P_n}{\Or, \wedge, \neg}$ называется
\odmkey{алгеброй булевых функций}{Алгебра (algebra)!булевых
функций (of boolean functions)}.


\end{frame}

\begin{frame}
\frametitle{3.2.4. Алгебра булевых функций (2/3)}
\begin{theorem}
Алгебра
булевых функций является булевой алгеброй.
\end{theorem}
\begin{proof}
Действительно, пусть
равносильности, 
перечисленные в \DMref{п.~3.2.3}{3.2.3.},
проверены путем построения таблиц истинности.
Ясно, что эти таблицы не зависят от того, откуда
взялись значения $a$, $b$, $c$. Таким образом, вместо $a$, $b$, $c$
можно подставить любые функции, а значит, любые реализующие их формулы,
поскольку выполнено правило подстановки.
Следовательно, аксиомы
булевой алгебры выполнены в алгебре $\Alg{P_n}{\Or, \wedge, \neg}$.
\end{proof}


Пусть $[{\cal F}]$~--- множество формул, равносильных
${\cal F}$ (то~есть класс эквивалентности по отношению равносильности).
Рассмотрим множество $\cal K$ классов эквивалентности по
отношению равносильности.
Пусть операции
$\Map{\Or,\wedge}{\cal K\times \cal K}{\cal K}$,
$\Map{\neg}{\cal K}{\cal K}$
определены
 следующим образом:
$
[{\cal F_1}] \Or [{\cal F_2}] \Def [{\cal F_1} \Or {\cal F_2}], \
[{\cal F_1}] \wedge [{\cal F_2}] \Def [{\cal F_1} \wedge {\cal F_2}], \
\neg[{\cal F_1}] \Def [\neg{\cal F_1}].
$

\smallskip
Тогда функция $\func$, сопоставляющая формуле булевскую функцию, которую она реализует,
является гомоморфизмом относительно операций $\wedge, \vee, \neg$, и к ней можно применить теорему о гомоморфизме 
(\DMref{п.~1.7.6}{1.7.6.}).
В этом случае $\Im{\func} \subset P_n$ и 
$\Dom{\func} = \mathfrak{F}[\vee, \wedge, \neg]$. 

\end{frame}

\begin{frame}
\frametitle{3.2.4. Алгебра булевых функций (3/3)}



\begin{remark}
В пп.~\DMref{3.4.3}{3.4.3.} и \DMref{3.6.4}{3.6.4.} показано, 
что $\func(\mathfrak{F}[\vee, \wedge, \neg]) = P_n$, то есть достаточно
рассматривать только формулы, построенные
над базисом $\{\vee, \wedge, \neg\}$.   
\end{remark}

\smallskip
Применение теоремы о гомоморфизме гласит
$$
\mathfrak{F}[\vee, \wedge, \neg]/\ker\func
\sim\Alg{P_n}{\vee, \wedge, \neg}.
$$
Другими словами, алгебра классов равносильных формул 
(\odmkey{алгебра Линденбаума--Тарского}{Алгебра (algebra)!Линденбаума--Тарского \\(Lindebaum--Tarski)}) изоморфна
алгебре булевых функций и, следовательно, является булевой алгеброй. \\ Носитель этой
алгебры~--- множество классов формул.
\begin{digress}
На практике мы говорим о функциях, а пишем формулы, хотя формулы и
функции~--- разные вещи. Например, формул бесконечно много, а
функций (булевых функций $n$ переменных)~--- только конечное
число, и свободная алгебра формул \em{не изоморфна} алгебре
функций. Но алгебра функций \odmkey{изоморфна}{Изоморфизм
(isomorphism)!алгебр (of algebras)} алгебре классов равносильных
формул, что позволяет манипулировать формулами, имея в виду
функции.
\end{digress}


\end{frame}

\subsection{3.3. Двойственность и симметрия}\label{3.3.}
\begin{frame}
\frametitle{3.3. Двойственность и симметрия}


Понятия двойственности и симметрии с
успехом применяется в самых разных областях математики. Мы рассматриваем
двойственность и симметрию на простейшем примере булевых функций.

\medskip
\begin{tabular}{p{12mm}|p{65mm}|p{65mm}}
  № & Название параграфа
  & Ключевые термины и обозначения \\
  \hline
\DMref{3.3.1.}{3.3.1.}\RS & Двойственные и самодвойственные функции  &
{ Мажоритарная функция }\\  
\DMref{3.3.2.}{3.3.2.} & Реализация двойственной функции &
{ }\\  
\DMref{3.3.3.}{3.3.3.}\RS & Принцип двойственности  &
{  }\\  
\DMref{3.3.4.}{3.3.4.} & Симметрические функции  & {}\\
\end{tabular}  


\end{frame}


\begin{frame}
\label{3.3.1.}\subsubsection{3.3.1. Двойственные и самодвойственные функции }\frametitle{3.3.1. Двойственные и самодвойственные функции (1/3)}

%\subsection{Двойственная функция}

Пусть $f(\Tuple{x}{1}{n}) \in P_n$~--- булева функция. Тогда
функция ${f^*}(\Tuple{x}{1}{n})$, определённая следующим образом:

\vspace{-16pt}
$$
{f^*}(\Tuple{x}{1}{n}) \Def  \overline{f(\Tuple{\overline
x}{1}{n})},
$$

\vspace{-4pt}
называется
\odmkey{двойственной}{Функция (function)!двойственная (dual)} к функции
$f$.
Из определения  видно, что двойственность инволютивна: $f^{**}=f$,
и по этой причине отношение <<быть двойственной к>>
 на множестве булевых функций
симметрично, то~есть если $f^*=g$, то $g^*=f$.
Если в таблице истинности булевой функции $f$ инвертировать
\em{все} значения, то получим таблицу истинности
двойственной функции $f^*$.


\begin{example}

\begin{tabular}{cc|c}
 $x_1$ & $x_2$ & $x_1\wedge x_2$\\
\hline
0 & 0 & 0  \\
0 & 1 & 0  \\
1 & 0 & 0  \\
1 & 1 & 1  \\
\end{tabular}
\hspace{6pt}
\begin{tabular}{cc|c}
 $x_1$ & $x_2$ & $(x_1\wedge x_2)^*$\\
\hline
1 & 1 & 1  \\
1 & 0 & 1  \\
0 & 1 & 1  \\
0 & 0 & 0  \\
\end{tabular}
$ = $
\begin{tabular}{cc|c}
 $x_1$ & $x_2$ & $(x_1\wedge x_2)^*$\\
\hline
0 & 0 & 0  \\
0 & 1 & 1  \\
1 & 0 & 1  \\
1 & 1 & 1  \\
\end{tabular}
\end{example}


Таким способом можно определить двойственную функцию для любой булевой
функции.
\end{frame}

\begin{frame}
\frametitle{3.3.1. Двойственные и самодвойственные функции (2/3)}
\begin{remark}
Если функции двойственны, то их отрицания также двойственны.
\end{remark}


\begin{example}
В таблице ниже приведены все пары двойственных функций двух переменных. Для каждой функции указан вектор значений при установленном порядке перечисления кортежей и обозначение, если оно введено в таблице \DMref{п.~3.1.4}{3.1.4.}.

\medskip
\begin{center}
\begin{tabular}{|c|c|c|c|c|c|c|c|c|c|c|}
\hline
$f$ & 0000 & 0001 & 0010 & 0011 & 0100 & 0101 & 0110 & 1000 & 1010 & 1100 \\
    & 0 & $x\And y$ & & $x$ & & $y$ 
    & $x +_{2} y$ & $x\downarrow y$ & $\Not y$ & $\Not x$ \\
\hline
    & 1 & $x\Or y$ & & $x$ & & $y$ 
    & $x \equiv y$ & $x\mid y$ & $\Not y$ & $\Not x$
    \\
$f^*$ & 1111 & 0111 & 1011 & 0011 & 1101 & 0101 & 1001 & 1110 & 1010 & 1100 \\
\hline
\end{tabular}
\end{center}
\end{example}

\medskip
Функция называется
\odmkey{самодвойственной\/}{Функция (function)!самодвойственная
(self-daul)}, если $f^* =
f$.

\begin{example}
 Тождественная функция и отрицание самодвойственны, 
 все функции двух существенных переменных, а также константы не самодвойственны. 
\end{example}

\begin{remark}
Если функция самодвойственна, то ее отрицание также самодвойственно.
\end{remark}
\end{frame}

\begin{frame}
\frametitle{3.3.1. Двойственные и самодвойственные функции (3/3)}
\begin{theorem}
Существует $2^{2^{n-1}}$ самодвойственных функций $n$ переменных.
\end{theorem}
\begin{proof}
Из алгоритма построения двойственной функции следует, 
что функция самодвойственна тогда и только тогда, 
когда вектор значений функции в таблице 
истинности является зеркально инверсным относительно своей середины. 
Получается, что первая половина вектора значений 
самодвойственной функции однозначно задаёт вторую.
\end{proof}

\begin{example}
Функция <<голосования>> $m(x,y,z)$ 
(эту функцию называют также \odmkey{мажоритарной}{Функция (function)!мажоритарная (majority)}) 
принимает значение 1, 
если большинство переменных принимают значение 1.
Из приведённой таблицы истинности функции голосования видно, 
что она самодвойственна.

\vspace{-10pt}
\begin{center}
$\quad$
\begin{tabular}{|c|c|c|c|c|c|c|c|c|}
\hline
$x$   & 0 & 0 & 0 & 0 & 1 & 1 & 1 & 1 \\
$y$   & 0 & 0 & 1 & 1 & 0 & 0 & 1 & 1 \\
$z$   & 0 & 1 & 0 & 1 & 0 & 1 & 0 & 1 \\
\hline
$m$   & 0 & 0 & 0 & 1 & 0 & 1 & 1 & 1 \\
\hline
\end{tabular}
\end{center}
\end{example}

\end{frame}

\begin{frame}
\label{3.3.2.}\subsubsection{3.3.2. Реализация двойственной функции}
\frametitle{3.3.2. Реализация двойственной функции}

%\subsection{Реализация двойственной функции}
\vspace{-4pt}
Формула, реализующая двойственную функцию, определённым образом
связана с формулой, реализующей исходную функцию.

\begin{remark}
Далее используется обозначение ${\overline f}(\cdots) \Def
\overline {f(\cdots)}$.
\end{remark}

\begin{theorem}
Если функция $\varphi(\Tuple{x}{1}{n})$ реализована формулой
$$
f(f_1(\Tuple{x}{1}{n}), \dots, f_n(\Tuple{x}{1}{n})),
$$
то формула
$$
f^*(f_1^{*}(\Tuple{x}{1}{n}),\dots,f_n^{*}(\Tuple{x}{1}{n}))
$$
реализует функцию $\varphi^*(\Tuple{x}{1}{n})$.
\end{theorem}

\begin{proof}
\begin{align*}
\varphi^*(\Tuple{x}{1}{n})&=
  \overline\varphi(\Tuple{\overline x}{1}{n})\!=\!
  \overline f(f_1(\Tuple{\overline x}{1}{n}),\dots,
  f_n(\Tuple{\overline x}{1}{n}))\!=\\
& =\overline f
   (\overline{\overline{f_1}}(\Tuple{\overline x}{1}{n}),\dots,
   (\overline{\overline{f_n}}(\Tuple{\overline x}{1}{n}))=\\
& =\overline f(\overline{f_1^{*}}(\Tuple{x}{1}{n}),\dots,
   \overline{f_n^{*}}(\Tuple{x}{1}{n}))=\\
& =f^*(f^{*}_1(\Tuple{x}{1}{n}),\dots,
   f_n^{*}(\Tuple{x}{1}{n})).
\end{align*}

\vspace{-18pt}
\end{proof}

\end{frame}


\begin{frame}
\label{3.3.3.}\subsubsection{3.3.3. Принцип двойственности }\frametitle{3.3.3. Принцип двойственности (1/2)}

\vspace{-2pt}
\odmkey{Принцип двойственности}{Принцип (principle)!двойственности
(of duality)} устанавливает связь между структурами формул,
реализующих пару двойственных функций. Рассмотрим две системы
булевых функций, $F=\{\Tuple{f}{1}{m}\}$ и
$F^{*}=\{f^*_1,\dots,f^*_m\}$, и введём обозначение
$$
\cal F^*[F^*]\Def \cal F[F][f_i^*/f_i^{}]_{i=1}^m.
$$



\begin{theorem} \NB{[Принцип двойственности]}
Пусть $F=\{\Tuple{f}{1}{m}\}$~{\upshape---} система булевых
функций, а $F^{*}=\{f^*_1,\dots,f^*_m\}$~{\upshape---} система
двойственных функций. Тогда если формула ${\cal F}$ над базисом
$F$ реализует функцию $f$, то формула ${\cal F^*}$ над базисом
$F^*$, полученная заменой функций $f_i$ двойственными функциями
${{f_i}^*}$, реализует функцию $f^*$:\\
$
\func{\cal F}[F] = f \Impl \func{\cal F^*}[F^*] = f^*.
$
\end{theorem}
%\vspace{-0.25\baselineskip}
%\pagebreak[3]
\begin{proof}
Индукция по структуре формулы $\cal F$.
База: если формула $\cal F$ имеет вид
$f(\Tuple{x}{1}{n})$,
где $f\in F$, то формула $\cal
F^*=f^*(\Tuple{x}{1}{n})$ реализует функцию ${f^*}$ по
определению. Индукционный переход по теореме \DMref{п.~3.3.2}{3.3.2.}.
\end{proof}
\begin{corollary}
${\cal F_1} = {\cal F_2} \Impl {{\cal F_1}^*} = {\cal F_2}^*$.
\end{corollary}

%\vspace{-.5\baselineskip}


\end{frame}


\begin{frame}
\frametitle{3.3.3. Принцип двойственности (2/2)}
\begin{example}
Из $\overline{x_1\wedge x_2}=\overline x_1\vee \overline x_2$ по
принципу двойственности сразу имеем $\overline {x_1\vee
x_2} =\overline x_1\wedge \overline x_2$.
\end{example}

\smallskip
\begin{digress}
Хорошо известен принцип математической индукции для натуральных
чисел:
$$
(P(1)\And (P(n)\Impl P(n+1)))\Impl \A{n\in\Bbb N}{P(n)}.
$$
Этот принцип является справедливым и для других множеств,
упорядоченных более сложным образом, нежели натуральные числа
(см.~\DMref{п.~1.8.9}{1.8.9.}).
Например, в доказательстве предыдущей теоремы был использован
принцип индукции в следующей форме.

Пусть задана
некоторая иерархия. Предположим, что
1) некоторое утверждение $P$ справедливо для всех узлов иерархии нижнего
уровня;
2) из того, что утверждение $P$ справедливо для всех узлов, подчинённых
данному узлу, следует, что утверждение $P$ справедливо для данного
узла.
Тогда утверждение $P$ справедливо для всех узлов иерархии.
\end{digress}
\end{frame}

\begin{frame}
\label{3.3.4.}\subsubsection{3.3.4. Симметрические функции }\frametitle{3.3.4. Симметрические функции (1/2)}

%\subsection{Симметрические функции}
Булева функция $n$ переменных называется
\odmkey{симметрической}{Функция (function)!симметрическая (symmetric)},
если её значение инвариантно относительно перестановок значений
аргументов. Заметим, что при перестановке значений в наборе общее
количество нулей и единиц сохраняется.
Понятие симметрической функции применимо только к функциям двух и большего числа переменных.


\begin{examples} Конъюнкция, сложение по модулю 2,
дизъюнкция, стрелка Пирса, эквивалентность и штрих Шеффера являются симметрическими булевыми функциями. Других
симметрических булевых функций двух существенных переменных нет.
\end{examples}


\begin{digress}
Симметрические функции изучаются и используются во многих 
разделах математики. Например, обычные сложение и умножение 
в арифметике уже являются симметрическими функциями. 
Формула Герона площади треугольника \\ $S(a,b,c)=
\sqrt{p(p-a)(p-b)(p-c)}$, где $p\Assign (a+b+c)/2$ и 
формула определения расстояния между точкой и началом 
координат $R(x_1,x_2,x_3)=\sqrt{x_1^2+x_2^2+x_3^2}$ 
являются примерами симметрических функций в элементарной геометрии.
\end{digress}



\end{frame}

\begin{frame}
\frametitle{3.3.4. Симметрические функции (2/2)}

\begin{theorem}
Существует $2^{n+1}$ различных симметрических булевых функций $n$ переменных.
\end{theorem}

\begin{proof} 
Ясно, что на любом наборе значений $n$ переменных, в котором содержится
$k$ значений 1 (и, соответственно, $n-k$ значений 0), симметрическая
функция принимает одно и то же значение (1 или 0).
Таким образом,
симметрическую функцию $f$ от $n$ переменных можно задать булевским
вектором $\left(f_0, f_1, \dots, f_n\right)$, где $f_i$ --- значение
функции на наборах значений, содержащих $i$ единиц.
Всего существует $2^{n+1}$ таких наборов.
\end{proof}

%%%%%%%%%%%%%%%%%%%%%%%%% Конец изменений 2012
\medskip
\begin{examples}
\begin{enumerate}
\item Функция голосования является симметрической.
\item Формула
$\neg x \wedge y \wedge z \Or x \wedge \neg y \wedge z \Or
x \wedge y \wedge \neg z$ реализует функцию, 
которая также является симметрической. 
\end{enumerate}
\end{examples}

\end{frame}

\begin{frame}
\frametitle{3.4. Нормальные формы (1/2)}

\subsection{3.4. Нормальные формы}\label{3.4.}

В данном разделе на примере
булевых функций обсуждается важное понятие <<нормальная форма>>,
то есть синтаксически однозначный способ записи формулы,
реализующей заданную функцию.

\medskip
\begin{tabular}{p{12mm}|p{45mm}|p{85mm}}
  № & Название параграфа
  & Ключевые термины и обозначения \\
  \hline
\DMref{3.4.1.}{3.4.1.}\RS & Разложение булевых функций по переменным  &
{ Мажоритарная функция }\\  
\DMref{3.4.2.}{3.4.2.} & Минимальные термы &
{Минтерм, ортогональная и выполнимая система булевых функций, макстерм }\\  
\DMref{3.4.3.}{3.4.3.}\RS & Совершенные нормальные формы &
{ Разрешимость, СДНФ, СКНФ }\\  
\DMref{3.4.4.}{3.4.4.} & Эквивалентные преобразования & {Правило склеивания/расщепления,
элиминация операций, протаскивание отрицаний,
дизъюнктивная форма, удаление нулей, расщепление 
переменных, системы подстановок термов, системы правил переписывания}\\
\end{tabular}  




\end{frame}

\begin{frame}
\frametitle{3.4. Нормальные формы (2/2)}


\medskip
\begin{tabular}{p{12mm}|p{45mm}|p{85mm}}
  № & Название параграфа
  & Ключевые термины и обозначения \\
  \hline
\DMref{3.4.5.}{3.4.5.} & Минимальные дизъюнктивные формы & {Ранг конъюнкции}
\\  
\DMref{3.4.6.}{3.4.6.} & Геометрическая интерпретация & {}
\\  
\DMref{3.4.7.}{3.4.7.} & Сокращённые дизъюнктивные формы & {Допустимая конъюнкция, максимальная конъюнкция}
\end{tabular}  
\end{frame}

\begin{frame}
\label{3.4.1.}
\subsubsection{3.4.1. Разложение булевых функций по переменным }
\frametitle{3.4.1. Разложение булевых функций по переменным (1/2)}


Положим $x^y \Def  x\wedge y \Or \overline x\wedge \overline y$.
Очевидно, что
$$x^y=
x \equiv y 
 = \textbf{ if } y \textbf{ then } x \textbf{ else } \overline{x} \textbf{ end if } =
\textbf{ if } x=y \textbf{ then } 1 \textbf{ else } 0 \textbf{ end if}.$$


\begin{theorem} \NB{ [О разложении булевой функции по переменным]}

$
f(\Tuple{x}{1}{m},\Tuple{x}{m+1}{n})
=\BigOr\limits_{(\Tuple{\sigma}{1}{m})}
x_1^{\sigma_1}\wedge \ldots \wedge x_m^{\sigma_m} \wedge
f(\Tuple{\sigma}{1}{m},\Tuple{x}{m+1}{n}),
$
где дизъюнкция берётся по всем возможным наборам
 $(\Tuple{\sigma}{1}{m})$.
\end{theorem}

\begin{proof}
Рассмотрим значение формулы в правой части на наборе значений \\ $a_1,\dots , a_n$. 
Имеем
$
\left(\BigOr\limits_{(\Tuple{\sigma}{1}{m})}x_1^{\sigma_1}\wedge\ldots\wedge
x_m^{\sigma_m}\wedge
f(\Tuple{\sigma}{1}{m},\Tuple{x}{m+1}{n})\right)(\Tuple{a}{1}{n})=
$\\
$
=\BigOr\limits_{(\Tuple{\sigma}{1}{m})}a_1^{\sigma_1}\wedge\ldots\wedge
a_m^{\sigma_m}\wedge
f(\Tuple{\sigma}{1}{m},\Tuple{a}{m+1}{n}).
$
Все конъюнкции, в которых \\
 $\E{i}{a_i \neq \sigma_i}$, равны 0,
и их можно опустить, поэтому остаётся только одно
слагаемое, для которого $\A{i\in 1..m}{a_i=\sigma_i}$,
значит \\
$ a_1^{a_1}\wedge\ldots\wedge
a_m^{a_m}\wedge f(\Tuple{a}{1}{m},\Tuple{a}{m+1}{n})
=f(\Tuple{a}{1}{n}).
$
\end{proof}
\end{frame}

\begin{frame}
\frametitle{3.4.1.Разложение булевых функций по переменным (2/2)}

%\vspace{-.5\baselineskip}

\begin{remark}
Здесь доказывается, что некоторая формула реализует заданную функцию.
Для этого достаточно взять \em{произвольный\/} набор значений аргументов
функции, вычислить на этом наборе значение формулы, и если оно окажется
равным значению функции на этом наборе аргументов, то из этого следует
доказываемое утверждение.
\end{remark}

%\vspace{-6pt}
\begin{corollary} \NB{ [Разложение булевой функции по одной переменной]}
$$
f(\Tuple{x}{1}{n-1},x_n) = \left( x_n \wedge f(\Tuple{x}{1}{n-1},1) \right) \Or
\left( \overline x_n \wedge f(\Tuple{x}{1}{n-1},0) \right).
$$
\end{corollary}


\begin{corollary} \NB{ [Разложение булевой функции по всем переменным]}
$$f(\Tuple{x}{1}{n}) =
\BigOr_{\left\{(\Tuple{\sigma}{1}{n}) |
f(\Tuple{\sigma}{1}{n})=1\right\}} x_1^{\sigma_1} \wedge \ldots
\wedge x_n^{\sigma_n}. $$
\end{corollary}

\begin{example}
Разложение функции голосования по переменной $x$:\\
$m(x,y,z) = (x \wedge m(1,y,z))\vee (\overline{x} \wedge m(0,y,z)) =
(x \wedge (y\vee z))\vee (\overline{x} \wedge (y\wedge z))$.


\end{example}

\end{frame}

\begin{frame}
\label{3.4.2.}\subsubsection{3.4.2. Минимальные термы }\frametitle{3.4.2. Минимальные термы (1/2)}
%\subsection{Минимальные термы}
%\vspace{-2pt}
Выражение вида $x_1^{\sigma_1}\wedge\dots\wedge x_n^{\sigma_n}$ имеет большое значение в теории булевых функций и носит
специальное название:
\odmkey{минимальный терм}{Терм (term)!минимальный  (minimal)},
или
\odmkey{совершенный одночлен}{Совершенный одночлен (perfect monomial)},
или
\odmkey{конституента единицы}{Конституента единицы (unit constituent)},
или, наиболее коротко,
\odmkey{минтерм}{Минтерм (minterm)}.
Нетрудно видеть, что минтерм --- это реализация булевой функции $n$
переменных, которая имеет значение 1 ровно на одном наборе значений
переменных, а~именно на наборе $\left(\Tuple{\sigma}{1}{n}\right)$.
Отсюда  вытекает еще одно следствие к теореме о разложении булевой функции по переменным.
%\setcounter{corollary}{2}
\begin{corollary}
Любую булеву функцию, кроме 0, можно представить как дизъюнкцию некоторых минтермов.
\end{corollary}

Ясно, что при фиксированном $n$ имеется ровно $2^n$ различных минтермов.
Обычно минтерм обозначают $m_i$, 
где $i\in 0..2^n-1$, $i$ называется
номером минтерма, или прямо указывают булевский вектор
$\left(\Tuple{\sigma}{1}{n} \right)$, на котором минтерм принимает
значение $1$. Из структуры формулы, задающей минтерм, 
следует, что двоичное представление номера минтерма задаёт тот набор
значений (в~установленном порядке),
на котором минтерм имеет значение 1.
%%%%%%%%%%%%%%%%%%%%%%%%% Изменения 2012
\end{frame}

\begin{frame}
\frametitle{3.4.2. Минимальные термы (2/2)}
\vspace{-2pt}
Говорят, что система булевых функций
\odmkey{ортогональна}{Ортогональная система булевых функций (Orthogonal system of Boolean functions)},
если их конъюнкция есть тождественный ноль.

\begin{example}
Система $\{ x, \overline{x}\}$ ортогональна.
\end{example}

 
%\setcounter{lemma}{0}
\begin{lemma} \NB{[1]}
Любое множество  различных минтермов ортого\-нально.
\end{lemma}
\begin{proof}
Не существует такого набора значений, 
на котором два  различных минтерма принимают значение 1.	
\end{proof}

Говорят, что система булевых функций
\odmkey{выполнима}{Выполнимость (satisfiability)!
системы булевых функций (system of Boolean functions)},
если их дизъюнкция не есть тождественный ноль. 

\begin{example}
Система $\{x, \overline{x}\}$ выполнима.
\end{example}


\begin{lemma} \NB{[2]}
Любое множество минтермов выполнимо.
\end{lemma}
\begin{proof}
Дизъюнкция  выполняется на всех наборах, на которых 
какой-либо  из  минтермов имеет значение 1.	
\end{proof}


Двойственным к минтерму является
\odmkey{макстерм}{Макстерм (maxterm)} --- булева функция, которая принимает значение 0 на
единственном наборе значений переменных. Ясно, что макстерм реализуется
формулой $x_1^{\sigma_1} \Or \dots \Or x_n^{\sigma_n}$.
Используя принцип двойственности, на макстермы нетрудно распространить все наблюдения для минтермов.
%%%%%%%%%%%%%%%%%%%%%%%%% Конец изменений 2012

\end{frame}

\begin{frame}
\label{3.4.3.}\subsubsection{3.4.3. Совершенные нормальные формы }\frametitle{3.4.3. Совершенные нормальные формы (1/4)}
%\subsection{Совершенные нормальные формы}
%\label{normforms}
%\vspace{-2pt}
Говорят, что некоторый класс формул $\cal K$ \odmkey{имеет
нормальную форму}{Форма (form)!нормальная (normal)}, если задан
другой класс формул $\cal K'$, которые называются нормальными
формами, такой, что любая формула класса $\cal K$ имеет
\emph{единственную\/} равносильную формулу из класса $\cal K'$.
Если задан алгоритм, позволяющий для любой формулы построить её
нормальную форму, то наличие у класса формул нормальной формы
обеспечивает \odmkey{разрешимость}{Разрешимость
(solvability)!класса формул (for formulas class)}, то есть наличие
алгоритма проверки равносильности. Действительно, в этом случае
достаточно сравнить нормальные формы двух формул.
Один и тот же класс $\cal K$ может иметь несколько различных
нормальных форм, то есть несколько различных классов $\cal K'$ 
(см. пп.~\DMref{2.2.3}{2.2.3.} и \DMref{2.2.4}{2.2.4.}).

\medskip
Рассмотрим понятие нормальной формы применительно к булевым функциям.
\odmkey{Cовершен\-ной дизъюнктивной нормальной формой}
{Форма (form)!совершенная нормальная (perfect normal)!дизъюнктивная
(disjunctive)} (СДНФ) называется реализация булевой функции
 $f(\Tuple{x}{1}{n})$ в виде формулы
$
\BigOr\limits_{\left\{(\Tuple{\sigma}{1}{n}) |
f(\Tuple{\sigma}{1}{n})=1\right\}} x_1^{\sigma_1} \wedge \ldots
\wedge x_n^{\sigma_n}.
$

\begin{remark}
СДНФ называется совершенной, потому что каждое слагаемое в
дизъюнкции включает все переменные; дизъюнктивной, потому что
главная операция~--- дизъюнкция;
нормальной, потому что СДНФ единственна,
как показано в следующей теореме.
\end{remark}

\end{frame}

\begin{frame}
\frametitle{3.4.3. Совершенные нормальные формы (2/4)}

%\vspace{-.5\baselineskip}

\begin{theorem}
Всякая булева функция имеет единственную СДНФ.
\end{theorem}

\begin{proof}
По следствию к теореме 
о разложении булевой функции по переменным имеем
\vspace{-4pt}
$$f(\Tuple{x}{1}{n})= 
\BigOr_{\left\{(\Tuple{\sigma}{1}{n}) |
f(\Tuple{\sigma}{1}{n})=1\right\}} x_1^{\sigma_1} \wedge \ldots
\wedge x_n^{\sigma_n}.$$
Если при записи СДНФ используется
установленный порядок, то СДНФ однозначно
определяет множество наборов значений переменных, на которых
функция, реализуемая \\СДНФ, принимает значение 1. Тем самым СДНФ
однозначно определяет таблицу истинности реализуемой функции, 
откуда любые две различные СДНФ над одним набором
переменных неравносильны.

\begin{examples}
\begin{enumerate}
\item СДНФ конъюнкции: 
$x\wedge y= x^1\wedge y^1$. 
\item СДНФ дизъюнкции: 
$x\vee y= x^0 \wedge y^1 \vee x^1\wedge 
y^0 \vee x^1\wedge y^1$.
\item СДНФ импликации: 
$x\to y = x^0 \wedge y^0 \vee x^0\wedge y^1 \vee x^1\wedge y^1$.
\end{enumerate}
\end{examples}

\end{proof}

\end{frame}

\begin{frame}
\frametitle{3.4.3. Совершенные нормальные формы (3/4)}

\begin{remark}
При рассмотрении дизъюнктивных и конъюнктивных
форм мы имеем дело с~формулами вида
$
\BigOr\limits_{i \in I} S_i$ и 
$ \bigwedge\limits_{i \in I} S_i$,
где $I$~--- некоторое множество индексов, а $S_i$~--- некоторые
формулы. Множество индексов может быть пусто. В этом случае удобно
считать, что пустая дизъюнкция имеет значение 0, а пустая
конъюнкция~--- значение 1:
\vspace{-0.25\baselineskip}
$$
I=\emptyset\Impl\BigOr_{i\in I} S_i =0,\quad
I=\emptyset\Impl\bigwedge_{i\in I} S_i =1.
$$
\end{remark}




\begin{corollary}
Всякая булева функция может быть выражена через дизъюнкцию, конъюнкцию и
отрицание:
$$
\A{f\in P_n}
{\E{\cal F\in \mathfrak{F}[\vee, \wedge,\neg]}{f = \func{\cal F}}}.
$$
\end{corollary}

\begin{remark}
Если $f=0$, то
$\Set{(\Tuple{\sigma}{1}{n})}{f(\Tuple{\sigma}{1}{n})=1}=\emptyset$.
В соответствии с соглашением предыдущего замечания пустая формула
считается СДНФ нуля.
\end{remark}


\end{frame}

\begin{frame}
\frametitle{3.4.3. Совершенные нормальные формы (4/4)}
%\vspace{-6pt}


\begin{theorem}
Всякая булева функция имеет единственную
\odmkey{совершенную конъюнктивную нормальную форму\/}
{Форма (form)!совершенная нормальная (perfect normal)!конъюнктивная
(conjunctive)} (СКНФ):
$$f(\Tuple{x}{1}{n}) =
\bigwedge_{\left\{(\Tuple{\sigma}{1}{n}) |
f^*(\Tuple{\sigma}{1}{n})=1\right\}} x_1^{\sigma_1} \Or \dots \Or
x_n^{\sigma_n}.$$
\end{theorem}

%\vspace{-6pt}
\begin{proof}
По принципу двойственности из предыдущей теоремы.
\end{proof}
\begin{examples}
Нормальные формы существуют у различных формул.

\begin{enumerate}
\item СДНФ и СКНФ являются нормальными формами для булевых формул
над любым базисом. 

\item Формулы вида
$\sum_{i=0}^{n} a_i x^i$ и $(\ldots(a_n x +
a_{n-1})x+\ldots+a_1)x + a_0$
являются нормальными формами для полиномов одной переменной
степени~$n$. 

\item Множество формул, построенных суперпозицией из рациональных
функций одной переменной, $e^x$, $\ln x$, прямых и обратных тригонометрических функций (то есть множество
формул, реализующих элементарные функции математического анализа),\\
нормальной формы не имеет. Доказательство последнего утверждения
выходит далеко за рамки этого курса.

\end{enumerate}
\end{examples}


\end{frame}

\begin{frame}
\label{3.4.4.}\subsubsection{3.4.4. Эквивалентные преобразования }\frametitle{3.4.4. Эквивалентные преобразования (1/8)}
%\subsection{Эквивалентные преобразования}
Используя уже доказанные равносильности, можно преобразовывать
по правилу замены одни формулы в другие, равносильные им.
Преобразование формулы в~равносильную ей называется
\odmkey{эквивалентным
преобразованием}{Преобразование (transformation)!эквивалентное
 (equivalent)}.

\begin{example} Используя равносильности из \DMref{п.~3.2.2}{3.2.2.},
покажем, что имеет место
\odmkey{правило склеивания/расщепления}
{Правило (rule)!склеивания/расщепления \\ (clutching/splitting)}:
$(x\And y)\Or (x\And\Not y)=x$.
Действительно,
$$
(x\And y)\Or(x\And\Not y)\ToE{5}
x\And(y\Or\Not y)\ToE{10}x\And 1\ToE{7}x
.$$
\end{example}

\begin{remark}
Если равносильность из предыдущего примера применяется для
уменьшения числа операций в формуле, то говорят, что производится
\odmkey{склеивание}{Склеивание (clutching)},
а если наоборот, то
\odmkey{расщепление}{Расщепление (splitting)}.
\end{remark}

\begin{theorem} Для любых двух равносильных формул ${\cal F_1}$ и ${\cal F_2}$
существует последовательность эквивалентных преобразований из
${\cal F_1}$ в ${\cal F_2}$, 
получаемая посредством равносильностей, указанных в
\DMref{п.~3.2.2}{3.2.2.}.
\end{theorem}


\end{frame}

\begin{frame}
\frametitle{3.4.4. Эквивалентные преобразования (2/8)}
%\vspace{-2pt}
\begin{proof} 
Любую формулу (кроме тех, которые реализуют~0) можно преобразовать
в СДНФ с помощью равносильностей из \DMref{п.~3.2.2}{3.2.2.} и
правила расщепления из предыдущего примера по следующему
алгоритму.
\proofsection{1. \odmkey{Элиминация операций}{Элиминация
(elimination)!операции
 (of operation)}}
Любая булева операция реализуется формулой над базисом\\
$\{\wedge,\Or,\neg\}$ 
(например, в виде СДНФ). 
Таким образом,
любая присутствующая в формуле подформула с главной операцией,
отличной от дизъюнкции, конъюнкции и отрицания, может быть
заменена подформулой, содержащей только три базисные операции.
Например, элиминация импликации выполняется с помощью
равносильности \\ $x_1\to x_2 = \neg x_1 \Or x_2$. В
результате первого шага в формуле остаются только базисные
операции. 
\proofsection{2. \odmkey{Протаскивание отрицаний}{Протаскивание
отрицаний (dragging of negations)}} 
С помощью инволютивности
отрицания и правил де~Моргана
 отрицание
<<протаскивается>> к переменным. 
В результате второго шага
отрицания могут присутствовать в формуле только непосредственно
перед переменными. 
\end{proof}

\end{frame}

\begin{frame}
\frametitle{3.4.4. Эквивалентные преобразования (3/8)}

\begin{proof} (Продолжение)
\proofsection{3. \odmkey{Раскрытие скобок}{Раскрытие
скобок (removal of brackets)}}
 По дистрибутивности конъюнкции
относительно дизъюнкции
раскрываются все скобки, являющиеся
операндами конъюнкции. 
В результате третьего шага формула
приобретает вид \odmkey{дизъюнктивной формы}{Форма
(form)!дизъюнктивная (disjunctive)}:
$\BigOr(A_i\wedge\dots\wedge A_j)$,
где $A_k$~--- это либо переменная, либо отрицание переменной.
\proofsection{4. \odmkey{Удаление нулей}{Удаление (removal)!нулей (of
zeros)}} 
Если в слагаемое дизъюнктивной формы входят переменная и
её отрицание 
($x\And \Not x$), 
то такое слагаемое удаляется.
Если при этом формула оказывается пустой, то процесс прерывается и
считается завершённым (исходная формула реализует 0). 
\proofsection{5. \odmkey{Расщепление
переменных}{Расщепление (splitting)!переменных (variables)}} 
По
правилу расщепления в каждую конъюнкцию, которая содержит не все
переменные, добавляются недостающие. 
В результате пятого шага
формула становится <<совершенной>>, то~есть в каждой конъюнкции
содержатся все переменные. 

\end{proof}

\end{frame}

\begin{frame}
\frametitle{3.4.4. Эквивалентные преобразования (4/8)}

%\vspace{-2pt}
\begin{proof} (Окончание) 
\proofsection{6. \odmkey{Приведение подобных}{Приведение подобных (collecting like 
terms)}} 
С помощью идемпотентности конъюнкции удаляются повторные
вхождения переменных в каждую конъюнкцию, а затем с помощью
идемпотентности дизъюнкции удаляются повторные вхождения
одинаковых конъюнкций в дизъюнкцию. 
В результате шестого шага
формула не содержит <<лишних>> переменных и <<лишних>>
конъюнктивных слагаемых.
\proofsection{7. \odmkey{Сортировка}{Сортировка
(sorting)!формулы (formula)}} 
С помощью коммутативности переменные
в каждой конъюнкции, а затем конъюнкции в дизъюнкции сортируются в
установленном порядке.
 (см.~пп.~\DMref{3.1.1}{3.1.1.} и~\DMref{3.4.3}{3.4.3.}). В результате
седьмого шага формула приобретает вид СДНФ.

\vspace{0.25\baselineskip}
Заметим, что все указанные преобразования обратимы.
Преобразуем ${\cal F_1}\mapsto {\cal S_1}$ и \\ ${\cal F_2}\mapsto {\cal S_2}$  этим
алгоритмом. Если ${\cal S_1}\ne {\cal S_2}$, значит, формулы 
\em{не равносильны} и эквивалентное преобразование одной в другую
невозможно. Если же ${\cal S_1}={\cal S_2}$, то, 
объединяя последовательность преобразований ${\cal F_1}\mapsto {\cal S_1}$ 
 и обратных преобразований   ${\cal S_2}\mapsto{\cal F_2}$,
имеем искомую последовательность.
\end{proof}
\end{frame}

\begin{frame}
\frametitle{3.4.4. Эквивалентные преобразования (5/8)}
\begin{example}
Покажем, что формулы
$\Not((y \Or \Not z)\rightarrow(\Not y \And \Not z) )\rightarrow x$ и \\
$(x\Or y)\rightarrow((\Not x\rightarrow y)  \And 
\Not(\Not x \And y) )$ равносильны. 
Для этого приведем первую формулу к СДНФ, используя шаги 1--7 из доказательства теоремы.

1. Элиминация импликации: $\Not\Not(\Not(y\Or\Not z)\Or(\Not y\And \Not z) )\Or x$

2. Протаскивание отрицаний: $((\Not y \And z)\Or(\Not y \And\Not z) )\Or x$

3. Раскрытие скобок: $(\Not y \And z)\Or(\Not y \And \Not z)\Or x$

4. Удаления нулей: нет

5. Расщепление переменных: 
$(\Not y \And z \And x)\Or(\Not y \And z \And \Not x)
\Or(\Not y \And \Not z \And x)\Or $ \\$\Or(\Not y \And\Not z \And \Not x)
\Or(x\And \Not y\And\Not z)\Or(x\And\Not y \And z)\Or(x\And y\And \Not z)
\Or(x\And y \And z)$

6. Приведение подобных: $(\Not y \And z \And x)\Or(\Not y \And z \And \Not x)
\Or(\Not y\And\Not z \And x)\Or$ \\
$\Or(\Not y\And\Not z\And \Not x)
\Or(x\And y\And\Not z)\Or(x\And y\And z)$

7. Сортировка: $(\Not x\And\Not y\And\Not z)
\Or(\Not x\And\Not y\And z)\Or
(x\And\Not y\And\Not z )\Or
(x\And\Not y\And z)\Or$\\
$\Or(x\And y \And\Not z)\Or(x\And y\And z)$
\end{example}
\end{frame}

\begin{frame}
\frametitle{3.4.4. Эквивалентные преобразования (6/8)}
\begin{example}
Аналогично обрабатывается вторая формула\\ 
$(x\Or y)\rightarrow((\Not x\rightarrow y)  \And 
\Not(\Not x \And y) )$. 

1. Элиминация импликации: $\Not(x\Or y)\Or
((\Not \Not x\Or y)  \And \Not(\Not x \And y) )$

2. Протаскивание отрицаний: $(\Not x\And\Not y)\Or((x\Or y)  \And(x \Or\Not y) )$

3. Раскрытие скобок: $(\Not x\And\Not y)
\Or(x\And x)  \Or(x\And\Not y)\Or(y\And x)\Or(y\And\Not y)$

4. Удаления нулей: $(\Not x\And\Not y)\Or(x\And x)  \Or(x\And\Not y)\Or(y\And x)$

5. Расщепление переменных: 
$(\Not x\And\Not y\And\Not z)
\Or(\Not x\And\Not y\And z)\Or
(x\And x\And \Not y\And \Not z)
\Or$ \\ $\Or (x\And x\And \Not y\And z)  
\Or(x\And x\And y\And\Not z)  
\Or(x\And x\And y\And z)  
\Or(x\And\Not y\And\Not z)
\Or$ \\ $\Or(x\And\Not y\And z)
\Or(y\And x\And\Not z)
\Or(y\And x\And z)$

6. Приведение подобных: 
$(\Not x\And\Not y\And\Not z)
\Or(\Not x\And\Not y\And z)
\Or(x\And\Not y\And\Not z)
\Or$ \\ $\Or(x\And\Not y\And z)
\Or(x\And y\And\Not z)
\Or(x\And y\And z)$  

7. Сортировка: нет.

Таким образом, формулы равносильны.
\end{example}
\end{frame}

\begin{frame}
\frametitle{3.4.4. Эквивалентные преобразования (7/8)}
Приведённый пример показывает, что, хотя построение последовательности эквивалентных преобразований за счёт приведения формул к нормальной форме всегда возможно, это не всегда рационально.

\begin{example}
В условиях предыдущего примера для формул
$\Not((y \Or \Not z)\rightarrow(\Not y \And \Not z) )\rightarrow x$\\ и 
$(x\Or y)\rightarrow((\Not x\rightarrow y)  \And 
\Not(\Not x \And y) )$
можно провести следующие эквивалентные преобразования:

$\Not((y \Or \Not z)\rightarrow(\Not y \And\Not z) )\rightarrow x =$ 
$\Not\Not(\Not(y\Or\Not z)\Or(\Not y\And\Not z) )\Or x =$ \\
$=((\Not y \And z)\Or(\Not y\And\Not z) )\Or x 
=(\Not y \Or( z\And\Not z) )\Or x = \Not y \Or x$. 

С другой стороны,
$(x\Or y)\rightarrow((\Not x\rightarrow y )\And\Not(\Not x\And y) )=$ \\
$=\Not(x\Or y)\Or((\Not\Not x\Or y)\And\Not(\Not x\And y) )=$ \\
$=(\Not x\And\Not y)\Or(x\And(y\Or\Not y) )=$
$(\Not x\And\Not y)\Or x=$ \\
$=(\Not x\Or x)\And(\Not y\Or x)=$ $\Not y\Or x$.
\end{example}
\end{frame}

\begin{frame}
\frametitle{3.4.4. Эквивалентные преобразования (8/8)}

\begin{digress}
Эквивалентные преобразования, рассмотренные здесь, 
являются частным случаем 
\odmkey{системы подстановок термов}{Система (system)!подстановок термов (of term rewriting)}, 
или \odmkey{системы правил переписывания}{Система (system)!правил переписывания (of rewriting rules)}, рассмотренных в \DMref{п.~2.2.3}{2.2.3.}.
В общем случае в такой системе задаётся класс формул и 
набор правил преобразования этих формул. 
Правила могут быть обратимыми и всегда применимыми 
(как в случае эквивалентных преобразований булевых формул), 
но могут не иметь обратных и быть обусловленными. 
С этими системами связан целый ряд проблем теории алгоритмов. 
Например, для заданной системы правил 
и двух заданных формул определить, можно ли преобразовать одну формулу в другую и найти кратчайшую последовательность преобразований. 
Системы подстановок термов нашли множество 
практически важных применений в математической логике, 
в теории алгоритмов и в практическом программировании. 
Заметим, однако, что построение \em{эффективного} алгоритма 
отыскания \em{кратчайшей} последовательности эквивалентных преобразований 
является непростой задачей даже для булевых формул, 
не говоря уже о других классах математических объектов 
и других системах переписывания. 
\end{digress} 

\end{frame}

\label{3.4.5.}\subsubsection{3.4.5. Минимальные дизъюнктивные формы}
\begin{frame}
\frametitle{3.4.5. Минимальные дизъюнктивные формы (1/3)}
%
%\subsection{Минимальные дизъюнктивные формы}

Булева функция может быть задана бесконечным числом различных, но
равносильных формул. Возникает естественная задача:
для данной булевой функции найти реализующую формулу,
обладающую теми или иными свойствами.

Практически наиболее востребованной оказалась задача минимизации:
найти минимальную формулу, реализующую функцию. 
Но и в этой постановке
задача имеет множество вариантов:
\begin{enumerate}
\item[1)] найти реализующую формулу, содержащую наименьшее количество
переменных;
\item[2)] найти реализующую формулу, содержащую наименьшее количество
опре\-де\-лённых операций;
\item[3)] найти реализующую формулу, содержащую наименьшее количество
подформул определённого вида.
\end{enumerate}


\end{frame}

\begin{frame}
\frametitle{ 3.4.5. Минимальные дизъюнктивные формы (2/3)}
Кроме того, могут быть наложены ограничения на синтаксический вид
искомой формулы, набор операций, которые разрешается использовать
в формуле, и~т.\;д. Из всего этого разнообразия наиболее детально,
по-видимому, изучена задача отыскания дизъюнктивных форм,
минимальных по числу вхождений переменных. 
Эту задачу мы бегло
рассматриваем в заключительных параграфах данного раздела.

\smallskip
\odmkey{Дизъюнктивной формой}{Форма (form)!дизъюнктивная (disjunctive)}
называется формула вида
$$
\BigOr_{i=1}^{k} K_i, \quad \text{где }
 K_i = \bigwedge_{p=1}^{m_i} x_{j_p}^{\sigma_p}.
$$


Формула $K_i$ называется элементарной конъюнкцией, переменные в
элементарную конъюнкцию входят в установленном порядке (хотя и не
обязательно все $n$). Количество переменных в конъюнкции
называется её \odmkey{рангом}{Ранг (rank)!конъюнкции (of
conjunction)} (обозначение $|K_i|$).
\end{frame}

\begin{frame}
\frametitle{ 3.4.5. Минимальные дизъюнктивные формы (3/3)}

\begin{theorem}
Число различных дизъюнктивных форм $n$ переменных равно $2^{3^n}$.
\end{theorem}
\begin{proof}
Переменная либо не входит в элементарную конъюнкцию,
либо входит с отрицанием,
либо входит без отрицания.
Таким образом, существует $3^n$ элементарных конъюнкций,
а каждое подмножество множества элементарных конъюнкций даёт
дизъюнктивную форму.
\end{proof}

\medskip
Из этой теоремы можно извлечь два практических вывода, важных для
рассматриваемой задачи минимизации дизъюнктивной формы.
Во-первых,
существует тривиальный алгоритм решения задачи: перебрать все формы
(их конечное число) 
и выбрать минимальную.
Во-вторых,
количество дизъюнктивных форм с ростом $n$ растёт \em{очень быстро},
и уже при небольших $n$ полный перебор практически неосуществим.

\end{frame}

\begin{frame}
\label{3.4.6.}\subsubsection{3.4.6. Геометрическая интерпретация }\frametitle{3.4.6. Геометрическая интерпретация (1/2)}
%\subsection{Геометрическая интерпретация}
При обсуждении задач, связанных с булевыми функциями, очень
полезна сле\-дую\-щая геометрическая интерпретация. 
%\vspace{-4pt}
Двоичным наборам
из $n$ элементов можно \\ взаимно-однозначно сопоставить вершины
$n$-мерного единичного гиперкуба~$E^n$.
При этом булева функция $f(\Tuple{x}{1}{n})$
задаётся подмножеством $N$ вершин гиперкуба,
на которых её значение равно единице:
 $N\Def\Set{(\Tuple{a}{1}{n})}{f(\Tuple{a}{1}{n})=1}$.


\begin{columns}[t] % contents are top vertically aligned
\begin{column}[T]{5.5cm}
\begin{example}
В этом и последующих примерах рассматривается булева функция
$g(x,y,z)$, заданная таблицей истинности, 
показанной в середине.
На рисунке справа выделены вершины куба, соответствующие функции
$g(x,y,z)$.
Для функции $g$ имеем $N=\{(0,0,0),(0,0,1),$ $(0,1,0),(1,1,0)\}$.
\end{example}
\end{column}
\begin{column}[T]{3.5cm} % each column can also be its own environment
\begin{tabular}{ccc|c}
%\hline
 $x$ &  $y$& $z$ & $g(x,y,z)$ \\
\hline
   0    &    0     &   0    &    1\\
   0    &    0     &   1    &    1\\
   0    &    1     &   0    &    1\\
   0    &    1     &   1    &    0\\
   1    &    0     &   0    &    0\\
   1    &    0     &   1    &    0\\
   1    &    1     &   0    &    1\\
   1    &    1     &   1    &    0\\
\hline
\end{tabular}
\end{column}
\begin{column}[T]{5cm}
\begin{figure}[h]
\begin{center}
%\hspace{2cm}
\includegraphics[height=40mm]{3_1.jpg}
\end{center}
%%\caption{Представление булевой функции в виде гиперкуба} \label{fig5-1}
\end{figure}
\end{column}
\end{columns}

\end{frame}

\begin{frame}
\frametitle{3.4.6. Геометрическая интерпретация (2/2)}
%\vspace{-4pt}
Элементарной конъюнкции $K=x_{i_1}^{\sigma_1} \wedge \ldots
\wedge x_{i_k}^{\sigma_k}$ соответствует множество вершин
гиперкуба, у которых $x_{i_1} = {\sigma_1}, \dots ,
x_{i_k}={\sigma_k}$, а значения остальных координат произвольны.
Другими словами, элементарной конъюнкции $K$ сопоставляется
множество вершин $\Set{(\Tuple{a}{1}{n})\in E^n}{K(\Tuple{a}{1}{n})=1}$. 
Мы
обозначаем одной и той же буквой и элементарную конъюнкцию,
и то множество вершин, на котором она принимает значение 1.
Легко видеть, что элементарная конъюнкция $k$ переменных образует
$(n-k)$-мер\-ную грань гиперкуба $E^n$.


\begin{examples}
{\color{blue} 1.} Элементарной конъюнкции
$\Not x\wedge\Not y$ соответствует ребро $((0,0,0),(0,0,1))$:
на рисунке это ребро выделено.

{\color{blue} 2.} Элементарной конъюнкции $z$
соответствует верхняя грань куба.
\end{examples}

\vspace{6pt}
Теперь ясен геометрический смысл \odmkey{задачи минимизации
дизъюнктивной формы}{Задача (problem)!минимизации дизъюнктивной
формы (minimum disjunctive form)}. Дано подмножество $N$ вершин
гиперкуба $E^n$. Требуется найти такой набор гиперграней $R_i$,
чтобы они в совокупности образовывали покрытие $N$ ($N  =
\bigcup_{i=1}^{k} K_i$)
 и сумма рангов $\sum_{i=1}^{k} |K_i|$ была бы минимальна.
\end{frame}

\begin{frame}
\label{3.4.7.}\subsubsection{3.4.7. Сокращённые дизъюнктивные формы }\frametitle{3.4.7. Сокращённые дизъюнктивные формы (1/3)}

%\subsection{Сокращённые дизъюнктивные формы}

Известно несколько различных методов решения задачи минимизации
дизъюнктивной формы. Все они используют одну и ту же идею:
уменьшить по возможности множество рассматриваемых элементарных
конъюнкций и затем найти минимальную дизъюнктивную форму,
перебирая оставшиеся конъюнкции. Рассмотрим один из самых простых
методов этого типа.

\medskip
Пусть функция $f$ задана множеством $N$
вершин единичного гиперкуба $E^n$.
Если $K\subset N$, то конъюнкция $K$ называется
\odmkey{допустимой}{Конъюнкция (conjunction)!допустимая
(admissible)} для функции $f$. Ясно, что в минимальную (да и любую
другую) дизъюнктивную форму функции $f$ могут входить только
допустимые конъюнкции. Если $K\cap(E^n\setminus N)\neq \emptyset$,
то $K$~--- недопустимая конъюнкция. Поэтому для каждого набора
$(\Tuple{a}{1}{n})$ из множества $E^n\setminus N$ следует удалить
из множества всех $3^n$ конъюнкций те $2^n$ конъюнкций, которые
можно построить из сомножителей $\{x_1^{a_1},\dots,x_n^{a_n}\}$.
\end{frame}

\begin{frame}
\frametitle{3.4.7. Сокращённые дизъюнктивные формы (2/3)}
\begin{example}
Для функции $g$ имеем

\begin{tabular}{c|rrrrrrrr}
\hline
 $g(x,y,z)=0$ &  \multicolumn{8}{c}{Недопустимые конъюнкции} \\
\hline
(0,\:1,\:1) & $1$ & $\Not x$ & $ y$ & $ z$ & $ \Not x\wedge y$ & $ \Not x\wedge z$ & $  y\wedge z$ & $ \Not x\wedge y\wedge z$\\
(1,\:0,\:0) & $1$ & $ x$ & $ \Not y$ & $ \Not z$ & $ x\wedge \Not y$ & $ x\wedge \Not z$ & $  \Not y \wedge \Not z$ & $  x\wedge \Not y\wedge \Not z$\\
(1,\:0,\:1) & $1$ & $x$ & $ \Not y$ & $ z$ & $ x\wedge \Not y$ & $ x\wedge z$ & $  \Not y\wedge z$ & $ x\wedge \Not y\wedge z$\\
(1,\:1,\:1) &
$1$ & $ x$ & $ y$ & $ z$ & $ x\wedge y$ & $ x\wedge z$ & $  y\wedge z$ & $ x\wedge y\wedge z$\\
\hline
\end{tabular}


\smallskip
Удаляя из множества всех 27 конъюнкций те, которые не являются
допустимыми, получаем следующий список допустимых конъюнкций:

\vspace{-\baselineskip}
$$
y\wedge \Not z,\ \Not x\wedge \Not y,\ \Not x\wedge \Not
z,\ x\wedge y\wedge \Not z,\ \Not x\wedge y\wedge \Not
z,\ \Not x\wedge \Not y\wedge z,\ \Not x\wedge \Not
y\wedge \Not z.
$$
\end{example}

\vspace{-4pt}
Конъюнкция $K$ называется
\odmkey{максимальной}{Конъюнкция (conjunction)!максимальная (maximal)}
для функции $f$, если

\vspace{-\baselineskip}

$$
K\subset N \And \left( K\subset K'\subset N\Impl K'=K \right).
$$


Очевидно, что всякая допустимая конъюнкция
содержится в некоторой максимальной. Поэтому
совокупность всех максимальных конъюнкций
образует покрытие множества $N$, то есть
дизъюнктивную форму, которая называется
\odmkey{сокращённой дизъюнктивной нормальной формой}
{Форма (form)!нормальная сокращённая (normal reduced)!дизъюнктивная
(disjunctive)}.

\end{frame}

\begin{frame}
\frametitle{3.4.7. Сокращённые дизъюнктивные формы (3/3)}

Сокращённая дизъюнктивная нормальная форма определена
для функции $f$ однозначно (с учётом установленного порядка переменных).

\smallskip
\begin{example}
Для функции $g$ сокращённая дизъюнктивная нормальная форма имеет
вид \\$(\Not x \wedge \Not y)\Or
(\Not x\wedge \Not z)\Or
(y\wedge \Not z)$,
что существенно короче, чем СДНФ этой функции:
$
(\Not x \wedge \Not y \wedge \Not z) \Or
(\Not x \wedge \Not y \wedge  z) \Or
(\Not x \wedge  y \wedge \Not z) \Or
(x \wedge  y \wedge \Not z)$.
\end{example}

\smallskip
\begin{theorem}
Минимальная дизъюнктивная форма является подформулой
сокращённой дизъюнктивной нормальной формы.
\end{theorem}
\begin{proof}
От противного. Пусть минимальная форма содержит не максимальную
конъюнкцию. Тогда эту конъюнкцию можно заменить соответствующей
максимальной, при этом покрытие сохранится, а сумма рангов
уменьшится, что противоречит минимальности формы.
\end{proof}

\begin{example}
Для функции $g$ минимальная дизъюнктивная форма имеет вид
$\Not x \wedge \Not y\Or  y\wedge \Not z$.
\end{example}

\begin{remark}
Хотя сокращённая дизъюнктивная нормальная форма 
единственна, 
минимальная дизъюнктивная форма не 
обязательно единственна, 
а потому её не следует называть нормальной.
Например,
$\Not x \wedge y \Or 
x \wedge \Not y \Or \Not x \wedge z =
\Not x \wedge y \Or 
x \wedge \Not y \Or \Not y \wedge z$ суть
две различных минимальных формы одной функции.
\end{remark}

\end{frame}

\begin{frame}
\frametitle{Онтология булевых функций}
\begin{center}
\includegraphics[height=7.8cm]{ODM3_1.png}
\end{center}
\end{frame}


\begin{frame}
\frametitle{3.5. Представление булевых функций в~программах}

\subsection{3.5. Представление булевых функций в~программах}
\label{3.5.}

Для представления булевых функций можно
использовать стандартные методы представления функций,
а также некоторые специальные приёмы, и эти идеи часто
оказываются применимыми для представления других,
более сложных объектов.

\medskip
\begin{tabular}{p{12mm}|p{65mm}|p{65mm}}
  № & Название параграфа
  & Ключевые термины и обозначения \\
  \hline
\DMref{3.5.1.}{3.5.1.}\RS & Табличные представления & {Алгоритм вычисления номера кортежа, схема Горнера}
\\  
\DMref{3.5.2.}{3.5.2.} & Строковые представления & {Алгоритм построения СДНФ}
\\  
\DMref{3.5.3.}{3.5.3.} & Алгоритм вычисления значения булевой функции & {}\\
\DMref{3.5.4.}{3.5.4.}\BS & Представление булевых функций арифметическими полиномами & {Буквальные замены, ортогонализация}
\\  
\DMref{3.5.5.}{3.5.5.} & Карты Карно и матрицы минтермов & {Диаграмма Вейча}
\\  
\DMref{3.5.6.}{3.5.6.} & Булевы функции  и карты Карно & {} \\
\DMref{3.5.7.}{3.5.7.} & Минимизация формул картами Карно & {Максимальная и избыточная компактная группа, тупиковое покрытие}
\\
\DMref{3.5.8.}{3.5.8.}\RS & Деревья решений & {Семантическое дерево}
\end{tabular}  

\end{frame}

\begin{frame}
\label{3.5.1.}\subsubsection{3.5.1. Табличные представления }\frametitle{3.5.1. Табличные представления (1/3)}

%\subsection{Табличные представления}

Область определения и область значений у булевой функции конечна,
поэтому булева функция может быть представлена массивом.
Самое бесхитростное представление прямо воспроизводит
таблицу истинности. Если условиться, что кортежи в таблице
всегда идут в установленном порядке, то для представления
булевой функции $f(\Tuple{x}{1}{n})$ достаточно хранить столбец значений:
\textbf{$F_1:$ array $[0..(2^n-1)]$ of $0..1$}

\begin{example}
Для булевой функции
$g(x,y,z)$ имеем \\
\textbf{$F_1(g):$ array $[0..7]$ of $0..1 \Assign$} $(1,1,1,0,0,0,1,0)$
\end{example}
%\end{column}
%
%\begin{column}[T]{5cm}
%
%\vspace{-16pt}
%\begin{center}
%\begin{tabular}{ccc|r}
%%\hline
% $x$ &  $y$& $z$ & $g(x,y,z)$ \\
%\hline
%   0    &    0     &   0    &    1\\
%   0    &    0     &   1    &    1\\
%   0    &    1     &   0    &    1\\
%   0    &    1     &   1    &    0\\
%   1    &    0     &   0    &    0\\
%   1    &    0     &   1    &    0\\
%   1    &    1     &   0    &    1\\
%   1    &    1     &   1    &    0\\
%\hline
%\end{tabular}
%\end{center}
%
%\end{column}
%\end{columns}

\medskip
Данное представление~--- не самое эффективное. В частности, если
набор аргументов задан массивом \textbf{$x:$ array $[1..n]$ of
$0..1$}, то для вычисления значения функции~$f$ с помощью
представления $F_1$ необходимо сначала вычислить индекс~$d$ (то
есть номер кортежа в установленном порядке), чтобы затем получить
значение $F_1[d]$. Индекс $d$ нетрудно вычислить, например, с
помощью следующего 
\odmkey{алгоритма вычисления номера кортежа}{Алгоритм (algorithm)!вычисления номера кортежа в установленном порядке (evaluation of tuple’s number in established order) }.



\end{frame}

\begin{frame}
\frametitle{3.5.1. Табличные представления (2/3)}


Алгоритм вычисления номера кортежа в
установленном порядке.

\begin{algorithmic}
\REQUIRE
кортеж \textbf{$x:$ array $[1..n]$ of $0..1$} значений
переменных.
\ENSURE номер $d$ кортежа $x$ при перечислении кортежей
в установленном порядке.
\STATE $d\Assign 0$
\COMMENT{\color{gray}\small начальное значение индекса }
\FOR{\textbf{$i$ from $1$ to $n$}}
  \STATE $d\Assign d * 2 + x[i]$
  \COMMENT {\color{gray}\small сдвигаем влево и добавляем разряд }
\ENDFOR
\end{algorithmic}

\begin{ground}
Если рассматривать кортеж булевых значений
как двоичную запись числа,
то это число~--- номер кортежа в установленном порядке.
Значение числа $d$, заданного  в
позиционной двоичной системе счисления цифрами $x_1\dots x_n$,
определяется
следующим образом:

\vspace{-16pt}
$$
d =x_1 2^{n-1}+\ldots + x_n 2^0 = \sum_{i=1}^{n} x_i 2^{n-i} =
2(2(\dots(2x_1+x_2)\dots)+x_{n-1})+x_n.
$$

Цикл в алгоритме непосредственно вычисляет последнюю формулу,
которая является частным случаем
\odmkey{схемы Горнера}{Схема Горнера (Horner's method)}.
\end{ground}

\end{frame}

\begin{frame}
\frametitle{3.5.1. Табличные представления (3/3)}

\begin{remark}
По схеме Горнера можно определить значение
числа $d$ по записи $x_1\dots x_n$
в \em{любой} позиционной системе счисления
с основанием $b$ следующим образом: 
$$d=b(b(\dots(bx_1+x_2)\dots)+x_{n-1})+x_n.$$
\end{remark}

\medskip
Более эффективным представлением таблицы истинности является использование\\
$n$-мерного массива:
\textbf{$F_2:$ array $[0..1,\dots,0..1]$ of $0..1$}

В случае использования представления $F_2$ значение функции
$f(\Tuple{x}{1}{n})$ задаётся выражением $F_2[\Tuple{x}{1}{n}]$.

\begin{example}
Для функции $g$ имеем\\
$
\textbf{$F_2:$ array $[0..1,0..1,0..1]$ of $0..1= (((1,1),(1,0)),((0,0),(1,0)))$}
$. 
\end{example}

\smallskip
\begin{remark}
Мы принимаем, что многомерные массивы располагаются в линейной памяти 
<<по строкам>>.
\end{remark}

%\addvspace{-3pt}
\smallskip
Представления булевой функции $n$ переменных в виде массива занимают память объёмом
$O(2^n)$.

\end{frame}

\begin{frame}
\label{3.5.2.}\subsubsection{3.5.2. Строковые представления }
\frametitle{3.5.2. Строковые представления (1/3)}

Булеву функцию можно представить с помощью реализующей её формулы. %\linebreak
Существует множество способов представления формул, но
наиболее удобной для человека является запись в виде цепочки символов
на некотором формальном языке.

Как уже указывалось, для любой булевой функции существует бесконечно много
различных реализующих её формул. Однако
булевы функции имеют нормальные формы, в частности СДНФ,
и при установленном порядке переменных СДНФ единственна.


СДНФ булевой функции может быть построена по заданной таблице
истинности с помощью следующего \odmkey{алгоритма построения СДНФ}{Алгоритм
(algorithm)!построения СДНФ (implementation of perfect DNF)}

%\par
%%\addvspace{-3pt}
%\begin{algorithm}
%\caption{Построение СДНФ}
%\addvspace{-4.3pt}%
%\end{algorithm}%

\begin{algorithmic}
\REQUIRE
вектор \textbf{$x:$ array $[1..n]$ of string} идентификаторов
переменных,\\
вектор \textbf{$F_1:$ array $[0..(2^n-1)]$ of $0..1$}
значений функции при установленном порядке кортежей.
\ENSURE последовательность символов, образующих
запись формулы СДНФ для заданной функции.

\end{algorithmic}
\end{frame}

\begin{frame}
\frametitle{3.5.2. Строковые представления (2/3)}

\begin{algorithmic}
\STATE $f\Assign\textbf{false}$
\COMMENT{\color{gray}\small признак присутствия левого операнда дизъюнкции }
\FOR{\textbf{$i$ from $0$ to $2^n-1$}}
  \IF{$F_1[i]=1$}
     \STATE \textbf{if $f$ then  yield $'\vee'$ else $f\Assign$ true end if}
     \COMMENT{\color{gray}\small  знак дизъюнкции }
     \STATE \textbf{$g\Assign$ false}
    \COMMENT{\color{gray}\small  признак присутствия левого операнда конъюнкции }
    \FOR{\textbf{$j$ from $1$ to $n$}}
     \STATE \textbf{if $\!g\!$ then 
      yield $\!'\wedge'\!$ else $\!g\Assign\!$ true end if}
     \COMMENT{\color{gray}\small знак конъюнкции }
     \STATE \textbf{ $v\Assign (i$ div $2^{n-j})$ mod $2$ }
      \COMMENT{\color{gray}\small $j$-й разряд $i$-го кортежа }
      \IF{$v=0$}
        \STATE \textbf{yield $'\neg'$}
        \COMMENT{\color{gray}\small  знак отрицания }
      \ENDIF
      \STATE \textbf{yield $x[j]$}
      \COMMENT{\color{gray}\small  добавление в формулу переменной }
    \ENDFOR
  \ENDIF
\ENDFOR
\end{algorithmic}


\end{frame}

%\end{algorithm}

%\addvspace{-10pt}%
%{\unitlength=1mm \line(1,0){130} }%
%\par%
%\addvspace{6pt}%
%
%}%baselinestretch
\begin{frame}
\frametitle{3.5.2. Строковые представления (3/3)}


\begin{ground}
Данный алгоритм буквально воспроизводит словесную запись
следующего правила:
для каждой строки таблицы истинности,
для которой значение функции равно 1,
построить дизъюнктивное слагаемое, включающее все переменные,
причём те переменные, которые имеют значение 0 в этой строке, входят
со знаком отрицания. Остальное в этом алгоритме~--- мелкие
программистские <<хитрости>>, которые полезно один раз посмотреть, но
не стоит обсуждать.
\end{ground}
%\vspace{-1\baselineskip}

\begin{example}
Для функции $g$, используемой  в примерах данного раздела,
алгоритм построит строку
\vspace{-6pt}
$$
\Not x \wedge \Not y \wedge \Not z \Or \Not x \wedge \Not y \wedge
z \Or \Not x \wedge  y \wedge \Not z \Or x \wedge  y \wedge \Not
z.
$$
\end{example}

\medskip
Представление функции в виде формулы, выраженной как строка символов,
совершенно необходимо при реализации интерфейса пользователя
в системах компью\-терной алгебры,
но крайне неудобно при выполнении других манипуляций с формулами.
В следующем параграфе рассматривается представление СДНФ, более
удобное, например, для вычисления значения функции.
\end{frame}


\begin{frame}
\label{3.5.3.}\subsubsection{3.5.3. Алгоритм вычисления значения булевой функции }
\frametitle{3.5.3. Алгоритм вычисления значения булевой функции (1/3)}


Некоторые классы формул допускают более эффективную интерпретацию по
сравнению  с алгоритмом Eval (\DMref{п.~3.2.1}{3.2.1.}).
Рассмотрим алгоритм
вычисления значения булевой
функции, заданной в виде СДНФ, для заданных значений переменных
$\Tuple{x}{1}{n}$. 
В этом алгоритме используется следующее
представление данных.

\smallskip
СДНФ задана массивом
\textbf{$F_3\!\!:$ array $[1..k,1..n]$ of $0..1$},
где строка
$F_3[i,*]$ содержит набор значений $\Tuple{\sigma}{1}{n}$, для которого
$f(\Tuple{\sigma}{1}{n}) = 1$, $i \in 1..k$.

\smallskip
\begin{example}
Для функции $g$
\textbf{$F_3:$ array $[1..4,1..3]$ of $0..1 = ((0,0,0),(0,0,1),(0,1,0),(1,1,0))$}.
\end{example}



\end{frame}


\begin{frame}
\frametitle{3.5.3. Алгоритм вычисления значения булевой функции (2/3)}

\begin{algorithmic}
\REQUIRE  массив, представляющий СДНФ: $f:$ \textbf{array} $[1..k,1..n]$ \textbf{of} $0..1$;\\
множество значений переменных $x:$ \textbf{array} $[1..n]$ \textbf{of} $0..1$.
\ENSURE
0..1~--- значение булевой функции.
\FOR{$i$ \textbf{from} $1$ \textbf{to} $k$}
  \FOR{$j$ \textbf{from} $1$ \textbf{to} $n$}
    \IF{$f[i,j] \not= x[j]$}
      \STATE \textbf{next for} $i$
      \COMMENT {\color{gray}\small   $x_j\neq\sigma_j\Impl x_j^{\sigma_j} = 0 \Impl {x_1}^{\sigma_1}
                 \wedge\ldots\wedge {x_n}^{\sigma_n} = 0$ }
    \ENDIF
  \ENDFOR
  \STATE \textbf{return} $1$
  \COMMENT {\color{gray}\small   ${x_1}^{\sigma_1}\!
             \And\!\dots\!\And\!{x_n}^{\sigma_n}\!=\!1 \Impl
             \BigOr_{(\Tuple{\sigma}{1}{n})}
             x_1^{\sigma_1}\wedge \ldots \wedge x_n^{\sigma_n} = 1$ }
\ENDFOR
\STATE \textbf{return} $0$ 
\COMMENT {\color{gray}\small   все слагаемые в дизъюнкции равны нулю }
\end{algorithmic}
%\end{algorithm}




\end{frame}


\begin{frame}
\frametitle{3.5.3. Алгоритм вычисления значения булевой функции (3/3)}
\begin{ground}
Алгоритм основан на следующих правилах. Можно
прекратить вычисление конъюнкции, как только получен конъюнктивный
сомножитель, равный 0 (вся конъюнкция имеет значение 0). Можно
прекратить вычисление дизъюнкции, как только получено
дизъюнктивное слагаемое, равное 1 (вся дизъюнкция имеет значение
1).
\end{ground}

%\vspace{-.5\baselineskip}

Этот алгоритм в худшем случае выполняет $k\cdot n$ сравнений, а в среднем~---
гораздо меньше. Таким образом, он существенно эффективнее общего алгоритма
интерпретации.

\smallskip
\begin{digress} Быстрое вычисление значения СДНФ имеет, помимо теоретического,
большое практическое значение. Например, во многих современных
программах с графическим интерфейсом для составления сложных
логических условий используется наглядный бланк в~виде таблицы: в
клетках записываются условия, причём клетки одного столбца
счита\-ются соединёнными конъюнкцией, а столбцы~--- дизъюнкцией,
то~есть образуют ДНФ (или наоборот, в таком случае получается
КНФ). В частности, так устроен графический интерфейс QBE
(Query-by-Example), применяемый для формулировки логических
условий при запросе к СУБД.
\end{digress}

%%%%%%%%%%%%%%%%%%%%%%%%% Изменения 2012
\end{frame}

\begin{frame}
\label{3.5.4.}\subsubsection{3.5.4. Представление булевых функций арифметическими полиномами }\frametitle{3.5.4. Представление булевых функций арифметическими полиномами (1/8)}

Существует удобное представление булевых функций, 
простое в понимании и эффективное в реализации, 
основанное на применении обычных арифметических операций 
к булевским значениям 0 и 1.

\begin{example}
Нетрудно убедиться, что $\Not x = 1 - x$.
\end{example}

\smallskip
Возникает вопрос: 
можно ли представить произвольную булеву функцию таким образом?

\smallskip
\odmkey{Арифметическим полиномом}{Полином (polynomial)!арифметический (arithmetic)} 
$P(\Tuple{x}{1}{n})$ для функции $f(\Tuple{x}{1}{n})$ 
называется полином $n$ переменных $\Tuple{x}{1}{n}$ c целыми коэффициентами такой, 
что  $$\A{\Tuple{x}{1}{n}}{P(\Tuple{x}{1}{n})=f(\Tuple{x}{1}{n})}.$$ 

Константы (булевские значения) являются своими арифметическими полиномами.
Ариф\-метический полином тождественной функции совпадает с ней, 
а арифметический полином отрицания построен в примере выше. 
\end{frame}

\begin{frame}

\frametitle{3.5.4. Представление булевых функций арифметическими полиномами (2/8)}


\begin{example}
Рассмотрим булеву функцию, реализуемую
формулой $(x \Or y)\And z$, и полином $xz + yz-xyz$. 
Прямое вычисление, проведенное в таблице, показывает, 
что этот полином является реализующим для данной функции.


\begin{tabular}{ccc|c|c}
$x$ & $y$ & $z$ & $(x \Or y)\And z$ & $xz + yz-xyz$ \\
\hline
0 & 0 & 0 & 0 & 0 \\
0 & 0 & 1 & 0 & 0 \\
0 & 1 & 0 & 0 & 0 \\
0 & 1 & 1 & 1 & 1 \\
1 & 0 & 0 & 0 & 0 \\
1 & 0 & 1 & 1 & 1 \\
1 & 1 & 0 & 0 & 0 \\
1 & 1 & 1 & 1 & 1 \\
\hline
\end{tabular}
\end{example}

Арифметические полиномы для всех булевых функций двух существенных переменных 
также построены и приведены в таблице
на следующем слайде.


\end{frame}

\begin{frame}

\frametitle{3.5.4. Представление булевых функций арифметическими полиномами (3/8)}

\begin{tabular}{l|c|c|c|c|c|l}
\multicolumn{2}{r|}{Переменная $x$}          & $0$ & $0$ & $1$ & $1$ & \\
%\hline
 \multicolumn{2}{r|}{Переменная $y$}         & $0$ & $1$ & $0$ & $1$ &   \\
\hline
Название             & Формула           &   &   &   &   & Полином   \\
\hline
%Нуль                &           $0$           & $0$ & $0$ & $0$ & $0$ & $x$, $y$   \\
Конъюнкция           &$x\wedge y$          & $0$ & $0$ & $0$ & $1$ &  $xy$     \\
                     & $x\wedge \Not y$    & $0$ & $0$ & $1$ & $0$ & $x-xy$        \\
%                     &                  & $0$ & $0$ & $1$ & $1$ & $y$      \\
                     &$\Not x\wedge y$     & $0$ & $1$ & $0$ & $0$ & $y-xy$       \\
%                     &                       & $0$ & $1$ & $0$ & $1$ & $x$     \\
Сложение по модулю 2 &$x +_2 y$         & $0$ & $1$ & $1$ & $0$ & $x+y-2xy$      \\
Дизъюнкция           & $x\vee y$          & $0$ & $1$ & $1$ & $1$ & $x+y-xy$       \\
Стрелка Пирса        & $x \downarrow y$  & $1$ & $0$ & $0$ & $0$ &  $1-x-y+xy$     \\
Эквивалентность      &    $x\equiv y$       & $1$ & $0$ & $0$ & $1$ & $1-x-y+2xy$      \\
%                    &                       & $1$ & $0$ & $1$ & $0$ & $x$      \\
                     & $x\vee \Not y$     & $1$ & $0$ & $1$ & $1$ &  $1-y+xy$      \\
%                    &                       & $1$ & $1$ & $0$ & $0$ & $y$       \\
Импликация           & $x \rightarrow y$ & $1$ & $1$ & $0$ & $1$ & $1-x+xy$       \\
Штрих Шеффера        &     $x\mid y$     & $1$ & $1$ & $1$ & $0$ & $1-xy$       \\
%Единица            &         $1$             & $1$ & $1$ & $1$ & $1$ &  $x$, $y$   \\
\hline
\end{tabular}

\end{frame}

\begin{frame}
\frametitle{3.5.4. Представление булевых функций арифметическими полиномами (4/8)}

\vspace{-2pt}
\begin{theorem}
Для любой булевой функции существует реализующий её арифметический полином.
\end{theorem}
\begin{proof}
Любую булеву функцию можно выразить через дизъюнкцию,
конъюнкцию и отрицание, см. \DMref{п.~3.4.3}{3.4.3.}.
Тогда полином можно построить, подставляя 
арифметические полиномы для каждой 
операции формулы, начиная с внешней,
а затем
упростить получившийся полином, 
раскрыв скобки, вычислив степени и приведя подобные.
\end{proof}


\begin{remark}
Поскольку $\A{n>0}{1^n=1 \And 0^n=0}$, вычисление 
степеней сводится к отбрасыванию показателей степени перед 
дальнейшими упрощениями.
\end{remark}


\begin{digress}
Представление булевых функций арифметическими полиномами может \\быть выгодно, 
так как вычисление значений линейных полиномов программно весьма просто, 
ибо сводится к суммированию коэффициентов при переменных не равных нулю.
\end{digress}

\end{frame}

\begin{frame}
\frametitle{3.5.4. Представление булевых функций арифметическими полиномами (5/8)}
\begin{example}
Представим арифметическим полиномом функцию 
$g$ (пример в \DMref{п.~3.4.6}{3.4.6.}), 
рассматриваемую в предыдущих параграфах. 
Из примера в \DMref{п.~3.4.7}{3.4.7.} известно, 
что её минимальная дизъюнктивная форма имеет вид:
$(\Not x \wedge \Not y) \vee (y \wedge \Not z)$.
Далее\\
$(1-x)(1-y) + y(1-z) -(1-x)(1-y)y(1-z) =$

$= 1-x -y +xy +y -yz -  (1-x -y +xy)(y -yz) =$ 

$= 1-x -y +xy +y -yz -  y +xy +y^2 - xy^2 + yz -xyz - y^2z +xy^2z =$ 

$= 1-x -y +xy +y -yz -  y +xy +y - xy + yz -xyz - yz +xyz =$ 

$= 1-x +xy -yz  .$ 

\end{example}

\smallskip
Рассмотрим ещё один способ реализации булевой функции, 
основанный на той же идее.
Пусть булева функция $f$ задана формулой $F$ в базисе 
$\{\wedge, \vee, \neg\}$ некоторой дизъюнктивной формой. 
Выполним в формуле следующие замены логических
операций арифметическими:
$\neg x \mapsto 1-x$; $x\wedge y\mapsto xy$; $x\vee y \mapsto x + y $.

\end{frame}

\begin{frame}
\frametitle{3.5.4. Представление булевых функций арифметическими полиномами (6/8)}

Такие замены уместно назвать \odmkey{буквальными}{Замена (replacement)!буквальная (literal)}.
Получится арифметическое выражение $A(F)$, которое реализует 
некоторую функцию $f_1$. В общем случае функции $f$
и $f_1$ различны.

\begin{example}

\begin{tabular}{cc|c|c}
$x$ & $y$ & $F = x \vee y$ & $A(F) = x+y$ \\
\hline
0 & 0 & 0 & 0 \\
0 & 1 & 1 & 1 \\
1 & 0 & 1 & 1 \\
1 & 1 & 1 & 2 \\
\hline
\end{tabular}
\end{example}

\smallskip
Однако, если формула $F$ предварительно 
\odmkey{ортогонализирована}{Ортогонализация (orthogonalization)!формулы (formula)}, 
то есть эквивалентно преобразована к такому виду,
когда все элементарные конъюнкции ортогональны,
то значения $F$ и $A(F)$ при буквальных заменах совпадают.

\begin{remark}
\odmkey{Ортогонализация}{Ортогонализация (orthogonalization)} сводится к систематическому применению
тождества\\ 
$a \vee b = a \wedge b \vee \neg a\wedge b \vee a\wedge \neg b$. 
\end{remark} 


\end{frame}

\begin{frame}
\frametitle{3.5.4. Представление булевых функций арифметическими полиномами (7/8)}


\begin{example}

\begin{tabular}{cc|c|c}
$x$ & $y$ & $F =  x \wedge y \vee \neg x\wedge y \vee x\wedge \neg y$ &
 $A(F) = xy +(1-x)y + x(1-y)$ \\
\hline
0 & 0 & 0 & 0 \\
0 & 1 & 1 & 1 \\
1 & 0 & 1 & 1 \\
1 & 1 & 1 & 1 \\
\hline
\end{tabular}
\end{example}


\smallskip
Ортогонализация усложняет исходную формулу, но позволяет упростить 
замены при построении арифметического полинома.

\smallskip
Имеется способ избежать ортогонализации и при этом 
сохранить буквальную простоту преобразования логической формулы в арифметическую.
Для этого достаточно заметить, что ортогонализация, фактически,
позволяет отбросить в арифметическом полиноме
дизъюнкции $x+y-xy$ слагаемое $xy$, поскольку оно всегда
равно нулю в ортогонализированной формуле. 
\end{frame}

\begin{frame}
\frametitle{3.5.4. Представление булевых функций арифметическими полиномами (8/8)}

Отбрасываемое 
слагаемое неотрицательное, поэтому 
$F = (A(F)>0)$, 
то есть в качестве значения функции можно
брать \em{знак} значения арифметического выражения,
полученного буквальными заменами.

\begin{example}
Пусть некоторая функция задана неортогонализированной
дизъюнктивной  формой
$ x \wedge  y \vee  y \wedge  z$.
Общий алгоритм даёт арифметический полином
$xy+yz -xyz$.
Буквальная замена даёт более простое выражение 
$xy+yz$,
знак которого совпадает со значением 
арифметического полинома и функции.  
\end{example}  
   

%Однако если формула дял функции $f$   
%При этом f=sign(f1), где
%
%Sign(f1) = {█( 1  если f1>0;@0  если f1=0.)┤			(1)
%
%То есть:
%
%Sign f  = 
%IF   x  >  0 
%THEN	1
%ELSE	0
%END IF
%
%Обоснование алгоритма:
%Способ основан на приведённом выше алгоритме представления булевых функций арифметическими полиномами, однако тогда замена логических операций арифметическими производилась в ортогонализированных формулах (п. 3.4.2. a ˅ b ~ a ˅ b&˥a ). В алгоритме с использованием sign ортогонализации можно избежать ввиду условия (1). 
%	
%Пример.
%Реализуем функцию a˅˥b˄с
%
% f= sign(a + (1-b)c) = sign(a + c-bc)
%
%Выражение существенно проще чем соответствующий f арифметический полином:
%
%a+c-bc-ac+abc 

\begin{remark}
Если исходная булева функция симметрическая, то и реализующий 
арифметический полином оказывается симметрическим.
\end{remark}

\end{frame}

\begin{frame}
\label{3.5.5.}\subsubsection{3.5.5. Карты Карно и матрицы минтермов }\frametitle{3.5.5. Карты Карно и матрицы минтермов (1/7)}
\odmkey{Карта Карно}{Карта Карно (Karnaugh map)}
(используется также термин
\odmkey{диаграмма Вейча}{Диаграмма (diagram)!Вейча (Veitch)})
 ---
это замечательное представление булевой функции, позволяющее наглядно и
эффектно описать и реализовать многие, в том числе сложные, операции с
булевыми функциями и~реализующими их формулами. В картах Карно
применяются сразу несколько понятий, рассматриваемых в этом курсе:
код Грея (\DMref{п.~1.3.4}{1.3.4.}), булевы матрицы (\DMref{п.~1.4.9}{1.4.9.}), таблицы истинности (\DMref{п.~3.1.1}{3.1.1.}), минтермы (\DMref{п.~3.4.2}{3.4.2.}),
единичный гиперкуб (\DMref{п.~3.4.6}{3.4.6.}),
табличные представления булевой функции (\DMref{п.~3.5.1}{3.5.1.}).

\medskip
Карта Карно для представления функций $n$ переменных является прямоугольной таблицей,
которая содержит $2^n$ ячеек. Положим $k\Assign n\DIV 2$. Если $n$ чётное, то таблица
содержит $2^k$ строк и $2^k$ столбцов, а если нечётное, то
 для удобства отображения обычно принимают, что
таблица содержит $2^k$ строк и $2^{k+1}$ столбцов. 
Обозначим число строк в таблице $r(n)\Assign 2^{k}$,
 число столбцов в таблице 
 $c(n)\Assign 2^{k +n \MOD 2}$.
\end{frame}


\begin{frame}
\frametitle{3.5.5. Карты Карно и матрицы минтермов (2/7)}

\vspace{-2pt}
 Поскольку всего существует $2^n$ минтермов для $n$ переменных,
  можно считать, что каждая ячейка карты Карно символизирует один 
  минтерм, а вся карта в целом символизирует  перечисление минтермов.

Вообще говоря, существует $ 2^n!$ способов сопоставить минтермы ячейкам карты Карно. 
Например, следующий 
%прямолинейный 
алгоритм заполняет матрицу минтермами в установленном порядке.
%Алгоритм 3.5. Заполнение карты Карно в порядке возрастания минтермов. 

\vspace{-2pt}
%\vspace{-\baselineskip}
\begin{algorithmic}
\REQUIRE число переменных $n$.
\ENSURE матрица минтермов \\ $M : $ \textbf{array} $[0..(r(n)-1), 0..(c(n)-1)]$  
 \textbf{of array} $ [1..n]$ \textbf{of} $0..1$  
\STATE{$m\Assign 0$} 
  \FOR{ $i$ \textbf{from} 0 \textbf{to} $r(n)-1$} 
  \FOR{ $j$ \textbf{from} 0 \textbf{to} $c(n)-1$} 
   \STATE{ $M[i, j]\Assign B(m)$}
   \COMMENT{\color{gray}\small $B(m)$~--- двоичный код числа $m$}
    \STATE{ $m\Assign m+1$}
  \ENDFOR 
  \ENDFOR
\end{algorithmic}


\end{frame}


\begin{frame}
\frametitle{3.5.5. Карты Карно и матрицы минтермов (3/7)}
\vspace{-2pt}
\begin{remark}
Полезно сопоставить этот алгоритм с алгоритмом 
перечисления подмножеств некоторого множества в 
\DMref{п.~1.3.3}{1.3.3.}.
В приведённом алгоритме не имеет значения, 
 индексировать массив начиная с $0$ или с  $1$.
Мы начинаем здесь с $0$, чтобы сохранить единство обозначений со следующим алгоритмом, где это существенно.  
\end{remark}

\medskip
Использование более сложных алгоритмов перечисления минтермов позволяет построить матрицу (таблицу) минтермов, обладающую двумя полезными свойствами.
 
{\color{blue}1)} Прямоугольную таблицу можно рассматривать как развертку тора, то есть считать верхнюю сторону смежной с нижней, а правую~--- с левой. В этом случае у \em{каждой} ячейки имеются четыре соседа. 

{\color{blue}2)} Можно расположить минтермы так, что любые два соседних отличаются ровно в одном разряде, как коды Грея и смежные вершины единичного гиперкуба. 

\medskip
Зафиксируем один из возможных способов построения матрицы минтермов с такими свойствами. 

%Алгоритм 3.6. Построение двумерного кода Грея. 
\end{frame}


\begin{frame}
\frametitle{3.5.5. Карты Карно и матрицы минтермов (4/7)}

%\vspace{-2pt}
Разобьём множество переменных на две группы, 
в одной $r(n)$ переменных, а в другой $c(n)$ переменных. 
Множество кортежей значений переменных каждой группы 
перечислим не в установленном порядке, а в порядке, 
задаваемом алгоритмом \DMref{п.~1.3.4}{1.3.4.} (код Грея). 
Эти коды отмечают строки и столбцы матрицы минтермов, 
а минтерм в ячейке получается конкатенацией кодов строки и столбца. 
%Легко видеть, что полученная матрица обладает требуемыми свойствами.

%\vspace{-2pt}
\begin{algorithmic}
\REQUIRE число переменных $n$.
\ENSURE матрица минтермов \\ $M : $ \textbf{array} $[0..(r(n)-1), 0..(c(n)-1)]$  
 \textbf{of array} $ [1..n]$ \textbf{of} $0..1$ 
  \FOR{ $i$ \textbf{from} 0 \textbf{to} $r(n)-1$} 
  \FOR{ $j$ \textbf{from} 0 \textbf{to} $c(n)-1$} 
   \STATE{$M[i, j]\Assign G(i) G(j)$}
   \COMMENT{\color{gray}\small конкатенация кодов }
  \ENDFOR 
  \ENDFOR
\end{algorithmic}

\vspace{-2pt}
Здесь $G(i)$~--- левый зеркальный код Грея 
(\DMref{п.~1.3.4}{1.3.4.}), 
а операция конкатенации кодов никак не обозначается (\DMref{п.~1.1.5}{1.1.5.}).  

\end{frame}


\begin{frame}
\frametitle{3.5.5. Карты Карно и матрицы минтермов (5/7)}
\begin{remark}
Удобнее индексировать матрицу начиная с $0$,
поскольку $G(0)$ как раз даёт нулевой код.
\end{remark}


Приведённый способ построения матрицы минтермов далеко не единственный. 
Можно варьировать приведённый алгоритм: 
разбивать переменные на две группы произвольным образом, 
а не по порядку, использовать различные коды Грея, 
а не только левый зеркальный, 
различным образом сцеплять разряды кодов строки и столбца, 
а не только конкатенировать их. Во всех случаях получаются матрицы, 
в которых соседние минтермы отличаются ровно в одном разряде.
Далее 
в этом курсе 
указанный способ заполнения 
считается фиксированным и приведены матрицы минтермов для $n=2$, $n=3$, $n=4$ и $n=5$.

\smallskip
\begin{tabular}{|c|c|c|c|}
\hline
\multicolumn{2}{|c|}{$n=2$} & \multicolumn{2}{|c|}{$x_2$}  \\
\cline{3-4}	
\multicolumn{2}{|c|}{}      & 0 & 1 \\
\hline
$x_1$ & 0 & 00 & 01 \\
\cline{2-4}
 &	1 &	10 & 11 \\
 \hline
\end{tabular}
\hspace{1cm}
\begin{tabular}{|c|c|c|c|c|c|}
\hline
\multicolumn{2}{|c|}{$n=3$} & \multicolumn{4}{|c|}{$x_2, x_3$}  \\
\cline{3-6}	
\multicolumn{2}{|c|}{}    & 00 & 01 & 11 & 10 \\
\hline
$x_1$ & 0 & 000 & 001 & 011 & 010 \\
\cline{2-6}
 &	1 & 100 & 101 & 111 & 110 \\
 \hline
\end{tabular}


%n=4	x_3,x_4
%	00	10	11	01
%x_1,x_2	00	0000	0010	0011	0001
%	10	1000	1010	1011	1001
%	11	1100	1110	1111	1101
%	01	0100	0110	0111	0101
\end{frame}

\begin{frame}
\frametitle{3.5.5. Карты Карно и матрицы минтермов (6/7)}
\vspace{6pt}
\begin{tabular}{|c|c|c|c|c|c|}
\hline
\multicolumn{2}{|c|}{$n=4$} & \multicolumn{4}{|c|}{$x_3, x_4$}  \\
\cline{3-6}	
\multicolumn{2}{|c|}{}    & 00 & 01 & 11 & 10 \\
\hline
$x_1,x_2$ & 00 & 0000 & 0001 & 0011 & 0010 \\
\cline{2-6}
 & 01 & 0100 & 0101 & 0111 & 0110 \\
\cline{2-6}	
& 11 & 1100 & 1101 & 1111 & 1110 \\
\cline{2-6}	
& 10 & 1000 & 1001 & 1011 & 1010 \\
 \hline
\end{tabular}


\vspace{6pt}

\begin{tabular}{|c|c|c|c|c|c|c|c|c|c|}
\hline
\multicolumn{2}{|c|}{$n=5$} & \multicolumn{8}{|c|}{$x_3, x_4, x_5$}  \\
\cline{3-10}	
\multicolumn{2}{|c|}{}    & 000 & 001 & 011 & 010 & 110 & 111 & 101 & 100 \\
\hline
$x_1,x_2$ & $00$ & $\underline{00000}$ & 
$\underline{00001}$ & $00011$ & $\underline{00010}$ & $00110$ & $00111$ & $00101$ & $\underline{00100}$
\\
\cline{2-10}	
 & $01^{\phantom{1}}$ & $\underline{01000}$ & $01001$ & $01011$ & $01010$ & $01110$ & $\vspace{2pt}\overline{01111}$ & $01101$ & $01100$ \\
\cline{2-10}
 & $11^{\phantom{1}}$ & $11000$ & $11001$ & $\overline{11011}$ & $11010$ & $\overline{11110}$ & $\overline{11111}$ & $\overline{11101}$ & $11100$ \\
\cline{2-10}
 & $10^{\phantom{1}}$ & $\underline{10000}$ & $10001$ & $10011$ & $10010$ & $10110$ & $\overline{10111}$ & $10101$ & $10100$ \\
 \hline
\end{tabular}


\end{frame}

\begin{frame}
\frametitle{3.5.5. Карты Карно и матрицы минтермов  (7/7)} 

Хранить в программе информацию, 
представленную в этих таблицах, 
совершенно не обязательно, 
поскольку при фиксированном способе заполнения матрицы 
легко вычислить минтерм, находящийся в клетке $M[i,j]$.
Другими словами, функция $K(i,j)$, 
доставляющая код минтерма в ячейке $M[i,j]$, 
вычисляется так:
$$K(i,j)\Assign G(i)G(j)=(B(i) +_2  (B(i)  \DIV 2 ) )(B(j)  +_2  (B(j)  \DIV 2 ) ),$$
где функция $G(i)$ доставляет $i$-й зеркальный код Грея 
(см. теорему \DMref{п.~1.3.4}{1.3.4.}), функция $B(i)$ доставляет двоичный код числа $i$, а операция конкатенации битовых шкал никак не обозначается.

\begin{remark}
Ячейка в матрице имеет четыре соседних, 
поэтому для $n=4$ 
матрица минтермов сохраняет ту же 
степень наглядности, что и гиперкуб 
для $n=3$: смежные ячейки являются соседними.
Однако уже для $n=5$ наглядность несколько утрачивается: соседние ячейки остаются смежными,
но не все смежные являются соседними.    
\end{remark}

\begin{example}
В таблице на предыдущем слайде минтерм 00000 и смежные с ними подчёркнуты, 
минтерм 11111 и смежные с ним надчёркнуты.
\end{example}


\end{frame}


\begin{frame}
\label{3.5.6.}\subsubsection{3.5.6. Булевы функции и карты Карно }\frametitle{3.5.6. Булевы функции и карты Карно (1/5)}

%\subsection{Представление булевой функции на карте Карно}
Поскольку всякая булева функция представляется дизъюнкцией минтермов 
(\DMref{п.~3.4.2}{3.4.2.}),
указание некоторого \em{подмножества} ячеек матрицы минтермов является 
\em{представлением} булевой функции. 

\smallskip
Рассмотрим булеву матрицу
$K : $ \textbf{array} $[0..(r(n)-1), 0..(c(n)-1)]$ \textbf{of} $0..1$  
 того же
размера, что и матрица минтермов. 
Каждой строке таблицы истинности булевой функции $f$
соответствует некоторый минтерм. Если значение в этой строке равно 1, то
записываем 1 в соответствующую ячейку матрицы, в противном случае
записываем 0. Фактически, элементы вектора значений таблицы истинности
булевой функции раскладываются по ячейкам булевой матрицы, а~правило
раскладывания задаётся матрицей минтермов.
Булеву матрицу $K$ назовём \em{картой Карно} булевой
функции $f$.


\begin{columns}[t]
\begin{column}[T]{12cm}
\begin{example}
Функция $g(x, y, z)$, заданная в примере 
\DMref{п.~3.5.1}{3.5.1.}, имеет
следующее представление картой Карно\\
$F_4(g) :$ \textbf{array} $[0..1, 0..3]$ \textbf{of} $0..1$ 
$\Assign ((1,1,0,1),(0,1,0,0))$.
\end{example}
\end{column}
\begin{column}[T]{2.5cm}

\vspace{-12pt}
\begin{center}
\begin{tabular}{|c|c|c|c|}
\hline
1 & 1 & 0 & 1 \\\hline
0 & 0 & 0 & 1 \\
\hline
\end{tabular}
\end{center}
\end{column}
\end{columns}
\end{frame}

\begin{frame}
\frametitle{ 3.5.6. Булевы функции и карты Карно (2/5)}
\vspace{-2pt}
Карта Карно действительно является представлением булевой функции, 
поскольку с её помощью можно выполнять основную присущую функциям операцию: 
вычислять значение функции по заданным значениям аргументов. 
Вообще говоря, если булева функция представлена картой Карно
 и задан набор значений переменных 
 $\Tuple{\sigma}{1}{n}$, то ответ получается за один шаг: 
 значение функции на заданном наборе содержится в ячейке, 
 соответствующей минтерму 
 $\left( x_1^{\sigma_1} \And \dots \And x_n^{\sigma_n} \right)$. 
%Алгоритм 3.7. Вычисление значения булевой функции по карте Карно.

%\vspace{-2pt}
\begin{algorithmic}
\REQUIRE карта Карно 
$K : $ \textbf{array} $[0..(r(n)-1), 0..(c(n)-1)]$ \textbf{of} $0..1$, 
массив значений переменных $x : $ \textbf{array} $[1..n]$ \textbf{of} $0..1$.
\ENSURE  значение булевой функции $0$ или $1$.
\STATE{$r\Assign n\DIV 2$}
\COMMENT{\color{gray}\small длина кода Грея номера строки }
\STATE{$c\Assign r + n\MOD 2$}
\COMMENT{\color{gray}\small длина кода Грея номера столбца }
\STATE{$i\Assign\GI (r, x[1..r])$}
\COMMENT{\color{gray}\small номер строки }
\STATE{$j\Assign\GI (c, x[(r+1)..n])$}
\COMMENT{\color{gray}\small номер столбца }
\STATE{ \textbf{return} $K[i, j]$}
\COMMENT{\color{gray}\small возвращаемое значение }
\end{algorithmic}

%\vspace{-2pt}

\end{frame}


\begin{frame}
\frametitle{ 3.5.6. Булевы функции и карты Карно (3/5)}

Рекуррентные поразрядные вычисления выполняются 
в функции $\GI$, которая вычисляет целое число~--- 
номер заданного кода Грея заданной длины.

%\vspace{-6pt}
\begin{algorithmic}
\REQUIRE $n$~--- длина кода Грея,
 $G : $ \textbf{array} $[1..n]$ \textbf{of} $0..1$~---
 код Грея.
% \vspace{-1pt}
\ENSURE  целое неотрицательное число~--- номер кода.
\STATE{$d \Assign 0$}
\COMMENT{\color{gray}\small вычисляемое значение числа }
\STATE{$p \Assign 0$}
\COMMENT{\color{gray}\small значение предыдущего разряда двоичного кода числа }
\FOR{ $i$ \textbf{from} $1$ \textbf{to} $n$} 	
	\STATE{$b \Assign G[i] +_2 p$}	
	\COMMENT{\color{gray}\small очередной разряд двоичного кода числа }
	\STATE{$p \Assign b$}
	\COMMENT{\color{gray}\small предыдущий разряд двоичного кода числа }
	\STATE{$d \Assign d * 2$}
	\COMMENT{\color{gray}\small сдвиг числа влево на 1 разряд }
	\IF{$ b = 1$}
       \STATE{$d \Assign d + 1$}
       \COMMENT{\color{gray}\small добавляем единицу в число }
     \ENDIF
\ENDFOR
\STATE{\textbf{return} $d$}
\end{algorithmic}


\end{frame}


\begin{frame}
\frametitle{ 3.5.6. Булевы функции и карты Карно (4/5)}
\begin{ground}
При выбранном способе заполнения матрицы минтермов код Грея номера строки 
располагается в левых $r$ разрядах, 
а код Грея номера столбца в правых $c$ разрядах. 
В цикле по следствию 2 теоремы \DMref{п.~1.3.4}{1.3.4.} 
определяется очередной разряд двоичного кода 
и по схеме Горнера вычисляется значение числа $d$ (см. \DMref{п.~3.5.1}{3.5.1.}).
\end{ground}

\medskip
Наряду с вычислением значений, 
карты Карно позволяют строить различные формульные представления 
для булевых функций и их комбинаций. 

\smallskip
Карта Карно для одной заданной функции 
очевидным образом задаёт её СДНФ~--- 
достаточно соединить знаками дизъюнкции минтермы, 
указанные ячейками, содержащими 1. 


\smallskip 
Столь же просто можно построить СДНФ для отрицания, 
дизъюнкции и конъюнкции заданных функций. 
Для получения СДНФ отрицания достаточно просто инвертировать матрицу.
\end{frame}


\begin{frame}
\frametitle{ 3.5.6. Булевы функции и карты Карно (5/5)}
 
Для получения СДНФ конъюнкции или дизъюнкции 
произвольного числа $k$ булевых функций можно поступить следующим образом. 
Взять матрицу, заполненную нулями, 
и добавить в каждую ячейку независимо 0 или 1 для каждой функции. 
При этом 
в некоторых ячейках окажется 0, в некоторых 1, 
а в некоторых какие-то числа вплоть до $k$. 
Все ненулевые ячейки образуют СДНФ дизъюнкции, 
все ячейки, равные $k$, образуют СДНФ конъюнкции.

\smallskip



 

\end{frame}

\begin{frame}
\label{3.5.7.}\subsubsection{3.5.7. Минимизация формул картами Карно }\frametitle{3.5.7. Минимизация формул картами Карно (1/8)}

Карты Карно являются наглядным средством решения задачи 
минимизации дизъюнктивных форм, 
обсуждаемой в пп. \DMref{3.4.5}{3.4.5.}--\DMref{3.4.7}{3.4.7.}.

\smallskip
Пусть булева функция представлена картой Карно. 
Если при этом в соседних ячейках оказались 1, 
то по основному свойству карт Карно соответствующие минтермы 
отличаются точно в одном разряде, 
и значит, к ним можно применить правило склеивания (\DMref{п.~3.4.4}{3.4.4.}). 
В таком случае дизъюнкцию данных минтермов можно заменить 
элементарной конъюнкцией (\DMref{п.~3.4.5}{3.4.5.}) ранга $n-1$.
 На карте Карно соответствующие ячейки можно объединить, 
 показывая, что они образуют 
 \odmkey{компактную группу}
{Группа (group)!компактная (compact)}. 
 Результат склеивания двух минтермов обычно также обозначают минтермом, 
 причём в позиции разряда, по которому проведено склеивание, 
 поставлен прочерк.
\begin{example}  
Результат склеивания минтермов (0101) и (0111) обозначается (01-1).
\end{example}

\end{frame}

\begin{frame}
\frametitle{3.5.7. Минимизация формул картами Карно (2/8)}

Операция склеивания является эквивалентным преобразованием (см.~\DMref{п.~3.4.4}{3.4.4.}),
сохраняет дизъюнктивность формы и при этом сокращает количество вхождений переменных,
поэтому операция склеивания и выделение компактных групп~--- 
основное средство минимизации дизъюнктивных форм.

\medskip
Если получились две смежные непересекающиеся компактные группы, 
образующие прямоугольник, их можно по правилу склеивания объединить 
в компактную группу, соответствующую элементарной конъюнкции ранга $n-2$ и т.д.
Компактная группа, которая не входит ни в какую другую компактную группу, 
называется \odmkey{максимальной}{Группа (group)!компактная (compact)!максимальная (maximal)}. 
Компактные группы могут включаться друг в друга, могут пересекаться, могут примыкать другу к другу, но не образовывать новой компактной группы при этом.
\end{frame}

\begin{frame}
\frametitle{3.5.7. Минимизация формул картами Карно (3/8)}

Каждая прямоугольная компактная группа из $2^k$ смежных ячеек, 
содержащих 1, соответствует элементарной конъюнкции ранга $n-k$. 
В частности, если ячейка содержащая 1 изолированная 
(нет смежных ячеек содержащих 1), 
то эта ячейка образует компактную группу, соответствующую минтерму, 
то есть элементарной конъюнкции ранга $n$. 
Таким образом, задача минимизации дизъюнктивной формы 
сводится к отысканию покрытия всех ячеек, 
содержащих 1 на карте Карно, 
минимальным количеством компактных групп. 

\begin{remark}
Элементарную конъюнкцию, соответствующую компактной группе, 
часто называют \odmkey{импликантой}{Импликанта (implicant)}, 
поскольку если $g$~--- элементарная конъюнкция, 
соответствующая некоторой компактной группе на карте Карно функции $f$, 
то имеет место импликация $g\rightarrow f$.
\end{remark}

\end{frame}

\begin{frame}
\frametitle{3.5.7. Минимизация формул картами Карно (4/8)}

\begin{example}
На рисунке показаны три компактных группы, 
которые имеются на карте Карно функции $g(x,y,z)$ из примера \DMref{п.~3.5.1}{3.5.1.}.
Компактные группы (00-) и (-10) не пересекаются, 
а компактные группы (0-0) и (-10), 
а также (00-) и (0-0), напротив, пересекаются. 
Все три компактные группы являются максимальными.
Хотя все четыре единицы на этом рисунке находятся в соседних ячейках, 
компактную группу они не образуют, поскольку форма не прямоугольная.
Минимальное покрытие на карте Карно содержит два элемента: 
(00-) и (-10), что соответствует дизъюнктивной форме 
$\Not x \And \Not y \Or y \And \Not z$.
\end{example}

\begin{figure}[h]
\begin{center}
\includegraphics[height=36mm]{3_2.jpg}
\end{center}
\end{figure}

\end{frame}



\begin{frame}
\frametitle{3.5.7. Минимизация формул картами Карно (5/8)}
\begin{example}
Рассмотрим карту Карно, приведённую на рисунке слева.
\begin{center}
\includegraphics[height=40mm]{3_3.jpg}
\end{center} 
%Рис. 3.3. Карта Карно и компактные группы различных рангов
На этой карте имеется 14 компактных групп ранга 3 (рисунок в центре)
 и еще 3 компактных группы ранга 2 (рисунок справа). 
 Других компактных групп нет.
 \end{example}
 
Определить компактные группы, покрывающие все ячейки, 
содержащие 1 на карте Карно, можно многими различными способами,
 в этом и состоит трудность задачи минимизации дизъюнктивной формы.


\end{frame}


\begin{frame}
\frametitle{3.5.7. Минимизация формул картами Карно (6/8)}
Известны различные алгоритмы нахождения минимальных дизъюнктивных форм 
с помощью карт Карно. 
Все эти алгоритмы переборные и используют методы сокращения перебора, 
основанные на фактах, устанавливаемых в 
следующих леммах.

\medskip
\begin{lemma} \NB{[1]}
Достаточно рассматривать только максимальные компактные группы.
\end{lemma}
\begin{proof} Максимальные компактные группы соответствуют максимальным конъюнкциям, 
то есть сокращённой ДНФ, далее по теореме
 \DMref{п.~3.4.7}{3.4.7.}.
\end{proof}
%\begin{example}
%Максимальные компактные группы для рассматриваемой функции 
%приведены на следующем рисунке в центре.
%\end{example}

\medskip
Если какая-то 1 на карте Карно покрывается только одной компактной группой, 
то говорят, что эта группа \em{входит в ядро} (данного покрытия).
\begin{lemma} \NB{[2]}
Всякая компактная группа, входящая в ядро сокращённой ДНФ, входит в минимальную ДНФ.
\end{lemma}  
\begin{proof}
 По лемме 1 достаточно рассматривать сокращённую ДНФ.	
 \end{proof}
%\begin{example}
% Компактные группы, показанные на предыдущем рисунке справа, 
% образуют ядро сокращённой ДНФ, показанной на следующем рисунке в центре.
%\end{example}

\end{frame}

\begin{frame}
\frametitle{3.5.7. Минимизация формул картами Карно (7/8)}
\begin{remark}
Если все максимальные компактные группы оказались в ядре, 
то сокращённая ДНФ является минимальной.
\end{remark}

\medskip
Компактную группу называют \odmkey{избыточной}{Группа (group)!компактная (compact)!избыточная (excess)} (в данном покрытии), 
если её можно удалить из покрытия без потери эквивалентности исходной СДНФ. 
Покрытие (и, соответственно, ДНФ) называется \odmkey{тупиковым}{Покрытие (covering)!тупиковое (irredundant)}, 
если оно не содержит избыточных компактных групп.

\medskip
\begin{lemma} \NB{[3]}
Минимальная ДНФ является тупиковой.
\end{lemma}
\begin{proof} В противном случае ДНФ не была бы минимальной. 		
\end{proof}

\medskip
Отсюда вытекает основная идея метода минимизации: построить сокращённую ДНФ, построить её ядро, перебрать все тупиковые ДНФ и найти минимальную. 

\end{frame}

\begin{frame}
\frametitle{3.5.7. Минимизация формул картами Карно (8/8)}

\begin{example}
Максимальные компактные группы для рассматриваемой функции 
приведены на рисунке в центре. 
Компактные группы, показанные на рисунке слева, образуют покрытие, 
но соответствующая ДНФ не является минимальной 
(она содержит $6\cdot 3=18$ литералов).
%на рис. 3.4. 
Справа приведена одна из трёх минимальных ДНФ данной функции, 
содержащая $3\cdot 2+2\cdot 3=12$ литералов.
\begin{center}
\includegraphics[height=40mm]{3_4.jpg}
\end{center} 
%Рис. 3.4. Некоторая ДНФ, сокращённая ДНФ и минимальная ДНФ
\end{example}
\end{frame}

\begin{frame}
\label{3.5.8.}\subsubsection{3.5.8. Деревья решений }\frametitle{3.5.8. Деревья решений (1/7)}

\vspace{-2pt}
%\subsection{Деревья решений}
Наиболее эффективными с точки зрения экономии памяти
и времени оказываются представления, которые не имеют
прямой связи с <<естественными>> представлениями функции
в виде графика (массива) или формулы (выражения),
но специально ориентированы на выполнение операций.


Начнём с простого наблюдения. 
Таблицу истинности булевой функции
$n$ переменных
можно представить в виде полного бинарного дерева
высоты $n+1$.
%\vspace*{-.25\baselineskip}
%

\begin{remark}
Бинарные деревья, равно как и другие виды деревьев, а также
способы их представления в программах и соответствующая
терминология, рассматриваются в главе~9.
\end{remark}
%
%\vspace*{-.25\baselineskip}
Ярусы дерева соответствуют переменным, дуги дерева соответствуют
значениям переменных, скажем, левая дуга~--- 0, а правая~--- 1.
Листья дерева на последнем ярусе хранят значение функции на кортеже,
соответствующем пути из корня в этот лист. Такое дерево называется
\odmkey{деревом решений}{Дерево (tree)!решений (decision)}
(или \odmkey{семантическим деревом}{Дерево (tree)!семантическое
(semantic)}).


\end{frame}

\begin{frame}
\frametitle{3.5.8. Деревья решений (2/7)}

Дерево решений можно сократить, если заменить корень каждого поддерева,
 все листья
которого имеют одно и то же значение, этим значением. Иногда такое сокращение
значительно уменьшает объём дерева.


\begin{example}
На рисунке слева представлено дерево решений для функции $g$, 
а справа представлено сокращённое дерево решений для функции $g$.
\end{example}

\begin{figure}[h]
\begin{center}
\includegraphics{03_02}
\end{center}
%\caption{Дерево решений и сокращённое дерево решений} \label{fig5-2}
\end{figure}
\end{frame}

\begin{frame}
\frametitle{3.5.8. Деревья решений (3/7)}

\vspace{-2pt}
\odmkey{Алгоритм вычисления значения функции}
{Алгоритм (algorithm)!вычисления значения функции по сокращённому дереву решений (evaluation of function using the decision tree)} осуществляется проходом по дереву решений.
При этом тип узла дерева 
$N= \textbf{ struct } \{ i:0..1;\ l,r:\:\uparrow N\}$.

%\begin{algorithm}
%\caption{Вычисление значения функции по сокращённому дереву
%решений}\label{decisiontree} \index{Алгоритм
%(algorithm)!вычисления значения функции по сокращённому дереву
%решений (evaluation of function on the base of decision tree)}
%\addvspace{-4.3pt}%
%\end{algorithm}%

\vspace{-2pt}
\begin{algorithmic}
\REQUIRE  указатель $T: \:\uparrow N$ на корень дерева решений,\\
массив $x:$ \textbf{array} $[1..n]$ \textbf{of} $0..1$  значений
переменных. 
\ENSURE 0..1~--- значение булевой функции. 
\STATE $i\Assign 0$
\COMMENT {\color{gray}\small номер переменной} 
\WHILE {\textbf{true}}
    \STATE{\textbf{if $T.l=$ nil $\And T.r =$ nil
    then return $T.i$ end if}}
     \STATE{} \COMMENT {\color{gray}\small листовой узел~--- возвращаем значение }
     \STATE $i\Assign i+1$
\COMMENT {\color{gray}\small следующая переменная} 
      \IF{$x[i]$}
      \STATE $T\Assign T.r$
      \COMMENT {\color{gray}\small 1~--- переход вправо  }
      \ELSE
      \STATE $T\Assign T.l$
      \COMMENT {\color{gray}\small 0~--- переход влево  }
      \ENDIF
\ENDWHILE
\end{algorithmic}
%\end{algorithm}

%\addvspace{-10pt}%
%{\unitlength=1mm \line(1,0){130} }%
%\par%
%\addvspace{6pt}%
\end{frame}

\begin{frame}
\frametitle{3.5.8. Деревья решений (4/7)}

Дерево решений можно сделать ещё компактнее, если отказаться от
древовидности связей, то есть допускать несколько дуг, входящих в
узел. В таком случае мы получаем 
\odmkey{бинарную диаграмму решений}{Диаграмма (diagram)!решений (decision)!бинарная (binary)}. 
Бинарная диаграмма решений получается из бинарного
дерева решений тремя последовательными преобразованиями:
\begin{enumerate}
\item[1)]
отождествляются листовые узлы, содержащие~0 и содержащие~1;
\item[2)]
в диаграмме выделяются изоморфные поддиаграммы и заменяются единственным
их экземпляром;
\item[3)]
исключаются узлы, обе исходящие дуги которых ведут в один узел.
\end{enumerate}

Интерпретация бинарной диаграммы решений
(вычисление значения функции) производится в точности так~же,
как и для дерева решений, то есть по указанному выше алгоритму.

На следующем слайде показан результат преобразования для функции $g(x,y,z)$.
\end{frame}

\begin{frame}
\frametitle{3.5.8. Деревья решений (5/7)}
\begin{figure}[p]
\begin{center}
\includegraphics{03_03}
\end{center}
%\caption{Построение бинарной диаграммы решений} \label{fig5-3}
\end{figure}
\end{frame}

\begin{frame}
\frametitle{3.5.8. Деревья решений (6/7)}


\begin{example}
На рисунке показаны последовательность преобразований и
окончательная диаграмма решений для функции $g$ в том случае,
когда переменные рассматриваются в следующем порядке: $y,x,z$.
\end{example} 

\vspace{-10pt}
\begin{figure}[p]
\begin{center}
\includegraphics[scale=0.85]{03_04}
\end{center}
%\caption{Бинарная диаграмма решений при другом порядке переменных} \label{fig5-4}
\end{figure}
\end{frame}


\begin{frame}
\frametitle{3.5.8. Деревья решений (7/7)}

Таким образом, результат преобразования
дерева решений в диаграмму решений
существенно зависит от того, в
каком порядке рассматриваются переменные при построении
исходного полного дерева решений.

Полученная в последнем примере диаграмма компактнее:
 вычисление значения
функции $g$ требует всего двух операций. Фактически, диаграмма 
 показывает, что функция $g$ может быть
реализована следующим условным выражением:

\begin{center}
\textbf{if $y$ then  $\Not z$ else $\Not x $ end if}.
\end{center}

\begin{digress}
Последний разобранный пример свидетельствует,
что иногда можно построить такое специальное представление функции,
которое позволяет хранить \em{меньше информации} и при этом
производить \em{вычисления быстрее},
чем это в принципе возможно при использовании
универсальных математических методов представления функций с помощью
графиков (массивов) и
формул (выражений).
Советуем читателям обдумать этот факт.
\end{digress}
\end{frame}

\begin{frame}
\frametitle{Онтология булевых функций}
\begin{center}
\includegraphics[width=12cm]{ODM3_2.png}
\end{center}
\end{frame}


\begin{frame}
\frametitle{3.6. Полные системы булевых функций}

\subsection{3.6. Полные системы булевых функций }
\label{3.6.}

В типичной современной цифровой вычислительной машине цифрами являются
0 и 1. Следовательно, команды, которые выполняет процессор, суть булевы
функции. Выше показано, что любая булева функция реализуется через
конъюнкцию, дизъюнкцию и отрицание. Следовательно, можно построить
нужный процессор, имея в распоряжении элементы, реализующие конъюнкцию,
дизъюнкцию и отрицание. Этот раздел посвящён ответу на
вопрос: существуют ли (и~если существуют, то какие) иные системы
булевых функций, обладающих тем свойством, что с их помощью можно выразить
все другие функции.

\medskip
\begin{tabular}{p{12mm}|p{65mm}|p{65mm}}
  № & Название параграфа
  & Ключевые термины и обозначения \\
  \hline
\DMref{3.6.1.}{3.6.1.} & Замкнутые классы & {Замыкание, классы монотонных, линейных, самодвойственных функций, классы функций,сохраняющих $1$ и $0$}
\\  
\DMref{3.6.2.}{3.6.2.} & Полные системы функций & {Полином Жегалкина}
\\  
\DMref{3.6.3.}{3.6.3.} & Полнота двойственной системы & {}\\
\DMref{3.6.4.}{3.6.4.} & 
Теорема Поста & {}
\end{tabular}  


\end{frame}

\begin{frame}
\label{3.6.1.}\subsubsection{3.6.1. Замкнутые классы }\frametitle{3.6.1. Замкнутые классы (1/5)}

Пусть $F = \{\Tuple{f}{1}{m}\}$, $\A{i\in 1.. m}{f_i \in P_n}$.
\odmkey{Замыканием\/}{Замыкание (closure)} $F$ (обозначается
$[F]$) называется множество всех булевых функций, реализуемых
формулами над $F$:
$$[F] \Def  \Set{f\in P_n}{f = \func {\cal F} \And \cal{F}\in\mathfrak{F}[F]}.
$$

Отметим следующие свойства 
 замыкания (см. также \DMref{2.1.2}{2.1.2.}).
\begin{enumerate}
\item  $F \subset [F]$.%
\item  $[[F]]=[F]$.%
\item  $F_1 \subset F_2 \Impl [F_1] \subset [F_2]$.%
\item  $([F_1] \cup [F_2]) \subset [F_1 \cup F_2]$.
\end{enumerate}


Класс  функций $F$ называется
\odmkey{замкнутым}{Класс (class)!замкнутый (closed)}, если $[F] = F$.

Рассмотрим следующие классы функций.

Класс функций, 
\odmkey{сохраняющих 0}{Класс (class)!функций (of function)!сохраняющих 0 (0-preserving)}: 
$T_0 \Def \Set{f}{f(0, \dots, 0) = 0}.$

 
Класс функций, \odmkey{сохраняющих 1}{Класс (class)!функций (of function)!сохраняющих 1 (1-preserving)}: 
$T_1 \Def  \Set{f}{f(1, \dots, 1) = 1}.$

Класс \odmkey{самодвойственных}{Класс (class)!функций (of function)!самодвойственных (self-dual)} функций: 
$T_* \Def  \Set{f}{f= f^*}.$
 



%\end{enumerate}

\end{frame}

\begin{frame}
\frametitle{3.6.1. Замкнутые классы (2/5)}
Класс \odmkey{монотонных}{Класс (class)!функций (of function)!монотонных (monotonic)} функций:
$T_{\le} \Def \Set{f}{\alpha\le\beta\Impl f(\alpha)\le f(\beta)}$, 
где  $\alpha\! =\! (\Tuple{a}{1}{n})$, 
$\beta\! =\! (\Tuple{b}{1}{n})$, 
${a_i, b_i}\! \in\! E_2$, 
$\alpha\! \le\! \beta 
\Def \A{i}{a_i\! \le\! b_i}$. 

Класс \odmkey{линейных}{Класс (class)!функций (of function)!монотонных (monotonic)} функций: 
${T_L}\Def  \Set{f}{f = c_0 + c_1x_1+ \dots +c_nx_n}$, 
где $+$ обоз\-начает сложение по модулю 2, а знак $\And$ опущен.

\begin{examples}
{\color{blue}\bf 1.} 
Рассмотрим отрицание и введём обозначение $\varphi(x)\Assign
\overline{x}$. Имеем $\varphi\not\in T_0$, так как
$\varphi(0)=1$, $\varphi\not\in T_1$, так как $\varphi(1)=0$,
$\varphi\not\in T_{\le}$, так как $0<1$, но
$\varphi(0)>\varphi(1)$. 
С другой стороны,
$\varphi\in T_*$, так как
$\varphi^*(x)=\overline{\varphi}(\overline{x})= \Not\varphi(\Not
x)= \Not \overline{\Not x} = \overline{x}= \varphi(x)$, и~$\varphi\in T_L$, так как $\varphi(x)=x+1$.

{\color{blue}\bf 2.}  Рассмотрим
конъюнкцию и введём обозначение $\psi(x,y)\Assign x\wedge y$.
Имеем: $\psi\in T_0$, так как $0 \wedge 0 = 0$, $\psi\in T_1$, так
как $1 \wedge 1 = 1$, $\psi\in T_{\le}$, так как $\psi(1,1)=1$, и \\$\A{(a,b) \ne(1,1)}{(a,b)\le(1,1)\And \psi(a,b)=0}$. С другой
стороны, $\psi\not\in T_*$, так как \\$\psi^*(x,y)=x\Or y$, и
$\psi\not\in T_L$. Действительно, от противного, пусть
$\psi(x,y) = ax +by +c$. Тогда имеем: если $x,y=0,0$,  то
$a0+b0+c=0\Impl c=0$; если $x,y=0,1$, то $a0+b1+0=0 \Impl b=0$; 
если $x,y=1,0$, то $a1+0\cdot0+0=0\Impl a=0$; если $x,y=1,1$, 
то $0\cdot 1+0\cdot 1+ 0=1$, и,
значит, $0=1$~--- противоречие.
%\end{enumerate}
\end{examples}

\end{frame}

\begin{frame}
\frametitle{3.6.1. Замкнутые классы (3/5)}


\begin{theorem}
Классы $T_0$, $T_1$, $T_{*}$, $T_\le$, $T_L$ замкнуты.
\end{theorem}
\begin{proof}
Чтобы доказать, что некоторый класс $F$ замкнут, достаточно показать,
что если функция реализована в виде формулы над $F$, то она принадлежит $F$.
Доказать, что произвольная формула обладает заданным свойством,
можно с помощью индукции по структуре формулы (см.~\DMref{п.~3.3.3}{3.3.3.}).
База индукции очевидна: функции из $F$ реализованы как тривиальные
формулы над $F$. Таким образом, осталось обосновать
индукционные переходы для пяти рассматриваемых классов.

\proofsection{$T_0$} Пусть $f,\Tuple{f}{1}{n}\in{T_0}$ и
$\Phi=f(f_1(\Tuple{x}{1}{n}),\dots,f_n(\Tuple{x}{1}{n}))$. \\ 
Тогда
$\Phi(0,\dots,0) =f(f_1(0,\dots,0),\dots,f_n(0,\dots,0))=
f(0,\dots,0)=0.$ \\ Следовательно, ${\Phi}\in{T_0}$.
\proofsection{$T_1$} Пусть $f,\Tuple{f}{1}{n}\in{T_1}$ и
$\Phi=f(f_1(\Tuple{x}{1}{n}),\dots,f_n(\Tuple{x}{1}{n}))$. \\Тогда $\Phi(1,\dots,1) =f(f_1(1,\dots,1),\dots,f_n(1,\dots,1))=
f(1,\dots,1)=1.$ \\ Следовательно, ${\Phi}\in{T_1}$.
\end{proof}

\end{frame}

\begin{frame}
\frametitle{3.6.1. Замкнутые классы (4/5)}

\begin{proof} (Продолжение)
\proofsection{$T_*$} Пусть $f,\Tuple{f}{1}{n}\in{T_*}$ и
$\Phi=f(f_1(\Tuple{x}{1}{n}),\dots,f_n(\Tuple{x}{1}{n}))$. Тогда \\
$\Phi^*\!=\! f^*(f_1^*(\Tuple{x}{1}{n}),\dots, f_n^*(\Tuple{x}{1}{n}))\!=\!
f(f_1(\Tuple{x}{1}{n}),\dots, f_n(\Tuple{x}{1}{n}))\!=\!\Phi$. \\
Следовательно,
$\Phi\in T_*$.
\proofsection{$T_{\le}$}
 Пусть $f,\Tuple{f}{1}{n}\!\in\!{T_\le}$ и
$\Phi=f(f_1(\Tuple{x}{1}{n}),\dots,f_n(\Tuple{x}{1}{n}))$.
\\ Тогда 
$\alpha\le\beta\Impl (f_1(\alpha),\dots,f_n(\alpha))\le
(f_1(\beta),\dots,f_n(\beta))\Impl$\\
$\Impl f(f_1(\alpha),\dots,f_n(\alpha))\le
f(f_1(\beta),\dots,f_n(\beta))\Impl \Phi(\alpha)\le \Phi(\beta)$. \\
Следовательно, ${\Phi}\in{T_\le}$.

\proofsection{$T_L$} Пусть 
$f,\Tuple{f}{1}{n} \in {T_L}$ и
$\Phi=f(f_1(\Tuple{x}{1}{n}),\dots,f_n(\Tuple{x}{1}{n}))$. \\Тогда 
$f = c_0+c_1x_1+\ldots +c_nx_n$,
$f_1 = c_0^1+c_1^1x_1+\ldots +c_n^1x_n$,
$\dots$, \\
$f_n = c_0^n+c_1^nx_1+\ldots+c_n^nx_n$.
Подставим эти формулы 
в формулу для $\Phi$. Имеем\\
$
\Phi(\Tuple{x}{1}{n})= c_0+c_1\left(c_0^1+c_1^1x_1+\ldots
+c_n^1x_n\right)+\ldots
 +c_n\left(c_0^n+c_1^n x_1+\ldots +c_n^n x_n\right)=$
\\$= d_0+d_1x_1+\ldots +d_nx_n$.
Следовательно, ${\Phi}\in{T_L}$.
\end{proof}

\end{frame}

\begin{frame}
\frametitle{3.6.1. Замкнутые классы (5/5)}

\begin{columns}[t] % contents are top vertically aligned
\begin{column}[T]{7cm} % each column can also be its own environment
\begin{tabular}{|c|c|c|c|c|c|}
\hline
                     & $T_0$ & $T_1$ & $T_*$ & $T_\le$ & $T_L$   \\
\hline
      $0$            &  $+$  &  $-$  &  $-$  &   $+$   &   $+$    \\
$x\And y$            &  $+$  &  $+$  &  $-$  &   $+$   &   $-$    \\
$x$                  &  $+$  &  $+$  &  $+$  &   $+$   &   $+$    \\
$y$                  &  $+$  &  $+$  &  $+$  &   $+$   &   $+$    \\
$x +_2 y$            &  $+$  &  $-$  &  $-$  &   $-$   &   $+$    \\
$x\Or y$             &  $+$  &  $+$  &  $-$  &   $+$   &   $-$    \\
$x\downarrow y$      &  $-$  &  $-$  &  $-$  &   $-$   &   $-$    \\
$x\equiv y$          &  $-$  &  $+$  &  $-$  &   $-$   &   $+$    \\
$\Not y$             &  $-$  &  $-$  &  $+$  &   $-$   &   $+$    \\
$\Not x$             &  $-$  &  $-$  &  $+$  &   $-$   &   $+$    \\
$x\rightarrow y$     &  $-$  &  $+$  &  $-$  &   $-$   &   $-$    \\
$x \mid y$           &  $-$  &  $-$  &  $-$  &   $-$   &   $-$    \\
      $1$            &  $-$  &  $+$  &  $-$  &   $+$   &   $+$    \\
\hline
\end{tabular}
\end{column}
\begin{column}[T]{7cm}
\begin{example}
Слева представлена таблица принадлежности булевых функций двух переменных рассмотренным выше замкнутым
классам. Плюс указывает, что функция принадлежит классу, минус указывает,
что функция не принадлежит классу. 
В каждом столбце таблицы 
есть как плюсы, так и минусы.
Таким образом, рассмотренные классы
$T_0$, $T_1$, $T_{*}$, $T_\le$, $T_L$ попарно различны,
не пусты и не совпадают с $P_n$.
\end{example}
\begin{remark}
Две функции, штрих Шеффера и стрелка Пирса, 
не принадлежат ни одному из рассматриваемых замкнутых классов.
\end{remark}
\end{column}
\end{columns}


\end{frame}

\begin{frame}
\label{3.6.2.}\subsubsection{3.6.2. Полные системы функций}
\frametitle{3.6.2. Полные системы функций (1/2)}

%
%\subsection{Полные системы функций}
%\label{full3}

%\vspace{-2pt}
Класс функций $F$ называется \odmkey{полным,}{Класс (class)!полный
(complete)} \odmkey{}{Полнота (completeness)!системы булевых
функций
 (of boolean functions set)} если его замыкание совпадает с $P_n$:
$
[F] = P_n.
$
Другими словами, множество функций $F$ образует полную систему, если
любая функция реализуема в виде формулы над $F$.



\begin{theorem}
Пусть заданы две системы функций:
$F = \{\Tuple{f}{1}{m}\}$ и $G = \{\Tuple{g}{1}{k}\}$.
Тогда, если система $F$ полна и все функции из $F$ реализуемы
формулами над $G$, то система $G$ также полна:

\vspace{-12pt}
$$
([F]=P_n \And \A{i\in 1..m}{f_i=\func{\cal G}_i[G]}) \Impl [G]=P_n.
$$
\end{theorem}



%\vspace{-6pt}
\begin{proof}
Пусть $h$~--- произвольная функция, $h\in P_n$. \\ Тогда
$[F]=P_n \Impl h=\func \cal F[F] \Impl \cal F[\cal
G_i/f_i]$~--- формула над $G$. \\ Следовательно, $h= \func\cal
G[G].$
\end{proof}

%
%{%
%\def\example{\pagebreak[2]\par\addvspace{9pt plus 1pt minus 3pt}{\sffamily\bfseries Примеры.}\quad\nopagebreak}
%\def\endexample{\par\addvspace{9pt plus 1pt minus 3pt}\pagebreak[2]}


\begin{examples}
Система $\{\Or,\wedge,\neg\}$~--- полная 
(\DMref{п.~3.4.3}{3.4.3.}). 
Следовательно:
\begin{enumerate}
\item[1)] система $\{\neg,\wedge\}$ полная, так~как $x_1\vee
x_2=\neg(\neg x_1\wedge \neg x_2)$; \item[2)] система $\{\neg,\vee
\}$ полная, так~как $x_1\wedge x_2=\neg(\neg x_1\vee \neg x_2)$;
\item[3)] система $\{\mid\}$ полная, так~как $\neg x ={x}\mid{x}$,
$x_1\wedge x_2=\neg({x_1}\mid{x_2})=
({x_1}\mid{x_2})\mid({x_1}\mid{x_2})$. 
\end{enumerate}
\end{examples}

\end{frame}

\begin{frame}
\frametitle{3.6.2. Полные системы функций (2/2)}
Cистема $\{0, 1,
\wedge, +\}$ полная, так~как $\neg x=x+1$ (здесь $+$ означает
сложение по модулю 2). Представление булевой функции над базисом
$\{0, 1, \wedge, +\}$ 
называется \odmkey{полиномом Жегалкина}{Полином Жегалкина (Zhegalkin polynomial)}. 

Таким образом,
всякая булева функция представима в виде
$
\sum\limits_{(\Tuple{i}{1}{s})}a_{\Tuple{i}{1}{s}}x_{i_1}\cdot\ldots\cdot
x_{i_s},
$
где $\sum$~--- сложение по модулю 2, знак $\cdot$ обозначает конъюнкцию и
$a_{\Tuple{i}{1}{s}}\in{E_2}$.



\begin{remark}
Фактически, в полиноме Жегалкина $\sum a_{\Tuple{i}{1}{s}}
x_{i_1}\cdot\ldots\cdot x_{i_s} = \sum x_{i_1}\cdots x_{i_s}$,
поскольку если $a_{\Tuple{i}{1}{s}}=1$, то этот коэффициент можно
опустить, а если $a_{\Tuple{i}{1}{s}}=0$, то можно опустить всё
слагаемое. Знак конъюнкции обычно также опускают.
\end{remark}

\begin{example}
$\Not x = x+1$, $x\vee y = xy+x+y$, $x\equiv y = x+y+1$.
\end{example}

%\vspace{-6pt}
\begin{remark}
Если в полиноме Жегалкина произвести 
следующие замены логических
операций арифметическими:
 $x\wedge y\mapsto xy$; $x+_2 y \mapsto x + y -2xy $, 
 то получится арифметический полином, который можно упростить и использовать для вычисления значения булевой функции.
 \end{remark}
 
\end{frame}

\begin{frame}
\label{3.6.3.}\subsubsection{3.6.3. Полнота двойственной системы}
\frametitle{3.6.3. Полнота двойственной системы}

%\subsection{Полнота двойственной системы}

\begin{theorem}
Если система $F=\{f_1,\dots, f_k\}$ полна, то система двойственных
функций $F^* =\{f_1^*,\dots, f_k^*\}$ также полна.
\end{theorem}

\begin{proof}
Пусть $h$~--- произвольная функция, $h\in P_n$. Рассмотрим
двойственную функцию~$h^*$. Система $F$ полна, так что $h^* =\func
\mathcal{H}[F]$.  По принципу двойственности $h= \func
\mathcal{H}^* [F^*]$.
\end{proof}

%\vspace*{-.75\baselineskip}

\begin{example}
Система $\{0, 1, \wedge, +\}$ полна, следовательно, система $\{1,
0, \vee, \equiv\}$ также полна.
\end{example}

%\vspace*{-.5\baselineskip}

\end{frame}


\begin{frame}
\label{3.6.4.}\subsubsection{3.6.4. Теорема Поста }\frametitle{3.6.4. Теорема Поста (1/5)}


%\subsection{Теорема Поста}

Теорема Поста устанавливает необходимые и достаточные условия
полноты системы булевых функций.\odmkey{}{Теорема (theorem)!Поста
 (Emil Leon Post)}

%\pagebreak[3]
\begin{theorem} Система булевых функций $F$ полна тогда и только
тогда, когда она содержит хотя бы одну функцию, не сохраняющую нуль, хотя
бы одну функцию, не
сохраняющую единицу, хотя бы одну несамодвойственную функцию, хотя бы
одну немонотонную функцию и хотя бы одну нелинейную функцию:
$$
[F] = P_n \Equiv \neg(F \subset T_0 \Or F \subset T_1 \Or F \subset T_* \Or
F \subset T_\le \Or F \subset T_L).
$$
\end{theorem}



\begin{proof}
\proofsection{Необходимость} От противного.
Пусть
$[F]=P_n$ и
$$
F\subset T_0\vee F\subset
T_1\vee F\subset T_*\vee F\subset T_\le \vee F\subset T_L.
$$
Введём обозначение: $i$~--- один из индексов, $0$, $1$, $*$, $\le$
или $L$.\\
Тог\-да $ T_i=[T_i]\Impl [F]\subset T_i\Impl P_n\subset T_i\Impl
P_n=T_i, $ но $P_n\not= T_i$ по таблице из
\DMref{п.~3.6.1}{3.6.1.}.
\end{proof}
\end{frame}


\begin{frame}
\frametitle{ 3.6.4. Теорема Поста (2/5)}
%\vspace{-4pt}
\begin{proof} (Продолжение 1)
%\addvspace{6pt}
\proofsection{Достаточность} Пусть $\neg(F\subset T_0\vee F\subset T_1\vee
F\subset T_*\vee F\subset T_\le\vee F\subset T_L)$. Тогда
$$
\E{F'=\langle f_0, f_1, f_*, f_\le, f_L\rangle}{f_0\not\in T_0\And
f_1\not\in T_1\And f_*\not\in T_*\And f_\le\not\in T_\le \And f_L\not\in T_L}.
$$
Функции $f_0, f_1, f_*, f_\le, f_L$  не обязательно различны и не
обязательно исчерпывают $F$.  Покажем, что отрицание и конъюнкция
реализуются в виде формул над~$F'$.
Тем самым теорема будет доказана (см.~\DMref{п.~3.6.2}{3.6.2.}).
Построение проводится в три этапа: на первом строятся формулы,
реализующие константы 0 и 1, которые нужны на третьем этапе. На
втором этапе строится формула, реализующая отрицание. На~третьем
этапе строится формула, реализующая конъюнкцию.



\proofsection{Константы} Построим формулу, реализующую 1. Пусть
$\varphi(x)\Assign f_0(\Tuple{x}{}{})$. Тогда\\
$\varphi(0)=f_0(0,\dots, 0)\not=0\Impl\varphi(0)=1$.
Возможны два случая: $\varphi(1)=1$ или
$\varphi(1)=0$.
Пусть $\varphi(1)=1$. В этом случае формула $\varphi$ реализует 1.
Пусть $\varphi(1)=0$. В этом случае формула $\varphi$ реализует отрицание.

\end{proof}
\end{frame}


\begin{frame}
\frametitle{3.6.4. Теорема Поста (3/5)}
\begin{proof} (Продолжение 2)
Тогда рассмотрим функцию $f_*$. Имеем
$$
f_*\not\in T_* \Impl \E{\Tuple{a}{1}{n}}{f_*(\Tuple{a}{1}{n})\ne
\overline{f_*}(\Tuple{\overline a}{1}{n})}.
$$
Следовательно, $f_*(\Tuple{a}{1}{n})=f_*(\Tuple{\overline{a}}{1}{n})$.
Пусть теперь  $\psi(x)\Assign f_*(x^{a_1}, \dots,
x^{a_n})$. Тогда
$$
\psi(0)=  f_*(0^{a_1},\dots, 0^{a_n})=
f_*(\Tuple{\overline {a}}{1}{n})=
 f_*(\Tuple{a}{1}{n})=f_*(1^{a_1},\dots,
1^{a_n})=\psi(1).
$$
Таким образом, $\psi(0)=\psi(1)$, откуда $\psi=1$  или $\psi=0$.
Если $\psi=1$, то требуемая константа 1 построена. В противном
случае $\psi$ реализует 0, и,~значит,
$\varphi(\psi(x)) =\overline\psi(x)$ реализует~1.

Построение $0$ аналогично, только вместо $f_0$ нужно использовать $f_1$.

\proofsection{Отрицание}
Построим формулу, реализующую отрицание.
Рассмотрим функцию $f_\le$. \\Имеем
$f_\le\not\in T_\le \Impl\E{\alpha=(\Tuple{a}{1}{n}),\beta=(\Tuple{b}{1}{n})}
{\alpha\le\beta\And f_\le(\alpha)>f_\le(\beta)}$.\\
Тог\-да $\alpha \le  \beta\Impl \A{i}{a_i = b_i\vee a_i = 0\And b_i = 1}$.
Но
$f_\le(\alpha)\not=f_\le(\beta)\Impl\alpha\not=
\beta\Impl$\\
$\Impl \E{J\subset 1..n}{j\in J \Impl a_j=0\And b_j=1}$.
Другими словами, $J$~--- это множество индексов $j$, для которых
$a_j\not=b_j$. 
\end{proof}
\end{frame}


\begin{frame}
\frametitle{3.6.4. Теорема Поста (4/5)}
\begin{proof} (Продолжение 3)
Пусть $\varphi(x)\Assign f_\le (\Tuple{c}{1}{n})$,
где $c_j\Assign x$, если $j\in J$, и $c_j\Assign a_j (=b_j)$, если
$j\notin J$. Тогда
$$\varphi(0) = f_\le
(\Tuple{c}{1}{n})[0/x] = f_\le(\alpha) > f_\le(\beta) = 
f_\le(\Tuple{c}{1}{n})[1/x] = \varphi(1).$$ 
Имеем
$\varphi(0)>\varphi(1)\Impl\varphi(0)=1\And \varphi(1)=0\Impl
\varphi(x)= \overline x$.
\proofsection{Конъюнкция}  Построим
формулу, реализующую конъюнкцию. Рассмотрим функцию $f_L$. Имеем
$
{f_L}\in{P_n}\Impl
f_L=\sum\limits_{\Tuple{i}{1}{s}} 
a_{\Tuple{i}{1}{s}}
 x_{i_1}\cdot\ldots\cdot x_{i_s}.
$
Но ${f_L}\notin{T_L}$, следовательно,
в полиноме Жегалкина существует нелинейное
слагаемое, содержащее конъюнкцию по крайней мере двух переменных.
Пусть, для определённости, это $x_1$ и $x_2$. Тогда\\
$
f_L= x_1\cdot x_2\cdot
f_a(\Tuple{x}{3}{n})+
x_1\cdot f_b(\Tuple{x}{3}{n}) 
+ x_2\cdot f_c(\Tuple{x}{3}{n})+f_d(\Tuple{x}{3}{n})$,\\
причём $f_a(\Tuple{x}{3}{n})\not= 0$. Следовательно,
$\E{\Tuple{a}{3}{n}} {f_a(\Tuple{a}{3}{n})=1}$. 
\\Пусть $b\Assign
f_b(\Tuple{a}{3}{n})$, $c\Assign f_c(\Tuple{a}{3}{n})$, $d\Assign
f_d(\Tuple{a}{3}{n})$ \\ и
$
\varphi(x_1, x_2)\Assign f_L(x_1, x_2, \Tuple{a}{3}{n})= x_1\cdot
x_2+b\cdot x_1+c\cdot x_2+d.
$
\end{proof}
\end{frame}


\begin{frame}
\frametitle{3.6.4. Теорема Поста (5/5)}
\begin{proof} (Окончание) Пусть далее $\psi(x_1, x_2)\Assign
\varphi(x_1+c, x_2+b)+b\cdot c+d$. Тогда
$
\psi(x_1, x_2)=  (x_1+c)\cdot (x_2+b)+ b\cdot (x_1+c)+c\cdot (x_2+b)
+d+b\cdot c+d=$\\
$=  x_1\cdot x_2+c\cdot x_2+b\cdot x_1+b\cdot c+b\cdot
x_1+b\cdot c+c\cdot x_2 +{b\cdot c}
 + d+b\cdot c+d=x_1\cdot
x_2.
$
Функции $x+a$ выразимы, так~как $x+1=\overline x$, $x+0=x$, а
константы $0$, $1$ и~отрицание уже построены.
\end{proof}
%
\medskip
%
\begin{example}
\begin{enumerate}
\item В системе $\{\neg, \wedge\}$ отрицание
 не сохраняет констант и немонотонно,
а~конъюнкция несамодвойственна и нелинейна.
\item Функция штрих Шеффера (также как функция стрелка Пирса)
не сохраняет констант, несамодвойственна,
немонотонна и нелинейна, а потому сама по себе образует полную систему. 
\end{enumerate}
\end{example}

%\vspace{-.5\baselineskip}
\end{frame}



% ===== tex/DM2022_4.tex =====
\begin{frame}
\frametitle{4. Логические исчисления (1/3)}
\section{4. Логические исчисления}
\vspace{-2pt}
С древнейших времён человечеству известна \em{логика}, 
или \em{искусство
правильно рассуждать}. Вообще говоря, способность к
рассуждениям~--- это именно искусство. Имея какие-то утверждения
(посылки), истинность которых проверена, скажем, на опыте, логик
путем умозрительных построений приходит к другому утверждению
(заключению), которое также оказывается истинным (в некоторых
случаях). Опыт древних (чисто наблюдательный) был систематизирован
Аристотелем. Он
рассмотрел конкретные виды рассуждений, которые назвал 
\em{силлогизмами\/}. А~именно, Аристотель рассмотрел так называемые
\em{категорические утверждения\/} четырех видов 
($A,B$~---  \em{категории}):
\begin{enumerate}
\item[1)] все $A$ обладают свойством $B$ (все $A$ суть $B$);%
\item[2)] некоторые $A$ обладают свойством $B$ (некоторые $A$ суть $B$);%
\item[3)] все $A$ не обладают свойством $B$ (все $A$ суть не $B$);%
\item[4)] некоторые\hspace{-1pt} $A$\hspace{-1pt}
 не \hspace{-1pt}обладают \hspace{-1pt}свойством 
 \hspace{-1pt}$B$ (некоторые $A$ суть не $B$)
\end{enumerate}
и зафиксировал все случаи, когда из посылок такого вида
выводятся заключения одного из этих же видов.
\end{frame}

\begin{frame}
\frametitle{4. Логические исчисления (2/3)}

\begin{examples}
\begin{enumerate}
\item Все люди смертны. Сократ~--- человек. Следовательно, Сократ смертен.
Это рассуждение правильно, потому что подходит под один из
силлогизмов Аристотеля.
\item
Все дикари раскрашивают свои лица.
Некоторые современные молодые люди раскрашивают свои лица.
Следовательно, некоторые современные молодые люди~--- дикари.
Рассуждение неправильно, хотя, видимо, все
входящие в~него утверждения истинны.

\end{enumerate}
\end{examples}

\smallskip
Логика Аристотеля~--- это \em{классическая} логика, то есть
наука, традиционно относящаяся к гуманитарному циклу и тем самым
находящаяся вне рамок данных лекций.

\smallskip
Предметом этой лекции являются некоторые элементы
логики \em{математической}, которая соотносится с логикой
классической примерно так, как язык Паскаль
соотносится с английским языком. 
Эта аналогия довольно точна
и по степени формализованности,
и по широте применимости в реальной жизни,
и по значимости для практического программирования.
\end{frame}

\begin{frame}
\frametitle{4. Логические исчисления (3/3)}


Математическая логика не является математической моделью классической логики, 
подобно тому, как язык Паскаль не является моделью английского языка. 
Для математической логики характерна формализация логических операций, 
абстрагирование от конкретного содержания предложений, 
выражающих какое-либо суждение. 
Истинность или ложность вывода зависит только от его формы, 
но не от конкретного содержания составляющих его предложений.

\medskip
План главы состоит в том, чтобы на основе небольшого
предварительного рассмотрения ввести
понятие <<формальная теория>>, или <<исчисление>>, в
наиболее общем виде, а затем конкретизировать это понятие
примерами двух наиболее часто используемых исчислений:
исчисления высказываний и исчисления предикатов.
Главу завершают два раздела, демонстрирующие связь
математической логики с программированием: 
наивная теория алгоритмов и автоматическое доказательство теорем.
\end{frame}


\subsection{4.1. Логические связки и кванторы}
\label{4.1.}
\begin{frame}
\frametitle{4.1. Логические связки и кванторы (1/2)}
Цель данного раздела~--- неформально ввести специфическую <<логическую>>
терминологию и~указать на её связь с материалом предшествующих глав.
В последующих разделах главы определения этого раздела уточняются
и конкретизируются.

\medskip
В конце раздела приведено краткое и неформальное 
обсуждение самого понятия <<формализация>>.

\medskip
\begin{tabular}{p{10mm}|p{40mm}|p{92mm}}
  № & Название параграфа
  & Ключевые термины и обозначения \\
  \hline
\DMref{4.1.1.}{4.1.1.} & Высказывания и связки  &
{ Простые высказывания, пропозициональные переменные, истинностные значения, логические связки, составные высказывания}\\  
\DMref{4.1.2.}{4.1.2.} & Формулы &
{ Пропозициональные формулы }\\  
\DMref{4.1.3.}{4.1.3.} & Интерпретация  &
{ Выполнимая формула, общезначимая формула, тавтология, невыполнимая формула, противоречие }\\ 
\DMref{4.1.4.}{4.1.4.} & 
{ Логические следование и эквивалентность } 
& { Алгебра высказываний }\\
\end{tabular}  
\end{frame}


\begin{frame}
\frametitle{4.1. Логические связки и кванторы (2/2)}
\begin{tabular}{p{10mm}|p{40mm}|p{92mm}}
  № & Название параграфа
  & Ключевые термины и обозначения \\
  \hline
\DMref{4.1.5.}{4.1.5.}\BS & Подстановка и замена  &
{ }\\  
\DMref{4.1.6.}{4.1.6.}\RS & Кванторы   &
{Одноместный и $n$-местный предикат,  универсальное
 высказывание, квантор всеобщности, 
  экзистенциальное высказывание, квантор существования, 
квантор существования и единственности, логический квадрат, контрарные, субконтрарные и контрадикторные высказывания }\\    
\DMref{4.1.7.}{4.1.7.} & Формализация   &
{ }\\  
\end{tabular}  

\end{frame}


%\subsection{Высказывания}
\begin{frame}
\label{4.1.1.}\subsubsection{4.1.1. Высказывания и связки }
\frametitle{4.1.1. Высказывания и связки (1/2)}
Элементами логических рассуждений являются утверждения, которые
либо истинны, либо ложны, но не то и другое вместе. Такие
утверждения называются \odmkey{(простыми)
высказываниями}{Высказывание (statement, proposition)! простое (atomic)}.
Простые высказывания обозначаются \odmkey{пропозициональными
переменными}{Переменная (variable)!пропозициональная
(propositional)}. Пропозициональная переменная может принимать
одно из двух \odmkey{истинностных значений}{Значение
(value)!истинностное (truth value)}, обозначаемых символами <<T>>
и <<F>>.

Из простых высказываний с помощью \odmkey{логических связок}{Связка (connective)!логическая (logical)}
могут быть построены
составные высказывания. Обычно рассматривают следующие
логические связки.

\medskip
\begin{tabular}{p{30mm}p{30mm}p{30mm}}
Название & Прочтение & Обозначение \\
\hline
Отрицание & не & $\neg$\\
Конъюнкция &и & $\And$ \\
Дизъюнкция &или & $\Or$ \\
Импликация & если\ldots то& $\to$ \\
\hline
\end{tabular}

\end{frame}

\begin{frame}
\frametitle{4.1.1. Высказывания и связки (2/2)}

\begin{remark}
Истинностные значения обозначают различными способами: латинскими
буквами T и F, как в этом курсе, русскими буквами И и Л (от
слов <<Истина>> и <<Ложь>>), цифрами 1 и 0, специальными значками
$\top$ и $\bot$. 
Каждый из способов обозначения истинностных
значений имеет свои достоинства и недостатки. 
Пара символов, T и F,
выбрана нами, во-первых, для удобства чтения, во-вторых, так как эти
символы не используются в курсе в другом смысле, и, в-третьих,
ради ассоциации с английскими словами <<true>> и <<false>>,
которые используются во многих языках программирования для
обозначения истинностных значений. При этом 
в программах мы иногда используем символы 1 и 0 ради краткости.
\end{remark}

\end{frame}


%\subsection{Формулы}
\begin{frame}
\label{4.1.2.}\subsubsection{4.1.2. Формулы }\frametitle{4.1.2. Формулы (1/2)}
Правильно построенные составные высказывания называются
\odmkey{(пропозициональными) формулами}{Формула (formula)!пропозициональная
(propositional)}.
%Формулы имеют следующий синтаксис:
Синтаксис формул:

\vspace{-1.2\baselineskip}
\begin{alignat*}{2}
&\llrle{формула}\dDef &\quad&   \text{T} \mid \text{F} \mid \\
&&&\llrle{пропозициональная переменная} \mid \\
&&&(\, \neg\llrle{формула}\, ) \mid \\
&&&(\, \llrle{формула} \And \llrle{формула}\,) \mid \\
&&& (\, \llrle{формула}\Or  \llrle{формула}\,) \mid \\
&&& (\,\llrle{формула}\to\llrle{формула}\,)
\end{alignat*}

\begin{remark}
Для описания синтаксиса языка
мы используем один из стандартных приёмов:
кон\-текст\-но"=сво\-бод\-ную порождающую грамматику
в форме Бэкуса--Наура. Здесь названия
синтаксических конструкций заключаются в угловые
скобки, метасимвол <<$\dDef$>> означает <<это>>,
метасимвол <<$\mid$>> означает <<или>>,
все остальные символы означают сами себя.
%Исчерпывающее изложение теории формальных грамматик
%применительно к программированию можно найти в ~\cite{Aho2}.
\end{remark}
\end{frame}


%\subsection{Формулы}
\begin{frame}
\frametitle{4.1.2. Формулы (2/2)}
Для упрощения записи вводится старшинство связок 
($\neg$, $\And$, $\Or$,
$\to$), и лишние скобки часто опускаются.

\medskip
Истинностное значение формулы определяется через истинностные значения её
составляющих в соответствии со следующей таблицей:

\medskip
\begin{center}
\begin{tabular}{c|c||c|c|c|c}
$A$ & $B$ & $\neg A$ & $A\And B$ & $A\Or B$ & $A\to B$ \\
\hline
 F  &   F  &      T    &       F    &      F    &      T \\
%\hline
 F  &   T  &      T    &       F    &      T    &      T  \\
%\hline
 T  &   F  &      F    &       F    &      T    &      F  \\
%\hline
 T  &   T  &      F    &       T    &      T    &      T  \\
\hline
\end{tabular}
\end{center}

\medskip
\begin{digress}
Важно понимать, что предыдущая таблица является соглашением, а не <<законом природы>>. 
\end{digress}
\end{frame}



\begin{frame}
\label{4.1.3.}\subsubsection{4.1.3. Интерпретация }\frametitle{4.1.3. Интерпретация (1/2)}
Пусть $A(x_1,\dots,x_n)$~--- пропозициональная формула, где
$x_1,\dots,x_n$~--- входящие  в неё пропозициональные переменные.
Конкретный набор истинностных значений, приписанных переменным
$x_1,\dots,x_n$, называется \odmkey{интерпретацией}{Интерпретация
(interpretation)!формулы (of formula)} формулы~$A$. Формула может
быть истинной (иметь значение T) при одной интерпретации и ложной
(иметь значение F) при другой интерпретации. Значение формулы $A$
в интерпретации $I$ будем обозначать $I(A)$. 

\medskip
Формула, истинная при
некоторой интерпретации, называется \odmkey{выполнимой}{Формула
(formula)!выполнимая (satisfiable)}. Формула, истинная при всех
возможных интерпретациях, называется \odmkey{общезначимой}{Формула
(formula)!общезначимая (valid)} (или
\odmkey{тавтологией}{Тавтология (tautology)}). Формула, ложная при
всех возможных интерпретациях, называется
\odmkey{невыполнимой}{Формула (formula)!невыполнимая
(unsatisfiable)} (или \odmkey{противоречием}{Противоречие
(contradiction)}).


\begin{examples}
$A\vee\neg A$~--- тавтология, $A\And\neg A$~--- противоречие,
$A\to\neg A$~--- выполнимая формула, она истинна при $I(A)=\text{F}$.
\end{examples}
\end{frame}


\begin{frame}
\frametitle{4.1.3. Интерпретация (2/2)}
%\setcounter{theorem}{0}
\begin{theorem}
Пусть $A$~{\upshape---} некоторая формула. Тогда:
\begin{enumerate}
\item[1)] если $A$~{\upshape---} тавтология, то $\neg
A$~{\upshape---} противоречие, и наоборот; \item[2)] если
$A$~{\upshape---} противоречие, то $\neg A$~{\upshape---}
тавтология, и наоборот; \item[3)] если $A$~{\upshape---}
тавтология, то неверно, что $A$~{\upshape---} противоречие, но не
наоборот; \item[4)] если $A$~{\upshape---} противоречие, то
неверно, что $A$~{\upshape---} тавтология, но не наоборот;
%\pagebreak
\item[5)] если $\neg A$ выполнима, то неверно, что
$A$~{\upshape---} тавтология; \item[6)] если  $A$ выполнима, то
неверно, что $A$~{\upshape---} противоречие.
\end{enumerate}
\end{theorem}


\begin{proof}
Очевидно из определений.
\end{proof}

\smallskip
\begin{theorem} Если формулы $A$ и $A\to B$~{\upshape---} тавтологии, то формула
$B$~{\upshape---} тавто\-логия.
\end{theorem}


\begin{proof}
От противного. Пусть $I(B)=\text{F}$.
Но $I(A)=\text{T}$, так~как $A$~--- тавтология.
Значит,~$I(A\to B)=\text{F}$,
что противоречит предположению о том, что $A\to B$~---
тавтология.
\end{proof}
\end{frame}


%\subsection{Логическое следование и логическая эквивалентность}

\begin{frame}
\label{4.1.4.}\subsubsection{4.1.4. Логические следование и эквивалентность }\frametitle{4.1.4. Логические следование и эквивалентность (1/3)}
Говорят, что формула $B$ \odmkey{логически следует}{Следствие логическое (logical consequence)} из формулы $A$ (обозначение
${A  \Impl B}$),
%\linebreak
если формула $B$ имеет значение T
при всех интерпретациях, при которых формула $A$ имеет значение T.
Говорят, что формулы $A$ и $B$
\odmkey{логически эквивалентны}{Эквивалентность логическая (logical equivalence)} (обозначается
$A\Equiv B$ или просто $A=B$), если они являются
логическими следствиями друг друга.
Логически эквивалентные формулы имеют
одинаковые значения при любой интерпретации.


%\setcounter{theorem}{0}
\begin{theorem}
$P\to Q = \neg P \Or Q.$
\end{theorem}



\begin{proof}
Для доказательства достаточно проверить, что формулы действительно
имеют одинаковые истинностные значения при всех интерпретациях.



%\addvspace{6pt}
%{\sffamily\small
%\def\arraystretch{1.1}
\begin{tabular}{p{10mm}|p{10mm}||p{15mm}||p{15mm}|p{15mm}p{45mm}}
$P$ & $Q$ & $P\to Q$ & $\neg P$ & $\neg P\vee Q$  \\
\cline{1-5}
F & F & T & T & T \\
F & T & T & T & T \\
T & F & F & F & F \\
T & T & T & F & T & \\
\cline{1-5}
\end{tabular}


\end{proof}
\end{frame}

\begin{frame}
\frametitle{4.1.4. Логические следование и эквивалентность (2/3)}
\begin{theorem}
Если $A$, $B$, $C$~--- любые формулы, то имеют место
следующие логические эквивалентности:
\end{theorem}

\vspace{-0.5\baselineskip}
\begin{tabbing}
14.4\= $\quad$ \= \kill
\>1. $A \Or A = A,$
  $A \And A = A.$ \\
\>2. $A \Or B = B \Or A,$
  $A \And B = B \And A.$ \\
\>3. $A \Or (B \Or C) = (A \Or B) \Or C,$
  $A\And (B \And C) = (A \And B) \And C.$ \\
\>4. $A \Or (B \And C) = (A \Or B) \And (A \Or C),$\\
\> \> $A \And (B \Or C) = (A \And B) \Or (A \And C).$ \\
\>5. $(A \And B) \Or A = A,$
  $(A \Or B) \And A = A.$ \\
\>6. $A \Or \text{F} =A,$
  $A \And \text{F} =\text{F}.$ \\
\>7. $A \Or \text{T} = \text{T},$
  $A \And \text{T} = A.$ \\
\>8. $\neg(\neg A) = A.$ \\
\>9. $\neg(A \And B) = \neg A \Or \neg B,$
  $\neg(A \Or B) = \neg A \And \neg B.$ \\
\>10. $A \Or \neg A = \text{T},$
  $A \And \neg A = \text{F}.$
\end{tabbing}

\vspace{-0.5\baselineskip}
\begin{proof}
Непосредственно проверяется построением таблиц
истиннос\-ти.
\end{proof}

\begin{remark}
Таким образом, алгебра $\Alg{\{\text{T,F}\}}{\Or,\And,\neg}$
является булевой алгеброй, которая называется \odmkey{алгеброй
высказываний}{Алгебра (algebra)!высказываний (propositional
algebra)}.
\end{remark}
\end{frame}

\begin{frame}
\frametitle{4.1.4. Логические следование и эквивалентность (3/3)}

\begin{theorem} \NB{[1]}
$P_1  \And  \dots  \And  P_n  \Impl  Q$ тогда и только
тогда, когда $(P_1 \And \dots \And  P_n) \to 
Q$~--- тавтология.
\end{theorem}


\begin{proof}
\proofsection{$\Longrightarrow$} 
Пусть $I(P_1 \And  \dots
  \And  P_n) = \text{T}$. Тогда $I(Q) = \text{T}$ и
$I(P_1 \And  \dots   \And  P_n \to  Q) =\text{T}$.
Пусть $I(P_1\And\dots\And P_n)=\text{F}$. Тогда
$I(P_1\And\dots\And P_n\to Q)=\text{T}$  при любой интерпретации
$I$. Таким образом, формула $P_1\&\dots\& P_n\to Q$ общезначима.
\proofsection{$\Longleftarrow$} 
 Пусть $I(P_1\And\dots\And P_n)=\text{T}$.
Тогда $I(Q)=\text{T}$,
иначе формула
$P_1\And\dots\And P_n\to Q$
не была бы тавтологией. Таким образом, формула $Q$~---
логическое следствие формулы
$P_1\And\dots\And P_n$.
\end{proof}


\begin{theorem} \NB{[2]}
$P_1  \And  \dots  \And  P_n  \Impl  Q$ тогда и только
тогда, когда $P_1 \And \dots \And  P_n \And   \Not
Q$~{\upshape---} противоречие.
\end{theorem}

\begin{proof}
По предыдущей теореме
$P_1 \And \dots \And P_n\Impl Q$ тогда и только тогда, когда
формула $P_1 \And \dots \And P_n\to Q$~--- тавтология. По первой теореме
\DMref{п.~4.1.3}{4.1.3.} формула \\
$P_1 \And \dots \And P_n\to Q$ является
тавтологией тогда и только тогда, когда формула\\
$\neg(P_1 \And \dots \And P_n\to Q)$ является противоречием. Имеем
$\neg(P_1 \And \dots \And P_n\to Q)=$ \\
$=\neg(\neg(P_1 \And \dots \And P_n)\vee Q) =\neg\neg(P_1 \And \dots \And P_n)\And \neg Q= P_1 \And \dots \And P_n \And \neg Q$.
\end{proof}

\end{frame}

\begin{frame}
\label{4.1.5.}\subsubsection{4.1.5. Подстановка и замена }\frametitle{4.1.5. Подстановка и замена (1/2)}

Пусть $A$~--- некоторая формула, в которую
входит переменная $x$ (обозначается $A(\dots x\dots)$) или входит некоторая
подформула $B$ (обозначается $A(\dots B \dots)$), и пусть $C$~---
некоторая формула.
Тогда
$A(\dots x\dots)[C/x]$
обозначает формулу,
полученную из формулы $A$ подстановкой формулы $C$ вместо {\it всех}
вхождений переменной~$x$, а $A(\dots B\dots)[C//B]$ обозначает любую формулу,
полученную из формулы $A$ подстановкой формулы $C$ вместо {\it
некоторых} (в частности, вместо одного) вхождений подформулы $B$.

%\setcounter{theorem}{0}
%\vspace{0.25\baselineskip}
\begin{theorem} \NB{[1]}
Если $A(\dots x\dots)$~{\upshape---} тавтология, а
$B$~{\upshape---} любая формула, \\то формула $A(\dots
x\dots)[B/x]$~{\upshape---} тавтология.
\end{theorem}

%\vspace{0.25\baselineskip}
\begin{proof}
Пусть $C\Assign A(\dots x\dots)[B/x]$.
Пусть $I$~--- интерпретация формулы $C$ (формула уже не
содержит переменной $x$).
Построим интерпретацию $I'$ формулы $A$, положив $x\Assign I(B)$,
а значения остальных переменных взяв такими же, как в интерпретации
$I$. Тогда
$I'(A)=I(C)$, но $I'(A)=\text{T}$, следовательно, $I(C)=\text{T}$.
\end{proof}
%\vspace{0.25\baselineskip}

\end{frame}

\begin{frame}
\frametitle{4.1.5. Подстановка и замена (2/2)}

\begin{theorem} \NB{[2]}
Если $A(\dots B\dots)$ и $B=C$, а $D\Assign A(\dots B\dots)[C//B]$, то
$A=D$.
\end{theorem}
%\vspace{0.25\baselineskip}
\begin{proof}
Пусть $I$~--- любая интерпретация, содержащая значения для всех
переменных в формулах $A$, $B$ и $C$. Тогда $I(B)=I(C)$, значит,
$I(A)\nolinebreak=\nolinebreak I(D).\nolinebreak\hspace{-5pt}$
\end{proof}

Теоремы~1 и 2 аналогичны правилам подстановки и замены, изложенным
в~\DMref{п.~3.2.3}{3.2.3.}.

\end{frame}

\begin{frame}
\label{4.1.6.}\subsubsection{4.1.6. Кванторы }\frametitle{4.1.6. Кванторы (1/7)}
\odmkey{Одноместным предикатом}{Предикат (predicate)!одноместный (unary)} (над множеством $M$) 
называется функция $\Map{P}{M}{\{\text{F},\text{T}\}}$.
Если $x\in M$, то выражение $P(x)$ обладает истинностным значением, 
и потому может рассматриваться как простое высказывание.

\medskip 
Аналогично \odmkey{$n$-местным предикатом}{Предикат (predicate)!$n$-местный ($n$-arity)} 
$P(\Tuple{x}{1}{n})$ называется функция\\ 
$\Map{P}{M_1\times\dots\times M_n}{\{\text{F},\text{T}\}}$,
причём $x_i\in M_i$. 

Поскольку прямое произведение множеств также является множеством, 
далее в этом параграфе в примерах и определениях рассматриваются 
для краткости только одноместные предикаты, поскольку для многоместных предикатов всё аналогично.

\begin{examples} 
Высказывание <<число пять больше числа два>> 
имеет истинностное значение $\text{T}$ 
и является простым высказыванием, но не является предикатом, 
поскольку не содержит переменных. 
Высказывание <<число $x$ больше двух>>
является одноместным предикатом и символически записывается $x>2$, 
а высказывание <<число $x$ больше числа $y$>> 
является двухместным предикатом и символически записывается $x>y$.
\end{examples}



\end{frame}

\begin{frame}
\frametitle{4.1.6. Кванторы (2/7)}

\vspace{-2pt}
\begin{remark}
Одноместный предикат соответствует \em{характеристической функции} подмножества, 
а $n$-местный предикат соответствует характеристической функции отношения 
(\DMref{п.~1.9.3}{1.9.3.}).
\end{remark}

Операция \odmkey{связывания квантором всеобщности}{Операция (operation)!связывания квантором (of binding by quantifier)!всеобщности (universal)} 
сопоставляет предикату $P$, определённому на множестве $M$, 
высказывание, обозначаемое $\A{x\in M}{P(x)}$, 
истинное в том и только в том случае, 
когда высказывание $P(x)$ имеет значение $\text{T}$ 
для всех $x\!\in\! M$, и ложное в противном случае. 
Высказывание $\A{x\!\in\! M}{P(x)}$ называется 
\odmkey{универсальным высказыванием}{Высказывание (statement, proposition)!универсальное (universal)} для предиката $P$.
\begin{remark}
Если множество $M$ подразумевается, то при записи его опускают: 
$\A{x}{P(x)}$.  
\end{remark} 
Символ $\forall x$ называется 
\odmkey{квантором всеобщности}{Квантор (quantifier)!всеобщности (universal)} по переменной $x$.

\begin{examples}
Предикат $x<10$, определенный на множестве $\Bbb{R}$, 
не является тождественно истинным, поэтому высказывание 
$\A{x\in\Bbb{R}}{x<10}$ имеет значение $\text{F}$. 
Предикат $x^2 \ge 0$, определенный на множестве $\Bbb{R}$, 
является тождественно истинным на множестве $\Bbb{R}$, 
поэтому высказывание $\A{x\in\Bbb{R}}{x^2 \ge 0}$ имеет значение $\text{T}$.
\end{examples}

\end{frame}

\begin{frame}
\frametitle{4.1.6. Кванторы (3/7)}
Если предикат $P(x)$ задан на конечном множестве 
$M = \{\Tuple{a}{1}{n}\}$, 
то операция связывания квантором всеобщности выражается через конъюнкцию:
$$\A{x\in M}{P(x)}\Equiv                  
P(a_1)\And P(a_2)\And \dots \And P(a_n).$$ 

\medskip
Операция \odmkey{связывания квантором существования}{Операция (operation)!связывания квантором (of binding by quantifier)!существования (existential)} 
сопоставляет предикату $P$, определённому на множестве $M$, 
высказывание, обозначаемое $\E{x\in M}{P(x)}$, 
истинное в том и только в том случае, 
когда высказывание $P(x)$ имеет значение $\text{T}$ 
хотя бы для одного $x\in M$, и ложное в противном случае. 
Высказывание $\E{x\in M}{P(x)}$ называется 
\odmkey{экзистенциальным высказыванием}{Высказывание (statement, proposition)!экзистенциальное (existential)} для предиката $P$.
\begin{remark}
Если множество $M$ подразумевается, то при записи его опускают: 
$\E{x}{P(x)}$.  
\end{remark}
 
Символ $\exists x$ называется 
\odmkey{квантором существования}{Квантор (quantifier)!существования (existential)} по переменной $x$.


\end{frame}


\begin{frame}
\frametitle{4.1.6. Кванторы (4/7)}

\vspace{-2pt}
\begin{examples}
Предикат $x=x+1$, определенный на множестве $\Bbb{N}$, 
является тождественно ложным, поэтому высказывание 
$\E{x\in\Bbb{N}}{x=x+1}$ имеет значение $\text{F}$. 
Высказывание $\E{x\in\Bbb{N}}{x+1 > 10}$ имеет значение $\text{T}$,
так как существует значение $x\in \Bbb{N}$,
например, 10, на котором предикат $x+1>10$, 
определенный на множестве $\Bbb{N}$, 
имеет значение $\text{T}$.
\end{examples}

Если предикат $P(x)$ задан на конечном множестве 
$M = \{\Tuple{a}{1}{n}\}$, 
то операция связывания квантором существования выражается через дизъюнкцию:
$$\E{x\in M}{P(x)}\Equiv 
P(a_1)\Or P(a_2)\Or \dots \Or P(a_n).$$


Иногда используется модификация квантора существования, 
при которой
операция \odmkey{связывания квантором существования и единственности}{Операция (operation)!связывания квантором (of binding by quantifier)!существования и единственности (existence and uniqueness)} 
сопоставляет предикату $P$, определенному на множестве $M$, 
высказывание, обозначаемое $\exists !\,x\in M\, (P(x))$, 
истинное в том и только в том случае, 
когда предикат $P(x)$ имеет истинностное значение $\text{T}$ 
в точности для одного элемента $x\in M$, 
и ложное в противном случае.
\end{frame}
 
%\end{document}
\begin{frame}
\frametitle{4.1.6. Кванторы (5/7)}

\begin{remark}
Если множество $M$ подразумевается, то при записи его опускают: 
$\exists\,!\,x\, (P(x))$.  
\end{remark}  
Символ $\exists !\, x$ называется 
\odmkey{квантором существования и единственности}{Квантор (quantifier)!существования и единственности (existence and uniqueness)} 
по переменной x.
\begin{examples}
Высказывание $\exists !\, x\in\Bbb{N}\, (x+1 = 10)$ имеет значение $\text{T}$,
так как существует единственное значение $9\in \Bbb{N}$,
на котором предикат $x+1=10$, 
определенный на множестве $\Bbb{N}$, 
имеет значение $\text{T}$.
\end{examples}


Квантор существования и единственности является не более чем 
<<синтаксическим сахаром>>, поскольку выражается через
кванторы существования, всеобщности и проверку равенства элементов:
$$
\exists !\,x\in M\, (P(x))\Equiv 
\E{x\in M}{P(x)\And \A{ y\in M}{P(y)\Impl x=y}}. 
$$  
 
\end{frame}
 
%\end{document}
\begin{frame}
\frametitle{4.1.6. Кванторы (6/7)}

%\vspace{-2pt}


\begin{columns}[t]
\begin{column}[T]{8cm}
Взаимосвязь между универсальными и экзистенциальными высказываниями 
удобно представлять в виде \odmkey{логического квадрата}{Квадрат (square)!логический (logical)}. Универсальные высказывания 
$\A{x}{P(x)}$ и $\A{x}{\Not P(x)}$, 
стоящие в двух верхних вершинах квадрата, 
называются \odmkey{контрарными}{Контрарность (contrary)}.
Экзистенциальные высказывания 
$\E{x}{P(x)}$ и $\E{x}{\Not P(x)}$, 
стоящие в двух нижних вершинах квадрата, 
 называются \odmkey{субконтрарными}{Субконтрарность (subcontrary)}.  
Высказывания, соединённые диагоналями квадрата, 
являются отрицаниями друг друга. 
Такие высказывания называются \odmkey{контрадикторными}{Контрадикторность (contradictory)}. 
Под каждым универсальным высказыванием 
стоит высказывание, логически  следующее из него.
\end{column}
\begin{column}[T]{7cm}
\includegraphics[height=60mm]{4_1.jpg}
\end{column}
\end{columns}

\end{frame}
 

\begin{frame}
\frametitle{4.1.6. Кванторы (7/7)} 
Контрарные высказывания 
$\A{x}{P(x)}$ и $\A{x}{\Not P(x)}$, 
стоящие в двух верхних вершинах квадрата, 
не могут быть одновременно истинными. 
Субконтрарные высказывания 
$\E{x}{P(x)}$ и $\E{x}{\Not P(x)}$, 
стоящие в двух нижних вершинах квадрата, 
не могут быть одновременно ложными. 

\begin{remark}
Контрарные высказывания могут быть одновременно ложными, 
субконтрарные высказывания могут быть одновременно истинными.
\end{remark}

\medskip
Контрадикторные высказывания всегда имеют
противоположные истинностные значения. 

\begin{example}
Высказывания 
<<неверно, что все студенты не понимают логический квадрат тогда и
только тогда, когда существуют студенты, 
которые понимают логический квадрат>>
и
<<неверно, что все студенты понимают логический квадрат тогда и
только тогда, когда существуют студенты, 
которые не понимают логический квадрат>>
имеют истинностное значение T.
\end{example}

\end{frame}

\begin{frame}
\label{4.1.7.}\subsubsection{4.1.7. Формализация }\frametitle{4.1.7. Формализация (1/4)}

\vspace{-2pt}
{\small
Рассмотрим текст на естественном языке~---
начальный фрагмент из замечательной книги   Д.~Гиль\-берта <<Основания геометрии>>
(с небольшими сокращениями и с упрощением форматирования).

{\rm\bf
2. Первая группа аксиом: аксиомы соединения (принадлежности)
}
{\rm
Аксиомы этой группы устанавливают отношения принадлежности 
между введёнными выше вещами~--- точками, прямыми и плоскостями~--- 
и гласят следующим образом.

$\text{I}_1$.
{ Для любых двух точек $A$, $B$ 
существует прямая $a$, принадлежащая каждой из этих двух точек $A$, $B$.
}

$\text{I}_2$.
{ Для двух точек $A$, $B$ существует не более одной прямой, 
принадлежащей каждой из точек $A$, $B$.
}


Здесь, как и в последующем, под двумя, тремя, ... 
точками (прямыми, плоскостями) всегда подразумеваются различные точки 
(прямые, плоскости).
Вместо термина <<принадлежат>> мы будем пользоваться также 
и другими формулировками. 
Например, вместо прямая $a$ принадлежит каждой из точек $A$ и $B$, 
мы будем говорить: прямая $a$ проходит через точки $A$ и $B$ 
или прямая $a$ соединяет точки $A$ и $B$; ...


$\text{I}_3$. {  
На прямой существуют по крайней мере две точки. 
Существуют по крайней мере три точки, не лежащие на одной прямой.}
}
}
\end{frame}


\begin{frame}
\frametitle{4.1.7. Формализация (2/4)}

\vspace{-2pt}
Здесь области изменения и обозначения предметных переменных заданы. 
Обозначим $P$~--- множество точек, $L$~--- множество прямых. 
Возможны различные варианты введения предикатов. 
Если ввести трёхместный предикат 
$\Map{R}{P\times P \times L}{\{\text{T},\text{F}\}}$, 
означающий принадлежность двух точек одной прямой, 
с ограничением, что два первых аргумента суть различные точки, 
то первые две аксиомы запишутся в виде одной простой формулы

\vspace{-12pt}
$$
\forall A\,\forall B\,\exists !\, a \ R(A,B,a),
$$

\vspace{-2pt}
а третья аксиома запишется в виде двух формул:

\vspace{-12pt}
$$\forall a \,\exists A \,\exists B \ R(A,B,a)\quad\text{и}$$ 
$$ 
\exists A\,\exists B\,\exists C\,\forall a 
\left(\Not\left(R(A,B,a) \And  R(B,C,a)  \And R(C,A,a)\right)\right).
$$

\vspace{-2pt}
Вторая часть третьей аксиомы позволяет определить новый предикат:

\vspace{-12pt}
$$S(A,B,C)\Assign \forall a 
\left(\Not\left(R(A,B,a) \And  R(B,C,a)  \And R(C,A,a)\right)\right),$$

\vspace{-2pt}
означающий, что три точки не лежат на одной прямой. 

\end{frame}


\begin{frame}
\frametitle{4.1.7. Формализация (3/4)}

В таком случае запись третьей аксиомы ещё более упростится:

\vspace{-12pt}
$$\forall a \,\exists A \,\exists B \ R(A,B,a) \quad\text{и}\quad
\exists A \,\exists B \,\exists C \ S(A,B,C).$$

Такой путь сокращает количество вхождений предикатов в формулы, но может затруднить понимание формул, поскольку предикаты получаются достаточно сложные и сопровождаются оговорками про различные значения переменных.

\smallskip
Можно пойти по другому пути, 
определив двухместный предикат 
$\Map{\in_{PL}}{P\times L}{\{\text{T},\text{F}\}}$, 
означающий, что точка лежит на прямой, явный предикат  
равенства точек  \\ $\Map{=_P}{P\times P}{\{\text{T},\text{F}\}}$
и предикат различия точек: $\ne_P(A,B)\Assign \Not(=_P(A,B))$. 
Аналогично можно ввести предикаты равенства и неравенства прямых: 
$\Map{=_L}{L\times L}{\{\text{T},\text{F}\}}$,\\
$\ne_L(a,b)\Assign \Not(=_L(a,b))$.
Поскольку предикаты двухместные, можно использовать
инфиксную запись, а индексы у значков опускать,
поскольку эти индексы указывают на типы аргументов,
которые можно восстановить из контекста.   

\end{frame}


\begin{frame}
\frametitle{4.1.7. Формализация (4/4)}

В этом случае аксиомы запишутся более длинными, 
но, может быть, более понятными формулами:

\vspace{-0.75\baselineskip}
$$
\A{A,B}{A\ne B \to \exists\,!\, {a}\ {(A\in a \And B\in a)}},
$$
%∀A ∀B (P_≠ (A,B)→∃a (Q(A,a)  & Q(B,a)) )
%$$
%\A{A,B}{\A{a,b}{A\in a \And B\in a \And A\in b \And B\in b \to a=b}},
%$$
%∀A ∀B ∀a ∀b ((Q(A,a)  & Q(B,a)  & Q(A,b)  & Q(B,b))→L_= (a,b))
$$
\A{a}{\E{A,B}{A\in a \And B\in a}},
$$
%∀a ∃A ∃B(Q(A,a)  & Q(B,a))
$$
\E{A,B,C}{\Not \E{a}{A\in a\And B\in a\And C\in a}}.
$$
%∃A ∃B ∃C (¬∃a(Q(A,a)& Q(B,a)& Q(C,a)) )


\medskip
{\large
Таким образом, при переводе высказываний с естественного языка 
на язык математической логики большое значение имеет выбор системы предикатов, 
который ни в коей мере не является однозначным и предопределённым, 
а всегда зависит от целей, ради которых проводится \odmkey{формализация}{Формализация (formalization)}.
}
\end{frame}


\begin{frame}
\frametitle{Онтология математической логики}
\begin{center}
\includegraphics[height=7.8cm]{ODM4_1.png}
\end{center}

\end{frame}


\subsection{4.2. Формальные теории}
\label{4.2.}

\begin{frame}
\frametitle{4.2. Формальные теории}
 
Исторически понятие формальной теории было
разработано в период интенсивных исследований в области оснований
математики для формализации собственно логики и теории
доказательства. Сейчас этот аппарат широко используется при
создании специальных исчислений для решения конкретных прикладных
задач. 
%Рамки курса не позволяют привести развёрнутые примеры
%подобных специальных исчислений (они довольно велики), но тем не
%менее такие примеры существуют.

\medskip
\begin{tabular}{p{10mm}|p{38mm}|p{95mm}}
  № & Название параграфа
  & Ключевые термины и обозначения \\
  \hline
\DMref{4.2.1.}{4.2.1.}\RS & Определение формальной теории  &
{Алфавит, формулы, аксиомы, правила вывода, язык, сигнатура, логические и собственные аксиомы }\\  
\DMref{4.2.2.}{4.2.2.} & Выводимость   &
{Непосредственная выводимость, посылка, заключение, вывод, гипотеза, теорема теории, метатеорема }\\    
\DMref{4.2.3.}{4.2.3.} & Интерпретация   &
{ Область интерпретации, выполнимость формулы, модель множества формул и формальной теории  }\\  
\DMref{4.2.4.}{4.2.4.} & Общезначимость и непротиворечивость   &
{Тавтология, противоречие, логическое следствие, семантически и формально непротиворечивая теория}\\  
\DMref{4.2.5.}{4.2.5.} & Полнота, независимость и разрешимость   &
{ Аксиоматизируемая система,  разрешимая и полуразрешимая теория }\\  
\end{tabular}  

\end{frame}


%\vspace{-0.25\baselineskip}
%\subsection{Определение формальной теории}
\begin{frame}
\label{4.2.1.}\subsubsection{4.2.1. Определение формальной теории }\frametitle{4.2.1. Определение формальной теории (1/2)}
\vspace{-2pt}
\odmkey{Формальная теория}{Теория (theory)!формальная (formal)} $\cal T$
имеет следующие составляющие:
\begin{enumerate}
\item[1)] множество $\cal A$ символов, образующих
\odmkey{алфавит}{Алфавит (alphabet)!формальной теории (of formal
theory)}; 
\item[2)] множество $\cal F$ слов в алфавите $\cal A$,
$\cal F\subset \cal A^*$, которые называются
\odmkey{формулами}{Формула (formula)!формальной теории (of formal
theory)}; 
\item[3)] подмножество $\cal B$ формул,
 $\cal B \subset \cal F$, которые называются
\odmkey{аксиомами}{Аксиома (axiom)!формальной теории (of formal
theory)}; \item[4)] множество $\cal R$ отношений $R$ на множестве
формул, $R\in \cal R$, $R\subset \cal{F}^{n+1}$, которые
называются \odmkey{правилами вывода}{Правило (rule)!вывода
(inference)!формальной теории (of formal theory)}.
\end{enumerate}

%\pagebreak[3]
 Множество символов $\cal A$ может быть конечным или
бесконечным. 
Обычно для образования множества символов $\cal A$ используют конечное
множество букв, к которым, если нужно, приписываются в качестве
индексов натуральные числа.

Множество формул $\cal F$ обычно задаётся индуктивным
определением, например, с~помощью порождающей формальной
грамматики. 
Как правило, это множество бесконечно. Множества $\cal
A$ и $\cal F$ в совокупности определяют \odmkey{язык}{Язык
(language)!формальной теории
 (of formal theory)}, или
\odmkey{сигнатуру}{Сигнатура (signature)!формальной теории (of
formal theory)}, формальной теории.
\end{frame}


%\vspace{-0.25\baselineskip}
%\subsection{Определение формальной теории}
\begin{frame}
\frametitle{4.2.1. Определение формальной теории (2/2)}
Множество аксиом $\cal B$ может быть конечным или бесконечным.
Если множество аксиом бесконечно, то обычно оно задаётся с помощью
конечного множества \odmkey{схем}{Схема (scheme)!аксиомы (of
axiom)} аксиом и правил порождения конкретных аксиом из схемы
аксиом. 
Часто аксиомы делятся на два вида:
\odmkey{логические}{Аксиома (axiom)!логическая (logical)}
аксиомы
 (общие для целого класса
формальных теорий) и \em{нелогические} (или
\odmkey{собственные}{Аксиома (axiom)!собственная (proper)})
аксиомы (определяющие специфику и содержание конкретной теории).

Множество правил вывода $\cal R$ чаще всего конечно.
\end{frame}


%\subsection{Выводимость}
%\label{prove} 
\begin{frame}
\label{4.2.2.}\subsubsection{4.2.2. Выводимость }\frametitle{4.2.2. Выводимость (1/2)}

\vspace{-2pt}
Пусть $F_1,\dots,F_n,G$~--- формулы теории $\cal T$,
то~есть $F_1,\dots,F_n,G\in \cal F$. Если существует такое правило
вывода $R \in \cal R$, что $(F_1,\dots,F_n,G)\in R$, то
говорят, что формула $G$ \odmkey{непосредственно
выводима}{Выводимость (inference)!непосредственная (direct)} из
формул $F_1,\dots,F_n$ по правилу вывода~$R$. Обычно этот факт
записывают следующим образом:

\vspace{-10pt}
$$
\Rule{F_1,\dots,F_n}{G}{(R)}.
$$

\vspace{-2pt}
Здесь формулы $F_1,\dots,F_n$ называются
\odmkey{посылками}{Посылка правила вывода (hypothesis of rule)},
формула $G$~--- \odmkey{заключением}{Заключение правила вывода
(conclusion of inference rule)}, а~$R$ является обозначением
правила вывода.

\begin{remark}
Обозначение правила вывода справа от черты, разделяющей посылки и
заключение, часто опускают.
\end{remark}

\odmkey{Выводом}{Выводимость (inference)} формулы $G$ из формул
$F_1,\dots,F_n$ в формальной теории $\cal T$ называется такая
последовательность формул $E_1,\dots,E_k$, что $E_k=G$, и любая %а -> и
формула $E_i$ ($i\le k$) является либо аксиомой ($E_i\in \cal B$), 
либо исходной формулой $F_j$ ($E_i=F_j$), либо непосредственно
выводима из ранее полученных формул $E_{j_1},\dots,E_{j_n}$
($j_1,\dots,j_n<i$).
\end{frame}

\begin{frame}
\frametitle{4.2.2. Выводимость (2/2)}
\vspace{-2pt}
Если в теории $\cal T$ существует вывод
формулы $G$ из формул $F_1,\dots,F_n$, то это записывают следующим
образом:
$
\Seq{F_1,\dots,F_n}{G}{\cal T}.
$
Формулы $F_1,\dots,F_n$ называются \odmkey{гипотезами}{Гипотеза (hypothesis)}
вывода. Если теория $\cal T$ подразумевается, то её обозначение обычно
опускают.
Если $\Seq{}{G}{\cal T}$, то формула $G$ называется
\odmkey{теоремой}{Теорема (theorem)!формальной теории (of formal
theory)} теории $\cal T$ (то~есть теорема~--- это формула,
выводимая только из аксиом, без гипотез).
Если $\Seq{\Gamma}{G}{\cal T}$, то очевидно,
что $\Seq{\Gamma,\Delta}{G}{\cal T}$,
где $\Gamma$ и $\Delta$~--- любые множества формул (то~есть  при
добавлении лишних гипотез выводимость сохраняется).

\begin{remark}
При изучении формальных теорий нужно различать
теоремы \em{самой} формальной теории
и теоремы \em{об этой} формальной теории,
или \odmkey{метатеоремы}{Метатеорема (metatheorem)}. Это
различие иногда не формализуется явно, но всегда является существенным.
В этой главе теоремы конкретной формальной теории, как правило, записываются
в виде формул, составленных из специальных знаков, а метатеоремы
формулируются на естественном языке, чтобы их легче было отличать
от теорем самой формальной теории.
\end{remark}

\end{frame}


%%\vspace{-6pt}
%\subsection{Интерпретация}
\begin{frame}
\label{4.2.3.}\subsubsection{4.2.3. Интерпретация}
\frametitle{4.2.3. Интерпретация}
%В этом подразделе вводится более общее и строгое определение
%понятия интерпретации по сравнению с п.~\ref{inter}.
\odmkey{Интерпретацией}{Интерпретация (interpretation)!формальной
теории (of formal theory)} формальной теории $\cal T$ в
\odmkey{область интерпретации}{Область (domain)!интерпретации (of interpretation)} $M$ называется функция $\Map{I}{\cal
F}{\cal M}$, которая каждой формуле формальной теории $\cal T$
однозначно сопоставляет некоторое содержательное высказывание
относительно объектов множества (алгебраической системы) $M$. Это
высказывание может быть истинным или ложным (или не иметь
истинностного значения). Если соответствующее высказывание
является истинным, то говорят, что формула
\odmkey{выполняется}{Выполнимость (satisfiability)!формулы (of
formula)} в~данной интерпретации.

\medskip
\begin{digress}
В конечном счёте нас интересуют такие формальные теории, которые
описывают какие-то реальные объекты и связи между ними. Речь идёт прежде
всего о математических объектах и математических теориях, которые
выбираются в качестве области интерпретации.
\end{digress}

\medskip
Интерпретация $I$ называется \odmkey{моделью}{Модель
(model)!множества формул (of a set of formulas)} множества формул
$\Gamma$, если все формулы этого множества выполняются в
интерпретации $I$. Интерпретация $I$ называется
\odmkey{моделью}{Модель
 (model)!формальной теории (of formal theory)}
формальной теории
$\cal T$, если все теоремы этой теории выполняются в~интерпретации $I$.
%(то~есть все выводимые формулы оказываются истинными в~данной интерпретации).

\end{frame}


%\subsection{Общезначимость и непротиворечивость}

\begin{frame}
\label{4.2.4.}\subsubsection{4.2.4. Общезначимость и непротиворечивость}
\frametitle{4.2.4. Общезначимость и непротиворечивость}
\vspace{-2pt}
Формула называется
\odmkey{общезначимой}{Формула (formula)!общезначимая (valid)} (или
\odmkey{тавтологией}{Тавтология (tautology)}), если она истинна в любой
интерпретации. Формула называется
\odmkey{противоречивой}{Формула (formula)!противоречивая (contradictory)},
если она ложна в любой
интерпретации.
Формула $G$ называется
\odmkey{логическим следствием}{Следствие логическое (logical consequence)}
множества формул $\Gamma$, если $G$ выполняется в любой модели $\Gamma$.
Формальная теория $\cal T$ называется
\odmkey{семантически непротиворечивой}{Непротиворечивость
(consistency)!семантическая (semantic)},
если ни одна её теорема не является противоречием. Таким образом,
формальная теория пригодна для описания тех множеств (алгебраических
систем), которые являются её моделями. Модель для формальной теории
$\cal T$ существует тогда и только тогда, когда $\cal T$ семантически
непротиворечива.
Формальная теория $\cal T$ называется
\odmkey{формально непротиворечивой}{Непротиворечивость
(consistency)!формальная (formal)},
если в ней не являются выводимыми одновременно формулы $F$ и $\Not F$.
Теория $\cal T$ формально непротиворечива тогда и только тогда, когда
она семантически непротиворечива.


\begin{remark}
Последние утверждения являются доказуемыми метатеоремами,
доказательства которых опускаются.
\end{remark}

\end{frame}


\label{4.2.5.}\subsubsection{4.2.5. Полнота, независимость и разрешимость}
\begin{frame}
\frametitle{4.2.5. Полнота, независимость и разрешимость}
%\vspace{-2pt}
Пусть интерпретация  $\Map{I}{\cal T}{M}$ является моделью
формальной теории $\cal T$. Тог\-да формальная теория $\cal T$
называется \odmkey{полной}{Полнота (completeness)!формальной
теории (of formal theory)} (или \odmkey{адекватной}{Адекватность
(adequacy)}), если каждому истинному высказыванию области
интерпретации $M$ соответствует теорема теории~$\cal T$.
Если для множества (алгебраической системы) $M$ существует
формальная полная непротиворечивая теория $\cal T$, то система $M$
называется \odmkey{аксиоматизируемой}{Аксиоматизируемость
(axiomatizability)} (или \odmkey{формализуемой}{Формализуемость
(formalizability)}).

\medskip
Система аксиом (или аксиоматизация) формально непротиворечивой
теории $\cal T$ называется \odmkey{независимой}{Множество
(set)!аксиом (of axiom)!независимое (independent)}, если никакая
из аксиом не выводима из остальных по правилам вывода теории $\cal
T$.

\medskip
Формальная теория $\cal T$ называется
\odmkey{разрешимой}{Разрешимость (solvability)!формальной теории
(of formal theory)}, если существует алгоритм, который для любой
формулы теории определяет, является ли эта формула теоремой
теории. Формаль\-ная теория $\cal T$ называется
\odmkey{полуразрешимой}{Полуразрешимость формальной теории (semi-solvability of formal theory)}, если
существует алгоритм, который для любой формулы $F$ теории выдает
ответ <<ДА>>, если $F$~является теоремой теории, и, может быть, не
выдает никакого ответа, если $F$ не является теоремой
 (то~есть
алгоритм применим не ко всем формулам).
\end{frame}

\begin{frame}
\frametitle{Онтология математической логики}
\begin{center}
\includegraphics[width=12cm]{ODM4_2.png}
\end{center}
\end{frame}

\subsection{4.3. Исчисление высказываний}
\label{4.3.}

\begin{frame}
\frametitle{4.3. Исчисление высказываний (1/3)}
В этом разделе дано описание формальной теории исчисления высказываний.
Поскольку исчисление высказываний уже знакомо 
по материалам разделов~\DMref{3.1}{3.1.} и~\DMref{4.1}{4.1.},
здесь сделан сильный акцент именно на \em{формальной\/} технике.
Знакомство с чисто механическим характером
формальных выводов подготавливает рассмотрение
методов автоматического доказательства теорем в
разделе~\DMref{4.6}{4.6.}.

\end{frame}

\begin{frame}
\frametitle{4.3. Исчисление высказываний (2/3)}
\medskip
\begin{tabular}{p{10mm}|p{38mm}|p{95mm}}
  № & Название параграфа
  & Ключевые термины и обозначения \\
  \hline
\DMref{4.3.1.}{4.3.1.}\RS & Определение исчисления высказываний  &
{Связки, пропозициональные переменные, схема аксиом, правило Modus ponens (правило отделения) }\\  
\DMref{4.3.2.}{4.3.2.} & Частный случай формулы   &
{ Унификатор, совместный частный случай формулы, унифицируемая формула, наиболее общий унификатор, частный случай набора формул, совместный частный случай набора формул }\\    
\DMref{4.3.3.}{4.3.3.} & Алгоритмы унификации    &
{ Алгоритм Робинсона  }\\  
\DMref{4.3.4.}{4.3.4.} & Конструктивное определение исчисления высказываний   &
{ }\\  
\DMref{4.3.5.}{4.3.5.} & Производные правила вывода    &
{ Правило введения импликации }\\  
\DMref{4.3.6.}{4.3.6.} & Дедукция   &
{  Теорема дедукции, правило транзитивности, правило сечения  }\\  
\end{tabular}  

\end{frame}

\begin{frame}
\frametitle{4.3. Исчисление высказываний (3/3)}
\medskip
\begin{tabular}{p{10mm}|p{38mm}|p{95mm}}
  № & Название параграфа
  & Ключевые термины и обозначения \\
  \hline
\DMref{4.3.7.}{4.3.7.} & Некоторые теоремы исчисления высказываний    &
{  }\\  
\DMref{4.3.8.}{4.3.8.} & Множество теорем исчисления высказываний    &
{ }\\  
\DMref{4.3.9.}{4.3.9.} & Другие аксиоматизации исчисления высказываний    &
{ }\\  
\end{tabular}  

\end{frame}

\begin{frame}
\label{4.3.1.}\subsubsection{4.3.1. Определение исчисления высказываний }\frametitle{4.3.1. Определение исчисления высказываний (1/2)}
\odmkey{Исчисление высказываний}{Исчисление
(calculus)!высказываний (propositions calculus)}~--- это
формальная теория $\cal L$, в которой определены
\begin{tabbing}
2.2\=Формулы: $\ \ $ \= \kill
\>1.\' Алфавит:
 \> $\Not$ и $\to$ ~---  \odmkey{связки}{Связка (connective)!логическая (logical)}; \\
 \> \>$( \ , \ )$ ~---  служебные символы;\\
 \> \>$a,b,\dots,a_1,b_1,\dots$~---
    \odmkey{пропозициональные переменные}{Переменная
(variable)!пропозициональная (propositional)}.\\
\>2.\' Формулы:
 \> 1) переменные суть формулы; \\
 \> \>2) если $A$, $B$~--- формулы, \\
 \> \> \ \ \ \ то $(\Not A)$ и $(A\to B)$~--- формулы.\\
\>3.\' Аксиомы:
 \> $A_1: (A \to (B \to A))$; \\
 \> \>$A_2: ((A\to (B\to C))\to ((A \to B)\to(A \to C)))$; \\
 \> \>$A_3: ((\Not B \to \Not A)\to ((\Not B \to A)\to B))$. \\
\>4.\' Правило:
 \> $\displaystyle\Rule{A,\ A\to B}{B}{\text{(Modus ponens)}}$.
\end{tabbing}



Здесь $A$, $B$ и $C$~--- \em{любые} формулы. Таким образом,
множество аксиом теории $\cal L$ бесконечно, хотя задано тремя
\odmkey{схемами}{Схема (scheme)!аксиомы (of axiom)} аксиом.
Множество правил вывода также бесконечно, хотя оно задано только
одной \odmkey{схемой}{Схема (scheme)!правила вывода (of inference
rule)}.



\end{frame}

\begin{frame}
\frametitle{4.3.1. Определение исчисления высказываний (2/2)}

\begin{remark}
Латинское словосочетание \em{Modus ponens\/} обычно переводят на
русский как \odmkey{правило отделения}{Правило (rule)!отделения
(of detachment)} и обозначают MP.
\end{remark}

При записи формул лишние скобки опускаются, если это не вызывает  
недоразумений. 
Другие связки вводятся определениями:
$$
A \And B \Def \Not(A\to\Not B), \qquad
A \Or B \Def \Not A \to B.
$$

Любая формула, содержащая эти связки,
рассматривается как синтаксическое сокращение
собственной формулы теории
$\cal L$.

\begin{digress}
Классическое задание теории $\cal L$ с помощью схем аксиом над
любыми формулами вполне соответствует математической традиции (мы
пишем 
${(a\!+\!b)}^2=a^2\!+\!2ab\!+\!b^2$, 
подразумевая под $a$ и $b$ не
только любые числа, но и любые выражения!). В то же время это не
очень удобно для представления формальных теорий в программах. В
следующих трёх подразделах вводится определение исчисления
высказываний, более удобное для программной реализации.
\end{digress}
\end{frame}

\begin{frame}
\label{4.3.2.}\subsubsection{4.3.2. Частный случай формулы }\frametitle{4.3.2. Частный случай формулы (1/3)}

Если в формулу $A$ вместо переменных
$\Tuple{x}{1}{n}$ подставить формулы
$\Tuple{B}{1}{n}$, то получится формула $B$, которая называется
\odmkey{частным случаем}{Частный случай (instance)!формулы
(formula)} формулы $A$:
$$
B \Def A(\dots,x_i^{},\dots)\left[B_i^{}/x_i^{}\right]^{n}_{i=1}.
$$
Каждая формула $B_i$ подставляется вместо \em{всех} вхождений
переменной $x_i$. 

\smallskip
Набор подстановок
$\left[B_i^{}/x_i^{}\right]_{i=1}^n$ называется
\odmkey{унификатором}{Унификатор (unifier)}.
Формула $C$ называется \odmkey{совместным}{Частный случай
 (instance)!совместный (shared)}
частным случаем формул $A$ и $B$, если $C$ является частным случаем
формулы $A$ и одновременно частным случаем формулы~$B$ при одном и том
же наборе подстановок, то~есть
$
C=A(\dots,x_i^{},\dots)\left[X_i^{}/x_i^{}\right]^{n}_{i=1}$
и $C=B(\dots,x_i^{},\dots)\left[X_i^{}/x_i^{}\right]^{n}_{i=1}.
$

\smallskip
Формулы, которые имеют совместный частный случай,
называются \odmkey{унифицируемыми}{Формула (formula)!унифицируемая
(unificated)}, а набор
подстановок $\left[X_i^{}/x_i^{}\right]^{n}_{i=1}$,
с помощью которого получается совместный частный случай, 
называется
\odmkey{общим}{Унификатор (unifier)!общий (general)}
 унификатором. Наименьший возможный
унификатор называется \odmkey{наиболее общим
\-унификатором}{Унификатор (unifier)!наиболее общий (most
general)}.

\end{frame}

\begin{frame}
\frametitle{4.3.2. Частный случай формулы (2/3)}
%%%%%%%%%%%%%%%%%%%%%%%%% Изменения 2012
Более точно, $\sigma$ называется наиболее общим унификатором (НОУ), если
\begin{enumerate}
\item[1)] $\sigma$ является унификатором; 
\item[2)] для любого другого унификатора $\delta$ существует 
такая подстановка $\lambda$, 
что унификатор $\delta$ выражается через 
суперпозицию $\sigma$ и $\lambda$, $\delta = \sigma \circ \lambda$
(см.~пример 2 в \DMref{п.~2.2.3}{2.2.3.}).

\end{enumerate}

\smallskip
Набор формул $\Tuple{B}{1}{n}$ называется \odmkey{частным случаем
набора формул}{Частный случай (instance)!набора формул (set of
formulas)} $\Tuple{A}{1}{n}$, если каждая формула $B_i$ является
частным случаем формулы $A_i$ при одном и~том же наборе
подстановок. Набор формул $\Tuple{C}{1}{n}$ называется
\odmkey{совместным}{Частный случай (instance)!наборов формул (set
of formulas)!совместный (shared)} частным случаем наборов формул
$\Tuple{A}{1}{n}$ и $\Tuple{B}{1}{n}$, если каждая формула $C_i$
является частным случаем формул $A_i$ и $B_i$ при одном и том же
наборе подстановок.
\end{frame}


%\subsection{Частный случай формулы}
\begin{frame}
\frametitle{4.3.2. Частный случай формулы (3/3)}

\begin{examples}
\begin{enumerate}
\item Формулы: $\Not a \to y, x \to \Not b$.
НОУ $[\Not a/x, \Not b/y]$.
Совместный частный случай: \\$\Not a \to \Not b$. 
\item Формулы: $\Not a \to y, x \to \Not b$.
Унификатор $[z/a, \Not w/y, \Not z/x, w/b]$.
Совместный частный случай $\Not z \to \Not w$, %общий -> совместный
но при этом
$[z/a, \Not w/y, \Not z/x, w/b] =\nobreak
[\Not a/x, \Not b/y] {\circ}[z/a, w/b]$. %b//y -> b/y
\item Формулы: $\Not a \to x, x \to \Not b$. Формулы не унифицируемы, так как нужно одновременно $[\Not a/x]$ и $[\Not b/x]$, что невозможно.
\item Формулы: $x, \Not x$, а также $x, x\to y$ или $x, y\to x$ и тому подобные не унифицируемы, потому что конечное дерево не может быть изоморфно своему собственному поддереву.
\end{enumerate}
\end{examples}

\begin{remark}
Хотя подстановка вида $[f(x)/x]$ синтаксически допустима, она никогда не может появиться в унификаторе, поскольку, как уже отмечено в примере, конечное дерево не может быть
изоморфно своему собственному поддереву, а потому такие подстановки можно не рассматривать.
\end{remark}
%%%%%%%%%%%%%%%%%%%%%%%%% Конец изменений 2012


\end{frame}


%%%%%%%%%%%%%%%%%%%%%%%%% Изменения 2012
%\subsection{Алгоритмы унификации}
%\label{unify}
%{\doublehyphendemerits=100000
%\index{Алгоритм (algorithm)!унификации (unification)}
\begin{frame}
\label{4.3.3.}\subsubsection{4.3.3. Алгоритмы унификации }\frametitle{4.3.3. Алгоритмы унификации (1/6)}
Сначала рассмотрим общий \odmkey{алгоритм унификации}{Алгоритм (algorithm)!унификации (unification)},
предложенный Робинсоном вместе с методом резолюций (см. раздел \DMref{4.6}{4.6.}).
В этом алгоритме унифицируются два терма (см. \DMref{п.~2.1.3}{2.1.3.}).
При этом ничего не предполагается относительно сигнатуры функциональных символов, входящих в термы. Предполагается только, что можно отличить
функциональные символы от символов переменных и~символов констант. Также считается, что задана функция $\firstNeqSubTerms$, которая отыскивает в заданных термах самые левые вхождения несовпадающих подтермов.

\medskip
\begin{remark}
Знак $\sim$ в следующем алгоритме означает, что та 
\em{пара} подтермов, которую
возвращает функция $\firstNeqSubTerms$, образует подстановку, то есть
состоит из переменной $v$ и некоторого терма $t$. При этом дополнительно
накладывается условие, что $v \not\in t$
(см. замечание на предыдущем слайде).
\end{remark}
\end{frame}


\begin{frame}
\frametitle{4.3.3. Алгоритмы унификации (2/6)}

\vspace{-6pt}
%\begin{algorithm}
%\caption{Оригинальный алгоритм Робинсона} \label{robisonalg}
\begin{algorithmic}
\REQUIRE формулы $A$ и $B$ в термальной форме.
\ENSURE \textbf{false}, если формулы не унифицируемы;
\textbf{true} и наиболее общий унификатор $\sigma$ в~противном случае.
\STATE $\sigma \Assign \emptyset$
\WHILE{$A \ne B$ }
  \STATE $[A_1, B_1] \Assign\firstNeqSubTerms (A, B)$
  \IF{ $\Not\left([A_1, B_1] \sim [v, t] \And v\not\in t\right)$}
    \STATE \textbf{return false} 
    \COMMENT{\color{gray}\small формулы не унифицируемы}
  \ENDIF
  \STATE $\sigma\Assign\sigma\circ[t/v]; A\Assign A[t/v]; B\Assign B[t/v]$
\ENDWHILE
\STATE \textbf{return true}
\end{algorithmic}
%\end{algorithm}

\vspace{-4pt}

%%%%%%%%%%%%%%%%%%%%%%%%% Конец изменений 2012
\end{frame}


\begin{frame}
\frametitle{4.3.3. Алгоритмы унификации (3/6)}

\begin{example}
В примере рассматривается унификация двух термов, при этом несовпадающие подтермы, найденные функцией firstNeqSubTerms, подчеркиваются.

\vspace{-2pt}
\begin{tabbing}
$P(a, f(a), f(g(y)))$\quad\= $P(a, f(a), f(g(y)))$\quad\= термы совпадают\kill
$P(\underline{a}, x, f(g(y)))$
    \> $P(\underline{z}, f(z), f(u))$
    \> $\sigma = [a/z]$ \\
$P(a, \underline{x}, f(g(y)))$
    \> $P(a, \underline{f(a)}, f(u))$
    \> $\sigma = [a/z, f(a)/x]$ \\
    $P(a, f(a), f(\underline{g(y)}))$ \> %\underline{f(g(y))} -> f(\underline{g(y)})
    $P(a, f(a), f(\underline{u}))$\>
    $\sigma = [a/z, f(a)/x, g(y)/u]$\\
$P(a, f(a), f(g(y)))$ \> $P(a, f(a), f(g(y)))$ \> термы совпадают
\end{tabbing}
\end{example}

%\medskip
Если язык формул
удовлетворяет определённым условиям,
то можно проверить, являются ли две заданные формулы
унифицируемыми, и если это так, то~найти
наиболее общий унификатор. В частности, это возможно для
типичного языка формул, описанного в \DMref{п.~4.3.1}{4.3.1.}.

\smallskip
В следующем рекурсивном алгоритме $\Unify$ предполагается, 
что функция $f$ осуществляет
синтаксический разбор формулы и возвращает знак главной операции
(то~есть $\to$, $\neg$ или признак того, что формула является
переменной),
а~функции $l$ и $r$ возвращают левый и правый операнды
главной операции, то~есть подформулы (считаем, что отрицание имеет
только правый операнд).
\end{frame}

\begin{frame}
\frametitle{4.3.3. Алгоритмы унификации (4/6)}
 Набор подстановок (унификатор) представлен
в виде глобального массива $S$, значениями индекса которого являются имена
переменных, а~значениями элементов~--- формулы, которые
подставляются вместо соответствующих переменных.
\begin{algorithmic}
\REQUIRE
формулы $A$ и $B$ исчиcления высказываний.
\ENSURE $0$, если формулы не унифицируемы;
$1$ и наиболее общий
унификатор $S$ в~противном случае.
%%%%%%%%%%%%%%%%%%%%%%%%% Изменения 2012
    \STATE {$a\Assign f(A); b\Assign f(B)$}
    \COMMENT{\color{gray}\small $a$ и  $b$~--- главные операции }
\STATE{} 
    \COMMENT{\color{gray}\small  обе формулы --- переменные }
\IF{$\operatorname{var}(a)\And\operatorname{var}(b)$}
   \STATE \textbf{if $a=b$ then return $1$ end if}
    \STATE \textbf{if $S[a]= \emptyset \And S[b] = \emptyset$ then
    $S[a] \Assign b;$ return $1$ end if}
    \STATE \textbf{if $S[a]=\emptyset \And S[b] \ne \emptyset$ then
    $S[a]\Assign S[b]$; return $1$ end if}
    \STATE \textbf{if $S[b]=\emptyset \And S[a] \ne \emptyset$ then
    $S[b]\Assign S[a]$; return $1$ end if}
    \STATE \textbf{if $S[a]\ne \emptyset \And S[b] \ne \emptyset$ then
    return $S[a]=S[b]$ end if}
\ENDIF
\end{algorithmic}

\end{frame}

\begin{frame}
\frametitle{4.3.3. Алгоритмы унификации (5/6)}
\begin{algorithmic}
\STATE\COMMENT{\color{gray}\small   одна формула --- переменная, а другая --- нет }
\IF{$\operatorname{var}(a)\Or\operatorname{var}(b)$}
  \IF{$\operatorname{var}(a) \And a\in B$}
         \STATE \textbf{return $0$}
  \ELSE
        \STATE \textbf{if $S[a] = \emptyset$ then $S[a] \Assign B$; return $1$}
        \STATE \textbf{else return $(S[a] = B)$ end if}
  \ENDIF
  \IF{$\operatorname{var}(b) \And b\in A$}
         \STATE \textbf{return $0$}
  \ELSE
    \STATE \textbf{if $S[b] = \emptyset$ then $S[b] \Assign A$; return $1$}
    \STATE \textbf{else return $(S[b] = A)$ end if}
  \ENDIF
\ENDIF
\end{algorithmic}


\end{frame}

\begin{frame}
\frametitle{4.3.3. Алгоритмы унификации (6/6)}
\begin{algorithmic}
\STATE\COMMENT{\color{gray}\small   обе формулы не переменные } 
\STATE{\textbf{if} $a\ne b$ \textbf{then return $0$ end if}}
\COMMENT{\color{gray}\small   главные операции различны }
\STATE{\textbf{if} $ a = \Not$ \textbf{then return} $\Unify(r(A), r(B))$ 
\textbf{end if}}
\STATE{\textbf{if} $ a = \to$ \textbf{then return} 
$\Unify(l(A), l(B)) \And \Unify(r(A), r(B))$ \textbf{end if}}
\end{algorithmic}

\medskip
\begin{ground}
При любых подстановках формул вместо переменных главная операция
(связка) формулы остаётся неизменной. Поэтому если главные
операции формул различны, то формулы заведомо не унифицируемы. В
противном случае, то~есть если главные операции совпадают, формулы
унифицируемы тогда и только тогда, когда унифицируемы подформулы,
являющиеся операндами главной операции. Эта рекурсия
заканчивается, когда сопоставление доходит до переменных.
Переменная унифицируется с любой формулой (в частности, с другой
переменной) простой подстановкой этой формулы вместо переменной.
Но подстановки для всех вхождений одной переменной должны
совпадать.
\end{ground}
\end{frame}



%\subsection{Конструктивное определение исчисления высказываний}


\begin{frame}
\label{4.3.4.}\subsubsection{4.3.4. Конструктивное определение исчисления высказываний}
\frametitle{4.3.4. Конструктивное определение исчисления высказываний}
Алфавит
и множество формул~--- те же (см. \DMref{п.~4.3.1}{4.3.1.}).
Аксиомы~--- три конкретные формулы:

%{
%\settowidth{\leftmargini}{$A_{1} \colon$}
%\addtolength{\leftmargini}{\labelsep}
\begin{enumerate}
%\itemsep 0pt
%\parsep 0pt
%\parskip 1pt plus 1pt
\item [$A_{1} \colon$] $(a \to (b \to a))$;
\item [$A_{2} \colon$] $((a\to (b\to c))\to ((a \to b)\to(a \to c)))$;
\item [$A_{3} \colon$] $((\Not b \to \Not a)\to ((\Not b \to a)\to b))$.
\end{enumerate}



Правила вывода:

\begin{enumerate}
\item Правило подстановки: если формула $B$ является частным случаем
формулы~$A$, то $B$ непосредственно выводима из $A$.



\item Правило Modus ponens: если набор формул $A,B,C$ является частным
случаем набора формул $a$, $(a\to b)$, $b$, то формула $C$ является
непосредственно выводимой из формул $A$ и $B$.
\end{enumerate}

%\vspace*{-.5\baselineskip}

\begin{remark}
Здесь $a$, $(a\to b)$, $b$~--- это три конкретные формулы, построенные с
помощью переменных $a$, $b$ и связки $\to$.
\end{remark}
\end{frame}


%\subsection{Производные правила вывода}
%\label{drules}

\begin{frame}
\label{4.3.5.}\subsubsection{4.3.5. Производные правила вывода }\frametitle{4.3.5. Производные правила вывода (1/3)}

%\vspace{-4pt}
Исчисление высказываний $\cal L$~--- это достаточно
богатая формальная теория, в~которой выводимы
многие важные теоремы.

\medskip
\begin{remark}
Выводимость формул в теории $\cal L$ доказывается путём
предъявления конкретного вывода, то~есть последовательности
формул, удовлетворяющих определению, данному в~\DMref{п.~4.2.2}{4.2.2.}. 
Для удобства чтения формулы
последовательности вывода выписываются друг под другом в столбик,
слева указываются их номера в последовательности, а справа
указывается, на каком основании формула включена в вывод (то~есть
является ли она гипотезой, или получена из схемы аксиом указанной
подстановкой, или получена из предшествующих формул по указанному
правилу вывода и~т.~д.).
\end{remark}

\medskip
Всякую доказанную выводимость можно использовать как новое
(\odmkey{производное}{Правило (rule)!вывода
(inference)!производное (derived)}) правило вывода. Например,
следующая доказанная выводимость называется \odmkey{правилом
введения импликации}{Правило (rule)!введения импликации (introduction of implication)}:
%\vspace{-4pt}
$$
\Rule{A}{B\to A}{(\to^+)}.
$$

%\vspace*{-.5\baselineskip}
%
%\setcounter{theorem}{0}

\end{frame}

\begin{frame}
\frametitle{4.3.5. Производные правила вывода (2/3)}

%\vspace{-2pt}
\begin{theorem} \NB{[1]} $\Seq{A}{B\to A}{\cal L}$.
\end{theorem}

%\vspace{-2pt}
\begin{proof} Вывод:
\vspace{-0.2\baselineskip}
\begin{tabbing}
2.2\=$ A\to (B\to A)\quad $ \= \hfill \=\kill
\>1.\' $A$ \> гипотеза. \\
\>2.\' $A\to (B\to A)$\> $A_1$. \\
\>3.\' $B\to A$\> МР; 1, 2.\`\qed
\end{tabbing}
\renewcommand{\qed}{\vspace{-\baselineskip}}\end{proof}

\begin{remark}
Малое количество аксиом и правил вывода
не влечёт лаконичности выводов в формальной теории.
Например, в аксиоматике исчисления высказываний 
всего три схемы аксиом и два правила вывода,
однако даже простые теоремы требуют 
длинных выводов.
\end{remark}
\end{frame}

\begin{frame}
\frametitle{4.3.5. Производные правила вывода (3/3)}


\begin{theorem} \NB{[2]} $\Seq{}{A\to A}{\cal L}$.
\end{theorem}

%\vspace{-2pt}
\begin{proof} Вывод:
%\vspace{-0.2\baselineskip}
\begin{tabbing}
2.2\= $((A\to(A\to A))\to (A\to A))\quad  $ \= \kill
1. \> $(A\to((A\to A)\to A))$
 \> $A_1; [A\to A/B]$.  \\
2. \>$((A\to((A\to A)\to A))\to$ \> \\
\> $((A\to (A\to A))\to (A\to A)))$
 \> $A_2; [A\to A/B, A/C]$.  \\
3.\> $((A\to(A\to A))\to (A\to A))$
 \> MР; 1, 2. \\
4. \>\' $A\to (A\to A)$
 \> $A_1; [A/B]$. \\
5.\> $A\to A$
 \> МР; 4, 3.\`\qed
\end{tabbing}
\renewcommand{\qed}{\vspace{-\baselineskip}}\end{proof}

\end{frame}

%\subsection{Дедукция}
\begin{frame}
\label{4.3.6.}\subsubsection{4.3.6. Дедукция }\frametitle{4.3.6. Дедукция (1/4)}
\vspace{-4pt}
В теории $\cal L$ импликация очень тесно связана с выводимостью.

\begin{theorem} \NB{[Дедукции]}  $\Seq{\Gamma,A}{B}{\cal L}$
если и только если
 $\Seq{\Gamma}{A\to B}{\cal L}$.
\end{theorem}
\begin{proof}
\proofsection{Необходимость}
Пусть $\Tuple{E}{1}{n}$~--- вывод $B$ из ${\Gamma},A$. $E_n=B$.\\
Индукцией по $i$ $(1\le i\le n)$ покажем, что $\Seq{\Gamma}{A\to E_i}{\cal L}$.
Тем самым теорема будет доказана.
База: $i=1$. Возможны три случая.

\vspace{-2pt}
\begin{enumerate}
\item Пусть $E_1$~--- аксиома. Тогда рассмотрим
вывод $E_1$, $E_1\to (A\to E_1)$, $A\to E_1$. Имеем
$\Seq{}{A\to E_1}{\cal L}$.
\item Пусть  $E_1\in \Gamma$. Тогда рассмотрим вывод
$E_1$, $E_1\to (A\to E_1)$, $A\to E_1$.
Имеем $\Seq{\Gamma}{A\to E_1}{\cal L}$.
\item Пусть $E_1=A$. Тогда
по второй теореме предыдущего параграфа $\Seq{}{E_1\to E_1}{\cal L}$,
а~значит, $\Seq{}{A\to E_1}{\cal L}$.
\end{enumerate}
В любом случае $\Seq{\Gamma}{A\to E_1}{\cal L}$. Таким образом,
база индукции доказана. Пусть теперь $\Seq{\Gamma}{A\to E_i}{\cal
L}$ для всех $i<k$. Рассмотрим $E_k$. Возможны четыре случая: либо
$E_k$~--- аксиома, либо $E_k\in \Gamma$, либо $E_k=A$, либо
формула $E_k$ получена по правилу Modus ponens из формул $E_i$ и
$E_j$, причём $i,j<k$ и $E_j=E_i\to E_k$. 
\end{proof}
\end{frame}

\begin{frame}
\frametitle{4.3.6. Дедукция (2/4)}
\begin{proof} (Продолжение)
Для первых трёх случаев
имеем $\Seq{\Gamma}{A\to E_k}{\cal L}$ аналогичным образом. Для
четвёртого случая по индукционному предположению имеется вывод
$\Seq{\Gamma}{A\to E_i}{\cal L}$  и вывод
$\Seq{\Gamma}{A\to(E_i\to E_k)}{\cal L}$. Объединим эти выводы и
достроим следующий вывод:
\vspace{-6pt}
\begin{tabbing}
$n+1.111$\= $(A\!\to\! (E_i\!\to\! E_k))\!\to\! ((A\!\to\! E_i)\!\to\! (A\!\to\! E_k))\quad$\= \kill
$n.$\>   $(A\!\to\! (E_i\!\to\! E_k))\!\to\! ((A\!\to\! E_i)\!\to\! (A\!\to\! E_k))$\> 
$A_2;[E_i/B, E_k/C]$. \\
$n+1.$\> $(A\!\to\! E_i)\!\to\! (A\!\to\! E_k)$ \> МР; $j$, $n$. \\
$n+2.$\> $A\!\to\! E_k$ \> MP; $i$, $n+1$.
\end{tabbing}

Таким образом, $\Seq{\Gamma}{A\to E_k}{\cal L}$ для любого $k$, в
частности для $k=n$. Но $E_n=B$, то~есть $\Seq{\Gamma}{A\to
B}{\cal L}$.

\proofsection{Достаточность} Имеем вывод $\Seq{\Gamma}{A\to
B}{\cal L}$, состоящий из $n$ формул. Достроим его следующим
образом:

\vspace{-6pt}
\begin{tabbing}
$n+1.111$\= $A \to B$\quad\= \kill
$n.$ \> $A \to B$. \\
$n+1.$ \> $A$ \> гипотеза. \\
$n+2.$ \> $B$ \> MP; $n+1$, $n$.
\end{tabbing}

\vspace{-6pt}
Таким образом, имеем $\Seq{\Gamma,A}{B}{\cal L}$.
\end{proof}
\end{frame}

\begin{frame}
\frametitle{4.3.6. Дедукция (3/4)}

%\begin{remark}
Схема аксиом $A_3$ теории $\cal L$ не использовалась в доказательстве,
поэтому теорема дедукции имеет место не только для теории
$\cal L$, но и для любых теорий, в которых выполнены аксиомы $A_1$
и $A_2$ и действует правило отделения.
%\end{remark}



%\setcounter{varcorollary}{0}
%\vspace{-6pt}
\begin{corollary} \NB{[1]}
$\Seq{A}{B}{\cal L}$ тогда и только тогда,
когда $\Seq{}{A\to B}{\cal L}$.
\end{corollary}

%\vspace{-18pt}
\begin{proof}
$\Gamma\Assign\emptyset$.
\end{proof}

%\vspace{\baselineskip}

%\vspace{-12pt}
\begin{corollary} \NB{[2]}
$\Seq{A\to B, B\to C}{A\to C}{\cal L}$.
\end{corollary}

%\vspace{18pt}
\begin{proof} Вывод:

\begin{tabbing}
7.7\= $\Seq{A\to (B\to C),B}{A\to C}{\cal L} \quad$ \= \kill
\>1.\' $A\to B$ \> гипотеза. \\%[3pt]
\>2.\' $B\to C$ \> гипотеза. \\%[3pt]
\>3.\' $A     $ \> гипотеза. \\%[3pt]
\>4.\' $B     $ \> MP; 3, 1. \\%[3pt]
\>5.\' $C     $ \> MP; 4, 2. \\%[3pt]
\>6.\' $A\to B, B\to C, \Seq{A}{C}{\cal L}$ \> 1--5. \\%[3pt]
\>7.\' $\Seq{A\to B, B\to C}{A\to C}{\cal L}$ \> по теореме дедукции.
\end{tabbing}
\vspace{-24pt}
\end{proof}


\end{frame}

\begin{frame}
\frametitle{4.3.6. Дедукция (4/4)}

\begin{remark}
Это производное правило называется \odmkey{правилом
транзитивности}{Правило (rule)!транзитивности (of transitivity)}.
\end{remark}

\smallskip
\begin{corollary} \NB{[3]}
$\Seq{A\to(B\to C), B}{A\to C}{\cal L}$.
\end{corollary}


\begin{proof} Вывод:
\begin{tabbing}
7.7\= $\Seq{A\to (B\to C),B}{A\to C}{\cal L} \quad$ \= \kill
\>1.\' $A\to (B\to C)$ \> гипотеза. \\[2pt]
\>2.\' $A$ \> гипотеза. \\[2pt]
\>3.\' $B\to C$ \> MP; 2, 1. \\[2pt]
\>4.\' $B$ \> гипотеза. \\[2pt]
\>5.\' $C$ \> MP; 4, 3. \\[2pt]
\>6.\' $\Seq{A\to (B\to C),B,A}{C}{\cal L}$ \> 1--5. \\[2pt]
\>7.\' $\Seq{A\to (B\to C),B}{A\to C}{\cal L}$ \> по теореме дедукции.
\end{tabbing}
\vspace{-\baselineskip}
\end{proof}

\begin{remark}
Это производное правило называется
\odmkey{правилом сечения}{Правило (rule)!сечения (cut)}.
\end{remark}
\end{frame}

\begin{frame}
\label{4.3.7.}\subsubsection{4.3.7. Некоторые теоремы исчисления высказываний }\frametitle{4.3.7. Некоторые теоремы исчисления высказываний (1/7)}
Множество теорем теории $\cal L$ бесконечно. Здесь приведены
некоторые теоремы, которые используются в дальнейших построениях.



%\vspace{-\baselineskip}
\begin{remark}
Каждая доказанная теорема (то~есть формула теории, для которой
построен вывод) может использоваться в дальнейших выводах
на правах аксиомы.
\end{remark}


\medskip
\begin{theorem} В теории ${\cal L}$ выводимы следующие теоремы:
{
\settowidth{\leftmargini}{\it{ж})}
\addtolength{\leftmargini}{\labelsep}
\begin{enumerate}
\itemsep 0pt
\parsep 0pt
\parskip 1pt plus 1pt
\item[a)]  $\Seq{}{\Not\Not A\to A}{\cal L}$.%
\item[б)]  $\Seq{}{A \to \Not\Not A}{\cal L}$.%
\item[в)]  $\Seq{}{\Not A \to (A\to B)}{\cal L}$.%
\item[г)]  $\Seq{}{(\Not B\to \Not A)\to(A\to B)}{\cal L}$.%
\item[д)]  $\Seq{}{(A\to B)\to(\Not B \to \Not A)}{\cal L}$.%
\item[е)]  $\Seq{}{A\to(\Not B\to \Not(A\to B))}{\cal L}$.%
\item[ж)]  $\Seq{}{(A\to B)\to ((\neg A\to B)\to B)}{\cal L}$.%
\end{enumerate}
}
\end{theorem}
%
%\begin{proof}
%Приведено в бумажной версии учебника.
%\end{proof}


\end{frame}

\begin{frame}
\frametitle{4.3.7. Некоторые теоремы исчисления высказываний (2/7)}


\begin{proof}
\proofsection{a) $\Seq{}{\Not\Not A\to A}{\cal L}$}

\begin{tabular}
{r  p{68mm} l}
%\No & Утверждение & Обоснование\\
1.&  $(\neg A\to\neg\neg A)\to ((\neg A\to\neg A)\to A)$ & $A_3;
     [A/B, \neg A/A]$.\\
2.&  $\neg A\to\neg A$ & вторая теорема 
\DMref{п.~4.3.5}{4.3.5.};  $[\neg A/A]$. \\
3.&  $(\neg A\to\neg\neg A)\to A$ &
     Сл.~3; 
$[\neg A\to\neg\neg A/A, \neg A\to\neg A/B, A/C]$.  \\
4.&  $\neg\neg A\to (\neg A\to\neg\neg A)$ & $A_1;
    [\neg\neg A/A, \neg A/B]$.\\
5.&  $\neg\neg A\to A$ & Сл.~2;  $[\neg\neg A/A, \neg
A\to\neg\neg A/B, A/C]$.
\end{tabular}

\medskip
\proofsection{б) $\Seq{}{A \to \Not\Not A}{\cal L}$}

\begin{tabular}{r  p{68mm} l}
%\No & Утверждение & Обоснование\\
1.&  $(\neg\neg\neg A\to \neg A)\to ((\neg\neg\neg A\to A)\to\neg\neg A)$ & $A_3; [\neg\neg A/B]$.\\
2.&  $\neg\neg\neg A\to \neg A$ & {\it а};  $[\neg A/A]$.\\
3.&  $(\neg\neg\neg A\to A)\to\neg\neg A$ & МР; 2, 1.\\
4.&  $A\to (\neg\neg\neg A\to A)$ & $A_1; [\neg\neg\neg A/B]$.\\
5.&  $A\to\neg\neg A$ & Сл.~2; $[\neg\neg\neg A/B, A/A, \neg\neg
A/C]$.
\end{tabular}

\end{proof}
\end{frame}

\begin{frame}
\frametitle{4.3.7. Некоторые теоремы исчисления высказываний (3/7)}


\begin{proof}  (Продолжение)

\proofsection{в) $\Seq{}{\Not A \to (A\to B)}{\cal L}$}

\begin{tabular}{r p{68mm} l}
%\No & Утверждение & Обоснование\\
1. & $\neg A$ & гипотеза. \\
2. & $A$ & гипотеза.\\
3. & $A\to (\neg B\to A)$ & $A_1; [\neg B/B]$.  \\
4. & $\neg A\to (\neg B\to\neg A)$ & $A_1; [\neg A/A, \neg B/B]$. \\
5. & $\neg B\to A$ & МР; 2, 3. \\
6. & $\neg B\to \neg A$ & МР; 1, 4. \\
7. & $(\neg B\to \neg A)\to ((\neg B\to A)\to B)$ & $A_3$.  \\
8. & $(\neg B\to A)\to B$ & МР; 6, 7. \\
9. & $B$ & МР; 5, 8. \\
10. & $\Seq{\neg A, A}{B}{\cal L}$ & 1--9.  \\
11. & $\Seq{\neg A}{A\to B}{\cal L}$ & теорема дедукции. \\
12. & $\Seq{}{\neg A\to (A\to B)}{\cal L}$ & теорема дедукции.
\end{tabular}

\end{proof}
\end{frame}

\begin{frame}
\frametitle{4.3.7. Некоторые теоремы исчисления высказываний (4/7)}

\begin{proof}  (Продолжение)

\proofsection{г) $\Seq{}{(\Not B\to \Not A)\to(A\to B)}{\cal L}$}


\begin{tabular}{r p{68mm} l}
%\No & Утверждение & Обоснование\\
1.&  $\neg B\to \neg A$  & гипотеза. \\
2.&  $A$ & гипотеза.\\
3.&  $(\neg B\to \neg A)\to ((\neg B\to A)\to B)$ & $A_3$.\\
4.&  $A\to (\neg B\to A)$ & $A_1;[\neg B/B]$. \\
5.&  $(\neg B\to A)\to B$ & МР; 1, 3.\\
6.&  $A\to B$ & Сл.~2;  $[\neg B\to A/B, B/C]$. \\
7.&  $B$ & MP; 2, 6. \\
8.&  $\Seq{\neg B\to\neg A, A}{B}{\cal L}$ & 1--7. \\
9.&  $\Seq{\neg B\to\neg A}{A\to B}{\cal L}$ & теорема дедукции. \\
10.&  $\Seq{}{(\neg B\to\neg A)\to (A\to B)}{\cal L}$ & теорема
дедукции.
\end{tabular}

\end{proof}
\end{frame}


\begin{frame}
\frametitle{4.3.7. Некоторые теоремы исчисления высказываний (5/7)}

\begin{proof}  (Продолжение)


\proofsection{д) $\Seq{}{(A\to B)\to(\Not B \to \Not A)}{\cal L}$}



\begin{tabular}{r  p{68mm} l}
%\No & Утверждение & Обоснование\\
1.&  $A\to B$ & гипотеза. \\
2.&  $\neg \neg A\to A$ & {\it а}. \\
3.&  $\neg\neg A\to B$ & Сл.~2;  $[A/B, \neg\neg A/A, B/C]$. \\
4.&  $B\to\neg\neg B$ & {\it б};  $[B/A]$.\\
5.&  $\neg\neg A\to\neg\neg B$ & Сл.~2;  $[\neg\neg A/A, \neg\neg B/C]$.\\
6.&  $(\neg\neg A\to\neg\neg B)\to (\neg B\to\neg A)$ & {\it г};  $[\neg A/B, \neg B/A]$.\\
7.&  $\neg B\to\neg A$ & МР; 5, 6.  \\
8.&  $\Seq{A\to B}{\neg B\to\neg A}{\cal L}$ & 1--7. \\
9.&  $\Seq{}{(A\to B)\to (\neg B\to \neg A)}{\cal L}$ & теорема дедукции.
\end{tabular}

\end{proof}
\end{frame}


\begin{frame}
\frametitle{4.3.7. Некоторые теоремы исчисления высказываний (6/7)}

\begin{proof}  (Продолжение)
\proofsection{е) $\Seq{}{A\to(\Not B\to \Not(A\to B))}{\cal L}$}


\begin{tabular}{r  p{55mm} l}
%\No & Утверждение & Обоснование\\
1.&  $A$ & гипотеза. \\
2.&  $A\to B$ & гипотеза. \\
3.&  $B$ & МР; 1, 2.  \\
4.&  $\Seq{A, A\to B}{B}{\cal L}$ & 1--3.  \\
5.&  $\Seq{A}{(A\to B)\to B}{\cal L}$ & теорема дедукции.\\
6.&  $\Seq{}{A\to((A\to B)\to B)}{\cal L}$ & теорема дедукции.\\
7.&  $\Seq{}{((A\to B)\to B)\to(\neg B\to\neg(A\to B))}{\cal L}$ & {\it д};  $[A\to B/A]$. \\
8.&  $\Seq{}{A\to(\Not B\to \Not(A\to B))}{\cal L}$ & Сл.~2; $[((A\to B)\to B)/ B(\neg B\to  \neg(A\to B))/C].$
\end{tabular}

\end{proof}
\end{frame}

\begin{frame}
\frametitle{4.3.7. Некоторые теоремы исчисления высказываний (7/7)}

\begin{proof}  (окончание)


\proofsection{ж) $\Seq{}{(A\to B)\to ((\neg A\to B)\to B)}{\cal L}$}

\begin{tabular}{r  p{68mm} l}
%\No & Утверждение & Обоснование\\
1.&  $A\to B$ & гипотеза. \\
2.&  $\neg A\to B$ & гипотеза. \\
3.&  $(A\to B)\to (\neg B\to\neg A)$ & {\it д}.  \\
4.&  $\neg B\to\neg A$ & МР; 1, 3.\\
5.&  $(\neg A\to B)\to (\neg B\to\neg\neg A)$ & {\it д};  $[\neg A/A]$.\\
6.&  $\neg B\to\neg\neg A$ & МР; 2, 5. \\
7.&  $(\neg B\to\neg\neg A)\to ((\neg B\to \neg A)\to B)$ & $A_3; [\neg A/A].$ \\
8.&  $(\neg B\to \neg A)\to B$ & МР; 6, 7. \\
9.&  $B$ & МР; 4, 8. \\
10.&  $\Seq{ A\to B, \neg A\to B}{B}{\cal L}$ & 1--9.\\
11.&  $\Seq{A\to B}{(\neg A\to B)\to B}{\cal L}$ & теорема дедукции. \\
12.&  $\Seq{}{(A\to B)\to ((\neg A\to B)\to B)}{\cal L}$ &
теорема дедукции.
\end{tabular} 

\end{proof}
\end{frame}


%\subsection{Множество теорем исчисления высказываний}
%\label{Lcomplete}
\begin{frame}
\label{4.3.8.}\subsubsection{4.3.8. Множество теорем исчисления высказываний }\frametitle{4.3.8. Множество теорем исчисления высказываний (1/3)}
Пусть формула $A$ содержит
переменные $\Tuple{a}{1}{n}$, и пусть задана некоторая
    интерпретация $I$
формулы $A$, то~есть переменным $\Tuple{a}{1}{n}$
приписаны истинностные значения.
Обозначим
$A'_i \Assign
\textbf{ if } a_i=\text{T} \textbf{ then } 
a_i \textbf{ else }
\neg a_i \textbf{ end if}$,\\
$A' \Assign \textbf{ if } I(A)= \text{T}
   \textbf{ then } A \textbf{ else } \neg A \textbf{ end if}$
в данной интерпретации.


\medskip
\begin{lemma}
$\Seq{\Tuple{A'}{1}{n}}{A'}{\cal L}$.
\end{lemma}



\begin{proof}
Индукция по структуре формулы $A$.

\proofsection{\!Переменная\!} Пусть $A=a$. Тогда $\Seq{a}{a}{\cal
L}$ и $\Seq{\neg a}{\neg a}{\cal L}$.

\proofsection{\!Отрицание\!} Пусть $A=\neg B$ и $I$~---
произвольная интерпретация формулы $A$. Возможны два случая:
$I(B)=\text{T}$ или $I(B)=\text{F}$. Пусть $I(B)=\text{T}$. Тогда
$I(A)=\text{F}$ и $A'=\neg A=\neg\neg B$. По индукционному
предположению $\Seq{A'_1,\dots,A'_n}{B}{\cal L}$. Но
$\Seq{}{B\to\neg\neg B}{\cal L}$ по теореме 
\DMref{п.~4.3.7}{4.3.7.}~{\it б}, 
следовательно, $\Seq{A'_1,\dots,A'_n}{\neg\neg B}{\cal L}=A'$.
Пусть теперь $I(B)=\text{F}$. Тогда \\$I(A)=\text{T}$ и $A'=A=\neg
B$. По индукционному предположению $\Seq{A'_1,\dots,A'_n}{\neg
B}{\cal L}=A'$.
\end{proof}
\end{frame}

\begin{frame}
\frametitle{4.3.8. Множество теорем исчисления высказываний (2/3)}

\vspace{-2pt}
\begin{proof} (Продолжение)
\proofsection{\!Импликация\!} Пусть $A=(B\to C)$ и $I$~---
произвольная интерпретация формулы $A$.
 По индукционному предположению
$\Seq{\Tuple{A'}{1}{n}}{B'}{\cal L}$ и
$\Seq{\Tuple{A'}{1}{n}}{C'}{\cal L}$.
Возможны три случая: $I(B)=\text{F}$,
или $I(B)=\text{T}$ и $I(C)=\text{T}$, или же
$I(B)=\text{T}$ и~$I(C)=\text{F}$.

Пусть $I(B)=\text{F}$. Тогда независимо от значения $I(C)$ имеем:
$I(A)=\text{T}$  и $B' =\neg B$, $A'=A$. Но
$\Seq{A'_1,\dots,A'_n}{\neg B}{\cal L}$, $\Seq{}{\neg B\to (B\to
C)}{\cal L}$ по теореме \DMref{п.~4.3.7}{4.3.7.}~{\it в}, следовательно,
$\Seq{A'_1,\dots,A'_n}{B\to C=A'}{\cal L}$.

Пусть $I(B)=\text{T}$ и $I(C)=\text{T}$.
Тогда $I(A)=\text{T}$ и $C'=C$, $A'=A=B\to C$.
Имеем
$\Seq{A'_1,\dots,A'_n}{C}{\cal L}$,
$\Seq{}{C\to (B\to C)}{\cal L}$ (аксиома $A_1$ с подстановкой
$[C/A]$), следовательно, $\Seq{A'_1,\dots,A'_n}{B\to C=A'}{\cal L}$. %\{C/A\} -> [C/A] 

Пусть теперь $I(B)=\text{T}$ и $I(C)=\text{F}$.
Тогда $A'=\neg A=\neg(B\to C)$, $B'=B$ и~$C' =\neg C$.
Имеем $\Seq{A'_1,\dots,A'_n}{B}{\cal L}$,
$\Seq{A'_1,\dots,A'_n}{\neg C}{\cal L}$,
$\Seq{}{B\to(\neg C\to \neg(B\to C))}{\cal L}$
по теореме \DMref{п.~4.3.7}{4.3.7.}~{\it е}.
Следовательно,
$\Seq{A'_1,\dots,A'_n}{\neg(B\to C)=A'}{\cal L}$.
\end{proof}


\end{frame}

\begin{frame}
\frametitle{4.3.8. Множество теорем исчисления высказываний (3/3)}

\begin{theorem}
Теоремами теории $\cal L$ являются общезначимые формулы и только они:
$$\Seq{}{A}{\cal L}\Equiv A {\text { {\upshape--- тавтология}}}.$$
\end{theorem}

\begin{proof}
\proofsection{$\Impll$} Пусть $A$~--- тавтология. Тогда
$\Seq{A'_1,\dots,A'_n}{A}{\cal L}$ в любой интерпретации. Таким образом,
имеется $2^n$ различных выводимостей $\Seq{A'_1,\dots,A'_n}{A}{\cal L}$.
Среди них есть две, которые различаются в $A'_n$:
$\Seq{A'_1,\dots,A'_{n-1},a_n}{A}{\cal L}$ и
$\Seq{A'_1,\dots,A'_{n-1},\neg a_n}{A}{\cal L}$. По теореме дедукции
$\Seq{A'_1,\dots,A'_{n-1}}{a_n\to A}{\cal L}$ и
$\Seq{A'_1,\dots,A'_{n-1}}{\neg a_n\to A}{\cal L}$. 

Но $\Seq{}{(a_n\to
A)\to ((\neg a_n\to A)\to A)}{\cal L}$ по
теореме \DMref{п.~4.3.7}{4.3.7.},~{\it ж}, с подстановкой $[a_n/A, A/B]$, %\ref{abcdefg} -> 4.3.7 ; $a_n // A$, $A//B$ -> $[a_n/A, A/B]$
следовательно, $\Seq{A'_1,\dots,A'_{n-1}}{A}{\cal L}$. Повторив этот
процесс ещё $n-1$ раз, имеем $\Seq{}{A}{\cal L}$.

\proofsection{$\Impl$} Аксиомы $A_1$, $A_2$, $A_3$ суть
тавтологии. Правило Мodus ponens сохраняет тавтологичность, 
следовательно, теоремы теории $\cal L$ суть
тавтологии.
\end{proof}

\begin{corollary} Теория $\cal L$ формально непротиворечива.
\end{corollary}

%\addvspace{-8pt}
\begin{proof}
Все теоремы $\cal L$ суть тавтологии. Отрицание тавтологии не есть
тавтология. Следовательно, в $\cal L$ нет теоремы и её отрицания.
\end{proof}

\end{frame}


\begin{frame}
\label{4.3.9.}\subsubsection{4.3.9. Другие аксиоматизации исчисления высказываний }\frametitle{4.3.9. Другие аксиоматизации исчисления высказываний (1/3)} 

Теория $\cal L$ не является единственной возможной аксиоматизацией
исчисления высказываний. Её основное достоинство~--- лаконичность
при сохранении определённой наглядности. Действительно, в теории
$\cal L$ всего две связки, три схемы аксиом и два правила.
Известны и многие другие аксиоматизации исчисления высказываний,
предложенные различными авторами.
%
%\footnotetext{\;Давид Гильберт (1862--1943).}
%\addtocounter{footnote}{+1} \footnotetext{\;Вильгельм Аккерман
%(1896--1962).} \addtocounter{footnote}{+1} \footnotetext{\;Джон
%Беркли Россер (1907--1989).} \addtocounter{footnote}{+1}
%\footnotetext{\;Стефан Клини (1909--1994).}

\medskip
\begin{tabbing}


1. \= Гильберт и Аккерман,  \= 1938.\kill
1. \> Гильберт и Аккерман, \> 1938.\\
\> Связки: \> $\Or, \Not; \qquad (A\to B \Def \Not A \Or B)$. \\[-1pt]
\> Аксиомы: \> $(A \Or A) \to A$, \\           %добавил скобки
\> \> $A \to (A \Or B)$, \\[-1pt]
\> \> $(A \Or B) \to (B \Or A)$, \\
\> \> $(B\to C) \to ( (A\Or B) \to (A \Or C))$. \\[-1pt]
\> Правило: \> Modus ponens. \\
\end{tabbing}

\end{frame}


\begin{frame}
\frametitle{4.3.9. Другие аксиоматизации исчисления высказываний (2/3)} 

\begin{tabbing}

1. \= Гильберт и Аккерман,  \= 1938.\kill
2. \>  Россер,\> 1953. \\
\> Связки: \> $\And, \Not; \qquad (A\to B \Def \Not(A\And\Not B))$.\\
\> Аксиомы: \> $A  \to (A \And A)$, \\
\> \> $(A \And B) \to A$, \\
\> \> $(A\to B) \to (\Not( B\And C) \to \Not (C \And A))$. \\[-1pt]
\> Правило:\> Modus ponens. \\
\end{tabbing}

\begin{tabbing}
1. \= Гильберт и Аккерман,  \= 1938.\kill
 3.\> Никод, \> 1917. \\
\> Связка:\> $\mid; \qquad (A \mid B = \Not A \Or \Not B)$. \\
\> Аксиома: \> $(A\mid(B\mid C))\mid((D\mid(D\mid D))\mid$ \\
\> \>          $\mid ((E\mid B)\mid((A\mid E)\mid (A\mid E))))$. \\
\> Правило:\> $\displaystyle \Rule{A,A\mid(B\mid C)}{C}{}.$\\[-1pt]
\end{tabbing}

\end{frame}

\begin{frame}
\frametitle{ 4.3.9. Другие аксиоматизации исчисления высказываний (3/3)} %добавлено (2/2)

\begin{tabbing}
1. \= Гильберт и Аккерман,  \= 1938.\kill
4. \>  Клини, \> 1952. \\
\> Связки: \> $\Not, \And, \Or, \to$. \\
\> Аксиомы: \> $A \to(B \to A)$, \\
\> \> $(A \to (B \to C)) \to ((A\to B)\to(A\to C))$, \\[-1pt]
\> \> $(A \And B) \to A$, \\
\> \> $(A \And B) \to B$, \\
\> \> $A \to (B\to (A \And B))$, \\
\> \> $A \to (A\Or B)$, \\
\> \> $B \to (A\Or B)$, \\
\> \> $(A\to C)\to((B\to C)\to((A\Or B)\to C))$, \\
\> \> $(A\to B) \to ((A\to \Not B)\to \Not A)$, \\
\> \> $\Not \Not A \to A$. \\
\> Правило:\> Modus ponens.\\
\end{tabbing}

\end{frame}

%\addvspace{-3pt}

\subsection{4.4. Исчисление предикатов}
\label{4.4.}
\begin{frame}
\frametitle{4.4. Исчисление предикатов (1/2)}
{\large
Описание формальной теории исчисления предикатов
в этом разделе носит конспективный характер,
в частности, некоторые технически сложные
доказательства опущены. В то же время
уделено внимание прагматическим аспектам теории,
то есть ответу на вопрос, что \em{выразимо} и что \em{невыразимо}
в исчислении предикатов.
}

\medskip
\begin{tabular}{p{10mm}|p{38mm}|p{95mm}}
  № & Название параграфа
  & Ключевые термины и обозначения \\
  \hline
\DMref{4.4.1.}{4.4.1.}\RS & Определение исчисления предикатов   &
{ Кванторы всеобщности и существования, 
предметные константы и переменные, 
предметные предикаты и функциональные символы, вместимость, терм, атомарная формула, 
свободные и связанные вхождения переменной, связанная переменная, 
замкнутая формула, литеральная формула, область действия квантора, терм, свободный для переменной, чистое исчисление предикатов, прикладное исчисление предикатов, исчисление первого порядка, исчисления высших порядков }\\  
\end{tabular}  
\end{frame}


\begin{frame}
\frametitle{4.4. Исчисление предикатов (2/2)}
{\small
\begin{tabular}{p{10mm}|p{48mm}|p{85mm}}
  № & Название параграфа
  & Ключевые термины и обозначения \\
  \hline
\DMref{4.4.2.}{4.4.2.} & Интерпретация &
{ Область интерпретации, выполнимость формулы, истинная формула, ложная формула, модель множества формул, открытая формула, замыкание }\\\DMref{4.4.3.}{4.4.3.} & Общезначимость &
{ }\\  
\DMref{4.4.4.}{4.4.4.} & Непротиворечивость и полнота чистого исчисления предикатов &
{ Пропозициональная структура }\\  
\DMref{4.4.5.}{4.4.5.} & Логическое следование и логическая эквивалентность  &
{ Предварённая форма, бескванторная формула }\\  
\DMref{4.4.6.}{4.4.6.} & Теория равенства &
{ Неотличимость элементов }\\  
\DMref{4.4.7.}{4.4.7.} & Формальная арифметика  &
{  }\\  
\DMref{4.4.8.}{4.4.8.}\RS & Неаксиоматизируемые теории  &
{ Теория групп, 
исчисление предикатов с равенством, 
абелева группа, конечный порядок, 
делимая абелева группа, 
периодическая абелева группа,
неаксиоматизируемые и конечно неаксиоматизируемые теории }\\
\DMref{4.4.9.}{4.4.9.} & 
Теоремы Гёделя о неполноте &
{ }\\  
\end{tabular}  
}

\end{frame}


\begin{frame}
\label{4.4.1.}\subsubsection{4.4.1. Определение исчисления предикатов }\frametitle{4.4.1. Определение исчисления предикатов (1/7)}
\vspace{-2pt}
\odmkey{(Чистое) исчисление предикатов (первого порядка)}
{Исчисление (calculus)!предикатов (predicates calculus)}
\odmkey{}{Исчисление (calculus)!предикатов (predicates
calculus)!чистое (pure)} \odmkey{}{Исчисление
(calculus)!предикатов (predicates calculus)!первого порядка (first
order)}~--- это формальная теория $\cal K$, в которой определены
следующие компоненты: алфавит, формулы, аксиомы и~правила вывода.

\vspace{-2pt}
{\color{blue}\bf 1}. Алфавит:
\begin{tabbing}
служебные символы  \= функциональные символы  \= \kill
связки \> основные \> $\Not,\ \to $\\%[-1pt]
  \> дополнительные  \> $\And,\ \Or$ \\%[-1pt]
служебные символы  \> \>  $(\ ,\ )$ \\%[-1pt]
кванторы \> \odmkey{всеобщности}{Квантор
(quantifier)!всеобщности (universal)} \> $\forall$ \\%[-1pt]
  \> \odmkey{существования}{Квантор (quantifier)!существования
(existential)} \> $\exists$ \\%[-1pt]
предметные \> \odmkey{константы}{Константа (constant)!предметная (object)}
\> $a,b,\dots,a_1,b_1,\dots$ \\%[-1pt]
  \> \odmkey{переменные}{Переменная
 (variable)!предметная (object)} \> $x,y,\dots,x_1,y_1,\dots$ \\%[-1pt]
предметные \> \odmkey{предикаты}{Предикат (predicate)} \>  $P,Q,\dots $\\%[-1pt]
 \> \odmkey{функциональные символы}{Функциональный символ
(functional symbol)} \> $f,g,\dots $
\end{tabbing}



С каждым предикатом и функциональным символом
связано некоторое натуральное число $n$,
которое называется \odmkey{арностью}{Предикат (predicate)!n-местный (n-place)}
\odmkey{}{Функциональный символ (functional symbol)!$n$-местный ($n$-place)},
\odmkey{местностью}{Предикат (predicate)!$n$-местный ($n$-place)}
\odmkey{}{Функциональный символ (functional symbol)!$n$-местный ($n$-place)}
или \em{вместимостью}.

\end{frame}

\begin{frame}
\frametitle{4.4.1. Определение исчисления предикатов (2/7)}

{\color{blue}\bf 2}. Формулы имеют следующий синтаксис:
\begin{tabbing}
$\llrle{список термов}$ \= $\dDef \ $ \= до$\ $ \= \kill
$\llrle{формула}$ \> $\dDef$  \> $\llrle{атом}\mid $\\
 \>  \> $\neg \llrle{формула}\mid$ \\
 \>  \> $(\llrle{формула}\to \llrle{формула}) \mid$ \\
 \>  \> $\A{\llrle{переменная}}{\llrle{формула}}\mid$ \\
 \>  \> $\E{\llrle{переменная}}{\llrle{формула}}$ \\[3pt]
$\llrle{атом}$ \> $\dDef$ \> $\llrle{предикат}(\llrle{список термов})$ \\[3pt]
$\llrle{список термов}$ \> $\dDef$ \> $\llrle{терм}\mid \llrle{терм},\llrle{список термов}$ \\[3pt]
$\llrle{терм}$ \> $\dDef$ \> $\llrle{константа}\mid$ \\
 \>  \> $\llrle{переменная}\mid$ \\
 \>  \> $\llrle{функциональный символ}(\llrle{список термов})$ \\
\end{tabbing}

\vspace{-\baselineskip}
При этом должны быть выполнены следующие \em{контекстные условия\/}: в
\odmkey{терме}{Терм (term)} $f(\Tuple{t}{1}{n})$ функциональный
символ $f$ должен быть $n$-местным. В \odmkey{атоме}{Атом (atom)}
(или \odmkey{атомарной формуле}{Формула (formula)!атомарная (atomic)}) 
$P(\Tuple{t}{1}{n})$ предикат $P$ должен быть
$n$-местным.


\end{frame}

\begin{frame}
\frametitle{4.4.1. Определение исчисления предикатов (3/7)}
%\vspace{-2pt}
Любая переменная может входить в формулу несколько раз. 
При  этом  различные  вхождения  одной  и  той  же  переменной в 
формулу имеют разный  смысл,  в  зависимости  от  того, 
в  какой  синтаксической  позиции рассматривается вхождение.

\smallskip
Прежде всего, необходимо различать 
\odmkey{свободные}{Вхождение
 (occurrence)!переменной (variable)!свободное (free)}
и \odmkey{связанные}{Вхождение
(occurrence)!переменной
 (variable)!связанное (bound)} вхождения переменной.  
Вхождения переменных  в  атомарную  формулу  считаются  свободными.  
Свободные вхождения переменных в формулах $A$ и $B$ по определению
считаются свободными в формулах $\Not A$ и $A \to B$. В формулах
$\forall x\, (A)$ и $\exists x\, (A)$ формула $A$,  как  правило,  имеет  свободные 
вхождения  переменной  $x$.  Вхождения  переменной  $x$  в  формулы
$\forall x\, (A)$  и  $\exists x\, (A)$ считаются  связанными. 
Вхождения  других  переменных  (отличных  от  $x$), которые были
свободными в формуле $A$, по определению остаются свободными и в
формулах $\forall x\, (A)$ и $\exists x\, (A)$. Одна и та же переменная может
иметь в одной и той же формуле как свободные, так  и  связанные
вхождения.  

\smallskip
Если все вхождения переменной в формулу связаны, то
переменная называется \odmkey{связанной}{Переменная
 (variable)!связанная (bounded)}
(в данной формуле).

\end{frame}

\begin{frame}
\frametitle{4.4.1. Определение исчисления предикатов (4/7)}
%\vspace{-2pt}
Формула, не содержащая
свободных вхождений переменных, называется
\odmkey{замкнутой}{Формула (formula)!замкнутая (closed)}.
В замкнутой формуле все переменные связаны.

\smallskip
\begin{example}
Рассмотрим формулу $\A{x}{P(x)\to\E{y}{Q(x,y)}}$ и её % $\A{x}{(P(x)\to\E{y}{Q(x,y)})}$ -> $\A{x}{P(x)\to\E{y}{Q(x,y)}}$ (лишняя пара скобок)
подформулы. В подформулу $\E{y}{Q(x,y)}$ переменная $x$ входит
свободно, а все вхождения переменной $y$ связаны (квантором
существования). Таким образом, эта подформула не замкнута. С
другой стороны, то же самое вхождение переменной $x$ в подформулу
$Q(x,y)$ является связанным в формуле
$\A{x}{P(x)\to\E{y}{Q(x,y)}}$. В~этой формуле все вхождения всех
переменных связаны, а потому формула замкнута.
\end{example}

\smallskip
\begin{remark}
Язык теории $\cal L$ не содержит кванторов, поэтому понятия
свободного и связанного вхож\-дения переменных не применимы
непосредственно. Можно расширить исчисление высказываний, допуская
замкнутые формулы исчисления предикатов наравне с
пропозициональными переменными. В таком случае для удобства обычно
полагают, что все формулы теории $\cal L$ замкнуты.
\end{remark}

\end{frame}

\begin{frame}
\frametitle{4.4.1. Определение исчисления предикатов (5/7)}

%\vspace{-2pt}
Формулы вида $A$ и $\Not A$, где $A$~--атом, называются
\odmkey{литеральными}{Литерал (literal)} формулами (или
\em{литералами}).
В формулах $\A{x}{A}$ и $\E{x}{A}$ подформула $A$ называется
\odmkey{областью действия}{Область действия квантора
(scope of quantifier)} квантора по $x$.
Cвязки и кванторы упорядочивают по приоритету следующим образом:
$\Not, \forall, \exists, \And, \Or, \to$.
Лишние скобки при этом опускают.

\medskip
Теpм $t$ называется
\odmkey{свободным}{Терм (term)!свободный для переменной в формуле
 (free for variable in formula)}
для переменной $x$ в формуле $A$, если никакое свободное вхождение
переменной $x$ в формулу $A$ не лежит в области действия никакого
квантора по переменной $y$, входящей в терм $t$. В частности, терм
$t$ свободен для любой переменной в формуле $A$, если никакая
переменная терма не является связанной переменной формулы~$A$.


\begin{example}
 Терм $y$ свободен для переменной $x$ в формуле $P(x)$,
но тот же терм $y$ не свободен для переменной $x$ в формуле
$\A{y}{P(x)}$.
 Терм $f(x,z)$ свободен для переменной $x$ в формуле
$\A{y}{P(x,y)\to Q(x)}$,
но тот же терм $f(x,z)$ не свободен для переменной $x$ в формуле
$\E{z}{\A{y}{P(x,y)\to Q(x)}}$.
\end{example}

\end{frame}

\begin{frame}
\frametitle{4.4.1. Определение исчисления предикатов (6/7)}

{\color{blue}\bf 3}. Аксиомы (логические): любая система аксиом исчисления
высказываний плюс
\begin{eqnarray*}
P_1\colon\quad \A{x}{A(x)}\to A(t), \\
P_2\colon\quad A(t)\to \E{x}{A(x)},
\end{eqnarray*}
где терм $t$ свободен для переменной $x$ в формуле $A$.




{\color{blue}\bf 4}. Правила вывода:
$$
\Rule{A,\ A\to B}{B}{\text{(Modus ponens)}},
$$
$$  
\Rule{B\to A(x)}{B\to \A{x}{A(x)}}{(\forall^+)},$$ 
$$\Rule{A(x)\to B}{\E{x}{A(x)}\to B}{(\exists^+)},  
$$
где формула $A$ содержит свободные вхождения переменной $x$, а формула
$B$ их не содержит.

\end{frame}

\begin{frame}
\frametitle{4.4.1. Определение исчисления предикатов (7/7)}
Исчисление предикатов, которое не содержит предметных констант,
функциональных символов, предикатов и собственных аксиом,
называется \odmkey{чистым}{Исчисление (calculus)!предикатов
(predicates calculus)!чистое (pure)}. Исчисление предикатов,
которое содержит предметные константы и/или функцио\-нальные
символы и/или предикаты и связывающие их собственные аксиомы,
называется \odmkey{прикладным}{Исчисление (calculus)!предикатов
(predicates calculus)!прикладное (applied)}.

\smallskip
Исчисление предикатов, в котором кванторы могут связывать только
предметные переменные, но не могут связывать функциональные
символы или предикаты, называется исчислением \odmkey{первого
порядка}{Исчисление (calculus)!предикатов (predicates
calculus)!первого порядка (first order)}. Исчисления, в которых
кванторы могут связывать не только предметные переменные, но и
функциональные символы, предикаты или иные множества объектов,
называются исчислениями \odmkey{высших порядков}{Исчисление
(calculus)!предикатов
 (predicates calculus)!высшего порядка (higher order)}.

\smallskip
Практика показывает, что прикладного исчисления предикатов первого порядка
оказывается достаточно для
формализации содержательных теорий во всех разумных случаях.


\end{frame}

\begin{frame}
\label{4.4.2.}\subsubsection{4.4.2. Интерпретация }\frametitle{4.4.2. Интерпретация (1/3)}
%\subsection{Интерпретация}



\odmkey{Интерпретация}{Интерпретация (interpretation)!исчисления
предикатов (of predicate calculus)} $I$ (прикладного) исчисления
предикатов $\cal K$ с областью интерпретации (или
\odmkey{носителем}{Носитель (support)!интерпретации
 (of interpretation)}) $M$~--- это набор
функций, которые сопоставляют:

{ \settowidth{\leftmargini}{3)}
\addtolength{\leftmargini}{\labelsep}
\begin{enumerate}
\item[1)] каждой предметной константе $a$
элемент носителя $I(a)$,  $I(a)\in M$;
\item[2)] каждому $n$-местному функциональному символу
$f$  операцию $I(f)$ на носителе, \\
$\Map{I(f)}{M^n}{M}$;
\item[3)] каждому $n$-местному предикату $P$
отношение $I(P)$ на носителе,  $I(P)\subset M^n$.
\end{enumerate}
}


Пусть $x=(x_1,\dots)$~--- набор
(последовательность) переменных (входящих в формулу), а $s=(s_1,\dots)$
--- набор значений из $M$.
Тогда всякий терм $f(\Tuple{t}{1}{n})$ имеет значение на $s$ (из области
$M$), то~есть существует функция $\Map{s^*}{\{t\}}{M}$, определяемая
следующим образом:
$$
s^*(a)\Def I(a),\ 
s^*(x_i)\Def s_i, \ 
s^*(f(\Tuple{t}{1}{n}))\Def I(f)(s^*(t_1),\dots,s^*(t_n)).
$$


\end{frame}

\begin{frame}
\frametitle{4.4.2. Интерпретация (2/3)}
Всякий атом $P(\Tuple{t}{1}{n})$ имеет на $s$ истинностное значение $s^*(P)$,
определяемое следующим образом:
$$
s^*(P(\Tuple{t}{1}{n}))\Def (s^*(t_1),\dots,s^*(t_n))\in I(P).
$$


Если $s^*(P)=\text{T}$, то говорят, что формула $P$
\odmkey{выполнена}{Выполнимость (satisfiability)!формулы (of
formula)} на $s$.


\smallskip
Формула $\Not A$ выполнена на $s$ тогда и только тогда, когда  формула $A$ не
выполнена на $s$.


\smallskip
Формула $A \to B$ выполнена на $s$ тогда и только тогда, когда формула $A$ не
выполнена на $s$ или формула $B$ выполнена на $s$.


\smallskip
Формула $\A{x_i}{A}$ выполнена на $s$ тогда и только тогда, когда формула $A$
выполнена на любом наборе $s'$, отличающемся от $s$ не более чем своим
$i$-м компонентом. 

\smallskip
Формула $\E{x_i}{A}$ выполнена на $s$ тогда и только тогда, когда формула $A$
выполнена на каком-либо наборе $s'$, отличающемся от $s$
не более чем своим $i$-м компонентом.

\end{frame}

\begin{frame}
\frametitle{4.4.2. Интерпретация (3/3)}
Формула называется \odmkey{истинной}{Формула (formula)!истинная (true)}
в данной
интерпретации $I$, если она выполнена на любом наборе $s$ элементов $M$.

\smallskip
Формула называется \odmkey{ложной}{Формула (formula)!ложная (false)} в данной
интерпретации $I$, если она не выполнена ни на одном наборе $s$ элементов $M$.


\smallskip
Интерпретация называется \odmkey{моделью}{Модель (model)!множества
формул (of set of formulas)} множества формул $\Gamma$, если все
формулы из $\Gamma$ истинны в данной интерпретации.


\smallskip
Всякая замкнутая формула истинна или ложна в
данной интерпретации.

\smallskip
\odmkey{Открытая}{Формула (formula)!открытая (open)}
 (то~есть \em{не} замкнутая)
формула $A(x,y,z,\dots)$ истинна в данной интерпретации тогда  и
только тогда, когда её \odmkey{замыкание}{Замыкание
(closure)!формулы (of formula)} $\A{ x}{\A{ y}{\A{ z\dots}{
A(x,y,z,\dots)}}}$ истинно в данной интерпретации.

\smallskip
\begin{remark}
Это
обстоятельство объясняет, почему собственные аксиомы прикладных теорий
обычно пишутся в открытой форме.
\end{remark}

\end{frame}

\label{4.4.3.}\subsubsection{4.4.3. Общезначимость}
\begin{frame}
\frametitle{4.4.3. Общезначимость}

%\vspace{-2pt}
Формула (исчисления предикатов)
\odmkey{общезначима}{Формула (formula)!общезначимая (valid)},
если она истинна в любой интерпретации.

%\vspace{-6pt}
\begin{theorem}
Формула $\A{x}{A(x)}\to A(t)$, где терм $t$ свободен для переменной
$x$ в~формуле $A$, общезначима.
\end{theorem}
%\pagebreak[3]
%\vspace{-12pt}
\begin{proof} Рассмотрим произвольную интерпретацию $I$,
произвольную последовательность значений из области интерпретации
$s$ и соответствующую функцию $s^*$. 

Заметим следующее. Пусть
$t_1$~--- терм и $s^*(t_1)=a_1$. Пусть $t(x)$~--- некоторый другой
терм, а $t'\Assign t(\dots x\dots)[t_1/x]$. Тогда
$s^*(t')=s_1^*(t)$, где $s_1$ имеет значение $a_1$ на месте $x$.

Пусть теперь $A(x)$~--- формула, а терм $t$ свободен для $x$ в~$A$. 
Положим $A(t) {\Assign} A(x)[t/x]$. Имеем $
s^*(A(t))=\text{T}\Equiv s_1^*(A(t))=\text{T}$, где $s_1$ имеет
значение $s^*(t)$ на месте $x$. Если $s^*(\A{x}{A(x)})=\text{T}$ и
терм $t$ свободен для $x$ в $A$, то $s^*(A(t))=\text{T}$.
Следовательно, формула $\A{x}{A(x)}\to A(t)$ 
выполнена на всех
последовательностях в произвольной интерпретации.
\end{proof}


%\vspace{-12pt}
\begin{remark}
Можно показать, что формула $A(t) \to \E{x}{A(x)}$,
где терм $t$ свободен для переменной
$x$ в~формуле $A$, общезначима.
\end{remark}
\end{frame}



%\vspace{-6pt}
%\subsection{Непротиворечивость и полнота чистого~исчисления предикатов}
\begin{frame}
\label{4.4.4.}\subsubsection{4.4.4. Непротиворечивость и полнота чистого исчисления предикатов }\frametitle{4.4.4. Непротиворечивость и полнота чистого исчисления предикатов (1/2)}
%Нам потребуется следующее преобразование $h$ формулы исчисления
%предикатов в~формулу исчисления высказываний:
%в формуле опускаются все кванторы и термы вместе с соответствующими
%скобками и запятыми, остаются предикатные символы (которые рассматриваются
%как пропозициональные переменные) и~связки.
%
%\begin{example}
%%\centering
%$h\left(\A{x}{P(x)\to \E{y}{Q(x,y)}}\right)=  P\to Q.
%$
%\end{example}

Рассмотрим преобразование $h$, которое <<забывает>>
всё, что связано с исчислением предикатов,
оставляя только
\em{пропозициональную структуру},
которая является формулой исчисления
высказываний.
При преобразовании $h$ 
в формуле исчисления предикатов опускаются все кванторы и термы вместе с соответствующими
скобками и запятыми, остаются только предикатные символы (которые рассматриваются
как пропозициональные переменные) и~связки.

\begin{example}
%\centering
$h\left(\A{x}{P(x)\to \E{y}{Q(x,y)}}\right)=  P\to Q.
$
\end{example}

Нетрудно видеть, что $h(\Not A) = \Not h(A)$
и $h(A\to B) =  h(A) \to h(B)$.

%\pagebreak[4]
%\setcounter{theorem}{0}
\begin{theorem}
Формальная теория $\mathcal{K}$ непротиворечива.
\end{theorem}

\begin{proof}
Рассмотрим пропозициональную структуру аксиом исчисления
предикатов $\mathcal{K}$, то есть применим к аксиомам
преобразование $h$. Ясно, что все полученные формулы суть
тавтологии, поскольку аксиомы $A_1,A_2,A_3$ сами по
себе совпадают со своими пропозициональными
структурами, которые являются тавтологиями. 
\end{proof}
\end{frame}

\begin{frame}
\frametitle{4.4.4. Непротиворечивость и полнота чистого исчисления предикатов (2/2)}
\begin{proof} (Продолжение)
Для
аксиом исчисления предикатов имеем \\
$h(P_1)=h(\A{x}{A(x)}\to A(t)) = A\to A$ и
$h(P_2)=h(A(t)\to\E{x}{A(x)} ) = A\to A$. 
Далее, правило Modus
ponens сохраняет тавтологичность пропозициональной структуры, то
есть если $h(A)$ и $h(A\to B)$~--- тавтологии, то $h(B)$~---
тавтология. Правила $\exists^+$ и $\forall^+$ также сохраняют
тавтологичность пропозициональной структуры,
поскольку в этих правилах пропозицио\-нальная структура заключения
совпадает с пропозициональной структурой посылки.
 Таким образом,
если формула $A$ является теоремой теории $\mathcal{K}$, то
формула $h(A)$ является тавтологией. Далее от противного. Если
$\Seq{}{A}{\mathcal{K}}$ и $\Seq{}{\Not A}{\mathcal{K}}$, то
$h(A)$ и  $\Not h(A)$~--- тавтологии, что невозможно.
\end{proof}

Следующие две метатеоремы устанавливают свойства полноты исчисления
предикатов. Теоремы приводятся без доказательств.

\begin{theorem}
Всякая теорема чистого исчисления предикатов первого порядка
общезначима.
\end{theorem}


\begin{theorem}
Всякая общезначимая формула является теоремой чистого исчисления предикатов
первого порядка.
\end{theorem}

%\begin{proof}
%Без доказательства.
%\end{proof}
\end{frame}



%\subsection{Логическое следование и логическая эквивалентность}
\begin{frame}
\label{4.4.5.}\subsubsection{4.4.5. Логическое следование и логическая эквивалентность }\frametitle{4.4.5. Логическое следование и логическая эквивалентность (1/2)}

Формула $B$ является
\odmkey{логическим следствием}{Следствие логическое (logical consequence)} 
формулы $A$
(обозначение ${A\Impl B}$), если формула $B$ выполнена на любом
наборе в любой интерпретации, на котором выполнена формула $A$.
Формулы $A$ и $B$ \odmkey{логически эквивалентны}{Эквивалентность логическая (logical equivalence)} (обозначение $A=B$ или $A\Leftrightarrow B$), если они
являются логическим следствием друг друга. 
Можно показать, что
имеют место следующие логические следования и~эквивалентности:


~1. $\Not \A{x}{A(x)} = \E{x}{\Not A(x)}$,

~2. $\Not \E{x}{A(x)} = \A{x}{\Not A(x)}$, 

~3. $\A{x}{(A(x)\And B(x))} =\A{x}{A(x)}\And\A{x}{B(x)}$,

~4. $\E{x}{A(x)\Or B(x)} = \E{x}{A(x)}\Or\E{x}{B(x)}$, %\Impl -> = (заменил более жёстким утверждением)

~5. $\E{x}{A(x)\And B(x)} \Impl \E{x}{A(x)}\And\E{x}{B(x)}$,

~6. $\A{x}{A(x)}\Or \A{x}{B(x)} \Impl \A{x}{A(x) \Or B(x)}$, 

~7. $\A{x}{\A{y}{A(x,y)}} = \A{y}{\A{x}{A(x,y)}}$,

~8. $\E{x}{\E{y}{A(x,y)}} = \E{y}{\E{x}{A(x,y)}}$,  



\end{frame}

\begin{frame}
\frametitle{4.4.5. Логическое следование и логическая эквивалентность (2/2)}

~9. $\A{x}{(A(x)\And C)} = \A{x}{A(x)}\And C$,

10. $\A{x}{A(x)\Or C} = \A{x}{A(x)}\Or C$, 

11. $\E{x}{A(x)\And C} = \E{x}{A(x)}\And C$,% \E{x}{(A(x)\And C)} -> \E{x}{A(x)\And C} (убрал лишние скобки)

12. $\E{x}{A(x)\Or C} = \E{x}{A(x)}\Or C$, 

13. $C\to\A{x}{A(x)} = \A{x}{C\to A(x)}$,

14. $C\to\E{x}{A(x)} = \E{x}{C\to A(x)}$, 

15. $\A{x}{A(x)}\to C= \A{x}{A(x)\to C}$,

16. $\E{x}{A(x)}\to C = \E{x}{A(x)\to C}$, 

17. $\E{x}{\A{y}{A(x,y)}} \Impl \A{y}{\E{x}{A(x,y)}}$,

18. $\A{x}{A(x)\to B(x)} \Impl \A{x}{A(x)} \to \A{x}{B(x)}$,

где формула $C$ не содержит никаких вхождений
переменной $x$.

\smallskip
Для всякой замкнутой формулы $A$ существует логически
эквивалентная ей формула $A'$, имеющая  \odmkey{предварённую
форму}{Формула (formula)!в предварённой форме (prenex form)}:
$
A' = Q_1\,x_1\dots Q_n\,x_n\ \bar A(\Tuple{x}{1}{n}),
$
где $\Tuple{Q}{1}{n}$~--- некоторые кванторы, а $\bar A$~---
\odmkey{бескванторная}{Формула (formula)!бескванторная
(quantifier-free)} формула.

\end{frame}

%\subsection{Теория равенства}
\begin{frame}
\label{4.4.6.}\subsubsection{4.4.6. Теория равенства }\frametitle{4.4.6. Теория равенства (1/2)}
Этот параграф начинает серию примеров
прикладных исчислений предикатов.

\begin{remark}
Всякое прикладное исчисление предикатов содержит
в себе чистое исчисление предикатов, поэтому при описании
прикладного исчисления аксиомы и все теоремы чистого исчисления подразумеваются,
указываются только собственные функциональные символы, предикаты и аксиомы.
\end{remark}

\odmkey{Теория равенства}{Теория (theory)!равенства (of equality)}
$\cal E\!$~--- это прикладное исчисление предикатов, в котором
имеются:

\vspace{-3pt}
\begin{enumerate}
\item Собственный двухместный предикат $=$ (инфиксная запись $x=y$).
\item Собственные аксиомы (схемы аксиом):

%\vspace{-4pt}
{
\settowidth{\leftmarginii}{$E_1\colon$}
\addtolength{\leftmarginii}{\labelsep}
\begin{enumerate}
\itemsep 0pt
\parsep 0pt
\parskip 1pt plus 1pt
\item[$E_1\colon$] $\A{x}{x=x}$,
\item[$E_2\colon$] $(x=y)\to(A(x)\to A(x)[y/1/x])$.
\end{enumerate}
}
\end{enumerate}


\end{frame}

%\subsection{Теория равенства}
\begin{frame}
\frametitle{4.4.6. Теория равенства (2/2)}
%\vspace{-4pt}
\begin{theorem} В теории ${\cal E}$ выводимы следующие теоремы:
\begin{enumerate}
\item $\Seq{}{t=t}{\cal E}$ для любого терма $t$.%
\item $\Seq{}{x=y\to y=x}{\cal E}$.%
\item $\Seq{}{x=y\to(y=z\to x=z)}{\cal E}$.%
\end{enumerate}
\end{theorem}


%\pagebreak[3]
%\vspace{-3pt}
\begin{proof}
\proofsection{1} Из $E_1$ и $P_1$ по Мodus ponens.
\proofsection{2} Имеем: $(x=y)\to (x=x\to y=x)$ из $E_2$ при
подстановке $[x=x/A(x)]$. Значит, $\Seq{x=y, x=x}{y=x}{\cal E}$.
Но $\Seq{}{x=x}{\cal E}$, следовательно, $\Seq{x=y}{y=x}{\cal E}$.
По~теореме дедукции $\Seq{}{x=y\to y=x}{\cal E}$.
\proofsection{3}
Имеем: $y=x\to (y=z\to x=z)$ из $E_2$ при подстановке
$[x=z/A(x),y/x, x/y]$, значит,
$\Seq{y=x}{(y=z\to x=z)}{\cal E}$. Но $\Seq{x=y}{y=x}{\cal E}$,
по транзитивности \\$x=y \vdash_{\cal E}$ $(y= z\to x=z)$,
по теореме дедукции
$\Seq{}{(x=\nobreak y) \to (y=\nobreak z \to x=z)}{\cal E}$.
\end{proof}



Таким образом, равенство является эквивалентностью. Но равенство~--- это
более сильное свойство, чем эквивалентность. Аксиома $E_2$ выражает
\em{неотличимость} равных элементов.

\end{frame}

\begin{frame}
\frametitle{4.4.7. Формальная арифметика}
\subsubsection{4.4.7. Формальная арифметика}
\label{4.4.7.}
%\vspace{-4pt}
\odmkey{Формальная арифметика}{Теория (theory)!формальной
арифметики (formal arithmetic)} $\cal A$~--- это прикладное
исчисление предикатов, в котором есть:

\begin{enumerate}

%\vspace{-2pt}
\item Предметная константа 0.



%\vspace{-2pt}
\item Двухместные функциональные символы $+$ и $\cdot$,\\
одноместный функциональный символ $'$.



%\vspace{-2pt}
\item Двухместный предикат $=$.



%\vspace{-2pt}
\item Собственные аксиомы (схемы аксиом):



\begin{enumerate}
\item[$A_1\colon$]  $(P(0) \And \A{x}{P(x)\to P(x')})\to\A{x}{P(x)}$,%
\item[$A_2\colon$]  ${t_1}' = {t_2}' \to t_1 = t_2$,
\item[$A_3\colon$]  $\Not (t'=0)$,
\item[$A_4\colon$]  $t_1=t_2\to(t_1=t_3\to t_2=t_3)$,
\item[$A_5\colon$]  $t_1=t_2 \to {t_1}'={t_2}'$,
\item[$A_6\colon$]  $t+0=t$,
\item[$A_7\colon$]  $t_1+{t_2}'=(t_1+t_2)'$,
\item[$A_8\colon$]  $t\cdot 0=0$,
\item[$A_9\colon$]  $t_1\cdot{t_2}'=t_1\cdot t_2 + t_1$,
\end{enumerate}


где $P$~--- любая формула, а $t,t_1,t_2$~--- любые термы теории
$\cal A$.
\end{enumerate}

%\begin{remark}
%Впервые данная система аксиом для арифметики была предложена
%Пеано\footnote{\;Джузеппе Пеано (1858--1932).}, а $A_1$~--- это
%схема аксиомы математической индукции.
%\end{remark}

\end{frame}


\begin{frame}
\label{4.4.8.}\subsubsection{4.4.8. Неаксиоматизируемые теории }\frametitle{4.4.8. Неаксиоматизируемые теории (1/5)}
В этом параграфе на примере некоторых понятий
теории (абелевых) групп неформально
устанавливаются границы применимости
исчисления предикатов первого порядка.

\smallskip
\odmkey{Теория групп}{Теория (theory)!групп (group)} $\cal G$~---
это прикладное исчисление предикатов \odmkey{с
равенством}{Исчисление (calculus)!предикатов (predicates
calculus)!с равенством (with equality)}, в котором имеются:
\begin{enumerate}
\item Предметная константа 0.
\item Двухместный функциональный символ $+$.
\item Собственные аксиомы (схемы аксиом):
%\end{enumerate}



%\vspace{-6pt}
{ \settowidth{\leftmargini}{$G_1\colon$}
\addtolength{\leftmargini}{\labelsep}
\begin{enumerate}
\itemsep 0pt
\parsep 0pt
\parskip 1pt plus 1pt
\item[$G_1\colon$] $\A{x,y,z}{x + (y + z) = (x + y) + z }$,\\[-2pt]%
\item[$G_2\colon$] $\A{x}{0+x=x}$,\\[-2pt]%
\item[$G_3\colon$] $\A{x}{\E{y}{x+y=0}}$.\\[-2pt]%
\end{enumerate}
}
\end{enumerate}
%\begin{eqnarray*}
%G_1\colon\quad \A{x,y,z}{x + (y + z) = (x + y) + z },\\
%G_2\colon\quad \A{x}{0+x=x},\\
%G_3\colon\quad \A{x}{\E{y}{x+y=0}}.
%\end{eqnarray*}
%}


\smallskip
\begin{remark}
Выражение <<теория первого порядка с равенством>> означает, что
подразумевается наличие предиката $=$,
 аксиом $E_1$ и $E_2$ и всех их следствий.
\end{remark}


\smallskip
Группа называется \odmkey{абелевой}{Группа (group)!абелева
 (abelian)}, если имеет место
аксиома

{
\settowidth{\leftmargini}{$G_1\colon$}
\addtolength{\leftmargini}{\labelsep}
\begin{itemize}
\itemsep 0pt
\parsep 0pt
\parskip 1pt plus 1pt
\item[$G_4\colon$] $\A{x,y}{x+y=y+x}$.
\end{itemize}
}


%\vspace{-\baselineskip}
%$$G_4\colon\quad \A{x,y}{x+y=y+x}.$$

\end{frame}

\begin{frame}
\frametitle{4.4.8. Неаксиоматизируемые теории (2/5)}
Говорят, что в абелевой группе все элементы имеют
\odmkey{конечный порядок $n$}{Группа (group)!порядка $n$ ($n$-order)},
если выполнена
собственная аксиома
{
\settowidth{\leftmargini}{$G_1\colon$}
\addtolength{\leftmargini}{\labelsep}
\begin{itemize}
\itemsep 0pt
\parsep 0pt
\parskip 1pt plus 1pt
\item[$G_5\colon$] $\A{x}{\E{k\le n}{kx=0}}$,
где $kx$~--- это сокращение для $\underbrace{x+x+\dots+x}_{k \text{раз}}$.
\end{itemize}
}

Эта формула не является формулой теории $\cal G$, поскольку
содержит <<посторонние>> предметные предикаты и переменные. Однако
для любого конкретного конечного $n$ собственная аксиома может
быть записана в виде допустимой формулы теории $\cal G$:
{
\settowidth{\leftmargini}{$G_1\colon$}
\addtolength{\leftmargini}{\labelsep}
\begin{itemize}
\itemsep 0pt
\parsep 0pt
\parskip 1pt plus 1pt
\item[$G_{5}'\colon$] $\A{x}{(x=0\Or 2x=0 \Or\dots\Or nx=0)}$.
\end{itemize}
}
%$$G_{5}'\colon\quad \A{x}{(x=0\Or 2x=0 \Or\dots\Or nx=0)}.$$


Абелева группа называется \odmkey{делимой}{Группа (group)!делимая
(divisible)}, если выполнена собственная аксиома
{
\settowidth{\leftmargini}{$G_1\colon$}
\addtolength{\leftmargini}{\labelsep}
\begin{itemize}
\itemsep 0pt
\parsep 0pt
\parskip 1pt plus 1pt
\item[$G_6\colon$] $\A{n\ge 1}{\A{x}{\E{y}{ny=x}}}$.
\end{itemize}
}
%$$G_6\colon\quad \A{n\ge 1}{\A{x}{\E{y}{ny=x}}}.$$


Эта формула не является
формулой теории $\cal G$, поскольку содержит <<посторонние>>
предметные
константы и переменные. 
\end{frame}

\begin{frame}
\frametitle{4.4.8. Неаксиоматизируемые теории (3/5)}


Однако
собственная аксиома 
может быть записана в виде
бесконечного множества допустимых формул тео\-рии~$\cal G$:
{
\settowidth{\leftmargini}{$G_1\colon$}
\addtolength{\leftmargini}{\labelsep}
\begin{itemize}
\itemsep 0pt
\parsep 0pt
\parskip 1pt plus 1pt
\item[$G'_6\colon$] $\A{x}{\E{y}{2y=x}}$,\\
$\A{x}{\E{y}{3y=x}}$,\\
$\dots$
\end{itemize}
}
%\begin{eqnarray*}
%&G'_6\colon\quad &\A{x}{\E{y}{2y=x}},\\
%& &\A{x}{\E{y}{3y=x}},\\
%& &\dots
%\end{eqnarray*}

Но можно показать, что любое \em{конечное} множество формул,
истинное во всех делимых абелевых группах, истинно и в некоторой
неделимой абелевой группе, то~есть теория делимых абелевых групп не
является \odmkey{конечно}{Аксиоматизируемость
(axiomatizability)!конечная
 (finite)} аксиоматизи\-руемой.


%\pagebreak
Абелева группа называется \odmkey{периодической}{Группа
(group)!периодическая (periodic)},
если выполнена собственная %\linebreak
аксиома
{
\settowidth{\leftmargini}{$G_1\colon$}
\addtolength{\leftmargini}{\labelsep}
\begin{itemize}
\itemsep 0pt
\parsep 2pt
\parskip 4pt plus 1pt
\item[$G_7\colon$] $\A{x}{\E{n\ge 1}{nx=0}}$.
\end{itemize}
}
%$$G_7\colon\quad \A{x}{\E{n\ge 1}{nx=0}}.$$

\vspace{-2pt}
Эта формула также не является
формулой теории $\cal G$, поскольку содержит <<посторонние>> предметные
константы и переменные. 
%\vspace*{-.5\baselineskip}
\end{frame}

\begin{frame}
\frametitle{4.4.8. Неаксиоматизируемые теории (4/5)}
Если попытаться преобразовать формулу $G_7$ по
образцу формулы $G_5$, то получится бесконечная <<формула>>
{
\settowidth{\leftmargini}{$G_1\colon$}
\addtolength{\leftmargini}{\labelsep}
\begin{itemize}
\itemsep 0pt
\parsep 2pt
\parskip 4pt plus 1pt
\item[$G'_7\colon$] $\A{x}{(x=0 \Or 2x=0\Or\dots\Or nx=0 \Or\dots)},$
\end{itemize}
}
%$$G'_7\colon\quad \A{x}{(x=0 \Or 2x=0\Or\dots\Or nx=0 \Or\dots)},$$

%\vspace{-2pt}
которая не является допустимой формулой исчисления предикатов и
тем более теории $\cal G$. Таким образом, теория периодических
абелевых групп не является
\odmkey{аксиоматизируемой}{Аксиоматизируемость (axiomatizability)}
(если не включать в теорию групп и всю формальную арифметику).
\end{frame}

\begin{frame}
\frametitle{4.4.8. Неаксиоматизируемые теории (5/5)}

\begin{digress}
Наличие неаксиоматизируемых и конечно неаксиоматизируемых формальных
теорий не означает практической неприменимости аксиоматического метода
на основе использования языка исчисления
предикатов первого порядка.
Это означает, что аксиоматический метод не применим <<в чистом виде>>. На
практике формальные теории, описывающие 
содержательные объекты, задаются
с помощью собственных аксиом, которые наряду с~собственными предикатами и
функциональными символами содержат <<внелогические>> предикаты и
функциональные символы, свойства
которых аксиомами не описываются, а~считаются известными (в данной
теории). В рассмотренных примерах 
внелогическими являются
натуральные числа и операции над ними. Аналогичное обстоятельство имеет
место и в системах логического программирования типа Пролог. Реализация
такой системы всегда снабжается библиотекой внелогических (или
\em{встроенных}) предикатов и функций, которые и обеспечивают
практическую применимость системы логического программирования.
\end{digress}

%\vspace*{-.5\baselineskip}

\end{frame}


%\addvspace{6pt}
%\subsection{Теоремы Гёделя о неполноте}\label{Gedel}

\begin{frame}
\label{4.4.9.}\subsubsection{4.4.9. Теоремы Гёделя о неполноте }\frametitle{4.4.9. Теоремы Гёделя о неполноте (1/2)}
\vspace{-2pt}
Давно известны некоторые важные свойства формальных теорий, в
частности прикладных исчислений предикатов первого порядка,
которые существенным образом влияют на практическую применимость
формальных  теорий (аксиоматического метода).
%Полное
%доказательство этих фактов выходит далеко за рамки данного
%учебника. 
Понимание \em{границ применимости} формальных методов
является, по нашему мнению, необходимым для использования таких
методов. Поэтому ниже приводятся 
 два важнейших факта, известные как
теоремы Гёделя о~неполноте.
Аксиоматический метод обладает множеством достоинств
и с успехом применяется на практике для формализации самых
разнообразных предметных областей в математике, физике и других
науках. Однако этому методу присуще следующее принципиальное
ограничение.

\begin{theorem} \NB{[Первая теорема Гёделя о неполноте]}
Во всякой достаточно богатой непротиворечивой теории первого порядка (в частности, %Во всякой достаточно богатой теории -> Во всякой достаточно богатой непротиворечивой теории
во всякой теории, включающей формальную арифметику) существует
такая истинная формула $F$, что ни $F$, ни $\Not F$ не являются
выводимыми в этой теории.\odmkey{ }{Теорема (theorem)!Гёделя (Kurt
Friedrich G$\ddot{\text{o}}$del)}
\end{theorem}
\end{frame}


%\addvspace{6pt}
%\subsection{Теоремы Гёделя о неполноте}\label{Gedel}

\begin{frame}
\vspace{-2pt}
\frametitle{\large 4.4.9. Теоремы Гёделя о неполноте (2/2)}

%\vspace{-6pt}
Утверждения о теории первого порядка также могут быть
сформулированы в виде формул теории первого порядка. В частности,
утверждения о свойствах формальной арифметики (например,
утверждение о непротиворечивости формальной арифметики) могут
быть сформулированы как арифметические утверж\-дения.


\begin{theorem} \NB{[Вторая теорема Гёделя о неполноте]}
Во всякой достаточно богатой теории первого порядка (в частности, 
во всякой теории, включающей формальную арифметику) формула $F$, 
утверждающая непротиворечивость этой теории, не является выводимой
в ней. \odmkey{ }{Теорема (theorem)!Гёделя (Kurt Friedrich
G$\ddot{\text{o}}$del)}
\end{theorem}


%\pagebreak

\begin{digress}
Вторая теорема Гёделя не утверждает, что арифметика противоречива.
Фор\-мальная арифметика \DMref{п.~4.4.7}{4.4.7.} как раз
непротиворечива. Эта теорема утверждает, что непротиворечивость
достаточно богатой формальной теории не может быть установлена
средствами самой этой теории. При этом вполне может статься, что
непротиворечивость одной конкретной теории  может быть установлена
средствами другой, более мощной формальной теории. Но тогда
возникает вопрос о~непротиворечивости этой второй теории и так далее.
\end{digress}

\end{frame}


\begin{frame}
\frametitle{Онтология математической логики}
\begin{center}
\includegraphics[width=12cm]{ODM4_34.png}
\end{center}

\end{frame}


\subsection{4.5. Наивная теория алгоритмов}
\label{4.5.}
\begin{frame}
\frametitle{ 4.5. Наивная теория алгоритмов (1/2)}

%Дополнение
%К курсу «Дискретная математика для программистов»
%Наивная теория алгоритмов
%Изложение построено в форме дополнения 
%к учебнику «Дискретная математика» издание 2013 года. 
%
%В этом Дополнении изложено доказательство первой теоремы Гёделя о неполноте,
%основанное на понятии алгоритма. 
%Приводимое доказательство в своих идеях следует брошюре Успенского  [32], 
%и является достаточно далёким от оригинального доказательства Гёделя. 
%Приведенное изложение далеко и от Успенского, поскольку сокращено, упрощено,
%«запрограммировано» и выдержано в стиле учебника.
В этом разделе изложено доказательство первой теоремы Гёделя о неполноте,
основанное на интуитивном понятии алгоритма.
Приводимое доказательство в своих идеях следует брошюре В.А.~Успенского 
и является достаточно далёким от оригинального доказательства Гёделя. 

\medskip
Основная цель раздела состоит в демонстрации
применимости и действенности \em{программистских} подходов в 
\em{математическом} контексте.

%
%…Ты когда-нибудь видела, как рисуют множество?
%– Множество чего? – спросила Алиса.
%– Ничего, – отвечала Соня. – Просто множество!
%– Не знаю, – начала Алиса, – может...
%– А не знаешь — молчи, – оборвал ее Болванщик.
%Льюис Кэрролл, Алиса в стране чудес

\medskip
\begin{tabular}{p{10mm}|p{46mm}|p{87mm}}
  № & Название параграфа
  & Ключевые термины и обозначения \\
  \hline
\DMref{4.5.1.}{4.5.1.} & Алгоритмы &
{ Машина Тьюринга, нормальный алгорифм Маркова, тезис Тьюринга--Чёрча, выполнение алгоритмов,  заголовок и тело алгоритма }\\
\DMref{4.5.2.}{4.5.2.} & Вычислимые функции &
{ Аргумент, результат, реализация вычислимой функции, функционально эквивалентные алгоритмы, представление вычислимой функции, тип }\\\DMref{4.5.3.}{4.5.3.} & Перечислимые множества &
{ Перечисляющий алгоритм }\\
\DMref{4.5.4.}{4.5.4.} & Разрешимые множества &
{ }\\    
\end{tabular}  

\end{frame}

\begin{frame}
\frametitle{4.5. Наивная теория алгоритмов (2/2)}
\begin{tabular}{p{10mm}|p{48mm}|p{85mm}}
  № & Название параграфа
  & Ключевые термины и обозначения \\
  \hline
\DMref{4.5.5.}{4.5.5.} & Протокол выполнения алгоритма   &
{ Детерминированный и недетерминированный алгоритм }\\  
\DMref{4.5.6.}{4.5.6.} & Программы    &
{ Применимый алгоритм, условно равные алгоритмы, равносильные алгоритмы, всюду отличающиеся алгоритмы, проекция алгоритма, остаточная программа, представительный класс, алгоритм интерпретации, 
синтаксис, семантика}\\  
\DMref{4.5.7.}{4.5.7.} & Невычислимые функции и неразрешимые множества  &
{ Универсальный алгоритм, всюду определенное продолжение }\\  
\DMref{4.5.8.}{4.5.8.}\RS & Истинность и доказуемость  &
{ Доказательство, истина, функция извлечения доказанного, теорема, полная / непротиворечивая / адекватная дедуктика}\\  
\DMref{4.5.9.}{4.5.9.} & Множество истин арифметики     &
{ Арифметика, терм, параметр, константа, атом, элементарная формула, множество истин арифметики, арифметическое множество, сопряжённое множество }\\  
\end{tabular}  


\end{frame}

\begin{frame}
\label{4.5.1.}\subsubsection{4.5.1. Алгоритмы }\frametitle{4.5.1. Алгоритмы (1/5)}
%\vspace{-2pt}
Интуитивное понятие алгоритма как 
<<набора инструкций, описывающих порядок действий исполнителя 
для достижения результата решения задачи за конечное число действий>> 
или как <<точно определённого предписания, 
позволяющего во многих случаях из исходных данных получить результат>>, 
обладает в точности теми же достоинствами и недостатками, 
что и определение понятия множества в \DMref{п.~1.1.1}{1.1.1.}. 
Приведённые определения понятия <<алгоритм>> понятны всем программистам, 
они вызывают верные практические ассоциации с программами, 
и они совершенно неполны и не обладают математической строгостью.
Тем не менее, даже этого нестрогого понятия оказывается достаточно во многих важных
случаях, так же как и нестрогого понятия <<множество>> 
в наивной теории множеств оказывается достаточно для многих полезных построений, 
в частности тех, которые приведены в этом курсе. 

\medskip
Будем считать \odmkey{алгоритм}{Алгоритм (algorithm)} неопределяемым понятием, 
постулируя те или иные его свойства по мере необходимости.
\end{frame}

\begin{frame}
\frametitle{4.5.1. Алгоритмы (2/5)}

\begin{digress}
Обычно в \em{строгой} теории алгоритмов вводится одно или несколько
<<уточнений>> понятия алгоритма, 
например: 
\odmkey{машина Тьюринга}{Машина (Machine)!Тьюринга (Turing)}, 
\odmkey{машина Поста}{Машина (Machine)!Поста (Post)}, 
\odmkey{нормальный алгорифм Маркова}{Алгоритм (algorithm)!нормальный Маркова (Markov normal)}, и многие другие. 
В сущности, каждое такое уточнение понятия алгоритма 
определяет формальный способ записи предписаний, 
т.~е. программ, и неформальный алгоритм интерпретации программ, 
т.~е. правила применения программ к аргументам для получения результатов. 
Все предложенные к настоящему моменту уточнения интуитивного понятия алгоритма
оказались эквивалентными (в строгом математическом смысле), 
что дало основание выдвинуть 
\odmkey{тезис Тьюринга--Чёрча}{Тезис (thesis)!
Тьюринга--Чёрча (Church--Turing)}, 
неформально постулирующий, что все уточнённые понятия алгоритма 
равнообъёмны интуитивному понятию алгоритма. 
Формальные уточнения понятия <<алгоритм>> 
и построенная на их основе строгая теория алгоритмов 
обладают несомненными достоинствами и подлежат неукоснительному изучению.

Однако теория алгоритмов имеет две имманентные особенности, 
которые в практическом применении проявляют себя как неудобства.
\end{digress}
\end{frame}

\begin{frame}
\frametitle{4.5.1. Алгоритмы (3/5)}

%\vspace{-2pt}
\begin{digress} (Продолжение) 
Во-первых, введение формального способа записи алгоритмов 
резко усложняет и затрудняет \em{программирование}, 
поскольку программировать приходится в терминах очень низкого уровня 
(ячейки ленты машины Тьюринга и тому подобное). 
Во-вторых, способ интерпретации программ 
(\odmkey{выполнения}{Выполнение (execution)} алгоритмов) всё равно 
остается неформальным и потому допускает различные реализации. Поэтому 
затруднительно ввести независящее от реализации средство 
прямого анализа свойств выполнения алгоритмов, 
и приходится применять косвенные критерии, 
подобные асимптотическим оценкам трудоёмкости и так далее.
\end{digress}

\medskip
{\large
В этом курсе тезис Тьюринга–Чёрча принимается безоговорочно; 
активно используемый способ записи программ (псевдокод) 
принимается как достаточно выразительный способ записи алгоритмов;  
убедительные рассуждения относительно текстов алгоритмов 
принимаются как достаточные обоснования свойств этих алгоритмов. 
В целом изложение можно назвать \odmkey{наивной теорией алгоритмов}{Теория (theory)!алгоритмов (of algorithms)!наивная (naive)} 
по аналогии с наивной теорией множеств.
}
\end{frame}

\begin{frame}
\frametitle{4.5.1. Алгоритмы (4/5)}


Не ограничивая общности, далее считается, 
что зафиксирован некоторый конечный алфавит символов $U$, 
достаточный для записи всех необходимых объектов.

\smallskip
Алгоритмы записываются с помощью последовательностей символов из этого алфавита;
данные, к которым применяются алгоритмы, также записываются с помощью
последовательностей символов (возможно, других символов), 
элементы всех множеств записываются с помощью слов в этом алфавите, 
функции действуют из множества слов в множество слов в этом алфавите и т.~д.

\smallskip 
Например, можно считать, что в алфавит $U$ входят кириллические, 
латинские и греческие буквы, прописные и строчные, 
в различных начертаниях, арабские цифры, знаки препинания, скобки, 
математические знаки и т.~д.

\smallskip
Правила записи (кодирования) различных объектов 
словами в алфавите $U$ считаются известными.  
\end{frame}

\begin{frame}
\frametitle{4.5.1. Алгоритмы (5/5)} 

В частности, считаются известными правила записи натуральных чисел 
(например, с помощью десятичных цифр), 
арифметических выражений (например, с помощью знаков операций, 
чисел и переменных), алгоритмов (например, с помощью псевдокода).
Таким образом, принимается, что существуют алгоритмы, 
которые по любому слову в данном фиксированном алфавите $U$ 
позволяют надёжно определить, соответствует ли это слово правилам записи, 
и определить, какой именно объект записан. 
Доказывать существование подобных алгоритмов лексического и 
синтаксического анализа нет нужды~--- это эмпирический факт, 
доступный в повседневных ощущениях при работе с компьютерами.

\smallskip
Имена алгоритмов записываются прописными 
полужирными латинскими буквами в прямом начертании
(например, $\Bl{A}$), 
чтобы легко отличать их от имен множеств и иных объектов. 
Сами алгоритмы записываются на псевдокоде. 
Если алгоритмы имеют входные данные и результаты, 
то они указываются отдельно и называются 
\odmkey{заголовком}{Заголовок (header)!алгоритма (of algorithm)} алгоритма, 
а операторы, задающие алгоритм, называются 
\odmkey{телом}{Тело (body)!алгоритма (of algorithm)} алгоритма.

\end{frame}

\begin{frame}
\label{4.5.2.}\subsubsection{4.5.2. Вычислимые функции }\frametitle{4.5.2. Вычислимые функции (1/5)}

%\vspace{-2pt}
Алгоритмы бывают разные. 
В частности, алгоритмы могут обладать разными 
возможностями в части вырабатывания результата. 
Рассмотрим случай, когда алгоритм $\Bl{A}$ 
принимает на вход некоторые данные 
(\odmkey{аргумент}{Аргумент (argument)}) и перерабатывает их в некоторые другие 
или те же самые данные (\odmkey{результат}{Результат (result)}). 
Тем самым алгоритм $\Bl{A}$ определяет отношение $A$ 
между множеством возможных входных данных (обозначение $\Dom \Bl{A}$) 
и множеством возможных результатов (обозначение $\Im \Bl{A}$),
$A \subset \Dom \Bl{A} \times \Im \Bl{A}.$
Если отношение $A$ обладает свойством функциональности (см. \DMref{п.~1.6.1}{1.6.1.}), 
то говорят, что алгоритм $\Bl{A}$ 
\odmkey{реализует}{Реализация функции (realization of function)! алгоритмом (by algorithm)} \odmkey{вычислимую функцию}{Функция (function)!вычислимая (computable)}. 
Если функция является вычислимой, 
то есть если предъявлен реализующий её алгоритм, 
то может быть построено бесконечно много других алгоритмов, 
также реализующих эту же функцию. 
Действительно, для этого достаточно вставлять 
в тело любого из уже построенных алгоритмов 
произвольные ненужные операторы, 
выполнение которых не влияет на результат. 
Все алгоритмы, реализующие одну и ту же функцию, 
называются \odmkey{функционально эквивалентными}{Алгоритм (algorithm)!функционально эквивалентный \\
 (functionally equivalent)}.

\end{frame}

\begin{frame}
\frametitle{4.5.2. Вычислимые функции (2/5)}

Функциональная эквивалентность, очевидно, 
 является отношением эквивалентности 
на классе алгоритмов и определяет фактормножество
 (\DMref{п.~1.7.3}{1.7.3.})
классов функционально эквивалентных алгоритмов. 
Это фактор-множество взаимно однозначно 
соответствует классу вычислимых функций, 
поэтому можно выбрать какой-либо из функционально эквивалентных алгоритмов, 
объявить его \odmkey{представлением}
{Представление (representation)!функции (of function)} 
вычислимой функции и использовать в качестве функции везде, 
где это необходимо, в частности в других алгоритмах 
для вычисления значения этой функции.

\smallskip 
Это обстоятельство оправдывает используемые 
далее вольности в обозначениях и в словесных оборотах, 
когда вычислимые функции и их представления не различаются, 
подобно тому, как мы говорим о функции, 
предъявляя запись реализующей её формулы (\DMref{п.~3.2.2}{3.2.2.}). 

\smallskip
Для целей наивной теории алгоритмов совершенно не важно, 
какой именно из функционально эквивалентных алгоритмов 
представляет вычислимую функцию. 

\end{frame}



\begin{frame}
\frametitle{4.5.2. Вычислимые функции (3/5)} 
Для практического программирования это, напротив, один из главных вопросов, 
потому что функционально эквивалентные алгоритмы могут 
драматически отличаться друг от друга по длине записи, 
по эффективности работы, по степени понятности и т. д.

\begin{remark}
 Для алгоритмов не известно нормальных форм 
(\DMref{п.~3.4.3}{3.4.3.}), к сожалению. %3.4.2 -> 3.4.3
\end{remark}

\smallskip
В используемом псевдокоде в теле вычислимой функции,
 как правило, присутствует оператор \textbf{return}. 
 
\smallskip
Для записи вычислимых функций применяются те же обозначения, 
что и для записи функций:
$\Map{\Bl{A}}{\Dom \Bl{A}}{\Im \Bl{A}}$, 
$\Im \Bl{A} = \Bl{A}(\Dom \Bl{A})$,
$r = \Bl{A}(a),$
где $a\in \Dom \Bl{A}, r\in \Im \Bl{A}.$

\begin{examples}
\begin{enumerate}
\item Суперпозиция вычислимых функций, очевидно, вычислима. 
Тем самым любая формула над базисом вычислимых функций 
задает вычислимую функцию.
\item Алгоритмы унификации (\DMref{п.~4.3.3}{4.3.3.}) и многие другие 
алгоритмы в этом курсе задают вычислимые функции.

\end{enumerate}
\end{examples}


\end{frame}

\begin{frame}
\frametitle{4.5.2. Вычислимые функции (4/5)}

\begin{remark}
Вычислимая функция может быть частичной, 
т.~е. в записи алгоритма $\Bl{A}$ может быть указано, 
что алгоритм принимает на вход значения типа $T_1$ 
и выдает на выход значения типа $T_2$ 
($\Map{\Bl{A}}{T_1}{T_2}$), 
но при этом, на самом деле, 
$\Dom \Bl{A} \SubsetNE T_1$ и 
$\Im \Bl{A} \SubsetNE T_2$, 
т.~е. область отправления и область прибытия 
шире области определения и области значений соответственно. 
В этом случае считается, что при применении к неподходящему 
аргументу вычислимая функция \em{не вырабатывает никакого результата}
и её поведение не определено.
Тем самым наивная теория алгоритмов 
приближается к суровой программистской практике.
\end{remark} 
\end{frame}

\begin{frame}
\frametitle{4.5.2. Вычислимые функции (5/5)}
\vspace{-2pt}
\begin{digress}
Понятие <<тип>>, или <<тип данных>>, использованное в предыдущем замечании, 
имеет важное значение в программировании. 
С точки зрения наивной теории алгоритмов 
\odmkey{тип}{Тип (type)}~--- это некоторое перечислимое 
(и разрешимое, см. \DMref{п.~4.5.4}{4.5.4.}) множество слов в алфавите $U$. %4.5.3 -> 4.5.4
В программировании понятие типа (и родственное понятие класса) 
применяется, в частности, для проверки правильности применения 
(частичных) вычислимых функций и алгоритмов. 
А именно, указывая типы параметров функции, 
программист задает (до некоторой степени) область определения функции, 
и система программирования может проверить, 
принадлежит переданный аргумент указанному типу, 
или нет (это возможно, поскольку тип~--- разрешимое множество). 
Результаты этой проверки могут быть использованы 
для диагностирования или предотвращения ошибки и для других полезных действий. % и других полезных действий -> и для других полезных действий
В этом разделе нам не понадобилась теория типов, 
но это важное и практически полезное ответвление наивной теории алгоритмов.
\end{digress}

\end{frame}


\begin{frame}
\label{4.5.3.}\subsubsection{4.5.3. Перечислимые множества }\frametitle{4.5.3. Перечислимые множества (1/7)}

%\vspace{-3pt}
Рассмотрим теперь случай, когда алгоритм $\Bl{A}$
не принимает на вход никаких данных, а будучи запущенным, 
вырабатывает некоторую последовательность результатов (см. \DMref{п.~1.1.2}{1.1.2.}). 
Тем самым алгоритм $\Bl{A}$ определяет некоторое множество $A$, 
являющееся множеством слов в алфавите $U$, 
которое называется \odmkey{перечислимым}{Множество (set)!перечислимое (enumerable)}:

\vspace{-13pt}
$$A = \Set{a\in U^*}{a\Assign\Bl{A}()}.$$

\vspace{-4pt}
Совершенно аналогично вычислимым функциям, 
перечислимые множества могут быть реализованы бесконечным множеством алгоритмов,
которые все эквивалентны в том смысле, что перечисляют одно и то же множество. Условимся для каждого перечислимого множества выбирать
\odmkey{представление}{Представление (representation)} в классе эквивалентных алгоритмов
и не различать перечислимое множество и его представление без нужды. 
Алгоритм, представляющий перечислимое множество, 
для краткости речи называется просто \odmkey{перечисляющим}{Алгоритм (algorithm)!перечисляющий (enumerating)}.

В используемом псевдокоде в теле алгоритма, перечисляющего множество, 
как правило, присутствует оператор \textbf{yield} $a$.
\end{frame}


\begin{frame}
\frametitle{4.5.3. Перечислимые множества (2/7)}
%\vspace{-2pt}
\begin{remark} {\color{blue} 1.} 
 Иногда оператор \textbf{yield} $a$ отсутствует в явном виде, 
 но тогда он всегда моделируется оператором $A\Assign A+a$ 
 добавления элемента $a$ в изначально пустое множество $A$.
\end{remark}

\medskip
\begin{remark} {\color{blue} 2.}  
Всякое конечное множество может быть задано перечислением элементов 
$\{\Tuple{a}{1}{n}\}$ (\DMref{п.~1.1.2}{1.1.2.}) и тем самым является перечислимым.
Действительно, алгоритм 
\textbf{begin~yield}~$a_1$;~...;~\textbf{yield}~$a_n$~\textbf{end} 
перечисляет это множество. 
В частности, алфавит $U$ перечислим, 
перечисляющий его алгоритм обозначим $\Bl{U}$.
\end{remark}

\medskip
\begin{remark} {\color{blue} 3.}
Всякое перечислимое множество $A$ линейно упорядочено (\DMref{п.~1.8.2}{1.8.2.}). 
Действи\-тельно, порядок выполнения операторов 
\textbf{yield} (или оператора добавления элемента) линеен.
\end{remark}

\medskip
\begin{remark} {\color{blue} 4.}
Перечислимые множества (и только они!) 
могут появляться в заголовке цикла по множеству 
\textbf{for $a\in A$ do $\Bl{B}(a)$ end for}, 
где $\Bl{B}$~--- некоторый алгоритм. 
Действительно, чтобы реализовать этот цикл, 
достаточно в теле перечисляющего алгоритма 
$\Bl{A}$ заменить оператор \textbf{yield $a$} вызовом алгоритма $\Bl{B}(a)$.
\end{remark}
\end{frame}


\begin{frame}
\frametitle{4.5.3. Перечислимые множества (3/7)}

%\vspace{-2pt}
\begin{remark} {\color{blue} 5.}
Конкретные перечисляющие алгоритмы, приводимые в примерах, 
генерируют каждый элемент ровно один раз, как это и требуется в обычных множествах.
Однако явно накладывать такое ограничение на перечисляющие алгоритмы необязательно.
Действительно, если имеется перечисляющий алгоритм, который перечисляет какие-то
элементы несколько раз, то в его тело можно вставить проверку, 
исключающую повторные добавления элементов. 
Перечисляющие алгоритмы, генерирующие элементы по не\-сколько раз, 
можно использовать для работы с мультимножествами, 
которые в этом разделе не понадобились.
\end{remark}


\begin{examples}
\begin{enumerate}
\item
Множество натуральных чисел перечислимо:\\
$\Bbb{N} =\Set{n}{ n\Assign 0;\ \textbf{while true do}\ n:=n+1; 
\textbf{yield}\ n\ \textbf{end while}}.$
\item
Алгоритмы пп.~\DMref{1.3.2}{1.3.2.}--\DMref{1.3.4}{1.3.4.} задают перечислимые множества.
\item
Пустое множество $\emptyset$ и множество всех слов $U^*$ перечислимы. 
Пустое множество перечислимо пустым алгоритмом \textbf{skip}, 
а множество всех слов перечислимо  алгоритмом, 
который вначале перечисляет все слова длины 1, 
потом все слова длины 2 ...
\end{enumerate}
\end{examples}
\end{frame}


\begin{frame}
\frametitle{4.5.3. Перечислимые множества (4/7)}

%\vspace{-2pt}
\begin{theorem}
Существуют невычислимые функции и неперечислимые множества.
\end{theorem}

\begin{proof}
Множество различных слов в непустом конечном алфавите счётно. \\
Всякий алгоритм~--- это слово в конечном алфавите, 
следовательно, множество различных алгоритмов не более чем счётно. 
Фактор-множество (\DMref{п.~1.7.3}{1.7.3.}) не может превосходить исходное множество по мощности, %1.7.2 -> 1.7.3
поэтому множество классов функционально эквивалентных алгоритмов не более чем счётно.
С другой стороны, множество всех подмножеств счётного множества несчётно.
Значит число различных множеств слов и 
число различных функций из множеств слов в множества слов несчётно. 
Следовательно, невозможно взаимно-однозначное 
соответствие из множества алгоритмов в 
множество перечислимых множеств и в множество вычислимых функций.
\end{proof}

\begin{corollary} 
Существуют перечислимые множества и вычислимые функции, причём множества вычислимых функций и перечислимых множеств не более чем счётны.
\end{corollary}

\end{frame}

\begin{frame}
\frametitle{4.5.3. Перечислимые множества (5/7)}

\begin{remark}
Наряду с вычислимыми функциями и перечислимыми множествами 
бывает полезно рассматривать и другие типы алгоритмов, 
например: итераторы, трансформации и др. 
\end{remark}

\smallskip
\begin{example}
Множество $\Bbb{N}$ перечислимо, 
множество $\Bbb{N}^2 = \Bbb{N}\times\Bbb{N}$ также перечислимо. \\
Например, следующий алгоритм $\Bl{N2}$ 
перечисляет пары натуральных чисел $(a, b)$:\\
\textbf{ $n\Assign 0;$ while true do $n \Assign n+1;$
 for $b$ from $1$ to $n$ do yield $(n-b+1, b)$ end for end while}.
\end{example}

\smallskip
Перечислимость множеств тесно связана 
с вычислимыми функциями натурального ряда.
Действительно, перечислимое множество $\Bbb{N}$ 
(и его отрезки, \DMref{п.~1.2.5}{1.2.5.})~--- это удобное и естественное средство \em{нумерации},
то есть перечисления множеств слов. 
Пусть $A$~--- перечислимое множество, 
заданное алгоритмом $\Bl{A}$. 
Построим  вычислимую функцию 
$\Map{F}{\Bbb{N}}{A}$, возвращающую элемент перечислимого множества 
по его номеру, и обратную функцию $\Map{F^{-1}}{A}{\Bbb{N}}$, 
возвращающую номер заданного элемента: 

\end{frame}

\begin{frame}
\frametitle{4.5.3. Перечислимые множества (6/7)}

\vspace{-1.5\baselineskip}
\begin{columns}[t]
\begin{column}{6.5cm}
\begin{algorithmic}
\STATE{ \textbf{func} $F (n: \Bbb{N}) : A$} 
\STATE{ $k\Assign 0$} 
\FOR {$a\in A$} 
  \STATE{ $k \Assign k+1$} 
  \STATE{ \textbf{if $k = n$ then return $a$ end if}}
\ENDFOR
\STATE{\textbf{return fail}}
\end{algorithmic}
\end{column}
\begin{column}{6.5cm}
\begin{algorithmic}
\STATE{ \textbf{func} $F^{-1} (x: A ) : \Bbb{N}$} 
\STATE{ $k\Assign 0$} 
\FOR {$a\in A$} 
  \STATE{ $k \Assign k+1$} 
  \STATE{ \textbf{if $x = a$ then return $k$ end if}}
\ENDFOR
\STATE{\textbf{return fail}}
\end{algorithmic}
\end{column}
\end{columns}

\medskip
Перечисляющую функцию $F$ всегда можно и иногда удобно использовать 
в качестве представления перечислимого множества наряду с генерирующим алгоритмом.

\smallskip
\begin{remark}
В этих алгоритмах показан один из способов работы
с частичными вычислимыми функциями в используемом псевдокоде.
Если значение функции не определено, то
возвращается специальное значение \textbf{fail}.
\end{remark}

\begin{example}
Формула $2^a (2b + 1)-1$ задает номер пары в
некотором перечислении натуральных пар $(a, b)$.
\end{example}


\end{frame}



\begin{frame}
\frametitle{4.5.3. Перечислимые множества (7/7)}

%\vspace{-2pt}

\begin{theorem}
 Объединение и пересечение перечислимых множеств перечислимы.
\end{theorem}
 
\begin{proof}
Пусть $X$ и $Y$~--- перечислимые множества
с перечисляющими функциями $\Map{\Bl{X}}{\Bbb{N}}{X}$ и
$\Map{\Bl{Y}}{\Bbb{N}}{Y}$.
{\color{blue}\proofsection{$\cup$}}
Если хотя бы одно из множеств пусто, 
то утверждение тривиально. 
Иначе вычислимая функция $\Map{\Bl{Z}}{\Bbb{N}}{X\cup Y}$
перечисляет объединение:\\
\textbf{$\Bl{Z}(n) \Assign$ if $n = (n\DIV 2)*2$ then return $\Bl{X}(n\DIV 2)$
else return $\Bl{Y}((n+1)\DIV 2)$ end if}.
{\color{blue}\proofsection{$\cap$}}
Если пересечение пусто, то оно перечислимо по определению. 
Иначе рассмотрим алгоритм $\Bl{N2}$, 
перечисляющий пары натуральных чисел (см. пример выше), 
и построим алгоритм, перечисляющий пересечение:\\
\textbf{for $(a, b)\in \Bl{N2}$ 
do if $\Bl{X}(a)=\Bl{Y}(b)$ then yield $\Bl{X}(a)$ end if end for}.
\end{proof}

\begin{remark}
Другие примеры перечисляющих алгоритмов 
в форме итераторов приведены в \DMref{п.~1.3.9}{1.3.9.}. 
\end{remark}
\end{frame}

\begin{frame}
\label{4.5.4.}\subsubsection{4.5.4. Разрешимые множества }\frametitle{4.5.4. Разрешимые множества (1/4)}
%\vspace{-0.25\baselineskip}
Подмножество слов $S$ перечислимого множества $M$ называется 
\odmkey{разрешимым}{Подмножество (subset)!разрешимое (solvable)}, если существует такой алгоритм
$\Map{\Bl{S}}{M}{\{\FF,\TT\}}$,  что $\Bl{S}(m)=\TT$ тогда и только тогда, 
когда $m\in S$.
Другими словами, множество разрешимо, если существует алгоритм,
который распознаёт принадлежность элементов этому множеству. 
Ясно, что если множество $S$ разрешимо относительно $M$, 
то множество $M \setminus S$ также разрешимо~--- 
достаточно в теле алгоритма $\Bl{S}$ 
заменить оператор \textbf{return}~$\TT$ 
оператором \textbf{return}~$\FF$ и наоборот.

\medskip
\begin{lemma} \NB{[1]}
Для любого $X$, $\emptyset$ и $X$ разрешимы относительно $X$. 
\end{lemma}

\begin{proof}
Для пустого множества искомый алгоритм~--- 
константа $\FF$, а для множества $X$~--- константа $\TT$. 							\end{proof}

\medskip
\begin{lemma} \NB{[2]}
Разрешимое подмножество $S$ перечислимого множества $M$ перечислимо.
\end{lemma}

\begin{proof} 
Пусть $\Bl{S}$~--- разрешающий алгоритм для $S$, 
а $\Bl{M}$~--- перечисляющий алгоритм для $M$. 
Построим перечисляющий алгоритм $\Bl{S}_1$ для $S$.\\ 
%Если $S=, то S := F, если же S, то существует элемент sS. Тогда 
\textbf{$\Bl{S}_1 \Assign$ for $m\in M$ do if
$\Bl{S}(m)$ then yield $m$ end if end for}.
\end{proof}


\end{frame}


\begin{frame}
\frametitle{4.5.4. Разрешимые множества (2/4)}
Отсюда немедленно следует, что всякое разрешимое подмножество натурального ряда перечислимо.

\medskip
\begin{lemma} \NB{[3]}
Подмножество $S$ перечислимого множества $M$ 
разрешимо тогда и только тогда, когда перечислимо как множество $S$, 
так и его дополнение $M \setminus S$.
\end{lemma}

\begin{proof} 
{\color{blue}\proofsection{$\Rightarrow$}}
Если множество $S$ разрешимо относительно $M$, 
то множество $M \setminus S$ также разрешимо, 
и остаётся применить лемму 2.
{\color{blue}\proofsection{$\Leftarrow$}} 
Если $S$ или $M \setminus S$ пусто, 
то по лемме 1 $S$ разрешимо. 
Пусть теперь оба множества 
$S$ и $M \setminus S$ не пусты и перечислимы перечисляющими алгоритмами
$\Bl{S}$ и $\Bl{S}_1$ соответственно, 
а вычислимые функции 
$\Map{F, F_1}{\Bbb{N}}{M}$ доставляют элементы этих множеств по номерам. 
Тогда следующий алгоритм разрешает множество $S$:
\end{proof}

\end{frame}


\begin{frame}
\frametitle{4.5.4. Разрешимые множества (3/4)}

\begin{proof} (Продолжение)
\begin{algorithmic}
\STATE{ \textbf{func} $SS (s: M ) : \{\FF,\TT\}$} 
\STATE{ $n\Assign 0$} 
\WHILE {\textbf{true}}
  \STATE{$n\Assign n+1$; $x \Assign F(n)$; $y \Assign F_1(n)$}  
  \STATE{ \textbf{if $x = s$ then return $\TT$ end if}} 
  \STATE{ \textbf{if $y = s$ then return $\FF$ end if}} 
\ENDWHILE
\end{algorithmic}
\end{proof}



\begin{lemma} \NB{[4]}
 Если функция вычислима, то образ перечислимого множества перечислим. 
\end{lemma}
%\vspace{-0.2\baselineskip} 
\begin{proof}
 Пусть функция $\Map{\Bl{F}}{\Dom \Bl{F}}{\Im \Bl{F}}$ вычислима, 
 а множество $A \subset \Dom \Bl{F}$ перечислимо, 
 тогда множество $\Bl{F}(A) \subset \Im \Bl{F}$
  перечисляется следующим алгоритмом:\\
\textbf{for $x\in A$ do yield $\Bl{F}(x)$ end for}.	
\end{proof}
\end{frame}


\begin{frame}
\frametitle{4.5.4. Разрешимые множества (4/4)}



\begin{lemma} \NB{[5]}
Если функция вычислима и определена на перечислимом множестве, то прообраз перечислимого множества перечислим. 
\end{lemma}
%\vspace{-0.2\baselineskip} 
\begin{proof} Пусть функция $\Map{\Bl{F}}{\Dom \Bl{F}}{\Im \Bl{F}}$ вычислима, 
а множество $B$ перечислимо вычислимой функцией $h$, 
и множество $\Dom \Bl{F}$  перечислимо вычислимой функцией $g$. 
Тогда множество $\Bl{F}^{-1}(B)$ перечисляется следующим алгоритмом:

\textbf{for $(a, b)\in \Bl{N2}$ do if $\Bl{F}(g(a))=h(b)$ then yield $g(a)$ 
end if end for}.
\end{proof}

\medskip
\begin{examples}
Конечные множества разрешимы.
Бесконечные множества чётных и нечётных, 
а также простых и составных натуральных чисел разрешимы.
\end{examples}

\medskip
{\large
Существуют неразрешимые множества, в том числе неразрешимые 
подмножества натуральных чисел, поскольку множество алгоритмов счётно,
а множество подмножеств натуральных чисел несчётно.
}
\end{frame}


\begin{frame}
\label{4.5.5.}\subsubsection{4.5.5. Протокол выполнения алгоритма }\frametitle{4.5.5. Протокол выполнения алгоритма (1/4)}

Программистская практика свидетельствует, 
что программы иногда удаётся отладить, 
то есть найти и устранить ошибки, содержащиеся в алгоритмах. 
Процесс отладки основан на том, 
что за работой алгоритма возможно наблюдение, в частности, 
для любого алгоритма и для любого случая выполнения этого алгоритма 
можно составить \odmkey{протокол выполнения}{Протокол выполнения (execution protocol)}, 
в котором исчерпывающим образом отражается история выполнения алгоритма 
(для заданных входных данных). Другими словами, 
в самих алгоритмах и в их выполнении нет ничего <<чудесного>>, 
ничего такого, чего нельзя записать словами в алфавите $U$.  

\medskip
Следующий факт мы принимаем без доказательства, 
как априорное свойство интуитивного понятия алгоритма.

{\bf Аксиома {\color{blue}[протокола]}.} 
{\it Для каждого алгоритма $\Bl{A}$ существует разрешимое множество
$H$ и существуют вычислимые функции $\Map{a}{H}{U^*}$ и $\Map{z}{H}{U^*}$, 
такие что $\Bl{A}(x)=y$ тогда и только тогда, 
когда существует такое $h\in H$, что $a(h)=x$ и $z(h)=y$.
}


\end{frame}


\begin{frame}
\frametitle{4.5.5. Протокол выполнения алгоритма (2/4)}
Множество $H$~--- это множество всех возможных протоколов 
выполнения алгоритма $\Bl{A}$, $h$~--- это протокол выполнения
с входом $x$ и выходом $y$, а  функции $a$ и $z$~--- 
это функции определения входов и выходов соответственно.

\smallskip
\begin{digress}
Если для каждой пары $(x, y)$, для которой 
$\Bl{A}(x)=y$ слово $h\in H$, 
такое что  $a(h)=x$ и $z(h)=y$ не только существует, 
но и единственно, то алгоритм называется 
\odmkey{детерминированным}{Алгоритм (algorithm)!детерминированный (deterministic)}. 
В протоколе детерминированного алгоритма 
порядок действий строгий линейный (\DMref{п.~1.8.2}{1.8.2.}). % 1.8.1 -> 1.8.2
В сформулированной аксиоме протокола единственность не требуется, 
тем самым допускается, что выполнение алгоритма может быть \odmkey{недетерминированным}{Алгоритм (algorithm)!недетерминированный (non-deterministic)}. 
В протоколе недетерминированного алгоритма порядок действий 
нестрогий частичный. В недетерминированности нет ничего страшного, 
более того, это путь к параллельному программированию 
и другим современным информационным технологиям. 
Поэтому накладывать требование детерминированности 
без необходимости не следует.
\end{digress}

\end{frame}


\begin{frame}
\frametitle{4.5.5. Протокол выполнения алгоритма (3/4)}

\begin{corollary} \NB{[1]}
Область применимости и множество результатов любого алгоритма перечислимы.
\end{corollary}
\begin{proof}
Область применимости~--- это $a(H)$, 
а множество результатов~--- это $z(H)$. 
Оба эти множества перечислимы по лемме 4.
\end{proof}

\smallskip
\begin{corollary} \NB{[2]}
Область определения и множество значений любой вычислимой функции перечислимы.
\end{corollary}
\begin{proof} 
Частный случай cледствия 1.
\end{proof}

\smallskip
\begin{corollary} \NB{[3]}
 График любой вычислимой функции перечислим.
\end{corollary}
\begin{proof} 
По аксиоме протокола для функции существует множество 
протоколов $H$ и вычислимые функции входа/выхода $a(h)$ и $z(h)$. 
Тогда следующий алгоритм перечисляет график функции:
\textbf{  for $h\in H$ do yield $(a(h), z(h))$ end for}
\end{proof}
\end{frame}


\begin{frame}
\frametitle{4.5.5. Протокол выполнения алгоритма (4/4)}
\begin{corollary} \NB{[4]}
 Функция с перечислимым графиком вычислима.
\end{corollary}
\begin{proof} 
Пусть алгоритм $\Bl{F}$ перечисляет график функции $f$, 
то есть множество $\Set{(a,b)}{b=f(a)}$. 
Тогда следующий алгоритм вычисляет значение функции $f(x)$:

\textbf{for $(a,b)\in \Bl{F}$ do if $x=a$ then return $b$ end if end for}
\end{proof}
\end{frame}


\begin{frame}
\label{4.5.6.}\subsubsection{4.5.6. Программы }\frametitle{4.5.6. Программы (1/6)}
%\vspace{-2pt}
Алгоритмы могут быть записаны как слова в универсальном алфавите $U$, 
причем такая запись заведомо не является единственной. 
Возникает естественный вопрос: если даны две записи, 
причём обе синтаксически распознаются как записи алгоритмов, 
то что можно сказать о свойствах алгоритмов по их записям? 
Например, это две записи одного алгоритма или разных алгоритмов? 
Ясно, что ответить на эти вопросы не так просто. 
Необходимо внести уточнения в описание способа записи алгоритмов.

\medskip 
Введем обозначения. 
Напомним, что мы рассматриваем алгоритмы над словами в алфавите $U$, 
т.~е. алгоритмы принимают в качестве аргументов слова 
(возможно, не принимают аргументов) 
и вырабатывают в качестве результата также слова 
(возможно, много слов). 

\medskip
Алгоритмы, в том числе вычислимые функции, не обязательно тотальны. 
Если аргументы принадлежат области определения алгоритма, 
то говорят, что алгоритм \odmkey{применим}{Алгоритм (algorithm)!применимый (applicable)}.
\end{frame}


\begin{frame}
\frametitle{4.5.6. Программы (2/6)}
%\vspace{-2pt} 
Говорят, что два алгоритма $\Bl{A}$ и $\Bl{B}$ 
\odmkey{условно равны}
{Алгоритм (algorithm)!условно равный (conventionally equal)} для аргумента $x$, 
обозначение \\ $\Bl{A}(x) \simeq \Bl{B}(x)$, 
если либо оба алгоритма применимы к аргументу $x$,
 и тогда их результаты совпадают, 
 либо оба алгоритма неприменимы к $x$.
 
\begin{remark}  
Алгоритмы в этом определении являются не только алгоритмами вычислимых функций, 
но любыми алгоритмами. Поэтому результатом работы может быть 
не только возвращаемое значение, но и перечислимое множество, 
или последовательность действий в случае итераторов, 
или изменение состояния объекта в случае трансформаций и т.д. 
Совпадение результатов для каждого типа алгоритмов 
определяется специфическим образом. 
Для вычислимых функций это равенство слов, 
для перечислимых множеств~--- равенство множеств слов.
\end{remark}

Два алгоритма $\Bl{A}$ и $\Bl{B}$ называются 
\odmkey{равносильными}{Алгоритм (algorithm)!равносильный (equivalent)}, если у них совпадают области применимости, 
и для любого слова из этой области, взятого в качестве аргумента, 
совпадают результаты работы алгоритмов, 
то есть $\A{x}{\Bl{A}(x) \simeq \Bl{B}(x)}$. 
\end{frame}


\begin{frame}
\frametitle{4.5.6. Программы (3/6)}

Равносильность обычно обозначается следующим образом: 
$\Bl{A}\sim\Bl{B}$. 
Ясно, что равносильность алгоритмов является эквивалентностью.



\begin{remark}
Все эти определения применимы к любым алгоритмам, 
в частности, к вычислимым функциям. 
Поэтому равносильность и функциональная эквивалентность~--- 
это равнообъёмные понятия.
\end{remark}

\vspace{0.25\baselineskip}
Говорят, что два алгоритма $\Bl{A}$ и $\Bl{B}$ 
\odmkey{всюду отличаются}{Алгоритм (algorithm)!всюду отличающийся (everywhere different)}, если $\Not\E{x}{\Bl{A}(x) \simeq \Bl{B}(x)}$. 
Другими словами, алгоритмы всюду отличаются, если дают разные результаты там, 
где применимы, и нет таких аргументов, при которых они оба неприменимы.

\vspace{0.25\baselineskip} 
Пусть имеется алгоритм с двумя параметрами 
$\Map{\Bl{A}}{X\times Y}{Z}$. 
\odmkey{Проекцией}{Проекция (projection)!алгоритма (algorithm)} этого алгоритма на первый (второй) параметр, 
или \odmkey{остаточной программой}{Программа (programme)!остаточная (residual)} по первому (второму) аргументу, 
называется алгоритм $\Map{\Bl{A}_x}{Y}{Z}$  ($\Map{\Bl{A}_y}{X}{Z}$) такой,\\ 
что $\A{y\in Y}{\Bl{A}_x(y) \simeq \Bl{A}(x, y)}$
(соответственно $\A{x\in X}{\Bl{A}_y(x) \simeq \Bl{A}(x, y)}$).

\end{frame}


\begin{frame}
\frametitle{4.5.6. Программы (4/6)}
 
\begin{remark}
Термин \em{остаточная программа} заимствован из теории смешанных вычислений.
Концепцию смешанных вычислений предложил А.~П.~Ершов.
\end{remark} 
  
Пусть есть некоторый класс алгоритмов $\Bl{\cal{A}}$. 
Класс алгоритмов $\Bl{\cal{B}}$ называется 
\odmkey{представительным}{Класс (class)!алгоритмов (of algorithm)!представительный (representative)} для класса
 $\Bl{\cal{A}}$, если для любого алгоритма $\Bl{A}\in \Bl{\cal{A}}$ 
существует равносильный алгоритм $\Bl{B}\in \Bl{\cal{B}}$,
$\Bl{A} \sim \Bl{B}$. 

%\vspace{0.25\baselineskip} 
\begin{remark}
Отношение представительности между классами алгоритмов 
не является эквивалентностью, но является рефлексивным и транзитивным, 
то есть является отношением предпорядка 
(\DMref{п.~1.8.1}{1.8.1.}). %добавлено (п. 1.8.1)
\end{remark}

\vspace{0.25\baselineskip}  
Идея состоит в том, чтобы для любого класса алгоритмов $\Bl{\cal{A}}$ 
подобрать более чётко очерченный (но, возможно, более широкий!) 
класс алгоритмов $\Bl{\cal{B}}$, 
который будет представительным по отношению к исходному классу. 
Элементы класса $\Bl{\cal{B}}$ называются 
\odmkey{программами}{Программа (programme)}, которые 
\odmkey{представляют}{Представление (representation)!алгоритма (of algorithm)} (или \odmkey{реализуют}{Реализация (implementation)!алгоритма (of algorithm)}) элементы класса
 $\Bl{\cal{A}}$.


 \end{frame}


\begin{frame}
\frametitle{4.5.6. Программы (5/6)}

%\vspace{-4pt}
Осталось определить, в каком смысле класс программ должен быть чётко очерчен.\\ 
Во-первых, класс программ должен быть разрешимым~--- требуется надёжно отличать
 программы от не программ. 
 Во-вторых, должен быть задан \odmkey{алгоритм интерпретации}{Алгоритм (algorithm)!интерпретации (interpretation)} $\Bl{I}$, 
 позволяющий по программе $\Bl{B}$ и исходным данным $x$ 
 определить результаты представляемого алгоритма $\Bl{A}(x)$~--- 
 программы необходимо уметь \em{выполнять}.

\smallskip
Всё сказанное приводит к следующей аксиоме.

{\bf Аксиома {\color{blue}[программы]}.}
{\it Существует представительное разрешимое множество 
программ $\Bl{\cal{P}}$ такое, что для любого алгоритма 
$\Map{\Bl{A}}{T_1}{T_2}$ существует равносильная программа 
$\Bl{P}\!\in\!\Bl{\cal{P}}$, $\Bl{A}\sim \Bl{P}$ 
и существует алгоритм интерпретации 
$\Map{\Bl{I}}{\Bl{\cal{P}}\!\times\!T_1}{T_2}$, 
причём $\Bl{A}\sim \Bl{I}_{\Bl{P}}$.
}

\begin{remark}
 Алгоритм $\Bl{I}$, как и все алгоритмы, может быть запрограммирован, 
 то есть для него может быть предъявлена равносильная программа. 
 Полученная программа является \em{самоприменимой}. 
 В этом нет ничего парадоксального.
\end{remark}


\end{frame}


\begin{frame}
\frametitle{4.5.6. Программы (6/6)}
\begin{digress} 
Этой схеме следуют все упомянутые выше теоретические уточнения понятия алгоритма.
Каждое такое уточнение описывает свой класс <<программ>> (более или менее формально),
неформально описывает алгоритм $\Bl{I}$ и в качестве недоказываемой догмы 
постулирует представительность выбранного класса программ.

\smallskip
Но этой же схеме следуют все практические системы программирования! 
Каждая такая система описывает свой класс программ 
(задает \odmkey{синтаксис}{Синтаксис (syntax)} языка программирования), 
полуформально описывает алгоритм $\Bl{I}$ 
(операционная \odmkey{семантика}{Семантика (semantics)} языка программирования), 
а <<выразительная сила>> языка программирования рассматривается как предмет веры. 

\smallskip
В том, что класс программ \em{шире} класса алгоритмов 
(если программа синтаксически правильна, это ещё не значит, 
что она закончит работу и даст правильный результат) 
программисты очень быстро убеждаются на собственном печальном опыте.
\end{digress}

\end{frame}


\begin{frame}
\label{4.5.7.}\subsubsection{4.5.7. Невычислимые функции и неразрешимые множества }\frametitle{4.5.7. Невычислимые функции и неразрешимые множества (1/5)}  

Рассмотрим более подробно важное понятие 
\odmkey{универсальной функции}{Функция (function)!универсальная (universal)}, 
или \odmkey{универсального алгоритма}{Алгоритм (algorithm)!универсальный (universal)}. 
Рассмотрим произвольный класс $F$ вычислимых функций 
$\Map{f}{X}{Y}$. 
Функция $\Map{U_F}{F\times X}{Y}$
называется \odmkey{универсальной}{Функция (function)!универсальная (universal)} для класса $F$, 
если \\$\A{f}{\A{x}{U_F(f,x)\simeq f(x)}}$. 
Другими словами, там, где исходная функция определена, 
универсальная функция вычисляет значение по функции и по аргументу, 
а там где функция не определена, универсальная функция также не определена.

\begin{examples}
\begin{enumerate}
\item Алгоритм в \DMref{п.~3.5.3}{3.5.3.} является универсальным алгоритмом
 в классе булевых функций $n$ переменных.
\item Алгоритм  в \DMref{п.~3.2.1}{3.2.1.} является универсальным алгоритмом в классе функций,
 заданных формулами в базисе $F$.
\item Интерпретатор $\Bl{I}$ является универсальным 
алгоритмом в классе программ.  
\end{enumerate}
\end{examples} 

\end{frame}

\begin{frame}
\frametitle{4.5.7. Невычислимые функции и неразрешимые множества (2/5)}  
Множество программ, представляющих вычислимые функции, 
перечислимо и их можно перенумеровать. 
Интуитивно ясно, что возможно построить программу 
(например, в форме <<очень длинного>> оператора \textbf{switch}), 
которая по номеру функции и значению аргумента 
будет вычислять значение этой функции. Таким образом, из аксиомы программы имеем следствие 1.

\vspace{0.2\baselineskip}  
\begin{corollary} \NB{[1]}
Каков бы ни был класс $F$ вычислимых функций 
$T_1 \to T_2$, существует вычислимая функция $\Map{U_F}{\Bbb{N}\times T_1}{T_2}$, универсальная в этом классе.
\end{corollary}

\vspace{0.2\baselineskip}  
В данном разделе наибольший интерес представляют вычислимые функции с натуральными аргументами и результатами.

\vspace{0.2\baselineskip}  
\begin{corollary} \NB{[2]} 
Существует вычислимая функция $\Map{U_N}{\Bbb{N}\times \Bbb{N}}{\Bbb{N}}$, универсальная в классе вычислимых функций $\Bbb{N}\to\Bbb{N}$.
\end{corollary}
\begin{proof}
 Положим $T_1\Assign \Bbb{N}$ и $T_2 \Assign \Bbb{N}$ и применим следствие 1.
\end{proof}

\end{frame}


\begin{frame}
\frametitle{4.5.7. Невычислимые функции и неразрешимые множества (3/5)}  
\begin{corollary} \NB{[3]}  
 Существует такая вычислимая функция $\Map{d}{\Bbb{N}}{\Bbb{N}}$, 
 что никакая вычислимая функция не может отличаться от неё всюду.
\end{corollary}
\begin{proof} 
Рассмотрим <<диагональную>> функцию 
$d(n) \Assign U_N(n, n)$. Пусть $f_n$~--- функция, 
имеющая номер $n$. Если $f_n(n)$ определено, 
то $f_n(n) = d(n)$ и $f_n(n) \simeq d(n)$, 
если же $f_n(n)$ не определено, то и $d(n)$ не определено, 
и все равно $f_n(n)\simeq d(n)$.
\end{proof}

\medskip
Предыдущее доказательство кажется парадоксом: 
действительно, если у нас есть функция, 
которая претендует на то, что, неформально говоря, её график 
<<пересекается>> с графиками всех других функций, 
то, казалось бы, возьмём функцию, 
график которой <<параллелен>> графику <<претендента>>, 
и получится контрпример к следствию 3. 
Но никакого парадокса и контрпримера нет! 
Дело в том, что мы рассматриваем \em{все} вычислимые функции, 
в том числе и частичные, и <<пересечение>> получается там, 
где функции не определены.
\end{frame}


\begin{frame}
\frametitle{4.5.7. Невычислимые функции и неразрешимые множества (4/5)}  
Но если вычислимая функция определена частично, 
то можно доопределить её произвольным образом, 
назначив какие-то значения функции для тех значений аргументов, 
где функция не определена. Такое <<доопределение>> называется 
\odmkey{всюду определённым продолжением}{Продолжение (extension)!всюду определённое (completely defined)}.

\medskip 
 \begin{corollary} \NB{[4]}   
Существует вычислимая функция $\Map{d_1}{\Bbb{N}}{\Bbb{N}}$, 
не имеющая всюду определённого вычислимого продолжения.
\end{corollary}
\begin{proof}  
Рассмотрим <<параллельную>> функцию $d_1(n) \Assign d(n)+1$. 
Эта функция определена частично, потому что функция $d(n)$ определена частично.
Рассмотрим тогда $D_1(n)$~--- всюду определённое продолжение $d_1(n)$. Как именно
продолжается $d_1$ несущественно.
Функция $D_1$ всюду определена и всюду отличается от $d(n)$. 
По следствию 3 заключаем, что функция $D_1(n)$ невычислима.
\end{proof}							

\medskip
Таким образом, построена всюду определённая, но невычислимая функция. 
\end{frame}


\begin{frame}
\frametitle{4.5.7. Невычислимые функции и неразрешимые множества (5/5)} 
\vspace{-2pt}
\begin{corollary} \NB{[5]} 
 Существует неразрешимое перечислимое подмножество натурального ряда.
\end{corollary}
\begin{proof}   
Рассмотрим множество
$ N\Assign \Dom d_1$. Множество $N$ перечислимо по лемме 5. 
Допустим, что множество $N$ разрешимо некоторым алгоритмом $\Bl{N}$. 
Тогда построим функцию \textbf{$d_2(x) \Assign$ if $\Bl{N}(x)$ then $d_1(x)$ 
else $1$ end if}. Функция $d_2(x)$ оказывается всюду определённым 
вычислимым продолжением функции $d_1$,
что невозможно в силу следствия 4, 
а значит алгоритма $\Bl{N}$ не существует.
\end{proof}

Итак, построено перечислимое, но неразрешимое множество.

\begin{corollary} \NB{[6]} 
Существует перечислимое подмножество натурального ряда с неперечислимым дополнением.
\end{corollary}
\begin{proof}  
Рассмотрим построенное перечислимое неразрешимое множество $N$. 
Его дополнение $\Bbb{N}\setminus N$ неперечислимо, 
в противном случае $N$ было бы разрешимо по лемме 3.
\end{proof}

Таким образом, построено перечислимое множество с неперечислимым дополнением. 


\end{frame}


\begin{frame}
\label{4.5.8.}\subsubsection{4.5.8. Истинность и доказуемость }\frametitle{4.5.8. Истинность и доказуемость (1/6)}  

%\vspace{-4pt}
К данному моменту накоплен достаточный запас фактов, 
чтобы показать применимость и полезность наивной теории алгоритмов 
в форме доказательства первой теоремы Гёделя о неполноте (\DMref{п.~4.4.9}{4.4.9.}),
в которой  
речь идёт об <<истинности>> и <<доказуемости>>. 
Чтобы доказывать какие-либо утверждения относительно этих понятий, 
необходимо придать им математический смысл. 

\smallskip
Понятие \odmkey{доказательства}{Доказательство (proof)} является едва ли не самым главным в математике, 
но формально определить его даже труднее, чем понятия множества и алгоритма. 
Прекрасное определение: \em{доказательство~--- это рассуждение, 
убеждающее нас настолько, что с его помощью мы готовы убеждать других}. 
Однако это определение принадлежит скорее психологии, нежели математике. 
Ничуть не проще для математического определения и понятие \odmkey{истины}{Истина (truth)}. 
Попытки раскрыть это понятие во всей полноте уводят нас в философию, языкознание,
теологию и в другие очень интересные и поучительные области мысли, которые, однако,
далеки от изучаемого предмета: дискретной математики для программистов.


\end{frame}


\begin{frame}
\frametitle{4.5.8. Истинность и доказуемость (2/6)}
Выход состоит в том, чтобы зафиксировать неопределяемые понятия 
в виде несложно устроенных конструктивных объектов 
и далее манипулировать этими объектами средствами наивной теории множеств и 
наивной теории алгоритмов. 
То, что получится, можно назвать \odmkey{наивной теорией доказательства}{}.
 
\smallskip
В данном случае множество истинных утверждений, 
или \odmkey{множество истин}{Множество (set)!истин (of truths)}, 
рассматривается просто как некоторое множество слов $T$ 
в  универсальном алфавите $U$. 
Далее, множество доказательств также рассматривается 
просто как некоторое \em{разрешимое} множество слов $D$ 
в универсальном алфавите $U$. 
Элементы множества $D$ называются \odmkey{доказательствами}{Доказательство (proof)}. 
Других ограничений на множества $T$ и $D$ пока не накладывается. 
Наряду с разрешимым множеством $D$ постулируется наличие вычислимой функции 
$\Map{E}{D}{U^*}$, называемой 
\odmkey{функцией извлечения доказанного}{Функция (function)!извлечения доказанного (proven extraction)}. 
Если $s=E(d)$, где $d\in D$, то $s$ называется \odmkey{теоремой}{Теорема (theorem)}. 
Таким образом, $S = E(D)$~--- множество теорем. 
Пара $\Alg{D}{E}$ называется \odmkey{дедуктикой}{Дедуктика (deduction)}.


\end{frame}

\begin{frame}
\frametitle{4.5.8. Истинность и доказуемость (3/6)}
%\vspace{-4pt}
\begin{example}  
Формальные теории, рассматриваемые в предыдущих разделах, 
являются дедуктиками. Множество $D$ в них строится как последовательности слов,
получаемых из <<аксиом>> по <<правилам вывода>> 
(проверка корректности вывода~--- это и есть разрешающий алгоритм), 
а функция извлечения доказанного $E$ очень проста: 
последнее слово в доказательстве является теоремой. 
\end{example}

\smallskip
Нас интересует, как соотносятся теоремы и истины в зависимости от того, 
какими свойствами обладают множества $T$ и $D$.

\begin{examples}
Рассмотрим некоторые примеры формулировок свойств интересующих 
нас множеств на естественном языке.
\begin{enumerate}
\item {\it Существуют недоказуемые истины}. 
Это утверждение тривиально и неинтересно. 
Достаточно положить $D\Assign\emptyset$, 
и при \em{любом} множестве $T$ все истины окажутся недоказуемыми. 
\item {\it Не существует дедуктики, в которой все истины доказуемы}.
 А это утверждение просто неверно. 
 Достаточно положить $D\Assign T$, $E(t)\Assign t$, 
 и при \em{любом} множестве истин $T$ все истины окажутся легко доказуемыми.
\end{enumerate}
\end{examples}


\end{frame}


\begin{frame}
\frametitle{4.5.8. Истинность и доказуемость (4/6)}
%\vspace{-2pt}
Дедуктика $\Alg{D}{E}$ называется \odmkey{непротиворечивой}{Дедуктика (deduction)!непротиворечивая (consistent)} 
относительно множества истин $T$, если $E(D)\subset T$ 
(т.~е. $S\subset T$, всё доказуемое истинно). 
Дедуктика $\Alg{D}{E}$ называется \odmkey{полной}{Дедуктика (deduction)!полная (complete)} 
относительно множества истин $T$, если $T\subset E(D)$ 
(т.~е. $T\subset S$, всё истинное доказуемо). 

Непротиворечивая и одновременно полная дедуктика 
называется \odmkey{адекватной}{Дедуктика (deduction)!адекватная (adequate)} множеству истин $T$, 
в такой дедуктике множество теорем совпадает с множеством истин, $T = S$. 


\begin{remark}
Разрешимое множество доказательств $D$ 
является подмножеством множества всех слов, 
которое перечислимо. Значит, по лемме 2 множество $D$ перечислимо. 
Далее, множество теорем $S$ является образом перечислимого 
множества при применении вычислимой функции. 
Значит, по лемме 4 множество $S$ также перечислимо.
\end{remark}

\begin{example}
Адекватные дедуктики для интересных множеств истин существуют: 
построение таблиц истинности является адекватной дедуктикой 
для тавтологий в исчислении высказываний 
(\DMref{п.~4.3.8}{4.3.8.}).
\end{example}
\end{frame}


\begin{frame}
\frametitle{4.5.8. Истинность и доказуемость (5/6)}  
%\vspace{-2pt}
Вопрос состоит в том, при каких условиях, 
налагаемых на множество истин $T$, 
может не найтись адекватной дедуктики?

\begin{theorem}
Если множество истин $T$ перечислимо, 
то существует дедуктика $\Alg{D}{E}$, адекватная $T$.
\end{theorem} 
\begin{proof}
Если $T=\emptyset$, то достаточно положить $D\Assign \emptyset$. 
Пусть $T\ne\emptyset$ и $\Map{\Bl{T}}{\Bbb{N}}{T}$ 
перечисляющий алгоритм в форме вычислимой функции, 
вычисляющей истину по её номеру. Тогда положим $D \Assign \Bbb{N}$ 
и $E\Assign \Bl{T}$.
\end{proof}

Другими словами, если множество истин перечислимо, 
и может быть расположено в определённой последовательности, 
то доказательством истины является указание номера истины 
в этой последовательности.

\begin{corollary}
Если множество истин $T$ неперечислимо, 
то не существует дедуктики $\Alg{D}{E}$, адекватной $T$.
\end{corollary}
\begin{proof}
 Действительно, множество теорем $S$ любой дедуктики перечислимо, 
 а потому не может совпадать с неперечислимым множеством истин $T$.	 
\end{proof}


\end{frame}


\begin{frame}
\frametitle{4.5.8. Истинность и доказуемость (6/6)}

{\large
Итак, ответ на основной вопрос получен~--- 
если множество истин можно задать алгоритмом, 
то адекватная дедуктика может быть построена и всё в порядке: 
всё истинное доказуемо и всё доказуемое истинно. 
Если же множество истин невозможно задать алгоритмически,
то и нет надежды на построение адекватной дедуктики (системы доказательств). 
}

\vspace{12pt}
{\large
Но, может быть, всё не так плохо и неалгоритмические множества истин 
нам просто не понадобятся на практике? 
Увы, ответ отрицательный~--- 
неалгоритмические множества истин нам совершенно необходимы.  
}

\end{frame}


\begin{frame}
\label{4.5.9.}\subsubsection{4.5.9. Множество истин арифметики }\frametitle{4.5.9. Множество истин арифметики (1/7)}
%\vspace{-2pt}
Сначала установим важный факт, относящийся к числовым множествам. 
Говорят, что множество натуральных чисел 
$N\subset \Bbb{N}$ \odmkey{выразимо}{Выразимость множества (expressiveness of set)} с помощью множества истин $T$, 
если существует такая вычислимая функция 
$\Map{\Bl{F}}{\Bbb{N}}{U^*}$, 
что $n\in N$ тогда и только тогда, когда $\Bl{F}(n)\in T$. 

\begin{theorem}
 Множество истин $T$ выражает хотя бы одно неперечислимое множество
 натуральных чисел $N$ тогда и только тогда, 
 когда не существует дедуктики $\Alg{D}{E}$, адекватной $T$.
\end{theorem}
\begin{proof}   
\proofsection{$\Rightarrow$}
Пусть множество $N$ неперечислимо и выразимо множеством истин $T$ 
с помощью вычислимой функции $\Map{\Bl{F}}{\Bbb{N}}{U^*}$. 
Имеем $N = \Bl{F}^{-1}(\Im \Bl{F} \cap T)$. 
Но множество значений $\Im \Bl{F}$ 
вычислимой функции натурального ряда перечислимо по лемме 4. 
Значит, множество $T$ неперечислимо, 
иначе множество $\Im \Bl{F} \cap T$ было бы перечислимо и 
множество $N$ было бы перечислимо по лемме 5. 
Для неперечислимого множества истин адекватной дедуктики не существует.
\end{proof}

\end{frame}


\begin{frame}
\frametitle{4.5.9. Множество истин арифметики (2/7)}
%\vspace{-2pt}
\begin{proof} {(Продолжение)}
\proofsection{$\Leftarrow$}
Пусть для множества истин $T$ адекватной дедуктики не существует. 
Тогда множество $Т$ неперечислимо. Но множество $U^*$ перечислимо, 
пусть оно перечисляется вычислимой функцией  $\Map{\Bl{F}}{\Bbb{N}}{U^*}$. 
Имеем $T=\Bl{F}(\Bl{F}^{-1}(T))$. Положим $N \Assign \Bl{F}^{-1}(T)$. 
Заключаем по лемме 5, 
что множество $N$ неперечислимо и выразимо множеством истин $T$.
\end{proof}
 
Данный факт сводит несуществование адекватной дедуктики 
к выразимости неперечислимых числовых множеств. 
Но, может быть, ещё не всё потеряно и такие множества истин нам просто не нужны, 
и про них можно забыть? К сожалению, это не так. 
Неприятные случаи совсем рядом.

Рассмотрим в качестве примера множество истин арифметики.

\odmkey{Арифметика}{Арифметика (arithmetic)}~--- понятие натуральных чисел, 
операции сложения, умножения, отношения между натуральными числами~--- 
это такой элемент общей культуры (не только раздел математики), 
без которого существование современного человечества невозможно.


\end{frame}


\begin{frame}
\frametitle{4.5.9. Множество истин арифметики (3/7)} 
\vspace{-2pt}
Как и ранее, обсуждение вопросов о том, является ли арифметика откровением, 
данным Богом, или же это итог развития науки, является ли способность к счёту
врожденной или благоприобретённой, увело бы нас далеко от рассматриваемого предмета: дискретной математики для программистов.
 
\medskip
Зафиксируем язык арифметики, определённый в 
\DMref{п.~4.4.7}{4.4.7.}. 
В этом языке используются: цифры для записи натуральных чисел, 
буквы для записи переменных, операции сложения $+$ и умножения $\cdot$, 
отношение $=$, логические связки $\Not$, $\to$ 
и кванторы $\forall$, $\exists$. 
Далее из этих символов в языке арифметики допускается 
строить следующие конструкции.


{\color{blue}\bf 1.} \odmkey{Термы}{Терм (term)}. 
Все числа и переменные суть термы. 
Выражения вида $(t+u)$ и ($t\cdot u)$, где $t$ и $u$ суть термы, 
также суть термы. Все переменные в терме являются его 
\odmkey{параметрами}{Параметр (parameter)!терма (of term)}. 
Терм без параметров называется \odmkey{константой}{Константа (constant)}.

{\color{blue}\bf 2.} \odmkey{Атомы}{Атом (atom)}. 
Выражения вида $(t=u)$, где $t$ и $u$ суть термы, 
называются атомами, или \odmkey{элементарными формулами}{Формула (formula)!элементарная (elementary)}. Параметрами атомов являются объединение параметров термов. 


\end{frame}


\begin{frame}
\frametitle{4.5.9. Множество истин арифметики (4/7)} 

 
{\color{blue}\bf 3.} \odmkey{Формулы}{Формула (formula)}. 
Если $A$ и $B$ суть формулы, а $x$~--- переменная, 
то $\Not A$, $A\to B$, $\A{x}{A}$, $\E{x}{A}$~--- формулы. 
Кванторы исключают свои переменные из числа параметров подкванторной формулы.


\smallskip
Далее формулам без параметров (\odmkey{суждениям}{Суждение (judgement)}) 
сопоставляются истинностные значения $\TT$ и $\FF$ по правилам \DMref{п.~4.4.2}{4.4.2.}. 
Все суждения арифметики, имеющие истинностное значение $\TT$, 
образуют \odmkey{множество истин арифметики}{Множество (set)!истин (of truths)!арифметики (of arithmetic)} $T$.
Например, $2+2=4$.

\smallskip
Охарактеризуем класс \odmkey{арифметических множеств}{Класс (class)!множеств (of sets)!арифметических  (arithmetic)}. 
Рассмотрим класс формул языка арифметики 
не более чем с одним параметром. 
Будем обозначать такие формулы $\alpha(x)$. 
Тогда для любого конкретного $x\in\Bbb{N}$ формула $\alpha(x)$~--- 
суждение, и имеет истинностное значение. 
Множество $M \Assign \Set{x\in \Bbb{N}}{\alpha(x)=\TT}$
называется \odmkey{сопряжённым}{Множество (set)!сопряжённое (conjugate)} с формулой $\alpha(x)$.
Множество, сопряжённое с некоторой формулой, 
называется \odmkey{арифметическим}{Множество (set)!арифметическое (arithmetic)}. 

\end{frame}


\begin{frame}
\frametitle{4.5.9. Множество истин арифметики (5/7)} 
Арифметические множества обладают рядом полезных свойств.

\smallskip
{\bf Свойство 1.}
{\it Дополнение к арифметическому множеству есть арифметическое множество.}
\begin{proof}
Если $M$ сопряжено с $\alpha(x)$, то $\Bbb{N} \setminus M$ сопряжено 
с $\Not\alpha(x)$.
\end{proof}

\smallskip		 
{\bf Свойство 2.}
{\it Объединение и пересечение арифметических множеств 
суть арифметические множества.}
\begin{proof}
Если $M$ сопряжено с $\alpha(x)$, 
$N$ сопряжено с $\beta(x)$, 
то $M\cup N$ сопряжено с $\alpha(x)\vee\beta(x)$, 
$M\cap N$ сопряжено с $\alpha(x)\wedge\beta(x)$.
\end{proof}

\smallskip
{\bf Свойство 3.}
{\it Арифметическое множество выразимо множеством истин арифметики.}
\begin{proof}
Пусть $M$ сопряжено с $\alpha(x)$,
определим $\Map{\Bl{F}}{\Bbb{N}}{\{\FF,\TT\}}$
следующим образом: $\Bl{F}(x)\Assign \alpha(x)$. 									 \end{proof}
\end{frame}


\begin{frame}
\frametitle{4.5.9. Множество истин арифметики (6/7)} 
Арифметика, как система операций с целыми (и натуральными) числами, 
реализована в подавляющем большинстве компьютеров. 
Опыт программирования этих компьютеров приводит к убеждению, 
что любую вычислимую функцию натурального ряда можно \em{запрограммировать}, 
если отвлечься от того обстоятельства,
 что память компьютера конечна и слишком большие числа могут не поместиться. 
 Впрочем, память можно увеличить и считать потенциально бесконечной. 
 Следовательно, всякое перечислимое множество натуральных чисел 
 (как множество значений вычислимой функции) 
 может быть запрограммировано с помощью машинных операций. 
 Основываясь на этих соображениях, мы принимаем, что имеет место
  
{\bf Аксиома {\color{blue} [арифметичности]}.}
{\it Всякое перечислимое множество натуральных чисел является арифметическим.}

Принятая аксиома позволяет в несколько шагов закончить доказательство первой теоремы Гёделя о неполноте.
 \end{frame}


\begin{frame}
\frametitle{4.5.9. Множество истин арифметики (7/7)} 
{\bf Шаг 1.}
{\it Существует неперечислимое арифметическое множество.}
\begin{proof}
Рассмотрим множество $N$, построенное в доказательстве следствия 5 \DMref{п.~4.5.7}{4.5.7.}. 
Это множество перечислимо и имеет неперечислимое дополнение по следствию 6.
 Но само множество $N$ арифметическое по аксиоме арифметичности 
 и его дополнение арифметично по свойству 1.
\end{proof}

\medskip
{\bf Шаг 2.}
{\it Существует неперечислимое множество, выразимое множеством истин арифметики.}
\begin{proof}
По свойству 3.
\end{proof}

\medskip
{\bf Шаг 3.}
{\it Не существует дедуктики, адекватной множеству истин арифметики.}
\begin{proof}
По теореме.
\end{proof}

\medskip
Таким образом, для \em{любого точно сформулированного понятия 
доказательства} найдется либо доказуемое, но ложное утверждение, 
формулируемое на языке арифметики, 
либо истинное утверждение того же языка, не являющееся доказуемым.

%Список литературы
%. . .
%32. Успенский В.А. Теорема Гёделя о неполноте. – М.: Наука, 1982. – 112с.
%

\end{frame}

\begin{frame}
\frametitle{Онтология математической логики}
\begin{center}
\includegraphics[height=7.8cm]{ODM4_5.png}
\end{center}
\end{frame}


\subsection{4.6. Автоматическое доказательство теорем}
\label{4.6.}
\begin{frame}
\frametitle{4.6. Автоматическое доказательство теорем} 
Автоматическое доказательство теорем~--- это краеугольный камень
логического программирования, искусственного интеллекта и других
современных направлений в программировании. Здесь излагаются
основы \odmkey{метода резолюций}{Метод (method)!резолюций (of
resolutions)}~--- классического (и в то же время до сих пор
популярного) метода автоматического доказательства теорем,
предложенного Робинсоном
%\footnote{\;Джон А. Робинсон (р. 1930).} 
в 1965 году.

\medskip
\begin{tabular}{p{10mm}|p{48mm}|p{85mm}}
  № & Название параграфа
  & Ключевые термины и обозначения \\
  \hline
\DMref{4.6.1.}{4.6.1.} & Постановка задачи автоматического доказательства теорем   &
{ Алгоритм автоматического доказательства теорем, метод резолюций }\\  
\DMref{4.6.2.}{4.6.2.} & Доказательство от противного  &
{ Пустая формула }\\  
\DMref{4.6.3.}{4.6.3.} & Сведение к дизъюнктам  &
{ Дизъюнкт, элиминация импликации, протаскивание отрицаний, разделение связанных переменных, сколемизация, элиминация кванторов всеобщности, элиминация конъюнкции }\\  
\end{tabular}  
\end{frame}

\begin{frame}
\frametitle{4.6. Автоматическое доказательство теорем (2/2)} 

\begin{tabular}{p{10mm}|p{48mm}|p{85mm}}
  № & Название параграфа
  & Ключевые термины и обозначения \\
  \hline
\DMref{4.6.4.}{4.6.4.} & Правило резолюции для исчисления высказываний &
{ Родительские дизъюнкты, правило резолюции, резольвента, контрарные литералы }\\  
\DMref{4.6.5.}{4.6.5.} & Правило резолюций для исчисления предикатов &
{  }\\  
\DMref{4.6.6.}{4.6.6.} & Опровержение методом резолюций &
{ Полуразрешимая теория }\\  
\DMref{4.6.7.}{4.6.7.} & Дерево вывода &
{  }\\  
\DMref{4.6.8.}{4.6.8.} & Алгоритм метода резолюций &
{  }\\  
\DMref{4.6.9.}{4.6.9.} & Стратегии метода резолюций &
{ Полные и неполные стратегии, уровень опровержения, стратегия насыщения уровня, линейная стратегия, стратегия единичной резолюции, хорновский дизъюнкт, факт, правило, цель, отрицательная стратегия }\\  
\end{tabular}  
\end{frame}

\begin{frame}
\label{4.6.1.}\subsubsection{4.6.1. Постановка задачи   автоматического доказательства теорем}
\frametitle{4.6.1. Постановка задачи   автоматического доказательства теорем (1/3)} 
Алгоритм, который проверяет отношение
$$ \Seq{\Gamma}{S}{\cal T}$$
для формулы $S$, множества формул $\Gamma$ и теории $\cal T$,
называется алгоритмом \odmkey{автоматического доказательства
теорем}{Автоматическое доказательство теорем (automatic theorem
proving)}. В общем случае такой алгоритм невозможен, то~есть не
существует алгоритма, который для любых $S$, $\Gamma$ и $\cal T$
выдавал бы ответ <<Да>>, если $\Seq{\Gamma}{S}{\cal T}$,  и ответ
<<Нет>>, если неверно, что $\Seq{\Gamma}{S}{\cal T}$. Более того,
известно, что нельзя построить алгоритм автоматического
доказательства теорем даже для большинства конкретных достаточно
богатых формальных теорий $\cal T$. В некоторых случаях удаётся
построить алгоритм автоматического доказательства теорем, который
применим не ко всем формулам теории (то~есть частичный алгоритм).
%см. п.~\ref{Solvable}).



\end{frame}

\begin{frame}
\frametitle{4.6.1. Постановка задачи  
автоматического доказательства теорем (2/3)}

Для некоторых простых формальных теорий
(например, для исчисления высказываний) и
некоторых простых классов формул
(например, для прикладного исчисления предикатов с
одним одноместным предикатом) алгоритмы автоматического доказательства
теорем известны.


\begin{example}
Поскольку для исчисления высказываний известно, что теоремами
являются общезначимые формулы, 
то достаточно вычислить истинностное значение
формулы при всех возможных интерпретациях (их конечное число).
Если во всех случаях получится значение $\TT$, то проверяемая
формула~--- тавтология и, следовательно, является теоремой теории
$\cal L$. Если же хотя бы в одном случае получится значение $\FF$, 
то проверяемая формула не является тавтологией и, следовательно, не
является теоремой теории~$\cal L$.
\end{example}

%\vspace*{-\baselineskip}
\end{frame}

\begin{frame}
\frametitle{4.6.1. Постановка задачи  
автоматического доказательства теорем (3/3)}
\begin{remark}
Приведённый выше пример является алгоритмом автоматического
доказательства теорем в теории $\cal L$, хотя и не является
алгоритмом автоматического \em{поиска вывода} 
теорем из аксиом теории $\cal L$.
\end{remark}

Наиболее известный классический алгоритм автоматического
доказательства теорем называется \odmkey{методом резолюций}{Метод
(method)!резолюций (of resolutions)}. Для любого прикладного
исчисления предикатов первого порядка $\cal T$, любой формулы $S$
и множества формул $\Gamma$ теории $\cal T$ метод резолюций выдаёт
ответ <<Да>>, если $\Seq{\Gamma}{S}{\cal T}$,  и выдаёт ответ
<<Нет>> или не выдаёт никакого ответа (то~есть <<зацикливается>>),
если неверно, что $\Seq{\Gamma}{S}{\cal T}$.

\end{frame}

\begin{frame}
\label{4.6.2.}\subsubsection{4.6.2. Доказательство от противного}
\frametitle{4.6.2. Доказательство от противного} 

%\vspace{-2pt}
%\subsection{Доказательство от противного}
В основе метода резолюций лежит идея <<доказательства от противного>>.
\begin{theorem} Если $\Seq{\Gamma,\neg S}{F}{}$, где $F$~{\upshape---} любое
противоречие (тождественно ложная формула), то
$\Seq{\Gamma}{S}{}$.
\end{theorem}

%{%
%\def\proof{\pagebreak[2]\par\addvspace{0pt plus 3pt minus 3pt}{\sffamily\bfseries\small Д{\scriptsize\uppercase{оказательство}}} (для случая $\cal L$).\quad\nopagebreak}
\begin{proof}
(для случая $\cal L$)
$\Gamma$, $\Seq{\neg S}{F}{}$ $\Equiv$
$\Gamma\And\neg S\to F$~--- тавтология. \\Но
$
\Gamma\And\neg S\to F=$ $\neg
(\Gamma\And\neg S)\vee F=$ $\neg(\Gamma\And\neg S)=$ $\neg\Gamma\vee
S=\Gamma\to S.
$
\\Имеем $\Gamma\to S$~--- тавтология $\Equiv
\Seq{\Gamma}{S}{}$.
\end{proof}



\odmkey{Пустая}{Формула (formula)!пустая (empty)}
формула не имеет никакого значения ни в
какой интерпретации, в~частности, не является истинной ни в
какой\vadjust{\penalty-1000} интерпретации и по определению является противоречием. В
качестве формулы $F$ при доказательстве от противного по методу
резолюций принято использовать пустую формулу.
Полезно сравнить использование
пустой формулы в методе резолюций с~\DMref{п.~3.4.3}{3.4.3.}.


\begin{remark}
При изложении метода резолюций пустая формула традиционно
обозначается~$\square$. Однако значок $\square$ используется 
в слайдах также для
обозначения конца доказательства теоремы или обоснования алгоритма.
Необходимо их различать по контексту.
% , 
%в каком смысле использован значок $\square$ в каждом случае.
\end{remark}


%\subsection{Свед$\acute{\text{е}}$ние к предложениям}

\end{frame}

\begin{frame}
\label{4.6.3.}\subsubsection{4.6.3. Свед$\acute{\text{е}}$ние к дизъюнктам }\frametitle{4.6.3. Свед$\acute{\text{е}}$ние к дизъюнктам (1/4) } 


Метод резолюций работает с особой стандартной формой формул,
которые называются \odmkey{дизъюнктами}{Дизъюнкт (clause)},
или \odmkey{предложениями}{Предложение (clause)}.
Дизъюнкт~--- это бескванторная дизъюнкция литералов. Любая
формула исчисления предикатов может быть преобразована в~множество
дизъюнктов следующим образом (здесь знак $\mapsto$ используется
для обозначения способа преобразования формул):


\smallskip
{\color{blue}\bf 1.} \odmkey{Элиминация импликации}{Элиминация
(elimination)!импликации (of implication)}. 
Преобразование:\\
$A\to B\mapsto \Not A \Or B$.\\
После первого этапа формула содержит только
$\Not,\Or,\And,\forall,\exists$.


\smallskip
{\color{blue}\bf 2.} \odmkey{Протаскивание отрицаний}{Протаскивание
отрицаний (dragging of negations)}.
Преобразования:\\
$\Not \A{x}{A}\mapsto\E{x}{\Not A},\quad
 \Not \E{x}{A}\mapsto\A{x}{\Not A},\quad
 \Not \Not A\mapsto A$,\\
$\Not (A\Or B)\mapsto \Not A \And \Not B,\quad
 \Not (A \And B)\mapsto\Not A \Or \Not B$.\\
После второго этапа формула содержит отрицания только
перед атомами.

\end{frame}

\begin{frame}
\frametitle{4.6.3. Свед$\acute{\text{е}}$ние к дизъюнктам (2/4) } 

\vspace{-2pt}
{\color{blue}\bf 3.} \odmkey{Разделение связанных переменных}{Разделение связанных переменных (splitting of bounded variables)}.
Преобразование:\\
$Q_1\,x\ A(\dots Q_2\,x\ B(\dots x\dots)\dots)\mapsto
 Q_1\,x\ A(\dots Q_2\,y\ B(\dots y\dots)\dots)$,\\
где $Q_1$ и $Q_2$~--- любые кванторы.
После третьего этапа формула не содержит
случайно совпадающих связанных переменных.

{\color{blue}\bf 4.} \em{Приведение к предварённой форме\/}.
Преобразования:\\
$Q\,x\ A \Or B \mapsto
 Q\,x\ (A \Or B), \quad
Q\,x\ A \And B \mapsto
 Q\,x\ (A \And B)$,\\
где $Q$~--- любой квантор. После четвёртого этапа формула
находится в~предварённой форме.

{\color{blue}\bf 5.} \odmkey{Элиминация кванторов существования}{Элиминация
(elimination)!квантора существования (of existential quantifier)}
(\odmkey{сколемизация}{Сколемизация
(sk$\ddot{\text{o}}$lemization)}).
Преобразования:\\

{\small
$\E{x_1}{Q_2\,x_2\ \dots Q_n\,x_n\ A(x_1,x_2,\dots,x_n)}\mapsto
Q_2\,x_2 \dots Q_n\,x_n\ A(a,x_2,\dots,x_n)$,\\
$\A{x_1}{\dots\A{x_{i-1}}
                {\E{x_i}{Q_{i+1}\,x_{i+1}\dots Q_n\,x_n\ 
                A(x_1,\dots,x_i,\dots,x_n)}}\dots} \mapsto$\\
$\A{x_1}{\dots\A{x_{i-1}}{Q_{i+1}\,x_{i+1} \dots Q_n\,x_n
A(x_1,\dots,f(x_1,\dots,x_{i-1}),\dots,x_n)}\dots}$,\\
}

где $a$~--- новая предметная константа,
$f$~--- новый функциональный символ, а
$Q_1,Q_2,\dots,Q_n$~--- любые кванторы.
После пятого этапа формула содержит только кванторы всеобщности.


\end{frame}

\begin{frame}
\frametitle{4.6.3. Свед$\acute{\text{е}}$ние к дизъюнктам (3/4) } 


{\color{blue}\bf 6.} \odmkey{Элиминация кванторов всеобщности}{Элиминация
(elimination)!квантора всеобщности (of universal quantifier)}.
Преобразование:\\
$\A{x}{A(x)}\mapsto A(x)$.\\
После шестого этапа формула не содержит кванторов.

{\color{blue}\bf 7.} \em{Приведение к конъюнктивной нормальной форме\/}. 
Преобразования:\\
$A\Or(B\And C)\mapsto (A\Or B) \And (A\Or C)$, 
$(A\And B)\Or C\mapsto (A\Or C) \And (B\Or C)$.\\
После седьмого этапа формула находится в конъюнктивной нормальной
форме.


{\color{blue}\bf 8.} \odmkey{Элиминация конъюнкции}{Элиминация
(elimination)!конъюнкции (of conjunction)}. 
Преобразование:\\
$A\And B \mapsto A,\, B$.\\
После восьмого этапа формула распадается на множество дизъюнктов.

Не все преобразования на этапах 1--8 являются логически эквивалентными.

\end{frame}

\begin{frame}
\frametitle{4.6.3. Свед$\acute{\text{е}}$ние к дизъюнктам (4/4) } 
%

\begin{theorem}
Если $\Gamma$~{\upshape---} множество дизъюнктов, полученных из
формулы $S$, то $S$ является противоречием тогда и только тогда,
когда множество $\Gamma$ невыполнимо.\parfillskip0pt
\end{theorem}

\begin{proof}
В доказательстве нуждается шаг 5~--- сколемизация.\\
 Пусть
$F\Assign\A{{x_1}}{\dots\A{x_{i-1}}{\E{x_i(Q_{i+1}x_{i+1}\dots
Q_nx_n} {A(\Tuple{x}{1}{n}))}\dots}\dots}$. Положим
$F' \Assign 
\forall{x_1} ( \dots
\forall x_{i-1} ( Q_{i+1}x_{i+1}\dots
Q_nx_n ( $ \\ $\quad\quad A(\Tuple{x}{1}{i-1}, f(\Tuple{x}{1}{i-1}),
\Tuple{x}{i+1}{n} )) \dots ) \dots )\dots )$.\\ Пусть $F$~--- противоречие, а
$F'$~--- не противоречие. Тогда существуют интерпретация $I$ и
набор значений $s=(\Tuple{a}{1}{i-1},\Tuple{a}{i+1}{n})$, такие,
что $s^*(F')=\text{T}$.
 Положим
$a_i\Assign f(\Tuple{a}{1}{i-1})$,
 $s_1\Assign (\Tuple{a}{1}{i-1}, a_i,
\Tuple{a}{i+1}{n})$. Тогда $s_1^*(F)=\text{T}$, и $F$~--- выполнимая формула.
\end{proof}
%
%\addvspace{16pt}

\begin{remark}
Невыполнимость множества формул $\Gamma$ означает,
что множество $\Gamma$ не имеет модели, то~есть не существует
интерпретации, в которой все формулы $\Gamma$ имели бы значение T.
\end{remark}

\end{frame}

\begin{frame}
\label{4.6.4.}\subsubsection{4.6.4. Правило резолюции для исчисления высказываний }\frametitle{4.6.4. Правило резолюции для исчисления высказываний (1/2) } 
%\addvspace{16pt}
%\subsection{Правило резолюции для исчисления высказываний}

Метод резолюций основан на использовании
специального правила вывода.

Пусть $C_1$ и $C_2$~--- два дизъюнкта в исчислении высказываний
и пусть $C_1=P\Or C_1'$, а $C_2=\Not P\Or C_2'$, где
$P$~--- пропозициональная переменная, а $C_1'$, $C_2'$~--- любые
дизъюнкты (в частности, может быть, пустые или состоящие только
из одного литерала). Правило вывода
$$
\Rule{C_1,\ C_2}{C_1'\Or C_2'}{(R)}
$$
называется \odmkey{правилом резолюции}{Правило
(rule)!резолюции (resolution)}, дизъюнкты
$C_1$ и  $C_2$ называются
\odmkey{резольвируемыми}{Дизъюнкт (clause)!резольвируемый (resolvable)}
 (или
\odmkey{родительскими}{Дизъюнкт (clause)!родительский
(parent)}),
  дизъюнкт
$C_1'\Or C_2'$~--- \odmkey{резольвентой}{Резольвента (resolvent)}, а формулы
$P$ и $\Not P$~--- \odmkey{контрарными}{Литерал (literal)!контрарный
 (contrary)}
литералами.

\begin{remark}
Резольвента является дизъюнктом, то~есть
множество дизъюнктов замкнуто
относительно правила резолюции.
\end{remark}

%\setcounter{theorem}{0}
\begin{theorem}{\bf\color{blue}[1]}
Правило резолюции логично, то~есть резольвента является логическим
следствием резольвируемых дизъюнктов.
\end{theorem}

\end{frame}

\begin{frame}
\frametitle{4.6.4. Правило резолюции для исчисления высказываний (2/2) }

\begin{proof}
Пусть $I(C_1)=\text{T}$ и $I(C_2)=\text{T}$. Тогда если $I(P)=\text{T}$, то
$C'_2\not=\emptyset$ и $I(C'_2)=\text{T}$,
а значит, $I(C'_1\vee C'_2)=\text{T}$. Если
же $I(P)=\text{F}$, то $C'_1\not= \emptyset$ и $I(C'_1)=\text{T}$, а значит,
$I(C'_1\vee C'_2)=\text{T}$.
\end{proof}


Правило резолюции~--- это очень мощное правило вывода.
Многие ранее рассмотренные правила являются его частными случаями:



\begin{align*}
\Rule{A, A\to B}{B}&{\text{Modus ponens}} & \Rule{A, \Not A \Or B}{B} & {R},\\
\Rule{A\to B, B\to C}{A\to C}&{\text{Транзитивность}} & \Rule{\Not A \Or B, \Not B\Or C}{\Not A \Or C}&{R}, \\
\Rule{A\Or B, A\to B}{B}&{\text{Слияние}} & \Rule{A\Or B, \Not
A\Or B}{B}&{R}.
\end{align*}


\end{frame}

\begin{frame}
\label{4.6.5.}\subsubsection{4.6.5. Правило резолюции для исчисления предикатов }\frametitle{4.6.5. Правило резолюции для исчисления предикатов (1/3) }
%\subsection{Правило резолюции для исчисления предикатов}
\begin{theorem} {\bf\color{blue}[2]}
Правило резолюции полно, то~есть всякая тавтология может быть выведена
с использованием только правила резолюции.
\end{theorem}

\begin{proof}
Modus ponens является частным случаем правила резолюции.
\end{proof}

Для применения
правила резолюции нужны контрарные литералы в резольвируемых
дизъюнктах.
В случае исчисления высказываний
контрарные литералы определяются \\ очень
просто~---
это пропозициональная переменная и её отрицание.
Для исчисления предикатов определение чуть
сложнее.

Пусть $C_1$ и $C_2$~--- два дизъюнкта в
исчислении предикатов. Правило вывода
$$
\Rule{C_1,\ C_2}{(C_1'\Or C_2')\sigma}{(R)}
$$
называется правилом резолюции в исчислении предикатов, если в
дизъюнктах $C_1$ и $C_2$ существуют унифицируемые контрарные
литералы $P_1$ и $P_2$, то~есть $C_1 =P_1\Or C_1'$ и $C_2=\Not
P_2 \Or C_2'$, причём атомарные формулы $P_1$ и $P_2$
унифицируются наиболее общим унификатором $\sigma$. В этом случае
резольвентой дизъюнктов $C_1$ и $C_2$ является дизъюнкт
$(C_1'\Or C_2')\sigma$, полученный из дизъюнкта $C_1'\Or C_2'$
применением унификатора $\sigma$.

\end{frame}

\begin{frame}
\frametitle{4.6.5. Правило резолюции для исчисления предикатов (2/3) }

\begin{example}
Пусть заданы дизъюнкты 
$P(x,a) \Or Q(x,y)$ и $\Not P(b,y) \Or R(y,x)$, где 
$P, Q, R$~--- двухместные предикаты, 
а $a$ и $b$~--- константы.
Тогда применение правила резолюции имеет вид
$$
\Rule{P(x,a) \Or Q(x,y),\qquad \Not P(b,y) \Or R(y,x)}
{Q(b,a)\Or R(a,b)}{(R)},
$$
при этом наиболее общий унификатор $\sigma = [b/x, a/y]$.
\end{example}

Если нет контрарных литералов, то правило резолюции неприменимо.

\begin{example}
Пусть заданы дизъюнкты 
$P(x,a) \Or Q(x,y)$ и $ P(b,y) \Or \Not R(y,x)$, где 
$P, Q, R$~--- двухместные предикаты, 
а $a$ и $b$~--- константы.
Тогда  правило резолюции  не применимо,
поскольку литералы  $P(x,a)$ и $P(b,y)$ 
хотя и унифицируемы, но не контрарны, и других 
контрарных литералов нет.
\end{example}

\end{frame}

\begin{frame}
\frametitle{4.6.5. Правило резолюции для исчисления предикатов (3/3) }

Следует ещё раз подчеркнуть, что наличия
атома и его отрицания в разных дизъюнктах ещё недостаточно:
необходимо, чтобы контрарные литералы были унифицируемыми.

\begin{example}
Пусть заданы дизъюнкты 
$P(x,a) \Or Q(x,y)$ и $\Not P(x,b) \Or R(y,x)$, где 
$P, Q, R$~--- двухместные предикаты, 
а $a$ и $b$~--- константы.
Тогда  правило резолюции  не применимо,
поскольку литералы  $P(x,a)$ и $P(x,b)$ 
не унифицируемы.
\end{example}

\end{frame}

\begin{frame}
\label{4.6.6.}\subsubsection{4.6.6. Опровержение методом резолюций }\frametitle{4.6.6. Опровержение методом резолюций (1/4) }

\odmkey{Опровержение методом резолюций}{Опровержение методом резолюций (refutation by resolution method)}~--- 
это алгоритм автоматического доказательства теорем в
прикладном исчислении предикатов, который сводится к~следующему.

\smallskip
Пусть нужно установить выводимость $\Seq{S}{G}{}.$
Каждая формула множества $S$ и формула $\Not G$ (отрицание целевой
теоремы) \em{независимо} преобразуются в множества дизъюнктов. В
полученном совокупном множестве дизъюнктов $C$ отыскиваются
резольвируемые дизъюнкты, к~ним применяется правило резолюции, и
резольвента добавляется в множество дизъюнктов до тех пор, пока
не будет получен пустой дизъюнкт. При этом возможны три
случая.



\end{frame}

\begin{frame}
\frametitle{4.6.6. Опровержение методом резолюций (2/4) }
\begin{enumerate}
\item 
Среди текущего множества дизъюнктов нет резольвируемых,
или же все  %\linebreak
резольвенты уже присутствуют в текущем множестве дизъюнктов. Это
означает, что теорема опровергнута, то~есть формула $G$ не выводима из
множества формул~$S$.


\item  В результате очередного применения правила резолюции получен пустой
дизъюнкт. Это означает, что теорема доказана, то~есть $\Seq{S}{G}{}$.
\item  Процесс не заканчивается, то~есть множество дизъюнктов пополняется всё
новыми резольвентами, среди которых нет пустых. Это ничего не означает.
\end{enumerate}
\begin{remark}
Таким образом, исчисление предикатов является
\odmkey{полуразрешимой}{Теория (theory)!полуразрешимая
(semisolvable)} теорией, а метод резолюций является 
\em{частичным} алгоритмом автоматического доказательства теорем.
\end{remark}


\end{frame}


\begin{frame}
\frametitle{4.6.6. Опровержение методом резолюций (3/4) }


\begin{example}
Рассмотрим выводимость (см. \DMref{п.~4.4.5}{4.4.5.}):
$\Seq{\E{x}{\A{y}{A(x,y)}}}{\A{y}{\E{x}{A(x,y)}}}{\cal{K}}$.
Сначала необходимо исходную формулу $\E{x}{\A{y}{A(x,y)}}$ и
отрицание целевой формулы 
$\Not\left(\A{y}{\E{x}{A(x,y)}}\right) $
преобразовать в дизъюнкты.

\begin{enumerate}
\item [1.] Формулы без изменений.
\item [2.] $\E{x}{\A{y}{A(x,y)}}$, $\E{y}{\A{x}{\Not A(x,y)}}$.
\item [3--4.] Формулы без изменений.
\item [5.] $\A{y}{A(a,y)}$, $\A{x}{\Not A(x,b)}$.
\item [6.] $A(a,y)$, $\Not A(x,b)$.
\item [7--8] Формулы без изменений.
\end{enumerate}

После этого проводится резольвирование имеющихся
дизъюнктов 1--2.
\begin{enumerate}
\item $A(a,y)$.
\item $\Not A(x,b)$.
\item $\square\qquad (1,2),\quad\sigma=[a/x,b/y]$.
\end{enumerate}

Таким образом, теорема доказана.

\end{example}
\end{frame}


\begin{frame}
\frametitle{4.6.6. Опровержение методом резолюций (4/4) }

%\vspace{-2pt}
Рассмотрим теперь обратную теорему 
$\Seq{\A{y}{\E{x}{A(x,y)}}}{\E{x}{\A{y}{A(x,y)}}}{\cal{K}}$.
Сначала необходимо исходную формулу $\A{y}{\E{x}{A(x,y)}}$ и
отрицание целевой формулы \\
$\Not\left(\E{x}{\A{y}{A(x,y)}}\right)$ преобразовать в дизъюнкты.

\vspace{-2pt}
\begin{enumerate}
\item [1.] Формулы без изменений.
\item [2.] $\A{y}{\E{x}{A(x,y)}}$, $\A{x}{\E{y}{\Not A(x,y)}}$.
\item [3--4.] Формулы без изменений.
\item [5.] $\A{y}{A(f(y),y)}$, $\A{x}{\Not A(x,g(x))}$.
\item [6.] $A(f(y),y)$, $\Not A(x,g(x))$.
\item [7--8] Формулы без изменений.
\end{enumerate}


\vspace{-2pt}
Попытка резольвирования имеющихся
дизъюнктов 1--2:
\begin{enumerate}
\item $A(f(y),y)$.
\item $\Not A(x,g(x))$.
\end{enumerate}

%\vspace{-0.2\baselineskip}
\vspace{-2pt}
Резольвенты нет, так как 
литералы $A(f(y),y)$, $A(x,g(x))$ не унифицируемы.
Действитель\-но, подстановка $[f(y)/x]$ 
даст $A(f(y),y)$, $A(f(y),g(f(y)))$,
а подстановка $[g(f(y))/y]$
не может входить в унификатор (см. \DMref{п.~4.3.2}{4.3.2.}).
Теорема опровергнута.
\end{frame}


\begin{frame}
\label{4.6.7.}\subsubsection{4.6.7. Дерево вывода }\frametitle{4.6.7. Дерево вывода (1/4) }

При применении метода резолюций (и не только),
часто используют графическое представление
вывода, которое называется \odmkey{деревом вывода}{Дерево (tree)!вывода (derivation)}. 

В узлах дерева помещаются дизъюнкты.
Дерево вывода растет снизу вверх. 
В верхних узлах (листьях) расположены исходные дизъюнкты, 
а в промежуточных узлах~--- резольвенты. 
В корне дерева~--- пустое предложение. 
На дугах можно показать резольвируемые литералы и используемые подстановки.

\begin{example}
Докажем методом резолюций теорему
$$\Seq{A\to B, B\to C}{A\to C}{\cal L}.$$

Сначала нужно преобразовать в дизъюнкты исходные формулы 
$A\to B$, $B\to C$ и
отрицание целевой формулы $\Not(A\to C)$.

\begin{enumerate}
\item [1.] $\Not A \Or B$, $\Not B \Or C$, $\Not(\Not A \Or C)$.
\item [2.] $\Not A \Or B$, $\Not B \Or C$, $ A \And \Not C$.
\item [3--7.] Формулы без изменений.
\item [8.] $\Not A \Or B$, $\Not B \Or C$, $ A$, $\Not C$.
\end{enumerate}
\end{example}
\end{frame}


\begin{frame}
\frametitle{4.6.7. Дерево вывода (2/4) }
 После этого проводится резольвирование имеющихся
 дизъюнктов 1--4.

\begin{columns}[t]
\begin{column}[T]{3.5cm}
\begin{tabbing}
{\color{blue} 1.} \= $\Not A \Or C\qquad$ \= $(1,2)$  \kill
{\color{blue} 1.} \> $\Not A\Or B$\> \\
{\color{blue} 2.} \> $\Not B\Or C$\> \\
{\color{blue} 3.} \> $A$ \> \\
{\color{blue} 4.} \> $\Not C$ \> \\
{\color{blue} 5.} \> $\Not A \Or C$ \> $(1,2)$ \\  
{\color{blue} 6.} \> $C$ \> $(3,5)$ \\  
{\color{blue} 7.} \> $\square$ \> $(4,6)$ \\  
\end{tabbing}
\end{column}
\begin{column}[T]{10cm}
\begin{center}
\includegraphics[height=45mm]{4_2.jpg}
\end{center}
\end{column}
\end{columns}

\medskip
Таким образом, теорема доказана.
Дерево вывода приведено на  рисунке справа. 


\end{frame}


\begin{frame}
\frametitle{4.6.7. Дерево вывода (3/4) }

Рассмотрим ещё один пример.
Докажем методом резолюций теорему\\
$\Seq{}{(((A\to B)\to A)\to A)}{\cal L}$.
Сначала нужно преобразовать в дизъюнкты
отрицание целевой формулы $\Not(((A\to B)\to A)\to A)$.


\begin{enumerate}
\item [1.] $\Not(\Not(\Not(\Not A\Or B)\Or A)\Or A)$.
\item [2.] $(((A\And\Not B)\Or A)\And \Not A)$.
\item [3--6.] Формула без изменений.
\item [7.] $(A\Or A)\And(\Not B\Or A)\And\Not A$.
\item [8.] $A\Or A, \Not B\Or A, \Not A$.
\end{enumerate}

\end{frame}


\begin{frame}
\frametitle{4.6.7. Дерево вывода (4/4) }

Затем проводится резольвирование имеющихся
дизъюнктов 1--3.

\vspace{-18pt}
\begin{columns}[t]
\begin{column}[T]{3.5cm}
\begin{tabbing}
{\color{blue} 1.} \= $\Not A \Or C\qquad$ \= $(1,2)$  \kill
{\color{blue} 1.} \> $ A\Or A$\> \\
{\color{blue} 2.} \> $\Not B\Or A$\> \\
{\color{blue} 3.} \> $\Not A$ \> \\
{\color{blue} 4.} \> $A$ \> $(1,3)$ \\  
{\color{blue} 5.} \> $\square$ \> $(3,4)$ \\  
\end{tabbing}
\end{column}
\begin{column}[T]{10cm}
\begin{center}
\includegraphics[height=40mm]{4_3.jpg}
\end{center}
\end{column}
\end{columns}

\medskip
\begin{remark}
Из этого примера видно, 
что устоявшийся термин <<дерево вывода>> не совсем корректен. 
Если дизъюнкт участвует в резольвировании несколько раз, 
то полученный граф может не быть деревом. 
Если же некоторые исходные дизъюнкты не участвуют в выводе, 
то граф может быть даже несвязным. 
\end{remark} 


\end{frame}


\begin{frame}
\label{4.6.8.}\subsubsection{4.6.8. Алгоритм метода резолюций }\frametitle{4.6.8. Алгоритм метода резолюций (1/3)}

Алгоритм\odmkey{}{Алгоритм (algorithm)!метода резолюций (of a resolution technique)}
поиска опровержения методом резолюций проверяет выводимость
формулы $G$ из множества формул $S$.


\begin{algorithmic}
\REQUIRE множество дизъюнктов $C$, полученных из множества формул
$S$ и~формулы $\Not G$. 
\ENSURE $1$~--- если $G$
выводимо из $S$, $0$~--- в противном случае. \STATE
$M\Assign C$ 
\COMMENT{\color{gray}\small  $M$~--- текущее множество дизъюнктов }
\WHILE{$\square \notin M$} 
\STATE %\boxed{\square} -> \square
$\text{Choose}(M, c_1, c_2, p_1, p_2, \sigma)$ 
\COMMENT{\color{gray}\small выбор
родительских дизъюнктов } \IF {$c_1, c_2 = \emptyset$}
  \STATE \textbf{return $0$}
  \COMMENT{\color{gray}\small нечего резольвировать~--- теорема опровергнута }
\ENDIF
\STATE $c\Assign R(c_1, c_2, p_1, p_2, \sigma)$
\COMMENT{\color{gray}\small вычисление резольвенты }
\STATE  $M\Assign M +c$
\COMMENT{\color{gray}\small пополнение текущего множества }
\ENDWHILE
\STATE \textbf{return $1$} \COMMENT{\color{gray}\small теорема доказана }
\end{algorithmic}

\end{frame}


\begin{frame}
\frametitle{4.6.8. Алгоритм метода резолюций (2/3)}

%\vspace{-0.25\baselineskip}
\begin{ground}
Если алгоритм заканчивает свою работу, то
правильность результата очевидным образом следует из
теорем предшествующих разделов.
Вообще говоря, завершаемость этого алгоритма
не гарантирована.
\end{ground}
%


\medskip
\begin{digress}
При автоматическом доказательстве теорем методом резолюций
львиная доля вычислений приходится на поиск унифицируемых
дизъюнктов. Таким образом, эффективная реализация алгоритма
унификации критически важна для практической применимости метода резолюций.
\end{digress}
\end{frame}


\begin{frame}
\frametitle{4.6.8. Алгоритм метода резолюций (3/3)}

\vspace{-2pt}
В алгоритме поиска опровержения 
методом резолюций использованы две вспомогательные функции.
Функция $R$ вычисляет резольвенту двух дизъюнктов, $c_1$ и $c_2$,
содержащих унифицируемые контрарные литералы $p_1$ и $p_2$ соответственно;
при этом $\sigma$~---
наиболее общий унификатор.
Результатом работы функции является
резольвента.
Процедура Choose выбирает в текущем множестве дизъюнктов $M$ два
резольвируемых дизъюнкта, то~есть два дизъюнкта, которые
содержат унифицируемые контрарные литералы. Если таковые есть, то
процедура их возвращает, в противном случае возвращается пустое
множество. 
В множестве дизъюнктов может существовать \em{много} пар дизъюнктов, к которым можно применить правило резолюций. Поэтому при автоматическом доказательстве теорем методом резолюций значительная часть вычислений приходится на поиск родительских дизъюнктов. Неограниченное применение правила резолюций может вызвать порождение \em{большого} количества дизъюнктов. 

\end{frame}

\begin{frame}
\label{4.6.9.}\subsubsection{4.6.9. Стратегии метода резолюций }\frametitle{4.6.9. Стратегии метода резолюций (1/6)}

\vspace{-2pt}
Порядок, в котором процедура
Choose выдаёт резольвируемые пары, имеет большое значение для применимости метода резолюций.
Если порядок удачный, то пустой дизъюнкт может быть получен быстро, а если порядок неудачный,
то резольвирование может продолжаться безнадёжно долго.  
Конкретные реализации процедуры Choose называются
\odmkey{стратегиями}{Стратегия (strategy)!метода резолюций (of
resolution method)} метода резолюций.

\begin{remark}
Следует иметь в виду, что автоматическое доказательство теорем
методом резолюций имеет, по существу, переборный характер,
и этот перебор столь велик, что может быть практически
неосуществим за приемлемое время на любом 
современном компьютере.
\end{remark}

В настоящее время предложено множество различных стратегий метода
резолюций. 

\smallskip
Среди них различаются \odmkey{полные} {Стратегия (strategy)!метода резолюций (of
resolution method)!полная
(complete)} и \odmkey{неполные}{Стратегия   (strategy)!метода резолюций (of
resolution method)!неполная (incomplete)}
стратегии. Полные стратегии~--- это такие, которые
гарантируют нахождение доказательства теоремы, если оно
вообще существует. Неполные стратегии могут в некоторых
случаях не находить доказательства, зато они работают быстрее.

\end{frame}

\begin{frame}
\frametitle{4.6.9. Стратегии метода резолюций (2/6)}

\vspace{-2pt}
Пусть $K_0$~--- исходное множество дизъюнктов. 
Тогда множество $K_i$ состоит из объединения множества $K_{i-1}$ и множества резольвент, порождённых из $K_{i-1}$. Индекс $i$ называется \odmkey{уровнем опровержения}{Уровень опровержения (rebuttal level)}. 
\odmkey{Стратегия насыщения уровня }{Стратегия (strategy)!насыщения уровня (level saturation)}
предполагает последовательное порождение всех резольвент уровня 1, затем уровня 2 и так далее до получения пустого дизъюнкта.

\begin{theorem}
Стратегия насыщения уровня является полной стратегией.
\end{theorem}
\begin{proof}
На каждом шаге строятся все возможные резольвенты, а значит, если какое-то опровержение и существует, то оно будет найдено ввиду полноты перебора.
\end{proof}

\begin{example}
Рассмотрим пример из \DMref{п.~4.6.7,}{4.6.7.}   уровень $K_0$
состоит из четырёх дизъюнктов:
\end{example}

\vspace{-4pt}
\begin{tabbing}
{\color{blue} 1.} \= $\Not A \Or C\qquad$ \= $(1,2)$  \kill
{\color{blue} 1.} \> $\Not A\Or B$\> \\
{\color{blue} 2.} \> $\Not B\Or C$\> \\
{\color{blue} 3.} \> $ A$ \> \\
{\color{blue} 4.} \> $\Not C$ \>  \\  
\end{tabbing}

%Тогда K_0 согласно стратегии:
%	¬A ⋁ B
%	¬B ⋁ C
%	A
%	¬C
\end{frame}

\begin{frame}
\frametitle{4.6.9. Стратегии метода резолюций (3/6)}

Уровень $K_1$ составляют дизъюнкты из $K_0$, а также:

%	¬A ⋁ С	(1, 2)
%	B		(1, 3)
%	¬B		(2, 4)
\begin{tabbing}
{\color{blue} 19.\ } \= $\Not A \Or C\qquad$ \= $(1,2)$  \kill
{\color{blue} 5.} \> $\Not A\Or C$\> $(1,2)$\\
{\color{blue} 6.} \> $B$\> $(1,3)$ \\
{\color{blue} 7.} \> $\Not B$ \> $(2,4)$\\
\end{tabbing}

\vspace{-6pt}
Уровень $K_2$ составляют дизъюнкты из $K_1$, а также:

%Уровень второго порядка K_2 содержит помимо K_1:
%	¬A		(1, 7)
%	 C		(2, 6)
%	 C		(3, 5)
%	 ¬A		(4, 5)

\begin{tabbing}
{\color{blue} 18.\ } \= $\Not A \Or C\qquad$ \= $(1,2)$  \kill
{\color{blue} 8.} \> $\Not A$\> $(1,7)$\\
{\color{blue} 9.} \> $C$\> $(2,6)$ \\
{\color{blue} 10.} \> $C$ \> $(3,5)$\\
{\color{blue} 11.} \> $\Not A$ \> $(4,5)$\\
\end{tabbing}

\vspace{-6pt}
И, наконец, уровень $K_3$:
%	 □		(3, 8)

\begin{tabbing}
{\color{blue} 18.\ } \= $\Not A \Or C\qquad$ \= $(1,2)$  \kill
{\color{blue} 12.} \> $\square$ \> $(3,8)$ \\  
\end{tabbing}

\vspace{-6pt}
Таким образом, пустой дизъюнкт получен на 12-ом шаге.

\end{frame}

\begin{frame}
\frametitle{4.6.9. Стратегии метода резолюций (4/6)}
\vspace{-2pt}
В \DMref{п.~4.6.7}{4.6.7.} пустой дизъюнкт 
для того же исходного множества дизъюнктов получен за
7 шагов. При этом для каждой резольвенты один из родительских дизъюнктов получен на предыдущем шаге. Второй дизъюнкт определяется простым линейным 
просмотром всех ранее полученных дизъюнктов и выбором первого подходящего.   
Такая стратегия называется \odmkey{линейной}{}.
Можно показать, что линейная стратегия также 
является полной.
 
\begin{remark}
Стратегия насыщения уровня соответствует поиску в ширину, а линейная стратегия~--- поиску в глубину. 
\end{remark}

%Замечание. Сравнение показывает, что данная стратегия уступает той, с помощью которой решался этот пример в 4.6.7., называемой линейной стратегией и также являющейся полной.

Стратегия \odmkey{единичной}{} резолюции состоит в том, что на каждом шаге используется единичное выражение~--- положительный литерал. Считается, что эта стратегия  эффективней и быстрее вышеуказанных полных, однако, она не всегда находит решение.

\begin{theorem}
 Единичная резолюция является неполной стратегией.
\end{theorem}
\begin{proof}
Достаточно привести пример, в котором единичная стратегия не может найти опровержение, несмотря на его существование. 
\end{proof}
\end{frame}

\begin{frame}
\frametitle{4.6.9. Стратегии метода резолюций (5/6)}
\begin{proof} (Продолжение)
Рассмотрим 
$K_0 = \{\Not A \Or B, A \Or B, A \Or \Not B,
\Not B \Or \Not A\}$.
%	¬A ⋁ B
%	A ⋁ B
%	A ⋁ ¬B
%	¬B ⋁ ¬A
%	B		(1, 2)
%	A		(3, 5)
%	¬B		(4, 6)
%	□		(5, 7)
%

\vspace{-6pt}
\begin{tabbing}
{\color{blue} 18.\ } \= $\Not A \Or C\qquad$ 
\= $(1,2)$  \kill
{\color{blue} 1.} \> $\Not A \Or B $\> \\
{\color{blue} 2.} \> $A \Or B$\> \\
{\color{blue} 3.} \> $A\Or \Not B$ \> \\
{\color{blue} 4.} \> $\Not A \Or \Not B$ \> \\
{\color{blue} 5.} \> $ B$ \> $(1, 2)$\\
{\color{blue} 6.} \> $A$ \> $(3, 5)$\\
{\color{blue} 7.} \> $\Not B$ \> $(4, 6)$\\
{\color{blue} 8.} \> $\square$ \> $(5, 7)$ \\
\end{tabbing}

\vspace{-\baselineskip}
То есть найдено опровержение линейной стратегией, однако, единичная резолюция его не найдёт, так как положительный литерал $B$ не входит в начальное множество резольвируемых дизъюнктов. 
\end{proof}

\smallskip
Возникает вопрос, при каких условиях эффективные стратегии,
подобные единичной резолюции, могут быть применимы?

\end{frame}

\begin{frame}
\frametitle{4.6.9. Стратегии метода резолюций (6/6)}

\vspace{-2pt}
\odmkey{Хорновским дизъюнктом}{Хорновский дизъюнкт 
(Horn clause)} называется дизъюнкт, содержащий не более одного положительного литерала. Всего возможно три формы хорновских дизъюнктов:

\vspace{-4pt}
\begin{tabbing}
\odmkey{Правило}{}: \= $\Not A_1 \Or \dots \Or\Not A_n \Or B$ \= 
импликация $A_1\And\dots\And A_n \to B$\kill
\odmkey{Факт}{}: \> $A$\> один положительный литерал\\
\odmkey{Правило}{}: \> $\Not A_1 \Or \dots \Or\Not A_n \Or B$ \> 
импликация $A_1\And\dots\And A_n  \to B$\\
\odmkey{Цель}{}: \> $\Not A_1 \Or \dots \Or  \Not A_n $ \> нет положительных литералов\\
\end{tabbing}
%	¬A_0  ⋁ ¬A_1…¬A_n⋁ B
%	A
%	¬A_0  ⋁ ¬A_1…¬A_n
%Первый называется правилом, второй – фактом, а третий – целью. 

\vspace{-16pt}
Заметим, что при применении правила резолюции к хорновским дизъюнктам получается хорновский дизъюнкт, потому что вычеркивается один положительный и один отрицательный литерал и в сумме получаем не более одного положительного литерала. 

\begin{digress}
Одной из немаловажных причин 
появления 
\em{логического программирования} как такового было открытие, что единичная резолюция оказывается 
полной в хорновском случае. 
Единичная резолюция (\odmkey{поиск от данных}{Поиск (search)!от данных (from data)}) является не единственной
эффективной стратегией в хорновском случае.
Например, \odmkey{отрицательная}{Стратегия (strategy)!отрицательная (negative)} стратегия
(\odmkey{поиск от целей}{Поиск (search)!от целей (from target)})
также является полной.   
На основе хорновского случая разработан 
\odmkey{язык Пролог}{Язык (language)!Пролог (Prolog)} и другие языки логического программирования. 
\end{digress}

\end{frame}



% ===== tex/DM2022_5.tex =====
\section{5. Комбинаторный анализ}

\begin{frame}
\frametitle{5. Комбинаторный анализ (1/3)}
Предмету \odmkey{комбинаторного анализа}{Анализ (analysis)!комбинаторный (combinatorial)} или, короче, \odmkey{комбинаторике}{Комбинаторика (combinatorics)} не так просто дать краткое исчерпывающее
определение. В некотором смысле слово <<комбинаторика>>
можно понимать как синоним термина <<дискретная математика>>,
то есть исследование дискретных конечных математических структур.
На школьном уровне с термином
<<комбинаторика>> связывают просто набор известных формул,
служащих для вычисления так называемых \em{комбинаторных чисел\/},
о которых идёт речь в первых разделах этой главы.
Может показаться, что эти формулы полезны
только для решения олимпиадных задач и не имеют отношения
к практическому программированию.
На самом деле это далеко не так.
Вычисления на дискретных конечных математических структурах,
которые часто называют \em{комбинаторными вычислениями\/},
требуют комбинаторного анализа для установления свойств
и выявления оценки применимости используемых алгоритмов.
Рассмотрим элементарный жизненный пример.
\end{frame}

\begin{frame}
\frametitle{5. Комбинаторный анализ (2/3)}
\begin{example}
Предположим, что забыт пароль на каком-то сайте. Так как в базе данных он хранится в виде кодированного образа, то единственный вариант восстановления --- полный перебор всех возможных вариантов. 
В наше время большинство сайтов требуют создать пароль из 8-ми символов.
Если известно, что пароль состоит только из цифр и строчных символов латиницы, то необходимо перебрать $(10+26)^8 \approx 2,82\cdot10^{12}$ вариантов. Современные программы для подбора забытого пароля осуществляют порядка 200 проверок паролей в секунду. Таким образом, на подбор такого простого варианта пароля в худшем случае уйдет около $1,41\cdot10^{10}$ секунд, что составляет около $160000$ дней или примерно $447$ лет. Если же в пароле использованы 
прописные буквы и специальные знаки (так рекомендуется делать), 
то подбор пароля <<в лоб>> практически неосуществим.

%Даже имея хеш пароля у себя на компьютере, вы сможете развить скорость перебора порядка $10^5$ паролей в секунду, %то есть вы потратите $70$ лет своей жизни на перебор всех вариантов. 
%А ведь умные пользователи добавляют в пароль спецсимволы и русские буквы, если это возможно...
%Пусть некоторое агентство недвижимости располагает базой данных из
%$n$ записей, причём каждая запись содержит одно предложение
%(чт$\acute{\text{о}}$ имеется) и один запрос (чт$\acute{\text{о}}$
%требуется) относительно объектов недвижимости. Требуется найти все
%такие пары записей, в которых предложение первой записи совпадает
%с~запросом второй и одновременно предложение второй записи
%совпадает с~запросом первой. (На бытовом языке это называется
%подбором вариантов обмена.) Допустим, что используемая СУБД
%позволяет проверить вариант за одну миллисекунду. Нетрудно
%сообразить, что при <<лобовом>> алгоритме поиска вариантов (каждая
%запись сравнивается с каждой) потребуется $n(n-1)/2$ сравнений.
%Если $n=100$, то все в порядке~--- ответ будет получен за
%4,95~секунды. Но если $n=100~000$ (более реальный случай), то
%ответ будет получен за 4~999~950 секунд, что составляет без малого
%1389 часов и вряд ли может считаться приемлемым. Обратите внимание
%на то, что мы оценили только трудоёмкость подбора \em{прямых\/}
%вариантов, а~существуют ещё варианты, когда число участников
%сделки больше двух\ldots
\end{example}
\end{frame}

\begin{frame}
\frametitle{5. Комбинаторный анализ (3/3)}
\vspace{-2pt}
Этот пример показывает, что практические задачи и
алгоритмы требуют предварительного комбинаторного анализа применимости.
%анализа и количественной оценки. 
Задачи обычно оцениваются с точки зрения \em{размера\/},
то есть общего количества вариантов, среди которых нужно
найти решение, а алгоритмы оцениваются с точки зрения
 \em{сложности\/}. При этом различают \odmkey{сложность по
времени}{Сложность (complexity)!по времени (time)} (или
\odmkey{временн$\acute{\text{у}}$ю сложность}{Сложность
 (complexity)!временн$\acute{\text{а}}$я (time)}),
то есть количество необходимых шагов алгоритма,
и \odmkey{сложность по памяти}{Сложность
(complexity)!по памяти (space)} (или
\odmkey{емкостн$\acute{\text{у}}$ю сложность}{Сложность
(complexity)!емкостн$\acute{\text{а}}$я (space)}),
то есть объём памяти, необходимый для работы алгоритма.
\odmkey{Худшим случаем}{Случай (case)!худший (worst)} называется вариант входных данных, в котором наблюдаются максимальные затраты вычислительного ресурса (времени или памяти).
\odmkey{Лучшим случаем}{Случай (case)!лучший (best)}, соответственно, называется вариант входных данных, в котором наблюдаются минимальные затраты вычислительного ресурса.
Таким образом появляются понятия сложности в лучшем случае (оценка снизу)
и сложности в худшем случае (оценки сверху).

\medskip
Во многих практических случаях возникает необходимость подсчитать
количество возможных комбинаций объектов, удовлетворяющих
определённым условиям. Такие задачи называются
\odmkey{комбинаторными}{Задача (problem)!комбинаторная
(combinatorial)}. 
Во всех случаях основным инструментом такого анализа
оказываются формулы и методы, рассматриваемые в этой главе.


\end{frame}

\begin{frame}
\frametitle{5.1. Комбинаторные задачи}
\subsection{5.1. Комбинаторные задачи}\label{5.1.}
%\label{7.1}

Разнообразие
комбинаторных задач не поддаётся исчерпывающему описанию, но среди
них есть целый ряд особенно часто встречающихся, для которых
известны способы подсчёта.


\bigskip
\begin{tabular}{p{11mm}|p{38mm}|p{97mm}}
  № & Название параграфа
  & Ключевые термины и обозначения \\
  \hline
\DMref{5.1.1.}{5.1.1.}\RS & Комбинаторные конфигурации &
{Правило суммы, правило произведения, принцип Дирихле }
\\  
\DMref{5.1.2.}{5.1.2.} & Размещения &
{ $U(m,n)$~--- число размещений}\\  
\DMref{5.1.3.}{5.1.3.} & Размещения без повторений&
{ $A(m,n)$~--- число размещений без повторений;

убывающий факториал
}\\  
\DMref{5.1.4.}{5.1.4.} & Перестановки &
{ $P(n)$~--- число перестановок;
матрёшка}\\  
\DMref{5.1.5.}{5.1.5.} & Сочетания &
{ $C(m,n)$~--- число сочетаний}\\  
\DMref{5.1.6.}{5.1.6.} & Сочетания с повторениями &
{ $V(m,n)$~--- число сочетаний с повторениями;

возрастающий факториал}\\  
\DMref{5.1.7.}{5.1.7.} & Дискретная вероятность &
{ Распределение вероятности, равномерное распределение, случайная величина,
математическое ожидание}\\  
\end{tabular}  


\end{frame}

\begin{frame}
\label{5.1.1.}\subsubsection{5.1.1. Комбинаторные конфигурации }\frametitle{5.1.1. Комбинаторные конфигурации (1/4)}

Для формулировки и решения комбинаторных задач используются
различные модели \odmkey{комбинаторных конфигураций}{Конфигурация (configuration)!комбинаторная (combinatorial)}. Рассмотрим
следующие две наиболее популярные:

\begin{enumerate}
\item  Дано $n$ мячей. Их нужно разместить по $m$ ящикам так, чтобы
выполнялись заданные ограничения. Сколькими способами это можно сделать?
\item
Рассмотрим множество функций
$\Map{f}{1..n}{1..m}$.
Сколько существует функций $f$, удовлетворяющих заданным ограничениям?
\end{enumerate}


Большей частью соответствие конфигураций, описанных на <<языке
мячей и ящиков>> и на <<языке функций>>, очевидно, поэтому доказательство
правильности способа подсчёта (вывод формулы) можно провести на
любом языке. Если свед$\acute{\text{е}}$ние одной модели к другой
не очевидно, то оно включается в доказательство.

\begin{remark}
Здесь рассматриваются именно функции,
то есть однозначные отношения 
в смысле \DMref{п.~1.6.1}{1.6.1.},
поскольку в большинстве случаев
этого оказывается достаточно,
а выводить формулы для функций проще.
Однако встречаются комбинаторные задачи,
в которых для построения адекватных
конфигураций приходится использовать
многозначные отношения.
\end{remark}

\end{frame}

\begin{frame}
\frametitle{5.1.1. Комбинаторные конфигурации (2/4)} 


При решении комбинаторных задач требуется вычислять 
мощность некоторых 
\em{множеств} (множества функций, множества 
размещений мячей по ящикам). При этом часто оказываются полезными следующие три наблюдения из теории множеств, которым даны специальные названия. 

\smallskip
Неформально \odmkey{правило суммы}{Правило (rule)!суммы (sum)} (или \odmkey{правило сложения}{Правило (rule)!сложения (addition)}) 
формулируется следующим образом: 
если возможности построения комбинаторной конфигурации взаимно исключают друг друга,
то их количества следует складывать. 
На языке теории множеств правило суммы~--- это следствие теоремы \DMref{п.~1.2.7.}{1.2.7.}
\smallskip
\begin{example}
Сколько возможных ходов имеет ферзь, 
стоящий на одном из центральных полей пустой шахматной доски? 
Решение: ферзь может пойти либо по горизонтали (7 ходов), 
либо по вертикали (7 ходов), либо по главной диагонали (7 ходов), 
либо по ортогональной диагонали (6 ходов). Ответ: $7+7+7+6=27$.
\end{example}

\smallskip
Неформально \odmkey{правило произведения}{Правило (rule)!произведения (product)} (или \odmkey{правило умножения}{Правило (rule)!умножения (multiplication)}) 
формулируется следующим образом: если возможности построения комбинаторной конфигурации независимы, то их количества следует умножать. 
На языке теории множеств правило произведения~--- это следствие 1 теоремы \DMref{п.~1.4.2.}{1.4.2.}

\end{frame}

\begin{frame}
\frametitle{5.1.1. Комбинаторные конфигурации (3/4)} 


\begin{example}
Сколькими способами можно на шахматной доске поставить две ладьи 
разного цвета так, 
чтобы они не били друг друга? 
Решение: первую ладью можно поставить на любое из 64 полей, 
при этом она будет бить 14 полей. 
Вторую ладью можно поставить на любое из оставшихся $64-15=49$ полей. 
Ответ: 
$64\cdot49= 3136 $ способов.
\end{example}


Правила суммы и произведения могут применяться совместно.

\begin{example}
Сколькими способами можно на шахматной доске поставить два короля 
разного цвета так, 
чтобы они не били друг друга? 
Решение: король может стоять либо в углу доски (4 занятых поля), 
либо на краю доски (6 занятых полей), либо не на краю (9 занятых полей). 
Ответ: $4\cdot(64-4) + 24\cdot(64-6) + 36\cdot(64-9)= 240 + 1392 + 1980 = 3612$ способов.  
\end{example}

Неформально \odmkey{принцип Дирихле}{Принцип (principle)! Дирихле (Dirichlet)} часто излагают в следующей форме: 
если $n\ne m$, то $n$ мячей невозможно положить в $m$ ящиков так, 
чтобы в каждом ящике было по одному мячу. 
На языке теории множеств принцип Дирихле~--- это следствие теоремы~2 \DMref{п.~1.2.5}{1.2.5.} 
о том, что различные отрезки натурального ряда неравномощны.

\end{frame}

\begin{frame}
\frametitle{5.1.1. Комбинаторные конфигурации (4/4)}


\begin{examples}
\begin{enumerate}
\item
Пусть в равностороннем треугольнике со стороной 1 расположены 5 точек. 
Тогда существует по крайней мере две точки, расстояние 
между которыми не больше ${}^1\!/\!{}_2$. 
Действительно, 
средние линии треугольника разбивают его на четыре равносторонних треугольника 
со стороной ${}^1\!/\!{}_2$. 
По принципу Дирихле две точки попадут в один из этих треугольников. 
Расстояние между ними будет не больше ${}^1\!/\!{}_2$.
\item Пусть грани куба окрашены в два цвета. 
Тогда найдутся соседние одноцветные грани. 
Действительно, рассмотрим три грани куба, имеющие общую вершину. 
Назовём их <<кроликами>>, а данные цвета~--- <<клетками>>. 
По принципу Дирихле, найдутся две соседние грани, окрашенные в один цвет. 
\end{enumerate}
\end{examples}

\begin{digress}
Понятия комбинаторного анализа 
применяются в самых разных областях,
при этом зачастую с использованием специфической
терминологии и традиций в обозначениях. 
Например, в статистике и в теории 
вероятности используются термины
\odmkey{$n$-выборка}{Выборка (sample)} из
\odmkey{$m$-генеральной совокупности}{Генеральная совокупность (general population)},
а в программировании рассматриваются
\odmkey{массивы}{Массив (array)}
\textbf{array $[1..n]$ of $1..m$}.
Для подсчёта
количества выборок или массивов,
обладающих определёнными свойствами, 
то применяются в точности те же формулы,
которые выведены на <<языке функций>>
или на <<языке ящиков>>.
\end{digress}
\end{frame}


\begin{frame}
\label{5.1.2.}\subsubsection{5.1.2. Размещения}\frametitle{5.1.2. Размещения}

Число всех функций $\Map{f}{1..n}{1..m}$ (при отсутствии ограничений), или
число всех возможных способов разместить $n$ мячей по $m$ ящикам,
называется \odmkey{числом размещений}{Число
 (number)!размещений (of menage)} и обозначается
$U(m,n)$.


\begin{theorem}
$U(m,n) = m^n$.
\end{theorem}

\begin{proof}
Поскольку ограничений нет, мячи размещаются независимо друг от друга, то есть каждый из $n$ мячей можно разместить $m$ способами. Таким образом, по правилу произведения получаем $m^n$ возможных размещений.
\end{proof}

\begin{examples}
\begin{enumerate}
\item	Число всех возможных бинарных последовательностей длины $n$ равно $U(2,n)=2^n$. Действительно, на каждую из $n$ позиций в последовательности приходится 2 варианта (1 или 0), значения в разных позициях независимы. 
\item
У кодового замка пароль состоит из трёх позиций, на каждой из них может стоять любая из цифр от 0 до  9. Сколько возможных паролей на данном замке?
Решение: $U(10,3) = 10^3=1000$.
\item Число всех возможных массивов 
\textbf{array $[1..n]$ of $1..m$} равно $m^n$,
поскольку в каждый из элементов массива
может иметь любое значение.
\end{enumerate}
\end{examples}
%При игре в кости бросаются две кости и выпавшие на верхних гранях
%очки складываются. Какова вероятность выбросить 12 очков?
%Каждый возможный исход соответствует функции
%$\Map{F}{\{1,2\}}{\{1,2,3,4,5,6\}}$
%(аргумент~--- номер кости, результат~--- очки на верхней грани).
%Таким образом, всего возможно
%$U(6,2) = 6^2 = 36$
%различных исходов. Из них только один ($6+6$) даёт
%двенадцать очков. Вероятность $1/36$.



%\vspace{-12pt}
\end{frame}

\begin{frame}
\label{5.1.3.}\subsubsection{5.1.3. Размещения без повторений }\frametitle{5.1.3. Размещения без повторений (1/2)}

%\vspace{-4pt}
Число инъективных функций $\Map{f}{1..n}{1..m}$, 
или число способов разместить $n$
мячей по $m$ ящикам, не более чем по одному в ящик, называется
\odmkey{числом размещений без повторений}{Число
(number)!размещений без повторений (of simple menage)} и
обозначается $A(m,n)$, или $[m]_n$, или $(m)_n$.



\begin{theorem}
$A(m,n)  = \displaystyle{\frac{m!}{(m-n)!}}.$
\end{theorem}

\begin{proof}
Ящик для первого мяча можно выбрать $m$ способами, для
второго мяча~--- $m-1$ способами и~т.~д. Имеем
$
A(m,n) =  m (m-1) \ldots (m-n+1)={m!}/{(m-n)!}.$
\end{proof}


Из определения ясно, что
$A(m,n) \Def  0$ при $n > m$ и
$A(m,0) \Def  1$.

\begin{example}
В некоторых видах спортивных соревнований исходом
является определение участников, занявших первое,
второе и третье места.
Сколько возможно различных
исходов, если в соревновании участвуют $n$ участников?
Каждый возможный исход соответствует
функции $\Map{F}{\{1,2,3\}}{1..n}$
(аргумент~--- номер призового места, результат~--- номер участника).
Всего возможно
$A(n,3)=n(n-1)(n-2)$ различных исходов.
\end{example}

\begin{remark}
Функцию $(x)_n\Def x (x-1) \ldots (x-n+1) =
\prod_{i=0}^{n-1} (x- (i)) $   
называют \odmkey{убывающим факториалом}
{Факториал (factorial)! убывающий (falling) } (часто обозначают $x^{\underline{n}}$).
Имеем $(m)_n=A(m,n)$. 
\end{remark}

\end{frame}

\begin{frame}
\frametitle{5.1.3. Размещения без повторений (2/2)}

\begin{example}
Число массивов
\textbf{$a:$ array $[1..n]$ of $1..m$},
таких что $\A{i\ne j}{a[i]\ne a[j]}$,
является числом инъективных функций 
и вычисляется убывающим факториалом.
\end{example}


\begin{digress}
Простые формулы, выведенные для числа размещений без повторений,
дают повод поговорить об элементарных, но весьма важных вещах.
Рассмотрим две формулы:
%\vspace{-10pt}
$$
A(m,n)=m\cdot(m-1)\cdot\ldots\cdot(m-n+1) \quad \text{и} \quad
A(m,n)=\frac{m!}{(m-n)!}.
$$

%\vspace{-6pt}
Формула слева выглядит сложной и незавершённой, формула справа~---
изящной и <<математичной>>. Но формула~--- это частный случай
алгоритма. При практическом вычислении левая формула оказывается
намного предпочтительнее правой. Во-первых, для вычисления по
левой формуле потребуется $n$ умножений, а по правой~---
$2m\!-\!n\!-\!2$ умножений и одно деление. Поскольку $n<m$,
левая формула вычисляется существенно быстрее. 
Во-вторых, число
$A(m,n)$ растёт довольно быстро и при больших $m$ и $n$ может не
поместиться в разрядную сетку. Левая формула работает правильно,
если результат помещается в разрядную сетку. При вычислении же по
правой формуле переполнение может наступить <<раньше времени>>
(то~есть промежуточные результаты не помещаются в разрядную сетку,
в то время как окончательный результат мог бы поместиться),
поскольку факториал~--- 
\em{очень} быстро растущая функция.
\end{digress}

\end{frame}

\begin{frame}\label{5.1.4.}\subsubsection{5.1.4. Перестановки }\frametitle{5.1.4. Перестановки (1/2)}
%\label{7.1.4}


Если $|X|=|Y|=n$, то существуют взаимно-однозначные функции
$\Map{f}{X}{Y}.$ 
Число взаимно-однозначных функций $\Map{f}{1..n}{1..n}$, или
\odmkey{число перестановок}{Число (number)!перестановок (of
permutations)} $n$ предметов, обозначается $P(n)$.

\begin{theorem}
$P(n) = n!.$
\end{theorem}

\begin{proof} $P(n)\!=\!A(n,n)\!=\!n\cdot(n\!-\!1)\cdot\ldots\cdot(n\!-\!n\!+\!1)\!=\!
n\cdot(n\!-\!1)\cdot\ldots\cdot 1\!=\!n!$.
\end{proof}

\medskip
\begin{example}
Число массивов
\textbf{$a:$ array $[1..n]$ of $1..n$},
таких что $\A{i\ne j}{a[i]\ne a[j]}$,
является числом биективных функций 
и вычисляется факториалом.
\end{example}

\medskip
В качестве развернутого примера рассмотрим понятие цепочки, введённое в \DMref{п.~1.8.2.}{1.8.2.}
Напомним, что последовательность $\cal E=(\Tuple{E}{1}{m})$ непустых подмножеств
множества $E$ %\linebreak
($\cal E\subset 2^E$, $E_i\subset E$, $E_i\not=\emptyset$)
называется \odmkey{цепочкой}{Цепочка (chain)!множеств (of sets)} в
$E$, если
$$
\A{i\in 1..(m-1)}{E_i\subset E_{i+1}\And E_i\not= E_{i+1}}.
$$

Цепочка $\cal E$ называется \odmkey{полной}{Цепочка (chain)!полная
(complete)} цепочкой в $E$, если $\vert\cal E\vert=\vert E\vert$.
Сколько существует полных цепочек? 
\end{frame}

\begin{frame}\frametitle{5.1.4. Перестановки (2/2)}
%\label{7.1.4}

\begin{remark}
Полную цепочку иногда образно называют
\odmkey{матрёшкой}{Матрёшка (nest)}.
\end{remark}

\medskip
Очевидно, что в полной цепочке
каждое следующее подмножество $E_{i+1}$ получено из предыдущего
подмножества $E_i$ добавлением ровно одного элемента из $E$ и
таким образом,
 $$\vert E_1\vert=1,
\vert E_2\vert=2, \dots, \vert E_m\vert=m.$$

Следовательно, полная цепочка определяется порядком, в котором
элементы $E$ добавляются для образования очередного элемента
полной цепочки. Отсюда количество полных цепочек~---
это количество перестановок элементов множества~$E$, равное~$m!$.

\medskip
\begin{remark}
В примере рассматриваются конечные цепочки. Понятие полной цепочки применимо и для бесконечных множеств.
\end{remark}




\end{frame}


\begin{frame}
\label{5.1.5.}\subsubsection{5.1.5. Сочетания }\frametitle{5.1.5. Сочетания (1/3)}
%\label{5.1.5}


Число строго монотонно возрастающих функций $\Map{f}{1..n}{1..m}$, или
число размещений $n$ неразличимых предметов по $m$ ящикам не более чем по
одному в ящик, то~есть число способов выбрать из $m$ ящиков $n$ ящиков с
предметами, называется \odmkey{числом сочетаний}{Число
 (number)!сочетаний (of combination)}
и обозначается $C(m,n)$, или $C_m^n$, или $\binom{m}{n}$.

По определению $C(m,n)\Def 0$ при $n>m$.
%\pagebreak[3]
\begin{theorem}
$\displaystyle{C(m,n) = \frac{m!}{n! (m-n)!}}.$
\end{theorem}


%\pagebreak
\begin{proof}
	\proofsection{\!Обоснование формулы\!} Сочетание является
	размещением без повторений неразличимых предметов. Следовательно,
	число сочетаний определяется тем, какие ящики заняты предметами,
	поскольку перестановка предметов при тех же занятых ящиках
	считается одним сочетанием. Таким образом,
	$
	\displaystyle{C(m,n)=\frac{A(m,n)}{P(n)}=\frac{m!}{n! (m-n)!}}.
	$
\end{proof}

\begin{remark}
Для подсчёта числа сочетаний выделяют граничные случаи: 
	$C(m,1) = m$, $C(m,m-1) = m$, $C(m,m) = 1$, $C(m,0) = 1$,
каждый из которых очевидным образом получается прямой подстановкой параметров в формулу. 
	
\end{remark}

\end{frame}


\begin{frame}
\frametitle{5.1.5. Сочетания (2/3)}
\begin{proof} (Окончание)	
\proofsection{\!Свед$\acute{\text{е}}$ние моделей\!} Число
сочетаний является количеством строго монотонно возрастающих функций,
потому что любая строго монотонно возрастающая функция
$\Map{f}{1..n}{1..m}$ определяется набором своих значений, причём
$1 \le f(1)<\ldots <f(n) \le m$. Другими словами, каждая
строго монотонно возрастающая функция определяется выбором $n$
чисел из диапазона $1..m$. Таким образом, число строго монотонно
возрастающих функций равно числу $n$-элементных подмножеств
$m$-элемент\-ного множества, которое
равно числу способов выбрать $n$ ящиков с предметами из $m$
ящиков.
\end{proof}


\begin{example}
	В начале игры в домино каждому играющему выдаётся 7 костей из
	имеющихся 28~различных костей. Сколько существует различных
	комбинаций костей, которые игрок может получить в начале игры?
	Очевидно, что искомое число равно числу 7-элементных подмножеств
	28-элементного множества. Имеем
	
	
	$$
	C(28,7)=\frac{28!}{7!(28-7)!}=
	\frac{28\cdot 27\cdot 26\cdot 25\cdot 24\cdot 23\cdot 22}
	{7\cdot 6\cdot 5\cdot 4\cdot 3\cdot 2\cdot 1}=
	1\:184\:040.
	$$
\end{example}

\end{frame}
\begin{frame}
\frametitle{5.1.5. Сочетания (3/3)}



\begin{example}
Сколько существует способов записать натуральное число $n$ в виде суммы $m$ натуральных чисел:
$a_1 + a_2 + … + a_m = n$?
Любое натуральное число можно представить в виде суммы единиц. Тогда решение задачи сводится к разделению $n$ выстроенных  в ряд единиц на $m$ непустых групп с помощью $m-1$ перегородки. Для перегородки существует $n-1$ позиция между единицами. Таким образом, число способов расставить перегородки: $C(n-1,m-1)$.
\end{example}

\vspace{-6pt}
\begin{center}
\includegraphics[scale=0.4]{5_1_5.jpg}
\end{center}	

\vspace{-4pt}
\begin{example}
Число массивов
\textbf{$a:$ array $[1..n]$ of $1..m$},
таких что $\A{i< j}{a[i]< a[j]}$,
является числом строго монотонных функций 
и вычисляется числом сочетаний.
\end{example}

\end{frame}

\begin{frame}
\label{5.1.6.}\subsubsection{5.1.6. Сочетания с повторениями }\frametitle{5.1.6. Сочетания с повторениями (1/2)}

Число монотонно возрастающих функций $\Map{f}{1..n}{1..m}$, 
или число размещений $n$
\em{неразличимых} мячей по $m$ ящикам, называется \odmkey{числом
сочетаний с повторениями}{Число (number)!сочетаний с повторениями
(of combination with repetition)} и~обозначается $V(m,n)$.

\smallskip

\begin{theorem}
$V(m,n) = C(n+m-1,n)$.
\end{theorem}



%\vspace{-12pt}
\begin{proof}
Монотонно возрастающей функции $\Map{f}{1..n}{1..m}$
однозначно соответствует
строго монотонно возрастающая функция $\Map{g}{1..n}{1..(n+m-1)}$.
Это соответствие устанавливается следующей формулой:
$g(k)=f(k)+k-1$.
\end{proof}

\smallskip
\begin{remark}
	$V(0,n)=C(n-1,n)=0$ при $n>0$ по определению для $C(m,n)$.
	Однако, обычно считают, что $V(0,0)\Def 1$.
\end{remark}




\begin{example}
Сколькими способами можно рассадить $n$ вновь
прибывших гостей среди $m$ гостей, уже сидящих за круглым столом?
Очевидно, что  имеется $m$
промежутков, в которые можно рассаживать вновь прибывших,
и число способов равно
$$V(m,n)=C(m+n-1,n)=\frac{(m+n-1)!}{n!(m-1)!}.$$
\end{example}

\end{frame}


\begin{frame}
\frametitle{5.1.6. Сочетания с повторениями (2/2)}
\begin{theorem}
	$V(m,n)=V(n+1,m-1)$.
\end{theorem}

\begin{proof}
	$\displaystyle{V(m,n)=C(m+n-1,n)=\frac{(m+n-1)!}{n!(m-1)!}=\frac{(n+m-1)!}{ (m-1)! n!}=}$ \\
	$\displaystyle{=C(n+m-1,m-1)=V(n+1,m-1).}$
\end{proof}
 
\smallskip
\begin{example}
Сколько существует способов записать натуральное число $n$ в виде суммы $m$ неотрицательных целых чисел: $a_1 + a_2 + … + a_m = n$?
\newline
Сведем задачу к примеру из \DMref{п.~5.1.5}{5.1.5.} с натуральными числами. 
Пусть $y_i=a_i+1$. Тогда $y_1 + y_2 + … + y_m=n+m$. 
Ответ: $C(m+n-1,m-1) = V(m,n)$.
\end{example}

\begin{remark}
Функцию $x^{(n)}\Def x (x+1) \ldots (x+n-1) =
\prod\limits_{i=0}^{n-1} (x+i) $   
называют \odmkey{возрастающим факториалом}
{Факториал (factorial)! возрастающий (raising)}
(часто обозначают $x^{\overline{n}}$).
Имеем $m^{(n)}=(m+n-1)!/(m-1)!=V(m,n)\cdot n!$. 
\end{remark}

\medskip
\begin{example}
Число массивов
\textbf{$a:$ array $[1..n]$ of $1..m$},
таких что $\A{i< j}{a[i]\le a[j]}$,
является числом монотонно возрастающих функций 
$V(m,n)=m^{(n)}/n!$.
\end{example}
\end{frame}

\begin{frame}
\frametitle{Онтология комбинаторного анализа}
\begin{center}
\includegraphics[height=7.8cm]{ODM5_1_1.png}
\end{center}

%{\color{red} Контрольные вопросы:
%что не так? Как должно быть? Почему?}
\end{frame}

\begin{frame}
\subsubsection{5.1.7. Дискретная вероятность }
\label{5.1.7.}
\frametitle { 5.1.7. Дискретная вероятность (1/7)}
\vspace{-2pt}
Пусть задано конечное множество $X=\{\Tuple{x}{1}{n}\}$,
$|X|=n$, и тотальная функция $\Map{\mathrm{P}}{X}{[0;1]}$,
такая, что $\sum\limits_{i=1}^n \mathrm{P}(x_i) =1$.
Тогда число $p_i\Assign\mathrm{P}(x_i)$ 
называется \odmkey{вероятностью}{Вероятность (probability)}
элемента $x_i$, а набор чисел 
$\mathrm{P}_X = (\Tuple{p}{1}{n})$ называется 
\odmkey{распределением вероятности}{Вероятность (probability)!распределение (distribution)} на множестве $X$.
Распределение вероятности $\mathrm{P}_X$, 
в котором $\A{i \in 1..n}{p_i = \frac{1}{n}}$, называется \odmkey{равномерным}{Вероятность (probability)!распределение (distribution)!равномерное (uniform)}. 

Элементы множества $X$ могут иметь произвольную 
природу: это могут быть результаты измерения 
числового значения некоторой физической величины,
способы разложить мячи по ящикам в комбинаторной
конфигурации, альтернативные события реального мира,
которые могут произойти в будущем и~т.~д.
	
\begin{example}
Одновременно бросаются два игральных кубика. Какова вероятность того, что на двух кубиках выпадут шестёрки?
Всего возможно 36 различных комбинаций, 
распределение естественно \em{считать} равномерным
(кубики <<честные>>),  значит вероятность двух шестёрок равна $\frac{1}{36}$.

\end{example}

\end{frame}

\begin{frame}
\frametitle { 5.1.7. Дискретная вероятность (2/7)}



\begin{digress}
\odmkey{Вероятность}
{Вероятность (probability)}~--- весьма общее понятие, которое изучается в специальном разделе математики~--- \odmkey{теории вероятности}{Теория (theory)!вероятности (of probability)}. Теория вероятности тесно связана с другими разделами математики: математической статистикой, теорией игр, теорией меры и др. Также теория вероятности нашла множество приложений в физике, экономике и биологии. 
Здесь используется \em{упрощённое} понятие вероятности, достаточное для проведения комбинаторных вычислений.
В настоящее время чаще употребляется термин 
\em{теория вероятностей}, подобный терминам
\em{теория множеств},
\em{теория графов}  и так далее.
Такое словоупотребление подчёркивает
однородность класса объектов, свойства которых изучаются в данной теории. Мы же используем
несколько устаревший термин, чтобы подчеркнуть,
что рассматриваем только \em{одно} 
свойство~--- вероятность,
не претендуя на охват всего класса 
объектов, могущих обладать этим свойством.
\end{digress}

\medskip
Рассмотрим множество $A \subseteq  X$. Если на $X$ задано распределение вероятности, то \odmkey{вероятностью подмножества}{Вероятность (probability)!подмножества (of subset)} $A$ называется число, равное $\sum\limits_{a \in A}{\mathrm{P}(a)}$.

Из определения ясно, что $ \mathrm{P}(\emptyset) = 0, \mathrm{P}(X) = 1, \forall A \; 0 \le \mathrm{P}(A) \le 1$.

\begin{example}
В игральной колоде 52 карты четырёх мастей, в каждой из которых по 13 карт разного достоинства.
Вероятность быть тузом равна $\frac{1}{52}+\frac{1}{52}+\frac{1}{52}+\frac{1}{52} = \frac{4}{52} = \frac{1}{13}$.
\end{example}

\end{frame}

\begin{frame}
\frametitle { 5.1.7. Дискретная вероятность (3/7)}
Стоит заметить, что в случае равномерного распределения вычисление вероятности подмножества $A$ упрощается: $\mathrm{P}(A) = \displaystyle\frac{|A|}{|X|}$.

\begin{example}
В игре <<покер>> комбинация из пяти карт одинаковой масти называется <<флеш>>. Используется обычная игровая колода из 52-ух карт.
Какова вероятность комбинации <<флеш>>? 
Множество $X$ содержит в качестве элементов всевозможные комбинации из пяти карт, порядок которых значения не имеет. Количество таких элементов равно $C(52,5)$.
В одной масти 13 карт, из которых в комбинацию входят пять любых. Мастей всего 4, значит количество элементов в множестве $A$ равно $4 \cdot C(13,5)$.
Таким образом, вероятность комбинации <<флеш>> равна $\frac{4 \cdot C(13,5)}{C(52,5)} = \frac{33}{16660} \approx 0,002$.
\end{example}

\end{frame}


\begin{frame}
\frametitle { 5.1.7. Дискретная вероятность (4/7)}
%\vspace{-4pt}
\odmkey{Случайная величина}{Величина (variable)!случайная (aleatory, random)}~--- это тотальная функция $\Map{\varphi}{X}{{\Bbb R}}$, причём на множестве $X$ заданы вероятности элементов.
\odmkey{Математическим ожиданием}{Математическое ожидание (expected value)} случайной величины при распределении вероятности $\mathrm{P}_X = (\Tuple{p}{1}{n})$ называется число
\vspace{-6pt}
$$\mathrm{E}_{\varphi} = \sum_{k=1}^{n}{\varphi(x_k) p_k} = \varphi(x_1) p_1 + \ldots + \varphi(x_n) p_n.$$

\vspace{-6pt}
Легко видеть, что среднее арифметическое~--- частный случай математического ожидания для равномерного распределения.
Поэтому математическое ожидание можно назвать средним значением случайной величины $\varphi(X)$.
\begin{remark}
Если элементы множества $X$ являются числами, то в качестве функции $\varphi$ может рассматриваться тождественная функция. 
Формула математического ожидания примет вид $\sum_{k=1}^{n}{x_k p_k}$.
\end{remark}

\begin{example}
При броске игрального кубика
$X=\{1,2,3,4,5,6\}$. Математическое ожидание числа очков на верхней грани 
$E = \frac{1}{6} \sum_{k=1}^{6}{x_k} = \frac{1}{6} (1 + 2 + 3 + 4 + 5 + 6) = 3,5$.
В среднем при одном подбрасывании кубика выпадает 3,5 очка.
\end{example}

\end{frame}


\begin{frame}
\frametitle { 5.1.7. Дискретная вероятность (5/7)}
\begin{remark}
Математическое ожидание~--- одно из важнейших понятий в теории вероятности. Имеет важное прикладное значение при оценке рисков, при торговле на финансовых рынках и во многих других областях.
\end{remark}

\begin{example}
В казино предлагают сыграть в игру: игрок ставит некоторую сумму на число от 1 до 6, а потом бросает три игральных кубика. Если загаданное число не выпадет ни разу, то его ставка сгорает. Если выпадет~--- ставку возвращают, а казино платит поставленную сумму столько раз, сколько раз выпало число.
Рассмотрим средний выигрыш при ставке в 1 доллар. Игрок может потерять доллар или выиграть 1, 2 или 3 доллара.
Тогда множество $X = \{-1, 1, 2, 3\}$,
а математическое ожидание: $$-\frac{125}{216} + \frac{ C(3,1) \cdot 5^2}{216} + 2 \cdot \frac{C(3,2) \cdot 5}{216} + 3 \cdot \frac{1}{216} = -\frac{17}{216} \approx -0,08$$
В среднем при каждой ставке в один доллар игрок потеряет восемь центов. 
\end{example}

\end{frame}


\begin{frame}
\frametitle { 5.1.7. Дискретная вероятность (6/7)}
К подмножествам множества $X$ можно применять обычные
операции: пересечение, объединение и другие.   
\begin{theorem} Если $A,B\subset X$, то
$\mathrm{P}(A\cup B)=\mathrm{P}(A) + 
\mathrm{P}(B)-\mathrm{P}(A\cap B)$.
\end{theorem}
\begin{proof}
Вероятности тех элементов, которые входят в пересечение, надобно включить в сумму только один раз.
\end{proof}

\medskip
По Колмогорову
\odmkey{условная вероятность}{Вероятность (probability)!условная (conditional)}
$\mathrm{P}(A|B)$ определяется так:
$\displaystyle{\ \mathrm{P}(A|B)\Def\frac{\mathrm{P}(A\cap B)}{\mathrm{P}(B)}}$. 
\begin{digress}
Существует и с успехом используется большое количество
других определений, интерпретаций и трактовок понятия условной вероятности и связанных с ним
понятий \odmkey{априорной}{Вероятность (probability)!априорная (prior)}
и \odmkey{апостериорной}{Вероятность (probability)!апостериорная (posterior)} 
вероятности. Все эти вариации оправданы и 
даже необходимы при использовании вероятностных
подходов в конкретных предметных областях.
Здесь мы ограничиваемся простейшей
возможной интерпретацией,
которой, однако, достаточно, чтобы вывести
формулы, необходимые для комбинаторных вычислений.
\end{digress}

\begin{remark}
Множества, входящие в определение операции условной вероятности, \\равноправны: $\quad\mathrm{P}(B|A)=\mathrm{P}(A\cap B)/\mathrm{P}(A)$. 
\end{remark}

\end{frame}


\begin{frame}
\frametitle { 5.1.7. Дискретная вероятность (7/7)}

\begin{theorem} {\NB{[формула Байеса]}}
$\mathrm{P}(A|B)\mathrm{P}(B)=\mathrm{P}(B|A)\mathrm{P}(A)$.
\end{theorem} 

\begin{proof}
$\mathrm{P}(A|B)\mathrm{P}(B)=
\mathrm{P}(A\cap B)=
\mathrm{P}(B|A)\mathrm{P}(A)$.
\end{proof}

\begin{remark}
Формулу Байеса часто записывают в несимметричной
форме 
$$\mathrm{P}(A|B)=
\frac{\mathrm{P}(B|A)\mathrm{P}(A)}{\mathrm{P}(B)}
$$
и называют $A$ событием, а $B$ --- гипотезой.
На применение формулы для вычислений это не влияет.  
\end{remark}

\begin{theorem} {\NB{[формула полной вероятности]}}
Если семейство множеств $\Tuple{B}{1}{n}$ является разбиением множества $X$,
то есть, если $B_1\cup\dots\cup B_n = X$ и 
$i\neq j\Impl B_i\cap B_j=\emptyset$, то
$$
\mathrm{P}(A) = \sum_{i=1}^{n} \mathrm{P}(A|B_i)\mathrm{P}(B_i).
$$
\end{theorem}

\begin{proof}
$\mathrm{P}(A) =\sum_{i=1}^{n} \mathrm{P}(A\cap B_i)= \sum_{i=1}^{n} \mathrm{P}(A|B_i)\mathrm{P}(B_i)$.
\end{proof}

\end{frame}



\begin{frame}
\frametitle{Онтология комбинаторного анализа}
\begin{center}
\includegraphics[height=7.8cm]{ODM5_1_2.png}
\end{center}

%{\color{red} Контрольные вопросы:
%что не так? Как должно быть? Почему?}
\end{frame}

\begin{frame}
\frametitle{5.2. Перестановки  (1/2)}
\subsection{5.2. Перестановки}\label{5.2.}

Для вычисления количества перестановок 
в~\DMref{п.~5.1.4}{5.1.4.} установлена
очень простая формула: $P(n)=n!$. Применяя эту формулу при решении
практических задач, не следует забывать, что факториал~--- это
быстро растущая функция, в частности, факториал
растёт \em{быстрее} экспоненты. 

%Действительно, используя известную из
%математического анализа \odmkey{формулу Стирлинга}{Формула
%(formula)!Стирлинга (James Stirling)}
%$$
%n^n \sqrt{2\pi n}\, e^{-n} < n! < n^n \sqrt{2\pi n}\,
%e^{-n+\frac{1}{12n}},
%$$
%нетрудно показать, что
%$$
%\lim_{n\to+\infty}\frac{n!}{2^n}=+\infty.
%$$
\vspace{-14pt}
\begin{figure}[C]
\includegraphics[height=65mm]{5_2.jpg}
\end{figure}

\end{frame}

\begin{frame}
\frametitle{5.2. Перестановки (2/2)}

\begin{tabular}{p{11mm}|p{45mm}|p{88mm}}
  № & Название параграфа
  & Ключевые термины и обозначения \\
  \hline
\DMref{5.2.1.}{5.2.1.} & Циклы в перестановках &
{Таблица подстановки; подвижная и неподвижная точка; петля; транспозиция; орбита }
\\  
\DMref{5.2.2.}{5.2.2.} & Порядок перестановки &
{ }\\  
\DMref{5.2.3.}{5.2.3.} & Инверсии &
{ $I(f)$~--- число инверсий; задача сортировки; сортировка методом пузырька
}\\  
\DMref{5.2.4.}{5.2.4.} & Генерация перестановок &
{ Лексикографический и антилексекографический порядок }\\  
\DMref{5.2.5.}{5.2.5.} & Двойные факториалы &
{ $n!!$ }\\  
\DMref{5.2.6.}{5.2.6.} & Быстрая сортировка &
{ Опорный элемент}\\  
\DMref{5.2.7.}{5.2.7.} & Трудоемкость в среднем &
{ Алгоритм бинарного поиска}\\  
\DMref{5.2.8.}{5.2.8.} & Гармонические числа &
{ Постоянная Эйлера}\\  
\DMref{5.2.9.}{5.2.9.} & Оценка в среднем для быстрой сортировки &
{ }\\  
\end{tabular}  

\end{frame}


\begin{frame}
\label{5.2.1.}\subsubsection{5.2.1. Циклы в перестановках }\frametitle{5.2.1. Циклы в перестановках (1/5)}

%\vspace{-2pt}
В~\DMref{п.~2.2.8}{2.2.8.} рассматриваются 
взаимно-однозначные функции
на отрезках натурального ряда
(перестановки),
которые удобно задавать \odmkey{таблицами подстановки}{Подстановка (substitution)}.
В таблице подстановки нижняя строка (значения функции) является
перестановкой элементов верхней строки (значения аргумента). Если
принять соглашение, что элементы верхней строки (аргументы) всегда
располагаются в определённом порядке (например, по возрастанию),
то верхнюю строку можно не указывать~--- подстановка определяется
одной нижней строкой. Таким образом, подстановки (таблицы)
взаимно-однозначно соответствуют перестановкам (функциям).

\medskip
Напомним, что множество перестановок образует
группу относительно суперпозиции. Нейтральный элемент, то есть тождественную подстановку,
обычно обозначают $e$.


%%%%%%%%%%%%%%%%%%%%%%%%%%%%%%%%%%%%%%%%%%%%%%%%%%%%%%%%%%%%
% перестановку в круглых скобках, это согласуется
% а последовтельность неравенств - в угловых. т.к. путь!
%
%%%%%%%%%%%%%%%%%%%%%%%%%%%%%%%%%%%%%%%%%%%%%%%%%%%%%%%%%%%%

\medskip
Перестановка $f$ (и соответствующая ей подстановка) элементов
$1,\dots,n$ обозначается \\$(\Tuple{a}{1}{n})$,
где  $a_i=f(i)$ и $\Tuple{a}{1}{n}$~--- 
\em{все} различные числа из диапазона $1..n$.
Если $f(i)=i$, то $i$ называется 
\odmkey{неподвижной точкой}{Неподвижная точка (fixed point)!перестановки (of permutation)} перестановки, 
если же $f(i)\ne i$, то $i$ называется \odmkey{подвижной точкой}{Подвижная точка (moving point)!перестановки (of permutation)} перестановки. 
\end{frame}

\begin{frame}
\frametitle{5.2.1. Циклы в перестановках (2/5)}
Если задана перестановка $f$, то \odmkey{циклом}{Цикл (cycle)}
называется \em{последовательность} 
элементов\\ $\Tuple{x}{1}{k}$, такая,
что
$f(x_i) = x_{i+1}$ при $1 \le i < k$ и
$f(x_k) =x_1$

Количество элементов в цикле называется \odmkey{длиной}{Длина (length)!цикла (of cycle)} цикла. 
Цикл длины 1 называется \odmkey{тривиальным}{Длина (length)!цикла (of cycle)} циклом или \odmkey{петлёй}{Петля (loop)}. 
Цикл длины 2 называется \odmkey{транспозицией}{Транспозиция (transposition)}.
Иногда перестановки удобно представлять в графической форме,
проводя стрелки от каждого элемента $x$ к элементу
$f(x)$.

\begin{columns}[t]
\begin{column}[T]{9.5cm}
\begin{example}
Графическое представление перестановки 
%(см.~п.~\ref{3.1.6})

\vspace{-8pt}
$$
f =
\begin{vmatrix}
1 & 2 & 3 & 4 & 5 \\
2 & 3 & 1 & 4 & 5
\end{vmatrix}
$$

\vspace{-2pt}
показано на рисунке.
%.~\ref{fig51}.
\end{example}
\end{column}
\begin{column}[T]{5.5cm}
\begin{figure}[h]
\begin{center}
\includegraphics[height=20mm]{5_2_1.jpg}
\end{center}
\end{figure}
\end{column}
\end{columns}

\medskip
\begin{remark}
Из графического представления перестановки наглядно видно
происхождение термина <<цикл>>.
\end{remark}

Перестановка может содержать несколько нетривиальных циклов. 
\end{frame}

\begin{frame}
\frametitle{5.2.1. Циклы в перестановках (3/5)}

Перестановка длины $n$, которая содержит ровно один нетривиальный цикл длины $k$,
называется \odmkey{циклической}{Перестановка (permutation)!циклическая (cyclic)} (перестановкой длины $k$), $k \le n$. 
В циклической перестановке все элементы цикла~--- подвижные точки, 
а остальные элементы~--- неподвижные точки.

\begin{example} 
В перестановке на рисунке ранее точки $1, 2, 3$~--- подвижные точки, 
образующие цикл $(1, 2, 3)$, а точки $4$ и $5$~--- неподвижные точки. 
Эта перестановка циклическая.
\end{example}

Множество элементов, входящих в цикл, называется \odmkey{орбитой}{Орбита (orbit)}. 
Орбита тривиального цикла состоит из одного элемента. 

\begin{lemma}
Орбиты различных циклов в перестановке не пересекаются.
\end{lemma}
\begin{proof}
От противного. Допустим, существует элемент $x$, 
который входит в два цикла. Тогда <<следующий>> элемент $f(x)$ 
также входит в оба цикла и т.~д. Ввиду конечности циклов, 
заключаем, что циклы совпадают.
\end{proof}

\begin{corollary}
 Множество орбит образует разбиение перестановки, 
 а отношение <<принадлежать одной орбите>> является эквивалентностью.
\end{corollary} 

\begin{example}    
В перестановке  на рисунке ранее $\{1,2,3\}$~--- орбита. 
\end{example}


\end{frame}

\begin{frame}
\frametitle{5.2.1. Циклы в перестановках (4/5)}

\begin{theorem}
Всякая неединичная перестановка представляется как суперпозиция циклических перестановок с непересекающимися орбитами, причём представление единственно.
\end{theorem}
\begin{proof} {\color{brown}\sffamily\small  [\;
\sffamily\small\color{brown} существование \sffamily\small\color{brown}\;]}
Построим требуемое разложение:
возьмём произвольный элемент и выделим орбиту, 
которой он принадлежит. Построим циклическую перестановку с циклом, 
соответствующим выделенной орбите. 
Затем возьмём любой элемент из оставшихся (не вошедших в найденные орбиты)
 и повторим выделение орбиты, соответствующей элементу. 
Ввиду конечности перестановок процесс закончится.
\proofsection{единственность}
 Допустим, что существует такая перестановка $\sigma$, что ее разложение на суперпозицию непересекающихся циклов не единственно, т.~е. $\sigma=\sigma_1\dots\sigma_m=\tau_1\dots\tau_n$. Тогда выберем некоторые циклы $\sigma_i$, $\tau_j$ ($i,j \in \mathbb{N}, 
 i \leq m, j \leq n$) такие, что для любого $x$, являющегося подвижной точкой, 
 было выполнено $\sigma_i(x) = \tau_j(x)$ 
 (такие $i, j$ обязательно существуют, иначе $\sigma(x) \neq \sigma(x)$~--- противоречие). Тогда
$\sigma_i (x) = \tau_j(x)$,
$\sigma_i^2(x) =  \tau_j^2(x)$,
$\dots$,
$\sigma_i^s(x) = \tau_j^s(x) = x$.
Повторяя рассуждение для всех оставшихся циклов в разложении перестановки, убеждаемся, что $m=n$ и разложение перестановки как суперпозиции непересекающихся циклических перестановок, единственно вплоть до порядка применения циклических перестановок.
\end{proof}

\end{frame}

\begin{frame}
\frametitle{5.2.1. Циклы в перестановках (5/5)}

\begin{remark}
Если орбита состоит из одного элемента, то цикл тривиален,
 а циклическая подстановка является единичной. 
 Единичную подстановку можно не включать в композицию. 
\end{remark}

\begin{example}
Перестановка
$$       
       f =
            \begin{vmatrix}
                1 & 2 & 3 & 4 & 5 & 6 & 7 & 8 & 9 \\
                7 & 6 & 1 & 3 & 9 & 8 & 4 & 2 & 5
            \end{vmatrix}
$$        

        раскладывается в три независимых цикла:

        $$f =  (1743)\circ(59)\circ(682)$$
\end{example}
\end{frame}

\begin{frame}
\label{5.2.2.}\subsubsection{5.2.2. Порядок перестановки }\frametitle{5.2.2. Порядок перестановки (1/3)}


Перестановка~--- это взаимно-однозначная функция. 
Что будет, если переставить элементы так, как задаёт функция, 
а потом переставить ещё раз и ещё раз? 
Другими словами, как ведёт себя степень перестановки?

\medskip
\begin{theorem}
$ (f^k)^{-1} = (f^{-1})^k$.
\end{theorem}

\begin{proof}
По индукции. База~--- по определению. 
Пуcть $ (f^k)^{-1} = (f^{-1})^k$. Тогда 
$(f^{k+1})^{-1} =(f f^k)^{-1} = (f^k)^{-1} (f)^{-1} =(f^{-1})^k f^{-1} = f^{-1} (f^{-1})^k = (f^{-1})^{k+1}$. 
Первый переход по определению степени 
\DMref{п.~1.6.5}{1.6.5.}, второй по свойству групп 
\DMref{п.~2.2.6}{2.2.6.}, 
третий по индукционному предположению, четвертый по лемме \DMref{п.~1.6.5}{1.6.5.},
 а пятый (последний) снова по определению степени.
\end{proof}

\begin{lemma}
Для любой перестановки $f$ существует такое 
натуральное число $k$,  что $f^k=e$.
\end{lemma}

\begin{proof} 
Поскольку различных перестановок лишь конечное число, а именно, $n!$,  
то в бесконечной  последовательности $f, f^2, f^3, f^4, \dots$ 
есть совпадающие элементы. 
Пусть для некоторых натуральных чисел $p, q$ ($p < q$) 
выполняется равенство $f^p = f^q$. 
Легко видеть, что $f^q = f^p\circ f^{q-p}$ 
и в данном случае $f^p = f^p\circ f^{q-p}$.
Домножим обе части на $(f^p)^{-1}$ слева.
Имеем
$(f^p)^{-1}\circ f^p = (f^p)^{-1}\circ f^p\circ f^{q-p}$, 
откуда  $e = f^{q-p}$. 
\end{proof}



\end{frame}


\begin{frame}
\frametitle{5.2.2. Порядок перестановки (2/3)}


Наименьшее число $k$ такое, что $f^k=e$, называется \odmkey{порядком перестановки}{Порядок (order)!перестановки (of permutation)}.

\begin{remark}
	Порядок перестановки~--- это индивидуальная характеристика отдельной перестановки. 
	В то же время порядок группы перестановок $S_n$~--- 
	это общая характеристика всего множества перестановок. 
	Эти два порядка не следует путать!
\end{remark}


 	\begin{example}  %Действительно:
 		
 		$$
 		f =
 		\begin{vmatrix}
 		1 & 2 & 3 & 4 & 5 \\
 		2 & 3 & 1 & 4 & 5	
 		\end{vmatrix},
 		\quad
 		f^2 =
 		\begin{vmatrix}
 		1 & 2 & 3 & 4 & 5 \\
 		3 & 1 & 2 & 4 & 5	
 		\end{vmatrix},
 		\quad
 		f^3 =
 		\begin{vmatrix}
 		1 & 2 & 3 & 4 & 5 \\
 		1 & 2 & 3 & 4 & 5	
 		\end{vmatrix}.
 		$$
 		Порядок перестановки $f$ равен 3.
 	\end{example}  

\begin{corollary}
	Симметрическая группа $S_n$ является периодической.
\end{corollary}
\begin{proof}
По лемме порядок любого элемента
$x\in S_n$ не превосходит некоторого 
натурального $k$, то
есть конечен. Значит $S_n$~--- 
периодическая группа (см. также \DMref{п.~2.2.8}{2.2.8.}, \DMref{п.~4.4.8}{4.4.8.}).
\end{proof}
\end{frame}
\begin{frame}
	\frametitle{5.2.2. Порядок перестановки (3/3)}

	
	\begin{lemma}
		Порядок циклической перестановки равен длине цикла.
	\end{lemma}
	\begin{proof}  
		Пусть цикл содержит $k$ элементов. 
		Все неподвижные точки останутся неподвижными при любом количестве применений
		перестановки, а подвижные точки после $k$ применений вернутся на своё место.
	\end{proof}
	
\medskip	
	\begin{theorem}
		Порядок перестановки равен наименьшему общему кратному длин всех циклов.
	\end{theorem}
	
	\begin{proof}
		Пусть перестановка $f$ разлагается в циклические перестановки по теореме \DMref{п.~5.2.1}{5.2.1.}:
		$f= f_1 \circ f_2 \circ\dots\circ f_s$. 
		Если перестановки $\Tuple{f}{1}{s}$ содержат циклы длины
		$\Tuple{k}{1}{s}$ соответственно, т.~е. имеют порядки $\Tuple{k}{1}{s}$, 
		то наименьшее число $m$, 
		для которого одновременно выполняются все равенства  
		$f_1^m = e$, $f_2^m = e$, ..., $f_s^m = e$, 
		равняется наименьшему общему кратному чисел $k_1, k_2,\dots, k_s$. 
	\end{proof} 
	
	
\end{frame}


\begin{frame}
\label{5.2.3.}\subsubsection{5.2.3. Инверсии }\frametitle{5.2.3. Инверсии (1/3)}

Если в перестановке $f =({\Tuple{a}{1}{n}})$ для элементов $a_i$ и $a_j$
имеет место неравенство \\ $a_i > a_j$ при
$i < j$, то пара $(a_i,a_j)$
называется \odmkey{инверсией}{Инверсия (inversion)}. Обозначим
$I(f)$~--- \odmkey{число инверсий}{Число
 (number)!инверсий (of inversions)} в перестановке $f$.

\begin{theorem}
Произвольную перестановку $f$ можно представить в виде суперпозиции
$I(f)$ транспозиций соседних элементов.
\end{theorem}

\begin{proof}
Пусть $f=({a_1,\dots,1,\dots,a_n})$. Переставим $1$ на первое
место, меняя её местами с соседними слева элементами. Обозначим
последовательность этих транспозиций через $t_1$. При этом все
инверсии (и только они), в которых участвовала $1$, пропадут.
Затем переставим $2$ на второе место и т.~д. Таким образом,
$f\circ t_1\circ\ldots\circ t_n=e$ и по свойству группы
$f=t^{-1}_n\circ\ldots\circ t^{-1}_1$, причём $\vert
t_1\vert+\vert t_2\vert+\ldots +\vert t_n\vert=I(f)$.
\end{proof}

\medskip
Рассмотрим входную последовательность, состоящую из $n$ элементов: 
$({a_1, a_2, \ldots, a_n})$.
	\odmkey{Задача сортировки}{Задача (problem)!сортировки (of sorting)} состоит в поиске такой перестановки $({b_1, b_2, \ldots, b_n })$ входной последовательности, 
	что: $\A{i \in 1..n-1}{b_i \le b_{i+1}}$.

\end{frame}

\begin{frame}
\frametitle{5.2.3. Инверсии (2/3)}

\begin{remark}
Множество, из которого берутся сортируемые элементы, предполагается
линейно упорядоченным.
\end{remark}

\begin{corollary}
Всякая сортировка может быть выполнена перестановкой соседних элементов.
\end{corollary}

\begin{digress}
Доказанная теорема утверждает, что произвольную перестановку можно
представить в~виде композиции определённого числа транспозиций,
но не утверждает, что такое представление является эффективным.
\end{digress}
	
\medskip
	Метод \odmkey{сортировки}{Сортировка (sorting)},
	основанный на предшествующей теореме,
	известен как \odmkey{метод пузырька}{Сортировка
		(sorting)!пузырьком (bubble)}. Заметим, что при
	перемещении элемента на своё место транспозициями соседних
	элементов, все элементы остаются на своих местах, кроме двух
	соседних элементов, которые, возможно, меняются местами.
	
\end{frame}

\begin{frame}
\frametitle{5.2.3. Инверсии (3/3)}

Таким образом, метод пузырька может быть выражен в форме следующего
алгоритма. Этот алгоритм прост, но является далеко не
самым эффективным алгоритмом сортировки.



%\addvspace{-6pt}
%\begin{algorithm}
%\caption{Сортировка методом пузырька}
%\label{bubble}
\begin{algorithmic}
\REQUIRE массив $A: \textbf{array } [1..n] \textbf{ of } B$, 
%где значения элементов массива расположены в произвольном порядке.
 где для значений типа $B$ задано отношение~$<$. 
\ENSURE массив
\textbf{$A:$ array $[1..n]$ of $B$}, отсортированный по возрастанию.%в котором значения расположены в порядке возрастания. 
\FOR{\textbf{$i$ from $2$ to
$n$}}
  \FOR{\textbf{$j$ from $n$ downto $i$}}
    \IF{$A[j]<A[j-1]$}
      \STATE $A[j]\leftrightarrow A[j-1]$
      \COMMENT {\color{gray}\small транспозиция соседних элементов }
    \ENDIF
  \ENDFOR
\ENDFOR
\end{algorithmic}
%\end{algorithm}

\end{frame}

\begin{frame}
\label{5.2.4.}\subsubsection{5.2.4. Генерация перестановок }\frametitle{5.2.4. Генерация перестановок (1/5)}

На множестве перестановок естественным образом можно определить
порядок на основе упорядоченности элементов. А именно, говорят,
что перестановка $({\Tuple{a}{1}{n}})$
\odmkey{лексикографически}{Отношение (relation)!порядка
(order)!лексикографического (lexicographical)} предшествует
перестановке $({\Tuple{b}{1}{n}})$, 
если $\E{k\!\le\!n\!\!}{a_k\!<\!b_k\!\And\! \A{i\!<\!k\!\!}{a_i\!=\!b_i}}$.
Аналогично, говорят, что перестановка $({\Tuple{a}{1}{n}})$
\odmkey{антилексикографически}{Отношение (relation)!порядка
(order)!антилексикографического (antilexico\-graphical)}
предшествует перестановке $({\Tuple{b}{1}{n}})$, если $\E{k\le
n}{a_k>b_k\And \A{i>k}{a_i=b_i}}$.

Следующий алгоритм генерирует \odmkey{}{Алгоритм
(algorithm)!генерации перестановок (generation of permutations)}
все перестановки элементов $1,\dots,n$ в антилексикографическом
порядке, при этом массив 
\textbf{$P:$ array $[1..n]$ of $1..n$} является
глобальным и предназначен для хранения перестановок.

%\addvspace{-3pt}
%\begin{algorithm}
%\caption{Генерация перестановок в антилексикографическом порядке}
%\label{alg52}
\begin{algorithmic}
\REQUIRE $n$~--- количество элементов
\ENSURE последовательность перестановок элементов $1,\dots, n$
в антилексикографическом порядке.

\STATE \textbf{for $i$ from $1$ to $n$ 
do $P[i]\Assign i$ end for} 
\COMMENT{\color{gray}\small инициализация }
\STATE $\Antilex(n)$ 
\COMMENT{\color{gray}\small вызов рекурсивной процедуры $\Antilex$ }
\end{algorithmic}
%\end{algorithm}
\end{frame}

\begin{frame}
\frametitle{5.2.4. Генерация перестановок (2/5)}

%\vspace{-4pt}
Генерация перестановок выполняется
рекурсивной процедурой $\Antilex$. 
%\pagebreak[3]


\begin{algorithmic}
%\small
\REQUIRE $m$~--- параметр процедуры~--- количество
первых элементов массива $P$, для которых генерируются перестановки.
\ENSURE
последовательность перестановок $1,\dots, m$  в антилексикографическом
порядке.
\IF{ $m = 1$}
  \STATE \textbf{yield $P$} 
  \COMMENT{\color{gray}\small  очередная перестановка }
\ELSE
  \FOR{\textbf{$i$ from $1$ to $m$}}
     \STATE $\Antilex ( m - 1 )$ 
     \COMMENT{\color{gray}\small  рекурсивный вызов }
     \IF{ $i< m$ }
        \STATE $P[i] \leftrightarrow P[m]$ 
        \COMMENT{\color{gray}\small  следующий элемент }
        \STATE $\Reverse(m-1)$ 
        \COMMENT{\color{gray}\small  изменение порядка элементов }
     \ENDIF
   \ENDFOR
\ENDIF
\end{algorithmic}

\end{frame}

\begin{frame}
\frametitle{5.2.4. Генерация перестановок (3/5)}
%\vspace{-4pt}
Вспомогательная процедура $\Reverse$ переставляет элементы
заданного отрезка массива $P$ в обратном порядке.



%\vspace{1pt}
\begin{algorithmic}

\REQUIRE $k$~--- номер элемента, задающий отрезок
массива $P$, подлежащий перестановке в~обратном порядке.
\ENSURE первые $k$ элементов массива $P$ переставлены
в обратном порядке
\STATE      $j \Assign  1$ 
\COMMENT{\color{gray}\small нижняя граница обращаемого диапазона }
\WHILE{$j < k$}
  \STATE $P[j] \leftrightarrow P[k]$ 
  \COMMENT{\color{gray}\small меняем местами элементы }
  \STATE $j \Assign  j + 1$ 
  \COMMENT{\color{gray}\small увеличиваем нижнюю границу }
  \STATE $k \Assign  k - 1$ 
  \COMMENT{\color{gray}\small уменьшаем верхнюю границу }
\ENDWHILE
\end{algorithmic}

\medskip
\begin{example}
Последовательность перестановок в антилексикографическом порядке
\\для $n=3$: $(1,2,3)$, $(2,1,3)$, $(1,3,2)$, $(3,1,2)$, $(2,3,1)$, $(3,2,1)$.
\end{example}

\end{frame}

\begin{frame}
\frametitle{5.2.4. Генерация перестановок (4/5)}

\begin{ground}
Заметим, что искомую последовательность перестановок $n$ элементов
можно получить из последовательности перестановок $n-1$ элементов
следующим образом. Нужно выписать $n$ блоков по $(n-1)!$
перестановок в каждом, соответствующих последовательности
перестановок $n-1$ элементов в антилексикографическом порядке.
Затем ко всем перестановкам в первом блоке нужно приписать справа
$n$, во втором~--- $n-1$ и т.~д. в убывающем порядке. Затем в
каждом из блоков (кроме первого), к перестановкам которого справа
приписан элемент $i$, нужно в перестановках блока заменить все
вхождения элемента $i$ на элемент $n$. В полученной
последовательности все перестановки различны, и их $n(n-1)!=n!$,
то~есть перечислены \em{все} перестановки. При этом
антилексикографический порядок соблюдён для последовательностей
внутри одного блока, потому что этот порядок был соблюдён в
исходной последовательности, а для последовательностей на границах
двух блоков~--- потому что происходит уменьшение самого правого
элемента. 
\end{ground}
\end{frame}

\begin{frame}
\frametitle{5.2.4. Генерация перестановок (5/5)}

\begin{ground} (Окончание) 
Обратимся к процедуре $\Antilex$~--- легко видеть, что в
ней реализовано указанное построение. В основном цикле сначала
строится очередной блок~--- последовательность перестановок первых
$m-1$ элементов массива $P$ (при этом элементы $P[m],\dots,P[n]$
остаются неизменными). Затем элемент $P[m]$ меняется местами с
очередным элементом $P[i]$. Вызов вспомогательной процедуры
$\Reverse$ необходим, поскольку последняя перестановка в блоке
является обращением первой, а~для генерации следующего блока на
очередном шаге цикла нужно восстановить исходный порядок.
\end{ground}


%\vspace{-8pt}


\end{frame}

\begin{frame}\label{5.2.5.}\subsubsection{5.2.5. Двойные факториалы }\frametitle{5.2.5. Двойные факториалы (1/2)}
\vspace{-2pt}
Содержание этого параграфа не имеет отношения к основной теме
раздела, т.~е. перестановкам, и включено в раздел по сходству
обозначений и с целью привести две формулы, которые иногда
используются в практических комбинаторных расчётах и требуются в
последующих разделах.

\odmkey{Двойным факториалом}{Факториал двойной (double factorial)} натурального
числа $n$ (обозначается $n!!$) называется
произведение числа $n$ и всех меньших натуральных чисел той же чётности.
\begin{example}
$10!!=10\cdot 8\cdot 6 \cdot 4 \cdot 2 = 3840$.
\end{example}
\begin{theorem}
\begin{enumerate}
\item
$(2k)!!=2^k\cdot k!$.
\item
$\displaystyle{(2k+1)!!=\frac{(2k+1)!}{2^k\cdot k!}}$.
\end{enumerate}
\end{theorem}
\begin{proof}
\proofsection{1} $(2k)!! = 2\cdot 4\cdot \cdots \cdot (2k-2) \cdot
2k  =
          (2\cdot 1)\cdot(2\cdot 2)\cdot\cdots\cdot(2\cdot(k-1))\cdot(2\cdot k)\rmb=
 \underbrace{(2\cdot2\cdot\cdots\cdot2\cdot 2)}_{k\ \text{раз}}\cdot
           (1\cdot 2\cdot \cdots \cdot \cdot(k-1)\cdot k) =2^k\cdot k!$.
\proofsection{2} Достаточно заметить, что $(2k+1)!=(2k+1)!!(2k)!!$.
\end{proof}


\end{frame}

\begin{frame}\frametitle{5.2.5. Двойные факториалы (2/2)}

По определению $0!!\Def 1$.

	
\begin{example}
\odmkey{Перестановками Стирлинга}{Перестановка (permutation)!Стирлинга (Stirling)} называются перестановки набора чисел \\$\{1, 1, 2, 2, …, k, k\}$ такие, что между любыми двумя равными числами могут находиться лишь б$\acute{\text{о}}$льшие по значению.  
\begin{example}
 $(1, 2, 3, 3, 2, 1, 4, 4)$.
\end{example}

\medskip
\begin{theorem}Число перестановок Стирлинга равно $(2k-1)!!$.
\end{theorem} 
\begin{proof}
База: при $k\! =\! 1$ существует единственная перестановка
$1, 1$,  $(2k\!-\!1)!! = 1$.
Индукционный переход: 
пусть число перестановок Стирлинга $1, 1, 2, 2, \ldots, k-1, k-1$ равно $(2(k-1)-1)!!=(2k - 3)!!$.
Так как в перестановке нет чисел б$\acute{\text{о}}$льших $k$, то элементы $k, k$ могут стоять только рядом. Мы можем их отбросить, получив перестановку с числами $1..k-1$,  учитывая, что у нас существует $2k-1$ место, куда можно поставить пару $k, k$. 
Добавив $k,k$ %, для которых в последовательности есть $2k-1$ место, 
получим $(2k - 3)!! \cdot (2k-1)= (2k-1)!!$.
\end{proof}

\end{example}
	
\end{frame}

\begin{frame}
	\label{5.2.6.}\subsubsection{5.2.6. Быстрая сортировка }\frametitle{5.2.6. Быстрая сортировка (1/6)}
	
\odmkey{Быстрая сортировка}{Сортировка (sorting)!быстрая (quick)} 
является наиболее широко применяемым и одним из самых эффективных алгоритмов сортировки. 
	
	Приведём неформальное описание идеи алгоритма быстрой сортировки.
	
	Пусть задан числовой массив, элементы которого требуется расположить в порядке возрастания.
	
	\begin{enumerate}
		
		\item {}
		В массиве выбирается элемент $x$, называемый \odmkey{опорным}{Элемент (element)!опорный (pivot)}.
		
		\item {}
		Массив делится на два подмассива относительно опорного элемента, и далее элементы меняются местами таким образом, чтобы элементы левого подмассива были меньше либо равны $x$, элементы правого
		 подмассива были больше $x$.
		
		\item {}
		Первые два шага рекурсивно применяются к левому и правому подмассивам, если они состоят более чем из одного  элемента.
		
	\end{enumerate}
	
\begin{digress}
Алгоритм быстрой сортировки придумал
сэр Чарльз Хоар во время его работы 
в МГУ в 1960 году. 
\end{digress}
	
	
\end{frame}
\begin{frame}
	\frametitle{5.2.6. Быстрая сортировка (2/6)}
\vspace{-4pt}	
	Массив \textbf{$A:$ array $[a..b]$ of $B$}, подлежащий сортировке, является глобальным. Основная работа выполняется функцией $\Quicksort$.
	
	
%\vspace{-4pt}	
	\begin{algorithmic}
		\REQUIRE числа $l$ и   $r$~--- индексы начала и конца части массива для сортировки соответственно.
		\ENSURE отсортированный по возрастанию массив $A[a..b]$.
		\IF{$l < r$}
		\STATE{ $i \Assign \Partition(l, r) $ }
		\COMMENT {\color{gray}\small разбиение массива  }
		\STATE{ $\Quicksort(l, i-1)$ }
		\COMMENT {\color{gray}\small рекурсивный вызов для левой части  }
		\STATE{ $\Quicksort(i+1, r)$ }
		\COMMENT {\color{gray}\small рекурсивный вызов для правой части  }
		\ENDIF
	\end{algorithmic}


%\vspace{-4pt}		
	Сортировка массива $A[a..b]$ выполняется вызовом функции $\Quicksort(a, b)$. 

\begin{remark}
Известно несколько способов выбора опорного
элемента $x$. Вообще говоря, можно выбирать любой элемент,
например, случайный. 
Теоретически лучше всего выбирать 
\odmkey{медиану}{Медиана (median)}, но это требует дополнительных вычислений.
Здесь используется простейший вариант: 
опорным выбирается последний элемент.  
\end{remark}	
\end{frame}

\begin{frame}
	\frametitle{5.2.6. Быстрая сортировка (3/6)}
	Функция  $\Partition$ находит опорный элемент.
	
	\begin{algorithmic}
		\REQUIRE числа $l$ и  $r$~--- индексы начала и конца части массива для поиска соответственно.
		%\COMMENT{$\Partition$}
		\ENSURE индекс $i$ опорного элемента.
		\STATE{ $x \Assign A[r] $ }
		\COMMENT{\color{gray}\small опорным является последний элемент массива }
		\STATE{ $i \Assign l-1 $ }
		\COMMENT{\color{gray}\small инициализация за пределами массива }
		\FOR{\textbf{$j$ from $l$ to $r-1$}}
		\IF{ $A[j] \le x$ }
		\STATE $i \Assign i+1$ 
		\STATE $A[i] \leftrightarrow A[j]$ 
		\ENDIF
		\ENDFOR
		\STATE{ $A[i] \leftrightarrow A[r]$} 
		%\textbf{return $i$}
		\COMMENT{\color{gray}\small опорный элемент помещается в нужное место }
		\STATE{\textbf{return $i$}}
	\end{algorithmic}
	
\end{frame}

\begin{frame}
	\frametitle{5.2.6. Быстрая сортировка (4/6)}
	
	Ключевой частью алгоритма <<быстрой сортировки>> является разбиение массива. На каждой итерации цикла выполняются четыре условия:
	
	\begin{enumerate}
		
		\item[1)] Отрезок массива
		$A[l.. i]$ содержит элементы, которые меньше или равны опорному;
		
		\item[2)] Отрезок массива
		$A[(i+1) .. (j-1)]$ содержит элементы, которые больше опорного;
		
		\item[3)] Отрезок массива
		$A[j .. (r-1)]$ содержит и те, и другие элементы;
		
		\item[4)]
		На позиции $r$ находится опорный элемент.
		
	\end{enumerate}
	
	\begin{center}
		\includegraphics[width=120mm]{5_2_6.jpg}
	\end{center}
	
\end{frame}


\begin{frame}
	\frametitle{5.2.6. Быстрая сортировка (5/6)}
	\begin{ground}
		\proofsection{Корректность разбиения}
		При инициализации цикла между $l$ и $i$ нет элементов, а, значит, первые три условия выполняются. Элемент с индексом $r$ принят опорным.
				На каждой итерации цикла необходимо рассмотреть два случая, которые определяются условием $A[j] \le A[r]$. 
		Если $A[j] > A[r]$, тогда следует добавить элемент $A[j]$ в множество элементов б$\acute{\text{о}}$льших опорного, значит, 
		следует увеличить индекс $j := j + 1$.
		Если же $A[j] \le x$, тогда элемент $A[j]$ подлежит добавлению в множество элементов меньших или равных опорному. Чтобы это сделать необходимо увеличить индекс $i := i + 1$, выполнить обмен элементов $A[i] \leftrightarrow A[j]$, а затем, чтобы сохранить множество элементов б$\acute{\text{о}}$льших опорного, следует увеличить индекс $j := j + 1$.
		
		После выполнения цикла массив разобьётся на три непересекающихся множества. Первое множество состоит из элементов меньших или равных опорному, второе~--- больших опорного, третье~--- одноэлементное множество из опорного элемента.
	\end{ground}
	
\end{frame}

\begin{frame}
	\frametitle{5.2.6. Быстрая сортировка (6/6)}
	
\begin{ground} (Продолжение)
\proofsection{Корректность сортировки}
Пусть $n$~--- длина входного массива.
База: $n = 1$~--- массив уже отсортирован. Индукционный переход: входной массив $А$ размера $n$ после вызова процедуры $\Partition$  разбивается на два подмассива размера меньше $n$ (опорный элемент не рассматривается, так как уже находится  на своём месте), которые отсортированы по индукционному предположению.
\end{ground}
	
\begin{example}
Пусть дан массив $A = \left\{ 8, 15, 24, 18, 9, 49, 20 \right\}$.
Полужирным шрифтом отмечены опорные элементы, а в прямоугольниках находятся элементы, которые уже встали на своё место.
		
		\begin{center}
			\begin{tabular}{c|ccccccc}
				Опорный элемент & $A[1]$  & $A[2]$ & $A[3]$ & $A[4]$ & $A[5]$ & $A[6]$ & $A[7]$  \\
				20    & 8 & 15 & 24 & 18 & 9 & 49 & \textbf{20}  \\
				9, 24   & 8 & 15 & 18 & \textbf{9} & \boxed{20} & 49 & \textbf{24}  \\
				8, 15, 49   & \textbf{8} & \boxed{9} & 18 & \textbf{15} & \boxed{20} & \boxed{24} & \textbf{49} \\
				18   & \boxed{8} & \boxed{9} & \boxed{15} & \textbf{18} & \boxed{20} & \boxed{24} & \boxed{49}  \\
				& \boxed{8} & \boxed{9} & \boxed{15} & \boxed{18} & \boxed{20} & \boxed{24} & \boxed{49}  \\  
			\end{tabular}
		\end{center}
		
		
	\end{example}
	
\end{frame}

\begin{frame}
	
	\label{5.2.7.}\subsubsection{5.2.7. Трудоёмкость в среднем }\frametitle{5.2.7. Трудоёмкость в среднем (1/5)}
	
\odmkey{Трудоёмкость алгоритма в среднем}{Трудоёмкость алгоритма в среднем (average labor intensity)}~--- математическое ожидание количества вычислительного ресурса (памяти или времени), используемого алгоритмом при заданном распределении вероятности появления тех или иных входных данных $P_X$ (\DMref{п.~5.1.7}{5.1.7.}).

\smallskip	
\begin{remark}
О трудоёмкости в среднем имеет смысл говорить только для тех алгоритмов, время работы которых зависит от $P_X$.
\end{remark}

\smallskip		
\begin{example}
Время работы алгоритма Уоршалла 
(\DMref{п.~1.5.2}{1.5.2.}) не зависит от 
распределения входных данных, а зависит только от размера данных.
С другой стороны, время работы алгоритма быстрой сортировки существенно
зависит от входных данных. 
Если исходный массив отсортирован в обратном порядке,
то приведённый алгоритм будет работать так же медленно,
как метод пузырька, поскольку на каждой итерации на своё место будет вставать только один опорный элемент. 
\end{example}


%%%%%%%%%%%%%%%%%%%%%%%%%%%%%%%%%%%%%%%%%%%%%%%%%%%%%%%%%%
% Не понял. Если понимаете, то верните. ФАН
%\begin{example}
%	Сортировка пузырьком закончит свою работу всего за один проход, если на вход придёт уже отсортированный массив.
%\end{example}

\medskip
Рассмотрим получение оценок в среднем на примере 	операции поиска записи по ключу с помощью следующего алгоритма, известного как алгоритм
\odmkey{}{Алгоритм (algorithm)!поиска (search)!бинарного (binary)}
\odmkey{бинарного}{Поиск (search)!бинарный (binary)} (или
\odmkey{двоичного}{Поиск (search)!двоичный (binary)}) поиска.
	
\end{frame}

\begin{frame}
\frametitle{5.2.7. Трудоёмкость в среднем (2/5)}
\begin{algorithmic}
\REQUIRE упорядоченный массив
\textbf{$A\Is$ array $[1..n]$ of struct $\{k\Is key; i\Is info\}$};
ключ  $a\Is key$.
\ENSURE
индекс записи с искомым ключом $a$ в массиве $A$ или 0,
если записи с~таким ключом нет.
\STATE $b \Assign  1$ 
\COMMENT{\color{gray}\small начальный индекс части массива для поиска }
\STATE $e \Assign  n$ 
\COMMENT{\color{gray}\small  конечный индекс части массива для поиска }
\WHILE{$b\le e$}
\STATE $c \Assign  (b+e)\DIV 2$ 
\COMMENT{\color{gray}\small  индекс проверяемого элемента }
\IF{$A[c].k > a$}
\STATE $e \Assign  c-1$ 
\COMMENT{\color{gray}\small  продолжаем поиск в первой половине }
\ELSIF{$A[c].k < a$}
\STATE $b \Assign  c+1$ 
\COMMENT{\color{gray}\small  продолжаем поиск во второй половине }
\ELSE
\STATE \textbf{return $c$} 
\COMMENT{\color{gray}\small  нашли искомый ключ }
\ENDIF
\ENDWHILE
\STATE \textbf{return $0$} \COMMENT{\color{gray}\small искомого ключа нет в массиве }

\end{algorithmic}
\end{frame}


\begin{frame}
\frametitle{5.2.7. Трудоёмкость в среднем (3/5)}
\begin{ground}
Достаточно заметить,
что на каждом шаге основного цикла искомый элемент
массива (если он есть) находится между
(включительно)
элементами с индексами $b$ и $e$. Поскольку
диапазон поиска на каждом шаге уменьшается вдвое,
общая трудоёмкость не превосходит $\log n$.
\end{ground}

\smallskip
\begin{digress}
Упорядоченные списки проигрывают упорядоченным массивам при реализации
операций, поскольку для списков нет аналога 
алгоритма двоичного поиска.
Это частное наблюдение является проявлением
весьма общего факта, состоящего в том,
что время доступа к элементу массива не зависит от
номера элемента и от размера массива.
Говорят, что массив является 
структурой данных \odmkey{прямого доступа}{Прямой доступ (direct access)},
это характеристическое свойство массива.
Во всех же <<динамических>> структурах данных, в частности, в списках,
время доступа \em{зависит} от положения элемента и от размера структуры данных. 
\end{digress}

\end{frame}

\begin{frame}
	
	\frametitle{5.2.7. Трудоёмкость в среднем (4/5)}

Рассмотрим оценки вычислительной сложности для алгоритма бинарного поиска на входном массиве длины $n$.

\medskip
\em{Оценка снизу}. В лучшем случае искомый элемент находится на позиции с индексом $\frac{n}{2}$. Значит потребуется только одно сравнение.

\medskip	
\em{Оценка в среднем}. Для упрощения вычислений положим $n = 2^k - 1$, 
$k \in \mathbb{N}$. Если искомый элемент имеет индекс $i = \frac{n}{2}$, то понадобится одно сравнение. Если же $i = \frac{n}{4}$ или $i = \frac{3n}{4}$, то понадобится уже два сравнения. Таким образом, общее количество сравнений равно $\sum\limits_{i=1}^{k}{i 2^{i-1}}$. Так как искомый элемент может находиться на любой из позиций с вероятностью $\frac{1}{n}$, среднее количество сравнений равно

$$\frac{1}{2^k - 1} \sum\limits_{i=1}^{k}{i 2^{i-1}} = \sum\limits_{i=1}^{k}{\frac{i 2^{i}}{2(2^k - 1)}} = \frac{2((k-1)2^k+1)}{2(2^k-1)}.$$

	
\end{frame}
\begin{frame}
	\frametitle{5.2.7. Трудоёмкость в среднем (5/5)}
Сокращая дробь и переходя к переменной $n$, имеем: 


$$\frac{(\log(n+1)-1)(n+1) + 1}{n} = \log(n+1) - 1 + \frac{ \log(n+1) }{n}.$$

\medskip	
	\em{Оценка сверху}. Основываясь на рассуждениях для оценки в среднем, можно увидеть, что в худшем случае количество сравнений, выполненных алгоритмом бинарного поиска, равно $\log(n+1)$, при $n=2^k-1, k \in \mathbb{N}$.
	
\medskip	
	
	Оценки в среднем и худшем случаях для бинарного поиска имеют величины порядка $O(\log n)$. Таким образом, алгоритм бинарного поиска хорош тем, что для него нет плохих случаев.
	
	\begin{remark}
		Распределение входных данных, как и трудоёмкость в среднем, зависит от конкретной задачи, которую необходимо исследовать.
	\end{remark}
	
\end{frame}

\begin{frame}
	\label{5.2.8.}\subsubsection{5.2.8. Гармонические числа }\frametitle{5.2.8. Гармонические числа (1/2)}
	\odmkey{Гармонические числа}{Число (number)!гармоническое (harmonic)} $H(n)$ определяются следующим образом:
	$$H(0)\Def 0,\quad H(n)\Def H(n-1)+\frac{1}{n}.$$
	Применив определение $n-1$ раз, получим ($n\ge 1$): 
	$$H(n)=H(n-1)+\frac{1}{n}=H(n-2)+\frac{1}{n}+\frac{1}{n-1}= \dots =
	\sum_{k=1}^n \frac{1}{k}$$
	
\begin{remark}
Из курса математического анализа известна 
формула для приближенного вычисления гармонических
чисел, которая в упрощённом виде записывается
следующим образом:
$$H_n = \ln n + \gamma + \frac{1}{2n} + O\left(\frac{1}{n^2} \right), \lim_{n\rightarrow \infty} (H_n - \ln n)  = \gamma \approx  0.57722$$
Число $\gamma$ называется \odmkey{постоянной Эйлера}{Постоянная (constant)!Эйлера (Euler)}.

Гармонические числа могут использоваться для вывода оценки в среднем для многих алгоритмов.
\end{remark}
\end{frame}

\begin{frame}
	\frametitle{5.2.8. Гармонические числа (2/2)}
По аналогии можно ввести понятие 
	\odmkey{обобщённых гармонических чисел}{Числа (numbers)!обобщенные гармонические (generalized harmonic)} порядка $m$:
	$$
	H(m,0)\Def0, \quad H(m,n)\Def H(m,n-1)+1/n^m,\quad 
	H(m,n)= \sum_{k=1}^n \frac{1}{k^m}.
	$$ 


\begin{digress}
C обобщёнными гармоническими числами 
связаны многие
важные и трудные открытые задачи.
В частности, 
\odmkey{дзета-функция Римана}{Дзета-функция Римана \\ (Riemann zeta-function)}
$\zeta(m)=\sum\limits_{k=1}^{\infty} \frac{1}{k^m}.$

\odmkey{Гипотеза Римана}{Гипотеза (hypothesis)!Римана (Riemann)}
о распределении нетривиальных нулей дзета-функции
тесно связана с распределением простых чисел,
и поэтому сложность многих алгоритмов над
целыми числами доказана в предположении
истинности гипотезы Римана.
Если гипотеза Римана окажется неверна,
то это обстоятельство радикально 
повлияет на криптографическую практику.
Гипотеза Римана входит
в число семи математических
<<проблем тысячелетия>>.  
\end{digress}	
\end{frame}

\begin{frame}
	\label{5.2.9.}\subsubsection{5.2.9. Оценка в среднем для быстрой сортировки }\frametitle{5.2.9. Оценка в среднем для быстрой сортировки (1/3)}
	
	\begin{theorem}
Время работы алгоритма <<быстрой сортировки>> в среднем $O(n \log n)$.
	\end{theorem}
	
	\begin{proof}
		Пусть $Q_n$~--- среднее количество сравнений, выполняемых алгоритмом, когда входные данные~--- это массив размера $n$. Положим $n \geq 2$. Имеем	$Q_0 = Q_1 = 0$. Известно, что после выбора опорного элемента необходимо провести $n-1$ сравнение со всеми оставшимися элементами. Опорный элемент разобьёт массив на две части. Первая часть может быть любого размера от $0$ до $n-1$. Пусть первая часть имеет размер $k$, значит вторая часть имеет размер $n-1-k$. 
Более того, $k \in \left\{0, \ldots , n-1 \right\}$ 		равновероятны. Тогда имеем следующее равенство для математического ожидания 
(\DMref{п.~5.1.7}{5.1.7.}):
		
$$Q_n = (n-1) + \frac{1}{n} \sum_{k=0}^{n-1}\left(Q_k + Q_{n-1-k}\right).$$
		
	\end{proof} 
	
\end{frame}

\begin{frame}
	\frametitle{5.2.9. Оценка в среднем для быстрой сортировки (2/3)}
\begin{proof} (Продолжение)	
	Заметим, что: 
	$\sum\limits_{k=0}^{n-1}{Q_{n-1-k}} = \sum\limits_{k=0}^{n-1}{Q_k}$, тогда:
	
	$$Q_n = (n-1)+\frac{2}{n} \sum_{k=0}^{n-1}{Q_k}.$$
	
Домножим на $n$, имеем: $n Q_n = n(n-1) + 2\sum\limits_{k=0}^{n-1}{Q_k}$.
	Заменим $n := n - 1$, что возможно, поскольку $ (n \geq 2 )$.
Получаем	$(n-1) Q_{n-1} = (n-1)(n-2) + 2 \sum\limits_{k=0}^{n-2}{Q_k}$.
	Вычтем выражения, получим:
	$n Q_n - (n-1)Q_{n-1} = n(n-1) - (n-1)(n-2) + 2Q_{n-1}$,\\
	$ nQ_n = 2(n-1)+(n+1)Q_{n-1} $.
	Разделим на $n(n+1)$:
	
	$$\frac{Q_n}{n+1} = 2\frac{n-1}{n(n+1)} + \frac{Q_{n-1}}{n}.$$
\end{proof}	
\end{frame}


\begin{frame}
	\frametitle{5.2.9. Оценка в среднем для быстрой сортировки (3/3)}
\begin{proof} (Окончание)	
	Введём новую переменную:	
	$F_n \Assign \frac{Q_n}{n+1}$.
	Имеем   
	$$F_n = 2 \frac{n-1}{n(n+1)} + F_{n-1}, \quad F_n = 2 \sum\limits_{k=1}^{n}{\frac{k-1}{k(k+1)}}.$$
	Разобьём на простейшие дроби: $\frac{k-1}{k(k+1)} = \frac{2}{1+k} - \frac{1}{k}$,
		$F_n = 4 \sum\limits_{k=1}^{n}{\frac{1}{1+k}} - 2 \sum\limits_{k=1}^{n}{\frac{1}{k}}$ или \\$F_n = 4(H_{n+1} - 1) - 2H_n = 4H_{n+1} - 2H_n - 4$.
	
	Получаем выражение в терминах гармонических чисел:
	
	$$Q_n = (n+1)(4H_{n+1} - 2H_n - 4).$$
	
	После разложения гармонических чисел в ряд получаем:
	
	$$Q_n = 2n(\ln n + \gamma - 2) + 2\ln n + 2\gamma + 1 + O\left(\frac{1}{n} \right).$$
\end{proof}	
\end{frame}

\begin{frame}
\frametitle{Онтология комбинаторного анализа}
\begin{center}
\includegraphics[width=12cm]{ODM5_2.png}
\end{center}
\end{frame}

\begin{frame}
\frametitle{5.3. Биномиальные коэффициенты (1/2)}
\subsection{5.3. Биномиальные коэффициенты}\label{5.3.}


Число сочетаний $C(m,n)$~--- это число различных $n$-элементных
подмножеств  $m$-элементного множества. 
Числа $C(m,n)$ обладают целым рядом
свойств, рассматриваемых в~этом разделе, которые
оказываются очень полезными при решении различных комбинаторных
задач.

\bigskip
\begin{tabular}{p{11mm}|p{45mm}|p{88mm}}
  № & Название параграфа
  & Ключевые термины и обозначения \\
  \hline
\DMref{5.3.1.}{5.3.1.} & Свойства биномиальных коэффициентов &
{$C(m,n)=m!/(n!(m-n)!)$; тождество Коши}
\\  
\DMref{5.3.2.}{5.3.2.} & Бином Ньютона &
{ Биномиальные коэффициенты}\\  
\DMref{5.3.3.}{5.3.3.} & Треугольник Паскаля &
{ }\\  
\DMref{5.3.4.}{5.3.4.} & Применение рекуррентного соотношения &
{  }\\  
\DMref{5.3.5.}{5.3.5.} & Генерация подмножеств &
{ Переборные задачи, пространство поиска}\\  
\DMref{5.3.6.}{5.3.6.} & Расширенные биномиальные коэффициенты &
{ Расширенный треугольник Паскаля}\\  
\end{tabular}  

\end{frame}

\begin{frame}
\frametitle{5.3. Биномиальные коэффициенты (2/2)}

\begin{tabular}{p{11mm}|p{45mm}|p{88mm}}
  № & Название параграфа
  & Ключевые термины и обозначения \\
  \hline
\DMref{5.3.7.}{5.3.7.} & Мультимножества и последовательности  &
{ Состав последовательности}\\  
\DMref{5.3.8.}{5.3.8.} & Мультиномиальные коэффициенты &
{ Постоянная Эйлера}\\  
\DMref{5.3.9.}{5.3.9.} & Биномиальная система счисления  &
{ }\\  
\end{tabular}  

\end{frame}



\begin{frame}
\label{5.3.1.}\subsubsection{5.3.1. Свойства биномиальных коэффициентов }\frametitle{5.3.1. Свойства биномиальных коэффициентов (1/5)}

Основная формула для числа сочетаний $$ C(m,n)=\frac{m!}{n!(m-n)!}
$$ позволяет получить следующие элементарные тождества.
\begin{theorem} {\NB{ [1] }}
\begin{enumerate}
\item $C(m,n) = C(m,m-n)$.%
\item $C(m,n) = C(m-1,n) + C(m-1,n-1)$.%
\item $C(n,i)C(i,m)=C(n,m)C(n-m,i-m)$.
\item $C(m,n) = \frac{m}{n}C(m-1,n-1)$.%
\end{enumerate}
\end{theorem}

\begin{proof}
\proofsection{1}
$\displaystyle{C(m,m-n)=\frac{m!}{(m-n)!\,\left(m-(m-n)\right)!}=\frac{m!}{(m-n)!\,n!}=C(m,n).}$
\end{proof}
\end{frame}

\begin{frame}
\frametitle{5.3.1. Свойства биномиальных коэффициентов (2/5)}
\begin{proof}
\proofsection{2}
$\displaystyle{C(m-1,n)+C(m-1,n-1)=}$\\[4pt]
$\displaystyle{= \frac{(m-1)!}{n!(m-n-1)!}+
\frac{(m-1)!}{(n-1)!\left(m-1-(n-1)\right)!}=}$\\
$\displaystyle{=
\frac{(m-1)!}{n(n-1)!\,(m-n-1)!}+\frac{(m-1)!}{(n-1)!\,(m-n)(m-n-1)!}=}$\\
$\displaystyle{=
\frac{(m-n)(m-1)!+n(m-1)!}{n(n-1)!\,(m-n)(m-n-1)!}=
\frac{(m-n+n)(m-1)!}{n!\,(m-n)!}}$ 
$\displaystyle{=\frac{m!}{n!\,(m-n)!}=C(m,n).}$

\proofsection{3}
$\displaystyle{C(n,i)C(i,m)= \frac{n!}{i!(n-i)!}\frac{i!}{m!(i-m)!}=
\frac{n!}{m!(i-m)!(n-i)!}=}$\\
$\displaystyle{=\frac{n!(n-m)!}{m!(i-m)!(n-i)!(n-m)!}=
\frac{n!}{m!(n-m)!}\frac{(n-m)!}{(i-m)!(n-i)!}=}$\\
$\displaystyle{=C(n,m)C(n-m,i-m).}$

\proofsection{4}
$\displaystyle{C(m,n)  
=\frac{m!}{n!(m-n)!} 
=\frac{m(m-1)!}{n(n-1)!((m-1)-(n-1))!} =
\frac{m}{n}C(m-1,n-1).}$

\end{proof}
\end{frame}

\begin{frame}
\frametitle{5.3.1. Свойства биномиальных коэффициентов (3/5)}

%\vspace{-4pt}
Биномиальные коэффициенты обладают целым рядом
других замечательных свойств.

%\addvspace{-5pt} \setcounter{theorem}{0}
\begin{theorem} {\NB{ [2] }}
$\sum\limits^m_{n=0} nC(m,n)  = m2^{m-1}$.
\end{theorem}

\begin{proof}
Рассмотрим следующую последовательность, составленную из чисел \\$1,\dots, m$.
Сначала выписаны все подмножества длины $0$, потом все подмножества длины $1$ и т.~д.
Имеется $C(m,n)$ подмножеств мощности $n$, и каждое из них имеет длину $n$,
таким образом, всего в этой последовательности $\sum^m_{n=0}nC(m,n)$
чисел. С другой стороны, каждое число $x$ входит в эту
последовательность $2^{\vert\{1,\dots, m\}\backslash\{x\}\vert}=2^{m-1}$ раз,
а всего чисел $m$.
\end{proof}

\begin{theorem} {\NB{ [3] }}
$C(m+n,k) =  \sum\limits^k_{i=0} C(m,i)C(n,k-i)$.
\end{theorem}

\begin{proof}
$C(n+m,k)$~--- это число способов выбрать $k$ шаров из
$m+n$ шаров. Шары можно выбирать в два приема:
сначала выбрать $i$ шаров из первых $m$ шаров,
а затем выбрать недостающие $k-i$ шаров из оставшихся $n$ шаров.
Отсюда общее число способов выбрать $k$ шаров составляет
$\sum^k_{i=0} C(m,i)C(n,k-i)$.
\end{proof}
\begin{remark}
Последнее свойство известно как \odmkey{тождество Коши}{Формула
(formula)!Коши (Augustin Louis Cauchy)}.
%\footnote{\;Огюстен Луи
%Коши (1789--1857).}.
\end{remark}

\end{frame}


\begin{frame}
\frametitle{5.3.1. Свойства биномиальных коэффициентов (4/5)}



\begin{theorem} \NB{[4]}
$\sum\limits^m_{n=0} (C(m,n))^2  = C(2m, m)$.
\end{theorem}

\begin{proof}
Рассмотрим множество из $2m$ объектов. Число способов выбрать из него $m$ объектов равно $C(2m, m)$. Разобьём это множество на 2 равномощных подмножества. Теперь применим тождество Коши при $n = m, k = m$. Тогда общее число способов выбрать $m$ предметов будет $\sum\limits^m_{n=0} (C(m,n))^2$ в силу того, что $C(m,n) = C(m,m-n)$.
\end{proof}


\begin{theorem} \NB{[5]}
$C(m,n) = \sum\limits^n_{k=0} C(m-(k+1),n-k))$.
\end{theorem}

\begin{proof}
Рассмотрим множество из $m$ элементов
и представление подмножеств
этого множества с помощью 
битовых шкал, \DMref{п.~1.3.1}{1.3.1.}. 
Число $n$-элементных подмножеств равно $C(m, n)$. При этом число битовых шкал,
в которых первый слева ноль
стоит в разряде $k+1$, равно 
$C(m -(k+1), n-k)$, поскольку в таких шкалах
все разряды определены до разряда $k+1$ включительно, и остаётся выбрать $n-k$
разрядов для размещения единиц
среди $m-(k+1)$ неопределённых разрядов. 
В битовой шкале для $n$-элемент\-ного 
подмножества ровно $m-n$ нулей,
значит первый слева ноль может находиться
в разряде с номером 
из диапазона $1..(n+1)$.  
\end{proof}


\end{frame}


\begin{frame}
\frametitle{5.3.1. Свойства биномиальных коэффициентов (5/5)}

\begin{digress}
Указанные свойства биномиальных коэффициентов нетрудно доказать,
проводя алгебраические выкладки обычным образом. Например, теорему 2 
можно получить так:
\begin{align*}
\sum_{n=0}^m nC(m,n) &= \sum_{n=1}^m \frac{n m!}{(m-n)!n!}
          =  \sum_{n=1}^m \frac{m (m-1)!}{(m-n)!(n-1)!} =\\
          &=  m \sum_{n=0}^{m-1} \frac{(m-1)!}{(m-n-1)!n!}
          =  m \sum_{n=0}^{m-1} C(m-1,n) = m2^{m-1}.
\end{align*}
В доказательстве использованы рассуждения <<в комбинаторном
духе>>, чтобы проиллюстрировать <<неалгебраические>> способы
решения комбинаторных задач.

\smallskip
Иногда наблюдается и обратный эффект:
комбинаторные рассуждения оказываются
короче даже 
простейших алгебраических выкладок.
 Например, теорему 1 пункт 2 
можно получить следующим рассуждением.
Зафиксируем произвольный элемент $x$.
Тогда количество способов выбрать $n$ элементов из $m$ 
элементов, то есть число $C(m,n)$,
складывается
из количества способов выбрать $n$ элементов, включающих
элемент $x$, то есть числа $C(m-1,n-1)$,
 и количества способов выбрать $n$ элементов не включающих элемент $x$,
то есть $C(m-1,n)$.
\end{digress}

\end{frame}



\begin{frame}
\label{5.3.2.}\subsubsection{5.3.2. Бином Ньютона }\frametitle{5.3.2. Бином Ньютона (1/2)}

Числа сочетаний $C(m,n)$ называются также \odmkey{биномиальными
коэффициентами}{Коэффициент (coefficient)!биномиальный
(binomial)}. Смысл этого названия устанавливается следующей
теоремой, известной также как формула \odmkey{бинома
Ньютона}{Бином Ньютона (binomial theorem)}.
%\footnote{\;Исаак Ньютон
%(1643--1727).}.



\begin{theorem}
$ \displaystyle{(x+y)^m   =  \sum^m_{n=0} C(m,n)x^n y^{m-n}}$.
\end{theorem}


%{\renewcommand{\baselinestretch}{0.94}\selectfont
%\begin{proof}
{\bf\color{brown} Доказательство.}
По индукции. База, $m=1$:
$(x+y)^1=x+y=1x^1y^0+1x^0y^1=$\\ $=C(1,1)x^1y^0+C(1,0)x^0y^1=\sum\limits^1_{n=0}C(1,n)x^ny^{1-n}$.

Индукционный переход:
\begin{align*}
(x+y)^m&=(x+y)(x+y)^{m-1} =
(x+y)\sum^{m-1}_{n=0}C(m-1,n)x^ny^{m-n-1}=\\
&=
\sum^{m-1}_{n=0}xC(m-1,n)x^ny^{m-n-1}+\sum^{m-1}_{n=0}yC(m-1,n)x^ny^{m-n-1}=\\
\end{align*}
\end{frame}

\begin{frame}
\frametitle{5.3.2. Бином Ньютона (2/2)}

\vspace{-12pt}
\begin{align*}
&=
\sum^{m-1}_{n=0}C(m-1,n)x^{n+1}y^{m-n-1}+\sum^{m-1}_{n=0}C(m-1,n)x^ny^{m-n}=\\
&=  \sum_{n=1}^m C(m-1,n-1)x^n y^{m-n}
          + \sum_{n=0}^{m-1} C(m-1,n)x^n y^{m-n}=\\
&= C(m-1,0)x^0 y^m
          + \sum_{n=1}^{m-1} (C(m-1,n-1)+C(m-1,n))x^n y^{m-n}\:
          + \\
&+ C(m-1,m-1)x^m y^0
= \sum^m_{n=0}C(m,n)x^ny^{m-n}.
\end{align*}

\vspace{-10pt}
%\addvspace{-5pt} \setcounter{corollary}{0}
\begin{corollary}
$\sum^m_{n=0} C(m,n)  =  2^m$.
\end{corollary}

%%{\bf Доказательство.}
%$2^m=(1+1)^m=\sum\limits^n_{n=0}C(m,n)1^n1^{m-n}=\sum\limits^m_{n=0}C(m,n)$

\begin{proof}
Достаточно взять $x=y=1$.
\end{proof}

\begin{corollary}
$\sum^m_{n=0} {(-1)}^{n} C(m,n)  = 0$.
\end{corollary}

%{\bf Д-во.}
%$0=(-1+1)^m=
%\sum\limits^m_{n=0}C(m,n)(-1)^n 1^{m-n}
%=\sum\limits^m_{n=0}(-1)^n C(m,n)$.
%

\begin{proof}
Достаточно взять $x=-1, y=1$.
\end{proof}

\end{frame}


\begin{frame}\label{5.3.3.}\subsubsection{5.3.3. Треугольник Паскаля }\frametitle{5.3.3. Треугольник Паскаля (1/7)}

Из второй формулы теоремы \DMref{п.~5.3.1}{5.3.1.} вытекает эффективный способ
рекуррентного вычисления значений биномиальных коэффициентов,
который можно представить в графической форме, известной как
\odmkey{треугольник Паскаля}{Треугольник Паскаля (Pascal's triangle)}
%\footnote{\;Блез Паскаль (1623--1662).}.
%\vspace{-0.5\baselineskip}
\begin{center}
\begin{tabular}{ccccccccccc}
   &   &   &   &   & 1 &   &   &   &   &    \\
   &   &   &   & 1 &   & 1 &   &   &   &    \\
   &   &   & 1 &   & 2 &   & 1 &   &   &    \\
   &   & 1 &   & 3 &   & 3 &   & 1 &   &    \\
   & 1 &   & 4 &   & 6 &   & 4 &   & 1 &    \\
 . &   & . &   & . &   & . &   & . &   &  . \\
\end{tabular}
\end{center}


В этом равнобедренном треугольнике каждое число
(кроме единиц на боковых сторонах) является суммой двух чисел,
стоящих над ним. Число сочетаний $C(m,n)$ находится в $(m+1)$-м
ряду на $(n+1)$-м месте.



\end{frame}


\begin{frame}\frametitle{5.3.3. Треугольник Паскаля (2/7)}

Рассмотрим вопрос реализации процесса вычислений биномиальных коэффициентов. 
Если просто 
по рекуррентной формуле $C(m,n) = C(m-1,n) + C(m-1,n-1)$ составить рекурсивную
процедуру, 
то можно столкнуться с переполнением стека вызовов из-за большой глубины рекурсии.
Кроме того, прямая рекурсивная реализация сугубо неэффективна, 
поскольку при рекурсивной реализации один и тот же биномиальный коэффициент 
будет вычисляться несколько раз. Например, 
$$C(5,3) = C(4,3)+C(4,2) = C(3,3) + \underline{C(3,2)} + 
\underline{C(3,2)} + C(3,1) = \dots.$$  
Для случая рекурсивного вычисления $C(m, n)$, получаем $2^n$  сложений, 
так как на каждом шаге рекурсии количество вычислений увеличивается в два раза, 
а всего шагов $n$. А для нерекурсивного метода имеем $n$ сложений, 
то есть рекурсивный метод экспоненциально хуже нерекурсивного.

\end{frame}


\begin{frame}\frametitle{5.3.3. Треугольник Паскаля (3/7)}

Замечательные свойства треугольника Паскаля позволяют организовать 
процесс вычислений намного эффективнее. 
Представим треугольник в другой форме, повернув его на 45 градусов против часовой стрелки. 
Получится бесконечная матрица, в которой верхняя строка и левый столбец 
содержат единицы, а каждый внутренний элемент является суммой двух элементов: соседнего слева и соседнего сверху. 
\begin{center}
\begin{tabular}{cccccc}
1 & 1 & 1 & 1 & 1 & 1 \\
1 & 2 & 3 & 4 & 5 & \\
1 & 3 & 6 & 10 &  & \\
1 & 4 & 10  &  &  & \\
1 & 5  &  &  &  & \\ 
1 &  &  &  &  & \\
\end{tabular}
\end{center}

\end{frame}

\begin{frame}\frametitle{5.3.3. Треугольник Паскаля (4/7)}

%\vspace{-4pt}
\begin{algorithmic}
\REQUIRE{ целые неотрицательные числа $m$ и $n$}
\ENSURE{ число сочетаний $C(m, n)$}
\STATE{\textbf{if $n = 1 \Or m = n + 1$ then return $m$ end if}}  
\COMMENT{\color{gray}\small граничные случаи }
\STATE{\textbf{if $n > m$ then return $0$ end if}} 
\COMMENT{\color{gray}\small по определению}
\STATE{\textbf{if  $n = m \Or n = 0$ then return $1$ end if}} 
\COMMENT{\color{gray}\small граничные случаи  }
\STATE{\textbf{if $n > m - n$ then $n \Assign m - n$ end if}}
\COMMENT{\color{gray}\small в силу симметрии } 
\STATE{\textbf{$A:$ array$[0..n]$ of int}}
\COMMENT{\color{gray}\small рабочий массив}
\FOR{\textbf{$i$ from $0$ to $n$}}
       \STATE{$A[i] \Assign 1$}
       \COMMENT{\color{gray}\small первая строка треугольника Паскаля }
\ENDFOR
\FOR{\textbf{$i$ from $1$ to $m-n$}}
   \FOR{\textbf{$j$ from $1$ to $n$}}
        \STATE{$A[j] \Assign A[j] + A[j - 1]$}
        \COMMENT{\color{gray}\small согласно  рекуррентной формуле }
\ENDFOR
\ENDFOR
\STATE{\textbf{return $A[n]$}}
\end{algorithmic}
\end{frame}


\begin{frame}\frametitle{5.3.3. Треугольник Паскаля (5/7)}
\begin{ground}
Сначала в алгоритме проверяются граничные случаи, для которых значения известны заранее (\DMref{п.~5.1.5}{5.1.5.}).  
Значение $C(m, n)$ находится в $(m-n+1)$-ой строке  на $(n + 1)$-м месте. 
Весь треугольник хранить в памяти нет нужды, 
достаточно хранить только одну строку. 
Первая строка содержит все единицы, а каждая следующая получается 
перевычислением слева направо по формуле $A[j] := A[j] + A[j - 1]$. 
При этом в правой части оператора $A[j]$~---
 это <<старое>>  значение, то есть элемент из <<верхней>> строки, 
 а $A[j-1]$~--- только что вычисленное значение слева.
 Перевычислив таким образом  массив $m-n$ раз, получаем $C(m,n)=A[n]$. 
\end{ground}

\begin{example}
Рассмотрим вычисление $C(5,2)$. Здесь $m=5$, $n=2$, $m-n=3$.
\vspace{-6pt}
\begin{center}
\begin{tabular}{c|ccc}
$i$ &   & $A$ &  \\
\hline
    & 1 & 1 & 1  \\
1   & 1 & 2 & 3  \\
2   & 1 & 3 & 6 \\
3   & 1 & 4 & 10  \\ 
\end{tabular}
\end{center}
\end{example}
\end{frame}


\begin{frame}\frametitle{5.3.3. Треугольник Паскаля (6/7)}

\begin{remark}
В силу симметрии треугольника Паскаля можно было бы использовать 
массив большей длины $m-n+1$ и сократить количество повторений основного цикла по $i$.
Однако это не даёт сокращения вычислений: необходимое количество сложений остаётся равным $n(m-n)$. 
\end{remark}


\end{frame}


\begin{frame}\frametitle{5.3.3. Треугольник Паскаля (7/7)}
\begin{digress} 
Для доказательства применимости и полезности приведённого алгоритма оценим максимальные
 значения числа сочетаний, которые можно вычислить по основной формуле 
  и по алгоритму. 
  Предположим, что используется 32-разрядная арифметика, 
  беззнаковые целые числа, максимально представимое число $2^{32}-1$. 
  При фиксированном $m$ биномиальный коэффициент $C(m, n)$ 
  имеет максимум для $n = m / 2$. Найдём максимально допустимые 
  $C(m, n)$ и $m$, для которых возможно вычисление. 
  Для формулы это $C(12, 6) = 924$, поскольку вычислить $m!$ уже не представляется возможным при $m > 12$ для используемого представления чисел в компьютере. 
  Использование формулы 
  $$C(m,n)= ((m-n+1)\cdot(m-n+2)\cdot\dots\cdot m)/(m-n)!$$ 
  не приведёт к существенному расширению диапазона вычисляемых значений. 
  A для приведённого алгоритма максимальное $m = 34$, $C(34, 17) = 2333606220$. 
  Если же варьировать $n$  от $0$ до $m$, 
  то получим максимальное значение $C(222, 5) = 4294249674$, 
  что близко к $2^{32}$, то есть, к максимальному представимому числу.  
  Значит, алгоритм применим для гораздо более широкого набора чисел $m$ и $n$ 
  по сравнению с замкнутой формулой из параграфа \DMref{5.1.5}{5.1.5.}, 
  что делает его весьма полезным.
\end{digress}   
\end{frame}


\begin{frame}
\label{5.3.4.}\subsubsection{5.3.4. Применение рекуррентного соотношения }\frametitle{5.3.4. Применение рекуррентного соотношения (1/5)}

Использование 
\odmkey{рекуррентных соотношений}{Рекуррентное соотношение (recurrence ratio)}
является одним из основных приёмов организации эффективных вычислений как в комбинаторике,
так и в других областях математики.


\begin{remark}
Однопараметрическое 
рекуррентное соотношение
$a_i=f(a_{i-1})$ позволяет построить итеративный алгоритм вычисления чисел $a_n$
без использования дополнительной памяти.
\end{remark}
 
 
Построим однопараметрическое
рекуррентное соотношение
для вычисления биномиальных коэффициентов. 


Положим
$a_1\Assign m-n+1$
и положим
$\displaystyle{a_i \Assign  \frac{m-n+i}{i} a_{i-1}}$, $i=2, 3, ..., n$.

\begin{lemma} В указанных обозначениях
$a_i=C(m-n+i,i)$.
\end{lemma}

\begin{proof} Действительно,
$\displaystyle{a_1 = m-n+1 =  \frac{m-n+1}{1} = C(m-n+1,1)}$,\\
$\displaystyle{a_2 = \frac{m-n+2}{2} a_1 = \frac{m-n+2}{2} C(m-n+1,1) = C(m-n+2,2)}$ 
по пункту 4 теоремы \DMref{п.~5.3.1}{5.3.1.},
и так далее вплоть до
$\displaystyle{a_n = \frac{m}{n} a_{n-1} = \frac{m}{n} C(m-1,n-1) = C(m,n)}$. 
\end{proof}

\end{frame}


\begin{frame}
\frametitle{5.3.4. Применение рекуррентного соотношения (2/5)}
\begin{remark}
Число $a_i$ при любом допустимом $i$~--- целое, поскольку является числом сочетаний. Значит, мы можем в процессе вычислений смело осуществлять целочисленное деление произведения величин
$ (m-n+i) a_{i-1}$ на $i$.
\end{remark}

%Из изложенного вытекает алгоритм
%прямого итеративного вычисления числа сочетаний.  
\begin{algorithmic}
\REQUIRE{ целые неотрицательные числа $m$ и $n$}
\ENSURE{ число сочетаний $C(m, n)$}
\STATE{\textbf{if $n = 1 \Or m = n + 1$ then return $m$ end if}}  
\COMMENT{\color{gray}\small граничные случаи }
\STATE{\textbf{if $n > m$ then return $0$ end if}} 
\COMMENT{\color{gray}\small по определению}
\STATE{\textbf{if  $n = m \Or n = 0$ then return $1$ end if}} 
\COMMENT{\color{gray}\small граничные случаи  }
\STATE{\textbf{if $n > m - n$ then $n \Assign m - n$ end if}}
\COMMENT{\color{gray}\small в силу симметрии } 
\STATE{$a \Assign m - n + 1$}
\COMMENT{\color{gray}\small первый элемент рекуррентной схемы}
\STATE{$k \Assign m - n + 2 $}
\COMMENT{\color{gray}\small для реализации рекуррентной схемы}
\FOR{\textbf{$i$ from $2$ to $n$}}
       \STATE{$a \Assign a \DIV i * k + a \MOD i * k \DIV i$}
       \COMMENT{\color{gray}\small рекуррентная схема }
        \STATE{$k \Assign k + 1$}
\ENDFOR
\STATE{\textbf{return $a$}}
\COMMENT{\color{gray}\small последний элемент рекуррентной схемы: $C(m,n)$ }
\end{algorithmic}
\end{frame}


\begin{frame}
\frametitle{5.3.4. Применение рекуррентного соотношения (3/5)}

\begin{ground}
 Формула 
 $a \Assign a \DIV i * k + a \MOD i * k \DIV i$
 получается из следующих соображений: 
 представляем число $a_{i-1} / i$ в виде неполного частного и остатка $a_{i-1} = ti+r$, раскрываем скобки и получаем: 
$\displaystyle{a_i=  \frac{(ti+r)(m-n+i)}{i} = 
t(m-n+i)+\frac{r(m-n+i)}{i}}$.
В силу того, что числа 
 $a_i$ целые, формула даёт правильный результат вне зависимости от того,
 какая именно из величин произведения 
$ (m-n+i) a_{i-1}$ делится на $i$ нацело.
\end{ground}

\begin{example}
 Рассмотрим нахождение биномиального коэффициента $C(6,3)$.
 
Начальный элемент $a_1=m-n+1=4$, $k=5$,  $i=2$.

Первая итерация: $a_2=4 \DIV 2 * 5 + 4 \MOD 2  * 5 \DIV 2= 2*5 +0*5 \DIV 2 = 10$. 

Вторая итерация: $a_3=10 \DIV 3*6+ 10 \MOD 3 * 6 \DIV 3 =3*6 + 1* 6 \DIV 3 = 18+2=20$.
\end{example}

\begin{remark}
При вычислении числа $a_i$ 
растут, и промежуточные результаты их не превосходят.
Поэтому построенный алгоритм имеет
такую же область определения и
сохраняет споcобность работать
без переполнения, как и предыдущий алгоритм.
Алгоритм содержит один цикл и имеет трудоёмкость
$O(n)$.   
\end{remark}


На следующих двух слайдах проведено сравнение двух представленных алгоритмов
вычисления биномиальных коэффициентов
по области применения и по быстродействию.

\end{frame}



\begin{frame}
\frametitle{5.3.4. Применение рекуррентного соотношения (4/5)}

\begin{columns}
\begin{column}{7cm}
Области определения этих алгоритмов совпадают.
На рисунке представлена область определения алгоритмов, то есть множество пар $(m,n)$
для которых величина $C(m,n)$ 
\em{может быть} вычислена в 32-х разрядной арифметике. На данном графике разными цветами, в зависимости от величины, обозначены значения биномиальных коэффициентов. В незакрашенной части
значения не могут быть вычислены. Участки, близкие к границам ($n=0$ или $n=m$), обрезаны.
\end{column}
\begin{column}{8cm}
\includegraphics[width=7.8cm]{5_3_4_1.jpg}
\end{column}
\end{columns}
\end{frame}



\begin{frame}
\frametitle{5.3.4. Применение рекуррентного соотношения (5/5)}

\smallskip
\begin{columns}
\begin{column}{7.8cm}
Асимптотическая сложность первого алгоритма равна 
$O(n*(m-n))$, а второго~---  $O(n)$. 
Но не стоит слепо доверять асимптотической 
оценке~--- на реальное быстродействие алгоритмов влияют и другие факторы, в частности, скорость выполнения арифметических операций.  
Известно, что, в сравнении со сложением, операции умножения и особенно деления являются более затратными по времени выполнения.
График справа отражает минимальное (из двух) время, затраченное для нахождения значений биномиальных коэффициентов.
Черная кривая обозначает границу, ниже которой квадратичный алгоритм справляется с поставленной задачей быстрее линейного! 
\end{column}
\begin{column}{7.2cm}
\includegraphics[width=7cm]{5_3_4_2.jpg}
\end{column}
\end{columns}

\end{frame}


\begin{frame}
\label{5.3.5.}\subsubsection{5.3.5. Генерация подмножеств }\frametitle{5.3.5. Генерация подмножеств (1/4)}
\label{gensubsets}


Элементы множества $\{1, \dots, m\}$ естественным образом
упорядочены. Поэтому каждое $n$-элементное подмножество этого
множества также можно рассматривать как упорядоченную
последовательность. На множестве таких последовательностей
определяется \em{лексикографический порядок}.
Следующий простой алгоритм \odmkey{}{Алгоритм
(algorithm)!генерации подмножеств (generation of subsets)}
генерирует все $n$-элементные подмножества $m$-эле\-мент\-ного
множества в лексикографическом порядке.

\begin{algorithmic}
\REQUIRE $n$~--- мощность подмножества, $m$~--- мощность множества, $m\ge n>0$.
\ENSURE
последовательность всех $n$-элементных подмножеств
$m$-элементного множества в лексикографическом порядке.
\FOR{\textbf{$i$ from $1$ to $m$}}
  \STATE $A[i] \Assign  i$ 
  \COMMENT{\color{gray}\small инициализация исходного множества }
\ENDFOR
\STATE {\textbf {if $m=n$ return $A[1..n]$ end if} }
%\IF{$m=n$}
%\STATE \textbf{return $A[1..n]$} \COMMENT{ единственное подмножество }
%\ENDIF
\end{algorithmic}

\end{frame}
\begin{frame}
\frametitle{5.3.5. Генерация подмножеств (2/4)}

%\pagebreak
%\begin{algorithm}
%\caption{Генерация $n$-элементных подмножеств $m$-элементного множества}
%\label{alg53}

\begin{algorithmic}
\STATE      $p \Assign n$ 
\COMMENT{\color{gray}\small $p$~--- номер первого изменяемого элемента }
\WHILE{$p \ge 1$}
\STATE \textbf{yield $A[1..n]$}
\COMMENT{\color{gray}\small очередное подмножество в первых $n$ элементах массива $A$ }
\IF{$A[n] = m$}
\STATE $p \Assign  p - 1$ \COMMENT{\color{gray}\small нельзя увеличить последний элемент }
\ELSE
  \STATE $p \Assign  n$ \COMMENT{\color{gray}\small можно увеличить последний элемент }
\ENDIF
\IF{$p \ge 1$}
  \FOR{\textbf{$i$ from $n$ downto $p$}}
    \STATE $A[i] \Assign  A[p] + i - p + 1$ 
    \COMMENT{\color{gray}\small увеличение элементов }
  \ENDFOR
\ENDIF
\ENDWHILE
\end{algorithmic}
%\end{algorithm}


\end{frame}


\begin{frame}
\frametitle{5.3.5. Генерация подмножеств (3/4)}
\vspace{-0.25\baselineskip}
\begin{ground}
Заметим, что в искомой последовательности $n$-элементных
подмножеств (каждое из которых является возрастающей
последовательностью $n$ чисел из диапазона $1..m$) вслед за
последовательностью $({\Tuple{a}{1}{n}})$ следует
последовательность\\
$({\Tuple{b}{1}{n}})=({a_1,\dots,a_{p-1},a_p+1,a_p+2,\dots,a_p+n-p+1})$,
где $p$~--- максимальный индекс, для которого $b_n=a_p+n-p+1\le
m$. Другими словами, следующая последовательность получается из
предыдущей заменой некоторого количества элементов в хвосте
последовательности идущими подряд натуральными числами, но так,
чтобы последний элемент не превосходил $m$, а первый изменяемый
элемент был на 1 больше, чем соответствующий элемент в предыдущей
последовательности. Таким образом, индекс $p$, начиная с которого
следует изменить <<хвост последовательности>>, определяется по
значению элемента $A[n]$. Если $A[n]<m$, то следует изменять
только\vadjust{\penalty-1000} $A[n]$, и при этом $p \Assign n$. Если же уже $A[n]=m$, то
нужно уменьшать индекс $p\Assign p-1$, увеличивая длину
изменяемого хвоста.
\end{ground}


\end{frame}


\begin{frame}
\frametitle{5.3.5. Генерация подмножеств (4/4)}

%\vspace{-2pt}
\begin{example}
 Последовательность $n$-элементных подмножеств
$m$-элементного множества \\в~лексикографическом порядке для $n=3$ и
$m=4$: $(1,2,3)$, $(1,2,4)$, $(1,3,4)$, $(2,3,4)$. 
\end{example}

\begin{digress}
Часто встречаются задачи, в которых требуется найти в
некотором множестве элемент, обладающий определёнными свойствами.
При этом исследуемое множество может быть велико, а его элементы
устроены достаточно сложно. Если неизвестно заранее, как именно
следует искать элемент, то остаётся перебирать все элементы, пока
не попадётся нужный (поэтому такие задачи называют
\odmkey{переборными}{Задача (problem)!переборная (search)}).
При решении переборных задач совсем необязательно и,
как правило, нецелесообразно строить всё исследуемое множество
(\odmkey{пространство поиска}{Пространство (space)!поиска
 (search)}) в явном виде.
Достаточно иметь итератор, 
%(см. п.~\ref{iterators}),
перебирающий элементы множества в подходящем порядке.
Многие рассмотренные алгоритмы
%~\ref{gray}, \ref{alg52} и \ref{alg53}
являются примерами таких итераторов
для некоторых типичных пространств поиска.
Чтобы использовать эти алгоритмы в качестве итераторов
при переборе, достаточно вставить
вместо оператора {\bf yield} проверку того,
что очередной элемент является искомым.
\end{digress}


\end{frame}

\begin{frame}
\label{5.3.6.}\subsubsection{5.3.6. Расширенные биномиальные коэффициенты }\frametitle{5.3.6. Расширенные биномиальные коэффициенты (1/3)}
Сократим множитель $(m-n)!$ в формуле для числа сочетаний:
$$
 C(m,n)=\frac{m (m-1)\dots(m-n+1)}{n!}.
$$ 
При определении чисел  $C(m,n)$ 
  по смыслу комбинаторной конфигурации параметры натуральны: $m,n \in \Bbb{N}$. 
Однако данная формула сохраняет свою вычислимость 
и при отрицательных 
значениях $m$, что позволяет ввести в 
рассмотрение \odmkey{расширенные биномиальные коэффициенты}{Коэффициент (coefficient)!биномиальный (binomial)!расширенный (extended)}
для $m \in \Bbb{Z}$ и $n\in \Bbb{N}_0$.  

\begin{remark} 
В качестве области определения расширенных
биномиальных коэффициентов можно 
рассматривать также рациональные $m\in \Bbb{Q}$
и вещественные значения  $m\in \Bbb{R}$.
\end{remark}

Пусть $m\in \Bbb{Z}, n\in \Bbb{N}_0$.
Положим $C(m,n)\Def 0$ при $0< m < n$, и положим
$C(m,0)\Def 1$ при $n=0$. В остальных случаях   
$C(m,n)\Def ( m (m-1)\dots(m-n+1))/n!$.
 
\begin{remark}
Можно доопределить $C(m,n)\Def 0$ при $n<0$.  
Тогда
$ \A{m\in \Bbb{Z}}{C(m,0)=1} $ и
$ \A{ n\in\Bbb{Z}, n\ne 0}{ C(0,n)=0} $.
Однако расширение области определения
$n$ не имеет существенного вычислительного смысла. 
\end{remark}
\end{frame}

\begin{frame}
\frametitle{5.3.6. Расширенные биномиальные коэффициенты (2/3)}

\begin{theorem}
 $C(m,n)=(-1)^n  C(n-m-1,n)$.
\end{theorem}
\begin{proof}
Заметим, что $m(m-1)\dots(m-n+1) =$ \\ 
$=(-1)^n (-m)(-m+1)\dots(-m+n-1) =(-1)^n (n-m-1)\dots(n-m-1 -(n-1))$.
Откуда
$C(m,n)= m\dots(m-n+1)/n!=(-1)^n (n -m-1)\dots (-m)/n! =$ \\ $=
(-1)^n C(n-m-1,n).
$
\end{proof}	


Доказанное тождество можно использовать для вычисления биномиальных 
коэффициентов с отрицательным первым параметром 
по уже вычисленным значениям коэффициентов с положительными параметрами.

\begin{remark}  
В теореме \DMref{п.~5.3.1}{5.3.1.} при доказательстве тождества 
$$C(m,n)=C(m-1,n-1)+C(m-1,n)$$  
 нигде не использовано то обстоятельство, что число $m>0$, 
 что позволяет организовать рекуррентное вычисление 
 биномиальных коэффициентов и
  для случая $m<0$.
\end{remark} 

Используя эти тождества, можно построить 
\odmkey{расширенный треугольник Паскаля}{Треугольник Паскаля (Pascal’s triangle)!расширенный (extended)},
в котором положительные значения $m$ распространяются вправо, 
отрицательные значения $m$ распространяются влево, 
а $n$ распространяется вниз. 
\end{frame}

\begin{frame}
\frametitle{5.3.6. Расширенные биномиальные коэффициенты (3/3)}
Поскольку 
при $m=0$ имеем $C(m,n)=0$, средний столбец,
обозначенный $n$,
содержит сами значения $n$, подобно тому, как верхняя строка содержит значения $m$. 

\begin{center}
\begin{tabular}{rrrr||c||rrrrr}
 $-4$ & $-3$ & $-2$ & $-1$ & $\frac{m}{n}$ & $\phantom{-}0$ & $\phantom{-}1$ & $\phantom{-}2$ & $\phantom{-}3$ & 
$\phantom{-}4$ \\
\hline
\hline
$1$   & $1$   & $1$  & $1$  & $0$ & $1$ & $1$ & $1$ & $1$ & $1$ \\
$-4$  & $-3$  & $-2$ & $-1$ & $1$ & $0$ & $1$ & $2$ & $3$ & $4$ \\
$10$  & $6$   & $3$  & $1$  & $2$ & $0$ & $0$ & $1$ & $3$ & $6$ \\
$-20$ & $-10$ & $-4$ & $-1$ & $3$ & $0$ & $0$ & $0$ & $1$ & $4$ \\
$35$  & $15$  & $5$  & $1$  & $4$ & $0$ & $0$ & $0$ & $0$ & $1$ \\
\hline 
\end{tabular}

%По вертикальной оси отложены -- n, по горизонтальной -- m.
\end{center}

\begin{example}
	Расширенные биномиальные коэффициенты применяются для разложения дробей в ряд:
	$$\frac{1}{(1+x)^m} = (1+x)^{-m} = \sum_{n=0}^\infty  C(-m,n) x^n = \sum_{n=0}^\infty (-1)^n C(m+n-1,n)x^n.$$
\end{example}	
 

\end{frame}

%\begin{frame}
%\frametitle{5.3.6. Расширенные биномиальные коэффициенты (4/4)}
%На область применения расширенных биномиальных коэффициентов 
%накладывается \\меньше ограничений, 
%что может облегчить доказательство теорем и тождеств. 
%
%
%\begin{example}
% Докажем утверждение 
% $\displaystyle{C(m,n)= \frac{m}{n} C(m-1,n-1)}$.
%При использовании формулы $\displaystyle{C(m,n) =\frac{m!}{n!(m-n)!}}$ 
%доказательство требует разбора случаев.
%
%Первый случай. При $m\ge n$ имеем $$\displaystyle{C(m,n)= \frac{m!}{n!(m-n)!}
%=\frac{m}{n}  \frac{(m-1)!}{(n-1)!((m-1)-(n-1) )!}=}
%{\frac{m}{n} C(m-1,n-1)}.$$
%
%Второй случай. При $ m<n$ имеем $C(m,n)=0$ и $C(m-1,n-1)=0$.
%
%При использовании расширенных биномиальных коэффициентов сразу имеем \\
%$C(m,n)= m(m-1)…(m-n+1)/n!=$ \\ $=(m/n)
%(m-1)…(m-n+1)/(n-1)!=(m/n)C(m-1,n-1)$.
%\end{example}
%
%%{\color{blue} А в чём облегчение, стесняюсь спросить? Ф.А.Н. }
%%
%%Также подобное обобщение позволяет 
%%расширить область применения некоторых теорем, 
%%например, биномиальную теорему становится 
%%возможным расширить на случай 
%%отрицательных и дробных степеней. 
%%
%%{\color{blue} Набор примеров слабоват. Есть ещё? Ф.А.Н. }
%%
%
%\end{frame}

\begin{frame}
\label{5.3.7.}\subsubsection{5.3.7. Мультимножества и последовательности }\frametitle{5.3.7. Мультимножества и последовательности (1/2)}

%\vspace{-2pt}
В \DMref{п.~5.1.5}{5.1.5.} показано, что число $n$-элементных подмножеств
$m$-элементного множества равно $C(m,n)$. Поскольку $n$-элементные
подмножества~--- это $n$-эле\-мент\-ные индикаторы
(\DMref{п.~1.1.4}{1.1.4.}), ясно, что число  $n$-элемент\-ных
индикаторов над $m$-элементным множеством также равно $C(m,n)$.

\medskip
Общее число  $n$-элементных мультимножеств $\left[ x_1^{n_1},
\dots, x_m^{n_m}\right]$, в которых $n_1 +\dots +\, n_m =
n$, над $m$-элементным множеством $X=\{\Tuple{x}{1}{m}\}$
нетрудно определить, если заметить, что это число способов
разложить $n$ неразличимых мячей по $m$ ящикам, то есть число
сочетаний с повторениями $V(m,n)=C(n +m -1,n)$.

\medskip
Пусть задано множество $X=\{\Tuple{x}{1}{m}\}$ и
последовательность $Y=({\Tuple{y}{1}{n}})$, где
$\A{i}{y_i\in X}$, причём в последовательности $Y$ элемент $x_i$
встречается $n_i$ раз. Тогда мультимножество $\widehat{X}=\left[
x_1^{n_1}, \dots, x_m^{n_m}\right]$ называется
\odmkey{составом}{Состав (compound)!последовательности (of
sequence)} последовательности $Y$. Ясно, что состав любой
последовательности определяется однозначно, но разные
последовательности могут иметь один и тот же состав.

\begin{example}
Пусть $X=\{1,2,3\}$, $Y_1=\{1,2,3,2,1\}$,
$Y_2=\{1,1,2,2,3\}$.
Тогда последовательности $Y_1$ и $Y_2$
имеют один и тот же состав
$\widehat{X}=\left[ 1^2,  2^2, 3^1 \right]$.
\end{example}

\end{frame}

\begin{frame}
\frametitle{5.3.7. Мультимножества и последовательности (2/2)}


\begin{digress}
Из органической химии известно, что существуют
различные вещества, имеющие один и~тот же химический состав.
Они называются изомерами.
\end{digress}

\medskip
Очевидно, что все последовательности над множеством
$X=\{\Tuple{x}{1}{m}\}$, имеющие один и тот же состав
$\widehat{X}=\left[ x_1^{n_1}, \dots, x_m^{n_m}\right]$, имеют
одну и ту же длину $n  = n_1 +\dots +\:n_m$. Обозначим число
последовательностей одного состава $C(n;\Tuple{n}{1}{m})$.

\begin{theorem}
$\displaystyle{C(n;\Tuple{n}{1}{m})=\frac{n!}{n_1!\dots n_m!}}$.
\end{theorem}

\begin{proof}
Рассмотрим мультимножество
$\widehat{X}=\left[ x_1^{n_1}, \dots, x_m^{n_m}\right]$,
которое можно записать в форме последовательности
${\Tuple{x}{1}{1},\Tuple{x}{2}{2},\dots,\Tuple{x}{m}{m}}$,
где элемент $x_i$ встречается $n_i$ раз.
Ясно, что все последовательности одного состава $\widehat{X}$
получаются из указанной последовательности перестановкой элементов.
При этом, если переставляются одинаковые элементы, например,
$x_i$, то получается та же самая последовательность.
Таких перестановок $n_i!$. 
Таким образом, общее число
перестановок $$n_1!\dots n_m!C(n;\Tuple{n}{1}{m}).$$
С другой стороны, число перестановок $n$ элементов равно $n!$.
\end{proof}

%\begin{example} Сколько различных слов можно составить из букв слова <<математика>>? 
%$C(10; 3,2,2,1,1,1) = 10!/(3!2!2!1!1!1!)=151200$ слов.
%\end{example}
\end{frame}


\begin{frame}
\label{5.3.8.}\subsubsection{5.3.8. Мультиномиальные коэффициенты }\frametitle{5.3.8. Мультиномиальные коэффициенты (1/3)}
%\vspace{-2pt}
Числа $C(n;n_1,\dots,n_m)$ встречаются в практических
задачах довольно часто, поскольку обобщают случай
индикаторов, которые описываются биномиальными коэффициентами,
на случай произвольных мультимножеств.
В частности, справедлива формула, обобщающая формулу
бинома Ньютона, в которой участвуют числа $C(n;n_1,\dots,n_m)$,
поэтому их иногда называют
\odmkey{мультиномиальными коэффициентами}{Коэффициент
 (coefficient)!мультиномиальный (multinomial)}.

\begin{theorem}
$ \displaystyle{(x_1+\ldots +x_m)^n   = \sum_{n_1+\ldots+n_m=n}
C(n;\Tuple{n}{1}{m}) x_1^{n_1}\cdots x_m^{n_m} }$.
\end{theorem}

\begin{proof}
Индукция по $m$. База: при $m=2$ по формуле бинома Ньютона имеем
${(x_1+x_2)}^n=\displaystyle{\sum_{i=0}^{n} C(n,i)x_1^{i} x_2^{n-i} =
\sum_{i=0}^{n} \frac{n!}{i!(n-i)!} x_1^{i} x_2^{n-i} =
\sum_{n_1+n_2=n} \frac{n!}{n_1! n_2!} x_1^{n_1} x_2^{n_2}.
}$
Пусть теперь $ \displaystyle{(x_1+\ldots +x_{m-1})^n   =
\sum_{n_1+\ldots+n_{m-1}=n} C(n;\Tuple{n}{1}{m-1})
 x_1^{n_1}\cdots x_{m-1}^{n_{m-1}} }.$
\end{proof}

\end{frame}

\begin{frame}
\frametitle{5.3.8. Мультиномиальные коэффициенты (2/3)}
\begin{proof} (Продолжение)
Рассмотрим ${(x_1+\ldots+x_m)}^n$. Имеем
$
(x_1+\ldots+x_m)^n= $\\
$\displaystyle{ =\left((x_1+\ldots+x_{m-1})+x_m \right)^n =
\sum^{n}_{i=0}C(n,i)(x_1+\ldots+x_{m-1})^i x_m^{n-i}=}$\\
$\displaystyle{=
\sum^{n}_{i=0} \frac{n!}{i!(n-i)!}
   \left(\sum_{n_1+\ldots+n_{m-1} = i}
   \frac{i!}{n_1!\cdots n_{m-1}!} x_1^{n_1}\cdots x_{m-1}^{n_{m-1}} \right)
  x_m^{n-i}  =}$\\
$\displaystyle{= \sum^{n}_{i=0} \sum_{n_1+\ldots+n_{m-1} = i}
   \frac{n!i!}{i!(n-i)!n_1!\cdots n_{m-1}!} x_1^{n_1}\cdots x_{m-1}^{n_{m-1}}
  x_m^{n-i}  =}$\\
$\displaystyle{= \sum_{n_1+\ldots+n_{m-1}+n_m = n}
   \frac{n!}{n_1!\cdots n_{m-1}!n_m!} x_1^{n_1}\cdots x_{m-1}^{n_{m-1}}
   x_m^{n_m}}$,\\
где $n_m\Assign n-i$.
\end{proof}

\end{frame}

\begin{frame}
	\frametitle{5.3.8. Мультиномиальные коэффициенты (3/3)}
	\begin{corollary}
		$\displaystyle{\sum_{n_1+\ldots+n_m = n}
			\frac{n!}{n_1!\cdots !n_m!} = m^n}$.
	\end{corollary}
	\begin{proof}
		$\displaystyle{m^n = (\underbrace{1+\ldots+1}_{m\ \text{раз}})^n =
\sum_{n_1+\ldots+n_m = n} \frac{n!}{n_1!\cdots !n_m!}1^{n_1}\dots 
1^{n_m}}=$ \\$=\displaystyle{
			\sum_{n_1+\cdots+n_m = n} \frac{n!}{n_1!\ldots !n_m!} }$.
	\end{proof}


\begin{example}
При игре в преферанс трём играющим
сдаётся по 10 карт из 32 карт и 2 карты остаются в <<прикупе>>.
Сколько существует различных раскладов?
$$
C(32;10,10,10,2)=\frac{32!}{10!10!10!2!}=
2\,7\,5\,3\,2\,9\,4\,4\,0\,8\,5\,0\,4\,6\,4\,0.
$$
\end{example}

\end{frame}

\begin{frame}
	\label{5.3.9.}\subsubsection{5.3.9. Биномиальная система счисления }\frametitle{ 5.3.9. Биномиальная система счисления (1/3)}
	\odmkey{$k$-биномиальной системой счисления}{Система счисления (number system, number language, notation, numeration)!k-биномиальная (k-binomial)}
	называется непозиционная система счисления,
	где каждое натуральное число $n$  представляется в виде:
	$$n=C(b_1,1)+ C(b_2,2)+\dots +C(b_k,k)=\sum\limits_{i=1}^k C(b_i,i),\quad 
	 k\in \Bbb{N},  
	0 \le b_1 < b_2< \dots <b_k.$$
	
	\begin{example}
		Пусть k = 2.\\
		$1=C(0,1)+ C(2,2) = 0 +1$\\
		$2=C(1,1)+ C(2,2)= 1+1$\\
		$3=C(0,1)+ C(3,2)= 0+3$\\
		$4=C(1,1)+ C(3,2)= 1+3$\\
		$5=C(2,1)+ C(3,2)=2+3$\\
		$6=C(0,1)+ C(4,2)=0+6$\\
		и т. д.
	\end{example}
	\begin{remark}  
		По определению,~$C(p,q)=0$ при целых $p$ и $q$ таких, что
		$0 \le p < q$.
	\end{remark}

\end{frame}



\begin{frame}
\frametitle{ 5.3.9. Биномиальная система счисления (2/3)}
	Покажем состоятельность данной системы счисления.

\smallskip
\begin{theorem}
При заданном натуральном значении $k$
для любого натурального числа $n$ 
существуют и единственны натуральные числа $\Tuple{b}{1}{k}$,
такие что $0 \le b_1 < b_2< \dots <b_k$ и
$n=C(b_1,1)+ C(b_2,2)+\dots +C(b_k,k)=\sum\limits_{i=1}^k C(b_i,i).
$

\end{theorem}

%\vspace{-0.5\baselineskip}
\begin{proof}
\proofsection{Существование} 
Всякой паре натуральных чисел $n$ и $k$ сопоставим набор \\
$0 \le b_1 < b_2< \dots <b_k$
строго возрастающих натуральных чисел 
по следующему правилу. Сначала выберем $b_k \ge 0$ так, чтобы
удовлетворялись неравенства\\ $C(b_k,k) \le n < C(b_k+1,k)$.
Такой выбор однозначен, потому что 
по основному рекуррентному соотношению
$C(b_k + 1,k)  =  C(b_k,k) + C(b_k,k - 1)$ и существует
единственное число $b_k:$
$C(b_k,k) \le n<C(b_k,k)+C(b_k,k-1)$.
Затем определим $b_{k-1}$: \\
$C(b_{k-1},k - 1)  \le  n - C(b_k,k)<C(b_{k-1} + 1,k - 1)$.
При таких условиях обязательно выполняется неравенство 
$b_k>b_{k-1}$.
Действительно,
$C(b_{k-1},k - 1)  \le  n  -  C(b_k,k)  <$ \\ $<  C(b_k + 1,k) - C(b_k,k)= C(b_k,k-1)$, 
откуда $b_{k-1}<b_k$.
\end{proof}
\end{frame}

\begin{frame}
\frametitle{ 5.3.9. Биномиальная система счисления (3/3)}
\vspace{-2pt}
\begin{proof} (Продолжение) 
Продолжая действовать подобным образом, определим числа
$b_{k-2}, b_{k-3}, \dots , b_1$. Имеем
$C(b_1,1) \le n-C(b_k,k)- \dots -C(b_2,2)<C(b_1+1,1)$.
Откуда
$n=C(b_1,1)+ C(b_2,2)+ \dots +C(b_k,k)$.
Таким образом, построили требуемый набор чисел $0 \le b_1 < b_2< \dots <b_k$,
дающий искомое представление числа $n$. 

\proofsection{Единственность}
Индукция по $k$.
При $k=1$ имеем $C(n,1)=n$.
Пусть для $k-1$ представление единственно и 
существует отличный от построенного набор 
натуральных чисел $0 \le c_1 < c_2< \dots <c_k$ такой, что
$n=C(c_1,1)+ C(c_2,2)+\dots +C(c_k,k)$.
В силу условия на $b_k$ имеем $c_k \le b_k$, причем равенство $c_k = b_k$ противоречит предположению индукции. Значит, $c_k \le b_k-1$, и имеем цепочку неравенств
$n  -  C(c_k,k)  \ge $ \\    $\ge C(b_k,k)  -  C(c_k,k)  \ge  C(b_k,k) -  C(b_k - 1,k) =C(b_k-1,k-1) \ge C(c_k,k-1)$.
С другой стороны, в силу предположения индукции, представление числа $n-C(b_k,k)$ в виде   
$n=C(c_1,1)+ C(c_2,2)+\dots +C(c_{k-1},k-1)$ единственно, 
причем $n-C(c_k,k)<$ \\$<C(c_{k-1}+1,k-1)$. 
Сравнивая два последних неравенства, приходим к противоречию
$c_{k-1}+1>c_k$.
\end{proof}
%\odmkey{k-биномиальной системой счисления}.
%называется непозиционная система счисления,
%где каждое натуральное число $n$  представляется в виде:
%
%\vspace{-\baselineskip}
%$$n=C(b_1,1)+ C(b_2,2)+\dots +C(b_k,k)=\sum_{i=1}^k C(b_i,i).$$

\end{frame}

\begin{frame}
\frametitle{Онтология комбинаторного анализа}
\begin{center}
\includegraphics[height=7.8cm]{ODM5_3.png}
\end{center}
\end{frame}

\begin{frame}
\frametitle{5.4. Разбиения}
\subsection{5.4. Разбиения}
\label{5.4.}

Разбиения не рассматривались среди типовых комбинаторных
конфигураций в разделе~\DMref{5.1}{5.1.}, потому что получить для них
явную формулу не так просто, как для остальных. В этом разделе
отмечаются основные свойства разбиений, а окончательные формулы
приведены в \DMref{п.~5.6.3}{5.6.3.}.

\bigskip
\begin{tabular}{p{11mm}|p{45mm}|p{88mm}}
  № & Название параграфа
  & Ключевые термины и обозначения \\
  \hline
\DMref{5.4.1.}{5.4.1.} & Числа Стирлинга  &
{ Блоки разбиения; $S(m,n)$~--- число разбиений}\\  
\DMref{5.4.2.}{5.4.2.} & Числа Моргана &
{ $M(m,n)=n!S(m,n)$}\\  
\DMref{5.4.3.}{5.4.3.} & Вычисление чисел Стирлинга  &
{ }\\  
\DMref{5.4.4.}{5.4.4.} & Числа Белла  &
{$B(m)$~--- число всех разбиений  }\\  
\DMref{5.4.5.}{5.4.5.} & Треугольник Белла  &
{ }\\  
\end{tabular}  



\end{frame}


\begin{frame}
\label{5.4.1.}\subsubsection{5.4.1. Числа Стирлинга }\frametitle{5.4.1. Числа Стирлинга (1/3)}

Пусть $\cal B = \{ B_1, \dots,B_n \}$~---
разбиение множества $X$ из $m$
элементов на $n$
подмножеств:
$$
B_i \subset X, \qquad
\bigcup_{i=1}^{n}B_i = X, \qquad
B_i \not= \emptyset, \qquad
B_i \cap B_j = \emptyset \quad\text{при} 
i \not= j.
$$
Подмножества $B_i$ называются \odmkey{блоками\/}{Блок
(block)!разбиения (of partition)} разбиения.


\medskip
Между разбиениями с непустыми блоками и отношениями
эквивалентности существует взаимно-однозначное соответствие
(\DMref{п.~1.7.1}{1.7.1.}). 

Пусть $E_1$ и $E_2$~--- два разбиения $X$. Говорят, что разбиение $E_1$ есть
\odmkey{измельчение\/}{Измельчение разбиения (refinement of
partition)} разбиения $E_2$, если каждый блок $E_2$ есть
объединение блоков $E_1$. Измельчение является частичным порядком
на множестве разбиений.





\sloppypar Число разбиений $m$-элементного множества на $n$ блоков, или число способов разложить $m$
мячей по $n$ неразличимым ящикам,  
называется
\odmkey{числом
%\linebreak
Стирлинга}{}
%\footnote{\;Джеймс Стирлинг (1699--1770).} второго рода}
%{Число (number)!Стирлинга второго рода (Stirling second kind)} 
и
обозначается $S(m,n)$. По определению положим
$S(m,m) \Def 1$ при $m \ge 0$,
$S(m,0) \Def 0$ при $m > 0$, 
$S(m,n) \Def 0$ при $n > m$.
%\begin{align*}
% S(m,m) &\Def  1, \quad S(m,0) \Def  0  \quad \text{при $m > 0$,}\\
% S(0,0) &\Def  1, \quad S(m,n) \Def  0 \quad  \text{при $n > m$.}
%\end{align*}
\end{frame}

\begin{frame}
\frametitle{5.4.1. Числа Стирлинга  (2/3)}

\begin{theorem} {\NB{[1]}}
$ S(m,n) = S(m-1,n-1) + nS(m-1,n)$.
\end{theorem}
\begin{proof}
Пусть $\cal B$~--- множество всех разбиений множества
$M\Assign\{1, \dots,m\}$ на $n$ блоков. Положим
$$\cal B_1\Assign\Set{X\in\cal B}{\E{B\in X}{B=\{m\}}},\quad
\cal B_2\Assign\Set{X\in\cal B}{\neg\E{B\in X}{B=\{m\}}},$$
то~есть в $\cal B_1$ входят разбиения, в которых элемент $m$
образует отдельный блок, а в $\cal B_2$~--- все остальные
разбиения. Заметим, что $\cal B_2=\Set{X\in\cal B}{m\in X
\Impl\vert X\vert>1}$. Тогда \\$\cal B=\cal B_1\cup \cal B_2$, $\cal
B_1\cap\cal B_2=\emptyset$. Имеем 
$$\vert\cal B_1\vert=S(m-1,n-1)\quad
\vert\cal B_2\vert=n S(m-1,n),$$
 так как все разбиения 
 $\cal B_2$ 
 получаются следующим образом: берём все разбиения множества
$\{1,\dots, m-1\}$ на $n$ блоков (их $S(m-1,n)$) и в каждый блок
по очереди помещаем элемент $m$. Следовательно, 
$
S(m,n)=\vert\cal B\vert= \vert\cal B_1\vert+\vert\cal
B_2\vert=S(m-1,n-1)+nS(m-1,n).
$
(\DMref{п.~1.2.6}{1.2.6.})
\end{proof}

\end{frame}

\begin{frame}
\frametitle{5.4.1. Числа Стирлинга  (3/3)}


\begin{theorem} {\NB{[2]}}
$ S(m,n) = \sum\limits^{m-1}_{i=n-1} C(m-1,i)S(i,n-1)$.
\end{theorem}


\begin{proof}
Пусть $\cal B$~--- множество всех разбиений множества
$M\Assign\{1,\dots, m\}$ на $n$ блоков. Рассмотрим семейство
$\overline B\Assign\Set{B\subset 2^M}{m\in B}$. Тогда $\cal
B=\bigcup_{B\in\overline B} {\cal B}_B$, где ${\cal
B}_B\Assign\Set{X}{X\in\cal B\And B\in X}$, причём ${\cal
B}_{B'}\cap{\cal B}_{B''}=\emptyset$, если $B'\not= B''$. Пусть
${B\in\overline B}$ и $b\Assign\vert B\vert$. Тогда $\vert\cal
B_B\vert=S(m-b,n-1)$. Заметим, что $\left\vert\Set{B\in\overline
B}{b=|B|}\right\vert =C(m-1,b-1)$. Имеем

$S(m, n)=
\vert\cal B\vert =
\sum\limits^{m-(n-1)}_{b=1}
\left\vert\bigcup\limits_{B\in\overline{B}\And
\vert B\vert=b} {\cal B}_B\right\vert =\sum\limits^{m-(n-1)}_{b=1}C(m-1,b-1)S(m-b,n-1)=$\\
$=\sum\limits^{n-1}_{i=m-1} C(m-1,m-i-1) S(i,n-1)=
  \sum\limits^{m-1}_{i=n-1} C(m-1,i) S(i,n-1),
$
где $i\Assign m-b$.
\end{proof}

\end{frame}


\label{5.4.2.}\subsubsection{5.4.2. Числа Моргана}
\begin{frame}\frametitle{5.4.2. Числа Моргана (1/2)}

Число сюръективных функций $\Map{f}{1..m}{1..n}$, то~есть число  
таких размещений $m$ мячей
по $n$ ящикам, что все ящики заняты, называется
\odmkey{числом Моргана}{Число (number)!Моргана (Morgan)} и обозначается $M(m,n)$. 


\medskip
\begin{example}
Число массивов
\textbf{$a:$ array $[1..m]$ of $1..n$} таких,
что $\A{x\!\in\!1..n}{\E{i\!\in\! 1..m}{a[i]=x}}$,
является числом сюръективных функций 
и вычисляется числом Моргана.
\end{example}



\medskip
\begin{remark}
Числа Стирлинга называют также 
\odmkey{числами Стирлинга второго рода}{Числа (numbers)!Стирлинга (Stirling)!второго рода (of second kind)}, 
а числа Моргана называют 
\odmkey{числами Стирлинга первого рода}{Числа (numbers)!Стирлинга (Stirling)!первого рода (of first kind)} и обозначают $s(m,n)$. 
\end{remark}

\medskip
\begin{theorem} {\NB{[1]}}
$M(m,n)=n!\,S(m,n)$.
\end{theorem}

\begin{proof}
Каждое разбиение множества $\{1,\dots,m\}$ соответствует семейству
множеств уровня сюръективной функции и обратно (\DMref{п.~1.7.4}{1.7.4.}).
Таким образом, число различных семейств множеств уровня
сюръективных функций~--- это число Стирлинга 
$S(m,n)$. Всего сюръективных функций $M(m,n)=n!\,
S(m,n)$, так~как число сюръективных функций с заданным семейством
множеств уровня равно числу перестановок множества значений
функции.
\end{proof}

\end{frame}

\begin{frame}\frametitle{5.4.2. Числа Моргана (2/2)}

\begin{theorem} {\NB{[2]}}
    $ M(m, n) = n \left(M(m-1, n-1) +  M(m-1, n)\right)$.
\end{theorem}
\begin{proof}
Используя рекуррентную формулу \DMref{п.~5.4.1}{5.4.1.}, имеем цепочку равенств:\\
$M(m, n) = n!S(m, n) = n!\left(S(m-1, n-1)+nS(m-1, n)\right) =$
\\ $=n!S(m-1, n-1) + n n!S(m-1, n)=n (n-1)!S(m-1, n-1)+n n!S(m-1, n)$. \\Применив теорему [1], получаем: \\$n (n-1)!S(m-1, n-1)+n n!S(m-1, n) = nM(m-1, n-1) + nM(m-1, n)$.    
\end{proof}

\begin{theorem} {\NB{[3]}}
$M(m, n) = \sum^{m-n+1}_{k=1} C(m, k) M(m-k, n-1).$   
\end{theorem}

\begin{proof}
Рассмотрим множество сюръективных функций 
$S=\{\Map{f}{1..m}{1..n}\}$. 
Разобьем всё множество сюръекций $S$ 
на дизъюнктные подмножества:\\ 
$S_1 = \Set{\Map{f}{1..m}{1..n}}{|f^{-1}(n)| = 1}$, 
 $\dots$, 
$S_{m-n+1} = \Set{\Map{f}{1..m}{1..n}}{|f^{-1}(n)| = m-n+1}$. Множества с  индексами большими, чем $m-n+1$ пусты. Заметим, что \\
$\left|S_k\right| =C(m, k) M(m-k, n-1)$, 
поскольку $k$ элементов из $X$ переходят
в элемент $n$, а значит 
только $m-k$ элементов из множества $X$ переходят
в оставшиеся $n-1$ элементов множества $Y$. 
Такие $k$ элементов можно выбрать $C(m, k)$ способами.  
Множества
$\Tuple{S}{1}{m-n+1}$ дизъюнктны и образуют покрытие по построению, а значит\\
$M(m, n) = |S| = \sum^{m-n+1}_{k=1}\left|S_k\right| = $
$=\sum^{m-n+1}_{k=1}C(m, k) M(m-k, n-1)$.               
\end{proof}

\end{frame}

\begin{frame}
\label{5.4.3.}\subsubsection{5.4.3. Вычисление чисел Стирлинга }\frametitle{5.4.3. Вычисление чисел Стирлинга (1/4)}

В \DMref{п.~5.4.1}{5.4.1.}  выведены две формулы для вычисления чисел Стирлинга.
Вторая формула 
$$ S(m,n) = \sum\limits^{m-1}_{i=n-1} C(m-1,i)S(i,n-1)$$  требует вcпомогательного вычисления биномиальных коэффициентов и выполнения большого количества сложений и умножений.

\smallskip 
При использовании первой рекуррентной формулы 
$$S(m,n) = S(m-1,n-1) + nS(m-1,n)$$
встаёт вопрос о хранении 
уже посчитанных чисел Стирлинга, необходимых для промежуточных вычислений, 
а также многократного перевычисления одних и тех же чисел в случае прямого использования рекурсии. 

\smallskip

\end{frame}

\begin{frame}\frametitle{5.4.3. Вычисление чисел Стирлинга (2/4)}

\begin{example}
Вычислим $S(7, 3)$. Для этого построим таблицу начальных значений \\чисел Стирлинга, пользуясь первой рекуррентной формулой.

\begin{center}
\begin{tabular}{c|cccccccc}
$m\setminus n$ & 0 & 1 & 2 & 3 & 4 & 5 & 6 & 7 \\
\hline
0 & 1 &  &  &  &  &  &  & \\
1 & \boxed{0} & \boxed{1} &  &  &  &  &  & \\
2 & \boxed{0} & \boxed{1} & \boxed{1} &  &  &  &  & \\
3 & \boxed{0} & \boxed{1} & \boxed{3} & \boxed{1} &  &  &  & \\
4 & \boxed{0} & \boxed{1} & \boxed{7} & \boxed{6} & 1 &  &  & \\
5 & 0 & \boxed{1} & \boxed{15} & \boxed{25} & 10 & 1 &  & \\
6 & 0 & 1 & \boxed{31} & \boxed{90} & 65 & 15 & 1 & \\
7 & 0 & 1 & 63 & \boxed{301} & 350 & 140 & 21 & 1 \\
\end{tabular}
\end{center}

%\vspace{-0.25\baselineskip}
Значения, которые необходимы 
для вычисления числа $S(7, 3)$,
обведены прямоугольниками. 

\end{example}

\end{frame}

\begin{frame}\frametitle{5.4.3. Вычисление чисел Стирлинга (3/4)}
%Алгоритм вычисления чисел Стирлинга.

\vspace{-6pt}

\begin{algorithmic}
\REQUIRE неотрицательные целые числа $m$ и $n$ 
\ENSURE число Стирлинга $S(m, n)$ 
\STATE{\textbf{if $m = n$ then return $1$ end if}}
\COMMENT {\color{gray}\small можем сразу вернуть известные значения }
\STATE{\textbf{if $n = 0$ then  return $0$ end if}}
\STATE{\textbf{if $n > m$ then return $0$ end if}}
\STATE{\textbf{if $n =1$ then return $1$ end if}}
\STATE{$d \Assign \min(n,m-n+1)$}
\COMMENT{\color{gray}\small  длина рабочего массива }
\STATE{\textbf{$D:$ array$[1..d]$ of integer}}
\COMMENT{\color{gray}\small  рабочий массив }
\STATE{$s \Assign \max(n,m-n+1)$}
\COMMENT{\color{gray}\small  количество повторений }
\STATE{\textbf{for $i$ from $1$ to $d$ do $D[i]\Assign 1$ end for}}
\COMMENT{\color{gray}\small инициализация }
\FOR{\textbf{$i$ from $2$ to $s$} }
     \FOR{\textbf{ $j$ from $2$ to $d$} }
%              \STATE{}
%              \COMMENT{\color{gray}\small перевычисление рабочего массива }
              \STATE\textbf{$D[j] \Assign $
              if $d=n$ then  $D[j-1]+j*D[j]$
              else  $D[j]+i*D[j-1]$ end if }
     \ENDFOR
\ENDFOR
\STATE{\textbf{return  $D[d]$ }}
\COMMENT{\color{gray}\small число Стирлинга } 
\end{algorithmic}
\end{frame}

\begin{frame}\frametitle{5.4.3. Вычисление чисел Стирлинга (4/4)}



\begin{ground}
Чтобы вычислить число Стирлинга, необходимо
вычислить все значения, показанные в 
<<параллелограмме>> предыдущего примера.
При этом хранить достаточно лишь рабочий массив,
равный длине более короткой стороны <<параллелограмма>>. 
При этом количество повторений 
соответствует более длинной стороне
<<параллелограмма>>.

Таким образом, рабочий массив соответствует либо 
диагонали, которая на каждой итерации опускается вниз,
либо столбцу, который на каждой итерации сдвигается вправо-вниз.  

В любом случае вначале рабочий массив инициализируется значениями $1$.
Далее на каждом шаге цикла по $i$ перевычисляется рабочий массив. Таким образом, 
вычисляются только необходимые промежуточные элементы, причём по одному разу.

В любом случае перевычисляется 
$(n-1)(m-n)$ элементов.  
\end{ground}

\medskip
\begin{remark}
При оценке области применимости алгоритма
следует обратить внимание на то, 
что числа Стирлинга для фиксированного $m$ 
с ростом $n$ сначала возрастают, а потом убывают. 
\end{remark}

\medskip
Формула $M(m,n)=n!\,S(m,n)$ позволяет использовать
алгоритм вычисления чисел Стирлинга для вычисления 
чисел Моргана.
\end{frame}


\begin{frame}
\label{5.4.4.}\subsubsection{5.4.4. Числа Белла}
\frametitle{5.4.4. Числа Белла}

Число всех разбиений
$m$-элементного множества называется
\odmkey{числом Белла}{Число
 (number)!Белла (Eric Temple Bell)}
% \footnote{\;Эрик Темпл Белл (1883--1960).}
и~обозначается $B(m)$,

\vspace{-10pt}
$$
B(m) \Def \sum^{m}_{n=0} S(m,n), \quad B(0) \Def  1.
$$


\begin{theorem}
$B(m+1) =   \sum\limits^{m}_{i=0} C(m,i) B(i)$.
\end{theorem}


\begin{proof}
Пусть $\cal B$~--- множество всех разбиений множества
$M_1=1..(m +1)$. Рассмотрим множество подмножеств множества
$M_1$, содержащих элемент $m+1$:\\
$\overline B\Assign{\left \{ \left.{B\subset 2^{ M_1}} \right| {m+1\in B} \right \}}$.
Тогда $\cal B=\bigcup\limits_{B\in\overline B}\cal B_B$,
где
$\cal B_B\Assign\Set{X\in\cal B}{B\in X}$.
Пусть $B\in\overline B$ и $b=\vert B\vert$.
Тогда $\vert \cal B_B\vert=B(m+1-b)$.
Заметим, что 
$\vert\Set{B\in\overline B}{\vert B\vert =b}\vert=$ $=C(m,b-1)$.
Следовательно,
%\begin{align*}
$B(m+1)=\vert\cal B\vert =\sum^{m+1}_{b=1}C(m,b-1)B(m-b+1)=$
$=\sum^0_{i=m}C(m,m-i)B(i)=\sum\limits^m_{i=0}C(m,i)B(i)$,
%\end{align*}
где $i\Assign m-b+1$.
\end{proof}

\end{frame}


\begin{frame}
\label{5.4.5.}\subsubsection{5.4.5. Треугольник Белла }\frametitle{5.4.5. Треугольник Белла (1/4)}
Вычисление числа Белла по рекуррентной формуле, 
выведенной в предыдущем параграфе, 
даже при условии хранения всех предыдущих вычисленных значений чисел Белла,
имеет трудоёмкость  $O(n^3)$. 
Рассмотрим алгоритм, использующий специальную структуру данных,
который в неявном виде хранит
все предыдущие вычисленные значения для биномиальных коэффициентов
и имеет трудоёмкость $O(n^2)$. 
\odmkey{Треугольник Белла}{Треугольник Белла (Bell's triangle)}~--- это бесконечная треугольная матрица, 
где каждое число (кроме чисел на левой боковой стороне) является суммой двух чисел: 
стоящего слева от него в этой же строке и слева над ним. Первое число в каждой строке 
(кроме первой) является последним числом из предыдущей строки. 
Число в первой строке~--- единица.


\begin{center}
\begin{tabular}{ccccccccccc}
   &    &    &    &    &  1  &    &    &   &    &   \\[-2pt] 
   &    &    &    &  1 &     & 2  &    &   &    &   \\[-2pt] 
   &    &    & 2  &    &  3  &    & 5  &   &    &   \\[-2pt]
   &    &  5 &    & 7  &     & 10 &    & 15&    &   \\[-2pt]
   & 15 &    & 20 &    &  27 &    & 37 &   & 52 &   \\[-4pt]
.  &    &  . &    & .  &    &  .  &    & . &    & .  \\[-1pt]
\end{tabular}
\end{center}

\begin{remark}
Треугольник Белла также называют треугольником Пирса.
\end{remark}
\end{frame}


\begin{frame}
\frametitle{5.4.5. Треугольник Белла (2/4)}
Для вычислений  
треугольник Белла удобнее хранить в  виде бесконечной нижней треугольной матрицы:

\begin{center}
\begin{tabular}{cccccc}
 1  &    &    &      &     &     \\ 
 1  & 2  &    &      &     &     \\ 
 2  & 3  &  5  &     &     &      \\
 5  & 7  &  10 &  15 &     &   \\
 15 & 20 &  27 &  37 &  52 &   \\
.   &  . &  .  & .   &  .  &  .   \\
\end{tabular}
\end{center}


Обозначим $A_{n,k}$~--- $k$-ый элемент в $n$-ой строке треугольника Белла.
По построению треугольника $A_{n,k} = A_{n,k-1}+ A_{n-1,k-1}$ при $k>1, n>1$,
$A_{n,1} =  A_{n-1,n-1}$ при $n>1$ 
и $A_{1,1}=1$.

\begin{lemma}  
$A_{n,k} =\sum\limits_{i=n-k+1}^{n} C(k-1,i-n+k-1) A_{i,1}.$
\end{lemma}
\begin{proof} Индукция по $k$. База. 
$A_{n,1} =\sum _{i=n}^{n}C(0,i-n)A_{i,1} = C(0,0)A_{n,1} = A_{n,1} $.
\end{proof}

\end{frame}


\begin{frame}
\frametitle{5.4.5. Треугольник Белла (3/4)}

\vspace{-2pt}
\begin{proof} (Продолжение)
Индукционный переход.
$A_{n,k} =  A_{n,k-1} + A_{n-1,k-1} = $\\[-2pt]
$=\sum\limits_{i=n-k+2}^{n} C(k-2,i-n+k-2)A_{i,1}  + 
\sum\limits_{i=n-k+1}^{n-1} C(k-2,i-n+k-1)A_{i,1} =$\\[-2pt]   
$=\sum\limits_{i=n-k+2}^{n-1}\left(C(k-2,i-n+k-2) +C(k-2,i-n+k-1) \right)A_{i,1} +$\\[-2pt] 
$ \phantom{\sum_{i=n-k+2}^{n-1}} +C(k-2,k-2)A_{n,1} +C(k-2,0)A_{n-k+1,1} = $\\[-2pt] 
$=\sum\limits_{i=n-k+2}^{n-1}      C(k-1,i-n+k-1)A_{i,1}+C(k-1,k-1)A_{n,1} +C(k-1,0)A_{n-k+1,1} =$\\[-2pt] 
$=\sum\limits_{i=n-k+1}^{n}C(k-1,i-n+k-1)A_{i,1}$.
\end{proof}

\begin{theorem}
$B(m)=A_{m+1,1} =A_{m,m} $.
\end{theorem}
 
\begin{proof}
Индукция по $m$. База $B(0) = 1 =A_{1,1}$ по построению треугольника. 
Пусть $B(m)=A_{m+1,1}$. Тогда
по лемме 
$A_{m+2,1} = A_{m+1,m+1} = 
\sum\limits_{i=0}^{m} C(m,i) A_{i+1,1} =$ \\ $= 
\sum\limits_{i=0}^{m}C(m,i) B(i) = B(m + 1)$. 
\end{proof} 

\end{frame}




\begin{frame}
\frametitle{5.4.5. Треугольник Белла (4/4)}
\vspace{-2pt}
Из теоремы непосредственно следует алгоритм.

\vspace{-2pt}
\begin{algorithmic}
\REQUIRE{ число $m$}
\ENSURE{ число Белла $B(m)$}
\STATE{\textbf{if $m=0$ then return $1$ end if}}
\COMMENT{\color{gray}\small  по определению }
\STATE{\textbf{$A:$ array $[1..m]$ of integer}}
\COMMENT{\color{gray}\small  рабочий массив, $A[k]$ хранит $A_{n,k}$ }
\STATE{$A[1] :=  1$}
\COMMENT{\color{gray}\small  $A_{1,1}= B(1) =1 $ }
\FOR{\textbf{$n$ from $2$ to $m$}}
\STATE{$t \Assign A[1]$}
\COMMENT{\color{gray}\small   переменная $t$ хранит $A_{n-1, k-1}$ }
\STATE{$ A[1] \Assign A[n - 1]$}
\COMMENT{\color{gray}\small  $A_{n,1} = A_{n-1,n-1} $ }
\FOR{\textbf{$k$ from $2$ to $n$}}
  \STATE{$s \Assign A[k]$ }
  \COMMENT{\color{gray}\small  переменная $s$ хранит $A_{n-1, k}$ }
  \STATE{$A[k] \Assign A[k-1] + t$}
  \COMMENT{\color{gray}\small   $A_{n,k} = A_{n,k-1}+ A_{n-1,k-1}$}
  \STATE{$ t \Assign s$} 
  \COMMENT{\color{gray}\small  $k\Assign k+1$}
\ENDFOR
\ENDFOR
\STATE{\textbf{return $A[m]$}}
\COMMENT{\color{gray}\small  число Белла }
\end{algorithmic}

\end{frame}

\begin{frame}
\frametitle{Онтология комбинаторного анализа}
\begin{center}
\includegraphics[height=7.8cm]{ODM5_4.png}
\end{center}
\end{frame}
\begin{frame}
\frametitle{5.5. Включения и исключения}
\subsection{5.5. Включения и исключения}
\label{5.5.}

Приведённые в предыдущих четырёх разделах формулы и алгоритмы дают
способы вычисления комбинаторных чисел для некоторых
распространённых комбинаторных конфигураций. Практические задачи
не всегда прямо сводятся к известным
комбинаторным конфигурациям. В таких случаях используются различные
методы свед$\acute{\text{е}}$ния одних комбинаторных конфигураций
к другим. В этом и двух следующих разделах рассматриваются три
наиболее часто используемых метода.


Мы начинаем с самого простого и прямолинейного,
но имеющего ограниченную область применения метода
\odmkey{включений и исключений}{Метод (method)!включений  и исключений \\(of inclusion and exclusion)}.

\medskip
\begin{tabular}{p{11mm}|p{65mm}|p{65mm}}
  № & Название параграфа
  & Ключевые термины и обозначения \\
  \hline
\DMref{5.5.1.}{5.5.1.} & Объединение конфигураций  &
{ }\\  
\DMref{5.5.2.}{5.5.2.}\RS & Формула включений и исключений &
{ }\\  
\DMref{5.5.3.}{5.5.3.} & Число булевых функций, существенно зависящих от всех своих переменных  &
{ }\\  
\DMref{5.5.4.}{5.5.4.} & Комбинации свойств  &
{ }\\  
\DMref{5.5.5.}{5.5.5.} & Беспорядки и субфакториал  &
{ $!n$}\\  
\end{tabular}  

\end{frame}


\begin{frame}
\label{5.5.1.}\subsubsection{5.5.1. Объединение конфигураций }\frametitle{5.5.1. Объединение конфигураций (1/2)}

\vspace{-0.25\baselineskip}

Часто комбинаторная конфигурация является объединением других,
число комбинаций в которых вычислить проще. В таком случае
требуется уметь вычислять число комбинаций в объединении.
В простых случаях формулы для вычисления очевидны:

$$\vert A\cup B\vert
 = \vert A\vert + \vert B \vert - \vert A \cap B \vert,$$ 
$$\vert A \cup B \cup C \vert = \vert A \vert + \vert B \vert +
\vert C \vert - \vert A \cap B \vert - \vert B \cap C \vert -
\vert A \cap C \vert + \vert A \cap B \cap C \vert.
$$

\medskip

\begin{center}
\includegraphics[width=130mm]{5_5_1.png}
\end{center}

\end{frame}

\begin{frame}\frametitle{5.5.1. Объединение конфигураций(2/2)}
	\begin{example}
		Сколько существует натуральных чисел меньших 1000,
		которые не делятся ни на 3,
		ни на 5, ни на 7? Всего чисел меньших тысячи, 999. Из них:
\label{key}
\begin{itemize}
\item[ ] $999:3=333$ делятся на 3,
\item[ ] $999:5=199$ делятся на 5,
\item[ ] $999:7=142$ делятся на 7,
\item[ ] $999:(3*5)=66$ делятся на 3 и на 5,
\item[ ] $999:(3*7)=47$ делятся на 3 и на 7,
\item[ ] $999:(5*7)=28$ делятся на 5 и на 7,
\item[ ] $999:(3*5*7)=9$ делятся на 3, на 5 и на 7.
\end{itemize}

Имеем: $999 - (333+199+142-66-47-28+9) = 457$.
\end{example}
\end{frame}


\begin{frame}\label{5.5.2.}\subsubsection{5.5.2. Формула включений и исключений }\frametitle{5.5.2. Формула включений и исключений (1/3)}


Следующая формула, известная как
\odmkey{формула включений и исключений}{Формула
(formula)!включений и исключений \\(inclusion-exclusion)},
позволяет вычислить мощность объединения нескольких множеств, если
известны их мощности и~мощности всех возможных пересечений.
\begin{theorem}
\begin{align*}
\left\vert\bigcup^{n}_{i=1} A_i\right\vert &=
\sum^{n}_{i=1} \vert A_i\vert-\sum_{1\le i<j\le n}\vert A_i\cap A_j\vert
+ \sum_{1\le i<j<k\le n}\vert A_i\cap A_j\cap A_k\vert - \dots\\
&+(-1)^{n-1}\vert A_1\cap\ldots\cap A_n\vert.
\end{align*}
\end{theorem}


\begin{proof}
Индукция по $n$. Для $n=2,3$ теорема проверена в предыдущем параграфе. Пусть
$$\left\vert\bigcup\limits_{i=1}^{n-1}A_i\right\vert=\sum\limits^{n-1}_{i=1}\vert
A_i\vert-\sum\limits_{1\leq i<j\leq n-1}\vert A_i\cap
A_j\vert+\ldots +(-1)^{n-2}\vert A_1\cap\ldots\cap A_{n-1}\vert.$$

\end{proof}
\end{frame}




\begin{frame}\frametitle{5.5.2. Формула включений и исключений (2/3)}
%{\small
\begin{proof} (Продолжение)
Заметим, что $\left(\bigcup\limits^{n-1}_{i=1}A_i\right)\cap A_n=
\bigcup\limits^{n-1}_{i=1}(A_i\cap A_n)$, и по индукционному
предположению
\begin{align*}
\left\vert\bigcup\limits^{n-1}_{i=1}(A_i\cap
A_n)\right\vert &=\sum\limits^{n-1}_{i=1}\vert A_i\cap A_n\vert-
\sum\limits_{1\leq i<j\leq n-1}\vert A_i\cap
A_j\cap A_n\vert+\dots+ \\
&+(-1)^{n-2}\vert A_1\cap\ldots\cap A_{n-1}\cap A_n\vert .
\end{align*}



Тогда
\begin{align*}
\left\vert\bigcup\limits^n_{i=1}A_i\right\vert  &=
\left\vert\left(\bigcup\limits^{n-1}_{i=1}A_i\right)\cup A_n\right\vert=
\left\vert\bigcup\limits^{n-1}_{i=1}A_i\right\vert+
\vert A_n\vert-
\left\vert\left(\bigcup\limits^{n-1}_{i=1}A_i\right)\cap A_n\right\vert= \\
%&=\left(\sum\limits^{n-1}_{i=1}\vert A_i\vert- \sum\limits_{1\leq
%i<j\leq n-1}\vert A_i\cap A_j\vert+\ldots+
%(-1)^{n-2}\vert A_1\cap\ldots\cap A_{n-1}\vert\right)+\\
%&+\vert A_n\vert-\left(\sum\limits^{n-1}_{i=1}\vert A_i\cap A_n\vert\right.
%   -\sum\limits_{1\leq i<j\leq n-1}\vert A_i\cap A_j\cap A_n\vert+\ldots+ \\
%&+(-1)^{n-2}\vert A_1\cap\ldots\cap A_{n-1}\cap A_n\vert\Biggr)=\\
%&=\left(\sum\limits^{n-1}_{i=1}\vert A_i\vert+
%\vert A_n\vert\right)-
%\left(\sum\limits_{1\leq i<j\leq n-1}\vert A_i\cap A_j\vert+
%\sum\limits^{n-1}_{i=1}\vert A_i\cap A_n\vert\right)+\ldots- \\
%&-(-1)^{n-2}\vert A_1\cap\ldots\cap A_{n-1}\cap A_n\vert \vphantom{\Biggr)}=\\
&=\sum\limits^n_{i=1}\vert A_i\vert- \sum\limits_{1\leq i<j\leq
n}\vert A_i\cap A_j\vert+\ldots+ (-1)^{n-1}\vert A_1\cap\ldots\cap
A_n\vert.
\end{align*}
\end{proof}
%}
\end{frame}

\begin{frame}\frametitle{5.5.2. Формула включений и исключений (3/3)}

%\vspace{-3pt}
\begin{example}
Докажем формулу для функции Эйлера $\varphi (n)$ (см.~\DMref{п.~2.4.9}{2.4.9.}) 
с помощью метода включений и исключений.

Пусть $n = p_1^{\alpha _1}p_2^{\alpha _2} \dots p_k^{\alpha _k}$ 
и $\varphi (n)$~--- количество чисел меньших 
$n$ и взаимно простых с ним. 
Пусть $A_i$ --- множество чисел, не превосходящих $n$ и 
делящихся на $p_i$. Тогда \\$\varphi (n) = n - |\bigcup\limits_{i = 1}^k A_i |$.
Заметим, что $|A_{i}| = \frac{n}{p_{i}}$
и
$|A_{i_1} \cap A_{i_2}\cap \dots \cap A_{i_s}| = \frac{n}{p_{i_1}p_{i_2}\dots p_{i_s}}$,
где подразумевается деление  нацело. 
Отсюда получаем\\ 
$\varphi (n) = n - (\frac{n}{p_1} + \frac{n}{p_2}+ \dots +\frac{n}{p_k}) + (\frac{n}{p_1p_2} + \frac{n}{p_2p_3}+ \dots +\frac{n}{p_{k-1} p_k})- \dots + (-1)^{k}\frac{n}{p_1p_2 \dots p_k}=$\\ $= n(1 - \frac{1}{p_1})(1 - \frac{1}{p_2})\dots(1 - \frac{1}{p_k})$, что согласуется с формулой, выведенной в \DMref{п.~2.4.9}{2.4.9.}.
\end{example}

\medskip
\begin{remark}
Обозначения сумм с неравенствами в пределах суммирования,
использованные в формулировке и доказательстве теоремы, являются
не более чем сокращённой формой записи кратных сумм. Например,
$
\sum_{1\leq i < j\leq n}$ означает 
$\sum_{i=1}^{n-1}\sum_{j=i+1}^{n}.
$
\end{remark}

\end{frame}




\begin{frame}
\label{5.5.3.}\subsubsection{5.5.3. Число булевых функций, существенно зависящих от всех своих переменных}
\frametitle{5.5.3. Число булевых функций, существенно зависящих от всех своих переменных}


%\vspace{-4pt}
Рассмотрим применение формулы включений и исключений.
 Пусть $p_n\Def\vert P_n\vert=2^{2^n}$~---
число всех булевых функций $n$~переменных,
а $\widetilde p_n$~--- число булевых функций, существенно
зависящих от всех $n$~переменных 
(\DMref{п.~3.1.2}{3.1.2.}).
Пусть $P^i_n$~--- множество булевых функций,
у которых переменная $x_i$ фиктивная (кроме $x_i$ могут быть и другие
фиктивные переменные). Имеем

\vspace{-14pt}
$$
\widetilde p_n=\vert P_n\setminus(P^1_n\cup\ldots\cup
P^n_n)\vert=\left\vert
P_n\setminus\bigcup\limits^n_{i=1}P^i_n\right\vert.
$$

\vspace{-2pt}
С другой стороны, $\vert P^i_n\vert=2^{2^{n-1}}$, более того,
$\vert P^{i_1}_n\cap\ldots\cap P^{i_k}_n\vert=2^{2^{n-k}}$.
 Следовательно,
$\widetilde p_n = $\\ 
$ = 2^{2^n} - \left(\sum\limits^n_{i=1}\vert
P^i_n\vert-   \sum\limits_{1\leq i<j\leq n}   \vert P^i_n\cap
P^j_n\vert+\ldots +(-1)^{n-1}\vert P^1_n\cap\ldots\cap
P^n_n\vert\right)=$\\
$=2^{2^n}-\left(C(n,1)2^{2^{n-1}}-C(n,2)2^{2^{n-2}}+\ldots
+(-1)^{n-1}C(n,n)2\right) =$\\
$=\sum\limits^n_{i=0}(-1)^iC(n,i)2^{2^{n-i}}.
$

\end{frame}


\begin{frame}
\label{5.5.4.}\subsubsection{5.5.4. Комбинации свойств }\frametitle{5.5.4. Комбинации свойств }

Достаточно часто комбинаторная задача
формулируется как задача подсчёта количества объектов, обладающих или не
обладающих определёнными \em{свойствами}.

Уточним постановку задачи. 
Пусть имеется множество объектов $M$, $|M|=m$,
и множество \odmkey{свойств}{Свойство (property)} $\Tuple{a}{1}{n}$,
то есть одноместных предикатов
$\Map{a_i}{M}{0..1}$. 
Обозначим $N(\Tuple{a}{i_1}{i_s})$~---
количество элементов, обладающих свойствами $\Tuple{a}{i_1}{i_s}$,\\
$N(\Tuple{\overline{a}}{j_1}{j_t})$~---
количество элементов, не обладающих свойствами $\Tuple{a}{j_1}{j_t}$. 	 

\begin{theorem}
$N(\Tuple{\overline{a}}{1}{n})= m-\sum\limits_{i=1}^n N(a_i) +
\sum\limits_{1\leq i<j\leq n} N(a_i,a_j) +\ldots +
(-1)^n N(\Tuple{a}{1}{n})$.
\end{theorem}
%\end{frame}
%
%
%\begin{frame}\frametitle{5.5.4. Комбинации свойств (2/2) }

\begin{proof}
Каждый одноместный предикат является характеристическим 
предикатом подмножества $M_i\Assign \Set{x\in M}{a_i(x)}$.
Поэтому $$N(\Tuple{\overline{a}}{1}{n})=
\left\vert \bigcap\limits_{i=1}^n \overline{M_i}\right\vert =
\vert M \vert - \left\vert \bigcup\limits_{i=1}^n M_i\right\vert $$
и далее по формуле включений и исключений.  
\end{proof}

%\begin{example}
%Симметричные пин-коды $xyyx$ и периодические пин-коды 
%$xyxy$ считаются ненадёжными.
%Сколько существует надёжных пин-кодов? Ответ:
%$10000 -(100+100) +10 = 9810$.
%\end{example}

\end{frame}


\begin{frame}\label{5.5.5.}\subsubsection{5.5.5. Беспорядки и субфакториал }\frametitle{5.5.5. Беспорядки и субфакториал (1/3)}
Перестановка без неподвижных точек называется \odmkey{беспорядком}{Беспорядок (disorder)}. 
Количество всех беспорядков $n$ предметов называется 
\odmkey{субфакториалом}{Субфакториал (sub-factorial)} числа $n$ и обозначается $!n$. 

\begin{theorem}
$!n=n!\sum\limits_{k=0}^n  (-1)^k/k! .$
\end{theorem}

\begin{proof}
Обозначим через $D(n,k)$~--- количество перестановок из $n$ элементов, 
в которых $k$ элементов~--- неподвижные точки. 
Существует $P(n-k)$ перестановок, в которых $k$ элементов остаются 
на своей позиции. Выбрать $k$ неподвижных элементов 
из $n$ можно $C(n,k)$ способами. 
Окончательно получаем $D(n,k) = C(n,k)P(n-k)$. 
Общее количество перестановок равняется $n!$. 
Воспользуемся формулой включений и исключений. Имеем:
$!n= $\\ $=n! -C(n,1)P(n-1)+C(n,2)P(n-2)-...+(-1)^k C(n,k)P(n-k)+\dots+(-1)^n C(n,n)= $ \\ $= n! -  n!(n-1)!/(n-1)! + n!(n-2)!/(n-2)!2! -\dots$ $+ 
(-1)^k  n!(n-k)!/(n-k)!k!  -
\dots$ $+(-1)^n  n!/(0!\cdot n!) = 
\sum_{k=0}^n(-1)^k   n!/k!$.
\end{proof}
\end{frame}


\begin{frame}\frametitle{5.5.5. Беспорядки и субфакториал (2/3)}
\vspace{-2pt}
\begin{remark}
Cумма в правой части 
формулы $!n=n!\sum_{k=0}^n \frac{(-1)^k}{k!}$
является частичной суммой 
разложения $e^{-1} =1-1/1!+1/2!-\ldots+(-1)^n/n!+\dots$. 
Поэтому субфакториал является ближайшим целым числом к числу $n!/e$.
\end{remark}

\begin{theorem}
$!n=(n-1)(!(n-1)+!(n-2))$.
\end{theorem}

\begin{proof}
$(n - 1)\left(!(n - 1) + !(n - 2)\right) =$ \\
$=(n - 1)\left((n - 1)!\sum\limits_{k=0}^{n-1} \frac{(-1)^k}{k!} + 
(n - 2)!\sum\limits_{k=0}^{n-2} \frac{(-1)^k}{k!} \right)=$ \\
$=(n-1)(n-2)!\left((n - 1)\sum\limits_{k=0}^{n-1} \frac{(-1)^k}{k!} + 
\sum\limits_{k=0}^{n-2} \frac{(-1)^k}{k!} \right)=$ \\
$=(n-1)!\left(n \sum\limits_{k=0}^{n-1} \frac{(-1)^k}{k!}
-\sum\limits_{k=0}^{n-1} \frac{(-1)^k}{k!} + 
\sum\limits_{k=0}^{n-2} \frac{(-1)^k}{k!} \right)=$ \\
$=(n-1)!\left( n \sum\limits_{k=0}^{n-1} \frac{(-1)^k}{k!}
- \frac{(-1)^{n-1}}{(n-1)!}  \right)=$ \\
$=(n-1)!\left(n \sum\limits_{k=0}^{n-1} \frac{(-1)^k}{k!}
+ \frac{n(-1)^n}{n(n-1)!}  \right)=n(n-1)!\sum\limits_{k=0}^{n} \frac{(-1)^k}{k!}=$ $!n$.
\end{proof}
\end{frame}


\begin{frame}\frametitle{5.5.5. Беспорядки и субфакториал (3/3)}

\vspace{-2pt}
\begin{corollary}
$ !n= !(n-1)\cdot n+ (-1)^n$.
\end{corollary}
\begin{proof}
По индукции. База: $!1=0$. Пусть $ !(n-1)= !(n-2)\cdot (n-1)+ (-1)^{n-1}$.
Тогда $!n=(n-1)\left(!(n - 1) + !(n - 2)\right) =
!(n-1)\cdot n - !(n-1)+!(n-2)\cdot n - !(n-2) =$ \\
$=!(n-1)\cdot n - (!(n-2)\cdot (n-1)+ (-1)^{n-1})+!(n-2)\cdot n - !(n-2) = $\\
$=!(n-1)\cdot n - !(n-2)\cdot n +!(n-2) + (-1)^{n}+!(n-2)\cdot n - !(n-2) =$\\  
$=!(n-1)\cdot n+ (-1)^n$.
\end{proof}

\begin{example}
Профессор собирается поручить четырём студентам $A, B, C, D$ 
проверить друг у друга
контрольную работу. Ни один студент не должен проверять свою работу. Сколько у профессора вариантов раздать работы на проверку?
 Из всех $P(4) = 4! = 24$ перестановок, подходят только $!4 = 9:$
$BADC$, $BCDA$, $BDAC$, $CADB$, $CDAB$, $CDBA$, $DABC$, $DCAB$, $DCBA$.
Любая другая перестановка будет иметь как минимум одну неподвижную точку.
\end{example}

%%%%%%%%%%%%%%%%%%%%%%%%%%%%%%%%%%%%%%%%%%%%%%%
% неплохо бы дать на графике рост субфакториала
\medskip
Для справки:
$!1 = 0$,
$!2 = 1$,
$!3 = 2$,
$!4 = 9$,
$!5 = 44$,
$!6 = 265$,
$!7 = 1~854$,
$!8 = 14~833$,
$!9 = 133~496$.


\end{frame}


\begin{frame}\frametitle{5.6. Формулы обращения}
\subsection{5.6. Формулы обращения}\label{5.6.}



Полезную, но очень специфическую группу приёмов
образуют различные способы преобразования
уже полученных комбинаторных выражений.
В этом разделе рассматривается один частный,
но важный случай.

\medskip
\begin{tabular}{p{11mm}|p{60mm}|p{70mm}}
  № & Название параграфа
  & Ключевые термины и обозначения \\
  \hline
\DMref{5.6.1.}{5.6.1.} & Теорема обращения  &
{ Комбинаторное уравнение}\\  
\DMref{5.6.2.}{5.6.2.} & Формулы обращения &
{ }\\  
\DMref{5.6.3.}{5.6.3.} & Формулы для чисел Стирлинга и Моргана  &
{ }\\  
\DMref{5.6.4.}{5.6.4.} & Полиномы с рациональными коэффициентами  &
{}\\  
\end{tabular} 

\end{frame}

\begin{frame}
\label{5.6.1.}\subsubsection{5.6.1. Теорема обращения}\frametitle{5.6.1. Теорема обращения}



Пусть $a_{n,k}$ и $b_{n,k}$~--- некоторые (комбинаторные) числа,
зависящие от параметров $n$ и $k$, причём $0\le k\le n$. Если
известно выражение чисел $a_{n,k}$ через числа $b_{n,k}$, то в
некоторых случаях можно найти и выражение чисел $b_{n,k}$ через
числа $a_{n,k}$, то~есть решить комбинаторное уравнение.

\begin{theorem} Пусть
$\A{ n}{\A{ k \le n}{ a_{n,k} =
\sum\limits^{n}_{i=0}\lambda_{n,k,i} b_{n,i}}}
$
и пусть\\
$
\E{\mu_{n,k,i}}{\A{ k\le n}{\A{ m\le n}
{\sum\limits^{n}_{i=0}\mu_{n,k,i} \lambda_{n,i,m} =
\begin{cases}
1, \quad m = k,\\
0,  \quad m\not= k
\end{cases}
}}}.
$\\
Тогда
$
\A{k\le n}{b_{n,k} = \sum\limits^{n}_{i=0}\mu_{n,k,i} a_{n,i}}.
$
\end{theorem}



\begin{proof}
$\sum\limits^n_{i=0}\mu_{n,k,i}a_{n,i}=$ $
\sum\limits^n_{i=0}\mu_{n,k,i}\left(\sum\limits^n_{m=0}
\lambda_{n,i,m}b_{n,m}\right)=$\\
$=\sum\limits^n_{m=0}\left(\sum\limits^n_{i=0}\mu_{n,k,i}
\lambda_{n,i,m}\right)b_{n,m}=b_{n,k}$.
\end{proof}



\end{frame}


\begin{frame}
\label{5.6.2.}\subsubsection{5.6.2. Формулы обращения  }\frametitle{5.6.2. Формулы обращения  (1/3)}



\vspace{-2pt}
Применение теоремы обращения предполагает отыскание для заданных
чисел $\lambda_{n,k,i}$ (коэффициентов комбинаторного уравнения)
соответствующих чисел $\mu_{n,k,i}$, удовлетворяющих условию
теоремы обращения. Особенно часто числами $\lambda_{n,k,i}$
являются биномиальные коэффициенты.

\begin{lemma}
$\displaystyle{
\sum\limits_{i=0}^{n} (-1)^{i-m} C(n,i)C(i,m) =
\begin{cases}
1, \quad m = n,\\
0,  \quad m<n.
\end{cases}}
$
\end{lemma}


\begin{proof}
Используя формулу 3 из \DMref{п.~5.3.1}{5.3.1.} и тот факт,
что $C(n-m,i-m)=0$ при $i<m$, имеем

\vspace{-18pt}
\begin{align*}
\sum\limits_{i=0}^{n} (-1)^{i-m} C(n,i)C(i,m) &=
\sum\limits_{i=0}^{n} (-1)^{i-m} C(n,m)C(n-m,i-m) =\\[-2pt]
&=
\sum\limits_{i=m}^{n} (-1)^{i-m} C(n,m)C(n-m,i-m) =\\[-2pt]
&=
C(n,m)\sum\limits_{i=m}^{n} (-1)^{i-m} C(n-m,i-m).
\end{align*}


\end{proof}

\end{frame}


\begin{frame}\frametitle{5.6.2. Формулы обращения (2/3)}

\begin{proof} (Продолжение)
Но при $m<n$ имеем

\vspace{-12pt}
$$
\sum\limits_{i=m}^{n} (-1)^{i-m} C(n-m,i-m) =
\sum\limits_{j=0}^{n-m} (-1)^{j} C(n-m,j) = 0,$$

\vspace{-6pt}
где $j\Assign i-m$.
С другой стороны, при $m=n$ имеем
$
C(n,m)\sum\limits_{i=m}^{n} (-1)^{i-m} C(n-m,i-m)=$\\
$=C(n,n)(-1)^{n-n}C(0,0) = 1$.
\end{proof}


%\pagebreak
%\setcounter{theorem}{0}
\begin{theorem}
Если $ a_{n,k} = \sum\limits^{k}_{i=0}C(k,i)b_{n,i}$, то
$b_{n,k} = \sum\limits^{k}_{i=0}(-1)^{k-i}C(k,i)a_{n,i}$.
\end{theorem}



\begin{proof}
Здесь $\lambda_{n,k,i} = C(k,i)$ и
$\mu_{n,k,i} = (-1)^{k-i}C(k,i)$.
При $k\le n$, $m\le n$ имеем:
$\sum\limits^{n}_{i=0}\mu_{n,k,i} \lambda_{n,i,m} =\sum\limits^{n}_{i=0} (-1)^{k-i}C(k,i) C(i,m) =$
(так как $C(k, i)=0$ при $i>k$) 
$= \sum\limits^{k}_{i=0} (-1)^{k - i + m - m}C(k,i) C(i,m)=
 (-1)^{k-m}\sum\limits^{k}_{i=0} (-1)^{i - m}C(k,i) C(i,m)=$\\
\textbf{$=$ if $k=m$ then $1$ else $0$ end if}.
\end{proof}
%

\end{frame}


\begin{frame}\frametitle{5.6.2. Формулы обращения 
 (3/3)}
\begin{theorem}
Если $a_{n,k} = \sum\limits^{n}_{i=k}C(i,k)b_{n,i}$, то
$b_{n,k} = \sum\limits^{n}_{i=k}(-1)^{i-k}C(i,k)a_{n,i}$.
\end{theorem}



\begin{proof}
Здесь $\lambda_{n,k,i} = C(i,k)$ и
$\mu_{n,k,i} = (-1)^{i-k}C(i,k)$.
При $k\le n$, $m\le n$ имеем:
$\sum\limits^{n}_{i=0}\mu_{n,k,i} \lambda_{n,i,m} =
\sum\limits^{n}_{i=0} (-1)^{i-k}C(i,k) C(m,i)=$ \\
$=\sum\limits^{n}_{i=0} (-1)^{i-k}C(m,i) C(i,k)=$
\textbf{if $k=n$ then $1$ else $0$ end if}.
\end{proof}


    \begin{digress}
Вообще говоря, формулы, подобные приведённым формулам обращения, имеют широкое использование в области 
алгебры,анализа и теории чисел 
(например, \em{формулы обращения Мёбиуса}), но такие применения выходят 
за рамки этого курса, поэтому здесь мы
ограничиваемся примерами, результаты
которых можно было бы получить и более элементарными способами.
Разбирая приведённые примеры,
следует иметь в виду,
что этими примерами область 
применения формул обращения 
отнюдь не исчерпывается.
    \end{digress}

\end{frame}


\begin{frame}
\label{5.6.3.}\subsubsection{5.6.3. Формулы для чисел Стирлинга и Моргана}
\frametitle{5.6.3. Формулы для чисел Стирлинга и Моргана}


В качестве примера использования формул обращения рассмотрим
получение явных формул для чисел Стирлинга и Моргана.
Рассмотрим множество функций $\Map{f}{A}{B}$, где $\vert A\vert=n$
и $\vert B\vert=k$. Число всех таких функций равно $k^n$. С другой
стороны, число функций $f$, таких, что $\vert f(A)\vert=i$, равно
$M(n,i)$, поскольку $M(n,i)$~--- это число сюръективных функций
$\Map{f}{1..n}{1..i}$. Но множество значений функции (при заданном
$i$) можно выбрать $C(k,i)$ способами. Поэтому
$
k^n=\sum\limits^k_{i=0}C(k,i)M(n,i).
$


Обозначив $a_{n,k}\Assign k^n$ и
$b_{n,i}\Assign M(n,i)$, имеем по первой теореме
\DMref{п.~5.6.2}{5.6.2.} 

$$M(n,k)=\sum\limits^k_{i=0}(-1)^{k-i}C(k,i)i^n.$$
Учитывая связь чисел Стирлинга и Моргана, окончательно имеем
$$
S(n,k)=\frac{1}{n!}\sum\limits^k_{i=0}(-1)^{k-i}C(k,i)i^n.
$$

\end{frame}

\begin{frame}
\label{5.6.4.}\subsubsection{5.6.4. Полиномы с рациональными коэффициентами}
\frametitle{5.6.4. Полиномы с рациональными коэффициентами (1/2)}

Комбинаторные числа~--- это целые числа по определению.
Заметим, что формулы, которыми выражаются комбинарные
комбинаторные числа, зачастую остаются вычислимыми
и при других, не обязательно целых значениях аргумента.
\begin{example} Функция
$\displaystyle{C(x,n)\Def\frac{x(x-1)\ldots(x-n+1)}{n!}}$ определена для любых вещественных $x$ и выражается
полиномом степени $n$ с рациональными коэффициентами.
\end{example} 

\medskip
Любой полином с целыми коэффициентами
принимает целые значения при целых значениях
аргумента~--- это очевидно. Если же коэффициенты
не целые, то значения полинома при целых аргументах
не обязаны быть целыми.
\begin{example} Полином $\frac{1}{2} x$ принимает целые значения
только при чётных целых $x$, а при 
нечётных~--- не принимает. 
\end{example} 

В тоже время существуют полиномы с нецелыми
коэффициентами, которые принимают
целые значения при целых аргументах.

\begin{example} Рассмотренная функция
$C(x,n)$ для любых целых $x$ принимает целые значения (\DMref{п.~5.3.6}{5.3.6.}).
\end{example} 
\end{frame}

\begin{frame}
\frametitle{5.6.4. Полиномы с рациональными коэффициентами (2/2)}
\begin{theorem}
Пусть $f(x)$~--- полином степени не более $n$ с рациональным коэффициентами. Тогда
    $$\A{a \in \Bbb{Z}}{f(a) \in \Bbb{Z}}   \Equiv
    f(0)\in \Bbb{Z} \And
    f(1)\in \Bbb{Z} \And
 \dots f(n) \in \Bbb{Z}.$$ 
  \end{theorem}

  \begin{proof}
\proofsection{$\Rightarrow$} Очевидно.
\proofsection{$\Leftarrow$}
 Заметим, что любой полином $f(x)$ 
 степени не больше $n$ с рациональными коэффициентами
     можно представить в виде 
$$f(x) = \sum_{i=0}^{n} g_{i}C(x,i) =
 \sum_{i=0}^{n} g_{i}\frac{x(x-1)\dots(x-i+1)}{i!}.$$ Тогда $\A{m\in 0..n}{f(m) = \sum_{i=0}^{m}
    g_{i}C(m,i)}$, и по теореме обращения\\
  $g_{m} = \sum_{i=0}^{m}(-1)^{m-i}f(i)C(m,i)$~--- целое число (т.~к. это линейная
  комбинация целых чисел). 
  Поскольку $\A{x \in \Bbb{Z}}{C(x,i) \in \Bbb{Z}}$ и
  $g_i \in \Bbb{Z}$, имеем 
  $ f(x)=\sum_{i=0}^{n}g_{i}C(x,i) \in \Bbb{Z}$
  для любого целого $x$.
  \end{proof}

\end{frame}


\begin{frame}\frametitle{ 5.7. Производящие функции}
\subsection{5.7. Производящие функции}\label{5.7.}

Для решения комбинаторных задач в некоторых случаях можно
использовать методы математического анализа. Разнообразие
применяемых здесь приёмов весьма велико и не может быть в полном
объёме рассмотрено в рамках этого курса. В данном разделе
рассматривается только основная идея метода производящих функций,
применение которой иллюстрируется несколькими простыми примерами. 
%Более
%детальное рассмотрение можно найти в литературе, например
%в~\cite{Sachkov2}.

\medskip
\begin{tabular}{p{11mm}|p{60mm}|p{70mm}}
  № & Название параграфа
  & Ключевые термины и обозначения \\
  \hline
\DMref{5.7.1.}{5.7.1.} & Формальные степенные ряды  &
{ Производящая функция; гармонический ряд}\\  
\DMref{5.7.2.}{5.7.2.} & Метод неопределённых коэффициентов &
{ }\\  
\DMref{5.7.3.}{5.7.3.} & Числа Фибоначчи  &
{ Формула Бине, золотое сечение}\\  
\DMref{5.7.4.}{5.7.4.} & Числа Каталана  &
{}\\  
\end{tabular} 


\end{frame}

\begin{frame}
\label{5.7.1.}\subsubsection{5.7.1. Формальные степенные ряды }\frametitle{5.7.1. Формальные степенные ряды (1/4)}

Пусть есть последовательность комбинаторных чисел $a_i$ и
последовательность функций $\varphi_i(x)$. Рассмотрим формальный
ряд
$$
\cal F(x)\Def\sum_{i}a_i\varphi_i(x).
$$



Выражение $\cal F(x)$ называется
\odmkey{производящей функцией}{Функция
(function)!производящая (course-of-value)}
(для заданной последовательности
комбинаторных чисел $a_i$ относительно
заданной последовательности функций $\varphi_i(x)$).



Обычно используют $\varphi_i(x)\Def x^i$ или
$\varphi_i(x)\Def x^i/i!$.



\begin{example}
Из формулы бинома Ньютона при $y = 1$ имеем
$$
(1+x)^n=\sum\limits_{i=0}^{n} C(n,i)x^i.
$$
Таким образом,
$(1+x)^n$ является производящей функцией для биномиальных коэффициентов.
\end{example}
\end{frame}

\begin{frame}
\frametitle{5.7.1. Формальные степенные ряды (2/4)}

\begin{lemma}
$V (m, n) = V (m-1, n) + V (m, n-1)$.
\end{lemma}

\begin{proof}
$V (m, n) = C (n + m - 1, n) = C (n + m- 2, n) + C (n + m - 2, n- 1) =$ 
$ = V (m - 1, n) + V (m, n - 1)$.
\end{proof}

\begin{example}
 Производящая функция для $V(m,n)$  при заданном $m$.

$\varphi_m(x) = \sum_{n=0}^\infty V(m,n) x^n  =$
$ 1 + \sum_{n=1}^\infty V(m,n) x^n   =$\\[2pt]
$= 1 + \sum_{n=1}^\infty V(m-1,n) x^n  + 
\sum_{n=1}^\infty V(m,n-1) x^n  =$\\[2pt] 
$=1+ \sum_{n=1}^\infty V(m-1,n) x^n  
+ x\sum_{n=1}^\infty V(m,n-1) x^{n-1}  =$\\[2pt]
$=\sum_{n=0}^\infty V(m-1,n) x^n  + x\sum_{n=0}^\infty V(m,n) x^n  =
\varphi_{m-1}(x) + x \varphi_m(x)$.

\smallskip
Получили, что 
$\varphi_m(x)  =$ $1/(1-x) \varphi_{m-1}(x) =$\\[2pt]
$= 1/ (1-x)^{2} \varphi_{m-2}(x)= \dots$
$= 1/ (1-x)^m  \varphi_0 (x) = 1/ (1-x)^m$   
так как\\[6pt]
$\varphi_0 (x) = \sum_{n=0}^\infty V(0,n) x^n 
= V(0,0) + \sum_{n=1}^\infty V(0,n) x^n   = 1 + 0 = 1$.

\smallskip
В итоге, $\varphi_m(x)  = 1/ (1-x)^m$. 
\end{example}
\end{frame}

\begin{frame}
\frametitle{5.7.1. Формальные степенные ряды (3/4)}
Степенные ряды рассматриваются здесь именно
 как \em{формальные} суммы,
как конечные, так и бесконечные, 
причём вопросы сходимости рядов
 и аналитичности функций не ставятся.
 
 \begin{example}
 Пусть 
 $\cal{F}(x)=1+x+x^2+\dots=\sum_{i=0}^\infty x^i$.
 Тогда
  $\cal{F}(x)=1+x+x^2+\dots=$ \\$=1+x(1+x+x^2+\dots)=1 + x\cal{F}(x)$.
Решая полученное функциональное уравнение
относительно $\cal{F}$ заключаем,
 что $\displaystyle{\cal{F}(x)=1+x+x^2+\dots=\sum_{i=0}^\infty x^i=
  \frac{1}{1-x}}$.
В полученное выражение бессмысленно подставлять
числовые значения --- это \em{не} числовое равенство!
\end{example} 

\medskip
Производящие функции используются как
 удобный <<контейнер>>, хранящий информацию о комбинаторных числах в их совокупности.

 \begin{example}
 Определим производящую функцию натурального ряда.
Пусть $\cal{F}(x)=1+2x+$ 
$+3x^2+\dots=
\sum_{i=0}^\infty (i+1)x^i$. 
Тогда
$\cal{F}(x)=\sum_{i=0}^\infty (i+1)x^i
=\sum_{i=0}^\infty ix^i+\sum_{i=0}^\infty x^i$.
Но $\sum_{i=0}^\infty x^i=
  \frac{1}{1-x}$, 
  а $\sum_{i=0}^\infty ix^i=\sum_{i=1}^\infty ix^i=x\sum_{i=1}^\infty ix^{i-1}=
  x\sum_{i=0}^\infty (i+1)x^i$, 
  поэтому \\$\displaystyle{\cal{F}(x)=x\cal{F}(x)+1/(1-x)}$, 
  и значит $\displaystyle{1+2x+3x^2+\dots = 1/(1-x)^2}$. 
 \end{example} 
\end{frame}

\begin{frame}
\frametitle{5.7.1. Формальные степенные ряды (4/4)}
Производящие функции для последовательности единиц (или любых других констант) и натурального ряда найдены составлением и решением 
функциональных уравнений относительно
неизвестной производящей функции. 
Составление и решение функциональных уравнений~--- это отнюдь не единственный способ нахождения производящих функций. Не возбраняется использовать любые средства математического анализа. 

\begin{example}
Определим производящую
функцию \odmkey{гармонического ряда}{Ряд (series)!гармонический (harmonic)}
$\displaystyle{\cal{F}(x)=\sum_{n=1}^\infty \frac{1}{n} x^n}$. \\
Очевидно, что 
$\displaystyle{\sum_{n=1}^\infty \frac{1}{n} x^n =
\sum_{n=0}^\infty \frac{1}{n+1} x^{n+1}}$.
Заметим, что функция
$\frac{1}{n+1} x^{n+1}$
является первообразной функции
$ x^{n}$.
Напомним, что
производящая функция суммы степеней
определяется следующим образом: 
$\displaystyle{\sum_{n=0}^{\infty} x^n =
 1/(1-x)}$.
Поэтому
$$
\cal{F}(x)=
\sum_{n=0}^\infty \frac{1}{n+1} x^{n+1} =
\sum_{n=0}^\infty \int x^{n}dx =
\int \left(\sum_{n=0}^\infty x^{n}\right)dx =
\int \frac{dx}{1-x} = \ln\frac{1}{1-x}.  
$$
\end{example}   

\end{frame}


\begin{frame}\label{5.7.2.}\subsubsection{5.7.2. Метод неопределённых коэффициентов }\frametitle{5.7.2. Метод неопределённых коэффициентов (1/2)}

%\vspace{-4pt}
Из математического анализа известно, что если

%\vspace{-12pt}
$$
\cal F(x) = \sum_{i}a_i\varphi_i(x)\quad\mbox{и}\quad
\cal F(x) = \sum_{i}b_i\varphi_i(x),
$$

%\vspace{-4pt}
то $\A{i}{a_i=b_i}$ (для рассматриваемых здесь систем функций $\varphi_i$).

В качестве примера применения производящих функций докажем теорему 4 из \DMref{п.~5.3.1}{5.3.1.}.




\begin{theorem}
$ \displaystyle{
C(2n,n) = \sum^{n}_{k=0}C(n,k)^2.}
$
\end{theorem}
\begin{proof}
Имеем: $(1+x)^{2n}=(1+x)^n(1+x)^n$. Следовательно,

\vspace{-12pt}
$$\sum\limits^{2n}_{i=0}C(2n,i)x^i=
\sum\limits^n_{i=0}C(n,i)x^i\cdot\sum\limits^n_{i=0}C(n,i)x^i.$$

%\vspace{-4pt}
Приравняем коэффициент при $x^n$:

\vspace{-12pt}
$$
C(2n,n)=\sum\limits^n_{k=0}C(n,k)C(n,n-k)=\sum\limits^n_{k=0}C(n,k)^2.
$$

\vspace{-2\baselineskip}
\end{proof}


\end{frame}

\begin{frame}\frametitle{5.7.2. Метод неопределённых коэффициентов (2/2)}
	
	\begin{example} Какие числа можно представить суммой степеней числа $2$  
	и сколькими способами (\DMref{п.~1.3.2}{1.3.2.})?
Рассмотрим произведение $$G(x)= \prod_{n=0}^k (1+x^{2^n})=(1+x)(1+x^2)\dots(1+x^{2^k}).$$ 
После раскрытия скобок и приведения подобных $G(x)$ представляется в виде суммы 
$G(x) = \sum_{n=0}^{2^{k+1}-1} g_n x^n$, в которой $x^n$ получается как произведение каких-то одночленов $x^{2^m}$, то есть коэффициент $g_n$~---  это количество различных представлений числа $n$ в виде суммы некоторых степеней $2$ (не исключая случая $g_n=0$).
Умножим теперь $G(x)$ на $(1-x)$. По формуле для разности квадратов получаем 
$(1-x)(1+x)(1+x^2)\dots(1+x^{2^k}) =$\\ $=(1-x^2)(1+x^2)\dots(1+x^{2^k})  =...= 1-x^{2^{k+1}}$. Имеем 
$G(x)=(1-x^{2^{k+1}})/(1-x)$. Cумма геометрической прогрессии:
$1+x+x^2+\dots+x^{2^{k+1}-1}=(1-x^{2^{k+1}})/(1-x) =G(x)=$ \\$ = \sum_{n=0}^{2^{k+1}-1} g_n x^n$.
По методу неопределённых коэффициентов 
получаем $\A{n} {g_n = 1}$, что означает, любое число $n$ можно представить в виде линейной комбинации степеней числа 2 и притом единственным образом. 
\end{example}
	
	
\end{frame}



\begin{frame}\label{5.7.3.}\subsubsection{5.7.3. Числа Фибоначчи }\frametitle{5.7.3. Числа Фибоначчи (1/4)}

\vspace{-2pt}
\odmkey{Числа Фибоначчи}{Число (number)!Фибоначчи (Fibonacci)}
%\label{s573}%
%\footnote{\;Фибоначчи (Ленардо Пизанский, 1180--1228).} 
$F(n)$
определяются следующим образом:
$$F(0)\Def 1,\; F(1)\Def 1,\; F(n+2)\Def F(n+1)+F(n).
$$
Найдём выражение для $F(n)$ через $n$. Для этого найдём
производящую функцию $\varphi(x)$ для последовательности чисел
$F(n)$. Имеем
\begin{align*}
\varphi(x)
     & = \sum_{n=0}^{\infty} F(n)x^n = F(0)x^0 + F(1)x^1 +
         \sum_{n=2}^{\infty}\left(F(n-2)+F(n-1)\right)x^n = \\
     & = 1 + x + \sum_{n=2}^{\infty} F(n-2)x^n + \sum_{n=2}^{\infty} F(n-1)x^n = \\
     & = 1 + x + x^2\sum_{n=2}^{\infty} F(n-2)x^{n-2}+x\sum_{n=2}^{\infty} F(n-1)x^{n-1} =\\
     & = 1\! +\! x\! +\! x^2\sum_{n=0}^{\infty} F(n)x^{n}+
                 x\left(\sum_{n=0}^{\infty} F(n)x^{n}-F(0)\right) = 
      1\! +\! x\! +\! x^2 \varphi(x) + x\left(\varphi(x)-1\right).
\end{align*}
\end{frame}


\begin{frame}\frametitle{5.7.3. Числа Фибоначчи (2/4)}
Решая это функциональное уравнение относительно $\varphi(x)$,
получаем, что
$$
\varphi(x) = \frac{1}{1-x-x^2}.
$$
Последнее выражение нетрудно разложить в ряд по степеням $x$.
Действительно, уравнение $1-x-x^2=0$ имеет корни $x_1= - (1+\sqrt{5})/2$
и $x_2= - (1-\sqrt{5})/2$, причём, как нетрудно убедиться,
$$
1-x-x^2= (1- ax)(1-bx), \quad \text{где}\ a = \frac{1+\sqrt{5}}{2},\quad
b = \frac{1-\sqrt{5}}{2}.
$$
Далее,
$$
\frac{1}{(1- ax)(1-bx)} = \frac{\alpha}{1-ax} + \frac{\beta}{1-bx},
\quad \text{где}\ \alpha = \frac{a}{a-b},\quad \beta = \frac{-b}{a-b}.
$$
\end{frame}


\begin{frame}\frametitle{5.7.3. Числа Фибоначчи (3/4)}
Из математического анализа известно, что для малых
(по абсолютной величине) $x$
$$
\frac{1}{1-\gamma x} =\sum_{n=0}^{\infty} \gamma^n x^n.
$$
Таким образом,
\begin{align*}
\varphi(x) & = \frac{\alpha}{1-ax}+\frac{\beta}{1-bx} =
           \alpha\sum_{n=0}^{\infty} a^n x^n +
               \beta\sum_{n=0}^{\infty} b^n x^n = \\
           & = \sum_{n=0}^{\infty} \frac{a}{a-b}a^n x^n -
               \sum_{n=0}^{\infty} \frac{b}{a-b} b^n x^n =
           \sum_{n=0}^{\infty} \frac{a^{n+1} - b^{n+1}}{a-b} x^n.
\end{align*}
Окончательно получаем соотношение, известное как \odmkey{формула Бине}{Формула (formula)!Бине (Binet)}
$$
F(n) = \frac{a^{n+1} - b^{n+1}}{a-b} = \frac{1}{\sqrt{5}}
\left(\left({\frac{1+\sqrt{5}}{2}}\right)^{n+1}-
      \left({\frac{1-\sqrt{5}}{2}}\right)^{n+1}
\right).
$$
%\footnote{\;Бине, Жак Филипп Мари
%(1786--1856).}.
\end{frame}

\begin{frame}\frametitle{5.7.3. Числа Фибоначчи (4/4)}
	
\begin{example}
Число $F_{n+2}=F_{n+1}+F_n$ является количеством двоичных последовательностей длины n, в которых нет двух $1$, находящихся на соседних позициях.
Доказательство по индукции.
База: $n = 1$, $F_3 = 2$, 
последовательности $0$ и $1$;
$n = 2$, $F_4 = 3$, последовательности $00$, $01$, $10$.
Индукционный переход: пусть $F_{n+2}$~--- число способов построить требуемые последовательности длины $n$.
Построим последовательность длины $n+1$: возьмем все такие последовательности длины $n$ и допишем к ним $0$ слева, не нарушая условие. Также мы можем взять все последовательности длины $n-1$ и дописать к ним $10$ слева. Таким образом, количество способов получить последовательность длины $n+1$ будет равно
  $F_{n+2}+F_{n+1} =F_{n+3}$.
\end{example}

\medskip
\begin{digress}
Число $\Phi \Def (1+\sqrt{5})/2\approx 1,618$ принято называть
\odmkey{золотым сечением}{Золотое сечение (golden ratio)},
поскольку оно является решением пропорции
$\displaystyle{\frac{a}{b}=\frac{b}{a+b}}$~---
меньшее так относится к большему, как большее к сумме. 
Это число известно ещё с древних времен и занимает важное место в науке и искусстве.
\end{digress}

%
%\medskip
%{\color{red} Контрольный вопрос: что такое Фибоначчиева система счисления и чем она может быть полезна?   
%}
	
\end{frame}


\begin{frame}\label{5.7.4.}\subsubsection{5.7.4. Числа Каталана }\frametitle{5.7.4. Числа Каталана (1/5)}
\odmkey{Числа Каталана}{Число (number)!Каталана (Catalan)} $C(n)$ можно
определить следующим образом:

\vspace{-8pt}
$$
C(0)\Def 1, \quad C(n)\Def \sum_{k=0}^{n-1} C(k)C(n-k-1).
$$

\vspace{-4pt}
Числа Каталана используются при решении различных
комбинаторных задач.

\begin{example}
Пусть $\langle S,\ast\rangle$~--- полугруппа (\DMref{п.~2.2.1}{2.2.1.}) и нужно вычислить
элемент $s_0\ast\dots\ast s_n$. Сколькими способами это можно
сделать (см. замечание \DMref{п.~3.2.2}{3.2.2.})? То есть сколькими способами можно расставить скобки,
определяющие порядок вычисления выражения?\\ Обозначим число
способов $C(n)$. Ясно, что $C(0)=1$. При любой схеме вычислений на
каждом шаге выполняется некоторое вхождение операции $\ast$ над
соседними элементами и результат ставится на их место. Пусть
последним выполняется то вхождение операции $\ast$, которое имеет
номер $k$ в исходном выражении. При этом слева от выбранного
вхождения знака $\ast$ находилось и было выполнено $k-1$ знаков
операции $\ast$, а справа, соответственно, $(n-1)-(k-1)=n-k$
знаков операции $\ast$. Тогда ясно, что\\
$
C(n)=\sum_{k=1}^{n}C(k-1)C(n-k)=\sum_{k=0}^{n-1} C(k)C(n-k-1),
$
и ответом на поставленный вопрос является число Каталана $C(n)$.
\end{example}

\end{frame}


\begin{frame}\frametitle{5.7.4. Числа Каталана (2/5)}
Числа Каталана выражаются через биномиальные коэффициенты.
Получим это выражение, используя метод производящих функций.


\begin{theorem}
$\displaystyle{C(n)=\frac{C(2n,n)}{n+1}}$.
\end{theorem}
\begin{proof}
Найдём производящую функцию для чисел Каталана
$\varphi(x)=\sum\limits_{n=0}^{\infty} C(n) x^n $, $x \neq 0.$
Для этого рассмотрим квадрат этой функции:
\begin{align*}
\varphi^2(x) &= \left(\sum_{n=0}^{\infty} C(n) x^n\right)^2 =
  \sum_{m=0}^{\infty} C(m) x^m \cdot \sum_{n=0}^{\infty} C(n) x^n =\\
&= \sum_{m,n=0}^{\infty} C(m)C(n) x^{m+n} =
   \sum_{n=0}^{\infty}\sum_{k=0}^{n} C(k) C(n-k) x^n =
 \sum_{n=0}^{\infty} C(n+1) x^n=\\
&=   \sum_{n=0}^{\infty} C(n+1) \frac{x^{n+1}}{x} = 
 \frac{1}{x}\left(\sum_{n=0}^{\infty} C(n) x^n - C(0)\right)=
   \frac{1}{x}\left(\varphi(x)- 1\right).
\end{align*}

\vspace{-16pt}
\end{proof}
\end{frame}

\begin{frame}
\frametitle{5.7.4. Числа Каталана (3/5)}
\begin{proof} (продолжение)
Решая уравнение $\varphi(x)=x\varphi^2(x)+1$ относительно функции
$\varphi(x)$, имеем
$$
\varphi(x)=\frac{1\pm\sqrt{1-4x}}{2x}.
$$
Обозначим $f(x)\Assign\sqrt{1-4x}$ и разложим
$f(x)$ в ряд по формуле Тейлора:
$$
f(x)=f(0)+\sum_{k=1}^{\infty}\frac{f^{(k)}(0)}{k!}x^k.
$$
Имеем
$\displaystyle{\frac{d^k\hfill}{dx^k}(1-4x)^{\frac{1}{2}} =\frac{1}{2}\cdot\left(\frac{1}{2}-1\right)\cdot\cdots\cdot
\left(\frac{1}{2}-k+1\right)\cdot(1-4x)^{\frac{1}{2}-k}\cdot (-4)^k =}$\\
$\displaystyle{= - 2^k\cdot 1\cdot 3\cdot\cdots\cdot(2k-3)\cdot(1-4x)^{\frac{1}{2}-k} = - 2^k\cdot (2k-3)!!\cdot(1-4x)^{\frac{1}{2}-k}=}$\\
$\displaystyle{= - 2^k\frac{(2k-3)!}{2^{k-2}(k-2)!} (1-4x)^{\frac{1}{2}-k}= -2^2 \frac{(2k-2)!(k-1)(k-1)!}{(2k-2)(k-1)!(k-1)!}(1-4x)^{\frac{1}{2}-k} =}$
\\$\displaystyle{= -2(k-1)!\frac{(2k-2)!}{(k-1)!(k-1)!}(1-4x)^{\frac{1}{2}-k} =} -2(k-1)!C(2k-2,k-1)(1-4x)^{\frac{1}{2}-k}.$
\end{proof}
\end{frame}

\begin{frame}\frametitle{5.7.4. Числа Каталана (4/5)}
\begin{proof}  (продолжение)
Таким образом,
$$
f^{(k)}(0)=-2(k-1)!C(2k-2,k-1), \quad f(x)=1-2\sum_{k=1}^{\infty}\frac{1}{k}C(2k-2,k-1)x^k.
$$
Подставляя выражение $f(x)$ в формулу для $\varphi(x)$, следует
выбрать знак <<минус>> перед корнем, чтобы удовлетворить условию
$C(0)=1$. Окончательно имеем
\begin{align*}
\sum_{n=0}^{\infty}C(n)x^n &= \varphi(x) 
= \frac{1-f(x)}{2x} =\frac{1-1+2\sum_{n=1}^{\infty} \frac{1}{n} C(2n-2,n-1) x^n}{2x} =\\
&= \sum_{n=1}^{\infty} \frac{1}{n} C(2(n-1),n-1) x^{n-1} =
   \sum_{n=0}^{\infty} \frac{C(2n,n)}{n+1} x^n,
\end{align*}
 и по методу неопределённых коэффициентов $C(n)=C(2n,n)/(n+1)$.
\end{proof}
\end{frame}


\begin{frame}\frametitle{5.7.4. Числа Каталана (5/5)}
В примере в начале параграфа показано, что число способов составить скобочную последовательность длины $2n$ равно $C(n)$. Рассмотрим взаимосвязь с другими задачами, где используются числа Каталана.
\begin{example}
Сколько существует монотонных маршрутов из левого нижнего угла квадрата $n \times n$ в правый верхний угол, которые идут по линиям сетки вверх или вправо и не заходят выше диагонали?
Решение: проведём биекцию с правильными скобочными последовательностями длины $2n$: открывающая скобка~--- горизонтальная линия, 
закрывающая скобка~--- вертикальная линия.
	 Очевидно, что соответствующий правильной скобочной последовательности путь не пересечет диагональ.
	На рисунке представлены возможные монотонные пути для случая $n = 3$.
	%\includegraphics[height = 20mm]{5_7_4.png}
	\begin{figure}
	\includegraphics[height = 20mm]{5_7_4(2).png}
	%\caption{Возможные монотонные пути для случая $n=3$}{}
	\end{figure}
\end{example}
	
\end{frame}

\begin{frame}
\frametitle{Онтология комбинаторного анализа}
\begin{center}
\includegraphics[height=7.8cm]{ODM5_5.png}
\end{center}
\end{frame}


% ===== tex/DM2022_6.tex =====
\section{6. Кодирование информации}

\begin{frame}
\frametitle{6. Кодирование информации (1/6)}
Что есть \odmkey{информация}{Информация (information)}
объяснить столь же нелегко,
как исчерпывающим образом ответить на вопрос 
<<что есть \em{материя},
\em{вещество} или \em{энергия}>>.
Неспособность ответить на этот вопрос
не препятствует нам жить в материальном мире,
пользоваться веществом и энергией.
Подобным образом мы с успехом обрабатываем информацию,
хотя и не можем однозначно определить, что же это такое. 

\smallskip
Устройства для обработки информации
известны человечеству с незапамятных 
времён~--- абак, счёты, ...
За исторически кратчайшее время 
компьютеры и информационные технологии стремительно
прошли путь от экспериментальных
академических экспонатов до устройств,
которые составляют основу современной
цивилизации.

\smallskip
Здесь 
информация трактуется с прикладных,
инженерных позиций. 
Для технического использования информации
необходимо, во-первых, уметь её численно измерять,
и, \\во-вторых, выработать эмпирически подтвержденные
приёмы и методы манипулирования информацией.
Эти приёмы и методы имеют собирательное название
\odmkey{теория информации}{Теория (theory)!информации (of information)}.

\end{frame}

\begin{frame}
\frametitle{6. Кодирование информации (2/6)}
Таким образом, центральными понятиями обсуждаемой теории информации
являются \odmkey{количество информации}{Количество (quantaty)!информации (of information)} и 
\odmkey{кодировние информации}{Кодирование (coding)!информации (of information)}, то есть представление информации в форме, пригодной для обработки. 

\smallskip
Вопросы кодирования информации издавна играли заметную роль как в математике, так и в жизни.

\begin{examples}
\begin{enumerate}
\item \odmkey{Десятичная позиционная система счисления} {Система
счисления (number system)!позиционная (positional)!десятичная
(decimal)}~--- это универсальный способ кодирования чисел, в том
числе натуральных. 
\item \odmkey{Римские цифры}{Цифры (numerals)!римские
(roman)}~--- другой способ кодирования небольших натуральных
чисел, причём гораздо более наглядный и естественный: палец~---~I,
пятерня~---~V, две пятерни~---~X.
%\footnote{\;Плотники часто
%маркируют брёвна сруба римскими цифрами, потому что их легко
%вырубить топором.}
Однако при этом способе кодирования трудно
выполнять арифметические операции над большими числами, поэтому он
был вытеснен десятичной позиционной  системой. 
\item \odmkey{Декартовы
координаты}{Координаты (coordinates)!декартовы (Cartesian)}~--- способ кодирования геометрических объектов
числами.
\end{enumerate}
\end{examples}

\end{frame}

\begin{frame}
\frametitle{6. Кодирование информации (3/6)}
Ранее средства кодирования играли
вспомогательную роль и не рассматривались
как отдельный предмет математического изучения,
но с появлением компьютеров ситуация радикально изменилась.
Кодирование буквально пронизывает
информационные технологии и является центральным
вопросом при решении самых разных
задач программирования. 

\begin{examples}
\begin{enumerate}
\item
Представление данных произвольной природы
 (например чисел, текста, графики, звуков, видео) в памяти компьютера.
\item
Защита информации от несанкционированного доступа, искажения или непреднамеренной потери.
\item
обеспечение помехоустойчивости при передаче
данных по каналам связи с естественными или искусственными помехами;
\item сжатие информации в базах данных.
\end{enumerate}
\end{examples}

\smallskip
\begin{remark}
Само составление текста программы
часто и совершенно справедливо
называют кодированием.
\end{remark}
 
\end{frame}

\begin{frame}
\frametitle{6. Кодирование информации (4/6)}

Не ограничивая общности, задачу кодирования можно сформулировать
следующим образом. 

Пусть заданы алфавиты $A=\{\Tuple{a}{1}{n}\}$,
$B=\{\Tuple{b}{1}{m}\}$ и~функция $\Map{F}{A^*}{B^*}$, причём
$\Dom F = S$, где $S$~--- некоторое множество слов в алфавите $A$,
$S\subset A^*$. Тогда функция $F$ называется
\odmkey{кодированием}{Кодирование (coding)}, элементы множества
$S$~--- \odmkey{сообщениями}{Сообщение (message)}, а~элементы
$\beta=F(\alpha)$, $\alpha\in S$, $\beta\in B^*$~---
\odmkey{кодами}{Код (code)!сообщения (of message)}
(соответствующих сообщений). Обратная функция $F^{-1}$ (если она
существует!) называется \odmkey{декодированием}{Декодирование
(decoding)}. Если $|B|=m$, то $F$ называется \odmkey{m-ичным
кодированием}{Кодирование (coding)!$m$-ичное ($m$-ary)}. Наиболее
распространенный случай $B = E_2 = \{0,1\}$~---
\odmkey{двоичное}{Кодирование (coding)!двоичное (binary)}
кодирование. 
Именно этот случай рассматривается в последующих
разделах; слово <<двоичное>> опускается.

\smallskip
Типичная задача теории
кодирования формулируется следующим образом: при заданных
алфавитах $A$, $B$ и множестве сообщений $S$ найти 
кодирование $F$, которое обладает определёнными свойствами (то
есть удовлетворяет заданным ограничениям) и оптимально в некотором
смысле. Критерий оптимальности, как правило, связан с минимизацией
длин кодов. 
\end{frame}

\begin{frame}
\frametitle{6. Кодирование информации (5/6)}

Свойства, которые требуются от кодирования, бывают
самой разнообразной природы. Здесь рассматриваются три наиболее 
востребованных на практике свойства.
\begin{itemize}
\item Существование декодирования, или
\odmkey{однозначность}{Кодирование (coding)!однозначное (unique)}
кодирования: функция кодирования $F$ обладает тем свойством, что
$\alpha_1\ne \alpha_2\Impl F(\alpha_1)\ne F(\alpha_2)$. Это очень
естественное свойство, однако даже оно требуется не всегда.
Например, трансляция программы на языке высокого уровня в машинные
команды~--- это кодирование, для которого не требуется
однозначного декодирования. \item
\odmkey{Помехоустойчивость}{Помехоустойчивость (noise-immunity)},
или исправление ошибок: продолжение функции декодирования $F^{-1}$
обладает тем свойством, что $F^{-1}(\beta)= F^{-1}(\beta^\prime)$,
где $\beta \in \Im F$, \\$\beta^\prime \in B^*\setminus\Im F$,
 если $\beta^\prime$
в определённом смысле близко к $\beta$ (см. раздел~\DMref{6.4}{6.4.}).
\item
Заданная \em{сложность} (или простота) кодирования и декодирования.
Например, в криптографии изучаются такие способы кодирования, при
которых функция $F$ вычисляется просто, но определение значения
функции $F^{-1}$ требует очень сложных вычислений
(см. \DMref{п.~6.5.7}{6.5.7.}).
\end{itemize}
\end{frame}

\begin{frame}
\frametitle{6. Кодирование информации (6/6)}

Большое значение для задач кодирования имеет природа множества
сообщений~$S$. При одних и тех же алфавитах $A$, $B$ и требуемых
свойствах кодирования $F$ оптимальные решения для разных множеств сообщений $S$ могут
кардинально различаться. Для описания множества $S$ (как правило,
очень большого или бесконечного) применяются различные методы:
\begin{itemize}
\item
теоретико-множественное описание, например,
$S=\Set{\alpha}{\alpha\in A^*\And\vert\alpha\vert=n}$;
\item
вероятностное описание, например, $S=A^*$, и заданы
вероятности $p_i$ появления букв $a_i$ в сообщении,
$\sum_{i=1}^{n} p_i = 1$;
\item
логико-комбинаторное описание, например,
$S$ задано порождающей формальной грамматикой.
\end{itemize}

В этой главе сначала вводятся базовые понятия теории информации,
а затем рассматривается несколько наиболее важных задач
теории кодирования и демонстрируется применение большей части
вышеупомянутых методов.
\end{frame}

\subsection{6.1 Информация и энтропия}
\label{6.1.}

\begin{frame}
\frametitle{6.1. Информация и энтропия }

Обороты <<получить информацию о событии>>, 
<<принять информацию к сведению>>,
<<оценить достоверность информации>> и другие,
наталкивают на мысль, что информация 
является величиной переменной и измеримой.
В этом разделе вводятся общепринятые меры информации.

\bigskip
\begin{tabular}{p{11mm}|p{25mm}|p{108mm}}
  № & Название параграфа
  & Ключевые термины и обозначения \\
  \hline
\DMref{6.1.1.}{6.1.1.}\RS & Количество информации &
{дискретный источник, бит, нат, информационная энтропия источника, энтропия, формула Шеннона, формула Хартли, избыточность источника, относительная избыточность, двоичный источник, функция Шеннона}
\\  
\DMref{6.1.2.}{6.1.2.} & Кодирование источника &
{Кодовое слово, схема кодирования, алфавитное кодирование, цена кодирования, эффективность кодирования, кодовое дерево, бинарное дерево, разделимая схема, неравенство Крафта, неравенство Макмиллана }  
\end{tabular}  
\end{frame}

\begin{frame}
\frametitle{6.1. Информация и энтропия (2/2)}
\begin{tabular}{p{11mm}|p{25mm}|p{108mm}}
  № & Название параграфа
  & Ключевые термины и обозначения \\
  \hline
\DMref{6.1.3.}{6.1.3.} & Взаимная и условная вероятность &
{ условная вероятность, формула Байеса, взаимная информация, априорная вероятность, апостериорная вероятность, условная информация, совместная энтропия, условная энтропия, относительная частота, стохастическая матрица}\\  
\DMref{6.1.4.}{6.1.4.} & Дискретный канал связи &
{передатчик, приёмник,матрица переходных вероятностей,
диаграмма переходов, диаграмма информационных потоков, утечка, шум}\\  
\DMref{6.1.5.}{6.1.5.} & Пропускная способность &
{стирание символа}\\  
\end{tabular}  
\end{frame}


\subsubsection{6.1.1. Количество информации}
\label{6.1.1}
\begin{frame}
\frametitle{6.1.1. Количество информации (1/7)}
В компьютере всё дискретно, недаром в старой
традиции компьютер назывался <<электронная
\em{цифровая} вычислительная машина>>.
Поэтому при определении понятия количества информации
для технических приложений оказывается
достаточно следующей простой дискретной модели.  

\smallskip
\odmkey{Дискретный источник информации}{Источник (source)!информации (of information)!дискретный (discrete)} 
выдаёт символ $x_i$ из конечного алфавита $X$ с заданным
на нём распределением вероятности $\mathrm{P}_X$ (\DMref{п.~5.1.7}{5.1.7.})

\vspace{-12pt}
$$
X=\{\Tuple{x}{1}{n}\},\quad
\mathrm{P}_X = \{\Tuple{p}{1}{n}\},\quad
\A{i\in 1..n}{0<p_i<1},\quad
\sum_{i=1}^n p_i = 1. 
$$ 

\vspace{-2pt}
В этом случае также говорят: 
происходит \odmkey{событие}{Событие (event)} $x_i$ из конечного множества возможных событий $X$ с вероятностью $p_i$. Наша первая цель состоит в том, чтобы ответить на вопрос, \em{сколько}
информации несёт единичное событие $x_i$?

\begin{remark}
Не ограничивая общности,
можно исключить из рассмотрения
невозможные события и достоверные события.
Если вероятность события равна нулю, то его можно 
не принимать в расчёт, а если вероятность 
некоторого события равна
единице, то вероятность остальных событий равна нулю
и их все можно не принимать в расчёт.
\end{remark}

\end{frame}

\begin{frame}
\frametitle{6.1.1. Количество информации (2/7)}
Для определения \odmkey{количества информации}{Количество (amount)!информации (of information)} события $x_i$, подберём функцию $I(x_i)$ так, что:
\begin{itemize}
\item  $\A{i}{I(x_i)\geqslant 0}$~--- количество информации о событии
неотрицательно;
\item $\A{i}{I(x_i)\sim 1/p_i}$~--- количество информации о событии находится в обратной
зависимости с вероятностью события;
\item $\A{i,j}{I(x_i\And x_j)= I(x_i)+I(x_j)}$~--- количество информации 
о нескольких независимых событиях равно сумме количеств информации
об отдельных событиях.
\end{itemize}

Этим требованиям удовлетворяет, например, \em{логарифм}
по основанию большему единицы величины, обратной к вероятности, откуда получается
\em{определение} количества информации,
доставляемой единичным событием $x_i$: 

$$
I(x_i)\Def \log_a \frac{1}{p_i} = - \log_a p_i.
$$

Приведённое определение количества информации является не более,
но и не менее, чем \em{общепринятым соглашением}. 

\smallskip
Как известно, $\log_a a =1$, т.~е. единица количества 
информации определяется выбором основания логарифма.
Напомним, что в наших обозначениях $\log x =\log_2 x$. 

\end{frame}

\begin{frame}
\frametitle{6.1.1. Количество информации (3/7)}

Наиболее употребительной в современной технике
единицей количества информации 
является \odmkey{бит}{Бит (bit)}, 
то есть количество 
информации для каждого из двух равновероятных событий
(обычно эти события обозначают символами 0 и 1), иначе говоря, количество информации, которое помещается в один двоичный разряд:  
$$
1\ \text{бит} = -\log \frac{1}{2} = \log 2. 
$$ 

\begin{remark}
По определению, количество информации~--- это любое
неотрицательное
\em{вещественное число}, совершенно не обязательно целое.
Однако, поскольку в один \em{двоичный разряд}
помещается ровно один бит информации,
очень часто допускают языковую небрежность,
и слово <<разряд>>  подменяют словом <<бит>>,
подразумевая не количество информации, а объём памяти.
В этом нет ничего страшного, если 
не забывать о том, что число разрядов непременно натуральное или ноль, а число битов может быть любым неотрицательным числом.   
\end{remark}

\smallskip
Если использовать натуральный логарифм, 
то единица информации называется \odmkey{нат}{Нат}.
$$
1\ \text{нат} = \ln e = \log e \ln 2\approx 1,443\ \text{бит}.
$$ 

\end{frame}

\begin{frame}
\frametitle{6.1.1. Количество информации (4/7)}
\begin{columns}[t]
\begin{column}[T]{9cm}
Зависимость количества информации от вероятности события:
чем больше вероятность,
тем меньше информации.
Достоверное событие не несёт никакой информации.
Невозможное событие даёт бесконечно много информации,
но оно невозможно. 
\end{column}
\begin{column}[T]{6cm}

\vspace{-12pt}
\includegraphics[width=4.5cm]{I2p.png}
\end{column}
\end{columns}


\smallskip
В теории информации \odmkey{информационная энтропия источника}{Энтропия (entropy)! информационная ((information))} определяется как \em{математическое ожидание} 
информации источника:
$$
H(X)\Def \sum_{i=1}^n p_i I(x_i) = -\sum_{i=1}^n p_i \log p_i.
$$ 

\begin{remark}
Эта формула называется \odmkey{формулой Шеннона}{Формула (formula)!Шеннона (Shannon)}.
\end{remark}

\begin{digress}
В физике (в термодинамике) \odmkey{энтропия}{Энтропия (entropy)}
определяется как мера неупорядоченности системы.
\end{digress}


\smallskip
\begin{theorem}
Энтропия источника с равномерным распределением вероятностей
максимальна и равна логарифму числа событий: 
$\max_{X} H(X)=\sum_{i=1}^{n}\frac{1}{n}I(x_i) =\log_a n.$
\end{theorem}

\end{frame}


\begin{frame}
\frametitle{6.1.1. Количество информации (5/7)}
\begin{proof} 
Энтропию можно считать в любых единицах, в битах или в натах.
Для упрощения выкладок воспользуемся натуральным логарифмом и
известной оценкой \\$\ln x \le x-1$. 
%Покажем, что энтропия любого 
%источника не превосходит энтропии источника с равномерным
%распределением вероятностей.
Рассмотрим источник $X=\{\Tuple{x}{1}{n}\}$ 
с распределением вероятностей $\{\Tuple{p}{1}{n}\}$ и
источник $Y=\{\Tuple{y}{1}{n}\}$ 
с равномерным распределением вероятностей $\{\Tuple{q}{1}{n}\}$,
где $\A{i\in 1..n}{q_i=\frac{1}{n}}$.
Используя оценку, имеем
$$
\ln q_i - \ln p_i = \ln \frac{q_i}{p_i} \le \frac{q_i}{p_i} -1.
$$
Умножая каждое неравенство на $p_i$ и суммируя по всем 
событиям, имеем 
$$
\sum\nolimits_{i=1}^n p_i(\ln q_i -\ln p_i) \le
\sum\nolimits_{i=1}^n p_i\left(\frac{q_i}{p_i}-1\right).
$$
Учитывая, что $H(X)= - \sum_{i=1}^n p_i(\ln p_i)$ и $\sum\nolimits_{i=1}^n p_i = 1 = \sum\nolimits_{i=1}^n q_i$, получаем
$$
H(X) + \sum\nolimits_{i=1}^n p_i\ln q_i \le \sum\nolimits_{i=1}^n q_i - \sum\nolimits_{i=1}^n p_i =0.
$$
Следовательно,
$
H(X) \le -\sum\nolimits_{i=1}^n p_i\ln q_i = -\sum\nolimits_{i=1}^n p_i\ln \frac{1}{n} = \ln n \sum\nolimits_{i=1}^n p_i = \ln n.
$
\end{proof}

\end{frame}

\begin{frame}
\frametitle{6.1.1. Количество информации (6/7)}
Из доказанной теоремы немедленно следует, что энтропия любого источника событий не превосходит
энтропии $H_0$ источника, в котором все события 
равновероятны:
$$
H_0= -\sum_{i=1}^n \frac{1}{n}\log \frac{1}{n} = 
\frac{1}{n}\sum_{i=1}^n \log n =\log n.
$$
\begin{remark}
Это формула называется \odmkey{формулой Хартли}{Формула (formula)!Хартли (Hartly)}. 
\end{remark}

\smallskip
\begin{example}
Как найти одну фальшивую (лёгкую) монету среди 27 монет
на рычажных весах? Если принять, что фальшивой 
может быть любая монета с равной вероятностью,
то по формуле Хартли для решения задачи необходимо получить
$\log 27 = 3\log 3$ бит информации. 
Одно взвешивание на рычажных весах даёт $\log 3$ 
бит информации, поэтому задачу невозможно решить 
меньше, чем за три взвешивания. Как это
сделать за три взвешивания, предоставляется догадаться читателю.  
\end{example}

\smallskip
Разность между максимальной возможной энтропией 
и энтропией источника называется \odmkey{избыточностью источника}{Избыточность!источника}:
$
R(X)\Def H_0 -H(X).
$
\odmkey{Относительная избыточность}{Избыточность!относительная}
определяется так:
$
r(X) \Def \frac{R(X)}{H_0} = 1-\frac{H(X)}{H_0} .
$

%\begin{example}
%Относительная избыточность <<читерского>> 
%игрального кубика (вместо единицы вторая шестерка) равна
%\end{example}
%   
\end{frame}

\begin{frame}
\frametitle{6.1.1. Количество информации (7/7)}
Особого внимания заслуживает 
\odmkey{двоичный источник}{Источник (source)!двоичный (binary)},
поскольку с помощью таких источников в 
большинстве случаев возможно описать процессы 
передачи данных.
В двоичном источнике есть два символа, 0 и 1, 
которые имеют вероятность $p$ и $1-p$ соответственно.
Энтропия двоичного источника зависит только от 
$p$ и выражается \odmkey{функцией Шеннона}{Функция (function)!Шеннона (Shannon)}:
$$
H_2(p)\Def -p\log p - (1-p)\log (1-p).
$$

Найдем производную функции Шеннона
$\displaystyle{
 (-p\log p - (1-p)\log (1-p))^\prime =}$\\
$\displaystyle{-\log p - \frac{p}{p\ln 2} +
 \frac{1}{(1-p)\ln 2} +
\log  (1-p) - \frac{p}{(1-p)\ln 2} = \log \frac{1-p}{p}}
$.  
Приравнивая \\производную нулю $\displaystyle{\log \frac{1-p}{p}=0}$,
находим точку максимума функции Шеннона $1/2$ из уравнения $1-p=p$.

\begin{columns}[t]
\begin{column}[T]{9cm}
Функция Шеннона выпукла вверх и 
симметрична относительно $1/2$. 
\end{column}
\begin{column}[T]{6cm}

\vspace{-22pt}
\includegraphics[width=4cm]{Sh2.png}
\end{column}
\end{columns}
\end{frame}

\subsubsection{6.1.2. Кодирование источника}
\label{6.1.2.}
%\begin{frame}
%\frametitle{6.1.2. Кодирование источника (1/7)}
%
%В современных условиях информация из различных источников
%по большей части обрабатывается с помощью компьютеров,
%а значит,
%должна быть \odmkey{закодирована}{Кодирование}, то есть представлена в
%виде последовательности символов 0 и 1.
%
%\medskip
%В этом разделе устанавливаются границы возможного
%для кодирования источников без потери информации. 
%
%\medskip
%\begin{remark}
%Здесь рассматриваются источники без памяти и
%двоичное кодирование, как наиболее важные
%на практике частные случаи. 
%\end{remark}
%\end{frame}

\begin{frame}
\frametitle{6.1.2. Кодирование источника (1/7)}

При кодировании источника $X=\{\Tuple{x}{1}{n}\}$ с 
распределением вероятности \\$\mathrm{P}=\{\Tuple{p}{1}{n}\}$,
необходимо каждому событию (символу) $x_i$
сопоставить некоторый $m$-ичный код $\alpha_i$,
который называется \odmkey{кодовым словом}{Кодовое слово (code word)},
а само сопоставление называют \odmkey{схемой кодирования}{Кодирование (coding)!Схема (scheme)} 
$\sigma(X)\Def\langle x_1\to\alpha_1,\dots, x_n\to\alpha_n \rangle$. 
Такое кодирование называют \odmkey{алфавитным}{Кодирование (coding)!алфавитное (alphabetic)}. 

\medskip
Пусть $l_i = |\alpha_i|$~--- длина кодового слова
$\alpha_i$.
Тогда математическое ожидание длины кода события источника
$$
l_* (\sigma(X)) \Def \sum_{i=1}^n p_i l_i
$$    
называют \odmkey{ценой кодирования}{Цена кодирования (price of coding)}.

\medskip
Цена кодирования прямо влияет на 
\odmkey{эффективность кодирования}{Эффективность кодирования (coding efficiency)},
то есть на расход памяти, скорость обработки и т.д.
Поэтому важно по возможности уменьшить цену кодирования
при сохранении других требований к кодированию, необходимых
в каждом конкретном случае.
\end{frame}

\begin{frame}
\frametitle{6.1.2. Кодирование источника (2/7)}
Алфавитное кодирование, то есть сопоставление символам $x_i$
$m$-ичных последовательностей $\alpha_i$ можно и удобно
представлять с помощью \odmkey{кодового дерева}{Дерево (tree)!кодовое (coding)}.
Кодовое дерево~--- это $m$-ичное дерево, в котором нумерация связей идёт слева направо от $0$ до $m-1$ 
(\DMref{п.~9.2.4}{9.2.4.}). 
В таком дереве каждый узел на $i$-м уровне однозначно представляет 
$m$-ичную последовательность длины $i$,
а именно последовательность, прочитанную вдоль пути
от корня дерева к данному узлу.

\begin{remark} Кодовое дерево для двоичного кодирования называется \odmkey{бинарным деревом}{дерево (tree)!бинарное (binary)}.
\end{remark}

\begin{theorem}
В кодовом дереве на $i$-м уровне 
имеется в точности $m^i$ узлов, а
всего в кодовом дереве высоты $h$ содержится 
$\frac{m^h − 1}{m - 1}$ узлов.
\end{theorem}

\begin{proof}
По формулам для геометрической прогрессии.
\end{proof}  

\begin{columns}[t]
\begin{column}[T]{5cm}
\begin{example}
На рисунке представлено 
кодовое дерево высоты 4,
в котором выделено красным
кодовое слово 101. 
\end{example}
\end{column}
\begin{column}[T]{10cm}

\vspace{-20pt}
\begin{center}
\includegraphics[width=9cm]{D2.png}
\end{center}
\end{column}
\end{columns}

\end{frame}

\begin{frame}
\frametitle{6.1.2. Кодирование источника (3/7)}
Говорят, что \odmkey{схема алфавитного кодирования}{Схема (scheme)!кодирования (coding)!алфавитная (alphabetic)}
 $\langle x_1 \to \alpha_1, 
\dots , x_n \to \alpha_n \rangle$ является \odmkey{префиксной}{Схема (scheme)!префиксная (prefix)}, если никакое кодовое слово $\alpha_i$
не является собственным началом никакого кодового 
слова $\alpha_j$.
В терминах кодовых деревьев это означает, 
что ни один узел, которому назначено кодовое слово,
не лежит ни в каком поддереве, корнем которого является другое
кодовое слово (помечено зелёным на предыдущем рисунке).

\begin{theorem}
Префиксное кодирование имеет
однозначное декодирование.
\end{theorem}   

\begin{proof}
При декодировании достаточно слева направо
однозначно выделять кодовые слова из декодируемой последовательности. 
\end{proof}

\smallskip
\begin{theorem} Если  $\langle x_1 \to \alpha_1, 
\dots , x_n \to \alpha_n \rangle$ является префиксной схемой
алфавитного кодирования, то

\vspace{-20pt}
$$
\sum_{i=1}^n \frac{1}{m^{|\alpha_i|}} \leq 1.
$$
\end{theorem}

\begin{proof}
Пусть $k\Assign \max |\alpha_i|$. 
Тогда заданная схема может быть представлена кодовым деревом
высоты $k$, в котором отмечены $n$ узлов. При этом, 
поскольку схема префиксная, ни один 
отмеченный узел в дереве не находится
в поддереве, определяемым другим отмеченным узлом.
\end{proof}
\end{frame}

\begin{frame}
\frametitle{6.1.2. Кодирование источника (4/7)}
\begin{proof}
(продолжение)
Заметим, что каждый отмеченный узел делает <<запретными>> 
для других узлов в точности $m^{k-|\alpha_i|}$ узлов из 
$m^k$
узлов на последнем $k$-м уровне кодового дерева. Поскольку <<запретные>> участки на последнем уровне не пересекаются, имеем в сумме
$$\sum_{i=1}^n m^{k-|\alpha_i|}\leq m^k,$$
откуда, поделив на $m^k$, получаем требуемое неравенство.   
\end{proof}

\smallskip
\begin{remark}
Доказанное неравенство называют
\odmkey{неравенством Крафта}{Неравенство!Крафта}.
В теории кодирования показано
(см. \DMref{6.2.4}{6.2.4}), что 
данное неравенство является не только необходимым, но и  
достаточным условием
существования однозначного декодирования,
и в этом случае обычно
говорят про
\odmkey{неравенство Макмиллана}{Неравенство!Макмиллана}.
\end{remark}
   
\smallskip
\begin{example}
Схема кодирования 
$x_1\to 0$,
$x_2\to 10$,
$\dots$,
$x_{n-1}\to 1\dots 10$,
$x_n\to 1\dots 11$
является префиксной, имеет кодовые слова длиной 
$2^0, 2^1, \dots, 2^{n-1},2^{n-1}$,
удовлетворяет неравенству Крафта и имеет
кодовое дерево высотой $n-1$. 
\end{example}

\end{frame}

\begin{frame}
\frametitle{6.1.2. Кодирование источника (5/7)}
Доказанное неравенство позволяет оценить 
границы возможного при построении различных схем алфавитного
 кодирования
для источника с заданной энтропией.


\begin{lemma} \NB{[Неравенство Гиббса]} Для двух распределений вероятности $\mathrm{P}_X = \{\Tuple{p}{1}{n}\}$ и
 $\mathrm{Q}_X = \{\Tuple{q}{1}{n}\}$ для источника $X$, выполнено:
$$-\sum_{i=1}^{n} p_i \log_a p_i \le 
-\sum_{i=1}^{n} p_i \log_a q_i
$$
\end{lemma} 

\begin{proof}
От логарифма по любому можно перейти к
натуральным логарифмам: $\log_a x = \frac{\ln x}{\ln a}$.
 Используя известное неравенство 
 $\ln x\le x - 1$, имеем 
 $\ln \frac{q_i}{p_i}\le \frac{q_i}{p_i} - 1$, 
 откуда $p_i \ln \frac{q_i}{p_i}\le q_i - p_i$.
  Суммируем по всем i:
$$
\sum_{i=1}^{n} p_i \ln \frac{q_i}{p_i}\le \sum_{i=1}^{n}q_i - \sum_{i=1}^{n}p_i = 1 -1 =0.$$
 Из чего получаем: 
 $-\displaystyle\sum_{i=1}^{n}p_i \ln p_i \le -\displaystyle\sum_{i=1}^{n}p_i \ln q_i$.
\end{proof}

\end{frame}

\begin{frame}
\frametitle{6.1.2. Кодирование источника (6/7)}
Теорема Шеннона устанавливает границы возможного
для алфавитного кодирования источников без потери информации. 
 
\begin{theorem}
\NB{[Шеннона]}
 Для любого источника $X = \{\Tuple{x}{1}{n}\}$
  с энтропией $H(X)$ cуществует префиксная схема алфавитного кодирования $\sigma(X)$, такая что: 
$$\frac{H(X)}{\log m} \le l_{\ast}(\sigma(X)) < \frac{H(X)}{\log m} + 1.$$
\end{theorem} 
\begin{proof} Положим $\l_i \Assign|\alpha_i|$ в схеме
$\sigma(X)=\langle x_1\to\alpha_1,\dots, x_n\to\alpha_n \rangle$.
\proofsection{$\frac{H(X)}{\log m} \le l_{\ast}(\sigma(X))$}
Обозначим $c \Assign \sum_{i=1}^{n} m^{-l_i} \le 1$.
По неравенству Крафта $c\le 1$. 
Пусть $q_i \Assign m_i^{-l_i}/c$, 
заметим что $\sum_{i=1}^{n}q_i=1$. 
По неравенству Гиббса имеем:
\begin{align*}
H(X) &= - \sum_{i=1}^{n}p_i \log p_i \le 
- \sum_{i=1}^{n}p_i \log q_i = 
\sum_{i=1}^{n}p_i (l_i \log m + \log c) \le \\
&\le\sum_{i=1}^{n}p_i l_i \ln m = l_{\ast}(\sigma(X)) \log m.
\end{align*} 
\end{proof}

\end{frame}

\begin{frame}
\frametitle{6.1.2. Кодирование источника (7/7)}

\begin{proof} (продолжение)
\proofsection{$l_{\ast}(\sigma(X)) < \frac{H(X)}{\log m} + 1$}
Положим $l_i = \lfloor - \log_a p_i \rfloor$, 
тогда $- \log_a p_i \le l_i < - \log_a p_i + 1$. 
Имеем $a^{- l_i} \le p_i \Impl 
\displaystyle\sum_{i=1}^{n} a^{- l_i} \le 1$. По неравенству Макмилана (\DMref{п.~6.2.4}{6.2.4.}) существует префиксная схема с соответствующими длинами кодовых слов в ней, поэтому для наших длин $\Tuple{l}{1}{n}$ существует префиксная схема
 $l_{\ast}$ такая, что:
 $$(\sigma(X)) = \displaystyle\sum_{i=1}^{n}  p_i l_i < \displaystyle\sum_{i=1}^{n} p_i (-\log_a p_i + 1) = \displaystyle\sum_{i=1}^{n} \left(\frac{-p_i \log p_i}{\log a}\right) + 1  = \frac{H(X)}{\log a} + 1.$$
\end{proof}

\begin{remark}
Схема кодирования, допускающая однозначное декодирование,
называется \odmkey{разделимой}{Схема (scheme)!разделимая}.
 Первая часть неравенства теоремы Шеннона верна для любой разделимой схемы.
\end{remark}

\end{frame}


\subsubsection{6.1.3. Взаимная и условная информация}
\label{6.1.3.}
\begin{frame}
\frametitle{6.1.3. Взаимная и условная информация (1/8)}

На практике информация получается из 
различных источников, и информация получаемая из
одного источника, может влиять
на оценку количества информации, получаемой 
из другого источника, причём это влияние взаимно.

\medskip
Даже в случае одного источника, 
информация, полученная из этого источника ранее,
может влиять на оценку количества информации,
получаемой из источника.

\medskip
\begin{example}
Если из источника кириллических букв
получена буква <<ъ>>, то неразумно считать,
что вероятности следующих букв равны 1/33,
поскольку в современном русском языке нет слов, 
оканчивающихся на твёрдый знак, и вслед за твёрдым
знаком следует одна из гласных е, ё, ю, я.  
\end{example}

\medskip
Определим информацию и энтропию для двух источников
$X=\{\Tuple{x}{1}{n}\}$ и $Y=\{\Tuple{y}{1}{m}\}$,
 для этого определим наш новый
источник такой, что он будет выдавать \em{пару} событий (символов) $(x_i, y_j)$.

\begin{remark}
Для упрощения обозначений в этом параграфе индексы можно
опустить, и рассматривать пару $(x, y)$,
подразумевая, что $x=x_i, y=y_j$.
\end{remark}

\end{frame}

\begin{frame}
\frametitle{6.1.3. Взаимная и условная информация (2/8)}

Введем понятие  \odmkey{условной вероятности}{Вероятность!условная}
события $a$ при условии события $b$ 
(обозначается $P(a|b)$) 
 следующим образом:
$$
P(a|b)\Def \frac{P(a, b)}{P(b)},
$$
где $P(a, b)$~--- совместная вероятность событий 
$a$ и $b$, иначе говоря, вероятность пары событий 
$(a, b)$
 
\begin{remark}
Определение условной вероятности дано в \DMref{п.~5.1.7}{5.1.7.},
однако сейчас наибольший интерес для нас представляют элементарные
события, поэтому здесь приведённого определения (без использования понятия события как множества) \em{достаточно}.
\end{remark} 


Из определения условной вероятности немедленно следует 
\odmkey{формула Байеса}{Формула (formula)!Байеса (Bayes)} (\DMref{п.~5.1.7}{5.1.7.}):
$$
P(a|b)P(b)=P(a,b)=P(b|a)P(a).
$$

Логарифмируя формулу Байеса, получаем выражение для информации пары событий $I(x,y)$:
$$
-I(x,y)=\log_2 P(x,y)= \log_2 P(x|y) - I(y)=
\log_2 P(y|x) - I(x).
$$ 
\end{frame}

\begin{frame}
\frametitle{6.1.3. Взаимная и условная информация (3/8)}
Несложные преобразования дают следующие выражения:
$$
I(x,y)=I(x)+I(y) - \log_2 \frac{P(x|y)}{P(x)}= 
I(x)+I(y) - \log_2 \frac{P(y|x)}{P(y)}.
$$ 

Таким образом, информация пары событий 
оказывается суммой информаций этих событий за вычетом некоторой неотрицательной величины, которая снижает неопределённость, то есть сама в свою очередь является информацией.

\odmkey{Взаимная информация}{Информация (information)!взаимная (mutual)}
пары событий $I(x;y)$ определяется так:
$$
I(x;y)\Def \log_2 \frac{P(x|y)}{P(x)}= 
\log_2 \frac{P(y|x)}{P(y)}.
$$

\begin{remark}
Взаимная информация событий положительна и
симметрична.
\end{remark}

\begin{digress} В теории вероятностей и в статистике
в дроби $P(x|y)/P(x)$ числитель $P(x|y)$ называют
\odmkey{апостериорной}{Вероятность!апостериорная}
вероятностью события $x$ относительно события $y$, а знаменатель $P(x)$
называют
\odmkey{априорной}{Вероятность!априорная}
вероятностью события $x$ (см. также \DMref{п.~5.1.7}{5.1.7.}).
\end{digress}
\end{frame}

\begin{frame}
\frametitle{6.1.3. Взаимная и условная информация (4/8)}
Рассмотрим два граничных случая.

\smallskip
1) Источники независимы, то есть $P(x,y)=P(x)P(y)$.
В этом случае $I(x;y)= 0$ и источники не обмениваются
информацией.

\smallskip
2) Источники жёстко связаны, то есть 
$P(x|y)=P(y|x)=1$. В этом случае\\
$I(x;y)= I(x)=I(y)$ и источники обмениваются
всей информацией.

\medskip
Введём понятие \odmkey{условной информации}{Информация (information)!условная (conditional)}: $I(x|y)\Def - \log_2 P(x|y)$. Заметим, что\\ 
$I(x;y)=I(x|y)-I(y)=I(y|x)-I(x)$.
Тогда из основного уравнения 
$$
I(x,y)=I(x)+I(y)-I(x;y)
$$
следует, что
$$
I(x,y)=I(x)+I(y|x)=I(y)+I(x|y).
$$
Другими словами, информация пары событий 
$(x,y)$ есть сумма информации события $x$ и
условной информации события $y$ при условии,
что событие $x$ уже известно, или, в другом порядке,
сумма информации события $y$ и
условной информации события $x$ при условии,
что событие $y$ уже известно.
\end{frame}

\begin{frame}
\frametitle{6.1.3. Взаимная и условная информация (5/8)}
Имея определение информации пар событий, определим 
\odmkey{совместную энтропию}{Энтропия (entropy)!совместная (joint)}
двух источников $X$ и $Y$ как математическое ожидание информации
пары событий:
$$
H(X,Y)\Def -\sum_{x\in X} \sum_{y\in Y} P(x,y)\log_2 P(x,y).
$$
Здесь подразумевается, что рассматриваются все пары 
совместных событий, то есть
$$
\sum_{x\in X} \sum_{y\in Y} P(x,y) =
\sum_{i=1}^n \sum_{j=1}^m P(x_i,y_j)=1.
$$

Определим 
\odmkey{условную энтропию}{Энтропия (entropy)!условная (conditional)}
источника $X$ относительно источника $Y$ как математическое ожидание условной информации:
$$
H(X|Y)\Def -\sum_{x\in X} \sum_{y\in Y} P(x,y)\log_2 P(x|y).
$$
Тогда $H(X,Y) = H(Y)+H(X|Y)=H(X)+H(Y|X)$.
\end{frame}

\begin{frame}
\frametitle{6.1.3. Взаимная и условная информация (6/8)}

Для вычисления взаимной и условной энтропии
необходимо знать вероятности всех пар событий.
В некоторых случаях, даже
если точные вероятности не известны, их
можно оценить и вычислить энтропию \em{приближённо}. 

\medskip
Пусть есть два дискретных источника
$X$ и $Y$.
Пусть также есть \em{выборка} пар совместных
событий этих источников, т.~е. \em{мультимножество} 
$\left[(x_i,y_j)^{a_{ij}}\right]$, причём
мощность выборки $k=\sum a_{ij}$. 
Если число $k$ достаточно велико,
т.~е. выборка \em{представительная},
то \odmkey{относительная частота}{Частота (frequency)!относительная (relative)}
пары событий может быть использована как оценка
совместной вероятности событий $(x_i,y_j)$ 
$$P(x_i,y_j)\approx a_{ij}/k.$$

\begin{example}
Пусть есть два дискретных источника
$X=\{x_1,x_2,x_3,x_4\}$ и $Y=\{y_1,y_2,y_3,y_4\}$,
алфавит каждого источника содержит по четыре символа.
Пусть также есть представительная выборка пар совместных
событий этих источников, причём значения 
относительной частоты пар событий подсчитаны 
и сведены в таблицу на следующем слайде. 
\end{example}
\end{frame}

\begin{frame}
\frametitle{6.1.3. Взаимная и условная информация (7/8)}
\begin{center}
\begin{tabular}{|c||c|c|c|c||c|}
\hline
 & $y_1$ & $y_2$ &$y_3$ &$y_4$ & $P(x_i)$ \\
\hline
\hline 
$x_1$ & $0,10$ & $0,05$ &$0,05$ & $0$ & $0,20$ \\
\hline
$x_2$ & $0,05$ & $0,15$ &$0,15$ &$0$ & $0,35$ \\
\hline
$x_3$ & $0$ & $0,10$ & $0,10$ & $0,10$ & $0,30$ \\
\hline
$x_4$ & $0$ & $0,05$ &$0,05$ &$0,05$ & $0,15$ \\
\hline
\hline
$P(y_j)$ & $0,15$ & $0,35$ &$0,35$ &$0,15$ & $1$ \\
\hline
\end{tabular}    
\end{center}

При таких исходных данных имеем:

\begin{itemize}
\item две энтропии истоников
$$H(X)=-\sum_{i=1}^4 P(x_i)\log_2 P(x_i) \approx 1,92,\quad
H(Y)=-\sum_{j=1}^4 P(y_j)\log_2 P(y_j) \approx 1,88 $$
\item совместная энтропия истоников
$$H(X,Y)=-\sum_{i=1}^4\sum_{j=1}^4 P(x_i,y_j)\log_2 P(x_i,y_j) \approx 3,45 $$
\end{itemize}
\end{frame}

\begin{frame}
\frametitle{6.1.3. Взаимная и условная информация (8/8)}
\begin{itemize}
\item две условные энтропии истоников
$$H(X|Y)=H(X,Y)-H(Y) \approx 1,57,\qquad
H(Y|X)=H(X,Y)-H(X) \approx 1,52 $$
\item условные вероятности $P(y_j|x_i)$
(\odmkey{стохастическая матрица}{Матрица (matrix)!стохастическая (stochastic)})
\begin{center}
\begin{tabular}{|c||c|c|c|c|}
\hline
 & $y_1$ & $y_2$ & $y_3$ & $y_4$ \\
\hline
\hline 
$x_1$ & $1/2$ & $1/4$ & $1/4$ & $0$ \\
\hline
$x_2$ & $1/7$ & $3/7$ & $3/7$ & $0$  \\
\hline
$x_3$ & $0$ & $1/3$ & $1/3$ & $1/3$  \\
\hline
$x_4$ & $0$ & $1/3$ & $1/3$ &$1/3$  \\
\hline
\end{tabular}    
\end{center}

\item условная энтропия
$$H(Y|X)=-\sum_{i=1}^4 \sum_{j=1}^4 P(x_i, y_j)\log_2 P(y_j|x_i) \approx 1,55. $$
\end{itemize}

\end{frame}


\subsubsection{6.1.4. Дискретный канал связи}
\label{6.1.4.}
\begin{frame}
\frametitle{6.1.4. Дискретный канал связи (1/9)}
Понятие \em{передачи информации}
допускает множество различных
интерпретаций. Здесь рассматривается
простейший, но одновременно 
наиболее актуальный случай
передачи информации по дискретным,
то есть \em{цифровым}, каналам связи.  

\smallskip
\odmkey{Дискретный канал связи без памяти}
{Канал (channel)!связи (communication)!дискретный (discrete)}
\odmkey{}
{Канал (channel)!связи (communication)!без памяти (memoryless)}
 мыслится как
пара взаимосвязанных дискретных источников.
Источник $X$, который называется \odmkey{передатчиком}{Передатчик (transmitter)}, передает в канал входной символ $x_i$.
символ $y_j$ источника $Y$, который называется
\odmkey{приёмником}{Приёмник (receiver)}. Оборот <<без памяти>> 
применительно к каналу означает,
что входной и выходной символы связаны в пары,
то есть каждый выходной символ зависит только от 
соответствующего входного символа.   

\medskip
Влияние входного символа на выходной символ
определятся свойствами канала. 
На практике свойства реальных каналов
зависят от множества разнообразных 
физических факторов.
В теоретических моделях каналов 
используется 
\odmkey{матрица переходных вероятностей}
{Матрица (matrix)!переходных вероятностей} $P(y_j|x_i)$.
Каждый элемент этой матрицы задаёт вероятность того,
что выходным символом является $y_j$ при условии,
что входным символом является $x_i$.
Матрица переходных вероятностей, очевидно, стохастическая.

\end{frame}


\begin{frame}
\frametitle{6.1.4. Дискретный канал связи (2/9)}
В терминах теории графов матрицу переходных вероятностей
можно интерпретировать как матрицу весов дуг 
двудольного орграфа. В этом случае орграф называют
\odmkey{диаграммой  переходов}{Диаграмма (diagram)!переходов (of transitions)}
канала.  
Двоичный симметричный канал описывается следущей диаграммой переходов и матрицей переходных вероятностей: 

\begin{columns}[t]
\begin{column}[T]{7cm}
\vspace{-20pt}
\begin{center}
\includegraphics[height=20mm]{2sch.png}
\end{center}
\end{column}
\begin{column}[T]{7cm}
$$\begin{pmatrix}
1-\varepsilon & \varepsilon \\
\varepsilon & 1-\varepsilon 
\end{pmatrix}$$
\end{column}
\end{columns}

Если распределение вероятностей входных символов равномерное,
то распределение вероятностей выходных символов также равномерное.
Отсюда определяются условные вероятности:
\begin{align*}
P(y_0 | x_0) =  & 1-\varepsilon,\\
P(y_1 |x_0) =& \varepsilon,\\
P(y_0 |x_1) =& \varepsilon,\\
P(y_1 | x_1) =& 1-\varepsilon.
\end{align*}

\end{frame}

\begin{frame}
\frametitle{6.1.4. Дискретный канал связи (3/9)}
Из условных (апостериорных) вероятностей,
учитывая, что априорная вероятность каждого символа равна $1/2$,
имеем взаимные информации пар символов: 
\begin{align*}
I(y_0;x_0) =& \log_2\frac{P(y_0|x_0)}{P(y_0)} =
\log_2\frac{1-\varepsilon}{1/2}=1+\log_2(1-\varepsilon),\\
I(y_1;x_0) =& \log_2\frac{P(y_1|x_0)}{P(y_1)} =
\log_2\frac{\varepsilon}{1/2}=1+\log_2(\varepsilon),\\
I(y_0;x_1) =& \log_2\frac{P(y_0|x_1)}{P(y_0)} =
\log_2\frac{\varepsilon}{1/2}=1+\log_2(\varepsilon),\\
I(y_1;x_1) =& \log_2\frac{P(y_0|x_0)}{P(y_0)} =
\log_2\frac{1-\varepsilon}{1/2}=1+\log_2(1-\varepsilon).
\end{align*}
\end{frame}

\begin{frame}
\frametitle{6.1.4. Дискретный канал связи (4/9)}
Рассмотрим три важных частных случая.

\begin{enumerate}
\item $\varepsilon=0$~--- передача без ошибок, в этом случае
взаимная информация\\$I(y_0;x_0)=I(y_1;x_1)=1$, других взаимных информаций не существует, и передача информации от источника к приёмнику происходит без потерь.
\item $\varepsilon=1/2$~--- передача c равновероятными ошибками, в этом случае
взаимная информация $I(y_0;x_0)=I(y_1;x_0) =I(y_1;x_1)=I(y_0;x_1)=0$ и передачи информации от источника к приёмнику не происходит.
\item $\varepsilon=1$~--- передача c постоянными ошибками. Инвертируя принятые символы, приходим к первому случаю.
\end{enumerate}

В общем случае, если $0<\varepsilon<1$, то
все взаимные информации строго меньше 1,
а значит невозможно передать информацию без потерь в двоичном симметричном канале с ошибками, даже если вероятность ошибки мала, но не ноль.
\end{frame}

\begin{frame}
\frametitle{6.1.4. Дискретный канал связи (5/9)}

Описание передачи информации с помощью переходных вероятностей
сводится, в конечном счёте, к совместным вероятностям пар событий $P(x,y)$, а значит различать приёмник $Y$ и передатчик $X$ нет нужды,
и тот и другой можно считать источником событий. 

\medskip
В таком случае определена 
\odmkey{совместная энтропия}{Энтропия (entropy)!совместная (joint)} двух источников:
$$
H(X,Y) = -\sum_{x\in X} \sum_{y\in Y} P(x,y)\log_2 P(x,y).
$$

Аналогично определяются \odmkey{условные энтропии}{Энтропия (entropy)!условная (conditional)} двух источников:
$$
H(X |Y) = -\sum_{x\in X} \sum_{y\in Y} P(x,y)\log_2 P(x |y),
$$
$$
H(Y |X) = -\sum_{x\in X} \sum_{y\in Y} P(x,y)\log_2 P(y |x).
$$

\vspace{-8pt}
Из этого следует, что
$$
H(X,Y)=H(Y)+H(X|Y)=H(X)+H(Y|X), \quad H(X,Y)\le H(X)+H(Y),
$$
причём равенство достигается, только если источники независимы.
\end{frame}

\begin{frame}
\frametitle{6.1.4. Дискретный канал связи (6/9)}
В определении совместной энтропии можно рассматривать любое 
количество источников, поскольку любую группу источников можно рассматривать как новый источник. В частности,  
$
H(X,Y,Z)= H(Z|X,Y) + H(X,Y)
$. Это наблюдение позволяет вывести неравенство, связывющее 
энтропии трёх источников.

\smallskip  
\begin{lemma}
$
H(X|Y) \le H(X|Y,Z) + H(Z).
$
\end{lemma}
\begin{proof}
По уже полученному имеем:
\begin{align*}
H(X,Y,Z) & = H(Z|X,Y)+ H(X,Y) = H(Z|X,Y) + H(X|Y) + H(Y), \\
H(X,Y,Z) & = H(X|Y,Z)+ H(Y,Z).
\end{align*} 
А так как $H(Z|X,Y)\ge 0$, получаем 
$$H(X|Y) + H(Y) \le H(X,Y,Z)= H(X|Y,Z) + H(Y,Z) \le H(X|Y,Z) + H(Y) + H(Z),$$ поскольку $H(Z,Y)\le H(Z)+H(Y)$.
Сокращая $H(Y)$, имеем требуемое.
\end{proof}

\smallskip
\begin{theorem}
\NB{[Неравенство Фано]}
Пусть $X$ и $Y$ два $m$-ичных источника. 
Тогда, если $\varepsilon$~--- вероятность того, 
что $x_i\ne y_i$ в паре $(x_i, y_i)$, то $ H(X|Y) \le H_2(\varepsilon) + \varepsilon \log (m-1)$,
где $H_2(\varepsilon)$~--- функция Шеннона.
\end{theorem}


\end{frame}

\begin{frame}
\frametitle{6.1.4. Дискретный канал связи (7/9)}
\begin{proof}
Положим, двоичный источник $Z$ выдаёт символ $1$, в случае если $x_i \ne y_i$, а в противном случае  выдаёт символ $0$. 
Тогда распределение вероятностей на 
источнике $Z$~--- это $(1-\varepsilon, \varepsilon)$. По предыдущей лемме $ H(X|Y) \le H(X|Y,Z) + H(Z) $, причём $H(Z)=H_2(\varepsilon)$ (\DMref{п.~6.1.1}{6.1.1.}).
Рассмотрим два случая для $H(X|Y,Z)$:
$H(x, y, 0) = 0$, так как тогда $x=y$ и мы не получаем никакой информации;
$H(x, y, 1) \le \log(m-1)$, поскольку у нас есть еще по крайней мере $m-1$ символов, не равных $y$.
Получаем:
$$
H(X|Y,Z) = \sum_{x\in X} \sum_{z \in Z} P(y, z) H(X|Y,Z) \le \sum_{y \in Y} P(y, 1) \log (m-1)
$$
Суммируем по $y$: $\displaystyle\sum_{y \in Y} P(y, 1) = \varepsilon$. Получаем $H(X|Y) \le H_2(\varepsilon) + \varepsilon \log (m-1)$
\end{proof}

\begin{remark}
Это неравенство часто применяется при доказательстве многих теорем в теории кодирования (\DMref{п.~6.1.5}{6.1.5.}). А также помогает, оценить степень связности источника и приёмника.
\end{remark}
\end{frame}


\begin{frame}
\frametitle{6.1.4. Дискретный канал связи (8/9)}
В случае двух источников, свзяанных каналом,
совместная неопределённость снижается,
поскольку событие одного источника позволяет
предсказывать событие другого.
Сниже\-ние неопредённости 
означает обмен информацией.
Отсюда приходим к определению:
среднее значение информации,
передаваемой по каналу, равно
математическому ожиданию взаимных информаций 
всех пар событий.

\begin{align*}
I(X;Y) &=
\sum_{x\in X} \sum_{y\in Y} P(x,y)I(x;y)=
 -\sum_{x\in X} \sum_{y\in Y} P(x,y)\log_2 
\frac{P(y |x)}{P(y)}= \\
& =
\sum_{x\in X} \sum_{y\in Y} P(x,y)\log_2 
P(y |x)-
\sum_{x\in X} \sum_{y\in Y} P(x,y)\log_2 
P(y).
\end{align*}

По определению энтропий и учитывая, что
$\sum_{x\in X} P(x,y)=P(y)$, имеем:
 $$\sum_{x\in X} \sum_{y\in Y} P(x,y)\log_2 
P(y |x) = -H(X|Y), \quad \sum_{x\in X} \sum_{y\in Y} P(x,y)\log_2 
P(y) = -H(Y),$$ 
 и поэтому
$$
I(X;Y)= H(X)-H(X|Y)=H(Y)- H(Y|X). 
$$
\end{frame}

\begin{frame}
\frametitle{6.1.4. Дискретный канал связи (9/9)}
Найденные зависимости энтропий можно наглядно
показать с помощью 
\odmkey{диаграммы информационных потоков}{Диаграмма (diagram)!информационных потоков (of information flow)}:
\begin{center}
\includegraphics[height=20mm]{DiaC.png}
\end{center}

Условная энтропия $H(X|Y)$ определяет меру неопределённости
посланного символа источника $X$, 
когда символ источника  $Y$ известен.
Энтропию $H(X|Y)$ часто называют  величиной 
\odmkey{утечкой информации}{Утечка (leak)!информации (of information)}.

\medskip
Условная энтропия $H(Y|X)$ определяет меру неопределённости
принятого символа источника $Y$, 
когда символ источника  $X$ известен.
Энтропию $H(Y|X)$ часто называют  величиной  
\odmkey{информационного шума}{Шум (noise)!информационный (infomation)}.

\medskip
Энтропия $H(X)$~--- это энтропия на входе канала,
энтропия $H(X)$~--- это энтропия на выходе канала.
В любом случае количество 
информации, передаваемой через канал,
описывается основным уравнением:
$
H(X)-H(X|Y) = I(X;Y) = H(Y)- H(Y|X). 
$

\end{frame}



\subsubsection{6.1.5. Пропускная способность}
\label{6.1.5.}
\begin{frame}
\frametitle{6.1.5. Пропускная способность (1/4)}
Величина $I(X;Y)$ играет важную роль в теории
информации, поскольку описывает передачу
информации по каналу связи. 
Эта величина зависит как от переходных вероятностей
канала, так и от распределения вероятностей 
символов на входе канала. Возникает вопрос:
какое максимальное количество информации
возможно передать по каналу с заданными 
переходными вероятностями?
Ответ даёт 
\odmkey{пропускная способность}{Пропускная способность} канала:
$$
C(X,Y)\Def \max_X I(X;Y),
$$
где максимум берётся по всем возможным распределениям 
вероятностей символов источника $X$.


\begin{remark}
В общем случае такая экстремальная задача может быть
решена \em{методом множителей Лагранжа},
который имеет высокую вычислительную трудоёмкость.
\end{remark}

\medskip
В дискретном случае помогает теорема,
гласящая, что в симметричном канале без памяти 
пропускная способность 
достигается при равномерном распределении 
вероятностей символов источника.
\end{frame}

\begin{frame}
\frametitle{6.1.5. Пропускная способность (2/4)}
Рассмотрим двоичный симметричный канал без памяти с
вероятностью ошибки $\varepsilon$.
Из равномерности распределения входных символов и 
симметричности переходов следует равномерное распределение
выходных символов:
$$
P(x_0)=P(x_1)=P(y_0)=P(y_1)= \frac{1}{2}.
$$
Используя формулу Байеса и определение величины $I(X;Y)$,
имеем:
\begin{align*}
C =& \sum_{i=0}^1 \sum_{j=0}^1 P(x_i, y_j) 
\log_2 \frac{P(y_j|x_i)}{P(y_j)} =
\sum_{i=0}^1 \sum_{j=0}^1 P(x_i)P(y_j| x_i) 
\log_2 \frac{P(y_j|x_i)}{P(y_j)} = \\
=&  \frac{1}{2}(1-\varepsilon)\log_2 (2(1-\varepsilon)) +
\frac{1}{2}\varepsilon\log_2 (2\varepsilon) +
\frac{1}{2}\varepsilon\log_2 (2\varepsilon) +
\frac{1}{2}(1-\varepsilon)\log_2 (2(1-\varepsilon)) = \\
=& 1 + (1-\varepsilon)\log_2 (1-\varepsilon)
+ \varepsilon\log_2 \varepsilon = 1 - H_2(\varepsilon),
\end{align*}
где $H_2(\varepsilon)$~--- функция Шеннона 
(см. \DMref{п.~6.1.1}{6.1.1.}).

Таким образом, пропускная способность в канале без помех ($\varepsilon=0$)
 равна 1,
а в зашумлённом канале ($\varepsilon=1/2$) равна 0. 
\end{frame}

\begin{frame}
\frametitle{6.1.5. Пропускная способность (3/4)}

\begin{columns}[t]
\begin{column}{11cm}
В реальных аналоговых каналах из-за шумов 
иногда бывает так, что невозможно определить,
был ли принят какой-либо символ. 
Другими словами, алфавит приёмника состоит 
из трёх символов 0, 1 и символа \_, 
который означает, что произошло \odmkey{стирание}{Стирание (erasing)!символа (of symbol)}  переданного символа.
Такой канал характеризуются двумя параметрами:
вероятность ошибки $p$ и вероятность стрирания $q$, как показано на рисунке.
\end{column}
\begin{column}{4cm}
\vspace{-40pt} 
\begin{center}
\includegraphics[height=32mm]{2s3.png}
\end{center}
\end{column}
\end{columns}

Так как канал симметричен, пропускная способность
достигается при равномерном распределении вероятностей
входных символов, и значит:
$$
P(x_0) = 1/2, \quad
P(x_1) = 1/2, \quad
P(y_0) = (1-q)/2, \quad
P(y_{\_}) = q, \quad
P(y_1) = (1-q)/2.
$$
Подставляя известные верятности в формулу пропускной способности, получаем
$$
C =\sum_{i=0}^1 \sum_{j=0}^2 
 P(x_i)P(y_j| x_i) 
\log_2 \frac{P(y_j|x_i)}{P(y_j)} =
1-q+(1-p-q)\log_2\frac{1-p-q}{1-q}+p\log_2\frac{p}{1-q}.
$$
\end{frame}

\begin{frame}
\frametitle{6.1.5. Пропускная способность (4/4)}

Размерность пропускной способности канала $C$ бит/символ.
Если известно, что символы передаются с постоянной скоростью,
каждый символ передаётся $t$ секунд, то величина
$C/t$ получает размерность бит/сек. 

\begin{remark}
В инженерной практике именно
величину $C/t$, имеющую размерность бит/сек часто
называют пропускной способностью канала.
\end{remark}

\medskip
Пределы возможного при передаче информации по
дискретному каналу без памяти устанавливают 
прямая и обратная теоремы Шеннона о кодировании.

\begin{theorem}
Если $H(X)<C$,
то существует кодирование, 
которое позволяет передать информацию
по каналу со сколь угодно малой вероятностью ошибки $\varepsilon$. 
\end{theorem}

Смысл этой теоремы состоит в том, что если передать 
достаточно много символов, то вероятность ошибки при
декодировании можно сделать сколь угодно малой.
К сожалению, доказательство не конструктивное, 
оно не содержит метода построения кодирования,
а потому здесь не приводится. 

\begin{theorem} 
Если $H(X)>C$,
то не существует никакого метода кодирования, 
который позволяет передать информацию
по каналу со сколь угодно малой вероятностью ошибки $\varepsilon$. 
\end{theorem}
  
\end{frame}


\subsection{6.2. Алфавитное кодирование}
\label{6.2.}
\begin{frame}
\frametitle{6.2. Алфавитное кодирование}
Кодирование $F$ может сопоставлять код всему сообщению
из множества $S$ как единому целому или же
строить код сообщения из кодов его частей.
Элементарной частью сообщения
является одна буква алфавита $A$.
Этот простейший случай рассматривается в этом
и следующих двух разделах.

\bigskip
\begin{tabular}{p{11mm}|p{38mm}|p{97mm}}
  № & Название параграфа
  & Ключевые термины и обозначения \\
  \hline
\DMref{6.2.1.}{6.2.1.} & Таблица кодов &
{Алфавитное кодирование, схема кодирования, 
множество элементарных кодов, двоично-десятичное кодирование }
\\  
\DMref{6.2.2.}{6.2.2.} & Разделимые схемы &
{ }\\  
\DMref{6.2.3.}{6.2.3.} & Префиксные и постфиксные схемы &
{ }\\  
\DMref{6.2.4.}{6.2.4.}\RS & Неравенство Макмиллана &
{}\\  
\end{tabular}  
\end{frame}


\begin{frame}
\label{6.2.1.}\subsubsection{6.2.1.  Таблица кодов}
\frametitle{6.2.1. Таблица кодов}
\odmkey{Алфавитное}{Кодирование (coding)!алфавитное (alphabetic)}
(или \odmkey{побуквенное}{Кодирование (coding)!побуквенное
(alphabetic)}) кодирование задается \odmkey{схемой}{Схема
(scheme)!кодирования (coding)} (или \odmkey{таблицей
кодов}{Таблица (table)!кодов (of codes)}) $\sigma$:\\
$\sigma\Def \lrle{a_1\to \beta_1,\dots, {a_n}\to \beta_n},\quad 
a_i \in A, \beta_i\in {B^*}.$
Множество кодов букв $V\Def \{\beta_i\}$ называется множеством
\odmkey{элементарных
кодов}{Код (code)!элементарный (elementary)}
(множеством \odmkey{кодовых слов}{Кодовое слово (code word)}).
Алфавитное кодирование пригодно для любого множества сообщений, поскольку
$\Map{F}{A^*}{B ^*}$,  и если \\
$\dTuple{a}{i_1}{i_k}=\alpha\in {A^*}$, 
то 
$F(\alpha)\Def \dTuple{\beta}{i_1}{i_k}$.

\begin{remark}
Ради общности в определениях 
указывают множества $A^*$ и $B^*$, 
содержащие пустое слово $\varepsilon$, хотя на практике,
чтобы передать информацию об отсутствии сообщения,
обычно используют какое-то условленное непустое сообщение,
а потому никогда не возникает нужды в пустых кодовых словах.  
\end{remark}

\begin{example}
Рассмотрим алфавиты $A\Assign \{0,1,2,3,4,5,6,7,8,9\}$, $B\Assign \{0,1\}$ и
схему\\
$\sigma_1\Assign \lrle{0\to
0, 1\to 1, 2\to 10, 3\to 11, 4\to 100, 5\to 101, 
6\to 110, 7\to 111, \break 8\to 1000, 9\to 1001}$.
Эта схема однозначна, но кодирование не
является взаимно-одно\-значным:
$F_{\sigma_1}(333)=111111=F_{\sigma_1}(77)$,
а значит, декодирование невозможно. С другой стороны, схема
$
\sigma_2\Assign \lrle{0\to
0000, 1\to 0001, 2\to 0010, 3\to 0011,
 4\to 0100, \linebreak 5\to 0101, 
6\to 0110, 7\to 0111, 8\to 1000, 9\to 1001}$,
известная под названием
<<\odmkey{двоично-десятичное кодирование}{Кодирование
 (coding)!двоично-десятичное (binary coded decimal)}>>,
допускает однозначное декодирование.
\end{example}

\end{frame}


\begin{frame}
\label{6.2.2.}\subsubsection{6.2.2.  Разделимые схемы}
\frametitle{6.2.2. Разделимые схемы}

Рассмотрим схему алфавитного кодирования $\sigma$ и различные слова,
составленные из элементарных кодов.
Схема $\sigma$ называется
\odmkey{разделимой}{Схема (scheme)!разделимая (separable)}, если
$$
\beta_{i_1}\dots \beta_{i_k} = \beta_{j_1}\dots \beta_{j_l}\Impl
k = l \And \A{t\in 1..k}{i_t = j_t},
$$
то~есть любое слово, составленное из элементарных кодов,
\em{единственным} образом разлагается на
элементарные коды.
Алфавитное кодирование с разделимой схемой 
допускает однозначное декодирование. 

\smallskip
Если таблица кодов содержит одинаковые элементарные
коды, то~есть если
$$
\E{i,j}{i \not= j\And\beta_i = \beta_j},
$$
где
${\beta_i},{\beta_j}\in V$, то схема заведомо не является
разделимой. Такие схемы далее не рассматриваются, 
то~есть $$ \A{i
\not= j}{\beta_i,\beta_j\in V \Impl \beta_i \not= \beta_j}. $$

\begin{remark}
Основные утверждения этого раздела установлены также в \DMref{п.~6.1.2}{6.1.2.} с использованием 
бинарного кодового дерева.
Здесь рассуждения проведены повторно с использованием
алгебраических выкладок и без аппеляции к геометрической 
интерпретации.  
\end{remark}

\end{frame}

\subsubsection{6.2.3.  Префиксные и постфиксные схемы}
\label{6.2.3.}
\begin{frame}
\frametitle{6.2.3. Префиксные и постфиксные схемы (1/2)}

Схема $\sigma$ называется \odmkey{префиксной}{Схема (scheme)!префиксная
 (prefix)}, если
элементарный код одной буквы не является собственным префиксом элементарного кода
другой буквы:
$$
\Not\E{\beta_i,\beta_j\in V,\beta\in B^*}{i\ne
j\And\beta_i=\beta_j\beta}.
$$
\begin{theorem} Префиксная схема является разделимой.
\end{theorem}
\begin{proof}
От противного. Пусть кодирование с префиксной
схемой $\sigma$ не является
разделимым. Тогда существует такое слово
$\beta\in F_\sigma (A^*)$,
что
$$
\beta=\beta_{i_1}\dots\beta_{i_k}=
\beta_{j_1}\dots\beta_{j_l}\And \E{t}{\A {s<t}
{\beta_{i_s}=\beta_{j_s}\And \beta_{i_t}\not=\beta_{j_t}}}.
$$
Поскольку $\beta_{i_t}\dots\beta_{i_k}=\beta_{j_t}\dots\beta_{j_l}$,
значит,
$\E{\beta'}{\beta_{i_t}=\beta_{j_t}\beta'\vee\beta_{j_t}=\beta_{i_t}\beta'}$,
но это противоречит тому, что схема префиксная.
\end{proof}

\begin{remark}
Свойство быть префиксной  достаточно, но не необходимо для
разделимости схемы.
\end{remark}

\begin{example}
Разделимая, но не префиксная схема: $$A=\{a,b\},\quad
B=\{0,1\},\quad \sigma={\langle a\to 0, b\to 01\rangle}.$$
\end{example}
\end{frame}

\begin{frame}
\frametitle{6.2.3. Префиксные и постфиксные схемы (2/2)}
Схема $\sigma$ называется \odmkey{постфиксной}{Схема (scheme)!постфиксная (postfix)}, если
элементарный код одной буквы не является собственным постфиксом элементарного кода
другой буквы:
$$
\Not\E{\beta_i,\beta_j\in V,\beta\in B^*}{i\ne
j\And\beta_i=\beta\beta_j}.
$$
\begin{theorem} Постфиксная схема является разделимой.
\end{theorem}
\begin{proof}
Аналогично предыдущему, только 
рассматривать слово $\beta$
надлежит справа, а не слева.
\end{proof}

\smallskip
Очень существенным на практике 
является случай, когда все кодовые слова имеют 
одинаковую длину $l$. 
Такая схема алфавитного кодирования
называется 
\odmkey{равномерной}{Схема (scheme)!равномерная
 (uniform)}.
 
\begin{corollary}
Равномерная схема алфавитного кодирования разделима.
\end{corollary} 
\begin{proof}
Совпадения кодовых слов запрещены,
поэтому всякая равномерная схема является как префиксной, так и постфиксной, поскольку 
никакое слово длины $l$ не может быть собственным префиксом или 
постфиксом другого слова длины $l$. 
\end{proof}

\begin{remark}
Таким образом, класс разделимых схем алфавитного кодирования
является объединением классов префиксных и постфиксных схем,
а класс равномерных схем является пересечением классов
префиксных и постфиксных схем.  
\end{remark}

\end{frame}

\begin{frame}
\label{6.2.4.}\subsubsection{6.2.4.  Неравенство Макмиллана }\frametitle{6.2.4. Неравенство Макмиллана (1/6)}

Чтобы схема алфавитного кодирования была разделимой, необходимо,
чтобы длины элементарных кодов удовлетворяли определённому
соотношению, известному как \odmkey{неравенство
Макмиллана}{Неравенство (inequality)!МакМиллана
(Kraft--MacMillan)}.

\begin{remark}
Употребляется также название <<неравенство Крафта--Макмиллана>>.
\end{remark}

\begin{theorem} Если схема
$\sigma={\lrle{a_i\to\beta_i}}_{i=1}^n$ разделима, то
$\sum\limits_{i=1}^n 2^{-|\beta_i|} \leq 1$.
\end{theorem}
\begin{proof} Обозначим $l_i\Assign |\beta_i|$,
$l\Assign\max\limits_{i\in 1..n}l_i$.
Рассмотрим некоторую $k$-ю степень левой части
неравенства $\left(\sum\limits^{n}_{i=1}2^{-l_i}\right)^k$.
Раскрывая скобки, имеем сумму
$\sum\limits_{(\Tuple{i}{1}{k})}({2^{l_{i_1}+\ldots+l_{i_k}}})^{-1}$,
где $\Tuple{i}{1}{k}$~--- различные наборы номеров элементарных
кодов. Обозначим через $\nu(k,t)$ количество входящих в эту сумму
слагаемых вида $1/2^t$, где $t=l_{i_1}+\ldots+l_{i_k}$. Для
некоторых $t$ может быть, что $\nu(k,t)=0$. Приводя подобные,
имеем сумму $\sum\limits^{kl}_{t=1}\frac{\nu(k,t)}{2^t}$.
Каждому слагаемому вида $({2^{l_{i_1}+\ldots +l_{i_k}}})^{-1}$
можно сопоставить код (слово в алфавите $B$) вида
$\beta_{i_1}\dots\beta_{i_k}$. Это слово состоит из $k$
элементарных кодов и имеет длину $t$.
\end{proof}
\end{frame}

\begin{frame}
\frametitle{6.2.4. Неравенство Макмиллана (2/6)}
\begin{proof} (Продолжение)
Таким образом, $\nu(k,t)$~--- это число некоторых слов вида
$\beta_{i_1}\dots\beta_{i_k}$, таких, что
$|\beta_{i_1}\dots\beta_{i_k}|=t$. В силу разделимости схемы
$\nu(k,t)\le 2^t$, в противном случае заведомо существовали бы два
одинаковых слова
$\beta_{i_1}\dots\beta_{i_k}=\beta_{j_1}\dots\beta_{j_k}$,
допускающих различное разложение. Имеем
$\sum\limits^{kl}_{t=1}\frac{\nu(k,t)}{2^t}\le\sum\limits^{kl}_{t=1}\frac{2^t}{2^t}=
kl$.
Следовательно, $ \A{k}{\left(\sum\limits^n_{i=1}2^{-l_i}\right)^k\le kl}$,
и, значит,
$\A{k}{\sum\limits^n_{i=1}2^{-l_i}\le\root k \of {kl}}$,
откуда
$
\sum\limits^{n}_{i=1}2^{-l_i}\le
\lim\limits_{k\to \infty}\root k \of
{kl}=1$.
\end{proof}
Неравенство Макмиллана является не только необходимым, но и
в некотором смысле достаточным условием разделимости схемы алфавитного
кодирования.
\begin{theorem}
Если числа $\Tuple{l}{1}{n}$ удовлетворяют неравенству
$\sum\limits_{i=1}^n 2^{-l_i}\leq 1$,
то существует такая разделимая (даже префиксная) схема алфавитного кодирования
$\sigma={\lrle{a_i\to\beta_i}}_{i=1}^n$,
что
$\A{i}{\vert\beta_i\vert = l_i}$.
\end{theorem}
\end{frame}

\begin{frame}
\frametitle{6.2.4. Неравенство Макмиллана (3/6)}

\begin{proof}
Без ограничения общности можно считать, что $l_1\le
l_2\le\dots \le l_n$. Разобьем множество $\{\Tuple{l}{1}{n}\}$
на классы эквивалентности по равенству $\{\Tuple{L}{1}{m}\}$,
$m\le n$. Пусть $\lambda_i\in L_i,\; \mu_i\Assign |L_i|$. Тогда $
 \sum\limits^m_{i=1}\mu_i=n,\;
\lambda_1 <\lambda_2 <\dots  <\lambda_m. $
\pagebreak
В этих обозначениях неравенство Макмиллана
можно записать так:
$\sum\limits^m_{i=1}\frac{\mu_i}{2^{\lambda_i}}\le 1$.
Из этого неравенства следуют $m$ неравенств для частичных сумм:

$\frac{\mu_1}{2^{\lambda_1}}\le 1\Impl\mu_1\le 2^{\lambda_1}$,

$\frac{\mu_1}{2^{\lambda_1}}+\frac{\mu_2}{2^{\lambda_2}}\le1
\Impl\mu_2\le 2^{\lambda_2}-\mu_12^{\lambda_2-\lambda_1}$,

$\dots$

$\frac{\mu_1}{2^{\lambda_1}}+\frac{\mu_2}{2^{\lambda_2}}+\ldots
+\frac{\mu_m}{2^{\lambda_m}}\le 1\Impl$ 
$ \mu_m 
 \leq 2^{\lambda_m}-\mu_12^{\lambda_m-\lambda_1}-
\mu_22^{\lambda_m-\lambda_2}-\ldots
-\mu_{m-1}2^{\lambda_m-\lambda_{m-1}}$.

\vspace{0.5\baselineskip}
Рассмотрим слова длины $\lambda_1$ в алфавите $B$. Поскольку
$\mu_1\le 2^{\lambda_1}$, из этих слов можно выбрать $\mu_1$
различных слов $\Tuple{\beta}{1}{\mu_1}$ длины $\lambda_1$.
Исключим из дальнейшего рассмотрения все слова из $B^*$,
начинающиеся со слов $\Tuple{\beta}{1}{\mu_1}$. 
\end{proof}

\end{frame}

\begin{frame}
\frametitle{6.2.4. Неравенство Макмиллана (4/6)}
\begin{proof} (Продолжение)
Далее рассмотрим
множество слов в алфавите $B$ длиной $\lambda_2$, не начинающихся
со слов $\Tuple{\beta}{1}{\mu_1}$. Таких слов будет
$2^{\lambda_2}-\mu_12^{\lambda_2-\lambda_1}$. 

Но $\mu_2\le
2^{\lambda_2}-\mu_12^{\lambda_2-\lambda_1}$, значит, можно выбрать
$\mu_2$ \em{различных} слов. Обозначим их\\
$\Tuple{\beta}{\mu_1+1}{\mu_1+\mu_2}$. Исключим слова,
начинающиеся с $\Tuple{\beta}{\mu_1+1}{\mu_1+\mu_2}$, из
дальнейшего рассмотрения. И~далее, используя неравенства для
частичных сумм, мы будем на $i$-м шаге выбирать $\mu_i$ слов длины
$\lambda_i$, $\Tuple{\beta}{\mu_1+\mu_2+\ldots
+\mu_{i-1}\!+\!1}{\mu_1+\mu_2+\ldots +\mu_{i-1}+\mu_i}$, причём эти
слова не будут начинаться с тех слов, которые были выбраны раньше.
В то же время длины этих слов всё время растут (так~как
$\lambda_1<\lambda_2<\ldots <\lambda_m$), поэтому они не могут
быть префиксами тех слов, которые выбраны раньше. Итак, в конце
имеем набор из $n$ слов $\Tuple{\beta}{1}{\mu_1+\ldots
+\mu_m}=\beta_n$, $|\beta_1|=l_1,\dots,|\beta_n|=l_n$, коды
$\Tuple{\beta}{1}{n}$ не являются префиксами друг друга, а значит,
схема $\sigma=\lrle{a_i\to\beta_i}^n_{i=1}$ будет префиксной и, по
теореме предыдущего подраздела, разделимой.
\end{proof}

\begin{remark}
Неравенство Макмиллана справедливо
не только для двоичного, но и для 
$m$-ичного кодирования.
\end{remark}

\end{frame}

\begin{frame}
\frametitle{6.2.4. Неравенство Макмиллана (5/6)}

\begin{example}
\odmkey{Азбука Морзе
%\footnote{\;Самуэль Морзе
%(1791--1872).}
}{Азбука Морзе (Morse alphabet)}~--- это схема
алфавитного кодирования


\begin{tabular}{lllll}
$\langle A\to 01$, & $B\to 1000$,  & $C\to 1010$,  & $D\to 100$,  & $E\to 0$, \\
$F\to 0010$,       & $G\to 110$,   & $H\to 0000$,  & $I\to 00$,   & $J\to 0111$, \\
$K\to 101$,        & $L\to 0100$,  & $M\to 11$,    & $N\to 10$, & $O\to 111$, \\        $P\to 0110$,       & $Q\to 1101$,  & $R\to 010$,   & $S\to 000$,  & $T\to 1$, \\
$U\to 001$,        & $V\to 0001$,  & $W\to 011$,   & $X\to 1001$, & $Y\to 1011$, \\
$Z\to 1100\rangle$,&               &               &              & \\   
\end{tabular}

где по историческим и техническим
причинам 0 называется \em{точкой\/} и обозначается
знаком <<$\cdot$>>, а 1 называется \em{тире\/} и обозначается знаком~<<--->>.
Имеем

\begin{tabular}{lllllllllll}
$\phantom{+}$
       & $1/4$  & $+$ & $1/16$ & $+$ & $1/16$ & $+$ & $1/8$  & $+$ & $1/2$  & $+$ \\
$+$    & $1/16$ & $+$ & $1/8$  & $+$ & $1/16$ & $+$ & $1/4$  & $+$ & $1/16$ & $+$ \\
$+$    & $1/8$  & $+$ & $1/16$ & $+$ & $1/4$  & $+$ & $1/4$  & $+$ & $1/8$  & $+$ \\ 
$+$    & $1/16$ & $+$ & $1/16$ & $+$ & $1/8$  & $+$ & $1/8$  & $+$ & $1/2$  & $+$ \\
$+$    & $1/8$  & $+$ & $1/16$ & $+$ & $1/8$  & $+$ & $1/16$ & $+$ & $1/16$ & $+$ \\
$+$    & $1/16$ & $=$ &  $2/2$  & $+$ & $4/4$  & $+$ & $8/8$  & $+$ & $12/16$ &$=$  \\
$=$    & $3$    & $+$ & $3/4$  & $>$ & $1$.   &     &         &     &   & \\ 
\end{tabular}
\end{example}
\end{frame}

\begin{frame}
\frametitle{6.2.4. Неравенство Макмиллана (6/6)}

Таким образом, неравенство Макмиллана для азбуки Морзе не
выполнено, и эта схема не является разделимой. 
Более того, в азбуке Морзе многие кодовые слова являются префиксами других
слов, а потому некоторые комбинации букв неразличимы.

\begin{examples}
$AM=0111=J,\quad TO=1111=MM,\quad WMW=01111011=JY$.
\end{examples}

\smallskip
На самом деле в азбуке
Морзе имеются дополнительные элементы~--- \em{паузы} между буквами (и словами),
которые позволяют декодировать сообщения. Эти дополнительные
элементы определены неформально, поэтому приём и передача сообщений
с помощью азбуки Морзе, особенно с высокой скоростью,
являлись некоторым искусством, а не
простой технической процедурой.


\begin{corollary}
Если схема алфавитного кодирования
$\sigma={\lrle{a_i\to\beta_i}}_{i=1}^n$ разделима,
то существует префиксная схема
$\sigma^\prime={\lrle{a_i\to\beta_i^\prime}}_{i=1}^n$,
причём $\A{i}{\vert\beta_i\vert=\vert\beta_i^\prime\vert}$.
\end{corollary}

\begin{example}
Схема $\langle a\to 0, b\to 01\rangle$~--- разделимая, но не
префиксная, а схема ${\langle a\to 0, b\to 10\rangle}$~---
префиксная (и разделимая).
\end{example}

\end{frame}


\subsection{6.3. Кодирование с минимальной избыточностью}
\label{6.3.}
\begin{frame}
\frametitle{6.3. Кодирование с минимальной избыточностью}

Для практики важно, чтобы коды сообщений имели
по возможности наименьшую длину.
Алфавитное кодирование пригодно для любых сообщений,
то есть для $S=A^*$. 
Если больше про множество $S$
ничего не известно, то точно сформулировать задачу
оптимизации затруднительно. Однако на практике часто доступна
дополнительная информация.
Например, для текстов на естественных языках
известно распределение вероятности появления
букв в сообщении. Использование такой информации
позволяет строго поставить
и решить задачу построения оптимального алфавитного кодирования.

\bigskip
\begin{tabular}{p{11mm}|p{38mm}|p{97mm}}
  № & Название параграфа
  & Ключевые термины и обозначения \\
  \hline
\DMref{6.3.1.}{6.3.1.} & Минимизация длины кода сообщения &
{ }
\\  
\DMref{6.3.2.}{6.3.2.} & Цена кодирования &
{Равномерное кодирование, оптимальное кодирование}\\  
\DMref{6.3.3.}{6.3.3.} & Алгоритм Фано &
{Медиана }\\  
\DMref{6.3.4.}{6.3.4.}\RS & Оптимальное кодирование &
{}\\  
\DMref{6.3.5.}{6.3.5.}\RS & Алгоритм Хаффмена &
{}\\ 
\end{tabular}  
\end{frame}


\begin{frame}
\label{6.3.1.}
\subsubsection{6.3.1. Минимизация длины кода сообщения}
\frametitle{6.3.1. Минимизация длины кода сообщения (1/2)}

Если задана разделимая схема алфавитного кодирования
$\sigma={\lrle{a_i\to\beta_i}}_{i=1}^n$, 
то любая схема
$\sigma'={\lrle{a_i\to{\beta_i}'}}_{i=1}^n$, где
последовательность $\lrle{\Tuple{\beta'}{1}{n}}$ является
перестановкой последовательности $\lrle{\Tuple{\beta}{1}{n}}$,
также будет разделимой. Если длины элементарных кодов равны, то
перестановка элементарных кодов в схеме не влияет на длину кода
сообщения. Но если длины элементарных кодов различны, то длина
кода сообщения зависит от состава букв в сообщении и от того,
какие элементарные коды каким буквам назначены. Если заданы
конкретное сообщение и конкретная схема кодирования, то нетрудно
подобрать такую перестановку элементарных кодов, при которой длина
кода сообщения будет минимальна. Пусть $\Tuple{k}{1}{n}$~---
количества вхождений букв $\Tuple{a}{1}{n}$ в сообщение $s\in S$,
а~$\Tuple{l}{1}{n}$~--- длины элементарных кодов
$\Tuple{\beta}{1}{n}$ соответственно. Тогда,
если $k_i\le k_j$ и $l_i \ge l_j$, то
$k_i l_i + k_j l_j \le k_i l_j + k_j l_i $.
Действительно, пусть $k_j=k+a$, $k_i=k$ и
$l_j=l$, $l_i=l+b$, где $a,b\ge 0$. Тогда
$ (k_i l_j + k_j l_i) - (k_i l_i + k_j l_j) =$\\
$=(kl+(k+a)(l+b)) -(k(l+b)+l(k+a)) =$\\ 
$=(kl+al+bk+ab+kl) - (kl+al+kl+bk)= ab \ge 0$.
\end{frame}


\begin{frame}
\frametitle{6.3.1.Минимизация длины кода сообщения (2/2)}
Отсюда вытекает алгоритм назначения
элементарных кодов, при котором длина кода конкретного
сообщения $s\in S$ будет минимальна: нужно отсортировать
буквы сообщения $s$ в порядке убывания количества вхождений,
элементарные коды отсортировать
в~порядке возрастания длины и назначить коды
буквам в~этом порядке.

\medskip
\begin{example}
Для алфавита из трёх букв кратчайшая разделимая схема
имеет вид 
$$\langle x\to 0, y\to 10 , z\to 11\rangle
\text{или} \langle x\to 1, y\to 01 , z\to 00\rangle,$$
поскольку схемы с более короткими кодовыми словами не 
удовлетворяют неравенству Макмиллана.
Если необходимо передать сообщение <<$SOS-SOS-SOS$>>,
то самой выгодной является схема 
$$\langle S\to 0, O\to 10 , - \to 11\rangle.$$
\end{example}

\begin{remark}
Этот простой метод решает задачу минимизации
длины кода только для фиксированного сообщения $s\in S$
и фиксированной схемы $\sigma$.
\end{remark}

\end{frame}

\begin{frame}
\label{6.3.2.}
\subsubsection{6.3.2. Цена кодирования}
\frametitle{6.3.2. Цена кодирования (1/3)}

Пусть заданы алфавит $A=\{\Tuple{a}{1}{n}\}$ и вероятности
появления букв в сообщении  $P=\lrle{\Tuple{p}{1}{n}}$ ($p_i$~---
вероятность появления буквы ${a_i}$). Не ограничивая общности,
можно считать, что ${p_1}+\ldots+{p_n}=1$ и
${p_1}\geq\ldots\geq{p_n}>0$ (то~есть можно сразу исключить буквы,
которые не могут появиться в сообщении, и упорядочить буквы по
убыванию вероятности их появления). Для каждой (разделимой) схемы
$\sigma={\lrle{a_i\to\beta_i}}_{i=1}^n$ алфавитного кодирования
математическое ожидание коэффициента увеличения длины сообщения
при кодировании $\sigma$ (обозначается $l_\sigma$) определяется
следующим образом:
$$ l_{\sigma}(P)\Def \sum_{i=1}^n{p_i}{l_i},
\qquad\text{где}\quad {l_i}\Def |\beta_i|,
$$
и называется средней
\odmkey{ценой}{Кодирование (coding)!цена (price)} (или
\odmkey{длиной}{Длина (length)!кодирования (of coding))})
кодирования $\sigma$ при
распределении вероятностей $P$.

\begin{example}
Для разделимой схемы $A=\{a,b\}$, $B=\{0,1\}$, $\sigma =\lrle{a\to 0,
b\to 01}$ 
при рас\-пределении вероятностей $\lrle{0,5;0,5}$ цена
кодирования составляет ${0,5*1} +{0,5 *2} =1,5$, а при
распределении вероятностей $\lrle{0,9;0,1}$ она равна
$0,9 *1 +0,1 *2 =1,1$.
\end{example}
\end{frame}


\begin{frame}
\frametitle{6.3.2. Цена кодирования (2/3)}
Введём обозначения:
$$
{l_*}(P)\Def \inf_\sigma{l_\sigma}(P), \qquad p_*\Def \min_{i\in 1..n}{p_i}, \qquad
L\Def \Lrf{\log(n-1)}+1.
$$
%\addvspace{-6pt}
Очевидно, что всегда существует разделимая схема
$\sigma={\lrle{a_i\to\beta_i}}_{i=1}^n$, такая, что\\
$\A{i}{|\beta_i|=L}$. Такая схема называется схемой
\odmkey{равномерного}{Кодирование (coding)!равномерное (uniform)}
кодирования. Следовательно, \\$1\leq{l_*}(P)\leq L$, и достаточно
учитывать только такие схемы, для которых $\A{i}{p_i l_i\leq
L}$, где $l_i$~--- целое и 
$l_i\leq L/{p_*}$. Таким образом,
имеется лишь конечное число схем $\sigma$, для которых
${l_*}(P)\leq {l_\sigma}(P)\leq L$. Значит, существует схема
$\sigma_*$, на которой инфимум достигается:
$l_{\sigma_*}(P)=l_*(P)$.

\medskip
Алфавитное (разделимое) кодирование $\sigma_*$, для
которого $l_{\sigma_*}(P)=l_*(P)$, называется кодированием с
\odmkey{минимальной избыточностью}{Кодирование (coding)!с
минимальной избыточностью (with minimal redundancy)}, или
\odmkey{оптимальным}{Кодирование (coding)!оптимальное (optimal)}
кодированием, для распределения вероятностей $P$.
\begin{example}
Пусть в алфавите 8 букв. 
Рассмотрим цену равномерного кодирования $C_1$ и 
<<наиболее неравномерного>> кодирования $C_2$, 
при равномерном $P_1$ и неравномерном $P_2$ распределении вероятностей. 
\end{example}
\end{frame}


\begin{frame}
\frametitle{6.3.2. Цена кодирования (3/3)}

\vspace{0.5\baselineskip}

\begin{tabular}{lrrrrr|lrrrrr}
$C_1$  & $l_1$ & $P_1$ & $L_{11}$ & $P_2$ & $L_{12}$ &
$C_2$  & $l_2$ & $P_1$ & $L_{21}$ & $P_2$ & $L_{22}$ \\
\hline
000 & 3 & 0,125 & 0,375 & 0,31 & 0,93 & 0 & 1 & 0,125 & 0,125 & 0,31 & 0,31 \\
001 & 3 & 0,125 & 0,375 & 0,24 & 0,72 & 10 & 2 & 0,125 & 0,25 & 0,24 & 0,48 \\
010 & 3 & 0,125 & 0,375 & 0,17 & 0,51 & 110 & 3 & 0,125 & 0,375 & 0,17 & 0,51 \\
011 & 3 & 0,125 & 0,375 & 0,11 & 0,33 & 1110 & 4 & 0,125 & 0,5 & 0,11 & 0,44 \\
100 & 3 & 0,125 & 0,375 & 0,09 & 0,27 & 11110 & 5 & 0,125 & 0,625 & 0,09 & 0,45 \\
101 & 3 & 0,125 & 0,375 & 0,05 & 0,15 & 111110 & 6 & 0,125 & 0,75 & 0,05 & 0,3 \\
110 & 3 & 0,125 & 0,375 & 0,02 & 0,06 & 1111110 & 7 & 0,125 & 0,875 & 0,02 & 0,14 \\
111 & 3 & 0,125 & 0,375 & 0,01 & 0,03 & 11111110 & 8 & 0,125 & 1 & 0,01 & 0,08 \\
\hline
 &  & 1 & 3 & 1 & 3 &  &  & 1 & 4,5 & 1 & 2,71 \\
\end{tabular}

\vspace{0.5\baselineskip}
Из таблицы видно, что при равномерном кодировании 
распределение вероятностей не влияет на цену кодирования,
в то время как неравномерное кодирование 
по сравнению с равномерным кодированием 
существенно проигрывает при равномерном распределении вероятностей
и несколько выигрывает при неравномерном распределении вероятностей.  

\end{frame}

\begin{frame}
\label{6.3.3.}
\subsubsection{6.3.3. Алгоритм Фано }
\frametitle{6.3.3. Алгоритм Фано (1/4)}

Рекурсивный \odmkey{алгоритм Фано}{Алгоритм (algorithm)!Фано (Fano)}
строит разделимую префиксную схему алфавитного кодирования,
близкого к оптимальному.

\begin{algorithmic}
\REQUIRE \textbf{$P:$ array $[1..n]$ of real}~--- массив
вероятностей появления букв в сообщении, упорядоченный по
невозрастанию;
\ENSURE \textbf{$C:$ array $[1..n,1..L]$ of
$0..1$}~--- массив элементарных кодов. 
\STATE $\Fano(1,n,0)$
\COMMENT{ вызов рекурсивной процедуры $\Fano$ }
\end{algorithmic}

\medskip
Алгоритм Фано использует функцию $\Med$, 
которая находит \odmkey{медиану}{Медиана (median)!множества
(of set)} указанной части массива $P[b..e]$, то~есть определяет
такой индекс $m$ ($b\le m<e$), что сумма элементов $P[b..m]$
наиболее близка к сумме элементов $P[(m+1)..e]$, то~есть
$\displaystyle{\min\limits_{m\in b..(e-1)}\left\vert\sum\limits_{i=b}^{m}P[i]-
\sum\limits_{i=m+1}^{e}P[i]\right\vert}$.
\end{frame}

\begin{frame}
\frametitle{6.3.3. Алгоритм Фано (2/4)}


Рекурсивная процедура $\Fano$ описана на этом слайде.
\begin{algorithmic}
\REQUIRE {$b$~--- индекс начала обрабатываемой части массива $P$,
$e$~--- индекс конца обрабатываемой части массива $P$,
$k$~--- длина уже построенных кодов в обрабатываемой части
массива $C$.}
\ENSURE {заполненный массив $C$.}
\IF {$e > b$}
  \STATE $k \Assign k + 1$ \COMMENT{\color{gray}\small  место для очередного разряда в коде }
  \STATE $m \Assign \Med(b,e)$ \COMMENT{\color{gray}\small  деление массива на две части }
  \FOR{\textbf{$i$ from $b$ to $e$}}
    \STATE $C[i,k] \Assign i > m$
    \COMMENT {\color{gray}\small  в первой части добавляем 0, во второй~--- 1 }
  \ENDFOR
  \STATE $\Fano(b,m,k)$ \COMMENT{\color{gray}\small  обработка первой части }
  \STATE $\Fano(m+1,e,k)$ \COMMENT{\color{gray}\small  обработка второй части }
\ENDIF
\end{algorithmic}

Функция $\Med$ описана на следующем слайде.
\end{frame}


\begin{frame}
\frametitle{6.3.3. Алгоритм Фано (3/4)}

\begin{algorithmic}
\REQUIRE $b$~--- индекс начала обрабатываемой части массива $P$,
$e$~--- индекс конца обрабатываемой части массива $P$.
\ENSURE $m$~--- индекс медианы
\STATE $S_b \Assign 0$
\COMMENT{\color{gray}\small сумма элементов первой части }
\FOR{\textbf{$i$ from $b$ to $e-1$}}
  \STATE $S_b \Assign S_b + P[i]$
  \COMMENT{\color{gray}\small вначале все, кроме последнего }
\ENDFOR
\STATE $S_e \Assign P[e]$
\COMMENT{\color{gray}\small сумма элементов второй части }
\STATE $m \Assign e$
\COMMENT{\color{gray}\small начинаем искать медиану с конца }
\REPEAT
  \STATE $d \Assign \vert S_b - S_e\vert$
  \COMMENT{\color{gray}\small разность сумм первой и второй частей }
  \STATE $m \Assign m-1$
  \COMMENT{\color{gray}\small сдвигаем границу медианы вниз }
  \STATE $S_b \Assign S_b - P[m]; S_e \Assign S_e + P[m]$
  \COMMENT{\color{gray}\small перевычисляем суммы }
\UNTIL {$\vert S_b - S_e\vert \ge d$}
\STATE \textbf{return $m$}
\end{algorithmic}

\end{frame}

\begin{frame}
\frametitle{6.3.3. Алгоритм Фано (4/4)}
\vspace{-2pt}
\begin{ground}
При каждом удлинении кодов в одной части коды удлиняются нулями, а в
другой~--- единицами. Таким образом, коды одной части не могут быть
префиксами другой. Удлинение кода заканчивается тогда и только
тогда, когда длина части равна 1, то есть остается единственный код.
Таким образом, схема по построению префиксная, а потому
разделимая.
\end{ground}

\begin{example}
Коды, построенные алгоритмом Фано для заданного 
неравномерного распределения
вероятностей ($n=8$).

\addvspace{-2pt}
{\renewcommand{\baselinestretch}{0.9}\selectfont
\begin{tabular}{r|lllllll|rr}
$p_i$ & &  &  &  &  &  &  & $l_i$ & $p_i l_i$ \\
\hline 
0,31 & 0,31 &  &  &  &  &  & 0 & 1 & 0,31 \\
0,24 & 0,69 &  & 0,24 &  &  &  & 100 & 3 & 0,72 \\
0,17 &  & 0,41 & 0,17 &  &  &  & 101 & 3 & 0,51 \\
0,11 &  & 0,28 & 0,11 &  &  &  & 110 & 3 & 0,33 \\
0,09 &  &  & 0,17 & 0,09 &  &  & 1110 & 4 & 0,36 \\
0,05 &  &  &  & 0,08 & 0,05 &  & 11110 & 5 & 0,25 \\
0,02 &  &  &  &  & 0,03 & 0,02 & 111110 & 6 & 0,12 \\
0,01 &  &  &  &  &  & 0,01 & 111111 & 6 & 0,06 \\
\hline
 &  &  &  &  &  &  &  &  & 2,66 \\
\end{tabular}
}
\end{example}

\end{frame}


\begin{frame}
\label{6.3.4.}\subsubsection{6.3.4. Оптимальное кодирование }\frametitle{6.3.4. Оптимальное кодирование (1/5)}
Оптимальные схемы алфавитного кодирования обладают определёнными
структурными свойствами, устанавливаемыми в следующих леммах.
Нарушение этих свойств для какой-либо схемы означает,
что схема не оптимальна.

\begin{lemma} {\NB{[1]}}
Пусть $\sigma = {\lrle{a_i\to\beta_i}}_{i=1}^n$~{\upshape---}
схема оптимального кодирования для распределения вероятностей
$P={p_1}\geq\ldots\geq{p_n}>0$. Тогда если $p_i > p_j$, то
$l_i~\le~l_j$.
\end{lemma}
\begin{proof}
От противного. Пусть $i<j$, $p_i>p_j$ и $l_i>l_j$. Тогда
рассмотрим схему $ \sigma'=\{a_1\to\beta_1,\dots
,a_i\to\beta_j,\dots,a_j\to\beta_i,\dots ,a_n\to\beta_n\}. $
Имеем $$
l_\sigma-l_{\sigma'} =(p_il_i+p_jl_j)-(p_il_j+p_jl_i)=(p_i-p_j)(l_i-l_j)>0
,$$ что противоречит тому, что схема $\sigma$~--- оптимальна.
\end{proof}

Таким образом, не ограничивая общности, можно считать, что
$l_1\leq\ldots\leq l_n$.

\begin{remark}
Фактически, лемма 1 повторяет для вероятностей то, что было
обнаружено для частот в~\DMref{п.~6.3.1}{6.3.1.}.
\end{remark}

\end{frame}


\begin{frame}
\frametitle{6.3.4. Оптимальное кодирование (2/5)}

\begin{lemma} {\NB{[2]}}
Если $\sigma={\lrle{a_i\to\beta_i}}_{i=1}^n$~{\upshape---} схема
оптимального префиксного кодирования для распределения
вероятностей $P={p_1}\geq\ldots\geq{p_n}>0$, то среди элементарных
кодов имеются два кода максимальной длины, которые различаются
только в~последнем разряде.
\end{lemma}

\begin{proof}
По предыдущей лемме кодовые слова максимальной длины являются
последними в схеме. Далее от противного. 
\proofsection{1} Пусть
кодовое слово максимальной длины одно и имеет вид $\beta_n=\beta
b$, \\где ${b=0\Or b=1}$. Имеем $\A{i\in 1 .. n-1}{l_i\le|\beta|}$.
Так~как схема префиксная, то слова $\Tuple{\beta}{1}{n-1}$ не
являются префиксами $\beta$. С другой стороны, $\beta$ не является
префиксом слов $\Tuple{\beta}{1}{n-1}$, иначе было бы
$\beta=\beta_j$, а значит, $\beta_j$ было бы префиксом $\beta_n$.
Тогда схема \\$\sigma'\Assign \lrle{a_1\to\beta_1,\dots
,a_n\to\beta}$ тоже префиксная, причём
$l_{\sigma'}(P)=l_\sigma(P)-p_n$, что противоречит оптимальности
$\sigma$. 
\end{proof}

\end{frame}

\begin{frame}
\frametitle{6.3.4. Оптимальное кодирование (3/5)}

\begin{proof} (Окончание)
\proofsection{2} Пусть теперь два кодовых слова
$\beta_{n-1}$ и $\beta_n$ максимальной длины различаются не в
последнем разряде, то~есть $\beta_{n-1}=\beta'b'$,
$\beta_n=\beta''b''$, $\beta'\not=\beta''$, причём $\beta'$,
$\beta''$ не являются префиксами для $\Tuple{\beta}{1}{n-2}$ и
наоборот.
 Тогда схема \\$\sigma' \Assign
{\lrle{a_1\to\beta_1,\dots
,a_{n-2}\to\beta_{n-2},a_{n-1}\to\beta'b',a_n\to\beta''}}$ также
является префиксной, причём $l_{\sigma'}(P)=l_\sigma(P)-p_n$, что
противоречит оптимальности $\sigma$. 
\proofsection{3} Пусть
кодовых слов максимальной длины больше двух. Тогда среди них
найдутся два, которые различаются не только в последнем разряде,
что противоречит уже доказанному пункту 2.
\end{proof}

\begin{remark}
В примере \DMref{п.~6.3.3}{6.3.3.} имеется два элементарных
кода максимальной длины, различающиеся в 
последнем разряде, но в примере \DMref{п.~6.3.5}{6.3.5.}
показано, что кодирование Фано не оптимально.  
Это говорит о том, что условие леммы 2
необходимо, но не достаточно для оптимальности.
\end{remark}
\end{frame}


\begin{frame}
\frametitle{6.3.4. Оптимальное кодирование (4/5)}

\begin{theorem} Если
$\sigma_{n-1}={\lrle{a_i\to\beta_i}}_{i=1}^{n-1}$~{\upshape---}
схема оптимального префиксного кодирования для распределения
вероятностей $P={p_1}\geq \ldots\geq {p_{n-1}} > 0$ и ${p_j =
{q'}+{q''}}$,\\
причём
$
{p_1}\geq\ldots\geq p_{j-1} \geq p_{j+1}\geq\ldots \geq p_{n-1}
\geq q'\geq q'' >0,
$
то кодирование со схемой

\vspace{-\baselineskip}
\begin{align*}
\sigma_n = \langle &a_1 \to \beta_1, \dots,
a_{j-1} \to \beta_{j-1},a_{j+1} \to \beta_{j+1}, \dots , \\
 & a_{n-1} \to
\beta_{n-1}, a_j \to \beta_j0, a_n \to \beta_j1\rangle
\end{align*}


является оптимальным
префиксным кодированием для распределения вероятностей\\
$P_n=\Tuple{p}{1}{j-1},\Tuple{p}{j+1}{n-1},q',q''$.
\end{theorem}

\begin{proof}
Заметим следующее. 
\proofsection{1} Если $\sigma_{n-1}$ было
префиксным, то $\sigma_n$ тоже будет префиксным по построению.
\proofsection{2} Цена кодирования для схемы $\sigma_n$ больше цены
кодирования для схемы $\sigma_{n-1}$ на величину
$p_j$: $
l_{\sigma_n}(P_n)=p_1l_1+p_2l_2+\ldots
+p_{j-1}l_{j-1}+p_{j+1}l_{j+1}+\ldots 
+p_{n-1}l_{n-1}+$\\ $+q'(l_j+1)+q''(l_j+1)=
p_1l_1+\ldots
+p_{n-1}l_{n-1}+l_j(q'+q'')+(q'+q'')=$\\
$=p_1l_1+\ldots
+p_{n-1}l_{n-1}+l_jp_j+p_j=l_{\sigma_{n-1}}(P_{n-1})+p_j$.
\end{proof}
\end{frame}


\begin{frame}
\frametitle{6.3.4. Оптимальное кодирование (5/5)}
\begin{proof} (Продолжение)
\proofsection{3} Пусть некоторая схема
$\sigma'_n\Assign\lrle{a_i\to\beta_i}^n_{i=1}$ оптимальна для $P_n$.
Тогда по лемме 2 два кода максимальной длины различаются в
последнем разряде: \\$\sigma'_n ={\langle a_1\to\beta'_1,\dots
,a_{n-2}\to{\beta}_{n-2}',a_{n-1}\to\beta 0, a_n\to\beta
1\rangle}$. Положим $l'\Assign |\beta|$ и рассмотрим схему
$\sigma'_{n-1}\Assign \langle a_1\to\beta'_1,\dots
,a_j\to\beta,\dots ,a_{n-2}\to{\beta}_{n-2}'\rangle$, где $j$
определено так, чтобы $p_{j-1}\ge q'+q''\ge p_{j+1}$. Цена
кодирования для схемы $\sigma'_n$ больше цены кодирования для
схемы $\sigma'_{n-1}$ также на величину $p_j$:
\begin{align*}
l_{\sigma'_n}(P_n)&=l_1p_1+\ldots
+l_{n-2}p_{n-2}+(l'+1)q'+(l'+1)q''=\\
&= l_1p_1+\ldots
+l_{n-2}p_{n-2}+l'(q'+q'')+(q'+q'')=
l_{\sigma'_{n-1}}(P_{n-1})+p_j.
\end{align*}
\proofsection{4}
Схема $\sigma'_n$~--- префиксная, значит,
схема $\sigma'_{n-1}$~--- тоже префиксная.
\proofsection{5}
Схема $\sigma_{n-1}$~--- оптимальна, значит,
$l_{\sigma'_{n-1}}(P_{n-1})\ge l_{\sigma_{n-1}}(P_{n-1})$.
\proofsection{6}
Имеем $l_{\sigma_n}(P_n)=l_{\sigma_{n-1}}(P_{n-1})+p_j\le
l_{\sigma'_{n-1}}(P_{n-1})+p_j=l_{\sigma'_n}(P_n)$. Но схема $\sigma'_n$
оптимальна, значит, схема $\sigma_n$ также оптимальна.
\end{proof}

\end{frame}

\begin{frame}
\label{6.3.5.}
\subsubsection{6.3.5. Алгоритм Хаффмена}
\frametitle{6.3.5. Алгоритм Хаффмена (1/5)}
Рекурсивный  \odmkey{алгоритм Хаффмена}{Алгоритм (algorithm)!Хаффмена (Huffman)}
%\footnote{\;Дэвид А. Хаффмен
%(1925--1999).} 
строит схему оптимального префиксного алфавитного
кодирования для заданного распределения вероятностей появления
букв.

%%\addvspace{-6pt}
%\begin{algorithm}
%\caption{Построение оптимальной схемы~--- рекурсивная процедура
%$\Huffman$} \index{Алгоритм (algorithm)!Хаффмена (David Albert
%Huffman)}
\begin{algorithmic}
\REQUIRE $n>1$~--- количество букв,
\textbf{$P:$ array $[1..n]$ of real}~--- глобальный массив вероятностей букв,
упорядоченный по убыванию.
\ENSURE
\textbf{$C:$ array $[1..n,1..L]$ of $0..1$}~--- глобальный массив элементарных кодов,\\
\textbf{$\ell:$ array $[1..n]$ of $1..L$}~--- глобальный массив длин элементарных кодов 
%схемы оптимального префиксного кодирования.
\IF{$n=2$}
 \STATE $C[1,1]\Assign 0; \ell[1]\Assign 1$ 
 \COMMENT {\small\color{gray} первый элемент }
 \STATE $C[2,1]\Assign 1; \ell[2]\Assign 1$ 
 \COMMENT {\small\color{gray} второй элемент }
\ELSE
\STATE $q\Assign P[n-1]+P[n]$ 
\COMMENT{\small\color{gray} сумма двух последних вероятностей }
\STATE $j\Assign \Up(n,q)$ 
\COMMENT{\small\color{gray} поиск места и вставка суммы }
\STATE $\Huffman(n-1)$ 
\COMMENT{\small\color{gray} рекурсивный вызов }
\STATE $\Down(n,j)$ 
\COMMENT{\small\color{gray} достраивание кодов }
\ENDIF
\end{algorithmic}
%\end{algorithm}

\end{frame}

\begin{frame}
\frametitle{6.3.5. Алгоритм Хаффмена (2/5)}

Функция $\Up$ находит в массиве $P$ место, в
котором должно находиться число~$q$ (см.~пре\-дыдущую теорему), и
вставляет это число, сдвигая вниз остальные элементы.


\begin{algorithmic}
\small
\REQUIRE $n$~--- длина обрабатываемой части массива $P$,
$q$~--- вставляемая сумма.
\ENSURE измененный массив $P$.
\STATE $j\Assign 1$
\COMMENT{\small\color{gray} место вставляемого элемента }
\FOR{\textbf{$i$ from $n-1$ downto $2$}}
  \IF {$P[i-1]< q$}
   \STATE $P[i]\Assign P[i-1]$ 
   \COMMENT {\small\color{gray} сдвиг элемента массива }
  \ELSE
   \STATE $j\Assign i$ 
   \COMMENT {\small\color{gray} определение места вставляемого элемента }
   \STATE \textbf{exit for $i$}
   \COMMENT{\small\color{gray} всё сделано~--- цикл не нужно продолжать }
\ENDIF
\ENDFOR
\STATE $P[j]\Assign q$ 
\COMMENT{\small\color{gray} запись вставляемого элемента }
\STATE\textbf{return $j$}
\end{algorithmic}
\end{frame}

\begin{frame}
\frametitle{6.3.5. Алгоритм Хаффмена (3/5)}
\vspace{-3pt}
{\renewcommand{\baselinestretch}{0.9}\selectfont
Процедура $\Down$ строит оптимальный код для $n$ букв на основе
построенного оптимального кода для $n-1$ буквы. Для этого код
буквы с номером $j$ временно исключается из массива $C$ путем
сдвига вверх кодов букв с номерами, большими $j$, а затем в конец
обрабатываемой части массива $C$ добавляется пара кодов,
полученных из кода буквы с номером $j$ удлинением на 0 и 1
соответственно. Здесь $C[i,*]$ означает вырезку из массива,
то~есть $i$-ю строку массива $C$.
}
%\vspace{6pt}
\begin{algorithmic}

\vspace{-3pt}
\REQUIRE $n$~--- длина обрабатываемой части массива $P$,
$j$~--- номер <<разделяемой>> буквы.
\ENSURE оптимальные коды в первых $n$ элементах массивов $C$ и
$\ell$.
\STATE {$c\Assign C[j,*]$, 
$l\Assign\ell[j]$} \COMMENT{\small\color{gray} запоминание кода буквы $j$ и длины этого кода }
\FOR{\textbf{$i$ from $j$ to $n-2$}}
\STATE $C[i,*]\Assign C[i+1,*]$ 
\COMMENT{\small\color{gray} сдвиг кода }
\STATE $\ell[i]\Assign\ell[i+1]$ 
\COMMENT{\small\color{gray} и его длины }
\ENDFOR
\STATE $C[n-1,*]\Assign c; C[n,*]\Assign c$ \COMMENT{\small\color{gray} копирование кода буквы $j$ }
\STATE $C[n-1,l+1]\Assign 0; C[n,l+1]\Assign 1$ \COMMENT{\small\color{gray} наращивание кодов }
\STATE $\ell[n-1]\Assign l+1; \ell[n]\Assign l+1$ \COMMENT{\small\color{gray} и увеличение длин }
\end{algorithmic}
\end{frame}

\begin{frame}
\frametitle{6.3.5. Алгоритм Хаффмена (4/5)}
\begin{ground}
Для пары букв при любом распределении вероятностей оптимальное
кодирование очевидно: первой букве нужно назначить код 0, а
второй~--- 1. Именно это и делается в первой части оператора
\textbf{if} основной процедуры $\Huffman$. Рекурсивная часть
алгоритма в точности следует доказательству теоремы предыдущего
подраздела. С помощью функции $\Up$ в исходном упорядоченном
массиве $P$ отбрасываются две последние (наименьшие) вероятности,
и~их сумма вставляется в массив $P$ так, чтобы массив (на единицу
меньшей длины) остался упорядоченным. Заметим, что при этом место
вставки сохраняется в локальной переменной~$j$. Так происходит до
тех пор, пока не останется массив из двух элементов, для которого
оптимальный код известен. После этого в~обратном порядке строятся
оптимальные коды для трёх, четырёх и т.~д. элементов. Заметим, что
при этом массив вероятностей $P$ уже не нужен~--- нужна только
последовательность номеров кодов, которые должны быть
изъяты из массива кодов и~продублированы в конце с добавлением
разряда. А эта последовательность хранится в экземплярах локальной
переменной $j$, соответствующих рекурсивным вызовам процедуры
$\Huffman$.
\end{ground}

\end{frame}

\begin{frame}
\frametitle{6.3.5. Алгоритм Хаффмена (5/5)}
\begin{example}
Построение оптимального кода Хаффмена для $n=8$. В левой части
таблицы показано изменение массива $P$, а в правой части~---
массива $C$. Позиция, соответствующая текущему значению
переменной $j$, выделена полужирным начертанием шрифта.

{\small
\begin{center}
\begin{tabular}{@{ }l@{ }l@{ }l@{ }l@{ }l@{ }l@{ }l@{ }|
@{  }l@{ }l@{ }l@{ }l@{ }l@{ }l@{ }l@{ }|l@{ }r}
$p_i$& & & & & & & & & & & & & $C$ & $l_i$ & $p_i l_i$ \\                  
\hline
0,31 & 0,31 & 0,31 & 0,31 & 0,31 & \bf{0,41} & \bf{0,59} & 
\bf{0} & \bf{1} & 1 & 00 & 00 & 00 & 00 & 2 & 0,62 \\
0,24 & 0,24 & 0,24 & 0,24 & \bf{0,28} & 0,31 & 0,41 & 
1 & 00 & \bf{01} & 10 & 10 & 10 & 10 & 2 & 0,48 \\
0,17 & 0,17 & 0,17 & 0,17 & 0,24 & 0,28 &  
&  & 01 & 000 & 11 & 11 & 11 & 11 & 2 & 0,34 \\
0,11 & 0,11 & 0,11 & \bf{0,17} & 0,17 &  &  & 
 &  & 001 & \bf{010} & 011 & 011 & 011 & 3 & 0,33 \\
0,09 & 0,09 & 0,09 & 0,11 &  &  &  & 
 &  &  & 011 & 0100 & 0100 & 0100 & 4 & 0,36 \\
0,05 & 0,05 & \bf{0,08} &  &  &  &  & 
 &  &  &  & \bf{0101} & 01010 & 01010 & 5 & 0,25 \\
0,02 & \bf{0,03} &  &  &  &  &  & 
 &  &  &  &  & \bf{01011} & 010110 & 6 & 0,12 \\
0,01 &  &  &  &  &  &  &  &  &  &  &  &  & 010111 & 6 & 0,06 \\
\hline
 &  &  &  &  &  &  &  &  &  &  &  &  &  &  & \bf{2,56} \\
\hline
\end{tabular}
\end{center}
}

Цена кодирования составляет
$2,56$,
что несколько лучше, чем в кодировании, полученном алгоритмом
Фано.
\end{example}
\end{frame}


\subsection{6.4. Помехоустойчивое кодирование}
\label{6.4.}
\begin{frame}
\frametitle{6.4. Помехоустойчивое кодирование (1/2)}
Надёжность электронных устройств по мере их
совершенствования всё время возрастает, но тем
не менее в их работе возможны
как систематические, так и случайные  ошибки.
Сигнал  в канале связи может быть искажён помехой,
поверхность магнитного носителя может быть повреждена,
в разъёме может быть потерян контакт 
и т.~п.
Ошибки аппаратуры ведут к искажению или
потере передаваемых или хранимых данных.
При определённых условиях,
некоторые из которых рассматриваются в этом разделе,
можно применять методы кодирования, позволяющие
правильно декодировать исходное сообщение, несмотря на
ошибки в данных кода.
В качестве исследуемой модели достаточно рассмотреть
\odmkey{канал связи с помехами}{Канал (channel)!связи
 (communication)!с помехами (noisy)},
потому что к этому случаю
легко сводятся остальные. Например, запись на диск
можно рассматривать как передачу данных в канал,
а чтение с~диска~--- как приём данных из канала.

\begin{tabular}{p{11mm}|p{38mm}|p{95mm}}
  № & Название параграфа
  & Ключевые термины и обозначения \\
  \hline
\DMref{6.4.1.}{6.4.1.} & Кодирование с исправлением ошибок &
{Помехоустойчивое кодирование; характеристика ошибок; обнаружение и исправление ошибок; бит чётности }
\end{tabular}

\end{frame}


\begin{frame}
\frametitle{6.4. Помехоустойчивое кодирование (2/2)}

\begin{tabular}{p{10mm}|p{42mm}|p{93mm}}
  № & Название параграфа
  & Ключевые термины и обозначения \\
  \hline
\DMref{6.4.2.}{6.4.2.} \RS & Возможность исправления всех ошибок &
{Двоичный симметричный канал; пропускная способность канала; достоверность декодирования; сжатие с потерями; теория информации}\\  
\DMref{6.4.3.}{6.4.3.} & Расстояния Левен\-штейна и Хэмминга &
{Метрика; аксиомы метрики: тождества, симметрии, неравенство треугольника; евклидово расстояние; шар; дистанция редактирования; окрестность}\\  
\DMref{6.4.4.}{6.4.4.} & Кодовое расстояние схемы &
{}\\  
\DMref{6.4.5.}{6.4.5.} & Мощность кодирования &
{Безызбыточное кодирование; контрольный разряд}\\ 
\DMref{6.4.6.}{6.4.6.} & Виды кодирования &
{Порождающая матрица кода; проверочная матрица}\\  
\DMref{6.4.7.}{6.4.7.} & Систематическая форма кодовых сообщений &
{Систематическое кодирование; информационный разряд}\\  
\DMref{6.4.8.}{6.4.8.} & Коды Хэмминга &
{Битовая шкала ошибок; синдром}\\  
\DMref{6.4.9.}{6.4.9.} & Исправление одного замещения &
{}\\ 
\end{tabular}  
\end{frame}

\begin{frame}
\label{6.4.1.}\subsubsection{6.4.1. Кодирование с исправлением ошибок}
\frametitle{6.4.1. Кодирование с исправлением ошибок (1/4)}

Пусть имеется канал связи, содержащий источник помех.
Помехи проявляются в том, что принятое сообщение может отличаться от
переданного, то есть имеется отношение 
$C\subset K\times K'$, $K , K' \subset B^*$,
где $K$~---множество переданных, а $K'$~--- соответствующее
множество принятых по каналу сообщений.
Используя запись $s' =C(s)$, $s\in K$, $s'\in K'$ подразумеваем, что разные
<<вызовы>> отношения $C$ с одним и тем же аргументом могут
возвращать различные результаты. Кодирование $\Map{F}{A^*}{B^*}$ (вместе с
декодированием $\Map{F^{-1}}{B^*}{A^*}$), обладающее тем свойством, что
$$
S\To{F} K\To{C} K'\To{F^{-1}} S,\ \A{s\in S\subset A^*}{F^{-1}(C(F(s))) = s},
$$
называется \odmkey{помехоустойчивым}{Кодирование
(coding)!помехоустойчивое (antinoise)}, или
\odmkey{самокорректирующимся}{Кодирование
(coding)!самокорректирующееся (self-correcting)}, или 
\odmkey{кодированием с исправлением ошибок}{Кодирование (coding)!с исправлением
ошибок (error correction)}. Без ограничения общности можно
считать, что $A  = B  = E_2 = \{0,1\}$. Кроме того, естественно
предположить, что содержательное кодирование (вычисление $F$ и
$F^{-1}$) выполняется на устройстве, свободном от помех, то есть
$F$ и $F^{-1}$ являются функциями в обычном смысле.
\begin{remark}
Функция $F^{-1}$ является декодированием для кодирования $F$,
но она \em{не} является обратной функцией для функции $F$
в смысле определений пп.~\DMref{1.4.4}{1.4.4.},~\DMref{1.6.3}{1.6.3.}.
\end{remark}
\end{frame}

\begin{frame}
\frametitle{6.4.1. Кодирование с исправлением ошибок (2/4)}

Вообще говоря, ошибки в канале могут быть следующих
типов:

\begin{itemize}
\item $0 \mapsto 1, \; 1 \mapsto 0$ --- ошибка типа замещения разряда;
\item $0 \mapsto \varepsilon, \;  1 \mapsto \varepsilon$ --- ошибка типа выпадения разряда;
\item $\varepsilon \mapsto 1, \;  \varepsilon \mapsto 0$ --- ошибка типа вставки разряда.
\end{itemize}

Если про источник помех $C$ ничего не известно,
то функции $F$ и $F^{-1}$ не могут быть определены,
за исключением тривиального случая $S=\emptyset$.

\begin{example}
Если канал физически разорвать, то все разряды 
выпадут, и никакое декодирование невозможно. 
\end{example}
 
Таким образом, для решения поставленной задачи
необходимо иметь описание возможных ошибок (проявлений помех).
Вообще говоря, помехи могут иметь различные характеристики
в зависимости от физической природы канала:

\begin{itemize}
\item ошибки могут зависеть от самого передаваемого сообщения;
\item ошибки могут случайно возникать с некоторой вероятностью;
\item ошибки могут определяться физическим устройством канала.
\end{itemize}

\end{frame}


\begin{frame}
\frametitle{6.4.1. Кодирование с исправлением ошибок (3/4)}

Здесь рассматривается простейшая математическая модель канала,
практически не зависящая от его физической природы. Предполагается, что:

\begin{itemize}
\item сообщения передаются 
\odmkey{блоками}{Блок (block)!сообщения (code)} определённой длины $n$,
причём ошибки в различных блоках независимы;
\vspace{-2pt}
\item значения разрядов $0$ и $1$ равноправны: 
если возможна ошибка $0 \mapsto 1$, то возможна и ошибка $1 \mapsto 0$;
\vspace{-2pt}
\item для каждого типа ошибок известна верхняя граница количества ошибок 
в блоке длины $n$.
\vspace{-2pt}
\end{itemize}
%\vspace{-2pt}

Общая \odmkey{характеристика}{Характеристика канала (channel characteristic)}
 ошибок канала (то~есть их количество и типы)
обозначается \\$\delta=\lrle{\delta_1,\delta_2,\delta_3}$, где
$\delta_1,\delta_2,\delta_3$~--- верхние оценки количества ошибок
замещения, выпадения, вставки разряда, соответственно.
%\vspace{-2pt}
\begin{example}
Допустим, что имеется канал с характеристикой
$\delta=\lrle{1,0,0}$, то~есть в канале возможно не более одной
ошибки типа замещения разряда при передаче блока длины~$n$.
Кодирование: $F(a)\Assign
aaa$ (то~есть каждый разряд в сообщении утраивается) и
декодирование $F^{-1}(abc) \Assign {a+b+c >1}$ (то~есть разряд
восстанавливается методом <<голосования>>) кажется помехоустойчивым для данного канала.  
\end{example}
\end{frame}


\begin{frame}
\frametitle{6.4.1. Кодирование с исправлением ошибок (4/4)}
\begin{example} (Продолжение)
Однако метод голосования очень расточителен 
(потребуется передать три блока вместо одного) и не вполне надёжен~--- 
если длина блока $n$ не кратна трём, 
то граница блоков пройдёт по коду одного разряда, 
и если ошибка произойдёт в последнем разряде первого блока 
и в соседнем первом разряде второго блока, то восстановление пройдёт
неверно. 
Так при $n =2$, $F(00)=000000$,
при отправке сообщений блоками $00$, $00$ и $00$, могут произойти ошибки в каждом блоке, так что мы получим 
$01$, $10$ и $01$, очевидно, что декодирование
$F^{-1}(011001)=10$ будет неверным.     
\end{example}

Таким образом, при построении помехоустойчивого
кодирования необходимо найти баланс между
надёжностью и эффективностью.   

\begin{digress} 
Следует различать \em{обнаружение} ошибок  и 
\em{исправление} ошибок.
Например, ранее при передаче данных широко использовался 
\odmkey{бит чётности}{Бит (bit)!чётности (parity)}, то есть 
дополнительный разряд, который добавлялся к каждому блоку так,
чтобы общее количество разрядов, имеющих значение $1$, 
было чётным (или нечётным).
%, то есть
%сумма разрядов блока по модулю два равна $0$ (или $1$).
Очевидно, что этот простой метод 
позволяет обнаружить одиночную ошибку типа замещения, 
но не позволяет её исправить.
В настоящее время бит чётности применяется реже,
поскольку имеет невысокую помехоустойчивость:
если в одном блоке произойдут две ошибки типа замещения, то они не будут обнаружены.  
\end{digress}

%
%
%\begin{example}
%Пусть $\delta=\lrle{1,0,0}$ для слов длины $n$. Метод голосования,
%рассмотренный в примере предыдущего подраздела, является
%кодированием с исправлением одной ошибки типа замещения, но не
%является кодированием с исправлением всех ошибок при $n>1$.
%\end{example}
%

\end{frame}


\begin{frame}
\frametitle{6.4.2. Возможность исправления всех ошибок (1/6)}

Рассмотрим некоторое кодирование 
$\Map{F}{A^*}{B^*}$, $A =B = \{0,1\}$ 
и канал $C$ с характеристикой $\delta$.    
Пусть $E_F^\delta (s)$~--- множество слов, которые могут быть получены
из сообщения $s\in S$ в результате кодирования $F$ и 
всех возможных комбинаций  допустимых в
канале ошибок типа $\delta$.
Возможность исправления ошибок определяется 
свойствами семейства множеств $E_F^\delta (s)$.
\begin{remark}
Иногда множество $E_F^\delta (s)$ 
называют \odmkey{окрестностью ошибок}{Окрестность (neighborhood)!ошибок (of errors)}, и поэтому для обозначения
традиционно используют букву $E$ от
английского Error (ошибка).
\end{remark}

\begin{example}
Пусть $S=\{00,11\}$, 
кодирование $F$ добавляет к сообщению третий разряд (бит чётности), 
а $\delta=\lrle{1,0,0}$ для блоков длины $3$. 
Тогда $F(00) = 000$ и $E^\delta_F (00) = \{ 000, 001,$ \\$ 010, 100 \}$,
$F(11)=110$ и $E^\delta_F (11) =\{ 110, 111, 100, 010 \}$.
Имеем $E^\delta_F (00)\cap E^\delta_F (11) = $\\
$= \{ 100, 010 \}$,
и если из канала получен один из кодов $100$ или $010$,
то однозначно определить исходное сообщение невозможно.

\smallskip
Если же в качестве третьего разряда при кодировании взять конъюнкцию,
то $F(00)=000$ и $E^\delta_F (00)  = \{ 000, 001, 010, 100 \}$,
$F(11) = 111$ и $E^\delta_F (11)  = \{ 111, 110, 101, 011 \}$. 
При этом
$E^\delta_F (00) \cap  E^\delta_F (11)  =  \emptyset$
и исходное сообщение определяется однозначно.
\end{example}

\end{frame}


\begin{frame}
\frametitle{6.4.2. Возможность исправления всех ошибок (2/6)}
\begin{remark}
Множество сообщений $S$ является существенным фактором.
Например, если в условиях предыдущего примера
множество сообщений $S=\{01,10\}$,
то, как нетрудно проверить, никакое добавление 
третьего разряда не является помехоустойчивым кодированием.    
\end{remark}

\medskip
\begin{theorem}
Чтобы существовало помехоустойчивое кодирование $F$ с исправлением
всех ошибок типа $\delta$, необходимо и достаточно, чтобы
$\A{s_1,s_2 \in S}{E^\delta_F (s_1) \cap E^\delta_F (s_2)= \emptyset}$.
\end{theorem}

\begin{proof}
\proofsection{$\Impl$} Если кодирование помехоустойчивое, то
$E^\delta_{F} (s_1)\cap E^\delta_F (s_2)=\emptyset$, иначе $F^{-1}$
невозможно было бы определить как функцию. 
\proofsection{$\Impll$}
Из условия 
$\A{s_1,s_2 \in S}{E^\delta_F (s_1) \cap E^\delta_F (s_2) = \emptyset}$ 
по теореме \DMref{п.~1.7.1}{1.7.1.} следует, что
существует разбиение $\cal B =\{B_s\}_{s\in S}$ множества $\{0,1\}^*$,
причём $\A{s \in S}{E^\delta_F (s) \subset B_s}$. 
По  разбиению $\cal B$  требуемая функция 
$\Map{F^{-1}}{\{0,1\}^*}{S}$ строится следующим
образом: \linebreak
\textbf{if $s'\in B_s$ then
$F^{-1}(s')\Assign s$ end if}.
\end{proof}

\end{frame}


\begin{frame}
\frametitle{6.4.2. Возможность исправления всех ошибок (3/6)}

Таким образом, сама возможность исправления ошибок зависит
от трёх факторов: множества сообщений, применяемого кодирования
и характеристики канала. 
Как правило, множество сообщений и характеристика канала 
считаются заданными, и задача состоит в том, 
чтобы подобрать по возможности эффективное 
помехоустойчивое кодирование. 

\smallskip
Однако в некоторых практических случаях исправление 
\em{всех} ошибок невозможно.

\begin{example}
Рассмотрим канал, в котором в любом передаваемом разряде
может происходить ошибка типа замещения с вероятностью $p$, причём
замещения различных разрядов статистически независимы. Такой канал
называется \odmkey{двоичным симметричным}{Канал (channel)!двоичный симметричный (binary symmetric)}. 
В~этом случае любое слово 
может быть преобразовано в любое другое слово той же длины
замещениями разрядов. Таким образом, 
исправить \em{все} ошибки в двоичном симметричном канале невозможно
даже при сколь угодно малом $p>0$.
\end{example}

\smallskip
Таким образом, в случае использования моделей со случайными
ошибками, в которых количество ошибок неограничено,
построить <<абсолютно>> помехоустойчивое кодирование
невозможно. Возможно лишь попытаться исправить 
ошибки с некоторой вероятностью успеха, но без полной гарантии.
Поэтому далее рассматриваются модели,
в которых количество ошибок ограничено.  

\end{frame}


\begin{frame}
\frametitle{6.4.2. Возможность исправления всех ошибок (4/6)}

\begin{digress}
Для различных моделей со случайными ошибками 
(в том числе для двоичного симметричного канала) 
средствами теории вероятностей  
рассматривается и решаются, в частности, такие задачи, как:

\vspace{-3pt}
\begin{itemize}
\item определение максимального количества информации, которая 
может быть передана через канал при наличии помех 
(\odmkey{пропускная способность канала}{Способность (capacity)!пропускная (сhannel)}, \DMref{п.~6.1.5}{6.1.5.});\\[-2pt] 
\item определение математического 
ожидания вероятности исправления ошибок в канале
(\odmkey{достоверность декодирования}{Достоверность декодирования (decoding fidelity)});\\[-2pt]
\item определение минимального количества информации,
которую необходимо передать через канал 
с сохранением возможности последующего использования данных
(\odmkey{сжатие с потерями}{Сжатие (compression)!с потерями (lossy)}).
\end{itemize}

%\vspace{-3pt}
%Все эти вопросы являются предметом 
%\odmkey{теории информации}{Теория (theory)!информации (of information)}, некоторые вопросы которой 
%рассматривается в разделе \DMref{6.1}{6.1.}.
\end{digress}

\medskip
Рассмотрим возможность исправления всех ошибок с практической точки зрения.   
Пусть в канале возможны только ошибки замещения разрядов
(выпадение и вставка разрядов исключены), сообщения 
передаются блоками длины $n$ и в каждом разряде возможна
ошибка с вероятностью $p$. Это допущение соответствует
устройству шины передачи данных в современном компьютере.

\end{frame}


\begin{frame}
\frametitle{6.4.2. Возможность исправления всех ошибок (5/6)}


При указанных предположениях вероятность того, что произойдёт ровно $k$
ошибок равна $v=p^k(1-p)^{n-k}$.        
Рассмотрим численное значение этой величины
при некоторых типичных значениях величин $n$, $k$ и $p$.

\begin{center}
\begin{tabular}{rrr|r||rrr|r}
$n$ & $k$ & $p$ & $v$ & $n$ & $k$ & $p$ & $v$ \\
\hline
$32$ & $1$ & $10^{-03}$ & $9,695\cdot 10^{-04}$ &
$64$ & $1$ & $10^{-03}$ & $9,389\cdot 10^{-04}$ \\
$32$ & $1$ & $10^{-06}$ & $1,000\cdot 10^{-06}$ &
$64$ & $1$ & $10^{-06}$ & $9,999\cdot 10^{-07}$ \\
$32$ & $2$ & $10^{-03}$ & $9,704\cdot 10^{-07}$ &
$64$ & $2$ & $10^{-03}$ & $9,399\cdot 10^{-07}$ \\
$32$ & $3$ & $10^{-03}$ & $9,714\cdot 10^{-10}$ &
$64$ & $3$ & $10^{-03}$ & $9,408\cdot 10^{-10}$ \\
$32$ & $2$ & $10^{-06}$ & $1,000\cdot 10^{-12}$ &
$64$ & $2$ & $10^{-06}$ & $9,999\cdot 10^{-13}$ \\
$32$ & $3$ & $10^{-06}$ & $1,000\cdot 10^{-18}$ &
$64$ & $3$ & $10^{-06}$ &	$9,999\cdot 10^{-19}$ \\
\hline
\end{tabular}
\end{center}

Таким образом, 
возможностью одновременного возникновения ошибок большой кратности
можно пренебречь, и достаточно рассматривать только каналы 
с ограниченной характеристикой
$\delta<<n$, 
для которых задача исправления всех ошибок оказывается разрешимой.

\end{frame}


\begin{frame}
\frametitle{6.4.2. Возможность исправления всех ошибок (6/6)} 

По определению $s \in S \subset
A^*$, $E_F^\delta (s) \subset B^*$. 
Пусть $E^\delta\lrle{s,s'}$ обозначает
кратчайшую последовательность ошибок типа $\delta$, которая позволяет получить
из слова $s$ слово $s'$.
Заметим, что таких кратчайших последовательностей может быть несколько.
Количество ошибок в кратчайшей последовательности
$E^\delta\lrle{s,s'}$ обозначим 
$\vert E^\delta\lrle{s,s'}\vert$. 

\begin{example}
Пусть $\delta=\lrle{2,1,1}$ для слов длины $2$. Тогда $\vert
E^\delta\lrle{01,10}\vert =2$, причём существует несколько
различных кратчайших последовательностей ошибок, в частности:\\ $01
\stackrel{0\mapsto 1}{\longrightarrow} 11 \stackrel{1\mapsto
0}{\longrightarrow} 10$; $01 \stackrel{0\mapsto
\varepsilon}{\longrightarrow} 1 \stackrel{\varepsilon\mapsto
0}{\longrightarrow} 10$; $01 \stackrel{\varepsilon\mapsto
0}{\longrightarrow} 010 \stackrel{0\mapsto
\varepsilon}{\longrightarrow} 10$.
\end{example}

\begin{remark}
Если характеристика канала
подразумевается, то индекс $\delta$ не указывается.
\end{remark}

\medskip
Окончательно приходим к следующему определению.
 $F$ является
кодированием с исправлением $k$ ошибок типа $\delta$,
если существует декодирование $F^{-1}$ такое, что
$$\A{s\in S}{\A{s'\in E^\delta_F(s)}{|E^\delta\lrle{F(s),s'}|\leq k \Impl F^{-1}(s')=s}}.$$
Кодирование с исправлением всех ошибок существует не всегда,
а только при определённых соотношениях между величинами 
$n$, $k$ и $\delta$.
\end{frame}


\begin{frame}
\label{6.4.3.}\subsubsection{6.4.3.  Расстояния Левенштейна и Хэмминга }\frametitle{6.4.3. Расстояния Левенштейна и Хэмминга (1/4)}

Неотрицательная функция $\Map{d(x,y)}{M \times M}{\Bbb{R}_+}$ называется
\odmkey{расстоянием}{Расстояние (distance)}
(или \odmkey{метрикой}{Метрика (metric)}) на
множестве $M$, если выполнены  три условия, которые называются
 \odmkey{аксиомами метрики}{Аксиома (axiom)!метрики (of metric)}.
\begin{enumerate}
\item \odmkey{Аксиома тождества}{Аксиома (axiom)!тождества (of identity)}: 
$d(x,y) = 0 \Equiv x = y$. 
\item \odmkey{Аксиома симметрии}{Аксиома (axiom)!симметрии (of symmetry)}:
$d(x,y) = d(y,x)$. 
\item \odmkey{Неравенство треугольника}{Неравенство (inequality)!треугольника (triangle)}: $d(x,y) \leq d(x,z) + d(y,z)$.
\end{enumerate}

\begin{example}
Функция $\sqrt{(x_1-x_2)^2+(y_1-y_2)^2+(z_1-z_2)^2}$ 
является метрикой и называется 
\odmkey{евклидовым расстоянием}{Расстояние (distance)!евклидово (Euclidean)} между точками $(x_1,y_1,z_1)$ и
$(x_2,y_2,z_2)$
в трёхмерном пространстве $\Bbb{R}^3$.
\end{example}


Множество $\Set{y}{d(x,y)\leq p}$ называется \odmkey{шаром}{Шар
(ball)}, или \odmkey{метрической окрестностью}{Окрестность
(neighborhood)!метрическая (metric)}, с центром в $x$ и радиусом
$p$.


Ошибки замещения, выпадения и вставки символов можно 
рассматривать как операции с двоичными словами.
Применяя эти операции, возможно одно слово $\alpha\in \{0,1\}^*$ 
преобразовать в другое слово $\beta\in \{0,1\}^*$.


\begin{example} 
$01\stackrel{0\mapsto 1}{\longrightarrow} 11 
\stackrel{1\mapsto 0}{\longrightarrow} 10$; $01 \stackrel{0\mapsto
\varepsilon}{\longrightarrow} 1 \stackrel{\varepsilon\mapsto
0}{\longrightarrow} 10$; $01 \stackrel{\varepsilon\mapsto
0}{\longrightarrow} 010 \stackrel{0\mapsto
\varepsilon}{\longrightarrow} 10$.
\end{example}
 
\end{frame}

\begin{frame}
\frametitle{6.4.3.  Расстояния Левенштейна и Хэмминга  (2/4)}
Минимальное число операций, переводящее слово $\alpha$ в слово $\beta$
называется \odmkey{расстоянием Левенштейна}{Расстояние (distance)!Левенштейна (Levenshtein)}
 и обозначается $L(\alpha,\beta)$.  
\begin{example} $L(01,10)=2$.
\end{example}
\begin{lemma}
Расстояние Левенштейна $L$ является метрикой.
\end{lemma}
\begin{proof}
\proofsection{1} $L(\alpha,\beta) = 0 \Equiv \alpha = \beta$,
поскольку преобразование отсутствует
тогда и только тогда, когда слова совпадают.
\proofsection{2} $L(\alpha,\beta) = L(\beta,\alpha)$,
поскольку ошибки симметричны и, значит,  преобразования слов обратимы.
\proofsection{3} $L(\alpha,\beta) \leq L(\alpha,\gamma) +
L(\beta,\gamma)$, поскольку
конкатенация последовательностей преобразования
 $\alpha\mapsto\gamma$ и
$\gamma\mapsto\beta$
является \em{некоторой} последовательностью, преобразующей
$\alpha$ в $\beta$, а $L(\alpha,\beta)$
является длиной кратчайшей из таких последовательностей.
\end{proof}

\begin{remark}
Расстояние Левенштейна нетрудно определить 
для слов в произвольном алфавите,
а не только в алфавите $\{0,1\}$.
В таком случае это расстояние часто называют
\odmkey{дистанцией редактирования}{Дистанция (distance)!редактирования (edit)}.
\end{remark}


\end{frame}

\begin{frame}
\frametitle{6.4.3.  Расстояния Левенштейна и Хэмминга  (3/4)}
\begin{digress}
На операции редактирования могут быть наложены различные ограничения.
При этом, если преобразования слов остаются обратимыми, 
а конкатенация преобразований является преобразованием,
то расстояние Левенштейна остаётся метрикой.
Например, в популярной интеллектуальной игре
требуется преобразовать одно слово в другое,
заменяя по одной букве. В этой игре $3\leq L(\text{коза},\text{волк})\leq 5$,
поскольку \\коза$\mapsto$поза$\mapsto$пола$\mapsto$полк$\mapsto$волк.
\end{digress}

\medskip
Пусть некоторый канал имеет характеристику 
$\delta=\lrle{\delta_1, \delta_2, \delta_3}$.
Пусть $E^\delta (\alpha)$~--- множество слов, которые могут быть получены
из слова $\alpha$ в результате всех возможных комбинаций  допустимых в
канале ошибок типа $\delta$. 
Поскольку рассматриваются симметричные ошибки,
ясно что $\beta\in E^\delta (\alpha)\Equiv \alpha \in E^\delta (\beta)$. 
\end{frame}

\begin{frame}
\frametitle{6.4.3.  Расстояния Левенштейна и Хэмминга  (4/4)}
Положим
\textbf {$H^\delta(\alpha, \beta) \Assign$
if $\alpha\in E^\delta (\beta)$ then
$L(\alpha, \beta)$ else $+\infty$ end if}.

\medskip
Функция $H^\delta$ называется
\odmkey{расстоянием Хэмминга}{Расстояние
 (distance)!Хэмминга (Hamming)}.
%\footnote{\;Ричард Весли Хэмминг (1915--1998).}.
Расстояние Хэмминга является метрикой
по тем же причинам, что и расстояние Левенштейна,
а потому множество $E^\delta(\alpha)$ 
уместно называть \odmkey{$\delta$-окрестностью}{Дельта-окрестность (delta neighborhood)}
слова $\alpha$. 


\begin{example}
Расстояние Хэмминга в $E_2^n$.  Пусть $\delta=\lrle{n,0,0}$, то
есть допускаются только ошибки типа замещения разрядов. Рассмотрим
слова $\alpha = a_1\dots a_n$ и $\beta =b_1\dots b_n$.
Тогда $$H^\delta(\alpha,\beta)\Def\sum\limits^n_{i=1}(a_i\ne b_i).$$
То есть расстоянием Хэмминга между битовыми шкалами одной длины
является число несовпадающих разрядов. Именно этот 
частный случай часто называют расстоянием Хэмминга в $E_2^n$ и обозначают
 $H^n$. 
\end{example}
\begin{remark}
В последней формуле, так же как и в других местах этой книги,
используется неявное приведение \textbf{false $\mapsto 0,$ true
$\mapsto 1$}.
\end{remark}
\end{frame}

\begin{frame}
\label{6.4.4.}\subsubsection{6.4.4.  Кодовое расстояние схемы }\frametitle{6.4.4. Кодовое расстояние схемы (1/3)}
Рассмотрим применение введённых понятий к простейшему случаю алфавитного кодирования.
Пусть $\sigma = {\lrle{a_i\to\beta_i}}_{i=1}^n$~--- схема
алфавитного кодирования, а $d$~--- некоторая метрика на $B^*$.
Тогда минимальное расстояние между элементарными кодами
$$d(\sigma)\Def \min\limits_{1\leq i<j \leq n} d(\beta_i,\beta_j) $$
называется \odmkey{кодовым расстоянием}{Расстояние (distance)!кодовое
(code)} схемы $\sigma$ в метрике $d$.

\begin{theorem} Алфавитное кодирование со схемой
$\sigma = {\lrle{a_i\to\beta_i}}_{i=1}^n$ и с кодовым расстоянием Хэмминга
$H(\sigma) = \min\limits_{\beta',\beta''\in V}{H^\delta(\beta',\beta'')}$
является кодированием с исправлением $h$ ошибок типа $\delta$ тогда и
только тогда, когда $H(\sigma) > 2h$.
\end{theorem}
\begin{proof}
Если можно исправить ошибки в каждом из кодовых слов, то можно
исправить ошибки и во всем сообщении. Поэтому достаточно
рассматривать только кодовые слова. Заметим, что множество
$E^\delta(\alpha)$ --- это шар радиуса $h$ с центром в $\alpha$. 
Таким образом, $E^\delta (\beta_i)\cap E^\delta (\beta_j)= \emptyset\Equiv
H(\sigma)>2h$. 

По теореме \DMref{п.~6.4.2}{6.4.2.} условие
$E^\delta (\beta_i)\cap E^\delta (\beta_j)=\emptyset$ равносильно существованию декодирования.
\end{proof}
\end{frame}

\begin{frame}
\frametitle{6.4.4. Кодовое расстояние схемы (2/3)}

Другими словами, исправление ошибок типа $\delta$ возможно тогда и только тогда,
когда кодовые слова схемы не лежат в $\delta$-окрестностях друг друга.
\begin{corollary}
Рассмотрим схему равномерного кодирования
$\sigma=\lrle{a_i\to\beta_i}_{i=1}^{n}$, \\где
$\A{i}{|\beta_i|=m}$, $m\ge \Lrf{\log(n-1)}+1$.
Такая схема может исправлять $h$ ошибок
типа замещения в $m$ разрядах, если, и только если
$\min\limits_{1\le i < j \le n}\left(\sum\limits_{k=1}^{m} \beta_i[k]\ne \beta_j[k]\right)
>2h$,
то есть если, и только если любые два кодовых слова различаются не менее чем в $2h$ разрядах.
\end{corollary}

\begin{proof}
Кодовое расстояние схемы равномерного кодирования
является расстоянием Хэмминга.
\end{proof}

\begin{remark}
Отсюда получается, что для блоков малой длины
$n=1, n=2$ исправление ошибок вообще невозможно,
и содержательное кодирование с исправлением
ошибок возможно только для блоков длины $n=3$ и более.
\end{remark}
\end{frame}

\begin{frame}
\frametitle{6.4.4. Кодовое расстояние схемы (3/3)}
\begin{columns}[t] % contents are top vertically aligned
\begin{column}[T]{9cm} % each column can also be its own environment
\begin{example}
На рисунке представлен
гиперкуб для пространства $E_2^3$ 
и две $1$-окрестности точек $101$ и $010$ 
(то есть два шара радиуса $1$ с центрами в точках
$101$ и $010$). Очевидно,
что можно взять два любых инверсных кода, 
и шары с центрами в этих кодах не пересекаются,
но все другие шары радиуса $1$ с центрами в других точках 
пересекаются с обоими выделенными шарами.
Из этого наблюдения и теоремы следует,
что трёхразрядный код 
может исправлять одну ошибку типа замещения 
для алфавитного кодирования только двухбуквенного алфавита.
\end{example}

\end{column}
\begin{column}[T]{6cm} % alternative top-align that's better for graphics
\begin{center}
\includegraphics[height=50mm]{6_3_3.jpg}
\end{center}
\end{column}
\end{columns}

\end{frame}

\begin{frame}
\label{6.4.5.}\subsubsection{6.4.5.  Мощность кодирования }\frametitle{6.4.5. Мощность кодирования (1/4)}
\odmkey{Мощностью}{Мощность (power)!кодирования (of encoding)} кодирования называется 
мощность множества сообщений $S$, 
которые удаётся закодировать с выполнением
заданных условий. Чем мощнее кодирование,
тем больше различных сообщений 
возможно закодировать с его помощью. 

Из теоремы \DMref{п.~6.4.2}{6.4.2.} следует,
что помехоустойчивое кодирование максимальной мощности
достигается в том случае,
когда разбиение множества принятых сообщений 
имеет максимальную мощность.  

\begin{example}
Помехоустойчивое кодирование
с исправлением не более одной ошибки типа замещения
в блоке длины $3$ имеет мощность не более $2$.   
\end{example}

Рассмотрим случай канала, в котором возможно не более $h$ 
ошибок замещения разрядов в блоках
длины $n$. Какова максимальная возможная мощность
помехоустойчивого кодирования? 
В рассматриваемом случае 
в метрическом пространстве $E_2^n$ 
с расстоянием Хэмминга $H^n$ 
множество возможных ошибок в сообщении 
$s$ определяется как шар радиуса $h$ 
с центром в коде сообщения $F(s)$:

$$E_F^{\lrle{h,0,0}} (s) =V^n(F(s),h), \quad 
\text{где} \  V^n(x,h) \Def \Set{y\in E_2^n}{H^n(x,y)\leq h}.$$    
      

\end{frame}

\begin{frame}
\frametitle{6.4.5. Мощность кодирования (2/4)}

\begin{lemma}
$\A{x\in E_2^n}{\A{h\geq 0}{\left\vert V^n(x,h)\right\vert =
\sum\limits_{i=0}^h C(n,i)}}$.
\end{lemma}
\begin{proof}
$ V^n(x,h) = \Set{y\in E_2^n}{H^n(x,y)\leq h} 
=\bigcup\limits_{i=0}^h \Set{y\in E_2^n}{H^n(x,y)= i}
$,\\ 
$ \vert \Set{y\in E_2^n}{H^n(x,y)= i}\vert = C(n,i)
$. 
\end{proof}

\begin{remark}
Доказанную лемму можно обобщить для любого дискретного метрического пространства,
в частности, пространства слов над любым алфавитом 
с расстоянием Левенштейна.
\end{remark}

\begin{theorem}
Мощность помехустойчивого кодирования для канала, 
в котором возможно не более $h$ 
ошибок замещения разрядов в блоках
длины $n$, не превосходит числа шаров,
которые можно уложить в пространстве $E^n_2$ 
без пересечений:
$
\displaystyle{|S| \leq \frac{2^n}{\sum\limits_{i=0}^h C(n,i)}.}
$ 
\end{theorem}
\begin{proof}
Для каждого сообщения множество возможных ошибок является шаром
радиуса $h$, для помехустойчивости шары не должны пересекаться, 
а всего в пространстве $2^n$ точек.
\end{proof}

\end{frame}


\begin{frame}
\frametitle{6.4.5. Мощность кодирования (3/4)}
\vspace{-3pt}
\begin{example}
В следующей таблице показаны верхние оценки для мощности
помехоустойчивовго кодирования
в зависимости от $h$ и $n$.

\vspace{-4pt}
\begin{center}
\begin{tabular}{l|rrrrrrrrrrr}
${}_h\setminus{}^n$ & 2 & 3 & 4 & 5 & 6 & 7 & 8 & 9 & 10 & 11 & 12 \\
\hline
0 & 4 & 8 & 16 & 32 & 64 & 128 & 256 & 512 & 1024 & 2048 & 4096  \\
1 & 1 & 2 & 3 & 5 & 9 & 16 & 28 & 51 & 93 & 170 & 315  \\
2 & 1 & 1 & 1 & 2 & 2 & 4 & 6 & 11 & 18 & 30 & 51  \\
3 & 1 & 1 & 1 & 2 & 2 & 2 & 3 & 5 & 7 & 11 & 17  \\
\end{tabular}
\end{center}
\end{example}

\vspace{-4pt}
\begin{remark}
Случай $h=0$ соответствует максимальной
мощности кодирования при отсутствии ограничений.
Такое кодирование называется \odmkey{безызбыточным}{Кодирование (coding)!безызбыточное (breakeven)}.
\end{remark}

Специфические особенности множества сообщений $S$ 
можно использовать при построении кодирования,
однако это требует специальных ухищрений и не 
является универсальным решением. 
Для практики 
наиболее важен случай $S=E_2^m$, $|S|=2^m$, то есть 
сообщением может быть любое слово длины $m$. 

\begin{corollary}
Если $S=E_2^m$ и существует помехоустойчивое кодирование
для канала, 
в котором возможно не более одной ошибки 
замещения  в блоке длины $n$, то $2^m\leq 2^n / (n+1)$.
\end{corollary}
\begin{proof}
По формуле теоремы $C(n,0)+C(n,1)=1+n$.
\end{proof}
\end{frame}


\begin{frame}
\frametitle{6.4.5. Мощность кодирования (4/4)}
Помехоустойчивое кодирование не может быть безызбыточным.
Обозначим $k\Assign n -m$~--- минимальное 
количество \odmkey{контрольных}{Разряд (digit)!контрольный (control)} разрядов,
то есть избыточных разрядов, необходимых для помехоустойчивости. 

\begin{example}
В следующей таблице показано соотношение для $m$, $n$ и $k$
при $h=1$. Здесь $n'\Assign \lfloor 2^n/(n+1)\rfloor$. 
Для $n=1,2$ помехоустойчивого кодирования не существует. 


\begin{tabular}{l|r|rrrr|rrrrrrrr|r}
$n$ & 3 & 4 & 5 & 6 & 7 & 8 & 9 & 10 & 11 & 12 & 13 & 14 & 15 & 16 \\
$n'$ & 2 & 3 & 5 & 9 & 16 & 28 & 51 & 93 & 170 & 315 & 585 & 1092 & 2048 & 3855 \\
$2^m$ & 2 & 2 & 4 & 8 & 16 & 16 & 32 & 64 & 128 & 256 & 512 & 1024 & 2048 & 2048 \\
$m$ & 1 & 1 & 2 & 3 & 4 & 4 & 5 & 6 & 7 & 8 & 9 & 10 & 11 & 11 \\
$k$ & 2 & 3 & 3 & 3 & 3 & 4 & 4 & 4 & 4 & 4 & 4 & 4 & 4 & 5 \\
\end{tabular}
\end{example}

\vspace{-2pt}
Можно заметить, например, что $7$-разрядные блоки передают 
столько же информации и с такой же надёжностью, что и $8$-разрядные,
аналогичная ситуация для $15$-разрядных и $16$-разрядных кодов и т.~д. 
Доля контрольных разрядов, 
необходимых для существования помехоустойчивого кодирования,
начиная с $n=6$ не превышает половины всех разрядов блока и с ростом $n$ 
постепенно уменьшается. В частности, при $n\in 27..57$ имеем 
$k=6$, а при $n=64$ имеем $k=7$. 

\end{frame}


\begin{frame}
\label{6.4.6.}\subsubsection{6.4.6.  Виды кодирования  }\frametitle{6.4.6. Виды кодирования  (1/5)}
Предложено множество разнообразных способов 
помехоустойчивого кодирования. Все эти способы 
подчиняются выведенным количественным ограничениям,
но отличаются используемым математическим аппаратом,
удобством применения, наглядностью и другими качественными
характеристиками.

\begin{example}   
Пусть $m=2, h=1$. Тогда минимально $n=5, k=3$. 
Можно \em{подобрать} подходящие четыре пятиразрядных кода так,
чтобы их $1$-окрестности в метрике Хэмминга не пересекались. В частности,
с помощью карты Карно (\DMref{п.~3.5.5}{3.5.5.}), как показано на рисунке.
Выбранные коды (центры шаров) выделены полужирным начертанием и подчёркнуты,
сами шары выделены различной штриховкой. 


\begin{center}
\includegraphics[height=30mm]{6_3_6_1.jpg}
\end{center}

\end{example}  
\end{frame}

\begin{frame}
\frametitle{6.4.6. Виды кодирования  (2/5)}
\vspace{-0.25\baselineskip}
\begin{example} (Продолжение) Затем можно произвольно назначить
коды сообщениям:\\
$00\mapsto 00001$, $01\mapsto 01010$, $11\mapsto 11111$, $10\mapsto 10100$. 
Декодирование можно реализовать с помощью дерева решений (\DMref{п.~3.5.8}{3.5.8.}), 
как показано на рисунке. 

\vspace{-0.5\baselineskip}
\begin{center}
\includegraphics[height=35mm]{6_3_6_2.jpg}
\end{center}

\vspace{-0.5\baselineskip}
Варианты с пометкой <<err>> означают появление ошибки,
при которой произошло более одного замещения разряда. 
Такую ошибку данное кодирование исправить не может, но может обнаружить.   
\end{example}

\begin{remark}
Выбор системы непересекающихся шаров в $E_2^n$,
назначения кодов словам из $E^m$ является произвольным.
Если сохранить этот выбор в тайне,
то кодирование, наряду с помехоустойчивостью,
станет обладать и свойствами тайнописи.    
\end{remark}

\end{frame}

\begin{frame}
\frametitle{6.4.6. Виды кодирования  (3/5)}

%\vspace{-0.25\baselineskip}
Среди пестроты методов помехоустойчивого кодирования 
своей простотой и ясностью выделяется группа методов,
основанных на применении математического аппарата и обозначений 
линейной алгебры. 

\smallskip
Напомним, что двоичная арифметика
 $\Alg{E_2}{+_2, \&}$ является полем (\DMref{п.~2.3.3}{2.3.3.})
относительно сложения по модулю 2 и конъюнкции, 
множество битовых шкал $E_2^n$ образует 
абелеву группу (\DMref{п.~2.2.6}{2.2.6}) относительно 
операции поразрядного логического сложения шкал,
и, как нетрудно проверить,
множество битовых шкал $E_2^n$ 
образует векторное пространство (\DMref{п.~2.5.1}{2.5.1.})
над двоичной арифметикой.

\begin{remark}
В контексте этого раздела знак + означает
сложение по модулю 2, а знак умножения опускается,
что позволяет использовать обычные обозначения
линейной алгебры, имея в виду, что все векторы и матрицы~--- булевы,
и операции выполняются по правилам, изложенным в \DMref{п.~1.4.9}{1.4.9.}.
Размеры матриц и векторов, если нужно, указываются в нижнем индексе. 
\end{remark} 

\smallskip
Введём обозначения: $\Bl{0}_m$ означает вектор из $m$ нулей, 
$\Bl{0}_{m\times n}$ означает матрицу размера $m\times n$ из нулей,
а $I_m$ означает единичную матрицу размера $m\times m$.
\end{frame}

\begin{frame}
\frametitle{6.4.6. Виды кодирования  (4/5)}
%\vspace{-0.25\baselineskip}
Рассмотрим
простейший случай кодирования, когда исходными 
сообщениями являются все слова длины $m$, 
а кодами являются блоки длины $n$, $n\ge m$.   

Пусть $\alpha=a_1\dots a_m \in E_2^m$~---  
исходное сообщение, а $\beta =b_1\dots b_n \in E_2^n$~--- кодовый блок. 
Рассмотрим случай, 
когда каждый разряд $\beta$ можно представить в виде \em{линейной} функции от
разрядов $\alpha$ (см.~\DMref{п.~3.6.1}{3.6.2.}). 
Тогда кодирование записывается следующим образом:
$\A{j \in 1..n }{b_j= \sum_{i=1}^{m} a_i  g_{i,j} }$.
Используя обозначения линейной алгебры, код $\beta$ можно представить 
в виде произведения вектора $\alpha$ и булевой
матрицы $G_{m \times n}$, составленной из коэффициентов $g_{i,j}$,
или $\beta={\alpha}{G}$.
Матрица $G$ называется \odmkey{порождающей матрицей}{Матрица (matrix)!порождающая (generator)} кода. 
 
\begin{example}
Рассмотрим сообщение $\alpha = \begin{pmatrix} 0 & 1 & 1 \end{pmatrix}$. 
Пусть матрица $G =\begin{pmatrix}
1 & 0 & 0 & 1 & 0\\
0 & 1 & 0 & 0 & 1 \\
0 & 0 & 1 & 1 & 1
\end{pmatrix}$. Тогда 
$\beta =
\begin{pmatrix}
0 & 1 & 1
\end{pmatrix}
\begin{pmatrix}
1 & 0 & 0 & 1 & 0\\
0 & 1 & 0 & 0 & 1 \\
0 & 0 & 1 & 1 & 1
\end{pmatrix}
 = 
 \begin{pmatrix}
0 & 1 & 1 & 1 & 0
\end{pmatrix}
$. Здесь $m=3, k=2, n=5$. 
\end{example}

\end{frame}

\begin{frame}
\frametitle{6.4.6. Виды кодирования  (5/5)}
Если $\E{G^{-1}_{n\times m}}{GG^{-1}=I_m}$,  
то 
 $\beta G^{-1} = (\alpha G) G^{-1} = \alpha (GG^{-1}) = \alpha I_m = \alpha$,
и декодирование выполняется умножением кода на матрицу $G^{-1}$.

\medskip
\begin{theorem}
Если $\E{H_{k\times n}}{GH^{\mathrm{T}}=\Bl{0}_{m\times k} }$, 
то $\beta=\alpha G \Equiv \beta H^{\mathrm{T}}=\Bl{0}_k$.
\end{theorem}
\begin{proof}
$\beta=\alpha G \Equiv \beta H^{\mathrm{T}}= \alpha G H^{\mathrm{T}} =
\alpha \Bl{0}_{m\times k} = \Bl{0}_k$.
\end{proof}

Матрица $H$ называется \odmkey{проверочной}{Матрица (matrix)!проверочная (parity check)},
поскольку с помощью равенства $\beta H^{\mathrm{T}}=\Bl{0}_k$
позволяет проверить, является ли принятый 
через канал код $\beta$ действительно кодом некоторого сообщения 
$\alpha$ 
при кодировании $G$. 
%{\small
\begin{example}
Матрица $H =\begin{pmatrix}
0 & 1 & 1 & 1 & 0\\
1 & 0 & 1 & 0 & 1 \\
\end{pmatrix}$ является проверочной для матрицы
$G$ предыдущего примера:
$
GH^\mathrm{T} =\begin{pmatrix}
1 & 0 & 0 & 1 & 0\\
0 & 1 & 0 & 0 & 1 \\
0 & 0 & 1 & 1 & 1
\end{pmatrix}
\begin{pmatrix}
1 & 0\\
0 & 1\\
1 & 1\\
1 & 0\\
0 & 1\\
\end{pmatrix} =
\begin{pmatrix}
0 & 0\\
0 & 0\\
0 & 0\\
\end{pmatrix}
$. 
\end{example}
%}

\end{frame}

\begin{frame}
\label{6.4.7.}\subsubsection{6.4.7.  Систематическая форма кодовых сообщений }\frametitle{6.4.7. Систематическая форма кодовых сообщений (1/3)}
\vspace{-4pt}
Известен очень простой способ определения матриц $G$, $G^{-1}$ и $H$,
который    
из-за своей формы называется
\odmkey{систематическим кодированием}{Кодирование (encoding)!систематическое (systematic)}. 
При систематическом кодировании блок из
$n$ разрядов состоит из $m$ \odmkey{информационных}{Разряд (digit)!информационный (information)} разрядов, 
за которыми следуют
 $k=n-m$ \odmkey{проверочных}{Разряд (digit)!проверочный (check)} (или \odmkey{контрольных}{Разряд (digit)!контрольный (control)}) разрядов,
причём информационные разряды копируются в код без изменения,
а проверочные разряды вычисляются как линейные функции 
от информационных разрядов.  
В этих соглашениях 
порождающая матрица имеет вид $G=[I_m|P_{m\times k}]$, где $I_m$~--- 
единичная матрица, формирующая
информационные разряды, 
а $P_{m\times k}$~--- матрица, определяющая проверочную часть.
  
\begin{digress}
В терминах линейной алгебры, если строки 
порождающей матрицы $G$ 
линейно независимы, то коды различных сообщений 
 образуют линейное подпространство
векторного пространства $E_2^n$.
В этом случае проверочная матрица $H$ задаёт
ортогональное дополнение векторного
пространства, заданного матрицей $G$.
Все положения распространяются на случай
векторных пространств над любым конечным полем,
а не только $E_2$.
Использование понятий линейной алгебры 
(в том числе понятий, выходящих за рамки изложенных в главе 2)
позволяет сделать описание кодирования 
лаконичным и общим,
но не меняет алгоритмической сути дела.%
\end{digress}

\end{frame}

\begin{frame}
\frametitle{6.4.7. Систематическая форма кодовых сообщений (2/3)}
\begin{theorem}
Если $G=[I_m|P_{m\times k}]$, то: 
\begin{enumerate}
\item[1)]
$G^{-1}_{n\times m} \Assign \left[I_m |\Bl{0}_{m\times k} \right]^\mathrm{T}$ и 
при этом $GG^{-1}=I_m$;
\item[2)] 
$H_{k\times n} \Assign \left[ P_{k\times m}^{\mathrm{T}} | {I_k} \right]$
и при этом $GH^{\mathrm{T}}=\Bl{0}_{m\times k}$.
\end{enumerate}
\end{theorem}

\begin{proof}
\proofsection{1}
Положим $C_{m\times m} \Assign GG^{-1}$. 
Имеем \textbf{$G_{i,l}=$ if $i=l$ then $1$ else $0$ end if} или, короче,
$G_{i,l}= (i=l)$ при $1\leq i,l \leq m$, и также $G_{l,j}^{-1}= (l=j)$.  
Тогда $C_{i,j} = \sum\limits_{l=1}^{n} G_{i,l} G^{-1}_{l,j} = (i=j)$,
то есть $C=I_m$.
\proofsection{2} 
Положим $C_{m\times k} \Assign GH^{\mathrm{T}}$. 
Имеем $G_{i,l}= (i=l)$ при $1\leq i,l\leq m$, 
$H^{\mathrm{T}}_{l,j}= (l-m=j)$ при $m+1 \leq l \leq m+k, 1\leq j\leq k$, 
а также $G_{i,l} = P_{i,l-m}$ при $l>m$ и 
$H_{l,j}^{\mathrm{T}} = P_{l,j}$ при $1\leq l \leq m$.  
Тогда $C_{i,j} = \sum\limits_{l=1}^{n} G_{i,l} H_{l,j}^{\mathrm{T}} =
\sum\limits_{l=1}^{m} G_{i,l} H_{l,j}^{\mathrm{T}}+
\sum\limits_{l=m+1}^{n} G_{i,l} H_{l,j}^{\mathrm{T}} = P_{i,j}+P_{i,j} = 0$.
\end{proof}

\begin{remark}
Напомним, что матрицы булевы,
и в операциях с матрицами знак $+$
означает сложение по модулю 2. 
\end{remark}

\end{frame}

\begin{frame}
\frametitle{6.4.7. Систематическая форма кодовых сообщений (3/3)}

%\begin{remark}
%Систематический код с параметрами $n$ и $m$ обозначается как $(n, m)$–код.
%\end{remark}

Систематическая форма кодовых сообщений позволяет удобно 
кодировать и декодировать сообщения, и позволяет обнаруживать 
некоторые ошибки, но не гарантирует исправления ошибок.

\vspace{-12pt}
\begin{example} В условиях предыдущих примеров
$G=\begin{pmatrix}
1 & 0 & 0 & 1 & 0\\
0 & 1 & 0 & 0 & 1\\
0 & 0 & 1 & 1 & 1\\
\end{pmatrix},
H =\begin{pmatrix}
0 & 1 & 1 & 1 & 0\\
1 & 0 & 1 & 0 & 1\\
\end{pmatrix}$.
Пусть $\alpha= 
\begin{pmatrix}
0 & 1 & 1
\end{pmatrix}$. Тогда 
$\beta=\alpha G = \begin{pmatrix}
0 & 1 & 1 & 1 & 0\\
\end{pmatrix}$.
Пусть теперь произошла ошибка 
в третьем разряде и из
канала принято слово
$\beta'=\begin{pmatrix}
0 & 1 & 0 & 1 & 0\\
\end{pmatrix}$.
Тогда 
$\beta'G^{-1}=\begin{pmatrix}
0 & 1 & 0
\end{pmatrix}
$ и сообщение декодировано
неверно, но 
$\beta'H^{\mathrm{T}} =
\begin{pmatrix}
0 & 1 & 0 & 1 & 0 
\end{pmatrix}
\begin{pmatrix}
1 & 0\\
0 & 1\\
1 & 1\\
1 & 0\\
0 & 1
\end{pmatrix} =
\begin{pmatrix}
1 & 1  
\end{pmatrix}
$ 
и ошибка обнаружена.

\end{example} 

\begin{remark}
По теореме \DMref{п.~6.4.5}{6.4.5.} кодирование с параметрами $m=3$, $k=2$, $n=5$,
%рассматриваемое в примере, 
не может быть помехустойчивым.
Если $m=3$, то необходимо $k\geq 3$.  
\end{remark}
\end{frame}

\begin{frame}
\label{6.4.8.}\subsubsection{6.4.8.  Коды Хэмминга }\frametitle{6.4.8. Коды Хэмминга (1/4)}
При помехоустойчивом кодировании основная задача состоит в том, 
чтобы определить, какую дополнительную
информацию надо передать, чтобы её объём был минимальным, а декодирование по возможности несложным.


Видимо, наиболее известным 
систематическим самокорректирующимся кодом 
в настоящее время является \odmkey{код Хэмминга}{Код (code)!Хэмминга (Hamming)}.

%  Далее идёт детская ботва, которую я выкинул ФАН 
%\newline
%Считаем, что канал допускает только ошибки типа замещения. Допустим, для кода со схемой кодирования 
%$\sigma = {\lrle{a_i\to\beta_i}}_{i=1}^n$, кодовое расстояние Хэмминга $d(\sigma) = 3$ и проверочная 
%матрица содержит $k$ строк. Всего существует $2^k - 1$ различных ненулевых столбцов длины $k$.
%Значит, для данного кода $n = 2^k - 1$, а $m = 2^k - 1 - k$. Получаем $(2^k - 1, 2^k - 1 - k)$ - код.
%Такие коды называются \odmkey{кодами Хэмминга}{Код (code)!Хэмминга (Hamming)} и по теореме из п.6.4.3 исправляют одиночные ошибки типа замещения.
%\begin{remark}
%На самом деле код Хэмминга способен исправлять одиночные или обнаруживать две ошибки типа замещения.
%\end{remark}
%\end{frame}
%
%\begin{frame}
%\frametitle{6.4.5. Код Хэмминга для исправления одного замещения (2/7)}
%При заданном $n$ количество дополнительных разрядов $k$ подбирается таким образом, чтобы $2^k\ge n+1$.
%Имеем
%$$
%2^k\ge n+1\Impl \frac{2^n}{n+1}\ge 2^{n-k}\Impl
%\frac{2^n}{n+1}\ge 2^{m}.
%$$
%\begin{example}
%Для сообщения длины  $m=4$ достаточно $k=3$ дополнительных разрядов, поскольку $8= 2^3\ge4+3+1$.
%$$
%G =
%\begin{pmatrix}
%1 & 0 & 0 & 0 & 1 & 1 & 0\\
%0 & 1 & 0 & 0 & 1 & 0 & 1\\
%0 & 0 & 1 & 0 & 0 & 1 & 1\\
%0 & 0 & 0 & 1 & 1 & 1 & 1
%\end{pmatrix}
%, H =
%\begin{pmatrix}
%1 & 1 & 0 & 1 & 1 & 0 & 0\\
%1 & 0 & 1 & 1 & 0 & 1 & 0\\
%0 & 1 & 1 & 1 & 0 & 0 & 1
%\end{pmatrix}
%$$
%\end{example}
%\end{frame}

%%\begin{frame}
%%\frametitle{6.4.5. Код Хэмминга (3/7)}
Введём обозначения.
Пусть $\alpha\in E_2^m$~--- исходный блок,
$\beta \in E_2^n$~--- кодовый блок, 
полученный с помощью порождающей матрицы $G_{m\times n}$,
для которой известна проверочная матрица $H_{k\times n}$,
такая, что $GH^{\mathrm{T}} = \Bl{0}_{m\times k}$.
Пусть   
$\beta' \in E_2^n$~--- блок с возможной ошибкой,
полученный через канал, 
а $\epsilon \in E_2^n$~--- \odmkey{битовая шкала ошибок}{Битовая шкала (bit scale)!ошибок (of error)}, 
которая содержит $1$
там, где произошла ошибка, и $0$ в остальных разрядах.
Результат применения проверочной матрицы к блоку с ошибкой 
называется \odmkey{синдромом}{Синдром (syndrome)} и обозначается 
$\zeta \Assign \beta' H^{\mathrm{T}}$.
\begin{theorem}
Синдром зависит только от битовой шкалы ошибок: $\zeta = \epsilon H^{\mathrm{T}}$.
\end{theorem}
\begin{proof}
По определению $\beta' = \beta + \epsilon$,
где $+$ означает поразрядное сложение битовых шкал по модулю 2.
Далее 
$\zeta = \beta' H^{\mathrm{T}} =
(\beta + \epsilon) H^{\mathrm{T}} =
\beta H^{\mathrm{T}} + \epsilon H^{\mathrm{T}} =\epsilon H^{\mathrm{T}}$,
поскольку $\beta H^{\mathrm{T}}=\Bl{0}$.
\end{proof}
\end{frame}

\begin{frame}
\frametitle{6.4.8. Коды Хэмминга (2/4)}
{%\small
\renewcommand{\baselinestretch}{0.8}\selectfont
Отсюда вытекает метод кодирования и декодирования
с исправлением одиночной ошибки замещения разряда
в блоке длины $n$.

\begin{enumerate}
\item По формулам или таблицам \DMref{п.~6.4.5}{6.4.5.} подобрать число $m$ так,
чтобы $2^m\leq 2^n/(n+1)$ и положить $k\Assign n-m$.

%\vspace{-1pt}
\item Определить булевы матрицы: 
порождающую матрицу $G_{m\times n}$, строки которой линейно
независимы;
матрицу декодирования $G^{-1}_{n\times m}$,
такую что $GG^{-1}=I_m$;
проверочную матрицу $H_{k\times n}$, 
такую что $GH^{\mathrm{T}}=\Bl{0}_{m\times k}$.
Все эти матрицы определяются неоднозначно.

%\vspace{-1pt}
\item Разбить исходное сообщение на блоки длины $m$.  

%\vspace{-1pt}
\item Каждый блок $\alpha$ длины $m$ независимо 
умножить на порождающую матрицу,
получая блок $\beta \Assign\alpha G$ длины $n$.

%\vspace{-1pt}
\item Блок $\beta$ передать через канал, 
получая блок $\beta' = b_1\dots b_n$, возможно содержащий 
одну ошибку замещения.  

%\vspace{-1pt}
\item Вычислить синдром $\zeta \Assign \beta'H^{\mathrm{T}}$.

%\vspace{-1pt}
\item Если синдром $\zeta =\Bl{0}_k$, то ошибок нет,
а если $\zeta \ne \Bl{0}_k$,  то определить $j$~ --- номер 
того столбца проверочной матрицы,
 который совпадает с 
синдромом, после чего инвертировать $j$-й разряд: $b_j\Assign 1+b_j$.

%\vspace{-1pt}
\item Декодировать блок $\alpha' \Assign \beta' G^{-1}$. 
\end{enumerate}
}
\end{frame}

\begin{frame}
\frametitle{6.4.8. Коды Хэмминга (3/4)}
\begin{theorem}
В указанных обозначениях $\alpha = \alpha'$.
\end{theorem}
\begin{proof}
Если ошибок не случилось, и $\beta'=\beta$,
то $\alpha' = \alpha$ по свойству матриц 
$G$ и $G^{-1}$. Пусть произошла ошибка в разряде $j$.    
Известно, что $\zeta=\epsilon H^{\mathrm{T}}$. 
Поскольку битовая шкала $\epsilon$ 
содержит $1$ в $j$-м разряде и $0$ во всех остальных разрядах,
синдром равен $j$-му столбцу матрицы $H$.
Инверсия $j$-го разряда устраняет ошибку в коде 
и позволяет произвести обычное декодирование.  
\end{proof}
 
\begin{remark}
Шаги 1 и 2 выполняются только один раз, и их трудоёмкость
некритична. Прочие шаги выполняются для каждого блока,
и скорость их выполнения важна.
Особого внимания заслуживают шаги 6 и 7,
поскольку это <<накладные расходы>> ради помехоустойчивости на кодирование
и декодирования.     
\end{remark}

Основная идея кода Хэмминга состоит в том, 
чтобы в проверочной матрице $H$ 
использовать к качестве $j$-го столбца 
двоичную запись номера этого столбца.
Тогда вычисление синдрома сразу даёт
номер разряда, в котором произошла ошибка.
Если синдром равен нулю, то ошибки нет.
Соотношение между $n$ и $k$ таково, что
$k$ разрядов как раз достаточно для записи чисел в
диапазоне $1..n$.
  
\end{frame}

\begin{frame}
\frametitle{6.4.8. Коды Хэмминга (4/4)}
\begin{example}
Пусть $m=4$, тогда $k=3$ и $n=7$. 
Матрица $
H =
\begin{pmatrix}
0 & 0 & 0 & 1 & 1 & 1 & 1\\
0 & 1 & 1 & 0 & 0 & 1 & 1\\
1 & 0 & 1 & 0 & 1 & 0 & 1
\end{pmatrix}
$ является проверочной.
Пусть $\alpha=1100$ и используется систематическое кодирование. 
Тогда $\beta=1100{b_5}{b_6}{b_7}$.
Соотношения для проверочных символов получаем из условия 
$\beta H^{\mathrm{T}} = \Bl{0}_k$, то есть
$\A{i \in 1..k}{\sum_{j=1}^{n}{b_j}{H_{j, i}} = 0}$.
Имеем $b_5+b_6+b_7=0$, $1+b_6+b_7=0$, $1+ b_5+b_7=0$. 
Получаем $b_5=1$, $b_6=1$, $b_7=0$. Значит $\beta=1100110$.
Пусть во время прохождения через канал произошла ошибка в третьем символе: $\beta'=1110110$.
 Найдём синдром в данном случае:
$
\zeta =
\begin{pmatrix}
1 & 1 & 1 & 0 & 1 & 1 & 0
\end{pmatrix}
\begin{pmatrix}
0 & 0 & 1\\
0 & 1 & 0\\
0 & 1 & 1\\
1 & 0 & 0\\
1 & 0 & 1\\
1 & 1 & 0\\
1 & 1 & 1
\end{pmatrix}
=
\begin{pmatrix}
0\\
1\\
1
\end{pmatrix}
$. Имеем
$j = 3$, в этом разряде как раз и произошла ошибка.
\end{example}

\end{frame}



\begin{frame}
\label{6.4.9.}
\subsubsection{6.4.9. Исправление одного замещения }
\frametitle{6.4.9. Исправление одного замещения (1/3)}
%{\color{blue} Тема для сочинения: 
%общее представление кодов Хэмминга с помощью матриц.
%Пример с исправлением одного замещения.
%}
%Очевидно, что для помехоустойчивости вместе с
%%основным сообщением нужно передавать какую-то дополнительную
%%информацию, с помощью которой в случае ошибок можно восстановить
%%исходное сообщение. При этом нельзя считать, что дополнительная
%%информация не подвержена ошибкам. Задача состоит в том, чтобы
%%определить, какую дополнительную информацию следует передавать,
%%чтобы её объём был минимальным, а процедуры кодирования и
%%декодирования по возможности простыми.
%%
%%Рассмотрим построение \odmkey{кода Хэмминга}{Код (code)!Хэмминга (Hamming)},
%%который позволяет исправлять одиночные ошибки типа замещения.
%%Пусть сообщение $\alpha=a_1\dots a_m$ кодируется
%%словом $\beta=b_1\dots b_n$, $n>m$. Обозначим
%%$k\Assign n-m$.
%%Пусть канал допускает не более одной ошибки типа замещения
%%в слове длины $n$.
%%
%%{\renewcommand{\baselinestretch}{0.95}\selectfont
%%\begin{digress}
%%Рассматриваемый случай простейший, но одновременно практически
%%очень важный. Таким свойством, как правило, обладают внутренние
%%шины передачи данных в современных компьютерах. При этом часто
%%применяют аппаратно реализуемый метод повышения надёжности,
%%который называется \odmkey{контролем чётности}{Контроль чётности
%%(parity check)}. Суть метода заключается в том, что к шине
%%добавляется дополнительный разряд, значение которого определяется
%%при передаче как двоичная сумма остальных разрядов. Таким образом,
%%общее количество единиц, передаваемых по шине, всегда чётно. Если
%%при приёме переданное количество единиц оказывается нечётным,
%%диагностируется произошедшая ошибка. Контроль чётности позволяет
%%обнаружить одиночную ошибку типа замещения разряда, но не
%%позволяет её исправить.
%%\end{digress}
%%
%%%\addvspace{-3pt}
%%При заданном $n$ количество дополнительных разрядов $k$
%%подбирается таким образом, чтобы $2^k\ge n+1$.
%%Имеем
%%$$
%%2^k\ge n+1\Impl \frac{2^n}{n+1}\ge 2^{n-k}\Impl
%%\frac{2^n}{n+1}\ge 2^{m}.
%%$$
%%
%%%\addvspace{-12pt}
%%\begin{example}
%%Для сообщения длиной $m=32$ достаточно $k=6$ дополнительных разрядов,
%%поскольку $64=2^6>32+6+1=39$.
%%\end{example}

%\addvspace{-5pt}

Существуют другие способы линейного кодирования, 
отличные от систематических кодов, 
при  которых проверочные разряды находятся на позициях, 
облегчающих кодирование. Рассмотрим в качестве примера случай 
одного замещения. 


Определим последовательности натуральных чисел в соответствии с их
представлениями в двоичной системе счисления следующим образом:\\
$V_1 \Assign 1,3,5,7,9,11,\dots$\hspace{4pt}--- все числа,
у которых разряд 1 равен 1;\hfill\\ %
$V_2 \Assign 2,3,6,7,10,\dots$\quad~--- все числа, у которых
разряд 2
равен 1;\hfill\\ %
$V_3 \Assign 4,5,6,7,12,\dots$\quad~--- все числа, у которых
разряд 3
равен 1,\hfill\\ %
и т.~д. По определению последовательность $V_k$
начинается с числа $2^{k-1}$. Также рассмотрим последовательности
$V'_k$, где $V'_k$ получается из $V_k$ отбрасыванием первого
элемента. Эти последовательности определены статически
и могут быть вычислены заранее. 


Пусть заданы числа $n$, $m$, 
причём $2^m\leq 2^n/(n+1)$ и $k=n-m$.
Рассмотрим $k$ разрядов с
номерами $2^0=1$, $2^1=2$, $2^2=4$,~\ldots $2^{k-1}$ 
в слове $b_1\dots b_n$. Эти разряды
считаются контрольными.
Остальные разряды, а их ровно $m$, считаются
информационными. Поместим в информационные разряды все разряды
слова $a_1\dots a_m$ как они есть. 

\end{frame}



\begin{frame}
\frametitle{6.4.9. Исправление одного замещения (2/3)}

Контрольные
разряды определим следующим образом:\\
$b_1 \Assign b_3+ b_5+ b_7+\dots, $ то~есть сумма всех разрядов с
номерами из $V'_1$;\hfill\\
$b_2 \Assign b_3+ b_6+ b_7+\dots,$ то~есть  сумма всех разрядов с
номерами из $V'_2$;\hfill\\
$b_4 \Assign b_5+ b_6+ b_7+\dots,$ то~есть  сумма всех разрядов с
номерами из $V'_3$;\hfill\\
и вообще, $ b_{2^{j-1}} \Assign
\sum\limits_{i\in V'_j}b_i, $ причём сложение выполняется по
модулю 2. Пусть после прохождения через канал получен код
$c_1\dots c_n$ и пусть 
$I$~--- номер разряда, в котором, возможно, произошла ошибка
замещения. Пусть это число имеет следующее двоичное представление:
$I=i_l\dots i_1$. Определим число $J=j_l\dots j_1$ следующим
образом:\\ %
$j_1 \Assign c_1 + c_3 + c_5 + c_7+\dots,$ то~есть сумма всех
разрядов с номерами из $V_1$;\\
$j_2 \Assign c_2 + c_3+ c_6 + c_7+\dots,$ то~есть  сумма всех
разрядов с номерами из $V_2$;\\
$j_3 \Assign c_4 + c_5 + c_6 + c_7 + \dots,$ то~есть  сумма всех
разрядов с номерами из $V_3$;\\
и вообще, $ j_{p} \Assign
\sum\limits_{q\in V_p}c_q, $ где сложение выполняется по модулю 2.



\end{frame}



\begin{frame}
\frametitle{6.4.9. Исправление одного замещения (3/3)}

\begin{theorem}
В указанных обозначениях $I=J$.
\end{theorem}

\begin{proof}
Эти числа равны, потому что поразрядно равны их двоичные
представления. Действительно, пусть $i_1=0$. Тогда $I\notin V_1$
и, значит, $$j_1=c_1+ c_3+ c_5+\ldots = b_1+ b_3+ b_5+\ldots = 0$$
по определению
разряда $b_1$.
Пусть теперь $i_1=1$. Тогда $I\in V_1$
и, значит,
$$j_1=c_1+ c_3+ c_5+\ldots =
b_1+ b_3+ b_5+\ldots+\overline{b_x}+\ldots = 1,
$$
так~как если в сумме по модулю 2 изменить ровно один разряд, то
изменится и~значение всей суммы. Итак, $i_1=j_1$. Аналогично
(используя последовательности $V_2$ и т.~д.) имеем $i_2=j_2$ и
т.~д. Таким образом, $I=J$.
\end{proof}

%%\vspace{-6pt}
%Отсюда вытекает метод декодирования с исправлением ошибки:
%нужно вычислить число $J$. Если $J=0$, то ошибки нет,
%иначе $c_J\Assign\overline{c_J}$.
%После этого из исправленного сообщения извлекаются
%информационные разряды, которые уже не содержат ошибок.
%
%\begin{remark}
%Универсальный метод построения наилучшего
%помехоустойчивого кодирования для произвольного
%$m$ и произвольного типа канала $\delta$ неизвестен.
%Однако для большинства
%встречающихся на практике комбинаций $m$ и $\delta$
%задача построения кодирования с исправлением всех ошибок решена.
%\end{remark}
\vspace{6pt}
\end{frame}



\subsection{6.5. Сжатие и шифрование данных}
\label{6.5.}
\begin{frame}
\frametitle{6.5. Сжатие и шифрование данных}
В этом разделе обсуждаются
способы кодирования, которые применяются в тех
случаях, когда на коды
накладываются дополнительные требования,
помимо однозначности декодирования 
и снижения цены кодирования. 

\begin{tabular}{p{10mm}|p{42mm}|p{93mm}}
  № & Название параграфа
  & Ключевые термины и обозначения \\
  \hline
\DMref{6.5.1.}{6.5.1.} & Сжатие текстов &
{Сжатие данных; коэффициент сжатия}\\  
\DMref{6.5.2.}{6.5.2.} & Предварительное построение словаря &
{Слово; словарь; задача полнотекстового поиска}\\  
\DMref{6.5.3.}{6.5.3.} & Алгоритм Лемпела-Зива &
{Адаптивное сжатие}\\  
\DMref{6.5.4.}{6.5.4.} & Криптография &
{Шифрование; ключ; надежность шифра}\\ 
\DMref{6.5.5.}{6.5.5.} & Шифрование с помощью случайных чисел &
{Псевдослучайные числа; гамма шифра; симметричный шифр}\\  
\DMref{6.5.6.}{6.5.6.} & Криптостойкость &
{DES}\\  
\DMref{6.5.7.}{6.5.7.}\RS & Шифрование с открытым ключом &
{Открытый, закрытый ключ}\\  
\DMref{6.5.8.}{6.5.8.} &Цифровая подпись &
{}\\ 
\end{tabular}  

\end{frame}


\begin{frame}
\label{6.5.1.}
\subsubsection{6.5.1. Сжатие текстов }
\frametitle{6.5.1. Сжатие текстов (1/2)}
Материал раздела~ \DMref{6.3}{6.3.} показывает,
что при кодировании наблюдается некоторый баланс
между временем и памятью.
Затратив дополнительные усилия при кодировании и
декодировании, можно сэкономить память, и наоборот,
если пренебречь оптимальным использованием памяти,
можно существенно выиграть во времени кодирования и
декодирования. Конечно, этот баланс имеет место
только в определённых пределах, и
нельзя сократить расход памяти до нуля или
построить мгновенно работающие алгоритмы кодирования.
Для алфавитного кодирования пределы возможного
установлены в разделе~\DMref{6.3}{6.3.}. Для
достижения дальнейшего прогресса
нужно рассмотреть неалфавитное кодирование.

\medskip

Допустим, что имеется некоторое сообщение, которое закодировано
каким-то общепринятым способом (для текстов это, например, код
ASCII) и хранится в памяти компьютера. Заметим, что равномерное
кодирование (в частности, ASCII) не является для текстов
оптимальным. Действительно, в текстах обычно используется
существенно меньше, чем 256 символов (в зависимости от языка~---
примерно 60--80 с учётом знаков препинания, цифр, строчных и
прописных букв). Кроме того, вероятности появления букв различны,
и для каждого естественного языка известны (с некоторой точностью)
частоты появления букв в~тексте. 
\end{frame}


\begin{frame}
\frametitle{6.5.1. Сжатие текстов (2/2)}

Таким образом, можно задаться
некоторым набором букв и частотами их появления в тексте и с
помощью алгоритма Хаффмена построить оптимальное алфавитное
кодирование текстов (для заданного алфавита и языка). Простые
расчёты показывают, что такое кодирование для распространенных
естественных языков будет иметь цену кодирования несколько меньше
6, то~есть даст выигрыш по сравнению с кодом ASCII на 25\% или
несколько больше.

\begin{remark}
Методы кодирования, позволяющие построить (без потери информации!)
коды сообщений меньшей длины по сравнению с исходным сообщением,
называются методами \odmkey{сжатия}{Сжатие (compression)} (или
\odmkey{упаковки}{Упаковка (packing)}) данных. Качество сжатия
определяется \odmkey{коэффициентом сжатия}{Коэффициент
(coefficient)!сжатия (of compression)}, который обычно измеряется
в процентах и показывает, на сколько процентов кодированное
сообщение короче исходного.
\end{remark}

\medskip
Известно, что практические программы сжатия (rar, zip и др.) имеют
гораздо лучшие показатели, чем код Хаффмена: при сжатии текстовых
файлов коэффициент сжатия достигает 70\% и более. Это означает,
что в таких программах используется \em{не} алфавитное
кодирование.

\end{frame}


\begin{frame}
\label{6.5.2.}
\subsubsection{6.5.2. Предварительное построение словаря}
\frametitle{6.5.2. Предварительное построение словаря (1/3)}


Рассмотрим следующий способ кодирования.

\begin{enumerate}
\item Исходное сообщение по некоторому алгоритму разбивается на
последовательности букв, называемые \odmkey{словами}{Слово (word)}
(слово может иметь одно или несколько вхождений в исходный текст
сообщения). 
\item Полученное множество слов считается буквами
нового алфавита. Для этого алфавита строится разделимая схема
алфавитного кодирования (равномерного кодирования или оптимального
кодирования, если для каждого слова подсчитать число вхождений в
текст). Полученная схема обычно называется
\odmkey{словарем}{Словарь (dictionary)}, так~как она сопоставляет
слову код. 
\item Далее код сообщения строится как пара~--- код
словаря и последовательность кодов слов из данного словаря. 
\item
При декодировании исходное сообщение восстанавливается путем
замены кодов слов на слова из словаря.
\end{enumerate}

\end{frame}


\begin{frame}
\frametitle{6.5.2. Предварительное построение словаря (2/3)}
\begin{example} 
Допустим, что требуется кодировать тексты на русском языке. В
качестве алгоритма деления на слова примем естественные правила
языка: слова отделяются друг от друга пробелами или знаками
препинания. Можно принять допущение, что в каждом конкретном
тексте имеется не более $2^{16}$ различных слов (обычно гораздо
меньше).
Таким образом, каждому слову можно сопоставить номер~---
целое число из двух байтов (равномерное кодирование). Поскольку в
среднем слова русского языка состоят более чем из двух букв, такое
кодирование даёт существенное сжатие текста (около 75\% для
обычных текстов на русском языке). 
Если текст достаточно велик
(сотни тысяч или миллионы букв, то есть сотни и тысячи страниц),
то дополнительные затраты на хранение словаря оказываются
сравнительно небольшими.
\end{example}

\medskip
\begin{remark}
Данный метод попутно позволяет решить
задачу \odmkey{полнотекстового поиска}{Поиск (search)!полнотекстовый
 (full text)},
то~есть определить, содержится ли заданное слово
(или слова) в данном тексте, причём для этого
не нужно просматривать весь текст
(достаточно просмотреть только словарь).
\end{remark}
\end{frame}


\begin{frame}
\frametitle{6.5.2. Предварительное построение словаря (3/3)}


Указанный метод можно усовершенствовать следующим
образом. На шаге~2 следует применить алгоритм Хаффмена для
построения оптимального кода, а на шаге~1~--- решить экстремальную
задачу разбиения текста на слова таким образом, чтобы среди всех
возможных разбиений выбрать то, которое даёт наименьшую цену
кодирования на шаге~2. Такое кодирование будет <<абсолютно>>
оптимальным. К сожалению, указанная экстремальная задача очень
трудоёмка, поэтому на практике не используется~--- время на
предварительную обработку большого текста оказывается чрезмерно
велико даже при использовании современных быстродействующих
компьютеров.
\end{frame}

\begin{frame}
\label{6.5.3.}
\subsubsection{6.5.3. Алгоритм Лемпела--Зива }
\frametitle{6.5.3. Алгоритм Лемпела--Зива (1/7)}

На практике используется следующая идея, которая известна также
как \odmkey{адаптивное сжатие}{Сжатие (compression)!адаптивное
(adaptive)}. За один проход по тексту одновременно динамически
строится словарь и кодируется текст. При этом нет необходимости
хранить словарь: за счёт того, что при декодировании используется
тот же самый алгоритм построения словаря, словарь динамически
восстанавливается.

\medskip
Ниже приведена простейшая реализация этой идеи,\\ известная как
\odmkey{алгоритм Лемпела--Зива}{Алгоритм (algorithm)!Лемпела-Зива (Lempel-Ziv)}.

\medskip
Вначале словарь \textbf{$D:$ array $[$int$]$ of string} содержит
пустое слово, имеющее код~0. Далее в тексте последовательно
выделяются слова. Выделяемое слово~--- это максимально длинное
слово из уже имеющихся в словаре плюс ещё один символ. В сжатое
представление записываются найденный код слова и~расширяющая
буква, а словарь пополняется расширенной комбинацией.

\begin{remark}
В алгоритмах этого раздела знак <<$+$>>, применённый к целым числам,
означает сложение, а применённый к словам или символам~---
конкатенацию.
\end{remark}
\end{frame}

\begin{frame}
\frametitle{6.5.3. Алгоритм Лемпела--Зива (2/7)}
\vspace{-6pt}
%\begin{algorithm}
%\caption{Упаковка по методу Лемпела--Зива}
%\addvspace{-4pt}%
%\end{algorithm}%
\begin{algorithmic}
\REQUIRE исходный текст, заданный массивом кодов символов
\textbf{$f:$ array $[1..n]$ of char}. 
\ENSURE сжатый текст,
представленный последовательностью пар $\lrle{p,q}$, где $p$~---
номер слова в словаре, $q$~--- код дополняющей буквы. 
\STATE
$D[0]\Assign \varepsilon; d\Assign 0$ 
\COMMENT{ начальное
состояние словаря } 
\STATE $k\Assign 1$ 
\COMMENT{ номер текущей
буквы в исходном тексте } 
\WHILE{$k\le n$} 
\STATE $p\Assign \FD (k)$ 
\COMMENT{ $p$~--- индекс найденного слова в словаре } 
\STATE $l\Assign \Length(D[p])$ 
\COMMENT{ $l$~--- длина найденного слова
в словаре } 
\STATE \textbf{yield $\lrle{p,f[k+l]}$} 
\COMMENT{ код
найденного слова и ещё одна буква } 
\STATE $d\Assign d+1;
D[d]\Assign D[p] + f[k+l]$ \COMMENT{ пополнение словаря } \STATE
$k\Assign k+l+1$ \COMMENT{ продвижение вперед по исходному тексту
} \ENDWHILE
\end{algorithmic}



%Слово в словаре ищется с помощью несложной функции $\FD$, приведённой на следующем слайде.

\end{frame}

\begin{frame}
\frametitle{6.5.3. Алгоритм Лемпела--Зива (3/7)}

%\vspace{-3pt}
Слово в словаре ищется с помощью несложной функции $\FD$.

%\vspace{3pt}
\begin{algorithmic}
\REQUIRE $k$~--- номер символа в исходном тексте, начиная с
которого нужно искать в тексте слова из словаря. \ENSURE $p$~---
индекс самого длинного слова в словаре, совпадающего с символами\\
$f[k]..f[k+l]$. Если такого слова в словаре нет, то $p=0$.
\STATE $l\Assign 0; p\Assign 0$ \COMMENT{ начальное состояние }
\FOR{\textbf{$i$ from $1$ to $d$}} \STATE $m\Assign \Length(D[i])$
\COMMENT{ длина слова в словаре } \IF{$D[i]=f[k..k+m-1] \And m >
l$} \STATE $p\Assign i; l\Assign m$ \COMMENT{ нашли более
подходящее слово } \ENDIF \ENDFOR \STATE\textbf{return $p$}
\end{algorithmic}


\end{frame}

\begin{frame}
\frametitle{6.5.3. Алгоритм Лемпела--Зива (4/7)}
\vspace{-3pt}
Распаковка осуществляется следующим алгоритмом.

\vspace{-3pt}
%\begin{algorithm}
%\caption{Распаковка по методу Лемпела--Зива}
\begin{algorithmic}
\REQUIRE {сжатый текст, представленный
массивом пар \\ \textbf{$g:$ array $[1..m]$
of record $p:$ int;
$q:$ char end record},\\
где $p$~--- номер слова в словаре,
$q$~--- код дополняющей буквы.}
\ENSURE
исходный текст, заданный
последовательностью строк и символов.
\STATE $D[0]\Assign \varepsilon; d\Assign 0$
\COMMENT{ начальное состояние словаря }
\FOR{\textbf{$k$ from $1$ to $m$}}
\STATE $p\Assign g[k].p$ \COMMENT{ $p$~--- индекс слова в словаре }
\STATE $q\Assign g[k].q$ \COMMENT{ $q$~--- дополнительная буква }
%\pagebreak
\STATE \textbf{yield $D[p] + q$} \COMMENT{ вывод слова и ещё одной
буквы } \STATE $d\Assign d+1; D[d]\Assign D[p] + q$ \COMMENT{
пополнение словаря } \ENDFOR
\end{algorithmic}
\end{frame}

\begin{frame}
\frametitle{6.5.3. Алгоритм Лемпела--Зива (5/7)}
\begin{example}
Упаковать по методу Лемпела – Зива
строку  $abcdecf\ abcdac$.

\smallskip

\begin{tabular}{@{ }l@{ }|@{ }c@{ }cccccccccc}
\hline
Вход & & $a$ & $b$ & $c$ & $d$ & $e$ & $cf$ & $ $ &
$ab$ & $cd$ & $ac$ \\
\hline
Код $\divC{p}{q}$ & & $\divC{0}{a}$ & $\divC{0}{b}$ & $\divC{0}{c}$ &
$\divC{0}{d}$ & $\divC{0}{e}$ & $\divC{3}{f}$ & $\divC{0}{}$ & $\divC{1}{b}$ &
$\divC{3}{d}$ & $\divC{1}{c}$ \\
\hline
Индекс $d$  & 0 & 1 & 2 & 3 & 4 & 5 & 6 & 7 & 8 & 9 & 10 \\
\hline
\end{tabular}

\smallskip

Получили сжатый текст, представленный массивом пар:

$\divC{0}{a}$,
$\divC{0}{b}$,
$\divC{0}{c}$,
$\divC{0}{d}$,
$\divC{0}{e}$,
$\divC{3}{f}$,
$\divC{0}{ }$,
$\divC{1}{b}$,
$\divC{3}{d}$,
$\divC{1}{c}$.
\end{example}

\medskip


\begin{example}
Распаковать по методу Лемпела – Зива  

$\divC{0}{a}$,
$\divC{0}{b}$,
$\divC{0}{c}$,
$\divC{0}{d}$,
$\divC{0}{e}$,
$\divC{3}{f}$,
$\divC{0}{ }$,
$\divC{1}{b}$,
$\divC{3}{d}$,
$\divC{1}{c}$.

\smallskip


\begin{tabular}{@{ }l@{ }|@{ }c@{ }cccccccccc}
\hline
Код $\divC{p}{q}$ & & $\divC{0}{a}$ & $\divC{0}{b}$ & $\divC{0}{c}$ &
$\divC{0}{d}$ & $\divC{0}{e}$ & $\divC{3}{f}$ & $\divC{0}{}$ & $\divC{1}{b}$ &
$\divC{3}{d}$ & $\divC{1}{c}$ \\
\hline
Словарь & & $a$ & $b$ & $c$ & $d$ & $e$ & $cf$ & $ $ &
$ab$ & $cd$ & $ac$ \\
\hline

Индекс $d$  & 0 & 1 & 2 & 3 & 4 & 5 & 6 & 7 & 8 & 9 & 10 \\
\hline
\end{tabular}

\smallskip
Получили исходный текст $abcdecf\ abcdac$.
\end{example}
%\end{algorithm}
\end{frame}

\begin{frame}
\frametitle{6.5.3. Алгоритм Лемпела--Зива (6/7)}
%\pagebreak
На практике применяют различные усовершенствования этой схемы:
\begin{enumerate}
\item
Словарь можно сразу инициализировать, например,
кодами символов (то~есть считать, что однобуквенные слова
уже известны).
\item
В текстах часто встречаются регулярные последовательности:
пробелы, табуляции в~таблицах и т.~п. Сопоставлять каждой
подпоследовательности такой последовательности отдельное
слово в словаре нерационально.
В таких случаях лучше применить специальный
приём, например, закодировать последовательность пробелов
парой $\lrle{k,s}$, где $k$~--- количество пробелов,
а $s$~--- код пробела.
\end{enumerate}

\begin{remark}
В данном параграфе представлен алгоритм Лемпела--Зива LZ78 (1978~г.), 
не следует путать его с алгоритмом LZ77 (Лемпела--Зива, 1977~г.), 
который имеет другой принцип работы.
\end{remark}
%{\color{blue} Тема для сочинения: 
%всё сказанное было сделано в конце 70-х.
%А как сейчас обстоят дела? Быстрые алгоритмы сжатия видео с потерями,
%обработка физических и астрономических экспериментов с большими данными и т.д.
%Слова для блога в Интернете не нужны. Нужны формулы,
%которых нет в учебнике, но которые можно вывести на основе того, что есть.
%}
\end{frame}

\begin{frame}
\frametitle{6.5.3. Алгоритм Лемпела--Зива (7/7)}



\begin{digress}
Основная идея LZ77 состоит в том, 
что второе и последующие вхождения некоторой строки символов 
в сообщении заменяются ссылками на её первое вхождение. 
LZ77 использует уже просмотренную часть сообщения как словарь. 
Чтобы добиться сжатия, он пытается заменить очередной фрагмент 
сообщения на указатель в содержимое словаря. 
LZ77 использует <<скользящее>> по сообщению <<окно>>, 
разделённое на две неравные части. 
Первая, большая по размеру, включает уже просмотренную часть сообщения. 
Вторая является буфером, содержащим ещё незакодированные символы входного потока.
Алгоритм пытается найти в словаре (большей части окна) фрагмент, 
совпадающий с содержимым буфера. 
LZ77 выдаёт коды, состоящие из трёх элементов: 
смещение в словаре относительно его начала подстроки, 
совпадающей с началом содержимого буфера; 
длина этой подстроки; первый символ буфера, 
следующий за подстрокой. 
Недостатки: невозможность кодирования подстрок, 
отстоящих друг от друга на расстоянии, большем длины словаря; 
длина подстроки, которую можно закодировать, 
ограничена размером буфера. 
Рассмотренный Алгоритм Лемпела--Зива LZ78 не имеет этих недостатков.
\end{digress}

\end{frame}


\begin{frame}
\label{6.5.4.}
\subsubsection{6.5.4.  Криптография }
\frametitle{6.5.4. Криптография (1/3)}

Защита информации, хранящейся в компьютере, от
несанкционированного доступа, искажения и уничтожения в настоящее
время является серьёзной социальной проблемой. Применяются
различные подходы к решению этой проблемы.

\medskip
{\color{blue}\bf [1]}
Поставить между злоумышленником и данными
в компьютере непреодолимый <<физический>> барьер, то есть исключить
саму возможность доступа к~данным
путем изоляции компьютера с данными в охраняемом помещении,
применения аппаратных ключей защиты и т.~п.
Такой подход надёжен, но он затрудняет доступ
к данным для законных пользователей,
а потому постепенно уходит в прошлое.

\medskip
{\color{blue}\bf [2]}
Поставить между злоумышленником и данными
в компьютере <<логический>> %\linebreak
барьер, то есть проверять
наличие прав на доступ к данным и
блокировать доступ при отсутствии таких прав.
Для этого применяются
различные системы паролей, регистрация и идентификация
пользователей, разграничение прав доступа и т.~п.
Практика показывает, что борьба между
<<хакерами>> и~разработчиками
модулей защиты операционных систем
идёт с переменным успехом.

\end{frame}


\begin{frame}
\frametitle{6.5.4. Криптография (2/3)}

{\color{blue}\bf [3]}
Хранить данные таким образом, чтобы
они могли <<сами за себя постоять>>.
Другими словами, так закодировать данные,
чтобы, даже получив к ним доступ,
злоумышленник не смог бы нанести ущерба.

\medskip
Этот раздел посвящён обсуждению методов
кодирования, применяемых в последнем случае.

\medskip
\odmkey{Шифрование}{Шифрование (enciphering)}~---
это кодирование данных с целью защиты
от несанкционированного доступа.
Процесс кодирования сообщения называется
\odmkey{шифрованием}{Шифрование (enciphering)}
(или \odmkey{зашифровкой}{Зашифровка (enciphering)}),
а процесс декодирования~---
\odmkey{расшифровыванием}{Расшифровывание (deciphering)}
(или \odmkey{расшифровкой}{Расшифровка (deciphering)}).
Само кодированное сообщение называется
\odmkey{шифрованным}{Сообщение (message)!шифрованное (coded)}
(или просто \odmkey{шифровкой}{Шифровка (coded message)}),
а применяемый метод называется
\odmkey{шифром}{Шифр (cipher)}.

Основное требование к шифру состоит в том, чтобы расшифровка (и,
может быть, зашифровка) была возможна только при наличии санкции,
то~есть некоторой дополнительной информации (или устройства),
которая называется \odmkey{ключом}{Ключ (key)!шифра (of
encryption)} шифра. Процесс декодирования шифровки {\it без} ключа
называется \odmkey{дешифрованием}{Дешифрование (decoding)} (или
\odmkey{дешифрацией}{Дешифрация (decoding)}, или просто
\odmkey{раскрытием}{Шифр (cipher)!раскрытие (decoding)} шифра).

\end{frame}

\begin{frame}
\frametitle{6.5.4. Криптография (3/3)}

Область знаний о  методах создания
и раскрытия шифров называется
\odmkey{криптографией}{Криптография (cryptography)}
(или \odmkey{тайнописью}{Тайнопись (cryptogram)}).
Свойство шифра противостоять раскрытию называется
\odmkey{криптостойкостью}{Криптостойкость (cryptostability)} (или
\odmkey{надёжностью}{Шифр (cipher)!надёжный (robust)}) и обычно
оценивается сложностью алгоритма дешифрации.

\begin{digress}
В практической криптографии криптостойкость шифра
оценивается из экономических соображений.
Если раскрытие шифра стоит (в денежном выражении,
включая необходимые компьютерные ресурсы,
специальные устройства, услуги специалистов и т.~п.) больше,
чем сама зашифрованная информация,
то шифр считается достаточно надёжным.
\end{digress}


\medskip
Криптография известна с глубокой древности и использует самые
разнообразные шифры, как чисто информационные, так и механические.

\medskip
{\large
В настоящее время наибольшее практическое значение имеет защита
данных в компьютере, поэтому далее рассматриваются программные
шифры для сообщений в алфавите $\{0,1\}$.
}
\end{frame}

\begin{frame}
\label{6.5.5.}\subsubsection{6.5.5.  Шифрование с помощью случайных чисел }\frametitle{6.5.5. Шифрование с помощью случайных чисел (1/2)}

Пусть имеется датчик \odmkey{псевдослучайных}{Число
(number)!псевдослучайное (pseudorandom)} чисел, работающий по
некоторому определённому алгоритму. Часто используют следующий
алгоритм:
$
T_{i+1} \Assign (a\cdot T_i+b)\mod c,
$
где $T_i$~--- предыдущее
псевдослучайное число, $T_{i+1}$~---
следующее псевдослучайное число,
а коэффициенты $a$, $b$, $c$ постоянны и хорошо
известны. Обычно $c=2^n$, где~$n$~--- разрядность процессора,
$a\bmod 4 =1$, а $b$~--- нечётное.
В этом случае последовательность псевдослучайных чисел
имеет период $c$.
%(см.~\cite{Knuth2}).

Процесс шифрования определяется следующим образом. Шифруемое
сообщение представляется в виде последовательности слов
$S_0,S_1,\dots$, каждое длины $n$, которые поразрядно складываются
по модулю 2 со словами последовательности $T_0,T_1,\dots$, то~есть
$
C_i\Assign S_i +_{2} T_i.
$

\begin{remark}
Последовательность $T_0,T_1,\dots$ называется \odmkey{гаммой
шифра}{Гамма шифра (cypher gamma)}.
\end{remark}

Процесс расшифровывания заключается в том, чтобы ещё
раз сложить шифрованную последовательность с той же
самой гаммой шифра:
$
S_i\Assign C_i +_{2} T_i.
$
Ключом шифра является начальное значение $T_0$,
которое является секретным и~должно быть известно
только отправителю и получателю шифрованного сообщения.
\end{frame}

\begin{frame}
\frametitle{6.5.5. Шифрование с помощью случайных чисел (2/2)}
\begin{remark}
Шифры, в которых для зашифровки и расшифровки используется
один и тот же ключ, называются
\odmkey{симметричными}{Шифр (cipher)!симметричный (symmetric)}.
\end{remark}

Если период последовательности псевдослучайных чисел достаточно
велик, чтобы гамма шифра была длиннее сообщения, то дешифровать
сообщение можно только подбором ключа. При увеличении $n$
экспоненциально
увеличивается криптостойкость шифра.
\begin{digress}
Этот очень простой и эффективный метод часто
применяют <<внутри>> программных систем, например,
для защиты данных на локальном диске.
Для защиты данных, передаваемых
по открытым каналам связи,
особенно в случае многостороннего обмена
сообщениями,
этот метод применяют не так часто,
поскольку возникают трудности
с надёжной передачей секретного
ключа многим пользователям.
\end{digress}

\end{frame}

\begin{frame}
\label{6.5.6.}\subsubsection{6.5.6.  Криптостойкость }\frametitle{6.5.6. Криптостойкость (1/2)}

Описанный в предыдущем подразделе метод
шифрования обладает существенным недостатком.
Если известна хотя бы
часть исходного сообщения,
то всё сообщение может быть легко дешифровано.
Действительно, пусть известно одно
исходное слово $S_i$. Тогда
$$
T_i\Assign C_i +_{2} S_i,
$$
и далее вся правая часть гаммы шифра определяется
по указанной формуле датчика псевдослучайных чисел.

%\vspace*{-.5\baselineskip}

\begin{remark}
На практике часть сообщения вполне может быть
известна злоумышленнику. Например, многие
текстовые редакторы помещают в начало файла
документа одну и ту же служебную информацию.
Если злоумышленнику известно, что исходное
сообщение подготовлено в данном редакторе,
то он сможет легко дешифровать сообщение.
\end{remark}
\end{frame}


\begin{frame}
\frametitle{6.5.6. Криптостойкость (2/2)}
Для повышения криптостойкости симметричных шифров
применяют различные приёмы:
\begin{enumerate}
\item[1)] вычисление гаммы шифра по ключу более сложным (или
секретным) способом; 
\item[2)] применение вместо $ +_{2}$ более
сложной (но обратимой) операции для вычисления шифровки; 
\item[3)]
предварительное перемешивание битов исходного сообщения по
фиксированному алгоритму.
\end{enumerate}

Наиболее надёжным симметричным шифром считается
DES (Data Encryption Standard), в котором
используется сразу несколько методов повышения
криптостойкости.

%{\color{blue} Тема для сочинения: 
%Рассказать современную правду про DES и про ГОСТ.
%Эллиптические функции ...
%Это не очень просто, потому что правда прикрыта
%толстым слоем информационного шума в Интернете.
%}
\end{frame}

%\begin{frame}
%\frametitle{6.5.4. Модулярная арифметика}
%
%Для объяснения метода шифрования с открытым ключом,
%описанного в следующем подразделе,
%нужны некоторые факты из теории чисел,
%изложенные здесь без доказательств.
%
%{\color{blue} Темы для сочинений: 
%а нужно с хорошими доказательствами и примерами.
%Необходимо написать \em{раздел} во вторую главу
%в котором охватить, по меньшей мере:
%делимость + алгоритм Евклида + вычеты + функция Эйлера + всё что еще нужно, 
%чтобы выйти на следующую теорему и подготовить
%качественное изложение шифрования с открытым ключём.    
%}

%В этом подразделе рассматриваются только целые числа. Говорят, что
%число $a$ \odmkey{сравнимо по модулю}{Сравнение по
%модулю (modulo comparison)} $n$ с числом $b$ (обозначение $a\equiv b \bmod
%n$), если $a$ и $b$ при делении на $n$ дают один и тот же остаток:
%$$
%a\equiv b \bmod n\Def a\bmod n = b\bmod n.
%$$
%
%
%\begin{remark}
%В левой части этого определения слово $\bmod$
%является частью обозначения трехместного отношения,
%а в правой части слово $\bmod$ является обозначением
%двухместной операции определения остатка от деления.
%\end{remark}
%Отношение
%\odmkey{сравнимости}{Отношение (relation)!сравнимости чисел (congruence)}
%рефлексивно, симметрично и транзитивно
%и является отношением эквивалентности.
%Классы эквивалентности по отношению сравнимости
%(по модулю $n$)
%называются \odmkey{вычетами}{Вычет (residual)}
%(по модулю $n$). Множество вычетов по модулю $n$
%обозначается $\Bbb Z_n$.
%Обычно из каждого вычета выбирают одного представителя~---
%неотрицательное число,
%которое при делении на $n$
%дает частное~0.
%Это позволяет считать, что
%$\Bbb Z_n=\{0,1,2,\dots,n-1\}$,
%и упростить обозначения.
%Над вычетами (по модулю $n$)
%определены операции сложения и умножения
%по модулю $n$, обозначаемые,
%соответственно, $+_n$ и $\cdot_n$ и
%определяемые следующим образом:
%$$
%a +_n b \Def (a+b)\bmod n, \qquad
%a \cdot_n b \Def (a\cdot b)\bmod n.
%$$
%
%\begin{remark}
%Если из контекста ясно,
%что подразумеваются операции
%по модулю $n$, то индекс $n$ опускается.
%\end{remark}
%
%Легко видеть, что
%$\Alg{\Bbb Z_n}{+_n}$ является
%абелевой группой, а
%$\Alg{\Bbb Z_n }{+_n ,\cdot_n }$~--- коммутативным кольцом с единицей.
%
%\begin{remark}
%Если $n$~--- простое число, то $\Alg{{\Bbb Z}_n}{+_n , \cdot_n}$
%является полем.
%\end{remark}
%
%
%Рассмотрим ${\Bbb Z}_n^*$~--- подмножество множества
% вычетов $\Bbb Z_n$, взаимно простых с $n$.
%
%\begin{remark}
%Числа $a$ и $b$ называются
%\odmkey{взаимно простыми}{Число (number)!взаимно простое (mutually prime)},
%если их наибольший общий делитель равен~1.
%\end{remark}
%
%Можно показать, что $\Alg{{\Bbb Z}_n^*}{\cdot_n}$~---
%абелева группа. Таким образом, для элементов
%множества ${\Bbb Z}_n^*$ существуют обратные по умножению
%по модулю $n$.
%
%Функция $\varphi(n)\Def\vert{\Bbb Z}_n^*\vert$ называется
%\odmkey{функцией Эйлера}{Функция (function)!Эйлера (Euler)}.
%
%\begin{remark}
%Если $p$~--- простое число, то $\varphi(p)=p-1$, и вообще,
%$\varphi(n)<n$.
%\end{remark}
%
%Можно показать, что
%$$
%\varphi(n)=n\left(1-\frac{1}{p_1}\right)\cdots
%\left(1-\frac{1}{p_k}\right),
%$$
%где $\Tuple{p}{1}{k}$~--- все простые делители $n$.
%
%Имеет место следующая теорема. \pagebreak[3]
%\begin{theorem}[Эйлера]
%Если $n>1$, то
%$$
%\A{a\in {\Bbb Z}_n^*}{a^{\varphi(n)}\equiv 1 \bmod n}.
%$$
%\end{theorem}
%Отсюда непосредственно вытекает
%\begin{theorem}[малая теорема Ферма\footnote{\;Пьер де Ферма (1601--1665).}]
%Если $p>1$~{\upshape---} простое число, то
%$$
%\A{a\in {\Bbb Z}_p^*}{a^{p-1}\equiv 1 \bmod p}.
%$$
%\end{theorem}
%
%Имеет место следующее утверждение.
%\begin{theorem}
%Если числа $\Tuple{n}{1}{k}$ попарно взаимно просты, $n=n_1
%n_2\cdots n_k$~{\upshape---} их произведение, $x$ и
%$a$~{\upshape---} целые числа, то
%$$
%x\equiv a \bmod n \Equiv
%\A{i\in 1..k}{ x\equiv a \bmod n_i}.
%$$
%\end{theorem}

%\vspace*{-.5\baselineskip}

%\begin{remark}
%Последнее утверждение известно как <<китайская теорема об остатках>>.
%\end{remark}
%
%\end{frame}
%
%\begin{frame}
%\frametitle{6.5.5. Шифрование с открытым ключом}
%%\label{6.5.5}
%{\color{blue} Темы для сочинений: 
%этот вопрос откладывается на следующий семестр,
%но альтернативное изложение на основе предшествующих
%сочинений приветствуется.    
%}
%В настоящее время широкое распространение получили
%шифры \odmkey{с открытым ключом}{Шифр (cipher)!с открытым ключом (public key)}.
%Эти шифры не являются симметричными~--- для зашифровки и расшифровки
%используются разные ключи. При этом ключ, используемый для %\linebreak
%зашифровки,
%является открытым (не секретным) и может быть сообщен
%всем желающим отправить шифрованное сообщение,
%а ключ, используемый для расшифровки, является закрытым и
%хранится в секрете получателем шифрованных сообщений.
%Даже знание всего зашифрованного сообщения и открытого ключа,
%с помощью которого оно было зашифровано,
%не позволяет дешифровать сообщение
%(без знания закрытого ключа).
%
%
%Шифрование с открытым ключом производится следующим образом:
%
%\begin{enumerate}
%\item Получателем сообщений производится
%генерация открытого ключа (пара чисел $n$ и $e$) и
%закрытого ключа (число $d$). Для этого
%\begin{itemize}
%\item
%выбираются два простых числа, $p$ и $q$;
%
%\item
%вычисляется первая часть открытого ключа $n\Assign pq$;
%
%\item%
%определяется вторая часть открытого ключа~--- выбирается небольшое
%не\-чётное число $e$, взаимно простое с числом $(p-1)(q-1)$
%(заметим, что \linebreak$(p-1)(q-1) =
%pq(1-1/p)(1-1/q)=\varphi(n)$);
%
%\item
%определяется закрытый ключ: $d\Assign e^{-1}\bmod(p-1)(q-1)$.
%\end{itemize}
%
%После чего открытый ключ (числа $n$ и $e$) сообщается всем отправителям
%сообщений.
%\item
%Отправитель шифрует сообщение
%(разбивая его, если нужно,
%на слова $S_i$ длиной
%менее $\log n$ разрядов):
%$$
%C_i\Assign (S_i)^e\bmod n,
%$$
%и отправляет получателю.
%\item
%Получатель расшифровывает сообщение с помощью закрытого
%ключа $d$:
%$$
%P_i\Assign (C_i)^d\bmod n.
%$$
%\end{enumerate}
%
%\begin{theorem}
%Шифрование с открытым ключом корректно,
%то~есть в предыдущих обозначениях
%$P_i=S_i$.
%\end{theorem}
%
%\begin{proof}
%Легко видеть, что $P_i=(S_i)^{ed}\bmod n$. Покажем, что
%$$\A{{M<n}}{M^{ed}\equiv M\bmod n}.$$ Действительно, числа $d$ и $e$
%взаимно обратны по модулю ${(p-1)(q-1)}$, то~есть
%$$
%ed=1+k(p-1)(q-1)
%$$
%при некотором $k$.
%Если $M\not\equiv 0\bmod p$, то по малой теореме Ферма имеем
%$$
%M^{ed}\equiv M\left( M^{(p-1)}\right)^{k(q-1)}\equiv M\cdot 1^{k(q-1)}
%\equiv M \bmod p.
%$$
%Если $M\equiv 0 \bmod p$,
%то сравнение $M^{ed}\equiv M \bmod p$, очевидно, выполняется.
%Таким образом,
%$$
%\A{0\le M<n}{ M^{ed}\equiv M \bmod p}.
%$$
%Совершенно аналогично имеем
%$$
%\A{0\le M<n}{ M^{ed}\equiv M \bmod q},
%$$
%%\pagebreak
%и по китайской теореме об остатках
%$$\A{M<n}{M^{ed}\equiv M\bmod n}.$$
%Поскольку $S_i<n$ и $P_i<n$, заключаем,
%что $\A{i}{P_i=S_i}$.
%\end{proof}
%
%\begin{example}
%Генерация ключей:
%\begin{enumerate}
%\item
%$p\Assign 3$, $q\Assign 11$.%
%\item
%$n\Assign pq= 3*11 =33$.%
%\item
%$(p-1)(q-1)=2*10=20$, $e\Assign 7$.%
%\item
%$d\Assign 7^{-1}\bmod 20=3$, ($7*3\bmod 20 =1$).
%\end{enumerate}
%Пусть $S_1\Assign 3$, $S_2\Assign 1$, $S_3\Assign 2$
%($S_1,S_2,S_3<n=33$). Тогда код определяется следующим образом:
%\begin{enumerate}
%\item $C_1\Assign 3^7\bmod 33 = 2187\bmod 33 = 9$.%
%\item $C_2\Assign 1^7\bmod 33 = 1\bmod 33 = 1$.%
%\item $C_3\Assign 2^7\bmod 33 = 128\bmod 33 = 29$.
%\end{enumerate}
%При расшифровке имеем
%\begin{enumerate}
%\item $P_1\Assign 9^3\bmod 33 = 729\bmod 33 = 3$.%
%\item $P_2\Assign 1^3\bmod 33 = 1\bmod 33 = 1$.%
%\item $P_3\Assign 29^3\bmod 33 = 24389\bmod 33 = 2$.
%\end{enumerate}
%\end{example}
%
%\begin{digress}
%Шифры с открытым ключом сравнительно просты в реализации,
%очень практичны (поскольку нет необходимости пересылать
%по каналам связи закрытый ключ и можно
%безопасно хранить его в одном месте) и в то же время
%обладают высочайшей криптостойкостью.
%Кажется, что дешифровать сообщение несложно:
%достаточно разложить открыто
%опубликованное число $n$ на множители,
%восстановив числа $p$ и $q$, и далее можно
%легко вычислить секретный ключ $d$.
%Однако дело заключается в следующем.
%В настоящее время
%известны эффективные алгоритмы
%определения простоты чисел, которые позволяют за
%несколько минут подобрать пару очень
%больших простых чисел (по 100 и больше цифр
%в десятичной записи).
%В то же время неизвестны эффективные
%алгоритмы разложения очень больших
%чисел на множители. Разложение на множители
%числа в 200 и больше цифр потребовало бы
%сотен лет работы самого лучшего
%суперкомпьютера.
%При практическом применении
%шифров с открытым ключом
%используют действительно
%большие простые числа (не менее 100 цифр
%в десятичной записи, а обычно значительно больше).
%В~результате вскрыть этот шифр оказывается
%невозможно, если не существует
%эффективных алгоритмов разложения на множители
%(что очень вероятно, хотя и не доказано строго).
%\end{digress}

%\end{frame}

\begin{frame}

\label{6.5.7.}\subsubsection{6.5.7.  Шифрование с открытым ключом }\frametitle{6.5.7. Шифрование с открытым ключом (1/5)}

%При изложении метода шифрования с открытым ключом воспользуемся терминологией
%раздела 2.4, без которой корректное изложение этого метода невозможно.
%можно текст получше, но смысл такой
%он очень нужен, чтобы на mod не ссылать много раз
%дальше просто скопирован текст из учебника 6.5.5 Шифрование... со вставками ссылок
% 6.5.4 не нужен!

В настоящее время широкое распространение получили
\odmkey{шифры с открытым ключом}{Шифр (cipher)!с открытым ключом (public key)}.
Эти шифры не являются симметричными~--- для зашифровки и расшифровки
используются разные ключи. При этом ключ, используемый для %\linebreak
зашифровки,
является открытым (не секретным) и может быть сообщен
всем желающим отправить шифрованное сообщение,
а ключ, используемый для расшифровки, является закрытым и
хранится в секрете получателем шифрованных сообщений.
Даже знание всего зашифрованного сообщения и открытого ключа,
с помощью которого оно было зашифровано,
не позволяет дешифровать сообщение
(без знания закрытого ключа).

\end{frame}

\begin{frame}
\frametitle{6.5.7 Шифрование с открытым ключом (2/5)}

Шифрование с открытым ключом производится следующим образом.

\begin{enumerate}
\item Получателем сообщений производится
генерация открытого ключа (пара чисел $n$ и $e$) и
закрытого ключа (число $d$). Для этого:
\begin{itemize}
\item
выбираются два простых числа (\DMref{п.~2.4.4}{2.4.4.}), $p$ и $q$;

\item
вычисляется первая часть открытого ключа $n\Assign pq$;

\item%
определяется вторая часть открытого ключа~--- выбирается небольшое
не\-чётное число $e$, взаимно простое 
(\DMref{п.~ 2.4.2}{2.4.2.}) с числом $(p-1)(q-1)$
(заметим, что \linebreak$(p-1)(q-1) =
pq(1-1/p)(1-1/q)=\varphi(n)$, \DMref{п.~2.4.9}{2.4.9.});

\item
определяется закрытый ключ: $d\Assign e^{-1}\bmod(p-1)(q-1)$.
\end{itemize}

После чего открытый ключ (числа $n$ и $e$) сообщается всем отправителям
сообщений.
\item
Отправитель шифрует сообщение
(разбивая его, если нужно,
на слова $S_i$ длиной
менее $\log n$ разрядов): $C_i\Assign (S_i)^e\bmod n$,
и отправляет получателю.
\item
Получатель расшифровывает сообщение с помощью закрытого
ключа $d$: \\ $P_i\Assign (C_i)^d\bmod n$.
\end{enumerate}

\vspace{-2pt}
\begin{theorem}
Шифрование с открытым ключом корректно,
то~есть в предыдущих обозначениях
$P_i=S_i$.
\end{theorem}

\end{frame}

\begin{frame}
\frametitle{6.5.7 Шифрование с открытым ключом (3/5)}


\begin{proof}
Легко видеть, что $P_i=(S_i)^{ed}\bmod n$. Покажем, что
$$\A{{M<n}}{M^{ed}\equiv M\bmod n}.$$ Действительно, числа $d$ и $e$
взаимно обратны по модулю ${(p-1)(q-1)}$, то~есть
$$
ed=1+k(p-1)(q-1)
$$
при некотором $k$.
Если $M\not\equiv 0\bmod p$, то по малой теореме Ферма (\DMref{п.~2.4.9}{2.4.9.}) имеем
$$
M^{ed}\equiv M\left( M^{(p-1)}\right)^{k(q-1)}\equiv M\cdot 1^{k(q-1)}
\equiv M \bmod p.
$$
Если $M\equiv 0 \bmod p$,
то сравнение $M^{ed}\equiv M \bmod p$, очевидно, выполняется.
Таким образом,
$
\A{0\le M<n}{ M^{ed}\equiv M \bmod p}.
$
Совершенно аналогично имеем\\
$
\A{0\le M<n}{ M^{ed}\equiv M \bmod q}
$
и по следствию из китайской теоремы об остатках (\DMref{п.~2.4.7}{2.4.7.})
$\A{M<n}{M^{ed}\equiv M\bmod n}.$
Поскольку $S_i<n$ и $P_i<n$, заключаем,\\
что $\A{i}{P_i=S_i}$.
%\pagebreak
\end{proof}

\end{frame}

\begin{frame}
\frametitle{6.5.7. Шифрование с открытым ключом (4/5)}

\begin{example}
Генерация ключей:
\begin{enumerate}
\item
$p\Assign 3$, $q\Assign 11$.%
\item
$n\Assign pq= 3*11 =33$.%
\item
$(p-1)(q-1)=2*10=20$, $e\Assign 7$.%
\item
$d\Assign 7^{-1}\bmod 20=3$, ($7*3\bmod 20 =1$).
\end{enumerate}
Пусть $S_1\Assign 3$, $S_2\Assign 1$, $S_3\Assign 2$
($S_1,S_2,S_3<n=33$). Тогда код определяется следующим образом:
\begin{enumerate}
\item $C_1\Assign 3^7\bmod 33 = 2187\bmod 33 = 9$.%
\item $C_2\Assign 1^7\bmod 33 = 1\bmod 33 = 1$.%
\item $C_3\Assign 2^7\bmod 33 = 128\bmod 33 = 29$.
\end{enumerate}
При расшифровке имеем
\begin{enumerate}
\item $P_1\Assign 9^3\bmod 33 = 729\bmod 33 = 3$.%
\item $P_2\Assign 1^3\bmod 33 = 1\bmod 33 = 1$.%
\item $P_3\Assign 29^3\bmod 33 = 24389\bmod 33 = 2$.
\end{enumerate}
\end{example}

\end{frame}

\begin{frame}

\frametitle{6.5.7 Шифрование с открытым ключом (5/5)}

\vspace{-2pt}
%{\small
\begin{digress}
Шифры с открытым ключом сравнительно просты в реализации,
очень практичны (поскольку нет необходимости пересылать
по каналам связи закрытый ключ и можно
безопасно хранить его в одном месте) и в то же время
обладают высочайшей криптостойкостью.
Кажется, что дешифровать сообщение несложно:
достаточно разложить открыто
опубликованное число $n$ на множители,
восстановив числа $p$ и $q$, и далее можно
легко вычислить секретный ключ $d$.
Однако дело заключается в следующем.
В настоящее время
известны эффективные алгоритмы
определения простоты чисел, которые позволяют за
несколько минут подобрать пару очень
больших простых чисел (по 100 и больше цифр
в десятичной записи).
В то же время неизвестны эффективные
алгоритмы разложения очень больших
чисел на множители. Разложение на множители
числа в 200 и больше цифр потребовало бы
сотен лет работы самого лучшего
суперкомпьютера.
При практическом применении
шифров с открытым ключом
используют действительно
большие простые числа (не менее 100 цифр
в десятичной записи, а обычно значительно больше).
В результате вскрыть этот шифр оказывается
невозможно, если не существует
эффективных алгоритмов разложения на множители
(что очень вероятно, хотя и не доказано строго).
\end{digress}
%}
\end{frame}


\begin{frame}
\label{6.5.8.}\subsubsection{6.5.8.  Цифровая подпись }\frametitle{6.5.8. Цифровая подпись (1/3)}

Шифр с открытым ключом позволяет выполнять и многие другие
полезные операции, помимо шифрования и посылки сообщений в
одну сторону. Прежде всего, для организации
многосторонней секретной связи каждому из участников
достаточно сгенерировать свою пару ключей (открытый и закрытый),
а затем сообщить всем партнёрам свой открытый ключ.

Заметим, что операции зашифровки и расшифровки, по существу,
одинаковы, и~различаются только показателем степени,
а потому коммутируют:
$$
M=({M^e})^d \bmod n = M^{ed} \bmod n =
M^{de} \bmod n = ({M^d})^e\bmod n = M.
$$
Это обстоятельство позволяет применять
различные приёмы, известные как
\odmkey{цифровая}{Цифровая подпись (digital signature)}
(или \odmkey{электронная}{Электронная подпись (digital signature)}) подпись.

Рассмотрим следующую схему взаимодействия
корреспондентов $X$ и $Y$.
Отправитель $X$ кодирует сообщение $S$
своим закрытым ключом ($C\Assign S^d \bmod n$)
и посылает получателю $Y$ пару $\lrle{S,C}$,
то~есть подписанное сообщение. Получатель~$Y$,
получив такое сообщение, кодирует подпись сообщения
открытым ключом отправителя $X$. 
\end{frame}

\begin{frame}
\frametitle{6.5.8. Цифровая подпись (2/3)}
\vspace{-0.5\baselineskip}
То~есть вычисляет $S^\prime\Assign C^e \bmod n$.
Если оказывается, что $S=S^\prime$,
то это означает, что (нешифрованное!)
сообщение $S$ действительно было отправлено
корреспондентом $X$. Если же $S\ne S^\prime$,
то сообщение было искажено при передаче
или фальсифицировано.

\begin{digress}
В подобного рода схемах возможны различные проблемы, которые носят
уже не математический, а социальный характер. Например, допустим,
что злоумышленник $Z$ имеет техническую возможность контролировать
всю входящую корреспонденцию получателя $Y$ незаметно для
последнего. Тогда, перехватив сообщение отправителя $X$, в котором
сообщался открытый ключ~$e$, злоумышленник $Z$ может подменить
открытый ключ отправителя $X$ своим собственным открытым ключом.
После этого злоумышленник сможет фальсифицировать все сообщения
отправителя $X$, подписывая их своей цифровой подписью, и, таким
образом, действовать от имени отправителя $X$. Другими словами,
цифровая подпись удостоверяет, что сообщение $S$ пришло из того же
источника, из которого был получен 
открытый ключ $e$, но не более того.
\end{digress}

\end{frame}

\begin{frame}
\frametitle{6.5.8. Цифровая подпись (3/3)}
Можно подписывать и шифрованные сообщения.
Для этого отправитель $X$ сначала кодирует
своим закрытым ключом сообщение $S$,
получая цифровую подпись $C$, а затем
кодирует полученную пару $\lrle{S,C}$
открытым ключом получателя~$Y$.
Получив такое сообщение, получатель $Y$ сначала
расшифровывает сообщение своим закрытым ключом,
а потом убеждается в подлинности полученного сообщения,
сравнив его
с результатом применения открытого ключа отправителя
$X$ к подписи $C$.

\begin{remark}
К сожалению, даже эти меры не смогут защитить
от злоумышленника $Z$, сумевшего подменить
открытый ключ отправителя $X$. Конечно, в этом случае
злоумышленник не сможет дешифровать исходное сообщение,
но он сможет подменить исходное сообщение
фальсифицированным, а получатель не сможет
обнаружить подмену сообщения.
\end{remark}

%{\color{blue} Тема для сочинения: 
%снабдить всё это диаграммами последовательности про 
%Алису и Боба.    
%}

\end{frame}




% ===== tex/DM2022_7.tex =====

\section{7. Графы}

\begin{frame}
\frametitle{7. Графы (1/3)}
{\large
Эта глава открывает заключительную часть курса, целиком
посвящённую графам и алгоритмам на них. Среди дисциплин и методов
дискретной математики теория графов, и особенно 
\em{алгоритмы на графах},
находят наиболее широкое применение в программировании.

\smallskip 
Как
показано в разделе~\DMref{7.5}{7.5.}, между понятием графа и понятием
бинарного отношения, рассмотренным в главе~1, имеется
глубокая связь~--- можно сказать, что это равнообъёмные понятия.
Возникает естественный вопрос: почему же тогда графам оказывается
столь заметное предпочтение при изучении дискретной математики и
программирования? 
}


\end{frame}

\begin{frame}
\frametitle{7. Графы (2/3)}
{\large
Дело в том, что теория графов предоставляет
очень удобный \em{язык\/} для описания программных (да и многих
других) моделей. Этот тезис можно пояснить следующей аналогией.
Понятие отношения также можно полностью выразить через понятие
множества (см. замечание в \DMref{п.~1.4.1}{1.4.1.} и далее). 


Однако
независимое определение понятия отношения \em{удобнее\/}~---
введение специальных терминов и обозначений упрощает изложение
теории и делает её более понятной. То же относится и к теории
графов. 

\smallskip
Стройная система специальных терминов и обозначений теории
графов позволяет просто и доступно описывать сложные и тонкие
вещи. Особенно важна возможность наглядной графической
интерпретации понятия графа (\DMref{п.~7.1.4}{7.1.4.}). Само название
<<граф>> подразумевает наличие графической интерпретации. Картинки
часто позволяют сразу <<усмотреть>> суть дела на интуитивном
уровне, дополняя и украшая утомительные текстовые доказательства и
сложные формулы. 
}
\end{frame}

\begin{frame}
\frametitle{7. Графы (3/3)}
{\large
Представление графов и подобных им объектов в программах даёт
убедительный пример решающей роли представления данных
в программировании. Для графов множественность альтернативных представлений
в программах
не только теоретически возможна, но и практически неизбежна.
Многие алгоритмы на
графах достаточно эффективны в том смысле, что имеют невысокую
теоретическую оценку трудоёмкости.
Такие алгоритмы бывают чувствительны к аддитивным и мультипликативным 
константам, которыми пренебрегают при получении асимптотической оценки.
Значе\-ние этих констант определяется представлением графа.
На практике, особенно для больших графов,
неудачно выбранное представление может воспрепятствовать
решению конкретной задачи на имеющемся оборудовании.
%
%\smallskip
Эта глава практически полностью посвящена
описанию \em{языка} теории графов.
Конкретные методы и алгоритмы более подробно 
рассматриваются далее.
}

\end{frame}


\subsection{7.1. Определения графов}\label{7.1.}
\begin{frame}
\frametitle{7.1. Определения графов (1/2)}
Как это ни удивительно, но у понятия <<граф>>
нет общепризнанного единого определения.
Разные авторы, особенно применительно к разным приложениям,
называют словом <<граф>> очень похожие, но всё-таки
различные объекты.
Здесь используется терминология из книги Харари
<<Теория графов>>,
которая была выбрана из соображений максимального
упрощения определений и доказательств.

\medskip
Главной целью этого раздела 
является явное определение и разграничение понятий:
диаграмма, граф и абстрактный граф.
Эти понятия тесно связаны, но их не следует смешивать.  
\end{frame}



\begin{frame}
\frametitle{7.1. Определения графов (2/2)}
\begin{tabular}{p{9mm}|p{50mm}|p{84mm}}
  № & Название параграфа
  & Ключевые термины и обозначения \\
  \hline
\DMref{7.1.1.}{7.1.1.} & История теории графов 
  & 
  Задачи о Кёнигсбергских мостах, трёх домах и трёх колодцах, четырёх красках \\
\DMref{7.1.2.}{7.1.2.}\RS & Основное определение 
  & 
  {$G(V,E)$~--- граф; $V$~--- множество вершин;
   $E$~--- множество рёбер; $p$~--- число вершин;
   $q$~--- число рёбер; смежность, инцидентность}
  \\
\DMref{7.1.3.}{7.1.3.} & Инцидентность, смежность и дополнение &
{$\Gamma(v)$~--- множество смежности;
 $\Lambda(v)$~--- множество инцидентности;
$\overline{G}$~--- дополнение графа }\\
\DMref{7.1.4.}{7.1.4.} & Диаграмма графа &
{\includegraphics[scale=1]{7_1.jpg}
}
\\
\DMref{7.1.5.}{7.1.5.} & Орграфы, псевдографы, мульти\-графы и гиперграфы &
{Узел, дуга, петля, гиперребро, помеченные и нумерованные вершины и рёбра }\\   
\DMref{7.1.6.}{7.1.6.}\RS & Изоморфизм графов &
{$\sim$~--- отношение изоморфизма графов;
 инварианты; абстрактный граф, рёберный изоморфизм}\\   
\end{tabular}
\end{frame}


\begin{frame}
\label{7.1.1.}\subsubsection{7.1.1. История теории графов }\frametitle{7.1.1. История теории графов (1/4)}
%\label{7.1.1}
Теория графов неоднократно переоткрывалась разными авторами при решении
различных прикладных задач.

{\color{blue}\bf 1.}
\odmkey{Задача о Кёнигсбергских мостах}{Задача (problem)!о
Кёнигсбергских мостах 
(K$\ddot{\text{o}}$nigsberg bridge problem)}. На
рисунке представлен схематический план центральной части
города Кёнигсберг (ныне Калининград), включающий два берега реки
Преголя, два острова на ней и семь соединяющих их мостов. Задача
состоит в том, чтобы обойти все четыре участка суши 
(обозначены $A$, $B$, $C$, $D$), пройдя по
каждому мосту один раз, и вернуться в исходную точку. В 1736 году
Эйлером
%\footnotemark \footnotetext{\;Леонард Эйлер (1707--1783).}
было показано, что решения этой задачи не существует. Точнее
говоря, Эйлер  получил необходимое и достаточное условие
существования решения для всех задач типа задачи о Кёнигсбергских
мостах (см. \DMref{п.~10.2.1}{10.2.1.}).
%\begin{figure}[h]

\vspace{-12pt}
\begin{center}
\includegraphics[height=30mm]{7_1_1_1.jpg}
\end{center}
%\caption{Иллюстрация к задаче о Кёнигсбергских мостах}
%\label{Euler}
%\end{figure}
\end{frame}

\begin{frame}
\frametitle{7.1.1. История теории графов (2/4)}
{\color{blue}\bf 2.}
\odmkey{Задача о трёх домах и трёх колодцах} {Задача
(problem)!о трёх домах и трёх колодцах (three houses, three
utilities)}. Имеется три дома и три колодца, каким-то образом
расположенные на плоскости. Требуется провести от каждого дома к
каждому колодцу тропинку так, чтобы тропинки не пересекались. 
Эта задача также не имеет решения. В 1930 году
Куратовским
%\footnote{\;Казимир Куратовский (1896--1979).} 
было
доказано намного более сильное утверждение, а именно достаточность
условия существования решения для всех задач типа задачи о трёх
домах и трёх колодцах (необходимость этого условия была известна и
ранее, см. \DMref{п.~10.8.2}{10.8.2.}).
%\begin{figure}[h]
\begin{center}
\includegraphics{07_02}
\end{center}
%\caption{Иллюстрация к задаче о трёх домах и трёх колодцах}
%\label{Kura}
%\end{figure}
\end{frame}

\begin{frame}
\frametitle{7.1.1. История теории графов (3/4)}
{\color{blue}\bf 3.} 
\odmkey{Задача о четырёх красках}{Задача (problem)!о четырёх
красках (four color)}. Разделение плоскости на неперекрывающиеся
области называется \odmkey{картой}{Карта (map)}. Области на карте
называются соседними, если они имеют общую границу. Задача состоит
в раскрашивании карты таким образом, чтобы никакие две соседние
области не были закрашены одним цветом. \\ 
С конца XIX века известна гипотеза, что для этого достаточно четырёх
красок. В 1976 году Аппель и Хейкен опубликовали решение задачи о
четырёх красках, которое базировалось на переборе вариантов с
помощью компьютера.
%\begin{figure}[h]
\begin{center}
\includegraphics[height=35mm]{7_1_1_3.jpg}
\end{center}
%\caption{Иллюстрация к задаче о четырёх красках}
%\label{4paint}
%\end{figure}
%\end{enumerate}
\end{frame}

\begin{frame}
\frametitle{7.1.1. История теории графов (4/4)}
\begin{digress}
Решение задачи о четырёх красках Аппелем и Хейкеном
<<программным путем>> явилось прецедентом, породившим
бурную дискуссию, которая отнюдь не закончена.
Суть опубликованного решения состоит в том, чтобы перебрать
большое, но конечное число (около 2000) типов
потенциальных контрпримеров к теореме о четырёх красках
и показать, что ни один случай
контрпримером не является. Этот перебор был выполнен
программой  примерно за тысячу часов работы суперкомпьютера.
Проверить <<вручную>> полученное решение невозможно~---
объём перебора выходит далеко за рамки человеческих возможностей.
Многие математики ставят вопрос: можно ли считать
такое <<программное доказательство>> действительным
доказательством? Ведь в программе могут  быть
ошибки... Методы доказательства правильности
программ не применимы к программам такой сложноcти, как
обсуждаемая. Тестирование не может гарантировать отсутствие
ошибок и~в~данном случае вообще невозможно.
Таким образом, остаётся уповать
на программистскую квалификацию авторов.
% и верить,
%что они всё сделали правильно.
\end{digress}

\end{frame}


\begin{frame}
\label{7.1.2.}\subsubsection{7.1.2. Основное определение}
\frametitle{7.1.2. Основное определение}
\vspace{-4pt}
\odmkey{Графом}{Граф (graph)} $G(V,E)$ называется совокупность
двух множеств~--- непустого множества~$V$ (множества
\odmkey{вершин}{Вершина (vertex)!графа (of graph)}) и
множества~$E$ двухэлементных подмножеств множества~$V$ ($E$~---
множество \odmkey{рёбер}{Ребро (edge)!графа (of graph)}),
$
G(V,E) \Def \Alg{V}{E\/}, \quad V \not= \emptyset, \quad
E\subset 2^V\And \A{e\in E}{|e|=2}.
$

\begin{remark}
Вершины графа считаются различными, им можно давать имена или номера. Чтобы подчеркнуть это обстоятельство,
иногда к слову <<граф>> добавляют прилагательное
и говорят \odmkey{помеченный}{Граф (graph)!помеченный (labeled)}. 
\end{remark}

Два графа \odmkey{равны}{Равенство (equality)!графов (of graphs)!}, если равны множества вершин и рёбер:\\
 $G_1(V_1,E_1)=G_2(V_2,E_2)\Def V_1=V_2\And E_1=E_2$.

\begin{remark}
Легко видеть, что любое множество $E$ двухэлементных подмножеств
множества $V$ определяет симметричное бинарное отношение на
множестве $V$. Поэтому можно считать, что
$E \subset V \times V$, $E=E^{-1}$,
и трактовать ребро не только как множество $\{v_1,v_2\}$,
но и как пару $(v_1,v_2)$.
\end{remark}

 Число вершин графа $G$ обозначим $p$, а число
рёбер~--- $q$:
$\ p(G)\Def  \vert V\vert, q(G)\Def  \vert E\vert.
$

Если хотят явно упомянуть числовые характеристики графа, то
говорят: $(p,q)$-граф. Про граф с числом вершин $p$ 
неформально говорят <<граф \em{на} $p$ вершинах>>.
Множества вершин и рёбер считаются конечными, если явно не оговорено противное.

\end{frame}

\begin{frame}
\label{7.1.3.}\subsubsection{7.1.3. Инцидентность, смежность и дополнение }\frametitle{7.1.3. Инцидентность, смежность и дополнение (1/3) }
Пусть $v_1$, $v_2$~--- вершины, $e=(v_1,v_2)$~--- соединяющее их
ребро. Тогда вершина~$v_1$ и ребро $e$
\odmkey{инцидентны}{Инцидентность (incidence)}, ребро $e$ и
вершина $v_2$ также инцидентны. Два ребра, инцидентные одной
вершине, называются \odmkey{смежными}{Смежность (adjacency)!рёбер
(edge)}; две вершины, инцидентные одному ребру, также называются
\odmkey{смежными}{Смежность (adjacency)!вершин (vertex)}.

\smallskip
Множество вершин, смежных с вершиной $v$, называется
\odmkey{множеством смежности}{Множество (set)!смежности
(adjacency)} (или \odmkey{окрестностью}{Окрестность
(neighborhood)!вершины (of vertex)}) вершины $v$
 и~обозначается
$\Gamma^{+}(v)$. По определению
$$\Gamma^{+}(v) \Def \Set{u\in V}{(u,v)\in E},\quad 
\Gamma^{*}(v) \Def  \Gamma^{+}(v) + v .$$


\begin{remark}
Если не оговорено противное, то символ
$\Gamma$  подразумевает $\Gamma^+$,
то есть саму вершину в окрестность не включают.
Важно отметить, что в данном случае 
множество вершин $A$ может пересекаться с 
$\Gamma(A)$.
\end{remark}

Очевидно, что $u\in \Gamma(v) \Equiv v \in \Gamma(u)$.
Если $A \subset V$~--- множество вершин, то $\Gamma(A)$~--- множество
всех вершин, смежных с вершинами из $A$:
$$
\Gamma(A)\Def \Set{u\in V}{\E{v\in A}{u\in\Gamma(v)}}=
\bigcup_{v\in A}\Gamma(v).
$$

\end{frame}

\begin{frame}
\frametitle{7.1.3.  Инцидентность, смежность и дополнение (2/3) }

Множество рёбер, инцидентных вершине $v$, 
называется \odmkey{множеством инцидентности}{Множество (set)!инцендентности вершин (of vertex incidence)} вершины $v$ 
и обозначается $\Lambda(v)$. По определению
$\Lambda(v)\Def \Set{e\in E}{v\in e}$.

\medskip
Рассмотрим граф $G(V,E)$ с некоторым множеством вершин $V$.
\odmkey{Дополнением}{Дополнение (complement)!графа (of graph)} $\overline{G}(V,\overline{E})$ называется другой граф,
с тем же множеством вершин $V$, 
в котором вершины смежны тогда и только тогда, 
когда они не смежны в исходном графе.
Другими словами, если ребро $e=(u,v)$ принадлежит графу,
то оно не принадлежит дополнению графа,
 и обратно,
если ребро $e = (u,v)$ принадлежит дополнению графа,
то оно не принадлежит графу.

\medskip
\begin{theorem}
Суммарное количество рёбер в графе и в его дополнении
определяется числом вершин $p$ и равно $p(p-1)/2$.
\end{theorem}
\begin{proof}
Число двухэлементных подмножеств множества вершин~--- это число
сочетаний $C(p,2)=\frac{p!}{2!(p-2)!}=p(p-1)/2$.  
\end{proof}

\medskip
\begin{corollary} 
Для любой вершины суммарное количество вершин смежных и не смежных с ней
 равно $p-1$.
\end{corollary}

\end{frame}

\begin{frame}
\frametitle{7.1.3.  Инцидентность, смежность и дополнение (3/3)}
Помеченный граф на $p$ вершинах определяется тем,
какие рёбра входят, а какие не входят в граф.
Из этого наблюдения немедленно следует

\begin{lemma}
Число помеченных графов на $p$ вершинах равно
$2^{p(p-1)/2}$.
\end{lemma} 


\medskip
Дополнение можно рассматривать и как операцию на графах, и как отношение
между графами.
Как операция,
дополнение инволютивно: дополнением к дополнению графа является сам граф.
Как отношение,
дополнение симметрично: если $\overline G$~--- дополнение $G$,
то и $G$~--- дополнение $\overline G$. 

\medskip
\begin{example}
Пусть 
$G\Assign(\{1, 2, 3, 4\}, \{(1, 2), (3, 4)\})$.
\\Тогда
$\overline{G}=(\{1, 2, 3, 4\}, 
\{(1, 3), (2, 4),(1, 4),(2, 3) \})$.
\end{example} 

\medskip
\begin{digress}
Для совместного рассмотрения графа и его дополнения
иногда оказывается полезной следующая метафора.
Раскрасим \em{все} возможные рёбра в два цвета:
скажем, рёбра графа покрасим в красный цыет,
а рёбра дополнения~--- в синий. Тогда
на одной диаграмме (\DMref{п.~7.1.4}{7.1.4.})
можно одновременно отобразить граф и его дополнение. 
\end{digress}   
\end{frame}


\begin{frame}
\label{7.1.4.}\subsubsection{7.1.4. Диаграмма графа}
\frametitle{7.1.4. Диаграмма графа}
%\label{8.1.4}
Обычно граф изображают \odmkey{диаграммой}{Диаграмма
(diagram)!графа (of graph)}: вершины~--- точками (или кружками),
рёбра~--- линиями.
\begin{example}
На рисунке приведен пример диаграммы графа, имеющего четыре вершины
и~пять рёбер. В этом графе
вершины $v_1$ и~$v_2$, $v_2$ и~$v_3$, $v_3$ и~$v_4$, $v_4$ и~$v_1$, $v_2$ и~$v_4$ смежны,
а~вершины $v_1$ и $v_3$ не смежны.
Смежные рёбра: $e_1$ и~$e_2$,
$e_2$ и~$e_3$, $e_3$ и~$e_4$, $e_4$ и~$e_1$,
$e_1$ и~$e_5$, $e_2$ и~$e_5$, $e_3$ и~$e_5$, $e_4$ и~$e_5$.
Несмежные рёбра: $e_1$ и~$e_3$, $e_2$ и~$e_4$.

%\begin{figure}[h]
\begin{center}
\includegraphics[width=45mm]{odm07-04}
\end{center}
%\caption{Диаграмма графа}
%\label{G}
%\end{figure}
\end{example}

\end{frame}

\begin{frame}

\label{7.1.5.}\subsubsection{7.1.5. Орграфы, псевдографы, мультиграфы и~гиперграфы }\frametitle{7.1.5. Орграфы, псевдографы, мультиграфы и~гиперграфы (1/3)}
Часто рассматриваются следующие родственные графам объекты:
%\vspace{-0.2\baselineskip}
%\begin{enumerate}
%\item 

\smallskip
{\color{blue} 1.}
Если элементами множества $E$ являются \em{упорядоченные}
пары (то есть $E\subset V\times V$), то граф называется
\odmkey{ориентированным}{Граф (graph)!ориентированный (directed)}
(или \odmkey{орграфом}{Орграф (digraph)}). В этом случае элементы
множества $V$ называются \odmkey{узлами}{Узел (node)}, а элементы
множества $E$~--- \odmkey{дугами}{Дуга (arc)}.
%{\parfillskip0pt\par}

\smallskip
{\color{blue} 2.}
Если
элементом множества $E$ может быть пара \em{одинаковых} (не
различных) элементов $V$, то такой элемент множества $E$
называется \odmkey{петлёй}{Петля (loop)}, а граф называется
\odmkey{графом с петлями}{Граф (graph)!с петлями (with loops)}
(или \odmkey{псевдографом}{Псевдограф (pseudograph)}).

\smallskip
{\color{blue} 3.}
Если
$E$ является не множеством, а \em{мультимножеством}, содержащим
некоторые элементы по нескольку раз, то эти элементы называются
\odmkey{кратными рёбрами}{Ребро (edge)!кратное (multiple)}, а граф
называется \odmkey{мультиграфом}{Мультиграф (multigraph)}.


\end{frame}

\begin{frame}

\frametitle{7.1.5. Орграфы, псевдографы, мультиграфы и~гиперграфы (2/3)}


{\color{blue} 4.}
Если элементами множества $E$ являются не обязательно
двухэлементные, а~\em{любые\/} (непустые) подмножества множества
$V$, то есть, если $E\subset 2^V$, то такие элементы множества $E$ называются
\odmkey{гиперрёбрами}{Гиперребро (hyperedge)}, а граф называется
\odmkey{гиперграфом}{Гиперграф  (hypergraph)}.

\smallskip
{\color{blue} 5.}
Если каждый элемент множества $E$ действует из множества 
<<начал>> в множество <<концов>>, 
$E\subset 2^V\times 2^V$,
то элементы множества $E$ называются 
\odmkey{гипердугами}{Гипердуга (hyperarc)}, а граф называется 
\odmkey{гиперорграфом}{Гиперорграф (hyperorgraph)}.

\smallskip
{\color{blue} 6.}
Если задана
функция $\Map{F}{V}{M}$ и (или) $\Map{F}{E}{M}$, то множество $M$
называется множеством \odmkey{пометок}{Множество (set)!пометок (of
labels)}, а граф называется \odmkey{помеченным}{Граф
(graph)!помеченный (labeled)} (или \odmkey{нагруженным}{Граф
(graph)!нагруженный (weighted)}). В качестве множества пометок
обычно используются буквы или целые числа. Если функция $F$
инъективна, то~есть разные вершины (рёбра) имеют разные пометки,
то граф называют \odmkey{нумерованным}{Граф (graph)!нумерованный
(numbered)}.
%\end{enumerate}
\end{frame}

\begin{frame}

\frametitle{7.1.5. Орграфы, псевдографы, мультиграфы и~гиперграфы (3/3)}

\begin{remark}
Отношение смежности является в некотором смысле определяющим для
графов и подобных им объектов.
При этом следует учитывать особенности каждого типа объектов.
% В
%орграфе узел $v$ смежен с узлом $u$, если существует дуга
%$(u,v)$. При этом узел $u$  может быть несмежен с узлом $v$.
Отношение смежности в графе симметрично, а в орграфе оно вовсе не
обязано быть симметричным. В графе обычно считают отношение
смежности рефлексивным, то~есть полагают, что вершина смежна сама
с собой. В псевдографе, напротив, вершину не считают смежной с
собой, если у~неё нет петли. В гиперграфе две вершины считаются
смежными, если они принадлежат одному гиперребру. В гиперорграфе
гипердуга обычно проводится из одного узла в множество узлов
(возможно, пустое). В таком случае отношение смежности оказывается
уже не бинарным отношением на $V$, а отношением из $V$ в $2^V$.
Эти и подобные естественные вариации определений обычно считают
ясными из контекста.
\end{remark}

\medskip
{\large
Далее выражение <<граф $G(V,E)$>> означает неориентированный непомеченный граф
без петель и кратных рёбер, с множеством вершин $V$ и множеством рёбер~$E$.
}
\end{frame}

\begin{frame}
\label{7.1.6.}\subsubsection{7.1.6. Изоморфизм графов }\frametitle{7.1.6. Изоморфизм графов (1/6)}
%\vspace{-2pt}
Говорят, что два графа, $G_1(V_1,E_1)$ и
$G_2(V_2,E_2)$, (вершинно)
\odmkey{изоморфны}{Изоморфизм (isomorphism)!графов (of graphs)} (обозначается %\\
$G_1 
\sim G_2$), если существует изоморфизм
$\Map{h}{V_1}{V_2}$, уважающий отношение смежности вершин 
(\DMref{п.~1.7.6}{1.7.6.}):
$
e_1=\divI{(u,v)}{E_1} \Longleftrightarrow e_2=\divI{\bigl(h(u),h(v)\bigr)}{E_2}$.

\begin{theorem}
Изоморфизм графов есть отношение эквивалентности.
\end{theorem}
\begin{proof}
\proofsection{\!Рефлексивность\!} Имеем
 $G\stackrel{h}{\sim} G$, где $h$~---  тождественная функция.\\[-3pt]
\proofsection{\!Симметричность\!} Если $G_1\stackrel{h}{\sim} G_2$,
то $G_2\stackrel{h^{-1}}{\sim} G_1$.\\[-3pt]
\proofsection{\!Транзитивность\!} Если $G_1\stackrel{h}{\sim} G_2$
 и $G_2\stackrel{g}{\sim} G_3$, то $G_1\stackrel{g\circ h}{\sim }G_3$.
\end{proof}


\begin{example}
Три внешне различные диаграммы, приведённые на рисунке,
являются диаграммами одного и того же графа $K_{3,3}$ (полный двудольный граф, 
см. \DMref{п.~7.3.2}{7.3.2.}).

\vspace{-8pt}
\begin{center}
\includegraphics[width=110mm]{odm07-05}
\end{center}
\end{example}
\end{frame}

\begin{frame}
\frametitle{7.1.6. Изоморфизм графов (2/6)}


Числовая характеристика, одинаковая для всех изоморфных графов,
называется \odmkey{инвариантом}{Инвариант (invariant)!графа (of
graph)} графа. Так, $p(G)$ и $q(G)$~--- инварианты графа $G$.
Неизвестно никакого простого набора инвариантов, определяющих граф
с точностью до изоморфизма.

\smallskip
\begin{example}
Количество вершин, рёбер и количество смежных вершин для каждой вершины
не определяют граф даже в простейших случаях! На рисунке
представлены диаграммы графов, у которых указанные инварианты
совпадают, но графы при этом не изоморфны.
%\begin{figure}[h]
\begin{center}
\includegraphics[width=80mm]{odm07-06_new}
\end{center}
%\caption{Диаграммы неизоморфных графов с совпадающими инвариантами}
%\label{nosim}
%\end{figure}


Представленные графы не изоморфны, 
потому что в  графе справа есть два треугольника,
а в графе слева треугольников нет.
\end{example}

\smallskip
Изоморфизм графа в себя называется \odmkey{автоморфизмом}{Автоморфизм (automorphism)!графов (of graphs)} графа.
\begin{example}
Треугольник автоморфен себе при любой перестановке вершин.
\end{example}
\end{frame}

\begin{frame}
\frametitle{7.1.6. Изоморфизм графов (3/6)}

Графы рассматриваются \em{с точностью
до изоморфизма}, то~есть рассматриваются
\em{классы эквивалентности по отношению изоморфизма}.
\begin{remark}
Чтобы подчеркнуть это обстоятельство,
иногда говорят 
\odmkey{абстрактный граф}{Граф (graph)!абстрактный (abstract)}, подразумевая
элемент фактормножества графов по отношению изоморфности.
\end{remark}

Смежными могут быть не только вершины, но и рёбра. 
Аналогично понятию вершинного изоморфизма, 
вводится понятие рёберного изоморфизма. 
Два графа $G_1(V_1, E_1)$ и $G_2(V_2,E_2)$
\odmkey{рёберно изоморфны}{Изоморфизм (isomorphism)!графов (of graphs)!рёберный (line)}, если существует изоморфизм $\Map{h}{E_1}{E_2}$, уважающий отношение смежности рёбер.

\begin{example}
Рассмотрим графы 
$G_1(V_1, E_1)$ и $G_2(V_2,E_2)$, где 
$V_1=V_2=\{v_1,v_2,v_3\}$, \\$E_1=\{(v_1,v_2),
(v_1,v_3)\}$, $E_2=\{(v_1,v_2),(v_2,v_3)\}$. По определению эти графы являются различными помеченными графами, так как $E_1\neq E_2$. В тоже время функция $h$, 
такая что $h(v_1)\Assign v_2,
h(v_2)\Assign v_1, h(v_3)\Assign v_3$ устанавливает изоморфизм этих графов.
\end{example}

\begin{remark}
Абстрактный граф, который является классом простых цепей
(см. \DMref{п.~7.2.3}{7.2.3.}) на $n$ вершинах, обозначается $P_n$. Графы $G_1, G_2$ в примере суть представители абстрактного графа $P_3$.
\end{remark}
%\smallskip
%{\color{red}
%Тема для сочинения: методы проверки изоморфизма.
%Очень трудный, но очень ценный параграф. 
%Гипотеза Улама. Труднорешаемая задача. 
%Быстрое установление отсутствия изоморфизма. 
%Проблема изоморфизма подграфов. Частные случаи.
%Приближенные алгоритмы. Примеры приложений.
%}
\end{frame}

\begin{frame}
\frametitle{7.1.6. Изоморфизм графов (4/6)}
В этом курсе слово <<граф>> означает как 
индивидуальный помеченный граф, так и представителя
класса изоморфных графов, то есть представителя 
абстрактного графа. Какая именно интерпретация подразумевается, надобно определять из контекста.


\begin{lemma} 
Пусть графы $G$ и $H$ рёберно изоморфны
 с изоморфизмом $h$, 
тогда если для рёбер $\Tuple{e}{1}{n+1}$,
где $n>2$, образы $h(e_1), \ldots, h(e_n)$
 инцидентны вершине $v_1$, а образы
 $h(e_2), \ldots, h(e_{n+1})$ инцидентны
 вершине $v_2$, то $v_1=v_2$. 
\end{lemma}
\begin{proof}
Поскольку две вершины полностью определяют ребро, то в случае, 
если $v_1 \neq v_2$, получаем 
$h(e_2)=\ldots=h(e_n)=(v_1, v_2) $, что противоречит биективности $h$.
\end{proof}
\begin{theorem}
Если графы вершинно изоморфны,
то они и рёберно изоморфны.
\end{theorem}

\begin{proof}
Пусть графы $G_1(V_1,E_1)$ и $G_2(V_2,E_2)$ вершинно изоморфны
 с изоморфизмом $\Map{h_1}{V_1}{V_2}$. 
Тогда построим функцию $\Map{h_2}{E_1}{E_2}$
следующим образом:
$h_2\left((u,v)\right)\Assign\left(h_1(u),h_1(v)\right)$.
Эта функция биективна и уважает смежность рёбер.  
\end{proof}

%\vspace{-12pt}
\begin{columns}[t]
\begin{column}{8cm}
Обратное, вообще говоря, неверно.

\begin{example}
На рисунке представлены 
диаграммы двух графов,
которые изоморфны рёберно,
но не изоморфны вершинно.
\end{example}
\end{column}
\begin{column}{7cm}

\vspace{-12pt}
\begin{center}
\includegraphics[height=16mm]{7_1_6_3.jpg}
\end{center}
\end{column}
\end{columns}

\end{frame}

\begin{frame}
\frametitle{7.1.6. Изоморфизм графов (5/6)}

\begin{digress}
Изоморфизм графов~--- это очень трудная, и в то
же время весьма насущная проблема.
Даже установить отсутствие или наличие 
изоморфизма у двух помеченных графов с совпадающими
инвариантами не всегда просто,
как видно из предшествующих примеров. 

\smallskip
На практике очень часто проблема изоморфизма возникает в более сложной постановке: 
даны два абстрактных графа, требуется определить,
имеются ли вхождения одного графа в другой
в качестве подграфа 
(как правильного, так и неправильного,
см.~\DMref{п.~7.2.1}{7.2.1.}), и если такие вхождения
имеются, то найти их (какое-нибудь одно или все).
Обычно эта проблема называется 
\em{изоморфизм подграфов}.

\smallskip
К решению проблемы изоморфизма подграфов
сводятся самые разные задачи в области химии и биоинформатики (например, входят ли в молекулу полимера те или иные характерные составные части), в области шифрования и кодирования информации (содержит ли некоторая 
информационная структура те или иные 
подструктуры) и в множестве других 
областей. Фактически, проблема 
изоморфизма подграфов встаёт везде,
где в качестве языка описания 
модели предметной области применяется теория 
графов (а язык теории графов~--- самый распространённый формализм дискретной математики).
\end{digress}
\end{frame}

\begin{frame}
\frametitle{7.1.6. Изоморфизм графов (6/6)}
\begin{digress} (Продолжение)
В такой самой общей постановке
проблема изоморфизма подграфов
оказывается трудоёмкой задачей и
решается, фактически, перебором.
В практически значимых моделях
(в биоинформатике, в криптографии)
 вершины в графах
исчисляются тысячами и миллионами,
поэтому проблема изоморфизма подграфов
в общей постановке
оказывается практически неразрешимой даже с применением
суперкомпьютеров.
Однако в каждой конкретной предметной области
удаётся отыскать такие особенности,
которые позволяют ограничить постановку 
задачи проверки изоморфизма так,
что задача становится практически разрешимой.
Эти специфические особенности и приёмы не
имеют универсальной силы и неприменимы в 
общем случае.

\smallskip
Поэтому в этом курсе нет описания
алгоритма проверки изоморфизма графов в общем случае~---
такой алгоритм не имеет практического смысла.
Зато в курсе имеется несколько
замечаний относительно того,
как можно эффективно проверить изоморфизм
в том или ином сугубо частном случае.
Комбинируя эти идеи, 
иногда удаётся построить 
эффективный алгоритм для решения конкретной 
частной задачи.  
\end{digress}

\end{frame}

\begin{frame}
\frametitle{Онтология теории графов}
\begin{center}
\includegraphics[height=7.8cm]{ODM7_1.png}
\end{center}

\end{frame}



\subsection{7.2. Элементы графов}\label{7.2.}
\begin{frame}
\frametitle{7.2. Элементы графов}
После рассмотрения определений, относящихся к графам
как к цельным объектам, естественно дать определения
различным элементам графов.

\begin{tabular}{p{9mm}|p{56mm}|p{79mm}}
  № & Название параграфа
  & Ключевые термины и обозначения \\
  \hline
\DMref{7.2.1.}{7.2.1.} & Подграфы 
  & 
  Остовные, собственные, правильные \\
\DMref{7.2.2.}{7.2.2.}\RS & Валентность 
  & 
  {$d(v)$~--- степень; 
  $\delta(G)$~--- минимальная степень;
  $\Delta(G)$~--- максимальная степень;
  $r(G)$~--- степень регулярности;
  $d^+(v)$~--- полустепень захода;
  $d^-(v)$~--- полустепень исхода;
  изолированные и висячие вершины}
  \\
\DMref{7.2.3.}{7.2.3.} & Маршруты, цепи, циклы &
{$\divC{u}{v}$~--- цепь; $\lrle{\stackrel{\rightarrow}{u,v}}$~--- путь; 
$z(G)$~--- число циклов; $z(G,k)$~--- число циклов
длины $k$}\\
\DMref{7.2.4.}{7.2.4.} & Связные и несвязные графы &
{$k(G)$~--- число компонент связности
}
\\
\DMref{7.2.5.}{7.2.5.} & Расстояние между вершинами, ярусы, эксцентриситет, диаметр, радиус, центр &
{$d(u,v)$~--- расстояние;
$D(v, n)$~--- ярус;

$e(v)$~--- эксцентриситет;
$D(G)$~--- диаметр; 
$R(G)$~--- радиус; $C(G)$~--- центр}\\   
\DMref{7.2.6.}{7.2.6.}\BS & Степенные последовательности &
{Откладывание вершины}\\   
\end{tabular}
\end{frame}

\begin{frame}
\label{7.2.1.}\subsubsection{7.2.1. Подграфы}
\frametitle{7.2.1. Подграфы}
Граф $G'(V',E')$ называется \odmkey{подграфом}{Подграф (subgraph)}
(или \odmkey{частью}{Часть графа (section of graph)})
 графа
$G(V,E)$  (обозначается \\$G'\subset G$), если $V'\subset V \And
E'\subset E$. Если $V'=V$, то $G'$ называется \odmkey{остовным
подграфом}{Подграф (subgraph)!остовный (spanning)} $G$. Если
$V'\subset V \And E'\subset E \And (V'\not= V\Or E'\not= E)$, то
граф $G'$ называется \odmkey{собственным}{Подграф
 (subgraph)!собственный (proper)} подграфом
графа $G$.
Подграф $G'(V',E')$ называется \odmkey{правильным}{Подграф
 (subgraph)!правильный (regular)}
подграфом графа $G(V,E)$, если $G'$ содержит все возможные рёбра $G$:
$\A{\divI{u,v}{V'}}{\divI{(u,v)}{E}\Impl \divI{(u,v)}{E'}}$.

Правильный подграф $G'(V',E')$ графа $G(V,E)$ определяется подмножеством
вершин $V'$.

\begin{example}
На рисунке показаны диаграммы графа, правильного подграфа, неправильного подграфа и остовного подграфа.

\vspace{-6pt} 
\begin{center}
\includegraphics[height=18mm]{7_2_1.jpg}
\end{center}
\end{example}

\vspace{-6pt}
Другими словами, правильный подграф получается из
исходного удалением вершин вместе с инцидентными рёбрами,
а остовный подграф получается удалением рёбер.
\begin{remark}
Иногда подграфами называют только правильные подграфы, 
а неправильные подграфы называют
\odmkey{изграфами}{Изграф (izgraph)}.
\end{remark}
\end{frame}

\begin{frame}
\label{7.2.2.}\subsubsection{7.2.2. Валентность }\frametitle{7.2.2. Валентность (1/3)}
Количество рёбер, инцидентных вершине $v$, называется
\odmkey{степенью}{Степень (degree)!вершины (of vertex)} (или
\odmkey{валентностью}{Валентность вершины (vertex valence)})
вершины $v$ и обозначается $d(v)$:
$$
\A{\divI{v}{V}}{0\leq d(v) \leq p-1}, \qquad
d(v) = \vert\Gamma^+(v)\vert.
$$
%
Таким образом, степень $d(v)$ вершины $v$ совпадает с количеством
смежных с ней вершин. Количество вершин, не смежных с $v$,
обозначают $\overline{d}(v)$. Ясно, что
$$\A{v\in V}{d(v)+\overline{d}(v)=p-1}.$$

Обозначим \odmkey{минимальную}{Степень (degree)!вершины (of
vertex)!минимальная (minimal)} степень вершины графа $G$ через
$\delta(G)$, а \odmkey{максимальную}{Степень (degree)!вершины (of
vertex)!максимальная (maximal)}~--- через $\Delta(G)$:
$$
\delta\left(G(V,E)\right)\Def \min_{\divI{v}{V}}d(v), \qquad
\Delta \left(G(V,E)\right)\Def \max_{\divI{v}{V}}d(v).
$$
Ясно, что $\delta(G)$ и $\Delta(G)$ являются инвариантами. Если
степени всех вершин равны $k$, то граф называется
\odmkey{регулярным}{Граф (graph)!регулярный (regular)}
степени~$k$:
$$
\delta(G) = \Delta(G) = k, \quad \A{v\in V}{d(v)=k}.
$$
Степень регулярности обозначается
$r(G)$. Для нерегулярных графов $r(G)$ не определено.
\end{frame}

\begin{frame}
\frametitle{7.2.2. Валентность (2/3)}
\begin{examples}
На рисунке приведены диаграммы трёх примечательных регулярных графов степени 3.
%На рис.~\ref{nosim} приведены диаграммы двух регулярных, но
%неизоморфных графов степени 3.
\begin{figure}[h]
\begin{center}
\includegraphics[width=100mm]{7_2_2.jpg}
\end{center}
%\caption{Диаграмма регулярного графа степени 3}
%\label{r3}
\end{figure}
\end{examples}
Если степень вершины равна нулю (то~есть $d(v) = 0$), то вершина
называется \odmkey{изолированной}{Вершина (vertex)!изолированная
(isolated)}. Если степень вершины равна единице (то~есть $d(v) =
1$), то вершина называется \odmkey{концевой}{Вершина
(vertex)!концевая (leaf)}, или \odmkey{висячей}{Вершина
(vertex)!висячая (pendant)}.

\vspace{\baselineskip}
Для орграфа число дуг, исходящих из
узла $v$, называется \odmkey{полустепенью
исхода}{Полустепень исхода (out-degree)}, а число входящих~---
\odmkey{полустепенью захода}{Полустепень захода (in-degree)}.
Обозначаются эти числа, соответственно,  $d^-(v)$ и $d^+(v)$.
\end{frame}

\begin{frame}
\frametitle{7.2.2. Валентность (3/3)}
\begin{theorem} \NB{[Лемма о рукопожатиях]}
Сумма степеней вершин графа (мультиграфа) равна удвоенному количеству рёбер:
\vspace{-0.2\baselineskip}
$$\sum_{\divI{v}{V}} d(v) = 2q.
$$
\end{theorem}

\vspace{-2pt}
\begin{proof}
При подсчёте суммы степеней вершин каждое ребро учитывается два
раза.
\end{proof}
%\setcounter{corollary}{0}
\begin{corollary}
Число вершин нечётной степени чётно.
\end{corollary}
\begin{proof}
Сумма степеней вершин чётной степени чётна, значит, сумма степеней
вершин нечётной степени также чётна, следовательно, их чётное число.
\end{proof}
\begin{corollary}
Суммы полустепеней узлов орграфа равны  количеству дуг:
\vspace{-0.25\baselineskip}
$$
\sum_{\divI{v}{V}}d^{-}(v) = q = \sum_{\divI{v}{V}}d^{+}(v) .
$$
\end{corollary}
\begin{proof}
Сумма полустепеней узлов орграфа равна сумме степеней вершин графа
(мультиграфа), полученного из орграфа забыванием ориентации
дуг.
\end{proof}
\end{frame}

\begin{frame}
\label{7.2.3.}\subsubsection{7.2.3. Маршруты, цепи, циклы }\frametitle{7.2.3. Маршруты, цепи, циклы (1/3)}
\odmkey{Маршрутом}{Маршрут (walk)} в графе называется чередующаяся
последовательность вершин и~рё\-бер, начинающаяся и кончающаяся вершиной,
$
v_0,e_1,v_1,e_2,v_2,\dots,e_k,v_k,
$
в которой любые два соседних элемента инцидентны,
причём однородные элементы (вершины, рёбра)
через один смежны или совпадают.
\begin{remark} Это определение подходит также для псевдо-, мульти- и
орграфов. При этом в графе (орграфе) достаточно указать
только последовательность вершин (узлов)
или только последовательность рёбер (дуг).
\end{remark}

\smallskip
Если $v_0=v_k$, то маршрут \odmkey{замкнут}{Маршрут (walk)!замкнутый (closed)},
иначе~--- \odmkey{открыт}{Маршрут (walk)!открытый (open)}.

\smallskip
Если все рёбра различны, то маршрут называется
\odmkey{цепью}{Цепь (circuit)}.
Если все вершины (а значит, и рёбра!) различны,
то маршрут называется \odmkey{простой}{Цепь (circuit)!простая (simple)} цепью.
В цепи $v_0, e_1,\dots,e_k,v_k$ вершины $v_0$ и $v_k$ называются
\odmkey{концами}{Конец цепи (circuit end)} цепи. Цепь с концами
$u$ и $v$ \em{соединяет} вершины $u$ и $v$ 
и обозначается $\divC{u}{v}$. Если нужно указать граф $G$,
которому принадлежит цепь, то добавляют индекс: ${\divC{u}{v}}_G$.


\end{frame}

\begin{frame}
\frametitle{7.2.3. Маршруты, цепи, циклы (2/3)}
\begin{lemma}\label{723}
Если есть какая-либо цепь, соединяющая вершины $u$ и $v$, то есть
и простая цепь.
\end{lemma}
\begin{proof}
Если имеющаяся цепь не простая, то будем удалять часть цепи между повторяющимися вхождениями одной вершины, получая новую цепь, соединяющую вершины $u$ и $v$, с меньшим количеством вершин,
пока цепь не станет простой.    
\end{proof}

Замкнутая цепь называется \odmkey{циклом}{Цикл (cycle)}; замкнутая простая цепь
называется \odmkey{простым циклом}{Цикл (cycle)!простой (simple)}.
Число циклов в графе $G$ обозначается $z(G)$. Граф без циклов называется
\odmkey{ациклическим}{Граф (graph)!ациклический (acyclic)}.
Число рёбер в цикле называется \odmkey{длиной цикла}{Длина (length)!цикла (of cycle)}.
Число циклов длины $k$ в графе $G$  обозначается  $z(G, k)$.

\begin{lemma}
Количество простых циклов длины $l$ одинаково для изоморфных графов.
\end{lemma}

\begin{proof}
Пусть графы $G_1(V_1,E_1)$ и $G_2(V_2,E_2)$ 
изоморфны с биекцией $\Map{h}{V_1}{V_2}$, 
сохраняющей смежность.
Количество простых циклов длины $l$ в $G_2$ не меньше, чем в $G_1$:
$ z(G_1,l)\le z(G_2,l)$. От противного. 
Пусть  $v_1,v_2,\dots,v_l$~--- простой цикл графа $G_1$, 
а $h(v_1),h(v_2),\dots,h(v_l)$ не является простым циклом графа $G_2$. 
Значит, найдётся такая пара смежных вершин $v_i$, $v_{i+1}$, 
для которой $(h(v_i ),h(v_{i+1}))\not\in E_2$, 
что противоречит тому, что $h$ сохраняет смежность.
Аналогично, количество простых циклов длины $l$ в $G_2$ не больше, чем в $G_1$: достаточно рассмотреть $\Map{h^{-1}}{V_2}{V_1}$.
\end{proof}
\end{frame}

\begin{frame}
\frametitle{7.2.3. Маршруты, цепи, циклы (3/3)}
\begin{remark}
Для псевдографов обычно особо
оговаривают, считаются ли петли циклами.
\end{remark}
Для орграфов цепь называется \odmkey{путём}{Путь (path)}, а цикл~---
\odmkey{контуром}{Контур (contour)}. Путь в орграфе из узла $u$ в узел $v$
обозначают $\lrle{\stackrel{\rightarrow}{u,v}}$.
\begin{example}
В графе, диаграмма которого приведена на рисунке:
\begin{enumerate}
\item[1)] $v_1$, $v_3$, $v_1$, $v_4$~--- маршрут, но не цепь;%
\item[2)] $v_1$, $v_3$, $v_5$, $v_2$, $v_3$, $v_4$~--- цепь, но не
простая цепь;%
\item[3)] $v_1$, $v_4$, $v_3$, $v_2$, $v_5$~--- простая цепь;%
\item[4)] $v_1$, $v_3$, $v_5$, $v_2$, $v_3$, $v_4$, $v_1$~--- цикл, но не простой цикл;%
\item[5)] $v_1$, $v_3$, $v_4$, $v_1$~--- простой цикл.%
\end{enumerate}
\end{example}

\vspace{-4pt}
\begin{figure}[h]
\begin{center}
\includegraphics[width=45mm]{odm07-08}
\end{center}
%\caption{Маршруты, цепи, циклы} \label{cicles}
\end{figure}
\end{frame}


\begin{frame}
\label{7.2.4.}\subsubsection{7.2.4. Связные и несвязные графы }\frametitle{7.2.4. Связные и несвязные графы (1/2)}
Говорят, что две вершины в графе \odmkey{связаны}{Вершина
(vertex)!связанная (bound)}, если существует соединяющая их
цепь. Граф, в котором все вершины связаны, называется
\odmkey{связным}{Граф (graph)!связный (connected)}.

\begin{remark}
Если две вершины связаны, то по первой лемме 
\DMref{п.~7.2.3}{7.2.3.} существует соединяющая их простая цепь. 
\end{remark}
 
\medskip
Нетрудно показать, что отношение связанности вершин
является эквивалентностью. 

\medskip
Классы эквивалентности по отношению
связанности называются \odmkey{компонентами связности}{Компонента
(component)!связности
 (of connectivity)} графа. Число
компонент связности графа $G$ обозначается $k(G)$. Граф $G$
является связным тогда и только тогда, когда $k(G)=1$. Если
$k(G)>1$, то $G$~--- \odmkey{несвязный}{Граф (graph)!несвязный
(disconnected)} граф. 

\medskip
Граф, состоящий только из изолированных
вершин в~котором $k(G)=p(G)$ и $q(G)=0$, иногда называется
\odmkey{вполне несвязным}{Граф (graph)!вполне несвязный
 (totally disconnected)}.
 
\begin{remark}
Вполне несвязный граф $G$ является регулярным нулевой степени, $r(G)=0$.
\end{remark} 
\end{frame}


\begin{frame}
\frametitle{7.2.4. Связные и несвязные графы (2/2)}
\begin{theorem}
Любой граф или его дополнение являются связным графом. 
\end{theorem}

\begin{proof} 
Если исходный граф связен, то теорема доказана. 
Пусть исходный граф несвязен и имеет несколько компонент связности. 
Рассмотрим любые две вершины. 
Если они принадлежат разным компонентам связности исходного графа, 
то они соединены в дополнении ребром, 
а потому связаны.
Если они принадлежат одной компоненте связности,
то в дополнении они связаны маршрутом из двух рёбер, 
проходящим через любую вершину из другой компоненты связности.
\end{proof}

\medskip
\begin{remark}
Как сам граф, так и его дополнение
могут быть связными одновременно,
но граф и его дополнение не могут быть 
несвязными одновременно!
\end{remark}

\end{frame}

\begin{frame}
\label{7.2.5.}\subsubsection{7.2.5. Расстояние между вершинами, ярусы, эксцентриситет, диаметр, радиус, центр    }\frametitle{7.2.5. Расстояние между вершинами, ярусы, эксцентриситет, диаметр, радиус, центр    (1/5)}
\odmkey{Длиной маршрута}{Длина (length)!маршрута (of walk)}
называется количество рёбер в нём (с учётом повторений).
Если~маршрут $M=v_0,e_1,v_1,e_2,v_2,\dots, e_k,v_k$, то длина $M$
равна $k$ (обозначается $\vert M\vert=k$).

\odmkey{Расстоянием}{Расстояние (distance)} между вершинами $u$ и
$v$ (обозначается $d(u,v)$) называется длина кратчайшей цепи
$\divC{u}{v}$, а сама кратчайшая цепь называется
\odmkey{геодезической}{Цепь (line)!Геодезическая (geodesic)},
$$
d(u,v)\Def \min\limits_{\{\langle u,v\rangle\}} \vert\langle u,v \rangle\vert.
$$

Если для любых двух вершин графа существует единственная
геодезическая цепь, то граф называется \odmkey{геодезическим}{Граф
(graph)!геодезический (geodesic)}. 
%\vspace{-0.25\baselineskip}

\begin{remark} Если вершины $u$ и
$v$ не связаны, то $d(u,v)\Def +\infty$.
\end{remark}

\medskip
Легко видеть, что расстояние между вершинами является метрикой (\DMref{п.~6.3.3}{6.3.3.})
в пространстве вершин графа.
Действительно, $\A{u,v\in V}{d(u,v)\ge 0}$,
так как $d(u,v)$~--- это количество рёбер в
кратчайшей цепи.
\end{frame}

\begin{frame}
\frametitle{7.2.5. Расстояние между вершинами, ярусы, эксцентриситет, диаметр, радиус, центр   (2/5)}

Проверим выполнение аксиом метрики.

\begin{enumerate}
\item $\A{u,v \in V}{d(u,v)=0\Equiv v=u}$, поскольку 
из нуля рёбер можно составить маршрут, 
который состоит только из одной вершины.
\item $\A{u,v \in V}{d(u,v)=d(v,u)}$, поскольку 
пересчёт рёбер цепи в обратном порядке не меняет их
количества.
\item $\A{u,v,w \in V}{d(u,v) \le 
d(u,w)+ d(w,v)}$, поскольку
$d(u,w)+ d(w,v)$~--- это длина некоторой
 цепи $\langle u,\ldots,w,\ldots,v\rangle$,
а $d(u,v)$~--- это длина кратчайшей
 цепи $\langle u,v\rangle$.   
\end{enumerate}

\smallskip
Множество вершин, находящихся на
заданном расстоянии $n$ от вершины $v$ (обозначение $D(v,n)$),
называется \odmkey{ярусом}{Ярус (tier)}:

\vspace{-8pt}
$$
D(v,n)\Def\Set{u\in V}{d(v,u)=n}.
$$

Ясно, что множество вершин $V$ всякого связного графа однозначно
разбивается на ярусы относительно данной вершины.


\end{frame}

\begin{frame}
\frametitle{7.2.5. Расстояние между вершинами, ярусы, эксцентриситет, диаметр, радиус, центр   (3/5)}

\odmkey{Эксцентриситетом}{Эксцентриситет (eccentricity)!вершины
(of vertex)}
 $e(v)$
вершины $v$ в связном графе $G(V,E)$ называется
максимальное расстояние от вершины $v$ до
других вершин графа $G$:

%\vspace{-8pt}
$$
e(v)\Def \max_{u\in V} d(v,u).
$$

\odmkey{Диаметром}{Диаметр графа (graph diameter)}  графа 
 называется длиннейшая
геодезическая цепь. Длина диаметра графа $G$ обозначается $D(G)$:

\vspace{-8pt}
$$
D(G)\Def \max\limits_{u,v\in V} d(u,v)
=\max\limits_{v\in V} e(v).
$$

%\vspace{-2pt}

%\vspace{-2pt}
Заметим, что наиболее эксцентричные вершины~--- это концы
диаметра. 

\smallskip
\odmkey{Радиусом}{Радиус графа (radius of graph)}
графа называется простая цепь, длина которой равна наименьшему из эксцентриситетов вершин. Длина радиуса графа $G$ обозначается $R(G)$:
$$
R(G) \Def \min_{v\in V} e(v).
$$

\end{frame}



\begin{frame}
\frametitle{7.2.5. Расстояние между вершинами, ярусы, эксцентриситет, диаметр, радиус, центр    (4/5)}
Вершина $v$
называется \odmkey{центральной}{Вершина (vertex)!центральная
(central)}, если её эксцентриситет совпадает с длиной радиуса графа,
$e(v)=R(G)$. Множество центральных вершин $C(G)$ называется
\odmkey{центром}{Центр (center)!графа (of graph)} графа:
$$
C(G)\Def \Set{v\in V}{e(v)=R(G)}.
$$

\begin{example}
На диаграммах ниже указаны эксцентриситеты вершин и центры двух графов.
Вершины, составляющие центр, на диаграммах выделены жирными точками.
\end{example}

%\vspace{-12pt}
\begin{figure}[h]
\begin{center}
\includegraphics[width=120mm]{07_09}
\end{center}
%\caption{Эксцентриситеты вершин и центры графов}
\label{f10-5}
\end{figure}
\end{frame}



\begin{frame}
\frametitle{7.2.5. Расстояние между вершинами, ярусы, эксцентриситет, диаметр, радиус, центр    (5/5)}

\begin{theorem}
Если граф $G$ связен, то $D(G)/2\le R(G)\le D(G)$.
\end{theorem}
%
\begin{proof}
\proofsection{$R(G)\le D(G)$} Следует из определения.
\proofsection{$D(G)/2\le R(G)$}
Рассмотрим центральную вершину $w\in C(G)$, 
$e(w)=R(G)$.
Имеем 
$\A{u,v \in V}{d(u,v) \le 
d(u,w)+ d(w,v)}$ по неравенству треугольника для метрики
$d(u,v)$.
Но $d(u,w)=d(w,u)\le e(w)=R(G)$ и
$d(w,v)\le e(w)=R(G)$, значит \\
$\A{u,v \in V}{d(u,v) \le 2R(G)}$,
откуда $D(G)\le 2R(G)$.   
\end{proof}

\begin{example}
В левой диаграмме на предыдущем слайде имеем
$D(G)=4, R(G)=2$.
В правой диаграмме
$D(G)=3, R(G)=2$, неравенства выполняются.
\end{example}
\end{frame}

\begin{frame}
\label{7.2.6.}\subsubsection{7.2.6. Степенные последовательности }
\frametitle{7.2.6. Степенные последовательности (1/4)}
Последовательность степеней вершин 
$(\Tuple{d}{1}{p})$ 
некоторого графа называется \odmkey{степенной последовательностью графа}{Последовательность (sequence)!степенная (of degrees)} 
или \odmkey{графовой последовательностью}{Последовательность (sequence)!графовая (of graph)}. 
Ясно, что всякая степенная последовательность является представлением
мультимножества, поэтому, не ограничивая общности, 
можно считать, что
$d_1\le d_2\le \dots\le d_p$.

\begin{example}
Последовательности $(2,2,2,2,3,3,4)$ и $(1,1,1,1,2,4)$ являются степенными последовательностями графов, 
диаграммы которых представлены на предыдущем рисунке.  
\end{example} 

\medskip
Не всякая последовательность неотрицательных целых чисел является степенной последовательностью.

\begin{example}
Последовательность $(0,1,2,\dots, p-1)$ не является степенной ни 
при каком \\$p>1$. Действительно, первая вершина должна быть изолированной, а последняя должна быть смежна со всеми остальными, что противоречит изолированности первой вершины.   
\end{example}


\end{frame}

\begin{frame}
\frametitle{7.2.6. Степенные последовательности (2/4)}

Если последовательность $D=(\Tuple{d}{1}{p})$ является степенной последовательностью 
графа $G$, то говорят, что граф $G$ \em{реализует} последовательность $D$. 


Легко заметить, что если два графа имеют различные степенные последовательности, 
то они не изоморфны. Но из рисунка \DMref{п.~7.1.6}{7.1.6.}. видно, что графы с одинаковыми степенными последовательностями также могут быть не изоморфными. Таким образом, степенная последовательность может иметь множество 
\em{различных} реализаций.

\smallskip
Отметим несколько очевидных свойств 
степенных последовательностей. 

\begin{lemma} \NB{[1]}
Сумма степенной последовательности~--- чётное число.
\end{lemma}
\begin{proof}
По лемме о рукопожатиях
$\sum_{i=1}^p d_i =2q$. 
\end{proof}

\smallskip
\begin{lemma} \NB{[2]}
Если последовательность $(\Tuple{d}{1}{p})$ степенная,
то последовательность \\$(0, \dots, 0, \Tuple{d}{1}{p})$ 
также степенная, и наоборот. 
\end{lemma}
\begin{proof}
Изолированные вершины можно не учитывать.
\end{proof}

\begin{lemma} \NB{[3]}
В степенной последовательности $(\Tuple{d}{1}{p})$ 
каждый элемент меньше длины последовательности:  
$\A{i\in 1..p}{d_i<p}$. 
\end{lemma}
\begin{proof}
Степень любой вершины меньше числа вершин.
\end{proof}

\end{frame}

\begin{frame}
\frametitle{7.2.6. Степенные последовательности (3/4)}

Таким образом, не умаляя общности, достаточно рассматривать только
последовательности $(\Tuple{d}{1}{p})$, 
такие, что $$0< d_1\le d_2\le \dots\le d_p < p.$$ 

Пусть $D=(\Tuple{d}{1}{p})$~---  неубывающая последовательность 
натуральных чисел, причём $\A{i}{d_i<p}$. 
Построим последовательность  
$D^\prime=(\Tuple{d^\prime}{1}{p-1})$ следующим образом:
\begin{align*}
d^\prime_{p-1} & \Assign d_{p-1}-1,\ \dots, \
d^\prime_{p-d_p+1}\Assign d_{p-d_p+1}-1, \\
d^\prime_{p-d_p} & \Assign d_{p-d_p}-1, \ 
d^\prime_{p-d_p-1}\Assign d_{p-d_p-1},
 \ \dots, \ 
d^\prime_1\Assign d_1.
\end{align*}

Другими словами, последний (наибольший) элемент последовательности $d_p$ 
удаляется, а $d_p$ предшествующих элементов уменьшаются на единицу. 
Оставшиеся вершины (они есть, если $d_p < p-1$) 
остаются без изменения.

\medskip
Такая операция называется   
\odmkey{откладыванием}{Вершина (vertex)!откладывание  (laying off)} вершины $d_p$, а $D^\prime$
называется \odmkey{остаточной последовательностью}{Последовательность (sequence)!остаточная (residual)}.
\end{frame}

\begin{frame}
\frametitle{7.2.6. Степенные последовательности (4/4)}

\begin{example}
Отложим последний элемент в последовательности: 
$$(2, 2, 3, 3, 3\, |\, 3)\mapsto (2, 2, 2, 2, 2).$$ При откладывании элемента упорядоченность может
нарушиться: $$(2, 2, 2, 2\, |\, 2)\mapsto (2, 2, 1, 1).$$
\end{example}

Графически откладывание вершины можно отобразить, соединяя рёбрами вершину, соответствующую удаляемому элементу, с вершинами, соответствующими 
уменьшаемым элементам.  
Фактически, при откладывании вершин постепенно прорисовывается 
диаграмма графа, начиная с вершины с максимальной степенью.



\begin{example}
На рисунке показано откладывание вершин для последовательности $(2,3,3,3,3)$. Откладываемые вершины выделены.
\begin{center}
\includegraphics[height=20mm]{7_2_7_1.jpg}
\end{center}
\end{example}

\end{frame}

\begin{frame}
\frametitle{7.2.7. Критерий степенной последовательности (1/6)}

\begin{theorem}
Последовательность $D=(\Tuple{d}{1}{p})$ является степенной последовательностью
 тогда и только тогда, когда последовательность $D^\prime=(\Tuple{d^\prime}{1}{p-1})$,
полученная откладыванием последнего элемента $d_p$, является степенной.
\end{theorem}
\begin{proof}
\proofsection{$\Rightarrow$} Пусть последовательность $D=(\Tuple{d}{1}{p})$
степенная, а $G$~--- реализующий граф. Рассмотрим в графе $G$ вершину  максимальной
степени $w$, $d(w)=d_p$.
Пусть $S$~--- множество вершин со степенями 
$d_{p-d_p},\dots,d_{p-1}$. 
Тогда если $\Gamma(w)=S$, то $G-w$~--- искомый граф, реализующий
последовательность $D^\prime=(\Tuple{d^\prime}{1}{p-1})$.
Пусть теперь $\Gamma(w)\ne S$, 
построим тогда такой граф $G^\prime$,
 в котором $\Gamma(w)=S$. 
Поскольку $\Gamma(w)\ne S$, существует вершина $z\in \Gamma(w)$,
которая не принадлежит $S$,  и существует вершина $x\in S$,
не смежная с вершиной $w$. 
По построению $S$ имеем  $d(z)\le d(x)$,
следовательно, существует
вершина $y$ смежная с $x$, но не смежная с $z$. 
Удалим рёбра $(w,z)$ и $(y,x)$, добавим рёбра $(w,x)$ и $(y,z)$. 
\end{proof}

\vspace{-6pt}
\begin{center}
\includegraphics[height=20mm]{7_2_77_1.jpg}
\end{center} 

\vspace{-6pt}
\end{frame}

\begin{frame}
\frametitle{7.2.7. Критерий степенной последовательности (2/6)}
\begin{proof} (Продолжение)
Такими построениями получаем граф $G^\prime$ такой, что \\$\Gamma(w)=S$, 
причём граф $G^\prime -w$ реализует 
последовательность $D^\prime=(\Tuple{d^\prime}{1}{p-1})$. 							

\vspace{-12pt}
\begin{center}
\includegraphics[height=18mm]{7_2_7_2.jpg}
\end{center} 

\vspace{-6pt}
 

 \proofsection{$\Leftarrow$} Пусть последовательность 
 $D^\prime=(\Tuple{d^\prime}{1}{p-1})$
степенная, а $G$~--- реализующий граф, вершины 
которого перенумерованы в порядке неубывания степеней
$1,\dots,p-1$. Добавим в граф $G$ вершину 
степени $d_p$ и соединим её с вершинами с номерами $p-d_p .. p-1$.
 Полученный граф реализует последовательность $D=(\Tuple{d}{1}{p})$.
\end{proof}

\medskip
Доказанная теорема позволяет проверить, 
является ли заданная последовательность натуральных
чисел степенной и построить граф, реализующий
эту последовательность, с помощью алгоритма, приведённого на следующих трёх слайдах.
\end{frame}

\begin{frame}
\frametitle{7.2.7. Критерий степенной последовательности (3/6)}

  
\vspace{-2pt}

%Алгоритм построения графа по степенной последовательности.

\begin{algorithmic} 
\REQUIRE{ натуральное число $p$ и 
двусвязный упорядоченный список пар натуральных чисел $(d,n)$,
представленный двумя указателями: на начало списка $s$ 
и на конец списка $t$.}  
\ENSURE{ \textbf{false}, 
если последовательность $D$ не степенная, 
иначе \textbf{true} и граф \textbf{$G:$ array $[1..q]$ of record $b,e: 1..p$} (см. 
\DMref{п.~7.4.5}{7.4.5.})}
\STATE{$q\Assign 0$}
\COMMENT{\color{gray}\small счётчик рёбер }\STATE{\textbf{if $t.d\ge p$ then return false end if}}
\COMMENT{\color{gray}\small $D$ не степенная }
\FOR{\textbf{$i$ from $1$ to $p$}}
   \STATE{$d\Assign t.d;\ n\Assign t.n$}
   \COMMENT{\color{gray}\small  $d$~--- степень, $n$~--- номер вершины }
   \STATE{\textbf{if $s.d<0$ then return false end if}}
   \COMMENT{\color{gray}\small граф нельзя построить }
   \STATE{\textbf{if $d=0$ then return true end if}}
   \COMMENT{\color{gray}\small граф уже построен}
   \STATE{$u\Assign t.l;\ u.r\Assign\textbf{nil}; 
   \ \sim t;\ t\Assign u$ }
   \COMMENT{\color{gray}\small откладываем вершину $t$}  
   \FOR{\textbf{$j$ from $1$ to $d$}}
   \STATE{$u.d\Assign u.d-1;\ q\Assign q+1;\
    G[q]\Assign (u.n, n);\ u\Assign u.l$ }
   \COMMENT{\color{gray}\small ребро } 
   \ENDFOR
   \STATE{\textbf{if $u\ne$ nil 
   then $\Ordering(s, t, u.r)$ end if}}
   \COMMENT{\color{gray}\small восстановление }
\ENDFOR    
\end{algorithmic}
\end{frame}

\begin{frame}
\frametitle{7.2.7. Критерий степенной последовательности (4/6)}

\vspace{-4pt}
Процедура $\Ordering$ восстанавливает упорядоченность двусвязного списка после откладывания вершины,
используя тот факт, что отрезки слева и справа  упорядочены.  

%\vspace{-12pt}
\begin{center}
\includegraphics[width=125mm]{7_2_7_3.jpg}
\end{center}

\vspace{-10pt}
\begin{algorithmic}
\REQUIRE{ указатель на начало списка $s$, 
указатель на конец списка $t$, 
указатель на тот элемент $u$, где может быть нарушен порядок.}  
\ENSURE{ переупорядоченный список }
\IF{$u.l.d > u.d$}
\STATE{$v\Assign u.l$}
\COMMENT{\color{gray}\small $v.d > u.d$ --- нарушение порядка }
   \STATE{\textbf{$z\Assign u;$ 
   while $z\ne t\And v.d>z.r.d$ do $z\Assign z.r$ 
   end while}}
   \COMMENT{\color{gray}\small $v.d\leq z.r.d$}
   \STATE{\textbf{$w\Assign v;$ 
   while $w\ne s\And w.l.d>u.d$ do $w\Assign w.l$ 
   end while}}
   \COMMENT{\color{gray}\small $ w.l.d \leq u.d$}

   \STATE{\textbf{if $z\ne t$  
   then $z.r.l\Assign v$ end if}}
   \COMMENT{\color{gray}\small справа что-то осталось }
   \STATE{\textbf{if $w\ne s$  
   then $w.l.r\Assign u$ end if}}
   \COMMENT{\color{gray}\small слева что-то осталось }
      \STATE{$v.r\Assign z.r;$  
          $u.l\Assign w.l;$ 
          $z.r\Assign w;$ 
          $w.l\Assign z$}
    \COMMENT{\color{gray}\small зашили разрыв}
\ENDIF    
\end{algorithmic}
\end{frame}

\begin{frame}
\frametitle{7.2.7. Критерий степенной последовательности (5/6)}
График является иллюстрацией изменения порядка элементов данного двусвязного списка. Более толстыми стрелками обозначены действия производимые в первую очередь. Тонкой пунктирной стрелкой, соответственно, обозначено <<зашивание разрыва>>. Случай, когда переменная $z$ и/или $w$ <<упирается>> в <<край>> списка, предполагается понятным для читателей после рассмотрения предложенного на графике случая.
\begin{center}
\includegraphics[width=100mm]{7_2_77_2.jpg}
\end{center}

\end{frame}
\begin{frame}
\frametitle{7.2.7. Критерий степенной последовательности (6/6)}
\begin{ground} Алгоритм следует доказательству теоремы (последовательно откладывает вершины) с учётом следующих замечаний.

1. Двусвязный список можно представить одним указателем, но здесь для удобства используются два: $s$ и $t$.

2. Если последнее число $t.d$ в последовательности $D$ больше или равно числу вершин $p$, то алгоритм нельзя применять.

3. Если в результате успешных откладываний 
в конце списка $t.d$ появился 0, то граф уже построен.

4. Если в результате откладывания в начале списка
$s.d$ появляются
отрицательные числа, то граф построить невозможно.

5. Деструктор элемента списка $t$ обозначен знаком 
$\sim$.

6. Элементу массива $G$ присваивается сразу пара вершин (ребро).

7. Восстанавливать упорядоченность списка необходимо, только если следующая вершина существует. 

Алгоритм заканчивает работу,
поскольку на каждом шаге основного цикла один
элемент из списка удаляется.
\end{ground}
\end{frame}



\begin{frame}
\frametitle{Онтология теории графов}
\begin{center}
\includegraphics[height=8cm]{ODM7_2.png}
\end{center}

\end{frame}


\subsection{7.3. Виды графов и операции над графами}
\label{7.3.}
\begin{frame}
\frametitle{7.3. Виды графов и операции над графами}
В данном разделе рассматриваются различные частные случаи графов и
вводятся операции над графами и их элементами.
Заметим, что не все используемые нами  обозначения
операций над графами являются традиционными и~общепринятыми.

\begin{tabular}{p{9mm}|p{56mm}|p{79mm}}
  № & Название параграфа
  & Ключевые термины и обозначения \\
  \hline
  & & \\
\DMref{7.3.1.}{7.3.1.} & Полные и пустые графы 
  & 
  {$C_k$~--- цикл; 
  $K_p$~--- полный граф; $\overline{K_p}$~--- пустой граф; клика, плотность}\\
\DMref{7.3.2.}{7.3.2.}\RS & Двудольные графы 
  & 
  {$K_{m, n}$~--- полный двудольный}
  \\
\DMref{7.3.3.}{7.3.4.} & Звёздные графы &
{$S_n=K_{1,n}$~--- звезда}\\
\DMref{7.3.4.}{7.3.4.} & Направленные орграфы и сети &
{Турнир; сеть; источник; сток
}
\\
\DMref{7.3.5.}{7.3.5.}\RS & Операции над графами &
{$\cup$~--- объединение;
$+$~--- соединение;
$-$~--- удаление;
$+$~--- добавление;
$/$~--- стягивание;
$\uparrow$~--- размножение
}\\   
\DMref{7.3.6.}{7.3.6.} & Свойства простых циклов &
{}\\ 
\DMref{7.3.7.}{7.3.7.}\BS & Рёберные графы &
{Матрица смежности рёбер; $I(G)$~--- рёберный граф}\\  
\end{tabular}

\end{frame}


\begin{frame}
\label{7.3.1.}\subsubsection{7.3.1. Полные и пустые графы }\frametitle{7.3.1. Полные и пустые графы (1/2)}

Граф, состоящий из одной вершины, называется
\odmkey{тривиальным}{Граф (graph)!тривиальный (trivial)}.

\medskip
Граф, состоящий из простого цикла с $k$ вершинами, обозначается~$C_k$.
\begin{example} $C_3$~---  треугольник,
$C_4$~---  квадрат.
\end{example}
\begin{remark}
В классе графов треугольники --- минимальные возможные циклы. 
В классах орграфов 
и псевдографов возможны контуры длиной $1$ 
(\em{петли}) и $2$ (\em{встречные дуги}).    
\end{remark}

\medskip
Граф, в котором любые две вершины смежны, называется
\odmkey{полным}{Граф (graph)!полный (complete)}.
Полный граф с $p$ вершинами обозначается $K_p$, он имеет максимально
возможное число рёбер: 
$$q(K_p) = \frac{p(p-1)}{2}.$$

\medskip
Граф, в котором нет смежных вершин, называется
\odmkey{пустым}{Граф (graph)!пустой (empty)}.
Пустой граф с $p$ вершинами обозначается $\overline{K_p}$, 
он имеет минимально
возможное число рёбер: $q(\overline{K_p}) = 0$.

\begin{remark}
Граф может относиться к различным классам.
Например, $C_3=K_3$.  
\end{remark}


\medskip
Дополнением полного графа является пустой, и обратно:
дополнением пустого графа является полный.
\end{frame}


\begin{frame}
\frametitle{7.3.1. Полные и пустые графы (2/2)}
Максимальный полный подграф (некоторого графа) называется
\odmkey{кликой}{Клика (clique)} (этого графа).
Мак\-симальным подграф является в том смысле, 
что любой другой подграф графа, включающий множество вершин клики, 
не является полным. Граф может иметь несколько клик.  
Максимальное число вершин среди клик данного графа 
называется \odmkey{кликовым числом}{Число (number)!кликовое (clique)} графа (или \odmkey{плотностью}{Число (number)!кликовое (clique)}).


\begin{example}
На рисунке изображён граф и две его клики,
рёбра клик на диаграмме   
отмечены утолщёнными и пунктирными линиями. 
Кликовое число в данном случае равно четырём.
\end{example}

\begin{center}
\includegraphics[height=30mm]{7_3_1.jpg}
\end{center}
 
\end{frame}

\begin{frame}
\label{7.3.2.}\subsubsection{7.3.2. Двудольные графы }\frametitle{7.3.2. Двудольные графы (1/4)}

%\looseness1
Граф $G(V,E)$ называется \odmkey{двудольным}{Граф
(graph)!двудольный (bipartite)} (или \odmkey{биграфом}{Биграф
(bigraph)}, или \odmkey{чётным}{Граф (graph)!чётный (even)}
графом), если множество $V$ может быть разбито на два
непересекающихся множества $V_1$ и $V_2$ (${V_1\cup V_2}=V$,
$V_1\cap V_2=\emptyset$, $V_1\ne\emptyset$, $V_2\ne\emptyset$), причём всякое ребро из $E$ инцидентно
вершине из $V_1$ и вершине из $V_2$ (то есть соединяет вершину из
$V_1$ с вершиной из $V_2$). Множества $V_1$ и~$V_2$ называются
\odmkey{долями}{Доля (part)} двудольного графа. Если двудольный
граф содержит все рёбра, соединяющие множества $V_1$ и $V_2$, то
он называется \odmkey{полным двудольным графом}{Граф
(graph)!двудольный (bipartite)!полный (complete)}. Если $\vert
V_1\vert=m$ и $\vert V_2\vert=n$, то полный двудольный граф
обозначается~$K_{m,n}$.

\begin{columns}[t]
\begin{column}[T]{9.2cm}
\begin{example}
На рисунке в \DMref{п.~7.1.6}{7.1.6.} приведена диаграмма графа $K_{3,3}$, который является полным двудольным,
а на рисунке справа изображена 
диаграмма двудольного графа, который не является полным двудольным графом.
\end{example}
\end{column}
\begin{column}[T]{6cm}

\vspace{-10pt}
\begin{figure}[h]
\begin{center}
\includegraphics[height=20mm]{7_11.jpg}
\end{center}
%\caption{Диаграмма графа $K_{3,3}$}
%\label{K33}
\end{figure}
\end{column}
\end{columns}

\medskip
\begin{remark}
Обсуждая двудольность, достаточно рассматривать только 
связные графы, поскольку каждую компоненту связности
можно рассматривать отдельно.
\end{remark}
\end{frame}

\begin{frame}
\frametitle{7.3.2. Двудольные графы (2/4)}
\begin{theorem}
Граф является двудольным тогда и только тогда, когда в нём нет
простых циклов нечётной длины.
\end{theorem}


\begin{proof}

\proofsection{$\Impll$}  Разобьём множество $V$ на
подмножества $V_1$ и $V_2$ следующей про\-цедурой.

\begin{algorithmic}
\REQUIRE граф $G(V,E)$.
\ENSURE Множества $V_1$ и $V_2$~--- доли графа.
\STATE \textbf{select $v\in V$} \COMMENT{\color{gray}\small  произвольная вершина }
\STATE {$V_1\Assign \{v\}$; 
$V_2\Assign\emptyset$} 
\COMMENT{\color{gray}\small  вначале первая доля содержит $v$, а вторая пуста }
\FOR{$u \in V - v$}
  \IF {$d(v,u)$~--- чётно}
    \STATE $V_1\Assign V_1 + u$
    \COMMENT{\color{gray}\small помещаем вершину $u$ в первую долю }
  \ELSE
    \STATE $V_2\Assign V_2 + u$
    \COMMENT{\color{gray}\small  помещаем вершину $u$ во вторую долю }
\ENDIF
\ENDFOR
\end{algorithmic}
\end{proof}
\end{frame}

\begin{frame}
\frametitle{7.3.2. Двудольные графы (3/4)}

\begin{proof} (Продолжение)
Далее от противного. Пусть есть две вершины в одной доле,
соединённые ребром. Пусть для определённости $u,w\in V_2$ и
$(u,w)\in E$. Рассмотрим геодезические $\divC{v}{u}$ и
$\divC{v}{w}$ (здесь $v$~--- та произвольная вершина, которая
использовалась в~алгоритме построения долей графа). Тогда длины
$\vert\divC{v}{u} \vert$ и $\vert\divC{v}{w} \vert$ нечётны. Эти
геодезические имеют общие вершины (по крайней мере, вершину $v$).
Рассмотрим наиболее удалённую от $v$ общую вершину геодезических
$\divC{v}{u}$ и $\divC{v}{w}$ и~обозначим её $v'$ (может оказаться
так, что $v=v'$).

Имеем $\vert\divC{v'}{u}
\vert+\vert\divC{v'}{w} \vert =\vert\divC{v}{u} \vert 
+
\vert\divC{v}{w} \vert-2\vert\divC{v}{v'} \vert$~--- чётно и
$v',\dots, u,w,\dots,v'$~--- простой цикл нечётной длины, что
противоречит условию. Если же $u,w\in V_1$, то длины
$\vert\divC{v}{u} \vert$ и $\vert\divC{v}{w} \vert$ чётны, и
 имеем $v',\dots,u,w,\dots,v'$~--- простой цикл
нечётной длины.
\proofsection{$\Impl$} \!От противного. Пусть
$G(\!V_1,\!V_2;\!E)$~--- двудольный граф и $v_1,v_2,\dots,
v_{2k+1},$ $v_1$~--- простой цикл нечётной длины. Пусть $v_1\in V_1$,
тогда $v_2\in V_2$, $v_3\in V_1$, ${v_4\in
V_2},\ldots,v_{2k+1}\in V_1$. Имеем: $v_1,v_{2k+1}\in
V_1$ и $(v_1,v_{2k+1})\in E$, что противоречит двудольности.
\end{proof}
\end{frame}

\begin{frame}
\frametitle{7.3.2. Двудольные графы (4/4)}
%\addvspace{-16pt}
\begin{corollary}
Ациклические графы двудольны.
\end{corollary}

\begin{example}
На рисунке представлены диаграммы всех ациклических связных двудольных графов с пятью вершинами.


\begin{figure}[h]
\begin{center}
\includegraphics[height=15mm]{7_3_2_2.jpg}
\end{center}
\end{figure}
\end{example}

\begin{remark}
Если двудольный граф связен, то вершины по долям 
распределяются единственным образом. Если же граф не является связным,
то распределение вершин по долям неоднозначно.
\end{remark} 
\begin{example}
В пустом графе
\em{любое} разбиение множества вершин на два непустых подмножества 
определяет доли двудольного графа.  
\end{example}

\begin{digress}
Среди вершин, входящих в одну долю графа, нет смежных,
то есть доли двудольного графа образуют \em{независимые множества
вершин} (см. \DMref{п.~10.4.1}{10.4.1.}). Исходя из этого
наблюдения, можно определить 
\odmkey{$k$-дольный граф}{Граф (graph)!к-дольный (k-partite)}, 
как граф, вершины которого можно 
разбить на $k$ независимых множеств. 
\end{digress}
\end{frame}

\begin{frame}
\label{7.3.3.}\subsubsection{7.3.3. Звёздные графы}
\frametitle{7.3.3. Звёздные графы (1/3)}
Полный двудольный граф вида $K_{1,n}$, где 
$n > 1$ называется 
\odmkey{звёздным графом}{Граф (graph)!звёздный (star)} 
(или кратко \odmkey{звёздной}{Граф (graph)!звёздный (star)}) и
обозначается  $S_n$. Единственная вершина, входящая в первую долю $V_1$
называется \odmkey{центральной вершиной звезды}{Центр (center)!звезды (of star)}.
%Пример. Звезды 𝑆3 и 𝑆4.
%Теорема

\begin{remark}
В определении мы полагаем $n>1$, чтобы
исключить из рассмотрения вырожденные случаи
тривиального графа 
$K_1=S_1$ и \odmkey{диполя}{Диполь (dipole)} $K_2=K_{1,1}=S_1$.
\end{remark}

\begin{theorem}
Центр звёздного графа и 
его центральная вершина совпадают.
\end{theorem}
\begin{proof}
Эксцентриситет центральной вершины равен 1,
а эксцентриситеты всех остальных вершин равны 2.
\end{proof}

\begin{remark}
Доказанная теорема
позволяет называть центральную вершину 
\em{центром}. 
\end{remark}


\medskip
Звёзды играют особую роль 
в рёберном изоморфизме.

\smallskip
\begin{lemma}  
Пусть $G_1(V_1, E_1)$ и $G_2(V_2, E_2)$~---
связные рёберно изоморфные графы c изоморфизмом $h$, причём $p(G_1)>3$, $p(G_2)>3$. Тогда если для некоторого $r>2$ ребра $\Tuple{e}{1}{r}$ составляют звезду $S_r$ в графе $G_1$, то их образы $h(e_1), 
\ldots, h(e_r)$
составляют такую же звезду $S_r$ в графе $G_2$.
\end{lemma}
\begin{remark}
То есть рёберный изоморфизм уважает звёзды. 
\end{remark}

\end{frame}

\begin{frame}
\frametitle{7.3.3. Звёздные графы (2/3)}

\begin{proof}
\begin{columns}[t]
\begin{column}{8cm}
Индукция по $r$. База $r=3$.
Возможны два изоморфизма, показанные на рисунке.
Покажем, что <<левый>> изоморфизм $g$ невозможен.
Действительно, поскольку $p(G_1)>3$, то в исходном графе должна быть ещё одна вершина, которая в силу связности должна иметь инцидентное ей ребро $e_4$. 
\end{column}
\begin{column}{7.5cm}
\vspace{-12pt}
\begin{center}
\includegraphics[width=7cm]{7_3_3_1.jpg}
\end{center}
\end{column}
\end{columns}

\vspace{4pt}
\begin{columns}[t]
\begin{column}{8cm}
Ребро $e_4$ смежно либо только одному ребру из 
множества $\{e_1, e_2, e_3\}$, либо всем трём, но при изоморфизме $g$ в графе $G_2$ ребро $g(e_4)$ может быть смежно только ровно двум ребрам из
множества $\{g(e_1),g(e_2),g(e_3)\}$.
Таким образом, изоморфизм $g$ невозможен и база индукции доказана.
\end{column}
\begin{column}{7.5cm}
\vspace{-12pt}
\begin{center}
\includegraphics[width=7cm]{7_3_3_2.jpg}
\end{center}
\end{column}
\end{columns}
\end{proof}
\end{frame}

\begin{frame}
\frametitle{7.3.3. Звёздные графы (3/3)}
\begin{proof} (Продолжение)
Пусть теперь лемма выполнена для 
всех звёзд с числом рёбер $r$ или менее.
Тогда если $r+1$ рёбер $\Tuple{e}{1}{r+1}$  образуют звезду $S_{r+1}$, а за вычетом
ребра $e_1$ или $e_{r+1}$ образуют звезду $S_r$ в графе $G_1$, то по предположению индукции рёбра 
$h(e_1),\ldots,h(e_r)$ образуют звезду $S_r$ с центром $v_1$ в графе $G_2$, аналогично рёбра 
$h(e_2),\ldots,h(e_{r+1})$ образуют звезду $S_r$ с центром $v_2$ в графе $G_2$. Значит по лемме  \DMref{п.~7.1.6}{7.1.6.} $v_1=v_2$ и рёбра $h(e_1),\dots,h(e_{r+1})$ действительно образуют
звезду $S_{r+1}$.
\end{proof}




%\begin{remark}
%Условие $p(G_1)\!>\!3$, $p(G_2)\!>\!3$ существенно,
%поскольку при $G_1\!=\!S_3, p(G_1)\!=\!4$, $G_2=C_3, p(G_2)=3$ лемма не имеет места, что видно из примера в п.~1.7.6.  
%\end{remark}

Любое подмножество рёбер звезды из трёх
или более элементов также образует звезду,
причём с тем же центром. 
Для вершины $v$ в графе $G$ можно 
определить \em{максимальную} звезду как звезду, 
центром которой является $v$, 
такую что она не является 
собственным подмножеством никакой другой звезды. Обозначим её как $S_G(v)$.

\begin{remark}
Фактически, максимальная звезда для вершины~---
это её окрестность, если степень больше двух.
\end{remark}

\begin{corollary}
Пусть $G_1(V_1, E_1)$ и $G_2(V_2, E_2)$~---
связные рёберно изоморфные графы c изоморфизмом $h$, причём $p(G_1)>3$, $p(G_2)>3$. Тогда максимальной звезде $S_{G_1}(v)$, образованной рёбрами  $\Tuple{e}{1}{d(v)}$, изоморфизм $h$ взаимно однозначно составляет максимальную звезду $S_{G_2}(v\prime)$, образованную рёбрами 
$h(e_1), \ldots, h(e_{d(v)})$
 в графе $G_2$. 
\end{corollary}

\end{frame}


\begin{frame}
\label{7.3.4.}\subsubsection{7.3.4. Направленные орграфы и сети}
\frametitle{7.3.4. Направленные орграфы и сети}
Если в графе ориентировать все рёбра, то получится орграф, который
называется \odmkey{направленным}{Орграф (digraph)!направленный
(directed)}, или \odmkey{антисимметричным}{Орграф
(digraph)!антисимметричный (antisymmetric)}. Направленный орграф,
полученный из полного графа, называется \odmkey{турниром}{Турнир
(tournament)}.

\begin{remark}
В направленном орграфе не может быть <<встречных>> дуг $(u,v)$ и $(v,u)$,
а в произвольном орграфе это допустимо.
\end{remark}

\begin{digress}
Название <<турнир>> имеет следующее происхождение.
Рассмотрим спортивное соревнование для пар участников
(или пар команд), где не предусматриваются ничьи.
Пометим вершины орграфа участниками и проведем дуги
от победителей к побеждённым. В таком случае
турнир в смысле
теории графов~--- это как раз
результат однокругового турнира в~спортивном смысле.
\end{digress}

\smallskip
Если в орграфе полустепень захода некоторого узла равна нулю
(то~есть $d^+(v) 
=0$), то такой узел называется
\odmkey{источником}{Источник (source)}, если же нулю равна
полустепень исхода
 (то~есть $d^-(v)=0$), то узел называется
\odmkey{стоком}{Сток (sink)}.
Направленный слабосвязный (см. 
\DMref{п.~8.5.1}{8.5.1.})
орграф с одним источником и одним стоком
называется \odmkey{сетью}{Сеть (network)}.

\end{frame}

\begin{frame}
\label{7.3.5.}\subsubsection{7.3.5. Операции над графами }\frametitle{7.3.5. Операции над графами (1/6)}

%Введем следующие операции над графами:
\begin{itemize}

%\item \odmkey{Дополнением графа}{Дополнение (complement)!графа (of
%graph)} $G_1(V_1,E_1)$ (обозначение --- $\overline G_1(V_1,E_1)$)
%называется граф
%%\linebreak
%$G_2(V_2,E_2)$, где
%$
%V_2= V_1\And
%E_2= \overline{E_1} =
%\{\divI{e}{V_1} \times V_1\mid e\not\in E_1\} = V\times V\setminus E.
%$
%\begin{example}
%$\overline{K_1}=K_1$.
%\end{example}
\item \odmkey{Объединение (дизъюнктное) графов}{Объединение
(union)!графов (of graphs)} $G_1(V_1,E_1)$ и $G_2(V_2,E_2)$ при условии
$V_1\cap V_2=\emptyset$
(обозначение --- $G_1(V_1,E_1)\cup G_2(V_2,E_2)$) даёт граф $G(V,E)$, \\где 
{$ V=V_1\cup V_2\And E= E_1\cup E_2.$}
\begin{example}
${\overline K_{3,3}}=C_3\cup C_3$.
\end{example}
\item \odmkey{Соединение графов}{Соединение графов (join of
graphs)} $G_1(V_1,E_1)$ и $G_2(V_2,E_2)$
при условии $V_1\cap
V_2=\emptyset$ \\(обозначение ---
$G_1(V_1,E_1)+ G_2(V_2,E_2)$) даёт граф $G(V,E)$,\\ где $ V = V_1\cup
V_2\And E= E_1\cup E_2 \cup \Set{e=(v_1,v_2)}{\divI{v_1}{V_1}\And
 \divI{v_2}{V_2}}.
$
\begin{example}
${K_{3,3}}=\overline C_3+\overline C_3$.
\end{example}
\item \odmkey{Удаление вершины}{Удаление (removal)!вершины (of
vertex)} $v$ из графа $G_1(V_1,E_1)$ при условии $\divI{v}{V_1}$ \\(обозначение ---
$G_1(V_1,E_1)-v$) даёт граф $
G_2(V_2,E_2)$, \\где
$V_2= V_1 - v \And
E_2= E_1\setminus
\Set{e=(v_1,v_2)}{v_1=v\Or v_2=v}.
$
\begin{example}
$C_3-v=K_2$.
\end{example}
\item \odmkey{Удаление ребра}{Удаление (removal)!ребра (of edge)}
$e$ из графа $G_1(V_1,E_1)$ при
условии $\divI{e}{E_1}$ \\(обозначение --- $G_1(V_1,E_1)-e$) даёт граф $ G_2(V_2,E_2)$, 
где $ V_2 = V_1\And E_2= E_1 - e. $
\begin{example}
$K_2-e=\overline K_2$.
\end{example}

\end{itemize}
\end{frame}

\begin{frame}
\frametitle{7.3.5. Операции над графами (2/6)}


\begin{itemize}


\item \odmkey{Добавление вершины}{Добавление (adding)!вершины (of
vertex)} $v$ в граф $G_1(V_1,E_1)$ при условии $v\notin V_1$ \\(обозначение ---
$G_1(V_1,E_1)+v$) даёт граф
$G_2(V_2,E_2)$, где $ V_2= V_1 +v \And E_2= E_1. $
\begin{example}
$K_2+v=K_2\cup K_1$.
\end{example}
\item \odmkey{Добавление ребра}{Добавление (adding)!ребра (of
edge)} $e$ в граф $G_1(V_1,E_1)$ при условии $e\not\in E_1$ \\(обозначение --- $G_1(V_1,E_1)+e$,
) даёт граф $ G_2(V_2,E_2)$, где $ V_2=
V_1\And E_2= E_1 + e. $ \\
\begin{example} $P_3+e=C_3$. \end{example}
\item \odmkey{Стягивание правильного
подграфа}{Стягивание подграфа (shrinkage of subgraph)} с
множеством вершин $A$ графа $G_1(V_1,E_1)$ к вершине $v$ при условии $A\subset V_1$
(обозначение ---
$G_1(V_1,E_1)/A$) даёт граф
$G_2(V_2,E_2)$, \\где
$V_2 = (V_1\setminus A) + v$,\\
$E_2 = E_1 \setminus \Set{e=(u,w)}{\divI{u}{A}\Or\divI{w}{A}}\cup
\Set{e=(u,v)}{\divI{u}{\Gamma(A)}\setminus A}$.
\begin{example}
$K_4/C_3=K_2$.
\end{example}
\begin{remark}
Если $A=\{u, v\}$, причём $e=(u,v)\in E$,
то допустима запись $G/e$.  
\end{remark}
\item \odmkey{Размножение вершины}{Размножение вершины (duplication of vertex)}
$v$ графа $G_1(V_1,E_1)$ при условии
$\divI{v}{V_1}, v'\not\in V_1$ \\
(обозначение --- $G_1(V_1,E_1) \uparrow v $)
даёт граф $G_2(V_2,E_2)$, \\где
$V_2= V_1+ v'\And E_2= E_1\cup \{(v,v')\}\cup
\Set{e=(u,v')}{\divI{u}{\Gamma^+(v)}}.
$
\begin{example}
$K_2\uparrow v=C_3$.
\end{example}
\end{itemize}


\end{frame}

\begin{frame}
\frametitle{7.3.5. Операции над графами (3/6)}
Некоторые из примеров, приведённых в определениях операций,
нетрудно обобщить. В частности, легко показать, что имеют место
следующие соотношения:
\begin{align*}
K_{m,n} &= \overline{K}_m+\overline{K}_n,
\quad &K_{p-1} &= K_p - v, \\
G+v &= G\cup K_1,
&K_{p-1} &= K_p / K_2, \\
K_p/K_{p-1} &=K_2, &K_{p} &=K_{p-1} \uparrow v.
\end{align*}
Введённые операции обладают рядом простых свойств, которые
легко вывести из определений. В частности:
\begin{align*}
G_1\cup G_2 &= G_2\cup G_1, &
G_1 + G_2 &= G_2 + G_1, \\
G_1 = G_2 &\Equiv \overline{G}_1 = \overline{G}_2, &
\overline{G_1\cup G_2} &= \overline{G}_1 + \overline{G}_2.
\end{align*}
\begin{remark}
Операции добавления и удаления ребра взаимно обратны:
$$
\A{e\in E, e'\notin E}{(G-e)+e=(G+e')-e' =G}.
$$
В то же время $\A{v'\notin V}{(G+v')-v'=G}$, но если вершина
$v$ графа $G$ не изолированная, то $(G-v)+v\ne G$,
потому что в этом случае $q((G-v)+v)<q(G)$.
\end{remark}
\end{frame}

\begin{frame}
\frametitle{7.3.5. Операции над графами (4/6)}
Результат выполнения нескольких последовательных операций удаления
или добавления вершин или рёбер не зависит от порядка выполнения операций:
\begin{align*}
(G + v_1) + v_2 &= (G + v_2) + v_1, &
(G - v_1) - v_2 &= (G - v_2) - v_1, \\
(G + e_1) + e_2 &= (G + e_2) + e_1, & (G - e_1) - e_2 &= (G - e_2)
- e_1.
\end{align*}
Поэтому операции удаления и добавления вершин и рёбер можно обобщить,
и~допускать в качестве второго аргумента \em{множества} вершин и рёбер.

\medskip
\begin{digress}
Приведённые определения операций над графами и примеры к ним дают
повод затронуть один тонкий вопрос, связанный с различиями в
традициях математических и программных обозначений. Рассмотрим
пример к операции дизъюнктного объединения графов
($\overline{K}_{3,3}=C_3\cup C_3$). Если введённые обозначения
понимать буквально, то этот пример кажется противоречащим
определению. 
\end{digress}

\end{frame}

\begin{frame}
\frametitle{7.3.5. Операции над графами (5/6)}

\begin{digress} (Продолжение)
Действительно, определение требует, чтобы множества
вершин объединяемых графов не пересекались. Если считать, что
$C_3$ обозначает конкретный треугольник, то придётся признать, что
в выражении $C_3\cup C_3$ множества вершин не только пересекаются,
но и совпадают, а значит, операция объединения неприменима. 
На
самом деле $C_3$ обозначает как класс треугольников (все
треугольники изоморфны как графы), так и отдельный объект~---
\odmkey{экземпляр}{Экземпляр класса (instance of class)} этого
класса. Как именно следует понимать обозначение, считается ясным
из контекста. Если стремиться к (излишней в данном случае)
строгости и однозначности обозначений, то приведённый пример можно
было бы записать, например, так:

$$
\forall\, G_1(V_1,E_1)\in C_3, 
G_2(V_2,E_2) \in C_3 $$ 
$$\big( V_1\cap V_2 = \emptyset\Impl\exists\, G_3(V_3,E_3)\in K_{3,3} (G_1\cup G_2\sim\overline{G}_3)\big).
$$
\end{digress}
%$$
%\A{G_1(V_1,E_1)\in C_3,G_2(V_2,E_2)\in C_3}{V_1\cap V_2 \ne \emptyset\Impl
%\E{G_3(V_3,E_3)\in K_{3,3}}{G_1\cup G_2\sim \overline{G}_3}}.
%$$
\end{frame}

\begin{frame}
\frametitle{7.3.5. Операции над графами (6/6)}
\begin{digress} (Окончание)
Подобная неоднозначность обозначений присуща и программированию,
хотя и в меньшей степени. Например, ключевое слово \textbf{int} в
различных контекстах может обозначать встроенный тип данных (класс
объектов), операцию порождения нового объекта этого типа
(экземпляра класса), операцию преобразования другого объекта в
объект типа \textbf{int} (явное приведение). Следуя стилю
объектно-ориентированных языков программирования, приведённый
пример можно было бы записать так:
$$
\overline{\textbf{new} K_{3,3}} = \textbf{new} C_3 \cup \textbf{new} C_3.
$$
Ради краткости и простоты изложения здесь принят
значительно менее строгий стиль обозначений в надежде на то, что
программистская интуиция и здравый смысл позволят читателю
избежать заблуждений, несмотря на вольности в обозначениях.
\end{digress}

\end{frame}



\begin{frame}
\label{7.3.6.}\subsubsection{7.3.6. Свойства простых циклов }\frametitle{7.3.6. Свойства простых циклов (1/2)}
\begin{theorem} \NB{[1]}
Регулярный граф степени 2 состоит из простых циклов.
\end{theorem}

\begin{proof}
Рассмотрим сначала случай связного
графа. Пусть $G(V,E)$~--- связный граф
и $\A{v\in V}{d(v)=2}$. 
Рассмотрим произвольную вершину $v\in V$.
У неё есть две смежные вершины. 
Пусть одна из них~--- вершина $u$.
Перейдём из вершины $v$ 
в вершину $u$ по ребру $(v,u)$. 
Тогда у вершины $u$ есть ровно одна смежная
и ранее не пройденная вершина $w$. 
Перейдем в неё по ребру $(u,w)$.
Таким образом, переходим единственным образом каждый раз по новому 
ребру либо в ранее непройденную вершину,
либо в начальную вершину $v$.
Действительно, поскольку переходы однозначные,
недостижение какой-либо вершины противоречит связности графа. Таким образом,
связный регулярный граф степени 2~--- 
это простой цикл. Рассмотрим теперь
несвязный граф. Он состоит из $k$ 
компонент связности, каждая из которых является простым циклом.
\end{proof}

\begin{theorem} \NB{[2]}
В ациклическом графе рёбер меньше, чем вершин:
$z(G)=0 \Impl q(G)< p(G)$.
\end{theorem}
\begin{proof}
Рассмотрим сначала случай, когда $G(V,E)$~--- связный ациклический граф. Индукция по $p$. 
База: при $p=1$ имеем $q=0$,
при $p=2$ имеем $q=1$.
Пусть утверждение теоремы 
верно для всех графов с числом вершин 
меньше $p$. 
\end{proof}
\end{frame}

\begin{frame}
\frametitle{7.3.6. Свойства простых циклов (2/2)}

\begin{proof} (продолжение)
Рассмотрим 
связный ацикличеcкий граф с $p$ вершинами.
Если в графе есть висячая вершина $v$,
то удалим её, $G\Assign G-v$. Тогда 
$q-1 < p-1$ по индукционному предположению,
и $q<p$. Если висячих вершин нет,
то $\A{v\in V}{d(v)\geq 2}$.
По образцу предыдущей теоремы,
возьмём произвольную вершину и 
перейдём в смежную вершину по ранее
не пройденному ребру. Это возможно,
поскольку степень любой 
вершины не меньше двух.
Ввиду конечности множества вершин,
рано или поздно попадём в одну
из ранее пройденных вершин, то есть
выявим цикл~--- противоречие.
Рассмотрим теперь несвязный ациклический граф.
Его компонентами являются связные ациклические графы,
в каждом из которых $q<p$.  
\end{proof}
\begin{corollary}
$q(G)\ge p(G) \Impl z(G)>0$.
\end{corollary}

\begin{theorem} \NB{[3]}
Если в графе нет изолированных и висячих вершин, 
то в нём есть цикл.
\end{theorem}

\begin{proof}
По лемме о рукопожатиях, если нет изолированных и висячих вершин, то 
$2q = \sum_{v\in V} d(v)\ge 2p$, а значит по предыдущему следствию $z(G)>0$.
\end{proof}

\smallskip
\begin{remark}
Указанное справедливо как для связных,
так и для несвязных графов. 
\end{remark} 
\end{frame}


\begin{frame}
\label{7.3.7.}\subsubsection{7.3.7. Рёберные графы }\frametitle{7.3.7. Рёберные графы (1/6)}
Пусть $G(V, E)$~--- граф. 
Составим квадратную матрицу \textbf{$I(G):$ array $[1..q,1..q]$ of $0..1$}, 
в которой и строки, и столбцы отвечают рёбрам $G$, и определим её элементы 
следующим образом:
$I(G)[x,y]\Assign 
\E{v\in V}{v\in x \And v\in y}.$
Другими словами, если рёбра $x$ и $y$ смежны и различны, 
то в соответствующей клетке матрицы ставим $1$, иначе $0$. 
Такую матрицу называют \odmkey{матрицей смежности рёбер}{Матрица (matrix)!смежности рёбер (adjacency of edges)}.

\smallskip
\odmkey{Рёберным графом}{Граф (graph)!рёберный (line)} для $(p,q)$-графа $G(V,E)$ 
(обозначение $I(G)$) называется граф с $q$ вершинами, 
соответствующими рёбрам исходного графа, 
причём вершины рёберного графа смежны тогда и только тогда, 
когда смежны соответствующие рёбра исходного графа.
Таким образом, матрица смежности рёбер является матрицей смежности 
(см. \DMref{п.~7.4.2}{7.4.2.}) рёберного графа.


\begin{remark}
Построить диаграмму рёберного графа $I(G)$ 
по диаграмме графа $G$ очень просто. 
Достаточно поставить по одной точке в произвольном месте каждой линии, 
изображающей ребро $G$, и затем соединить попарно все точки, 
которые соответствуют смежным ребрам. 
\end{remark}

\end{frame}


\begin{frame}
\frametitle{7.3.7. Рёберные графы (2/6)}
\begin{examples} Ниже приведены диаграммы 
графов $I(P_3)=K_2, I(C_3)=C_3, 
I(S_3)=C_3$,\\ $I(C_4)=C_4, I(S_4)=C_4$.

%\vspace{-8pt}
\begin{center}
\includegraphics[width=14cm]{7_3_7_1.jpg}
\end{center} 
\end{examples}
Мы видим, что для всякого графа существует соответствующий ему рёберный граф.

\medskip
Возникает естественный вопрос: 
как узнать, является ли данный граф рёберным графом 
некоторого другого графа? Ответ на вопрос дает следующая теорема.

\begin{theorem}
Граф $I(G)$ является рёберным графом для некоторого графа $G$ 
тогда и только тогда, когда он представим в виде множества рёберно 
непересекающихся клик $K_{d(v)}$ с $d(v)$ вершинами, таких, 
что $K_{d(v)}$ имеет не более одной общей вершины с любой другой кликой, причём вершина может входить не более чем в две клики.
\end{theorem}


\end{frame}


\begin{frame}
\frametitle{7.3.7. Рёберные графы (3/6)}

\begin{proof} 
\proofsection{$\Rightarrow$}
Пусть $I(G)$~--- рёберный граф для графа $G(V, E)$. 
Рассмотрим произвольную вершину $v\in V$. 
Ей инцидентны $d(v)$ рёбер. 
Рассмотрим любую вершину $u \in \Gamma(v)$
 и пусть $e = (v, u)$. 
 Тогда в $I(G)$ вершина $e$, соответствующая ребру $e$ графа $G$, 
 соединяется ребрами с $d(v)-1$ вершинами, 
 соответствующими остальным ребрам $G$, инцидентным $v$.
То есть, в $I(G)$ имеем полный подграф $K_{d(v)}$, 
порождённый вершиной $v$ графа $G$. 
Рассмотрим теперь вершину $u$. 
Проведя для неё аналогичные рассуждения, 
приходим к тому, что $e$ соединена ребрами также с $d(u)-1$ вершинами, 
которые вместе с $e$ образуют другой полный граф $K_{d(u)}$. 
Два графа, $K_{d(v)}$ и $K_{d(u)}$, имеют ровно одну общую вершину, 
определяемую единственным ребром $e$, связывающим $v$ и $u$.

\vspace{-6pt}
\begin{center}
\includegraphics[height=30mm]{7_3_6.jpg}
\end{center}
%\vspace{-6pt}

\vspace{-20pt} 
\end{proof}

\end{frame}


\begin{frame}
\frametitle{7.3.7. Рёберные графы (4/6)}
\begin{proof} (Окончание)
\proofsection{$\Leftarrow$}
Предположим, что, наоборот, 
граф $G_1$ представим в виде множества рёберно непересекающихся 
клик $U(e_i)$, и любая пара этих клик $(U(e_i), U(e_j))$ 
имеет не более одной общей вершины, 
причём вершина может входить не более чем в две клики. 
Тогда $G_1$ можно рассматривать как рёберный граф $I(G)$ 
некоторого графа $G$. Ясно, что каждая клика $U(e_i)$ 
имеет $d_i \le d(e_i) + 1$ общих вершин с другими кликами $U(e_j)$. 
Каждому $U(e_i)$ поставим в соответствие одну вершину $v_i$ 
и соединим $v_i$ и $v_j$ ребром в $G$ тогда и только тогда, 
когда $U(e_i)$ и $U(e_j)$ имеют общую вершину. 
К этим рёбрам добавим $d(e_i)-d_i$ рёбер $(v_i, u)$, 
идущих к новым вершинам $u$, в которых существует только одно это ребро. 
\end{proof}
%
%\medskip
%Говорят, что два графа $G_1(E_1, V_1)$ и $G_2(E_2, V_2)$  
%\odmkey{рёберно изоморфны}{}, 
%если существует такое взаимно однозначное соответствие 
%$E_1 \leftrightarrow E_2$ между их рёбрами, 
%что если $x_1$ и $y_1$~--- смежные рёбра в $G_1$, 
%то соответствующие им рёбра $x_2$ и $y_2$ смежны в $G_2$, 
%и наоборот. 

\smallskip
Заметим, что два графа рёберно изоморфны тогда и только тогда, 
когда их рёберные графы изоморфны.
% На рисунке представлены два графа, 
%изоморфные рёберно, но не изоморфные в обычном смысле.
%
%\begin{center}
%\includegraphics[height=20mm]{7_3_6__.jpg}
%\end{center}

\smallskip
Как мы выяснили ранее, любой граф $G$ определяет, причём единственным образом, соответствующий рёберный граф $I(G)$. Возникает вопрос, однозначно ли определяет рёберный граф $I(G)$ исходный граф $G$?
Рассмотренный пример~--- треугольник $C_3$ и трёхконечная звезда  
$S_3$~--- показывает, 
что это не всегда так, поскольку 
$I(C_3)=C_3=I(S_3)$. 
Исчерпывающий ответ  даёт следующая теорема.


\end{frame}


\begin{frame}
\frametitle{7.3.7. Рёберные графы (5/6)}

\begin{theorem} \NB{[Уитни]}
Любой связный  граф, 
отличный от $C_3$ и $S_3$, однозначно задаётся своим рёберным графом.
\end{theorem}

\begin{proof}
Рассмотрим графы $G_1$ и $G_2$ такие,
 что $I(G_1) \stackrel{h}{\sim} I(G_2)$ 
 и покажем, что $G_1 \stackrel{g}{\sim} G_2$. 
 Если 
 $p(G_1)\leq 4, p(G_2)\leq 4$, то
рассмотрев конечное число всех таких графов, нетрудно убедиться, что среди них только пара
$C_3$ и $S_3$ имеет
изоморфные рёберные графы,
но именно эта пара исключена 
из рассмотрения условием теоремы.
Далее рассмотрим случай 
$p(G_1)> 4, p(G_2)> 4$. Для
каждой вершины 
$v$ из $G_1$ найдём 
соответствующую ей вершину $u=g(v)$ 
из $G_2$, рассмотрев два случая.

\proofsection{$d(v)>1$}
Если вершина $v$ не висячая, 
то она является центром максимальной звезды. 
По следствию \DMref{п.~7.3.3}{7.3.3.}. ей соответствует
максимальная звезда $S_{G_2}(u)$ с центром 
$u$ в графе $G_2$. Мы построили частичную функцию $g_1(v)\Assign u$ на
множестве вершин графа $G_1$,
причём функция $g_1$ инъективна.  
Действительно, если
$h(S_{G_1}(v_1)) = h(S_{G_1}(v_2)) = S_{G_2}(u)$, то из биективности $h$ имеем 
$S_{G_1}(v_1) = S_{G_1}(v_2) = h^{-1}(S_{G_2}(u))$.
Поскольку это максимальные звезды, то в каждой звезде как минимум
два ребра. В таком случае центр 
единственный и $v_1 =v_2 =v$.
Значит $g_1$~--- инъективная функция, а по соображениям симметрии $g_1$~--- 
биекция на своей области определения.
\end{proof}
\end{frame}


\begin{frame}
\frametitle{7.3.7. Рёберные графы (6/6)}
\begin{proof} (Продолжение)
\proofsection{$d(v)=1$}
Если вершина $v$ висячая, то в силу связности графа и того, что
$p(G_1) > 4$, 
существует вершина $w$
смежная с $v$,  такая что $d(w) > 1$.
Положим $x\Assign(v,w)$.
Вершина $w$ является центром звезды и для неё уже определена вершина
$u=g_1(w)$. Выберем в звезде 
 $S_{G_2}(u) = h(S_{G_1}(w))$ такое ребро
$y = (u, t)$, что $h(x) = y$.
 Покажем, что $d(t) = 1$. 
 Пусть это не так, тогда в
звезде $S_{G_2}(t)$ 
существует ребро $z \ne y$
 не входящее в $S_{G_2}(u)$. Следовательно, 
в звезде $h^{−1}(S_{G_2}(t))$ также будет ребро
$h^{−1}(z)$ смежное с $e$, 
но не входящее в $S_{G_1}(w)$. Значит его концом является $v$,
но тогда $d(v) \ne 1$~--- противоречие. Получаем инъекцию множества висячих
вершин $G_1$ в такое же множество для $G_2$, обозначим её как $g_2$. По
соображениям симметрии $g_2$~--- биекция
на своей области определения.
Наконец покажем, что биекция
$
g(v)\Assign \textbf{ if } d(v)>1 
\textbf{ then } g_1(v) \textbf{ else }
g_2(v) \textbf{ end if }
$
является вершинным изоморфизмом. 
%При этом
%сопоставляемая вершине $v$ вершина $g(v)$ полностью определяется
%отображением $h$ между их максимальными звездами. 
Расширим
обозначение $S_G(v)$ и на случай $d(v) = 1$.
Вершины $v_1$ и $v_2$ смежны в $G_1$ тогда
и только тогда, когда $S_{G_1}(v_1)$
 и $S_{G_1}(v_2)$ имеют общее ребро. В случае, если $e$
является общим ребром $S_{G_1}(v_1)$
 и $S_{G_1}(v_2)$, то $h(e)$ является общим ребром $S_{G_2}(g(v_1)) = h(S_{G_1} (v_1))$
 и $S_{G_2}(g(v_2)) = h(S_{G_1} (v_2))$.
Аналогично и в обратную сторону.
Таким образом $(v_1, v_2) \in G_1 
\Equiv (g(v_1), g(v_2)) \in G_2$. 
Графы $G_1$ и $G_2$
изоморфны.
\end{proof}
\end{frame}




\begin{frame}
\frametitle{Онтология теории графов}
\begin{center}
\includegraphics[height=8cm]{ODM7_3.png}
\end{center}

\end{frame}
\subsection{7.4. Представление графов в программах}
\label{7.4.}
\begin{frame}
\frametitle{7.4. Представление графов в программах}
%\label{7.4}
Следует ещё раз подчеркнуть, что
конструирование структур данных
для представления в программе
объектов математической модели~---
это основа искусства практического программирования.
Мы приводим четыре различных базовых представления графов.
Выбор наилучшего представления
определяется требованиями конкретной задачи.
Более того, на практике используются, как правило,
некоторые комбинации или модификации
указанных представлений,
общее число которых необозримо.
Но все они, так или иначе, основаны на тех базовых идеях,
которые описаны в этом разделе.

\begin{tabular}{p{10mm}|p{56mm}|p{77mm}}
  № & Название параграфа
  & Ключевые термины и обозначения \\
  \hline
\DMref{7.4.1.}{7.4.1.} & Требования к представлению графов 
  & 
  {$O(n)$~--- асимптотическая оценка}\\
\DMref{7.4.2.}{7.4.2.} & Матрица смежности 
  & 
  {$
M:\; \mathbf{array} \; [1..p,1..p] \;  \mathbf{of} \;  0..1
$}
  \\
\DMref{7.4.3.}{7.4.3.} & Матрица инциденций &
{\textbf{$H:$~array~$[1..p,1..q]$~of~$-1..1$}}\\
\DMref{7.4.4.}{7.4.4.} & Списки смежности &
{\textbf{$\Gamma:$ array $[1..p]$ of $\uparrow N$},

\textbf{$N=$ struct $\{ v:1..p\,;\; n:\, \uparrow N\}$},
}
\\
\DMref{7.4.5.}{7.4.5.} & Массив дуг &
{\textbf{$E:$ array $[1..q]$ of struct
$\{ b, e: 1..p \}$}
}\\   
\DMref{7.4.6.}{7.4.6.}\RS & Обходы графов &
{Поиск в ширину и поиск в глубину}\\ 
\end{tabular}

\end{frame}

\begin{frame}
\label{7.4.1.}\subsubsection{7.4.1. Требования к представлению графов }\frametitle{7.4.1. Требования к представлению графов (1/2)}
Известны различные способы представления графов в памяти
компьютера, которые различаются объёмом занимаемой
памяти и скоростью выполнения операций над графами.
Представление выбирается исходя из потребностей
конкретной задачи. Далее приведены четыре наиболее
часто используемых представления с указанием
характеристики
$n(p,q)$~--- объёма памяти для каждого представления.
Здесь $p$~--- число вершин, а $q$ ~--- число рёбер.

\begin{remark}
Значение характеристики $n(p,q)$ указывается с помощью символа
$O$, который означает совпадение по порядку величины (или
равенство с точностью до мультипликативной константы $c$, см.
замечание в \DMref{п.~1.3.2}{1.3.2.}).
Применительно к измерению занимаемой памяти
использование символа $O$ связано с тем, что память может быть
измерена в битах, байтах, машинных словах или иных единицах.
Коэффициент $c$ при этом меняется, а~порядок величины остаётся.
\end{remark}


\end{frame}

\begin{frame}
\frametitle{7.4.1. Требования к представлению графов (2/2)}
Представления иллюстрируются на конкретных примерах графа $G$
и орграфа $D$, диаграммы которых представлены
на рисунке. 
%рис.~\ref{GD}.
\begin{figure}[h]
\begin{center}
\includegraphics[width=100mm]{odm07-10}
\end{center}
%\caption{Диаграммы графа ({\it слева}) и орграфа ({\it справа}),
%используемых в~качестве примеров} \label{GD}
\end{figure}
\end{frame}


\begin{frame}
\label{7.4.2.}\subsubsection{7.4.2. Матрица смежности}
\frametitle{7.4.2. Матрица смежности}
Представление графа с помощью квадратной булевой матрицы
$$
M:\; \mathbf{array} \; [1..p,1..p] \;  \mathbf{of} \;  0..1,
$$
отражающей смежность вершин, называется
\odmkey{матрицей смежности}{Матрица (matrix)!смежности (adjacency)}, \\где
$M[i,j]\Assign (v_i\in\Gamma(v_j))$, то есть вершины $v_i$ и $v_j$
смежны.  

Для матрицы смежности $n(p,q) = O(p^2)$.
\begin{example}
%\vspace{-24pt}
$$
G: \quad
\begin{vmatrix}
0 & 1 & 0 & 1\\
1 & 0 & 1 & 1\\
0 & 1 & 0 & 1\\
1 & 1 & 1 & 0
\end{vmatrix}\qquad
D: \quad
\begin{vmatrix}
0 & 1 & 0 & 0\\
0 & 0 & 1 & 1\\
0 & 0 & 0 & 0\\
1 & 0 & 1 & 0
\end{vmatrix}
$$
\end{example}
\begin{remark}
Матрица смежности графа симметрична относительно
главной диагонали,
поэтому достаточно хранить только верхнюю (или нижнюю)
треугольную матрицу.
\end{remark}
\end{frame}


\begin{frame}
\label{7.4.3.}\subsubsection{7.4.3. Матрица инциденций}
\frametitle{7.4.3. Матрица инциденций}
Представление графа матрицей
\textbf{$H:$~array~$[1..p,1..q]$~of~$ 0..1$},
отражающей инцидентность вершин и рёбер, называется
\odmkey{матрицей инциденций}{Матрица (matrix)!инциденций (incidence)}, где
$H[i,j]\Assign \E{e_j\in E}{v_i\in e_j}$,
то есть вершина $v_i$ инцидентна ребру $e_j$.
Представление орграфа
матрицей \\ \textbf{$H:$~array~$[1..p,1..q]$~of~$-1..1$},
отражающей инцидентность узлов и дуг, также называется
\odmkey{матрицей инциденций}{Матрица (matrix)!инциденций (incidence)}, где
\textbf{$H[i,j]\Assign$
if $\E{e_j\in E}{e_j=(x, v_i)}$ then $1$ \\ elseif
$\E{e_j\in E}{e_j=(v_i, x)}$ then $-1$ else $0$ end if},
то есть если узел $v_i$ инцидентен дуге $e_j$ и является её концом,
то $H[i,j]=1$, а если узел $v_i$ инцидентен дуге $e_j$ и является её началом,
то $H[i,j]=-1$, если же
 узел $v_i$ и ребро $e_j$ не инцидентны,
  то $H[i,\!j]=0$.
 Для  матрицы инциденций  $n(p,q)=O(pq)$.

%\vspace{-3pt}%
\begin{example}
%\addvspace{-12pt}

\vspace{-24pt}
$$G: \quad
\begin{vmatrix}
1 & 0 & 0 & 1 & 0\\
1 & 1 & 0 & 0 & 1\\
0 & 1 & 1 & 0 & 0\\
0 & 0 & 1 & 1 & 1
\end{vmatrix}\qquad
D: \quad
\begin{vmatrix}
-1 &  0 &  0 &  1 &  0\\
 1 & -1 &  0 &  0 & -1\\
 0 &  1 &  1 &  0 &  0\\
 0 &  0 & -1 & -1 &  1
\end{vmatrix}
$$
\end{example}

%\vspace{-6pt}
\begin{remark}
Для связных графов $q\geq p-1$, поэтому матрица смежности
несколько компактнее матрицы инциденций.
\end{remark}
\end{frame}


\begin{frame}
\label{7.4.4.}\subsubsection{7.4.4. Списки смежности }
\frametitle{7.4.4. Списки смежности }
%\label{s744}
Представление графа с помощью списочной структуры,
отражающей смежность вершин и состоящей из массива указателей
$\Gamma: \textbf{array $[1..p]$ of $\uparrow N$}
$
на списки смежных вершин, где
элемент списка представлен структурой
\textbf{$N=$ struct $\{ v:1..p\,;\; n:\, \uparrow N\}$},
называется
\odmkey{списком смежности}{Список (list)!смежности (adjacency)}.
В случае представления неориентированных графов
списками смежности $n(p,q)\!=\!O(p\!+\!2q)$,
а в случае ориентированных графов $n(p,q)\!=\!O(p\!+\!q)$.

\begin{remark}
Массив $\Gamma$ также можно представить списком.
\end{remark}



\begin{example}
Списки смежности для графа $G$ и орграфа $D$ представлены на
рисунке.
\begin{figure}[h]
\begin{center}
\includegraphics[height=35mm]{odm07-11}
\end{center}
%\caption{Списки смежности для графа {\itshape G\/} ({\it слева}) и
%орграфа {\itshape D\/} ({\it справа})} \label{Gamma}
\end{figure}
\end{example}

%\vspace{-3pt}
%{\color{blue} Тема для сочинения: чистые списки.}
\end{frame}


\begin{frame}
\label{7.4.5.}\subsubsection{7.4.5. Массив дуг}
%\frametitle{7.4.4. Списки смежности }
\frametitle{7.4.5. Массив дуг}

Массив структур
$
E: \textbf{array $[1..q]$ of struct
$\{ b, e: 1..p \}$}$,
отражающий список пар смежных вершин (или, для орграфов, узлов),
называется \odmkey{массивом рёбер}{Массив (array)!рёбер (of
edges)} (\odmkey{массивом дуг}{Массив (array)!дуг (of arcs)}). Для
массива рёбер (или дуг) $n(p,q)=O(2q)$.

\begin{remark}
Для представления графов с изолированными вершинами может понадобиться
хранить ещё и число $p$, если только система программирования
не позволяет извлечь значение $p$ из массива структур $E$.
\end{remark}

\begin{columns}[t]
\begin{column}[T]{5cm}
\begin{example}
Представление с помощью массива рёбер (дуг) показано в 
таблице (для графа $G$~--- слева, а для орграфа $D$~--- справа).
\end{example}
\end{column}
\begin{column}[T]{9cm}
%\addvspace{6pt}
%{\sffamily\small
\vspace{-26pt}
\begin{center}
\def\arraystretch{1.1}
\begin{tabular}{p{10mm}p{10mm}||p{10mm}p{10mm}}
$b$ & $e$ & $b$ & $e$ \\
\hline
1 & 2 & 1 & 2 \\
1 & 4 & 2 & 3 \\
2 & 3 & 2 & 4 \\
2 & 4 & 4 & 1 \\
3 & 4 & 4 & 3 \\
\hline
\end{tabular}
\end{center}
\end{column}
\end{columns}

\medskip
\begin{remark}
Указанные представления пригодны для графов и орграфов,
а после некоторой модификации~--- также и для псевдографов,
мультиграфов и гиперграфов.
\end{remark}

\end{frame}


\begin{frame}
\label{7.4.6.}\subsubsection{7.4.6. Обходы графов }\frametitle{7.4.6. Обходы графов (1/9)}
%\label{traversegraph}
Обход графа~--- это некоторое систематическое
перечисление его вершин и (или) рёбер. Наибольший интерес представляют
обходы, использующие локальную информацию (списки смежности). Среди всех
обходов наиболее известны поиск \odmkey{в~ширину}{Алгоритм
(algorithm)!поиска (search)!в ширину (breadth first)}
и \odmkey{в~глубину}{Алгоритм (algorithm)!поиска
(search)!в глубину (depth first)}. Алгоритмы поиска в
ширину и в глубину лежат в основе многих конкретных алгоритмов на графах
и изложены во множестве источников. Определение поиска в ширину и в глубину обычно даётся не
постулированием требуемых свойств обходов, но алгоритмически,
предъявлением конкретных алгоритмов. Это обстоятельство, а также
многочисленные технические детали изложения несколько затуманивают
основную идею, и может сложиться впечатление, что поиск в ширину и~в~глубину~--- это принципиально различные алгоритмы: например, поиск в~ширину~--- итеративный алгоритм, а~поиск в~глубину~--- рекурсивный.
На самом деле, хотя поиск в~ширину и~в~глубину действительно дают
различные обходы, по существу это две реализации одного алгоритма,
применяющие различные структуры данных.

Если структура данных $T$~--- это
стек (LIFO~--- Last In First Out), то обход называется
\odmkey{поиском в глубину}{Поиск (search)!в глубину (depth
first)}. Если $T$~--- это очередь (FIFO~--- First In First Out),
то обход называется \odmkey{поиском в ширину}{Поиск (search)!в
ширину (breadth first)}.

% как показано в~алгоритме 7.1
%%%%%%%%%%%%%%%%%%%%%%%%% Конец изменений 2012
\end{frame}


\begin{frame}
\frametitle{7.4.6. Обходы графов (2/9)}
%\begin{algorithm}
%\caption{Поиск в ширину и в глубину}
%\label{graphsearch}
\vspace{-4pt}
\begin{algorithmic}
\REQUIRE граф $G(V,E)$, представленный списками смежности $\Gamma$.
\ENSURE последовательность вершин обхода.
\STATE \textbf{for $v\in V$ do $x[v]\Assign 0$ end for} 
\COMMENT{\color{gray}\small вначале все вершины не отмечены }
\STATE\textbf{select $v\in V$} 
\COMMENT{\color{gray}\small  начало обхода~--- произвольная вершина }
\STATE $v\rightarrow T$ 
\COMMENT{\color{gray}\small  помещаем $v$ в структуру данных $T$\ldots }
\STATE $x[v]\Assign 1$ 
\COMMENT{\color{gray}\small  \ldots и отмечаем вершину $v$ }
\REPEAT
\STATE $u\leftarrow T$ 
\COMMENT{\color{gray}\small  извлекаем вершину из структуры данных $T$\ldots }
\STATE \textbf{yield $u$} 
\COMMENT{\color{gray}\small  \ldots и возвращаем её в качестве очередной пройденной }
%вершины }
\FOR{$w\in \Gamma(u)$}
  \IF{$x[w]=0$}
    \STATE $w\rightarrow T$ 
    \COMMENT{\color{gray}\small  помещаем $w$ в структуру данных $T$\ldots }
    \STATE $x[w]\Assign 1$ 
    \COMMENT{\color{gray}\small  \ldots и отмечаем вершину $w$ }
  \ENDIF
\ENDFOR
\UNTIL{$T=\emptyset$}
\end{algorithmic}
%\end{algorithm}

\end{frame}


\begin{frame}
\frametitle{7.4.6. Обходы графов (3/9)}
%%%%%%%%%%%%%%%%%%%%%%%%% Изменения 2012


\begin{remark}
Это изложение алгоритмов заметно отличается от обычного изложения.
Здесь не используется три цвета вершин, и при буквальном сравнении
протоколов работы алгоритмов поиска в глубину получаются различные
обходы. В результате у внимательного читателя может создаваться ошибочное
впечатление, что приведённый алгоритм <<неправильный>>.
На самом деле
это не так.
В случае поиска в ширину различий между приведённым алгоритмом 
и обычным
изложением алгоритма поиска в ширину нет.
Различие между приведённым алгоритмом в случае поиска в
глубину и обычным рекурсивным алгоритмом поиска в глубину заключается
только в порядке рассмотрения вершин $w$, смежных с текущей вершиной $u$.
В нашем варианте вершины $w$ заносятся в структуру данных $T$ в том
порядке, в каком вершины встречаются в списке смежности $\Gamma(u)$, и,
соответственно, рассматриваться будут в обратном порядке, если $T$
--- это стек. В рекурсивном варианте вершины будут рассматриваться в порядке следования в списке смежности. Таким образом, различие в том, как просматривать списки смежности: от начала к концу или от конца к началу. На свойства алгоритмов и получаемых обходов это различие не влияет.
\end{remark}

\end{frame}


\begin{frame}
\frametitle{7.4.6. Обходы графов (4/9)}
\vspace{-2pt}
\begin{example}
В следующей таблице показаны протоколы поиска
в глубину и~в~ширину для графа, диаграмма которого
приведена на рисунке в начале раздела.
Предполагается, что начальной является вершина 1.
%%%%%%%%%%%%%%%%%%%%%%%%% Изменения 2012
В первой паре столбцов в таблице приведён
протокол поиска в ширину, а во второй паре столбцов~--- в глубину.
Кроме того, для сравнения, в третьем столбце приведен протокол поиска в
глубину для случая, когда смежные вершины заносятся в стек в обратном порядке. Тем самым третий столбец показывает,
как ведёт себя рекурсивный вариант алгоритма поиска в глубину.
На рисунке сплошные стрелки с номерами
рядом с рёбрами показывают движение по графу при поиске в глубину, 
а  пунктирные~---  в ширину.
\end{example}

\vspace{-1.2\baselineskip}
\begin{columns}[t]
\begin{column}[T]{7cm}
\begin{center}
\begin{tabular}[t]{p{7mm}p{7mm}||p{7mm}p{7mm}|| p{7mm}p{7mm}}
$u$ & $T$           & $u$ & $T$      & $u$ & $T$\\
\hline
1 & 2,4         & 1 & 2,4         & 1 & 4,2        \\
2 & 4,3         & 4 & 2,3         & 2 & 4,3        \\
4 & 3           & 3 & 2           & 3 & 4          \\
3 & $\emptyset$ & 2 & $\emptyset$ & 4 & $\emptyset$ \\
\hline
\end{tabular}
\end{center}
%%%%%%%%%%%%%%%%%%%%%%%%% Конец изменений 2012
\end{column}
\begin{column}[T]{5cm}
\begin{figure}[h]
\begin{center}
\includegraphics[height=30mm]{odm07-12-new}
\end{center}
%{\caption{\hfil Диаграмма \hfil графа к~примеру обхода в ширину и~в~глубину}%
%\label{f0712}}
\end{figure}\hfil
%\end{tabular}
\end{column}
\end{columns}

%\vspace{-6pt}
\end{frame}


\begin{frame}
\frametitle{7.4.6. Обходы графов (5/9)}
%%%%%%%%%%%%%%%%%%%%%%%%% Изменения 2012

\begin{theorem}
Если граф G связен и конечен, то поиск в ширину и поиск в глубину
обходят все вершины по одному разу за время,
пропорциональное не более чем суммарному числу вершин и рёбер.
\end{theorem}

\begin{proof}
\proofsection{Единственность обхода вершины} Обходятся только
вершины из $T$. В $T$ попадают неотмеченные
вершины, при этом они отмечаются. Следовательно,
любая вершина будет обойдена не более одного раза.
\proofsection{Завершаемость алгоритма} Всего в $T$ может
попасть не более $p$ вершин. На каждом шаге одна вершина удаляется
из $T$. Следовательно, алгоритм завершит работу не более чем через
$p$ шагов.
\proofsection{Обход всех вершин} От противного.
Пусть алгоритм закончил работу и вершина $w$ не обойдена. Значит,
$w$ не попала в $T$. Следовательно, она не была отмечена. Отсюда
следует, что все вершины, смежные с $w$, не были обойдены и
отмечены. Аналогично, любые вершины, связанные с неотмеченными,
сами не отмечены (после завершения алгоритма). Но $G$ связен,
значит, существует путь $\lrle{v,w}$. Следовательно, исходная вершина $v$
не отмечена --- но она отмечена на первом шаге! 
\end{proof}


\end{frame}


\begin{frame}
\frametitle{7.4.6. Обходы графов (6/9)}
%\setcounter{corollary}{0}
\vspace{-2pt}
\begin{proof} (Продолжение)
\proofsection{Трудоёмкость алгоритма}
Инициализация (вторая строчка) выполняется за $O(p)$. В цикле каждая
вершина просматривается не более одного раза,
и для каждой просматриваемой
вершины её список смежности просматривается один раз. Следовательно, цикл
выполняется за $O(q)$. Значит трудоёмкость алгоритма $O(p+q)$.
\end{proof}

\smallskip
Таким образом, алгоритмы обхода графа в ширину и в глубину
обладают одними и теми же теоретическими свойствами.
Возникает вопрос: чем же отличаются алгоритмы \em{практически}
и в каких случаях их целесообразно применять.
Если целью применения алгоритмов является
построение обходов графа,
то алгоритмы практически не отличаются.
Если же целью применения алгоритмов является
поиск в графе какой-то конкретной вершины, 
то алгоритмы существенно различны с практической
точки зрения, причём эти различия 
зависят от особенностей графа поиска, и от того,
как соотносятся начальная и искомая вершины.  
Мы даём частичный ответ на вопрос 
о сравнении алгоритмов поиска в ширину и в глубину в
форме серии следствий. В следствиях 
приводятся оценки номера шага алгоритма,
на котором достигается некоторая искомая вершина.
Номер шага алгоритма можно трактовать как 
\em{время поиска}.  


\end{frame}


\begin{frame}
\frametitle{7.4.6. Обходы графов (7/9)}
\begin{corollary} \NB{[1]}
Пусть $(u_1,\dots ,u_i,\dots ,u_j,\dots ,u_p)$~{\upshape---} обход
(то~есть последовательность вершин) при поиске в ширину. Тогда
$
\A{i < j}{d(u_1,u_i) \leq d(u_1,u_j)}.
$
\end{corollary}
Другими словами, расстояние текущей вершины от начальной является
монотонно возрастающей функцией времени поиска в ширину,
 вершины обходятся в~порядке
возрастания расстояния от начальной вершины.

\begin{corollary} \NB{[2]}
Пусть $(u_1,\dots ,u_i,\dots ,u_p)$~{\upshape---} обход при поиске
в глубину. \\Тогда
$
\A{i\ge 1}{d(u_1,u_i)< i\leq p}.
$
\end{corollary}
%\addvspace{-6pt}
Другими словами, время поиска в глубину любой вершины не менее
расстояния от начальной вершины и не более общего числа вершин.
В худшем случае время поиска в глубину может быть максимальным,
независимо от расстояния до начальной вершины.

\begin{corollary} \NB{[3]}
Пусть $(u_1,\dots ,u_i,\dots ,u_p)$~{\upshape---} обход при поиске
в ширину, \\а $D(u_1,1)$, $D(u_1,2),\dots$~{\upshape---} ярусы графа
относительно вершины $u_1$. Тогда
$$
\A{i\ge 1}{\sum_{j=0}^{d(u_1,u_i)-1}\vert D(u_1,j)\vert < i\leq
\sum_{j=0}^{d(u_1,u_i)}\vert D(u_1,j)\vert}.
$$
\end{corollary}


\end{frame}


\begin{frame}
\frametitle{7.4.6. Обходы графов (8/9)}
Другими словами, время поиска в ширину ограничено снизу
количеством вершин во всех ярусах, находящихся на расстоянии меньшем,
чем расстояние от начальной вершины до текущей,
и ограничено сверху количеством вершин в~ярусах,
начиная с яруса текущей вершины и включая все меньшие ярусы.

%\medskip
\begin{corollary} \NB{[4]}
Пусть $(u_1,\dots ,u_i,\dots ,u_p)$~{\upshape---} обход при поиске
в ширину, \\а $(v_1,\dots $, $v_j,\dots ,v_p)$~{\upshape---} обход
при поиске в глубину, где $u_i=v_j$.
Пусть также вершина 
$u_i=v_j$~--- одна из вершин
самого внешнего яруса относительно 
вершины $u_1=v_1$, то есть
расстояние $d(u_1,u_i)$ максимально. Тогда математическое ожидание 
$\mathrm{E}(\frac{i}{j})\approx 2$.
\end{corollary}

Другими словами, поиск в глубину отдалённой вершины в среднем вдвое быстрее, чем поиск в
ширину.

\begin{remark}
В формулировке следствия используется
знак приближённого равенства,
поскольку граф может иметь особенности,
влияющие на применение алгоритма.
Например, к искомой вершине может вести 
цепь из промежуточных вершин степени два.
В таком случае поиск в глубину, если 
<<повезёт>>, найдёт искомую вершину за
минимальное время. Математическое ожидание 
равно двум, если особенностей нет, например,
граф регулярный. 
\end{remark}

\end{frame}


\begin{frame}
\frametitle{7.4.6. Обходы графов (9/9)}
\begin{example}


\vspace{-1.5\baselineskip}
\begin{center}
\hspace{1cm} \includegraphics[height=45mm]{7_4_6.jpg}
\end{center}
\end{example}

%%%%%%%%%%%%%%%%%%%%%%%%% Изменения 2012
\vspace{-0.75\baselineskip}
\begin{remark}
Алгоритм поиска в ширину и в глубину 
может быть применен также к несвязным графам.
Действительно, для этого достаточно применить алгоритм к одной компоненте
связности, потом к следующей и так далее. Алгоритм может быть
применён также к любым орграфам, в том числе к орграфам, которые не
являются сильно связными 
(см. \DMref{п.~8.5.1}{8.5.1.}). Для этого достаточно применить
алгоритм, начиная с~произвольного узла,
удалить все обойденные узлы с~инцидентными дугами, а затем
применить алгоритм к оставшейся части орграфа и так далее.
\end{remark}
%%%%%%%%%%%%%%%%%%%%%%%%% Конец изменений 2012
\end{frame}

\begin{frame}
\frametitle{Онтология теории графов}
\begin{center}
\includegraphics[height=8cm]{ODM7_4.png}
\end{center}

\end{frame}

\subsection{7.5. Графы и отношения}\label{7.5.}
%\label{7.5}
\begin{frame}
\frametitle{7.5. Графы и отношения }
Целью  заключительного раздела данной главы является установление связи теории
графов с другими разделами дискретной математики.

\medskip
\begin{tabular}{p{10mm}|p{60mm}|p{70mm}}
  № & Название параграфа
  & Ключевые термины и обозначения \\
  \hline
\DMref{7.5.1.}{7.5.1.} & Орграфы и бинарные отношения 
  & 
  {}\\
\DMref{7.5.2.}{7.5.2.} & Граф инциденций 
  & 
  {}
  \\
\DMref{7.5.3.}{7.5.3.}\BS & Гиперграфы и многоместные отношения &
{Деревья И/ИЛИ}
\\
\DMref{7.5.4.}{7.5.4.}\RS & Достижимость и частичное упорядочение &
{Достижимый узел}
\\   
\DMref{7.5.5.}{7.5.5.}\RS & Достижимость и транзитивное замыкание &
{}\\ 
\end{tabular}


\end{frame}


\begin{frame}
\label{7.5.1.}\subsubsection{7.5.1. Орграфы и бинарные отношения }\frametitle{7.5.1. Орграфы и бинарные отношения (1/6)}


Любой орграф $G(V,E)$, возможно, с петлями, но  без  кратных  дуг,
задаёт бинарное отношение~$E$ на множестве $V$. 
Обратно, любое бинарное отношение~$E$ на множестве $V$
определяет орграф $G(V,E)$:
пара элементов принадлежит отношению $(a,b)\in E\subset V\times V$
тогда и только тогда, когда в графе $G$ есть дуга $(a,b)$.

\begin{example}
% это определение должно быть в главе 2
%Рассмотрим отношение делимости $a\mid b$ на множестве натуральных чисел.
%Известно, что $\A{n\in\Bbb{N}}{1\mid n}$ и $\A{n\in\Bbb{N}}{n\mid n}$,
%то есть любое число делится на 1 и себя. Такие делители называются 
%\odmkey{тривиальными}{}. Прочие делители, если они есть,
%называются \odmkey{нетривиальными}{}, или \odmkey{собственными}{}.
На рисунке представлена диаграмма нагруженного орграфа, 
соответствующего отношению собственной 
делимости (см. \DMref{п.~2.4.1}{2.4.1.}), 
заданному на множестве
\\ $\{ 2, 3, 4, 5, 6, 7, 8, 9 \}$. 

\begin{center}
\includegraphics[scale=0.75]{7_5_1_1.jpg}
\end{center}
  
\end{example}

\end{frame}


\begin{frame}
\frametitle{7.5.1. Орграфы и бинарные отношения (2/6)}

Орграф, соответствующий рефлексивному отношению, в каждой вершине имеет петлю, а значит является псевдографом. Если же отношение антирефлексивно, то орграф вовсе не имеет петель. 
Может оказаться так, что для рефлексивного отношения рисовать петлю на диаграмме \em{каждой} вершины утомительно и неинформативно. Поэтому петли на диаграммах графов и
орграфов показывают, только если эта информация существенна.


\begin{columns}[t]
\begin{column}[T]{8cm}
\begin{example}
Рассмотрим отношение 
$\Set{(a,b)\in\Bbb{N}\times\Bbb{N}}{(a + b) \mod 4 = 0}$
на множестве $\{2, 3, 4, 5, 6, 7, 8, 9\}$.  
Это отношение симметрично, поэтому ему удобнее ставить в соответствие граф, а не орграф. 
Петли на диаграмме этого графа являются существенными, 
так как отношение не рефлексивно и не антирефлексивно.
\end{example}
\end{column}
\begin{column}[T]{6cm}
\begin{center}
\vspace{-1.3\baselineskip}
\includegraphics[width=5.5cm]{7_5_1_2.jpg}
\end{center}
\end{column}
\end{columns}


\end{frame}


\begin{frame}
\frametitle{7.5.1. Орграфы и бинарные отношения (3/6)}
Дополнение орграфа  $G(V,E)$ (обозначение $\overline{G}$)~--- это орграф, который соответствует
дополнению отношения $E$, то есть $\overline{G(V,E)}=
G(V,\overline{E})$.
Изменение направления всех дуг в орграфе соответствует
обратному отношению. Другими словами,
\odmkey{обратный орграф}{Орграф (digraph)!обратный (reverse)}~--- это орграф, соотвествующий обратному отношению, то есть $(G(V,E))^{-1}=
G(V,E^{-1})$. 

\smallskip
Степени бинарного отношения $E\subset V\times V$ соответствуют
путям в орграфе. Степень $E^1$ соответствует путям длины $1$, 
то есть дугам,
$E^2$ соответствует путям длины $2$ и т.~д.

\smallskip
Если отношение антисимметрично, 
то соответствующий ему орграф не имеет встречных дуг.
Антисимметричными являются, в частности, отношения порядка, 
которые рассматриваются в следующем параграфе с точки зрения теории графов.
Если же отношение симметрично, то всякая дуга орграфа имеет встречную,
и направление дуги не имеет значения. Таким образом, всякому симметричному отношению соответствует неориентированный граф, и обратно, 
граф (неориентированный) соответствует симметричному отношению.

\begin{remark}
Понятие обратного графа не используется,
поскольку симметричное отношение на множестве
совпадает со своим обратным отношением.
\end{remark}

\end{frame}


\begin{frame}
\frametitle{7.5.1. Орграфы и бинарные отношения (4/6)}
\begin{columns}[t]
\begin{column}[T]{8cm}
\begin{example}
Напомним, что отношение взаимной простоты  
$\Set{(a,b)\in \Bbb{N}\times\Bbb{N}}{\gcd(a,b)=1}$ на множестве натуральных чисел
симметрично.
На рисунке представлена диаграмма графа, соответствующего отношению взаимной 
простоты на множестве $\{ 2, 3, 4, 5, 6, 7, 8, 9 \}$.
\end{example}
\end{column}
\begin{column}[T]{7cm}
\begin{center}
\vspace{-2.0\baselineskip}
\includegraphics[width=4.5cm]{7_5_1_3.jpg}
\end{center}
\end{column}
\end{columns}

\smallskip
Полный граф соответствует универсальному отношению.
 
\medskip
Исходя из названий <<граф>> и <<ориентированный граф>>
в первый момент может показаться, 
что класс графов шире класса орграфов.
Однако на самом деле, поскольку симметричные отношения являются
подклассом класса всех отношений,
класс графов является подклассом класса орграфов.
Так исторически сложилась терминология.
  

\end{frame}


\begin{frame}
\frametitle{7.5.1. Орграфы и бинарные отношения (5/6)}
В то же время, следует обратить внимание, что
мультиорграф  не может задавать бинарное отношение,
так как пары в подмножестве прямого произведения двух множеств не повторяются.
Иначе говоря, между языками мультиграфов  
и бинарных отношений нет прямого соответствия.

\medskip
Функция является частным случаем отношения 
и может быть представлена орграфом. 
Рассмотрим функцию $\Map{f}{A}{B}$, 
ей соответствует двудольный орграф $G(A,B,E)$, 
для которого $V=A\cup B$ и $E=\Set{(a,b)\in A \times B}{ b = f(a)}$. 
Множество всех вершин орграфа $G$, 
имеющих исходящие дуги, есть область определения функции, 
а множество всех вершин, имеющих хотя бы одну входящую дугу,~--- область значений функции. 
При этом в силу однозначности, 
любой узел орграфа $G$ имеет не более одной исходящей дуги.
Если функция является функцией \em{на} множестве $V$,
$\Map{f}{V}{V}$, то диаграмма орграфа $G(V,E)$
может иметь довольно причудливую форму.
\end{frame}


\begin{frame}
\frametitle{7.5.1. Орграфы и бинарные отношения (6/6)}

\begin{examples}
На рисунке слева представлена диаграмма орграфа для 
функции $x\MOD 9 +1$, а справа~--- 
функции  $x\MOD 4 +1$, действующих 
на множестве $\{1, 2, 3, 4, 5, 6, 7, 8, 9 \}$.
Так как функции тотальные, каждый узел 
 имеет ровно одну исходящую дугу.

%\vspace{-2pt} 
\begin{center}
\includegraphics[width=15cm]{7_5_1_3(2).jpg}
\end{center}
\end{examples}

\begin{remark}
Диаграммы орграфов пригодны 
для представления только таких функций,
у которых область определения и область значений существенно конечны. 
Для бесконечно заданных функций,
в частности, непрерывных функций вещественной переменной, 
традиционным и наиболее используемым графическим средством
является представление графика функции в виде линии, 
которая к тому же наглядно иллюстрирует многие свойства этой функции.
\end{remark}



\end{frame}


\begin{frame}
\label{7.5.2.}\subsubsection{7.5.2. Граф инциденций }\frametitle{7.5.2. Граф инциденций (1/2)} 
Отношение инцидентности между вершинами и рёбрами
(узлами и дугами) также может быть представлено графом
(орграфом) \odmkey{инциденций}{Граф (graph)!инциденций (of incidence)}.
Для графов отношение инцидентности симметрично,
поэтому граф инциденций является двудольным графом
(а не орграфом).

\begin{example}
На рисунке представлены диаграммы графа (слева)
и его графа инциденций (справа), 
для примера из \DMref{п.~7.4.1}{7.4.1.}. 

\vspace{-0.5\baselineskip}
\begin{center}
\includegraphics[scale=0.7]{7_5_1_4.jpg}
\end{center}
  
\end{example}

Чтобы не потерять информацию о направлении дуг,
для орграфов обычно считают,
что узел инцидентен исходящим дугам
и входящие дуги инцидентны узлу.


\end{frame}


\begin{frame}
\frametitle{7.5.2. Граф инциденций (2/2)}  

Таким образом, для орграфов отношение инциденции задаётся орграфом.

\begin{example}
На рисунке представлены диаграммы орграфа (слева)
и его орграфа инциденций (справа), для примера из \DMref{п.~7.4.1}{7.4.1.}. 

\vspace{-0.5\baselineskip}
\begin{center}
\includegraphics[scale=0.7]{7_5_1_5.jpg}
\end{center}
  
\end{example}  
   
\begin{remark}
В графе инциденций степени вершин сохраняются, а степени 
рёбер равны 2. В орграфе инциденций полустепени узлов 
сохраняются, а у дуг полустепени равны 1.
\end{remark}

Граф инциденций оказывается особенно полезен при 
рассмотрении отношений, устроенных  более сложно, нежели 
бинарные отношения на конечном множестве.
\end{frame}


\begin{frame}
\label{7.5.3.}\subsubsection{7.5.3. Гиперграфы и многоместные отношения }\frametitle{7.5.3. Гиперграфы и многоместные отношения (1/3)}  
В гиперграфах рёбра связывают произвольные 
непустые подмножества множества вершин. Для 
гиперграфов нет единой общепринятой нотации диаграмм.
На рисунке приведены две наиболее распространенные.
Слева рёбра изображены в виде овалов, охватывающих
смежные вершины. При большом количестве вершин и рёбер
такая диаграмма трудно читается. Можно
рисовать рёбра, имеющие несколько концов (см. рисунок в центре), 
с помощью специальных <<ромбиков>>, которые считаются частью
линии, а не фигурой. Фактически, этот способ 
задаёт диаграмму графа инциденций гиперграфа (см. рисунок справа).

\vspace{-4pt}      
\begin{center}
\includegraphics[scale=0.65]{7_5_1_6.jpg}
\end{center}


\end{frame}


\begin{frame}
\frametitle{7.5.3. Гиперграфы и многоместные отношения (2/3)}

Таким образом, граф инциденций гиперграфа~---
это двудольный граф без ограничений на степени вершин.

\smallskip
Аналогичные построения легко провести для гиперорграфов.

\medskip
\begin{digress}
Не следует думать, что гиперграфы и гиперорграфы~---
это какая-то экзотика. \\ Напротив, такие конструкции часто встречаются в
реальных приложениях. Например, 
\odmkey{игровые деревья}{Дерево (tree)!игровое (game)}, известные также как 
\odmkey{ деревья И/ИЛИ}{Дерево (tree)! И/ИЛИ (and-or)}, которые являются 
основной структурой данных в игровых программах
в \odmkey{символическом искусственном интеллекте}{Искусственный интеллект (artificial intelligence)!символический (symbolic)},
фактически являются гиперорграфами специального вида. 
   
Другим примером являются блок-схемы, диаграммы UML,
диаграммы <<сущность-связь>> и тому подобное.   
\end{digress}

\end{frame}


\begin{frame}
\frametitle{\large 7.5.3. Гиперграфы и многоместные отношения (3/3)}

\large
Таким образом, имеется полная аналогия между орграфами и бинарными
отношениями~--- фактически, это один и тот же класс объектов,
только описанный разными средствами. Отношения (в частности,
функции) являются базовым средством для построения подавляющего
большинства математических моделей, используемых при решении
практических задач. С другой стороны, графы допускают наглядное
представление в виде диаграмм. Этим обстоятельством объясняется
широкое использование графов в программировании.
%\vspace{-6pt}

  
\end{frame}


\begin{frame}
\label{7.5.4.}\subsubsection{7.5.4. Достижимость и частичное упорядочение }\frametitle{7.5.4. Достижимость и частичное упорядочение (1/3)}
%\label{8.5.1} 
В качестве примера связи между орграфами и бинарными
отношениями рассмотрим отношения частичного порядка с точки зрения
теории графов. Узел $u$ в орграфе $G(V,E)$
\odmkey{достижим}{Вершина (vertex)!достижимая (reachable)} из
узла  $v$, если  существует путь из $v$ в $u$. Путь из $v$ в $u$
обозначим $\lrle{\stackrel{\rightarrow}{v,u}}$. Отношение
достижимости  можно представить   матрицей \textbf{$T:$ array
$[1..p,1..p]$  of $0..1$}, где $T[i,j]=1$, если узел $v_j$
достижим из узла $v_i$, и~$T[i,j]=0$, если узел $v_j$ недостижим
из узла $v_i$. 

\smallskip
Рассмотрим отношение строгого частичного порядка
$\succ$, которое характеризуется следующими аксиомами:

\begin{tabbing}
  1.\quad \= антирефлексивность:\quad  \= $\A{v\in V}{\neg(v\succ v)}$;  \kill
  % \> for next tab, \\ for new line...
  1. \> Антирефлексивность:  \> $\A{v\in V}{\neg(v\succ v)}$. \\
  2. \> Транзитивность:  \> $\A{u,v,w}{(v\succ w\And w\succ u\Impl v\succ u)}$. \\
  3. \> Антисимметричность:  \> $\A{u,v}{\neg(u\succ v\And v\succ u)}$.
\end{tabbing}
Отношению строгого частичного порядка $\succ\subset V\times V$
можно сопоставить орграф $G(V,E)$, в котором $a\succ b \Equiv
(a,b)\in E$.

\end{frame}


\begin{frame}
\frametitle{\large 7.5.4. Достижимость и частичное упорядочение (2/3)}

\begin{theorem} \NB{[1]}
Если отношение $E$ есть строгое  частичное  упорядочение,
то орграф $G(V,E)$ не имеет контуров.
\end{theorem}
\begin{proof}
От противного. Пусть в $G$ есть контур.
Рассмотрим любую дугу $(a,b)$ в этом контуре.
Тогда имеем $a\succ b$, но $b\succ a$ по
транзитивности, что противоречит антисимметричности упорядочения.
\end{proof}


\smallskip
\begin{theorem} \NB{[2]}
Если орграф $G(V,E)$ не  имеет  контуров,  то  отношение достижимости
есть
строгое частичное упорядочение.
\end{theorem}
\begin{proof}
\proofsection{\!Антирефлексивность\!} Нет контуров, следовательно,
нет петель. \proofsection{\!Транзитивность\!} Если существуют пути
из $v$ в $w$ и из $w$ в $u$, то существует и~путь из $v$ в $u$.
\proofsection{\!Антисимметричность\!} От противного.
 Пусть
$\E{u,v}{u\succ v\And v\succ u}$,
то~есть существует путь $\lrle{\stackrel{\rightarrow}{v,u}}$ из $v$ в $u$
 и путь $\lrle{\stackrel{\rightarrow}{u,v}}$ из $u$ в
$v$.
Следовательно, существует контур вида
$\lrle{\stackrel{\rightarrow}{u,v}}+\lrle{\stackrel{\rightarrow}{v,u}}$,
что противоречит условию.
\end{proof}
\end{frame}


\begin{frame}
\frametitle{7.5.4. Достижимость и частичное упорядочение (3/3)}
\begin{theorem} \NB{[3]}
Если орграф не имеет контуров, то в нем
есть узел, полустепень исхода которого равна 0.
\end{theorem}
\begin{proof}
От противного. Пусть такого узла нет, тогда для любого узла найдётся
исходящая дуга. Пойдём по этой дуге.
Ввиду конечности числа узлов, рано или поздно 
попадём в узел, в котором уже были.
 Следовательно, имеем
контур~--- противоречие.
\end{proof}
\begin{remark}
Эта теорема позволяет найти минимальный элемент в конечном частично
упорядоченном множестве, который требуется в алгоритме топологической
сортировки (\DMref{п.~1.8.4}{1.8.4.}).
А именно, минимальный элемент~--- это сток, то~есть
узел, которому в матрице смежности соответствует нулевая строка.
\end{remark}
\end{frame}


\begin{frame}
\subsubsection{7.5.5. Достижимость и транзитивное замыкание}\label{7.5.5.}
\frametitle{7.5.5. Достижимость и транзитивное замыкание}
Если $E$~--- бинарное отношение на $V$, то транзитивным
замыканием  $E^+$ на $V$ будет
отношение достижимости на орграфе $G(V,E)$.
\begin{theorem}
Пусть $M$~{\upshape---} матрица смежности орграфа $G(V,E)$. Тогда
$M^k [i,j]  
= 1$ в том и только в том случае, если существует
путь длиной $k$ из узла $v_i$ в~узел $v_j$.
\end{theorem}
\begin{proof}
Индукция по $k$. База: $k=1$, $M^1=M$~--- пути длины 1. Пусть $M^{k-1}$
содержит пути длины $k-1$. Тогда
$
M^k[i,j]=\BigOr\limits^p_{l=1} \left(M^{k-1}[i,l]\And M[l,j]\right),
$
то~есть путь длины $k$ из узла $i$ в узел $j$ существует тогда и
только тогда, когда найдётся узел $l$ такой, что  существует путь
длины $k-1$ из $i$ в $l$ и~дуга $(l,j)$, то~есть
\\
$\E{\lrle{\stackrel{\rightarrow}{i,j}}}
{\vert\lrle{\stackrel{\rightarrow}{i,j}}\vert = k}
\Equiv \E{l}{\E{\lrle{\stackrel{\rightarrow}{i,l}}}
{\vert\lrle{\stackrel{\rightarrow}{i,l}}\vert=k-1 \And (l,j) \in E}}$.
\end{proof}

Если $T$~--- матрица достижимости, то
очевидно, что $ T = {\BigOr^{p-1}_{k=1}}{M^k}$. 
Трудоёмкость прямого вычисления по этой формуле составит $O(p^4)$.
Матрица достижимости $T$ может быть вычислена по матрице смежности
$M$ алгоритмом Уоршалла 
(\DMref{п.~1.5.2}{1.5.2.}) за $O(p^3)$.
\end{frame}


%\begin{frame}
%\frametitle{7.5.6. Транзитивная редукция}
%
%{\color{blue}
%Тема для сочинения: связь с диаграммами Хассе и прочее.
%Возможно, алгоритм, построения подграфа с минимальным 
%достаточным числом рёбер (дуг).
%
%Особая тема для сочинения: 
%транзитивное замыкание на гиперграфах и структурный синтез программ. 
%}
%
%\end{frame}

%\printindex




% ===== tex/DM2022_8.tex =====
\section{8. Связность}

\begin{frame}
\frametitle{8. Связность}

В этой главе обсуждается важное для приложений понятие связности, доказывается фундаментальная 
теорема~--- теорема Менгера
и рассматриваются
наиболее популярные алгоритмы поиска кратчайших путей.                                                                                                                                                                                                                                               
\end{frame}                                                                                                       


\subsection{8.1. Компоненты связности}
\label{8.1.}

\begin{frame}
\frametitle{8.1. Компоненты связности}
В русском языке есть как слово <<компонент>> мужского рода,
так и слово <<компонента>> женского рода, оба варианта допустимы.
В современной языковой практике чаще используется слово мужского
рода.
Однако исторически сложилось так, что <<компонента связности>> имеет
женский род, и в данном случае мы подчиняемся традиции.

\medskip
\begin{tabular}{p{9mm}|p{50mm}|p{84mm}}
  № & Название параграфа
  & Ключевые термины и обозначения \\
  \hline
\DMref{8.1.1.}{8.1.1.} & Объединение графов и компоненты связности 
  & 
  {} \\
\DMref{8.1.2.}{8.1.2.}\RS & Точки сочленения, мосты и блоки 
  & 
  {}
  \\
\DMref{8.1.3.}{8.1.3.} & Свойства блоков &
{ }\\
\DMref{8.1.4.}{8.1.4.} & Вершинная и рёберная связность &
{$\varkappa(G)$~--- вершинная связность;
$\lambda(G)$~--- рёберная связность 
}
\\
\DMref{8.1.5.}{8.1.5.}\RS & Оценка числа рёбер &
{}\\   
\end{tabular}
\end{frame}

\begin{frame}
\label{8.1.1.}\subsubsection{8.1.1. Объединение графов и компоненты связности }\frametitle{8.1.1. Объединение графов и компоненты связности (1/2)}

\begin{lemma}
В нетривиальном связном графе $G(V, E)$
для любого разбиения множества вершин $V$ на непустые блоки $U$ и $W$
существует ребро, инцидентное верши\-нам из разных блоков.
% Эта хитрая формула сейчас не нужна. Пусть полежит. ФАН
\begin{align*}
\forall\, U,W \subset V\; (
    U\ne\emptyset \And
    W\ne\emptyset &\And
    U\cap W= \emptyset \And
    U\cup W=V \Impl\\
&\Impl \E{(u, w)\in E}{u\in U \And w\in W}).
\end{align*}
\end{lemma}

\begin{proof}
Рассмотрим вершины $u\in U$, $w\in W$ 
и цепь $\langle u, w \rangle$, 
которые существуют по условиям леммы. 
Если цепь $\langle u, w \rangle$
состоит из одного ребра, то это ребро~--- требуемое.
Если цепь состоит из нескольких рёбер,
то рассмотрим вершину $v$ на цепи $\langle u, w \rangle$,
смежную с вершиной $u$. Если $v \in W$, то ребро $(u,v)$~--- требуемое.
В противном случае отщепим от цепи $\langle u, w \rangle$ ребро $(u,v)$,
положим $u\Assign v$ и снова рассмотрим цепь $\langle u, w \rangle$,
которая удовлетворяет условиям. Будем продолжать процесс отщеплений,
пока не получим требуемое ребро.
Процесс отщеплений закончится, потому что исходная цепь конечна.
\end{proof}
%%%%%%%%%%%%%%%%%%%%%%%%% Конец изменений 2012
\end{frame}

\begin{frame}
\frametitle{8.1.1. Объединение графов и компоненты связности (2/2)}
\begin{theorem}
Граф связен тогда и только тогда, когда его нельзя представить в~виде
объединения двух графов.
\end {theorem}


\begin{proof}
{[$\Impl$]} От противного. Пусть $k(G)=1$ 
 и граф $G$ является объединением двух графов, то~есть $G=G_1(V_1,E_1)\cup
G_2(V_2,E_2)$, где $V_1\cap V_2=\emptyset$, $E_1\cap
E_2=\emptyset$, $V_1\not=\emptyset$, \\$V_2\not=\emptyset$.
Возьмём $v_1\in V_1$, $v_2\in V_2$. Тогда
$\E{\lrle{v_1,v_2}}{v_1\in V_1\And v_2\in V_2}$. В этой цепи
по лемме
$\E{e=(a,b)}{a\in V_1\And b\in V_2}$. Но тогда
$e\not\in E_1$, $e\not\in E_2$, следовательно, $e\not\in E$.

\proofsection{$\Impll$} От противного. Пусть $k(G)>1$. Тогда
существуют две вершины, $u$ и~$v$, не соединённые цепью.
Следователь\-но, вершины $u$ и $v$ принадлежат разным компонентам
связности. Положим $G_1$~--- компонента связности для $u$, а
$G_2$~--- все остальное. 
Имеем $G=G_1\cup G_2$.
\end{proof}

%\vspace*{-1.0\baselineskip}

\begin{remark}
Несвязный граф всегда можно представить как объединение
связных компонент. Эти компоненты можно рассматривать независимо.
Поэтому во многих случаях можно без ограничения общности предполагать,
что рассматриваемый граф связен.
\end{remark}
\end{frame}

%В учебнике этого нет.
%%\begin{frame}
%%\label{8.1.2.}\subsubsection{8.1.2. Связность и простые циклы }\frametitle{8.1.2. Связность и простые циклы (1/3)}
%%
%%Очевидно, что все вершины в любом цикле связаны, и всякий цикл принадлежит
%%одной компоненте связности.
%%Простые циклы дополнительно обладают рядом полезных характеристических свойств. 
%%
%%\begin{theorem}
%%Граф $G(V,E)$ является простым циклом тогда и только тогда, 
%%когда $G(V,E)$ ~--- связный граф и $\A{ v\in V} {d(v)=2}$.
%%\end{theorem}
%%
%%\begin{proof}
%%\proofsection{$\Longleftarrow$}
%%Пусть $G$ связен и все вершины имеют степень 2. 
%%%Покажем, что цепь, содержащая все его вершины, существует, простая, замкнутая и единственная с точностью до выбора начала и направления (а значит $G$ ~--- простой цикл).
%%Построим простую замкнутую цепь, содержащую все вершины и все рёбра графа.
%%Рассмотрим произвольную вершину $v_1\in V$, ей инцидентны два ребра. 
%%Включим в цепь вершину $v_1$, 
%%выйдем из вершины по любому ребру и попадём в вершину $v_2$ по ребру $(v_1,v_2)$. 
%%Включим в цепь ребро $(v_1,v_2)$ и вершину $v_2$.
%%Далее выйдем из вершины $v_2$ через единственное непройденное ребро в вершину $v_3$.
%%Так будет продолжаться единственным образом, пока из $v_p$ мы не вернёмся в $v_1$.
%%Это произойдет именно из вершины $v_p$, 
%%поскольку иное противоречило бы связанности $v_1$ и $v_p$.
%%%,
%%% и обязательно произойдет, поскольку у всех остальных вершин уже были использованы все их ребра.
%%%Простота получившейся цепи получается приведенного алгоритма автоматически, а её единственность достаточно очевидна.
%%\proofsection{$\Longrightarrow$} При обходе  простого цикла 
%%ровно один раз выйдем из каждой вершины и ровно один раз войдём, 
%%то есть $\A{ v\in V} {d(v)=2}$.
%%\end{proof}
%%
%%
%%
%%\end{frame}
%%
%%
%%\begin{frame}
%%\frametitle{8.1.2. Связность и простые циклы (2/3)}
%%\begin{remark}
%%Связность существенна, поскольку
%%любая совокупность двух и более простых циклов
%%удовлетворяет условию $\A{ v\in V} {d(v)=2}$, 
%%но не является простым циклом. 
%%\end{remark}
%%
%%\begin{theorem}
%%Граф $G(V,E)$ является простым циклом  тогда и только тогда, 
%%когда граф $G(V,E)$~--- связный граф с одним циклом и без висячих вершин.
%%\end{theorem}
%%
%%\begin{proof}
%%\proofsection{$\Longleftarrow$} 
%%В графе $G$ есть цикл, значит $p(G)\ge 3$. 
%%Граф $G$ связен, значит в нём нет вершин нулевой степени. 
%%Так как в графе $G$ нет висячих вершин, в нём нет вершин степени 1.
%%Имеем $\A{ v\in V}{d(v)\ge 2}$.
%%Единственный цикл является простым,
%%поскольку всякий не простой цикл состоит по меньшей мере из двух простых циклов.
%%
%%Предположим, что $\exists u\in V (d(u)>2)$. Обозначим инцедентные $u$ рёбра $e_1...e_n, n\ge 3$.
%%Так как степени всех вершин в графе $G$ больше единицы, операцию перехода от одной из них к другой можно будет повторить бесконечное количество раз, но $p(G)<\infty$, значит в какой-то момент мы попадём в цикл. То есть начав путь из вершины $u$, мы независимо от выбора  начального ребра придём к циклу. Обозначим первые повторившиеся вершины на каждом из путей $w_1...w_n$ соответственно. Так как по условию в графе $G$ ровно один цикл, все пути приведут к одному и тому же циклу $Z$, содержащему вершины $w_1...w_n$.
%%\end{proof}
%%
%%\end{frame}
%%
%%
%%\begin{frame}
%%\frametitle{8.1.2. Полезные теоремы о простом цикле (2*/2)}
%%
%%\begin{proof}(Продолжение) 
%%Если вершина $u$ содержится в цикле $Z$, то она находится между двумя вершинами $w_i, w_j, i,j\in[1, n], i \ne j$ и от неё идёт путь к некоторой вершине $w_k, k\in[1, n], k\not= i \And k\not= j$. Если вершина $u$ не принадлежит циклу $Z$, то подойдут любые различные номера $i,j,k$ меньшие $n+1$. Рассмотрим два цикла: $Z$ и $\langle u,w_i,w_k,u \rangle$. Они различны, так как только один из них содержит вершину $w_j$. Но по условию цикл в графе $G$ один, значит предположение существования вершины $u, d(u)\ge 2$ ведёт к противоречию. Имеем $\forall v\in V (d(v)\ge 2\And d(v)<3)$. То есть $\forall v\in V (d(v)= 2)$, по теореме 1 граф $G$ ~--- простой цикл.
%%
%%$[\Longleftarrow]$ Граф $G(V,E)$ ~--- простой цикл, значит он связный и в нём содержится ровно один цикл. По теореме 1: $\forall v\in V (d(v)=2)$, то есть висячих вершин в графе $G$ нет (степень висячей вершины ~--- 1, а такие  в $G$ отсутствуют).
%%
%%
%%\end{proof}
%%
%%\end{frame}
%%
%%
\begin{frame}
\label{8.1.2.}\subsubsection{8.1.2. Точки сочленения, мосты и блоки }\frametitle{8.1.2. Точки сочленения, мосты и блоки (1/5)}
%\label{s821}

%\vspace{-2pt}
Вершина графа называется \odmkey{точкой сочленения}{Точка
сочленения (articulation point)}, или \odmkey{разделяющей
вершиной}{Вершина (vertex)!разделяющая (cut)}, если её
удаление увеличивает число компонент связности.
\odmkey{Мостом}{Мост (bridge)} называется ребро, удаление которого
увеличивает число компонент связности. \odmkey{Блоком}{Блок
(block)!графа (of graph)} называется связный граф,
 не имеющий точек сочленения.

\begin{example}
В графе, диаграмма которого представлена
на рисунке,

\begin{enumerate}
\item[1)] вершины $u,v$~--- точки сочленения, и других точек сочленения нет;%
\item[2)] ребро $x$~--- мост, и других мостов нет;%
\item[3)] правильные подграфы%
$\{a,b,w\}$, $\{b,u,w\}$, $\{a,b,u,w\}$, $\{c,d,v\}$,
$\{e,f,v\}$~--- блоки, и других 
нетривиальных блоков нет.
Связные подграфы с двумя вершинами являются тривиальными блоками (\DMref{п.~8.1.3}{8.1.3.}).
\end{enumerate}

\vspace{-0.5\baselineskip}
\begin{figure}[h]
\begin{center}
\includegraphics[width=75mm]{odm08-01}
\end{center}
%\caption{Точки сочленения, мосты и блоки}
%\label{f0801}
\end{figure}
\end{example}

\end{frame}

\begin{frame}
\frametitle{8.1.2. Точки сочленения, мосты и блоки (2/5)}

\begin{lemma}
В любом нетривиальном графе есть по крайней мере две вершины,
которые не являются точками сочленения.
\end{lemma}

\begin{proof}
Например, это концы диаметра.
\end{proof}

\begin{theorem}
Пусть $G(V,E)$~{\upshape---} связный граф и $v\in V$. Тогда
следующие утверждения эквивалентны:

\vspace{-2pt}
\begin{enumerate}
\item $v$~{\upshape---} точка сочленения.

\item $\E{u,w\in V}{{u}\not= {w}\And \A{\divC{u}{w}_G}{v\in
\divC{u}{w}_G}}$.

\item 
$\exists\,U,W\;
(U\cap W = \emptyset \And U\cup W = V-v \And$
$\forall\,u\in U,w\in W\;(\forall\,{\divC{u}{w}}_{G}\;(v\in {\divC{u}{w}}_G)))
$.
%$\E{U,W}
%{U\cap W = \emptyset \And U\cup W = V-v \And
%\A{u\in U,w\in W}{\A{{\divC{u}{w}}_{G}}{v\in {\divC{u}{w}}_G}}}
%$.
\end{enumerate}
\end{theorem}


%\vspace{-6pt}
\begin{proof}
\proofsection{1$\Impl$3}
Рассмотрим $G'\Assign G-v$. Граф $G'$ не связен, $k(G')>1$,
следовательно, $G'$ имеет
по крайней мере две компоненты связности,
то~есть
$V-v = U\cup W$,
где  $U$~--- множество вершин одной из
компонент связности, а $W$~--- все остальные
вершины.
Пусть $u\in U$, $w\in W$.
Тогда $\Not\exists {\divC{u}{w}}_{G'}$.
Но
$k(G)=1$, следовательно, $\exists {\divC{u}{w}}_G$
и, значит,
$\A{{\divC{u}{w}}_G}{v\in{\divC{u}{w}}_G}$.
\proofsection{3$\Impl$2} Тривиально.
\end{proof}

\end{frame}

\begin{frame}
\frametitle{8.1.2. Точки сочленения, мосты и блоки (3/5)}

\begin{proof} (Продолжение)
\proofsection{2$\Impl$1} Рассмотрим $G'\Assign G-v$. Имеем
$\Not\exists {\divC{u}{w}}_{G'}$.
Следовательно, $k(G')>1$, откуда $v$~--- точка сочленения.%\qed
\end{proof}

\smallskip
\begin{example}
На рисунке в начале параграфа: $v$~--- точка сочленения, при этом
$U=\{a,b,u,w\}$, $W=\{c,d,e,f\}$.  
\end{example}

\smallskip
\begin{corollary}
Если вершина инцидентна мосту и не является висячей,
то она является точкой сочленения.
\end{corollary}
\begin{proof}
Пусть $(u,v)$~--- мост и $d(v)>1$. Тогда $v$ смежна с $w$, причём
$w \neq u$. Далее разбор случаев. Если $\E{\divC{u}{w}}{v\notin
\divC{u}{w}}$, то ребро $(u,v)$ принадлежит циклу, что
противоречит тому, что $(u,v)$~--- мост. Если же
$\Not\E{\divC{u}{w}}{v\notin \divC{u}{w}}$, то
$\A{\divC{u}{w}}{v\in \divC{u}{w}}$ и, значит, $v$~--- точка
сочленения по пункту 2 теоремы.
\end{proof}
%
%\vspace*{-.5\baselineskip}

\medskip
\begin{remark}
Отнюдь
не всякая точка сочленения является концом моста.
\end{remark}
\begin{example}
В условиях предыдущего примера, при удалении 
моста $x$, вершина $v$ остаётся точкой сочленения
и не является концом  моста.
\end{example}


\end{frame}

\begin{frame}
\frametitle{8.1.2. Точки сочленения, мосты и блоки (4/5)}

\begin{corollary}
Если граф $G(V, E)$ связен, а вершина $v$ не является точкой сочленения,
то для любых двух других вершин, $u$ и $w$, существует простая
$\langle u, w \rangle$-цепь, не содержащая вершину $v$.
\end{corollary}
\begin{proof}
Данное утверждение является отрицанием пункта 2 предыдущей теоремы.
\end{proof}

\smallskip
\begin{theorem}
Пусть $G(V,E)$~{\upshape---} связный граф и $x\in E$. Тогда
следующие утверждения эквивалентны:

\begin{enumerate}
\item $x$~{\upshape---} мост.
\item $x$ не принадлежит ни одному простому циклу.
\item $\E{u,w\in V}{\A{{\divC{u}{w}}_G}{x\in{\divC{u}{w}}_G}}$.
\item
$\exists\,U,W\,
(U\cap W = \emptyset\And U\cup W = V\And$
$\forall\, u\in U\, (\forall\, w\in W\,
(\forall\, \divC{u}{w}_G\, (x\in\divC{u}{w}_G))))$.
%$\E{U,W}
%{U\cap W = \emptyset\And U\cup W = V\And
%\A{u\in U}{\A{w\in W}
%{\A{{\divC{u}{w}}_G}{x\in\divC{u}{w}}_G}}}$.
\end{enumerate}
\end{theorem}

\begin{proof}
\proofsection{1$\Impl$2}
Рёбра любого цикла, очевидно, не являются мостами.
\proofsection{2$\Impl$1}
Если ребро $x$ не принадлежит ни одному (простому) циклу, то концы $x$
не связаны в графе $G'\Assign G-x$, а значит ребро $x$~--- мост.
\end{proof}
\end{frame}

\begin{frame}
\frametitle{8.1.2. Точки сочленения, мосты и блоки (5/5)}
\begin{proof}
(Продолжение) 
\proofsection{1$\Impl$4}
Рассмотрим $G'\Assign G-x$. Граф $G'$ не связен, $k(G')>1$,
следовательно, $G'$ имеет
по крайней мере две компоненты связности,
то~есть
$V = U\cup W$,
где  $U$~--- множество вершин одной из
компонент связности, а $W$~--- все остальные
вершины.
Пусть $u\in U$, $w\in W$.
Тогда $\Not\exists {\divC{u}{w}}_{G'}$.
Но
$k(G)=1$, следовательно, $\exists {\divC{u}{w}}_G$
и, значит,
$\A{{\divC{u}{w}}_G}{x\in{\divC{u}{w}}_G}$.
\proofsection{4$\Impl$3} Тривиально.
\proofsection{3$\Impl$1} Рассмотрим $G'\Assign G-x$. Имеем
$\Not\exists {\divC{u}{w}}_{G'}$.
Следовательно, $k(G')>1$, откуда $x$~--- мост.
\end{proof}

\begin{example}
В графе, диаграмма которого представлена
на рисунке, только ребро $x$ не входит в циклы и является мостом.

\vspace{-0.55\baselineskip}
\begin{figure}[h]
\begin{center}
\includegraphics[width=70mm]{odm08-01}
\end{center}
%\caption{Точки сочленения, мосты и блоки}
%\label{f0801}
\end{figure}
\end{example}
\end{frame}

\begin{frame}
\label{8.1.3.}\subsubsection{8.1.3. Свойства блоков }\frametitle{8.1.3. Свойства блоков (1/7)}
%%%%%%%%%%%%%%%%%%%%%%%%% Изменения 2012
Графы $K_1$ и $K_2$ формально можно считать блоками, поскольку у них нет
точек сочленения. Эти блоки уместно назвать
\odmkey{тривиальными}{Блок (block)!тривиальный (trivial)},
поскольку они удовлетворяют определению, но не обладают 
свойствами, присущими нетривиальным блокам.
В частности, нетрудно видеть, что в нетривиальном блоке не только нет
точек сочленения, но также нет мостов и нет висячих вершин. Более того,
все вершины блока связаны неединственным 
образом, что устанавливает следующая теорема.

\begin{theorem}
Если граф $G(V, E)$ связен и $p(G)>2$,
то следующие утверждения эквиваленты:
\begin{enumerate}\leftskip-8pt
\item[1)] $G$~{\upshape ---} нетривиальный блок;\\%[-2pt]
\item[2)] любые две вершины принадлежат одному простому циклу;\\%[-2pt]
\item[3)] любые вершина и  ребро принадлежат одному простому циклу;\\%[-2pt]
\item[4)] любые два ребра принадлежат одному простому циклу;\\%[-2pt]
\item[5)] для любых двух вершин и любого ребра существует простая цепь, соединяющая эти вершины и содержащая это ребро;\\%[-2pt]
\item[6)] для любых трёх вершин существует простая цепь, соединяющая две вершины и~содержащая третью.
\end{enumerate}
\end{theorem}
\end{frame}

\begin{frame}
\frametitle{8.1.3. Свойства блоков (2/7)}
\begin{proof}
\proofsection{1$\Impl$ 2}
Пусть $u$ и $w$~--- любые две вершины.
Рассмотрим $U$~--- множество всех вершин,
входящих во все простые циклы, содержащие вершину $u$. Ясно,
что $U\ne\emptyset$, поскольку вершина $u$ не является висячей,
поэтому для любых двух вершин, $u_1$ и $u_2$, смежных с вершиной $u$,
по следствию 2 к теореме 1 существует соединяющая их цепь, не содержащая
вершины $u$, и, значит, существует некоторый простой цикл, содержащий
вершину $u$. Положим $W \Assign V \setminus U$.
\begin{figure}[h]
\begin{center}
\includegraphics[scale=0.6]{08_02}
%%% здесь нужно сделать и вставить картинку
\end{center}
%\caption{К доказательству пункта [1$\Impl$2] теоремы о блоках}
\end{figure}

\vspace{-16pt}
\end{proof}
\end{frame}


\begin{frame}
\frametitle{8.1.3. Свойства блоков (3/7)}
\begin{proof} (Продолжение)
Далее от противного.
Пусть $W \ne\emptyset$ и $w\in W$. Граф связен, и, значит,
в нём есть ребро $(s, t)$ такое, что $s\in U$, $t\in W$
(лемма \DMref{п.~8.1.1}{8.1.1.}). Поскольку $u\in U$ и $s\in U$,
обе эти вершины принадлежат некоторому простому циклу $Z$.
По следствию 2 существует простая $\langle u, t\rangle$-цепь $P$,
не содержащая вершину $s$. Обозначим через $v$ первую вершину, считая от
вершины $t$, на цепи $P$, которая также принадлежит циклу $Z$.
Тогда рассмотрим следующую последовательность простых цепей:
отрезок цепи $P$ от вершины $t$ до вершины $v$,
отрезок цикла $Z$ от вершины $v$ до вершины $u$,
отрезок цикла $Z$ от вершины $u$ до вершины $s$,
ребро $(s, t)$. Это цикл, причём простой,
и вершины $u$ и $t$ обе принадлежат этому циклу
и входят в множество $U$,
что противоречит выбору вершины $t$ и предположению $w\in W$.


\proofsection{2$\Impl$3}
Пусть $u$~--- любая вершина и $(v, w)$~--- любое ребро.
Вершины $u$ и $v$ принадлежат некоторому простому циклу $Z$.
Если $w\in Z$, то требуемый цикл строится из ребра $(v, w)$
и из того отрезка цикла $Z$ от $v$ до $w$,
который содержит вершину $u$. 
\end{proof}
\end{frame}


\begin{frame}
\frametitle{8.1.3. Свойства блоков (4/7)}
\begin{proof} (Продолжение)
Если же $w\not\in Z$,
то по следствию 2 существует простая  $\langle u, w \rangle$-цепь $P$,
не содержащая вершину $v$. 
Обозначим через $s$ первую вершину,
считая от вершины $w$, на цепи $P$, которая также принадлежит циклу $Z$.
Тогда рассмотрим следующую последовательность простых цепей:
отрезок цепи $P$ от вершины $w$ до вершины $s$,
отрезок цикла $Z$ от вершины $s$ до вершины $v$,
содержащий вершину $u$, и ребро $(v, w)$.
Это цикл, причём простой, содержащий вершину $u$ и ребро $(v, w)$.

\begin{figure}[h]
\begin{center}
\includegraphics[scale=0.75]{08_03}
%%% здесь нужно сделать и вставить картинку
\end{center}
%\caption{К доказательству пункта [2$\Impl$3] теоремы о блоках}
\end{figure}
\end{proof}
\end{frame}


\begin{frame}
\frametitle{8.1.3. Свойства блоков (5/7)}
\begin{proof} (Продолжение)
\proofsection{3$\Impl$4}
Пусть $(u, w)$ и $(s, t)$~--- два любых ребра.
По условию, ребро $(u, w)$ и~вершина $s$
принадлежат некоторому циклу $Z_1$, и в то же время
ребро $(u, w)$ и~вершина $t$ принадлежат некоторому циклу $Z_2$.
%(см.~рис.~8.4).
Тогда рассмотрим следующую последовательность простых цепей:
ребро $(u, w)$, отрезок цикла $Z_1$ от вершины $u$ до вершины $s$,
ребро $(s, t)$, отрезок цикла $Z_2$ от вершины $t$ до вершины $w$.

Это цикл, причём простой, содержащий рёбра $(u, w)$ и $(s, t)$.

\begin{figure}[h]
\begin{center}
\includegraphics[scale=0.75]{08_04}
%%% здесь нужно сделать и вставить картинку
\end{center}
%\caption{К доказательству пункта [3$\Impl$4] теоремы о блоках}
\end{figure}

\end{proof}
\end{frame}


\begin{frame}
\frametitle{8.1.3. Свойства блоков (6/7)}
\begin{proof} (Продолжение)
\proofsection{4$\Impl$5}
Пусть $u$ и $w$~--- любые две вершины и $(s, t)$~--- любое ребро.
Граф связен, значит, существует простая $\langle u, w\rangle$-цепь $P$. Пусть $P = (u, x, \dots, w)$. По условию, рёбра $(u, x)$ и $(s, t)$ принадлежат некоторому простому циклу $Z$ (см.~рисунок).
Обозначим через $v$ первую вершину, считая от вершины $w$,
на цепи $P$, которая также принадлежит циклу $Z$
(может оказаться, что $v=x$). Тогда рассмотрим
следующую последовательность простых цепей:
отрезок цикла $Z$ от вершины $u$ до вершины $v$,
проходящий через ребро $(s, t)$,
отрезок цепи $P$ от вершины $v$ до вершины $w$.
Это цепь, причём простая, содержащая ребро $(s, t)$.

\vspace{-12pt}
\begin{figure}[h]
\begin{center}
\includegraphics[scale=0.7]{08_05}
%%% здесь нужно сделать и вставить картинку
\end{center}
%\caption{К доказательству пункта [4$\Impl$5] теоремы о блоках}
\end{figure}

\vspace{-12pt}
\end{proof}
\end{frame}



\begin{frame}
\frametitle{8.1.3. Свойства блоков (7/7)}
\begin{proof} (Окончание)
\proofsection{5$\Impl$6}
Пусть $u$, $v$, $w$~--- любые три вершины и $(w, x)$~--- некоторое ребро,
инцидентное вершине $w$. По условию существует простая
$\langle u, v\rangle$-цепь $P$, содержащая ребро $(w, x)$ и, значит, вершину $w$.
\proofsection{6$\Impl$1}
Пусть $v$ --- любая вершина. Покажем, что вершина $v$ не является точкой
сочленения, то есть что граф $G-v$ связен. Рассмотрим $u$ и $w$~---
любые две вершины. По условию существует простая
$\langle u, v\rangle$-цепь $P$, проходящая через вершину $w$.
Отрезок цепи $P$ от $u$ до $w$ обеспечивает связанность вершин
$u$ и $w$ в графе $G-v$.
\end{proof}
%%%%%%%%%%%%%%%%%%%%%%%%%% Конец изменений 2012

%\begin{example}
%{\color{blue} Диаграмма блока, иллюстрирующая все положения теоремы}
%\end{example}
\end{frame}


\begin{frame}
\label{8.1.4.}\subsubsection{8.1.4. Вершинная и рёберная связность }\frametitle{8.1.4. Вершинная и рёберная связность (1/3)}


При взгляде на диаграммы связных графов возникает естественное впечатление,
что связный граф может быть <<более>> или <<менее>> связен. В этом
параграфе рассматриваются некоторые количественные меры
связности.

\medskip
\odmkey{Вершинной связностью}{Связность (connectivity)!вершинная (vertex)}
 графа $G$
(обозначается $\varkappa(G)$) называется наименьшее~число
вершин, удаление которых приводит к несвязному или тривиальному графу.

%\pagebreak
\begin{example}
\begin{tabbing}
$\varkappa(K_p)=p-1,\quad$\=     \kill
$\varkappa(G)=0,$ \> если $G$  несвязен;\\[3pt]
$\varkappa(G)=1,$ \> если $G$  имеет точку сочленения;\\[3pt]
$\varkappa(K_p)=p-1,$ \> если $K_p$ ~--- полный граф.
\end{tabbing}
\end{example}



Граф $G$ называется \odmkey{$n$-связным}{Граф (graph)!$n$-связный ($n$-connected)},
если $\varkappa(G)=n$.


\medskip
\odmkey{Рёберной связностью}{Связность (connectivity)!рёберная (edge)}
 графа $G$
(обозначается $\lambda(G)$) называется наименьшее число
рёбер, удаление которых приводит к несвязному или тривиальному графу.


\end{frame}

\begin{frame}
\frametitle{8.1.4. Вершинная и рёберная связность (2/3)}
\begin{example}
\begin{tabbing}
$\lambda(K_p)=p-1,\quad$ \= если \kill
$\lambda(G)=0,$ \> если $G$ несвязен;\\[2pt]

$\lambda(G)=1,$ \> если $G$ имеет мост;\\[2pt]
$\lambda(K_p)=p-1,$ \> если $K_p$~--- полный граф.
\end{tabbing}
\end{example}



\begin{theorem} $\varkappa(G) \leq \lambda(G) \leq \delta(G)$.
\end{theorem}


\begin{proof}

\proofsection{$\varkappa\leq\lambda$} Если $\lambda=0$, то
$\varkappa=0$. Если $\lambda=1$, то есть мост, а значит, либо есть
точка сочленения, либо $G=K_2$. В любом случае $\varkappa=1$.
Пусть теперь $\lambda\geq 2$. Тогда, удалив $\lambda-1$ ребро,
получим граф $G'$, имеющий мост ($u,v$). Для каждого из удаляемых
$\lambda-1$ рёбер удалим инцидентную вершину, отличную от $u$ и
$v$. Если после такого удаления вершин граф несвязен, то
$\varkappa\leq\lambda-1<\lambda$. Если связен, то удалим один из
концов моста~--- $u$ или~$v$. Имеем $\varkappa\leq\lambda$.



\proofsection{$\lambda\leq\delta$}
Если $\delta=0$, то $\lambda=0$. Пусть вершина
$v$ имеет наименьшую степень: $d(v)=\delta$.
Удалим все рёбра, инцидентные $v$.
Тогда $\lambda\leq\delta$.
\end{proof}

%%%%%%%%%%%%%%%%%%%%%%%%% Изменения 2012
\end{frame}

\begin{frame}
\frametitle{8.1.4. Вершинная и рёберная связность (3/3)}
\begin{example}
Слева $\varkappa=\lambda=\delta=2$, 
в центре $\varkappa=1, \lambda=\delta=2$, 
справа $\varkappa=1, \lambda=2, \delta=3$. 
\end{example}


\begin{center}
\includegraphics[width=100mm]{8_1_4.png}
\end{center}

Особое значение имеют
\odmkey{двусвязные}{Граф (graph)!двусвязный (biconnected)}
графы по двум причинам.
Во-первых, они особенно часто встречаются в приложениях.
Во-вторых, двусвязные графы являются частным случаем блоков,
которые обладают многими полезными свойствами,
часть из которых приведена в предыдущем параграфе.

\begin{remark}
Легко видеть, что любой двусвязный граф является блоком,
поскольку не содержит точек сочленения.
В то же время всякий нетривиальный блок является,
по меньшей мере, двусвязным,
но может быть и трёхсвязным,
и вообще иметь произвольную величину вершинной связности,
поскольку, например, всякий полный граф является блоком.
\end{remark}
%%%%%%%%%%%%%%%%%%%%%%%%% Конец изменений 2012


\end{frame}


\begin{frame}
\label{8.1.5.}\subsubsection{8.1.5. Оценка числа рёбер }\frametitle{8.1.5. Оценка числа рёбер (1/3)}
%\label{pqk}

Легко заметить, что инварианты $p(G)$, $q(G)$ и $k(G)$
не являются независимыми.


\begin{theorem}
Пусть $p$~{\upshape---} число вершин, $q$~{\upshape---} число
рёбер, $k$~{\upshape---} число  компонент связности графа. Тогда
$p - k \le q \le (p-k)(p-k+1)/2$.
\end{theorem}


\begin{proof}

\proofsection{$p-k\leq q$} Индукция по $p$. База: $p=1$, $q=0$,
$k=1$. Пусть оценка верна для всех графов с числом вершин меньше
$p$. Рассмотрим граф $G$, $p(G)=p$. Удалим в $G$ вершину $v$,
которая не является точкой сочленения: $G'\Assign G-v$. Тогда если
$v$~--- изолированная вершина, то $p'=p-1$, $k'=k-1$, $q'=q$.
Имеем $p-k=p'-k' \leq q'=q$. Если $v$~--- не изолированная
вершина, то $p'=p-1$, $k'=k$, $q'<q$. Имеем $p-k=1+p'-k'\leq
1+q'\leq q$.



\proofsection{$q\leq (p-k)(p-k+1)/2$}
Метод выделения <<критических>> графов.
Число рёбер $q$ графа с $p$ вершинами и $k$ компонентами связности
не превосходит числа рёбер в графе с $p$ вершинами и $k$ компонентами
связности, в котором все компоненты связности суть полные графы.
\end{proof}
\end{frame}


\begin{frame}
\frametitle{8.1.5. Оценка числа рёбер (2/3)}
\begin{proof} (Продолжение)
Следовательно, достаточно рассматривать только графы, в которых все
компоненты связности полные.
Среди всех графов с $p$ вершинами и $k$ полными компонентами
связности наибольшее число рёбер имеет граф $\widetilde{K_p}$, в котором
все рёбра сосредоточены в одной компоненте связности:
$$
\widetilde{K_p}\Assign K_{p-k+1}\cup\bigcup\limits_{i=1}^{k-1}K_1.
$$
Действительно, пусть есть две компоненты связности, $G_1$ и $G_2$,
такие, что $$1 <p_1 =p(G_1)\leq p_2=p(G_2).$$
Если перенести одну
вершину из компоненты связности $G_1$ в компоненту связности
$G_2$, то число рёбер изменится на $\Delta
q ={p_2-(p_1-1)} ={p_2-p_1+1>0}$.

Следовательно, число рёбер
возросло.
Таким образом, достаточно рассмотреть граф $\widetilde{K_p}$. Имеем
$
q(\widetilde{K_p})={(p-k+1)(p-k+1-1)}/{2}+0=
{(p-k)(p-k+1)}/{2}$.
\end{proof}

\end{frame}


\begin{frame}
\frametitle{8.1.5. Оценка числа рёбер (3/3)}
%\setcounter{corollary}{0}
\begin{corollary} \NB{[1]}
Если $q > (p-1)(p-2)/{2}$, то граф связен.
\end{corollary}
\begin{proof}
От противного. Пусть $q>(p-1)(p-2)/2$, но граф не связен, то есть $k \ge 2$.
Но тогда $(p-1)(p-2)/2 < q \le (p-k)(p-k+1)/2 \le (p-2)(p-1)/2$~--- противоречие.
\end{proof}

\medskip
\begin{corollary} \NB{[2]}
В связном графе $p-1\leq q\leq p(p-1)/2$.
\end{corollary}
\begin{proof}
Следует из теоремы при $k=1$.
\end{proof}
%
%\medskip
%{\color{red} Запрос на удачу: 
%нет ли ещё красивых следствий из этой теоремы,
%полезных при решении задач и
%с внятными доказательствами?}

\end{frame}

\begin{frame}
\frametitle{Онтология теории графов}
\begin{center}
\includegraphics[height=8cm]{ODM_8_1.png}
\end{center}

\end{frame}

\subsection{8.2. Теорема Менгера}
\label{8.2.}
\begin{frame}
\frametitle{8.2. Теорема Менгера}

Интуитивно очевидно,
что граф тем более связен (при использовании различных мер связности),
чем больше существует различных цепей, соединяющих
одну вершину с другой, и тем менее связен,
чем меньше нужно удалить промежуточных
вершин, чтобы отделить одну вершину от другой.
В этом разделе устанавливается теорема Менгера,
которая придаёт этим неформальным наблюдениям
точный и строгий смысл.

\medskip
\begin{tabular}{p{9mm}|p{50mm}|p{84mm}}
  № & Название параграфа
  & Ключевые термины и обозначения \\
  \hline
\DMref{8.2.1.}{8.2.1.} & Непересекающиеся цепи и разделяющие множества 
  & 
  {$P(u,v)$~--- множество непересекающихся цепей;  $S(u,v)$~--- разделяющее множество вершин} \\
\DMref{8.2.2.}{8.2.2.}\RS & Теорема Менгера в <<вершинной форме>> 
  & 
  {}
  \\
\DMref{8.2.3.}{8.2.3.} & Варианты теоремы Менгера &
{ }\\
\end{tabular}
\end{frame}


\begin{frame}
\label{8.2.1.}\subsubsection{8.2.1. Непересекающиеся цепи и~разделяющие~множества }\frametitle{8.2.1. Непересекающиеся цепи и~разделяющие~множества (1/2)}

Пусть $G(V,E)$~--- связный граф, $u$ и $v$~--- две его несмежные
вершины. Две цепи $\divC{u}{v}$ называются
\odmkey{вершинно-непересекающимися}{Цепь
(circuit)!вершинно-непересекающаяся (pairwise
vertex-independent)}, если у них нет общих вершин, отличных от $u$
и $v$. Две цепи $\divC{u}{v}$ называются
\odmkey{рёберно-непересекающимися}{Цепь
(circuit)!рёберно-непересекающаяся (pairwise edge-independent)},
если у них нет общих рёбер. Если две цепи вершинно не
пересекаются, то они и~рёберно не пересекаются. Обозначим через
$P(u,v)$ множество вершинно-непере\-сека\-ющихся цепей
$\divC{u}{v}$.
% ЭТИ ФОРМУЛЫ ВРЕМЕННО УБРАНЫ
%$$
%P(u,v)\Def \left\{ \langle u,v\rangle\left|\bigcap_{\langle u,v\rangle\in P(u,v)}\langle u,v\rangle=\{ u,v\}\right. \right\}.
%
%\subset\{\divC{u}{v}\} \And {\divC{u}{v}}_1\in
%P\And{\divC{u}{v}}_2\in P\Impl{\divC{u}{v}}_1\cap{\divC{u}{v}}_2=\{u,v\}.
%$$
Множество $S$ вершин и (или) рёбер графа $G$ \em{разделяет} две
вершины, $u$ и $v$, если $u$ и $v$ принадлежат разным компонентам
связности графа $G-S$. 
\odmkey{Разделяющее множество
рёбер}
{Множество (set)!рёбер (of edges)!разделяющее (separate)}
называется \odmkey{разрезом}{Разрез (cut)}. \odmkey{Разделяющее
множество вершин}{Множество (set)!вершин (of vertex)!разделяющее
(dividing)} для вершин $u$ и $v$ обозначим $S(u,v)$.
%
%% ЭТИ ФОРМУЛЫ ВРЕМЕННО УБРАНЫ
%%$$
%%S(u,v)\Def\Set{w\in V}{{G-S=G_1\cup G_2}\And {v\in G_1}\And {u\in G_2}}\!\!.
%%$$
%
Для заданных вершин $u$ и $v$ множества $P(u,v)$ и $S(u,v)$ можно выбирать
различным образом.
Из определений следует, что при $\vert P(u,v)\vert>1$ любое
множество $P(u,v)$ вер\-шин\-но-не\-пе\-ре\-сека\-ющих\-ся цепей обладает
тем свойством, что
$
\bigcap_{\divC{u}{v}\in P(u,v)} \divC{u}{v} =\{ u,v\},
$
а любое разделяющее множество вершин $S(u,v)$ обладает тем свойством, что
$
G\!-\!S=G_1\!\cup\!G_2\!\And\!{v\!\in\!G_1}\!\And\!{u\!\in\!G_2}.
$
\end{frame}


\begin{frame}
\frametitle{8.2.1. Непересекающиеся цепи и~разделяющие~множества (2/2)}

%\vspace*{-0.5\baselineskip}
\begin{remark}
Если $u$ и $v$ принадлежат разным компонентам связности графа $G$, \\то
$\vert P(u,v)\vert=0$ и $\vert S(u,v)\vert=0$. Поэтому без ограничения
общности можно считать, что $G$~--- связный граф.
\end{remark}

%\vspace*{-0.5\baselineskip}
\begin{example}
В графе, диаграмма которого представлена на рисунке,
можно выбрать мно\-жества~вершинно-непересекающихся путей
$P_1=\left\{ \langle u,a,d,v\rangle,\langle
u,b,e,v\rangle\right\}$ и \\ $P_2 =\left\{ \langle
u,b,d,v\rangle,\langle u,c,e,v\rangle\right\}$. Заметим, что путь
$\langle u,c,b,d,v\rangle$ образует множество $P_3$\\
вершинно-непересекающихся путей, состоящее из одного элемента.
Множества вершин $S_1=\{a,b,c\}$ и $S_2=\{d,e\}$ являются
разделяющими для вершин~$u$~и~$v$.
\end{example}

\vspace{-6pt}
\begin{figure}[h]
\begin{center}
\includegraphics[width=90mm]{odm08-02_add}
\end{center}
%\caption{Вершинно-непересекающиеся пути и разделяющие множества вершин}
%\label{f0802add}
\end{figure}

\end{frame}


\begin{frame}
\label{8.2.2.}\subsubsection{8.2.2. Теорема Менгера в <<вершинной форме>> }\frametitle{8.2.2. Теорема Менгера в <<вершинной форме>> (1/4)}
%\label{Menger}



\begin{theorem} {\NB{[Менгера]}}
Пусть $u$ и $v$~{\upshape---} несмежные вершины в графе $G$.
Наименьшее число вершин в множестве, разделяющем $u$ и $v$, равно
наибольшему числу вершинно-непересекающих\-ся простых
$\divC{u}{v}$-цепей: 
\odmkey{}{Теорема (theorem)!Менгера (Menger)}
$$
\max\vert P(u,v)\vert = \min\vert S(u,v)\vert.
$$
\end{theorem}


\begin{remark}
Пусть $G$~--- связный граф и вершины $u$ и $v$ несмежны. Пусть
$P\Assign P(u,v)$, $S \Assign S(u,v)$. Легко видеть, что $\vert
P\vert\leq\vert S\vert$. Действительно, любая $\divC{u}{v}$-цепь
проходит через $S$. Если бы $\vert P\vert>\vert S\vert$, то в $S$
была бы вершина, принадлежащая более чем одной цепи из $P$, что
противоречит выбору $P$. Таким образом, $\A{P}{\A{S}{\vert
P\vert\leq\vert S\vert}}$. Следовательно, $\max\vert
P\vert \leq\min\vert S\vert$. Утверждение теоремы состоит в том,
что в любом графе \em{существуют\/} такие $P$ и $S$, что
достигается равенство $\vert P\vert=\vert S\vert$.
\end{remark}

%
\begin{proof}
Пусть $G$~--- связный граф, $u$ и $v$~--- несмежные вершины.
Совместная индукция по $p$ и $q$. База: наименьший граф,
удовлетворяющий условиям теоремы, состоит из трёх вершин, $u,w,v$, и
двух рёбер, $(u,w)$ и $(w,v)$. В нём
$P(u,v)=\left\{\lrle{u,w,v}\right\}$ и $S(u,v)=\{w\}$. Таким
образом, $\vert P(u,v)\vert=\vert S(u,v)\vert=1$. 
\end{proof}
\end{frame}


\begin{frame}
\frametitle{8.2.2. Теорема Менгера в <<вершинной форме>> (2/4)}
\begin{proof} (Продолжение)
Пусть
утверждение теоремы верно для всех графов с числом вершин меньше
$p$ и (или) числом рёбер меньше $q$. Рассмотрим граф $G$ с $p$
вершинами и $q$ ребрами. Пусть $u$, $v\in V$, причём $u$, $v$~---
не смежны и $S$~--- некоторое наименьшее множество вершин,
разделяющее $u$ и $v$. Обозначим $n\Assign\vert S\vert$.
Рассмотрим три случая.


\proofsection{1} Пусть в $S$ есть вершины, несмежные с $u$ и
несмежные с $v$. Тогда граф $G-S$ состоит из двух 
\em{нетривиальных\/} графов, $G_1$ и $G_2$. Образуем два новых графа,
$G_u$ и $G_v$, стягивая  нетривиальные графы $G_1$ и  $G_2$ к
вершинам $u$ и $v$: $G_u\Assign G/G_1$, $G_v\Assign
G/G_2$. 

\begin{figure}[h]
\begin{center}
\includegraphics[width=120mm]{odm08-02-lebedev}
\end{center}
\end{figure}


\vspace{-6pt}
$S$ по-прежнему является наименьшим разделяющим множеством для $u$
и $v$ как в~$G_u$,  так и в $G_v$. Так как $G_1$ и $G_2$
нетривиальны, то $G_u$ и $G_v$ имеют меньше вершин и (или) рёбер,
чем $G$. 
\end{proof}
\end{frame}


\begin{frame}
\frametitle{8.2.2. Теорема Менгера в <<вершинной форме>> (3/4)}
\begin{proof} (Продолжение)
Следовательно, по индукционному предположению в $G_u$ и~в~$G_v$ имеется $n$ вершинно-непересекающихся простых цепей.
Комбинируя отрезки этих цепей от $u$ до $S$ и от $S$ до $v$,
получаем $n$ непересекающихся простых цепей в $G$.



\proofsection{2} Пусть все вершины $S$ смежны с $u$ или с $v$ 
(пусть с $u$) и~среди вершин $S$ есть вершина $w$,
смежная одновременно с $u$ и с $v$. 
%(рис.~\ref{f0803}).



\begin{figure}[h]
\begin{center}
\includegraphics[width=45mm]{odm08-03-lebedev}
\end{center}
%\caption{К доказательству теоремы Менгера, случай 2}
\label{f0803}
\end{figure}


\vspace{-4pt}
Рассмотрим граф $G'\Assign G-w$. В нём $S-w$~--- разделяющее
множество для $u$ и $v$, причём наименьшее. По индукционному
предположению в $G'$ имеется ${\vert S-w\vert} ={n-1}$
вершинно-непересекаю\-щихся простых цепей. Добавим к ним цепь
$\lrle{u,w,v}$. Она простая и вершинно не пересекается с
остальными. Таким образом, имеем $n$ вершинно-непересе\-кающихся
простых цепей в $G$.
\end{proof}
\end{frame}


\begin{frame}
\frametitle{8.2.2. Теорема Менгера в <<вершинной форме>> (4/4)}


\begin{proof} (Окончание)
\proofsection{3} Пусть теперь все вершины $S$ смежны с $u$ или с
$v$ (пусть с $u$) и среди вершин $S$ нет
вершин, смежных одновременно с вершиной $u$ и~с вершиной $v$.
Рассмотрим {\it кратчайшую} $\divC{u}{v}$-цепь
$\lrle{u,w_1,w_2,\dots, v}$, $w_1\in S$, $w_2\not= v$.
%(рис.~\ref{f0804}).

\vspace{-6pt}

\begin{figure}[h]
\begin{center}
\includegraphics[width=40mm]{odm08-04-lebedev}
\end{center}
%\caption{К доказательству теоремы Менгера, случай 3}
\label{f0804}
\end{figure}


\vspace{-8pt}
Рассмотрим граф $G'\Assign G/\{w_1,w_2\}$, полученный из $G$
склейкой вершин $w_1$ и $w_2$ в вершину $w_1$. Имеем $w_2\not\in
S$, в противном случае цепь $\lrle{u,w_2,\dots, v}$ была бы ещё
более короткой. Следовательно, в графе $G'$ множество $S$
по-прежнему является наименьшим, разделяющим $u$ и $v$, и граф
$G'$ имеет по крайней мере на одно ребро меньше. По индукционному
предположению в $G'$ существуют $n$ простых непересекающихся
 цепей. Но цепи, которые не пересекаются в $G'$, не
пересекаются и в $G$. Получаем $n$
цепей в $G$.
\end{proof}



\end{frame}

\begin{frame}
\label{8.2.3.}\subsubsection{8.2.3. Варианты теоремы Менгера}
\frametitle{8.2.3. Варианты теоремы Менгера}



Теорема Менгера представляет собой весьма общий факт, который в разных
формулировках встречается в различных областях математики. Комбинируя
вер\-шин\-но- и рёберно\-непересекающиеся цепи, разделяя не отдельные
вершины, а множества вершин, используя инварианты $\varkappa$ и $\lambda$,
можно сформулировать и другие утверждения,
подобные теореме Менгера. Например:


%\setcounter{theorem}{0}

\begin{theorem}
Для любых двух несмежных вершин $u$ и $v$ графа $G$ наибольшее
число рё\-бер\-но-не\-пе\-ре\-се\-кающихся $\divC{u}{v}$-цепей равно
наименьшему числу рёбер в $(u,v)$"~раз\-резе.
\end{theorem}



\begin{theorem}
Чтобы граф $G$ был $n$-связным, необходимо и
достаточно, чтобы любые две несмежные вершины были соединены не менее
чем $n$ вершинно-не\-пе\-ре\-се\-ка\-ющи\-ми\-ся простыми цепями.
\end{theorem}



Другими словами, в любом графе $G$ любые две несмежные вершины соединены
не менее чем $\varkappa(G)$ вершинно-непересекающимися простыми цепями.

%Доказательство подобных утверждений требует места, и мы оставляем их
%за пределами книги. В последующих разделах приведены и некоторые другие
%результаты, в которых проявляется теорема Менгера.

%\medskip
%{\color{red}
%Запрос: красивое следствие из теоремы Менгера,
%полезное при решении задач и с доказательством на один слайд.
%}
\end{frame}

\begin{frame}
\frametitle{Онтология теории графов}
\begin{center}
\includegraphics[height=8cm]{ODM_8_2.png}
\end{center}

\end{frame}

\subsection{8.3. Паросочетания}
\label{8.3.}

\begin{frame}
\frametitle{8.3. Паросочетания}

\odmkey{Паросочетанием}{Паросочетание (matching)} (или
\odmkey{независимым множеством рёбер}{Множество (set)!рёбер (of
edges)!независимое (independent)}) называется множество рёбер, в
котором никакие два ребра не смежны. Паросочетание называется
\odmkey{максимальным}
{Паросочетание (matching)!максимальное (maximal)}, если никакое его
надмножество не является независимым.
Паросочетание называется
\odmkey{наибольшим}
{Паросочетание (matching)!наибольшее
(maximum)}, если оно включает наибольшее возможное
число рёбер.


\smallskip
Задача отыскания паросочетаний имеет множество различных
применений и интерпретаций,
часть которых приведена в этом разделе.

%\smallskip
%Рассмотрение  интерпретаций начинается c
%теоремы Холла
%<<о~свадьбах>> и продолжается другими примерами. 

\medskip
\begin{tabular}{p{9mm}|p{65mm}|p{70mm}}
  № & Название параграфа
  & Ключевые термины и обозначения \\
  \hline
\DMref{8.3.1.}{8.3.1.} & Задачи о свадьбах, трансверсалях и паросочетаниях 
  & 
  {Система различных представителей; 
  совершенное паросочетание} \\
\DMref{8.3.2.}{8.3.2.}\RS & Теорема Холла 
  & 
  {}
  \\
\DMref{8.3.3.}{8.3.3.} & Латинские квадраты &
{Латинский прямоугольник; ортогональные латинские квадраты }\\
\DMref{8.3.4.}{8.3.4.} & Устойчивые паросочетания &
{Неустойчивое паросочетание}
\\
\DMref{8.3.5.}{8.3.5.}\BS & Наибольшее паросочетание &
{Полное паросочетание}
\\
\DMref{8.3.6.}{8.3.6.}\BS & Алгоритм поиска наибольшего паросочетания &
{}
\\
\end{tabular}
\end{frame}

\begin{frame}
\label{8.3.1.}\subsubsection{8.3.1. Задачи о свадьбах, трансверсалях и паросочетаниях }\frametitle{8.3.1. Задачи о свадьбах, трансверсалях и паросочетаниях (1/2)}
%\label{s841}
{\small
Рассмотрим три внешне различных задачи.

\vspace{-6pt} 
\begin{itemize}\leftskip-12pt
\item
Пусть имеется конечное множество юношей, каждый из которых знаком с
некоторым подмножеством конечного множества девушек.
В каком случае всех юношей можно женить так, чтобы каждый
женился на знакомой девушке?
\odmkey{}{Задача (problem)!о свадьбах (marriage)}
\\[-6pt]
\item
Пусть $S=\{S_1,\dots,S_m\}$~--- семейство непустых подмножеств конечного
множества $E$. Подмно\-жества $S_k$ не обязательно различны и могут
пересекаться. \odmkey{Системой различных представителей}{Система
(system)!различных представителей (of distinct representatives)}
в семействе $S$ (или \odmkey{трансверсалью}{Трансверсаль
(transversal)} в семействе $S$) называется любое подмножество \\
$C=\{c_1,\dots,c_m\}$ из $m$ элементов множества $E$ таких, что
$\A{k\in 1..m}{c_k\in S_k}$. В каком случае существует
трансверсаль?
\begin{remark}
Все элементы множества $C$  различны,
откуда и происходит
название <<система различных представителей>>.
\end{remark}

\vspace{-4pt}
\item
Пусть $G(V_1,V_2,E)$~--- двудольный граф.
\odmkey{Совершенным паросочетанием}{Паросочетание
(matching)!совершенное (perfect)}
из $V_1$ в $V_2$ называется паросочетание, покрывающее вершины $V_1$.
В каком случае существует совершенное паросочетание из $V_1$ в $V_2$?
\\
\begin{remark}
Совершенное паросочетание является максимальным и наибольшим.
\end{remark}
\end{itemize}
}
\end{frame}

\begin{frame}
\frametitle{8.3.1. Задачи о свадьбах, трансверсалях и паросочетаниях (2/2)}

Вообще говоря, указанные задачи~---
это одна и та же задача. 
Задача о свадьбах  сводится к задаче 
о паросочетаниях следующим образом.
$V_1$~---
множество юношей, $V_2$~--- множество девушек, рёбра~--- знакомства
юношей с девушками. В~таком случае
совершенное паросочетание~--- искомый набор свадеб.
Задача о трансверсали
сводится к задаче о паросочетаниях следующим образом.
Положим $V_1\Assign S$, $V_2\Assign E$,
ребро $(S_k,e_i)$ существует, если $e_i\in S_k$.
В таком случае совершенное
паросочетание~--- искомая трансверсаль.
Все задачи 
имеют общий ответ:
в том и только том случае, когда

\vspace{-14pt}
\begin{equation*}
\begin{split}
&\mbox{любые } k
\left(\begin{array}{c}
\mbox{юношей}\\
\mbox{подмножеств}\\
\mbox{вершин из $V_1$}
\end{array}\right)
\mbox{ в совокупности }
\left(\begin{array}{c}
\mbox{знакомы с}\\
\mbox{содержат}\\
\mbox{смежны с}
\end{array}\right)
\\
&\mbox{не менее чем } k
\left(\begin{array}{c}
\mbox{девушками}\\
\mbox{элементов}\\
\mbox{вершинами из $V_2$}
\end{array}\right),
\end{split}
\end{equation*}

\vspace{-3pt}
что устанавливается теоремой Холла.

\end{frame}

\begin{frame}
\label{8.3.2.}\subsubsection{8.3.2. Теорема Холла }\frametitle{8.3.2. Теорема Холла (1/2)}


\begin{theorem} Пусть $G(V_1,V_2,E)$~--- двудольный граф.
Совершенное паросочетание из $V_1$ в $V_2$ существует тогда и
только тогда, когда $\A{A\subset V_1}{\vert A\vert  \leq\vert
\Gamma(A)\vert}$. 
\end{theorem}

\begin{proof}


\proofsection{$\Rightarrow$}
Пусть существует совершенное паросочетание из $V_1$ в $V_2$.
Тогда в~$\Gamma(A)$ входит $\vert A\vert$ вершин из $V_2$, парных к
вершинам из множества $A$, и, возможно, еще что-то. Таким образом,
$\vert A\vert\leq\vert\Gamma(A)\vert$.

\begin{columns}[t]
\begin{column}[T]{8cm}
\proofsection{$\Leftarrow$} Добавим  две новые вершины, $u$ и $v$,
так что вершина $u$ смежна со всеми вершинами из $V_1$, а вершина
$v$ смежна со всеми вершинами из $V_2$. 
Совершенное паросочетание
из $V_1$ в $V_2$ существует тогда и только тогда, когда существуют
$\vert V_1\vert$ вершинно-непересекающихся простых
$\divC{u}{v}$-цепей. 

Ясно, что $\vert
P(u,v)\vert \leq\vert V_1\vert$ (так~как $V_1$ разделяет вершины
$u$ и $v$). \qed
\end{column}
\begin{column}[T]{7cm}
\vspace{-12pt}
\begin{figure}[h]
\begin{center}
\includegraphics[width=60mm]{odm08-05}
\end{center}
\end{figure}
\end{column}
\end{columns}
\end{proof}
\end{frame}

\begin{frame}
\frametitle{8.3.2. Теорема Холла (2/2)}
\begin{proof}
(Окончание)
По теореме Менгера, 
если $S$~--- наименьшее множество,
разделяющее вершины $u$ и $v$, то
 $\max\vert P(u,v)\vert=\min\vert
S(u,v)\vert=\vert S\vert$. Имеем $\vert S\vert\leq\vert V_1\vert$. 
Покажем, что $\vert S\vert\geq\vert V_1\vert$. Пусть
$S=A\cup B$, $A\subset V_1$, $B\subset V_2$. Тогда
$\Gamma(V_1\setminus A)\subset B$, не считая вершин $u$, $v$.
Действительно, если бы $\Gamma(V_1\setminus A)\not\subset B$, то
существовал бы <<обходной>> путь $\lrle{u,v_1,v_2,v}$
 и $S$ не было бы разделяющим множеством для
$u$ и $v$. Имеем $\vert V_1\!\setminus\!
A\vert\leq\vert\Gamma(V_1\!\setminus\! A)\vert\leq\vert B\vert.$
Следовательно, 
$\vert S\vert=\vert A\vert+\vert
B\vert\geq\vert A\vert+\vert V_1\setminus
A\vert=\nobreak\vert V_1\vert.$
\end{proof}

\begin{corollary}
 В любом непустом регулярном двудольном графе 
 $G(V_1, V_2)$ существует совершенное паросочетание.
\end{corollary}

\begin{proof}
Пусть $\A{v}{d(v) =r > 0}$. Возьмем произвольные $k$ вершин первой доли
$\{\Tuple{u}{1}{k}\}$ и смежные с ними $n$ вершин второй доли $\{\Tuple{v}{1}{n}\}$. Рассмотрим порожденный этими $k+n$ вершинами подграф исходного графа. В полученном подграфе $\A{i}{d(u_i)=r}$ и  
$\A{j}{d(v_j)\le r}$. 
Число рёбер двудольного графа равно сумме степеней вершин любой из его долей. Поэтому
$ \sum_{i=1}^k d(u_i) = kr = \sum_{j=1}^n d(v_j)\le nr,$
откуда $k\le n$. 
\end{proof}
\end{frame}

\begin{frame}
\label{8.3.3.}\subsubsection{8.3.3. Латинские квадраты }\frametitle{8.3.3. Латинские квадраты (1/3)}

Любопытным приложением теоремы Холла
является известная комбинаторная головоломка о
латинских квадратах.
Матрица 
\textbf{$L:$ array $[1..m, 1..n]$ of $M$}, 
где $M\Assign\{\Tuple{a}{1}{n}\}$  называется \odmkey{латинским прямоугольником}{Прямоугольник (rectangle)!латинский (Latin)}, 
если элементы каждой строки этой матрицы образуют перестановку множества $M$, и в каждом столбце все элементы разные. По принципу Дирихле число строк латинского прямоугольника $m$ не превосхо­дит числа его столбцов $n$. В случае $m=n$ латинский прямоугольник называют \odmkey{латинским квадратом}{Квадрат (square)!латинский (Latin)} порядка $n$.

\medskip
\begin{remark}
Латинские квадраты известны с античных времен. В настоящее время в качестве множества $M$ обычно берётся отрезок $1..n$,  однако  Эйлер использовал буквы латинского алфавита, откуда латинские квадраты и получили своё название. 
\end{remark}

\medskip
Построить латинский квадрат можно, например, так. Заполнить первую строку
любой перестановкой, а каждую следующую строку получать из предыдущей сдвигом на одну позицию вправо, 
перенося последний элемент в первую позицию. 
\end{frame}

\begin{frame}
\frametitle{8.3.3. Латинские квадраты (2/3)}
\begin{theorem}
Любой латинский прямоугольник 
\textbf{array $[1..m, 1..n]$ of $1..n$}, можно дополнить  до латинского квадрата порядка $n$ приписыванием новых $n-m$ строк.
\end{theorem}

\begin{proof}
Построим двудольный граф $G(V_1, V_2 )$, 
в котором $V_1$ есть множество столбцов латинского прямоугольника, а $V_2=1..n$. 
Вершина $i \in V_1$ и вершина $j \in V_2$ 
соединены ребром тогда и только тогда, 
когда в $i$-м столбце прямоугольника нет числа $j$. Ясно, что $\A{v\in V_1}{d(v) = n-m}$.
С другой стороны, любое число $j$ 
встречается в $m$ строках исходного латинского прямоугольника $m$ раз, значит оно появляется в $m$ столбцах и отсутствует в $n-m$ столбцах. 
Отсюда следует, что
$\A{v\in V_2}{d(v)=n-m}$. 
Таким образом, двудольный граф $G(V_1, V_2 )$ 
является непустым и регулярным.  
По следствию к теореме Холла в непустом
регулярном двудольном графе можно выделить  совершенные паросочетание,
которое задаёт очередную новую строку латинского квадрата, и продолжить с латинским прямоугольником, в котором на одну строку больше. 
\end{proof}

\medskip
\begin{digress}
 Известная игра <<судоку>> тоже по сути является латинским квадратом.
\end{digress}
 
\end{frame}

\begin{frame}
\frametitle{8.3.3. Латинские квадраты (3/3)}


Два латинских квадрата $L_1$ и $L_2$ порядка $n$
называются \odmkey{ортогональными}{Квадрат (square)!латинский (Latin)!ортогональный (orthogonal)}, если все упорядоченные пары $(L_1[i,j],L_2[i,j])$ различны.

\medskip
\begin{example}
Среди 36 офицеров поровну уланов, драгунов, гусаров, кирасиров, кавалергардов и гренадеров и, кроме того, поровну генералов, полковников, майоров, капитанов,
лейтенантов и поручиков, причем каждый полк представлен офицерами шести рангов. Можно ли выстроить этих офицеров в каре 6 на 6 так, чтобы в любой колонне и в любой шеренге встречались офицеры всех полков и всех рангов? 
\end{example}

\medskip
\begin{digress}
Эйлер не смог построить ортогональных латинских квадратов порядков 6 и 10, затем высказал гипотезу, что не существует ортогональных квадратов для $n = 4r + 2$. В 1901 г. гипотеза была подтверждена для $n = 6$, а опровергнута в 1959 г. путём построения  ортогональных квадратов для $n = 22$. В 1960 г. с помощью ЭВМ было построено два ортогональных квадрата для $n = 10$. Задача про офицеров оказалась удивительным исключением! 
\end{digress}


\end{frame}

\begin{frame}
\label{8.3.4.}\subsubsection{8.3.4. Устойчивые паросочетания }\frametitle{8.3.4. Устойчивые паросочетания (1/8)}


Пусть имеется конечное множество мужчин и конечное множество женщин (равномощные). Элементы множества 
мужчин обозначим прописными\hspace{-1pt} буквами, 
\hspace{-1pt} а\hspace{-1pt} элементы\hspace{-1pt} множества\hspace{-1pt} женщин~--- 
строчными.  
Пусть также задано совершенное взаимно-однозначное 
паросочетание из множества мужчин в множество женщин и обратно 
(брачное отношение): \textbf{Aa}, \textbf{Bb} и т.~д.
Пусть теперь у каждого мужчины есть список женщин, 
упорядоченный по убыванию <<привлекательности>>
для этого мужчины. А у каждой женщины~--- список мужчин,
также упорядоченный по убыванию <<привлекательности>> для неё.

\medskip
Паросочетание называется \odmkey{неустойчивым}{Паросочетание (matching)!неустойчивое (unstable)}, если в нём найдутся мужчина \textbf{X} (женатый на женщине \textbf{x}) и женщина \textbf{y} (замужем за мужчиной \textbf{Y}), предпочитающие друг друга своим супругам, то есть:
\begin{enumerate}
\item[1)]	\textbf{X} считает \textbf{y} привлекательней, чем \textbf{x};
\item[2)]	\textbf{y} считает \textbf{X} привлекательней, чем \textbf{Y}.
\end{enumerate}

\medskip
Паросочетание называется \odmkey{устойчивым}{Паросочетание (matching)!устойчивое (stable)}, 
если таких пар в нём нет.

\end{frame}

\begin{frame}
\frametitle{8.3.4. Устойчивые паросочетания (2/8)}
\begin{example}
Четверо мужчин (\textbf{A, B, C, D}) состоят в браке с четырьмя женщинами (\textbf{a, b, c, d}).
Порядок предпочтений представлен в таблицах.
\begin{center}
\begin{tabular}{|c|c||c|c|}
\hline
Мужчины & Предпочтения & Женщины & Предпочтения \\
\hline
\textbf{A} &  \textbf{c b d a} & \textbf{a} & \textbf{A B D C} \\ 
\textbf{B} &  \textbf{b a c d} & \textbf{b} & \textbf{C A B D} \\
\textbf{C} &  \textbf{b d a c} & \textbf{c} & \textbf{C B D A} \\
\textbf{D} &  \textbf{c a d b} & \textbf{d} & \textbf{B A C D} \\
\hline
\end{tabular}
\end{center}
Паросочетание (\textbf{Aa, Bb, Cc, Dd}) неустойчиво, так как \textbf{C} и \textbf{b} предпочитают друг друга своим супругам.
\end{example}

\medskip
Задача заключается в нахождении устойчивого паросочетания по данным спискам предпочтений.
Иными словами, нужно в полном двудольном графе $K_{n,n}$ найти совершенное устойчивое паросочетание, если оно существует.

Решение задачи существует, как показывает следующий алгоритм,
который называют \em{алгоритмом отложенного согласия}.
\end{frame}


\begin{frame}

\frametitle{8.3.4. Устойчивые паросочетания (3/8)}

В алгоритме используются следующие структуры данных.

{\color{blue}\bf 1.}
Матрица предпочтений мужчин \textbf{$M:$ array $[1..n,1..n]$
of $1..n$}. Элемент  $M[m,p]$ означает \em{номер} женщины, стоящей
на месте $p$
в списке предпочтений мужчины $m$.

{\color{blue}\bf 2.}
Матрица предпочтений женщин \textbf{$W:$ array $[1..n,1..n]$
of $1..n$}. Элемент  $W[w,m]$ означает 
\em{рейтинг} мужчины $m$
в списке предпочтений женщины $w$.
Данные примера представляются 
следующими матрицами:
$$
M=\begin{pmatrix}
3 & 2 & 4 & 1 \\
2 & 1 & 3 & 4 \\
2 & 4 & 1 & 3 \\
3 & 1 & 4 & 2 \\
\end{pmatrix}, \quad 
W=\begin{pmatrix}
4 & 3 & 1 & 2 \\
3 & 2 & 4 & 1 \\
1 & 3 & 4 & 2 \\
3 & 4 & 2 & 1 \\
\end{pmatrix}.
$$


{\color{blue}\bf 3.}
Список женихов \textbf{$B:$ array $[1..n]$
of $0..n$}, 
где $B[m]$~--- номер невесты, помолвленной с женихом $m$. 

{\color{blue}\bf 4.}
Список предложений \textbf{$P:$ array $[1..n]$
of $0..n$},  
где $P[m]$~--- количество предложений, уже сделанных женихом $m$.

%{\color{blue}\bf 5.}
%Список невест \textbf{$G:$ array $[1..n]$
%of $0..n$}, 
%где $G[w]$~--- номер жениха, помолвленного с невестой $w$.
%
%{\color{blue}\bf 6.}
%Текущий рейтинг \textbf{$R:$ array $[1..n]$
%of $0..n$}, 
%где $R[w]$~--- рейтинг жениха, помолвленного с невестой $w$.

\end{frame}


\begin{frame}%[plain]
\frametitle{8.3.4. Устойчивые паросочетания (4/8)}
{\color{blue}\bf 5.}
Список невест \textbf{$G:$ array $[1..n]$
of $0..n$}, 
где $G[w]$~--- номер жениха, помолвленного с невестой $w$.

{\color{blue}\bf 6.}
Текущий рейтинг \textbf{$R:$ array $[1..n]$
of $0..n$}, 
где $R[w]$~--- рейтинг жениха, помолвленного с невестой $w$.


\begin{algorithmic}
\REQUIRE $n$  число мужчин и женщин, матрицы предпочтений $M, W$.
%\vspace{-1pt}
\ENSURE $n$ сформированных пар в списках $B, G$.
%\vspace{-1pt}
\STATE{$k\Assign 0$}
\COMMENT{\color{gray}\small количество сформированных пар }
\STATE{\textbf{for $i$ from $1$ to $n$ do 
$B[i]\Assign 0; G[i]\Assign 0; P[i]\Assign 0;$ end for}}
%\COMMENT{\color{gray}\small инициализация }
%\vspace{-1pt}
%\WHILE{$k < n$}
\STATE{\textbf{while $k < n$ do}}
%\vspace{-1pt}
    \FOR{\textbf{$m$ from $1$ to $n$}}
    	\IF{$B[m]=0$}
    	\STATE{}\COMMENT{\color{gray}\small  первый свободный жених $m$ делает предложение невесте $w$}
    	\STATE{$P[m]\Assign P[m]+1; w\Assign M[m,P[m]]$}
    	\STATE{\textbf{exit for $m$}}
    	\COMMENT{\color{gray}\small  цикл не нужно продолжать }
    	\ENDIF
%\vspace{-1pt}
    \ENDFOR
%\vspace{-1pt}    
\end{algorithmic}


\end{frame}

\begin{frame}
\frametitle{8.3.4. Устойчивые паросочетания (5/8)}
\begin{algorithmic}
    \IF{$G[w]=0$}
    	\STATE{}\COMMENT{\color{gray}\small  первое предложение всегда принимается~--- помолвка}
        \STATE{$k\Assign k+1; G[w]\Assign m, B[m]\Assign w;
        R[w]\Assign W[w,m]$}
        \STATE{\textbf{next while $k<n$}} 
        \COMMENT{\color{gray}\small  следующий жених} 
    \ENDIF
%\vspace{-1pt}    
    \IF{$W[w,m]>R[w]$}
        \STATE{}\COMMENT{\color{gray}\small  новый жених лучше~--- прежнему отставка} 
        \STATE{$ B[G[w]]\Assign 0; B[m]\Assign w; R[w]\Assign W[w,m]; G[w]\Assign m$}
%\vspace{-1pt}
    \ENDIF
%\vspace{-1pt}
%\ENDWHILE
\STATE{\textbf{end while}}
\end{algorithmic}
\vfill
\end{frame}

\begin{frame}
\frametitle{8.3.4. Устойчивые паросочетания (6/8)}
\begin{ground}
\proofsection{Завершаемость}
На каждой итерации рассматривается одно предложение,
а всего возможно не более $n^2$ предложений.
\proofsection{Паросочетаемость}
На каждой итерации для ненулевых элементов
инвариантно выполняется условие:
$G[w]=m\Equiv B[m]=w$, пары образуют паросочетание.
\proofsection{Совершенность}
Всего создано $n$ пар, значит паросочетание совершенное.
\proofsection{Устойчивость} От противного.
Пусть в найденном паросочетании \textbf{Xx, Yy}
существует \\неустойчивая комбинация \textbf{X--y}.
Возможны два случая. Мужчина \textbf{X} 
не делал предложения женщине \textbf{y}.
Но это означает, что \textbf{x} предшествует \textbf{y}
в списке предпочтений \textbf{X}, что противоречит 
неустойчивости. Мужчина \textbf{X} 
делал предложение женщине \textbf{y} и 
она ему отказала.
Но это означает, что \textbf{Y} предшествует \textbf{X}
в списке предпочтений \textbf{y}, что противоречит 
неустойчивости.      
\end{ground}
\end{frame}

\begin{frame}
\frametitle{8.3.4. Устойчивые паросочетания (7/8)}
\begin{example}
Для рассматриваемого примера 
$$
M=\begin{pmatrix}
3 & 2 & 4 & 1 \\
2 & 1 & 3 & 4 \\
2 & 4 & 1 & 3 \\
3 & 1 & 4 & 2 \\
\end{pmatrix}, \quad 
W=\begin{pmatrix}
4 & 3 & 1 & 2 \\
3 & 2 & 4 & 1 \\
1 & 3 & 4 & 2 \\
3 & 4 & 2 & 1 \\
\end{pmatrix},
$$
алгоритм последовательно построит 
следующие паросочетания:\\
$\emptyset\mapsto\{\textbf{Ac}\}\mapsto\{\textbf{Ac,Bb}\}
\mapsto\{\textbf{Ac, Cb}\}\mapsto\{\textbf{Ac,Cb,Ba}\}\mapsto\{\textbf{Cb, Ba, Dc}\}\mapsto\{\textbf{Cb, Ba, Dc, Ad}\}$.
\end{example}

\begin{remark}
Алгоритм находит одно из многих возможных решений задачи. В данном случае алгоритм строит наилучший вариант для мужчин 
(для каждого мужчины выбранная алгоритмом супруга
наиболее привлекательна для него из всех возможных
устойчивых паросочетаний) и наихудший для  женщин
(выбранный супруг нравится каждой женщине не больше, 
чем возможный супруг в любом другом устойчивом паросочетании).
\end{remark}
\end{frame}

\begin{frame}
\frametitle{8.3.4. Устойчивые паросочетания (8/8)}

\begin{digress}
Рассмотрим задачу об <<устойчивом зачислении>> $n$ абитуриентов в $m$ университетов.
Пусть университет с номером $k$ может принять $n_k$  абитуриентов. 
Не ограничивая общности можно считать, что $n_1 + n_2 + \ldots + n_m = n$,
рассматривая, если нужно, самый нежелательный <<никакой>> университет 
с неограниченным приёмом.
Каждый абитуриент имеет свой рейтинг университетов, 
каждый университет~--- свой рейтинг абитуриентов. 

Можно говорить об <<устойчивом зачислении>>, 
если никакой абитуриент \textbf{A}, не зачисленный в университет \textbf{a}, но предпочитающий его своему ВУЗу, не окажется предпочтительнее для университета \textbf{a}, чем кто-то из его студентов.
\end{digress}

%{\color{red}
%Запрос: развить этот пример,
%поставив и решив задачу 
%о <<справедливом зачислении>>. 
%Здесь возможны варианты от одного абзаца на слайд
%до статьи в приличном журнале.  
%}
\end{frame}




\begin{frame}
\label{8.3.5.}\subsubsection{8.3.5. Наибольшее паросочетание }\frametitle{8.3.5. Наибольшее паросочетание (1/5)}

Ясно, что наибольшее паросочетание является максимальным, 
но не всякое максимальное паросочетание является наибольшим.
\begin{columns}[t]
\begin{column}[T]{8cm}
\begin{example} 
На рисунке показаны две диаграммы графа, в котором 
паросочетание $\{(b,d)\}$ является максимальным, но не является наибольшим, а паросочетание $\{(a,d), (c,b)\}$ является наибольшим и максимальным.
\end{example}
\end{column}
\begin{column}[T]{7cm}
\vspace{-24pt}
\begin{figure}[h]
\begin{center}
\includegraphics[width=30mm]{8_3_4.jpg}
\end{center}
\end{figure}\end{column}
\end{columns}

\vspace{8pt}
Простая цепь называется \em{альтернирующей}
относительно паросочетания $M$, если она попеременно содержит 
рёбра из $M$ и не из $M$.

\begin{example}
На рисунке цепь $\langle (a,d), (d,b), (b,c)\rangle$
является альтернирующей относительно обоих паросочетаний.
\end{example}

Альтернирующая цепь называется \em{аугментальной}
относительно паросочетания $M$, если её крайние рёбра не принадлежат $M$.

\begin{example}
На рисунке альтернирующая цепь $\langle (a,d), (d,b), (b,c)\rangle$
является аугментальной  только относительно паросочетания
 $\{(b,d)\}$.
\end{example}


\end{frame}




\begin{frame}
\frametitle{8.3.5. Наибольшее паросочетание (2/5)}
Вершины, инцидентные рёбрам паросочетания $M$
естественно назвать
\em{занятыми} в этом паросочетании, 
а прочие вершины можно назвать \em{свободными}. 
\begin{example}
На рисунке на предыдущем слайде вершины 
$a$ и $c$ свободны в паросочетании
 $\{(b,d)\}$, а в паросочетании $\{(a,d), (c,b)\}$
 все вершины заняты.
\end{example}


\begin{lemma}
Если существует аугментальная цепь, то паросочетание не наибольшее.
\end{lemma}

\begin{proof}
Пусть $M$~--- паросочетание, а $P$~--- аугментальная цепь.
Здесь $P$ рассматривается как \em{множество} рёбер,
последовательность не имеет значения. Тогда никакое 
ребро в множестве $M\setminus(P\cap M)$ не может быть 
смежно ни с каким ребром в цепи $P$,
поскольку две концевые вершины свободны по определению,
и они не инцидентны никаким рёбрам из $M$,
а все остальные вершины заняты по построению,
и они не могут быть инциденты ещё каким-то рёбрам из 
паросочетания $M$.
Следовательно, множество рёбер
$$
M'\Assign (M\setminus(P\cap M))\cup(P\setminus(P\cap M))
=(M\cup P)\setminus(P\cap M)
$$
является паросочетанием.
Но так как оба концевых ребра из
$P$ не лежат в $M$, 
то $|P\setminus(P\cap M)|=|P\cap M|+1$ и
 $|M'|=|M|+1$, так что $M$ 
не является наибольшим паросочетанием.
\end{proof}

\end{frame}

\begin{frame}
\frametitle{8.3.5. Наибольшее паросочетание (3/5)}

\begin{remark}
Паросочетание $M'$ получается
из паросочетания $M$ изменением
<<ролей>> всех рёбер в аугментальной цепи:
те рёбра, которые входили в паросочетание,
исключаем из паросочетания,
а те, которые не входили, наоборот, включаем.  
\end{remark}

\begin{columns}[t]
\begin{column}[T]{7.2cm}
\begin{example}
В условиях предыдущих примеров
паросочетание $\{(a,d), (c,b)\}$ получается из паросочетания  $\{(b,d)\}$
изменением ролей рёбер в аугментальной цепи 
$\langle (a,d), (d,b), (b,c)\rangle$.
\end{example}
\end{column}
\begin{column}[T]{7cm}

\vspace{-16pt}
\begin{center}
\includegraphics[width = 70mm]{8_3_5_1.jpg}
\end{center}
\end{column}
\end{columns}

\medskip
Наибольшее паросочетание может быть \em{неединственным}.

\medskip
\begin{columns}[t]
\begin{column}[T]{8cm}
\begin{example} 
В верхней диаграмме  
наибольшее паросочетание $\{(a,e), (b,d), (c,f)\}$ единственно,  
а в нижней диаграмме существуют два наибольших паросочетания:
 $\{(a,d), (c,b)\}$  и  $\{(a,b), (c,d)\}$.
\end{example}
\end{column}
\begin{column}[T]{7cm}
\vspace{-20pt}
\begin{figure}[h]
\begin{center}
\includegraphics[width=50mm]{8_3_4_2.jpg}
\end{center}
\end{figure}\end{column}
\end{columns}


\end{frame}

\begin{frame}
\frametitle{8.3.5. Наибольшее паросочетание (4/5)}
\begin{theorem}
Паросочетание является наибольшим тогда и только тогда, когда
для него не существует аугментальной цепи.
\end{theorem}

\begin{proof}
\proofsection{Необходимость}
Если паросочетание наибольшее, то аугментальная цепь не существует
по лемме.

\proofsection{Достаточность}
Пусть в графе $G$ задано некоторое паросочетание $M$, 
для которого нет ни одной
аугментальной цепи, а $M^*$~--- наибольшее паросочетание 
в этом графе.
Рассмот\-рим остовный подграф $G'(V,E')$, в который включим 
рёбра симметрической разности паросочетаний,
$E'\Assign (M\cup M^*)\setminus (M\cap M^*)$.
В графе $G'$ степени всех вершин не больше двух, 
поскольку любая вершина инцидентна не более чем одному ребру
из $M$ и не более чем одному ребру
из $M^*$. Таким образом, граф $G'=\bigcup G_i$,
где каждый подграф $G_i$ либо
 изолированная вершина, либо простая цепь, 
либо простой цикл.  
При этом как цепи, так и циклы являются
альтернирующими как для $M$, так и для $M^*$.
Заметим, что изолированные вершины появляются, 
если $M\cap M^*\ne\emptyset$.
\end{proof}

\end{frame}

\begin{frame}
\frametitle{8.3.5. Наибольшее паросочетание (5/5)}
%\vspace{-2pt}
\begin{proof} (Продолжение)
Пусть $N(G_i)$~--- число рёбер из $M$, а
$N^*(G_i)$~--- число рёбер из $M^*$, входящих 
в  подграф $G_i$. 
Если $G_i$~--- изолированная вершина, то $N(G_i)=N^*(G_i)=0$.  
Если $G_i$~---  альтернирующая цепь, то она не может быть 
аугментальной для $M$ по условию теоремы, 
и не может быть 
аугментальной для $M^*$ по лемме. Значит, концевые
рёбра $G_i$ принадлежат различным альтернативам, 
стало быть их чётное число, причём $N(G_i)=N^*(G_i)$.
Если $G_i$~---  альтернирующий цикл, то он является чётным 
(иначе этот цикл не был бы альтернирующим),
следовательно, опять имеем $N(G_i)=N^*(G_i)$.
Таким образом,
$$
|M| = |M\cap M^*|+\sum_i N(G_i) = |M\cap M^*|+\sum_i N^*(G_i) =|M^*|
$$
и паросочетание $M$~--- наибольшее.   
\end{proof}

\begin{remark}
Паросочетание называется \odmkey{полным}{Паросочетание (matching)!полное (complete)},
если оно не оставляет свободных вершин.
Полное паросочетание, если оно существует,
является наибольшим.
\end{remark}

\end{frame}


\begin{frame}
\label{8.3.6.}\subsubsection{8.3.6. Алгоритм поиска наибольшего паросочетания }\frametitle{8.3.6. Алгоритм поиска наибольшего паросочетания (1/4) }
На основании доказанной теоремы  можно построить 
алгоритм отыскания наибольшего паросочетания.
Для этого достаточно взять какое-то паросочетание,
например, пустое, и 
далее попытаться найти аугментальную цепь.
Если аугментальная цепь найдена, 
то паросочетание перестраивается в паросочетание, содержащее
на одно ребро больше текущего. 
Если же аугментальной цепи нет, 
то текущее паросочетание наибольшее.  

\smallskip
\begin{remark}
В наибольшем паросочетании не 
более $p/2$ рёбер, так что искать
аугментальные цепи потребуется не более,
чем $p/2$ раз.
Однако, только простых цепей в графе 
может быть до $p(p-1)/2$, и каждая
цепь может содержать до $p-1$ рёбер,
так что поиск наибольшего паросочетания может оказаться трудоёмким.    
\end{remark}

\smallskip
Алгоритм отыскания наибольшего паросочетания 
оказывается наиболее простым в случае
двудольных графов. 
Пусть имеется двудольный граф $G(V_1,V_2,E)$, 
причём $|V_1|=p_1$, $|V_2|=p_2$, $p_1\leq p_2$
и для каждой вершины $v$ задано множество
смежных вершин $\Gamma(v)$.

 \end{frame}


\begin{frame}
\frametitle{8.3.6. Алгоритм поиска наибольшего паросочетания (2/4) }    

Алгоритм основан на идее поиска в глубину,
для отметки вершин используется массив
$x$: \textbf{array} $[1..p_1]$ \textbf{of bool}.
Паросочетание представлено массивом
$M$: \textbf{array} $[1..p_2]$ \textbf{of} $0..p_1$,
где $M[u]=0$, если вершина $u$ свободна
в паросочетании, и $M[u]=v$, если ребро $(v,u)$ входит
в паросочетание.
   

\begin{algorithmic}
\REQUIRE двудольный граф $G(V_1,V_2,E)$.
\ENSURE наибольшее паросочетание 
$M$: \textbf{array} $[1..p_2]$ \textbf{of} $0..p_1$.
%\STATE{}\COMMENT{\color{gray}\small инициализация }
\FOR{\textbf{$u$ from $1$ to $p_2$}} 
\STATE{$M[u]\Assign 0$}
\COMMENT{\color{gray}\small  паросочетание пусто }
\ENDFOR
\STATE{}\COMMENT{\color{gray}\small запускаем поиск в глубину из всех вершин первой доли }
\FOR{\textbf{$v$ from $1$ to $p_1$}}
    \FOR{\textbf{$w$ from $1$ to $p_1$}}
    \STATE{\textbf{$x[w]\Assign$ false}}
    \COMMENT{\color{gray}\small перед поиском всё не отмечено }
    \ENDFOR
  	\STATE{$\Aug(v)$}
   	\COMMENT{\color{gray}\small поиск в глубину от вершины $v$ }
\ENDFOR
\end{algorithmic}

 \end{frame}


\begin{frame}
\frametitle{8.3.6. Алгоритм поиска наибольшего паросочетания (3/4) }    

Поиск в глубину выполняется рекурсивной функцией 
$\Aug$, которая так названа потому, что ищет аугментальную 
цепь и перестраивает паросочетание, если находит такую цепь.

\begin{algorithmic}
\REQUIRE вершина $v$ первой доли.
\ENSURE \textbf{true}, если 
аугментальная цепь найдена и паросочетание перестроено,
\textbf{false} в противном случае.  
\IF{$x[v]$}
    	\STATE{\textbf{return false}}
    	\COMMENT{\color{gray}\small  вершина занята }
\ENDIF
\STATE{$x[v]\Assign$ \textbf{true}}
\COMMENT{\color{gray}\small занимаем вершину }
\FOR{$u\in\Gamma(v)$}
    \IF{$M[u]=0 \Or \Aug(M[u])$}
    	\STATE{$M[u]\Assign v$}
    	\COMMENT{\color{gray}\small  добавляем ребро в паросочетание }
    	\STATE{\textbf{return true}}
    	\COMMENT{\color{gray}\small  успех }
    \ENDIF
\ENDFOR
\STATE{\textbf{return false}}
\COMMENT{\color{gray}\small  неудача }
\end{algorithmic}

 \end{frame}


\begin{frame}
\frametitle{8.3.6. Алгоритм поиска наибольшего паросочетания (4/4) }

\begin{example}
Дан двудольный граф $G=(V_1, V_2, E)$. 
$V_1=\{1, 2, 3\}$, 
$V_2=\{1, 2, 3\}$, \\
$E=\{(1, 1), (1, 2), (1, 3), (2, 1), (2, 2), (3, 1)\}$.
Паросочетания  в процессе работы функции 
Aug перестраиваются следующим образом:
$\{(1, 1)\} \to$ $\{(1, 1), (1,2)\} \to$ 
$\{(2, 1), (1, 2)\} \to$ 
$\{(2, 1), (1, 2), (1, 3)\} \to$ 
$\{(2, 1), (2, 2),(1, 3)\} \to \{(2, 2), (1, 3), (3, 1)\}.
$
\end{example}

\begin{digress}
Задача отыскания наибольшего паросочетания 
имеет множество интерпретаций, 
которым соответствуют математические 
модели и алгоритмы различной сложности.
Например, рассмотренный алгоритм можно
применить к решению следующей задачи:
имеется множество работников, каждый
из которых может выполнять некоторое подмножество
множества работ. Как трудоустроить
наибольшее количество работников?
Рёбрам в графе могут быть назначены веса,
и тогда возникает задача 
найти паросочетание наибольшего суммарного веса,
известная как <<\em{задача о назначениях}>>.
Вес ребра можно интерпретировать
как эффективность работника при выполнении 
данной работы, и тогда решение задачи о назначениях
означает наиболее эффективное распределение работ.
Для решения задачи о назначениях был предложен
<<\em{венгерский алгоритм}>>, один из первых и самых
известных алгоритмов оптимизации.
В графах общего вида могут быть циклы 
нечётной длины, что привносит
дополнительные сложности в отыскание паросочетаний 
(\odmkey{алгоритм Эдмондса}
{Алгоритм (algorithm)!Эдмондса (Edmonds)}) и т.~д.   
\end{digress}  
   

\end{frame}



\begin{frame}
\frametitle{Онтология теории графов}
\begin{center}
\includegraphics[width=14cm]{ODM_8_3.png}
\end{center}

\end{frame}


%\vspace{-2pt}

\begin{frame}
\frametitle{8.4. Потоки в сетях (1/3)}
\subsection{8.4. Потоки в сетях}
\label{8.4.}

Рассмотрим некоторые примеры практических задач.

%\vspace{-2pt}

%\begin{examples}
%\begin{enumerate}
%\item
\medskip
{\color{blue}\bf 1.}
Пусть имеется сеть автомобильных
дорог, по которым можно проехать из пункта
$A$ в пункт $B$.
Дороги могут пересекаться в промежуточных пунктах.
Количество автомобилей, которые могут
проехать по каждому отрезку дороги за единицу времени,
не безгранично: оно определяется такими факторами,
как ширина проезжей части, качество дорожного
покрытия, действующие ограничения скорости движения и т.~д.
(обычно это называют <<пропускной способ\-ностью>> дороги).
Каково максимальное количество автомобилей,
которые могут проехать из пункта $A$ в пункт $B$
за единицу времени без образования
пробок на дорогах (обычно это называют <<автомобильным потоком>>)?
Можно поставить и другой вопрос: какие дороги
и насколько нужно расширить или улучшить,
чтобы увеличить максимальный
автомобильный поток на заданную величину?
\end{frame}


\begin{frame}
\frametitle{8.4. Потоки в сетях (2/3)}
%\item 
{\color{blue}\bf 2.}
Пусть имеется сеть трубопроводов, соединяющих пункт $A$
(скажем, нефтепромысел) с пунктом $B$ (скажем,
нефтеперерабатывающим заводом). Трубопроводы могут соединяться и
разветвляться в промежуточных пунктах. Количество нефти, которое
может быть перекачано по каждому отрезку трубопровода в единицу
времени, также не безгранично и определяется такими факторами, как
диаметр трубы, мощность нагнетающего насоса и т.~д. (обычно это
называют <<пропускной способностью>> или <<максимальным расходом>>
трубопровода). Сколько нефти можно прокачать через такую сеть за
единицу времени?

\medskip
{\color{blue}\bf 3.}
Пусть имеется система машин для производства готовых изделий
из сырья и~последовательность технологических операций может быть
различной (то есть операции могут выполняться на разном
оборудовании или в разной последовательности), но все допустимые
варианты заранее строго определены. Максимальная
производительность каждой единицы оборудования, естественно, также
заранее известна и постоянна. Какова максимально возможная
производительность всей системы в целом, и как нужно организовать
производство для достижения максимальной производительности?
%
%
%\item 
\end{frame}


\begin{frame}
\frametitle{8.4. Потоки в сетях (3/3)}

%\end{enumerate}
%\end{examples}



Изучение этих и многочисленных подобных им практических задач
приводит к~теории потоков в сетях. В данном разделе
рассматривается решение только одной (но самой существенной)
задачи этой теории, а именно: задачи определения максимального
потока. 
%Описание других родственных задач, например задачи
%определения критического пути, можно найти в литературе,
%перечисленной в~конце главы.

\medskip
\begin{tabular}{p{9mm}|p{55mm}|p{80mm}}
  № & Название параграфа
  & Ключевые термины и обозначения \\
  \hline
\DMref{8.4.1.}{8.4.1.}\RS & 
Определение потока 
  & 
  {$c(u,v)$~--- пропускная способность дуги;
  $w(f)$~--- величина потока в сети;
  $\operatorname{div}(f,u)$~--- дивергенция потока в узле; 
  насыщенная дуга} \\
\DMref{8.4.2.}{8.4.2.} & Разрезы 
  & 
  {$C(P)$~--- пропускная способность разреза}
  \\
\DMref{8.4.3.}{8.4.3.} & Теорема Форда--Фалкерсона &
{Аугментальная цепь}\\
\DMref{8.4.4.}{8.4.4.} & Алгоритм нахождения максимального потока &
{} \\
\DMref{8.4.5.}{8.4.5.}\BS & Связь между теоремой Менгера и теоремой Форда--Фалкерсона &
{} \\
\end{tabular}

\end{frame}


\begin{frame}
\label{8.4.1.}\subsubsection{8.4.1. Определение потока }\frametitle{8.4.1. Определение потока (1/2)}

Пусть $G(V,E)$~--- сеть, $s$ и $t$~--- источник и
сток сети соответственно. Дуги сети нагружены неотрицательными вещественными
числами, $\Map{c}{E}{\Bbb {R}_{+}}$. Если $u$ и $v$~--- узлы сети,
то число $c(u,v)$~--- называется \odmkey{пропускной
способностью}
{Пропускная способность (capacity)!дуги (of arc)}
дуги $(u,v)$.

\begin{remark}
Матрица пропускных способностей \textbf{$C:$ array $[1..p,1..p]$ of real}
является представлением сети, аналогичным матрице смежности.
Нулевое значение элемента $C[u,v]$ соответствует дуге с нулевой
пропускной способностью,
то~есть отсутствию дуги, а
положительное значение элемента $C[u,v]$
 соответствует дуге с ненулевой пропускной способностью,
то~есть дуга присутствует.
\end{remark}



Пусть задана функция $\Map{f}{E}{\Bbb{R}_+}$.
\odmkey{Дивергенцией}{Дивергенция (divergence)} функции $f$ в
узле $u$ называется число $\operatorname{div}(f,u)$, которое
определяется следующим образом:
$$
\operatorname{div}(f,u)\Def \sum_{\Set{v}{(u,v)\in E}} f(u,v)
   -   \sum_{\Set{v}{(v,u)\in E}} f(v,u).
$$
\end{frame}


\begin{frame}
\frametitle{8.4.1. Определение потока (2/2)}
\begin{remark}
В физике дивергенция обычно определяется наоборот:
то, что пришло, минус то, что ушло.
Но в данном случае удобнее,
чтобы дивергенция источника была положительна.
\end{remark}



Функция $\Map{f}{E}{\Bbb R}$ называется \odmkey{потоком}{Поток
(flow)} в сети $G$, если

\begin{enumerate}\leftskip2pt
\item[1)] $\A{(u,v)\in E}{ 0 \le f(u,v) \le c(u,v)}$,
то~есть поток через дугу
неотрицателен и не превосходит пропускной способности дуги;

\item[2)] $\A{u\in V\setminus\{s,t\}}{\operatorname{div}(f,u) = 0}$,
то~есть дивергенция потока равна нулю во всех узлах,
кроме источника и стока.
\end{enumerate}

Число $w(f)\Def\operatorname{div}(f,s)$ называется
\odmkey{величиной}{Величина потока (value of flow)} потока $f$.
Если $f(u,v)=c(u,v)$, то дуга $(u,v)$ называется
\odmkey{насыщенной}{Дуга (arc)!насыщенная (saturated)}.



\end{frame}


\begin{frame}
\label{8.4.2.}\subsubsection{8.4.2. Разрезы}
\frametitle{8.4.2. Разрезы}

Пусть $P$~--- $(s,t)$-разрез, $P\subset E$. Всякий разрез
определяется разбиением множества вершин $V$ на два подмножества,
$S$ и $T$, так, что $S\subset V$, $T\subset V$, $S\cup T=V$, $S\cap
T=\emptyset$, \\$s\in S$, $t\in T$, а в $P$ попадают все дуги,
соединяющие $S$ и $T$. Тогда $P=P^+\cup P^-$, где $P^+$~--- дуги
от $S$ к $T$,
 $P^-$~--- дуги от $T$ к $S$.
Сумма потоков через дуги разреза $P$ обозначается $F(P)$. Сумма
пропускных способностей дуг разреза~$P$ называется
\odmkey{пропускной способностью}{Пропускная способность
(capacity)!разреза (of cut)} разреза и обозначается $C(P)$:
$$
F(P)\Def\sum_{e\in P}f(e), \qquad
C(P)\Def\sum_{e\in P}c(e).
$$

Аналогично, $F(P^+)$ и $F(P^-)$~--- суммы потоков через
соответствующие части разрезов (<<положительную>> и
<<отрицательную>>).

\end{frame}


\begin{frame}
\label{8.4.3.}\subsubsection{8.4.3. Теорема Форда--Фалкерсона }\frametitle{8.4.3. Теорема Форда--Фалкерсона (1/4)}
%\label{s853}

В этом параграфе устанавливается
количественная связь между величиной
максимального потока и
пропускными способностями дуг сети.

\medskip
%\setcounter{varlemma}{0}
\begin{lemma}  \NB{[1]}
$w(f) = F(P^+)-F(P^-)$.
\end{lemma}

\begin{proof}
Рассмотрим сумму $W\Assign\sum_{v\in S}\operatorname{div}(f,v)$.
Пусть дуга $(u,v)\in E$. Если $u,v\in S$, то в сумму $W$ попадают два
слагаемых для этой дуги: $+f(u,v)$ при вычислении
$\operatorname{div}(f,u)$ и $-f(u,v)$ при вычислении
$\operatorname{div}(f,v)$, итого $0$. Если $u\in S$, $v\in T$, то в
сумму $W$ попадает одно слагаемое $+f(u,v)$, все такие слагаемые дают
$F(P^+)$. Если $u\in T$, $v\in S$, то в сумму $W$ попадает одно
слагаемое $-f(u,v)$, все такие слагаемые дают 
$-F(P^-)$. Таким образом,
$W=F(P^+)-F(P^-)$. С другой стороны,\\
$W=\sum_{v\in S}\operatorname{div}(f,v)=\operatorname{div}(f,s)=w(f)$.
\end{proof}
%
%\vspace*{-.5\baselineskip}

\medskip
\begin{lemma}  \NB{[2]}
$\operatorname{div}(f,s)=-\operatorname{div}(f,t)$.
\end{lemma}


%\addvspace{-6pt}
\begin{proof}
Рассмотрим разрез $P\Assign (S,T)$, где
$S\Assign V-t$, а $T\Assign\{t\}$. \\Имеем
$
\operatorname{div}(f,s)=w(f)=F(P^+)-F(P^-)=F(P^+)=
\displaystyle{\sum_{\Set{v}{(v,t)\in E}}f(v,t)} =
-\operatorname{div}(f,t).
$
\end{proof}
%
%
%\vspace*{-.5\baselineskip}
\end{frame}


\begin{frame}
\frametitle{8.4.3. Теорема Форда--Фалкерсона (2/4)}
\begin{lemma}  \NB{[3]}
$w(f)\le F(P)$.
\end{lemma}

\begin{proof}
$w(f)=F(P^+)-F(P^-)\leq F(P^+)\leq F(P)$.
\end{proof}

\smallskip
\begin{lemma}  \NB{[4]}
$\max\limits_f w(f) \le \min\limits_P C(P)$.
\end{lemma}

\begin{proof}
По предыдущей лемме
$w(f)\leq F(P)$, следовательно,\\
$\max\limits_f w(f) \leq\min\limits_P F(P)$. %\\
По определению
$F(P)\leq C(P)$, следовательно,\\
$\min\limits_P F(P)\leq\min\limits_P C(P)$.
Имеем
$\max\limits_{f}w(f)\leq\min\limits_{P}C(P)$.
\end{proof}

%
%\vspace*{-.75\baselineskip}
\smallskip
\begin{theorem}
Максимальный поток в сети равен пропускной способности
положительной части минимального разреза, то~есть существует поток $f^*$
такой, что \\
$$
w(f^*)=\max_f w(f)=\min_P C(P^+).
$$
\end{theorem}

%\addvspace{-6pt}
\begin{proof}
Пусть $f$~--- некоторый максимальный поток. Покажем, что
существует разрез $P$ такой, что $w(f)=C(P)$. Рассмотрим граф
$G'$, полученный из сети $G$ забыванием ориентации дуг. 
\end{proof}
\end{frame}


\begin{frame}
\frametitle{8.4.3. Теорема Форда--Фалкерсона (3/4)}
\begin{proof} (Продолжение)
Построим
множество вершин $S$ следующим образом:
\begin{equation*}
\begin{split}
S\Def &\{ u\in V \mid \exists \divC{s}{u}\in G'\;
                  \forall (u_i,u_{i+1})\in\divC{s}{u}\in G'\\
(&(u_i,u_{i+1})\in E\Impl f(u_i,u_{i+1})<C(u_i,u_{i+1})\And\\
      &\And(u_{i+1},u_i)\in E\Impl f(u_{i+1},u_i)>0 ) \},
\end{split}
\end{equation*}
то~есть вдоль пути $\divC{s}{u}$ дуги в направлении пути не
насыщены, а дуги против направления пути имеют положительный
поток. Такая цепь $\divC{s}{u}$ называется
\odmkey{аугментальной}
{Цепь (circuit)!аугментальная (augmenting)}.
Имеем $s\in S$ по построению. Следовательно, $S\not=\emptyset$.
Положим $T\Assign V\setminus S$. Покажем, что $t\in T$.
Действительно, пусть $t\in S$. Тогда существует аугментальная цепь
$\divC{s}{t}$, обозначим её $R$. Но тогда можно найти число
$\delta>0$:
$$
\delta\Assign \min\limits_{e\in R}\Delta(e),\text{ где }
\Delta(e)\Assign
\begin{cases}
c(e)-f(e)>0,& \text{если $e$  вдоль $R$,}\\
f(e)>0,         & \text{если $e$  против $R$.}
\end{cases}
$$
\end{proof}
\end{frame}


\begin{frame}
\frametitle{8.4.3. Теорема Форда--Фалкерсона (4/4)}
\begin{proof} (Продолжение)
В этом случае можно увеличить величину потока на $\delta$, изменив
поток для всех дуг аугментальной цепи:
$$
f(e)\Assign
\begin{cases}
f(e)+\delta,& \text{если $e$ ориентировано вдоль $R$,}\\
f(e)-\delta,& \text{если $e$ ориентировано против $R$.}
\end{cases}
$$
При этом условия потока выполняются: $0\le f(e)\leq C(e)$,
$\operatorname{div}(v)=0$.



Таким образом, поток $f$ увеличен на величину $\delta$, что
противоречит максимальности потока $f$. Имеем $t\in T$ и
$T\not=\emptyset$. Следовательно, $S$ и $T$ определяют разрез
$P=P^+\cup P^-$. В этом разрезе все дуги $e^+$ насыщены
$(f(e^+)=c(e^+))$, а все дуги $e^-$ не нагружены $(f(e^-)=0)$,
иначе $S$ можно было бы расширить. Имеем
$w(f)=F(P^+)-F(P^-)=C(P^+)$, таким образом, $P$~--- искомый
разрез.
\end{proof}


%\addvspace{-6pt}
\end{frame}


\begin{frame}
\label{8.4.4.}\subsubsection{8.4.4. Алгоритм нахождения максимального потока }\frametitle{8.4.4. Алгоритм нахождения максимального потока (1/5)}

%%\addvspace{-3pt}
Следующий алгоритм \odmkey{нахождения максимального потока}
{Алгоритм (algorithm)!нахождения
максимального потока (for
maximum flow problem)} определяет максимальный поток
в сети, заданной ма\-три\-цей пропускных способностей дуг. Этот
алгоритм использует идею доказательства теоремы Форда--Фалкерсона,
а именно, задавшись начальным приближением потока, определяет
множество вершин $S$, которые соединены аугментальными цепями с
источником $s$. Если оказывается, что $t\in S$, то это означает,
что поток не максимальный и его можно увеличить на величину
$\delta$. Для определения аугментальных цепей и одновременного
подсчёта величины $\delta$ в алгоритме использованы
 структуры данных:
%
%
%
%{
%\small
%\vspace{3pt}
\begin{algorithmic}
\STATE \textbf{$P:$ array $[1..p]$ of struct $\{$ }
\STATE \quad\textbf{$s:$ enum $(-,+)$} 
\COMMENT{\color{gray}\small <<знак>>, то~есть направление дуги }
\STATE \quad\textbf{$n: 1..p$} 
\COMMENT{\color{gray}\small предшествующий узел в аугментальной цепи }
\STATE \quad\textbf{$\delta:$ real} 
\COMMENT{\color{gray}\small величина возможного увеличения потока }
\STATE {\} }
\STATE \textbf{$N:$ array $[1..p]$ of $0..1$} 
\COMMENT{\color{gray}\small отметка узла }
\STATE \textbf{$S:$ array $[1..p]$ of $0..1$}
\COMMENT{\color{gray}\small признак принадлежности узла множеству $S$}
\end{algorithmic}

\end{frame}


\begin{frame}
\frametitle{8.4.4. Алгоритм нахождения максимального потока (2/5)}

\begin{algorithmic}
\REQUIRE
сеть $G(V,E)$ с источником $s$ и стоком $t$,
заданная матрица пропускных способностей
\textbf{$C:$ array $[1..p,1..p]$ of real}.
\ENSURE матрица максимального потока
\textbf{$F:$ array $[1..p,1..p]$ of real}.
\FOR{$u,v\in V$}
  \STATE $F[u,v]\Assign 0$
  \COMMENT{\color{gray}\small  вначале поток нулевой }
\ENDFOR
\STATE $M:$ \COMMENT{\color{gray}\small  итерация увеличения потока }
\FOR{$v\in V$}
  \STATE $S[v]\Assign  0; N[v]\Assign 0; P[v]\Assign(+,\textbf{nil},0)$
  \COMMENT{\color{gray}\small  инициализация }
\ENDFOR
\STATE $S[s]\Assign 1; P[s]\Assign (+, \textbf{nil}, \infty)$
\COMMENT{\color{gray}\small так~как $s\in S$ }
\STATE\textbf{repeat}
  \STATE $a\Assign 0$ 
  \COMMENT{\color{gray}\small  признак расширения $S$ }
\end{algorithmic}

\end{frame}


\begin{frame}
\frametitle{8.4.4. Алгоритм нахождения максимального потока (3/5)}

\begin{algorithmic}
  \FOR{$v\in V$}
    \IF{$S[v]=1 \And N[v]=0$}
       \FOR{$u\in \Gamma(v)$}
          \IF{$S[u]=0 \And F[v,u]<C[v,u]$}
             \STATE $S[u]\Assign 1$
             \STATE $P[u]\Assign (+, v, \min(P[v].\delta,
             C[v,u]-F[v,u])); a\Assign 1$           
          \ENDIF
        \ENDFOR
        \FOR{$u\in \Gamma^{-1}(v)$}
          \IF{$S[u]=0 \And F[u,v]>0$}
             \STATE $S[u]\Assign 1; P[u]\Assign (-, v, \min(P[v].\delta,F[u,v])); a\Assign 1$
          \ENDIF
        \ENDFOR
        \STATE $N[v]\Assign 1$
        \COMMENT{\color{gray}\small   узел $v$ отмечен }
    \ENDIF
  \ENDFOR
  \end{algorithmic}

\end{frame}


\begin{frame}
\frametitle{8.4.4. Алгоритм нахождения максимального потока (4/5)}

\begin{algorithmic}
  \IF{$S[t]$}
    \STATE $x\Assign t$
    \COMMENT{\color{gray}\small  текущий узел аугментальной цепи }
    \STATE $\delta\Assign P[t].\delta$
    \COMMENT{\color{gray}\small  величина изменения потока }
    \WHILE{$x\ne s$}
      \IF{$P[x].s=+$}
        \STATE $F[P[x].n,x]\Assign F[P[x].n,x]+\delta$
        \COMMENT{\color{gray}\small  увеличение потока }
      \ELSE
        \STATE $F[x,P[x].n]\Assign F[x,P[x].n]-\delta$
        \COMMENT{\color{gray}\small  уменьшение потока }
      \ENDIF
      \STATE $x\Assign P[x].n$
      \COMMENT{\color{gray}\small  предыдущий узел аугментальной цепи }
    \ENDWHILE
    \STATE \textbf{goto $M$}
  \ENDIF
\STATE\textbf{until $a=0$}
\end{algorithmic}
\end{frame}


\begin{frame}
\frametitle{8.4.4. Алгоритм нахождения максимального потока (5/5)}

\begin{ground}
В качестве первого приближения берётся нулевой поток, который по
определению является допустимым. В основном цикле, помеченном
меткой $M$, делается попытка увеличить поток. Для этого в цикле
\textbf{repeat} расширяется, пока это возможно, множество $S$
узлов, достижимых из источника $s$ по аугментальным цепям. При
этом, если в множество $S$ попадает сток $t$, то поток вдоль
найденной аугментальной цепи $\lrle{s,t}$ немедленно увеличивается
на величину $\delta$ и начинается новая итерация с целью увеличить
поток. Процесс расширения множества $S$ в цикле \textbf{repeat}
заканчивается, потому что множество узлов конечно, а отмеченные в
массиве $N$ узлы повторно не рассматриваются. Если процесс
расширения множества $S$ заканчивается и при этом сток $t$ не
попадает в множество $S$, то по теореме Форда--Фалкерсона
найденный поток $F$ является максимальным и работа алгоритма
завершается.
\end{ground}


\begin{remark}
Приведённый алгоритм не является самым эффективным
из известных. 
\end{remark}

\end{frame}



\begin{frame}
\label{8.4.5.}\subsubsection{8.4.5. Связь между теоремой Менгера и теоремой Форда--Фалкерсона }\frametitle{8.4.5. Связь между теоремой Менгера и теоремой Форда--Фалкерсона (1/3)}



Теорема Менгера является фундаментальным результатом, который
проявляется в различных формах. 
Теорема Форда--Фалкерсона также может
быть получена из теоремы Менгера. Далее приведена схема
доказательства теоремы Форда--Фалкерсона на основе теоремы
Менгера. Сначала нужно получить вариант теоремы Менгера в
ориентированной рёберной форме: наибольшее число
$\divC{s}{t}$-путей, не пересекающихся по дугам, равно наименьшему
числу дуг в ($s,t$)-разрезе. Это теорема Форда--Фалкерсона в том
случае, когда $\A{e\in E}{c(e)= 1}$. Действительно, пусть $Q$~---
множество дуг из максимального набора непересекающихся
$\divC{s}{t}$-путей. Назначим поток следующим образом:
$f(e)\Assign 1$, если $e\in Q$, и $f(e)\Assign 0$, если $e\not\in
Q$. Это поток, так~как дивергенция в любом узле равна $0$, и поток
через дугу не превосходит пропускной способности. Величина этого
потока равна $k\leq d^+(s)$, где $k$~--- число дуг, выходящих из
$s$, которые начинают пути из $Q$. Этот поток максимальный.

\end{frame}



\begin{frame}
\frametitle{8.4.5. Связь между теоремой Менгера и теоремой Форда--Фалкерсона (2/3)}

Действительно, если положить $f(e)\Assign a>0$ для некоторой дуги
$e\not\in Q$, то
\begin{enumerate}
\item[1)] если дуга $e$ входит в узел, входящий в пути из $Q$, или
выходит из такого узла, то нарушается условие потока (дивергенция
$-a$ или $+a$ соответственно);



\item[2)] если вновь назначенные дуги с ненулевыми потоками
образуют новый $\divC{s}{t}$-путь, то это противоречит выбору $Q$.
\end{enumerate}



Пусть теперь пропускные способности суть натуральные числа. Тогда можно
провести аналогичные рассуждения для мультиграфов, считая, что вместо
одной дуги с пропускной способностью $n$ имеется $n$ дуг с пропускной
способностью $1$.
Если пропускные способности~--- рациональные числа, то их можно привести
к~общему знаменателю и свести задачу к предыдущему случаю.
Для вещественных пропускных способностей заключение теоремы можно получить
путем перехода к пределу в условиях предыдущего случая.

\end{frame}



\begin{frame}
\frametitle{8.4.5. Связь между теоремой Менгера и теоремой Форда--Фалкерсона (3/3)}

\begin{remark}
Приведённое в \DMref{п.~8.4.3}{8.4.3.} доказательство
теоремы Форда--Фалкерсона
\em{конструктивно},
из него не только следует заключение
о величине максимального потока,
но и извлекается способ нахождения
максимального потока.
Набросок доказательства
в этом параграфе \em{неконструктивен}~---
из него нельзя извлечь алгоритм.
\end{remark}


\begin{digress}
Рассмотренная математическая модель 
неотрицательно нагруженной сети 
используется во многих приложениях.
Например, сеть можно использовать для 
календарного планирования работ в некотором проекте
(это называют \em{сетевой график}). 
Дуги соответствуют работам, узлы связывают зависимые 
работы: исходящие работы можно начать, только когда 
выполнены входящие работы, источник обозначает начало проекта, 
сток обозначает окончание проекта.
Нагрузка на дугу соответствует \em{минимальному}
времени, за которое может быть выполнена данная работа.
Как назначить даты узлам так, чтобы
завершить проект как можно раньше?
Оказывается, для этого надо найти \em{максимальный}
путь из источника в сток. Этот путь называется \em{критическим}.
\end{digress}

\end{frame}

%\begin{frame}
%\frametitle{8.4.6. Задача о критическом пути}
%
%
%{\color{blue}
%Тема для несложного, но длинного, а потому 
%дорогого (в баллах) сочинения.
%Теория потоков в сетях получилась слишком куцая --- 
%всего одна задача. Надобно рассмотреть вторую, но так, чтобы 
%было полное единство обозначений и стиля с 
%задачей о максимальном потоке.  
%}
%
%\end{frame}

%\pagebreak[3]
\subsection{8.5. Связность в орграфах}
\label{8.5.}

\begin{frame}
\frametitle{8.5. Связность в орграфах}
Связность является одним из немногих
понятий, которые не распространяются
непосредственно с графов на другие
родственные объекты и требуют отдельного
определения и рассмотрения в каждом случае.
В данном разделе рассматривается связность в орграфах,
поскольку именно это понятие чаще всего
применяется в~программировании.

\medskip
\begin{tabular}{p{9mm}|p{65mm}|p{70mm}}
  № & Название параграфа
  & Ключевые термины и обозначения \\
  \hline
\DMref{8.5.1.}{8.5.1.}\RS & 
Сильная, односторонняя и слабая связность 
  & 
  {Сильно, односторонне и слабо связанные узлы} \\
\DMref{8.5.2.}{8.5.2.} & Компоненты сильной связности 
  & 
  {$G^*$~--- конденсация; КСС~--- компонента сильной связности}
  \\
\DMref{8.5.3.}{8.5.3.} & Выделение компонент сильной связности &
{ }\\
\DMref{8.5.4.}{8.5.4.} & Базисы ориентированных графов &
{Порождающее множество узлов}
\end{tabular}

\end{frame}

\begin{frame}
\label{8.5.1.}\subsubsection{8.5.1. Сильная, односторонняя и слабая связность }\frametitle{8.5.1. Сильная, односторонняя и слабая связность (1/2)}
%\label{s851}


В неориентированном графе две вершины либо связаны (если
существует сое\-диняющая их цепь), либо не связаны. В
ориентированном графе отношение связанности узлов несимметрично, а
потому определение связности различается.


\smallskip
Пусть $G(V,E)$~--- орграф, $v_1$ и $v_2$~--- его узлы. Говорят,
что два узла, $v_1$ и $v_2$, \odmkey{сильно связаны}{Связанность
узлов (connectivity of nodes)!сильная (strong)} в орграфе $G$,
если существует путь (ориентированная цепь) как из $v_1$ в~$v_2$,
так и из $v_2$ в~$v_1$. Говорят, что два узла, $v_1$ и~$v_2$,
\odmkey{односторонне связаны}{Связанность узлов (connectivity of
nodes)!односторонняя (one-way)} в орграфе $G$, если существует
путь либо из $v_1$ в $v_2$, либо из $v_2$ в $v_1$. Говорят, что
два узла, $v_1$ и $v_2$, \odmkey{слабо связаны}{Связанность узлов (connectivity of nodes)!слабая (weak)} в орграфе~$G$, если
они связаны в графе $G'$, полученном из $G$ забыванием ориентации
дуг. 

\smallskip
Если все вершины в орграфе сильно (односторонне, слабо)
связаны, то орграф называется \odmkey{сильно}{Связность
(connectivity)!сильная (strong)} (\odmkey{односторонне}{Связность
(connectivity)!односторонняя (one-way)}, \odmkey{слабо}{Связность
(connectivity)!слабая (weak)}) {\it связным}. Сильная связность
влечет одностороннюю связность, которая влечет слабую связность.
Обратное, разумеется, неверно.
\end{frame}

\begin{frame}
\frametitle{8.5.1. Сильная, односторонняя и слабая связность (2/2)}
%\addvspace{-3pt}
\begin{example}
На рисунке показаны диаграммы
сильно, односторонне и слабо связных орграфов.


\vspace{-12pt}
\begin{figure}[h]
\begin{center}
\includegraphics[height=30mm]{8_5_1.jpg}
\end{center}
%\caption{Сильная ({\it слева}), односторонняя ({\it в центре}) и
%слабая ({\it справа}) связность} \label{f0806}
\end{figure}
\end{example}

\vspace{-12pt}
\begin{lemma}
В произвольном бесконтурном орграфе $G(V,E)$ узлы
можно перенумеровать так, что у любой дуги номер начала меньше
номера конца:  $\A{(v_i,v_j)\in E}{i<j}$.
\end{lemma}
\begin{proof}
По теореме \DMref{п.~7.5.4}{7.5.4.} %\ref{s752}
отношение достижимости в бесконтурном
графе  есть строгое частичное упорядочение на конечном множестве.
Применяя алгоритм топологической сортировки 
(\DMref{п.~1.8.3}{1.8.3.}), получаем
требуемую нумерацию узлов.
\end{proof}

\begin{example}
Левая диаграмма на рисунке содержит контур, 
и требуемая нумерация узлов невозможна.
Для других двух диаграмм приведены
нумерации в условиях леммы.
\end{example}

\end{frame}

\begin{frame}
\label{8.5.2.}\subsubsection{8.5.2. Компоненты сильной связности}
\frametitle{8.5.2. Компоненты сильной связности}

\smallskip
\odmkey{Компонента сильной связности}{Компонента
(component)!сильной связности (of strong connectivity)}
(употребляется аббревиатура КСС) орграфа $G$~--- это его
максимальный сильно связный подграф.

Каждый узел орграфа принадлежит только одной КСС. Если узел не
связан сильно с другими, то считаем, что он сам образует КСС.
\odmkey{Конденсацией}{Конденсация орграфа (condensation of
digraph)} $G^*$ орграфа $G$ (или \odmkey{графом Герца}{Граф
(graph)!Герца (Herz)}, или \odmkey{фактор-графом}{Фактор-граф
(quotient graph)}) называется орграф, который получается
стягиванием в один узел каждой КСС орграфа $G$.


%\addvspace{-5pt}
\begin{example}
На рисунке показаны диаграммы
орграфа и его конденсации.


%\addvspace{-10pt}
\begin{figure}[h]
\begin{center}
\includegraphics[width=100mm]{odm08-07}
\end{center}
%\caption{Орграф ({\it слева}) и его фактор-граф ({\it справа})}
%\label{f0807}
\end{figure}
\end{example}
%\pagebreak[3]


\begin{remark}
Фактор-граф не содержит контуров.
\end{remark}

%\vspace*{-.5\baselineskip}


\end{frame}


\begin{frame}
\label{8.5.3.}\subsubsection{8.5.3. Выделение компонент сильной связности }\frametitle{8.5.3. Выделение компонент сильной связности (1/7)}


Следующий \odmkey{алгоритм выделения компонент
сильной связности}{Алгоритм (algorithm)!выделения компонент сильной связности (finding strongly connectivity component)}, 
основанный на
методе поиска в глубину, находит все компоненты сильной связности
орграфа.


%
%\begin{algorithm}
%\caption{Выделение компонент сильной связности}
\begin{algorithmic}
\REQUIRE
орграф $G(V,E)$, заданный списками смежности $\Gamma(v)$.
\ENSURE
список $C$ компонент сильной связности,
каждый  элемент  которого есть
список узлов, входящих в компоненту сильной связности.
\STATE $C \Assign \emptyset$
\FOR{$v \in V$}
  \STATE $M[v] \Assign  \{v\}$
  \COMMENT{\color{gray}\small $M[v]$ список узлов, входящих в ту же КСС, что и $v$ }
  \STATE $e[v] \Assign 0$
  \COMMENT{\color{gray}\small  все узлы не рассмотрены }
\ENDFOR
\WHILE{$V\ne \emptyset$}
  \STATE \textbf{select $v\!\in\!V;\ v\!\rightarrow\!T;$}
  \COMMENT{\color{gray}\small  взять  $v$, положить  в стек, }
    \STATE $ e[v]\Assign 1; \KCC$
  \COMMENT{\color{gray}\small отметить и вызвать процедуру $\KCC$ }
 \ENDWHILE
\end{algorithmic}


Рекурсивная процедура без параметров $\KCC$:
\end{frame}


\begin{frame}
%\frametitle{8.5.3. Выделение компонент сильной связности (2/3)}
%%
%Основная работа выполняется рекурсивной процедурой без параметров
%$\KCC$, которая использует стек $T$ для хранения просматриваемых
%узлов. Процедура $\KCC$ выделяет все компоненты сильной связности,
%достижимые из узла,
%выбранного в основном алгоритме.


\vspace{-6pt}
%\vspace{6pt}
%{\small
\begin{algorithmic}
%\IF{$T=\emptyset$}
\STATE\textbf{if $T=\emptyset$ then return end if} 
\COMMENT{\color{gray}\small  негде выделять }
%\ENDIF
\STATE $v\Assign$ \textbf{top} $T$ 
\COMMENT{\color{gray}\small  $v$~--- верхний элемент стека}
\IF{$\Gamma[v] =\emptyset$}
  \STATE {$C\Assign C \cup M[v]$;  
   $V\Assign V-v$; 
  $v\leftarrow T$}  
  \COMMENT{\color{gray}\small  это КСС, удалить узел, снять узел $v$ со стека }
  \STATE $\KCC$ 
  \COMMENT{\color{gray}\small  возврат из тупика }
\ELSE
%\pagebreak
  \FOR{$u\in\Gamma[v]$}
    \IF{$e[u]=0$}
      \STATE $u \rightarrow T; e[u]\Assign 1$
       \COMMENT{\color{gray}\small  положить узел $u$ в стек и отметить }
%      \STATE $e[u]\Assign$ 1 \COMMENT{  отметить узел $u$ }
    \ELSE
      \REPEAT
         \STATE $w\leftarrow T; V\Assign V-w; \Gamma[u]\Assign \Gamma[u]\cup\Gamma[w]; M[u]\Assign M[u]\cup M[w]$ 
%         \COMMENT{ $w$~--- склеиваемый узел }
%         \STATE $V\Assign V-w$ \COMMENT{ удалить узел }
%         \STATE $\Gamma[u]\Assign \Gamma[u]\cup\Gamma[w]$
%          \COMMENT{ запомнить смежность }
%         \STATE $M[u]\Assign M[u]\cup M[w]$ \COMMENT{ склеивание узлов }
      \UNTIL {$u=w$}
      \STATE $w\rightarrow T; V\Assign V+w$ 
      \COMMENT{\color{gray}\small  чтобы не убрать узел, с которого начали}
    \ENDIF
    \STATE $\KCC$ \COMMENT{\color{gray}\small  поиск в глубину }
  \ENDFOR
  
  \vspace{-2pt}
\ENDIF
\end{algorithmic}
%}
\end{frame}


\begin{frame}
\frametitle{8.5.3. Выделение компонент сильной связности (3/7)}

\begin{ground}
Достаточно заметить, что любой контур принадлежит ровно одной КСС,
и, более того, если в КСС входит несколько узлов, то они
обязательно принадлежат некоторым контурам. Поэтому, если при
обходе в глубину мы попадаем в уже отмеченный узел $u$, то это
означает, что обнаружен контур, причём предшествующие узлы
найденного контура находятся на стеке, начиная от вершины стека и
до рассматриваемого отмеченного узла $u$, который также
присутствует в стеке. Все узлы найденного контура можно
<<склеить>>. При склеивании узла $w$
 его список смежности и список уже найденных узлов той КСС,
которой принадлежит узел $w$, объединяются с соответствующими
списками узла $u$. После этого можно продолжить поиск в глубину от
узла $u$, который для этой цели следует оставить в стеке.
\end{ground}

\begin{corollary}
Конденсация любого орграфа содержит источники. 
\end{corollary}
\begin{proof}
Конденсация не может содержать контуров по алгоритму,
и по теореме~3 \DMref{п.~7.5.4}{7.5.4.} имеет хотя бы один узел, 
полустепень захода которого равна $0$.  
\end{proof}

\medskip
Приведённый алгоритм не является единственным возможным.
Известен, например, 
\odmkey{алгоритм Косарайю}{Алгоритм (algorithm)!Косарайю (Kosaraju)} 
(иногда это имя транскрибируют как алгоритм \em{Косараджу} ), примечательный своеобразным использованием 
обратного орграфа \DMref{(п.~7.5.1)}{7.5.1.}.
\end{frame}


\begin{frame}
\frametitle{8.5.3. Выделение компонент сильной связности (4/7)}
Cначала с помощью 
рекурсивной процедуры $F$ орграф обходится
в глубину, и при этом узлы переупорядочиваются в 
порядке, обратном порядку выхода из узлов при обходе.
Затем с помощью рекурсивной процедуры $B$ обратный 
орграф обходится
в глубину, и при этом узлы рассматриваются 
в найденном порядке.

\begin{algorithmic}
\REQUIRE
орграф $G(V,E)$, заданный списками смежности 
\textbf{$\Gamma :$ array $[1..p]$ of set of $1..p$}.
\ENSURE
 массив разметки КСС
\textbf{$M :$ array $[1..p]$ of  $1..p$},
причём $M[u]=M[v]$ тогда и только тогда,
когда узлы $u$ и $v$ принадлежат одной КСС.
%\STATE{\textbf{for $v$ from $1$ to $p$ do $M[v] \Assign 0$ end for}}
%\COMMENT{\color{gray}\small вначале разметка неизвестна}
\STATE{\textbf{$X:$ array $[1..p]$ of $0..1$}}
\COMMENT{\color{gray}\small отметка посещения узла} 
\STATE{\textbf{for $v$ from $1$ to $p$ do $X[v] \Assign 0$ end for}}
\COMMENT{\color{gray}\small вначале все узлы не посещены}
\STATE{\textbf{$\Gamma^{-1}:$ array $[1..p]$ of set of  $0..1$}}
\COMMENT{\color{gray}\small списки смежности обратного орграфа} 
\STATE{\textbf{for $v$ from $1$ to $p$ do $\Gamma^{-1}[v] \Assign \emptyset$ end for}}
\COMMENT{\color{gray}\small вначале списки пусты}
\STATE{\textbf{$T:$ array $[1..p]$ of  $0..1$}}
\COMMENT{\color{gray}\small массив узлов в порядке, обратном порядку выхода из узла} 
\STATE{\textbf{for $v$ from $1$ to $p$ do $T[v] \Assign 0$ end for}}
\COMMENT{\color{gray}\small вначале последовательность выходов неизвестна}
\STATE{\textbf{$t\Assign p$}}
\COMMENT{\color{gray}\small величина, обратная времени выхода}
\STATE{\textbf{for $v$ from $1$ to $p$ do $F(v)$ end for}}
\COMMENT{\color{gray}\small обход исходного орграфа в глубину}
\STATE{\textbf{for $v$ from $1$ to $p$ do $B(T[v],T[v])$ end for}}
\COMMENT{\color{gray}\small обход обратного орграфа в глубину}
\end{algorithmic}

\end{frame}

\begin{frame}
\frametitle{8.5.3. Выделение компонент сильной связности (5/7)}

\begin{columns}[t]
\begin{column}{8cm}
Рекурсивная процедура $F(u)$ обхода в глубину орграфа
определяет порядок выхода и заодно строит списки смежности обратного орграфа.


\begin{algorithmic}
\IF{$X[u] = 0$}
\STATE{$X[u] \Assign 1$}
\COMMENT{\color{gray}\small отмечаем посещение узла}
\FOR{$w\in \Gamma[u]$}
\STATE{$F(w)$}
\COMMENT{\color{gray}\small рекурсивный вызов в глубину}
\STATE{$\Gamma^{-1}[w] \Assign \Gamma^{-1}[w] + u$}
\COMMENT{\color{gray}\small  записываем $u$ как смежный с $w$ в обратном орграфе}
\ENDFOR
\STATE{$T[t]\Assign u$}
\COMMENT{\color{gray}\small запись времени выхода в обратном порядке}		
\STATE{$t\Assign t-1$}
\COMMENT{\color{gray}\small уменьшение счётчика времени выхода}				
\ENDIF
\end{algorithmic}
\end{column}
\begin{column}{7cm}
Рекурсивная процедура $B(u,v)$ обхода в глубину  обратного орграфа относит узлы $u$ и $v$ к одной КСС.


\begin{algorithmic}
\IF{$X[u] = 1$}
\STATE{$X[u] \Assign 0$}
\COMMENT{\color{gray}\small отмечаем посещение узла}
\STATE{$M[u]\Assign v$}
\COMMENT{\color{gray}\small $u$ принадлежит 
той же КСС, что и $v$}
\FOR{$w\in \Gamma^{-1}[u]$}
\STATE{$B(w,v)$}
\COMMENT{\color{gray}\small рекурсивный вызов в глубину}
\ENDFOR
\ENDIF
\end{algorithmic}
\end{column}
\end{columns}

\end{frame}

\begin{frame}
\frametitle{8.5.3. Выделение компонент сильной связности (6/7)}

\begin{ground}
Совершенно ясно, что если начать 
поиск в глубину с любого узла $u$ в любой КСС, 
то алгоритм посетит \em{все}
узлы $v$ в данной КСС, для которых существует 
как путь из  $u$ в $v$, так и путь из
$v$ в $u$. Однако, если граф не сильно 
связен, то алгоритм, возможно,
посетит
ещё некоторые <<лишние>> узлы,
то есть такие, для которых существует 
путь из  $u$ в $v$, но \em{не} существует пути из
$v$ в $u$. 
Заметим, что если существует 
путь из  $u$ в $v$, но не существует пути из
$v$ в $u$, то при любом обходе в глубину
алгоритм выйдет из узла $u$ позже, чем из узла $v$.

  
\smallskip
Если узел $v$ пройден в результате выполнения поиска в глубину 
процедурой $B(v, u)$ с началом в $u$, то это значит, 
что: 
1) узел $u$ имеет более позднее время выхода, 
чем узел $v$, то есть
узел $u$ расположен левее, 
чем узел $v$, в массиве $T$;
2) существует путь из $u$ в $v$ в обратном графе,
то есть существует путь из $v$ в $u$ в исходном графе.

При этом в исходном орграфе существует путь из $u$ в $v$,
который вместе с обратным путем из $v$ в $u$ даёт сильную
связанность узлов и заключение, что
 $u$ и $v$ принадлежат одной КСС. 
 Вариант, при котором существует путь
  $v$ в $u$ в исходном графе, но не существует
  пути из $u$ в $v$, невозможен, 
  поскольку противоречит тому, что узел 
  $v$ расположен правее в массиве $T$.
%  
%  
% 
% 
%%2) a и b никак не были связаны в исходном орграфе (в ином случае, то есть при существовании пути только из b в a, b был бы пройден первым шагом позже и третий шаг алгоритма не мог бы начать поиск с a). Но это противоречит тому, что в инвертированном орграфе мы нашли путь из a в b.
% 
%Сложность: использует рекурсивный DFS, в который встроены построение инвертированного орграфа и запись порядка выхода из узлов в порядке, обратном порядку прохождения по ним. Рекурсивный DFS проходит по |V| узлов и вызывает суммарно |E| дуг. Следовательно, итоговая сложность |V| + |E|.
%
%Сравнение алгоритма из 8.5.3. с алгоритмом Косарайю: используют одну и ту же идею поиска контуров в орграфе с использованием DFS, но делают это разными методами - алгоритм из слайдов ищет контуры нахождением уже отмеченных узлов, в то время как Косарайю доказывает существование пути как из a в b, так и из b в a с помощью порядка выхода. Оба имеют сложность |V| + |E|. Алгоритм из слайдов имеет нерекурсивное решение для списков смежности, для алгоритма Косарайю такого решения автору найти не удалось (однако это не значит, что его не существует).

\end{ground}

\end{frame}

\begin{frame}
\frametitle{8.5.3. Выделение компонент сильной связности (7/7)}


\begin{remark}
Если два узла $u$ и $v$ сильно связаны в орграфе, то они сильно связаны и в обратном орграфе, просто потому, 
что паре путей, связывающих узлы в орграфе, 
соответствует пара обратных путей, связывающих эти
узлы в обратном орграфе. Таким образом, КСС в орграфе 
и обратном орграфе совпадают.
\end{remark}

\begin{digress}
Алгоритм изложен в соответствии с традицией,
согласно которой используются 
списки смежности для представления орграфов
и рекурсия для реализации поиска в глубину.
В данном случае следование традициям
упрощает изложение и спсобствует пониманию материала,
но никоим образом не повышает
эффективность реализации.
Списки смежности, по-видимому, ~--- это
наихудшее представление для построения
обратного орграфа.
Действительно, если 
используется другое представление,
то для построения обратного орграфа 
достаточно транспонировать матрицу смежности,
или заменить все $1$ на $-1$ и наоборот в матрице
инциденций, или поменять местами элементы в
каждой паре списка дуг. 
Рекурсия здесь также использована
не по существу, а для передачи 
текущих значений переменных 
через стек вызовов в форме параметров.
Ведь никаких новых значений 
переменных $u$, $v$ не вычисляется,
это изначально заданные номера узлов.
Из этого наблюдения следует вывод,
что целесообразно использовать различные
подходы при \em{публикации} 
и при \em{реализации} алгоритмов. 
\end{digress}

\end{frame}

\begin{frame}
\label{8.5.4.}\subsubsection{8.5.4. Базисы ориентированных графов }\frametitle{8.5.4. Базисы ориентированных графов (1/2)}

Множество узлов $B$ орграфа $G(V,E)$ называется 
\odmkey{порождающим}{Множество (set)!узлов (of nodes)!порождающее (generative)}, если любой узел орграфа достижим
из некоторого  узла,
принадлежащего этому множеству: 
$$\A{v\in V}{\E{u\in B}{\exists \Path{u}{v}\in G}}.$$
Порождаю\-щее множество узлов называется
\odmkey{базисом}{Базис (base)!орграфа (of digraph)} орграфа, если оно минимально по включению.

\begin{remark}
Узлы одного базиса недостижимы друг из друга.
\end{remark}

Ясно, что порождающие множества, а значит и базисы, существуют у любого орграфа. 

\begin{lemma}
Если орграф не имеет контуров, то
имеет единственный базис, состоящий   
из всех источников.
\end{lemma}
\begin{proof}
Любой источник входит в базис.
Поскольку в орграфе нет контуров,
любой узел достижим из некоторого
источника. Действительно, чтобы найти этот источник,
достаточно двигаться из узла против 
направления входящих дуг.   
\end{proof}

\begin{columns}[t]
\begin{column}[T]{8cm}
\begin{example}
На рисунке слева каждый узел образует базис, 
в центре базис состоит из одного узла, а справа
 базис состоит из двух узлов. 
\end{example}
\end{column}
\begin{column}[T]{7cm}
\vspace{-18pt}
\begin{figure}[h]
\begin{center}
\includegraphics[width=60mm]{odm08-06}
\end{center}
\end{figure}
\end{column}
\end{columns}
\end{frame}

\begin{frame}
\frametitle{8.5.4. Базисы ориентированных графов (2/2)}
\begin{corollary}
Конденсация любого орграфа имеет единственный базис,
состоящий из источников.
\end{corollary}

\em{Базисными} компонентами сильной связности орграфа
называются такие компоненты сильной связности, которым соответствуют источники в конденсации этого орграфа.

\begin{theorem}
Подмножество узлов орграфа образует базис тогда и только тогда,
когда оно является трансверсалью базисных компонент сильной связности. 
\end{theorem}   

\begin{proof}
[ $\Impl$ ] Пусть $B$~--- базис. 
Тогда узлы $B$ принадлежат различным базисным КСС,
поскольку узлы базиса недостижимы друг из друга.
В базис должен входить узел из каждой 
базисной КСС, в противном случае
узлы <<обделённой>> КСС оказались бы недостижимы из базиса.
Таким образом, базис является трансверсалью базисных КСС.      
\proofsection{$\Impll$}
Любая трансверсаль базисных 
КСС образует базис, поскольку все узлы
базисных КСС достижимы из своих представителей в базисе
ввиду сильной связности КСС,
а любой узел в другой КСС достижим из 
некоторой базисной КСС.
\end{proof}

\medskip
\begin{corollary}
Все базисы любого орграфа равномощны.
\end{corollary}

\end{frame}


%\vspace*{-.5\baselineskip}
%
%\begin{frame}
%\frametitle{Темы для сочинений}
%
%{\color{red} Тема для сочинения:
%в разделе про потоки в сетях остро недостаёт
%сравнительного разбора второй задачи,
%например, задачи о критическом пути.
%
%Разбор должен быть математический,
%но программистский.
%
%А.Н. Терехов?
%
%}
%
%\medskip
%{\color{red} Тема для сочинения:
%пересказать алгоритм Ахо, Гэри, Ульман 1972 про построение сокращения в общем случае: 
%граф Герца, каждая нетривиальная КСС --– один простой цикл. 
%
%http://dept-info.labri.fr/~thibault/tmp/0201008.pdf
%
%}
%
%\medskip
%{\color{red} Тема для сочинения:
%построение графа Герца, или фактор-графа какой-нибудь
%модели программы~---
%это основной метод снижения объёма вычислений
%при компиляции, оптимизации и всяком другом 
%потоковом анализе программ.
%Было бы желательно сделать ход в эту сторону.
%См. Евстигнеев, Касьянов.
%}
%\end{frame}
%%

%

%%%%%%%%%%%%%%%%%%%%%%%%%%%%%%%%%%%%%%%%%%%%%%%%%%%%%%%%%%%%
% Вариант Сергея Лебедева
% редактировал ФАН 13.04.2015
%%%%%%%%%%%%%%%%%%%%%%%%%%%%%%%%%%%%%%%%%%%%%%%%%%%%%%%%%%%%

\begin{frame}
\frametitle{Онтология теории графов}
\begin{center}
\includegraphics[width=12cm]{ODM_8_5.png}
\end{center}

\end{frame}


\subsection{8.6. Кратчайшие пути}
\label{8.6.}

\begin{frame}
\frametitle{8.6. Кратчайшие пути (1/2)}

Задача нахождения кратчайшего пути
в графе имеет столько практических
применений и интерпретаций, что
читатель, без сомнения, может сам
легко привести множество примеров.
Здесь рассматриваются несколько классических
алгоритмов, которые обязан знать
каждый программист.
\end{frame}

\begin{frame}
\frametitle{8.6. Кратчайшие пути (2/2)}

\begin{tabular}{p{9mm}|p{65mm}|p{70mm}}
  № & Название параграфа
  & Ключевые термины и обозначения \\
  \hline
\DMref{8.6.1.}{8.6.1.}\RS & 
Взвешенные орграфы 
  & 
  {$w(u,v)$~--- вес дуги; $W(u,w)$~--- длина пути;
  $\Omega(s,t)$~--- длина кратчайшего пути}; процедура релаксации \\
\DMref{8.6.2.}{8.6.2.} & Кратчайшие пути в бесконтурном графе 
  & 
  {}
  \\
\DMref{8.6.3.}{8.6.3.} & Алгоритм  Беллмана--Форда &
{ }\\
\DMref{8.6.4.}{8.6.4.}\RS & Алгоритм Дейкстры &
{}
\\
\DMref{8.6.5.}{8.6.5.}\RS & Область применения алгоритма Дейкстры  &
{$R(\Path{u}{v})$~--- рекорд пути}\\   
\DMref{8.6.6.}{8.6.6.}\RS & Анализ алгоритма Дейкстры  &
{}\\   
\DMref{8.6.7.}{8.6.7.} & Алгоритм Флойда--Уоршалла  &
{}\\
\DMref{8.6.8.}{8.6.8.}\BS & Алгоритмы с предобработкой  &
{$h(v)$~--- потенциал}\\
\end{tabular}
\end{frame}


\begin{frame}
\label{8.6.1.}\subsubsection{8.6.1. Взвешенные орграфы}\frametitle{8.6.1. Взвешенные орграфы (1/8)}

Орграф $G(V,E)$ называется 
\odmkey{взвешенным}{Граф (graph)!взвешенный (weighted)},
если задана числовая функция  $\Map{w}{V\times V}{\Bbb{R}}$,
называемая \odmkey{весом}{Вес (weight)!дуги (of arc)} или \odmkey{длиной}{Длина (length)!дуги (of arc)} дуги.

\begin{remark}
Если дуга $e=(u,v)\in E$ ведёт из узла $u$ в узел $v$,
то записи $w(e)$ и $w(u,v)$ обе считаются допустимыми и равноправными.
\end{remark}

Если $M = \Path{s}{t}$~--- путь из узла $s$ в узел $t$, то
\odmkey{длиной пути}{Длина (length)!пути (of path)} называется величина

\vspace{-6pt}
$$W(M)=W(\Path{s}{t}) = W(s,t)=\sum\limits_{e\in M}w(e),$$

а \odmkey{длиной кратчайшего пути}{Длина (length)!кратчайшего пути (of shortest path)} называется   
величина 

\vspace{-6pt}
$$ \Omega (s,t)=\min\limits_{\Path{s}{t}\subset E^*} W(s,t).$$

\vspace{-2pt}
\begin{remark}
Поскольку неориентированные графы являются частным случаем орграфов,
приведённые определения, утверждения и алгоритмы распространяются
также и на графы.
Поэтому изложение ведётся исключительно в терминах орграфов.
\end{remark}


\end{frame}


\begin{frame}
\frametitle{8.6.1. Взвешенные орграфы (2/8)}
\begin{theorem}
Для взвешенных орграфов выполняется неравенство 
$\Omega(s, v) \le \Omega(s,u)+\Omega(u,v)$ (\odmkey{неравенство треугольника}{Неравенство (inequality)!треугольника (of triangle)}).
\end{theorem}
\begin{proof}
Пусть существует кратчайший путь из $s$ в $v$.
Тогда его длина не превышает длины любого другого пути
$\Path{s}{v}$ и, 
в частности, пути, состоящего из кратчайшего пути из $s$ в $u$ и 
кратчайшего пути из $u$ в $v$.
\end{proof}

\smallskip
Представим взвешенный орграф  \odmkey{матрицей длин дуг}{Матрица (matrix)!длин дуг (of arc's length)}
\textbf{$W:$ array $[1..p,1..p]$ of $\Bbb{R}$}, 
\\где
\textbf{$W[i,j]\Assign$
if $i=j$ then $0$ elseif $(i,j)\in E$ then $w(i,j)$ else $+\infty$ end if}.

\smallskip
Для хранения информации о последовательности узлов в кратчайших путях  используется \odmkey{матрица предшествования}{Матрица (matrix)!предшествования (precedence)}
$\Pi$: \textbf{array} $[1..p,1..p]$ \textbf{of} $0..p$, где
$\Pi[i,j] \Assign 0$, если из узла $i$ в узел $j$ нет пути,
и $\Pi[i,j] \Assign k$, если $k\!$~--- узел, достигаемый
на кратчайшем пути из $i$ непосредственно перед узлом $j$.

\smallskip
Для хранения информации о длинах кратчайших путей 
 используется \odmkey{матрица длин путей}{Матрица (matrix)!длин путей (of path's length)}
$T$: \textbf{array} $[1..p,1..p]$ \textbf{of} 
$\Bbb{R}$, где
$T[i,j] =+\infty$, если из узла $i$ в узел $j$ нет пути, и $T[i,j] =
\Omega(i,j)$ в противном случае.


\end{frame}


\begin{frame}
\frametitle{8.6.1. Взвешенные орграфы (3/8)}

\begin{digress}
Матрица $\Pi$ размера $O(p^2)$ хранит информацию обо \em{всех}
(кратчайших) путях в графе. Заметим, что всего в графе $O(p^2)$
путей, состоящих из $O(p)$ узлов. Таким образом, непосредственное
представление всех путей потребовало бы памяти объема $O(p^3)$.
Экономия памяти при использовании матрицы предшествования 
достигается за счёт \odmkey{интерпретации
представления}
{Интерпретация (interpretation)!представления (of implementation)}, 
то~есть динамического вычисления некоторой части
информации вместо её хранения в памяти. 
В данном случае любой
конкретный путь $\lrle{u,v}$ легко извлекается из матрицы с
помощью следующего алгоритма.


\begin{algorithmic}
\STATE $w\Assign v$; \textbf{yield $w$} 
\COMMENT{\color{gray}\small последний узел }
\WHILE{$w\ne u$}
\STATE $w\Assign \Pi[u,w]$;
\STATE \textbf{yield $w$} 
\COMMENT{\color{gray}\small  предыдущий узел }
\ENDWHILE
\end{algorithmic}


Здесь узлы пути выдаются в обратном 
порядке. Полученную последовательность можно обратить, 
если это потребуется, и получить прямой порядок узлов пути.

\end{digress}

\end{frame}


\begin{frame}
\frametitle{8.6.1. Взвешенные орграфы (4/8)}
Если требуется хранить информацию только о путях из одного узла $s$ или, 
в частности, о единственном пути между двумя узлами $s$  и $t$, 
то матрицу предшествования можно сократить до вектора
 \textbf{$\Pi:$ array $[1..p]$ of $0..p$},
а матрицу длин путей можно сократить до вектора
$T$: \textbf{array} $[1..p]$ \textbf{of} 
$\Bbb{R}$.
Это возможно, поскольку первый индекс $s$ 
является фиксированным.

\medskip
Во многих алгоритмах данного раздела используются две процедуры~--- инициализации и релаксации. Здесь приводятся реализации
для случая векторов предшествования и длин путей,
реализации для случая матриц совершенно 
аналогичны.

\medskip
Процедура инициализации $\Init(s)$ строит начальное состояние
вектора длин путей и вектора предшествования.  

%\vspace{-2pt}
\begin{algorithmic}
\REQUIRE
узел $s$ %\COMMENT{\color{gray}\small  $\Init(s)$ }
\ENSURE заполненные векторы
\textbf{$T:$ array $[1..p]$ of $\Bbb{R}$} длин путей и
\textbf{$\Pi:$ array $[1..p]$ of $0..p$} самих путей.
\STATE\textbf{for $v$ from $1$ to $p$ do 
$T[v]\Assign +\infty;\Pi[v] \Assign 0$ end for}
\STATE\textbf{$T[s] \Assign 0$}
\end{algorithmic}
\end{frame}


\begin{frame}
\frametitle{8.6.1. Взвешенные орграфы (5/8)}

Процедура 
\odmkey{ослабления}{Ослабление (relaxation)} или
\odmkey{релаксации}{Релаксация (relaxation)}~--- 
$\Relax(s, v, u)$
проверяет, 
возможно ли улучшить известный путь из узла $s$ 
в узел $v$, 
проведя новый путь через узел $u$, и обновляет пути, если это возможно.


\begin{algorithmic}
\REQUIRE
узлы $s$, $v$, $u$, %\COMMENT{\color{gray}\small  $\Relax(s, v, u)$ }
матрица \textbf{$W:$ array $[1..p,1..p]$ of $\Bbb{R}$} длин дуг,
вектор \textbf{$T:$ array $[1..p]$ of $\Bbb{R}$} длин путей,
вектор \textbf{$\Pi:$ array $[1..p]$ of $0..p$} самих путей.
\ENSURE обновлённые векторы $T, \Pi$.
\IF {$T[v] > T[u] + W[u,v]$}
\STATE{ $T[v] \Assign T[u] + W[u,v]; 
\Pi[v] \Assign u$}
\COMMENT{\color{gray}\small  новый путь короче }
\ENDIF
\end{algorithmic}


Здесь параметр $s$ не используется в теле процедуры
релаксации, и оставлен для единообразия с общим случаем использования матриц. 

\begin{remark}
Для случая использования матриц,
тело процедуры релаксации достаточно
модифицировать следующим образом.

\vspace{-1pt}
\begin{algorithmic}
\IF {$T[s,v] > T[s,u] + W[u,v]$}
\STATE{ $T[s,v] \Assign T[s,u] + W[u,v]; \Pi[s,v] \Assign u$}
\COMMENT{\color{gray}\small  новый путь короче }
\ENDIF
\end{algorithmic}

\end{remark}
\end{frame}


\begin{frame}
\frametitle{8.6.1. Взвешенные орграфы (6/8)}

Процедура релаксации обладает рядом очевидных, но важных свойств.

\begin{enumerate}

\item \textbf{Верхняя граница (безопасность)}\\
Для всех пар узлов $T[s, v] \ge \Omega[s,v]$, причем $T[s, v]$ не изменяется после того, как станет равным  $\Omega[s, v]$.\\

\item \textbf{Отсутствие пути}\\
Если из узла $s$ в $t$ нет пути, то $T[s,t] = \Omega(s,t) = +\infty$.

\item \textbf{Сходимость}\\
%ЗАМЕЧАНИЕ ПО ЛЕКЦИИ:
%Это свойство ни в коем случае не эквивалентно первому.
%Как показала прошлая лекция, и студенты, и лектор могут упустить из виду чередование букв.
Если $\Path{s}{u}, (u,v)$ --- кратчайший путь
$\Path{s}{v}$, и до ослабления 
$\Relax(s,v,u)$ выполняется равенство $T[s,u] = \Omega(s,u)$, то после ослабления $T[s,v] = \Omega(s,v)$.



\item \textbf{Ослабление пути}\\
Если $M = \langle s, v_1, ... , v_k, t \rangle$ --- кратчайший путь 
$\Path{s}{t}$, и ослабление производится в порядке
$\Relax(s,s,v_1)$, $\Relax(s,v_1,v_2),\dots, \Relax(s,v_k,t)$,
то $T[s, t] = \Omega(s, t)$, и это свойство выполняется вне зависимости от других операций ослабления.

\end{enumerate}

\end{frame}


\begin{frame}
\frametitle{8.6.1. Взвешенные орграфы (7/8)}
\begin{remark}
Если в $G$ есть контур с отрицательным весом, то на
этом контуре можно <<накручивать>> сколь угодно короткий путь
(вплоть до~$-\infty$). В этой
ситуации решение задачи о поиске кратчайшего пути, вообще говоря, не
определено. Однако, в зависимости от конкретной прикладной задачи,
данная ситуация может быть интерпретирована по-разному, 
как полное отсутствие решения, или как решение, содержащее путь  бесконечного отрицательного веса, или ещё как-то. 

\smallskip
В любом случае, способность предсказуемо реагировать на контуры
отрицательного веса является важной характеристикой алгоритмов поиска
кратчайших путей.
\end{remark}

\smallskip
Совокупность всех путей от узла $s$ к узлу $t$ образует 
\em{сеть} в смысле \DMref{п.~7.3.3}{7.3.3.}, точнее говоря, \em{нагруженную 
$(s,t)$-сеть} в 
смысле \DMref{п.~8.4.1}{8.4.1.}.  В сети может встретиться контур $C$.
Если $W(C)>0$, то $C$ не входит
ни в какой кратчайший путь.
Если $W(C)=0$, то $C$ не влияет
на кратчайшие пути.
Если $W(C)<0$, то $C$ препятствует построению 
кратчайших путей.  



\end{frame}


\begin{frame}
\frametitle{8.6.1. Взвешенные орграфы (8/8)}
В любом случае контуры  
можно исключить из сети путей и не ограничивая общности
рассматривать только пути без самопересечений.
Обозначим $P(s,t)$~--- множество всех путей без самопересечений
из узла $s$ в узел $t$.

\smallskip
Источник $s$ и сток $t$ являются общими узлами для всех 
$(s,t)$-путей.
Если исключить из рассмотрения узлы $s$ и $t$, то пути 
в $(s,t)$-сети называются:
\begin{itemize}
\item \em{непересекающимися}, если они не имеют общих узлов;
\item \em{пересекающимися}, если они имеют общие узлы, но не имеют общих дуг; 
\item \em{слипающимися},
если они имеют общие дуги 
(а значит, и узлы). 
\end{itemize}

\smallskip
Совокупность кратчайших путей от узла $s$ к другим узлам образует
\em{корневое дерево} в смысле \DMref{п.~9.2.1}{9.2.1.}. 
Действительно, никакой определённый
 кратчайший путь
не содержит контуров.
Всякий узел
достижим из корня по построению. 
Если в узел ведет несколько кратчайших пересекающихся  или слипающихся путей одинаковой длины, 
всегда можно выбрать один из них. 
\end{frame}



\begin{frame}
\label{8.6.2.}\subsubsection{8.6.2. Кратчайшие пути в бесконтурном орграфе }\frametitle{8.6.2. Кратчайшие пути в бесконтурном орграфе (1/2)}
Существует эффективный \odmkey{алгоритм определения расстояний от источника}{Алгоритм (algorithm)!определения расстояний от источника (finding distances from source)}, который позволяет найти в орграфе
\odmkey{дерево кратчайших путей}
{Дерево (tree)!кратчайших путей (of shortest paths)}, растущее из заданного узла,
даже если длины дуг могут быть отрицательными,
но известно, что орграф не содержит контуров.

\begin{remark}
Так как орграф не имеет контуров, 
можно применить лемму \DMref{п.~8.5.1}{8.5.1.}
и перенумеровать узлы так, что 
номер начала каждой дуги меньше номера конца. 
\end{remark}

 
\begin{algorithmic}
\REQUIRE
взвешенный орграф $G(V,E)$,
списки предшествующих узлов $\Gamma^{-1}$, \\матрица весов
\textbf{$W:$ array $[1..p,1..p]$ of $\Bbb{R}$}, 
источник $s = 1$.
\ENSURE
вектор \textbf{$T:$ array $[1..p]$ of $\Bbb{R}$}
длин кратчайших путей от источника, \\вектор \textbf{$\Pi:$ array $[1..p]$ of $0..p$} самих путей.
\STATE $\Init(s)$
\FOR{\textbf{$v$ from $2$ to $p$}}
  \FOR{$u\in \Gamma^{-1}(v)$}
       \STATE $\Relax(1, v, u)$
  \ENDFOR
\ENDFOR
\end{algorithmic}

\end{frame}


\begin{frame}
\frametitle{8.6.2. Кратчайшие пути в бесконтурном орграфе (2/2)} 

\begin{ground}
Докажем индукцией по $v$, что основной цикл имеет инвариант
$$\A{u\in 1..v}{T[u]=\Omega(1,u)}.$$ База: если $v=1$, то
$T[1]=0=\Omega(1,1)$. Пусть  $\A{u<v}{T[u]=\Omega(1,v)}$. 
В кратчайшем
пути $\Path{1}{v}$ все промежуточные узлы имеют номера
больше $1$ и меньше $v$ по построению орграфа, в частности, узел
$u$, предшествующий $v$ на кратчайшем пути, будет выбран во
внутреннем цикле, для него по индукционному предположению
$T[u]=\Omega(1,u)$ и, значит, $T[v]=\Omega(1,v)$.
Оценка времени работы алгоритма складывается 
из времени топологической сортировки, 
которая выполняется за время $O(p+q)$, 
инициализации~--- $O(p)$ и итераций двойного цикла, 
которых будет ровно $q$, следовательно время будет $O(q)$. 
Так получаем $O(p\!+\!q)+O(p)+O(q)=O(p\!+\!q)$.
\end{ground}

\medskip
\begin{remark}
В алгоритме используется множество
<<предшествования>>
$$
\Gamma^{-1}(v)\Def \Set{u}{v\in\Gamma(u)},
$$
которое совпадает со списками смежности $\Gamma(v)$ для графов
и не пересекается с $\Gamma(v)$ для направленных орграфов (см. \DMref{п.~7.1.3}{7.1.3.}).
\end{remark}
\end{frame}


\begin{frame}
\label{8.6.3.}\subsubsection{8.6.3. Алгоритм Беллмана--Форда }\frametitle{8.6.3. Алгоритм Беллмана--Форда (1/3)} 

\odmkey{Алгоритм Беллмана--Форда}{Алгоритм (algorithm)!Беллмана--Форда (Bellman--Ford)} также решает задачу о построении дерева кратчайших путей, 
но на взвешенных орграфах общего вида.

\begin{algorithmic}
\REQUIRE
взвешенный орграф $G(V,E)$, источник $s$, \\матрица весов
\textbf{$W:$ array $[1..p,1..p]$ of real}.
\ENSURE
вектор \textbf{$T:$ array $[1..p]$ of real}
длин кратчайших путей от источника, \\вектор \textbf{$\Pi:$ array $[1..p]$ of $0..p$} самих путей.

\STATE $\Init(s)$

\FOR{\textbf{$i$ from $1$ to $p-1$}}
  \FOR{$(u, v) \in E$}
    \STATE $\Relax(s, v, u)$
  \ENDFOR
\ENDFOR
%\STATE{} \COMMENT{\color{gray}\small Проверка на отрицательные контуры}
\FOR{$(u, v) \in E$}
  \STATE \textbf{if $T[v] > T[u] + W[u,v]$ then
    stop end if }
    \COMMENT {\color{gray}\small найден контур отрицательного веса}
\ENDFOR
\end{algorithmic} 


\end{frame}


\begin{frame}
\frametitle{8.6.3. Алгоритм Беллмана--Форда (2/3)}

\begin{ground}
Очевидно, что кратчайший путь содержит не более чем $p-1$ дугу при условии отсутствия контуров отрицательного веса, доступных из источника. Соответственно, если провести релаксацию всех дуг взвешенного орграфа $p-1$ раз, то по свойствам релаксации (в особенности по свойствам ослабления пути и безопасности) результатом окажется дерево кратчайших путей.
Если после этого провести дополнительную релаксацию всех дуг, и при этом дерево кратчайших путей изменится, то этот факт по свойству безопасности 
войдёт в противоречие с отсутствием контуров отрицательного веса, и можно за счёт этого выдавать предупреждение об их наличии.
\end{ground}

\begin{example}
На рисунке показано состояние после инициализации и начало последовательности релаксаций некоторого орграфа.

\vspace{-6pt}
\begin{center}
\includegraphics[width=150mm]{8_6_3.jpg}
\end{center}
\end{example}

\end{frame}


\begin{frame}
\frametitle{8.6.3. Алгоритм Беллмана--Форда (3/3)}

\begin{remark}
В предыдущем примере дуги релаксировались в лексикографическом порядке:
$(s,u), (s,w), (u,v), (u,w), (w,u)$,
причём этот порядок \em{случайно} выбран
настолько удачным, что все релаксации оказались результативными,   
и дерево кратчайших путей 
построено уже на первой итерации, а последующие итерации ничего не меняют.

На основании такого наблюдения
иногда может иметь смысл добавить в алгоритм проверку того, что изменений дерева на прошедшей итерации внешнего цикла не произошло,
и заканчивать алгоритм <<досрочно>>.   
\end{remark}

\medskip
Оценка времени работы алгоритма тривиальна. 

Время инициализации~--- $O(p)$. 
Процедура ослабления вызывается $(p-1)q$ раз, 
всего имеем 
$O(p) + (p-1)q \cdot O(1) = O(p) + O((p-1)q) = O(p) + O(pq) = O(pq)$. Очевидно, это медленнее, чем оценка $O(p+q)$ алгоритма поиска на бесконтурном орграфе. Однако следует напомнить, что 
алгоритм Беллмана--Форда 
корректно работает на любых орграфах.
\end{frame}


\begin{frame}
\label{8.6.4.}\subsubsection{8.6.4. Алгоритм Дейкстры }\frametitle{8.6.4. Алгоритм Дейкстры (1/3)}

\odmkey{Алгоритм Дейкстры}
{Алгоритм (algorithm)!Дейкстры (Dijkstra)}, безусловно, самый знаменитый алгоритм поиска кратчайших путей. 
%Однако он имеет как преимущества, так и недостатки относительно
%алгоритма Беллмана-Форда, что  разобрано в следующем параграфе.

\begin{algorithmic}
\REQUIRE
взвешенный орграф $G(V,E)$, источник $s$,
 \\матрица весов
\textbf{$W:$ array $[1..p,1..p]$ of real}.
\ENSURE
вектор \textbf{$T:$ array $[1..p]$ of real}
длин кратчайших путей от источника, \\вектор \textbf{$\Pi:$ array $[1..p]$ of $0..p$} самих путей.

\STATE $\Init(s)$
\COMMENT {\color{gray}\small инициализация векторов $T$ и $\Pi$}
\STATE $Q \Assign V$
\COMMENT {\color{gray}\small вначале контейнер $Q$ содержит все узлы}
\WHILE{\textbf{$Q \neq \emptyset$}}
  \STATE $u \Assign \ExtractMin(Q)$
  \COMMENT {\color{gray}\small извлечение узла с минимальным $T[u]$}
    \STATE \textbf{if $T[u]=+\infty$ then stop end if}
  \COMMENT {\color{gray}\small можно прервать}
  \FOR {$v \in \Gamma(u)$}
  \IF{$v\in Q$}
    \STATE $\Relax(s, v, u)$
    \COMMENT {\color{gray}\small релаксация дуг, ведущих из $u$ в $Q$}
  \ENDIF  
  \ENDFOR
\ENDWHILE

\end{algorithmic}

\end{frame}


\begin{frame}
\frametitle{8.6.4. Алгоритм Дейкстры (2/3)}
%\vspace{-2pt}
Алгоритм Дейкстры работает итеративно.
На каждой итерации алгоритм выполняет операцию $\ExtractMin$,
 извлекая узел $u$  с наименьшей оценкой $T[u]$ из контейнера 
 $Q$, и выполняет релаксацию дуг, исходящих из
 узла $u$ и ведущих в узлы, ещё остающиеся в 
 контейнере $Q$.
Начнём с некоторых наблюдений относительно работы алгоритма 
Дейкстры.

{\color{blue}\bf [1]}
Алгоритм в любом случае начинает и заканчивает работу.
Действительно, по крайней мере одна итерация
выполняется в любом случае, поскольку при инициализации
 $T[s]\Assign 0 < +\infty$, и не более $p$ 
итераций выполняется всего, поскольку на каждой итерации
из контейнера $Q$ извлекается один элемент.

{\color{blue}\bf [2]}
Алгоритм может завершить работу досрочно.
Действительно, 
 если $T[u] = +\infty$, то  $\A{x\in Q}{T[x]=+\infty}$,
  а это означает, что 
 ни разу не проводилась релаксация дуг, направленных в узлы $Q$ из
 уже рассмотренных узлов, а это, в свою очередь, означает, 
 что оставшиеся узлы в $Q$ недостижимы из узла $s$,
возможная часть дерева кратчайших\hspace{-2pt} 
путей\hspace{-2pt} уже\hspace{-2pt} 
построена\hspace{-2pt} и\hspace{-2pt} 
работу\hspace{-1pt} алгоритма\hspace{-1pt} можно\hspace{-2pt} прервать. 

\end{frame}


\begin{frame}
\frametitle{8.6.4. Алгоритм Дейкстры (3/3)}
{\color{blue}\bf [3]}
Алгоритм в любом случае строит ориентированное дерево, 
представленное массивом $\Pi$ с корнем в узле $s$.
Действительно, поскольку при первой релаксации дуги, 
входящей в узел $v$, 
когда оценка 
$\infty$ меняется на конечную величину, к дереву добавляется лист,
а при последующих релаксациях, когда оценка уменьшается,
происходит переназначение предка для узла $v$,
причём вновь назначаемый предок не является потомком узла $v$
по построению, так что контуров появиться не может.

{\color{blue}\bf [4]}
Алгоритм не меняет оценку узлов, извлеченных из 
контейнера $Q$.
Действительно, релаксируются только дуги,
направленные в узлы контейнера $Q$, 
следовательно, могут меняться оценки только
узлов из $Q$.   

\begin{example}
На рисунке показана работа алгоритма Дейкстры.

\vspace{-6pt}
\begin{center}
\includegraphics[width=150mm]{8_6_4.jpg}
\end{center}
\end{example}

\end{frame}


\begin{frame}
\label{8.6.5.}\subsubsection{8.6.5. Область применимости алгоритма Дейкстры }\frametitle{8.6.5. Область применимости алгоритма Дейкстры (1/8)}

\begin{digress}
В своей оригинальной статье Дейкстра употребил
термин <<длина>>, применительно к дугам и путям,
подразумевая, что длина не может быть отрицательной. 
Действительно, в положительно взвешенном орграфе алгоритм Дейкстры
всегда применим и находит решение,
поскольку контуры отрицательного веса невозможны.
Поэтому в литературе  область 
применимости алгоритма Дейкстры не обсуждается, 
и рассматриваются 
только положительно взвешенные графы и орграфы.
На самом деле класс орграфов,
в которых применим алгоритм Дейкстры, существенно шире,
что видно из предыдущего примера.      
\end{digress}

\medskip
Назовём \odmkey{рекордом пути}{Рекорд пути (path record)} 
$M=\Path{u}{w}$ (обозначение $R(\Path{u}{w})$) максимальный вес начального отрезка этого пути, 
$R(\Path{u}{w})\Def\max\limits_{v\in M-u}W(u,v)$. 
Отдельно отметим, что $v\neq u$, то есть рекорд не может достигаться в начальном узле по определению.
Из определения следует, что 
$$\A{\Path{u}{w}}{R(\Path{u}{w})\geq W(\Path{u}{w}) \geq 
\Omega(u,w)}.$$


\end{frame}
 

\begin{frame}
\frametitle{8.6.5. Область применимости алгоритма Дейкстры (2/8)}
Рассмотрим сеть путей $P(s,t)$ из
узла $s$ в узел $t$. 
Заметим, что для любого промежуточного узла $v$
в сети $P(s,t)$ определены  сети $P(s,v)$ и $P(v,t)$,
причём они являются \em{сегментами} исходной сети:  

\vspace{-8pt}
$$
\A{v\in \Path{s}{t}}{P(s,v)\subset P(s,t)
\And P(v,t)\subset P(s,t)}.
$$

\vspace{-2pt}
Далее можно взять промежуточный узел в одном 
из сегментов и получить ещё два сегмента,
причём один из них уже не будет инцидентен
исходным узлам  $s$ и $t$. Таким образом можно
декомпозировать сеть на сегменты вплоть до отдельных дуг.   
 
Пути в сегментах могут не пересекаться, 
а могут пересекаться и даже слипаться. 
Если начальный узел некоторого сегмента имеет только одну
исходящую дугу, то такой сегмент называется
\em{безальтернативным}. 
Все пути безальтернативного сегмента слиплись в первой дуге.  
Если из начального узла исходит более одной дуги, то 
сегмент называется 
\em{альтернативным}, или \em{развилкой}.

\begin{remark}
При поиске кратчайших путей в сегменте 
безальтернативные участки путей можно пропускать,
и начинать поиск с первой развилки. 
\end{remark} 

\end{frame}


\begin{frame}
\frametitle{8.6.5. Область применимости алгоритма Дейкстры (3/8)}

\vspace{-1pt}
\begin{theorem}
Алгоритм Дейкстры на взвешенном 
орграфе без контуров отрицательного веса строит дерево кратчайших путей с корнем в узле $s$ тогда и только тогда, 
когда для каждого узла $t \in G -s$, 
достижимого из узла $s$, в сети $P(s,t)$ для каждой развилки
$P(w,t)$ рекорд кратчайшего пути
$\Path{w}{t}$ строго меньше, чем рекорд любого не кратчайшего пути 
$\Path{w}{t}$.
\end{theorem}

\vspace{-1pt}
\begin{proof} 
Заметим, что когда  узел $u$ извлечён из контейнера,
все узлы на пути  $\Path{s}{u}$ уже были извлечены ранее,
и поэтому оценка извлечённого узла не превосходит рекорда найденного пути
к нему:  $T[u]\leq R(\Path{s}{u})$. 
\proofsection{Достаточность}
Покажем по индукции,
что  на каждой итерации 
при извлечении узла $u$
путь $M_0=\Path{s}{u}$, определяемый значением $\Pi[u]$,
кратчайший, то есть 
$$W(\Path{s}{u})=\Omega(s,u).$$
Тем самым покажем, 
что на каждой итерации извлечённые узлы и пути к ним
образуют 
частичное дерево кратчайших путей из узла $s$,
и по завершении работы алгоритма построено
дерево кратчайших путей из узла $s$ ко всем достижимым узлам.


\end{proof}
\end{frame}


\begin{frame}
\frametitle{8.6.5. Область применимости алгоритма Дейкстры (4/8)}
\begin{proof} (Продолжение)
База: на первой итерации всегда выбирается
узел $s$, который образует корень дерева.
Длина кратчайшего пути $\Omega(s,s)=0$, 
поскольку контуры  запрещены.
Индукционное предположение:
пусть на предыдущих итерациях
построено частичное дерево кратчайших путей и пусть
 из контейнера $Q$
выбран узел $u$ с оценкой $T[u]<\infty$ и
  узел $u$ получил свою оценку в результате 
релаксации дуги $(v,u)$. Если путь $\Path{s,\dots,v}{u}$ 
не единственный, то
в сети $P(s,u)$ рассмотрим
первую (считая от узла $s$) развилку $P(w,u)$ и   
возможные альтернативные пути:  $M_0=\Path{w,\dots,v}{u}$,
  $M_1$, в котором все узлы, кроме последнего,
уже исключены из $Q$ на предыдущих итерациях,
$M_2$, в котором есть узлы, принадлежащие $Q$.

\vspace{-6pt}
\begin{center}
\includegraphics[width=100mm]{8_6_5_1.jpg}
\end{center}

\vspace{-12pt}
\end{proof}
 
\end{frame}


\begin{frame}
\frametitle{8.6.5. Область применимости алгоритма Дейкстры (5/8)}
\begin{proof} (Продолжение)
В случае пути $M_1$  дуги $(v,u)$
и $(z,u)$ уже участвовали в релаксации,
и в построенном дереве сохранена та,
которая даёт меньшую оценку, поэтому путь $M_1$
не оказывает влияния на построенное дерево
и его можно не учитывать. 
Далее от противного. Пусть  $W(\Path{s,\dots,w}{M_2})<W(\Path{s,\dots,w}{M_0})$.
По условию теоремы это означает, что $R(M_2)<R(M_0)$. 
Рассмотрим те узлы, на которых достигаются рекорды этих путей.
Пусть $u_0\in M_0$~--- место достижения рекорда пути $M_0$.
Тогда $T[u_0]=T[w]+R(M_0)$ (узел $u_0$ 
может быть узлом $u$, $v$ или любым промежуточным узлом).    
Рассмотрим ту итерацию, когда был выбран узел $u_0$,
и
пусть при этом $u_2$~--- первый ещё не извлечённый узел на пути $M_2$
(узел $u_2$ 
может быть узлом $y$, $x$ или любым промежуточным узлом).
Тогда 
$T[u_0]-T[w]=R(M_0)>R(M_2)\geq T[u_2]-T[w],$
то есть $T[u_0]> T[u_2]$,
что противоречит тому, что узел $u_0$
извлечён раньше узла $u_2$. 
Доказательство очевидным образом распространяется 
на все последующие развилки.
\end{proof}
 
\end{frame}


\begin{frame}
\frametitle{8.6.5. Область применимости алгоритма Дейкстры (6/8)}
%\vspace{-2pt}
\begin{proof} (продолжение)
\proofsection{Необходимость} 
Нужно показать, что если алгоритм построил дерево кратчайших путей,
то в любой развилке рекорд кратчайшего пути
меньше рекорда не кратчайшего пути, 
или, в контрапозитивной форме,
если существует развилка, 
в которой рекорд кратчайшего пути 
больше либо равен рекорду не кратчайшего пути,
то алгоритм <<ошибётся>> и 
построенное дерево не будет деревом кратчайших путей.    

\smallskip
Рассмотрим сначала случай, когда 
рекорд кратчайшего пути строго больше рекорда
не кратчайшего пути.

\smallskip
Пусть алгоритм Дейкстры построил дерево, 
и пусть сущеcтвует узел $u$, такой, что
в него ведут два пути: путь $M_{0}^{'}$~--- 
кратчайший,   и $M_{2}^{'}$~--- не кратчайший, 
$W(M_{0}^{'})<W(M_{2}^{'})$,  
и $P(w,u)$~--- первая развилка этих путей,
$M_0$ и $M_2$~--- продолжения путей после развилки
соответственно, 
причём $R(M_0)> R(M_2)$ (рисунок на следующем слайде).
\end{proof}

\end{frame}


\begin{frame}
\frametitle{8.6.5. Область применимости алгоритма Дейкстры (7/8)}


\begin{proof} (Продолжение)
\begin{center}
\includegraphics[width=80mm]{8_6_5_2.jpg}
\end{center}

Рассмотрим узел $v\in M_0$, в котором достигается рекорд,
то есть $T[v]=R(M_0)$. Тогда 
$\A{x\in M_2}{T[x]\le R(w,x)< R(M_0)=T[v]}$,
и алгоритм выберет все узлы 
пути $M_2$, включая узел $u$, раньше, 
чем узел $v$, а значит какие-то 
дуги на пути $\Path{v}{u}$ 
не будут участвовать в релаксации и кратчайший
путь не будет найден.

\smallskip
		 Рассмотрим теперь случай, когда рекорд кратчайшего пути \em{равен} рекорду не кратчайшего пути.
В такой ситуации корректность работы алгоритма Дейкстры зависит от воли случая, то есть от того,
какой из узлов с равными значениями оценки $T$ будет выбран. 
\end{proof}
\end{frame}


\begin{frame}
\frametitle{8.6.5. Область применимости алгоритма Дейкстры (8/8)}
\begin{proof} (Окончание)
Например, в орграфе на рисунке ниже с начальным узлом $s$ на первом шаге
оба узла $u$ и $v$ получают одинаковые оценки
$T[u]=T[v]=1$. Если на втором шаге будет выбран узел $v$, 
то произойдет релаксация дуги $(v, u)$ и будет
найден кратчайший путь $\Path{s,v}{u}$.
Если же на втором шаге будет выбран узел $u$,
то релаксации дуги $(v,u)$ не произойдет и
 кратчайший путь не будет найден.      	 

\begin{center}
\begin{figure}
\includegraphics[width=4cm]{fig_bad_choice.png}
\end{figure}
\end{center}
Мы считаем, что в тех случаях, когда мы не можем быть уверены в корректной работе алгоритма, следует считать работу алгоритма некорректной. 		
\end{proof}
\end{frame}


\begin{frame}
\label{8.6.6.}\subsubsection{8.6.6. Анализ алгоритма Дейкстры }\frametitle{8.6.6. Анализ алгоритма Дейкстры (1/3)}

\begin{example}
На рисунках приведены примеры работы алгоритма на орграфах
с положительными и отрицательными весами дуг.  
На первом рисунке  рекорд кратчайшего пути $s,u,t$ равен $-4$,
в то время как рекорд пути $s,v,t$ равен $0$, 
и алгоритм работает правильно.



\begin{figure}[h]
\begin{center}
\includegraphics[width=140mm]{8_6_5_3.jpg}
\end{center}
\end{figure}

На втором рисунке рекорд кратчайшего пути $s,u,t$ равен $+1$,
но рекорд пути $s,v,t$ равен $0$, 
и алгоритм работает неправильно.

\begin{figure}[h]
\begin{center}
\includegraphics[width=140mm]{8_6_5_4.jpg}
\end{center}
\end{figure}

\end{example}

\end{frame}


\begin{frame}
\frametitle{8.6.6. Анализ алгоритма Дейкстры (2/3)}

\begin{corollary}
Алгоритм Дейкстры корректен на классе неотрицательно взвешенных
орграфов.
\end{corollary}
\begin{proof}
Неотрицательно взвешенный орграф очевидно не содержит контуров отрицательного веса. Далее заметим, что из неотрицательности весов следует, что рекорд любого пути достигается в его конечном узле, поэтому рекорд кратчайшего пути меньше 
рекордов альтернативных или равен им в конечном узле.
\end{proof}

\medskip
Поскольку алгоритм Дейкстры не имеет в себе никакой разумной реакции на контуры отрицательного веса, а анализ орграфов на выполнение необходимого и достаточного условия обычно сложен,
алгоритм повсеместно применяется только для неотрицательно взвешенных орграфов, что покрывает подавляющее большинство прикладных задач поиска кратчайшего пути.
\end{frame}


\begin{frame}
\frametitle{8.6.6. Анализ алгоритма Дейкстры (3/3)}

\begin{remark}
Алгоритмом Дейкстры можно найти кратчайший путь между двумя узлами. Для этого достаточно после операции $\ExtractMin(Q)$ завершить работу, если извлеченный узел $u$ является целевым, поскольку $T[u]$ для него больше не изменится.
\end{remark}

\smallskip
Время работы алгоритма Дейкстры зависит от реализации контейнера $Q$ и вектора $T$ с операциями построения $Q \Assign V$, уменьшения $T[v]$ в процедуре релаксации и вызова $\ExtractMin(Q)$. При работе алгоритма каждая дуга ослабляется не более чем один раз, поскольку начальные узлы дуг $u$ не повторяются, что в сумме произойдёт не более чем $q$ раз. 
Пусть $Q$ --- массив, тогда одна вставка выполняется за $O(1)$, а все вставки --- за $O(p)$, 
уменьшения значений $T$ займут в сумме $qO(1)=O(q)$, а  $pO(p) = O(p^2)$.
Итого имеем 
$O(p)+O(q)+O(p^2) = O(p^2 + q) = O(p^2)$.
Пусть $\ExtractMin$ и уменьшение $T$ выполняются за $O(\log p)$, а присваивание $Q \Assign V$ --- за $O(p)$. Тогда \\ 
$O(p)+pO(\log p)+qO(\log p) = O((p+q) \log p))$, что будет меньше $O(p^2)$, \\если $q < p^2 / \log p$. Такими свойствами обладают \em{неубывающие пирамиды}.

 
\end{frame}
 
 



\begin{frame}
\label{8.6.7.}\subsubsection{8.6.7. Алгоритм Флойда--Уоршалла }\frametitle{8.6.7. Алгоритм Флойда--Уоршалла (1/3)}

В отличие от предыдущих алгоритмов, 
\odmkey{алгоритм Флойда--Уоршалла}
{Алгоритм (algorithm)!Флойда--Уоршалла (Floyd--Warshall)} решает задачу о поиске не дерева кратчайших путей, а \em{всех} кратчайших путей в орграфе,
и делает это за время $O(p^3)$.

\begin{algorithmic}
\REQUIRE матрица $W[1..p,1..p]$ длин дуг.
\ENSURE матрица $T[1..p,1..p]$ длин путей и матрица $\Pi [1..p,1..p]$ самих путей.

\STATE \COMMENT{\color{gray}\small инициализация }
\FOR{\textbf{$i$ from $1$ to $p$}}
  \FOR{\textbf{$j$ from $1$ to $p$}}
    \STATE $T[i,j]\Assign W[i,j]$ 
    \IF{$W[i,j]=\infty$}
      \STATE $\Pi[i,j]\Assign 0$ 
      \COMMENT{\color{gray}\small  нет дуги из $i$ в $j$ }
    \ELSE
      \STATE $\Pi[i,j]\Assign j$ 
      \COMMENT{\color{gray}\small  есть дуга из $i$ в $j$ }
    \ENDIF
  \ENDFOR
\ENDFOR
\end{algorithmic}

\end{frame}

\begin{frame}
\frametitle{8.6.7. Алгоритм Флойда--Уоршалла (2/3)}

\begin{algorithmic}
\FOR{\textbf{$i$ from $1$ to $p$}}
  \FOR{\textbf{$j$ from $1$ to $p$}}
    \FOR{\textbf{$k$ from $1$ to $p$}}
      \IF{$i\ne j \And T[j,i]\ne\infty \And i\ne k \And T[i,k]\ne\infty \And$ \\
      $\phantom{\textbf{if }} \And T[j,k] >T[j,i]+T[i,k]$}
        \STATE $T[j,k]\Assign T[j,i]+T[i,k]$ 
        \COMMENT{\color{gray}\small  запомнить длину нового пути}
        \STATE $\Pi[j,k]\Assign \Pi[j,i]$ 
        \COMMENT{\color{gray}\small   и сам путь }
      \ENDIF
    \ENDFOR 
  \ENDFOR 
  \FOR{\textbf{$j$ from $1$ to $p$}}
    \IF{$T[j,j]<0$}
      \STATE \textbf{stop} 
      \COMMENT{\color{gray}\small  узел $j$ входит в цикл отрицательной длины }
    \ENDIF
  \ENDFOR
\ENDFOR
\end{algorithmic}

\end{frame}


\begin{frame}
\frametitle{8.6.7. Алгоритм Флойда--Уоршалла (3/3)}

\begin{ground}
Алгоритм Флойда--Уоршалла является обобщением алгоритма Уоршалла (\DMref{п.~1.5.2}{1.5.2.}). 
Покажем по индукции, что после выполнения $i$-го
шага основного цикла элементы матриц $T[j,k]$ и $\Pi[j,k]$
содержат, соответственно, длину кратчайшего пути и первый узел
на кратчайшем пути из узла $j$ в~узел $k$, проходящем через
промежуточные узлы из диапазона $1..i$. База: $i=0$, то~есть до
начала цикла элементы матриц $T$ и $\Pi$ содержат информацию о
кратчайших путях (дугах!), не проходящих ни через какие
промежуточные узлы. Пусть теперь перед началом выполнения тела
цикла на $i$-м шаге $T[j,k]$ содержит длину кратчайшего пути от
$j$ к $k$, а $\Pi[j,k]$ содержит первый узел (если таковой есть) на кратчайшем пути
из узла $j$ в узел $k$. В таком случае,
если в результате добавления узла $i$ к диапазону промежуточных
узлов находится более короткий путь (в частности, если это 
первый найденный путь), то он записывается.
%Таким образом,
Таким образом,
после окончания цикла, когда $i=p$, матрицы содержат кратчайшие пути,
проходящие через промежуточные узлы $1..p$, то~есть искомые
кратчайшие пути. 
Алгоритм не всегда выдаёт решение, поскольку оно не всегда определено. 
Дополнительный цикл по $j$ служит для
прекращения работы в случае обнаружения
% в орграфе
контура с отрицательным весом.
\end{ground}


\end{frame}


\begin{frame}
\label{8.6.8.}\subsubsection{8.6.8. Алгоритмы с предобработкой }\frametitle{8.6.8. Алгоритмы с предобработкой (1/6)}

Как уже было сказано, алгоритм Дейкстры обычно применяется только к неотрицательно взвешенным орграфам. 

\smallskip
Однако, в некоторых случаях можно провести \em{предобработку},
то есть предварительное преобразование данных задачи в
более выгодный для решения вид.

\smallskip
В данном случае предобработка означает
преобразование исходного орграфа к неотрицательно взвешенному
орграфу, 
применение алгоритма Дейкстры и поиск кратчайшего пути, а потом преобразование найденного пути в путь в исходном 
орграфе. 

\smallskip
\begin{digress}
Предобработка имеет смысл, если предстоит решать не разовую, а \em{массовую} задачу. Скажем, искать кратчайшие пути на графе,
представляющем карту автомобильных дорог. 
Такая задача в навигационных приложениях
должна решаться многие тысячи раз в день, 
поэтому затрат времени на предобработку не жаль,
если предобработка позволяет ускорить
выполнение основной задачи поиска кратчайшего пути.    
\end{digress}
\end{frame}


\begin{frame}
\frametitle{8.6.8. Алгоритмы с предобработкой (2/6)}

Одним из наиболее показательных алгоритмов с предобработкой является
 \odmkey{алгоритм\\ Джонсона}{Алгоритм (algorithm)!Джонсона (Jonson)}, предложенный в 1977 году.
В этом алгоритме  
%(Дональд Б. Джонсон) 
поиска кратчайших путей между всеми узлами орграфа используется предобработка типа
\textit{<<reweighting>>} (<<перевзвешивание>>)~--- введение функции потенциала. 
Каждому узлу $v$ присваивается \em{потенциал} $h(v)$. 
Веса дуг переопределяются.
Для каждой дуги $e=(u,v)$ новый вес
определяется через старый вес и потенциалы 
концов: 
$w'(e)\Assign w(e)+h(u)-h(v).$ 
Потенциалы определяются так, чтобы новые веса были
неотрицательными. Можно показать, что
такое преобразование сохраняет отношение порядка весов путей $\Path{u,t_1,t_2,\dots,}{v}$ из 
узла $u$ в узел $v$:

% TODO не вышло нормально запилить эту формулу в $$ $$
\begin{align*}
W'(\langle u,v\rangle)= &\ w((u,p_{1}))+h(u)-h(p_{1})+w((p_{1},p_{2}))
+h(p_{1})-h(p_{2})+...\\&+w((p_{n},v))+h(p_{n})-h(v)=\\= & \ W(\langle u,v\rangle )+h(u)-h(v)=W(\langle u,v\rangle)+c
\end{align*}


Поскольку прибавление константы $c$ сохраняет отношение порядка весов путей,
по результатам работы алгоритма Дейкстры очевидным образом можно найти 
кратчайшие пути в исходном орграфе.
\end{frame}


\begin{frame}
\frametitle{8.6.8. Алгоритмы с предобработкой (3/6)}

Осталось приписать потенциалы узлам так, чтобы новые веса дуг стали неотрицательными.
В своем алгоритме Джонсон предлагает следующий
метод поиска требуемых 
потенциалов.

\begin{enumerate}
\item Фиктивный узел $s$ дугами нулевого веса присоединить
к реальным узлам орграфа.
\item Воспользоваться алгоритмом Беллмана--Форда для поиска кратчайших путей из узла $s$ во все реальные узлы орграфа.
Нетрудно понять, что для неотрицательно взвешенного орграфа веса всех таких путей будут нулевыми, в то время как для других
орграфов это, вообще говоря, не так.
\item Взять $\Omega(s,v)$ в качестве $h(v)$, то есть искомых потенциалов, после чего удалить фиктивный узел и дуги.
\end{enumerate}

\begin{remark}
По наследству от алгоритма Беллмана--Форда,
метод  работоспособен для всех взвешенных орграфов без циклов отрицательного веса и выдаёт предупреждение в случае наличия таких циклов.
\end{remark}

\end{frame}


\begin{frame}
\frametitle{8.6.8. Алгоритмы с предобработкой(4/6)}

\begin{ground}
Предложенный способ построения потенциальной функции корректен.\\
Действительно, если в неравенство треугольника 
$\Omega(s,v)\le \Omega(s,u)+w((u,v))$
подставить весовую функцию Джонсона  $h(v) = \Omega(s,v)$, $h(u) = \Omega(s,u)$, получим\\
$h(v)\le h(u)+w((u,v))$, и тогда
$w'((u,v))=w((u,v))+h(u)-h(v)\ge 0$.
\end{ground}

\begin{example}
На следующем рисунке приведены диаграммы
(слева направо): исходного орграфа с отрицательными дугами, дополнительного построения для предобработки, перевзвешенного орграфа. Кратчайшие пути 
$\Path{u}{t}$ и дополнительные построения алгоритма Джонсона выделены красным цветом.  
\end{example}

\vspace{-6pt}
\begin{figure}[h]
\begin{center}
\includegraphics[width=120mm]{8_6_5_5.png}
\end{center}
\end{figure}

\end{frame}

%%%%%%%%%%%%%%%%%%%%%%%%%%%%%%%%%%%%%
\begin{frame}
\frametitle{8.6.8. Алгоритмы с предобработкой(5/6)}
\begin{remark}
Из орграфа с отрицательными дугами
получитcя неотрицательно взвешенный орграф,
если просто увеличить веса всех дуг на величину
$c\Assign - \min_{e\in E} w(e)$.
Однако при таком тривиальном перевзвешивании 
алгоритм Дейкстры может дать неверный ответ.
Нетрудно показать, что при тривиальном перевзвешивании
алгоритм Дейкстры работает правильно, только если  
для любой пары путей 
$M_1\Assign \Path{s,v_1,\dots}{v_m,t}$ и \\
$M_2\Assign \Path{s,u_1,\dots}{u_n,t}$, 
$W(M_1)\le W(M_2)$ оказывается,
что если $W(M_1)< W(M_2)$, \\то $W(M_2)-W(M_1)>c(m-n)$,
а если $W(M_1)= W(M_2)$, то $m=n$.  
\end{remark}

\begin{example}
 Слева приведен пример неудачного, а справа~---
удачного перевзвешивания. Кратчайшие пути  выделены красным цветом.  
\end{example}

\vspace{-6pt}
\begin{figure}[h]
\begin{center}
\includegraphics[width=140mm]{8_6_5_6.png}
\end{center}
\end{figure}

\end{frame}

\begin{frame}
\frametitle{8.6.8. Алгоритмы с предобработкой (6/6)}

Рассмотрим оценки трудоёмкости.
Пусть алгоритм используется для поиска одного дерева кратчайших путей.
Предварительная обработка алгоритмом Беллмана--Форда~--- $O(pq)$, 
алгоритм Дейкстры~--- $O(p\log p+q)$, обратное преобразование дерева~--- $O(p)$.
В сумме имеем $O(pq)+O(p\log p+q) + O(p) = O(pq)$, 
что, очевидно, не дает никакого выигрыша в сравнении с алгоритмом 
Беллмана--Форда.

Но если применить алгоритм для поиска $p$ деревьев кратчайших путей, предварительную обработку потребуется провести только \em{один} раз.
В этом и есть главная идея предобработки!

\begin{enumerate}
\item Провести единственную предобработку. 
\item Найти $p$ деревьев кратчайших путей алгоритмом Дейкстры.
\item Сделать $p$ обратных преобразований, выдать ответ.
\end{enumerate}

Время работы алгоритма Джонсона (при <<сложной>> реализации алгоритма Дейкстры) ---
$O(pq) + O(p\cdot (p\log p+q)) + pO(p) $ $=O(p^{2}\log p + pq)$,
что получилось лучше, чем алгоритм 
Флойда--Уоршалла ($O(p^{3})$), если $q$ асимптотически меньше, чем $p^2$.

\end{frame}
\begin{frame}
\frametitle{Онтология теории графов}
\begin{center}
\includegraphics[height=8cm]{ODM_8_6.png}
\end{center}

\end{frame}

%%%%%%%%%%%%%%%%%%%%%%%%%%%%%%%%%%%%%%%%%%%%%%%%%%%%%%%%%%%%%%%%%%%%%%%%%%%%%
% Далее идёт вариант ДО работы Сергея Лебедева
%%%%%%%%%%%%%%%%%%%%%%%%%%%%%%%%%%%%%%%%%%%%%%%%%%%%%%%%%%%%%%%%%%%%%%%%%%%%
%%
%%\subsection{8.6. Кратчайшие пути}
%%
%%\begin{frame}
%%\frametitle{8.6. Кратчайшие пути}
%%
%%Задача нахождения кратчайшего пути
%%в графе имеет столько практических
%%применений и интерпретаций, что
%%читатель, без сомнения, может сам
%%легко привести множество примеров.
%%Здесь рассматриваются четыре классических
%%алгоритма, которые обязан знать
%%каждый программист.
%%
%%%\vspace*{1.5\baselineskip}
%%%{\color{blue} Тема для сочинения:
%%%а почему именно четыре? Может быть, надобно пять? или шесть?
%%%
%%%}
%%
%%\end{frame}
%%
%%
%%\begin{frame}
%%\frametitle{8.6.1. Длина дуг}
%%
%%Алгоритм Уоршалла  позволяет ответить на вопрос,
%%достижима ли вершина $v$ из вершины $u$, то есть
%%существует ли цепь $\divC{u}{v}$. Часто нужно не только определить,
%%существует ли цепь, но и найти эту цепь.
%%
%%
%%Если задан орграф $G(V,E)$, в котором дуги нагружены числами (эти
%%числа обычно называют \odmkey{весами}{Вес дуги (weight of arc)},
%%или \odmkey{длинами}{Длина (length)!дуги (of arc)}, дуг),  то этот
%%орграф можно представить в виде матрицы весов (длин) $C$:
%%$$
%%C[i,j]=
%%\begin{cases}
%%0,     & \text{для $i=j$},\\
%%c_{ij},& \text{конечная величина, если есть дуга из узла $i$ в узел  $j$,} \\
%%\infty,& \text{если нет дуги из узла $i$ в узел $j$.}
%%\end{cases}
%%$$
%%
%%
%%
%%\odmkey{Длиной пути}{Длина (length)!пути (of path)} называется
%%сумма длин дуг, входящих в этот путь. Наиболее часто на практике
%%встречается задача отыскания \odmkey{кратчайшего}{Путь
%%(path)!кратчайший (shortest)} пути.
%%
%%%\vspace*{-.5\baselineskip}
%%
%%\begin{remark}
%%Данное представление, равно как и последующие алгоритмы, применимо как к графам,
%%так и к орграфам.
%%\end{remark}
%%
%%%\vspace*{-.5\baselineskip}
%%
%%\end{frame}
%%
%%\begin{frame}
%%\frametitle{8.6.2. Алгоритм Флойда }
%%
%%Алгоритм Флойда находит кратчайшие
%%пути между всеми парами вершин (узлов) в (ор)графе. В этом
%%алгоритме для хранения информации о путях используется матрица
%%$H$: \textbf{array} $[1..p,1..p]$ \textbf{of} $1..p$, где
%%$$
%%H[i,j] = \!
%%\begin{cases}
%%k, & \!\!\!\!\text{если $k\!$~--- первая вершина, достигаемая}\\[-4pt]
%%    & \!\!\!\!\text{на кратчайшем пути из $i$ в $j$;}\\
%%0, & \!\!\!\!\text{если из вершины $i$ в вершину $j$ нет пути.}
%%\end{cases}
%%$$
%%
%%%\addvspace{-0.5\baselineskip}
%%%\begin{algorithm}
%%%\caption{Алгоритм Флойда}
%%%\begin{algorithmic}
%%%\REQUIRE матрица $C[1..p,1..p]$ длин дуг.
%%%\ENSURE
%%%матрица $T[1..p,1..p]$ длин путей и
%%%матрица $H[1..p,1..p]$ самих путей.
%%%\FOR{\textbf{$i$ from $1$ to $p$}}
%%%\FOR{\textbf{$j$ from $1$ to $p$}}
%%%\STATE $T[i,j]\Assign C[i,j]$ \COMMENT{ инициализация }
%%%\IF{$C[i,j]=\infty$}
%%%  \STATE $H[i,j]\Assign 0$ \COMMENT{ нет дуги из $i$ в $j$ }
%%%\ELSE
%%%  \STATE $H[i,j]\Assign j$ \COMMENT{ есть дуга из $i$ в $j$ }
%%%\ENDIF \ENDFOR \ENDFOR \FOR{\textbf{$i$ from $1$ to $p$}}
%%%\FOR{\textbf{$j$ from $1$ to $p$}} \FOR{\textbf{$k$ from $1$ to
%%%$p$}} \IF{$i\ne j \And T[j,i]\ne\infty \And i\ne k \And
%%%T[i,k]\ne\infty\And (T[j,k]=\infty \lor T[j,k] >T[j,i]+T[i,k])$}
%%%\STATE $H[j,k]\Assign H[j,i]$ \COMMENT{ запомнить новый путь }
%%%\STATE $T[j,k]\Assign T[j,i]+T[i,k]$ \COMMENT{ и его длину }
%%%\ENDIF \ENDFOR \ENDFOR \FOR{\textbf{$j$ from $1$ to $p$}}
%%%\IF{$T[j,j]<0$}
%%%  \STATE \textbf{stop} \COMMENT{ нет решения: вершина $j$ входит в цикл отрицательной
%%%длины }
%%%\ENDIF
%%%\ENDFOR
%%%\ENDFOR
%%%\end{algorithmic}
%%%\end{algorithm}
%%%
%%%\addvspace{-0.25\baselineskip}
%%\begin{digress}
%%Матрица $H$ размера $O(p^2)$ хранит информацию обо всех
%%(кратчайших) путях в графе. Заметим, что всего в графе $O(p^2)$
%%путей, состоящих из $O(p)$ вершин. Таким образом, непосредственное
%%представление всех путей потребовало бы памяти объема $O(p^3)$.
%%Экономия памяти достигается за счет \odmkey{интерпретации
%%представления}{Интерпретация (interpretation)!представления (of
%%representation)}, то~есть динамического вычисления некоторой части
%%информации вместо её хранения в памяти. 
%%%В данном случае любой
%%%конкретный путь $\lrle{u,v}$ легко извлекается из матрицы с
%%%помощью следующего алгоритма.
%%%
%%%%\pagebreak[3]{ \small
%%%\begin{algorithmic}
%%%\STATE $w\Assign u$; \textbf{yield $w$} \COMMENT{ первая вершина }
%%%\WHILE{$w\ne v$}
%%%\STATE $w\Assign H[w,v]$; \textbf{yield $w$} \COMMENT{ следующая вершина }
%%%\ENDWHILE
%%%\end{algorithmic}
%%%}
%%\end{digress}
%%
%%%}%baselinestretch
%%
%%%\vspace*{-.5\baselineskip}
%%
%%
%%\begin{remark}
%%Если в $G$ есть цикл с отрицательным весом, то решения поставленной задачи
%%не существует, так~как на этом цикле можно <<накручивать>>  сколь угодно
%%короткий путь.
%%\end{remark}
%%
%%
%%%
%%%\begin{ground}
%%%Алгоритм Флойда имеет много общего с алгоритмом Уоршалла (алгоритм
%%%\ref{warshall}). Покажем по индукции, что после выполнения $i$-го
%%%шага основного цикла по $i$ элементы матриц $T[j,k]$ и $H[j,k]$
%%%содержат, соответственно, длину кратчайшего пути и первую вершину
%%%на кратчайшем пути из вершины $j$ в~вершину $k$, проходящем через
%%%промежуточные вершины из диапазона $1..i$. База: $i=0$, то~есть до
%%%начала цикла элементы матриц $T$ и $H$ содержат информацию о
%%%кратчайших путях (если таковые есть), не проходящих ни через какие
%%%промежуточные вершины. Пусть теперь перед началом выполнения тела
%%%цикла на $i$-м шаге $T[j,k]$ содержит длину кратчайшего пути от
%%%$j$ к $k$, а $H[j,k]$ содержит первую вершину на кратчайшем пути
%%%из вершины $j$ в вершину $k$ (если таковой есть). В таком случае,
%%%если в результате добавления вершины $i$ к диапазону промежуточных
%%%вершин находится более короткий путь (в частности, если это вообще
%%%первый найденный путь), то он записывается. Таким образом, после
%%%окончания цикла, когда $i=p$, матрицы содержат кратчайшие пути,
%%%проходящие через промежуточные вершины $1..p$, то~есть искомые
%%%кратчайшие пути. Алгоритм не всегда выдаёт решение, поскольку оно
%%%не всегда существует. Дополнительный цикл по $j$ служит для
%%%прекращения работы в случае обнаружения в графе цикла с
%%%отрицательным весом.
%%%\end{ground}
%%%
%%%\vspace*{-.5\baselineskip}
%%
%%
%%\end{frame}
%%
%%
%%\begin{frame}
%%\label{8.6.3.}\subsubsection{8.6.3. Алгоритм Дейкстры }\frametitle{8.6.3. Алгоритм Дейкстры (1/4)}
%%
%%Алгоритм Дейкстры находит
%%кратчайший путь между  двумя данными вершинами (узлами) в
%%(ор)графе, если длины рёбер (дуг) неотрицательны.
%%%
%%%\begin{algorithm}
%%%\caption{Алгоритм Дейкстры}
%%%\label{algDijkstra}
%%%\vspace*{-4.3pt}
%%%\end{algorithm}
%%%\vspace*{-4.3pt}
%%\begin{algorithmic}
%%\REQUIRE
%%орграф $G(V,E)$, заданный матрицей длин дуг
%%\textbf{$C:$ array $[1..p,1..p]$ of real};\\
%%$s$ и $t$~--- вершины графа. \ENSURE векторы \textbf{$T:$ array
%%$[1..p]$ of real}; и \textbf{$H:$ array $[1..p]$ of $0..p$}. Если
%%вершина $v$ лежит на кратчайшем пути от $s$  к $t$, то $T[v]$~---
%%длина кратчайшего пути от $s$  к~$v$; $H[v]$~--- вершина,
%%непосредственно  предшествующая $v$ на кратчайшем пути.
%%\FOR{\textbf{$v$ from $1$ to $p$}} \STATE $T[v]\Assign \infty$
%%\COMMENT{ кратчайший путь неизвестен } \STATE $X[v]\Assign 0$
%%\COMMENT{ все вершины не отмечены } \ENDFOR \pagebreak[3] \STATE
%%$H[s] \Assign 0 $ \COMMENT{ $s$ ничего не предшествует } \STATE
%%$T[s] \Assign 0 $ \COMMENT{ кратчайший путь имеет длину 0\ldots }
%%\STATE $X[s] \Assign 1 $ \COMMENT{ \ldots и он известен } \STATE
%%$v \Assign s $ \COMMENT{ текущая вершина }
%%\end{algorithmic}
%%\end{frame}
%%
%%
%%\begin{frame}
%%\frametitle{8.6.3. Алгоритм Дейкстры (2/4)}
%%
%%\begin{algorithmic}
%% \STATE $M$: \COMMENT{обновление пометок }
%%
%%\FOR{$u\in\Gamma(v)$}
%%  \IF{$X[u]=0\And T[u]>T[v]+C[v,u]$}
%%     \STATE $T[u]\Assign T[v]+C[v,u]$
%%     \COMMENT{ найден более короткий путь из $s$ в $u$ через $v$ }
%%     \STATE $H[u]\Assign v$ \COMMENT{ запоминаем его }
%%  \ENDIF
%%\ENDFOR
%%
%%
%%\STATE $m\Assign\infty;v\Assign 0$
%%\COMMENT{ поиск конца кратчайшего пути }
%%\FOR{\textbf{$u$ from $1$ to $p$}}
%%  \IF{$X[u]=0\And T[u]<m$}
%%    \STATE $v\Assign u; m\Assign T[u]$
%%    \COMMENT{ вершина $v$ заканчивает кратчайший путь из $s$ }
%%  \ENDIF
%%\ENDFOR
%%\STATE{\textbf{ if $v=0$ then stop end if}}
%% \COMMENT{ нет пути из $s$ в $t$ }
%%\STATE{\textbf{ if $v=t$ then stop end if}}
%%\COMMENT{ найден кратчайший путь из $s$ в $t$ }
%%\STATE $X[v]\Assign  1$ \COMMENT{ найден кратчайший путь из $s$ в $v$ }
%%\STATE\textbf{goto $M$}
%%\end{algorithmic}
%%
%%%\end{algorithm}
%%\end{frame}
%%
%%
%%\begin{frame}
%%\frametitle{8.6.3. Алгоритм Дейкстры (3/4)}
%%
%%%\vspace*{-.25\baselineskip}
%%
%%\begin{remark}
%%Для применимости алгоритма Дейкстры достаточно выполнения
%%\odmkey{неравенства треугольника}{Неравенство
%%(inequality)!треугольника (triangle)}:
%%$$
%%\A{u,v,w\in V}{d(u,v)\le d(u,w)+d(w,v)},
%%$$
%%которое, очевидно, выполняется, если длины дуг неотрицательны.
%%Если же допускаются отрицательные длины дуг, то алгоритм Дейкстры
%%может оказаться неприменимым. Например, в графе с тремя узлами,
%%$s$, $t$ и $u$, если $C[s,t]=2$, $C[s,u]=3$ и $C[u,t]=-2$,
%%алгоритм не найдет кратчайшего пути длиной 1 и
%%выдаст путь длиной 2.
%%\end{remark}
%%\begin{ground}
%%Для доказательства корректности алгоритма Дейкстры достаточно
%%заметить, что при каждом выполнении\vadjust{\penalty-1000} тела цикла, начинающегося
%%меткой $M$, в качестве $v$ используется вершина, для которой
%%известен кратчайший путь из вершины $s$. Другими словами,
%%если $X[v]=1$, то $T[v]=d(s,v)$, и все вершины на пути
%%$\lrle{s,v}$,
%%определяемом вектором $H$, обладают тем же свойством,
%%то~есть
%%$$
%%\A{u}{X[u]=1\Impl T[u]=d(s,u)\And X[H[u]]=1}.
%%$$
%%\end{ground}
%%%\vspace*{-.25\baselineskip}
%%\end{frame}
%%
%%
%%\begin{frame}
%%\frametitle{8.6.3. Алгоритм Дейкстры (4/4)}
%%\begin{ground} Продолжение.
%%Действительно (по индукции), первый раз в качестве $v$
%%используется вершина~$s$, для которой кратчайший путь пустой и
%%имеет длину 0 (непустые пути не могут быть короче, потому что
%%длины дуг неотрицательны). Пусть $T[u]=d(s,u)$ для всех ранее
%%помеченных вершин $u$. Рассмотрим вновь помеченную вершину~$v$,
%%которая выбрана из условия $T[v]=\min\limits_{X[u]=0}T[u]$.
%%Заметим, что если известен путь, проходящий через помеченные
%%вершины, то тем самым известен кратчайший путь. Допустим (от
%%противного), что $T[v]>d(s,v)$, то~есть найденный путь, ведущий из
%%$s$ в $v$, не является кратчайшим. Тогда на этом пути должны быть
%%непомеченные вершины. Рассмотрим первую вершину $w$ на этом пути,
%%такую, что $X[w]=0$. Имеем $T[w]=d(s,w)\le d(s,v)<T[v]$, что
%%противоречит выбору вершины $v$.
%%\end{ground}
%%
%%%\vspace*{-\baselineskip}
%%
%%\begin{remark}
%%Известно, что применение алгоритма Флойда в среднем примерно вдвое
%%менее трудоёмко, чем применение алгоритма Дейкстры для всех пар
%%вершин.
%%\end{remark}
%%
%%%\addvspace{18pt}
%%\end{frame}
%%
%%%\begin{frame}
%%%\label{8.6.4.}\subsubsection{8.6.4. Алгоритм Беллмана-Форда }\frametitle{8.6.4. Алгоритм Беллмана-Форда (1/5)}
%%%Перед демонстрацией самого алгоритма необходимо обратить внимание на такое понятие, как динамическое программирование.
%%% 
%%%Динамическое программирование~--- способ решения сложных задач путём разбиения их на более простые подзадачи. Он применим к задачам, выглядящим как набор перекрывающихся подзадач, сложность которых чуть меньше исходной. В этом случае время вычислений, по сравнению с <<наивными>> методами, можно значительно сократить.
%%%Таким образом идею динамического программирования можно выразить простым алгоритмом:
%%%
%%%1.	Разбиение задачи на подзадачи меньшего размера $n$.
%%%
%%%2.	Нахождение оптимального решения подзадач рекурсивно, проделывая такой же трёхшаговый алгоритм.
%%%
%%%3.	Использование полученного решения подзадач для конструирования решения исходной задачи.
%%%
%%%Показательный пример такого подхода являет собой алгоритм Беллмана — Форда. 
%%%Алгоритм Беллмана — Форда находит кратчайший путь от данной вершины $s$ до всех остальных вершин в ориентированном взвешенном графе $G$ с $V$ вершинами и $E$ рёбрами.
%%%\end{frame}
%%%
%%%\begin{frame}
%%%\frametitle{8.6.4. Алгоритм Беллмана-Форда (2/5)}
%%%В отличие от алгоритма Дейкстры из предыдущего пункта, 
%%%этот алгоритм применим также и к графам, содержащим рёбра отрицательного веса.
%%%
%%%\begin{remark}
%%%Если граф $G$ содержит отрицательный цикл, 
%%%то кратчайшего пути до некоторых вершин может не существовать.
%%%\end{remark}
%%%
%%%Общая идея. Cчитаем, что граф не содержит цикла отрицательного веса. Заведём массив расстояний от данной точки s до всех остальных $d[0...V-1]$, который после отработки алгоритма будет содержать кратчайшие пути от данной вершины s до всех остальных вершин. В начале работы мы заполняем его следующим образом: $d[s]=0$, а все остальные элементы $d[...]=inf$.
%%%Алгоритм состоит из нескольких фаз. На каждой фазе просматриваются все рёбра графа, и алгоритм пытается произвести релаксацию вдоль каждого ребра $(a,b)$ стоимости $w$. Релаксация вдоль ребра - это попытка улучшить значение $d[b]$ значением $d[a]+w$. Достаточно $V-1$ фазы алгоритма, чтобы посчитать длины всех кратчайших путей в графе. Для недостижимых вершин расстояние $d[...]=inf$.
%%%\end{frame}
%%%
%%%\begin{frame}
%%%\begin{algorithmic}
%%%\REQUIRE
%%%орграф $G(V,E)$, заданный массивом структур.\\*
%%%$s$~--- вершины графа.\\*
%%%\ENSURE возрастает $false$, если есть цикл отрицательного веса, $true$ в противном случае.
%%%\FOR{$v\in G(v)$}
%%%\STATE $d[v]\Assign inf;$
%%%\ENDFOR
%%%\STATE $d[s]\Assign 0;$
%%%\FOR{\textbf{$i$ from $1$ to $V-1$}}
%%%	\FOR{$(u,v)\in E$}
%%%		\IF{$d[v]>d[u]+\omega(u, v)$}
%%%			\STATE $d[v]\Assign d[u]+\omega(u, v);$
%%%		\ENDIF
%%%	\ENDFOR
%%%\ENDFOR
%%%\FOR{$(u,v)\in E$}
%%%	\IF{$d[v]>d[u]+\omega(u, v)$}
%%%		\STATE $return$ $false;$
%%%	\ENDIF
%%%\ENDFOR
%%%\STATE $return$ $true;$
%%%\end{algorithmic}
%%%\end{frame}
%%%
%%%\begin{frame}
%%%\begin{theorem}
%%%Пусть $G(V,E)$ — взвешенный орграф, $s$ — стартовая вершина. 
%%%Тогда после завершения $V-1$ итераций цикла для всех вершин, 
%%%достижимых из $s$, выполняется равенство
%%%$d[v]=\delta(s,v)$, где $d[v]$ - оценка веса кратчайшего 
%%%пути из вершины $s$ в каждую вершину  $v\in V$. $\delta(s,v)$ - точный вес кратчайшего пути из $s$ в вершину $v$.%
%%%\end{theorem}
%%%
%%%\begin{proof}
%%%Рассмотрим произвольную вершину $v$, достижимую из $s$.
%%%Пусть $p=\langle v_0...v_k \rangle$, 
%%%где $v_0 = s$, $v_k=v$ — кратчайший ациклический путь из $s$ в $v$.
%%%Путь $p$ содержит не более $V-1$ рёбер. Поэтому $k\leq V-1$.
%%%
%%%Докажем следующее утверждение: 
%%%после $n(n \leq k)$ итераций первого цикла алгоритма,  $d[v_n]=\delta(s,v_n)$.
%%%Воспользуемся индукцией по $n$:
%%%
%%%База индукции. $n = 1.$ $d[v_0]=d[s]=\delta(s,s)=0$.
%%%  
%%%Индукционный переход. Пусть после $n(n \leq k)$ итераций, верно что $d[v_n]=\delta(s,v_n)$.
%%%Так как $(v_n,v_{n+1})$ принадлежит кратчайшему пути от $s$ до $v$, то $\delta(s,v_{n+1})=\delta(s,v_n)+\omega(v_n,v_{n+1})$.
%%%Во время $l+1$ итерации релаксируется ребро $(v_n,v_{n+1})$, следовательно, по завершении итерации будет выполнено $d[v_{n+1}] \leq d[v_n] + \omega(v_n,v_{n+1}) = \delta(s,v_n)+\omega(v_n,v_{n+1})=\delta(s,v_{n+1})$.
%%%Ясно, что $d[v_n+1] \ge \delta(s,v_{n+1})$, поэтому верно, что после $l+1$ итерации $d[v_{n+1}]=\delta(s,v_{n+1})$. Индукционный переход доказан. 
%%%Итак, выполнены равенства $d[v]=d[v_k]=\delta(s,v_k)=\delta(s,v)$.
%%%\end{proof}
%%%\end{frame}
%%%
%%%\begin{frame}
%%%\frametitle{8.6.4. Алгоритм Беллмана-Форда (5/5)}
%%%Нахождение отрицательного цикла. %выделить курсивом
%%%При наличии отрицательного цикла, расстояние до всех вершин на этом цикле, а также расстояния до достижимых из этого цикла вершин не определены — они должны быть равны минус бесконечности.
%%%Алгоритм Форда-Беллмана сможет бесконечно делать релаксации среди всех вершин этого цикла и вершин, достижимых из него. Значит, если не ограничивать число фаз числом $V-1$, то алгоритм будет работать бесконечно, постоянно улучшая расстояния до этих вершин.
%%%Отсюда получаем критерий наличия достижимого цикла отрицательного веса: если после $V-1$ фазы выполним ещё одну фазу, и на ней произойдёт хотя бы одна релаксация, то граф содержит цикл отрицательного веса, достижимый из s; в противном случае, такого цикла нет.
%%%\begin{remark}
%%%Ищется отрицательный цикл, достижимый из некоторой стартовой вершины $s$. Однако алгоритм можно модифицировать, чтобы он искал просто любой отрицательный цикл в графе. Для этого надо положить все расстояния равными нулю, а не бесконечности, так как будто бы мы ищем кратчайший путь изо всех вершин одновременно; на корректность обнаружения отрицательного цикла это не повлияет.
%%%\end{remark}
%%%%\begin{remark}
%%%%Алгоритм Форда-Беллмана можно модифицировать таким образом, чтобы он выводил сам отрицательный %цикл в виде последовательности вершин, входящих в него, если такой имеется.
%%%%\end{remark}
%%%\end{frame}
%%%
%%%
%%
%%\begin{frame}
%%\label{8.6.4.}\subsubsection{8.6.4. Дерево кратчайших путей }\frametitle{8.6.4. Дерево кратчайших путей (1/3)}
%%
%%Алгоритм Дейкстры можно использовать и в том случае, когда нужно найти
%%кратчайшие цепи (пути) от вершины $s$ до всех остальных вершин.
%%
%%\begin{remark}
%%В предыдущем параграфе мы отметили,
%%что алгоритм Дейкстры применим как к графам,
%%так и к орграфам в равной мере.
%%Действительно, в этом алгоритме достаточно знать,
%%какие вершины (узлы) являются соседними и какова цена
%%(длина) перехода. Для графа матрица длин $C$ и отношение
%%смежности $\Gamma$ симметричны, для орграфа это не так,
%%но для работы алгоритма Дейкстры данное обстоятельство
%%не имеет значения.
%%Таким образом, алгоритм Дейкстры, равно как и
%%предшествующий, применим как к графам, так и к орграфам.
%%Именно поэтому в этом разделе мы позволяем себе вольности,
%%используя в одном контексте термины, относящиеся к графам (<<вершина>>)
%%и к орграфам (<<путь>>).
%%\end{remark}
%%
%%Если достаточно знать только длины путей и все вершины достижимы из $s$,
%%а~длины дуг положительны,
%%то алгоритм можно записать в следующей компактной и~наглядной форме.
%%
%%%\addvspace{6pt}
%%\end{frame}
%%
%%\begin{frame}
%%%\frametitle{8.6.4. Дерево кратчайших путей (2/3)}
%%%\begin{algorithm}
%%%\caption{Алгоритм Дейкстры для вычисления дерева кратчайших путей}
%%%\index{Алгоритм (algorithm)!Дейкстры (Edsger Wybe Dijkstra)}
%%%\label{Dijkstra2}
%%%\vspace{-\baselineskip}
%%\begin{algorithmic}
%%\REQUIRE
%%орграф $G(V,E)$, заданный матрицей длин дуг
%%\textbf{$C:$ array $[1..p,1..p]$ of real} и списками смежности $\Gamma$;
%%$s$~--- исходная вершина графа.
%%\ENSURE
%%вектор \textbf{$T:$ array $[1..p]$ of real}
%%длин кратчайших путей от $s$.
%%\FOR{$u\in V$}
%%  \STATE $T[u]\Assign C[s,u]$ \COMMENT{ начальное приближение определяется матрицей }
%%  \STATE $X[u]\Assign 0$ \COMMENT{ все вершины не отмечены }
%%\ENDFOR
%%\FOR{\textbf{$i$ from $1$ to $p$}}
%%  \STATE $m\Assign\infty$ \COMMENT{ поиск конца нового кратчайшего пути }
%%  \FOR{$u\in V$}
%%    \IF{$X[u]=0\And T[u]<m$}
%%       \STATE $v\Assign u$,  $m\Assign T[u]$
%%       \COMMENT{ вершина $u$ заканчивает путь из $s$~}
%%    \ENDIF
%%  \ENDFOR
%%  \FOR{$u\in \Gamma(v)$}
%%       \STATE $T[u]\Assign \min(T[u],T[v]+C[v,u])$
%%       \COMMENT{ пересчет оценки длины пути }
%%  \ENDFOR
%%  \STATE $X[v]\Assign  1$ \COMMENT{ найден кратчайший путь из $s$ в $v$}
%%\ENDFOR
%%\end{algorithmic}
%%%\end{algorithm}
%%
%%\end{frame}
%%
%%\begin{frame}
%%\frametitle{8.6.4. Дерево кратчайших путей (3/3)}
%%\begin{ground}
%%Приведённый вариант алгоритма основан на той же идее, что и~предшествующий: на очередном шаге нужно выбрать вершину $v$, до
%%которой мы точно знаем кратчайший путь от $s$, и пересчитать
%%оценки путей для вершин, смежных с $v$. Такой вершиной на первом
%%шаге окажется вершина $s$ (для нее $T[s] =0$). На следующих
%%шагах новой вершиной $v$ является вершина, которая еще не
%%рассмотрена ($X[u]=0$) и для которой оценка пути минимальна.
%%Действительно, даже если в будущем мы обнаружим другой путь,
%%ведущий в $v$, его длина может быть разве что больше.
%%\end{ground}
%%
%%%\vspace*{-.5\baselineskip}
%%
%%Таким образом, грубая оценка трудоёмкости
%%алгоритма Дейкстры составляет $O(p^2)$, поскольку вложенные циклы
%%выполняются $O(p)$ раз.
%%Заметим,\vadjust{\penalty-1000} что во внутреннем цикле
%%явно совершается излишняя работа: вершины, далекие от $s$,
%%не могут повлиять на выбор $v$ на ранних шагах, и их не стоит
%%рассматривать раньше времени. Известно,
%%что, более тщательно подбирая структуры данных для представления
%%графа, можно уменьшить трудоёмкость алгоритма Дейкстры до $O(q\log_2 p)$.
%%
%%
%%%\begin{remark}
%%%Последнее замечание этого параграфа
%%%относится к его названию.
%%%Совокупность путей от вершины $s$ к другим вершинам образует
%%%корневое дерево в смысле п.~\ref{s921}.
%%%Действительно, никакой кратчайший путь,
%%%очевидно, не содержит контуров (напомним, что длины дуг положительны).
%%%Всякий узел
%%%достижим из корня по построению. Более того, этот путь единственный,
%%%поскольку, даже если в вершину ведет несколько кратчайших путей одинаковой длины,
%%%алгоритм выберет из них только один. Отсюда и из определений
%%%п.~\ref{s921} немедленно
%%%следует, что традиционное название,
%%%использованное для данного параграфа, оправданно.
%%%\end{remark}
%%
%%\end{frame}
%%
%%
%%\begin{frame}
%%\label{8.6.5.}\subsubsection{8.6.5. Кратчайшие пути в бесконтурном орграфе }\frametitle{8.6.5. Кратчайшие пути в бесконтурном орграфе (1/2)} 
%%%\label{8.7.5}
%%
%%Алгоритм Дейкстры применим, если длины дуг неотрицательны.
%%Существует эффективный алгоритм, который позволяет найти
%%дерево кратчайших путей в~орграфе,
%%даже если длины дуг могут быть отрицательными,
%%но известно, что орграф не содержит контуров.
%%
%%%\pagebreak[3] {%\renewcommand{\baselinestretch}{1.03}\selectfont
%%\begin{lemma}
%%В произвольном бесконтурном орграфе $G(V,E)$ узлы
%%можно перенумеровать так, что $\A{(v_i,v_j)\in E}{i<j}$.
%%\end{lemma}
%%\begin{proof}
%%По теореме п.~7.5.2 отношение достижимости в бесконтурном
%%графе  есть строгое частичное упорядочение на конечном множестве.
%%Применяя алгоритм топологической сортировки, получаем
%%требуемую нумерацию узлов.
%%\end{proof}
%%
%%
%%Допустим теперь, что узлы в бесконтурном орграфе перенумерованы
%%так, что каждая дуга ведёт из узла с меньшим номером в узел с
%%б$\acute{\text{о}}$льшим номером, а~источник этого орграфа, из
%%которого нужно построить дерево кратчайших путей, имеет номер 1.
%%Тогда следующий алгоритм находит кратчайшие пути от узла 1 до
%%всех достижимых узлов.
%%\end{frame}
%%
%%
%%\begin{frame}
%%\frametitle{8.6.5. Кратчайшие пути в бесконтурном орграфе (2/2)} 
%%%\begin{algorithm}
%%%\caption{Определение расстояний от источника в бесконтурном графе}
%%%\label{a8.5} \index{Алгоритм (algorithm)!определения расстояний от
%%%источника (distances from source)}
%%\begin{algorithmic}
%%\REQUIRE
%%орграф $G(V,E)$, заданный матрицей длин дуг
%%\textbf{$C:$ array $[1..p,1..p]$ of real}
%%и~списками предшествующих узлов $\Gamma^{-1}$;
%%источник = 1.
%%\ENSURE
%%вектор \textbf{$T:$ array $[1..p]$ of real}
%%длин кратчайших путей от источника.
%%\FOR{\textbf{$i$ from $1$ to $p$}}
%%  \STATE $T[i]\Assign C[1,i]$
%%  \COMMENT{ начальное приближение определяется матрицей }
%%\ENDFOR
%%\FOR{\textbf{$i$ from $2$ to $p$}}
%%  \FOR{$j\in \Gamma^{-1}(i)$}
%%       \STATE $T[i]\Assign \min(T[i],T[j]+C[j,i])$
%%       \COMMENT{ пересчёт оценки длины пути }
%%  \ENDFOR
%%\ENDFOR
%%\end{algorithmic}
%%%\end{algorithm}
%%
%%%\begin{ground}
%%%Докажем индукцией по $i$, что основной цикл имеет инвариант
%%%$\A{j\in 1..i}{T[j]=d(1,j)}$. База: если $j=1$, то
%%%$T[1]=0=d(1,1)$. Пусть теперь $\A{j<i}{T[j]=d(1,j)}$. В кратчайшем
%%%пути $\langle 1,i\rangle$ все промежуточные узлы имеют номера
%%%больше 1 и~меньше $i$ по построению орграфа, в частности, вершина
%%%$j$, предшествующая $i$ на кратчайшем пути, будет выбрана во
%%%внутреннем цикле, для неё по индукционному предположению
%%%$T[j]=d(1,j)$ и, значит, $T[i]=d(1,i).$\hfill\qed
%%%\renewcommand{\qed}{\vspace{-\baselineskip}}
%%%\end{ground}
%%
%%\begin{remark}
%%В алгоритме используется множество
%%<<предшествования>>
%%$$
%%\Gamma^{-1}(v)\Def \Set{u}{v\in\Gamma(u)},
%%$$
%%которое совпадает со списками смежности $\Gamma(v)$ для графов
%%и не пересекается с $\Gamma(v)$ для направленных орграфов.
%%\end{remark}
%%\end{frame}







% ===== tex/DM2022_9.tex =====

\section{9. Деревья}

\begin{frame}
\frametitle{9. Деревья}

Деревья заслуживают отдельного и подробного рассмотрения по двум
причинам.
\begin{itemize}
\item Деревья являются в некотором смысле простейшим классом
графов. Для них выполняются многие интересные утверждения, которые
не всегда выполняются для графов в общем случае. Применительно к
деревьям многие доказательства и рассуждения оказываются намного
проще. Выдвигая какие-то гипотезы при решении задач теории графов,
целесообразно сначала их проверять на деревьях. 
\item Деревья
являются самым распространённым классом графов, 
применяемых в~программировании, причём в самых разных ситуациях. 
Более половины
объёма этой главы посвящено рассмотрению конкретных применений
деревьев в~программировании.
\end{itemize}

\end{frame}


\subsection{9.1. Свободные деревья}
\label{9.1.}

\begin{frame}
\frametitle{9.1. Свободные деревья}
Изучение деревьев целесообразно начать
с самых общих определений и установления
основных свойств.

\medskip
\begin{tabular}{p{9mm}|p{50mm}|p{84mm}}
  № & Название параграфа
  & Ключевые термины и обозначения \\
  \hline
\DMref{9.1.1.}{9.1.1.} & Свободные деревья и их элементы 
  & 
  Лес, ацикличность, cвободное дерево, древочисленность, субцикличность, хорда дерева \\
\DMref{9.1.2.}{9.1.2.}\RS & Основные свойства свободных деревьев 
  & 
  \\
\DMref{9.1.3.}{9.1.3.} & Центр свободного дерева &
\\
\end{tabular}

\end{frame}

\begin{frame}
\label{9.1.1.}\subsubsection{9.1.1. Свободные деревья и их элементы }\frametitle{9.1.1. Свободные деревья и их элементы (1/2)}


Граф без циклов называется \odmkey{ациклическим}{Граф
(graph)!ациклический (acyclic)}, или \odmkey{лесом}{Лес (forest)}.
Связный ациклический граф называется (\odmkey{свободным}{Дерево
(tree)!свободное (free)}) \odmkey{деревом}{Дерево (tree)}. Таким
образом, компонентами связности леса являются деревья.


\begin{remark}
Слово <<свободное>> употребляется
в том случае, когда нужно подчеркнуть
отличие деревьев от других объектов, родственных деревьям:
ориентированных деревьев, упорядоченных деревьев и т.~д.
\end{remark}

\smallskip
В связном графе $G$ выполняется неравенство $q(G) \ge p(G)-1$
(\DMref{п.~8.1.5}{8.1.5.}). Граф $G$ (не обязательно связный!), в котором
$q(G) = p(G)-1$, называется \odmkey{древочисленным}{Граф
(graph)!древочисленный (tree-numerical)}.
В ациклическом графе $G$ \ $z(G) = 0$. Пусть $u$, $v$~---
несмежные вершины графа $G$, $x = (u,v) \not\in E$. Если граф $G +
x$ имеет ровно один простой цикл, $z(G+x) = 1$, то граф $G$ (не
обязательно ациклический!) называется \odmkey{субциклическим}{Граф
(graph)!субциклический (subcyclic)}.
%\pagebreak[3]


\begin{remark}
Графы $K_3\cup K_1$ и $K_3\cup K_2$ являются древочисленными и
субциклическими и в то же время не являются ни связными, ни
ациклическими.
\end{remark}

\smallskip
Ребро, соединяющее несмежные вершины свободного дерева, 
называется \odmkey{хордой}{Хорда (chord)}. Хорда дерева не принадлежит дереву!

\end{frame}



\begin{frame}
\frametitle{9.1.1. Свободные деревья и их элементы (2/2)}
\begin{example}
На рисунках последовательно показаны диаграммы \em{всех} различных
свободных деревьев с 4, 5 и 6 вершинами.
\begin{figure}[h]
\begin{center}
\includegraphics[height=8mm]{9_0.jpg}
\end{center}
%\caption{Свободные деревья с 4 вершинами} 
\end{figure}
\begin{figure}[h]
\begin{center}
\includegraphics[height=20mm]{odm09-01_new}
\end{center}
%\caption{Свободные деревья с 5 вершинами} \label{f0901}
\end{figure}
\begin{figure}[h]
\begin{center}
\includegraphics{odm09-02}
\end{center}
%\caption{Свободные деревья с 6 вершинами}
\label{f0902}
\end{figure}
\end{example}
\end{frame}


\begin{frame}
\label{9.1.2.}\subsubsection{9.1.2. Основные свойства свободных  деревьев }\frametitle{9.1.2. Основные свойства свободных деревьев (1/8)}

%Следующая теорема устанавливает,
%что любыедва из четырех свойств~--- связность, ацикличность, древочисленность и
%субцикличность~--- характеризуют граф как дерево.
%{\small
\begin{theorem}
Пусть $G(V,E)$~{\upshape---} граф с $p$ вершинами, $q$  ребрами,
$k$ компонентами связности и $z$ простыми циклами.  Пусть далее
$x$~{\upshape---} ребро, соединяющее любую пару несмежных  вершин
в $G$. Тогда следующие утверждения эквивалентны.
\end{theorem}


\begin{enumerate}
%\itemindent=0pt \listparindent=0pt
%\parsep=0pt plus 1pt
%\itemsep=\parsep
{\leftskip-6pt
\item $G$~--- дерево, то~есть связный граф без
циклов,
$k(G)=1\And z(G)=0.$\\[-2pt]
\item Любые две вершины соединены в $G$ единственной простой
цепью,
$\A{u,v}{\left\vert P(u,v)\right\vert\!=\!1}.$\\[-2pt]
\item $G$~--- связный граф, и любое ребро есть мост,
$k(G)=1 \And \A{e\in E}{k(G-e)>1}.$\\[-2pt]
\item $G$~--- связный и древочисленный,
$k(G)=1 \And q(G)=p(G)-1.$\\[-2pt]
\item $G$~--- ациклический и древочисленный,
$z(G)=0 \And q(G)=p(G)-1.$\\[-2pt]
\item $G$~--- древочисленный и субциклический (за двумя
исключениями),
$q(G)=p(G)-1$, $G\not= {K_1}\cup K_3$, 
$G \not= {K_2}\cup K_3$, $z(G+x)=1.$\\[-2pt]
\item $G$~--- ациклический и субциклический,
$z(G)=0 \And z(G+x)=1.$\\[-2pt]
\item $G$~--- связный, субциклический и неполный,
$k(G)=1$, $G\not= {K_p}$, $p\geq 3$, $z(G+x)=1.$ \\
}
\end{enumerate}
%\addvspace{-6pt}
%}
\end{frame}


\begin{frame}
\frametitle{9.1.2. Основные свойства свободных деревьев (2/8)}
\begin{proof}
\proofsection{1$\Impl$2} 
От противного. Пусть существуют две цепи
$\lrle{u,v}$. Некото\-рые вершины этих цепей различны. Обозначим
через $w_1$ первую вершину при перечислении вершин от $u$ к $v$,
такую, что следующие вершины в цепях различны, а через $w_2$
обозначим первую вершину при перечислении вершин от $v$ к $u$,
такую, что следующие вершины в цепях различны (рисунок слева). 
Рассмотрим отрезок $\lrle{w_1,w_2}$ первой цепи при
перечислении вершин от $u$ к $v$ и отрезок $\lrle{w_2,w_1}$ второй
цепи при перечислении вершин от $v$ к $u$. \\Тогда $\lrle{w_1,w_2} +\lrle{w_2,w_1}$~--\- цикл, что противоречит ацикличности графа
$G$.

\vspace{-12pt}
\begin{figure}[h]
\begin{center}
\includegraphics[width=90mm]{9_1_2_1.jpg}
\end{center}
\end{figure}

\vspace{-2pt}
\proofsection{2$\Impl$3} Любые две вершины соединены цепью
(единственной), следовательно, \\ $k(G) = 1$. Далее от противного.
Пусть ребро $x$~--- не мост. Тогда в $G-x$ концы этого ребра
связаны цепью. Само ребро $x$~--- вторая цепь.
\end{proof}

\end{frame}


\begin{frame}
\frametitle{9.1.2. Основные свойства свободных деревьев (3/8)}
\begin{proof} (Продолжение 1)
\proofsection{3$\Impl$4} Индукция по $p$. База: $p=1\Impl q=0$.
Пусть $q(G)=p(G)-1$ для всех связных графов $G$ с числом вершин
меньше $p$, у которых любое ребро является мостом. Тогда удалим из
графа $G$ некоторое ребро $x$ (которое является мостом). Получим
две компоненты, $G'$  и $G''$, удовлетворяющие индукционному
предположению. Имеем \\
$q'=p'-1$, $q''=p''-1$, $q=q'+q''+1=p'-1+p''-1+1=p-1$.
\proofsection{4$\Impl$5} От противного. Пусть есть цикл с $n$
вершинами и $n$ рёбрами. Остальные $p-n$ вершин имеют инцидентные
им рёбра, которые связывают их с циклом. Следовательно, $q\geq p$,
что противоречит условию $q=p-1$.
\proofsection{5$\Impl$6} Граф $G$~--- ациклический, следовательно,
$G\not= K_2\cup K_3$, $G\not= K_1\cup K_3$. Далее от противного.
Пусть $z(G+x)\ne 1$. Если $z(G+x)\ge 2$, 
то $x\in Z_1$, $x\in Z_2$, 
где $Z_1$ и $Z_2$~--- вновь образованные циклы
в графе $G+x$. Но тогда $Z_1\cup Z_2 \setminus \{x\}$~---
цикл в графе $G$, что противоречит ацикличности графа $G$. 
\end{proof}

\end{frame}


\begin{frame}
\frametitle{9.1.2. Основные свойства свободных деревьев (4/8)}
\begin{proof} (Продолжение 2)
Если же $z(G+x)=0$, то удалим из графа $G+x$ ребро $x$. 
Получим две компоненты, $G_1$  и $G_2$. Графы $G_1$ и 
$G_2$~--- ациклические. Убедимся, что $G_1$ и $G_2$ являются связными. Если не связны они оба, то $q(G_1)<p(G_1)-1$ и $q(G_2)<p(G_2)-1$, то есть $$q(G)=q(G_1)+q(G_2)<p(G_1)-1+p(G_2)-1=p(G)-2,$$ что противоречит древочисленности $G$. Если несвязным является только один из них, пусть $G_1$, то $G_2$ связный ($q(G_2)\geq p(G_2)-1$) и ациклический ($q(G_2)<p(G_2)$), что даёт  древочисленность. 
Тогда $p(G_1)=p(G)-p(G_2)=q(G)+1-(q(G_2)+1)=q(G_1)$, что противоречит условию ацикличности $G_1$ 
(\DMref{п.~7.3.5}{7.3.5.}). То есть $G_1$ и $G_2$~--- связные и ациклические, а значит древочисленные.
Имеем $q(G)=q(G_1)+q(G_2)=p(G_1)-1+p(G_2)-1=p(G)-2$ и получаем
$q(G)=p(G)-2$, что противоречит условию древочисленности $G$.    
\end{proof}


\end{frame}


\begin{frame}
\frametitle{9.1.2. Основные свойства свободных деревьев (5/8)}
\begin{proof} (Продолжение 3)
\proofsection{6$\Impl$7} От противного. Пусть в $G$ есть
цикл $Z=C_n$. Если $n>3$, то рассмотрим пары вершин 
цикла, которые не соединены рёбрами цикла. Если среди них уже есть смежные
вершины, имеем три цикла.
Если среди них нет смежных вершин, то,
соединив несмежные вершины в $Z$, получим три цикла.
Следовательно, в цикле нет вершин, не соединённых рёбрами цикла, значит  $Z=C_3$. Этот цикл $Z$ явля\-ется компонентой
связности $G$. Действительно, пусть это не так. Тогда существует
вершина $w$, смежная с $Z$. Если $w$ смежна более чем с одной
вершиной $Z$, то имеем больше одного цикла. Если $w$ смежна только
с одной вершиной $Z$, то, соединив её с другой вершиной, получим
два цикла
(см. рис. на слайде <<9.1.2. Основные свойства деревьев
(2/8)>>). Рассмотрим $G'\Assign G-Z$. Имеем $p=p'+3$, $q=q'+3$.
Но $q=p-1$, значит, $q'=p'-1$. Отсюда $z(G')=0$, так как
один цикл уже есть. Следовательно, компоненты $G'$~--- деревья.
Пусть их $k$. Имеем
$
q'=\sum_{i=1}^{k}q_i=\sum_{i=1}^{k}(p_i-1)=
\sum_{i=1}^{k}p_i-k=p'-k,
$
но $q'=p'-1$, следовательно, $k=1$, то~есть дерево одно. 
Если в этом
дереве соединить несмежные вершины, то получим второй цикл. Два
исключения:
деревья, которые не имеют несмежных вершин,~--- это $K_1$~и~$K_2$.

\end{proof}
\end{frame}


\begin{frame}
\frametitle{9.1.2. Основные свойства свободных деревьев (6/8)}
\begin{proof} (Продолжение 4)

\proofsection{7$\Impl$8} При $p\geq 3$ граф $K_p$ содержит цикл,
следовательно, $G\not=K_p$. Далее от противного. Пусть $G$
несвязен, тогда при соединении ребром двух компонент связности
цикл не возникнет, что противоречит субцикличности. 

\proofsection{8$\Impl$1} От противного. Пусть в графе $G$ есть 
единственный цикл (большее количество циклов противоречит
субцикличности графа $G$). Если в $G$ нет висячих вершин и $p=3$,
то $G=K_3$, что противоречит условию неполноты $G$. Если в связном $G$ нет висячих вершин, $p>3$ и есть единственный цикл,  то $G=C_p$.
Соединим любые
две несмежные вершины $u$ и $v$ ребром $x$ (рисунок слева). 

\begin{figure}[h]
\begin{center}
\includegraphics[width=140mm]{9_1_2_2.jpg}
\end{center}
\end{figure}

\vspace{-12pt} 
\end{proof}



\end{frame}


\begin{frame}
\frametitle{9.1.2. Основные свойства свободных деревьев (7/8)}
\begin{proof} (Окончание)

Имеем три цикла в $G+x$,
что противоречит условию субцикличности. Если в $G$ есть висячая
вершина $u$, соединим её c любой несмежной с ней вершиной $v$
ребром $x$. Вершина $v$ может принадлежать (рисунок в центре)
или не принадлежать циклу (рисунок справа). 
Теперь из $u$ в $v$ есть две или три простые цепи (одна
или две цепи были раньше, так как $G$~--- связный с одним простым
циклом, и появилась цепь из ребра $x$). Имеем два, три или более
циклов в графе $G+x$, что противоречит условию субцикличности.
\end{proof}

\begin{corollary} {\NB{[1]}}
В любом нетривиальном дереве имеются  по  крайней  мере
две висячие вершины.
\end{corollary}
\begin{proof}
Рассмотрим дерево $G(V,E)$. Дерево~--- связный граф,
следовательно, $\A{v_i\in V}{d(v_i)\geq 1}$.
По теореме 3 \DMref{п. 7.3.6}{7.3.6.} какие-то висячие вершины
есть.
Далее от противного. Пусть есть только одна висячая вершина
$\A{i\in 1..(p-1)}{d(v_i)>1}$. Тогда
$
2q=\sum^p_{i=1}d(v_i)\geq 2(p-1)+1=2p-1.
$
Но $q=p-1$, то~есть $2q=2p-2$.
Имеем противоречие: $2p-2\leq 2p-1$.
\end{proof}


%\setcounter{corollary}{0}
\end{frame}

\begin{frame}
\frametitle{9.1.2. Основные свойства свободных деревьев (8/8)}


\begin{remark}
Легко видеть, в частности, что висячими вершинами в дереве являются
концы любого диаметра.
\end{remark}
\begin{corollary} {\NB{[2]}}
Каждая не висячая вершина свободного
дерева является точкой сочленения.
\end{corollary}
\begin{proof}
Пусть $G(V,E)$~--- дерево, $v\in V$ и $d(v)>1$. Тогда по
определению $\E{u,w\in V}{u\ne w \And (u,v)\in E\, \And\,
{(v,w)\in E}}$. Граф $G$ связен, поэтому существует цепь $\langle
u , w\rangle$. Если $v\not\in \langle u , w\rangle$, то имеем цикл
$v,\langle u , w\rangle, v$, что противоречит тому, что $G$~---
дерево. Следовательно, $\E{u,w\in V}{\A{\langle u ,
w\rangle}{v\in\langle u , w\rangle}}$, и~по теореме \DMref{п.~8.1.2}{8.1.2.}
вершина $v$~--- точка сочленения.
\end{proof}

%%\pagebreak[3] \vspace{-\baselineskip}
%\begin{corollary} {\NB{[3]}}
%Если в связном графе нет висячих вершин,
%то в нём есть цикл.
%\end{corollary}
%%\vspace{-0.25\baselineskip}
%\begin{proof}
%От противного. Если связный граф не имеет циклов, то
%он является деревом и по следствию~1 обязан иметь
%висячие вершины.
%\end{proof}
%%\pagebreak[3] \vspace{-\baselineskip}
\end{frame}



\begin{frame}
\label{9.1.3.}\subsubsection{9.1.3. Центр свободного дерева}
\frametitle{9.1.3. Центр свободного дерева}
Свободные деревья выделяются из других графов тем, что их центр
всегда оправдывает своё название. 
%\vspace{-0.25\baselineskip}
\begin{theorem}
Центр свободного дерева состоит из одной вершины
или из двух смежных вершин:

\vspace{-12pt}
$$
(z(G)=0 \And k(G)=1) \Impl (C(G)=K_1 \Or C(G)= K_2).
$$
\end{theorem}
\vspace{-0.25\baselineskip}
\begin{proof}
Для деревьев $K_1$ и $K_2$ утверждение теоремы очевидно. Пусть
теперь $G(V,E)$~--- некоторое свободное дерево, отличное от $K_1$
и $K_2$. Рассмотрим граф $G'(V',E')$, полученный из $G$ удалением
всех висячих вершин. Заметим, что $G'$~--- дерево, так как
ацикличность и связность при удалении висячих вершин сохраняется.
Далее, если эксцентриситет $e_G(v)=d(v,u)$, то $u$~--- висячая
вершина в дереве $G$ (иначе можно было бы продолжить цепь <<за>>
вершину $u$). Поэтому $\A{v\in V'}{e_G(v)=e_{G'}(v)+1}$, и при
удалении висячих вершин эксцентриситеты оставшихся уменьшаются на
1. Следовательно, при удалении висячих вершин центр не меняется,
$C(G)=C(G')$. Но дерево $G$ конечно, и удаляя на каждом
шаге все висячие вершины, за несколько шагов придём
к $K_1$ или $K_2$.
\end{proof}

\end{frame}

\subsection{9.2. Ориентированные, упорядоченные и~бинарные деревья }
\label{9.2.}
\begin{frame}
\frametitle{9.2. Ориентированные, упорядоченные и~бинарные деревья }
Ориентированные (упорядоченные)
деревья являются абстракцией иерархических отношений, которые
очень часто встречаются как в практической жизни, так и в
математике и программировании. \em{Дерево (ориентированное) и
иерархия~--- это равнообъёмные понятия.} 

\medskip
\begin{tabular}{p{9mm}|p{50mm}|p{84mm}}
  № & Название параграфа
  & Ключевые термины и обозначения \\
  \hline
\DMref{9.2.1.}{9.2.1.} & Ориентированные деревья 
  & 
  Корень ордерева, поддерево узла, лист, крона, ветвь, уровень узла, ярус, потомок, предок, родитель, братья, генеалогические графы \\
\DMref{9.2.2.}{9.2.2.} & Эквивалентные определения ордерева
  & 
  \\
\DMref{9.2.3.}{9.2.3.}\RS  & Упорядоченные деревья & Прямой порядок обхода
\\
\DMref{9.2.4.}{9.2.4.}  & Бинарные деревья & $m$-ичное дерево
\\
\end{tabular}

\end{frame}

\begin{frame}
\label{9.2.1.}\subsubsection{9.2.1. Ориентированные деревья }\frametitle{9.2.1. Ориентированные деревья (1/5)}
%\vspace{-0.25\baselineskip}
%\label{s921} 
\odmkey{Ориентированным деревом}{Дерево
(tree)!ориентированное (directed)} (или
\odmkey{oрдеревом}{Ордерево (directed tree)}, или
\odmkey{корневым}{Дерево (tree)!корневое (rooted)} деревом)
называется  орграф со следующими свойствами:
%\vspace{-0.25\baselineskip}
\begin{enumerate}
\item[1)] существует единственный узел $r$, полустепень захода
которого равна 0, $d^{+}(r)=0$; он называется
\odmkey{корнем}{Корень ордерева (root of directed tree)} ордерева;
\item[2)] полустепень захода всех остальных узлов равна 1, 
$\A{v\in V-r}{d^{+}(v)=1}$;

\vspace{-6pt} 
\item[3)] каждый узел достижим из корня, 
$\A{v\in V-r} {\exists {\lrle{\stackrel{\rightarrow}{r,v}}}}$.
\end{enumerate}

{\renewcommand{\baselinestretch}{0.95}\selectfont
%\vspace{-0.25\baselineskip}
\begin{theorem} Ордерево обладает следующими свойствами:
%\vspace{-6pt}
\begin{enumerate}
\item[1)] ордерево древочисленно, $q=p-1$;\\[-2pt]
\item[2)] если в ордереве забыть ориентацию дуг, то  получится
свободное дерево;\\[-2pt]
\item[3)] в ордереве нет контуров;\\[-2pt]
\item[4)] для каждого узла существует единственный путь, ведущий
     в этот узел из \hbox{корня};\\[-2pt]
\item[5)] подграф, определяемый множеством  узлов,  достижимых из
      узла $v$, является ордеревом с корнем $v$
       (\odmkey{поддеревом}{Поддерево (subtree)} узла $v$);\\[-2pt]
\item[6)] если в свободном дереве любую вершину назначить корнем, то
      получится ордерево.
\end{enumerate}
\end{theorem}
}
\end{frame}

\begin{frame}
\frametitle{9.2.1. Ориентированные деревья (2/5)}
\begin{example}
На рисунках приведены диаграммы всех различных
ориентированных деревьев с~3 и 4 узлами.
\end{example}
%\pagebreak[3]

\begin{figure}[h]
\begin{center}
\includegraphics{odm09-05}
\end{center}
%\caption{Ориентированные деревья с 3 узлами}
\label{f0905}
\end{figure}
\begin{figure}[h]
\begin{center}
\includegraphics{odm09-06}
\end{center}
%\caption{Ориентированные деревья с 4 узлами}
\label{f0906}
\end{figure}
%\vspace{-11pt}
\end{frame}

\begin{frame}
\frametitle{9.2.1. Ориентированные деревья (3/5)}

\begin{proof}
\proofsection{1} Каждая дуга входит в какой-то узел. Из 
определения  имеем $\A{  v\in V-r}{d^+(v)=1}$,
где $r$~--- корень. Следовательно, $q=p-1$.
\proofsection{2} Пусть
$G$~--- ордерево, граф $G'$ получен из $G$ забыванием ориентации
рёбер, $r$~--- корень. Тогда любые две вершины в $G'$ связаны с корнем,
а значит связаны между собой и граф $G'$
связен. Таким образом, учитывая п.~4. теоремы \DMref{9.1.2}{9.1.2.}, $G'$~---
дерево. 
\proofsection{3} Следует из п. 2. 
\proofsection{4} От
противного. Если бы в $G$ существовали два пути из $r$ в $v$, то в
$G'$ имелся бы цикл. 
\proofsection{5} Пусть $G_v$~---
правильный подграф, определяемый множеством узлов, достижимых из
$v$. Тогда $d^+_{G_v}(v)=0$, иначе узел $v$ был бы достижим из
какого-то узла $v'\in G_v$ и, таким образом, в $G_v$, а значит, и
в $G$ имелся бы контур, что противоречит п.~3. Далее имеем:
$\A{v'\in G_v - v}{d^+(v')=1}$, так~как $G_v\subset G$. Все узлы
$G_v$ достижимы из $v$ по построению. По определению 
получаем, что $G_v$~--- ордерево. 
\end{proof}
\end{frame}

\begin{frame}
\frametitle{9.2.1. Ориентированные деревья (4/5)}

\begin{proof} (Окончание)
\proofsection{6} Пусть вершина
$r$ назначена корнем и дуги последовательно ориентированы <<от
корня>> обходом в глубину. Тогда $d^+(r)=0$ по построению;
$\A{v\in V -r}{d^+(v)=1}$, так~как входящая дуга появляется при
первом посещении узла; все узлы достижимы из корня, так~как обход
в глубину посещает все вершины связного графа. Таким образом, по
определению получаем ордерево.
\end{proof}

\smallskip
\begin{corollary}
Алгоритм поиска в глубину строит ордерево с корнем в начальном
узле.
\end{corollary}

\smallskip
\begin{remark}
Каждую вершину в свободном дереве с $p$ вершинами можно
назначить корнем и получить ордерево.
Некоторые из полученных ордеревьев могут оказаться изоморфными.
Следовательно, свободное дерево определяет не более
 $p$ различных ориентированных
деревьев. Таким образом, общее число различных ордеревьев с $p$ узлами
не более чем в $p$ раз превосходит общее число различных свободных
деревьев с $p$ вершинами.
\end{remark}


\end{frame}

\begin{frame}
\frametitle{9.2.1. Ориентированные деревья (5/5)}

Сток ордерева называется \odmkey{листом}{Лист ордерева
(leaf of directed tree)}. Множество листьев называется
\odmkey{кроной}{Крона ордерева (tree top)}. Путь из корня в лист
называется \odmkey{ветвью}{Ветвь ордерева (branch of directed
tree)}. Длина наибольшей ветви ордерева называется его
\odmkey{высотой}{Высота (height)!дерева (of tree)}.
\odmkey{Уровень}{Уровень узла (level of node)} узла ордерева~---
это расстояние от корня до узла. Сам корень имеет уровень $0$.
Узлы одного уровня образуют \odmkey{ярус}{Ярус (tier)!дерева (of
tree)} ордерева.


\begin{remark}
Наряду с <<растительной>> применяется ещё и <<генеалогическая>>
терминология. Узлы, достижимые из узла $u$, называются
\odmkey{потомками}{Потомок узла (uncestor node)} узла $u$
(потомки одного узла образуют дерево). Если узел $v$ является
потомком узла $u$, то узел $u$ называется \odmkey{предком}{Предок
узла (ancestor node)} узла $v$. Если в дереве существует дуга
$(u,v)$, то узел $u$ называется \odmkey{отцом}{Отец узла (node
parent)} (или \odmkey{родителем}{Родитель узла (parent node)})
узла $v$, а~узел $v$ называется \odmkey{сыном}{Сын узла (child
node)} узла $u$. Сыновья одного отца называются
\odmkey{братьями}{Брат узла (sibling node)}.
\end{remark}

\smallskip
\begin{digress}
Применительно к деревьям генеалогическая терминология неудачна,
поскольку отношение <<родительства>> не является иерархией.
Родственные отношения точнее описываются 
\odmkey{генеалогическими графами}{Граф (graph)!генеалогический (genealogical)}~--- 
 бесконтурными орграфами,
 в которых все узлы разбиты на два множества
 $M$ (мужчин) и $F$ (женщин) так, что в любой узел
 входит не более чем по одной дуге из классов 
 $M$ и $F$.  
\end{digress}

\end{frame}




\begin{frame}
\label{9.2.2.}\subsubsection{9.2.2. Эквивалентные определения ордерева}
\frametitle{9.2.2. Эквивалентные определения ордерева}
\label{s922} 
Ордерево $T$~--- это непустое конечное множество
узлов, на котором определено разбиение, обладающее следующими
свойствами:
\begin{enumerate}
\item Имеется один выделенный одноэлементный блок $\{r\}$,
называемый корнем данного ордерева.
\item Остальные узлы (исключая корень) содержатся в $k$ ($k\ge 0$)
блоках $T_1,\dots,T_k$, каждый из которых, в свою очередь,
      является ордеревом и называется \emph{поддеревом}:
\end{enumerate}


$$
T\Def\left\{\{r\},T_1,\dots,T_k\right\}.
$$

Нетрудно видеть, что данное определение эквивалентно
определению \DMref{п.~9.2.1}{9.2.1.}. Достаточно построить орграф, проводя дуги
от заданного корня ордерева $r$ к корням поддеревьев
$\Tuple{T}{1}{k}$ и далее повторяя рекурсивно этот процесс для
каждого из поддеревьев.

%------------------------------------------------------
% TO DO ФАН
%{\color{blue}
%Тема для сочинения: системы непересекающихся множеств, и представление 
%деревьев с указателями от листьев к корням. 
%Есть соображения и заготовки.}
%-------------------------------------------------------
\begin{remark}
В терминах \DMref{п.~8.5.4}{8.5.4.} можно определить 
ордерево как орграф, в котором базис состоит
из одного узла, причём путь от базиса единствен.  
\end{remark}
\end{frame}


\begin{frame}
\label{9.2.3.}\subsubsection{9.2.3. Упорядоченные деревья }\frametitle{9.2.3. Упорядоченные деревья (1/5)}
%\label{s923} 

Если относительный порядок поддеревьев
$T_1,\dots,T_k$ в эквивалентном определении ордерева фиксирован,
то ордерево называется \odmkey{упорядоченным}{Дерево
(tree)!упорядоченное (ordered)}.

%\vspace{-0.25\baselineskip}
\begin{examples}
Ориентированные и упорядоченные ориентированные деревья интенсивно
используются в программировании. 
%Несколько примеров представлены на рисунке.
\end{examples}


\addvspace{-10pt}
\begin{figure}[h]
\begin{center}
\includegraphics[width=130mm]{9_2_3.jpg}
\end{center}
%\caption{Примеры изображения деревьев в программировании}
%\label{f0907}
\end{figure}
%\pagebreak[3]

%\begin{enumerate}
%\item 

\addvspace{-16pt}         
{\color{blue}\bf 1.}
Для представления выражений языков
программирования, как правило, используются ориентированные
упорядоченные деревья. 
Пример представления выражения
 $a+b*c$ показан на рисунке {\it а}.

{\color{blue}\bf 2.}
Для представления
иерархической структуры вложенности элементов данных и (или)
операторов управления часто используется техника отступов,
показанная на 
рисунке {\it в}.


\end{frame}


\begin{frame}
\frametitle{9.2.3. Упорядоченные деревья (2/5)}
%\item 

{\color{blue}\bf 3.}
Для представления блочной структуры программы и связанной с
ней  структуры  областей  определения  идентификаторов  часто 
используется  ориентированное  дерево  (может  быть, неупорядоченное,
так как  порядок  определения  переменных  в  блоке  в  большинстве
языков программирования считается несущественным). 
На
рисунке {\it б} показана структура областей определения
идентификаторов $a,b,c,d,e$, причём  для  отображения  иерархии
использованы вложенные области. 



\smallskip 
%\item 
{\color{blue}\bf 4.}
Структура вложенности
 каталогов и файлов в современных операционных системах
является упорядоченным ориентированным деревом. Обычно для
изображения таких деревьев применяется способ, показанный на
рисунке {\it г}. 


{\color{blue}\bf 5.}
Различные <<правильные скобочные
структуры>> 
(например, $(a(b)(c(d)(e)))$) являются ориентированными
упорядоченными деревьями.
\begin{digress}
Тот факт, что большинство систем управления файлами использует
ориентированные деревья, отражается даже в терминологии, например:
<<корневой каталог диска>>.
\end{digress}
\end{frame}


\begin{frame}
\frametitle{9.2.3. Упорядоченные деревья (3/5)}




\begin{remark}
Общепринятой практикой при изображении деревьев является
соглашение о том, что корень находится наверху и все дуги
ориентированы сверху вниз, поэтому стрелки можно не изображать.
Диаграммы свободных, ориентированных и
упорядоченных деревьев оказываются графически неразличимыми, и
требуется дополнительное уточнение, дерево какого класса
изображено на диаграмме. 
\end{remark}


\begin{example}
На рисунке приведены три диаграммы деревьев,
которые внешне выглядят различными.
%%%%%%% Внимание!
% требуется на рис. 9.8. добавить номера (1) (2) и (3)
%Обозначим дерево слева~--- $(1)$,
%в центре~--- $(2)$ и справа~--- $(3)$.
Как упорядоченные деревья они действительно все различны:
 \textit{а}$\not= $\textit{б}, \textit{б}$\not=$\textit{в}, \textit{в}$\not=$\textit{а}.
Как ориентированные деревья \textit{а}$=$\textit{б}, но \textit{б}$\not=$\textit{в}.
Как свободные деревья они все изоморфны: \textit{а}=\textit{б}=\textit{в}.

\vspace{-6pt}
\begin{figure}[h]
\begin{center}
\includegraphics[height=30mm]{9_2_3_2.jpg}
\end{center}
%\caption{Диаграммы деревьев}
%\label{f0908}
\end{figure}

\end{example}

\end{frame}


\begin{frame}
\frametitle{9.2.3. Упорядоченные деревья (4/5)}
Указанное в \DMref{п.~9.2.2}{9.2.2.} построение позволяет определить
на упорядоченном ордереве единый линейный порядок узлов,
согласованный с уже заданными в дереве упорядочениями.
\begin{theorem}
В упорядоченном ордереве с $p$ узлами существует такая нумерация
узлов числами из диапазона $1..p$, что номера потомков больше
номеров предков и~номера старших братьев больше номеров младших
братьев.
\end{theorem}

%\pagebreak[3]%

\begin{proof}
Применим к упорядоченному ордереву алгоритм обхода в глубину, 
начиная с корня, причём узлы,
смежные с данным, посещаются в порядке, определяемом порядком
поддеревьев. Присвоим узлам номера в порядке
первого посещения. Тогда все узлы получат
уникальные номера из диапазона $1..p$, и корень получит номер 1.
Далее, старшие братья получат номера, б$\acute{\text{о}}$льшие,
чем номера младших братьев по условию обхода, а потомки получат
номера б$\acute{\text{о}}$льшие, чем номера предков, поскольку
алгоритм обхода в глубину посещает предков раньше, чем потомков.
\end{proof}

\end{frame}


\begin{frame}
\frametitle{9.2.3. Упорядоченные деревья (5/5)}
\begin{remark}
Указанный порядок обхода часто называют \odmkey{прямым}{Обход дерева (traverse of tree)!прямой (preorder)}.
\end{remark}
\begin{example}
Узлы упорядоченного ордерева на рисунке при прямом обходе
получат следующие номера:
$a\mapsto 1$,
$b\mapsto 2$,
$c\mapsto 5$,
$d\mapsto 3$,
$e\mapsto 4$,
$f\mapsto 6$,
$g\mapsto 7$,
$h\mapsto 8$,
$i\mapsto 9$.
\begin{figure}[h]
\begin{center}
\includegraphics[height=30mm]{9_2_3_3.jpg}
\end{center}
\end{figure}
\end{example}
\begin{remark}
Построенная  нумерация не
является единственной нумерацией, обладающей указанными свойствами.
\end{remark}
\begin{example}
Можно обойти упорядоченное ордерево по ярусам. При этом узлы
упорядоченного ордерева на рисунке получат
следующие номера: $a\mapsto 1$, $b\mapsto 2$, $c\mapsto 3$,
$d\mapsto 4$, $e\mapsto 5$, $f\mapsto 6$, $g\mapsto 7$, $h\mapsto
8$, $i\mapsto 9$.
\end{example}
\end{frame}

\begin{frame}
\label{9.2.4.}\subsubsection{9.2.4. Бинарные деревья }\frametitle{9.2.4. Бинарные деревья (1/2)}
\odmkey{Бинарное}{Дерево (tree)!бинарное (binary)} (или
\odmkey{двоичное}{Дерево (tree)!двоичное (binary)}) дерево~--- это
непустое конечное множество узлов, на котором определена
структура, обладающая следующими свойствами:
\begin{enumerate}\leftskip-2pt
\item[1)] имеется один выделенный узел  $r$, называемый корнем данного
бинарного дерева;
\item[2)] остальные узлы (исключая корень)
содержатся в двух непересекающихся множествах (поддеревьях)~---
левом и правом, каждое из которых, в свою очередь,
 либо пусто, либо является бинарным деревом.
\end{enumerate}
На первый взгляд может показаться, что бинарное дерево~---
это частный случай упорядоченного ориентированного дерева,
в котором у каждого узла не более двух смежных. Но это не так,
бинарное дерево \em{не является\/} упорядоченным ордеревом. Дело
в том, что даже если у некоторого узла бинарного дерева имеется только одно
непустое поддерево, то всё равно известно, какое именно это поддерево:
левое или правое.
\end{frame}

\begin{frame}
\frametitle{9.2.4. Бинарные деревья (2/2)}
\begin{example}
На рисунке приведены две диаграммы деревьев,
которые изоморфны как упорядоченные, ориентированные
и свободные деревья, но не изоморфны как
бинарные деревья.
\begin{figure}[h]
\begin{center}
\includegraphics{odm09-09}
\end{center}
%\caption{Два различных бинарных дерева}
%\label{f0909}
\end{figure}
%\pagebreak[3]

\end{example}
\begin{remark}
Понятие двоичного дерева допускает обобщение. \odmkey{$m$-ичным
деревом}{Дерево (tree)!$m$-ичное ($m$-ary)} называется непустое
конечное множество узлов, которое состоит из корня  и $m$
непересекающихся подмножеств, имеющих номера $1,\ldots,m$, каждое
из которых, в свою очередь, либо пусто, либо является $m$-ичным
деревом. Большая часть утверждений и алгоритмов для двоич\-ных
деревьев может быть сравнительно легко распространена и на
$m$-ичные деревья (с соответствующими модификациями). Поэтому, хотя
на практике $m$-ичные деревья и~используются достаточно часто,
здесь мы ограничиваемся только двоичными деревьями.
\end{remark}

\end{frame}

\begin{frame}
\frametitle{Онтология теории графов}
\begin{center}
\includegraphics[height=7cm]{ODM_9_1_2.png}
\end{center}

%{\color{red} Тема для сочинения: квадродерево. 
% Попытки есть, но слабые}
\end{frame}

\subsection{9.3. Представление деревьев в программах}
\label{9.3.}
\begin{frame}
\frametitle{9.3. Представление деревьев в программах (1/2)}
Обсуждению представлений деревьев можно
предпослать в точности те же рассуждения,
что предпосланы обсуждению представлений графов. Кроме того,
следует подчеркнуть, что задача представления
деревьев в~программе встречается гораздо чаще,
чем задача представления графов общего вида,
а~потому методы её решения оказывают ещё
большее влияние на практику программирования.

\end{frame}


\begin{frame}
\frametitle{9.3. Представление деревьев в программах (2/2)}
\medskip
\begin{tabular}{p{9mm}|p{50mm}|p{84mm}}
  № & Название параграфа
  & Ключевые термины и обозначения \\
  \hline
\DMref{9.3.1.}{9.3.1.} & Представление свободных деревьев 
  & 
  Код Прюфера, адгоритм упаковки, алгоритм распаковки \\
\DMref{9.3.2.}{9.3.2.} & Представление упорядоченных корневых деревьев
  & 
  \\
\DMref{9.2.3.}{9.2.3.} & Число упорядоченных ориентированных деревьев & Числа Каталана
\\
\DMref{9.2.4.}{9.2.4.}  & Проверка правильности скобочной структуры & 
\\
\DMref{9.2.5.}{9.2.5.}\RS  & Представление бинарных деревьев & 
Списочные структуры, упакованные массивы, размеченная степень, польская запись \\
\DMref{9.2.6.}{9.2.6.}\RS  & Обходы бинарных деревьев & 
Прямой (префиксный, левый) обход, внутренний (инфиксный, симметричный) обход, концевой (постфиксный, правый) обход, прямая польская запись, обратная польская запись, прошитые деревья \\
\end{tabular}

\end{frame}

\begin{frame}
\label{9.3.1.}\subsubsection{9.3.1. Представление свободных деревьев }\frametitle{9.3.1. Представление свободных деревьев (1/6)}
Для представления деревьев можно использовать те же приёмы, что и
для представления графов общего вида~--- матрицы смежности и
инциденций, списки смежности и др. Но, используя особенные свойства
деревьев, можно предложить существенно более эффективные
представления. 
Рассмотрим следующее представление свободного
дерева,
известное как \odmkey{код Прюфера}{Код (code)!Прюфера (Pr$\ddot{\text{u}}$fer)}. %\footnote{Прюфер (ХХХХ--ХХХХ)}
Допустим, что вершины дерева $T(V,E)$ пронумерованы числами из
интервала $1..p$. Построим последовательность $A:
\textbf{array}\,\, [1..p-1]\,\, \textbf{of}\,\, 1..p\,$.
в
соответствии со следующим 
\odmkey{алгоритмом упаковки}{Алгоритм (algorithm)!построения кода Прюфера (generating of Pr$\ddot{\text{u}}$fer’s code)}.

\begin{algorithmic}
\REQUIRE
дерево $T(V,E)$ в любом представлении,
вершины дерева пронумерованы числами $1..p$ произвольным образом.
\ENSURE
массив \textbf{$A:$ array  $[1..p - 1]$ of $1..p$}~--
код Прюфера дерева $T$.
\FOR{$i$ \textbf{ from } $1$ \textbf{ to } $p-1$}
  \STATE $v\Assign \min\Set{k\in V}{d(k)=1}$
  \COMMENT{\color{gray}\small выбираем висячую вершину $v$}
  \STATE $A[i]\Assign \Gamma(v)$
  \COMMENT{\color{gray}\small  заносим в код номер вершины, смежной с $v$ }
  \STATE $V\Assign V- v$
  \COMMENT{\color{gray}\small  удаляем вершину $v$ из дерева }
\ENDFOR
\end{algorithmic}


\end{frame}

\begin{frame}
\frametitle{9.3.1. Представление свободных деревьев (2/6)}

По построенному коду можно
восстановить исходное дерево с
помощью следующего \odmkey{алгоритма распаковки}{Алгоритм (algorithm)!распаковки кода Прюфера (unpacking Prüfer’s code)}.


\begin{algorithmic}
\REQUIRE
Массив $A: \textbf{ array } [1..p - 1] \textbf{ of } 1..p$~--
код Прюфера дерева $T$.
\ENSURE
Дерево $T(V,E)$, заданное множеством рёбер $E$,
вершины дерева пронумерованы числами $1..p$.
\STATE $E\Assign \emptyset$
\COMMENT{\color{gray}\small   в начале множество рёбер пусто }
\STATE $B\Assign 1..p$
\COMMENT{\color{gray}\small   множество неиспользованных номеров вершин }
\FOR{$i$ \textbf{ from } $1$ \textbf{ to } $p-1$}
  \STATE $v\Assign \min\Set{k\in B}{\A{j\geq i}{k\neq A[j]}}$
  \STATE{$\qquad$}
  \COMMENT{\color{gray}\small   выбираем вершину $v$~--- неиспользованную вершину }
  \STATE{$\qquad$}
  \COMMENT {\color{gray}\small  с наименьшим номером, }
  \STATE{$\qquad$}
  \COMMENT {\color{gray}\small   который не встречается в остатке кода Прюфера }
  \STATE $E\Assign E+(v,A[i])$
  \COMMENT{\color{gray}\small   добавляем ребро $(v,A[i])$ }
  \STATE $B\Assign B- v$
  \COMMENT{\color{gray}\small   удаляем вершину $v$ из списка неиспользованных }
\ENDFOR
\end{algorithmic}
%\end{algorithm}


\end{frame}

\begin{frame}
\frametitle{9.3.1. Представление свободных деревьев (3/6)}
%\addvspace{-12pt}


\begin{ground}
Код Прюфера действительно является представлением свободного
дерева. Чтобы убедиться в этом, покажем, что если $T'$~--- дерево,
построенное алгоритмом распаковки по коду $A$, который
построен алгоритмом упаковки по дереву $T$, то $T'\sim T$.
Для этого установим отображение
$\Map{f}{1..p}{1..p}$ между номерами вершин в деревьях $T$ и $T'$
по порядку выбора вершин в алгоритмах: если вершина $v$ выбрана на
$i$-м шаге алгоритма упаковки, то вершина $f(v)$ выбрана на
$i$-м шаге алгоритма распаковки. 
Заметим, что $\Dom f =
1..p$, поскольку висячие вершины есть в любом дереве и удаление
висячей вершины оставляет дерево деревом. Далее, 
$\Im f = 1..p$,
поскольку на $i$-м шаге алгоритма распаковки использовано
$i - 1$ число из $p$ чисел и остаётся $p - i + 1$ чисел, а в хвосте кода
Прюфера занято не более $p - i$ чисел. Более того, $\A{v}{f(v) = v}$.
Действительно, номера вершин, которые являются висячими в исходном
дереве, не появляются в коде Прюфера (кроме, может быть, одной
висячей вершины с наибольшим номером), а номера вершин, которые не
являются висячими, обязательно появляются.
\end{ground}

\end{frame}

\begin{frame}
\frametitle{9.3.1. Представление свободных деревьев (4/6)}
%\addvspace{-12pt}


\begin{ground} (Окончание)
 Поскольку при выборе
первой вершины $v$ в алгоритме упаковки все вершины с меньшими
номерами не являются висячими, их номера будут присутствовать в
коде и, значит, не могут быть использованы на первом шаге
алгоритма распаковки. Таким образом, на первом шаге
алгоритм распаковки выберет ту же вершину $v$. Но после
удаления вершины $v$ на втором шаге снова имеется дерево, к
которому применимы те же рассуждения. Итак, $f$~--- тождественное
и, значит, взаимно-однозначное отображение. Заметим теперь, что
для определения $i$-го элемента кода на $i$-м шаге
алгоритма упаковки используется, а затем удаляется ребро
$(v,A[i])$, и в точности то же ребро добавляется в~дерево $T'$ на
$i$-м шаге алгоритма распаковки. Следовательно, $f$~---
взаимно-однозначное отображение, сохраняющее смежность и $T\sim
T'$.
\end{ground}

\end{frame}

\begin{frame}
\frametitle{9.3.1. Представление свободных деревьев (5/6)}
\begin{example}
Для дерева, представленного на рисунке, код Прюфера $7,
9, 1, 7, 2, 2, 7, 1, 2, 5, 12$. На этом рисунке числа
в вершинах~--- это их номера, а числа на рёбрах указывают порядок,
в котором выбираются висячие вершины и удаляются рёбра при
построении кода Прюфера.
\end{example}

\vspace{-12pt}
\begin{figure}[h]
\begin{center}
\includegraphics[height=40mm]{09_10}
\end{center}
%\caption{Построение кода Прюфера}
%\label{f0991}
\end{figure}

\begin{remark}
Код Прюфера~--- наиболее экономное по памяти
представление дерева. Его можно немного улучшить,
если заметить, что существует всего одно
дерево с двумя вершинами~--- $K_2$,
а потому информацию о <<последнем ребре>> можно
не хранить, она восстанавливается однозначно.
\end{remark}
\end{frame}

\begin{frame}
\frametitle{9.3.1. Представление свободных деревьев (6/6)}


\begin{example}
Для дерева, представленного на рисунке, 
код Прюфера 7, 9, 1, 7, 2, 2, 7, 1, 2, 5, 12. 
Рассмотрим протокол выполнения распаковки кода Прюфера:


\begin{tabular}{lllll}
$i$ & $B$ & $A$ & $v$ & $   A[i]$ \\
\hline
$1$ & $\{1,2,3,4,5,6,7,8,9,10,11,12\}$ & $[| 7,9,1,7,2,2,7,1,2,5,12]$ & $3$ & $7$ \\ 
$2$ & $\{1,2,4,5,6,7,8,9,10,11,12\}$   & $[7,| 9,1,7,2,2,7,1,2,5,12]$ & $4$ & $9$ \\ 
$3$ & $\{1,2,5,6,7,8,9,10,11,12\}$     & $[7,9,| 1,7,2,2,7,1,2,5,12]$ & $6$ & $1$ \\ 
$4$ & $\{1,2,5,7,8,9,10,11,12\}$       & $[7,9,1,| 7,2,2,7,1,2,5,12]$ & $8$ & $7$ \\ 
$5$ & $\{1,2,5,7,9,10,11,12\}$         & $[7,9,1,7,| 2,2,7,1,2,5,12]$ & $9$ & $2$ \\ 
$6$ & $\{1,2,5,7,10,11,12\}$           & $[7,9,1,7,2,| 2,7,1,2,5,12]$ & $10$ & $2$ \\ 
$7$ & $\{1,2,5,7,11,12\}$              & $[7,9,1,7,2,2,| 7,1,2,5,12]$ & $11$ & $7$ \\ 
$8$ & $\{1,2,5,7,12\}$                 & $[7,9,1,7,2,2,7,| 1,2,5,12]$ & $7$ & $1$ \\ 
$9$ & $\{1,2,5,12\}$                   & $[7,9,1,7,2,2,7,1,| 2,5,12]$ & $1$ & $2$ \\ 
$10$ & $\{2,5,12\}$                    & $[7,9,1,7,2,2,7,1,2,| 5,12]$ & $2$ & $5$ \\ 
$11$ & $\{5,12\}$                      & $[7,9,1,7,2,2,7,1,2,5,| 12]$ & $5$ & $12$ \\ 
\hline
\end{tabular}

\end{example}
\end{frame}

\begin{frame}
\label{9.3.2.}\subsubsection{9.3.2. Представление упорядоченных корневых деревьев }\frametitle{9.3.2. Представление упорядоченных корневых деревьев (1/5)}
В упорядоченных ордеревьях выделен корень и поддеревья
упорядочены. Использование этой информации позволяет построить ещё
более компактное представление по сравнению с кодом Прюфера для
свободных деревьев. Пусть упорядоченное ордерево задано списками
смежности, причём узлы перенумерованы в
соответствии с порядком прямого обхода, 
определённым в~доказательстве теоремы \DMref{п.~9.2.3}{9.2.3.}. Тогда следующий \odmkey{алгоритм
 построения кода упорядоченного ордерева}{Алгоритм (algorithm)!построения кода упорядоченного ордерева (generating of shortest spanning tree) } сгенерирует  
битовую шкалу (код) из $2q$ разрядов,
где $q$~--- число дуг в ордереве.
%\begin{algorithm}уг в ордео
%\caption{Построение кода упорядоченного ордерева}
%\label{ortreecode} \index{Алгоритм (algorithm)!построения кода
%упорядоченного ордерева (ordered directed tree code)}

\begin{algorithmic}
\REQUIRE
Нетривиальное упорядоченное ордерево $T$, \\
заданное списками смежности
\textbf{$\Gamma:$ array $[1..p]$ of $\Pointer N$},\\
где \textbf{$N=$ struct $\{v\Is 1..p\,; n\Is \Pointer N\}$},
причём узлы перенумерованы в прямом порядке.
\ENSURE
Массив \textbf{$C:$ array $[1..2q]$ of $0..1$},
являющийся кодом ордерева $T$.
\STATE $i\Assign 0$
\COMMENT{\color{gray}\small количество заполненных разрядов кода }
\STATE $\TraverseTree(\Gamma[1])$
\COMMENT{\color{gray}\small  указатель на список смежности корня }
\end{algorithmic}
%\end{algorithm}
\end{frame}


\begin{frame}
\frametitle{9.3.2. Представление упорядоченных корневых деревьев (2/5)}
Основная работа выполняется рекурсивной процедурой
$\TraverseTree$, для которой переменная $i$ и массивы $\Gamma$ и $C$ глобальны.

%\pagebreak[3]
\begin{algorithmic}
\REQUIRE $d : \Pointer N$~--- указатель на список исходящих дуг.
\ENSURE
заполнение двух разрядов кода $C$.
\WHILE{\textbf{$d\ne$ nil}}
\STATE $i\Assign i+1;\ C [i]\Assign 1$
\COMMENT{\color{gray}\small   отмечаем вход в узел }
\STATE $\TraverseTree(\Gamma[d.v])$
\COMMENT{\color{gray}\small   построение кода поддерева  }
\STATE $i\Assign i+1;\ C [i]\Assign 0$
\COMMENT{\color{gray}\small   отмечаем выход из узла }
\STATE $d \Assign d.n$
\COMMENT{\color{gray}\small   выбор следующего узла в списке }
\ENDWHILE
\end{algorithmic}

\begin{remark}
Представленный алгоритм, в сущности, строит протокол
выполнения алгоритма обхода в глубину упорядоченного ордерева.
Запись 1 отмечает вход в узел по единственной входящей дуге,
запись 0 отмечает возврат из узла по этой же дуге.
\end{remark}

\medskip
\odmkey{Восстановление упорядоченного ордерева по коду}{Алгоритм (algorithm)!восстановления упорядоченного ордерева по коду (of construction of ordered directed tree by code)} представлено в следующем алгоритме.
\end{frame}

\begin{frame}
%\frametitle{9.3.2. Представление упорядоченных корневых деревьев (3/5)}
%По построенному коду нетрудно восстановить исходное представление
%упорядоченного ордерева 

%\begin{algorithm}


%\label{undoortreecode} \index{Алгоритм (algorithm)!восстановления
%упорядоченного ордерева по коду (ordered directed tree recovery)}
\begin{algorithmic}
\REQUIRE
массив \textbf{$C:$ array $[1..2q]$ of $0..1$}~---
код ордерева $T$.
\ENSURE
массив
\textbf{$\Gamma:$ array $[1..p]$ of $\Pointer N$},
\textbf{$N= $ struct $ \{v : 1..p; n : \Pointer N \}$ } 
%---
%списки смежности упорядоченного ордерева $T$.
\STATE $p\Assign 1$
\COMMENT{\color{gray}\small   счётчик узлов }
\STATE \textbf{$\Gamma [1]\Assign$ nil}
\COMMENT{\color{gray}\small   корень }
\STATE $n\Assign 1$
\COMMENT{\color{gray}\small   текущий узел }
\FOR{\textbf{$i$ from $1$ to $2q$}}
  \IF{$C[i]=1$}
    \STATE $p\Assign p+1$
    \COMMENT{\color{gray}\small   номер нового узла }
    \STATE $\Gamma [p]\Assign \textbf{nil}$
    \COMMENT{\color{gray}\small   новый узел пока листовой }
    \STATE $d\Assign \NewNode(p,\textbf{nil})$
    \COMMENT{\color{gray}\small  дуга от текущего узла к новому узлу }
    \STATE $\Append(\Gamma[n],d)$
    \COMMENT{\color{gray}\small   добавить узел в список смежности }
    \STATE $n\rightarrow S$
    \COMMENT{\color{gray}\small  положить номер текущего узла на стек }
    \STATE $n\Assign p$
    \COMMENT{\color{gray}\small   перейти в новый узел }
  \ELSE
    \STATE $n\leftarrow S$
    \COMMENT{\color{gray}\small   снять со стека и перейти в новый узел }
  \ENDIF
\ENDFOR
\end{algorithmic}
%\end{algorithm}
\end{frame}

\begin{frame}
%\frametitle{9.3.2. Представление упорядоченных корневых деревьев (4/5)}
В алгоритме используются вспомогательный стек $S$ для хранения
номеров узлов и две процедуры:
\begin{itemize}
\item $\NewNode(v\Is 1..p, n\Is\Pointer N)\Is \Pointer N$~---
функция, создающая новый элемент списка смежности с полями $v$ и
$n$ и возвращающая указатель на него; \item $\Append(L,e \Is
\Pointer N)$~--- процедура, присоединяющая элемент, указатель на
который задан параметром $e$, к списку, указатель на первый
элемент которого задан параметром $L$.
\end{itemize}
{\renewcommand{\baselinestretch}{0.95}\selectfont
\begin{ground}
Массив $C$ является представлением упорядоченного ордерева $T$, то
есть если к $T$ применить алгоритм кодирования, а затем к
результату применить алгоритм декодирования, то получится
то же самое упорядоченное ордерево $T$. Действительно,
алгоритм декодирования интерпретирует протокол работы
алгоритма кодирования. 
Каждый раз, когда
алгоритм упаковки входит по дуге в некоторый узел первый
раз, он записывает в~протокол 1, и в точности в том же порядке
алгоритм распаковки создаёт узлы. Каждый раз, когда
алгоритм упаковки возвращается на предыдущий уровень и
записывает в~протокол 0, алгоритм распаковки
возвращается к родительскому узлу, снимая его номер со стека.
\end{ground}
}
\end{frame}

\begin{frame}
\frametitle{9.3.2. Представление упорядоченных корневых деревьев (5/5)}
\begin{example}
Применение алгоритма к упорядоченному ордереву на
рисунке даст код\\ 1101001101011000.
\begin{figure}[h]
\begin{center}
\includegraphics[height=30mm]{9_2_3_3.jpg}
\end{center}
\end{figure}
\end{example}
\begin{remark}
Наложенное здесь требование прямого порядка нумерации узлов
не является ограничением и используется только для упрощения изложения.
Можно показать, что при \em{любой} нумерации узлов ордерева
алгоритм~кодирования построит код, а алгоритм~декодирования
восстановит по этому коду изоморфное ордерево, отличающееся разве что
нумерацией узлов.
\end{remark}

\end{frame}


\begin{frame}

\label{9.3.3.}\subsubsection{9.3.3. Число упорядоченных ориентированных деревьев }\frametitle{9.3.3. Число упорядоченных ориентированных деревьев (1/3)} 
%{\renewcommand{\baselinestretch}{0.95}\selectfont
Представление упорядоченных ордеревьев, построенное в предыдущем
подразделе, обладает замечательным характеристическим свойством,
позволяющим получить явную формулу для числа упорядоченных
ордеревьев. 

Введём обозначения. Пусть \textbf{$b:$ array $[1..n]$
of $0..1$}~--- любая битовая шкала. Обозначим через $N_0(b,k)$~---
количество нулей в отрезке шкалы $b[1..k], k\le n$, и через
$N_1(b,k)$~--- количество единиц в отрезке шкалы $b[1..k], k\le
n$. Пусть теперь \textbf{$c:$ array $[1..2q]$ of $0..1$}~---
битовая шкала, построенная по упорядоченному ордереву $T$
алгоритмом кодирования. Тогда по построению алгоритма имеем
$$
N_0(c,2q)=N_1(c,2q) \And \A{k<2q}{N_0(c,k)\le N_1(c,k)}.
\eqno{(*)}
$$
Из алгоритма декодирования следует, что
данное свойство является характеристическим, то есть
\em{любая} шкала, обладающая свойством $(*)$,
является кодом некоторого дерева,
причём по теореме \DMref{п.~9.2.3}{9.2.3.} соответствие между деревьями и кодами
взаимно-однозначно.
Пусть $S_n$~--- множество битовых шкал длины $2n$:
$$
S_n\Def\Set{c=(\Tuple{a}{1}{2n})}{\A{i\in 1..2n}{a_i\in\{0,1\}}}.
$$
%}
\end{frame}


\begin{frame}
\frametitle{9.3.3. Число упорядоченных ориентированных деревьев (2/3)}
Обозначим число тех битовых шкал $c\in S_n$,
которые обладают свойством $(*)$,
через $C(n)$ (по определению $C(0)\Def 1$):

\vspace{-12pt}

$$
C(n)\Def \left\vert
\Set{c\in S_n}{N_0(c,2n) = N_1(c,2n) \And \A{k < 2n}{N_0(c,k)\le N_1(c,k)}}
\right\vert.
$$


\vspace{-2pt}
\begin{example}
Ясно, что $C(1)=\vert\{(1,0)\}\vert=1$, $C(2)=2$,
$C(3)=4$~--- см. рисунки в \DMref{п.~9.2.1}{9.2.1.}.
\end{example}
\begin{lemma}
Для числа $C(n)$ выполняется следующее равенство:

\vspace{-12pt}
$$
C(n)=\sum_{k=0}^{n-1}C(k)C(n-k-1).
$$
\end{lemma}

\vspace{-2pt}
%{\renewcommand{\baselinestretch}{0.95}\selectfont
\begin{proof}
Рассмотрим подмножество кодов $c\in S_n$, удовлетворяющих более
сильному условию

\vspace{-12pt}
$$
N_0(c,2q)=N_1(c,2q) \And \A{k<2q}{N_0(c,k) < N_1(c,k)},
\eqno{(*')}
$$

\vspace{-2pt}

и обозначим число таких кодов через $C'(n)$. Ясно, что в 
\em{любом} коде, 
удовлетворяющем условию $(*')$, первый бит равен 1, а
последний~--- 0. Если их отбросить, то оставшаяся часть кода будет
удовлетворять условию $(*)$, и поэтому $C'(n)=C(n-1)$, причём
по определению $C'(1)=C(0)=1$. 
\end{proof}

\end{frame}


\begin{frame}
\frametitle{9.3.3. Число упорядоченных ориентированных деревьев (3/3)}
%\vspace{-2pt}
\begin{proof} (Продолжение)
Пусть теперь $s$~--- наименьшее
число, такое, что код $c[1..s]$ удовлетворяет условию $(*')$.
Сгруппируем все коды, удовлетворяющие условию $(*)$, по значениям
числа $s$. В группе, в которой $s=n$, имеется $C'(n)$ кодов, а в
группах, в которых $1\le s \le n-1$, имеется $C'(s)C(n-s)$ кодов.
Имеем
\begin{align*}
C(n) &= C'(n)+\sum_{s=1}^{n-1}C'(s)C(n-s) = C(n-1)\cdot 1+
\sum_{s=1}^{n-1}C(s-1)C(n-s) = \\[-2pt]
&= C(n-1)C(0)+ \sum_{k=0}^{n-2}C(k)C(n-k-1) = \sum_{k=0}^{n-1}C(k)C(n-k-1),
\end{align*}
где $k\Assign s-1$.
\end{proof}

Таким образом, для числа $C(n)$ битовых шкал, удовлетворяющих
условию $(*)$, выполняется характеристическое рекуррентное
соотношение чисел Каталана (\DMref{п.~5.7.4}{5.7.4.}),
откуда немедленно следует 
%теорема.
\begin{theorem}
Число ориентированных упорядоченных деревьев
с $q$ дугами равно
$
\displaystyle{\frac{C(2q,q)}{q+1}.}
$
\end{theorem}


\end{frame}



\begin{frame}
\label{9.3.4.}\subsubsection{9.3.4. Проверка правильности скобочной структуры }\frametitle{9.3.4. Проверка правильности скобочной структуры (1/2)}
Из характеристического свойства $(*)$ рассматриваемого
представления упорядоченных ордеревьев в качестве побочного,
но полезного наблюдения можно извлечь эффективный \odmkey{алгоритм
проверки правильности скобочной структуры}{Алгоритм (algorithm)!проверки правильности скобочной структуры (checking parenthesis structure accuracy)}.

\begin{algorithmic}
\REQUIRE строка \textbf{$s:$ array $[1..n]$ of char}, возможно,
содержащая скобки <<(>> и <<)>>. 
\ENSURE число $0$, если скобочная
структура правильна, или число в диапазоне $1..(n+1)$, указывающее
на позицию в строке, где скобочная структура нарушена. 
\STATE $p\Assign 0$ 
\COMMENT{\color{gray}\small число прочитанных <<(>> минус число прочитанных <<)>> }
\FOR{\textbf{$i$ from $1$ to $n$}}
  \STATE{\textbf{if $s[i]="("$ then $p\Assign p+1$ end if }}   
  \COMMENT{\color{gray}\small  прочли  <<(>> }
  \STATE{\textbf{if $s[i]=")"$ then $p\Assign p-1$ end if }}   
  \COMMENT{\color{gray}\small  прочли <<)>> }
  \STATE{\textbf{if $p<0$ then return $i$ end if }}   
  \COMMENT{\color{gray}\small  лишняя закрывающая скобка }
\ENDFOR
  \STATE{\textbf{if $p=0$ then return $0$  }}   
  \COMMENT{\color{gray}\small  скобочная структура правильна }
  \STATE{\textbf{else return $n+1$ end if }}  
  \COMMENT{\color{gray}\small  не хватает закрывающих скобок }
\end{algorithmic}

\end{frame}



\begin{frame}
\frametitle{9.3.4. Проверка правильности скобочной структуры (2/2)}
%\pagebreak[3]
\begin{ground}
Всякая правильная скобочная структура взаимно-однозначно
соответствует упорядоченному ордереву. Если трактовать
открывающую скобку как 1, а закрывающую~--- как 0,
то последовательность скобок, содержащихся в~строке,
образует код дерева. Алгоритм очевидным образом проверяет
выполнение условия $(*)$, то есть проверяет допустимость
этого кода.
\end{ground}
\begin{example}
Строка $a(b(d)(e))(c(f)(g)(h(i)))$ является естественной записью с
помощью правильной скобочной структуры дерева на рисунке.

\begin{figure}[h]
\begin{center}
\includegraphics[height=30mm]{9_2_3_3.jpg}
\end{center}
\end{figure}
\end{example}
%\addvspace{-3pt}
\end{frame}


\begin{frame}
\label{9.3.5.}\subsubsection{9.3.5. Представление бинарных деревьев }\frametitle{9.3.5. Представление бинарных деревьев (1/5)}
%\addvspace{-3pt}

Всякое свободное дерево можно ориентировать, назначив один из узлов корнем.
Всякое ордерево можно произвольно упорядочить.
Для потомков одного узла (братьев)
упорядоченного ордерева определено отношение старше--младше (правее--левее).
Всякое упорядоченное дерево можно представить бинарным деревом, например,
проведя правую связь к старшему брату,
а левую~--- к младшему сыну.
Таким образом, достаточно рассмотреть представление
в программе бинарных
де\-ревьев. 
%\addvspace{-6pt}
\begin{example}
На рисунке приведены диаграммы упорядоченного и
соответствующего ему бинарного деревьев.
\end{example}

\vspace{-1pt}
\begin{figure}[h]
\begin{center}
\includegraphics[width=140mm]{9_3_5.jpg}
\end{center}
%\caption{Упорядоченное и бинарное деревья}
%\label{f0910}
\end{figure}
\end{frame}




\begin{frame}
\frametitle{9.3.5. Представление бинарных деревьев (2/5)}

\begin{remark}
Из данного представления следует, что множество бинарных деревьев \\
взаимно-однозначно соответствует множеству упорядоченных лесов
упорядоченных ордеревьев. Действительно, в указанном представлении
одиночному упорядоченному ордереву всегда соответствует бинарное
дерево, у которого правая связь корня пуста, а упорядоченному
лесу~--- бинарное дерево, у которого правая связь корня не пуста.
\end{remark}


\smallskip
Обозначим через $n(p)$ объём памяти, занимаемой представлением
бинарного дерева, где $p$~--- число узлов. Наиболее часто
используются следующие представления бинарных деревьев.
%\begin{enumerate}
%\item

\smallskip
{\color{blue}\bf 1.}
\textbf{Списочные структуры}: каждый узел  представляется записью \\
типа \textbf{$N=$ struct $\{ i\Is info;\ l,r\Is\Pointer N \}$},
содержащей два поля ($l$ и $r$)
с указателями на левый и правый узлы и еще одно поле $i$
для хранения указателя на информацию об узле,
где тип $info$ считается заданным.
Дерево представляется
указателем на корень.
Для этого представления $n(p) = 3p$.
\end{frame}


\begin{frame}
\frametitle{9.3.5. Представление бинарных деревьев (3/5)}

\begin{remark}
Поскольку в бинарном дереве,
как и в любом другом, $q=p-1$,
то из $2p$ указателей,
отводимых для хранения дуг,
$p+1$ всегда хранит значение \textbf{nil},
 т.~е. половина связей не используется.
\end{remark}

%\begin{enumerate}\setcounter{enumi}{1}
%\item 

{\color{blue}\bf 2.}
\textbf{Упакованные массивы}: узлы располагаются в
массиве так, что все узлы поддерева данного узла располагаются
вслед за этим узлом. Вместе с каждым узлом хранится индекс узла,
который является первым узлом правого поддерева данного узла.
Дерево $T$ определяется так: \textbf{$T\Is$
array $[1..p]$ of struct $\{ i\Is info;$ $k\Is 1..p\}$}, где
тип $info$ считается заданным. Для этого представления $n(p) =
2p$. 


{\color{blue}\bf 3.}
\textbf{Размеченная степень}: аналогично упакованным массивам, 
но вместо связей
фиксируется <<размеченная степень>> каждого узла. Например, 0
означает, что это лист, 1~--- есть правая связь, но нет левой,
2~--- есть левая связь, но нет правой, 3~--- есть обе связи.
Для представления размеченной степени в случае бинарных деревьев достаточно битовой шкалы длины $2$.
Дерево $T$ определяется так: \textbf{$T:$ array
$[1..p]$ of struct $\{i:info; $ $d:0..3\}$}, где тип $info$
считается заданным. Для этого представления $n(p) = 2p$. 

\end{frame}


\begin{frame}
\frametitle{9.3.5. Представление бинарных деревьев (4/5)}

{\color{blue}\bf 4.}
\textbf{Польская запись}:
Если
степень узла известна из информации, хранящейся в самом узле, то
можно не хранить и степень. Такой способ представления деревьев
называется \odmkey{польской записью}{Польская запись (polish
notation)} и обычно используется для представления выражений. В
этом случае представление дерева оказывается наиболее компактным:
объём памяти $n(p) = p$.

\vfill
%------------------------------------------------------------
% TO DO
% репозиторная модель необходима, но в учебнике недостаточно наглядно
% ФАН 12 апреля 2015


%\begin{example}
%Покажем, как выглядят в памяти бинарные деревья в разных
%представлениях. В приводимых ниже таблицах первая группа столбцов
%соответствует полям списочных структур,
%вторая~--- упакованным массивам и третья~--- польской записи с
%явными размеченными степенями. Условные адреса~--- это целые
%числа, а пустой указатель~--- 0.
%
%\begin{enumerate}
%\item Бинарное дерево на рис.~\ref{f0909}, {\it слева}; считаем,
%что верхний узел содержит $a$, а~нижний~--- $b$.
%%\addvspace{6pt}
%%\begin{center}
%
%{\sffamily\small
%\def\arraystretch{1.1}
%\begin{tabular}{l||c|c|c||c|c||c|c}
%%\hline
%Адрес & $i$ & $l$ & $r$ & $i$ & $k$ & $i$ & $d$ \\
%\hline
%$1$     & $a$ & $2$ & $0$ & $a$ & $0$ & $a$ & $1$ \\
%$2$     & $b$ & $0$ & $0$ & $b$ & $0$ & $b$ & $0$ \\
%\hline
%\end{tabular}}
%
%%\addvspace{9pt}
%
%\item Бинарное дерево на рис.~\ref{f0909}, {\it справа}; считаем,
%что верхний узел содержит $a$, а~нижний~--- $b$.
%%\addvspace{6pt}
%%\begin{center}
%
%{\sffamily\small
%\def\arraystretch{1.1}
%\begin{tabular}{l||c|c|c||c|c||c|c}
%%\hline
%Адрес & $i$ & $l$ & $r$ & $i$ & $k$ & $i$ & $d$ \\
%\hline
%$1$     & $a$ & $0$ & $2$ & $a$ & $2$ & $a$ & $2$ \\
%$2$     & $b$ & $0$ & $0$ & $b$ & $0$ & $b$ & $0$ \\
%\hline
%\end{tabular}}
%%\addvspace{9pt}
%
%\item Бинарное дерево на рис.~\ref{f0907}, {\it слева}.
%%\addvspace{6pt}
%%\begin{center}
%
%{\sffamily\small
%\def\arraystretch{1.1}
%\begin{tabular}{l||c|c|c||c|c||c|c}
%%\hline
%Адрес & $i$ & $l$ & $r$ & $i$ & $k$ & $i$ & $d$ \\
%\hline
%$1$     & $+$ & $2$ & $3$ & $+$ & $3$ & $+$ & $3$ \\
%$2$     & $a$ & $0$ & $0$ & $a$ & $0$ & $a$ & $0$ \\
%$3$     & $*$ & $4$ & $5$ & $*$ & $5$ & $*$ & $3$ \\
%$4$     & $b$ & $0$ & $0$ & $b$ & $0$ & $b$ & $0$ \\
%$5$     & $c$ & $0$ & $0$ & $c$ & $0$ & $c$ & $0$ \\
%\hline
%\end{tabular}}
%%\addvspace{9pt}
%\vspace{9pt} \pagebreak[4] \item Бинарное дерево на
%рис.~\ref{f0910}, {\it справа}.
%%\addvspace{6pt}
%%\begin{center}
%
%{\sffamily\small
%\def\arraystretch{1.1}
%\begin{tabular}{l||c|c|c||c|c||c|c}
%%\hline
%Адрес & $i$ & $l$ & $r$ & $i$ & $k$ & $i$ & $d$ \\
%\hline
%$1$    & $a$ & $2$ & $0$ & $a$ & $0$ & $a$ & $1$ \\
%$2$    & $b$ & $3$ & $5$ & $b$ & $5$ & $b$ & $3$ \\
%$3$    & $d$ & $0$ & $4$ & $d$ & $4$ & $d$ & $2$ \\
%$4$    & $e$ & $0$ & $0$ & $e$ & $0$ & $e$ & $0$ \\
%$5$    & $c$ & $6$ & $0$ & $c$ & $0$ & $c$ & $1$ \\
%$6$    & $f$ & $0$ & $7$ & $f$ & $7$ & $f$ & $2$ \\
%$7$    & $g$ & $0$ & $8$ & $g$ & $8$ & $g$ & $2$ \\
%$8$     & $h$ & $9$ & $0$ & $h$ & $0$ & $h$ & $1$ \\
%$9$     & $i$ & $0$ & $0$ & $i$ & $0$ & $i$ & $0$ \\
%\hline
%\end{tabular}}
%
%\end{enumerate}
%\end{example}
%\addvspace{12pt}
\end{frame}


\begin{frame}
\frametitle{9.3.5. Представление бинарных деревьев (5/5)}


\begin{example}
Рассмотрим бинарное дерево на
рисунке и покажем, как выглядят в памяти бинарные деревья в разных
представлениях. В приводимой ниже таблице первая группа столбцов
соответствует полям списочных структур,
вторая~--- упакованным массивам и третья~--- польской записи с
явными размеченными степенями. Условные адреса~--- это целые
числа, а пустой указатель~--- $0$.

\begin{columns}[t]
\begin{column}[T]{8cm}
\vspace{-12pt}
\begin{center}
\begin{tabular}{|l||c|c|c||c|c||c|c|}
\hline
Адрес & $i$ & $l$ & $r$ & $i$ & $k$ & $i$ & $d$ \\
\hline
$1$    & $a$ & $2$ & $0$ & $a$ & $0$ & $a$ & $2$ \\
$2$    & $b$ & $3$ & $5$ & $b$ & $5$ & $b$ & $3$ \\
$3$    & $d$ & $0$ & $4$ & $d$ & $4$ & $d$ & $1$ \\
$4$    & $e$ & $0$ & $0$ & $e$ & $0$ & $e$ & $0$ \\
$5$    & $c$ & $6$ & $0$ & $c$ & $0$ & $c$ & $2$ \\
$6$    & $f$ & $0$ & $7$ & $f$ & $7$ & $f$ & $1$ \\
$7$    & $g$ & $0$ & $8$ & $g$ & $8$ & $g$ & $1$ \\
$8$     & $h$ & $9$ & $0$ & $h$ & $0$ & $h$ & $2$ \\
$9$     & $i$ & $0$ & $0$ & $i$ & $0$ & $i$ & $0$ \\
\hline
\end{tabular}
\end{center}
\end{column}
\begin{column}[T]{4cm}
\medskip
\includegraphics[scale=1]{9_3_5_2.jpg}
\end{column}
\end{columns}
\end{example}



\end{frame}

\begin{frame}
\label{9.3.6.}\subsubsection{9.3.6. Обходы бинарных деревьев }\frametitle{9.3.6. Обходы бинарных деревьев (1/5)}
Большинство алгоритмов работы с
деревьями основаны на \em{обходах}. 
Возможны следующие основные обходы бинарных
деревьев.
\begin{itemize}
\item
\odmkey{Прямой}{Обход дерева (traverse of tree)!прямой
(preorder)} (\odmkey{префиксный}{Обход дерева (traverse of tree)!левый (preorder)},
  \odmkey{левый}{Обход дерева (traverse of tree)!левый (preorder)}) обход:
                      попасть в узел,  
                       обойти левое поддерево, 
                       обойти правое поддерево. 
\item
\odmkey{Внутренний}{Обход дерева (traverse of tree)!внутренний
(inorder)} (\odmkey{инфиксный}{Обход дерева (traverse of tree)!инфиксный (inorder)},
  \odmkey{симметричный}{Обход дерева (traverse of tree)!симметричный (inorder)}) обход:
                     обойти левое поддерево,
                     попасть в узел,
                     обойти правое поддерево.
\item                     
\odmkey{Концевой}{Обход дерева (traverse of tree)!концевой
(postorder)} (\odmkey{постфиксный}{Обход дерева (traverse of tree)!постфиксный (postorder)},
  \odmkey{правый}{Обход дерева (traverse of tree)!правый (postorder)}) обход:
                      обойти левое поддерево,
                      обойти правое поддерево, 
                      попасть в узел.
\end{itemize}


Возможны еще три обхода,
отличающихся порядком рассмотрения левых и правых поддеревьев. Этим
исчерпываются обходы, если в представлении есть только дуги,
ведущие от каждого узла к левому и правому поддеревьям, и не используются дополнительные структуры данных.
\begin{example}
Прямой обход дерева выражения  $a+b*c$  даёт
\odmkey{прямую}{Польская запись (polish notation)!прямая
(prefix)} польскую запись этого выражения: $+a*bc$.
Концевой обход дерева выражения  $a+b*c$  даёт
\odmkey{обратную}{Польская запись (polish notation)!обратная
(postfix)} польскую запись этого выражения: $abc*+$.
Внутренний обход дерева выражения  $a+b*c$  восстанавливает исходную запись выражения: $a+b*c$.
\end{example}

\end{frame}

\begin{frame}
\frametitle{9.3.6. Обходы бинарных деревьев (2/5)}

%-------------------------------------------------------------
% TO DO вставить алгоритм внутреннего интерпретатора Форта ФАН
%-------------------------------------------------------------
%{\renewcommand{\baselinestretch}{0.95}\selectfont
\begin{digress}
Польская запись выражений (прямая или обратная) применяется в
некоторых языках программирования непосредственно и используется в
качестве внутреннего представления программ во многих трансляторах
и интерпретаторах. Причина заключается в том, что такая форма
записи допускает очень эффективную интерпретацию (вычисление
значения) выражений. Например, значение выражения в
обратной польской записи может быть вычислено при однократном
просмотре выражения слева направо с использованием одного стека. В
таких языках, как Forth и PostScript, обратная польская запись
используется как основная.
\end{digress}


\begin{remark}
%------------------------ ФАН 12.04.2105 -----------------------------
% Это то, во что трансформировалось замечание Ивана из файла remarks
Ордерево задаёт частичный порядок на множестве узлов, причём
минимальный элемент единственный (корень). 
Всякий обход дерева задаёт строгий линейный порядок на множестве узлов.
Прямой обход дерева задаёт линейный порядок на множестве узлов,
который дополняет исходный порядок. 
Некоторые другие обходы также сохраняют
исходный порядок (например, поиск в ширину или 
обход \odmkey{по уровням}{Обход дерева (traverse of tree)!по уровням (breadth-first search)}). В то же  время, 
часть обходов нарушает исходный порядок  
(например, концевой и симметричный обходы). 
\end{remark}
%}

\end{frame}

\begin{frame}
\frametitle{9.3.6. Обходы бинарных деревьев (3/5)}
Программные реализации различных обходов бинарных деревьев однотипны.

\smallskip
В качестве примера рассмотрим инфиксный обход.
Реализация обхода бинарного дерева с помощью рекурсивной процедуры
$\TraverseTree$
не вызывает затруднений. 

\begin{algorithmic}
\REQUIRE
     бинарное дерево, представленное списочной структурой,
     $q$~--- указатель на узел, при первом вызове это указатель на корень дерева.
\ENSURE
     последовательность узлов бинарного дерева в порядке симметричного обхода.
\STATE{\textbf{if $q.l  \ne $ nil then $\TraverseTree(q.l)$ end if}}
 \COMMENT{\color{gray}\small обойти левое поддерево}
\STATE\textbf{yield $q$} 
\COMMENT{\color{gray}\small  очередной узел при симметричном обходе }
\STATE{\textbf{if $ q.r  \ne $ nil then $ \TraverseTree(q.r) $ end if}}
 \COMMENT{\color{gray}\small  обойти правое поддерево}
\end{algorithmic}


В некоторых случаях из соображений
эффективности применение явной рекурсии оказывается нежелательным.
Следующий алгоритм реализует наиболее
популярный симметричный обход без явной рекурсии, но с
использованием стека.
\end{frame}

\begin{frame}
\frametitle{9.3.6. Обходы бинарных деревьев (4/5)}
\begin{algorithmic}
\REQUIRE
     бинарное дерево, представленное списочной структурой, $r$~--- указатель на корень.
\ENSURE последовательность узлов бинарного дерева в порядке симметричного обхода.
\STATE $T \Assign  \emptyset; p \Assign  r$
\COMMENT{\color{gray}\small  вначале стек пуст и $p$ указывает на корень дерева }
%\STATE $M:$ \COMMENT{\color{gray}\small  анализируем узел, на который указывает $p$ }
\REPEAT
\IF{\textbf{$p =$ nil}}
%  \IF{$T=\emptyset$}
%    \STATE\textbf{stop} \COMMENT{\color{gray}\small  обход закончен }
%  \ENDIF
  \STATE $p \leftarrow T $ \COMMENT{\color{gray}\small  левое поддерево обойдено }
  \STATE\textbf{ yield $p$} 
  \COMMENT{\color{gray}\small  очередной узел при симметричном обходе }
  \STATE $p\Assign p.r$ 
  \COMMENT{\color{gray}\small  начинаем обход правого поддерева }
\ELSE \STATE $p\rightarrow T$ 
\COMMENT{\color{gray}\small  запоминаем текущий узел
в стеке \ldots }
\STATE $p\Assign p.l$ \COMMENT{\color{gray}\small  \ldots и начинаем обход левого поддерева }
\ENDIF
\UNTIL\textbf{$T=\emptyset \And p=$ nil}
\end{algorithmic}


%\end{algorithm}
\end{frame}

\begin{frame}
\frametitle{9.3.6. Обходы бинарных деревьев (5/5)}
\begin{ground}
Заметим, что при каждом повторении тела основного цикла стек $T$ хранит все узлы, которые находятся выше узла $p$ и ещё не обойдены. 
Ясно, что условие \\$T=\emptyset$ является необходимым условием окончания обхода. Однако, если при этом\\
\textbf{$p\ne $ nil}, то условие пустоты стека
не является достаточным, обход ещё не завершен,
поскольку остались непройденные узлы, в частности 
узел $p$.
\end{ground}

\begin{remark}
Другие варианты обходов также могут быть
запрограммированы с использованием вспомогательных структур данных и без использования рекурсии. 
\end{remark}

\medskip
Как уже отмечено, в рассматриваемом 
представлении бинарных деревьев половина 
связей не используется.
Деревья, в
которых пустые поля $l$ и $r$ в структуре $N$ используются для
хранения дополнительных связей, называются \odmkey{прошитыми
деревьями}{Дерево (tree)!прошитое (threaded)}.

\medskip
\begin{remark}
Если кроме связей <<влево>> и <<вправо>> в представлении есть другие связи,
то возможны другие обходы.
В частности, деревья можно прошить 
таким образом, что обходы выполняются 
без рекурсии и без дополнительных структур данных. 
\end{remark}

\end{frame}

%\begin{frame}
%\frametitle{9.3.7. Изотропное представление бинарных деревьев}
%
%{\color{red}
%Тема для сочинения: изотропные обходы бинарных деревьев.
%Если в каждом узле хранить три указателя во все стороны, 
%то никакие стеки и рекурсии не нужны. 
%
%Алгоритм симметричного обхода с обоснованием.
%
%
%
%}
%\end{frame}

%\begin{frame}
%\frametitle{9.3.8. Прошитые бинарные деревья}
%
%{\color{red}
%
%Тема для сочинения: прошитые бинарные деревья.
%Если в каждом узле вместо пустых указателей вниз 
%хранить пунктирные указатели вверх, 
%то возможны фиксированные обходы без стека и рекурсии.
%См. Д.Кнут. Искусство программирования для ЭВМ.
%
% 
%}
%\end{frame}

\begin{frame}
\frametitle{Онтология теории графов}
\begin{center}
\includegraphics[height=8cm]{ODM_9_1_3.png}
\end{center}

\end{frame}

\subsection{9.4. Деревья сортировки}
\label{9.4.}
\begin{frame}
\frametitle{9.4. Деревья сортировки (1/2)}
В этом разделе обсуждается одно конкретное применение деревьев в
программировании, а именно \odmkey{деревья сортировки}{Дерево
(tree)!сортировки (sorting)} (также называемые \odmkey{деревьями
упорядочивания}{Дерево (tree)!упорядочивания (sorting)} или
\odmkey{деревьями поиска}{Дерево
	(tree)!поиска (search)}). При этом
рассматривается практическая
реализация алгоритмов работы с деревом сортировки,
а также целый ряд прагматических аспектов применения деревьев
сортировки.

\medskip
Деревья сортировки мы рассматриваем в контексте
общими методами организации хранения и поиска данных, сопоставляя их с 
\odmkey{таблицами расстановки}{Таблица (table)!расстановки (hash)}.

\end{frame}

\begin{frame}
\frametitle{9.4. Деревья сортировки (2/2)}
\medskip
\begin{tabular}{p{9mm}|p{70mm}|p{64mm}}
  № & Название параграфа
  & Ключевые термины и обозначения \\
  \hline
\DMref{9.4.1.}{9.4.1.} & Ассоциативная память 
  & 
 Ассоциативный масиив, 
 словарь, ключ, запись, многозначная ассоциация, таблица расстановки, хэш-таблица, дерево сортировки \\
\DMref{9.4.2.}{9.4.2.} & Таблицы расстановки
   & 
  Хэш-функция, хэш-значение, коллизия, парадокс дней рождения\\
\DMref{9.4.3.}{9.4.3.} & Эффективность хэширования & Метод деления, метод умножения
\\
\DMref{9.4.4.}{9.4.4.}  & Алгоритм поиска в дереве сортировки & 
\\
\DMref{9.4.5.}{9.4.5.}\RS & Алгоритм вставки в дереве сортировки & 
Обфускация \\
\DMref{9.4.6.}{9.4.6.}\RS  & Алгоритм удаления из дерева сортировки & \\
\DMref{9.4.7.}{9.4.7.} & Вспомогательные алгоритмы для дерева сортировки & 
 \\
\DMref{9.4.8.}{9.4.8.}  & Сравнение представлений ассоциативной памяти & 
 \\
\end{tabular}

\end{frame}



\begin{frame}
\label{9.4.1.}\subsubsection{9.4.1. Ассоциативная память }
\frametitle{9.4.1. Ассоциативная память (1/6)}
В практическом программировании для организации хранения данных и
доступа к ним часто используется механизм, 
который обычно называют
\odmkey{ассоциативной памятью}{Память ассоциативная \\
(content-addressable memory)} (или 
 \odmkey{ассоциативным массивом}{Массив (array)!ассоциативный массив (associative)}, или же просто \odmkey{словарем}{Словарь (dictionary)}). 

При использовании ассоциативной
памяти данные делятся на порции (называемые
\odmkey{записями}{Запись (record)}), и с каждой записью
ассоциируется \odmkey{ключ}{Ключ (key)!ассоциативной памяти
(content-addressable memory)}. Ключ~--- это значение из некоторого
линейно упорядоченного множества, а записи могут иметь
произвольную природу и различные размеры. Предполагается, что ассоциативная память не может хранить записи с одинаковыми ключами, то есть ключи записей различны. Доступ к данным
осуществляется по значению ключа, которое обычно выбирается
простым, компактным и~удобным для работы.
%\vspace{-6pt}
\begin{examples}
Ассоциативная  память используется во многих случаях:

\begin{enumerate}
\item Адресная книга: ключом является имя абонента,
а~записью~--- адресная информация (телефон(ы), почтовый адрес и
т.~д.).
\item Банковские счета: ключом является номер счета,
а записью~--- финансовая информация (которая может быть 
сложной).
\end{enumerate}
\end{examples}
\end{frame}

\begin{frame}
\frametitle{9.4.1. Ассоциативная память (2/6)}
Таким образом, ассоциативная память должна поддерживать
по меньшей мере три основные операции:
\begin{enumerate}
\item[1)] \em{добавить} (ключ, запись);%
\item[2)] \em{найти} (ключ): запись;%
\item[3)] \em{удалить} (ключ).%
\end{enumerate}
Указанные операции при необходимости дополняются и другими. Широко применяемые расширения включают в себя операции: 
\begin{enumerate}
\item[4)] \em{изменить} (ключ, запись);%
\item[5)] \em{найти максимум} (): запись;%	
\item[6)] \em{найти минимум} (): запись.%
\end{enumerate}
\begin{remark}
	Часто операции изменения и добавления реализуются совместно. В этом случае выполняемое действие зависит от наличия или отсутствия записи с данным ключом в ассоциативной памяти.
\end{remark}

%{\color{blue} Тема для простого и эффективного сочинения:
%примеры, когда кто у кого выигрывает по скорости в зависимости от 
%$n$ и от соотношения частот. 
%Возможно, эти примеры нужны в \em{последнем} параграфе раздела. }
\end{frame}

\begin{frame}
\frametitle{9.4.1. Ассоциативная память (3/6)}
%\medskip
Эффективность каждой операции зависит от структуры данных,
используемой для представления ассоциативной памяти,
и, как правило, существенно зависит от количества записей $n$.
Эффективность ассоциативной памяти в целом
зависит от соотношения частоты выполнения различных операций
в~данной конкретной программе, использующей ассоциативную память.

%\medskip
\begin{remark}
	Таким образом, невозможно указать способ
	организации ассоциативной памяти,
	который оказался бы наилучшим во всех возможных
	случаях.
\end{remark}

%\medskip
	Обобщением понятия ассоциативной памяти является механизм, называемый \odmkey{многозначная ассоциация}{Многозначная ассоциация (multimap)}, также ассоциирующий записи с ключами, но при этом допускается, чтобы одному ключу соответствовало множество записей.
	
%\medskip
\begin{examples}
\begin{enumerate}
		\item Толковый словарь: записью является толкование словарной статьи, а ключом --- слово или словосочетание (заголовок словарной статьи).
 Как известно, многие слова имеют несколько значений.
		\item Хранение взвешенных ребер мультиграфа $G(V, E)$: ключом является пара вершин, а записью --- вес ребра.
	\end{enumerate}
\end{examples}
\end{frame}


\begin{frame}
%\label{9.4.2.}\subsubsection{9.4.2. Способы реализации ассоциативной памяти }
\frametitle{9.4.1. Ассоциативная память (4/6)}

%{\renewcommand{\baselinestretch}{0.95}\selectfont
Для представления ассоциативной памяти используются разнообразные структуры данных:
\begin{enumerate}
\item[1)] неупорядоченный массив, в элементах которого хранятся значения
ключей актуальных записей; 
\item[2)] упорядоченный массив, в элементах которого хранятся значения
ключей актуальных записей в порядке возрастания (или убывания);
%\pagebreak[2] 
\item[3)] \odmkey{таблица расстановки}{Таблица
(table)!расстановки (hash)} (или \odmkey{хэш-таблица}{Таблица
(table)!хэш (hash)}),
в которой
адрес хранения записи определяется по значению 
ключа;
\item[4)] неупорядоченный список, в элементах которого хранятся значения
ключей актуальных записей; 
\item[5)] упорядоченный список, в элементах которого хранятся значения
ключей актуальных записей в порядке возрастания (или убывания);
\item[6)] \odmkey{дерево сортировки}{Дерево
(tree)!сортировки (sorting)}~--- бинарное дерево, каждый узел
которого содержит ключ (и указатель на запись) и обладает
следующим свойством: значения ключа во всех узлах левого поддерева
меньше, а во всех узлах правого поддерева~--- больше, чем значение
ключа в узле. 
\end{enumerate}
%}


\end{frame}


\begin{frame}
\frametitle{9.4.1. Ассоциативная память (5/6)}
%{\renewcommand{\baselinestretch}{0.92}\selectfont
%\small
При использовании неупорядоченного массива алгоритмы реализации
операций ассоциативной памяти очевидны.

\vspace{-1pt}
\begin{enumerate}
\item Операция <<добавить (ключ, запись)>> реализуется добавлением
записи в конец массива (трудоёмкость $O(1)$). 

\vspace{-1pt}
\item Операция <<найти (ключ) : запись>>
реализуется проверкой в цикле всех записей в массиве (трудоёмкость $O(n)$). 

\vspace{-1pt}
\item
Операция <<удалить (ключ)>> реализуется поиском удаляемой записи,
а затем перемещением последней записи на место удаляемой (трудоёмкость $O(n)$).
\end{enumerate}


Неупорядоченный список 
не имеет  преимуществ по сравнению с массивом.
\vspace{-1pt}
\begin{enumerate}
\item Операция <<добавить (ключ, запись)>> реализуется добавлением
записи в начало списка (трудоёмкость $O(1)$). 

\vspace{-1pt}
\item Операция <<найти (ключ): запись>>
реализуется проверкой в цикле всех записей в списке (трудоёмкость $O(n)$). 

\vspace{-1pt}
\item
Операция <<удалить (ключ)>> реализуется поиском удаляемой записи,
а затем исключением записи из списка (перемещать 
элементы при этом нет нужды, но трудоёмкость всё равно составляет $O(n)$).
\end{enumerate}
%}

\end{frame}

\begin{frame}
\frametitle{9.4.1. Ассоциативная память (6/6)}
%{\renewcommand{\baselinestretch}{0.92}\selectfont
%\small
Упорядоченный массив нередко используется
для реализации ассоциативной памяти,
поскольку является сравнительно эффективной структурой.

\vspace{-1pt}
\begin{enumerate}
\vspace{-1pt}
\item Операция <<найти (ключ) : запись>>
реализуется алгоритмом
бинарного поиска, см. \DMref{п.~5.2.7}{5.2.7.} 
(трудоёмкость $O(\log n)$). 

\item Операция <<добавить (ключ, запись)>> реализуется в два этапа: 
сначала с помощью операции <<найти>>
надобно найти место для вставки нового элемента, а затем <<подвинуть хвост>>
упорядоченного массива, чтобы освободить это место (общая трудоёмкость $O(n)$). 


\vspace{-1pt}
\item
Операция <<удалить (ключ)>> совершенно
аналогично реализуется поиском удаляемой записи
и последующее <<перемещением хвоста>> (трудоёмкость $O(n)$).
\end{enumerate}


Упорядоченный список 
не имеет  никаких преимуществ по сравнению с неупорядоченным, явно проигрывает всем остальным представлениям ассоциативной памяти и здесь более не рассматривается.

\medskip
Таким образом, осталось рассмотреть 
таблицы расстановки и деревья сортировки.

\end{frame}
%
%%\begin{frame}
%%\frametitle{9.4.1. Ассоциативная память (6/6)}
%%%{\renewcommand{\baselinestretch}{0.93}\selectfont\small
%%\begin{digress}
%%Таблицы расстановки являются важным практическим
%%приёмом программирования, подробное описание которого не включено
%%в курс из экономии места. Идея заключается в
%%том, что подбирается специальная функция, называемая
%%\odmkey{хэш-функцией}{Функция (function)!хэш (hash)}, переводящая
%%значение ключа в адрес хранения записи (адресом может быть индекс
%%в массиве, номер кластера на диске и т.~д.). Таким образом, по
%%значению ключа с помощью хэш-функции сразу определяется место
%%хранения записи и открывается доступ к ней. Хэш-функция
%%подбирается таким образом, чтобы разным ключам соответствовали, по
%%возможности, разные адреса из диапазона возможных адресов записей.
%%Как правило, мощность множества ключей существенно больше размера
%%пространства адресов, которое, в свою очередь, больше количества
%%одновременно хранимых записей. Поэтому при использовании
%%хэширования возможны \odmkey{коллизии}{Коллизия (collision)}~---
%%ситуации, когда хэш-функция сопоставляет один и тот же адрес двум
%%актуальным записям с различными ключами. Методы
%%хэширования отличаются друг от друга способами разрешения коллизий
%%и приёмами вычисления хэш-функций.
%%Тщательная программная реализация и соблюдение
%%ограничений на мощность множества ключей, адресов и записей
%%позволяют  избегать коллизий.
%%Если коллизий нет, то трудоёмкость всех трёх операций составляет 
%%$O(1)$.
%%\end{digress}
%%%}
%%
%%\end{frame}
%
%%\begin{frame}
%%\frametitle{9.4.1. Ассоциативная память (6/6)}
%\end{frame}
%
%\begin{frame}
%\frametitle{9.4.3. Алгоритм бинарного (двоичного) поиска (1/3)}
%
%В этом параграфе обсуждается 
%использование упорядоченного массива
%для реализации ассоциативной памяти. 
%
%
%\medskip
%При использовании упорядоченного массива для представления
%самой важной оказывается операция поиска, поскольку 
%чтобы удалить запись, её необходимо сначала найти,
%и при вставке записи необходимо найти место для этой записи.
%
%В упорядоченном массиве операция поиска записи по ключу может быть
%выполнена за время  $O(\log n)$ (где $n$~--- количество
%записей) с помощью следующего алгоритма, известного как алгоритм
%\odmkey{}{Алгоритм (algorithm)!поиска (search)!бинарного (binary)}
%\odmkey{бинарного}{Поиск (search)!бинарный (binary)} (или
%\odmkey{двоичного}{Поиск (search)!двоичный (binary)}) поиска.
%
%
%
%\end{frame}
%
%
%\begin{frame}
%\frametitle{9.4.3. Алгоритм бинарного (двоичного) поиска (2/3)}
%\begin{algorithmic}
%\REQUIRE упорядоченный массив
%\textbf{$A\Is$ array $[1..n]$ of struct $\{k\Is key; i\Is info\}$};
%ключ  $a\Is key$.
%\ENSURE
%индекс записи с искомым ключом $a$ в массиве $A$ или 0,
%если записи с~таким ключом нет.
%\STATE $b \Assign  1$ 
%\COMMENT{\color{gray}\small начальный индекс части массива для поиска }
%\STATE $e \Assign  n$ 
%\COMMENT{\color{gray}\small  конечный индекс части массива для поиска }
%\WHILE{$b\le e$}
%\STATE $c \Assign  (b+e)\DIV 2$ 
%\COMMENT{\color{gray}\small  индекс проверяемого элемента }
%\IF{$A[c].k > a$}
%\STATE $e \Assign  c-1$ 
%\COMMENT{\color{gray}\small  продолжаем поиск в первой половине }
%\ELSIF{$A[c].k < a$}
%\STATE $b \Assign  c+1$ 
%\COMMENT{\color{gray}\small  продолжаем поиск во второй половине }
%\ELSE
%\STATE \textbf{return $c$} 
%\COMMENT{\color{gray}\small  нашли искомый ключ }
%\ENDIF
%\ENDWHILE
%\STATE \textbf{return $0$} \COMMENT{\color{gray}\small искомого ключа нет в массиве }
%
%\end{algorithmic}
%\end{frame}
%
%
%\begin{frame}
%\frametitle{9.4.3. Алгоритм бинарного (двоичного) поиска (3/3)}
%\begin{ground}
%Достаточно заметить,
%что на каждом шаге основного цикла искомый элемент
%массива (если он есть) находится между
%(включительно)
%элементами с индексами $b$ и $e$. Поскольку
%диапазон поиска на каждом шаге уменьшается вдвое,
%общая трудоёмкость не превосходит $\log n$.
%\end{ground}
%
%\smallskip
%\begin{digress}
%Упорядоченные списки проигрывают упорядоченным массивам при реализации
%операций
%ассоциативной памяти, поскольку для списков нет аналога 
%алгоритма двоичного поиска.
%Это частное наблюдение является проявлением
%весьма общего факта, состоящего в том,
%что время доступа к элементу массива не зависит от
%номера элемента и от размера массива.
%Говорят, что массив является 
%структурой данных \odmkey{прямого доступа}{Прямой доступ (direct access)},
%это характеристическое свойство массива.
%Во всех же <<динамических>> структурах данных, в частности, в списках,
%время доступа \em{зависит} от положения элемента и от размера структуры данных. 
%\end{digress}
%
%
%\end{frame}


\begin{frame}
\label{9.4.2.}\subsubsection{9.4.2. Таблицы расстановки}
\frametitle{9.4.2. Таблицы расстановки (1/7)}
Если множество возможных ключей $U$ 
(обозначим $u\Assign  |U|$) 
сравнимо по мощности с множеством актуальных записей $R$ (обозначим $n\Assign  |R|$),
то есть если $n\approx u, n\le u$, 
то обычный массив \textbf{$T:$  array $[U]$
of  $\Pointer R$} является вполне
удовлетворительным представлением 
ассоциативной памяти. Действительно,
в этом случае индексами являются ключи,
а элементами являются адреса записей 
в некотором внешнем хранилище.
При этом обеспечивается \em{прямой доступ} к записи по ключу (см. \DMref{п.~5.2.7}{5.2.7.})
и сравнительно небольшой перерасход памяти на хранение пустых ячеек, если $n<u$.  

\smallskip
Однако на практике чаще бывает так, 
что  мощность множества ключей 
существенно больше допустимого размера
масcива $m$, который, в свою очередь, 
несколько больше количества
одновременно хранимых записей ($u \gg m \ge n$).

\smallskip
В таких случаях
применяют
 \odmkey{таблицу расстановки}{Таблица
	(table)!расстановки (hash)}, в которой вместо непосредственного использование ключа в качестве индекса, индекс вычисляется по ключу посредством \odmkey{хэш-функции}{Функция (function)!хэш (hash)}. Основная идея хэширования: уменьшить рабочий диапазон индексов массива $m$.

\begin{remark} 
Таблицы расстановки применяют не только для реализации ассоциативной памяти,
но и в ряде других областей, 
например, в шифровании. 
\end{remark}

\smallskip

\end{frame}

\begin{frame}
\frametitle{9.4.2. Таблицы расстановки (2/7)}
\odmkey{Хэш-функцией}{Функция (function)!хэш (hash)} $\Map{h}{U}{0.. m - 1}$ называется функция, отображающее множество ключей $U$ на индексы ячеек таблицы расстановки длины $m$.
%\textbf{$T:$  array $[0..m - 1]$}.
%of struct $\{ k: U; r : \Pointer R\}$}. 
Величина $h(k)$ называется \odmkey{хэш-значе\-нием}{Хэш (hash)} ключа $k$
(часто говорят просто <<\em{хэш $k$}>>). 

\begin{remark}
Индексы таблицы расстановки
принято нумеровать с нуля.
\end{remark}

Множество ключей актуальных записей обозначим
 $K$, $K\subset U$.  
Может случиться так,
что сужение функции $h\big|_K$ 
на множество $K$ является инъекцией.
Но по принципу Дирихле  
функция $h$ на всём множестве $U$
не может быть инъекцией, поскольку $u\gg m$.
Ситуация, при которой хэш-функция сопоставляет один и тот же хэш двум актуальным записям с различными ключами, называется  \odmkey{коллизией}{Коллизия (collision)}.

%\medskip
\begin{example}
Пусть всё множество ключей $U=0..19$, 
подмножество актуальных ключей \\
$K = \{0, 2, 3, 6, 8, 12, 15, 17 \}$,
$h(k)\Assign k \MOD 10$ и $m=10$.
Тогда массив $T$ имеет вид
\begin{center}
\includegraphics[width=120mm]{9_4_2_1.jpg}
\end{center}
и возникает коллизия, поскольку $h(2)=h(12)=2$.	
\end{example}

\end{frame}

\begin{frame}
\frametitle{9.4.2. Таблицы расстановки (3/7)}
%Так как $|U| > m$, полностью избежать коллизий невозможно, можно лишь минимизировать их количество. В любом случае, необходим метод разрешения возникающих коллизий. 
%
%\medskip
\begin{digress}
Подобрать инъективное сужение 
$h\big|_K$ в общем случае довольно трудно,
если только множество $K$ случайно.
При этом мала надежда на то, 
что коллизий не случится.
Например, известен 
\odmkey{парадокс дней рождения}{Парадокс (paradox)!дней рождения (birthday)},
состоящий в том, 
что если в комнате собрались не менее
23 человек, то с вероятностью 0,4927 у всех
собравшихся дни рождения различны.
Другими словами, уже при $n = 23, m=365$ 
вероятность коллизии оказывается более половины!  
\end{digress}

%Методы
%хэширования отличаются друг от друга способами разрешения коллизий
%и приёмами вычисления хэш-функций.
%
\smallskip
К счастью, существуют эффективные способы  разрешения конфликтов, вызываемых коллизиями. Здесь рассматриваются \em{метод цепочек} и \em{метод линейной открытой адресации}.
При анализе производительности методов используют \em{коэффициент заполнения} таблицы 
$\alpha \Assign n / m$. 


\smallskip
При разрешении коллизий по методу цепочек  все элементы, хэшированные в одну и ту же ячейку, объединяются в связный список. Ячейка таблицы расстановки содержит указатель на первую ячейку этого списка. Если таких элементов нет, список пуст, то ячейка таблицы расстановки содержит значение \textbf{nil}.


Поиск записи в такой таблице расстановки происходит следующим образом: после вычисления хеш-значения необходимо последовательно просмотреть список элементов. 
Вставка и удаление элементов выполняются как обычно в связных списках.

\end{frame}

\begin{frame}
\frametitle{9.4.2. Таблицы расстановки (4/7)}
\begin{example}
	В условиях предыдущего примера метод цепочек даёт такое решение.
	\begin{center}
		\includegraphics[width=120mm]{9_4_2_2.jpg}
	\end{center}	
\end{example}
Средняя длина одного списка равна 
$n/m$, то есть равна коэффициенту заполнения таблицы $\alpha$, поэтому доступ к записям оказывается эффективным, если хэш-функция достаточно равномерно распределяет ключи по ячейкам таблицы расстановки.    

\end{frame}

\begin{frame}
\frametitle{9.4.2. Таблицы расстановки (5/7)}
При использовании метода линейной открытой адресации все элементы хранятся непосредственно в таблице расстановки. 
Для выполнения вставки последовательно проверяются ячейки таблицы расстановки, начиная с той ячейки, на которую указывает хэш-функ\-ция, и до тех пор, пока не находится пустая ячейка.
При поиске элемента аналогично проверяются  ячейки таблицы, начиная с той ячейки, на которую указывает хэш-функция, пока не находится искомый элемент или пустая ячейка в случае его отсутствия. 


%%%%%%%%%%%%%%%%%%%%%%%%%%%%%%%%%%%%%%%%%%%%%%%%%
%  я не понял исходного замечания и написал своё
%  но выкинул его, потому что ботва 
%%%%%%%%%%%%%%%%%%%%%%%%%%%%%%%%%%%%%%%%%%%%

%\medskip
%\begin{remark}
%В методе открытой адресации, если просмотр в таблице расстановки делается только в одну строну, может не найтись пустой ячейки, что делает невозможным вставку новых элементов, хотя коэффициент заполнения $\alpha < 1$.
%\end{remark}
%\medskip

% В результате в хэш-функция включается номер исследования $i$:
%
%$$
%h: U \times \{0, 1, ..., m - 1\} \rightarrow \{0, 1, ..., m - 1\} 
%$$
%

\begin{remark}
Этот метод называется методом \em{линейной } открытой адресации, поскольку
просматриваются соседние ячейки в линейном порядке. Можно просматривать не обязательно соседние ячейки, а ячейки, 
задаваемые любой перестановкой
отрезка $0..m-1$ без неподвижных точек. При этом получаются другие методы открытой адресации,
которые также имеют право на существование. 
\end{remark}



\begin{example}
Работа метода открытой адресации в условиях предыдущих примеров.
\begin{center}
\includegraphics[width=120mm]{9_4_2_3.jpg}
\end{center}	
\end{example}

\end{frame}


\begin{frame}
\frametitle{9.4.2. Таблицы расстановки (6/7)}
В сильно заполненной таблице метод открытой адресации даёт несколько худшие результаты,
поскольку пустая ячейка может далеко отстоять 
от первого хеш-значения.

\medskip
Процедура удаления из таблицы расстановки с открытой адресацией вызывает некоторые трудности, так как при удалении записи недостаточно просто её ячейку пометить как свободную~---  процедура поиска после этого может работать некорректно!

\begin{example}
В условиях предыдущих примеров после удаления
записи с ключом 3 запись с ключом 12 не будет найдена, хотя и присутствует в таблице!
\begin{center}
\includegraphics[width=120mm]{9_4_2_4.jpg}
\end{center}	
\end{example}

\medskip
Одно из решений~--- ввести три состояния ячеек:
свободные, занятые и удалённые. При поиске 
удалённые ячейки пропускаются как занятые,
а при вставке новый ключ вставляется в первую
встретившуюся удалённую или свободную ячейку.

% помечать ячейку специальным значением  \textbf{DELETED}, при этом слегка изменив процедуру вставки, однако тогда время поиска перестанет зависеть от коэффициента заполнения $\alpha$.
%
%\medskip
%Преимущество открытой адресации заключается в неиспользование указателей. Вместо этого мы вычисляем последовательность проверяем ячеек, что позволяет создавать таблицу расстановки б$\acute{\text{о}}$льшего размера.
%
%\medskip
%Для получения последовательности исследования при открытой адресации обычно используются три метода, описанных в п. 9.4.5. 
%
\end{frame}

\begin{frame}
\frametitle{9.4.2. Таблицы расстановки (7/7)}
Эта простая идея работает только при небольшом количестве удалений. Если удалений (и вставок)
произведено достаточно много, то постепенно
все ячейки таблицы будут помечены как занятые 
и удалённые, поскольку однажды занятая ячейка никогда не может стать свободной.
Такая деградация таблицы приводит к тому, 
что поиск отсутствующего элемента 
требует просмотра всей таблицы. 

%\medskip
В случае линейной открытой адресации 
решить проблему удаления можно следующим алгоритмом.
Если ячейка, следующая за удаляемой, уже свободна, то просто удалить ключ и прекратить работу.
Иначе найти первую среди следующих занятых ячеек такую ячейку,
в которой хэш ключа не совпадает с индексом и при этом меньше хеша удалённого ключа.
Если такая ячейка найдётся (это источник ошибки при следующем поиске!), то переместить
из найденной ячейки ключ в удаляемую ячейку, а к найденной ячейке применить алгоритм удаления.

\begin{example}
В условиях предыдущих примеров результат удаления
записи с ключом 3.

\vspace{-4pt}
\begin{center}
\includegraphics[width=120mm]{9_4_2_5.jpg}
\end{center}	
\end{example}



\end{frame}

\begin{frame}
\label{9.4.3.}\subsubsection{9.4.3. Эффективность хэширования}
\frametitle{9.4.3. Эффективность хэширования (1/5)}

%\smallskip 
%
%Пункт 9.4.5 посвящен вопросу о составлении и выборе хэш-функций. 
%
%\medskip
%\medskip
%Оценки сложности в среднем имеют вероятностный характер и представлены в пункте 9.4.6.

Эффективность хэширования во многом определяется свойствами выбранной 
хэш-функ\-ции.
В реальных данных часто встречается много близких
ключей, поэтому желательно выбрать 
хэш-функцию так, чтобы близким ключам соответствовали
далёкие друг от друга хэш-значения, и при этом 
хэш-значения были распределены по отрезку 
$0..m-1$ равномерно.

%\smallskip
\begin{remark}
На этом основании 
таблицы расстановки или хэш-таблицы раньше называли 
\em{рассеянными таблицами}.
\end{remark}

Из великого множества предложенных
способов построения хэш-функций 
чаще всего реализуют два сравнительно простых
метода: 1) \em{метод деления}; 2) \em{метод умножения}. 

%\begin{digress}
%Описание построения универсального множества не включено в курс из соображений экономии места. В методе строится  множество хэш-функций, из которого используемая функция выбирается случайным образом, не зависящим от того, с каким ключом ей предстоит работать. Рандомизация гарантирует, что одни и те же входные данные не могут постоянно давать наихудшее поведение алгоритма. 
%
%\medskip
%Таким образом, данный подход позволяет избежать ситуаций, когда злоумышленник, имея информацию о виде хэш-функции, может специально подбирать такие ключи, которые будут хэшироваться в одну и ту же ячейку.

\smallskip
Построение хэш-функций методом деления состоит в
%отображении ключа $k$ в одну из ячеек путем
получении остатка от деления ключа $k$ на число $m$:

\vspace{-14pt}
$$
h(k) \Assign k \MOD m.
$$

\begin{remark}
Зачастую хорошие результаты можно получить, выбирая в качестве значения $m$ простое число,
причём так, чтобы $r^s \not\equiv \pm a \mod m$, где
$r$~--- основание системы счисления, 
а $s$ и $a$~--- небольшие натуральные числа.
Например, для таблиц до тысячи элементов хороший 
результат даёт число $m= 1009$. 
%Также стараются избегать значений $m$, являющихся степенью 2 и 10~--- лучше строить хэш-функцию так, чтоб ее результат зависел от всех битов ключа.	
\end{remark}


\end{frame}

\begin{frame}
\frametitle{9.4.3. Эффективность хэширования (2/5)}


Построение хэш-функций методом умножения выполняется в четыре этапа:
1)~умножение ключа $k$ на константу $0 < A < 1$;
2)~взятие дробной части от результата; 
3)~умножение полученного значения на $m$; 
4)~взятие целой части от результата:
%умножением $k$ на константу $0 < A < 1$ со взятием дробной части от результата и умножением полученного значения на $m$ с последующим округлением числа в меньшую сторону:
$$
h(k) \Assign \lfloor m
\left( kA - \lfloor kA \rfloor \right) \rfloor.
$$

\begin{remark}
Рекомендуется использовать в качестве $A$ золотое сечение: $A \Assign(\sqrt{5}-1)/2$ (\DMref{п.~5.7.3}{5.7.3.}).	
\end{remark}
\medskip
Достоинство метода умножения в том, что выбор числа $m$ перестает быть критичным.

\medskip
Если ключами являются не числа, имеющие 
одинаковую длину машинного представления, а произвольные данные различной длины,
то обычно их приводят к ключам одинаковой
длины. Для этого битовое представление ключа делят на блоки фиксированной длины 
и блоки складывают по модулю 2.

\begin{digress}
Одним из интересных методов является 
\em{метод универсального хеширования}. В методе строится  семейство хэш-функций, из которого используемая функция выбирается случайным образом. Рандомизация гарантирует, что построенные хэши имеют равномерное распределение
для любых исходных данных. 
\end{digress}
  
\end{frame}

%\begin{frame}
%\frametitle{9.4.3. Эффективность хэширования (3/5)}
%При использовании открытой адресации для вычисления последовательности исследований обычно используются методы:
%\vspace{-1pt}
%\begin{enumerate}
%	\item Основанные на \em{линейном исследовании}. 
%	
%	\vspace{-1pt}
%	\item Основанные на \em{квадратичном исследовании}. 
%	
%	\vspace{-1pt}
%	\item Метод \em{двойного хэширования}.
%\end{enumerate}	
%\medskip
%Эти методы гарантируют, что для каждого ключа $k$ последовательность исследований $\langle h(k,0), h(k,1), ..., h(k,m-1)\rangle$ представляет собой перестановку множества $\langle0,1,...,m-1\rangle$, чтобы в конечном счете могли быть просмотрены все ячейки таблицы.
%
%\medskip
%Однако эти методы не удовлетворяют предположению о равномерном хэшировании, так как не в состоянии сгенерировать необходимые $m!$ различных последовательностей исследований.
%\end{frame}
%
%\begin{frame}
%\frametitle{9.4.3. Эффективность хэширования (4/5)}
%\medskip
%Метод линейного исследования использует хэш-функцию 
%$$
%h(k,i) = (h'(k) + i) \: mod \: m
%$$
%где $h': U \rightarrow \{0, 1, ..., m - 1\} $ — вспомогательная функция.
%
%\medskip
%Метод легко реализуется, но имеет проблему появления длинных последовательностей занятых ячеек, что увеличивает среднее время поиска.
%
%\medskip
%Квадратичное исследование использует хэш-функцию вида
%$$
%h(k,i) = (h'(k) + c_1i+c_2i^2) \: mod \: m \qquad c_1^2+c_2^2 \neq 0
%$$
%\medskip
%Метод существенно лучше первого, но для того, чтоб он обхватывал все ячейки, необходим выбор специальных значений $c_1$, $c_2$  и $m$.
%
%\medskip
%Однако, поскольку в обоих методах исследуемая ячейка однозначно определяет всю последовательность исследований целиком, в них имеется только $m$ различных последовательностей.
%\end{frame}
%
%\begin{frame}
%\frametitle{9.4.3. Эффективность хэширования (5/5)}
%\medskip
%Двойное хэширование представляет собой один из наилучших способов использования открытой адресации, поскольку получаемые при этом перестановки обладают многими характеристиками случайных перестановок.
%$$
%h(k,i) = (h_1(k) + ih_2(k)) \: mod \: m 
%$$
%где $h_1$, $h_2$ — вспомогательные хэш-функции.
%\medskip
%\begin{remark}
%	Для того, чтоб метод мог обхватить всю таблицу, значение $h_2 (k)$ должно быть взаимно простым с размером таблицы расстановки размером $m$. Для этого можно выбрать в качестве $m$ простое число, а вспомогательные хэш-функции следующими:
%	$$
%	h_1 (k)=k\:mod\:m
%	$$
%	$$
%	h_2 (k)=1+(k\:mod\:m')
%	$$
%\end{remark}
%\begin{remark}
%	Так как каждая пара $(h_1 (k),h_2 (k))$ дает свою, отличную от других последовательность исследований, двойное хэширование имеет $m^2$ последовательностей исследования.
%\end{remark}
%\end{frame}

\begin{frame}
\frametitle{9.4.3. Эффективность хэширования (3/5)}
Трудёмкость выполнения операций в различных методах хеширования оценивается в количестве
сравнений, которые необходимо выполнить, чтобы получить доступ к записи по ключу.

\smallskip
Здесь приводятся оценки в предположении,
что хэш-значения поступающих ключей распределены 
по отрезку $0..m-1$ \em{равномерно}.

\smallskip
\begin{remark}
Предположение о 
равномерности распределения
правомерно,
поскольку многочисленные прямые вычислительные эксперименты на реальных данных
и фактически реализованных хэш-функциях
показывают распределения очень близкие к равномерным. 
\end{remark}

%\medskip
\begin{theorem}
	В таблице расстановки с разрешением коллизий методом цепочек математическое ожидание количества сравнений при поиске не превосходит $1+\alpha/2$. 
%в предположении, что все элементы хэшируются по ячейкам равномерно и независимо, 
%равно в среднем $O(1+\alpha)$.
\end{theorem}
\begin{proof}
	Докажем оценку в среднем для неудачного и удачного поиска отдельно.
\proofsection{нет ключа} 
Метод цепочек устанавливает отсутствие ключа за
одно сравнение (значение \textbf{nil} в ячейке). 	
\proofsection{ключ есть} 
Математическое ожидание количества элементов, проверяемых в процессе успешного поиска, равно половине средней длины списка. 
%Количество элементов, проверяемых в процессе успешного поиска элемента $x$, на 1 больше, чем количество элементов, находящихся в списке перед $x$. Тогда математическое ожидание числа проверяемых элементов равно:
%	$$
%	1/\alpha * \sum\limits_{i=1}^\alpha i = 1 / \alpha * (\alpha(\alpha+1))/2 = (1 + \alpha)/2
%	$$
%	Включая время на вычисление хэш-функции, полное время для проведения успешного поиска составляет $O(3/2+\alpha/2) = O(1+\alpha)$.
\end{proof}
\end{frame}

\begin{frame}
\frametitle{9.4.3. Эффективность хэширования (4/5)}
%\medskip
%Таким образом, если количество ячеек в таблице расстановки, как минимум, пропорционально количеству элементов, хранящихся в ней, то $n=O(m)  \Longrightarrow \alpha=O(m) /m= O(1)$, а значит поиск требует постоянного времени. Поскольку при использовании двухсвязного списка операции вставки и удаления в худшем случае занимают $O(1)$, все ассоциативные операции в среднем выполняются за константное время.
%\medskip

Таким образом, математическое ожидание $E$ трудоёмкости поиска в таблице расстановки с цепочками не зависит от размера 
таблицы и $1\le E \le 1,5$. 

\begin{theorem}
В таблице расстановки с линейной открытой адресацией математическое ожидание количества сравнений при вставке нового элемента  не превышает 
$1/(1-\alpha)$.
\end{theorem}
\begin{proof}
%	При неуспешном поиске каждая последовательность исследований завершается на пустой ячейке. В предположении равномерного хэшировании не возникает проблемы кластеризации — создания длинных последовательностей занятых ячеек. Имеет место распределение Бернулли.
В предположении равномерного кэширования
существует $C(m,n)$ вариантов заполнения 
таблицы, и все они равновероятны.
Для вставки нового элемента потребуется ровно $i$
сравнений, если среди исследуемых
ячеек ровно $i-1$ заняты и ещё одна свободна.
Таких конфигураций $C(m-i,n-(i-1))$.
Значит вероятность вставки за 
$i$ сравнений
$
P_i = C(m-i,n-i+1)/C(m,n)
$.
Отсюда математическое ожидание
\begin{align*}
E_n &=\sum_{i=1}^m iP_i = m\!+\!1-\!\sum_{i=1}^m 
(m\!+\!1\!-\!i)P_i =
m\!+\!1-\!\sum_{i=1}^m (m\!+\!1\!-\!i)\frac{C(m-i,m-n-1)}{C(m,n)} = \\
&=m\!+\!1-\!\sum_{i=1}^m (m-n)\frac{C(m\!+\!1\!-\!i,m\!-\!n)}{C(m,n)} = 
m\!+\!1-(m\!-\!n)\frac{C(m+1,m-n+1)}{C(m,n)} = \\
&=m\!+\!1\!-\!(m\!-\!n)(m\!+\!1)/(m\!-\!n\!+\!1) = 
(m\!+\!1)/(m\!-\!n\!+\!1)
\le m/(m\!-\!n)= 1/(1\!-\!\alpha).
\end{align*}
 
  
%для каждой исследуемой ячейки таблицы 
%расстановки вероятность быть занятой равна   равна $\alpha$, 
%а математическое ожидание числа исследований равно $1/(1-\alpha)$.
\vspace{-6pt}
\end{proof}

\end{frame}

\begin{frame}
\frametitle{9.4.3. Эффективность хэширования (5/5)}

\begin{corollary}
	Неуспешный поиск в таблице расстановки с открытой адресацией в предположении равномерного хеширования требует в среднем не более $1/(1-\alpha)$ сравнений.
\end{corollary}
\begin{proof}
При неуспешном поиске проводятся в точности те же сравнения, что и при вставке.
\end{proof}
\begin{theorem}
%	Математическое ожидание количества исследований при удачном поиске в таблице расстановки с открытой адресацией в предположении равномерного хэширования и равновероятного поиска любого из ключей не превышает $\frac{1}{\alpha} ln \frac{1}{1-\alpha}$.
%
	В таблице расстановки с разрешением коллизий методом  линейной открытой адресации математическое ожидание количества сравнений при успешном поиске не превосходит 
$\displaystyle{\frac{1}{\alpha} \ln \frac{1}{1-\alpha}}$. 	
\end{theorem}
\begin{proof}
	Поиск ключа выполняется с той же последовательностью исследований, что и его вставка. Вставка каждого нового ключа увеличивает коэффициент заполнения таблицы. Следовательно, математическое ожидание числа сравнений при успешном поиске 
$$ E= \frac{1}{n}\sum_{i=1}^n E_i =
	 \frac{1}{n}\sum\limits_{i=1}^{n}
	  \frac{1}{1-i/m} \approx \frac{1}{\alpha}\int\limits_1^\alpha \frac{1}{1-x}dx = \frac{1}{\alpha} \ln \frac{1}{1-\alpha}
	$$
\end{proof}
\end{frame}



\begin{frame}
\label{9.4.4.}\subsubsection{9.4.4. Алгоритм поиска в дереве сортировки}
\frametitle{9.4.4. Алгоритм поиска в дереве сортировки (1/2)}
Обратимся теперь к деревьям сортировки (поиска).
\medskip
\begin{remark}
Симметричный обход дерева поиска перечисляет множество узлов
в порядке возрастания значения ключа, то есть определяет
строгий линейный порядок на множестве узлов. 
Таким образом, любое конечное линейно упорядоченное множество 
можно представить в виде дерева поиска. Польза 
от такого представления в том, что это позволяет производить многие операции с множеством (в частности, поиск элемента) за время $O(h)$, 
где $h$~--- высота дерева.
В частности, в предельном случае, когда  $h=\log n$, поиск с помощью сортировки
работает также быстро, как двоичный поиск.
\end{remark}


\medskip
Следующий алгоритм находит в дереве сортировки
узел с указанным ключом, если он там есть.
\end{frame}




\begin{frame}
\frametitle{9.4.4. Алгоритм поиска в дереве сортировки (2/2)}
%\vspace{-4pt}
{\renewcommand{\baselinestretch}{0.98}\selectfont
\begin{algorithmic}
\REQUIRE
дерево сортировки $T$, заданное указателем на корень; ключ  $a\Is key$.
\ENSURE
указатель $p$ на найденный узел или \textbf{nil}, если в дереве нет такого
ключа.
\STATE $p \Assign  T$ 
\COMMENT{\color{gray}\small указатель на проверяемый узел }
\WHILE{\textbf{$p\ne$ nil}}
  \IF{$a<p.i$}
    \STATE $p \Assign  p.l$ 
    \COMMENT{\color{gray}\small  продолжаем поиск слева }
  \ELSIF{$a>p.i$}
    \STATE $p\Assign  p.r$ 
    \COMMENT{\color{gray}\small  продолжаем поиск справа }
   \ELSE
     \STATE \textbf{return $p$} 
     \COMMENT{\color{gray}\small  нашли узел }
  \ENDIF
\ENDWHILE
\STATE \textbf{return nil}
\COMMENT{\color{gray}\small  искомого ключа нет в дереве }
\end{algorithmic}
%\end{algorithm}

\begin{ground}
Этот алгоритм работает в точном соответствии с определением дерева
сортировки: если текущий узел не искомый, то в зависимости от
того, меньше или больше искомый ключ по сравнению с текущим, нужно
продолжать поиск слева или справа.
\end{ground}
}
\end{frame}

\begin{frame}
\label{9.4.5.}\subsubsection{9.4.5. Алгоритм вставки в дерево сортировки }\frametitle{9.4.5. Алгоритм вставки в дерево сортировки (1/4)}

%{\renewcommand{\baselinestretch}{0.93}\selectfont
Один из возможных способов вставки в дерево сортировки
узла с указанным ключом описан в следующем
\odmkey{алгоритме вставки узла}{Алгоритм (algorithm)!вставки узла в дерево сортировки (node insertion in sorting tree)}.

\begin{remark} Необходимо отметить два важных обстоятельства.\\
{\color{blue}\bf 1.} В дереве уже может быть узел с указанным ключом. Требуется определить, 
что надлежит делать в этом случае: выдавать сообщение об ошибке,
менять старую запись на новую с тем же ключом и т.~д.
В данной реализации ничего не делается.\\
{\color{blue}\bf 2.} Алгоритм хорошо работает,
если дерево сортировки <<динамичное>>,
то есть в него постоянно добавляются и удаляются ключи,
причём значения ключей случайны и распределены равномерно.
Если же значения ключей меняются регулярно,
например, если ключи всё время возрастают,
то алгоритм применять не следует. 
\end{remark}

\begin{ground}
Алгоритм вставки, в сущности, аналогичен алгоритму поиска:
в~дереве ищется такой узел, имеющий свободную связь для
подцепления нового узла, чтобы не нарушалось условие дерева
сортировки. А именно, если новый ключ меньше текущего,
то либо его можно подцепить слева (если левая связь свободна),
либо нужно найти слева подходящее место. Аналогично,
если новый ключ больше текущего.
\end{ground}
%}
\end{frame}

\begin{frame}
\frametitle{9.4.5. Алгоритм вставки в дерево сортировки (2/4)}
%{\renewcommand{\baselinestretch}{0.95}\selectfont
\begin{algorithmic}
\REQUIRE дерево сортировки $T$, заданное указателем на корень;
ключ  $a\Is key$.
\ENSURE
модифицированное дерево сортировки $T$.
\STATE\textbf{if $T =$ nil then $T \Assign  \NewNode(a);$ return $T$ end if}
  \COMMENT{\color{gray}\small новый узел}
\STATE $ p \Assign T$
\COMMENT{\color{gray}\small указатель на текущий узел }
\WHILE{\textbf{true}}
  \IF{$ a<p.i$}
  	\STATE{}\COMMENT{\color{gray}\small  создаём новый узел слева от $p$}
    \STATE{\textbf{if $p.l =$ nil then
           $p.l \Assign  \NewNode(a);$ 
           return $T$ else $p \Assign  p.l $ 
           end if}}
  \ELSIF{$ a>p.i$}
    \STATE{}\COMMENT{\color{gray}\small  создаём новый узел справа от $p$}
     \STATE{\textbf{if $p.r =$ nil then
      $p.r \Assign  \NewNode(a);$ 
     return $T$ else $p \Assign  p.r $ end if}}
 \ELSE
    \STATE \textbf{return $T$}
    \COMMENT{\color{gray}\small $a=p.i$~--- в дереве уже есть такой ключ }
  \ENDIF
\ENDWHILE
\end{algorithmic}
%}
\end{frame}


\begin{frame}
\frametitle{9.4.5. Алгоритм вставки в дерево сортировки (3/4)}
\begin{digress}
В последнее время в среде профессиональных программистов
наблюдается противоестественное и достойное сожаления явление~---
\odmkey{обфускация}{Обфускация (obfuscation)} программ, то есть
искусственное составление текстов программ таким образом, чтобы
сделать их как можно более непонятными для читателя. Обычно
обфускация обосновывается необходимостью защиты интеллектуальной
собственности, хотя на самом деле, по мнению автора, часто
является проявлением комплекса неполноценности, развивающегося на
фоне недостаточного уровня профессиональной подготовки. Для
иллюстрации последнего тезиса приведём пример того, как \em{не
надо} писать программы, проведя обфускацию
алгоритма вставки в дерево сортировки.

\medskip
Эта программа примерно вдвое короче, чуть медленнее (доступ к
элементу массива обычно медленнее доступа к полю структуры),
использует неявное приведение \textbf{false $\mapsto 0,$
true~$\mapsto~1$}, и её трудно понять без дополнительных
пояснений.
\end{digress}

\end{frame}


\begin{frame}
\frametitle{9.4.5. Алгоритм вставки в дерево сортировки (4/4)}
\begin{algorithmic}
\REQUIRE дерево сортировки, заданное указателем на корень
$T\Is\Pointer N$, \\где \textbf{$N = $ struct $\{ i\Is key;\;
l\Is$ array $[0..1]$ of $\Pointer N\}$ };
ключ  $a\Is key$.
\ENSURE
модифицированное дерево сортировки $T$.
\STATE{\textbf{if $T =$ nil then $T \Assign  \NewNode(a);$ return $T$ end if}}
\STATE $ p \Assign T$
\WHILE{\textbf{true}}
    \STATE \textbf{if $a=p.i$ then return $T$ end if}
  \STATE $b\Assign a>p.i$
  \IF{\textbf{$p.l[b] =$ nil}}
      \STATE $p.l[b] \Assign  \NewNode(a)$
      \STATE \textbf{return $T$}
  \ELSE
      \STATE $p \Assign  p.l[b] $
  \ENDIF
\ENDWHILE
\end{algorithmic}
%\addvspace{-3pt}
\end{frame}

\begin{frame}
\label{9.4.6.}\subsubsection{9.4.6. Алгоритм удаления из дерева сортировки }\frametitle{9.4.6. Алгоритм удаления из дерева сортировки (1/4)}

На следующем слайде приведён 
\odmkey{алгоритм удаления узла из дерева сортировки}{Алгоритм (algorithm)!удаления узла из дерева сортировки (node removal from sorting tree)}. Если узла с указанным ключом
нет в дереве, то ничего не делается. Вспомогательные процедуры
$\Find$ и $\Delete$ описаны в следующем параграфе.

\medskip
\begin{remark}
В замечательной книге Дональда Кнута, из которой заимствован данный
алгоритм, показано, что, хотя алгоритм <<выглядит несимметричным>>
(правые и левые связи обрабатываются по-разному), на самом деле
в среднем характеристики дерева сортировки не искажаются.
\end{remark}

\end{frame}

\begin{frame}
%\frametitle{9.4.6. Алгоритм удаления из дерева сортировки (2/4)}

{\renewcommand{\baselinestretch}{0.90}\selectfont%\small
%\addvspace{-3pt}
\begin{algorithmic}
\REQUIRE дерево сортировки $T$, заданное указателем на корень; ключ  $a\Is key$.
\ENSURE
модифицированное дерево сортировки $T$.
\STATE  $\Find(T,a,p,q,s)$ \COMMENT{ поиск удаляемого узла }
\STATE{\textbf{if $p = $ nil then return $T$ end if}}
\COMMENT{\color{gray}\small   нет такого узла~--- ничего делать не нужно }
\IF{\textbf{$p.r =$ nil}}
  \STATE $\Delete(p,q,p.l,s)$
  \COMMENT{\color{gray}\small   случай 1, рисунок {\it слева }}
\ELSE
  \STATE $u \Assign  p.r $
  \IF{\textbf{$ u.l = $ nil}}
    \STATE $u.l \Assign  p.l$
    \STATE $\Delete(p,q,u,s)$ \
    \COMMENT{\color{gray}\small   случай 2, рисунок {\it в центре} }
  \ELSE

    \STATE $w \Assign  u;  v \Assign  u.l$
    \STATE{\textbf{while $ v.l\ne$ nil do 
            $w \Assign  v ; v \Assign  v.l$ end while}}
    \STATE $p.i \Assign  v.i $
    \STATE $\Delete(v,w,v.r,-1)$
    \COMMENT{\color{gray}\small   случай 3, рисунок {\it справа} }
  \ENDIF
\ENDIF
\STATE\textbf{return $T$}
\end{algorithmic}
}

\end{frame}


\begin{frame}
\frametitle{9.4.6. Алгоритм удаления из дерева сортировки (3/4)}
\begin{figure}[h]
\begin{center}
\includegraphics[scale=0.85]{odm09-11}
\end{center}
%\caption{Иллюстрация к алгоритму удаления узла из дерева сортировки}
%\label{f0911}
\end{figure}

\vspace{-\baselineskip}
\begin{ground}
Удаление узла производится перестройкой дерева сортировки.
При этом возможны три случая (не считая тривиального случая,
когда удаляемого узла нет в дереве и ничего делать не нужно).
\proofsection{1} Правая связь удаляемого узла $p$ пуста (см. рисунок {\it слева}). В этом случае левое поддерево 1
узла $p$ подцепляется к родительскому узлу $q$ с той же стороны, с
которой был подцеплен узел $p$. Условие дерева сортировки,
очевидно, выполняется. 
\end{ground}


\end{frame}


\begin{frame}
\frametitle{9.4.6. Алгоритм удаления из дерева сортировки (4/4)}
\begin{ground} (продолжение)
\proofsection{2} Правая связь удаляемого
узла $p$ не пуста и ведёт в узел $u$, левая связь которого пуста
(см. рисунок {\it  в центре}). В этом случае левое
поддерево 1 узла $p$ подцепляется к узлу $u$ слева, а сам узел $u$
подцепляется к родительскому узлу $q$ с той же стороны, с которой
был подцеплен узел $p$. Нетрудно проверить, что условие дерева
сортировки выполняется и в этом случае. 
\proofsection{3} Правая
связь удаляемого узла $p$ не пуста и ведёт в узел $u$, левая связь
которого не пуста. Поскольку дерево сортировки конечно, можно
спуститься от узла $u$ до узла $v$, левая связь которого пуста
(см. рисунок {\it справа}). В этом случае выполняются
два преобразования дерева. Сначала информация в узле $p$
заменяется информацией узла $v$. Поскольку узел $v$ находится в
правом поддереве узла $p$ и в левом поддереве узла $u$, имеем
$p.i<v.i<u.i$. Таким образом, после этого преобразования условие
дерева сортировки выполняется. Далее правое поддерево 4 узла $v$
подцепляется слева к узлу $w$, а сам узел $v$ удаляется. Поскольку
поддерево 4 входило в левое поддерево узла $w$, условие дерева
сортировки также сохраняется.
\end{ground}


\end{frame}

\begin{frame}
\label{9.4.7.}\subsubsection{9.4.7. Вспомогательные алгоритмы для~дерева~сортировки }\frametitle{9.4.7. Вспомогательные алгоритмы для~дерева~сортировки (1/3)}

%{\renewcommand{\baselinestretch}{0.92}\selectfont\small
Алгоритмы двух предыдущих параграфов
используют вспомогательные функции:

{\NB 1.} Поиск узла~--- функция $\Find$. 
\begin{algorithmic}
\REQUIRE
дерево сортировки $T$, заданное указателем на корень; ключ  $a\Is key$.
\ENSURE
$p$~--- указатель на найденный узел или \textbf{nil}, если в дереве нет такого
ключа;
$q$~--- указатель на отца узла $p$;
$s$~--- способ подцепления узла $p$ к узлу $q$
($s=-1$, если $p$ слева от $q$;
$s=+1$, если $p$ справа от $q$; $s=0$, если $p$~--- корень).
\STATE $p \Assign  T; q \Assign  \textbf{nil}; s \Assign  0$
\COMMENT{\color{gray}\small инициализация }
\WHILE {\textbf{$p\ne$ nil}}
  \STATE\textbf{if $p.i = a$ then return $p,q,s$ end if}
  \STATE $q \Assign  p$ 
  \COMMENT{\color{gray}\small сохранение значения $p$ }
  \IF{$a<p.i$}
    \STATE $p \Assign  p.l; s \Assign  -1 $ 
    \COMMENT{\color{gray}\small  поиск слева }
  \ELSE \pagebreak[3]
    \STATE $p \Assign  p.r; s \Assign  +1 $ 
    \COMMENT{\color{gray}\small  поиск справа }
  \ENDIF
\ENDWHILE
\STATE \textbf{return $p,q,s$}
\end{algorithmic}
%}

\end{frame}

\begin{frame}
\frametitle{9.4.7. Вспомогательные алгоритмы для~дерева~сортировки (2/3)}

\begin{digress}
В этой простой функции ст$\acute{\text{о}}$ит обратить внимание на
использование указателя $q$, который отслеживает значение
указателя $p$ <<с запаздыванием>>, то~есть указатель $q$
<<помнит>> предыдущее значение указателя $p$. Такой приём полезен
при обходе однонаправленных структур данных, в которых невозможно
вернуться назад.
\end{digress}


{\NB 2.} Удаление узла~--- процедура $\Delete$.

%{\renewcommand{\baselinestretch}{0.92}\selectfont
\begin{algorithmic}
\REQUIRE $p1$~--- указатель на удаляемый узел;
$p2$~--- указатель на подцепляющий узел;
$p3$~--- указатель на подцепляемый узел;
$s$~--- способ подцепления.
\ENSURE преобразованное дерево.
\IF{$s = -1$}
    \STATE $p2.l \Assign  p3$ \COMMENT{\color{gray}\small   подцепляем слева }
\ENDIF
\IF{$s = +1$}
    \STATE $p2.r \Assign  p3$ \COMMENT{\color{gray}\small   подцепляем справа }
\ENDIF
\STATE $\Dispose(p1)$ \COMMENT{\color{gray}\small   удаляем узел }
\end{algorithmic}
%}

\end{frame}

\begin{frame}
\frametitle{9.4.7. Вспомогательные алгоритмы для~дерева~сортировки (3/3)}

{\NB 3.} Создание нового узла~--- конструктор $\NewNode$.
\begin{algorithmic}
\REQUIRE ключ $a$.
\ENSURE
указатель $p$ на созданный узел.
\STATE $\New(p);
p.i \Assign a;
p.l \Assign \textbf{nil};
p.r \Assign \textbf{nil}$
\STATE\textbf{return $p$}
\end{algorithmic}

\end{frame}


\begin{frame}
\label{9.4.8.}\subsubsection{9.4.8. Сравнение представлений ассоциативной памяти }\frametitle{9.4.8. Сравнение представлений ассоциативной памяти (1/2)}


Пусть $n$~--- количество элементов в
ассоциативной памяти. Рассмотрим:


%{\renewcommand{\baselinestretch}{0.7}\selectfont
\begin{columns}[t]
\begin{column}[T]{6cm} 
{\color{blue}\bf 1)} неупорядоченный массив;

{\color{blue}\bf 2)} упорядоченный массив;

{\color{blue}\bf3)} таблицу расстановки;
\end{column}
\begin{column}[T]{6cm} 

{\color{blue}\bf 4)} неупорядоченный список; 

{\color{blue}\bf 5)} упорядоченный список;

{\color{blue}\bf 6)} дерево сортировки. 

\end{column}
\end{columns}

\medskip
В таблице указаны оценки трудоёмкости операций  для 
представлений.
%\addvspace{6pt}

\begin{center}
\begin{tabular}{|l| c c c c c c |}
\hline
         & 1)      & 2) & 3) & 4) & 5)  &  6)          \\
\hline
Добавить & $O(1)$ & $O(n)$        &  $O(1)$ & $O(1)$ & $O(n)$ & $O(h)$ \\
Найти    & $O(n)$ & $O(\log n)$ &  $O(1)$ & $O(n)$ & $O(n)$  & $O(h)$ \\
Удалить  & $O(n)$ & $O(n)$        &  $O(1)$ & $O(n)$ & $O(n)$  & $O(h)$ \\
\hline
\end{tabular}
\end{center}
\end{frame}


\begin{frame}
\frametitle{9.4.8. Сравнение представлений ассоциативной памяти (2/2)}
\begin{remark} Необходимо сделать два существенных уточнения.

{\color{blue}\bf 1}.~Для таблиц расстановки указаны 
условные оценки в среднем, которые выполняются только при определённых
ограничениях. В худшем случае оценки для таблиц расстановки 
безнадёжно плохие. 

{\color{blue}\bf 2}.~Трудоёмкость операций с деревом сортировки
ограничена сверху высотой дерева, которая может меняться в пределах
от $\log n$ до $n$.
Дерево сортировки может расти неравномерно. Например, если при загрузке
дерева исходные данные уже упорядочены, то полученное дерево будет право-
или леволинейным и окажется даже менее эффективным, чем неупорядоченный массив.
\end{remark}
%}

%--------------------------------------------------------
% TO DO параграф подлежит пересмотру.
% Надобно рассмотреть лучшие и худшие случаи,
% верхние и нижние оценки для каждого представления и для каждой операции,
% а потом обсудить, как можно сравнивать
% представления, если заданы соотношения частот выполнения операций
% + примеры из жизни (малый массив, квазистационарный справочник). ФАН
%-------------------------------------------------------- 

%}%baselinestretch
%\pagebreak[3]
\end{frame}

\subsection{9.5. Специальные деревья}
\label{9.5.}
\begin{frame}
\frametitle{9.5. Специальные деревья (1/2)}
Высота дерева сортировки является критическим параметром эффективности.
В этом разделе обсуждаются методы уменьшения
высоты дерева сортировки, а также использование 
других структур данных, родственных бинарным деревьям,
для реализации функций ассоциативной памяти.
\end{frame}

\begin{frame}
\frametitle{9.5. Специальные деревья (2/2)}
\medskip
\begin{tabular}{p{9mm}|p{70mm}|p{64mm}}
  № & Название параграфа
  & Ключевые термины и обозначения \\
  \hline
\DMref{9.5.1.}{9.5.1.} & Выровненные и полные деревья 
  & 
 Выровненное, заполненное, полное бинарное дерево \\
\DMref{9.5.2.}{9.5.2.}\RS & Сбалансированные деревья
   & 
  АВЛ-дерево, подровненное, сбалансированное по высоте, cбалансированное по весу дерево сортировки\\
\DMref{9.5.3.}{9.5.3.} & Балансировка деревьев & Правило простого вращения, правило двойного вращения
\\
\DMref{9.5.4.}{9.5.4.}  & Двоичные кучи & 
\\
\DMref{9.5.5.}{9.5.5.} & Восстановление свойства кучи & 
\\
\DMref{9.5.6.}{9.5.6.}  & Реализация операций кучи & \\
\DMref{9.5.7.}{9.5.7.} & Построение графа по степенному ряду & 
 Откладываение вершины, очередь с приоритетом\\
\DMref{9.5.8.}{9.5.8.}  & Дерево отрезков & 
 Отрезок массива, предработка \\
\DMref{9.5.9.}{9.5.9.}  & Оценки эффективности дерева отрезков & 
 \\
\end{tabular}

\end{frame}

\begin{frame}
\label{9.5.1.}\subsubsection{9.5.1. Выровненные и полные деревья }\frametitle{9.5.1. Выровненные и полные деревья (1/3)}

Бинарное дерево называется \odmkey{выровненным}{Дерево
(tree)!выровненное (aligned tree)}, если все листья находятся на
одном (последнем) уровне. В выровненном дереве все ветви имеют
одну длину, равную высоте дерева, однако выровненное дерево не
всегда является эффективным деревом сортировки.

\begin{example}
Дерево, состоящее из корня и двух поддеревьев, леволинейного
и праволинейного, ведущих к двум листьям, является выровненным,
однако такое дерево имеет высоту $(p-1)/2$.
\end{example}

Бинарное дерево называется \odmkey{заполненным}{Дерево
(tree)!заполненное (thick)}, если все узлы, степень которых меньше
2, располагаются на одном или двух последних уровнях. Другими
словами, в~заполненном дереве $T$ все ярусы $D(r,i)$, кроме, может
быть, последнего, заполнены:


\begin{columns}[t]
\begin{column}[T]{8cm}
$$
\A{i\in 0..(h(T)-1)}{\vert D(r,i)\vert=2^i}.
$$
\begin{example}
На рисунке приведены диаграммы заполненного ({\it слева})
и незаполненного ({\it справа}) деревьев.
\end{example}
\end{column}
\begin{column}[T]{7cm}
\vspace{-12pt}
\begin{figure}[h]
\begin{center}
\includegraphics[scale=0.8]{odm09-12}
\end{center}
\end{figure}
\end{column}
\end{columns}

\end{frame}

\begin{frame}
\frametitle{9.5.1. Выровненные и полные деревья (2/3)}

%\begin{lemma}
%$\sum^k_{i=0}2^i=2^{k+1}-1.$
%\end{lemma}
%\begin{proof}
%По формуле для суммы геометрической прогрессии $\sum^{k}_{i=0}
%2^i = 1 + 2 + \ldots + 2^k = \frac{2^{k+1}-1}{2-1}= 2^{k+1}-1$.
%\end{proof}

Заполненное
дерево имеет наименьшую возможную для данного $p$ высоту~$h$.

\medskip
\begin{theorem}
Для заполненного бинарного дерева
$
\log(p+1)-1 \le h < \log(p+1).
$
\end{theorem}
\begin{proof}
На $i$-м уровне может быть самое большее $2^i$ вершин,
следовательно, $\sum^{h-1}_{i=0}2^i<p\leq \sum^{h}_{i=0}2^i$. 
Значит $2^h<p+1\le 2^{h+1}$, и, логарифмируя, имеем
$h<\log(p+1)$ и $h+1\ge\log(p+1)$. Имеем
${\log(p+1)-1} \le h <\log(p+\nobreak 1)$.
\end{proof}

\medskip
Выровненное заполненное бинарное дерево называется
\odmkey{полным}{Дерево (tree)!бинарное (binary)!полное (full)}. В
полном бинарном дереве все листья находятся на последнем уровне и
все узлы, кроме листьев, имеют полустепень исхода 2. Другими
словами, в полном дереве \em{все} ярусы заполнены.

\begin{corollary}
Полное бинарное дерево высотой $h$
имеет $2^{h+1}-1$ узлов.
\end{corollary}
\end{frame}

\begin{frame}
\frametitle{9.5.1. Выровненные и полные деревья (3/3)}
Иногда полным называют бинарное дерево,
в котором все нелистовые узлы имеют
полустепень исхода 2, но листья могут встречаться на любом уровне.
Высота дерева, полного в таком смысле,
может изменяться в широких пределах.

\medskip
\begin{remark}
Заполненные деревья дают наибольший возможный эффект при поиске.\\
Однако известно, что вставка/удаление в заполненное дерево может
потребовать полной перестройки всего дерева и, таким образом,
трудоёмкость операции в худшем случае составит $O(p)$.
\end{remark}

\medskip
Желательно выделить такой подкласс деревьев
сортировки, в котором высота ограничена и в то же время
операции по вставке и удалению элементов выполняются эффективно.
Было найдено несколько таких подклассов, наиболее популярный из
них описывается в следующем параграфе.
\end{frame}


\begin{frame}
\label{9.5.2.}\subsubsection{9.5.2. Сбалансированные деревья }\frametitle{9.5.2. Сбалансированные деревья (1/3)}
%{\small
(Бинарное) дерево называется \odmkey{подровненным деревом}{Дерево
(tree)!подровненное (balanced)}, или \odmkey{АВЛ- деревом}{Дерево
(tree)!АВЛ-дерево (AVL-tree)}\\
(Адельсон-Вельский и Ландис, 1962), 
или \odmkey{сбалансированным по высоте
деревом}{Дерево (tree)!сбалансированное по высоте (height
balanced)}, если для любого узла высоты левого и правого
поддеревьев различаются не более чем на~$1$.


\begin{columns}[t]
\begin{column}[T]{6cm}
\begin{example}
На рисунке приведена диаграмма наиболее несимметричного
сбалансированного по высоте дерева, в котором для всех узлов высота левого
поддерева ровно на 1 больше высоты правого поддерева.
\end{example}
\end{column}
\begin{column}[T]{9cm}

\vspace{-10pt}
\begin{figure}[h]
\begin{center}
\includegraphics[width=70mm]{odm09-13}
\end{center}
\end{figure}
\end{column}
\end{columns}

%-------------------------------------------------------------
% !!!!!!!!!!!!!!!!!!!!!!!!!!!!!!!!!!!!!!!!!!!!!!!!!!!!!!!!!!!!
% TO DO TO DO TO DO TO DO TO DO TO DO TO DO TO DO TO DO TO DO 
% Самый важный и трудный вопрос главы.
% Балансировка деревьев описана без подробностей,
% потому что проходили у Беляева и в других местах.
% А важные и, может быть нужные вариации не описаны:
% 2-3 деревья, красно-чёрные и т.д.
% Что включить в главу, а что ИСКЛЮЧИТЬ !!!!????
% Но для этого надобно решить, чего хотим:
% показать свою учёность или принести пользу.
%-------------------------------------------------------------  

\medskip
\begin{remark}
Кроме сбалансированных по высоте, рассматривают и другие классы
деревьев, например, \odmkey{сбалансированные по весу}{Дерево
(tree)!сбалансированное по весу (weight balanced)} деревья, в
которых для каждого узла количества узлов в левом и правом
поддеревьях находятся в определённом соотношении. Здесь такие
классы не рассматриваются, поэтому далее термин <<сбалансированное
дерево>> означает сбалансированное по высоте бинарное дерево
сортировки.
\end{remark}
%}
%\pagebreak[4]
%{\renewcommand{\baselinestretch}{0.90}\selectfont
\end{frame}


\begin{frame}
\frametitle{9.5.2. Сбалансированные деревья (2/3)}
\begin{theorem}
Для сбалансированного по высоте бинарного дерева $ h \leq 2\log
p. $
\end{theorem}
%\vspace{-3pt}
\begin{proof}
Рассмотрим наиболее несимметричные сбалансированные деревья,
скажем, такие, у которых
всякое левое поддерево на $1$ выше правого. Такие сбалансированные деревья имеют
максимальную возможную высоту
среди всех сбалансированных деревьев с заданным числом
вершин. Пусть $P_h$~--- число вершин в наиболее несимметричном
сбалансированном дереве высоты~$h$. По построению имеем
$p=P_h=P_{h-1}+P_{h-2}+1$.
Непосредственно проверяется, что $P_0=1$, $P_1=2$, $P_2=4$.
Покажем по индукции, что
$P_h\ge (\sqrt 2)^h$.
База: $P_0=1$, $(\sqrt 2)^0=1$, $1\ge 1; P_1=2$,
$(\sqrt 2)^1=\sqrt 2$, $2\ge \sqrt 2$.\\
Пусть $P_h\ge (\sqrt 2)^h$. Тогда
\begin{align*}
P_{h+1}  &= P_h+P_{h-1}+1\ge(\sqrt 2)^h+(\sqrt 2)^{h-1}+1 \\
 &=(\sqrt 2)^h\left(1+\frac{1}{(\sqrt 2)}+\frac{1}{(\sqrt 2)^h}\right)>(\sqrt 2)^h\left(1+\frac{1}{(\sqrt 2)}\right)>(\sqrt
2)^h \sqrt 2=(\sqrt 2)^{h+1}.
\end{align*}
Имеем $(\sqrt 2)^h=2^{h/2}\le P_h=p$, следовательно,
$\frac{h}{2}\le \log p$ и $h\le 2\log p$.
\end{proof}
%\vspace{-3pt}
\end{frame}


\begin{frame}
\frametitle{9.5.2. Сбалансированные деревья (3/3)}
Можно заметить, что рекуррентное соотношение для $P_h$ похоже на
определение чисел Фибоначчи $F(n)$ (\DMref{п.~5.7.3}{5.7.3.}),
откуда имеем
\begin{corollary}
$P_h=F(h+2)-1$.
\end{corollary}
%\vspace{-3pt}
%
\begin{proof}
По индукции. База: $P_0=1$, $F(0)=1$, $F(1)=1$, $F(2)=2$, $1=2-1$.
Пусть $P_{h}=F(h+2)-1$. Рассмотрим $P_{h+1}$. Имеем $P_{h+1} =
P_{h}+P_{h-1}+1 = $ \\$= F(h+2) - 1 + F(h-1+2) - 1 + 1 = F(h+2) +
F(h+1) -1 = F(h+3) -1 =$ \\$= F((h+1)+2) - 1$.
\end{proof}

%\vspace*{-.5\baselineskip}

\begin{remark}
Известна более точная оценка высоты сбалансированного дерева:
$$h < \log_{{\left(\sqrt{5}+{1}\right)/{2}}} 2\, \log_2 p.$$
\end{remark}


Сбалансированные деревья уступают заполненным деревьям по скорости
поиска (менее чем в два раза), однако их преимущество состоит в
том, что известны алгоритмы вставки узлов в сбалансированное
дерево и удаления их из него, которые сохраняют сбалансированность
\pagebreak[3] и в то же время при перестройке дерева затрагивают
только конечное число узлов (см. следующий параграф). Поэтому во
многих случаях сбалансированное дерево оказывается наилучшим
вариантом представления дерева сортировки.

\end{frame}


\begin{frame}
\label{9.5.3.}\subsubsection{9.5.3. Балансировка деревьев }\frametitle{9.5.3. Балансировка деревьев (1/4)}

При добавлении (или при удалении) узла в сбалансированном дереве
сортировки может возникнуть ситуация, в которой баланс нарушается.
В этом случае необходимо перестроить дерево, чтобы восстановить
баланс. Алгоритм балансировки дерева представлен на рисунках, 
при этом использованы
следующие обозначения: $x$~--- добавляемый узел, $v$~--- первый
узел, в котором нарушен баланс на пути от корня $r$ к новому узлу
$x$. Тогда возможны два варианта. Путь от $v$ к $x$ состоит либо
из дуг одной ориентации (скажем, только из левых дуг), либо из
левых и правых дуг. Варианты, в~которых путь состоит только из
правых дуг либо из правых и левых дуг, зеркально симметричны. В
первом варианте дерево перестраивается в соответствии с
\odmkey{правилом простого вращения}{Правило (rule)!вращения (rotation)!простого
(single)}.
Простое вращение восстанавливает баланс и сохраняет условие дерева
сортировки. Действительно, до преобразования имеем
$
u<v,p<v,q<v,p<u,q>u,w>v.
$
Следовательно,
$
p<u,u<v,u<q,u<w,q<v,w>v
$
и преобразование не нарушает условия дерева сортировки.
\end{frame}

\begin{frame}
\frametitle{9.5.3. Балансировка деревьев (2/4)}
\begin{figure}[h]
\begin{center}
\includegraphics[width=140mm]{9_5_3_1.jpg}
\end{center}
%\caption{Простое вращение}
%\label{simpleturn}
\end{figure}
\end{frame}



\begin{frame}
\frametitle{9.5.3. Балансировка деревьев (3/4)}

 Во втором
варианте дерево перестраивается в соответствии с \odmkey{правилом
двойного вращения}{Правило (rule)!вращения (rotation)!двойного (double)}.
Двойное вращение восстанавливает баланс и сохраняет условие
дерева сортировки. Действительно, до преобразования имеем 
$u<v,w>v, p>v,q>v,p<w, \break q>w,s>v,t>v,s<w,t<w,s<p,t>p.$
Следовательно,
$v<p, u<p, s<p, \break u<v, s>v, w>p, t> p, q>p, t<w, q>w$
и преобразование не нарушает условия дерева сортировки.

\end{frame}





\begin{frame}
\frametitle{9.5.3. Балансировка деревьев (4/4)}

\begin{figure}[h]
\begin{center}
\includegraphics[width=140mm]{9_5_3_2.jpg}
\end{center}
%\caption{Двойное вращение}
%\label{doubleturn}
\end{figure}


%\begin{remark}
%Детальное обсуждение алгоритмов работы с АВЛ-деревьями, включая
%оценки трудоёмкости в среднем и в худшем случае, можно найти в
% \cite{Lavrov} и  \cite{Knuth3}.
%\end{remark}
\end{frame}

\begin{frame}
\label{9.5.4.}\subsubsection{9.5.4. Двоичные кучи }\frametitle{9.5.4. Двоичные кучи (1/2)}

Довольно часто в приложениях 
оказывается востребованной модификация 
ассоциативной памяти, в которой записи добавляются
в произвольном порядке, а ищется и извлекается
всегда только запись с максимальным ключом
из хранящихся в данный момент в памяти.

\begin{example}
Выбор для обслуживания процесса с высшим приоритетом.
\end{example}

\begin{remark}
Случай минимального ключа совершенно аналогичен.
\end{remark} 
  
Такая ассоциативная память 
предполагает наличие следующих базовых операций:
\begin{enumerate}
\item[1)] добавить (ключ, запись);
\item[2)] извлечь по максимальному ключу (): запись;
\item[3)] изменить (ключ старый, ключ новый, запись).
\end{enumerate}  

%{\small
\begin{remark}
В некоторых случаях рассматривают несколько иной набор операций.
Вторую операцию делят на две: 
только найти и просмотреть запись с максимальным ключом и
удалить запись с максимальным ключом. 
Также иногда не используют третью операцию
или не разрешают задавать новое значение записи при изменении ключа.
Такие изменённые операции легко выражаются через базовые.
\end{remark} 
%}
    
\end{frame}


\begin{frame}
\frametitle{9.5.4. Двоичные кучи (2/2)}

\odmkey{Двоичная куча}{}~--- это нагруженное бинарное дерево со следующими свойствами:
\begin{enumerate}
\item[1)] значение в любом узле не меньше, чем значения в узлах поддеревьев;
\item[2)] глубина листьев (расстояние до корня) отличается не более чем на 1;
\item[3)] последний уровень заполняется слева направо.
\end{enumerate}

Двоичная куча интересна тем, что при её использовании в качестве
представления ассоциативной памяти реализация операций 
 оказывается весьма эффективной. 

Удобная структура данных для двоичной кучи~--- представление бинарного
дерева массивом $A$, у которого первый элемент,
$A[1]$~--- значение  в корне, а потомками элемента $A[i]$ являются 
элементы $A[2i]$ и $A[2i+1]$.
При таком способе хранения свойства двоичной кучи 2 и 3 выполнены автоматически.

\begin{columns}[t]
\begin{column}[T]{8cm}
\begin{example}
На рисунке представлена двоичная куча,
которая хранится в виде массива

\begin{center}
\begin{tabular}{|c|c|c|c|c|c|}
%1 & 2 & 3 & 4 & 5 & 6 & 7 \\
\hline
10 & 8 & 5 & 7 & 3 & 1 \\
\hline
\end{tabular}
\end{center}
\end{example}
\end{column}
\begin{column}[T]{6cm}
\begin{figure}[h]
\vspace{-2\baselineskip}
\begin{center}
\includegraphics[width=40mm]{9_5_4.jpg}
\end{center}
\end{figure}
\end{column}
\end{columns}

\end{frame}


\begin{frame}
\label{9.5.5.}\subsubsection{9.5.5. Восстановление свойства кучи }\frametitle{9.5.5. Восстановление свойства кучи (1/4)}
%\subsection{9.5.5. Восстановление свойства кучи}
Если в куче изменяется один из элементов, 
то она может перестать удовлетворять свойству 1 двоичной кучи.
Для восстановления этого свойства служит процедура $\Heapify$. 
%Эта процедура принимает на вход массив элементов $A$ и индекс $i$. 
%Она восстанавливает свойство кучи во всём поддереве, 
%корнем которого является элемент $A[i]$ в случае,
%когда поддеревья обладают совйством кучи.

\vspace{-2pt}
\begin{algorithmic}
\REQUIRE массив $A$, индекс $i$.
\ENSURE  восстанавливается свойство кучи в поддереве 
с корнем $A[i]$, причём
левое и правое поддеревья  имеют свойства кучи.
    \STATE {$l \Assign 2 * i$; $r \Assign 2 * i + 1$;  $m \Assign i$ }
    \COMMENT{\color{gray}\small инициализация }
    \IF {$l \leq |A| \And A[l] > A[i]$}
        \STATE $m := l$    
    \ENDIF
    \IF {$r\leq |A| \And A[r] > A[m]$}
        \STATE $m := r$    
    \ENDIF
    \IF {$i \neq m$}
        \STATE $A[i] \leftrightarrow A[m]$ 
        \COMMENT{\color{gray}\small  транспозиция элементов массива }
        \STATE{$\Heapify (A, m)$}
    \ENDIF
\end{algorithmic}
\end{frame}



\begin{frame}
\frametitle{9.5.5. Восстановление свойства кучи (2/4)}
\begin{ground}
В случае, если корень поддерева меньше хотя бы одного из сыновей, 
нужно поменять местами значения
корня с максимальным значением сыновей.
Тогда достаточно спуска по дереву с рекурсивным вызовом функции Heapify
 для гарантирования
восстановления свойств двоичной кучи.
Время работы процедуры прямопропорционально высоте дерева, 
то есть $O \left( \log{n} \right)$.
\end{ground}

\begin{example}
На рисунке слева выделен элемент, 
в котором нарушено условие кучи, а правее показан ход 
выполнения процедуры Heapify
  
\begin{center}
\includegraphics[height=35mm]{9_5_5.jpg}
\end{center}
\end{example}

\end{frame}

\begin{frame}
\frametitle{9.5.5. Восстановление свойства кучи (3/4)}
%\subsection{9.5.6. Построение кучи}

\begin{lemma}
Если выполнить $\Heapify$ для всех элементов массива $A$, начиная с последнего
и кончая первым, то массив станет кучей.
\end{lemma}

\begin{proof} 
Заметим, что $\Heapify(A, i)$ не делает ничего, если $i > \lfloor |A|/2 \rfloor$, 
так как узлы с такими индексами являются листьями, 
следовательно они изначально имеют свойства двоичной кучи. 
Далее по индукции, к моменту выполнения $\Heapify(A, i)$
все поддеревья, чьи корни имеют индекс больше $i$, уже кучи, 
и, следовательно, после выполнения 
$\Heapify(A, i)$ кучей будет вся часть массива $A$, начиная с $i$.
\end{proof}

\medskip
Указанное построение выполняет процедура $\MakeHeap$.

\begin{algorithmic}
\REQUIRE массив $A$
\ENSURE устанавливается свойство кучи в дереве, заданным массивом $A$
    \STATE{\textbf{for $i$ from $\lfloor |A|/2 \rfloor$ downto $1$ do
     \textnormal{Heapify($A$, $i$)} end for}}
\end{algorithmic}

Время работы равно $O \left( n \log{n} \right)$, 
так как функция $\Heapify(A, i)$
вызывается $O \left( n \right)$ раз.
\end{frame}

\begin{frame}
\frametitle{9.5.5. Восстановление свойства кучи (4/4)}
\begin{example}
В таблице показано изменение массива $A$ при построении кучи.

\begin{center}
\begin{tabular}{|c||c|c|c|c|c|c|}
%1 & 2 & 3 & 4 & 5 & 6 & 7 \\
\hline
$i$ & 1 & 3 & 5 & 7 & 8 & 10 \\
\hline
\hline
$3$ & 1 & 3 & 10 & 7 & 8 & 5 \\
\hline
$2$ & 1 & 8 & 10 & 7 & 3 & 5 \\
\hline
$1$ & 10 & 8 & 5 & 7 & 3 & 1 \\
\hline

\end{tabular}
\end{center}
\end{example}


\end{frame}

\begin{frame}
\label{9.5.6.}\subsubsection{9.5.6. Реализация операций кучи }\frametitle{9.5.6. Реализация операций кучи (1/3)}
%\subsection{9.5.7. Реализация операций кучи }
{\renewcommand{\baselinestretch}{0.95}\selectfont
Если при изменении значения элемента c индексом $i$ новый ключ $k$ меньше,
то достаточно вызвать процедуру $\Heapify$, разобранную в \DMref{п.~9.5.5}{9.5.5.}.
При изменении значения элемента $A[i]$ на большее, также могут нарушиться свойства кучи.
В этом случае достаточно выполнить следующий алгоритм $\IncElt(A,i,k)$.

\vspace{-2pt}
\begin{algorithmic}
\REQUIRE массив $A$, индекс изменяемого элемента $i$, новое значение $k$.
\ENSURE значение элемента с индексом $i$ изменяется на $k$ 
без нарушения свойств двоичной кучи.
    \IF{$k < A[i]$}
        \STATE{ \textbf{ return fail}} 
        \COMMENT {\color{gray}\small  новое значение не должно быть  меньше}
    \ENDIF
    \STATE $A[i] \Assign k$
    \WHILE{ $i > 1 \And A[\lfloor i/2 \rfloor] < A[i]$}
        \STATE $A[i] \leftrightarrow A[\lfloor i/2 \rfloor]$ 
        \COMMENT{\color{gray}\small  транспозиция элементов }
        \STATE $i \Assign \lfloor i/2 \rfloor$
    \ENDWHILE
\end{algorithmic}

\vspace{-2pt}
Время работы алгоритма $\IncElt$ равно $O \left( \log{n} \right)$, так как количество транспозиций
прямо пропорционально высоте дерева.
}
\end{frame}

\begin{frame}
\frametitle{9.5.6. Реализация операций кучи (2/3)}


Добавить элемент в двоичную кучу можно так: 
в конец массива добавляется новый элемент с 
заведомо минимальным значением,
затем его значение увеличивается с помощью  
только что описанной процедуры.
Время работы получается такое же~--- $O\left( \log{n} \right)$.

Максимальный элемент двоичной кучи находится в корне,
то есть в первом элементе массива, поэтому его поиск не является
проблемой. Нужно только правильно перестроить кучу. 
Для этого достаточно  поменять местами значения в корне
и в последней вершине дерева, после чего удалить последнюю вершину.
Для правильного перестроения кучи, нужно <<протащить>> 
новое значение в корне вниз, что
и делает процедура $\Heapify$.
\end{frame}

\begin{frame}
\frametitle{9.5.6. Реализация операций кучи (3/3)}
\begin{algorithmic}
\REQUIRE массив $A$
\ENSURE извлечение максимального элемента из кучи с сохранением свойств бинарной кучи
    \STATE{\textbf{if $|A| = 0$ then return fail end if}}
    \COMMENT{\color{gray}\small  куча пуста }
    \STATE $m \Assign A[1]$ 
    \COMMENT{\color{gray}\small  сохранение значения корневого элемента }
    \STATE $A[1] \Assign A[|A|]$ 
    \COMMENT{\color{gray}\small  копирование последнего элемента в корень дерева }
    \STATE $|A| \Assign |A|-1$ 
    \COMMENT{\color{gray}\small  уменьшение длины массива }
    \STATE {$\Heapify(A,1)$}
    \STATE \textbf{return $m$}
\end{algorithmic}

Время работы прямо пропорционально высоте дерева~--- $O\left( \log{n} \right)$.

\end{frame}



\begin{frame}
\label{9.5.7.}
\subsubsection{9.5.7. Построение графа по степенному ряду}
\frametitle{9.5.7. Построение графа по степенному ряду (1/3)}

%\subsection{9.5.8. Построение графа по степенному ряду}

В качестве примера использования двоичной 
кучи рассмотрим задачу 
построения графа по заданным степеням вершин, если такой граф существует.
В \DMref{п.~7.2.6}{7.2.6.} показано, что неубывающая последовательность 
$D = (\Tuple{d}{1}{p})$ является степенной последовательностью графа
тогда и только тогда, когда последовательность 
$D^\prime = (\Tuple{d^\prime}{1}{p-1})$,
полученная \em{откладыванием} последнего элемента $d_p$, является степенной.
Эта теорема определяет конструктивный способ ответа на вопрос, 
является ли данная последовательность
степенной.

\smallskip
Для решения этой задачи подходит структура данных 
\odmkey{очередь с приоритетом}{Очередь (queue)!с приоритетом (priority)}, реализуемая
на основе двоичной кучи.

\begin{remark}
Приоритет~--- это ключ, который ассоциируется с записью. 
Понятие <<приоритета>> прижилось,
так как изначально эта структура данных была придумана для распределения вычислительных
ресурсов в многозадачных операционных системах.
\end{remark}


\end{frame}



\begin{frame}
\frametitle{9.5.7. Построение графа по степенному ряду (2/3)}
Очередь с приоритетом позволяет хранить пары (ключ, значение) и поддерживает две операции:

\begin{enumerate}
\item[1)] добавить (ключ, значение); 
\item[2)] извлечь запись с максимальным ключом : (ключ, значение).
\end{enumerate}


Идея реализации очереди с приоритетом состоит в построении двоичной кучи,
ключами в которой являются приоритеты. Как хранить значения~--- не имеет
принципиального значения, можно считать, что элементами кучи являются
пары (ключ, значение).  


Добавление ключа $k$ со значением $v$ в очередь $Q$ записывается $(k,v)\rightarrow Q$,
а извлечение, соответственно, $(k,v)\leftarrow Q$.  



В таком случае граф по степенной последовательности строится 
с помощью очереди с приоритетом следующим алгоритмом.
В этой очереди ключом является степень вершины, а значением~--- её номер.

\begin{remark}
Для заданной последовательности чисел граф существует не всегда.
Если вершину наибольшей степени не удаётся \em{отложить},
то граф построить невозможно.
\end{remark} 

\end{frame}



\begin{frame}
%\frametitle{9.5.7. Построение графа по степенному ряду (3/3)}

{\small
\begin{algorithmic}
\REQUIRE $A$~--- массив степеней вершин. 
\ENSURE  массив рёбер $E$ графа с заданными степенями, если такой существует.  
\STATE $Q \Assign \emptyset$ 
\COMMENT{\color{gray}\small   инициализация очереди с приоритетом }
\STATE{\textbf{ for $i$ from $1$ to $|A|$ do $(A[i], i)\rightarrow Q$ end for }} 
     \STATE    \COMMENT{\color{gray}\small  номера вершин~--- индексы в начальном массиве}
    \STATE $E := \emptyset$ \COMMENT{\color{gray}\small  массив рёбер }
\WHILE {$|Q|>0$}
    \STATE $(d, n)\leftarrow Q$ 
    \COMMENT{\color{gray}\small  извлечение из очереди откладываемой вершины }
    \STATE{\textbf{ if $d>|Q|$ then return fail end if}}
    \COMMENT {\color{gray}\small  невозможно отложить вершину}    
    \STATE $S := \emptyset$ 
    \COMMENT{\color{gray}\small  инициализация временного контейнера }
    \FOR{\textbf{$k$ from $1$ to $d$}}
        \STATE $(d^\prime, n^\prime)\leftarrow Q$
        \COMMENT{\color{gray}\small  извлечение вершины, смежной с откладываемой }
        \STATE $(n, n^\prime)\rightarrow E$ 
        \COMMENT{\color{gray}\small  добавление ребра в граф}
        \STATE{\textbf{ if $d^\prime > 1$ then $(d^\prime-1, n^\prime)\rightarrow S$ end if}} 
            \COMMENT{\color{gray}\small  убрать в контейнер}
    \ENDFOR
    \STATE{\textbf{for $(d^\prime, n^\prime) \in S$ do $S(d^\prime, n^\prime)\rightarrow Q$
    end for}} 
         \COMMENT{\color{gray}\small  достать из контейнера}
\ENDWHILE
\STATE{\textbf{return $E$}}
\end{algorithmic}
}

\end{frame}






\begin{frame}
\label{9.5.8.}\subsubsection{9.5.8. Дерево отрезков }\frametitle{9.5.8. Дерево отрезков (1/8)}

{\renewcommand{\baselinestretch}{0.92}\selectfont
%\small
Пусть задано непустое конечное линейно упорядоченное множество (т.~е. массив)
некоторых элементов, и требуется производить различные 
\em{групповые} операции
над \em{отрезками} этого множества (то есть, операции с 
\em{отрезками массива}).
Например: найти минимум, максимум или сумму элементов отрезка массива;
присвоить заданную величину всем элементам отрезка 
или увеличить все элементы отрезка на заданную величину и т.д.

\begin{remark}
В языке Algol-68 часть массива называется \odmkey{вырезкой из массива}{Вырезка из массива (slice from array)},
но словосочетание <<дерево вырезок>> неблагозвучно,
поэтому здесь используется термин \odmkey{отрезок массива}{Отрезок (segment)!массива (of array)},
означающий часть массива, определяемую двумя индексами~--- 
начало отрезка и конец отрезка.
\end{remark}

Если требуется только один раз выполнить операцию
над одним отрезком, то дополнительные ухищрения
ничего не дадут, и разумнее всего просто выполнить требуемую
операцию в цикле. Однако, если операцию требуется выполнять
много раз над различными отрезками,
то можно провести   
\odmkey{предобработку}{Предобработка (pre-processing)} заданного массива,
построив дополнительную структуру данных,
которая обеспечит более эффективное выполнение операций.
Для реализации многократно повторяемых групповых
операций на линейно упорядоченном множестве используется структура данных, 
которая называется \odmkey{дерево отрезков}{Дерево (tree)!отрезков (of segments)}.
}


 
\end{frame}

\begin{frame}
\frametitle{9.5.8. Дерево отрезков (2/8)}

Пусть задан массив \textbf{$A:$ array $[1 .. n]$ of $M$},
причём $M$ является моноидом (\DMref{п.~2.2.5}{2.2.5.}) с
ассоциативной операцией $\Map{\varphi}{M \times M}{M}$ 
и нейтральным элементом $e$. 
Также на множестве $M$ задана произвольная
функция преобразования $\Map{f}{M}{M}$. 
Дерево отрезков позволяет поддерживать, по крайней мере, три основные операции:
\begin{enumerate}
\item[1)] построить дерево отрезков (массив $A$);
\item[2)] изменить в массиве элементы (индекс $l$, индекс $r$, функция  $f$);
\item[3)] применить операцию (индекс $l$, индекс $r$) : 
значение $\varphi(A[l..r])$.
\end{enumerate}

Дерево отрезков позволяет реализовать \em{любую} ассоциативную операцию
над отрезком, то есть групповую операцию, результат которой
не зависит от порядка применения операции
к индивидуальным элементам отрезка, 
и позволяет использовать \em{любую} функцию 
независимого преобразования элементов отрезка.
Не ограничивая общности, 
здесь в
примерах в качестве операции рассматривается суммирование элементов,\\
то есть
 $\varphi(A[l..r])\Assign \sum_{i=l}^r A[i]$,
а в качестве функции преобразования 
рассматривается присваивание конкретного значения $x$, 
т.~е.  $\A{i\in 1..n}{A[i] \Assign x}$.

\end{frame}

\begin{frame}
\frametitle{\vspace{-4pt} 9.5.8. Дерево отрезков (3/8)}

\vspace{-2pt}
{\renewcommand{\baselinestretch}{0.88}\selectfont
Дерево отрезков является нагруженным бинарным деревом, в котором
во всех листовых узлах хранятся элементы массива $A$
(или нейтральные элементы, если нужно),
а все нелистовые узлы 
хранят результат применения операции $\varphi$ к 
своим сыновьям.  
Дерево отрезков представляется 
массивом \textbf{$T:$ array $[1 .. (2n -1 )]$ of $M$}, 
в котором потомками элемента $T[i]$ являются элементы $T[2i]$ и $T[2i + 1]$, 
причём $T[i] = \varphi(T[2i], T[2i + 1])$. 
Если $T[i]=\varphi(A[l..r])$~--- значение операции на отрезке $A[l..r]$, 
то $T[2i]=\varphi(A[l..(l +r )\DIV 2])$  значение операции на левой <<половине>> отрезка,
а $T[2i + 1]= \varphi(A[(l+r)\DIV 2 + 1..r])$~--- значение операции на правой <<половине>> отрезка. Значением операции на отдельном элементе (когда $l=r$) 
является сам этот элемент $A[l]$.
Построение дерева отрезков производится следующим образом: 
$T[1]$~--- значение операции на отрезке $A[1..n]$. 
Для вычисления значения $T[1]$ вычисляются значения в потомках корня, 
т.е. $T[2]$ и $T[3]$, где $T[2]$~--- значение операции на отрезке $A[1..n\DIV 2]$,
а $T[3]$~--- на отрезке $A[n\DIV 2 + 1..n]$. 
Далее отрезки делятся рекурсивно. 
Окончание рекурсии: 
$T[i]$~--- значение операции на отрезке $A[j..j]=A[j]$. 
}


\end{frame}


\begin{frame}
\frametitle{9.5.8. Дерево отрезков (4/8)}
{\renewcommand{\baselinestretch}{0.96}\selectfont
\begin{example}
Пусть исходный массив $A = [3, 12, 6, 8, 4, 3, 5]$.
Тогда построенное дерево (для операции сложения)
имеет вид $T = [41, 29, 12, 15, 14, 7, 5, 3, 12, 6, 8, 4, 3, 5, 0]$

\vspace{-4pt}
\begin{center}
\includegraphics[width=80mm]{9_5_8.jpg}
\end{center}
\end{example}

\vspace{-4pt}
\begin{remark}
Нелистовые элементы с суммами 
находятся в массиве $T$ левее элементов
исходного массива $A$. Дополнительные
элементы с нулями (нейтральными элементами моноида)
находятся в массиве $T$ правее
элементов исходного массива $A$.
Вообще говоря, они излишни для вычислений,
и вставлены только для соблюдения
формального правила, 
что каждый нелистовой узел имеет ровно двух     
сыновей. 
\end{remark}
}
\end{frame}


\begin{frame}
\frametitle{9.5.8. Дерево отрезков (5/8)}

Построение дерева отрезков можно выполнить вызвав 
рекурсивную процедуру \\$\Build(1, 1, n)$. При этом массивы 
$A[1..n]$ и $T[1..2n-1]$ считаются глобальными.


\begin{algorithmic}
\REQUIRE{ номер узла $v$, границы $l,r$ отрезка, соответствующего узлу.}
\ENSURE{ построенное дерево $T$ }
    \IF{$l=r$}
        \STATE{ $T[v] \Assign A[l]$ } 
        \COMMENT {\color{gray}\small отрезок является листом }
    \ELSE
				\STATE{ $m \Assign (l + r) \DIV 2$ }
				\COMMENT {\color{gray}\small  середина отрезка }
				\STATE{ $\Build(2v, l, m)$ }
				\COMMENT {\color{gray}\small  левое поддерево }
				\STATE{ $\Build(2v + 1, m + 1, r)$ }
				\COMMENT {\color{gray}\small  правое поддерево }
				\STATE{ $T[v] \Assign \varphi(T[2v], T[2v + 1])$ }
				\COMMENT {\color{gray}\small  выполняем операцию в данном узле }
    \ENDIF
\end{algorithmic}

\begin{ground}
Построенный массив $T$ действительно
является представлением дерева отрезков,
поскольку порядок элементов в этом массиве в точности совпадает
с обходом бинарного дерева по уровням, а узлы 
нагружены нужными значениями по построению.
\end{ground}


\end{frame}

\begin{frame}
\frametitle{9.5.8. Дерево отрезков (6/8)}

\vspace{-1pt}
Если дерево отрезков $T$ для данного массива $A$ 
и для заданной операции $\varphi$ построено 
(и считается глобальным в следующих процедурах),
то вычислить $\varphi(A[l..r])$, \\где 
$1\leq l \leq r \leq n$ можно следующим образом.  

\vspace{-1pt}
\begin{algorithmic}
\REQUIRE границы $l$ и $r$ запроса.
\ENSURE значение операции на отрезке $\varphi(A[l..r])$.
    \STATE{ \textbf{return $\Phi(1, 1, n, l, r)$} } 
\end{algorithmic}

\vspace{-1pt}
Рекурсивная процедура $\Phi$ вычисления суммы на отрезке.

\vspace{-1pt}
\begin{algorithmic}
\REQUIRE индекс узла $v$, границы $p$ и $q$ отрезка, соответствующего узлу $v$, 
границы $l$ и $r$ запроса на вычисление значения $\varphi(A[l..r])$.
\ENSURE сумма на отрезке $[l..r]$
%    \IF{$l > r$}
%        \STATE{ \textbf{return $0$} } 
%    \ELSIF{$l = p \And r = q$}
%				\STATE{ \textbf{return $T[v]$} }
%		\ELSE 
% не удержался и переписал в своём стиле. ФАН 26.04ю2015
				\STATE{\textbf{if $l > r$ then return $e$ end if}}
				\COMMENT {\color{gray}\small  пустой запрос }
				\STATE{\textbf{if $l = p \And r = q$ return $T[v]$ end if}}
				\COMMENT {\color{gray}\small  знаем ответ }
				\STATE{ $m \Assign (p + q) \DIV 2$ }
				\COMMENT {\color{gray}\small  середина отрезка }
				\STATE{ $g \Assign \Phi(2v, p, m, l, \min(r, m))$ }
				\COMMENT {\color{gray}\small  в левом поддереве }
				\STATE{ $h \Assign \Phi(2v + 1, m + 1, q, \max(l, m + 1), r)$ }
				\COMMENT {\color{gray}\small  в правом поддереве }
				\STATE{ \textbf{return $\varphi(g, h)$} }
				\COMMENT {\color{gray}\small  ответ }
%		\ENDIF
\end{algorithmic}
\end{frame}


\begin{frame}
\frametitle{9.5.8. Дерево отрезков (7/8)}
\begin{ground}
При первом вызове 
функции $\Phi$ имеем $p\leq l \leq r \leq q$, 
поскольку \\$p=1\leq l \leq r \leq n=q$.
Далее либо границы запроса совпали с границей ответственности узла,
и тогда можно сразу вернуть ответ, либо нет, и тогда    
возможны три случая.

\proofsection{$p\leq l\leq r\leq m$} В этом случае ответ надобно искать 
в левом поддереве, и для него выполнено требуемое предусловие,
поскольку $p\leq l\leq \min(r,m) \leq m$, а в правом
поддереве запрос заведомо вернёт нейтральный элемент,
поскольку $\max(l,m+1)>r$.

\proofsection{$m+1\leq l\leq r\leq q$} В этом случае ответ надобно искать 
в правом поддереве, и для него выполнено требуемое предусловие,
поскольку $m+1\leq \max(l,m+1)\leq r \leq q$, а в левом
поддереве запрос заведомо вернёт нейтральный элемент,
поскольку $l>\min(r,m)$.

\proofsection{$p\leq l\leq m\leq r\leq q$} В этом случае ответ надобно искать 
в обоих поддеревьях, и для них выполнены требуемые предусловия,
поскольку $p\leq l\leq \min(r,m) \leq m$ и \\$m+1\leq \max(l,m+1)\leq r \leq q$.
\end{ground}



\end{frame}

\begin{frame}
\frametitle{9.5.8. Дерево отрезков (8/8)}

Изменение значения элемента и перестройка дерева отрезков.

\vspace{-2pt}
\begin{algorithmic}
\REQUIRE  $i$~--- индекс изменяемого элемента в исходном массиве $A$, 
$f$~--- функция изменения значения. 
\ENSURE изменённое дерево $T$
    \STATE{ $\Modify(1, 1, n, i, f)$ } 
\end{algorithmic}

Рекурсивная процедура $\Modify$ изменения значения элемента.

\begin{algorithmic}
\REQUIRE  индекс узла $v$, границы $p$ и $q$ отрезка, соответствующего узлу $v$,
индекс $i$ элемента в массиве $A$, $f$~--- функция изменения значения.
\ENSURE изменённое дерево $T$.
%   \IF{$l$ = $r$}
        \STATE{\textbf{if $p = q$ then $T[v] \Assign f();$ return end if }} 
        \COMMENT {\color{gray}\small  это лист }
%    \ELSE
				\STATE{ $m \Assign (p + q) \DIV 2$ }
				\COMMENT {\color{gray}\small  середина отрезка }
				\STATE{\textbf{if $ p \leq m $ then $ \Modify(2v,p,m,i,f) $
				 else $ \Modify(2v + 1, m + 1,q,i,f) $ end~if}}
%					\STATE{ $\ModifyImpl(T, 2v, l, m, p, q)$ }
%					\COMMENT { Левое поддерево }
%				\ELSE
%					\STATE{ $\ModifyImpl(T, 2v + 1, m + 1, r, p, q)$ }
%					\COMMENT { Правое поддерево }
%				\ENDIF
				\STATE{ $T[v] = \varphi(T[2v], T[2v + 1])$ }
				\COMMENT {\color{gray}\small  пересчитываем узел $v$ }
%		\ENDIF
\end{algorithmic}

\begin{ground}
Рекурсивно спускаемся по дереву отрезков до листа,
изменяем его, и, возвращаясь из рекурсии, пересчитываем предков.   
\end{ground}
\end{frame}

\begin{frame}
\label{9.5.9.}\subsubsection{9.5.9. Оценки эффективности дерева отрезков }\frametitle{9.5.9. Оценки эффективности дерева отрезков (1/3)}
\begin{remark}
Операция $\varphi$ не может являться параметром запроса,
поскольку узлы дерева хранят конкретные значения,
и при изменении операции все значения необходимо пересчитать.
Функция преобразования $f$, напротив, может быть параметром запроса.
Более того, эта функция может не зависеть от параметров и быть константой,
может зависеть от прежнего значения элемента,
или может зависеть от каких-то внешних параметров.
Таким образом, дерево отрезков можно эффективно использовать и в том случае,
когда массив $A$ динамически меняет свои элементы (но не длину).     
\end{remark}

\medskip
\begin{theorem}
Число узлов в дереве отрезков не превосходит $2n - 1$.
\end{theorem}

\begin{proof} Если $n=2^k$, то дерево отрезков является полным и количество
узлов в точности равно $2n - 1$, поскольку крона полного дерева содержит на 1 узел больше, чем все остальное дерево. В неполном дереве отрезков узлов разве что меньше. 
\end{proof}

\begin{theorem}
Время работы процедуры построения 
дерева отрезков $O(n)$.
\end{theorem}
\begin{proof}
Процедура построения дерева отрезков посещает каждый узел один раз, 
а количество узлов не более  $2n - 1$.
\end{proof}
\end{frame}

\begin{frame}
\frametitle{9.5.9. Оценки эффективности дерева отрезков (2/3)}

\begin{theorem}
Процедура запроса на каждом уровне дерева отрезков посещает не более четырёх узлов.
\end{theorem}
%{\small
\begin{proof}
Для удобства введём обозначение запроса: 
$(l, r, l', r')$, где $l, r$~--- границы запроса, 
а $l', r'$~--- границы отрезка массива для данного узла.
Для краткости записи знак $/$ здесь означает деление с остатком.  
Проведём разбор случаев.
\proofsection{$l=l' \And r=r'$} Запрос такого типа отработает мгновенно.
\proofsection{$r' \leq (l  +  r)/2 \Or l' > (l + r)/ 2$} 
В этом случае произойдет всего один рекурсивный вызов: 
либо от левого поддерева, либо от правого. 
\proofsection{$r' > (l  +  r)/2 \And l' \le (l + r)/2$}
В этом случае запрос разделится на два запроса:\\ 
$(l, (l + r)/2, l', (l + r)/2)$ и 
$((l + r)/2  +  1, r, (l  +  r)/2  +  1, r')$.
 Теперь, если $l' > (3l  +  r)/4$, либо $r' \le (3l  +  r)/4$, 
 то имеем случай 2. 
 В противном случае, если $l' \le (3l  +  r)/4$, 
 то левый запрос разделится на два запроса: 
 $(l, (3l + r)/4, l', (3l  +  4)/4)$ и \\
 $((3l + r)/4  +  1, (l + r)/2, (3l + r)/4  +  1,  (l + r)/2)$. 
 Видно, что второй из них отработает мгновенно. 
 Аналогично для правого запроса. 
Таким образом, в любой момент времени запрос посещает не более четырёх вершин 
на каждом уровне и имеется не более двух реально работающих ветвей рекурсии.
\end{proof}
%}

\end{frame}

\begin{frame}
\frametitle{9.5.9. Оценки эффективности дерева отрезков (3/3)}
\begin{theorem}
Время работы процедуры запроса $O(\log n)$.
\end{theorem}
\begin{proof}
По предыдущей теореме процедура запроса посещает не более 4-х вершин на каждом уровне,
а количество уровней в дереве отрезков $O(\log n)$.
\end{proof}

\begin{theorem}
Время работы процедуры изменения элемента в дереве отрезков $O(\log n)$.
\end{theorem}
\begin{proof}
Элемент $A[i]$ участвует в $O(\log n)$ узлах дерева отрезков~--- 
по одному на каждом уровне.
\end{proof}



\begin{example}
Работа процедуры запроса суммы на отрезке $A[1..6]$. 
Процедура посещает только выделенные вершины.
\begin{center}
\includegraphics[width=90mm]{9_5_9.jpg}
\end{center} 
\end{example}
\end{frame}
%\vspace{-0.5\baselineskip}
\subsection{9.6. Кратчайший остов}
\label{9.6.}

\begin{frame}
\frametitle{9.6. Кратчайший остов}

Поиск кратчайшего остова графа является классической 
задачей теории графов. Методы решения этой задачи послужили
основой для многих других важных результатов. В частности,
исследования \em{алгоритма Краскала}, 
описанного в \DMref{п.~9.6.3}{9.6.3.}, привели в
конечном счёте к теории жадных алгоритмов, изложенной в \DMref{п.~2.7.4}{2.7.4}.

\medskip
\begin{tabular}{p{9mm}|p{70mm}|p{64mm}}
  № & Название параграфа
  & Ключевые термины и обозначения \\
  \hline
\DMref{9.6.1.}{9.6.1.} & Стягивающие деревья и остовы 
  & 
Остовный подграф, остов, каркаc, стягивающее дерево, остовный лес, кратчайший остов \\
\DMref{9.6.2.}{9.6.2.}& Схема алгоритма построения кратчайшего остова
   & 
 \\
\DMref{9.6.3.}{9.6.3.}\RS  & Алгоритм Краскала & 
\\
\DMref{9.5.4.}{9.5.4.}  & Алгоритм Прима & 
Задача Штейнера, точки Штейнера\\
\end{tabular}

\end{frame}


\begin{frame}
\label{9.6.1.}\subsubsection{9.6.1. Стягивающие деревья и остовы  }\frametitle{9.6.1. Стягивающие деревья и остовы  (1/2)}
Пусть $G(V,E)$~--- граф. \odmkey{Остовный подграф}{Подграф
(subgraph)!остовный (spanning)} графа $G(V,E)$~--- это подграф,
содержащий все вершины. Остовный подграф, являющийся деревом,
называется \odmkey{остовом}{Остов (spanner)} или
\odmkey{каркасом}{Каркас (spanner)}.
Если подграф, являющийся деревом, в то же время является
остовным, то такое дерево называется 
\odmkey{стягивающим деревом}{Дерево (tree)!стягивающее (spanning)}.
Таким образом, стягивающее дерево и остов~--- 
это одно и то же. 

%\vspace{-0.25\baselineskip}
\begin{remark}
Остов определяется множеством рёбер, поскольку вершины остова суть
вершины графа. Таким образом, всякий остов определяет разбиение 
множества рёбер графа на два множества: 
множество рёбер остова и множество хорд остова. 
\end{remark}

Несвязный граф не имеет остова.  Связный  граф  может
иметь много остовов. Но несвязный граф имеет \odmkey{остовный лес}{Лес (forest)!остовный (spanning)}~---
множество остовов компонент связности.
Если задать длины рёбер, то можно поставить задачу нахождения
\odmkey{кратчайшего}{Остов (spanner)!кратчайший (shortest)}
остова. 

%\vspace{-0.5\baselineskip}
\begin{remark}
Существует множество различных способов найти какой-то остов графа.\\
Например, алгоритм поиска в глубину строит остов (по рёбрам возврата).
\end{remark}

\end{frame}


\begin{frame}
\frametitle{9.6.1. Стягивающие деревья и остовы (2/2) }
%\vspace{-0.5\baselineskip}
\begin{remark}
Множество кратчайших путей из заданной вершины ко всем остальным также
образует остов. Однако этот остов может не быть кратчайшим.
\end{remark}

\begin{example}
На рисунке показаны диаграммы ({\it слева направо})
графа, дерева кратчайших путей из вершины $v_1$ с суммарным весом $5$ и
двух кратчайших остовов 
этого графа с суммарным весом $4$. 
%\vspace{-0.25\baselineskip}
%\vspace{-0.25\baselineskip}
\begin{figure}[h]
\begin{center}
\includegraphics[scale=1]{odm09-14}
\end{center}
%\caption{Граф, дерево кратчайших путей и два кратчайших остова}
%\label{f0914}
\end{figure}
\end{example}
\begin{digress}
Задача нахождения кратчайшего остова имеет множество практических интерпретаций.
Например, пусть задан набор аэродромов и нужно определить
минимальный (по сумме расстояний) набор авиарейсов, который
позволил бы перелететь с любого аэродрома на любой другой.
Решением этой задачи будет кратчайший остов полного графа
расстояний между аэродромами.
\end{digress}

%\pagebreak[3]
\end{frame}


\begin{frame}
\label{9.6.2.}\subsubsection{9.6.2. Схема алгоритма построения кратчайшего остова }\frametitle{9.6.2. Схема алгоритма построения кратчайшего остова (1/2)}

Рассмотрим следующую схему алгоритма построения кратчайшего
остова\index{Алгоритм (algorithm)!построения кратчайшего остова
(minimum spanning tree)}. Пусть $T$~--- множество
непересекающихся деревьев, являющихся подграфами графа $G$.
Вначале $T$ состоит  из отдельных вершин графа $G$, в конце $T$
содержит единственный элемент~--- кратчайший остов графа $G$.
%\pagebreak

%\begin{algorithm}
%\caption{Построение кратчайшего остова}
%\label{MST}
\begin{algorithmic}
\REQUIRE граф $G(V,E)$, заданный матрицей длин рёбер $C$.
\ENSURE
кратчайший остов $T$.
              \STATE $T \Assign V$
              \WHILE{в $T$ больше одного элемента}
              \STATE взять любое поддерево из $T$
              \STATE найти к нему ближайшее
              \STATE соединить эти деревья в $T$
              \ENDWHILE
\end{algorithmic}
%\end{algorithm}

\medskip
\begin{remark}
Различные способы выбора поддерева для наращивания на первом шаге
тела цикла дают различные конкретные варианты алгоритма построения
$MST$ (\em{Minimum Spanning Tree}~---
стандартное обозначение для кратчайшего остова).
\end{remark}

\end{frame}

\begin{frame}
\frametitle{9.6.2. Схема алгоритма построения кратчайшего остова (2/2)}

%{\renewcommand{\baselinestretch}{0.96}\selectfont
\begin{ground}
Возьмём произвольное дерево $T_i$ и выберем ближайшее к нему дерево $T_j$ в $T$,
$T_i\cap T_j=\emptyset$.
Положим
$\displaystyle{
\Delta_{i,j}\Assign \min_{t_i\in T_i, t_j\in T_j}d(t_i,t_j).}
$
Так~как с самого начала все вершины $G$ покрыты деревьями из $T$,
а $T_j$ выбирается ближайшим к $T_i$, то
$\Delta_{i,j}$ всегда реализуется на некотором ребре с длиной $c_{i,j}$.
Далее индукцией по шагам алгоритма покажем,
что все рёбра, включенные в деревья из $T$,
принадлежат кратчайшему остову $MST$. Вначале
выбранных рёбер нет, поэтому их множество включается в кратчайший
остов. Пусть теперь все рёбра, добавленные в $T$, принадлежат
$MST$. Рассмотрим очередной шаг алгоритма. Пусть на этом шаге
добавлено ребро $(i,j)$, соединяющее поддерево $T_i$ с~поддеревом
$T_j$. Если $(i,j)\not\in MST$, то, поскольку $MST$ является
деревом и, стало быть, связен, $\exists (i^*,j^*)\in MST$,
соединяющее $T_i$ с остальной частью $MST$. Тогда удалим из $MST$
ребро $(i^*,j^*)$ и добавим ребро $(i,j)$: $MST^{\prime}\Assign
MST - (i^*,j^*) + (i,j)$. Полученный подграф $MST^{\prime}$
является остовом, причём более коротким, чем $MST$, что
противоречит выбору $MST$.
\end{ground}
%\vspace{-6pt}
%}
\end{frame}

\begin{frame}
\label{9.6.3.}\subsubsection{9.6.3. Алгоритм Краскала }\frametitle{9.6.3. Алгоритм Краскала (1/3)}

Следующий алгоритм известен как  \odmkey{алгоритм
Краскала}{Алгоритм (algorithm)!Краскала (Kruskal)}.
 
\begin{algorithmic}
\REQUIRE
список $E$ рёбер графа $G$ с длинами,
упорядоченный в порядке возрастания длин.
\ENSURE
множество $T$ рёбер кратчайшего остова.

\STATE $T \Assign  \emptyset; k \Assign 1$ 
\COMMENT{\color{gray}\small номер рассматриваемого ребра }
\FOR{\textbf{$i$ from $1$ to $p-1$}}
  \WHILE {$z(T+E[k])>0$}
     \STATE $k \Assign  k+1$ 
     \COMMENT{\color{gray}\small пропустить это ребро }
  \ENDWHILE
  \STATE $T \Assign  T +E[k]$ 
  \COMMENT{\color{gray}\small добавить это ребро в $MST$ }
  \STATE $k \Assign k+1$ 
  \COMMENT{\color{gray}\small и исключить его из рассмотрения }
\ENDFOR
\end{algorithmic}
%\end{algorithm}

\begin{ground}
Для обоснования алгоритма Краскала достаточно показать, что
выдерживается схема алгоритма предыдущего параграфа.
Действительно, поскольку в множестве рёбер $T$ нет циклов по
построению, это множество суть совокупность рёбер некоторого
множества деревьев. 
\end{ground}
\end{frame}

\begin{frame}
\frametitle{9.6.3. Алгоритм Краскала (2/3)}
\begin{ground} (Продолжение) 
Если множество $T$ содержит более одного
дерева, то существует ребро, при добавлении которого не возникает
цикла~--- оно соединяет два дерева в $T$. Добавляемое ребро~---
кратчайшее возможное, значит, на нём реализуется расстояние между
некоторыми деревьями в~$T$. По построению в конце работы алгоритма
множество рёбер $T$ содержит $p-1$ элемент, а значит, является
искомым остовом.
\end{ground}

%{\renewcommand{\baselinestretch}{0.93}\selectfont
\medskip
\begin{theorem}
Семейство всех таких подмножеств множества рёбер графа,
которые не содержат циклов, является матроидом.
\end{theorem}

\medskip
\begin{proof}
Пустое множество рёбер не содержит циклов, и аксиома $M_1$
 выполнена. Далее, если множество рёбер не
содержит циклов, то любое его подмножество также не содержит
циклов, и аксиома $M_2$ выполнена. 
\end{proof}
\end{frame}

\begin{frame}
\frametitle{9.6.3. Алгоритм Краскала (3/3)}
\begin{proof} (Продолжение)
Пусть теперь $E'\subset E$~---
произвольное множество рёбер, а $G'$~--- правильный подграф графа
$G$, образованный этими рёбрами. Очевидно, что любое максимальное
не содержащее циклов подмножество множества $E'$ является
объединением остовов компонент связности графа $G'$ и по
теореме \DMref{п.~9.1.2}{9.1.2.} содержит $p(G')-k(G')$ элементов. Таким
образом, все максимальные не содержащие циклов подмножества
произвольного множества рёбер содержат одинаковое количество
элементов, и по теореме \DMref{п.~2.7.2}{2.7.2.} семейство ациклических
подмножеств рёбер образует матроид.
\end{proof}

\medskip
Таким образом, алгоритм Краскала как жадный алгоритм, 
применённый к матроиду, находит
ациклическое подмножество рёбер наименьшего веса. По построению
алгоритма (основной цикл до $p-1$) это подмножество рёбер
древочисленно, а значит, является кратчайшим остовом.
\end{frame}

\begin{frame}
\label{9.6.4.}\subsubsection{9.6.4. Алгоритм Прима }\frametitle{9.6.4. Алгоритм Прима (1/5)}

В \odmkey{алгоритме Прима}{Алгоритм (algorithm)!Прима (Prim)} 
кратчайший остов порождается в процессе разрастания одного дерева,
к которому присоединяются ближайшие одиночные вершины.
Каждая одиночная вершина является деревом, и схема алгоритма
\DMref{п.~9.6.2}{9.6.2.} выполнена. 

  При этом для каждой
вершины $v$, кроме начальной, используются две пометки:
$\alpha[v]$~--- это ближайшая к $v$ вершина, уже включённая в
остов, а $\beta[v]$~--- это длина ребра, соединяющего $v$ с
остовом. Если вершину $v$ ещё нельзя соединить с остовом одним
ребром, то $\alpha[v]\Assign 0, \beta[v]\Assign\infty$.
%\pagebreak
%\begin{algorithm}%
%\caption{Алгоритм Прима}%
%
%\label{Prim} \index{Алгоритм (algorithm)!Прима (Robert Clay Prim)}
%
%\addvspace{-4.3pt}%
%\end{algorithm}%


\end{frame}

\begin{frame}
\frametitle{9.6.4. Алгоритм Прима (2/5)}
\begin{algorithmic}
\REQUIRE граф $G(V,E)$, заданный матрицей длин рёбер $C$. 
\ENSURE множество $T$ рёбер кратчайшего остова. 
\STATE \textbf{ select $u\in V$  } 
\COMMENT {\color{gray}\small выбираем произвольную вершину } 
\STATE $S\Assign\{u\}$ 
\COMMENT {\color{gray}\small  $S$~--- множество вершин, включённых в
кратчайший остов } 
\STATE  $T\Assign\emptyset$ 
\COMMENT {\color{gray}\small  $T$~---
множество рёбер, включённых в кратчайший остов }
 \FOR{$v\in V -u$}
  \IF{$v\in\Gamma(u)$}
    \STATE $\alpha[v] \Assign  u $
    \COMMENT {\color{gray}\small  $u$~--- ближайшая вершина остова }
    \STATE $\beta[v] \Assign  C[u,v]$
    \COMMENT {\color{gray}\small  $C[u,v]$~--- длина соответствующего ребра }
  \ELSE
    \STATE $\alpha[v] \Assign  0 $
    \COMMENT {\color{gray}\small  ближайшая вершина остова неизвестна }
    \STATE $\beta[v] \Assign  \infty$
    \COMMENT {\color{gray}\small  и расстояние также неизвестно }
  \ENDIF
\ENDFOR
\end{algorithmic}

\end{frame}

\begin{frame}
%\frametitle{9.6.4. Алгоритм Прима (3/5)}

\vspace{-2pt}
{\renewcommand{\baselinestretch}{0.92}\selectfont
\begin{algorithmic}
\FOR{\textbf{$i$ from $1$ to $p-1$}}
  \STATE $x\Assign\infty$
  \COMMENT{\color{gray}\small  начальное значение для поиска ближайшей вершины }
  \FOR{$v\in V\setminus S$}
     \IF{$\beta[v] < x$}
        \STATE $w \Assign  v$
        \COMMENT{\color{gray}\small  нашли более близкую вершину }
        \STATE $x \Assign \beta[v]$
        \COMMENT{\color{gray}\small  и расстояние до неё }
     \ENDIF
   \ENDFOR
   \STATE $S \Assign  S +w $
   \COMMENT{\color{gray}\small  добавляем найденную вершину в остов }
   \STATE  $T \Assign  T + (\alpha[w],w)$
   \COMMENT{\color{gray}\small  добавляем найденное ребро в остов }
   \FOR{$v\in\Gamma(w)$}
      \IF{$v\notin S$}
         \IF{$\beta[v] > C[v,w]$}
            \STATE $\alpha[v] \Assign  w$
            \COMMENT{\color{gray}\small  изменяем ближайшую вершину остова }
            \STATE $\beta[v] \Assign  C[v,w]$
            \COMMENT{\color{gray}\small  и длину ведущего к ней ребра }\pagebreak[3]
          \ENDIF
      \ENDIF
   \ENDFOR
\ENDFOR
\end{algorithmic}
}

\end{frame}

\begin{frame}
\frametitle{9.6.4. Алгоритм Прима (4/5)}


\begin{digress}
Задача о нахождении кратчайшего остова принадлежит к числу
немногих задач теории графов, которые можно считать полностью
решёнными. Между тем, если изменить условия задачи, на первый
взгляд даже незначительно, то она оказывается гораздо более
трудной. Рассмотрим следующую задачу. Пусть задано множество
городов на плоскости и нужно определить минимальный (по сумме
расстояний) набор железнодорожных линий, который позволил бы
переехать из любого города в любой другой. Кратчайший остов
полного графа расстояний между городами не будет являться решением
этой (практически, очевидно, очень важной) задачи,
известной как
\odmkey{задача Штейнера}{Задача (problem)!Штейнера (Jacob Steiner)}. %\footnote{Штейнер (ХХХХ--ХХХХ)}.
В задаче Штейнера допускается введение дополнительных вершин,
называемых \odmkey{точками Штейнера}{Точка Штейнера (Steiner
point)}.  Задача Штейнера на плоскости хорошо изучена, в
частности, известно много важных свойств точного решения.
Например, известно, что каждая точка Штейнера имеет степень 3 и
отрезки в ней встречаются под углом $120^\circ$. 
\end{digress}
\end{frame}

\begin{frame}
\frametitle{9.6.4. Алгоритм Прима (5/5)}

\begin{digress} (Продолжение)
Тем не менее, до
сих пор неизвестно достаточно эффективных алгоритмов решения
данной задачи, и она остаётся предметом интенсивных исследований.
\end{digress}

\medskip
\begin{example}
На рисунке приведены, соответственно, диаграммы
кратчайшего остова, наивного <<решения>> задачи Штейнера и
правильного решения для случая, когда города расположены в~вершинах квадрата.
\end{example}
\begin{figure}[h]
\begin{center}
\includegraphics[scale=1.1]{odm09-15}
\end{center}
%\caption{Кратчайший остов, приближённое и точное решения задачи
%Штейнера} \label{f0915}
\end{figure}


\end{frame}





% ===== tex/DM2022_10.tex =====
\section{10. Циклы, независимость и раскраска}

\begin{frame}
\frametitle{10. Циклы, независимость и раскраска}

%\addvspace{-12pt}
После рассмотрения ациклических связных
графов, то есть деревьев, естественно перейти к рассмотрению
графов с циклами. Некоторые задачи, связанные с~циклами,
в частности с гамильтоновыми циклами, оказываются
труднорешаемыми, так называемыми переборными задачами.
В этой главе на примере задачи отыскания
наибольшего независимого множества вершин
рассматриваются подходы к программному решению
таких задач.
К числу переборных задач относится и задача раскрашивания
графов, которая имеет много приложений в программировании.
С раскраской связана задача об укладывании графа
на заданной поверхности, имеющая
приложения в электронике.

Таким образом, в этой заключительной главе рассматриваются
некоторые известные задачи на графах и указываются связи между
ними. В частности, приводятся решения тех исторических задач, с
которых начато изложение теории графов в главе 7.
\end{frame}

\subsection{10.1. Фундаментальные циклы и разрезы}
\label{10.1.}

\begin{frame}
\frametitle{10.1. Фундаментальные циклы и разрезы}
Первый раздел главы посвящён установлению структуры множеств
циклов и~разрезов в графе с точки зрения векторных пространств.

\medskip
\begin{tabular}{p{12mm}|p{65mm}|p{65mm}}
  № & Название параграфа
  & Ключевые термины и обозначения \\
  \hline
\DMref{10.1.1.}{10.1.1.} & Циклы и разрезы  &
{ Простой цикл, правильный разрез, простой разрез}\\  
\DMref{10.1.2.}{10.1.2.} & Фундаментальная система циклов и циклический ранг &
{Кодерево, цикломатическое число, циклический вектор }\\  
\DMref{10.1.3.}{10.1.3.} & Фундаментальная система разрезов и коциклический ранг  &
{ Коцикломатическое число, коцикл }\\  
\DMref{10.1.4.}{10.1.4.} & Подпространства циклов и коциклов  &
{ Независимые циклы и разрезы, базис }\\  
\end{tabular}  

\end{frame}

\begin{frame}
\label{10.1.1.}\subsubsection{10.1.1. Циклы и разрезы }\frametitle{10.1.1. Циклы и разрезы (1/5)}
%\label{s10-1-1}
%{\renewcommand{\baselinestretch}{0.90}\selectfont
Цикл может входить только в одну компоненту связности графа $G(V,E)$,
а~в~не\-связном графе понятие разреза является вырожденным,
поэтому без ограничения общности
в этом разделе граф $G(V,E)$ считается связным.

\smallskip
Цикл не может содержать одно ребро более одного раза, поэтому в
этом разделе цикл рассматривается как множество рёбер.
 В связи с этим можно дать эквивалентное определение простого цикла:
\odmkey{простым}{Цикл (cycle)!простой (simple)} называется цикл,
никакое собственное  подмножество  которого  циклом  не является.

\smallskip
Напомним, что \odmkey{разрезом}{Разрез (cut)} связного  графа
называется  множество  рёбер,  удаление которых делает граф
несвязным. Заметим, что любое \em{разбиение} множества вершин $V$ на
два непустых подмножества, $V_1$ 
$V_2$, 
$V_1 \cap V_2 = \emptyset$, $V_1 \cup V_2 = V$,
определяет разрез $S\Assign\Set{(v_1,v_2) \in E}{v_1 \in V_1 \And
v_2 \in V_2}$, поскольку правильные подграфы $G_1$ и $G_2$,
определяемые подмножествами $V_1$ и $V_2$,
являются, очевидно, компонентами связности графа $G-S$. Заметим
далее, что множества $V_1$ и $V_2$ определяют друг друга:
$V_1=V\setminus V_2$, $V_2=V\setminus V_1$, поэтому достаточно
задать только одно из них. Естественно ввести обозначение
$\A{U\subset V}{\overline{U}\Def V\setminus U}.$
%}
\end{frame}

\begin{frame}
\frametitle{10.1.1. Циклы и разрезы (2/5)}
Введём обозначение $E(V_1, V_2)$ для множества рёбер,
соединяющих два дизъюнктных непустых подмножества вершин графа $G(V,E)$:
$$
E(V_1,V_2) \Def \Set{(v_1,v_2)\in E}{v_1\in V_1 \And v_2\in V_2},
$$
где $V_1\subset V$, $V_2\subset V$, $V_1\neq\emptyset$,
$V_2\neq\emptyset$, $V_1\cap V_2=\emptyset$. Заметим, что
$E(V_1,V_2) =E(V_2,V_1)$.

Разрез связного графа $G(V,E)$, определяемый непустым
подмножеством $U$ множества вершин $V$, называется
\odmkey{правильным}{Разрез (cut)!правильный (regular)} разрезом и
обозначается $S(U)$:
$$
S(U)\Def \Set{(v_1,v_2)\in E}{v_1\in U \And v_2\in \overline{U}}
= E(U,\overline{U}).
$$
Правильный разрез не содержит <<лишних>> рёбер,
то есть таких рёбер, включение или исключение которых
не меняет компонент связности, получаемых при удалении рёбер разреза.
Ясно, что всякий разрез содержит
некоторый правильный разрез. В этом разделе разрез считается \em{правильным}, если не оговорено обратное.
\end{frame}

\begin{frame}
\frametitle{10.1.1. Циклы и разрезы (3/5)}
\begin{lemma}
Симметрическая разность двух различных
правильных разрезов,  определяемых 
 множествами 
$V_1$  и  $V_2$,  является  правильным  разрезом, 
определяемым симметрической разностью множеств $V_1$ и 
$V_2$:

\vspace{-12pt}
$$
V_1\neq V_2 \Impl S(V_1)\vartriangle S(V_2) = S(V_1\vartriangle
V_2).
$$
\end{lemma}

\vspace{-2pt}
\begin{proof}
<<Пересечение>> правильных разрезов $S(V_1)$ и $S(V_2)$ образует разбиение
множества вершин $V$ на четыре подмножества:

\vspace{-12pt}
$$
V_{11} \Assign V_1\cap V_2, \quad
V_{10} \Assign V_1\cap \overline{V_2}, \quad
V_{01} \Assign \overline{V_1}\cap V_2, \quad
V_{00} \Assign \overline{V_1}\cap \overline{V_2}.
$$

\vspace{-2pt}
В этих обозначениях
\vspace{-4pt}
\begin{align*}
S(V_1) &=
E(V_{11},V_{01})\cup
E(V_{10},V_{01})\cup
E(V_{11},V_{00})\cup
E(V_{10},V_{00}),\\
S(V_2) &=
E(V_{11},V_{10})\cup
E(V_{01},V_{10})\cup
E(V_{11},V_{00})\cup
E(V_{01},V_{00}),
\end{align*}
откуда, учитывая, что $E(V_{10},V_{01})=E(V_{01},V_{10})$, имеем
$$
S(V_1)\vartriangle S(V_2) = E(V_{11},V_{01})\cup
E(V_{10},V_{00})\cup E(V_{11},V_{10})\cup E(V_{01},V_{00}).
$$
Заметим, что\\
$
V_1\vartriangle V_2 =  (V_1\cap \overline{V_2}) \cup
(\overline{V_1}\cap V_2) = V_{10}\cup V_{01}$,
$\overline{V_1\vartriangle V_2} = (V_1\cap V_2) \cup
(\overline{V_1}\cap \overline{V_2}) = V_{11}\cup V_{00}.
$
\end{proof}
\end{frame}

\begin{frame}
\frametitle{10.1.1. Циклы и разрезы (4/5)}
\begin{proof} (Продолжение)
Поэтому
$$
S(V_1\vartriangle V_2) = E(V_{10},V_{11})\cup E(V_{10},V_{00})\cup
E(V_{01},V_{11})\cup E(V_{01},V_{00}),
$$
и, учитывая, что $E(V_{10},V_{11})=E(V_{11},V_{10})$ и
$E(V_{01},V_{11})=E(V_{11},V_{01})$, окончательно имеем
$S(V_1)\vartriangle S(V_2) = S(V_1\vartriangle V_2)$.
\end{proof}

\begin{figure}[h]
\begin{center}
\includegraphics[scale=1.0]{10_1_1.jpg}
\end{center}
%\caption{К доказательству леммы п.~\ref{s10-1-1}}
%\label{f1000new}
\end{figure}


\odmkey{Простым}{Разрез (cut)!простой (simple)} разрезом
называется минимальный разрез, то есть такой разрез, никакое
собственное подмножество которого разрезом не является.

\end{frame}

\begin{frame}
\frametitle{10.1.1. Циклы и разрезы (5/5)}

\begin{remark}
Всякий простой разрез является правильным, но не всякий правильный разрез является простым.
\end{remark}

\begin{example}
Рассмотрим граф $G$, представленный на рисунке {\it а}. Разрез $S=\{(v_1,v_2),(v_1,v_3),\break(v_1,v_4),(v_4,v_5),(v_2,v_3)\}$ 
графа $G$ не является правильным разрезом (рисунок {\it б}). Исключив ребро $(v_2,v_3)$, получим 
$S1=\{(v_1,v_2),(v_1,v_3),(v_1,v_4),(v_5,v_4)\}$ ~--- правильный разрез графа $G$ (рисунок {\it в}), причём $S1$ не является простым разрезом. Правильный разрез $S1$ содержит в себе два простых разреза $S1_1=\{(v_1,v_2),(v_1,v_3),(v_1,v_4)\}$ и $S1_2=\{(v_5,v_4)\}$, (рисунки {\it г} и {\it д}).
\end{example}

\vspace{-6pt}
\begin{figure}[h]
\begin{center}
\includegraphics[width=150mm]{10_1_1_2.jpg}
\end{center}
%\caption{Фундаментальные циклы}
%\label{f1002new}
\end{figure}

\addvspace{-3pt}
\begin{remark}
Чем больше в графе циклов, тем труднее его
разр$\acute{\text{е}}$зать. В дереве, напротив, каждое ребро само
по себе является простым разрезом.
\end{remark}

\end{frame}

\begin{frame}
\label{10.1.2.}\subsubsection{10.1.2. Фундаментальная система циклов и~циклический~ранг }\frametitle{10.1.2. Фундаментальная система циклов и~циклический~ранг (1/4)}
%\label{ciclerank}

Пусть $T(V,E_T)$~--- некоторый остов графа $G(V,E)$.
\odmkey{Кодеревом}{Кодерево (cotree)} $T^*(V,E^*_T)$ остова $T$
называется остовный подграф, такой, что $E^*_T=E\setminus E_T$.
(Кодерево не является деревом!) Рёбра кодерева являются
\odmkey{хордами}{Хорда (chord)} остова. По теореме п~9.1.2 об
основных свойствах деревьев каждая хорда $e\in T^*$ остова $T$
порождает ровно один простой цикл, обозначим его $Z_e$. Таким
образом, имеем систему простых циклов
$\mathcal{Z}\Def\left\{Z_e\right\}_{e\in T^*}$,
определяемых выбранным остовом $T$, которая называется
\odmkey{фундаментальной системой циклов}{Система (system)!циклов
(of cycles)!фундаментальная (fundamental)}. Циклы фундаментальной
системы называются \odmkey{фундаментальными}{Цикл
(cycle)!фундаментальный (fundamental)}, а количество циклов в
(данной) фундаментальной системе называется \odmkey{циклическим
рангом}{Ранг (rank)!циклический (cycle)} (или
\odmkey{цикломатическим числом}{Число (number)!цикломатическое
(cyclomatic)}) графа $G$ и обозначается $m(G)$.

%\addvspace{-3pt}
\medskip
\begin{theorem}
Любой цикл в связном графе $G(V,E)$
можно представить как симметрическую разность
нескольких фундаментальных циклов
из системы $\mathcal{Z}$, определяемой
произвольным остовом $T$.
\end{theorem}
\end{frame}

\begin{frame}
\frametitle{10.1.2. Фундаментальная система циклов и~циклический~ранг (2/4)}
%\pagebreak[3]
%\addvspace{-3pt}
\begin{proof}
Если в графе $G$ нет циклов, то это дерево, $G=T$,
$T^*=\emptyset$, $\mathcal{Z} =\emptyset$, и утверждение теоремы
тривиально. 
Рассмотрим цикл $Z$ в графе $G$. Этот цикл содержит
хорды $\Tuple{e}{1}{n}\in T^*$. Такие хорды в $Z$ обязательно
есть, в~противном случае $Z \subset  T$, что невозможно, поскольку
$T$~--- дерево. Докажем индукцией по $n$, что
$Z = Z_{e_1} \vartriangle \dots \vartriangle Z_{e_n}$. База: пусть
$n=1$, тогда $e_1\not\in T$, $Z-e_1 \subset T$ и $Z =Z_{e_1}$,
так~как если бы $Z\neq Z_{e_1}$, то концы $e_1$ были бы соединены
в $T$ двумя цепями, что невозможно по теореме \DMref{п.~9.1.2.}{9.1.2.} Пусть
(индукционное предположение)
$Z=Z_{e_1} \vartriangle\ldots \vartriangle Z_{e_m}$ для всех
циклов $Z$ с числом хорд $m<n$. Рассмотрим цикл $Z$ с $n$ хордами
$\Tuple{e}{1}{n}\in T^*$ и цикл $Z_{e_n}$.

\begin{columns}[t]
\begin{column}[T]{9cm}
Имеем $Z'\Assign Z\vartriangle Z_{e_n} =
(Z-e_n)\cup(Z_{e_n}-e_n)$~--- тоже цикл (возможно, не простой).
Но~$Z'$ содержит только $n-1$ хорд $\Tuple{e}{1}{n-1}$. По
индукционному предположению
$Z'=Z_{e_1}\vartriangle{\ldots}\vartriangle Z_{e_{n-1}}$.
Имеем 
$ Z  =  Z' \vartriangle  Z_{e_n} = 
\left(Z_{e_1} \vartriangle \dots \vartriangle Z_{e_{n-1}}\right)
\vartriangle Z_{e_n} =  Z_{e_1} \vartriangle \dots \vartriangle 
Z_{e_{n-1}}  \vartriangle  Z_{e_n}.$
\end{column}
\begin{column}[T]{6cm}
\begin{figure}[h]
\begin{center}
\includegraphics[width=60mm]{10_1_2.jpg}
\end{center}
\end{figure}

\end{column}
\end{columns}
\end{proof}
\end{frame}

\begin{frame}
\frametitle{10.1.2. Фундаментальная система циклов и~циклический~ранг (3/4)}
\begin{corollary}
Количество циклов в фундаментальной
системе равно числу хорд остова:\\
$m(G)=q-p+1$.
\end{corollary}

\begin{proof}   $m(G) =  q(T^*) = q(G) - q(T)= q - (p - 1)=q - p + 1$.
\end{proof}

\begin{columns}[t]
\begin{column}[T]{7cm}
\begin{example}
На рисунке представлена система фундаментальных
циклов, определяемых некоторым остовом. Рёбра остова выделены
жирными линиями, а соответствующие фундаментальные циклы
 $Z_1=\{(v_1,v_2),(v_2,v_3),(v_3,v_1)\}$,
$Z_2=\{(v_2,v_5),(v_5,v_3),(v_3, v_2)\}$,\\
$Z_3=\{(v_2,v_5),(v_5,v_6),(v_6,v_3),(v_3,v_2)\}$, $Z_4
=\{(v_2,v_4),(v_4,v_5),(v_5,v_2)\}$,~--- 
пунк\-тирными
контурами.
\end{example}
\end{column}
\begin{column}[T]{8cm}

\vspace{-10pt}
\begin{figure}[h]
\begin{center}
\includegraphics[scale=0.9]{10_02}
\end{center}
\end{figure}
\end{column}
\end{columns}
\end{frame}

\begin{frame}
\frametitle{10.1.2. Фундаментальная система циклов и~циклический~ранг (4/4)}
\begin{remark}
Совокупность всевозможных симметрических разностей фундаментальных
циклов содержит множеств больше, чем нужно: в неё входят не только
все циклы графа $G$, но и~некоторые объединения таких циклов. В
частности, симметрическая разность двух (фундаментальных) циклов,
которые не имеют общих вершин, снова даёт два этих же цикла.
Иногда множество объектов, определяемое всевозможными циклическими
разностями фундаментальных циклов, называют множеством
\odmkey{циклических векторов}{Вектор (vector)!циклический (cyclic
vector)}.
\end{remark}

\begin{example}
На рисунке нет фундаментальных циклов, не имеющих
общих вершин. Симмет\-рическая разность трёх фундаментальных
циклов, $Z_1\vartriangle Z_3\vartriangle Z_4$, даёт цикл\\
$(v_1,v_2),(v_2,v_4),(v_4,v_5),(v_5,v_6),(v_6,v_3),(v_3,v_1)$~---
 <<внешнюю границу>> графа.
Симметри\-ческая разность всех фундаментальных циклов
$Z_1\vartriangle Z_2\vartriangle Z_3\vartriangle Z_4$ даёт эйлеров
цикл данного графа (см. раздел \DMref{10.2}{10.2.}).
\end{example}

%\addvspace{-3pt}
\end{frame}

\begin{frame}
\label{10.1.3.}\subsubsection{10.1.3. Фундаментальная система разрезов и~коциклический ранг }\frametitle{10.1.3. Фундаментальная система разрезов и~коциклический ранг (1/4)}
%\label{cocyclerank}

Пусть опять $T(V,E_T)$~--- некоторый остов графа $G(V,E)$.
Рассмотрим  ребро  остова\\  
$e \in  E_T$  и 
определим  разрез  $S_e$  следующим 
образом. Так как ребро $e$~--- разрез дерева $T$, 
множество вершин $V$ разбивается
на два непустых подмножества, $V_1$ и $V_2$,
так что\\
$V_1 \subset  V$, $V_1 \neq  \emptyset$,
$V_2 \subset  V$, $V_2 \neq \emptyset$,
$V_1\cup V_2=V$, $V_1\cap V_2=\emptyset$.
Включим в разрез $S_e$ ребро $e$ и все хорды остова $T$,
соединяющие вершины  $V_1$ с вершинами $V_2$:
$$
S_e\Def e+\Set{(v_1,v_2)\in T^*}{v_1\in V_1 \And v_2\in V_2}.
$$
Тогда $S_e$~--- это простой разрез.
Система разрезов
$
\mathcal{S}\Def\left\{S_e\right\}_{e\in T}
$
называется \odmkey{фундаментальной системой разрезов}{Система
(system)!разрезов (of cuts)!фундаментальная (fundamental)}.
Разрезы фундаментальной системы называются
\odmkey{фундаментальными}{Разрез (cut)!фундаментальный
(fundamental)}, а количество разрезов в (данной) фундаментальной
системе называется \odmkey{коциклическим рангом}{Ранг
(rank)!коциклический (cocycle)} (или
\odmkey{коцикломатическим числом}{Число (number)!коцикломатическое
(cocyclomatic)}) 
графа $G$ и обозначается $m^*(G)$.

\begin{remark}
Между циклами и правильными разрезами существует определённая двойственность,
поэтому правильные разрезы иногда называют \odmkey{коциклами}{Коцикл
(cocycle)}. Отсюда название <<фундаментальная система коциклов>> и~<<коциклический ранг>>. 
%Детальное рассмотрение этой двойственности
%выходит за рамки данной книги.
\end{remark}

\end{frame}

\begin{frame}
\frametitle{10.1.3. Фундаментальная система разрезов и~коциклический ранг (2/4)}
%\vspace{-4pt}

{\renewcommand{\baselinestretch}{0.97}\selectfont
\begin{theorem}
Любой правильный разрез в связном графе $G(V,E)$
можно представить как симметрическую разность
некоторых фундаментальных разрезов
из системы $\mathcal{S}$, определяемой
произвольным остовом $T$.
\end{theorem}


%\addvspace{-18pt}
\begin{proof}
Действительно, любой разрез $S$ содержит хотя бы одно ребро из
остова $T$, так~как $T$~--- связный и содержит все вершины $G$. Пусть
$S$~--- правильный разрез, который содержит рёбра
$\Tuple{e}{1}{n}\in T$. Докажем индукцией по $n$, что
$S=S_{e_1}\vartriangle\ldots\vartriangle S_{e_n}$. База: пусть
$n=1$, тогда $T - e_1 = T_1 \cup T_2$, где $T_1$ и $T_2$~---
компоненты, получаемые из остова $T$ удалением ребра $e_1$. Имеем
$S_{e_1}\subset S$, иначе $S$ не был бы разрезом, и $S\subset
S_{e_1}$, иначе $S$ не был бы правильным разрезом. Таким образом,
$S=S_{e_1}$. 

Пусть теперь
$S=S_{e_1} \vartriangle\ldots \vartriangle S_{e_m}$ для всех
правильных разрезов $S$ с числом рёбер остова $m<n$. Рассмотрим
правильный разрез $S$ с $n$ рёбрами $\Tuple{e}{1}{n}\in T$.
Положим $S'\Assign S\vartriangle S_{e_n}$, то есть исключим из
разреза $S$ все рёбра фундаментального разреза $S_{e_n}$.
 Разрез
$S'$~--- правильный по лемме \DMref{п.~10.1.1}{10.1.1.} и содержит рёбра
$\Tuple{e}{1}{n-1}$, а значит, по индукционному предположению
$S'=S_{e_1}\vartriangle\ldots \vartriangle S_{e_{n-1}}$. Имеем
$S = S'~\vartriangle~S_{e_n}$, и окончательно
$S =S_{e_1}~\vartriangle~\ldots~\vartriangle~S_{e_{n-1}}\vartriangle~S_{e_n}$.
\end{proof}
}
\end{frame}

\begin{frame}
\frametitle{10.1.3. Фундаментальная система разрезов и~коциклический ранг (3/4)}



\begin{corollary}
Количество разрезов в фундаментальной
системе равно числу рёбер остова:
$
m^*(G)=p-1.
$
\end{corollary}

%\vspace{-6pt}
\begin{proof}
По определению.
\end{proof}


\begin{example}
На рисунке представлены
все пять фундаментальных разрезов, соответствующих
графу и его остову, представленным на предыдущем рисунке.
\end{example}

\vspace{-6pt}
\begin{figure}[h]
\begin{center}
\includegraphics[width=150mm]{10_03}
\end{center}

\end{figure}

%\enlargethispage{\baselineskip}
%\addvspace{-3pt}
\end{frame}


\begin{frame}
\frametitle{10.1.3. Фундаментальная система разрезов и~коциклический ранг (4/4)}
\begin{remark}
Совокупность всевозможных симметрических разностей фундаментальных
разрезов содержит множеств меньше, чем может понадобиться: в неё входят только правильные разрезы; разрезы, не являющиеся правильными, в ней не содержатся. В частности, разрез может содержать циклы, в правильном разрезе или симметрической разности правильных разрезов циклов быть не может, поэтому разрезов с циклами нам не получить с помощью фундаментальной системы разрезов.
\end{remark}

\begin{columns}[t]
\begin{column}[T]{7cm}
\begin{example}
Обратимся ещё раз к приведённому на рисунке графу. Рассмотрим разрез
$S=\{(v_1,v_2)$, $(v_2,v_3)$, $(v_3,v_5)$, $(v_5,v_6)$,
$(v_6,v_3)$, 
$(v_3,v_1)\}$. Этот разрез графа не является правильным и его нельзя представить в виде симметрической разности фундаментальных разрезов.
\end{example}
\end{column}
\begin{column}[T]{7.5cm}

\vspace{-16pt}
\begin{figure}[h]
\begin{center}
\includegraphics[width=7cm]{10_02}
\end{center}
%\caption{Фундаментальные циклы}
%\label{f1002new}
\end{figure}
\end{column}
\end{columns}
\end{frame}

\label{10.1.4.}
\subsubsection{10.1.4. Подпространства циклов и коциклов}
\begin{frame}
\frametitle{10.1.4. Подпространства циклов и коциклов (1/3)}

Рассмотрим векторное пространство (\DMref{п.~2.5.1}{2.5.1.}) над двоичной
арифметикой (\DMref{п.~2.3.3}{2.3.3.}), натянутое
на множество рёбер $E$ графа $G(V,E)$. Элемент этого векторного
пространства 
(линейную комбинацию, \DMref{п.~2.5.2}{2.5.2.})
 можно
отождествлять с подмножеством рёбер: если коэффициент в линейной
комбинации равен 1, то ребро входит в подмножество, если
коэффициент равен 0, то не входит. При такой интерпретации
сложению векторов соответствует симметрическая разность множеств
рёбер.
Множество циклических векторов и множество правильных
разрезов замкнуты относительно симметрической разности,
а потому являются подпространствами  общего 
пространства  подмножеств  рёбер  графа.

%\vspace{-0.2\baselineskip}
\medskip
\begin{remark}
Строго говоря, множество циклов, равно как и множество разрезов,
не образует подпространства, поскольку не замкнуто относительно
симметрической  разности.  
Множество  циклических  векторов 
 шире 
множества  циклов,  а  множество правильных  разрезов 
$\acute{\text{у}}$же  множества разрезов. Однако интуитивно ясно,
что <<по потребительским свойствам>> циклические векторы и циклы,
равно как разрезы и правильные разрезы, весьма близки, поэтому
далее рассматриваются только циклические векторы
и правильные разрезы, которые для краткости называются циклами и
разрезами (коциклами).
\end{remark}

\end{frame}

%
\begin{frame}
\frametitle{10.1.4. Подпространства циклов и коциклов (2/3)}
%\pagebreak


Множество циклов $\{Z_i\}_{i=1}^{n}$ называется
\odmkey{независимым} {Множество (set)!циклов (of
cycles)!независимое (independent)}, если ни один из циклов $Z_i$
не является линейной комбинацией остальных.

Множество разрезов $\{S_i\}_{i=1}^{n}$ называется
\odmkey{независимым} {Множество (set)!разрезов (of
cuts)!независимое (independent)}, если ни один из разрезов~$S_i$
не является линейной комбинацией остальных.

Максимальное независимое множество циклов (или минимальное
множество циклов, от которых зависят все остальные, или
независимое порождающее множество циклов) является
\odmkey{базисом}{Базис (basis)!пространства циклов (of cycle
space)} пространства циклов. Максимальное независимое множество
разрезов (или минимальное множество разрезов, от которых зависят
все остальные, или независимое порождающее множество разрезов)
является \odmkey{базисом}{Базис (basis)!пространства коциклов (of
cocyclic space)} пространства разрезов.

Введенная  фундаментальная
система циклов является базисом подпространства циклов.
Действительно, циклы любой фундаментальной системы
$\mathcal{Z}=\{Z_e\}_{e\in T*}$
независимы, поскольку каждый из них содержит индивидуальную
хорду $e\in T^*$, и по теореме  фундаментальная
система порождает множество циклов.
Таким образом, циклический ранг~--- это размерность
пространства циклов.

%\vspace{-0.2\baselineskip}

\end{frame}


\begin{frame}
\frametitle{10.1.4. Подпространства циклов и коциклов (3/3)}
%\enlargethispage{\baselineskip}%\vspace{-0.1\baselineskip}

\begin{remark}
Фундаментальные системы циклов, порождаемые остовами,~--- это не
единственные базисы пространства циклов. Например, четыре
<<естественных>> треугольника на рисунке в \DMref{п.~10.1.2}{10.1.2.}~---
$\{(v_1,v_2),(v_2,v_3),(v_3,v_1)\}$,
$\{(v_2,v_4),(v_4,v_5),(v_5,v_2)\}$,\\
$\{(v_2,v_5),(v_5,v_3),(v_3,v_2)\}$, $\{(v_3,v_5),
(v_5,v_6),(v_6,v_3)\}$~--- являются базисом подпространства
циклов, но не являются фундаментальной системой ни для какого
остова.
\end{remark}

\medskip
Введенная  фундаментальная система
разрезов является базисом подпространства разрезов. Действительно,
разрезы любой фундаментальной системы $\mathcal{S}=\{S_e\}_{e\in
T}$ независимы, поскольку каждый из них содержит индивидуальное
ребро $e\in T$, и по теореме  фундаментальная
система порождает множество разрезов. Таким образом, коциклический
ранг~--- это размерность пространства разрезов.

%%\vspace{-0.2\baselineskip}
\begin{remark}
Все утверждения этого и следующего разделов распространяются
на случай м\-ультиграфов.
\end{remark}
\end{frame}


\begin{frame}
\frametitle{Онтология теории графов}
\begin{center}
\includegraphics[width=14cm]{ODM_10_1.png}
\end{center}

%{\color{red} Тема для 
%сочинения: двойственность циклов и коциклов. }
\end{frame}


\subsection{10.2. Эйлеровы циклы}\label{10.2.}

\begin{frame}
\frametitle{10.2. Эйлеровы циклы}
Здесь приведено исчерпывающее решение задачи о Кёнигсбергских мостах
(см. \DMref{п.~7.1.1}{7.1.1.}),
приведшей к исторически первой успешной попытке развития
теории графов как самостоятельного предмета.

\medskip
\begin{tabular}{p{12mm}|p{65mm}|p{65mm}}
  № & Название параграфа
  & Ключевые термины и обозначения \\
  \hline
\DMref{10.2.1.}{10.2.1.} & Эйлеровы графы  &
{ Эйлеров цикл, эйлерова цепь, полуэйлеров граф}\\  
\DMref{10.2.2.}{10.2.2.} & Алгоритм построения эйлерова цикла в эйлеровом графе  &
{ }\\  
\DMref{10.2.3.}{10.2.3.} & Оценка числа эйлеровых графов &
{  }\\  
\end{tabular}  

\end{frame}



\begin{frame}
\label{10.2.1.}\subsubsection{10.2.1. Эйлеровы графы }\frametitle{10.2.1. Эйлеровы графы (1/3)}
%\label{s10-2-1}

\vspace{-2pt}
Если связный граф имеет цикл (не обязательно простой), содержащий все
рёбра графа, то такой цикл называется \odmkey{эйлеровым}{Цикл
(cycle)!эйлеров (eulerian)} циклом, а граф называется
\odmkey{эйлеровым}{Граф (graph)!эйлеров (eulerian)} графом. Если связный граф
имеет цепь (не обязательно простую), содержащую все рёбра, то
такая цепь называется \odmkey{эйлеровой}{Цепь (circuit)!эйлерова
(eulerian)} цепью, а граф называется \odmkey{полуэйлеровым}{Граф
(graph)!полуэйлеров (semi-eulerian)} графом.
И эйлеров цикл и эйлерова цепь содержат не только все рёбра, но и все
вершины графа (возможно, по нескольку раз).

%\vspace{-2pt}

\begin{theorem}
Если граф $G$ связен и нетривиален, то следующие утверждения эквивалентны.
\begin{enumerate}
\item $G$~{\upshape---} эйлеров граф. \item Каждая вершина $G$
имеет чётную степень. \item Множество рёбер $G$ можно разбить на
простые циклы.
\end{enumerate}
\end{theorem}


%\vspace{-2pt}
\begin{proof}
%\vspace{-2pt}
\proofsection{1$\Impl$ 2} Пусть $Z$~--- эйлеров цикл в $G$.
Двигаясь по $Z$, подсчитаем степени вершин, полагая их до
начала подсчёта нулевыми. Прохождение каж\-дой вершины вносит 2
в степень этой вершины. 
Поскольку $Z$ содержит все рёбра, то,
когда обход $Z$ будет закончен, будут учтены все рёбра, а степени
всех вершин~--- чётные.
\end{proof}
\end{frame}



\begin{frame}
\frametitle{10.2.1. Эйлеровы графы (2/3)}

\begin{proof} (Продолжение)
\proofsection{2$\Impl$ 3}
$G$~--- связный и нетривиальный граф,
следовательно, $\A{v_i}{d(v_i)>0}$.
Степени вершин чётные, следовательно,
$\A{v_i}{d(v_i)\geq 2}$.
Имеем 
$$2q=\sum_{i=1}^{p} d(v_i)\geq 2p \Impl q\geq p \Impl q>p-1.$$
Следовательно, граф $G$~--- не дерево,
а значит, граф $G$
содержит (хотя бы один) простой цикл $Z_1$. ($Z_1$~--- множество рёбер.)
Тогда $G-Z_1$~--- остовный подграф,
в~котором опять все степени вершин
чётные. Исключим из рассмотрения изолированные вершины. Таким~образом, $G-Z_1$
тоже удовлетворяет условию~2, следовательно, существует простой цикл
$Z_2\subset(G-Z_1)$.
Далее выделяем циклы~$Z_i$,
пока граф не будет пуст.
Имеем $E=\bigcup Z_i$ и $\bigcap Z_i=\emptyset$.
\end{proof}
\end{frame}



\begin{frame}
\frametitle{10.2.1. Эйлеровы графы (3/3)}
\begin{proof} (Окончание)
\proofsection{3$\Impl$1} Возьмем какой-либо цикл $Z_1$ из данного
разбиения. Если $Z_1=E$, то теорема доказана. Если нет, то
существует цикл $Z_2$, не имеющий общих рёбер с~$Z_1$ (см.
рисунок), такой, что $\E{v_1}{(v_1\in Z_1\And v_1\in
Z_2)}$, так~как $G$ связен. Маршрут $Z_1\cup Z_2$ является циклом
и содержит все свои рёбра по одному разу.%\linebreak

%\pagebreak
Если $Z_1\cup Z_2=E$, то теорема доказана. Если нет, то существует
цикл $Z_3$, такой, что $\E{v_2}{v_2\in Z_1\cup Z_2\And v_2\in
Z_3}$. Далее будем наращивать эйлеров цикл, пока он не исчерпает
разбиения.
\end{proof}

\begin{figure}[h]
\begin{center}
\includegraphics[height=30mm]{10_04}
\end{center}
%\caption{К доказательству теоремы п.~\ref{s10-2-1}}
%\label{f1003new}
\end{figure}


\end{frame}

\begin{frame}
\label{10.2.2.}\subsubsection{10.2.2. Алгоритм построения эйлерова цикла в~эйлеровом графе }\frametitle{10.2.2. Алгоритм построения эйлерова цикла в~эйлеровом графе (1/2)}


{\renewcommand{\baselinestretch}{0.95}\selectfont
%\vspace{-2pt}
\begin{algorithmic}
\REQUIRE эйлеров граф $G(V,E)$, заданный списками смежности
($\Gamma[v]$~--- список вершин, смежных с вершиной $v$).
\ENSURE последовательность вершин эйлерова цикла.
\STATE $S\Assign \emptyset$ 
\COMMENT{\color{gray}\small стек для хранения вершин }
\STATE \textbf{select $v\in V$} 
\COMMENT{\color{gray}\small  произвольная вершина }
\STATE $v\rightarrow S$ 
\COMMENT{\color{gray}\small  положить $v$ в стек $S$ }
\WHILE{$S\ne\emptyset$}
  \STATE \textbf{$v\Assign$ top $S$}
  \COMMENT{\color{gray}\small  $v$~--- верхний элемент стека }
  \IF{$\Gamma[v]=\emptyset$}
    \STATE $v\leftarrow S$; \textbf{yield $v$}
    \COMMENT{\color{gray}\small  очередная вершина эйлерова цикла }
  \ELSE
    \STATE\textbf{select $u\in \Gamma[v]$}
    \COMMENT{\color{gray}\small  взять первую вершину из списка смежности }
     \STATE $u\rightarrow S$ 
     \COMMENT{\color{gray}\small  положить $u$ в стек }
     \STATE $\Gamma[v]\Assign\Gamma[v]-u;
      \Gamma[u]\Assign\Gamma[u]-v$
     \COMMENT{\color{gray}\small  удалить ребро $(v,u)$ }
  \ENDIF
\ENDWHILE
\end{algorithmic}
}
\end{frame}

\begin{frame}
\frametitle{10.2.2. Алгоритм построения эйлерова цикла в~эйлеровом графе (2/2)}

%{\renewcommand{\baselinestretch}{0.98}\selectfont

\begin{ground}
 Начав с
произвольной вершины $v$, строим путь, удаляя рёбра и запоминая
вершины в стеке, до тех пор, пока множество смежности очередной
вершины не окажется пустым, что означает, что путь удлинить
нельзя. Заметим, что при этом мы обязательно придём в ту
вершину, с которой начали. В противном случае вершина $v$ имеет
нечётную степень, что невозможно по условию. Таким образом, из
графа были удалены рёбра цикла, а вершины цикла были сохранены в
стеке $S$, при этом степени всех вершин остались
чётными. Далее вершина $v$ выводится в качестве первой вершины
эйлерова цикла, а процесс продолжается с вершины, стоящей на
вершине стека.
\end{ground}
%\vspace{-6pt}
%{\color{blue} Тема для сочинения: пример с картинкой.}
%
\begin{remark}
В предыдущем разделе был установлен эффективный способ проверки
наличия эйлерова цикла в графе. А именно, для этого достаточно
убедиться, что степени всех вершин чётные, что нетрудно сделать
при любом представлении графа. Приведённый алгоритм
\odmkey{}{Алгоритм (algorithm)!построения эйлерова цикла (for eulerian
cycle problem)} находит эйлеров цикл в графе, если известно, что граф
эйлеров, то есть цикл заведомо существует.
Если же граф не эйлеров и цикла не существует,
то алгоритм неприменим.
\end{remark}
%}

\end{frame}

\begin{frame}
\label{10.2.3.}\subsubsection{10.2.3. Оценка числа эйлеровых графов }\frametitle{10.2.3. Оценка числа эйлеровых графов (1/2)}


Пусть $\cal G(p)$~--- множество
всех графов с $p$ вершинами,
а ${\cal E}(p)$~--- множество эйлеровых графов с
$p$ вершинами.

\begin{remark}
В этом параграфе речь идёт о числе \odmkey{нумерованных}{Граф
(graph)!нумерованный (numbered)} графов, то есть о числе графов, в~которых вершины перенумерованы. Считается, что если перенумеровать
вершины в~другом порядке, то это будет другой граф. Очевидно, что
нумерованных графов (любого типа) больше, чем графов, определяемых
как классы эквивалентности по отношению изоморфизма, поэтому
приводимые здесь оценки являются достаточно грубыми.
\end{remark}


\begin{theorem}
Эйлеровых графов почти нет, то~есть $$
\lim\limits_{p\to\infty}\frac{\vert\cal E(p)\vert}{\vert\cal G(p)\vert} = 0.
$$
\end{theorem}


\begin{proof}
Пусть $\cal E'(p)$~--- множество графов с $p$ вершинами и чётными
степенями. Тогда по предыдущей теореме $\cal E(p)\subset\cal
E'(p)$ и $\vert\cal E(p)\vert\leq\vert\cal E'(p)\vert$. В любом
графе число вершин нечётной степени чётно, следовательно, любой
граф из $\cal E'(p)$ можно получить из некоторого графа $\cal
G(p-1)$, если добавить новую вершину и~соединить её со всеми
старыми вершинами нечётной степени. 
\end{proof}
\end{frame}

\begin{frame}
\frametitle{10.2.3. Оценка числа эйлеровых графов (2/2)}

\begin{proof} (Продолжение)
Следовательно, $|\cal
E'(p)|\leq |\cal G(p-1)|$. Но $|\cal G(p)|=2^{C(p,2)}$, поскольку
(нумерованный) граф с $p$ вершинами определяется подмножеством
включённых в него рёбер, выбранных из множества всех возможных
рёбер, а их $C(p,2)$ (лемма \DMref{п.~7.1.3}{7.1.3.}).
Заметим, что
$
\displaystyle C(k,2)-C(k-1,2) = \frac{k(k-1)}{2}-\frac{(k-1)(k-2)}{2}= k-1.
$
 


%\vspace{-12pt}
Далее имеем
$$
|\cal E(p)|\leq |\cal E'(p)|\leq |\cal
G(p-1)|=2^{C(p-1,2)}=2^{C(p,2)-(p-1)}=|\cal G(p)|2^{-(p-1)}
$$
%\addvspace{6pt}
и
$$\displaystyle{\frac{\vert\cal E(p)\vert}{\vert\cal G(p)\vert}\leq \frac{1}{2^{p-1}}},$$
откуда
$$\displaystyle{\lim\limits_{p\to\infty}\frac{\vert\cal E(p)\vert}{\vert\cal G(p)\vert}= 0}.$$
\end{proof}


\end{frame}


\subsection{10.3. Гамильтоновы циклы}
\label{10.3.}

\begin{frame}
\frametitle{10.3. Гамильтоновы циклы}

\begin{columns}[t]
\begin{column}{10cm}
Название <<гамильтонов цикл>> произошло от задачи <<Кругосветное
пу\-те\-шест\-вие>>, придуманной Гамильтоном в XIX веке: 
нужно обойти все вершины
графа, диаграмма которого показана на рисунке (в исходной
формулировке вершины были помечены названиями столиц различных
стран), по одному разу и~вернуться в исходную точку. Этот граф
представляет собой укладку додекаэдра.
\end{column}
\begin{column}{5cm}

\vspace{-20pt}
\begin{figure}[h]
\begin{center}
\includegraphics{odm10-01}
\end{center}
\end{figure}
\end{column}
\end{columns}


\medskip
\begin{tabular}{p{12mm}|p{60mm}|p{70mm}}
  № & Название параграфа
  & Ключевые термины и обозначения \\
  \hline
\DMref{10.3.1.}{10.3.1.} & Гамильтоновы циклы  &
{ Гамильтонова цепь, полугамильтонов граф}\\  
\DMref{10.3.2.}{10.3.2.} & Задача коммивояжёра  &
{ Полный перебор; $NP$-полные задачи}\\  
\DMref{10.3.3.}{10.3.3.} & Теорема Поша &
{  }\\  
\end{tabular}  

\end{frame}

\begin{frame}
\label{10.3.1.}\subsubsection{10.3.1. Гамильтоновы графы }\frametitle{10.3.1. Гамильтоновы графы (1/3)}



Если граф имеет  простой цикл, содержащий все вершины графа (по
одному разу), то такой цикл называется \odmkey{гамильтоновым}{Цикл
(cycle)!гамильтонов (hamiltonian)} циклом, а граф называется
\odmkey{гамильтоновым}{Граф (graph)!гамильтонов (hamiltonian)}
графом. Если граф имеет  простую цепь, содержащую все вершины графа (по
одному разу), то такая цепь называется \odmkey{гамильтоновой}{Цепь
(chain)!гамильтонова (hamiltonian)} цепью, а граф называется
\odmkey{полугамильтоновым}{Граф (graph)!полугамильтонов (semi-hamiltonian)}
графом. 


\medskip
Гамильтонов цикл не обязательно содержит все рёбра графа.
Ясно, что гамильтоновым может быть только связный граф.


%
%%\vspace{-2pt}
\begin{remark}
Любой граф $G$ можно превратить в гамильтонов, добавив достаточное
количество новых вершин и инцидентных им рёбер и не добавляя
рёбер, инцидентных только старым вершинам. Для этого, например,
достаточно к вершинам $\Tuple{v}{1}{p}$ графа $G$ добавить вершины
$\Tuple{u}{1}{p}$ и множество рёбер
$\{(v_i,u_i)\}\cup\{(u_i,v_{i+1})\}$, считаем $v_{p+1}=v_1$.
\end{remark}

\medskip
Простые необходимые и достаточные условия гамильтоновости графа неизвестны.\\
Известны только некоторые достаточные условия, одно из которых приведено
в~следующей теореме.

%%\vspace{-3pt}
%%\pagebreak[4]


\end{frame}

\begin{frame}
\frametitle{10.3.1. Гамильтоновы графы (2/3)}

\begin{theorem} %[Дирак]
Если $p (G)\ge 3$ и $\delta (G)\ge p/2$, то граф $G$ является
гамильтоновым.
\end{theorem}
%\vspace{-2pt}
\begin{proof}
От противного. Пусть $G$~--- не гамильтонов. Добавим к $G$ \em{минимальное}
количество новых вершин $\Tuple{u}{1}{n}$, соединяя их со \em{всеми} вершинами
$G$ так, чтобы граф $G'\Assign G+u_1+\dots +u_n$ был гамильтоновым.
Пусть $v,u_1,w,\dots, v$~--- гамильтонов цикл в графе $G'$, причём
$v\in G$, $u_1\in G'$, $u_1\not\in G$. Такая пара вершин, $v$ и
$u_1$, в гамильтоновом цикле обязательно найдется, иначе граф $G$
был бы гамильтоновым. Тогда $w\in G$, $w\not\in
\{\Tuple{u}{1}{n}\}$, иначе вершина $u_1$ была бы не нужна. Более
того, вершина $v$ не смежна с вершиной $w$, иначе вершина $u_1$
была бы не нужна. Далее, если в цикле $v,u_1,w,\dots, v',w',\dots,
v$ есть вершина~$w'$, смежная с вершиной~$w$, то вершина $v'$ не
смежна с вершиной $v$, так~как иначе можно было бы построить
гамильтонов цикл $v,v',\dots, w,w'\dots v$ без вершины $u_1$, взяв
последовательность вершин $w,\dots, v'$ в обратном порядке. 




\end{proof}

\end{frame}

\begin{frame}
\frametitle{10.3.1. Гамильтоновы графы (3/3)}


%\vspace{-2pt}
\begin{proof} (Продолжение)
\begin{figure}[h]
\begin{center}
\includegraphics[width=130mm]{10_3_1.jpg}
\end{center}
%\caption{К доказательству}
%\label{f1001}
\end{figure}
Отсюда
следует, что число вершин графа~$G'$, не смежных с $v$, не менее
числа вершин, смежных с $w$. Но для любой вершины $w$ графа $G$
в графе $G'$ имеем $d(w)\ge p/2+n$ по построению, в том числе $d(v)\geq p/2+n$. 
Общее
число вершин (смежных и не смежных с $v$, за исключением самой
вершины $v$) составляет $n+p-1$. Таким образом, имеем\\
$
n+p-1 = \overline{d}(v)+d(v)\ge
d(w)+d(v)\ge \frac{p}{2}+n+\frac{p}{2}+n=2n+p.
$\\
Следовательно, $0\ge n+1$, что противоречит тому,
что $n>0$.
\end{proof}

\medskip
\begin{corollary}
При $p\ge 3$ полный граф $K_p$ гамильтонов.
\end{corollary}

\end{frame}

\begin{frame}
\label{10.3.2.}\subsubsection{10.3.2. Задача коммивояжёра }\frametitle{10.3.2. Задача коммивояжёра (1/3)}

Рассмотрим следующую \odmkey{задачу
коммивояжёра}{Задача (problem)!коммивояжёра (travelling
salesman)}. Имеется $p$ городов, расстояния между которыми
известны. Коммивояжёр должен посетить все $p$ городов по одному
разу, вернувшись в тот, с которого начал. Требуется найти 
маршрут движения, при котором суммарное пройденное расстояние
будет минимальным.
Ясно, что задача коммивояжёра~--- это задача отыскания
кратчайшего гамильтонова цикла в нагруженном полном графе.
Можно предложить следующую простую схему решения
задачи коммивояжёра: сгенерировать все $p!$ возможных
перестановок вершин полного графа, подсчитать для каждой
перестановки длину маршрута и выбрать из них
кратчайший. Очевидно, такое вычисление
потребует не менее $O(p!)$ шагов.
Как известно, $p!$~--- быстро растущая функция. Таким образом,
решение задачи коммивояжёра описанным методом \odmkey{полного
перебора}{Полный перебор (exhaustive search)} оказывается
практически неосуществимым даже для сравнительно небольших $p$.
Более того, известно, что задача коммивояжёра принадлежит к числу
 \odmkey{NP-полных}{Задача
(problem)!\emph{NP}-полная (\emph{NP}-complete)} задач, подробное
обсуждение которых выходит за рамки этого курса.

\end{frame}

\begin{frame}
\frametitle{10.3.2. Задача коммивояжёра (2/3)}
Вкратце суть проблемы \emph{NP}-полноты сводится к следующему. В
различных областях дискретной математики, комбинаторики, логики и
т.~п. известно множество задач, принадлежащих к числу наиболее
фундаментальных, для которых, несмотря на все усилия, не удалось
найти алгоритмов решения, имеющих \pagebreak[3]полиномиальную
сложность. Более того, если бы удалось отыскать эффективный
алгоритм решения хотя бы одной из этих задач, то из этого
немедленно следовало бы существование эффективных алгоритмов для
всех остальных задач данного класса. На этом основано общепринятое
мнение, что таких алгоритмов не существует.

\begin{digress} 
Полезно сопоставить задачи отыскания эйлеровых и гамильтоновых
циклов. Внешне
формулировки этих задач очень похожи, однако они оказываются
принципиально различными с точки зрения практического применения.
 Эйлером получено просто проверяемое необходимое и
достаточное условие существования в графе эйлерова цикла. Что
касается гамильтоновых графов, то для них неизвестны простые
необходимые и достаточные условия. 
\end{digress}
\end{frame}

\begin{frame}
\frametitle{10.3.2. Задача коммивояжёра (3/3)}
\begin{digress} (Продолжение)
На основе необходимого и
достаточного условия существования эйлерова цикла можно построить
эффективные алгоритмы отыскания такого цикла. В то же время задача
проверки существования гамильтонова цикла оказывается\\
\emph{NP}-полной (так~же, как и задача коммивояжёра). Далее,
известно, что почти нет эйлеровых графов, и эффективный алгоритм
отыскания эйлеровых циклов редко оказывается применимым на
практике. 
С другой стороны, можно показать, что почти все
графы~--- гамиль\-то\-новы, то есть
$$
\lim_{p\rightarrow\infty} \frac{|\mathcal{H}(p)|}{|\mathcal{G}(p)|} = 1,
$$
где $\mathcal{H}(p)$~--- множество гамильтоновых графов с $p$ вершинами, а
$\mathcal{G}(p)$~--- множество всех графов с $p$ вершинами.
Таким образом, задача отыскания гамильтонова цикла
или задача коммивояжёра являются
практически востребованными,
но
эффективный алгоритм решения
для них неизвестен (и, скорее всего, не существует).
\end{digress}

\end{frame}



\begin{frame}
\label{10.3.3.}\subsubsection{10.3.3. Теорема Поша }\frametitle{10.3.3. Теорема Поша (1/4)}
В степенной последовательности гамильтонова графа не может быть 
много вершин с малыми степенями. Введём обозначение 
$N_{=}(G,n)$ числа вершин графа
$G$ степени $n$, $0\le n <p$:
$$N_{=}(G,n)\Def\left\vert\Set{v\in V}{d(v)=n}\right\vert.$$
В этих обозначениях
$$
N_{\le}(G,n)\Def\left\vert\Set{v\in V}{d(v)\le n}\right\vert
=\sum\limits_{i=0}^n N_{=}(G,i),
$$
$$
N_{>}(G,n)\Def\left\vert\Set{v\in V}{d(v)> n}\right\vert
=\sum\limits_{i=n+1}^{p-1} N_{=}(G,i),
$$
$$
N_{\le}(G,n)+N_{>}(G,n)=p.
$$
\begin{lemma}
$\A{G(V,E)}{\A{e\not\in E}{\A{0\le n <p}{N_{\le}(G+e,n)\le N_{\le}(G,n)}}}.$
\end{lemma}

Другими словами, 
добавление рёбер разве что уменьшает число вершин ограниченной степени.
\end{frame}



\begin{frame}
\frametitle{10.3.3. Теорема Поша (2/4)} 
Следующая теорема даёт достаточное условие того, 
что граф является гамильтоновым. 
Она обобщает результаты, полученные ранее Оре  и Дираком, 
которые оказываются её следствиями.

\begin{theorem} \NB{[Поша]}
Если $G (V, E)$~--- связный граф  с числом вершин $p\ge 3$,\\
и $\A{n,1\le n< (p-1)/2}{N_{\le}(G,n)<n}$,\\
и для нечётных $p$ выполняется неравенство $N_{\le}(G,(p-1)/2)\le (p-1)/2$,\\   
то граф $G$ гамильтонов.
\end{theorem}

\begin{proof}
От противного.
Пусть $G$~---  не гамильтонов граф с $p$ вершинами, 
удовлетворяющий условиям теоремы. Доказательство содержит три шага.
\proofsection{\color{blue}\bf 1} 
По лемме добавление рёбер не нарушает неравенств 
в условии теоремы.
Добавим в $G$ все возможные рёбра так,
чтобы он оставался не гамильтоновым,
то есть рассмотрим \em{максимальный} не гамильтонов надграф $G$. 
Тогда, поскольку добавление к перестроенному 
максимальному графу $G$ произвольного ребра 
приводит к гамильтонову графу, 
любые две несмежные вершины соединимы простой 
гамильтоновой (остовной, то есть   
содержащей все вершины графа) цепью.
\end{proof}
\end{frame}

\begin{frame}
\frametitle{10.3.3. Теорема Поша (3/4)}
\begin{proof} (Продолжение)
 
\proofsection{\color{blue}\bf 2}
Покажем, что всякая вершина, 
степень которой \em{не меньше} $(p-1)/2$, 
смежна с каждой вершиной со степенью, 
\em{большей чем} $(p-1)/2$. 
Допустим (не теряя общности), что $d(v_1) \ge (p-1)/2$ и 
$d(v_p) \ge p/2 > (p-1)/2$, 
но вершины $v_1$ и $v_p$ не смежны. 
Тогда существует гамильтонова цепь 
$\langle v_1, v_2, \dots, v_p\rangle$, 
соединяющая $v_1$ и $v_p$.
Все вершины, смежные с $v_1$, находятся 
в этой же цепи, поскольку цепь гамильтонова. 
Пусть $v_1$ смежна с $v_j$.
Тогда вершина $v_p$ не смежна с вершиной $v_{j-1}$, 
поскольку иначе в $G$ был бы гамильтонов цикл 
$v_1, v_2,\dots, v_{j-1}, v_p, v_{p-1},\dots, v_j, v_1$.
\begin{center}
\includegraphics[width=80mm]{10_3_3.jpg}
\end{center}
Положим $n\Assign d(v_1)\le (p-1)/2$.
Тогда среди всех $p$ вершин по меньшей мере $n$ вершин,
предшествующих на гамильтоновом пути вершинам, смежным с $v_1$,
не смежны с $v_p$.   
Имеем противоречие:
$p/2 \le d(v_p) \le p-1-n \le p-1-(p-1)/2 = (p-1)/2$. 
\end{proof}
\end{frame}

\begin{frame}
\frametitle{10.3.3. Теорема Поша (4/4)}
\begin{proof} (Окончание) 
{\color{blue}\bf [3]} 
Если $\A{v\in V}{d(v)\ge p/2}$, то граф гамильтонов. 
Значит в $G$ есть вершины со степенями меньше, чем  $p/2$. 
Пусть $U\Assign\Set{v\in V}{d(v)<p/2}\ne\emptyset$ и $W\Assign V\setminus U$.
Заметим, что $\A{w\in W}{d(w)\ge p/2}$.
Выберем вершину $v_1\in U$ так, что $m\Assign d(v_1) =\max_{u\in U} d(u)$.
По условию теоремы $\vert U\vert =N_{\le}(G,m) <m <p/2$,
и значит $\vert W \vert\ge p/2$ и среди вершин $W$ есть вершина 
$v_p$ не смежная с $v_1$.
Тогда существует гамильтонова цепь 
$\langle v_1, v_2, \dots, v_p\rangle$, 
соединяющая $v_1$ и $v_p$.
Как показано в п.~2 $v_{j-1}\not\in \Gamma(v_p)$,
в противном случае нашёлся бы гамильтонов цикл.
Но тогда $m<(p-1)/2$, потому что если 
$m\ge(p-1)/2$, то имеем 
$$p/2 \le d(v_p) \le p-1-m \le p-1-(p-1)/2 = (p-1)/2.$$
По условию теоремы $N_{\le}(G,m)<m$ и значит хотя бы одна из вершин
$v_{j-1}$, предшествующих вершинам, смежным с $v_1$, имеет степень больше $p/2$,
$d(v_{j-1})\ge p/2$.
Получили пару несмежных вершин, $v_{j-1}$ и $v_p$, таких, что 
$d(v_{j-1}) \ge p/2$ и $d(v_p)\ge p/2$. Это противоречит п. 2.  
\end{proof}


\begin{corollary}
Если $p\geq 3$ и $d(u)+d(v)\geq p$ для любой пары 
$u$ и $v$ несмежных вершин графа $G$, то $G$~--- гамильтонов граф.
\end{corollary}

\end{frame}

\begin{frame}
\frametitle{Онтология теории графов}
\begin{center}
\includegraphics[width=10cm]{ODM_10_3.png}
\end{center}

%\smallskip
%{\color{red} Тема для сочинения: перечисление графов.
%Любые простые формулы или оценки числа графов разных типов, помимо тех, что уже есть в курсе, 
%очень желательны.}

%\smallskip
%{\color{red} Тема для сочинения: эффективный 
%приближённый алгоритм решения задачи коммивояжёра.
%Понятия приближённого и частичного алгоритма.}
%
%\smallskip
%{\color{red} Тема для сочинения: $P=NP?$ 
%в стиле наивной теории алгоритмов.}

\end{frame}



\subsection{10.4. Независимые и покрывающие множества}
\label{10.4.}

\begin{frame}
\frametitle{10.4. Независимые и покрывающие множества (1/2)}
В этом разделе вводятся определения и рассматриваются
основные свойства независимых и покрывающих
множеств вершин и рёбер. Эти определения и свойства используются
в последующих разделах.

\medskip
Значительное количество \em{экстремальных} задач в различных 
предметных областях сводится к отысканию наибольших
независимых или наименьших покрывающих множеств вершин и рёбер
в графах, входящих в модели этих предметных областей.
Поэтому данный раздел теории графов имеет
множество практических приложений.

\medskip
Обсуждение \em{переборных} алгоритмов решения экстремальных задач
представляет особый интерес.      
\end{frame}

\begin{frame}
\frametitle{10.4. Независимые и покрывающие множества (2/2)}
\medskip
\begin{tabular}{p{12mm}|p{60mm}|p{70mm}}
  № & Название параграфа
  & Ключевые термины и обозначения \\
  \hline
\DMref{10.4.1.}{10.4.1.} & Независимые и покрывающие множества вершин и рёбер  &
{ Вершинное покрытие, число вершинного покрытия 
$\alpha_0$, рёберное покрытие, число рёберного покрытия $\alpha_1$, независимое множество вершин,
внутренне устойчивое множество вершин, вершинное число независимости $\beta_0$,
независимое множество рёбер, паросочетание, рёберное число независимости $\beta_1$}\\  
\DMref{10.4.2.}{10.4.2.} & Связь чисел независимости и покрытий  &
{ }\\  
\DMref{10.4.3.}{10.4.3.} & Оценка числа вершинной независимости &
{  Средняя степень графа}\\  
\DMref{10.4.4.}{10.4.4.} & Теорема Кёнига &
{  }\\  \end{tabular}  
\end{frame}

\begin{frame}
\label{10.4.1.}\subsubsection{10.4.1. Независимые и покрывающие множества вершин и рёбер }
\frametitle{10.4.1. Независимые и покрывающие множества вершин и рёбер (1/5)}


Говорят, что вершина \odmkey{покрывает}{Вершина
(vertex)!покрывающая (covering)} инцидентные ей рёбра, а ребро
\odmkey{покрывает}{Ребро (edge)!покрывающее (covering)}
инцидентные ему вершины. Множество вершин, которые в совокупности
покрывают все рёбра, называется \odmkey{вершинным
покрытием}{Покрытие (covering)!вершинное (vertex)}. Вершинное покрытие называется 
\odmkey{минимальным}{Покрытие (covering)!минимальное (minimal)!вершинное (vertex)}, если никакое его подмножество не является вершинным покрытием. Вершинное покрытие называется 
\odmkey{наименьшим}{Покрытие (covering)!наименьшее (minimum)!вершинное (vertex)}, если число элементов в нём наименьшее возможное, а мощность такого покрытия называется \odmkey{числом вершинного покрытия}{Число (number)!вершинного покрытия (of vertex covering)} и обозна\-ча\-ется~$\alpha_0$.

\begin{examples}
\begin{enumerate}
\item $\alpha_0(K_p)=p-1$, если $K_p$~--- полный граф. 
\item
$\alpha_0(K_{m,n})=\min(m,n)$, если $K_{m,n}$~--- полный
двудольный граф. 
\item $\alpha_0(\overline K_p)= 0$, если
$\overline K_p$~--- вполне несвязный граф. 
\item $\alpha_0(G_1\cup
G_2)=\alpha_0(G_1)+\alpha_0(G_2)$, если $G_1$ и $G_2$~---
компоненты несвязного графа.
\end{enumerate}
\end{examples}

\begin{digress}
Задача поиска наименьшего вершинного покрытия имеет наглядную практическую интерпретацию: пусть граф описывает план охраняемого объекта, при этом рёбрам соответствуют просматриваемые пространства, например, коридоры, а вершинам~--- возможные точки наблюдения (блокпосты). Необходимо расставить минимальное количество охранников на блокпостах так, чтобы ими контролировались все коридоры.
\end{digress}



\end{frame}


\begin{frame}
\frametitle{10.4.1. Независимые и покрывающие множества вершин и рёбер (2/5)}
Множество рёбер, которые в совокупности покрывают
все вершины, называется \odmkey{рёберным покрытием}{Покрытие (covering)!рёберное (edge)}. Рёберное покрытие называется 
\odmkey{минимальным}{Покрытие (covering)!минимальное (minimal)!рёберное (edge)}, если никакое его подмножество не является рёберным покрытием. Рёберное покрытие называется \odmkey{наименьшим}{Покрытие (covering)!наименьшее (minimum)!рёберное (edge)}, если число элементов в нём наименьшее возможное, а мощность такого покрытия называется \odmkey{числом рёберного
покрытия}{Число (number)!рёберного покрытия (of edge covering)} и
обозначается $\alpha_1$.



\begin{remark}
Для графа с изолированными вершинами
$\alpha_1$ не определено.
\end{remark}


\begin{examples}
\begin{enumerate}
\item
$\alpha_1(C_{2n})=n$, если $C_{2n}$~--- чётный цикл.
\item
$\alpha_1(C_{2n+1})=n+1$, если $C_{2n+1}$~--- нечётный цикл.
\item
$\alpha_1(K_{2n})=n$, если $K_{2n}$~--- полный граф с чётным числом вершин.
\item
$\alpha_1(K_{2n+1})=n+1$, если $K_{2n+1}$~---
полный граф с нечётным числом вершин.
\item
$\alpha_1(K_{m,n})=\max(m,n)$, если $K_{m,n}$~--- полный двудольный граф.
\end{enumerate}
\end{examples}
\end{frame}

%
\begin{frame}
\frametitle {10.4.1. Независимые и покрывающие множества вершин и рёбер (3/5)}
Множество вершин называется \odmkey{независимым}{Множество
(set)!вершин (of vertex)!независимое (independent)} (или
\odmkey{внутренне устойчивым}{Множество (set)!вершин (of
vertex)!внутренне устойчивое (internally stable)}), если никакие
две из них не смежны. Независимое множество вершин называется \odmkey{максимальным}{Множество
(set)!вершин (of vertex)!независимое (independent)! максимальное (maximal}, если никакое его надмножество не является независимым. 
Независимое множество вершин называется \odmkey{наибольшим}{Множество
(set)!вершин (of vertex)!независимое (independent)! наибольшее (maximum}, если число элементов в нём наибольшее возможное, а мощность такого множества называется \odmkey{вершинным числом
независимости}{Число (number)!независимости (of independence)!вершинное (vertex)} и обозначается~$\beta_0$.



%\vspace{-6pt}
\begin{examples}
\vspace{-3pt}
\begin{enumerate}
\item $\beta_0(K_p)= 1$, если $K_p$~--- полный граф. \item
$\beta_0(K_{m,n})=\max(m,n)$ , если $K_{m,n}$~--- полный
двудольный граф. \item $\beta_0(\overline K_p)= p$, если
$\overline K_p$~--- вполне несвязный граф. \item $\beta_0(G_1\cup
G_2)=\beta_0(G_1)+\beta_0(G_2)$ , если $G_1$ и $G_2$~---
компоненты несвязного графа.
\end{enumerate}
\end{examples}

\begin{remark}
Практической интерпретацией этого понятия может служить следующий пример: пусть имеется множество ресурсов и множество проектов, каждому из которых требуется некоторое подмножество ресурсов для выполнения. Построим граф, в котором вершины будут отвечать проектам, а ребра будут обозначать наличие общих ресурсов у проектов. Тогда наибольшее независимое множество вершин графа будет представлять наибольшее множество проектов, которые можно выполнить одновременно.
\end{remark}


\end{frame}


\begin{frame}
\frametitle {10.4.1. Независимые и покрывающие множества вершин и рёбер (4/5)}
Ясно, что множество вершин $S$ является независимым тогда и только
тогда, когда $$\A{v\in S}{\Gamma(v)\cap S = \emptyset}.$$


Множество рёбер называется \odmkey{независимым}{Множество
(set)!рёбер (of edges)!независимое (independent)} (или
\odmkey{паросочетанием}{Паросочетание (matching)}), если никакие
два из них не смежны. Наибольшее число рёбер в независимом
множестве рёбер называется \odmkey{рёберным числом
независимости}{Число (number)!независимости (of independence)!рёберное (edge)} и обозначается $\beta_1$. (Понятия \odmkey{максимального}{Паросочетание (matching)! максимальное (maximal)} и \odmkey{наибольшего}{Паросочетание (matching)! наибольшее (maximum)} паросочетания были введены в разделе \DMref{8.3}{8.3.}).



\begin{examples}
\begin{enumerate}
\item
$\beta_1(C_{2n})=n$, если $C_{2n}$~--- чётный цикл.
\item
$\beta_1(C_{2n+1})=n$, если $C_{2n+1}$~--- нечётный цикл.
\item
$\beta_1(K_{2n})=n$, если $K_{2n}$~--- полный граф с чётным числом вершин.
\item
$\beta_1(K_{2n+1})=n$, если $K_{2n+1}$~---
полный граф с нечётным числом вершин.
\item
$\beta_1(K_{m,n})=\min(m,n)$, если $K_{m,n}$~--- полный двудольный граф.
\end{enumerate}
\end{examples}


\begin{remark}
Индексы 0 и 1 в обозначениях $\alpha_0$,  $\alpha_1$, $\beta_0$,  $\beta_1$
имеют следующую мнемонику.
Индекс 0 соответствует вершинам,
так~как вершина нульмерна, а~1~--- рёбрам, так~как ребро одномерно.
\end{remark}

\end{frame}


\begin{frame}
\frametitle {10.4.1. Независимые и покрывающие множества вершин и рёбер (5/5)}
Ясно, что наибольшее независимое множество
вершин или рёбер является максимальным,
а наименьшее покрывающее множество 
вершин или рёбер 
является минимальным.
Обратное неверно:
независимое множество может быть 
максимальным, не будучи наибольшим,
а покрывающее множество может быть 
минимальным, не будучи наименьшим.

\begin{example}
На диаграмме слева 
вершины $\{v_1, v_4\}$ образуют наименьшее
покрытие, а вершины $\{v_2, v_3, v_5\}$
образуют наибольшее независимое множество.
Cправа 
вершины $\{v_1, v_4\}$ образуют максимальное, но не наибольшее независимое множество, 
а вершины \\$\{v_2, v_3, v_5\}$
образуют минимальное, но не наименьшее покрывающее множество.

\vspace{-2pt}
\begin{center}
\includegraphics[width=90mm]{10_4_1.jpg}
\end{center}
\end{example} 

\end{frame}

\begin{frame}
\label{10.4.2.}\subsubsection{10.4.2. Связь чисел независимости и покрытий }\frametitle{10.4.2. Связь чисел независимости и покрытий (1/3)}


%{\renewcommand{\baselinestretch}{0.98}\selectfont
Приведённые примеры наводят на мысль, что числа независимости и
покрытия связаны друг с другом и с количеством вершин $p$.


\begin{theorem}
Для любого нетривиального связного графа
$$
\alpha_0 + \beta_0 = p = \alpha_1 + \beta_1.
$$
\end{theorem}



\begin{proof}
Докажем четыре неравенства, из которых следуют
два требуемых равенства.

\proofsection{$\alpha_0+\beta_0\ge p$}
Пусть $M_0$~--- наименьшее вершинное покрытие,
то~есть $|M_0|=\alpha_0$. Рассмот\-рим $V\setminus M_0$.
Тогда $V\setminus M_0$~--- независимое множество,
так~как если бы в~множестве $V\setminus M_0$ были
смежные вершины, то $M_0$
не было бы покрытием.
Имеем $|V\setminus M_0|\le \beta_0$,
следовательно,
$p=|M_0|+|V\setminus M_0| \le \alpha_0 + \beta_0$.

\proofsection{$\alpha_0+\beta_0\le p$}
Пусть $N_0$~--- наибольшее независимое множество вершин,
то~есть $|N_0|=\beta_0$. Рассмотрим $V\setminus N_0$.
Тогда $V\setminus N_0$~--- вершинное покрытие,
так~как нет рёбер, инцидентных только вершинам из
$N_0$, стало быть, любое ребро инцидентно~вершине
(или вершинам) из
$V\setminus N_0$.
Имеем $|V\setminus N_0|\ge \alpha_0$,
следовательно, %\linebreak
$p=|N_0|+|V\setminus N_0| \ge \beta_0 + \alpha_0$.
\end{proof}
%}
\end{frame}

\begin{frame}
\frametitle{10.4.2. Связь чисел независимости и покрытий (2/3)}
\begin{proof} (Продолжение)
\proofsection{$\alpha_1+\beta_1\ge p$} Пусть $M_1$~--- наименьшее
рёберное покрытие, то~есть $|M_1|=\alpha_1$. Множество $M_1$ не
содержит цепей длиной больше 2. Действительно, если бы в~$M_1$
была цепь длиной 3, то среднее ребро этой цепи можно было бы
удалить из $M_1$ и это множество все равно осталось бы покрытием.
Следовательно, $M_1$ состоит из звёзд (\DMref{п.~7.3.3}{7.3.3.}). Пусть этих звёзд
$m$. Имеем $|M_1|=\sum_{i=1}^{m}n_i$, где $n_i$~--- число рёбер в
$i$-й звезде. Заметим, что звезда из $n_i$ рёбер покрывает $n_i+1$
вершину. \\Имеем
$p=\sum_{i=1}^{m}(n_i+1)=m+\sum_{i=1}^{m}n_i=m+|M_1|$. Возьмём по
одному ребру из каждой звезды и составим из них множество~$X$.
Тогда ${\vert X\vert=m}$, множество $X$~--- независимое, то~есть
$|X|\le\beta_1$. Следовательно, ${p=|M_1|+m} =|M_1|+|X|\le
\alpha_1+\beta_1$.
\proofsection{$\alpha_1+\beta_1\le p$} Пусть $N_1$~--- наибольшее
независимое множество рёбер, то~есть $|N_1|=\beta_1$. Построим
рёберное покрытие $Y$ следующим образом. Множество $N_1$ покрывает
$2|N_1|$ вершин. Добавим по одному ребру, инцидентному непокрытым
$p-2|N_1|$ вершинам, таких рёбер $p-2|N_1|$. Тогда множество
$Y$~--- рёберное покрытие,  $|Y| = |N_1| + p - 2|N_1| = p - |N_1|$ и
$|Y|\ge\alpha_1$. Имеем ${p = p - |N_1| + |N_1|}
  = {|Y| + |N_1| \ge \alpha_1 + \beta_1}$.
\end{proof}
\end{frame}

\begin{frame}
\frametitle{10.4.2. Связь чисел независимости и покрытий (3/3)}

%\pagebreak
\begin{remark}
Условия связности и нетривиальности
гарантируют, что все четыре инварианта
определены. Хотя
это условие является достаточным,
оно не является необходимым.
Например, для графа $K_n\cup K_n$ заключение
теоремы остается справедливым, хотя условие не выполнено.
\end{remark}


\begin{digress}
Задача отыскания экстремальных независимых  множеств
возникает  во  
многих  случаях.  
Например,  пусть  дано 
множество  процессов, использую\-щих неразделяемые ресурсы.
Соединим рёбрами вершины, соответствующие процессам, которым
требуется один и тот же ресурс. Тогда $\beta_0$  определяет
количество возможных параллельных процессов.
\end{digress}

\end{frame}

\begin{frame}
\label{10.4.3.}\subsubsection{10.4.3. Оценка числа вершинной независимости }\frametitle{10.4.3. Оценка числа вершинной независимости (1/3)}
\begin{theorem}
Для любого графа $G(V,E)$ справедливо неравенство: 
$\displaystyle{\beta_0(G)\geq \sum\limits_{v\in V} \frac{1}{1+d(v)}.}$ 
\end{theorem}

\begin{proof}
Заметим, что если $G=K_p$, то $\beta_0(K_p) = 1$ 
и \\ $\displaystyle{\sum_{v\in V} \frac{1}{1+d(v)}
= \sum_{v\in V} \frac{1}{1+p-1} = 1,}$ 
то есть достигается равенство. 
Далее индукция по $p$. 
База для $p=1$ проверена. 
Ввиду того, что для полных графов достигается равенство, 
можно считать, что $G\ne K_p$.
Пусть теперь $p>2$ и для всех графов с числом вершин меньше $p$
неравенство выполнено. Рассмотрим в графе $G$ 
вершину $u$ минимальной степени: $d(u)=\delta(G)$.
Так как $G\ne K_p$, имеем $\Gamma^*(u)\ne V$.
Удалим вершину $u$ и её окрестность из 
множества вершин $V'\Assign V\setminus \Gamma^*(u)$,
и рассмотрим граф $G'$, который является правильным подграфом 
графа $G$, определяемым множеством $V'$. Степень вершины 
в графе $G'$ обозначим $d'$. Заметим, что
$\A{v\!\in\! V'}{d'(v)\!\leq\! d(v)}$ и  
$\A{v\!\in\! V'}{(1\!+\!d'(v))^{-1}\!\geq\! (1\!+\!d(v))^{-1}}$.
Далее, если $M$~--- независимое множество вершин в графе $G'$,
то $M+u$~--- независимое множество вершин в графе $G$,
и $\beta_0(G)\geq \beta_0(G')+1$.   
\end{proof}
\end{frame}


\begin{frame}
\frametitle{10.4.3. Оценка числа вершинной независимости (2/3)}
\begin{proof} (Продолжение)
Заметим теперь, что в силу выбора вершины $u$ имеем
$\displaystyle{\sum\limits_{v\in \Gamma^*(u)} (1+d(v))^{-1}
\leq \sum\limits_{v\in \Gamma^*(u)} (1+d(u))^{-1} =1}$.
Окончательно получаем
\begin{align*}
\beta_0(G)&\geq \beta_0(G')+1 \geq
\sum\limits_{v\in V'} (1+d'(v))^{-1} +1 \geq \\
&\geq\sum\limits_{v\in V'} (1+d(v))^{-1} +1 \geq
\sum\limits_{v\in V'} (1+d(v))^{-1} + 
\sum\limits_{v\in \Gamma^*(u)} (1+d(v))^{-1}\geq \\
&\geq\sum\limits_{v\in V} (1+d(v))^{-1}.
\end{align*}
\end{proof}
\begin{remark}
Полезно выделить крайние случаи для пустого и полного графа: 
	$$ \beta_0(\overline{K_p})=\sum_{i=1}^{p} \frac{1}{1} = p,  \qquad \beta_0(K_p)=\sum_{i=1}^{p} \frac{1}{1+p-1} = 1 .$$
\end{remark}


\end{frame}


\begin{frame}
\frametitle{10.4.3. Оценка числа вершинной независимости (3/3)}

\smallskip
Рассмотрим \odmkey{среднюю степень}{} графа: 
$\displaystyle{\bar{d}(G)\Def\frac{1}{p}\sum_{v\in V} d(v)}$. 

\medskip
\begin{corollary}
Для любого графа $G(V,E)$ справедливо неравенство: 
$\displaystyle{\beta_0(G)\geq \frac{p}{1+\bar{d}}.}$ 
\end{corollary}
\begin{proof}
Известно \odmkey{неравенство Коши--Буняковского}{Неравенство (inequality)!Коши-Буняковского \\
(Cauchy--Bunyakovsky--Schwarz)}
$$\left(\sum\limits_{i=1}^n a_i b_i\right)^2 \leq 
\left(\sum\limits_{i=1}^n a_i^2\right)
\left(\sum\limits_{i=1}^n b_i^2\right).$$
 Положим
$a_i\Assign \sqrt{(1+d(v_i))^{-1}}$, $b_i\Assign \sqrt{1+d(v_i)}$.
Тогда
$$p^2 = \left(\sum\limits_{i=1}^p \sqrt{\frac{1+d(v_i)}{1+d(v_i)}}\right)^2
\leq \beta_0 \left(\sum\limits_{i=1}^p (1+d(v_i))\right),$$ и $p^2\leq\beta_0(p+p\bar{d})$, то есть $p\leq\beta_0(1+\bar{d})$.
\end{proof}

\end{frame}

\begin{frame}
\label{10.4.4.}\subsubsection{10.4.4. Теорема Кёнига}
\frametitle{10.4.4. Теорема Кёнига}

\begin{theorem}
В любом двудольном графе $\beta_1(G) = \alpha_0(G).$
\end{theorem}
\begin{proof}
Добавим к вершинам графа $G$ две вершины $u$ и $v$ и соединим с $u$ все
вершины левой доли, а с $v$~--- правой доли. В полученном 
графе $G'$ любое подмножество вершин разделяет вершины $u$ и $v$ тогда и только 
тогда, когда оно покрывает все рёбра исходного графа (белые вершины на рисунке). Следовательно, $\min\vert S(u,v)\vert = \alpha_0(G).$ Кроме того, любое 
множество вершинно-непересекающихся простых $\divC{u}{v}$-цепей образует 
некоторое паросочетание в графе $G$ (жирные рёбра), а любое паросочетание в графе $G$ может быть
продолжено до множества вершинно-непересекающихся простых $\divC{u}{v}$-цепей.
Следовательно, $\max\vert P(u,v)\vert = \beta_1(G)$. Тогда, по теореме Менгера:\\
$\beta_1(G) = \alpha_0(G).$
\end{proof}

\vspace{-4pt}
\begin{center}
\includegraphics[width=65mm]{10_4_4.png}
\end{center}

\end{frame}


\subsection{10.5. Построение независимых множеств вершин}
\label{10.5.}

\begin{frame}
\frametitle{10.5. Построение независимых множеств вершин}
Этот раздел фактически является вводным обзором методов решения
переборных задач. Методы рассматриваются на примере задачи
отыскания наибольшего независимого множества вершин графа. Наш
обзор не претендует на полноту, но описание основополагающих идей
и терминов в нём присутствует.

\medskip
\begin{tabular}{p{12mm}|p{60mm}|p{70mm}}
  № & Название параграфа
  & Ключевые термины и обозначения \\
  \hline
\DMref{10.5.1.}{10.5.1.} & Постановка задачи отыскания наибольшего независимого множества вершин  &
{ Эвристический алгоритм }\\  
\DMref{10.5.2.}{10.5.2.} & Поиск с возвратами  &
{ Бэктрекинг}\\  
\DMref{10.5.3.}{10.5.3.} & Улучшенный перебор &
{  Хранение контекста}\\  
\DMref{10.5.4.}{10.5.4.} & Алгоритм построения независимых множеств &
{  }\\ 
 \end{tabular}  
\end{frame}


\begin{frame}
\label{10.5.1.}\subsubsection{10.5.1. Постановка задачи отыскания наибольшего независимого множества вершин}
\frametitle{10.5.1. Постановка задачи отыскания наибольшего независимого множества вершин (1/3)}

Задача отыскания наибольшего независимого множества вершин (и тем
самым определения вершинного числа независимости $\beta_0$)
принадлежит к числу трудоёмких.
Эту задачу можно поставить следующим образом.
Пусть задан граф $G(V,E)$.
Найти такое множество вершин $X$, $X\subset V$, что
$|X|=\max\limits_{Y\in \cal E} |Y|$,
а $\cal E\Def\Set{Y\subset V}{\A{u,v\in Y}{(u,v)\not\in E}}$.
Если бы семейство $\cal E$ независимых множеств оказалось
матроидом, то поставленную задачу можно было бы
решить жадным алгоритмом.
Однако семейство $\cal E$ матроидом не является.
Действительно, хотя аксиомы $M_1$ и $M_2$ (см. \DMref{п.~2.7.1}{2.7.1.})
выполнены, так~как пустое множество вершин независимо
и всякое подмножество независимого множества независимо,
но аксиома $M_3$ может быть не выполнена.

\begin{example}
Рассмотрим полный двудольный граф $K_{n,n+1}$.
Доли этого графа образуют независимые множества,
и их мощности различаются на~1. Однако никакая вершина из большей
доли не может быть добавлена к~вершинам меньшей доли с сохранением
независимости.
\end{example}

\end{frame}


\begin{frame}
\frametitle{10.5.1. Постановка задачи отыскания наибольшего
независимого множества вершин (2/3)}


Нижняя оценка числа вершинной независимости $\beta_0$,
полученная в \DMref{п.~10.4.3}{10.4.3.}
 оказывается точным
значениям для крайних случаев~--- для полных и для пустых графов.
На первый взгляд кажется разумным предположить,
что эта оценка <<достаточно точна>> и для других графов.
В таком случае, на основе конструктивного
доказательства нижней оценки, возникает соблазн предложить эффективный
\odmkey{эвристический}{Алгоритм (algorithm)!эвристический (heuristic)} алгоритм поиска независимого
множества, <<близкого>> к наибольшему.  

\vspace{-2pt}
\begin{algorithmic}
\REQUIRE граф $G(V,E)$.
\ENSURE наибольшее независимое множество $X$.
\STATE $X\Assign \emptyset$
\COMMENT{\color{gray}\small вначале независимое множество пусто }
\WHILE{$G\ne \emptyset$}
\STATE {\textbf{select $u\in V \And d(u) = \delta(G)$}}
\COMMENT{\color{gray}\small $u$ имеет минимальную степень }
\STATE {$X\Assign X + u$}
\COMMENT{\color{gray}\small помещаем её в независимое множество }
\STATE {$G\Assign G\setminus \Gamma^*(u)$}
\COMMENT{\color{gray}\small удаляем вершину $u$ вместе с её окрестностью }
\ENDWHILE
\end{algorithmic}
\end{frame}


\begin{frame}
\frametitle{10.5.1. Постановка задачи отыскания наибольшего
независимого множества вершин (3/3)}

\begin{example}
На рисунке слева приведена диаграмма графа, 
для которого эвристический алгоритм находит точное решение
$\beta_0 = 2$, а справа приведена диаграмма графа,
для которого алгоритм даёт ответ $2$, 
в то время, как наибольшим независимым множеством
является доля в центре и $\beta_0 =3$.
\vspace{-12pt}
\begin{figure}[h]
\begin{center}
\includegraphics[width=90mm]{10_5_1.jpg}
\end{center}
%\caption{К доказательству}
%\label{f1001}
\end{figure}
\vspace{-12pt}
\end{example}

Приведённый пример легко обобщить.
Рассмотрим граф $\overline{K}_m + K_n$ и 
добавим к нему одну вершину, соединив её с вершинами доли 
$\overline{K}_m$. Тогда при $n\ge m > 2$ 
наибольшим независимым множеством является средняя доля $\overline{K}_m$
и $\beta_0 =m$, в то время как эвристический 
алгоритм по-прежнему выдаёт ответ $2$.
Эвристический алгоритм может дать ответ, сколь угодно 
далёкий от точного решения!
   


\end{frame}


\begin{frame}
\label{10.5.2.}\subsubsection{10.5.2. Поиск с возвратами }\frametitle{10.5.2. Поиск с возвратами (1/5)}

Даже если для решения задач, подобных поставленной
в предыдущем параграфе, не удаётся найти эффективного
алгоритма, остаётся возможность попробовать
найти решение <<полным перебором>> всех возможных
вариантов, просто в
силу конечности числа возможностей.
Например, наибольшее
независимое множество можно найти по следующей схеме.



%\vspace{6pt}

\begin{algorithmic}
\REQUIRE граф $G(V,E)$.
\ENSURE наибольшее независимое множество $X$.
\STATE $m\Assign 0$
\COMMENT{\color{gray}\small наилучшее известное значение $\beta_0$ }
\FOR{$Y\in 2^V$}
  \IF{$Y\in \cal E\And \vert Y\vert > m$}
    \STATE $m:=\vert Y\vert; X\Assign Y$
    \COMMENT{\color{gray}\small  наилучшее известное значение $X$ }
  \ENDIF
\ENDFOR
\end{algorithmic}


\begin{remark}
Для выполнения этого алгоритма потребуется $O(2^p)$ шагов.
\end{remark}




\end{frame}


\begin{frame}
\frametitle{10.5.2. Поиск с возвратами (2/5)}

При решении переборных задач большое значение имеет способ
организации перебора (в нашем случае~--- способ построения и
последовательность перечисления множеств $Y$). Наиболее популярным
является следующий способ организации перебора, основанный на идее
поиска в глубину и называемый \odmkey{поиском с~возвратами}{Алгоритм (algorithm)!поиска (search)!с возвратами
(backtracking)}.


\medskip
\begin{remark}
Иногда употребляется термин <<\odmkey{бэктрекинг}{Бэктрекинг
(backtracking)}>> (транслитерация английского названия этого
метода~--- {\it back\-tracking}). 
Буквальный перевод английского названия~---
обратное прослеживание~--- явно неудачен,
поскольку в методе нет ничего <<обратного>>,
и непонятно, что <<прослеживается>>.
\end{remark}

\medskip
Идея поиска с возвратами состоит в
следующем. Находясь в некоторой ситуации,
пробуем изменить её допустимым образом в надежде найти решение.
Если изменение не привело к успеху, то возвращаемся
в исходную ситуацию
(отсюда название <<поиск с возвратами>>)
и пробуем изменить её другим образом,
и так до тех пор, пока не будут перебраны все возможности.



\end{frame}


\begin{frame}
\frametitle{10.5.2. Поиск с возвратами (3/5)}
{\renewcommand{\baselinestretch}{0.95}\selectfont



\begin{algorithmic}
\REQUIRE граф $G(V,E)$.
\ENSURE наибольшее независимое множество $X$.
\STATE $m\Assign 0$
\COMMENT{\color{gray}\small  наилучшее известное значение $\beta_0$ }
\STATE $X\Assign \emptyset$
\COMMENT{\color{gray}\small  наибольшее известное независимое множество $X$ }
\STATE $\BT(\emptyset,V)$
\COMMENT{\color{gray}\small  вызов рекурсивной процедуры $\BT$ }
\end{algorithmic}
%\end{algorithm}
%
%\pagebreak[3]
\vspace{-3pt}
Основная работа выполняется рекурсивной процедурой $\BT$.

\vspace{-3pt}

\begin{algorithmic}
\REQUIRE $S$~--- текущее независимое множество вершин,
$T$~--- оставшиеся вершины графа.
\ENSURE изменение глобальной переменной $X$,
если текущее множество не может быть расширено
(является максимальным).
      \IF{$\vert S\vert > m$}
        \STATE $X:=S; m:=\vert S\vert$
        \COMMENT{\color{gray}\small  наибольшее известное независимое множество }
      \ENDIF
\FOR{$v\in T$}
  \IF{$S+v\in\cal E$}
    \STATE $\BT(S+v,T\setminus\Gamma^*(v))$
    \COMMENT{\color{gray}\small  пробуем добавить $v$ }
  \ENDIF
\ENDFOR
\end{algorithmic}
}


\end{frame}


\begin{frame}
\frametitle{10.5.2. Поиск с возвратами (4/5)}

\begin{ground}
Для рассматриваемой задачи отыскания наибольшего независимого
множества вершин метод поиска с возвратами реализован с
помощью указанного рекурсивного алгоритма.

По построению вершина $v$ добавляется в множество $S$ только при
сохранении независимости расширенного множества. В алгоритме это
обстоятельство указано в форме условия $S+v\in\cal E$. Проверить
сохранение условия независимости нетрудно, например, с помощью
следующей функции.

%\vspace{3pt}
%{
%\small
\begin{algorithmic}
\REQUIRE независимое множество $S$ и проверяемая вершина $v$.
\ENSURE \textbf{true}, если множество $S+v$ независимое, %\\
\textbf{false} --- в противном случае.
\FOR{$u\in S$}
\IF{$(u,v)\in E$}
\STATE \textbf{return false}
\COMMENT{\color{gray}\small  множество $S+v$ зависимое }
\ENDIF
\ENDFOR
\STATE \textbf{return true}
\COMMENT{\color{gray}\small множество $S+v$ независимое }
\end{algorithmic}
%}
\end{ground}

\end{frame}


\begin{frame}
\frametitle{10.5.2. Поиск с возвратами (5/5)}

\begin{ground} (Окончание)
Этот цикл не включен в явном виде в рекурсивную процедуру $\BT$,
чтобы не загромождать основной текст и не затуманивать идею поиска
с возвратами. 

Таким образом, множество $S$, а следовательно, и
множество $X$~--- независимые. В~тот момент, когда множество $S$
нельзя расширить, оно максимально по определению. Перемен\-ная $m$
глобальна, поэтому среди всех максимальных независимых множеств в
конце работы алгоритма построенное множество $X$ является
наибольшим независимым множеством вершин.
\end{ground}

\begin{digress}
Алгоритм, трудоёмкость которого (число шагов) ограничена полиномом
от характерного размера задачи, принято называть
\odmkey{эффективным}{Алгоритм (algorithm)!эффективный
(efficient)}, в противоположность \odmkey{неэффективным}{Алгоритм
(algorithm)!неэффективный (inefficient)} алгоритмам, трудоёмкость
которых ограничена функцией,  растущей быстрее, например,
экспонентой. Таким образом, жадный алгоритм эффективен, а полный
перебор~--- нет.
\end{digress}


\end{frame}


\begin{frame}
\label{10.5.3.}\subsubsection{10.5.3. Улучшенный перебор }\frametitle{10.5.3. Улучшенный перебор (1/3)}



Применение метода поиска с возвратами
не гарантирует эффективности~--- трудоёмкость поиска с возвратами
имеет тот же порядок,
что и другие способы перебора (в худшем случае).
Используя конкретную информацию о задаче, в некоторых случаях
можно существенно сократить трудоёмкость выполнения каждого шага
перебора или уменьшить количество перебираемых возможностей в
среднем (при сохранении оценки количества шагов в худшем случае).
Такие приёмы называются методами 
\odmkey{улучшения перебора}{Улучшение перебора (search improvement)}.
Например, может
оказаться, что некоторые варианты заведомо не могут привести к
решению, а потому их можно не рассматривать.



\begin{remark}
Рекурсивная форма метода поиска с возвратами удобна для понимания,
но не является самой эффективной. По сути, рекурсия здесь
используется для сохранения \odmkey{контекста}{Контекст
(context)}~--- то~есть информации, характеризующей текущий
рассматриваемый вариант. Если использовать другие методы
сохранения контекста, то поиск с возвратами может быть реализован
без явной рекурсии, а значит, более эффективно.
\end{remark}


\end{frame}

\begin{frame}
\frametitle{10.5.3. Улучшенный перебор (2/3)}
Рассмотрим методы улучшения перебора на примере
задачи отыскания всех максимальных независимых
множеств вершин.

Идея: начинаем с пустого множества и пополняем его
вершинами с сохранением независимости (пока возможно).


Пусть $S_k$~--- уже полученное множество из $k$ вершин,
$Q_k$~--- множество вершин, которое можно добавить к $S_k$,
то~есть $S_k \cap \Gamma(Q_k) = \emptyset$.
Среди вершин $Q_k$ будем различать те, которые уже использовались
для расширения $S_k$ (обозначим их множество $Q_k^-$),
и те, которые еще не
использовались ($Q_k^+$). Тогда
общая схема
нерекурсивной реализации поиска с возвратами будет состоять из следующих
шагов.



Шаг вперед от $k$ к $k+1$ состоит в выборе вершины $x\in Q_k^+$:
\begin{align*}
S_{k+1}   &\Assign S_k+x,\\
Q_{k+1}^- &\Assign Q_k^- \setminus \Gamma^+(x),\\
Q_{k+1}^+ &\Assign Q_k^+ \setminus\Gamma^*(x).
\end{align*}




\end{frame}

\begin{frame}
\frametitle{10.5.3. Улучшенный перебор (3/3)}

Шаг назад от $k+1$ к $k$:
\begin{align*}
S_k   &\Assign S_{k+1}-x,\\
Q_k^+ &\Assign Q_k^+ - x, \\
Q_k^- &\Assign Q_k^- + x.
\end{align*}
Если $S_k$~--- максимальное, то $Q_k^+ = \emptyset$.
Если $Q_k^- \ne \emptyset$, то
$S_k$ было расширено раньше и не является максимальным. Таким~образом,
проверка максимальности задается следующим условием:
$Q_k^+=Q_k^-=\emptyset$.



Перебор можно улучшить, если заметить следующее.



Пусть $x \in Q_k^-$ и $\Gamma(x) \cap Q_k^+ =\emptyset$. Эту
вершину $x$ никогда не удалить из $Q_k^-$, так~как из $Q_k^-$
удаляются только вершины, смежные с $Q_k^+$. Таким образом,
существование~$x$, такого, что $x\in Q_k^-$ и $\Gamma(x) \cap
Q_k^+ = \emptyset$, является достаточным условием для возвращения.
Кроме того, $k\le p-1$.


\end{frame}


\subsubsection{10.5.4. Алгоритм построения независимых множеств}
\label{10.5.4.}
\begin{frame}
\frametitle{10.5.4. Алгоритм построения независимых
множеств (1/3) }


%\addvspace{-4pt}

Приведённый \odmkey{алгоритм построения 
независимых множеств}{Алгоритм (algorithm)!построения максимальных независимых множеств (generating of maximal independent sets)},
обоснование которого дано в предыдущем параграфе,
строит все максимальные независимые множества вершин заданного графа.

%\vspace{-3pt}

%\addvspace{-8pt}
%{\scriptsize
%\begin{algorithm}
%\label{selectmaxalg} \caption{Построение максимальных независимых
%множеств} \index{Алгоритм (algorithm)!построения максимальных
%независимых множеств (maximal independent sets)}

%\addvspace{-4pt}%
%\end{algorithm}%

\begin{algorithmic}
%\addvspace{6pt}
\REQUIRE граф $G(V,E)$, заданный списками смежности $\Gamma[v]$.
\ENSURE
последовательность максимальных независимых множеств.
\STATE $k \Assign  0$
\COMMENT{\color{gray}\small  количество элементов в текущем независимом множестве }
\STATE $S[k] \Assign\emptyset$
\COMMENT{\color{gray}\small  $S[k]$~--- независимое множество из $k$ вершин }
\STATE $Q^{-}[k] \Assign \emptyset$
\COMMENT{\color{gray}\small $Q^{-}[k]$ уже использованы для расширения~$S[k]$ }
\STATE $Q^{+}[k] \Assign  V$
\COMMENT{\color{gray}\small  $Q^{+}[k]$ ещё не использованы для расширения~$S[k]$ }
\STATE $M1:$ \COMMENT{\color{gray}\small  шаг вперед }
\STATE \textbf{select $v\in Q^{+}[k]$} \COMMENT{\color{gray}\small  расширяющая вершина }
\STATE $X[k]\Assign v$ \COMMENT{\color{gray}\small  запоминаем расширяющую вершину }
\STATE $S[k+1] \Assign S[k]\cup\{v\}$  \COMMENT{\color{gray}\small  расширенное множество }
\STATE $Q^{-}[k+1] \Assign Q^{-}[k]\setminus \Gamma[v]$
\COMMENT{\color{gray}\small  вершина $v$ использована для расширения }
\STATE $Q^{+}[k+1] \Assign Q^{+}[k]\setminus (\Gamma[v]\cup\{v\})$
\COMMENT{\color{gray}\small  не могут быть использованы }
\STATE $k \Assign k+1$
\end{algorithmic}

\end{frame}



\begin{frame}
\frametitle{10.5.4. Алгоритм построения независимых
множеств (2/3) }
%\vspace{-6pt}

\begin{algorithmic}
\STATE $M2:$ \COMMENT{ проверка }
\FOR{$u\in Q^{-}[k]$}
  \STATE{\textbf{if $\Gamma[u]\cap Q^{+}[k]=\emptyset$ then goto $M3$ end if}} 
  \COMMENT{\color{gray}\small  возврат }
\ENDFOR
\IF{$Q^{+}[k]=\emptyset$}
      \STATE{\textbf{if $Q^{-}[k]=\emptyset$ then yield $S[k]$ end if} }
      \COMMENT{\color{gray}\small   $S[k]$ максимально }
  \STATE \textbf{goto $M3$} \COMMENT{\color{gray}\small  можно возвращаться }
  \ELSE
  \STATE\textbf{goto $M1$} \COMMENT{\color{gray}\small  можно идти вперед }
\ENDIF \STATE $M3:$ \COMMENT{\color{gray}\small  шаг назад }
\STATE $k \Assign k-1$
\STATE $S[k]\Assign S[k+1]-X[k+1]$
\STATE $Q^{-}[k] \Assign Q^{-}[k]+X[k+1]$
\COMMENT{\color{gray}\small  вершина $v$ уже добавлялась }
\STATE $Q^{+}[k] \Assign Q^{+}[k]-X[k+1]$ 
\end{algorithmic}

\end{frame}



\begin{frame}
\frametitle{10.5.4. Алгоритм построения независимых
множеств (3/3) }
%\vspace{-6pt}

\begin{algorithmic}
\IF{$k=0\And
Q^{+}[k]=\emptyset$} \STATE \textbf{stop} \COMMENT{\color{gray}\small  перебор
завершён } \ELSE \STATE \textbf{goto  $M2$}  \COMMENT{\color{gray}\small  переход на
проверку } \ENDIF
\end{algorithmic}

%\end{algorithm}


%\addvspace{-10pt}%
%{\unitlength=1mm \line(1,0){130} }%
%\par%
%\addvspace{6pt}%
%}

%\pagebreak[3]
\begin{example}
Известная \odmkey{задача о восьми ферзях}{Задача (problem)!о
восьми ферзях (eight queens)} (расставить на шахматной доске 8
ферзей так, чтобы они не били друг друга) является задачей об
отыскании максимальных независимых множеств. Действительно,
достаточно представить доску в виде графа с 64 вершинами
(соответствующими клеткам доски), которые смежны, если клетки
находятся на одной вертикали, горизонтали или диагонали.
\end{example}

%\addvspace{6pt}
%{\color{red} Тема для сочинения:
%придумать и реализовать способ наглядной публикации этого алгоритма.
%Например, с помощью диаграммы переходов состояний.}

\end{frame}


\subsection{10.6. Доминирующие множества}
\label{10.6.}

\begin{frame}
\frametitle{10.6. Доминирующие множества}

\odmkey{Задача о наименьшем покрытии}{Задача (problem)!о наименьшем покрытии (of least coverage)} (сокращённо ЗНП) является примером
общей экстремальной задачи, к которой прямо или косвенно сводятся
многие практические задачи. Эта задача является классической,
хорошо изучена и часто используется в качестве теста для сравнения
и оценки различных общих методов решения трудоёмких задач.


В этом разделе ЗНП формулируется в общем виде, а также на основе решения задач, связанных с покрытиями графа, в частности, поиска доминирующих множеств.

\medskip
\begin{tabular}{p{12mm}|p{60mm}|p{70mm}}
  № & Название параграфа
  & Ключевые термины и обозначения \\
  \hline
\DMref{10.6.1.}{10.6.1.} & Минимальные и наименьшие доминирующие множества  &
{Внешне устойчивое множество }\\  
\DMref{10.6.2.}{10.6.2.} & Доминирование и независимость   &
{ Ядро графа}\\  
\DMref{10.6.3.}{10.6.3.} & Задача о наименьшим покрытии &
{  }\\  
\DMref{10.6.4.}{10.6.4.} & Связь задачи о наименьшем покрытии с другими задачами &
{ Задача о выборе переводчиков, задача о развозке }
\\ 
 \end{tabular}  
%\vspace{-0.25\baselineskip}
\end{frame}

\begin{frame}
\label{10.6.1.}\subsubsection{10.6.1. Минимальные и наименьшие доминирующие множества}
\frametitle{10.6.1. Минимальные и наименьшие доминирующие множества (1/2)}
Множество вершин $S\subset V$ графа $G(V,E)$ называется
\odmkey{доминирующим множеством}{Множество (set)!вершин (of
vertex)!доминирующее (dominating)} (или \odmkey{внешне
устойчивым}{Множество (set)!вершин (of vertex)!внешне устойчивое
(externally stable)}), если $S\cup\Gamma(S)=V$, то~есть
$$
\A{v\in V}{v\in S \Or \E{s\in S}{(s,v)\in E}}.
$$

Очевидно, что множество $S$ доминирует тогда и только тогда, когда
$$\A{v \not\in S}{\Gamma(v) \cap S \ne \emptyset},$$ что
равносильно утверждению {\small$\A{v\in V}{\E{s\in S}{d(v,s)\le
1}}$}.

Доминирующее множество называется \odmkey{минимальным}{Множество
(set)!минимальное (minimal)}, если никакое его подмножество не
является доминирующим. Доминирующее множество называется
\odmkey{наименьшим}{Множество (set)!наименьшее (least)}, если
число элементов в нём наименьшее возможное.


\begin{example}
Известная \odmkey{задача о пяти ферзях}{Задача (problem)!о пяти
ферзях (five queens)} (расставить на шахматной доске 5 ферзей так,
чтобы они били всю доску) является задачей об отыскании наименьших
доминирующих множеств.
\end{example}

\end{frame}

\begin{frame}
\frametitle{10.6.1. Минимальные и наименьшие доминирующие множества (2/2)}
Множество рёбер $D\subset E$ графа $G(V,E)$ называется
\odmkey{доминирующим множеством}{Множество (set)!рёбер (of
edges)!доминирующее (dominating)} (или \odmkey{внешне
устойчивым}{Множество (set)!рёбер (of edges)!внешне устойчивое
(externally stable)}), если 
$$
\A{e\in E}{e\in D \Or \E{d = (v_1, v_2)\in D}{e\in (\Lambda(v_1) \cup \Lambda(v_2))}}.
$$

Аналогично доминирующему множеству вершин множество
ребёр $D$ доминирует тогда и только тогда, когда
$$\A{e = (v_1, v_2) \not\in D}{(\Lambda(v_1) \cup \Lambda(v_2)) \cap D \ne \emptyset}.$$

Понятия \odmkey{минимального}{Множество (set)!минимальное (minimal)} и \odmkey{наименьшего}{Множество (set)!наименьшее (least)} доминирующих множеств рёбер вводятся так же, как и для доминирующего множества вершин.
\end{frame}


\begin{frame}
\label{10.6.2.}\subsubsection{10.6.2. Доминирование и независимость }
\frametitle{10.6.2. Доминирование и независимость (1/3)}
Доминирование тесно связано с вершинной независимостью.


%\vspace{-0.25\baselineskip}
\begin{theorem}
Независимое множество вершин является максимальным тогда
%\linebreak
и~только тогда,
когда оно является доминирующим.
\end{theorem}

%}%baselinestretch
%\vspace{-0.25\baselineskip}
\begin{proof}
%%{\renewcommand{\baselinestretch}{0.89}\selectfont
%
\proofsection{$\Impl$} Пусть множество вершин $S\subset V$~---
максимальное независимое. Допустим, что оно не
доминирующее. Тогда существует вершина $v$, находящаяся на
расстоянии большем 1 от всех вершин множества $S$. Эту вершину
можно добавить к $S$ с сохранением независимости, что противоречит
максимальности.
%
%%}%baselinestretch
%
%%{\renewcommand{\baselinestretch}{0.98}\selectfont
%\pagebreak[3]
%
\proofsection{$\Impll$} Пусть $S$~--- независимое доминирующее
множество. Допустим, что оно не максимальное.
Тогда существует вершина $v$, не смежная ни с одной из вершин
множества $S$, то~есть находящаяся на расстоянии больше 1 от всех
вершин множества $S$. Это противоречит тому, что множество $S$~---
доминирующее.
%}%baselinestretch
\end{proof}

%{\renewcommand{\baselinestretch}{0.98}\selectfont

\medskip
Независимое доминирующее множество вершин~---
\odmkey{ядро}{Ядро (kernel)!графа (of graph)} графа.


\begin{example}
В полном графе $K_p$ каждая из $p$ вершин
является ядром и других ядер нет.
\end{example}

\end{frame}

\begin{frame}
\frametitle{10.6.2. Доминирование и независимость (2/3)}
%{\renewcommand{\baselinestretch}{0.98}\selectfont
\begin{corollary}
Любой граф имеет ядро.
\end{corollary}
\begin{proof}
Следующий простой алгоритм,
основанный на доказательстве предыдущей теоремы,
строит некоторое ядро $S$.

\vspace{-4pt}
\begin{algorithmic}
\STATE $S\Assign \emptyset$
\COMMENT{\color{gray}\small  ядро }
\WHILE{$V \ne \emptyset$}
   \STATE \textbf{select $v\in V$}
   \COMMENT{\color{gray}\small  любая нерассмотренная вершина }
   \STATE $S \Assign S + v$
   \COMMENT{\color{gray}\small  расширяем множество $S$ }
   \STATE $V \Assign V\setminus \Gamma^* (v)$
   \COMMENT{\color{gray}\small  удаляем вершины, которые не могут быть использованы для расширения }
\ENDWHILE 
\end{algorithmic}

\vspace{-12pt}
\end{proof}
%\end{proof}
%} \renewcommand{\qed}{\vspace{-\baselineskip}}

%\vspace{-4pt}
Понятия независимости, доминирования и ядра применимы как к
графам, так и к орграфам. Множество узлов $S$ орграфа называется
\odmkey{независимым}{Множество (set)!узлов (of nodes)!независимое
(independent)}, если $\A{{v\in S}}{\Gamma(v)\cap S =
\emptyset}$, и называется \odmkey{доминирующим}{Множество
(set)!узлов (of nodes)!доминирующее (dominating)}, если $\A{v
\not\in S}{\Gamma(v)\cap S \ne \emptyset}$. Независимое
доминирующее множество узлов в орграфе называется
\odmkey{ядром}{Ядро (kernel)!орграфа (of digraph)}.

\begin{remark}
Существуют орграфы, не имеющие ядра. Например, контур $C_3$ не
имеет ядра. Более того, задача выделения ядра в произвольном
орграфе оказывается \emph{NP}-полной.
\end{remark}

%}
\end{frame}

\begin{frame}
\frametitle{10.6.2. Доминирование и независимость (3/3)}

\begin{theorem}
Паросочетание является максимальным тогда и только тогда, когда множество его рёбер является доминирующим.
\end{theorem}

\begin{proof}
\proofsection{$\Impl$} Пусть паросочетание $M\subset E$~---
максимальное. Допустим, что оно не
доминирующее, то есть существует ребро $e\not\in M$, не смежное ни одному ребру из множества $M$. Следовательно при добавлении этого ребра в множество $M$ независимость сохранится, что противоречит максимальности паросочетания.


\proofsection{$\Impll$} Пусть $M$~--- доминирующее
множество рёбер, образующее паросочетание. 

Допустим, что оно не максимальное.
Тогда существует ребро $e\not\in M$, не смежное ни одному из рёбер множества $M$, что противоречит тому, что множество $M$~--- доминирующее.
\end{proof}

\end{frame}


\begin{frame}
\label{10.6.3.}\subsubsection{10.6.3. Задача о наименьшем покрытии }
\frametitle{10.6.3. Задача о наименьшем покрытии (1/2)}

\odmkey{Задача о наименьшем покрытии}{Задача (problem)!о наименьшем покрытии (minimum covering)} является весьма общей задачей, к которой сводятся многие задачи, в том числе и задача отыскания наименьших доминирующих множеств.

Данная задача получила такое название благодаря следующей формулировке на языке теории множеств: 

Пусть имеется конечное множество $V=\{\Tuple{v}{1}{p} \}$ и
семейство подмножеств этого множества 
$\cal{E}=\{E_1,\dots,E_t\}$, $\A{i\in 1..t}{E_i\subset V}$.
Каждому подмножеству $E_i$ приписан вес. Найти покрытие $\cal{E}'$
($\cal{E}'\subset \cal{E}$) наименьшего веса.

\medskip

Эту задачу можно сформулировать также как задачу целочисленного программирования. Пусть 
\textbf{$A:$~array~$[1..p,1..t]$~of~$ 0..1$}, где $A[i,j]\Assign (v_i\in E_j))$ ~--- матрица, отражающая принадлежность элементов подмножествам, а 
\textbf{$c:$~array~$[1..t]$~of integer}~--- вектор весов подмножеств. Тогда нужно найти такой вектор \textbf{$x:$~array~$[1..t]$~of~$ 0..1$},  где $x[i] := (E_i\subset E'),$ что 

$$
\min{(c x | (Ax)_j >= 1)}.
$$


\end{frame}


\begin{frame}
\frametitle{10.6.3. Задача о наименьшем покрытии (2/2)}
\begin{remark}
Условие $(Ax)_j >= 1$ вводится для того, чтобы каждый элемент множества попал хотя бы в одно из выбранных подмножеств~--- тогда подмножества, попавшие в $\cal{E}'$ будут образовывать покрытие. Таким образом обе формулировки оказываются полностью эквивалентными.
\end{remark}

Известно, что ЗНП относится
к числу трудоёмких задач,
и для её решения применяются
переборные алгоритмы с теми или иными улучшениями.


\end{frame}

\begin{frame}
\label{10.6.4.}\subsubsection{10.6.4. Связь задачи о наименьшем покрытии с другими задачами }
\frametitle{10.6.4. Связь задачи о наименьшем покрытии с другими задачами (1/7)}

Рассмотрим следующую задачу. Пусть каждой вершине сопоставлена
некоторая цена. Требуется выбрать доминирующее множество с
наименьшей суммарной ценой. Покажем, что это ЗНП, сформулированная на языке теории графов. Пусть $V=\{\Tuple{v}{1}{p} \}$~--- множество вершин графа, а $E_i =\Gamma^{*}(v_i)$.
Тогда, если $E_i \subset E'$, то вершина $v_i$ принадлежит искомому доминирующему множеству.

\medskip

Чтобы показать эквивалентность ЗНП на языке теории графов и задачи целочисленного программирования, нужно принять матрицу $A$ равной матрице смежности графа с единицами на главной диагонали. Тогда, если $x[i] = 1$, то вершина $v_i$ войдёт в доминирующее множество вершин.

\medskip

\begin{remark}
Если принять веса всех подмножеств тождественно равными единицами, то мы получим так называемую ЗНП с единичными весами, которая также является задачей отыскания обычного доминирующего множества вершин.
\end{remark}
\begin{remark}
Если в качестве $A$ взять матрицу смежности рёбер (т.е. $A[i][j] = 1$ в том случае, если ребро $e_i$ смежно ребру $e_j$), то в результате решения ЗНП мы найдём наименьшее доминирующее множество рёбер.
\end{remark}

\end{frame}

\begin{frame}

\frametitle{10.6.4. Связь задачи о наименьшем покрытии с другими задачами (2/7)}

 \odmkey{Задача о выборе переводчиков}{Задача
(problem)!о выборе переводчиков (translator choice)}. Организации
нужно нанять переводчиков для перевода с определённого множества
языков. Каждый из переводчиков владеет некоторыми
иностранными языками и требует определённую зарплату. Требуется
определить, каких переводчиков следует нанять, чтобы сумма
расходов на зарплату была минимальной. Эта задача сводится к ЗНП следующим образом.Рассмотрим двудольный граф с~долями $V_1$ и $V_2$, где доля
$V_1$ соответствует множеству переводчиков, а доля $V_2$~---
множеству языков. Вершины $v_1\in V_1$ и $v_2 \in V_2$ смежны,
если переводчик $v_1$ владеет языком $v_2$. Требуется найти
наименьшее доминирующее множество такое, что $S\subset V_1$.

\medskip

\odmkey{Задача о развозке}{Задача
(problem)!о развозке (transshipment)}. Поставщику нужно доставить
товары своим потребителям. Имеется множество возможных маршрутов,
каждый из которых позволяет обслужить определённое подмножество
потребителей и требует определённых расходов. Требуется
определить, какие маршруты следует использовать, чтобы минимизировать расходы и обслужить всех потребителей. Задача о развозке сводится к ЗНП аналогичным образом.
Здесь $V_1$~--- множество маршрутов, $V_2$~--- множество
потребителей и вершины $v_1\in V_1$ и $v_2 \in V_2$ смежны, если
маршрут $v_1$ обслуживает потребителя $v_2$.

\end{frame}

\begin{frame}

\frametitle{10.6.4. Связь задачи о наименьшем покрытии с другими задачами (3/7)}
Задача нахождения наибольшего независимого множества вершин так же может быть интерпретирована как частный случай ЗНП. Для этого будем составлять матрицу $A$ таким образом, что её строкам будут соответствовать рёбра графа, а столбцам~--- вершины: 
$$
A[i,j]\Assign (e_i\in \Lambda(v_j)).
$$
Тогда, если $X^*$~--- полученное при решении данной ЗНП множество вершин, то $V \setminus X^*$ будет наибольшим независимым множеством вершин. Это верно благодаря тому, что $X^*$ получается наименьшим вершинным покрытием, то есть каждое ребро графа инцидентно хотя бы одной вершине из $X^*$, а значит никакие две вершины из $V \setminus X^*$ не могут быть смежны. Более того $X^*$ имеет наименьшую мощность, а значит $V \setminus X^*$ имеет наибольшую мощность.

\end{frame}

\begin{frame}
\frametitle{10.6.4. Связь задачи о наименьшем покрытии с другими задачами (4/7)}
\begin{remark}
Легко заметить, что аналогично наименьшему вершинному покрытию 
можно находить наименьшее рёберное покрытие: для этого достаточно сопоставлять строкам матрицы $A$ вершины графа, а столбцам~--- рёбра.
\end{remark}
\medskip

Но неверно считать, что наибольшее паросочетание в графе можно найти аналогично наименьшему доминирующему множеству рёбер~--- просто вычитая из множества всех рёбер полученное наименьшее рёберное покрытие. На приведённом на рисунке графе пунктирными линиями изображено множество рёбер, полученное таким образом, и оно паросочетанием не является.


\begin{center}
\includegraphics[width=40mm]{10_6_4_4.png}
\end{center}


\end{frame}


\begin{frame}
\frametitle{10.6.4. Связь задачи о наименьшем покрытии с другими задачами (5/7)}

Есть ещё одна распространённая экстремальная задача на графах~--- поиск наибольшей клики. На самом деле данная задача эквивалентна задаче поиска наибольшего независимого множества вершин. Действительно: рассмотрим дополнение к графу $\overline{G}(V,\overline{E})$ и наибольшее независимое множество вершин $G$. Заметим, что в графе $\overline{G}$ это же множество вершин будет образовывать полный подграф, при этом мощность этого множества будет наибольшей. На рисунке приведён граф $G$ и его дополнение $\overline{G}$. Закрашенные вершины (наибольшее независимое множество в $G$) образуют в графе $\overline{G}$ клику наибольшего размера.

\begin{center}
\includegraphics[width=70mm]{10_6_4_5.png}
\end{center}

\end{frame}

\begin{frame}
\frametitle{10.6.4. Связь задачи о наименьшем покрытии с другими задачами (6/7)}

В данном разделе рассмотрены способы свести ЗНП к решению различных экстремальных задач. На схеме, приведённой на следующем слайде, показаны связи ЗНП и некоторых других задач на графах. На этой схеме стрелка от задачи $A$ к задаче $B$ означает, что решение задачи
$A$ влечет за собой решение задачи $B$.

\medskip

Можно заметить, что не все связи, изображённые на схеме, были разобраны. Наибольшее паросочетание в графе $G$ можно найти, как наибольшее независимое множество вершин в рёберном графе, соответствующем $G$. Также, задачи наименьшего рёберного покрытия и наибольшего паросочетания оказываются эквивалентными, и, решив одну из этих задач, можно получить решение другой.
\end{frame}


\begin{frame}
\frametitle{10.6.4. Связь задачи о наименьшем покрытии с другими задачами (7/7)}

\begin{figure}[h]
\begin{center}
\includegraphics[width=120mm]{10_6_4_6.png}
\end{center}
\end{figure}
\end{frame}


\begin{frame}
\frametitle{Онтология теории графов}
\begin{center}
\includegraphics[width=15cm]{ODM_10_6.png}
\end{center}

\end{frame}


\subsection{10.7. Раскраска графов}
\label{10.7.}
\begin{frame}
\frametitle{10.7. Раскраска графов}
Задача раскрашивания графов,
которая на первый взгляд кажется
просто праздной головоломкой,
имеет неожиданно широкое
применение в программировании,
особенно при решении фундаментальных
теоретических проблем
(см., например, книгу А.П.~Ершова <<Введение в теоретическое программирование>>).

\medskip
\begin{tabular}{p{12mm}|p{60mm}|p{70mm}}
  № & Название параграфа
  & Ключевые термины и обозначения \\
  \hline
\DMref{10.7.1.}{10.7.1.} & Хроматическое число графа и его оценки  &
{$\chi(G)$~--- хроматическое число; одноцветный класс}\\  
\DMref{10.7.2.}{10.7.2.} & Хроматические числа графа и его дополнения   &
{ }\\  
\DMref{10.7.3.}{10.7.3.} & Точный алгоритм раскрашивания &
{  }\\  
\DMref{10.7.4.}{10.7.4.} & Приближенные алгоритмы раскрашивания &
{ Алгоритм последовательного раскрашивания, улучшенный 
алгоритм последовательного раскрашивания }\\ 
\DMref{10.7.5.}{10.7.5.} & Хроматический полином &
{ }\\ \end{tabular}  
\end{frame}

\begin{frame}
\label{10.7.1.}\subsubsection{10.7.1. Хроматическое число графа и его оценки }
\frametitle{10.7.1. Хроматическое число графа и его оценки (1/3)}


\odmkey{Раскраской}{Раскраска графа (coloring of graph)} графа $G$
называется такое приписывание цветов (натуральных чисел) его
вершинам, что никакие две смежные вершины не получают одинаковый
цвет. 

\begin{remark}
Часто добавляют уточняющее прилагательное и говорят
\em{правильная} или \em{допустимая} раскраска.
Здесь никаких иных раскрасок не рассматривается, 
а потому прилагательное опускается. 
\end{remark}

Если задана допустимая раскраска графа, использующая $m$
цветов, то говорят, что граф $m$-раскраши\-ва\-емый. Наименьшее
возможное количество цветов в раскраске называется
\odmkey{хроматическим чис\-лом}{Число (number)!хроматическое
(chromatic)} и обозначается $\chi(G)$.

\begin{examples}
$\chi(\overline K_p)=1$,
$\chi(K_p)=p$, $\chi(K_{m,n})=2$, $\chi(C_{2n})=2$,
$\chi(C_{2n+1}) =3$, $\chi(T)=2$, где $T$~{\upshape---} свободное дерево.
\end{examples}


Очевидно, что существует $m$-раскраска графа $G$ для любого $m$ в
диапазоне \\$\chi(G)\leq m\leq p$. Множество вершин,
покрашенных в один цвет, называется
\odmkey{одноцветным классом}{Класс
(class)!одноцветный (unicolored)}. Одноцветные классы образуют
независимые множества вершин, то~есть никакие две вершины в
одноцветном классе не смежны.


\medskip
Способ явного выражения хроматического
числа через другие стандартные инварианты графа неизвестен.
Известны только некоторые оценки, часть из которых приведена
ниже.


\end{frame}

\begin{frame}
\frametitle{10.7.1. Хроматическое число графа и его оценки (2/3)}
\begin{theorem}
$\chi (G) \le 1 +  \Delta(G)$.
\end{theorem}

\begin{proof}
Индукция по $p$: \\$p = 1\Impl \chi(G) = 1\And\Delta(G) = 0$.
Пусть
$\A{G}{p(G)<p\Impl\chi(G)\leq\Delta(G)+1}$.
Рассмот\-рим граф $G$,  $p(G) = p$.  Тогда
$\A{v \in  V}{\chi(G - v) \leq \Delta(G - v) + 1 \le \Delta(G) + 1}$.
Но $d(v)\leq\Delta(G)$, значит, хотя бы один цвет в
$(\Delta(G)+1)$-раскраске графа $G-v$ свободен для $v$. Покрасив
вершину $v$ в этот цвет, получаем $(\Delta(G)+1)$-раскраску графа
$G$.
\end{proof}


%\vspace{-0.25\baselineskip}
\begin{theorem}
$p/\beta_0(G)\le\chi(G)\le p-\beta_0(G)+1$.
\end{theorem}

\begin{proof}
\proofsection{$p/\beta_0(G)\le\chi(G)$}
Пусть $\chi (G)=n$ и $V=V_1\cup\dots\cup V_n$, где $V_i$~--- одноцветные
классы. $V_i$~--- независимое множество вершин, следовательно,
$|V_i|\leq\beta_0(G)$. Имеем \\
$p=\sum\limits_{i=1}^n|V_i|\leq n\beta_0(G)\Impl p/\beta_0\leq \chi$.
\proofsection{$\chi(G)\le p-\beta_0(G)+1$} Пусть
$S\subset V$~--- наибольшее независимое множество,
$|S| =\beta_0(G)$. Тогда $\chi(G-S)\leq |V-S|=p-\beta_0(G)$. Из
$n$-раскраски графа $G-S$ можно получить $(n+1)$-раскраску графа
$G$, так~как все вершины из $S$ можно покрасить в один новый цвет.
Следовательно, $\chi(G)= \chi (G-S)+1\leq p-\beta_0+1$.
\end{proof}
\end{frame}

\begin{frame}
\frametitle{10.7.1. Хроматическое число графа и его оценки (3/3)}

Нахождение хроматического числа 
является нетривиальной и практически востребованной задачей.

\medskip

\begin{example}
Предположим, что у каждой студенческой группы свой список экзаменов, которые необходимо сдать
в сесиию. Если группы соответствуют разным направлениям подготовки,
то списки сдаваемых предметов не совпадают,
хотя, возможно, пересекаются. Нужно так составить расписание сессии, чтобы экзамены для каждой группы были назначены в разные дни, а число экзаменационных дней было бы минимальным. 

Для решения задачи построим такой граф $G$: множество вершин совпадает с множеством предметов, а вершины соединены ребром, если хотя бы одна группа сдаёт оба предмета. Тогда $\chi(G)$ определит минимальное количество экзаменационных дней.
\end{example}

\end{frame}

\begin{frame}
\label{10.7.2.}\subsubsection{10.7.2. Хроматические числа графа и его дополнения }\frametitle{10.7.2. Хроматические числа графа и его дополнения (1/2)}

Хроматические числа графа $G$ и его дополнения $\overline{G}$
связаны: если в $G$ сравнительно мало рёбер, то и $\chi(G)$ будет
невелико, но тогда в $\overline{G}$~--- много рёбер и
$\chi(\overline{G})$ будет близко к $p$.

\begin{theorem}
Пусть $\chi \Assign \chi(G)$, $\overline{\chi} \Assign \chi(\overline{G})$.
Тогда
$$
2\sqrt{p} \le  \chi+\overline{\chi} \le  p+1 , \quad p \le
\chi\overline{\chi}  \le \left(\frac{p+1}{2}\right)^2.
$$
\end{theorem}



\begin{proof} 
\proofsection{$p \le  \chi\overline{\chi}$} Пусть $\chi(G)=n$,
$\Tuple{V}{1}{n}$~--- одноцветные классы, $p_i\Assign |V_i|$.
Тогда $\sum\limits_{i=1}^np_i=p$, следовательно,
$\max\limits_{i=1}^np_i\geq p/n$. Но $V_i$~--- независимые
множества в $G$, следовательно, $V_i$~--- клики в $\overline{G}$.
Значит, $\overline{\chi}\geq\max\limits^n_{i=1}p_i\geq p/n$. Имеем
$\chi\overline{\chi}\geq n\cdot p/n\nolinebreak=\nolinebreak p$.


\proofsection{$2\sqrt{p} \le  \chi+\overline{\chi}$}
Известно, что среднее геометрическое двух чисел
не превосходит среднего арифметического:
$(a+b)/2\ge\sqrt{ab}$.
Следовательно,
$\chi+\overline{\chi}\ge 2\sqrt{\chi\overline{\chi}}\ge 2\sqrt{p}$.
\end{proof}
\end{frame}

\begin{frame}
\frametitle{10.7.2. Хроматические числа графа и его дополнения (2/2)}

\begin{proof} (Продолжение)
\proofsection{$\chi+\overline{\chi} \le  p+1 $} Индукция по $p$.
База: $p=1\Impl \chi=1\And\overline{\chi}=1$. Пусть
$\chi+\nobreak\overline{\chi}\le\nobreak p$ для всех графов с
$p-1$ вершинами. Рассмотрим граф $G$ с $p$ вершинами и~вершину %\linebreak
$v\in V$. Тогда, очевидно, $\chi(G)\le \chi(G-v)+1$ и
$\chi(\overline{G})\le \chi(\overline{G}-v)+1$. Если
$\chi(G)<\chi(G-v)+1$ или
$\chi(\overline{G})<\chi(\overline{G}-v)+1$, то
$\chi+\overline{\chi}=\chi(G)+\chi(\overline{G})
<\chi(G-v)+1+\chi(\overline{G}-v)+1< p+2\le p+1.$ Следовательно,
$\chi+\overline{\chi}\!\le\! p+1$. Пусть теперь
$\chi(G)\!=\!\chi(G-v) +1$ и
$\chi(\overline{G})\!=\!\chi(\overline{G}-v) +1$. Положим
$d\Assign d(v)$ в графе $G$, тогда $\overline{d}=p-d-1$~---
степень вершины $v$ в графе $\overline{G}$. Имеем
$d\ge\chi(G-v)$. Действительно, $\chi(G)=\chi(G-v)+1$, и если бы
$d<\chi(G-v)$, то вершину $v$ можно было бы покрасить в любой из
свободных $\chi(G-v)-d$ цветов и~получить $\chi(G-v)$-раскраску
графа $G$. Аналогично,
$\overline{d}=p-d-1\ge\chi(\overline{G}-v)$. Таким образом,

\vspace{-15pt}
{\small
$$\chi+\overline{\chi}=\chi(G)+\chi(\overline{G})=
\chi(G-v)+1+\chi(\overline{G}-v)+1\le d+1+p-d-1+1=p+1.$$}

\vspace{-15pt}
\proofsection{$\chi\overline{\chi} \le \left((p+1)/2\right)^2$\!} Имеем
$2\sqrt{\chi\overline{\chi}}\!\le\! \chi+\overline{\chi}\!\le\!
p+1$. 
Следовательно, $\left((p+1)/2\right)^2
\!\ge\! \chi\overline{\chi}$.
\end{proof}


\end{frame}

\begin{frame}
\label{10.7.3.}\subsubsection{10.7.3. Точный алгоритм раскрашивания }\frametitle{10.7.3. Точный алгоритм раскрашивания (1/2)}

%{\renewcommand{\baselinestretch}{0.95}\selectfont
Поскольку точная формула для хроматического числа неизвестна, задача
нахождения наилучшей раскраски графа оказывается, как и следовало
ожидать, труднорешае\-мой.

Рассмотрим следующую схему рекурсивной процедуры $P$:
\begin{enumerate}
\item
Выбрать в графе $G$  некоторое
 максимальное независимое множество вершин~$S$.

\item Покрасить вершины множества $S$ в очередной цвет.

\item Применить процедуру $P$  к графу $G-S$.
\end{enumerate}



На псевдокоде процедура $P$ может быть записана так.


\begin{algorithmic}
\REQUIRE граф $G(V,E)$, номер свободного цвета $i$.
\ENSURE раскраска, заданная массивом $C[V]$,~--- номера цветов, приписанные вершинам.
\STATE \textbf{if $V=\emptyset$ then return end if} 
\COMMENT{\color{gray}\small  раскраска закончена }
\STATE $S\Assign \Selectmax(G)$
\COMMENT{\color{gray}\small $S$~--- максимальное независимое множество }
\STATE $C[S]\Assign i$
\COMMENT{\color{gray}\small  раскрашиваем вершины множества $S$ в цвет $i$ }
\STATE $P(G-S,i+1)$
\COMMENT{\color{gray}\small  рекурсивный вызов }
\end{algorithmic}


\begin{remark}
Функция $\Selectmax$ может быть реализована, например, на основе алгоритма
10.5.4.
\end{remark}
%}

\end{frame}

\begin{frame}
\frametitle{10.7.3. Точный алгоритм раскрашивания (2/2)}
%\pagebreak[3]
%{\renewcommand{\baselinestretch}{0.92}\selectfont
%\vspace{-2pt}
\begin{theorem}
Если граф $G$~{\upshape---} $k$-раскрашиваемый, то существует
такая последовательность выборов множества $S$ на шаге 1 процедуры
$P$, что применение процедуры $P$ к графу $G$ построит не более
чем $k$-раскраску графа $G$.
\end{theorem}

%\vspace{-2pt}
\begin{proof}
Пусть имеется некоторая $k$-раскраска графа $G(V,E)$. Перестроим
её в не более чем $k$-раскраску, которая может быть получена
процедурой $P$. Пусть $V_1\subset V$~--- множество вершин в данной
$k$-раскраске, покрашенных в~цвет $1$. Множество $V_1$~---
независимое, но, может быть, не максимальное. Рассмотрим
множество ${V_1}'$,~такое, что \\$V_1\cup {V_1}'$~--- максимальное
независимое множество (может оказаться, что ${V_1}'=\emptyset)$.
Вершины из ${V_1}'$ не смежны с $V_1$, значит, вершины из ${V_1}'$
можно перекрасить в цвет $1$. Пусть далее $V_2 \subset 
V \setminus (V_1 \cup  {V_1}')$~--- множество вершин, покрашенных в
цвет 2. Аналогично рассмотрим множество ${V_2}'$, такое, что
$V_2\cup {V_2}'$~--- максимальное независимое в
$G\backslash(V_1\cup {V_1}')$, покрасим вершины $V_2\cup {V_2}'$ в
цвет $2$ и т.~д. Всего в исходной раскраске было $k$ независимых
множеств. При перекраске их число не возрастёт (но может
уменьшиться, если $\chi(G)<k$). На каждом шаге алгоритма
рассматривается одно из непустых множеств, следовательно, процесс
закончится.
\end{proof}
%}
\end{frame}

\begin{frame}
\label{10.7.4.}\subsubsection{10.7.4. Приближённые алгоритмы  раскрашивания }\frametitle{10.7.4. Приближённые алгоритмы  раскрашивания (1/5)}

В предыдущем параграфе некоторый алгоритм точного раскрашивания
был построен на основе алгоритма выделения максимальных
независимых множеств вершин, который имеет переборный характер.
Таким образом, предложенный алгоритм точного раскрашивания также
имеет переборный характер. Можно показать, что это не случайно и
задача построения минимальной раскраски является $NP$-полной. 
%
При
практическом решении $NP$-полных задач целесообразно рассматривать
\odmkey{приближённые}{Алгоритм (algorithm)!приближённый
(approximate)} алгоритмы, которые не всегда находят точное решение
задачи (иногда они находят только приближение к нему, и мы не
можем знать этого заранее), но зато достаточно~эффективны.
Здесь понятие приближённого алгоритма 
имеет более широкий смысл по сравнению с тем, который
обычно подразумевается в вычислительной математике
при проведении численных расчётов.
В вычислительной математике приближённый алгоритм, 
например, численное решение уравнения,
находит приближение к корню с заданной точностью: $x \pm \delta$.
Увеличивая число итераций, можно уменьшать $\delta$.

\end{frame}


\begin{frame}
\frametitle{10.7.4. Приближённые алгоритмы  раскрашивания (2/5)}
Приближённые алгоритмы решения переборных задач 
находят решение с той точностью, 
с какой умеют~--- дополнительные вычисления не улучшают решения.
Поэтому приближённый алгоритмы решения переборных задач
характеризуются двумя параметрами:
насколько далеко найденное решение может отстоять от точного решения,
и насколько меньшие трудозатраты по сравнению с полным перебором
при этом требуются.  
  


\smallskip
Рассмотрим следующий \odmkey{алгоритм последовательного
раскрашивания}{Алгоритм (algorithm)!последовательного
раскрашивания (sequential coloring)}.

Идея состоит в том, чтобы красить вершины последовательно,
выбирая среди допустимых 
цвет с наименьшим номером.

\begin{ground}
В основном цикле рассматриваются все вершины,
и каждая из них получает~допус\-тимую раскраску.
Таким образом, процедура строит допустимую раскраску.
\end{ground}
\end{frame}


\begin{frame}
\frametitle{10.7.4. Приближённые алгоритмы  раскрашивания (3/5)}
\begin{algorithmic}
\REQUIRE граф $G$.
\ENSURE раскраска графа~--- массив \textbf{$C:$ array $[1..p]$ of $1..p$}.
\FOR{$v \in V$}
\STATE $C[v] \Assign  0 $
\COMMENT{\color{gray}\small  все вершины не раскрашены }
\ENDFOR
\FOR{$v \in V$}
\STATE $A \Assign  \{1,\dots,p\}$
\COMMENT{\color{gray}\small  все цвета }
\FOR{$u \in\Gamma(v)$}
\STATE $A \Assign  A\setminus\{C[u]\}$
\COMMENT{\color{gray}\small  занятые для вершины $v$ цвета }
\ENDFOR
\STATE $C[v] \Assign  \min A$
\COMMENT{\color{gray}\small  свободный цвет с наименьшим номером }
\ENDFOR
\end{algorithmic}

\odmkey{Улучшенный алгоритм последовательного раскрашивания}{Алгоритм
(algorithm)!последовательного раскрашивания (улучшенный) (improved
sequential coloring)} строит допустимую раскраску, применяя
такую эв\-рис\-тику: начинать раскрашивать следует с вершин
наибольшей степени, поскольку если их раскрашивать в конце
процесса, то более вероятно, что для них не найдётся свободного
цвета и придётся использовать ещё один.


\end{frame}


\begin{frame}
\frametitle{10.7.4. Приближённые алгоритмы  раскрашивания (4/5)}



\begin{algorithmic}
\REQUIRE граф $G$.
\ENSURE раскраска графа~--- массив \textbf{$C:$ array $[1..p]$ of $1..p$}.
\STATE $\Sort(V)$ \COMMENT{\color{gray}\small  упорядочить вершины по невозрастанию степени }
\STATE $c \Assign  1$ \COMMENT{\color{gray}\small  первый цвет }
\STATE\textbf{for $v \in V$ do $C[v] \Assign  0$ end for} \COMMENT{\color{gray}\small  все не раскрашены }
\WHILE{$V \not=\emptyset$}
\FOR{$v \in V$}
\FOR{$u \in \Gamma^+(v)$}
\IF{$C[u] = c$}
\STATE \textbf{next for $v$}
\COMMENT{\color{gray}\small  вершину $v$ нельзя покрасить в цвет $c$ }
\ENDIF
\ENDFOR
\STATE {$C[v] \Assign  c$;
$V\Assign V\setminus \{v\}$} \COMMENT{\color{gray}\small красим вершину $v$ в цвет $c$  и удаляем её из рассмотрения }
\ENDFOR
\STATE$c \Assign c+1$ \COMMENT{\color{gray}\small  следующий цвет }
\ENDWHILE
\end{algorithmic}

\end{frame}


\begin{frame}
\frametitle{10.7.4. Приближённые алгоритмы  раскрашивания (5/5)}

\begin{ground}
Заметим, что данный алгоритм отличается от предыдущего тем, что
основной цикл идет не по вершинам, а по цветам: сначала всё, что
можно, красим в цвет~1, затем в оставшемся красим всё, что можно,
в цвет 2 и т.~д. В~остальном алгоритмы аналогичны, и данный
алгоритм заканчивает свою работу построением допустимой раскраски
по тем же причинам, что и предыдущий.
\end{ground}


\medskip
%\vspace{-0.5\baselineskip}
\begin{digress}
Улучшенный алгоритм последовательного раскрашивания
несколько сложнее первого алгоритма
и основан на остроумной эвристике. Можно было бы ожидать,
что он даст существенно лучшие результаты.
Однако прямые вычислительные эксперименты показывают,
что эти алгоритмы почти во всех случаях
работают одинаково хорошо (или плохо).
Таким образом, программистские <<хитрости>> далеко не всегда
дают практически значимые результаты.
\end{digress}

\end{frame}

\subsubsection{10.7.5. Хроматический полином }
\label{10.7.5.}
\begin{frame}
\frametitle{10.7.5. Хроматический полином (1/6)}
Раскраски (допустимые) нетривиальных графов не являются единственными.

\smallskip
Обозначим $P(G,n)$~--- число раскрасок графа $G$ в 
$n$ цветов. Положим $\A{G}{P(G,0)\Assign 0}$.
Ясно, что $\A{n\in 0..(\chi(G)-1)}{P(G,n) =0}$
и $\A{n\in \chi(G)..p}{P(G,n) >0}$.
Для $n>p$ число $P(G,n)$ не имеет 
естественного смысла, а потому обычно не рассматривается.     

\begin{examples}
\begin{enumerate}
\item $P(K_1, 1)=1$. Действительно,
единственную изолированную вершину можно
раскрасить единственным образом.
\item $P(K_2, 1)=0$, $P(K_2, 2)=2$.
Действительно,
пару смежных вершин можно
раскрасить в два цвета двумя способами.
\item  $P(K_3, 1)=0$, $P(K_3, 2)=0$, $P(K_3, 3)=6$.
Действительно,
первую вершину можно покрасить в любой из трёх цветов,
вторую вершину можно покрасить в один из двух
оставшихся цветов, а цвет третьей вершины определяется однозначно.
\end{enumerate}
\end{examples}

\begin{lemma}
$\A{G}{\max\Set{n\in 0..p}{P(G,n)}= P(G,p)= p!}$.
\end{lemma}

\begin{proof}
По принципу Дирихле при раскраске
$p$ вершин в $p$ цветов каждая вершина получает свой цвет, а число раскрасок равно
числу перестановок. 
При уменьшении числа цветов число 
раскрасок также уменьшается.  
\end{proof}


\end{frame}

\begin{frame}
\frametitle{10.7.5. Хроматический полином (2/6)}
\begin{example}
$P(K_p,p)=p!$ и $\A{i<p}{P(K_p,i)=0}$.
\end{example}

\begin{lemma}
$P(G_1\cup G_2, n)=P(G_1, n)P(G_2, n)$.
\end{lemma}
\begin{proof}
Раскраски компонент связности независимы,
и значит, применимо правило произведения
\DMref{п.~5.1.1}{5.1.1.}. 
\end{proof}

\begin{digress}
Здесь подразумеваются нумерованные графы, вершины которых отличимы. Если определить допустимую раскраску \em{с точностью до изоморфизма}, то подсчёт числа допустимых раскрасок существенно усложняется.
Представим себе, что диаграмма графа $K_2$ нарисована
на листе бумаги и левая вершина покрашена в красный цвет, а правая~--- в синий. Тогда следует ли считать, что это уже другая раскраска, если повернуть лист на 180 градусов? Ведь физически рисунок диаграммы не перекрашен!
Допустим, что ответ на предыдущий вопрос отрицательный и 
согласуется с <<физическим здравым смыслом>>.
При таком подходе треугольник $K_3$ 
можно раскрасить всего двумя способами, которые
устойчивы по отношению к поворотам изображения диаграммы.
Однако диаграмму треугольника можно нарисовать
на прозрачном стекле, и тогда оказывается,
что существует всего одна правильная раскраска,
и всё зависит от того, с какой стороны смотреть
на диаграмму.
Учёт подобных тонкостей не входит в
базовые концепции языка теории графов, 
и относится, скорее, к \em{вычислительной геометрии} .
\end{digress}

\end{frame}


\begin{frame}
\frametitle{10.7.5. Хроматический полином (3/6)}

Очевидно, что 
$\chi(G)=\min\Set{n\in 0..p}{P(G,n)>0}$,
поэтому, зная числа $P(G,n)$, нетрудно
найти хроматическое число $\chi(G)$.
В обратную сторону это рассуждение
не работает: по хроматическому числу
в общем случае невозможно определить даже
число минимальных раскрасок. 

\smallskip
Возникает вопрос: как найти компактное описание
множества $\Set{n\in 0..p}{P(G,n)}$?\\
Далее в этом параграфе рассматривается один
из популярных ответов на этот вопрос.

\smallskip
Пусть дан $(p,q)$-граф $G$. 
Полином $P_G(x)$ такой, что 
$\A{n\in 0..p}{P_G(n) =P(G,n)}$
называется
\odmkey{хроматическим полиномом}{Полином (polynomial)!хроматический (chromatic)} графа $G$.

\smallskip
Из математического анализа известно, 
что существует единственный 
\em{интерполяционный полином}  
степени не выше $k$, проходящий 
через множество точек
$\left\lbrace \left( x_i, y_i \right) \right\rbrace_{i=0}^k$, если $x_i$ различны.
Для хроматического полинома $P_G(x)$ имеем множество точек
  с целочисленными координатами 
$\left\lbrace \left( n, P(G,n) \right) \right\rbrace_{n=0}^p$, 
а значит для любого графа $G$ существует единственный 
хроматический полином $P_G(x)$.
В общем случае метод построения 
хроматического полинома для произвольного
графа неизвестен, однако в некоторых 
случаях хроматический полином удаётся построить. 

\end{frame}

\begin{frame}
\frametitle{10.7.5. Хроматический полином (4/6)}
\begin{example}
Для полного графа $K_p$ имеем
$\A{n\in 0..(p\!-\!1)}{P(K_p,n)=0}$ и \\$P(K_p,p)=p!$,
то~есть искомый полином $P_{K_p}(x)$ имеет корни $0,1,\dots,p-1$.
Из математического анализа известно,
что в таком случае 
$P_{K_p}(x)=x(x-1)\ldots(x-(p-1))$.
В частности, $P_{K_3}(x)=x(x-1)(x-2)$ и
$P_{K_3}(3)=3(3-1)(3-2) = 6$,
что согласуется с предыдущими наблюдениями. 
\end{example}

\begin{remark}
Поскольку значение хроматического полинома
требуется вычислять только при значениях
$n\in 0..p$, иногда записывают полином 
в следующей форме: 
$$P_{K_p}(x)=x(x-1)\ldots(x-(p-1))= \frac{x!}{(x-p)!}.$$
Следует быть осторожными при использовании
подобных вольностей в обозначениях: эта 
запись имеет смысл и даёт правильный результат только при значении $x=p$.
\end{remark}
\smallskip
Удаление рёбер в графе разве что увеличивает 
число раскрасок.
\begin{example}
Рассмотрим пустой граф $\overline{K_p}$.
Любую вершину можно покрасить в любой цвет,
поэтому число раскрасок~--- это число
размещений (\DMref{п.~5.1.1}{5.1.1.}) 
и $P(\overline{K_p}, n)=n^p$.
Значит $P_{\overline{K_p}}(x)= x^p 
 > x(x-1)\ldots(x-(p-1)) =  P_{K_p}(x)$.
\end{example}
\end{frame}

\begin{frame}
\frametitle{10.7.5. Хроматический полином (5/6)}
\begin{theorem} {\NB{[редукции]}}
Пусть $e=(u,v)$~--- любое ребро в графе $G$.\\
Тогда $P_{G}(x)= P_{G-e}(x)-P_{G/e}(x)$.
\end{theorem}

\begin{proof}
Все раскраски графа $G-e$ в $n$ цветов можно разбить на два блока:
в один блок входят такие раскраски, в которых 
вершины $u$ и $v$ окрашены в один цвет, 
в другой блок входят раскраски, в которых эти 
вершины покрашены в разные цвета.
Тогда в первом блоке $P(G/e, n)$ раскрасок,
а во втором блоке $P(G, n)$ раскрасок.
Это справедливо для любого $n\in 0..p$,
а значит
$P_{G-e}(x)= P_{G}(x)+P_{G/e}(x)$.
\end{proof}

\smallskip
\begin{example} Рассмотрим граф $K_2$.
Если редуцировать ребро, то
по теореме получаем \\
$P_{K_2}(x) = 
P_{K_2-e}(x)-P_{K_2/e}(x)$. 
Но $P_{K_2-e}(x)=P_{K_1\cup K_1}(x) =x^2$ и 
$P_{K_2/e}(x)= P_{K_1}(x)= x$.
Поэтому $P_{K_2}(x)=x^2-x=x(x-1)$.
\end{example}

\smallskip
\begin{corollary}
Если ребро $e=(u,v)$ не входит в граф $G$, то
 $P_{G}(x)= P_{G+e}(x)+P_{G/e}(x)$.
\end{corollary}

\begin{proof}
Пусть $G_1 = G + e$. Заметим, что $G_1/e = G/e$, 
так как при стягивании ребра $e$ оно исчезает. По теореме редукции 
$ P_{G_1}(x) = P_{G_1- e}(x)-P_{G_1/e}(x)$.
Тогда $ P_{G+e}(x)= P_G(x)- P_{G/e}(x)$. 
\end{proof}
Имея запас хроматических полиномов для простейших
графов, построим 
хроматические полиномы для более сложных графов по индукции, используя теорему редукции. 
\end{frame}

\begin{frame}
\frametitle{10.7.5. Хроматический полином (6/6)}

\begin{examples}
{\color{blue}\bf 1.} Для простой цепи $P_n$ из $n$ вершин
имеем по теореме редукции крайнего ребра
$P_{P_n}(x)= P_{P_n-e}(x)-P_{P_n/e} (x) =
P_{P_{n-1}\cup K_1}(x)-P_{P_{n-1}} (x)
= xP_{P_{n-1}}{x} - P_{P_{n-1}} (x) =$\\
$=(x-1)P_{P_{n-1}} (x)$.
Учитывая, что $P_{P_2}(x)=x(x-1)$,
получаем что
$P_{P_n}(x)= x(x-1)^{n-1}$.

{\color{blue}\bf 2.}
Для простого цикла $C_n$ из $n$ вершин 
$P_{C_n}(x)=(x-1)^n + (-1)^n(x-1)$.
Доказательство по индукции.
$P_{C_3}(x)=P_{K_3}(x)=x(x-1)(x-2) = (x-1)^3
+(-1)^3(x-1)$. 
По теореме редукции
$P_{C_{n+1}}(x) = P_{C_{n+1}-e}(x)-P_{C_{n+1}/e}(x)
= P_{P_{n+1}}(x)-P_{C_{n}}(x) = $\\
$=x(x-1)^n -(x-1)^n -(-1)^n(x-1) =
(x-1)^{n+1} +(-1)^{n+1}(x-1)$.

{\color{blue}\bf 3.} 
Для свободного дерева $T_n$ c $n$ вершинами
 $P_{T_n}(x) = x(x\!-\!1)^{n-1}$. 
Доказательство индукцией 
по числу рёбер. База $K_2=P_2=T_2$ проверена.
При удалении ребра в дереве оно разбивается на два дерева, а при стягивании ребра остаётся деревом с меньшим числом вершин. Имеем по теореме редукции 
 $P_{T_n}(x) = P_{T_k\cup T_{n-k}}(x) - P_{T_{n-1}}(x)=$\\ $= x(x\!-\!1)^{k-1}\cdot x(x\!-\!1)^{n-k-1} - x(x\!-\!1)^{n-2} = x^2(x\!-\!1)^{n-1}- x(x\!-\!1)^{n-2}  = x(x\!-\!1)^{n-1}$. 
\end{examples}


\begin{digress}
Помимо вершинной раскраски существует рёберная, т.е. при которой никакие два смежных ребра не имеют один и тот же цвет. Реберное хроматическое число графа $G$ равно вершинному хроматическому числу его реберного графа $I(G)$.
\end{digress}
\end{frame}

\subsection{10.8. Планарность}
\label{10.8.}


\begin{frame}
\frametitle{10.8. Планарность}
Обсуждение планарности в этом разделе
позволяет решить вторую историческую задачу из
перечисленных в \DMref{п.~7.1.1}{7.1.1.},
а также подготавливает результаты,
необходимые для доказательства теоремы о пяти красках.

\medskip
\begin{tabular}{p{12mm}|p{60mm}|p{70mm}}
  № & Название параграфа
  & Ключевые термины и обозначения \\
  \hline
\DMref{10.8.1.}{10.8.1.} & Укладка графов  &
{Планарный граф, плоский граф, $f(G)$~--- число граней графа}\\  
\DMref{10.8.2.}{10.8.2.} & Эйлерова характеристика   &
{ Операция подразбиения ребра,  гомеоморфизм
графов, теорема Куратовского }\\  
\DMref{10.8.3.}{10.8.3.} & Теорема о пяти красках &
{  }\\  
\end{tabular}  
\end{frame}

\begin{frame}
\label{10.8.1.}\subsubsection{10.8.1. Укладка графов }\frametitle{10.8.1. Укладка графов (1/2)}

Граф \odmkey{укладывается}{Укладка графа (graph flattering)} на
некоторой поверхности, если его можно нарисовать на этой
поверхности так, чтобы рёбра графа при этом не пересекались. Граф
называется \odmkey{планарным}{Граф (graph)!планарный (planar)},
если его можно уложить на плоскости. \odmkey{Плоский}{Граф
(graph)!плоский (flat)} граф~--- это граф, уже уложенный на
плоскости.


Область, ограниченная ребрами в плоском графе, называется
\odmkey{гранью}{Грань (face)!графа (of graph)}. Грань не содержит
других граней. Число граней плоского графа $G$
обозначается~$f(G)$.


%\vspace{-0.5\baselineskip}
\begin{remark}
Внешняя часть плоскости также образует грань.
\end{remark}


%\vspace{-0.5\baselineskip}
\begin{example}
На рисунке показаны диаграммы планарного графа $K_4$
и две его укладки на плоскости. Этот граф имеет 4 грани.
\end{example}


%\vspace{-0.5\baselineskip}
\begin{figure}[h]
\begin{center}
\includegraphics[height=30mm]{10_8_new.jpg}
\end{center}
%\caption{Планарный граф и его укладка}
%\label{f1201}
\end{figure}
\end{frame}

\begin{frame}
\frametitle{10.8.1. Укладка графов (2/2)}
%\pagebreak[3]
\begin{digress}
Точное определение некоторых понятий, используемых в этом разделе,
в частности, таких понятий, как <<поверхность>>, <<область>>,
<<граница>>, выходит далеко за рамки этого учебника и стандартного
курса дискретной математики. Мы полагаемся на геометрическую
интуицию читателя и не даём никаких определений. Для понимания
простейших рассуждений этого раздела достаточно неформальных
интуитивных представлений. Однако читателю следует иметь в виду,
что при решении сложных практических задач, например из
вычислительной геометрии, интуитивных представлений может
оказаться недостаточно.
\end{digress}
%
%%}%baselinestretch

%{\color{red} Роскошная тема для сочинения:
%ещё про планарность, например,
%математически грамотное отступление про эйлеровы характеристики 
%других поверхностей, в частности, тора. 
%Алгоритмы укладки также приветствуются.
%Новый материал должен быть безупречно встроен в имеющийся.
%}

\end{frame}

\begin{frame}
\label{10.8.2.}\subsubsection{10.8.2. Эйлерова характеристика }\frametitle{10.8.2. Эйлерова характеристика (1/3)}
%\label{s10-8-2}


Для графов, уложенных на некоторой
поверхности, справедливо определённое
соотношение между числом вершин, рёбер и граней
графов, которые укладываются на этой поверхности.



%\pagebreak[3]
\begin{theorem} \NB{ [формула Эйлера] }
Для связного планарного графа справедливо следующее
соотношение:\odmkey{}{Формула (formula)!Эйлера (Leonard Euler)}
$
p - q + f = 2.
$
\end{theorem}

\begin{remark}
Число в правой части этого соотношения называется
\odmkey{эйлеровой характеристикой}{Эйлерова характеристика \\(Euler
characteristic)} поверхности.
\end{remark}

\vspace{3pt}
\begin{proof}
Индукция по $q$. База: $q=0\Impl p=1\And f=1$. Пусть теорема верна для
всех графов с $q$ рёбрами: $p-q+f=2$. Добавим еще одно ребро. Если
добавляемое ребро соединяет существующие вершины, то
легко видеть, что $q'=q+1$, $p'=p$,
$f'=f+1$ и  %\linebreak
$p'-q'+f'=p-q-1+f+1=p-q+f=2$. Если добавляемое ребро соединяет
существующую вершину с новой, то $p'=p+1$, $q'=q+1$, $f'=f$ и\\
${p'-q'+f'} = {p+1-q-1+f} ={p-q+f=2}$.
\end{proof}



\end{frame}

\begin{frame}
\frametitle{10.8.2. Эйлерова характеристика (2/3)}
%\setcounter{varcorollary}{0}
\begin{corollary} \NB{[1]}
Если $G$~{\upshape---} связный планарный граф $(p>3)$, то $q \le
3p-6$.
\end{corollary}


\begin{proof}
Каждая грань ограничена по крайней мере тремя ребрами, каждое ребро
ограничивает не более двух граней, отсюда $3f\le 2q$. Имеем\\
$2=p-q+f\le
p-q+\frac{2q}{3} \Impl  3p-3q+2q\ge 6 \Impl  q\le 3p-6$.
\end{proof}

\begin{corollary} \NB{[2]} Графы
$K_5$ и $K_{3,3}$ не планарны.
\end{corollary}


\begin{proof}
\proofsection{\!$K_5$ не планарен\!} Имеем $p=5$, $q=p(p-1)/2=10$. Если
$K_5$ планарен, то по предыдущему следствию $q=10\le
3p-6= 3\cdot 5-6=9$~--- противоречие.\pagebreak[2]

\proofsection{\!$K_{3,3\!}$ не планарен} Имеем $p=6$, $q=9$. В
этом графе нет треугольников, значит, если этот граф планарен, то
в его плоской укладке каждая грань ограничена не менее чем
четырьмя ребрами и, следовательно, $4f\le 2q$ или $2f\le q$. По
формуле Эйлера $6 - 9 + f = 2$, откуда $f = 5$. Имеем
$2f = 2\cdot 5 = 10 \le  q = 9$~--- противоречие.
\end{proof}
\end{frame}

\begin{frame}
\frametitle{10.8.2. Эйлерова характеристика (3/3)}%

%\vspace{-3pt}
\begin{remark}
Операция \odmkey{подразбиения}{Операция (operation)!подразбиения (subpartition)}
ребра $x=(u,v)$ состоит в замене его двумя рёбрами, $(u,w)$ и~$(w,v)$, где $w$~--- новая вершина. Два графа называются
\odmkey{гомеоморфными}{Граф (graph)!гомеоморфный (homeomorphic)},
если они могут быть получены из одного графа подразбиением рёбер.
Граф планарен тогда и только тогда, когда он не содержит
подграфов, гомеоморфных $K_5$ или $K_{3,3}$. Доказательство
достаточности этого утверждения (\odmkey{теоремы
Куратовского}{Теорема (theorem)!Куратовского (Kazimierz
Kuratowski)}) выходит за рамки данного курса.
\end{remark}

%\pagebreak[3]
%
%{\renewcommand{\baselinestretch}{0.98}\selectfont
%
\begin{corollary} \NB{[3]}
В любом планарном графе существует вершина, степень
которой не больше 5.
\end{corollary}



\begin{proof}
От противного. Пусть $\A{v\in V}{d(v)\ge 6}$. Тогда
\vspace{-2pt}
$$
6p\le
\sum\limits_{v\in V}d(v)=2q\Impl 3p\le q,
$$
но $q\le 3p-6$. Имеем $3p\le 3p-6$~--- противоречие.
\end{proof}



\begin{digress}
Для случая многогранников формулу Эйлера вывел Декарт. Действительно,
каркас многогранника~--- это плоский граф, и обратно,
связному плоскому графу (при $f >3$) соответствует многогранник.
\end{digress}

\end{frame}


\begin{frame}
\label{10.8.3.}\subsubsection{10.8.3. Теорема о пяти красках }\frametitle{10.8.3. Теорема о пяти красках (1/2)}

\begin{theorem}
Всякий планарный граф можно раскрасить пятью красками.
\end{theorem}



%\vspace{-12pt}
\begin{proof}
Достаточно рассматривать связные графы, потому что
компоненты  связности  можно  раскрашивать  независимо.
$$
\chi\left(\bigcup\limits^n_{i=1}G_i\right)=\max\limits^n_{i=1}\chi\left(G_i\right) .
$$
Индукция по $p$. База: если $p\le 5$, то $\chi\le p\le 5$. Пусть
теорема верна для всех связных планарных графов с $p$ вершинами.
Рассмотрим граф $G$ с $p+1$ вершинами. По третьему следствию к
формуле Эйлера $\E{v\in V}{d(v)\le 5}$. По индукционному
предположению $\chi(G-v)\le 5$. Нужно раскрасить вершину $v$. Если
$d(v)<5$, то в $5$-раскраске графа $G-v$ существует цвет,
свободный для вершины $v$. Если $d(v)=5$ и для $\Gamma^+(v)$
использованы не все пять цветов, то в $5$-раскраске графа $G-v$
существует цвет, свободный для вершины $v$. 

Остался случай,
ког\-да $d(v)=5$ и все пять цветов использованы. 
Пусть $G_{13}$~---
правильный подграф графа $G-v$, порожденный всеми вершинами,
покрашенными в цвета $1$ или $3$ в $5$-рас\-краске графа $G-v$.
\end{proof}


\end{frame}



\begin{frame}
\frametitle{10.8.3. Теорема о пяти красках (2/2)}
\begin{proof} (Продолжение)
\begin{columns}[t]
\begin{column}[T]{8cm}

\vspace{-12pt}
\begin{figure}[h]
\begin{center}
\includegraphics[scale=1.0]{odm12-02}
\end{center}
\end{figure}
\end{column}
\begin{column}[T]{7cm}
Если $v_1$ и~$v_3$ принадлежат разным компонентам связности графа
$G_{13}$, то в той компоненте, в которой находится вершина $v_1$,
произведём перекраску $1\leftrightarrow 3$. При этом получится
$5$-раскраска графа $G-v$, но цвет $1$ будет свободен для вершины
$v$. 
\end{column}
\end{columns}

В противном случае существует простая цепь, соединяющая $v_1$
и $v_3$ и состоящая из вершин, покрашенных в цвета $1$ и $3$.
Тогда вершины $v_2$ и $v_4$ принадлежат разным компонентам
связности подграфа $G_{24}$ (так~как граф $G$~--- планарный).
Перекрасим вершины $2\leftrightarrow 4$ в той компоненте связности
графа $G_{24}$, которой принадлежит $v_2$, и получим $5$-раскраску
графа $G-v$, в которой цвет $2$ свободен для вершины~$v$.
\end{proof}

\end{frame}



